\documentclass[UTF8,12pt,a4paper,fontset=fandol]{ctexbook}


% 支持从项目根目录或当前笔记目录编译。
\makeatletter
\def\input@path{{../../}}
\makeatother
\usepackage{researchnotes}


\title{数学分析笔记}
\author{}
\date{\today}


\begin{document}
% 允许每一页的正文高度不完全相同，不强制把页面内容在垂直方向拉伸到页底。
% 密集公式页面保持正常的段间距，由页底承担余白。
\raggedbottom


\maketitle
\tableofcontents
\clearpage


\chapter{函数}
数学分析主要由微积分和级数理论组成，它所研究的主要对象是实函数，即以实数为自变量并且在实数中取值的函数。因此，在本章中我们首先简要地介绍一下实数系的连续性，然后介绍函数的概念和有关的基本知识。


\section{实数}
在近几个世纪中，科学技术之所以取得了辉煌的成就，在很大程度上是因为数学研究取得了重大进展，其中微积分的创立是现代数学的里程碑。微积分在物理、天文、技术、化学、生物等的研究中显示了强大的威力，解决了许多过去认为高不可攀的困难问题，促进了科学技术的发展。然而，由于在创建初期微积分是以几何直观和物理直觉为依据而进行演绎推理的，因此就形成了方法上有效但逻辑上不能自圆其说的矛盾局面。为了解决微积分在理论上面临的问题，许多著名数学家都投身于微积分理论基础的研究。人们后来发现，微积分的主要理论基础是严格的极限理论。到了 19 世纪初，柯西（Cauchy）以极限理论为微积分奠定了理论基础。但是柯西构筑的理论大厦起初并不完善，这是因为柯西并没有对实数给出严格的定义。而后来人们又发现，极限理论的某些基本原理依赖于实数系的连续性。为此，本节简要地介绍一下这方面的内容。


\subsection{数集}
在介绍实数理论前，我们简要地介绍一下集合的概念。在数学中，所谓\keyterm{集合}是将具有某种特性的对象的全体放在一起作为一个整体，通常用大写字母 $A,B,C,X,Y$ 等记之。这些对象称为集合中的\keyterm{元素}，通常用小写字母 $a,b,c,x,y$ 等记之。若 $a$ 是 $A$ 中的元素，则记为 $a\in A$，并读做 $a$ 属于 $A$。如果 $a$ 不是 $A$ 中的元素，则记为 $a\notin A$，并读做 $a$ 不属于 $A$。集合的概念是最基本的概念，它不能用别的概念来加以定义.


集合通常有列举法和描述法两种表示法。列举法是将集合中的元素全部列出，如：$A=\{1,2,3,\cdots,10\}$；而描述法是将集合的特性精确给出，如：$B=\{x:x\text{是 }x^2-1=0\text{ 的根}\}$；$E=\{x:x\text{是正有理数}\}$。一般来说，在数学分析中我们主要研究由数构成的集合，并且很少研究由有限个元素构成的集合；大部分集合是由无穷多个元素组成 (这种集合简称为无穷集合)，因此这些集合必须用描述法表示。


若集合 $A$ 中的每一个元素 $x$ 都属于集合 $B$，则称 $B$ 包含 $A$，记为 $A\subseteq B$，此时也称 $A$ 是 $B$ 的\keyterm{子集}。如果 $A\subseteq B$ 和 $B\subseteq A$ 同时成立，则认为 $A,B$ 是同一个集合，此时也记为 $A=B$。如上面的集合 $B=\{x:x\text{是 }x^2-1=0\text{ 的根}\}$ 与 $C=\{1,-1\}$ 是相等的。

若 $A\subseteq B$ 且 $A\ne B$，则称 $A$ 是 $B$ 的\keyterm{真子集}，记为 $A\subset B$。特别地，我们引入\keyterm{空集} $\varnothing$，即 $\varnothing$ 中不含有任何元素，因此 $\varnothing$ 是任何集合的子集.

集合有如下的运算：给定集合 $A,B$，有
\[
  A\cup B=\{x:x\in A\text{ 或 }x\in B\}
\]
称为 $A$ 与 $B$ 的\keyterm{并}；
\[
  A\cap B=\{x:x\in A\text{ 且 }x\in B\}
\]
称为 $A$ 与 $B$ 的\keyterm{交}；
\[
  A\setminus B=\{x:x\in A\text{ 且 }x\notin B\}
\]
称为 $A$ 与 $B$ 的\keyterm{差}.

读者容易将集合的并与交推广到多个集合的情形.

为了数学论述简便，习惯上常将逻辑推理中一些常用的词用符号表示。现将本书常用的两个数学记号 $\forall$ 和 $\exists$ 说明如下：
\[
  \forall \quad \text{代表``任意一个''；}
\]
\[
  \exists \quad \text{代表``存在''.}
\]
利用上述记号可以使得我们以后在行文上非常方便。例如，我们可以将 $A\subseteq B$ 的定义叙述为：$\forall x\in A$，有 $x\in B$.

\subsection{实数系的连续性}
\label{subsec:ma-1-1-2}

人类所认识的第一个数系就是自然数系，它的定义为 $\{0,1,2,\cdots\}$。虽然自然数对于计数来说是够用了，但是自然数系并不是一个完善的数系。首先作为量的描述手段，它只能表示一个单位量的整数倍，而无法去表示此单位量的部分。此外，作为量的运算手段，它只能自由地进行加、乘运算，而不能自由地进行加、乘运算的逆运算。自然数的这种离散性和运算的不完备性，促使人们去对它进行扩充。人们首先引进了负数，得到了整数系。对整数系人们可以自由进行加、减运算。人们为了得到可以自由地进行加、减、乘、除四则运算的数系，对整数系中的任意两个数进行加、减、乘、除 (除数不为零) 得到的数的全体记为一个新的数系，这就是随后得到的有理数系.

对于一个数集 $K$，若 $K$ 中至少有一个非零元素，且 $K$ 中任何两个元素的加、减、乘、除 (除数不为零) 运算后得到的数仍然属于 $K$，即 $K$ 关于四则运算封闭，则称 $K$ 为一个数域。容易看出，有理数集就是一个数域。从代数上来讲，有理数系已经是一个完美的数系，它可以在任何精度要求下对一个量进行表示和实施有效的运算，并且任何两个有理数之间必有有理数存在 (或说有理数有稠密性)。然而，早在 2500 年前，人们就发现有理数系也有缺陷。例如，若用 $c$ 来表示一个边长为 $1$ 的正方形的对角线的长度，则 $c$ 就无法用有理数表示。所以，有理数虽然在数轴上密密麻麻，但并没有布满整个数轴，留有许多``空隙''。这说明还有新的数存在，但它没有被严格地定义。由于有理数系已经对四则运算封闭，因此我们必须用不同于以前的数系的扩充才有可能得到新的数。我们前面曾经指出，微积分的理论基础是极限论，但在这种具有``空隙''的有理数系上实施极限运算，就会极为不便。因此，正如四则运算需要一个封闭的数域一样，极限运算也需要一个关于极限封闭的数域。这就需要将有理数进行扩充，使其能够填补有理数在数轴上留下的所有这些``空隙''。于是，这又归结到如何定义无理数的问题上了。这一历史任务，终于在 19 世纪后半叶，由戴德金 (Dedekind) 和康托尔 (Cantor) 等人完成。以下我们简要地介绍一下戴德金关于实数的构造方法.

\begin{definition}\label{definition:ma-1-1-1}
设 $S$ 是一个有大小顺序的非空数集，$A$ 和 $B$ 是它的两个子集，如果它们满足以下条件：
\begin{enumerate}
  \item $A\ne\varnothing,\ B\ne\varnothing$；
  \item $A\cup B=S$；
  \item $\forall a\in A,\ \forall b\in B$，都有 $a<b$；
  \item $A$ 中无最大数，
\end{enumerate}
则我们将 $A,B$ 称为 $S$ 的一个\keyterm{分划}，记为 $(A|B)$.
\end{definition}

现在我们来考虑有理数系 $\mathbb{Q}$ 的分划。对 $\mathbb{Q}$ 的任意分划 $(A|B)$，必有以下两种情形之一发生：
\begin{enumerate}
  \item $B$ 中存在最小数，此时称 $(A|B)$ 是一个\keyterm{有理分划}；
  \item $B$ 中不存在最小数，此时称 $(A|B)$ 是一个\keyterm{无理分划}.
\end{enumerate}

\begin{example}\label{example:ma-1-1-1}
记
\[
  A=\{x\in\mathbb{Q}:x\leqslant0\text{ 或 }x>0\text{ 且 }x^2<2\},\qquad B=\mathbb{Q}\setminus A,
\]
则 $(A|B)$ 是一个无理分划.
\end{example}

有理数系 $\mathbb{Q}$ 的所有分划构成了一个集合，我们称这个集合为\keyterm{实数系}，并记为 $\mathbb{R}$。显然，有理数集与 $\mathbb{R}$ 中的有理分划是一一对应的。因此 $\mathbb{R}$ 可以被认为是由有理数集加上无理分划所构成。我们称 $\mathbb{R}$ 中的这些无理分划为\keyterm{无理数}。换句话说，$\mathbb{R}$ 是由全体有理数与无理数所构成的集合。例如，例~\ref{example:ma-1-1-1} 中的无理分划对应的就是大家熟知的无理数 $\sqrt{2}$。我们可以在 $\mathbb{R}$ 上很自然地定义大小顺序关系及四则运算，经过冗长但不是很困难的论证，可以证明它是一个以有理数域为其子域的有序域.

前面我们已经指出有理数系有稠密性，但没有连续性，即有理数之间有许多``空隙''。下面的戴德金分割定理告诉我们，在有理数集合中加入无理数之后，就没有``空隙''了，也就是说实数系具有了连续性。对于一个数集 $A$，若对任意两个数 $a,b\in A$，$a,b$ 之间的所有的数都在 $A$ 中，则称 $A$ 是\keyterm{连通的}。换句话说，实数集是一个连通的集合.

\begin{theorem}[{戴德金分割定理}]\label{theorem:ma-1-1-1}
对 $\mathbb{R}$ 的任一分划 $(A|B)$，$B$ 中必有最小数.
\end{theorem}

定理~\ref{theorem:ma-1-1-1} 告诉我们，对 $\mathbb{R}$ 再进行分划就不可能产生新数了。换句话说，$\mathbb{R}$ 已将整个实轴填满了。因此 $\mathbb{R}$ 是一个有序的连通域。对于每一个实数 $x$，若 $x\geqslant0$，则它在原点 $O$ 的右边且它到 $O$ 的距离为 $x$；若 $x<0$，则它在原点 $O$ 的左边且它到 $O$ 的距离为 $-x$。根据上面的对应法则，显然，对于 $\mathbb{R}$ 中的每个数，在数轴上有且只有一个点与之对应。反之，任给数轴上的一个点 $A$，若 $A$ 就是原点 $O$，则它与 $0\in\mathbb{R}$ 对应；若 $A$ 在原点 $O$ 的右边，则它与 $A$ 到 $O$ 的距离 $x\in\mathbb{R}$ 对应；若 $A$ 在原点 $O$ 的左边，则它与 $-x$ 对应，其中 $x$ 为 $A$ 到 $O$ 的距离。这样，我们就建立了 $\mathbb{R}$ 中的数与数轴上的点之间的一一对应。因此，今后我们将不再区分数轴上的点与它所对应的实数.

\subsection{有界集与确界}
\label{subsec:ma-1-1-3}

设集合 $E\subseteq\mathbb{R}$，并且 $E\ne\varnothing$。如果存在 $M\in\mathbb{R}$，使得对 $\forall x\in E$，有 $x\leqslant M$，则称 $E$ 是\keyterm{有上界的}，并且说 $M$ 是 $E$ 的一个\keyterm{上界}；如果存在 $m\in\mathbb{R}$，使得对 $\forall x\in E$，有 $x\geqslant m$，则称 $E$ 是\keyterm{有下界的}，并且说 $m$ 是 $E$ 的一个\keyterm{下界}；如果 $E$ 既有上界又有下界，则称 $E$ 是\keyterm{有界的}。显然，$E$ 是有界的充分必要条件是：存在 $M>0$，使得对任意的 $x\in E$，有 $|x|\leqslant M$.

\begin{example}\label{example:ma-1-1-2}
集合
\[
  E=\left\{0,\frac32,\frac23,\cdots,1+\frac{(-1)^n}{n},\cdots\right\}
\]
是有界集，它的一个上界是 $\dfrac32$，一个下界是 $0$.
\end{example}


\begin{example}\label{example:ma-1-1-3}
集合
\[
  E=\left\{1,\frac12,\frac13,\cdots,\frac1n,\cdots\right\}
\]
是有界集，其中 $1$ 是它的一个上界，$0$ 是它的一个下界.
\end{example}

若 $E$ 是一个有界数集时，它的上 (下) 界必有无穷多个。是否有一个最小 (大) 的上 (下) 界呢？为了对此进行精确描述，我们引入：

\begin{definition}\label{definition:ma-1-1-2}
设 $E\subset\mathbb{R}$ 为一个非空数集。若有 $M\in\mathbb{R}$ 满足
\begin{enumerate}
  \item $M$ 是 $E$ 的一个上界，即 $\forall x\in E$，有 $x\leqslant M$；
  \item 对 $\forall\varepsilon>0$，存在 $x'\in E$，使得 $x'>M-\varepsilon$，
\end{enumerate}
则称 $M$ 为 $E$ 的\keyterm{上确界}，记为 $M=\sup E=\sup_{x\in E}\{x\}$.
\end{definition}

若有 $m\in\mathbb{R}$ 满足
\begin{enumerate}
  \item $m$ 是 $E$ 的一个下界，即 $\forall x\in E$，有 $x\geqslant m$；
  \item 对 $\forall\varepsilon>0$，存在 $x'\in E$，使得 $x'<m+\varepsilon$，
\end{enumerate}
则称 $m$ 为 $E$ 的\keyterm{下确界}，记为 $m=\inf E=\inf_{x\in E}\{x\}$.

从定义容易看出，$E$ 的上确界就是它的最小上界，而下确界就是它的最大下界。如果 $\sup E\in E$，则 $\sup E$ 可记为 $\max E$，这时上确界即为 $E$ 中的最大数；如果 $\inf E\in E$，则 $\inf E$ 可记为 $\min E$，这时下确界即为 $E$ 中的最小数.

此外，引入记号 $+\infty$ 和 $-\infty$，分别表示正无穷大和负无穷大，简称正无穷和负无穷。为了以后行文方便，我们约定：当数集 $E$ 无上界时，记做 $\sup E=+\infty$；当数集 $E$ 无下界时，记做 $\inf E=-\infty$。这里需特别指出的是，$+\infty$ 和 $-\infty$ 只是记号而不是实数。当数集 $E$ 满足 $\sup E=+\infty(\inf E=-\infty)$ (即 $E$ 是无上界 (下界)) 时，它的等价说法是：$\forall M>0,\exists x\in E$，并且有 $x>M(x<-M)$.

\begin{example}\label{example:ma-1-1-4}
设
\[
  E=\left\{\frac pq:p,q\text{ 为正整数，而且满足 }p\leqslant q\right\},
\]
则有 $\inf E=0\notin E,\ \sup E=1=\max E$。这表明，$E$ 中没有最小数，而有最大数 $1$.
\end{example}

下面的定理与戴德金分割定理等价，即如下定理是实数系连续性的另一种表述，但它使用起来特别方便.

\begin{theorem}[{确界存在定理}]\label{theorem:ma-1-1-2}
非空有上界的实数集必有上确界；非空有下界的实数集必有下确界.
\end{theorem}


\begin{proof}
我们只证明上确界的情形，下确界的情形完全类似。设 $E$ 是一个非空有上界的实数集。若 $E$ 中存在最大数 $M$ 时，则
\[
  \sup E=\max E=M.
\]
现假设 $E$ 中没有最大数。我们对 $\mathbb{R}$ 作分划：$B$ 是由 $E$ 的所有上界组成的集合，而 $A=\mathbb{R}\setminus B$。由 $E$ 的有界性，推出 $B\ne\varnothing$；而由 $E\ne\varnothing$，推出 $A\ne\varnothing$。显然，对任意 $a\in A$ 和 $b\in B$，有 $a<b$，而且 $A$ 中无最大数。因此，$(A|B)$ 是 $\mathbb{R}$ 的一个分划，从而由戴德金分割定理知，$B$ 中存在最小数 $M$，即 $M=\sup E$.
\end{proof}


\subsection{几个常用不等式}
\label{subsec:ma-1-1-4}

下面介绍几个常用的不等式。实数 $x$ 的绝对值定义为
\[
  |x|=\begin{cases}
    x, & x\geqslant0,\\
    -x, & x<0.
  \end{cases}
\]
由此可知，对任何 $x,y\in\mathbb{R}$，有
\[
  -|x|\leqslant x\leqslant |x|,\qquad -|y|\leqslant y\leqslant |y|,
\]
将此两个不等式相加，即得
\[
  -(|x|+|y|)\leqslant x+y\leqslant |x|+|y|,
\]
因此有
\[
  |x+y|\leqslant |x|+|y|.
\]
这一不等式通常称做\keyterm{三角不等式}。运用三角不等式，可以得到
\[
  \bigl||x|-|y|\bigr|\leqslant |x-y|.
\]

\begin{theorem}[{伯努利 (Bernoulli) 不等式}]\label{theorem:ma-1-source-10}
对任意的 $x\geqslant -1$ 和任意正整数 $n$，有
\[
  (1+x)^n\geqslant 1+nx.
\]
\end{theorem}


\begin{proof}
我们采用数学归纳法来证明这一不等式。当 $n=1$ 时，上述不等式显然成立。假设我们已经证明了
\[
  (1+x)^{n-1}\geqslant 1+(n-1)x
\]
对一切的 $x\geqslant -1$ 成立，则有
\[
\begin{aligned}
  (1+x)^n
  &= (1+x)^{n-1}(1+x)\\
  &\geqslant (1+(n-1)x)(1+x)\\
  &=1+nx+(n-1)x^2\\
  &\geqslant 1+nx.
\end{aligned}
\]
由数学归纳法原理知，伯努利不等式对一切的正整数 $n$ 成立.
\end{proof}


\begin{theorem}[{算术--几何平均不等式}]\label{theorem:ma-1-source-12}
对任意 $n$ 个非负实数 $x_1,x_2,\cdots,x_n$，有
\[
  \frac{x_1+x_2+\cdots+x_n}{n}\geqslant \sqrt[n]{x_1x_2\cdots x_n}.
\]
\end{theorem}


\begin{proof}
我们同样采用数学归纳法来证明这一不等式。当 $n=1$ 时，上述不等式显然成立。假设我们已经证明了，对任意 $n-1$ 个非负实数，算术--几何平均不等式成立。下面我们来考虑任意 $n$ 个非负实数 $x_1,x_2,\cdots,x_n$ 的情形。不妨假定 $x_n$ 是这 $n$ 个数中最大的一个，并且令
\[
  y=\frac{x_1+x_2+\cdots+x_{n-1}}{n-1},
\]
则有
\[
  x_n\geqslant y=\frac{x_1+x_2+\cdots+x_{n-1}}{n-1}
  \geqslant \sqrt[n-1]{x_1x_2\cdots x_{n-1}},
\]
从而有
\[
\begin{aligned}
  \left(\frac{x_1+x_2+\cdots+x_n}{n}\right)^n
  &=\left(y+\frac{x_n-y}{n}\right)^n\\
  &=y^n+ny^{n-1}\frac{x_n-y}{n}+\cdots+\left(\frac{x_n-y}{n}\right)^n\\
  &\geqslant y^n+ny^{n-1}\frac{x_n-y}{n}\\
  &=y^{n-1}x_n\\
  &\geqslant x_1x_2\cdots x_n,
\end{aligned}
\]
即有
\[
  \frac{x_1+x_2+\cdots+x_n}{n}\geqslant \sqrt[n]{x_1x_2\cdots x_n}.
\]
由数学归纳法原理知，算术--几何平均不等式对一切的正整数成立.
\end{proof}


\subsection{常用记号}
\label{subsec:ma-1-1-5}

全体正整数组成的集合、全体整数构成的集合、全体有理数组成的集合和全体实数组成的集合是本书最常遇到的数集，我们约定分别用空心字母 $\mathbb{N},\mathbb{Z},\mathbb{Q}$ 和 $\mathbb{R}$ 表示，即
\[
\begin{array}{ll}
  \mathbb{N} & \text{表示全体正整数组成的集合；}\\
  \mathbb{Z} & \text{表示全体整数组成的集合；}\\
  \mathbb{Q} & \text{表示全体有理数组成的集合；}\\
  \mathbb{R} & \text{表示全体实数组成的集合.}
\end{array}
\]
显然有
\[
  \mathbb{N}\subset\mathbb{Z}\subset\mathbb{Q}\subset\mathbb{R}.
\]
本书中常要用到以下形式的实数子集：
\[
  [a,b]=\{x\in\mathbb{R}:a\leqslant x\leqslant b\};
\]
\[
  (a,b)=\{x\in\mathbb{R}:a<x<b\};
\]
\[
  (a,b]=\{x\in\mathbb{R}:a<x\leqslant b\};
\]
\[
  [a,b)=\{x\in\mathbb{R}:a\leqslant x<b\}.
\]
这里 $a,b\in\mathbb{R}$，且 $a<b$.

推广区间的概念，把
\[
  (-\infty,b),\quad (-\infty,b],\quad (a,+\infty),\quad [a,+\infty),\quad (-\infty,+\infty)
\]
都称做\keyterm{无穷区间}。读者应该清楚知道上面每个无穷区间的精确定义，例如 $(-\infty,b]=\{x\in\mathbb{R}:x\leqslant b\}$.

今后，我们提及的区间必是以上形式的区间之一。特别地，对 $a\in\mathbb{R}$，我们把开区间 $(a-\varepsilon,a+\varepsilon)(\varepsilon>0)$ 称为 $a$ 的一个 $\varepsilon$\keyterm{邻域}，记为 $U(a,\varepsilon)$；而把 $(a-\varepsilon,a+\varepsilon)\setminus\{a\}$ 称为 $a$ 的一个\keyterm{去心邻域}，记为 $U_0(a,\varepsilon)$.

在本书第一册中，我们还要遇到平面内的点集。在平面直角坐标系中，平面上的点集 $E$ 可以通过描述 $E$ 中元素的坐标所具有的特性来定义。如 $E=\{(x,y):x^2+y^2=1\}$ 就是平面内以原点 $O$ 为圆心，以 $1$ 为半径的圆周上的点组成的集合.

\section{函数的概念}
\label{sec:ma-1-2}

\subsection{函数的定义}
\label{subsec:ma-1-2-1}

当我们观察或者研究客观世界的各种事物时，会遇到很多不同的量。这些量一般来说都是在不断变化的，这是客观世界不断变化、不断运动、不断发展在量的方面的体现。例如，当我们研究一个自由落体运动时，就会遇到地球的引力、空气的浮力、时间、物体的高度等各种量，这些量都是变化的，叫做变量。我们知道，客观世界中各个事物之间是相互联系、相互依赖与相互制约的。因此，从量的方面来看，各个变量之间也是相互联系和相互制约的。变量之间的这种相互依赖的关系就是所谓的函数关系.

例如，若在研究自由落体运动时设初始高度为 $h_0$，则在不计空气阻力等其他因素的前提下，该物体距地面的高度 $h$ 与时间 $t$ 有以下关系：
\begin{equation}
  h=h_0-\frac12gt^2,\qquad t\in[0,\sqrt{h_0/g}],
  \label{eq:ma-1-2-1}
\end{equation}
其中 $g$ 是重力加速度.

在以上自由落体运动过程中，$g$ 始终保持不变，因此它是一个常量；当时间 $t$ 在 $[0,\sqrt{h_0/g}]$ 中不断变化时，$h$ 也随着 $t$ 的变化而不断地变化，因此它们都是变量。在这一过程中，只要知道时间 $t$，根据关系 \eqref{eq:ma-1-2-1}，我们就可以知道物体距地面的高度 $h$。换句话说，对于每一个 $t\in[0,\sqrt{h_0/g}]$，根据 \eqref{eq:ma-1-2-1} 所规定的对应法则，就有唯一确定的一个高度 $h$ 与之对应。这就是自由落体运动中的函数关系。一般地，我们有

\begin{definition}\label{definition:ma-1-2-1}
对于给定的集合 $X\subseteq\mathbb{R}$，如果存在某种对应法则 $f$，使得对 $X$ 中的每一个数 $x$，在 $\mathbb{R}$ 中存在唯一的数 $y$ 与之对应，则称对应法则 $f$ 为从 $X$ 到 $\mathbb{R}$ 的一个\keyterm{函数}，记做
\[
  f:X\to\mathbb{R},
\]
\[
  x\mapsto y=f(x),
\]
其中 $y$ 称为 $f$ 在点 $x$ 的值，$X$ 称为函数 $f$ 的定义域，数集 $\{f(x):x\in X\}$ 称为函数 $f$ 的值域，记为 $f(X)$；$x$ 称做自变量，$y$ 称做因变量.
\end{definition}

简言之，函数就是由定义域到它的值域的一个单值对应。此外，构成一个函数必须具备三个基本要素：定义域、值域和对应法则。但其中最重要的两个基本的要素是定义域和对应法则，而值域可由定义域和对应法则所确定。通常，我们用 $y=f(x)(x\in X)$ 记定义在 $X$ 上的一个函数。有时，我们也用 $y=g(x),y=\varphi(x)$ 等记号来表示不同的函数.

在自由落体的运动中，$X=[0,\sqrt{h_0/g}]$ 是定义域，对应法则是 $f$：对每一个 $t\in[0,\sqrt{h_0/g}]$，由解析式 $h=h_0-\frac12gt^2$ 来确定对应的数 $h\in\mathbb{R}$。由于对每一个 $t$，此解析式可确定唯一的数 (即物体高度) 与之对应。因此，物体自由下落这个过程确定了其高度与时间的函数关系.

设 $y=f(x)$ 是定义在 $X$ 上的一个函数，通常称平面点集
\[
  T=\{(x,y):y=f(x),x\in X\}
\]
为该函数的\keyterm{图像}。一般来讲，它是平面内的一条曲线，它和任何一条平行于 $y$ 轴的直线至多只有一个交点.

在中学数学中，我们所遇到的主要是如下的六类函数：
\begin{enumerate}
  \item 常值函数 $y=C$；
  \item 幂函数 $y=x^\alpha\ (\alpha>0)$；
  \item 指数函数 $y=a^x\ (a>0)$；
  \item 对数函数 $y=\log_a x\ (a>0,\ a\ne1)$；
  \item 三角函数 $y=\sin x,\ y=\cos x,\ y=\tan x,\ y=\cot x$；
  \item 反三角函数 $y=\arcsin x,\ y=\arccos x,\ y=\arctan x,\ y=\operatorname{arccot} x$.
\end{enumerate}
这些函数统称为\keyterm{基本初等函数}，它们的定义域是使得其解析式有意义的数 $x$ 的全体。在本书中，它们仍将是我们研究的主要对象。基本初等函数的主要特点是函数关系有明显的解析表达式。值得注意的是，这些解析式有着各自特殊的意义。如 $y=\sin x$，我们是用以下方式来确定这个函数关系的.

对每个 $x\in[0,2\pi)$，在 $O\xi\eta$ 坐标平面中取射线 $\overrightarrow{OP}$，使得 $O\xi$ 轴正向逆时针旋转到 $\overrightarrow{OP}$ 的角度为 $x$。设 $\overrightarrow{OP}$ 与单位圆周的交点为 $P$，则 $P$ 的纵坐标即是 $y=\sin x$ 的值 (图 1.2.1)；再规定 $x$ 与 $x+2k\pi\ (k\in\mathbb{Z})$ 取相同的函数值，就得到了定义在 $(-\infty,+\infty)$ 上的函数 $y=\sin x$.

另外，在中学数学中，对某些函数关系也没有给出精确定义。如，在初等数学中就很难弄清楚 $(\sqrt2)^\pi$ (它是函数 $f(x)=x^\pi$ 在 $x=\sqrt2$ 处的取值) 是如何定义的。我们将在后续章节中给出它们的精确定义。今后对于具有解析表达式的函数，若无特殊说明，该函数的定义域即为使得该解析式有意义的自变量的全体所构成的集合。因此，有时候我们对具有解析式的函数不特别指出其定义域.

作为练习，请读者做出基本初等函数的图像.

函数的表示法有穷举法、描述法、列表法及图形法等，以上表示法在中学教科书中已有相应的例子，我们在这里不再赘述。下面我们给出几个中学数学中未曾出现的但在数学分析中经常用到的函数的例子.

\begin{example}[{符号函数}]\label{example:ma-1-2-1}
所谓符号函数，是指如下函数：
\[
  y=\operatorname{sgn}x=
  \begin{cases}
    1, & x>0,\\
    0, & x=0,\\
    -1, & x<0,
  \end{cases}
\]
其图像如图 1.2.2 所示。这个函数的表达式是分段给出的，通常称这样的函数为分段函数。读者应该注意它是一个函数，而不是三个函数.
\end{example}


\begin{example}[{狄利克雷 (Dirichlet) 函数}]\label{example:ma-1-2-2}
函数
\[
  y=D(x)=
  \begin{cases}
    1, & x\in\mathbb{Q},\\
    0, & x\in\mathbb{R}\setminus\mathbb{Q}
  \end{cases}
\]
称为狄利克雷函数。这一函数的图像，是分布在 $y=0$ 和 $y=1$ 这两条直线上的两个离散的点集 (参见图 1.2.3)。注意，这样函数的图像是无法精确画出的.
\end{example}


\begin{example}[{高斯 (Gauss) 取整函数}]\label{example:ma-1-2-3}
取整函数定义为 $y=[x]$，其中 $[x]$ 表示不超过 $x$ 的最大整数，即若 $n\leqslant x<n+1$，则 $[x]=n$。其图像如图 1.2.4 所示.
\end{example}


\begin{example}[{黎曼 (Riemann) 函数}]\label{example:ma-1-2-4}
所谓的黎曼函数，是指由下式确定的函数：
\[
  y=R(x)=
  \begin{cases}
    \dfrac1p, & x=\dfrac qp\in(0,1),\ p,q\text{ 为互素的正整数},\\
    0, & x\in[0,1]\setminus\mathbb{Q},\\
    1, & x=0\text{ 或 }1.
  \end{cases}
\]
它的图像也不是连续不断的曲线 (参见图 1.2.5)，也是无法精确画出的.
\end{example}


\begin{example}[{特征函数}]\label{example:ma-1-2-5}
设 $E\subset\mathbb{R}$，定义在 $\mathbb{R}$ 中的函数
\[
  y=\chi_E(x)=
  \begin{cases}
    1, & x\in E,\\
    0, & x\notin E
  \end{cases}
\]
称为集 $E$ 的特征函数.
\end{example}


\subsection{由已知函数构造新函数的方法}
\label{subsec:ma-1-2-2}

从已知函数出发，可以通过代数四则运算、限制与延拓、复合运算和取反函数等手段构造出许多新的函数.

\paragraph*{函数的四则运算}


设 $y=f_j(x),x\in X_j\subset\mathbb{R}\ (j=1,2)$ 为两个已知函数，且 $X=X_1\cap X_2\ne\varnothing$，则我们可以利用实数的四则运算构造新函数如下：
\[
  (f_1\pm f_2)(x)=f_1(x)\pm f_2(x),\qquad x\in X;
\]
\[
  (f_1f_2)(x)=f_1(x)f_2(x),\qquad x\in X;
\]
\[
  \frac{f_1}{f_2}(x)=\frac{f_1(x)}{f_2(x)}\ (f_2(x)\ne0),\qquad x\in X.
\]
显然，函数的加法、减法与乘法运算可推广到任意有限个函数的情形，而且对于加法和乘法运算，它们具有交换律与结合律.

\paragraph*{函数的限制与延拓}


设函数 $f(x),x\in X_1$ 和 $g(x),x\in X_2$ 满足：$X_1\subset X_2$，且 $f(x)=g(x),\forall x\in X_1$，则称 $f(x)$ 是 $g(x)$ 在 $X_1$ 上的\keyterm{限制}，而 $g(x)$ 是 $f(x)$ 在 $X_2$ 上的\keyterm{延拓}。这样，我们就可以从一个已知函数出发，通过限制和延拓来产生新函数.

这里须特别指出的是，对于两个函数而言，只有当它们的定义域及对应关系完全一样时，才能说它们是相等的！例如，$y_1=|x|,x\in(-\infty,+\infty)$ 和 $y_2=x\operatorname{sgn}x,x\in(-\infty,+\infty)$，虽然形式不同，但由于它们具有同样的定义域和对应关系，所以 $y_1$ 和 $y_2$ 是同一个函数；而 $f(x)=x^2,x\in(-\infty,+\infty)$ 和 $g(x)=x^2,x\in(0,+\infty)$ 就是两个不同的函数.

\paragraph*{函数的复合}


函数的复合运算是数学中一种重要的运算。设 $y=f_j(x),x\in X_j\subset\mathbb{R}\ (j=1,2)$ 为两个函数，若 $Y_1=f_1(X_1)\subseteq X_2$，则定义在 $X_1$ 上的函数 $y=f_2(f_1(x))$ 称为 $f_1$ 和 $f_2$ 的\keyterm{复合函数}，记做 $f_2\circ f_1:X_1\to\mathbb{R}$。通常称 $f_1$ 为该复合函数的\keyterm{内函数}，$f_2$ 为\keyterm{外函数}.

同样我们可以考虑多个函数的复合 (如果复合有意义的话)。一般来说，函数复合不具有交换律。例如，取 $f(x)=x^2+1,x\in\mathbb{R},g(x)=x^4,x\in\mathbb{R}$，则 $f(g(x))=x^8+1,x\in\mathbb{R}$，而 $g(f(x))=(x^2+1)^4,x\in\mathbb{R}$.

读者要特别注意是：函数的复合运算可以进行的前提条件是，外函数的定义域必须包含内函数的值域。当然，当两个函数的定义域都是 $\mathbb{R}$ 时，则复合运算可以自由地进行.

函数的复合运算是构造新函数的重要途径.

\begin{example}\label{example:ma-1-2-6}
设 $f(x)=\dfrac{x}{\sqrt{1+x^2}}$，试求出 $n$ 次复合函数 $\underbrace{f\circ f\circ\cdots\circ f}_{n\text{ 个}}(x)$ 的表达式.
\end{example}


\begin{solve}
显然，$f(x)$ 的定义域是 $(-\infty,+\infty)$，因此所要讨论的复合运算是可行的。从
\[
  f\circ f(x)=
  \frac{\dfrac{x}{\sqrt{1+x^2}}}{\sqrt{1+\dfrac{x^2}{1+x^2}}}
  =\frac{x}{\sqrt{1+2x^2}}
\]
我们猜出 $n$ 次复合的结果应该为
\[
  \underbrace{f\circ f\circ\cdots\circ f}_{n\text{ 个}}(x)=\frac{x}{\sqrt{1+nx^2}}.
\]
下面用归纳法证之。当 $n=1$ 时，结论显然成立。现在假定当 $n=k$ 时，有
\[
  \underbrace{f\circ f\circ\cdots\circ f}_{k\text{ 个}}(x)=\frac{x}{\sqrt{1+kx^2}},
\]
则当 $n=k+1$ 时，有
\[
  f\circ(\underbrace{f\circ\cdots\circ f}_{k\text{ 个}}\circ f)(x)
  =f\left(\frac{x}{\sqrt{1+kx^2}}\right)
  =\frac{x}{\sqrt{1+(k+1)x^2}}.
\]
这样，由数学归纳法原理知，对一切的 $n$，有
\[
  \underbrace{f\circ f\circ\cdots\circ f}_{n\text{ 个}}(x)=\frac{x}{\sqrt{1+nx^2}}.
\qedhere
\]
\end{solve}


\paragraph*{反函数}


在函数的定义中，我们要求对 $\forall x\in X$，在 $\mathbb{R}$ 中必须有唯一的数 $f(x)$ 与之对应，但我们并没有要求 $X$ 中的任何两个不同的数 $x_1$ 和 $x_2$，在 $\mathbb{R}$ 中与之对应的两个数 $f(x_1)$ 和 $f(x_2)$ 也必须不同。最极端的例子是常值函数，其所有点上的函数值取同一个数.

设 $f:X\to Y$ 是一个函数。若对任意的 $x_1,x_2\in X$，只要 $x_1\ne x_2$，就有 $f(x_1)\ne f(x_2)$ 成立，则称 $f(x)$ 是\keyterm{单的}；若 $Y=f(X)$，则称 $f(x)$ 为\keyterm{满的}；若 $f(x)$ 既是单的又是满的，则称它为一个\keyterm{一一对应}.

\begin{definition}\label{definition:ma-1-2-2}
设 $f:X\to Y$ 是一个一一对应。定义函数 $g:Y\to X$ 如下：对任意的 $y\in Y$，函数值 $g(y)$ 规定为由关系式 $y=f(x)$ 所唯一确定的 $x\in X$。这样定义的函数 $g(y)$ 称为是函数 $f(x)$ 的\keyterm{反函数}，记为 $g=f^{-1}$.
\end{definition}

反函数 $x=f^{-1}(y)$ 的定义域和值域恰为原来函数 $y=f(x)$ 的值域和定义域。函数 $f$ 和其反函数 $f^{-1}$ 满足：
\[
  f(f^{-1}(y))=y,\qquad \forall y\in Y,
\]
\[
  f^{-1}(f(x))=x,\qquad \forall x\in X.
\]
习惯上，我们总是把自变量记为 $x$，因变量记为 $y$。因此，为了遵从统一性，$y=f(x)$ 的反函数也记为 $y=f^{-1}(x)$。这样一来，从几何上来看，$y=f(x)$ 的图像与它的反函数 $y=f^{-1}(x)$ 的图像正好关于直线 $y=x$ 对称.

在中学数学中，我们曾经学习过反三角函数。回想一下，在定义这类函数时，我们必须将三角函数限制在某个区间上。这说明，一个函数在它的定义域中可能不是单的，但把它做适当的限制后则常常可以变成单的，从而可以定义它的反函数.

\begin{example}\label{example:ma-1-2-7}
求 $y=f(x)=\dfrac{e^x-e^{-x}}{2}$ 的反函数，这里 $e=2.7182\cdots$ 为一常数.
\end{example}


\begin{solve}
固定 $y$，为了从方程 $y=\dfrac{e^x-e^{-x}}{2}$ 中求出 $x$，我们令 $e^x=z$，则有
\[
  2zy=z^2-1.
\]
解之，得
\[
  z=y\pm\sqrt{y^2+1}.
\]
由于 $z=e^x>0$，故舍去 $y-\sqrt{y^2+1}$，从而有 $x=\ln(y+\sqrt{y^2+1})$。因此，所求的反函数为
\[
  y=f^{-1}(x)=\ln(x+\sqrt{x^2+1}),\qquad x\in(-\infty,+\infty).
\]

本题中的函数 $f(x)$ 称为双曲正弦函数，并记为 $f(x)=\sinh x$ (以后我们将给出所有双曲函数的定义)，它的反函数 $\sinh^{-1}x$ 称为反双曲正弦函数。双曲正弦函数和反双曲正弦函数的图像如图 1.2.6 所示.
\end{solve}


\begin{example}\label{example:ma-1-2-8}
设 $y=f(x)$ 为定义在 $X$ 上的一个函数，并且记 $Y=f(X)$。证明：若存在 $Y$ 上定义的函数 $g(y)$，使得 $g(f(x))=x$，则 $f(x)$ 的反函数存在，而且 $g=f^{-1}$.
\end{example}


\begin{proof}
容易验证，复合函数是单的充分必要条件是其内、外函数都是单的。由 $g(f(x))=x$ 是 $X$ 的恒等函数可知，$f(x)$ 是单的，从而 $f(x)$ 是 $X$ 到 $Y$ 的一一对应。于是，$f(x)$ 的反函数存在，而且对 $\forall y\in Y=f(X)$，有 $g(y)=g(f(x))=x=f^{-1}(y)$，即有 $g=f^{-1}$.

作为本节的结束，我们来讨论一下分式线性函数的反函数和复合函数。形如
\[
  y=\frac{ax+b}{cx+d}\qquad (ad-bc\ne0)
\]
的函数称为\keyterm{分式线性函数}，这里的条件 $ad-bc\ne0$ 是防止它变成一个常值函数.

由 $y=\dfrac{ax+b}{cx+d}$，可推出 $x=\dfrac{b-dy}{cy-a}$，而且有
\[
  (-d)(-a)-bc=ad-bc\ne0.
\]
这表明，任何一个分式线性函数都具有反函数，而且其反函数仍然为分式线性函数.

另外，设
\[
  f_j(x)=\frac{a_jx+b_j}{c_jx+d_j}\quad (a_jd_j-b_jc_j\ne0),\qquad j=1,2,
\]
则容易导出
\[
  f_1\circ f_2(x)=
  \frac{(a_1a_2+b_1c_2)x+a_1b_2+b_1d_2}
  {(c_1a_2+d_1c_2)x+c_1b_2+d_1d_2},
\]
而且
\[
\begin{aligned}
 &(a_1a_2+b_1c_2)(c_1b_2+d_1d_2)
 -(a_1b_2+b_1d_2)(c_1a_2+d_1c_2)\\
 &\qquad=(a_1d_1-b_1c_1)(a_2d_2-b_2c_2)\ne0.
\end{aligned}
\]
这表明，任意两个分式线性函数的复合还是分式线性函数.

分式线性函数在数学的许多分支中起着重要的作用。这个函数类有着许多有趣的性质。例如，每个分式线性函数 $y=\dfrac{ax+b}{cx+d}(ad-bc\ne0)$ 可对应一个 $2\times2$ 矩阵 $\begin{pmatrix}a&b\\ c&d\end{pmatrix}$。有趣的是，任何分式线性函数的反函数所对应的矩阵刚好是它对应矩阵的逆乘以它的行列式的 $-1$ 倍，而任意两个分式线性函数的复合函数所对应的矩阵则刚好是它们各自对应矩阵的乘积.

细心的读者可以发现，对 $y=f(x)=\dfrac{ax+b}{cx+d}$ 来说，当 $c\ne0$ 时，$f(x)$ 在 $x=-\dfrac dc$ 没有定义；另外 $f(x)$ 的值域也不包含 $\dfrac ac$。如果我们形式地令 $f(-\dfrac dc)=\infty$ 和 $f(\infty)=\dfrac ac$，而当 $c=0$ 时令 $f(\infty)=\infty$，则可以验证 $f(x)$ 是 $\mathbb{R}\cup\{\infty\}$ 到 $\mathbb{R}\cup\{\infty\}$ 的一个一一对应，从而以上的求反函数和复合的运算可以自由地进行.
\end{proof}


\section{函数的性质}
\label{sec:ma-1-3}

在本节中，我们主要介绍某些函数具有的一些特别性质。这些性质的讨论是我们以后经常要遇到的，因此要求读者能用准确的数学语言来叙述它们的定义和一些等价的命题.

\subsection{函数的有界性}
\label{subsec:ma-1-3-1}

设 $y=f(x)$ 是定义在 $X$ 上的函数。若存在常数 $M$，使得对 $\forall x\in X$，都有 $f(x)\leqslant M$，则称 $f(x)$ 在 $X$ 上有\keyterm{上界}，同时称 $M$ 是 $f(x)$ 的一个上界；若存在常数 $m$，使得对 $\forall x\in X$，都有 $f(x)\geqslant m$，则称 $f(x)$ 在 $X$ 上有\keyterm{下界}，同时称 $m$ 是 $f(x)$ 的一个下界；若 $f(x)$ 在 $X$ 上既有上界 $M$ 又有下界 $m$，则称 $f(x)$ 在 $X$ 上\keyterm{有界}。此时，对 $\forall x\in X$，有 $m\leqslant f(x)\leqslant M$。若 $|f(x)|\leqslant M(x\in X)$，则称 $M$ 是 $f(x)$ 的一个\keyterm{界}.

容易看出，$f(x)$ 在 $X$ 上有界的充分必要条件是存在 $M>0$，使得当 $x\in X$ 时，有 $|f(x)|\leqslant M$，即 $f(X)\subseteq[-M,M]$。换句话说，$f(x)$ 在 $X$ 有界也就是其值域 $f(X)$ 是一个有界集.

从几何上来看，若函数 $y=f(x)$ 在 $X$ 上有上界 $M$，则 $f(x)$ 的图像将位于直线 $y=M$ 的下面；若 $f(x)$ 有下界 $m$，则 $f(x)$ 的图像在直线 $y=m$ 的上面；若 $f(x)$ 有界，则必存在 $M>0$，使得 $f(x)$ 的图像位于直线 $y=M$ 和 $y=-M$ 之间.

\begin{example}\label{example:ma-1-3-1}
函数 $y=\sin x$ 和 $y=\cos x$ 在 $\mathbb{R}$ 上有界，$M=1$ 就是它们的一个界。而对函数 $y=x^n,\ x\in(-\infty,+\infty)$ ($n$ 为正整数) 来说，当 $n$ 为奇数时，它既无上界，又无下界；而当 $n$ 为偶数时，它无上界，但有下界 $m=0$.
\end{example}

若 $f(x)$ 不是 $X$ 上的有界函数，则称该函数在 $X$ 上\keyterm{无界}。如何用肯定的语气来叙述函数的无界性呢？我们分析一下，$f(x)$ 在 $X$ 上无界的等价说法是，任何正数 $M$ 都不是 $y=f(x)$ 在 $X$ 上的界。而 $M$ 不是 $y=f(x)$ 在 $X$ 上的界，就等价于，必然 $\exists x_0\in X$，使得 $|f(x_0)|>M$。因此，若用肯定的语气，函数 $f(x)$ 在 $X$ 上无界可以叙述为：$\forall M>0$，$\exists x_0\in X$，使得 $|f(x_0)|>M$.

同理我们也可以给出函数 $f(x)$ 在 $X$ 上无上界或无下界的定义，请读者自己给出这些定义的精确数学语言的描述.

\begin{example}\label{example:ma-1-3-2}
证明：$f(x)=\dfrac{\sin\frac1x}{x}$ 在 $(0,1]$ 区间上既无上界也无下界，但它却在 $[a,1]$ 上有界，这里 $1>a>0$.
\end{example}


\begin{proof}
$\forall M>0$，取正整数 $n_0>M$ 并取
\[
  x_0=\frac{1}{2n_0\pi+\pi/2}\in(0,1],
\]
则有
\[
  f(x_0)=\frac{1}{x_0}=2n_0\pi+\frac{\pi}{2}>n_0>M.
\]
这就证明了 $y=\dfrac{\sin\frac1x}{x}$ 在 $(0,1]$ 上是无上界的.

取
\[
  x_1=\frac{1}{2n_0\pi-\pi/2}\in(0,1],
\]
则有
\[
  f(x_1)=\frac{-1}{x_1}=-\left(2n_0\pi-\frac{\pi}{2}\right)<-n_0<-M.
\]
这就证明了 $y=\dfrac{\sin\frac1x}{x}$ 在 $(0,1]$ 上是无下界的.

对 $a>0$，取 $M=\dfrac1a>0$，则 $\forall x\in[a,1]$，有
\[
  \left|\frac{\sin\frac1x}{x}\right|\leqslant \left|\frac1x\right|\leqslant M,
\]
即 $y=\dfrac{\sin\frac1x}{x}$ 在 $[a,1]$ 上有界.
\end{proof}


\subsection{函数的单调性}
\label{subsec:ma-1-3-2}

设 $y=f(x)$ 是定义在 $X$ 上的一个函数。若对任意的 $x_1,x_2\in X$，只要 $x_1<x_2$，便有 $f(x_1)\leqslant f(x_2)(f(x_1)\geqslant f(x_2))$，则称 $f(x)$ 在 $X$ 上是\keyterm{单调上升 (下降) 函数}或\keyterm{单调递增 (递减) 函数}。在上述不等式中将``$\leqslant$'' (``$\geqslant$'') 换成``$<$'' (``$>$'')，则称 $f(x)$ 在 $X$ 上是\keyterm{严格单调上升 (下降) 函数}。单调上升函数和单调下降函数统称为\keyterm{单调函数}.

严格单调的函数是定义域到值域的一一对应，所以它存在反函数。但存在反函数的函数未必是单调的.

\begin{example}\label{example:ma-1-3-3}
$y=\sin x(x\in\mathbb{R})$ 不是单调函数，但当其定义域 $X$ 取为 $\left[-\dfrac\pi2,\dfrac\pi2\right]$ 时，它是严格单调递增函数；当 $X$ 取为 $\left[\dfrac\pi2,\dfrac{3\pi}2\right]$ 时，它是严格单调递减函数。因此该函数在这两个区间都分别存在反函数。在习惯上，我们用 $y=\arcsin x$ 表示 $y=\sin x$ 在 $\left[-\dfrac\pi2,\dfrac\pi2\right]$ 上的反函数.
\end{example}


\begin{example}\label{example:ma-1-3-4}
考查函数
\[
  y=f(x)=
  \begin{cases}
    x, & x\in[0,1]\cap\mathbb{Q},\\
    1-x, & x\in[0,1]\setminus\mathbb{Q}.
  \end{cases}
\]
容易看出，该函数的反函数就是它自己，但 $f(x)$ 在 $[0,1]$ 的任何子区间都不是单调的.
\end{example}


\subsection{函数的周期性}
\label{subsec:ma-1-3-3}

设 $y=f(x)$ 是在 $X$ 上有定义的函数。若存在 $T>0$，使得对任意 $x\in X$ 有 $f(x+T)=f(x)$，则称 $f(x)$ 为\keyterm{周期函数}，$T$ 称为 $f(x)$ 的一个\keyterm{周期}。若存在一个最小的周期 $T_0$，则称 $T_0$ 为 $f(x)$ 的\keyterm{基本周期}.

显然，以 $T$ 为周期的周期函数 $f(x)$ 的定义域 $X$ 必须满足条件：对一切的 $x\in X$，有 $T+x\in X$.

\begin{example}\label{example:ma-1-3-5}
函数 $y=\tan x$ 的定义域为
\[
  \mathbb{R}\setminus\left\{\frac\pi2+n\pi:n\in\mathbb{Z}\right\},
\]
它是以 $\pi$ 为基本周期的周期函数.

这里须特别指出的是，并非所有的周期函数都有基本周期.
\end{example}


\begin{example}\label{example:ma-1-3-6}
容易验证，任何正有理数都是狄利克雷函数 $D(x)$ (参看例~\ref{example:ma-1-2-2}) 的周期。因此，它没有基本周期.
\end{example}

若一个定义在闭区间 $[a,b]$ 上的函数 $f(x)$ 满足 $f(a)=f(b)$，则可将它延拓成 $(-\infty,+\infty)$ 上以 $T=b-a$ 为周期的周期函数 $\tilde f(x)$.

\begin{example}\label{example:ma-1-3-7}
设 $f(x)$ 为定义在 $\mathbb{R}$ 上的周期函数，并且有基本周期 $\tau>0$。证明：若 $\forall x\in(0,\tau)$，有 $f(x)\ne f(0)$，则 $g(x)=f(x^2)$ 不是周期函数.
\end{example}


\begin{proof}
用反证法。假定 $g(x)$ 是周期函数，而且 $T>0$ 是它的一个周期，则有
\begin{equation}
  f((x+T)^2)=g(x+T)=g(x)=f(x^2),\qquad \forall x\in\mathbb{R}.
  \label{eq:ma-1-3-1}
\end{equation}
以 $x=0$ 代入 \eqref{eq:ma-1-3-1} 式，得 $f(T^2)=f(0)$。而 $f(x)$ 是以 $\tau$ 为基本周期的周期函数，因此必有某一正整数 $k$，使得 $T^2=k\tau$，即 $T=\sqrt{k\tau}$.

再以 $x=\sqrt{(k+1)\tau}$ 代入 \eqref{eq:ma-1-3-1} 式，得
\[
  f\left(\left(\sqrt{(k+1)\tau}+\sqrt{k\tau}\right)^2\right)=f((k+1)\tau)=f(0),
\]
而
\[
  f\left(\left(\sqrt{(k+1)\tau}+\sqrt{k\tau}\right)^2\right)
  =f\left((2k+1)\tau+2\sqrt{(k+1)k}\tau\right)
  =f(2\sqrt{(k+1)k}\tau),
\]
所以
\[
  f(2\sqrt{(k+1)k}\tau)=f(0).
\]
这又蕴涵着，必有某一正整数 $n$，使得 $2\sqrt{(k+1)k}\tau=n\tau$。由此可得 $\left(\dfrac n2\right)^2=k(k+1)$。这一等式表明，$\dfrac n2$ 是正整数，而且满足 $k<\dfrac n2<k+1$。这显然是不可能的。这一矛盾表明，$g(x)$ 不可能是周期函数.
\end{proof}


\subsection{函数的奇偶性}
\label{subsec:ma-1-3-4}

设 $y=f(x)$ 是定义在 $X$ 上的一个函数，而且 $X$ 是关于原点对称的，即 $x\in X$ 蕴涵着 $-x\in X$。若 $f(x)=-f(-x)$ 对一切的 $x\in X$ 成立，则称 $f(x)$ 是 $X$ 上的\keyterm{奇函数}；若 $f(x)=f(-x)$ 对一切的 $x\in X$ 成立，则称 $f(x)$ 是 $X$ 上的\keyterm{偶函数}.

容易看出，偶函数的图像是关于 $y$ 轴对称的，而奇函数的图像是关于原点对称的.

\begin{example}\label{example:ma-1-3-8}
设 $y=f(x)$ 是 $X$ 上的奇函数且存在反函数。证明：它的反函数也是奇函数.
\end{example}


\begin{proof}
由奇函数的定义知，对 $\forall x\in X$，有 $f(-x)=-f(x)$。于是
\[
  -f^{-1}(f(x))=-x=f^{-1}(f(-x))=f^{-1}(-f(x)),\qquad \forall x\in X,
\]
即 $\forall y\in f(X)$，有 $-f^{-1}(y)=f^{-1}(-y)$。这就证明了 $f^{-1}$ 也是奇函数.
\end{proof}


\begin{example}\label{example:ma-1-3-9}
证明 $f(x)=\ln(x+\sqrt{1+x^2})$ 是奇函数.
\end{example}


\begin{proof}
由于 $f(x)$ 是奇函数 $y=\sinh x$ 的反函数，故由上例知，该函数是奇函数.
\end{proof}


\section{初等函数}
\label{sec:ma-1-4}

在 第~\ref{sec:ma-1-2}~节 中已经指出，常值函数、幂函数、指数函数、对数函数、三角函数和反三角函数统称为基本初等函数。从这些基本初等函数出发，经过有限多次加、减、乘、除和复合运算所能得到的所有函数统称为\keyterm{初等函数}。对基本初等函数及其基本性质，大家已经十分熟悉，这里我们就不再赘述。下面我们只简要地介绍一下双曲函数，这是一类十分有用的初等函数，在以后的学习中常常用到它们.

双曲函数定义如下：
\[
  \sinh x=\frac{e^x-e^{-x}}{2},\qquad
  \cosh x=\frac{e^x+e^{-x}}{2},
\]
\[
  \tanh x=\frac{e^x-e^{-x}}{e^x+e^{-x}},\qquad
  \coth x=\frac{e^x+e^{-x}}{e^x-e^{-x}}.
\]
我们把这四个函数分别称之为\keyterm{双曲正弦}、\keyterm{双曲余弦}、\keyterm{双曲正切}和\keyterm{双曲余切}。由定义不难看出，$\sinh x,\tanh x,\coth x$ 都是奇函数，而 $\cosh x$ 是偶函数 (参见图 1.4.1 (a),(b)).

由双曲函数定义出发，可以导出许多类似于三角函数的恒等式。例如，有
\[
  \sinh 2x=2\sinh x\cosh x,\qquad
  \cosh^2x-\sinh^2x=1,
\]
\[
  \cosh^2x+\sinh^2x=\cosh 2x,\qquad
  \cosh x=1+2\sinh^2\frac x2.
\]
这些恒等式正好对应于如下的三角函数恒等式：
\[
  \sin 2x=2\sin x\cos x,\qquad
  \cos^2x+\sin^2x=1,
\]
\[
  \cos^2x-\sin^2x=\cos 2x,\qquad
  \cos x=1-2\sin^2\frac x2.
\]
令 $X=\cosh x,\ Y=\sinh x$，则它们满足双曲方程
\[
  X^2-Y^2=1.
\]
这是双曲函数名称由来的原因之一。在单位圆盘上的非欧几何——双曲几何 (俄国人称之为罗巴切夫斯基几何，西方人称之为庞加莱 (Poincaré) 几何) 理论中，双曲函数是基本的研究工具.

\section*{习题一}
\phantomsection
\addcontentsline{toc}{section}{习题一}
\label{exercises:ma-1}

\begin{enumerate}
  \item 设集合 $A=\{1,2,4,8\},\ B=\{1,3,6,9\}$，试求 $A\cup B,\ A\cap B$ 和 $A-B$.
  \item 设集合 $A=\{2n-1:n\in\mathbb{N}\},\ B=\{2n:n\in\mathbb{N}\}$，试求 $A\cup B,\ A\cap B,\ A-B$ 和 $B-A$.
  \item $\forall n\in\mathbb{N}$，设 $A_n=\left[-1+\dfrac{1}{2n},1-\dfrac1n\right]$，试求 $\displaystyle\bigcup_{n=1}^{+\infty}A_n$.
\end{enumerate}

% 整合来源：math-2.tex
\chapter{序列的极限}
\label{chap:ma-2}

极限是构筑微积分坚实理论体系的基石。要想对数学分析这门学科的实质有一个真正的了解和掌握, 就必须准确掌握极限的概念和无穷小的分析方法。这一章我们将致力于序列极限的讨论, 从中建立极限的主要理论与方法。下一章我们来讨论函数的极限.

\section{序列极限的定义}
\label{sec:ma-2-1}

\subsection{序列}
\label{subsec:ma-2-1-1}

所谓的\keyterm{序列} (有时也称为\keyterm{数列}) 实质上就是一个从正整数集 $\mathbb{N}$ 到实数集 $\mathbb{R}$ 的一个函数 $f:\mathbb{N}\to\mathbb{R}$。但我们常将一个序列看做是按照一定顺序排列的一列数:
\[
  x_1=f(1),\ x_2=f(2),\ \cdots,\ x_n=f(n),\ \cdots .
\]
一个序列通常记为 $\{x_n\}$, 其中 $x_n$ 称为\keyterm{通项}。有时为了简单起见, 我们也将构成一个序列的所有数作成的集合记为 $\{x_n\}$, 即此时 $\{x_n\}=f(\mathbb{N})$。读者以后应注意``集合 $\{x_n\}$''和``序列 $\{x_n\}$''的区别.

例如, 序列
\[
  1,\ \frac12,\ \frac13,\ \cdots,\ \frac1n,\ \cdots
\]
的通项是 $\dfrac1n$。但在有的时候序列的通项却很难用公式给出, 如圆周率 $\pi$ 精确到小数点后 $n$ 位的有限小数组成的序列
\[
  3,\ 3.1,\ 3.14,\ 3.141,\ 3.1415,\ \cdots
\]
就不可能写出通项公式.

我们已经知道在任意两个实数之间必有有理数, 有理数集的这一性质称为有理数集在实数集中的稠密性。使人吃惊的是, 在数轴上密密麻麻分布的有理数的全体竟然可以排成一个序列。下面我们仅排列 $[0,1)$ 中的有理数, $\mathbb{R}$ 中有理数的排列留给读者作为练习.

由于 $(0,1)$ 中的任何一个有理数都可唯一地表示为 $\dfrac pq$, 其中 $p$ 和 $q$ 是两个互素的正整数且 $p<q$, 因此我们可以将 $[0,1)$ 中的所有有理数按照其分母的大小从小到大排列, 而相同分母的再按其分子的大小从小到大排列。这样, 我们就可得到序列:
\[
  0,\ \frac12,\ \frac13,\ \frac23,\ \frac14,\ \frac34,\ \frac15,\ \frac25,\ \frac35,\ \cdots .
\]
对此序列, 似乎也很难写出它的通项公式.

\subsection{序列极限的定义}
\label{subsec:ma-2-1-2}

对于给定的一个序列 $\{x_n\}$, 我们这里主要关心的是随着 $n$ 的增大, $x_n$ 的变化趋势。为此, 我们首先必须说清楚``变化趋势''是什么意思。下面我们先看几个具体的例子.

\begin{example}\label{example:ma-2-1-1}
序列 $x_n=\dfrac1n(n=1,2,\cdots)$ 随着 $n$ 的增大, $x_n$ 从 $0$ 的右边越来越趋近于 $0$.
\end{example}


\begin{example}\label{example:ma-2-1-2}
序列 $x_n=\dfrac{(-1)^n}{n}(n=1,2,\cdots)$ 随着 $n$ 的增大, $x_n$ 时而在 $0$ 点的右边, 时而在 $0$ 点的左边, 但 $x_n$ 与 $0$ 的距离越来越趋近于 $0$.
\end{example}


\begin{example}\label{example:ma-2-1-3}
序列 $\{x_n\}$ 由下列法则确定:
\[
  x_{2n}=\frac1{2n},\qquad
  x_{2n+1}=\frac1{2^{2n+1}}
  \quad(n=1,2,\cdots).
\]
对于这一序列, 虽然我们不能说, 随着 $n$ 的增大, $x_n$ 与 $0$ 的距离越来越趋近于 $0$, 但 $x_n$ 与 $0$ 的距离也是无限地接近于 $0$, 只不过奇数项与偶数项接近零的速度不同而已.
\end{example}


\begin{example}\label{example:ma-2-1-4}
设 $\{x_n\}$ 由下式给出:
\[
  x_{2n-1}=\frac1n,\qquad
  x_{2n}=0
  \quad(n=1,2,\cdots),
\]
则 $\{x_n\}$ 有无穷多项为 $0$, 而随着 $n$ 的增大, $x_{2n-1}$ 也无限地趋于 $0$。因此, 随着 $n$ 的增大, 整个序列也是趋近于 $0$ 的.
\end{example}

分析以上例子我们发现, 当 $n$ 无限变大时, 尽管它们的变化过程各具特点, 但四个序列最终都无限地趋于一个常数 $a=0$。这种随着 $n$ 的增大可以无限地趋于一个常数的序列正是我们感兴趣的。那么怎样来精确地描述一个序列 $\{x_n\}$ 无限地趋于一个常数 $a$ 呢? 也许大家会说, 这是指 $x_n$ 与 $a$ 的差 (随着 $n$ 的增大) 可以达到任意小的程度。显然, 这只是一种描述性的说法, 十分不宜于进行严格的数学推理和演绎。其实, 给出某精确的定义并非一件易事, 经过众多数学家的不懈努力和不断探索, 直到 19 世纪才有了数学上的如下定义:

\begin{definition}\label{definition:ma-2-1-1}
设 $\{x_n\}$ 是一个序列。若存在常数 $a\in\mathbb{R}$, 使得 $\forall\varepsilon>0$, $\exists N\in\mathbb{N}$, 当 $n>N$ 时, 有
\[
  |x_n-a|<\varepsilon,
\]
则称该序列是\keyterm{收敛的}, 并称 $a$ 为该序列的\keyterm{极限} (或者说序列 $\{x_n\}$ 收敛于 $a$), 记做 $\lim\limits_{n\to\infty}x_n=a$ 或 $x_n\to a(n\to\infty)$。若不存在 $a\in\mathbb{R}$, 使得 $\{x_n\}$ 收敛于 $a$, 则称之为\keyterm{发散序列}.
\end{definition}

在上述定义中, $\varepsilon$ 是事先给定的任意小的正数, 若对 $\varepsilon_1>0$, 我们已经找到了正整数 $N_1$, 使得当 $n>N_1$ 时, 有 $|x_n-a|<\varepsilon_1$, 则对 $\forall\varepsilon_2>\varepsilon_1$, 当 $n>N_1$ 时, 必有 $|x_n-a|<\varepsilon_1<\varepsilon_2$。所以, $\varepsilon$ 贵在``小''。此外, 一般来说, $N$ 是依赖于 $\varepsilon$ 的, $\varepsilon$ 越小, 所需的 $N$ 就会越大。对一个给定的正数 $\varepsilon$, 如果 $N$ 是满足定义中的要求的正整数, 则显然对任何比 $N$ 大的正整数也满足定义的要求.

现在我们来考查序列极限的几何意义。极限定义中的不等式
\[
  |x_n-a|<\varepsilon,\qquad \forall n>N,
\]
可以写成
\[
  a-\varepsilon<x_n<a+\varepsilon,\qquad \forall n>N,
\]
即
\[
  x_n\in U(a,\varepsilon)=(a-\varepsilon,a+\varepsilon),\qquad \forall n>N.
\]
因此, 如果采用几何的语言, 极限的定义可以表述为: $\forall\varepsilon>0$, 在 $a$ 的 $\varepsilon$ 邻域 $U(a,\varepsilon)$ 内包含了 $\{x_n\}$ 自某项之后的所有项。与之等价说法是: $\forall\varepsilon>0$, 在 $a$ 的 $\varepsilon$ 邻域 $U(a,\varepsilon)$ 外只有 $\{x_n\}$ 的有限项。如图 2.1.1 所示.

\paragraph*{思考题}
将序列 $\{x_n\}$ 看成是定义在正整数集上的函数, 试描述序列极限的几何意义.

\begin{example}\label{example:ma-2-1-5}
证明 $\lim\limits_{n\to\infty}\dfrac{n}{n+1}=1$.
\end{example}


\begin{proof}
对任意给定的 $\varepsilon>0$, 要使
\[
  \left|\frac{n}{n+1}-1\right|
  =
  \frac1{n+1}<\varepsilon,
\]
只需
\[
  n>\frac1\varepsilon-1.
\]
于是, 取大于 $\dfrac1\varepsilon-1$ 的任意正整数作为 $N$ (例如, 取 $N=\left[\dfrac1\varepsilon\right]$), 则当 $n>N$ 时, 就有
\[
  \left|\frac{n}{n+1}-1\right|
  =
  \frac1{n+1}<\varepsilon.
\]
由极限定义知, $\lim\limits_{n\to\infty}\dfrac{n}{n+1}=1$ 成立.
\end{proof}


\begin{example}\label{example:ma-2-1-6}
证明 $\lim\limits_{n\to\infty}q^n=0\ (|q|<1)$.
\end{example}


\begin{proof}
不妨设 $q\ne0$, 否则该序列为常序列 $(x_n=0,\ n=1,2,\cdots)$, 所述极限自然成立.

$\forall\varepsilon>0$ (不妨设 $\varepsilon<1$), 要使
\[
  |q^n-0|=|q|^n<\varepsilon,
\]
只要
\[
  n\ln |q|<\ln\varepsilon.
\]
注意, 这里 $\ln |q|<0,\ \ln\varepsilon<0$, 于是上式等价于
\[
  n>\frac{\ln\varepsilon}{\ln |q|}.
\]
因此, 取 $N=\left[\dfrac{\ln\varepsilon}{\ln |q|}\right]$, 则当 $n>N$ 时, 就有 $|q^n-0|<\varepsilon$, 即 $\lim\limits_{n\to\infty}q^n=0$ 成立.
\end{proof}


\begin{example}\label{example:ma-2-1-7}
证明 $\lim\limits_{n\to\infty}\dfrac{n^2-n+2}{3n^2+2n+4}=\dfrac13$.
\end{example}


\begin{proof}
对任意的正整数 $n$, 我们有
\[
  \left|\frac{n^2-n+2}{3n^2+2n+4}-\frac13\right|
  =
  \frac{5n-2}{3(3n^2+2n+4)}
  <
  \frac{5n}{9n^2}
  <
  \frac1n.
\]
于是, $\forall\varepsilon>0$, 只要取 $N=\left[\dfrac1\varepsilon\right]+1$, 则当 $n>N$ 时, 就有
\[
  \left|\frac{n^2-n+2}{3n^2+2n+4}-\frac13\right|
  <
  \frac1n<\varepsilon.
\]
这就证明了 $\lim\limits_{n\to\infty}\dfrac{n^2-n+2}{3n^2+2n+4}=\dfrac13$.
\end{proof}


\begin{example}\label{example:ma-2-1-8}
证明 $\lim\limits_{n\to\infty}\sqrt[n]{a}=1(a>1)$.
\end{example}


\begin{proof}[{证法 1}]
$\forall\varepsilon>0$, 要使
\[
  |\sqrt[n]{a}-1|=\sqrt[n]{a}-1<\varepsilon,
\]
只要 $\sqrt[n]{a}<1+\varepsilon$, 即 $\dfrac1n\lg a<\lg(1+\varepsilon)$, 这只要 $n>\dfrac{\lg a}{\lg(1+\varepsilon)}$。于是, 取 $N=\left[\dfrac{\lg a}{\lg(1+\varepsilon)}\right]$, 则当 $n>N$ 时, 就有 $|\sqrt[n]{a}-1|<\varepsilon$, 即
\[
  \lim_{n\to\infty}\sqrt[n]{a}=1.
\qedhere
\]
\end{proof}


\begin{proof}[{证法 2}]
令 $\sqrt[n]{a}-1=h_n$, 则 $h_n>0$, 而且
\[
  a=(1+h_n)^n=1+nh_n+\cdots+(h_n)^n>nh_n,
\]
从而
\[
  0<h_n<\frac an,\qquad \forall n\in\mathbb{N}.
\]
这样, $\forall\varepsilon>0$, 要使 $|\sqrt[n]{a}-1|=h_n<\varepsilon$, 只要 $\dfrac an<\varepsilon$。现在取 $N=\left[\dfrac a\varepsilon\right]$, 则只要 $n>N$, 就有 $|\sqrt[n]{a}-1|<\varepsilon$, 即 $\lim\limits_{n\to\infty}\sqrt[n]{a}=1$.

用定义来证明极限的过程, 实际上就是对于相对确定的正数 $\varepsilon$ 去找相应的正整数 $N$ 的过程。换句话说就是, 你任给我一个正数 $\varepsilon$, 我去给你找一个相应的正整数 $N$, 在找 $N$ 的过程中, $\varepsilon$ 是相对固定的。而找 $N$ 的过程实质就是解不等式的过程。因此, 为了使所解的不等式尽可能的简单, 我们可以将所估计量做适当的放大 (参见例~\ref{example:ma-2-1-7}), 但不能放得太大, 要保证从中能够找到所需要的正整数 $N$.

下面我们来讨论一下发散序列。如何用肯定的语气来表述一个序列是发散的呢? 从定义知, 一个序列 $\{x_n\}$ 是发散的等价于任何数 $a\in\mathbb{R}$ 都不是它的极限。这样, 上述问题就转化为怎样用肯定的语气来表述 $\{x_n\}$ 不收敛于 $a$。大家已经知道, $\{x_n\}$ 收敛于 $a$ 意味着 $a$ 的任意小的邻域的外面只有序列中的有限多项。因此, 若 $\{x_n\}$ 不收敛于 $a$, 则存在 $a$ 的一个邻域 $U(a,\varepsilon_0)$, 使得在它之外必有 $\{x_n\}$ 的无穷多项。于是, 若用肯定的语气, 发散序列可以叙述为:

设 $\{x_n\}$ 是一个序列。若 $\forall a\in\mathbb{R}$, $\exists\varepsilon_0>0$, 对 $\forall N\in\mathbb{N}$, 总存在 $n_0>N$, 使得 $|x_{n_0}-a|>\varepsilon_0$, 则称 $\{x_n\}$ 为\keyterm{发散序列}.

仔细观察上述的表述, 不难发现, 在发散序列的定义中, 只是将极限定义中的``$\forall$''换成了``$\exists$'', ``$\exists$''换成了``$\forall$''。尽管这样, 我们还是希望读者能理解这样表述的实质, 而不是形式地加以记忆.
\end{proof}


\begin{example}\label{example:ma-2-1-9}
证明如下定义的序列是发散的:
\[
  x_n=
  \begin{cases}
    \dfrac1n, & n=2k-1,\\
    n, & n=2k,
  \end{cases}
  \qquad k=1,2,\cdots .
\]
\end{example}


\begin{proof}
对任意给定的 $a\in\mathbb{R}$, 取 $\varepsilon_0=1$, 则对于任给的正整数 $N$, 选取 $n_0=2k_0>\max\{N,a+1\}$, 其中 $k_0$ 为正整数, 便有 $n_0>N$, 而且 $x_{n_0}-a=2k_0-a>1=\varepsilon_0$, 即 $a$ 不是 $\{x_n\}$ 的极限。由 $a$ 的任意性知序列 $\{x_n\}$ 发散.
\end{proof}


\subsection{无穷小量}
\label{subsec:ma-2-1-3}

作为序列极限的特例, 我们引入无穷小量的概念.

\begin{definition}\label{definition:ma-2-1-2}
设 $\{x_n\}$ 是一个序列。若 $x_n\to0(n\to\infty)$, 则称序列 $\{x_n\}$ 为\keyterm{无穷小量}, 记为 $x_n=o(1)(n\to\infty)$.

这里需特别指出的是, 无穷小量并不是一个很小的量, 而是极限为零的一个变量.
\end{definition}


\begin{example}\label{example:ma-2-1-10}
证明序列 $\left\{\dfrac{a^n}{n!}\right\}$ 是无穷小量, 这里 $a>1$.
\end{example}


\begin{proof}
注意到
\[
  \frac{a^n}{n!}
  =
  \frac a1\cdot\frac a2\cdots\frac a{[a]}\cdot
  \frac a{[a]+1}\cdots\frac an
  < C\frac an
  \quad(n>[a]),
\]
其中 $C=\dfrac{a^{[a]}}{[a]!}$ 是常数, $\forall\varepsilon>0$, 要使
\[
  \left|\frac{a^n}{n!}-0\right|<\varepsilon,
\]
只需
\[
  n>\frac{Ca}{\varepsilon}.
\]
于是, 取 $N=\left[\dfrac{Ca}{\varepsilon}\right]+1$, 则只要 $n>N$, 就有
\[
  \left|\frac{a^n}{n!}-0\right|<\varepsilon,
\]
即 $\lim\limits_{n\to\infty}\dfrac{a^n}{n!}=0$, 从而序列 $\left\{\dfrac{a^n}{n!}\right\}$ 是无穷小量.
\end{proof}


\begin{example}\label{example:ma-2-1-11}
证明序列 $\{x_n\}$ 是无穷小量, 其中
\[
  x_n=
  \begin{cases}
    \dfrac1{2^n}, & n=2k-1,\\[0.4em]
    \dfrac1{\sqrt n}, & n=2k,
  \end{cases}
  \qquad n=1,2,\cdots .
\]
\end{example}


\begin{proof}
因为对 $\forall n\in\mathbb{N}$, 有
\[
  \frac1{2^n}<\frac1{\sqrt n},
\]
所以, $\forall\varepsilon>0$, 取 $N=\left[\dfrac1{\varepsilon^2}\right]$, 则当 $n>N$ 时, 有
\[
  |x_n-0|\leqslant\frac1{\sqrt n}<\varepsilon,
\]
即 $\lim\limits_{n\to\infty}x_n=0$。这就证明了所述序列 $\{x_n\}$ 是无穷小量.

利用极限的定义容易证明无穷小量有如下的基本性质:
\end{proof}


\begin{theorem}\label{theorem:ma-2-1-1}
设 $\{x_n\}$ 是一个序列.

(1) $\{x_n\}$ 是无穷小量的充分必要条件是 $\{|x_n|\}$ 是无穷小量;

(2) 若 $\{x_n\}$ 是无穷小量, $M$ 是一个常数, 则 $\{Mx_n\}$ 是无穷小量;

(3) $\lim\limits_{n\to\infty}x_n=a$ 的充分必要条件是 $\{x_n-a\}$ 是无穷小量.
\end{theorem}

证明留作练习.

\subsection{无穷大量}
\label{subsec:ma-2-1-4}

对于序列 $x_n=(-1)^n$ 和 $x_n=(-1)^nn\ (n=1,2,\cdots)$ 来说, 虽然它们都不存在极限, 但这两者有着本质的区别: 前者中的 $x_n$ 随着 $n$ 的增加在 $-1$ 和 $1$ 这两点上跳来跳去; 而后者, 则随着 $n$ 的增加, 其绝对值目标一致地趋向 $+\infty$。因此我们有如下定义:

\begin{definition}\label{definition:ma-2-1-3}
设 $\{x_n\}$ 是一个序列。若 $\forall M>0$, $\exists N$, 当 $n>N$ 时, 有 $x_n>M$, 则称 $\{x_n\}$ 为\keyterm{正无穷大量} (有时也称 $\{x_n\}$ 的极限为 $+\infty$, 记为 $\lim\limits_{n\to\infty}x_n=+\infty$); 若 $\forall M>0$, $\exists N$, 当 $n>N$ 时, 有 $x_n<-M$, 则称 $\{x_n\}$ 为\keyterm{负无穷大量} (有时也称 $\{x_n\}$ 的极限为 $-\infty$, 记为 $\lim\limits_{n\to\infty}x_n=-\infty$); 若 $\{|x_n|\}$ 是正无穷大量, 则称 $\{x_n\}$ 为\keyterm{无穷大量} (有时也称 $\{x_n\}$ 的极限为 $\infty$, 记为 $\lim\limits_{n\to\infty}x_n=\infty$).
\end{definition}

显然, 正无穷大量和负无穷大量都是无穷大量。具有有穷极限和无穷极限的序列有着本质的区别, 必须加以区别。因此, 当序列有有穷极限 $a$ 时, 我们说它\keyterm{收敛于} $a$; 当序列有无穷极限 $a$ 时, 我们说它\keyterm{发散到} $+\infty$, $-\infty$ 或 $\infty$。在本书的后续部分, 如果没有特别说明, 凡提到极限存在, 均指有穷极限的情形; 若包括无穷极限时, 则说其\keyterm{广义极限存在}, 此时也说序列是\keyterm{广义收敛的}.

\begin{example}\label{example:ma-2-1-12}
设
\[
  x_n=1+\frac12+\frac13+\cdots+\frac1n\quad(n=1,2,\cdots),
\]
证明 $\lim\limits_{n\to\infty}x_n=+\infty$.
\end{example}


\begin{proof}
注意到, 当 $n>2^k$ 时, 有
\[
\begin{aligned}
  x_n
  &=1+\frac12+\frac13+\frac14+\cdots+\frac1{2^{k-1}+1}
    +\cdots+\frac1{2^k}+\cdots+\frac1n\\
  &>1+\frac12+\frac14+\frac14+\frac18+\frac18+\frac18+\frac18
    +\cdots+\frac1{2^k}+\frac1{2^k}+\cdots+\frac1{2^k}\\
  &=1+\frac12+\frac12+\frac12+\cdots+\frac12\\
  &>\frac k2.
\end{aligned}
\]
故 $\forall M>0$, 欲使
\[
  x_n>M,
\]
只需
\[
  k>2M.
\]
这样, 只要取 $N=2^{[2M]+1}$, 则当 $n>N$ 时, 有
\[
  x_n>\frac{[2M]+1}{2}>\frac{2M}{2}=M,
\]
从而由极限定义知 $\lim\limits_{n\to\infty}x_n=+\infty$ 成立.

无穷大量和无穷小量之间有如下的关系:
\end{proof}


\begin{theorem}\label{theorem:ma-2-1-2}
$\{x_n\}$ 是无穷小量的充分必要条件是 $\left\{\dfrac1{x_n}\right\}$ 是无穷大量, 这里假定对任意的正整数 $n$, 有 $x_n\ne0$.
\end{theorem}


\begin{proof}
先证必要性。设 $\lim\limits_{n\to\infty}x_n=0$, 则 $\forall M>0$ (即 $\varepsilon=\dfrac1M>0$), $\exists N$, 当 $n>N$ 时, 有
\[
  |x_n|<\frac1M,
\]
从而有
\[
  \left|\frac1{x_n}\right|>M,
\]
即 $\lim\limits_{n\to\infty}\dfrac1{x_n}=\infty$.

再证充分性。设 $\lim\limits_{n\to\infty}\dfrac1{x_n}=\infty$, 则 $\forall\varepsilon>0$ (即 $M=\dfrac1\varepsilon>0$), $\exists N$, 当 $n>N$ 时, 有
\[
  \left|\frac1{x_n}\right|>\frac1\varepsilon,
\]
从而有
\[
  |x_n|<\varepsilon,
\]
即 $\lim\limits_{n\to\infty}x_n=0$.

作为本节的结束, 我们再举一例.
\end{proof}


\begin{example}\label{example:ma-2-1-13}
设 $\lim\limits_{n\to\infty}x_n=a$, 试证 $\lim\limits_{n\to\infty}\dfrac{x_1+\cdots+x_n}{n}=a$, 这里 $a$ 可以是有限实数, $+\infty$ 或 $-\infty$.
\end{example}


\begin{proof}
先证 $a$ 是有限实数的情形。$\forall\varepsilon>0$, 由于 $\lim\limits_{n\to\infty}x_n=a$, 故 $\exists N_1$, 使当 $n>N_1$ 时, 有
\[
  |x_n-a|<\frac\varepsilon2.
\]
对于取定的正整数 $N_1$, 由于
\[
  |x_1-a|+|x_2-a|+\cdots+|x_{N_1}-a|
\]
就是一个不随 $n$ 的变化而变化的常数, 因此有
\[
  \lim_{n\to\infty}
  \frac{|x_1-a|+|x_2-a|+\cdots+|x_{N_1}-a|}{n}=0.
\]
于是, 又存在正整数 $N_2>0$, 当 $n>N_2$ 时, 有
\[
  \frac{|x_1-a|+|x_2-a|+\cdots+|x_{N_1}-a|}{n}<\frac\varepsilon2.
\]
现在取 $N=\max\{N_1,N_2\}$, 则当 $n>N$ 时, 有
\[
\begin{aligned}
  \left|\frac{x_1+\cdots+x_n}{n}-a\right|
  &\leqslant
  \frac{|x_1-a|+|x_2-a|+\cdots+|x_n-a|}{n}\\
  &=
  \frac{|x_1-a|+\cdots+|x_{N_1}-a|}{n}
  +
  \frac{|x_{N_1+1}-a|+\cdots+|x_n-a|}{n}\\
  &\leqslant \frac\varepsilon2+\frac{n-N_1}{2n}\varepsilon<\varepsilon,
\end{aligned}
\]
从而有 $\lim\limits_{n\to\infty}\dfrac{x_1+\cdots+x_n}{n}=a$.

下证 $a=+\infty$ 的情形。$\forall M>0$, 由 $\lim\limits_{n\to\infty}x_n=+\infty$ 知, $\exists N_1$, 当 $n>N_1$ 时, 有 $x_n>2M$, 且 $x_1+x_2+\cdots+x_{N_1}>0$, 从而有
\[
  \frac{x_1+\cdots+x_n}{n}
  >
  \frac{x_{N_1+1}+\cdots+x_n}{n}
  >
  \frac{2(n-N_1)}nM.
\]
现在取 $N=2N_1$, 则当 $n>N$ 时, 有
\[
  \frac{n-N_1}{n}=1-\frac{N_1}{n}>\frac12,
\]
从而有
\[
  \frac{x_1+\cdots+x_n}{n}
  >
  \frac{2(n-N_1)}nM>M.
\]
这就证明了 $\lim\limits_{n\to\infty}\dfrac{x_1+\cdots+x_n}{n}=+\infty$.

对 $a=-\infty$ 的情形, 证明是类似的, 请读者自己给出.
\end{proof}


\paragraph*{思考题}


(1) 若 $\lim\limits_{n\to\infty}x_n=\infty$, 是否有 $\lim\limits_{n\to\infty}\dfrac{x_1+\cdots+x_n}{n}=\infty$?

(2) 若 $\{x_n\}$ 发散, 是否有 $\left\{\dfrac{x_1+\cdots+x_n}{n}\right\}$ 亦发散?

\section{序列极限的性质}
\label{sec:ma-2-2}

在上一节中, 我们已经对极限有了初步的认识。对于一个序列, 我们要讨论的主要问题是:
\[
\begin{array}{ll}
  \text{(1)} & \text{它是否存在极限?}\\
  \text{(2)} & \text{如果极限存在, 如何来求出这一极限?}
\end{array}
\]
对一些简单的序列, 我们很容易看出它是否收敛, 甚至能看出它的极限是什么。但对一般的序列来说, 这是两个相当困难的问题。可以这么说, 整个数学分析自始至终都在讨论各种极限的存在性以及求法等问题。极限论可以看成是分析学与代数学的主要区别.

在这一节, 我们来进一步讨论序列极限的基本性质.

\begin{definition}\label{definition:ma-2-2-1}
设 $\{x_n\}$ 是一个序列。若 $\exists M>0$, 对 $\forall n$, 有 $|x_n|\leqslant M$ 成立, 则称 $\{x_n\}$ 是\keyterm{有界的}.
\end{definition}

显然, 以上定义等价于数集 $\{x_n\}$ 是一个有界集.

若一个序列 $\{x_n\}$ 是有界的, 则记为 $x_n=O(1)$ (读做大 $O$); 若存在 $M_2>M_1>0$ 和正整数 $N$, 使得当 $n>N$ 时, 有 $M_1<|x_n|<M_2$, 则以 $x_n=O_0(1)$ 表示之.

\begin{theorem}\label{theorem:ma-2-2-1}
(1) 改变一个序列 $\{x_n\}$ 的有限多项, 不改变其敛散性; 当 $\{x_n\}$ 收敛时, 则不改变其极限值.

(2) (\keyterm{唯一性}) 收敛序列的极限是唯一的.

(3) (\keyterm{有界性}) 收敛序列是有界的.
\end{theorem}


\begin{proof}
(1) 由极限定义立即可得.

(2) 用反证法证之。如果结论不真, 则存在收敛序列 $\{x_n\}$ 和实数 $b\ne a$, 使得 $\lim\limits_{n\to\infty}x_n=a$ 且 $\lim\limits_{n\to\infty}x_n=b$。不失一般性, 我们不妨设 $a<b$。现在取 $\varepsilon_0=\dfrac{b-a}{2}$, 则由极限定义知, 存在正整数 $N_1$ 和 $N_2$, 使得
\[
  |x_n-a|<\varepsilon_0,\qquad \forall n>N_1,
\]
\[
  |x_n-b|<\varepsilon_0,\qquad \forall n>N_2.
\]
令 $N=\max\{N_1,N_2\}$, 则当 $n>N$ 时, 上述两个不等式都成立。由第一个不等式, 得
\[
  x_{N+1}<a+\varepsilon_0=\frac{a+b}{2},
\]
而由第二个不等式, 得
\[
  \frac{a+b}{2}=b-\varepsilon_0<x_{N+1}.
\]
这显然是不可能的, 因此 (2) 成立.

(3) 设 $\lim\limits_{n\to\infty}x_n=a$。特别取 $\varepsilon_0=1$, 则 $\exists N$, 使得当 $n>N$ 时, 有 $|x_n-a|<1$, 从而
\[
  |x_n|\leqslant |a|+|x_n-a|<|a|+1,\qquad \forall n>N.
\]
令 $M=\max\{|a|+1,|x_1|,\cdots,|x_N|\}$, 则对 $\forall n$, 有 $|x_n|\leqslant M$, 即 $\{x_n\}$ 是有界的.
\end{proof}


\paragraph*{思考题}
试举例说明定理~\ref{theorem:ma-2-2-1}(3) 的逆命题不真.

\begin{theorem}[{保序性}]\label{theorem:ma-2-2-2}
给定两个序列 $\{x_n\}$ 和 $\{y_n\}$, 并且假定
\[
  \lim_{n\to\infty}x_n=a,\qquad
  \lim_{n\to\infty}y_n=b,
\]
则有:

(1) 若 $a<b$, 则对任意给定的 $c\in(a,b)$, $\exists N_0>0$, 使得当 $n>N_0$ 时, 有 $x_n<c<y_n$;

(2) 若 $\exists N_0>0$, 当 $n>N_0$ 时, 有 $x_n\leqslant y_n$, 则 $a\leqslant b$.
\end{theorem}


\begin{proof}
(1) 令 $\varepsilon_0=\min\{c-a,b-c\}$, 则由极限的定义知, $\exists N>0$, 使得当 $n>N$ 时, 有 $|x_n-a|<\varepsilon_0$ 且 $|y_n-b|<\varepsilon_0$ 成立。由此可得, 当 $n>N$ 时, 有
\[
  x_n<\varepsilon_0+a\leqslant(c-a)+a=c=b-(b-c)\leqslant b-\varepsilon_0<y_n.
\]

(2) 用反证法证之。如若不然, 则有 $a>b$。现取 $\varepsilon_0=\dfrac{a-b}{2}$, 由 (1) 所证知, $\exists N>0$, 使得当 $n>N$ 时, 有
\[
  y_n<\frac{a+b}{2}<x_n.
\]
这与已知条件 $x_n\leqslant y_n\ (n>N_0)$ 相矛盾, 故必有 $a\leqslant b$ 成立.
\end{proof}


\begin{corollary}\label{corollary:ma-2-source-35}
设 $\lim\limits_{n\to\infty}x_n=a>0$, 则对任给 $0<a'<a$, $\exists N_1>0$, 使得当 $n>N_1$ 时, 有 $x_n>a'$; 若 $\lim\limits_{n\to\infty}x_n=b<0$, 则对任给 $0>b'>b$, $\exists N_2>0$, 使得当 $n>N_2$ 时, 有 $x_n<b'$.
\end{corollary}


\begin{remark}\label{remark:ma-2-source-36}
注意, 在定理~\ref{theorem:ma-2-2-2} 的 (2) 中, 即使对一切的 $n$ 有 $x_n<y_n$, 也未必有 $a<b$。例如,
\[
  x_n=1-\frac1n<1+\frac1n=y_n,\qquad n=1,2,\cdots,
\]
而
\[
  \lim_{n\to\infty}x_n=\lim_{n\to\infty}y_n=1.
\]
\end{remark}


\begin{theorem}[{极限的四则运算}]\label{theorem:ma-2-2-3}
设 $\lim\limits_{n\to\infty}x_n=a$, $\lim\limits_{n\to\infty}y_n=b$, 则

(1) $\lim\limits_{n\to\infty}(x_n\pm y_n)=a\pm b$;

(2) $\lim\limits_{n\to\infty}(x_ny_n)=ab$;

(3) $\lim\limits_{n\to\infty}\dfrac{x_n}{y_n}=\dfrac ab$, 其中 $b\ne0,\ y_n\ne0$.
\end{theorem}


\begin{proof}
(1) $\forall\varepsilon>0$, 由 $\lim\limits_{n\to\infty}x_n=a$ 知, $\exists N_1>0$, 使得当 $n>N_1$ 时, 有
\[
  |x_n-a|<\frac\varepsilon2;
\]
又由 $\lim\limits_{n\to\infty}y_n=b$ 知, $\exists N_2>0$, 使得当 $n>N_2$ 时, 有
\[
  |y_n-b|<\frac\varepsilon2.
\]
取 $N=\max\{N_1,N_2\}$, 则当 $n>N$ 时, 有
\[
  |(x_n\pm y_n)-(a\pm b)|
  \leqslant |x_n-a|+|y_n-b|
  <\frac\varepsilon2+\frac\varepsilon2=\varepsilon,
\]
即 $\lim\limits_{n\to\infty}(x_n\pm y_n)=a\pm b$.

(2) 首先, 我们有
\[
\begin{aligned}
  |x_ny_n-ab|
  &=|x_ny_n-y_na+y_na-ab|\\
  &\leqslant |y_n||x_n-a|+|a||y_n-b|.
\end{aligned}
\]
(这里, 为了与已知的极限发生联系, 我们使用了加一项、减一项的插项方法, 它是一个常用的方法).

由有界性定理知, $\exists M_1>0$, $\forall n$, $|y_n|\leqslant M_1$。令
\[
  M=\max\{M_1,|a|\}>0.
\]
$\forall\varepsilon>0$, 由 $\lim\limits_{n\to\infty}x_n=a$ 知, $\exists N_1>0$, 使得当 $n>N_1$ 时, 有
\[
  |x_n-a|<\frac{\varepsilon}{2M};
\]
又由 $\lim\limits_{n\to\infty}y_n=b$ 知, $\exists N_2>0$, 使得当 $n>N_2$ 时, 有
\[
  |y_n-b|<\frac{\varepsilon}{2M}.
\]
取 $N=\max\{N_1,N_2\}$, 则当 $n>N$ 时, 有
\[
\begin{aligned}
  |x_ny_n-ab|
  &\leqslant |y_n||x_n-a|+|a||y_n-b|\\
  &\leqslant M(|x_n-a|+|y_n-b|)\\
  &\leqslant M\left(\frac{\varepsilon}{2M}+\frac{\varepsilon}{2M}\right)\\
  &=\varepsilon,
\end{aligned}
\]
即 $\lim\limits_{n\to\infty}(x_ny_n)=ab$.

(3) 由 (2) 成立, 要证 (3), 只需证 $\lim\limits_{n\to\infty}\dfrac1{y_n}=\dfrac1b$ 即可。首先, 我们有
\[
  \left|\frac1{y_n}-\frac1b\right|
  =
  \left|\frac{y_n-b}{y_nb}\right|.
\]
然后, 由 $\lim\limits_{n\to\infty}y_n=b\ne0$, 从定理~\ref{theorem:ma-2-2-2} 的推论知, $\exists N_1>0$, 使得当 $n>N_1$ 时, 有
\[
  |y_n|>\frac{|b|}{2}.
\]
这样, $\forall\varepsilon>0$, 由 $\lim\limits_{n\to\infty}y_n=b$ 知, $\exists N_2>0$, 使得当 $n>N_2$ 时, 有
\[
  |y_n-b|<\frac{|b|^2}{2}\varepsilon.
\]
取 $N=\max\{N_1,N_2\}$, 则当 $n>N$ 时, 有
\[
  \left|\frac1{y_n}-\frac1b\right|
  =
  \left|\frac{y_n-b}{y_nb}\right|
  <
  \frac2{|b|^2}\frac{|b|^2}{2}\varepsilon
  =
  \varepsilon,
\]
即 $\lim\limits_{n\to\infty}\dfrac1{y_n}=\dfrac1b$.
\end{proof}


\begin{example}\label{example:ma-2-2-1}
求下列极限:

(1) $\lim\limits_{n\to\infty}\dfrac{3n^2+4n-100}{4n^2+5n+10^5}$;

(2) $\lim\limits_{n\to\infty}(1+q+q^2+\cdots+q^n)\ (|q|<1)$.
\end{example}


\begin{solve}
(1) 由于
\[
  \lim_{n\to\infty}\frac{3n^2+4n-100}{4n^2+5n+10^5}
  =
  \lim_{n\to\infty}
  \frac{3+\dfrac4n-\dfrac{100}{n^2}}
       {4+\dfrac5n+\dfrac{10^5}{n^2}},
\]
而
\[
  \lim_{n\to\infty}\left(3+\frac4n-\frac{100}{n^2}\right)=3,
\]
\[
  \lim_{n\to\infty}\left(4+\frac5n+\frac{10^5}{n^2}\right)=4,
\]
故有
\[
  \lim_{n\to\infty}\frac{3n^2+4n-100}{4n^2+5n+10^5}=\frac34.
\]

(2) 由于
\[
  1+q+q^2+\cdots+q^n=\frac{1-q^{n+1}}{1-q},
\]
而 $|q|<1$ 蕴涵着 $\lim\limits_{n\to\infty}q^n=0$, 于是
\[
  \lim_{n\to\infty}(1+q+q^2+\cdots+q^n)
  =
  \lim_{n\to\infty}\frac{1-q^{n+1}}{1-q}
  =
  \frac1{1-q}.
\qedhere
\]
\end{solve}


\begin{theorem}[{夹逼收敛原理}]\label{theorem:ma-2-2-4}
设序列 $\{x_n\}$, $\{y_n\}$ 和 $\{z_n\}$ 满足
\[
  x_n\leqslant z_n\leqslant y_n,\qquad \forall n>N_0.
\]
若 $\lim\limits_{n\to\infty}x_n=\lim\limits_{n\to\infty}y_n=a$, 则 $\lim\limits_{n\to\infty}z_n=a$.
\end{theorem}


\begin{proof}
$\forall\varepsilon>0$, 由 $\lim\limits_{n\to\infty}x_n=\lim\limits_{n\to\infty}y_n=a$ 知, 存在正整数 $N_1$ 和 $N_2$, 使得
\[
  |x_n-a|<\varepsilon,\qquad \forall n>N_1,
\]
\[
  |y_n-a|<\varepsilon,\qquad \forall n>N_2.
\]
令 $N=\max\{N_2,N_1,N_0\}$, 则当 $n>N$ 时, 有
\[
  a-\varepsilon<x_n\leqslant z_n\leqslant y_n<a+\varepsilon,
\]
即
\[
  |z_n-a|<\varepsilon.
\]
这表明 $\lim\limits_{n\to\infty}z_n=a$.
\end{proof}


\begin{example}\label{example:ma-2-2-2}
求下列极限:

(1) $\lim\limits_{n\to\infty}\sqrt[n]{n}$;

(2) $\lim\limits_{n\to\infty}(a_1^n+a_2^n+\cdots+a_m^n)^{\frac1n}\ (a_i>0,\ i=1,2,\cdots,m)$.
\end{example}


\begin{solve}
(1) 令 $\sqrt[n]{n}=1+h_n$, 则 $h_n>0$, 而且
\[
  n=(1+h_n)^n
  =
  1+nh_n+\frac{n(n-1)}2h_n^2+\cdots+h_n^n
  >
  \frac{n(n-1)}2h_n^2\quad(n>3).
\]
于是
\[
  0<h_n<\frac{\sqrt2}{\sqrt{n-1}}.
\]
由于 $\lim\limits_{n\to\infty}\dfrac{\sqrt2}{\sqrt{n-1}}=0$, 由夹逼收敛原理知, $\lim\limits_{n\to\infty}h_n=0$, 从而 $\lim\limits_{n\to\infty}\sqrt[n]{n}=1$.

(2) 令 $a=\max\{a_1,a_2,\cdots,a_m\}$, 则有
\[
  a^n\leqslant a_1^n+a_2^n+\cdots+a_m^n\leqslant ma^n.
\]
于是
\[
  a\leqslant (a_1^n+a_2^n+\cdots+a_m^n)^{\frac1n}\leqslant \sqrt[n]{m}\,a.
\]
由于当 $n\geqslant m$ 时有
\[
  1\leqslant\sqrt[n]{m}\leqslant\sqrt[n]{n},
\]
由夹逼收敛原理和 (1) 知 $\lim\limits_{n\to\infty}\sqrt[n]{m}=1$。再由夹逼收敛原理, 得
\[
  \lim_{n\to\infty}(a_1^n+a_2^n+\cdots+a_m^n)^{\frac1n}
  =
  a
  =
  \max\{a_1,a_2,\cdots,a_m\}.
\qedhere
\]
\end{solve}


\begin{example}\label{example:ma-2-2-3}
设 $x_n\to0(n\to\infty)$, $p$ 为正整数。证明 $\lim\limits_{n\to\infty}\sqrt[p]{1+x_n}=1$.
\end{example}


\begin{proof}
由题设, 不妨假定 $|x_n|<1$ 对一切 $n$ 成立。这样, 我们有
\[
\begin{aligned}
  1-|x_n|
  &=\sqrt[p]{(1-|x_n|)^p}\\
  &\leqslant \sqrt[p]{1-|x_n|}
   \leqslant \sqrt[p]{1+x_n}
   \leqslant \sqrt[p]{1+|x_n|}\\
  &\leqslant \sqrt[p]{(1+|x_n|)^p}
  =1+|x_n|.
\end{aligned}
\]
而
\[
  \lim_{n\to\infty}(1-|x_n|)
  =
  \lim_{n\to\infty}(1+|x_n|)
  =
  1,
\]
所以, 由夹逼收敛原理知 $\lim\limits_{n\to\infty}\sqrt[p]{1+x_n}=1$ 成立.
\end{proof}


\begin{example}\label{example:ma-2-2-4}
设
\[
  a_1=2,\qquad a_{n+1}=2+\frac1{a_n}\quad(n=1,2,\cdots),
\]
试问 $\{a_n\}$ 是否收敛? 若收敛则求其极限.
\end{example}


\begin{analyze}
首先我们发现, 若 $\lim\limits_{n\to\infty}a_n=a$, 则必有
\[
  a=2+\frac1a.
\]
注意到 $a_n>0$, 从而得 $a=\sqrt2+1$.

令 $a_n=\sqrt2+1+h_n,\ (n=1,2,\cdots)$, 则 $\{a_n\}$ 收敛的充要条件是 $h_n=o(1)(n\to\infty)$.
\end{analyze}


\begin{solve}
令 $a_n=\sqrt2+1+h_n,\ n=1,2,\cdots$。我们有
\[
  1+\sqrt2+h_{n+1}=2+\frac1{1+\sqrt2+h_n},
\]
整理得
\[
  h_{n+1}=\frac{h_n(1-\sqrt2)}{1+\sqrt2+h_n}.
\]
显然 $0<|h_1|=\sqrt2-1<1$。对 $k\geqslant1$, 若 $|h_k|<1$, 则有
\[
  |h_{k+1}|<|h_k|\frac{\sqrt2-1}{\sqrt2}<1.
\]
因此对所有的正整数 $n$, 有
\[
  |h_n|<1.
\]
再利用 $h_n$ 的递推关系我们推知
\[
  |h_{n+1}|\leqslant\frac12|h_n|,\qquad n=1,2,\cdots,
\]
从而对任意的 $n$, 有
\[
  0<|h_n|<\frac1{2^{n-1}}.
\]
由于 $\lim\limits_{n\to\infty}\dfrac1{2^{n-1}}=0$, 由夹逼收敛原理得 $\lim\limits_{n\to\infty}h_n=0$。因此 $\{a_n\}$ 收敛, 并且 $\lim\limits_{n\to\infty}a_n=\sqrt2+1$.

作为本节的结束, 我们简要地讨论一下无穷小量和无穷大量的四则运算。对任意两个无穷小量 $\{x_n\}$ 和 $\{y_n\}$, 由无穷小量定义和极限的四则运算定理知, 它们的和、差以及积仍然是无穷小量, 但它们的商可能有各种情形发生, 请看下例.
\end{solve}


\begin{example}\label{example:ma-2-2-5}
设
\[
  x_n=\frac1{n+1},\qquad y_n=\frac1{n^2},\qquad z_n=\frac{(-1)^{n+1}}{n+1}\quad(n=1,2,\cdots).
\]
试考查:
\[
  \text{(1) }\lim_{n\to\infty}\frac{x_n}{y_n};\qquad
  \text{(2) }\lim_{n\to\infty}\frac{y_n}{x_n};\qquad
  \text{(3) }\lim_{n\to\infty}\frac{z_n}{x_n}.
\]
\end{example}


\begin{solve}
(1)
\[
  \lim_{n\to\infty}\frac{x_n}{y_n}
  =
  \lim_{n\to\infty}\frac{\dfrac1{n+1}}{\dfrac1{n^2}}
  =
  \lim_{n\to\infty}\frac{n^2}{n+1}
  =
  \infty;
\]
(2)
\[
  \lim_{n\to\infty}\frac{y_n}{x_n}
  =
  \lim_{n\to\infty}\frac{\dfrac1{n^2}}{\dfrac1{n+1}}
  =
  \lim_{n\to\infty}\frac{n+1}{n^2}
  =
  0;
\]
(3)
\[
  \lim_{n\to\infty}\frac{z_n}{x_n}
  =
  \lim_{n\to\infty}\frac{\dfrac{(-1)^{n+1}}{n+1}}{\dfrac1{n+1}}
  =
  \lim_{n\to\infty}(-1)^{n+1},
\]
极限不存在.

我们习惯把两个无穷小量的商形式地记为 $\dfrac00$, 上面例子告诉我们 $\dfrac00$ 是一个不定式, 即它的极限是不确定的, 可以有各种可能情况发生.

现在我们再来讨论无穷大量的情况。我们有:
\[
\begin{array}{ll}
  \text{(1)} & (\pm\infty)+(\pm\infty)=\pm\infty;\\
  \text{(2)} & \infty\cdot\infty=\infty;\\
  \text{(3)} & \infty\cdot O_0(1)=\infty.
\end{array}
\]
但 $(+\infty)-(+\infty)$, $\dfrac{\infty}{\infty}$ 等则是不定式。请读者作为练习, 举例来说明它们的不定性.

对于以上的形式记号, 它们的意义是明确的。例如, (3) 中表示的意思是一个无穷大量与一个绝对值有正下界的有界序列的积仍然是一个无穷大量.
\end{solve}


\begin{example}\label{example:ma-2-2-6}
设 $x_n=a_kn^k+a_{k-1}n^{k-1}+\cdots+a_0(a_k\ne0,k\geqslant1)$, 证明 $\lim\limits_{n\to\infty}x_n=\infty$.
\end{example}


\begin{proof}
设
\[
  x_n=n^k(a_k+a_{k-1}n^{-1}+a_{k-2}n^{-2}+\cdots+a_0n^{-k})=n^ky_n,
\]
由于
\[
  \lim_{n\to\infty}y_n=a_k\ne0,
\]
因此对所有充分大 $n$ 有
\[
  \frac{3|a_k|}{2}\geqslant |y_n|\geqslant\frac{|a_k|}{2},
\]
即 $y_n=O_0(1)$。再由 $\lim\limits_{n\to\infty}n^k=\infty$, 即知
\[
  \lim_{n\to\infty}x_n=\infty.
\qedhere
\]
\end{proof}


\paragraph*{思考题}
对于这一例子试问下面的证明是否正确?
\[
  \lim_{n\to\infty}x_n
  =
  \lim_{n\to\infty}a_kn^k+\lim_{n\to\infty}a_{k-1}n^{k-1}+\cdots+\lim_{n\to\infty}a_0
  =
  \infty+\infty+\cdots+\infty+a_0=\infty.
\]

\section{单调收敛原理}
\label{sec:ma-2-3}

解决一个序列的收敛性问题, 只从定义出发往往是不够的, 还需要建立一系列理论结果。本节我们就来介绍单调收敛原理和它的一些典型的应用范例.

\subsection{单调收敛原理}
\label{subsec:ma-2-3-1}

若序列 $\{x_n\}$ 满足:
\[
  x_n\leqslant x_{n+1},\qquad \forall n\in\mathbb{N},
\]
则称 $\{x_n\}$ 是\keyterm{单调递增} (上升) 的序列; 若满足:
\[
  x_n\geqslant x_{n+1},\qquad \forall n\in\mathbb{N},
\]
则称 $\{x_n\}$ 是\keyterm{单调递减} (下降) 的序列; 单调上升的序列和单调下降的序列统称为\keyterm{单调序列}。细心的读者不难发现, 若将序列看成定义在 $\mathbb{N}$ 上的函数 $f(x)$, 则序列的单调性即是函数 $f(x)$ 的单调性.

\begin{theorem}[{单调收敛原理}]\label{theorem:ma-2-3-1}
单调有界序列必收敛.
\end{theorem}


\begin{proof}
这里只对单调上升的情形给出证明, 单调下降的情形请读者自己证明.

设 $\{x_n\}$ 单调上升, 而且有上界, 即
\[
  x_1\leqslant x_2\leqslant\cdots\leqslant x_n\leqslant x_{n+1}\leqslant\cdots\leqslant M,
\]
其中 $M$ 是一个正数。现考虑集合 $E=\{x_n:n\in\mathbb{N}\}$, 则 $E$ 是一个非空有上界的集合, 故由确界定理, $E$ 必有上确界, 令 $a=\sup\{x_n\}$。由上确界的定义知, $a$ 满足:

(1) $x_n\leqslant a,\ \forall n\in\mathbb{N}$;

(2) $\forall\varepsilon>0$, $\exists x_N$, 使得 $a-\varepsilon<x_N$.

这样, $\forall\varepsilon>0$, 当 $n>N$ 时, 有
\[
  a-\varepsilon<x_N\leqslant x_n\leqslant a,
\]
因此有
\[
  |x_n-a|<\varepsilon.
\]
这就证明了 $\lim\limits_{n\to\infty}x_n=a=\sup\{x_n\}$.
\end{proof}


\begin{remark}\label{remark:ma-2-source-56}
容易证明: 若 $\{x_n\}$ 单调上升无上界, 则 $\lim\limits_{n\to\infty}x_n=+\infty$; 若 $\{x_n\}$ 单调下降无下界, 则 $\lim\limits_{n\to\infty}x_n=-\infty$。这样, 结合定理~\ref{theorem:ma-2-3-1}, 我们知单调序列总是广义收敛的, 并且有

(1) 若 $\{x_n\}$ 单调上升, 则 $\lim\limits_{n\to\infty}x_n=\sup\{x_n\}$;

(2) 若 $\{x_n\}$ 单调下降, 则 $\lim\limits_{n\to\infty}x_n=\inf\{x_n\}$.
\end{remark}

容易看出: 对一个单调递增序列 $\{x_n\}$, 若记 $A=\lim\limits_{n\to\infty}x_n$, 则或者存在 $N>0$ 使得 $x_n=x_N=A(\forall n>N)$, 或者 $x_n<A(\forall n>0)$.

\begin{example}\label{example:ma-2-3-1}
设 $A>0$, $x_1>0$, 定义
\[
  x_{n+1}=\frac12\left(x_n+\frac A{x_n}\right),\qquad n=1,2,\cdots.
\]
证明 $\lim\limits_{n\to\infty}x_n$ 存在, 并求出它的值.
\end{example}


\begin{proof}
易知这样产生的 $x_n$ 均为正数, 故有
\[
  x_{n+1}=\frac12\left(x_n+\frac A{x_n}\right)
  \geqslant \sqrt{x_n\frac A{x_n}}=\sqrt A
  \quad(n\geqslant2).
\]
这表明 $\{x_n\}$ 有下界。此外,
\[
  x_{n+1}-x_n
  =
  \frac12\left(x_n+\frac A{x_n}\right)-x_n
  =
  \frac{A-x_n^2}{2x_n}\leqslant0
  \quad(n\geqslant2).
\]
这又证明了 $\{x_n\}$ 是单调下降的。因此, $\lim\limits_{n\to\infty}x_n$ 存在, 设其为 $a$, 则在
\[
  x_{n+1}=\frac12\left(x_n+\frac A{x_n}\right)
\]
两边取极限, 得
\[
  a=\frac12\left(a+\frac Aa\right).
\]
解上述方程得 $a=\pm\sqrt A$。注意到 $a\geqslant\sqrt A$, 便有 $a=\sqrt A$, 即 $\lim\limits_{n\to\infty}x_n=\sqrt A$.
\end{proof}


\begin{remark}\label{remark:ma-2-source-59}
例~\ref{example:ma-2-3-1} 给出了一个求开方的近似方法。例如, 令 $x_1=A=2$, 直接计算可得 $x_4=1.414257$, 这已经是 $\sqrt2$ 的一个很好的近似, 其实 $\sqrt2=1.414213\cdots$.
\end{remark}

此外, 还需指出的是, 当我们知道了极限存在时, 可通过在序列的递推关系式中取极限而得到所求的极限值; 倘若我们并不知道极限是否存在, 则一般不能在递推式中取极限, 否则将得到荒谬的结果, 参见下面的例子.

\begin{example}\label{example:ma-2-3-2}
设 $x_1=1,\ x_{n+1}=-x_n,\ n=1,2,\cdots$, 易见 $\lim\limits_{n\to\infty}x_n$ 不存在。但是, 若假定 $\lim\limits_{n\to\infty}x_n=a$, 并且在 $x_{n+1}=-x_n$ 两边取极限, 得 $a=-a$, 由此推出 $\lim\limits_{n\to\infty}x_n=0$ 的荒谬结果.
\end{example}


\begin{example}\label{example:ma-2-3-3}
设序列 $\{a_n\}$ 满足
\[
  a_1=1,\qquad a_{n+1}=1+\frac1{a_n},\qquad n=1,2,\cdots,
\]
证明 $\{a_n\}$ 收敛, 并求其极限.
\end{example}


\begin{analyze}
直接计算, 得
\[
  a_1=1,\quad a_2=2,\quad a_3=\frac32,\quad a_4=\frac53,\quad a_5=\frac85.
\]
它们满足
\[
  a_1<a_3<a_5<a_4<a_2,
\]
由此我们推测: $\{a_{2k+1}\}$ 单调上升, 而 $\{a_{2k}\}$ 单调下降.
\end{analyze}


\begin{proof}
由 $a_n$ 满足的条件可导出
\begin{equation}
  a_{n+2}=2-\frac1{a_n+1}.
  \label{eq:ma-2-3-1}
\end{equation}
于是
\[
\begin{aligned}
  a_{n+2}-a_n
  &=
  \left(2-\frac1{a_n+1}\right)
  -
  \left(2-\frac1{a_{n-2}+1}\right)\\
  &=
  \frac{a_n-a_{n-2}}{(a_{n-2}+1)(a_n+1)}.
\end{aligned}
\]
这表明 $a_{n+2}-a_n$ 与 $a_n-a_{n-2}$ 同号。因此, 利用这一结果, 从 $a_1<a_3$ 和 $a_4<a_2$ 出发, 就可归纳地导出 $\{a_{2k+1}\}$ 单调上升和 $\{a_{2k}\}$ 单调下降。此外, \eqref{eq:ma-2-3-1} 式和 $a_n$ 的递推公式蕴涵着
\[
  1<a_{2k},\qquad a_{2k+1}<2,\qquad \forall k\in\mathbb{N}.
\]
这样, 由单调收敛原理知, $\{a_{2k+1}\}$ 和 $\{a_{2k}\}$ 都收敛。令
\[
  \lim_{k\to\infty}a_{2k+1}=a,\qquad
  \lim_{k\to\infty}a_{2k}=b,
\]
则有 $1<a,b<2$, 而且 \eqref{eq:ma-2-3-1} 式蕴涵着
\[
  a=2-\frac1{a+1},\qquad b=2-\frac1{b+1}.
\]
由此可导出
\[
  a=b=\frac12(1+\sqrt5).
\]
这就证明了 $\{a_n\}$ 收敛, 而且 $\lim\limits_{n\to\infty}a_n=\dfrac12(1+\sqrt5)$ (见本章习题 7).
\end{proof}


\subsection{\texorpdfstring{无理数 $e$ 和欧拉常数 $c$}{无理数 e 和欧拉常数 c}}
\label{subsec:ma-2-3-2}

作为单调收敛原理的应用, 我们给出无理数 $e$ 和欧拉常数 $c$ 的定义, 它们是数学中两个重要的常数.

考查序列
\[
  x_n=\left(1+\frac1n\right)^n,\qquad
  y_n=\left(1+\frac1n\right)^{n+1},\qquad n=1,2,\cdots.
\]
首先, 对任意的正整数 $n$, 有
\[
  y_n=\left(1+\frac1n\right)x_n>x_n.
\]
其次, 对任意的正整数 $n$, 有
\[
\begin{aligned}
  \frac{x_{n+1}}{x_n}
  &=
  \frac{(n+2)^{n+1}}{(n+1)^{n+1}}\frac{n^n}{(n+1)^n}\frac n{n+1}\\
  &=
  \left(1-\frac1{(n+1)^2}\right)^{n+1}\frac{n+1}{n}\\
  &>
  \left(1-\frac1{n+1}\right)\frac{n+1}{n}
  =
  1;
\end{aligned}
\]
\[
\begin{aligned}
  \frac{y_n}{y_{n+1}}
  &=
  \left(\frac{n+1}{n}\right)^{n+1}
  \left(\frac{n+1}{n+2}\right)^{n+2}\\
  &=
  \left(1+\frac1{n^2+2n}\right)^{n+1}\frac{n+1}{n+2}\\
  &>
  \left(1+\frac{n+1}{n^2+2n}\right)\frac{n+1}{n+2}\\
  &=
  \frac{n^3+4n^2+4n+1}{n^3+4n^2+4n}>1.
\end{aligned}
\]
这样我们就证明了: $\{x_n\}$ 单调上升有上界, $\{y_n\}$ 单调下降有下界。因此, 它们二者都收敛, 而且有相同的极限。我们约定用字母 $e$ 来表示其极限值, 即
\[
  \lim_{n\to\infty}\left(1+\frac1n\right)^n
  =
  \lim_{n\to\infty}\left(1+\frac1n\right)^{n+1}
  =
  e.
\]
$e$ 是数学中最重要的常数之一。以后我们可以证明它是一个无理数并计算出具有任何给定精确度的近似值, 如
\[
  e=2.718281828459045\cdots.
\]
在数学的理论研究和应用中, 以 $e$ 为底的对数起着重要的作用, 这种对数称为\keyterm{自然对数}, 记为 $\ln x$, 即
\[
  \ln x=\log_e x.
\]

此外, 上述讨论已经证明了如下的不等式
\begin{equation}
  \left(1+\frac1n\right)^n
  <
  e
  <
  \left(1+\frac1n\right)^{n+1},
  \qquad \forall n\in\mathbb{N}.
  \label{eq:ma-2-3-2}
\end{equation}
对 \eqref{eq:ma-2-3-2} 两边取对数, 得
\[
  n\ln\left(1+\frac1n\right)
  <
  1
  <
  (n+1)\ln\left(1+\frac1n\right).
\]
由此即得
\begin{equation}
  \frac1{n+1}
  <
  \ln\left(1+\frac1n\right)
  <
  \frac1n,\qquad \forall n\in\mathbb{N}.
  \label{eq:ma-2-3-3}
\end{equation}
从而
\[
  \sum_{k=1}^n\frac1{k+1}
  <
  \sum_{k=1}^n\ln\left(1+\frac1k\right)
  <
  \sum_{k=1}^n\frac1k.
\]
再注意到
\[
  \sum_{k=1}^n\ln\left(1+\frac1k\right)=\ln(n+1),
\]
便有
\begin{equation}
  \sum_{k=1}^n\frac1{k+1}
  <
  \ln(n+1)
  <
  \sum_{k=1}^n\frac1k,\qquad \forall n\in\mathbb{N}.
  \label{eq:ma-2-3-4}
\end{equation}
现考虑序列
\[
  z_n=1+\frac12+\cdots+\frac1n-\ln n,\qquad n=1,2,\cdots.
\]
由不等式 \eqref{eq:ma-2-3-4}, 可得
\[
  z_n>\ln(n+1)-\ln n>0
\]
和
\[
\begin{aligned}
  z_n-z_{n+1}
  &=
  \ln(n+1)-\ln n-\frac1{n+1}\\
  &=
  \ln\left(1+\frac1n\right)-\frac1{n+1}>0,
\end{aligned}
\]
其中最后一个不等式用到了不等式 \eqref{eq:ma-2-3-3}。这样, 我们就证明了 $\{z_n\}$ 单调下降有下界。因此它收敛。记
\[
  c=\lim_{n\to\infty}\left(1+\frac12+\cdots+\frac1n-\ln n\right),
\]
称 $c$ 为\keyterm{欧拉 (Euler) 常数}。我们可计算出
\[
  c=0.577216\cdots.
\]
欧拉常数也是一个重要的数学常数, 并且是一个不好研究的常数, 迄今为止人们还不知道它是有理数还是无理数.

由上面的讨论, 我们有
\[
  1+\frac12+\cdots+\frac1n=\ln n+c+\alpha_n,
\]
其中 $\alpha_n=o(1)(n\to\infty)$.

此外, 由不等式 \eqref{eq:ma-2-3-2}, 可得
\[
  \frac{(n+1)^n}{n!}
  =
  \prod_{k=1}^n\left(1+\frac1k\right)^k
  <
  e^n
  <
  \prod_{k=1}^n\left(1+\frac1k\right)^{k+1}
  =
  \frac{(n+1)^{n+1}}{n!}.
\]
由此可得
\[
  \left(\frac{n+1}{e}\right)^n<n!<\left(\frac{n+1}{e}\right)^n(n+1),
\]
从而
\[
  \frac{n+1}{e}<\sqrt[n]{n!}<\frac{n+1}{e}\sqrt[n]{n+1}.
\]
这说明 $\lim\limits_{n\to\infty}\sqrt[n]{n!}=+\infty$, 而且
\[
  \lim_{n\to\infty}\frac n{\sqrt[n]{n!}}=e.
\]

这样, 我们对无穷大量 $\left\{1+\dfrac12+\cdots+\dfrac1n\right\}$ 和 $\{n!\}$ 的``量级''就有了一个比较清楚的认识, 前者是一个``弱得惊人''的无穷大量, 而后者是一个``大得出奇''的无穷大量.

\begin{example}\label{example:ma-2-3-4}
求极限
\[
  \lim_{n\to\infty}\left(1-\frac12+\frac13-\cdots+\frac{(-1)^{n-1}}n\right).
\]
\end{example}


\begin{solve}
记
\[
  S_n=1-\frac12+\frac13-\cdots+\frac{(-1)^{n-1}}n,
\]
则有
\[
\begin{aligned}
  \lim_{n\to\infty}S_{2n}
  &=
  \lim_{n\to\infty}
  \left(1+\frac13+\cdots+\frac1{2n-1}
  -\left(\frac12+\cdots+\frac1{2n}\right)\right)\\
  &=
  \lim_{n\to\infty}
  \left(1+\frac12+\cdots+\frac1{2n}
  -2\left(\frac12+\cdots+\frac1{2n}\right)\right)\\
  &=
  \lim_{n\to\infty}
  \left(c+\ln2n+\alpha_{2n}-(c+\ln n+\alpha_n)\right)\\
  &=
  \ln2,
\end{aligned}
\]
其中 $c$ 是欧拉常数, $\alpha_n=o(1)(n\to\infty)$。注意到
\[
  \lim_{n\to\infty}S_{2n+1}
  =
  \lim_{n\to\infty}\left(S_{2n}+\frac1{2n+1}\right)
  =
  \ln2,
\]
因此有 (见本章习题 7)
\[
  \lim_{n\to\infty}S_n=\ln2.
\qedhere
\]
\end{solve}


\section{实数系连续性的基本定理}
\label{sec:ma-2-4}

这一节我们给出实数系连续性的几个常用的等价表述, 它们各有特点, 各自适合于不同的应用场合。本节中的很多定理是许多数学论证中的基本工具.

\subsection{闭区间套定理}
\label{subsec:ma-2-4-1}

\begin{theorem}[{闭区间套定理}]\label{theorem:ma-2-4-1}
设 $\{[a_n,b_n]\}$ 是一列闭区间, 并满足:

(1) $[a_n,b_n]\supseteq[a_{n+1},b_{n+1}],\ n=1,2,\cdots$;

(2) $\lim\limits_{n\to\infty}(b_n-a_n)=0$,

则存在唯一的一点 $c\in\mathbb{R}$, 使得 $c\in[a_n,b_n],\ n=1,2,\cdots$, 即
\[
  \{c\}=\bigcap_{n=1}^\infty [a_n,b_n].
\]
\end{theorem}


\begin{proof}
由条件 (1), 知
\[
  a_n\leqslant a_{n+1}\leqslant b_{n+1}\leqslant b_n,\qquad \forall n\in\mathbb{N}.
\]
这表明 $\{a_n\}$ 单调上升有上界, 并且 $\forall n\in\mathbb{N}$, $b_n$ 均是 $\{a_n\}$ 的一个上界; 另一方面, $\{b_n\}$ 单调下降有下界, 并且 $\forall n\in\mathbb{N}$, $a_n$ 均是 $\{b_n\}$ 的一个下界。因此, $\{a_n\}$ 和 $\{b_n\}$ 都收敛。令
\[
  \lim_{n\to\infty}a_n=a=\sup\{a_n\},\qquad
  \lim_{n\to\infty}b_n=b=\inf\{b_n\}.
\]
则有
\[
  a_n\leqslant a\leqslant b\leqslant b_n,\qquad \forall n\in\mathbb{N}.
\]
再由条件 (2), 知
\[
  a-b=\lim_{n\to\infty}a_n-\lim_{n\to\infty}b_n
  =
  \lim_{n\to\infty}(a_n-b_n)=0.
\]
记 $c=a=b$, 则有
\[
  a_n\leqslant\lim_{n\to\infty}a_n=c=\lim_{n\to\infty}b_n\leqslant b_n,\qquad \forall n\in\mathbb{N}.
\]
若另有 $d$ 满足
\[
  a_n\leqslant d\leqslant b_n,\qquad \forall n\in\mathbb{N},
\]
则由夹逼收敛原理, 有
\[
  d=\lim_{n\to\infty}a_n=\lim_{n\to\infty}b_n=c,
\]
故 $c$ 是所有区间 $[a_n,b_n]$ 的唯一公共点.
\end{proof}


\paragraph*{注 1}
闭区间的条件是重要的, 若区间是开的, 则定理的结论不一定成立。例如, 对于开区间列 $\left(0,\dfrac1n\right),\ n=1,2,\cdots$, 显然满足定理的条件 (1) 和 (2), 但
\[
  \bigcap_{n=1}^\infty\left(0,\frac1n\right)=\varnothing.
\]
另一方面, 若开区间列 $\{(a_n,b_n)\}$ 满足条件

(i) $a_n<a_{n+1}<b_{n+1}<b_n,\ n=1,2,\cdots$;

(ii) $\lim\limits_{n\to\infty}(b_n-a_n)=0$,

则定理~\ref{theorem:ma-2-4-1} 的结论仍然成立.

\paragraph*{注 2}
若定理~\ref{theorem:ma-2-4-1} 中条件 (2) 不满足, 即 $\lim\limits_{n\to\infty}(a_n-b_n)\ne0$, 则由定理的证明过程知, 此时 $a\ne b$, 而且容易导出, 此时有
\[
  \bigcap_{n=1}^\infty [a_n,b_n]=[a,b].
\]

\begin{example}\label{example:ma-2-4-1}
设 $\{[a_n,b_n]\}$ 是一列闭区间, 并且满足:
\[
  [a_i,b_i]\cap[a_j,b_j]\ne\varnothing,\qquad \forall i,j\in\mathbb{N},
\]
则存在一点 $c\in\mathbb{R}$, 使得 $c\in[a_n,b_n],\ n=1,2,\cdots$.
\end{example}


\begin{proof}
注意到
\[
  [a,b]\cap[c,d]\ne\varnothing\Rightarrow a\leqslant d,\ c\leqslant b,
\]
即知闭区间列所满足的条件蕴涵着任意的 $b_n$ 均是 $\{a_n\}$ 的上界。于是, 令
\[
  c=\sup_{n\in\mathbb{N}}\{a_n\},
\]
则有
\[
  a_n\leqslant c\leqslant b_n,\qquad n=1,2,\cdots .
\qedhere
\]
\end{proof}


\begin{remark}\label{remark:ma-2-source-70}
本例中的闭区间列并不一定是区间套, 只满足``两两有公共点''的条件, 同样推出它们有公共点的结论。如若再加条件 $b_n-a_n\to0(n\to\infty)$, 则可得到与闭区间套定理完全相同的结论.
\end{remark}


\subsection{有限覆盖定理}
\label{subsec:ma-2-4-2}

设 $A$ 是 $\mathbb{R}$ 中的一个子集, $\{E_\lambda\}_{\lambda\in\Lambda}$ 是 $\mathbb{R}$ 中的一族子集组成的集合, 其中 $\Lambda$ 是一个指标集。若 $A\subseteq\bigcup_{\lambda\in\Lambda}E_\lambda$, 则称 $\{E_\lambda\}_{\lambda\in\Lambda}$ 是 $A$ 的一个\keyterm{覆盖}。若 $\{E_\lambda\}_{\lambda\in\Lambda}$ 是 $A$ 的一个覆盖, 而且对每个 $\lambda\in\Lambda$, $E_\lambda$ 均是一个开区间, 则称 $\{E_\lambda\}_{\lambda\in\Lambda}$ 为 $A$ 的一个\keyterm{开覆盖}; 若 $\{E_\lambda\}_{\lambda\in\Lambda}$ 是 $A$ 的一个覆盖, 而且 $\Lambda$ 的元素只有有限多个, 则称 $\{E_\lambda\}_{\lambda\in\Lambda}$ 是 $A$ 的一个\keyterm{有限覆盖}.

\begin{theorem}[{有限覆盖定理}]\label{theorem:ma-2-4-2}
设 $[a,b]$ 是一个闭区间, $\{E_\lambda\}_{\lambda\in\Lambda}$ 是 $[a,b]$ 的任意一个开覆盖, 则必存在 $\{E_\lambda\}_{\lambda\in\Lambda}$ 一个子集构成 $[a,b]$ 的一个有限覆盖, 即在 $\{E_\lambda\}_{\lambda\in\Lambda}$ 中必有有限个开区间 $E_1,E_2,\cdots,E_N$, 使得 $[a,b]\subseteq\bigcup\limits_{j=1}^N E_j$.
\end{theorem}


\begin{proof}
如若不然, 我们将 $[a,b]$ 等分成两个闭区间, 则其中必有一个不能被 $\{E_\lambda\}_{\lambda\in\Lambda}$ 中的有限多个开区间所覆盖, 记其为 $[a_1,b_1]$; 再将 $[a_1,b_1]$ 等分成两个闭区间, 则必有其中的一个不能被 $\{E_\lambda\}_{\lambda\in\Lambda}$ 中的有限多个开区间所覆盖, 记其为 $[a_2,b_2]$。如此进行下去, 就可得到一列闭区间 $[a_n,b_n]$, 满足:

(1) $[a_{n+1},b_{n+1}]\subseteq[a_n,b_n],\ \forall n$;

(2) $b_n-a_n=\dfrac{b-a}{2^n}\to0(n\to\infty)$;

(3) $\forall n$, $[a_n,b_n]$ 不能被 $\{E_\lambda\}_{\lambda\in\Lambda}$ 中的有限多个开区间所覆盖.

由闭区间套定理知, 存在唯一的 $\xi$, 使得 $\xi\in[a_n,b_n],\ \forall n\in\mathbb{N}$。由于 $\xi\in[a,b]$, 而 $\{E_\lambda\}_{\lambda\in\Lambda}$ 是 $[a,b]$ 的一个开覆盖, 故必存在它的一个开区间 $E_{\lambda_0}=(c,d)$, 使得 $\xi\in(c,d)$, 即 $c<\xi<d$。由于 $\{a_n\}$ 单调上升, $\{b_n\}$ 单调下降, 而且
\[
  \lim_{n\to\infty}a_n=\xi=\lim_{n\to\infty}b_n,
\]
故必有充分大的正整数 $n$, 使得
\[
  c<a_n\leqslant\xi\leqslant b_n<d,
\]
即 $[a_n,b_n]\subseteq E_{\lambda_0}$。这与条件 (3) 矛盾。定理得证.

关于有限覆盖定理, 我们必须注意以下两点:

(1) 不能将闭区间改为开区间。如, $\left(\dfrac1n,\dfrac2n\right)(n=1,2,\cdots)$ 是 $(0,1)$ 的一个开覆盖, 但没有有限覆盖.

(2) 该定理是说, 对于闭区间的任意一个开覆盖必可从中找出有限多个开区间就可将该闭区间覆盖, 而并不是说, 任意一个闭区间都可被有限多个开区间覆盖。若是这样, $(a-1,b+1)$ 总是 $[a,b]$ 的一个有限开覆盖.

在今后的应用中, 常可根据函数的局部性质得到一些开区间, 然后通过有限覆盖定理而得到函数在区间的整体性质.
\end{proof}


\begin{example}\label{example:ma-2-4-2}
设函数 $f(x)$ 在 $[a,b]$ 上有定义, 且对任意的 $x\in[a,b]$, 都存在邻域 $U(x,\delta(x))(\delta(x)>0)$, 使得 $f(x)$ 在 $U(x,\delta(x))\cap[a,b]$ 上有界。试证明 $f(x)$ 在 $[a,b]$ 上有界.
\end{example}


\begin{proof}
对任意的 $x\in[a,b]$, 由题设存在 $U(x,\delta(x))$, 使得 $f(x)$ 在 $U(x,\delta(x))\cap[a,b]$ 上有界, 记其界为 $M_x>0$, 即
\[
  |f(t)|\leqslant M_x,\qquad \forall t\in U(x,\delta(x))\cap[a,b].
\]
令
\[
  \mathcal{F}=\{U(x,\delta(x))=(x-\delta(x),x+\delta(x)):x\in[a,b]\}.
\]
显然 $\mathcal{F}$ 是 $[a,b]$ 的一个开覆盖。由有限覆盖定理知, 存在 $\mathcal{F}$ 的有限个开区间
\[
  U(x_1,\delta(x_1)),\ U(x_2,\delta(x_2)),\ \cdots,\ U(x_N,\delta(x_N)),
\]
使得
\[
  [a,b]\subseteq\bigcup_{k=1}^N U(x_k,\delta(x_k)).
\]
令
\[
  M=\max\{M_{x_1},M_{x_2},\cdots,M_{x_N}\},
\]
则对任意的 $x\in[a,b]$, 必存在 $k(1\leqslant k\leqslant N)$, 使得 $x\in U(x_k,\delta(x_k))$, 从而
\[
  |f(x)|\leqslant M_{x_k}\leqslant M.
\]
这表明 $f(x)$ 在 $[a,b]$ 上有界.
\end{proof}


\begin{remark}\label{remark:ma-2-source-75}
由上题的逆否命题我们可以得到以下有趣的事实: 若 $f(x)$ 在 $[a,b]$ 上无界, 则存在 $\xi\in[a,b]$, 使得 $\forall\delta>0$, $f(x)$ 在 $U(\xi,\delta)\cap[a,b]$ 上均无界.
\end{remark}


\begin{example}\label{example:ma-2-4-3}
证明: 无论你用什么方法都不可能将 $[0,1]$ 区间中的全体实数排列成一个序列.
\end{example}


\begin{proof}
用反证法。假定
\[
  [0,1]=\{x_1,x_2,x_3,\cdots,x_n,\cdots\}.
\]
取 $0<\varepsilon<\dfrac14$, 则开区间簇
\[
  \mathcal{F}=\left\{U\left(x_n,\frac{\varepsilon}{2^n}\right):n\in\mathbb{N}\right\}
\]
是 $[0,1]$ 区间的一个开覆盖。由有限覆盖定理知, $\mathcal{F}$ 中存在有限个开区间
\[
  U\left(x_{n_1},\frac{\varepsilon}{2^{n_1}}\right),\
  U\left(x_{n_2},\frac{\varepsilon}{2^{n_2}}\right),\
  \cdots,\
  U\left(x_{n_\ell},\frac{\varepsilon}{2^{n_\ell}}\right),
\]
它们可将 $[0,1]$ 区间覆盖。既然它们覆盖了 $[0,1]$ 区间, 那么这 $\ell$ 个区间的长度和必大于 $1$。令
\[
  m=\max\{n_1,n_2,\cdots,n_\ell\},
\]
则有
\[
  1<
  2\sum_{k=1}^{\ell}\frac{\varepsilon}{2^{n_k}}
  <
  2\varepsilon\sum_{k=1}^m\frac1{2^k}
  =
  2\varepsilon\frac{1-\dfrac1{2^{m+1}}}{1-\dfrac12}
  <
  4\varepsilon<1.
\]
这一矛盾说明假设区间 $[0,1]$ 中的全体实数可以排列成一个序列是错误的.
\end{proof}


\begin{remark}\label{remark:ma-2-source-78}
若一个数集 $E$ 中只有有限个元素或可以将它的所有元素排成一个序列, 则称 $E$ 是一个\keyterm{可数集}。上述例子说明区间 $[0,1]$ 不是可数集.
\end{remark}


\subsection{聚点原理}
\label{subsec:ma-2-4-3}

\begin{definition}\label{definition:ma-2-4-1}
设 $E$ 是 $\mathbb{R}$ 中的一个子集。若 $x_0\in\mathbb{R}$ ($x_0$ 不一定属于 $E$) 满足: 对 $\forall\delta>0$, 有 $U_0(x_0,\delta)\cap E\ne\varnothing$, 则称 $x_0$ 是 $E$ 的一个\keyterm{聚点}.
\end{definition}


\begin{remark}\label{remark:ma-2-source-80}
(1) $x_0$ 是 $E$ 的聚点与 $x_0$ 是否属于 $E$ 无关;

(2) 由聚点的定义容易证明如下三个命题等价:

(i) $x_0$ 是 $E$ 的聚点;
\end{remark}

(ii) $\forall\delta>0$, 在 $U(x_0,\delta)$ 中有 $E$ 的无穷多个点;

(iii) 存在 $E$ 中互异的点组成的序列 $\{x_n\}$, 使得 $\lim\limits_{n\to\infty}x_n=x_0$.

(3) 若 $x_0\in E$, 但它不是 $E$ 的聚点, 则称 $x_0$ 是 $E$ 的一个\keyterm{孤立点}。由定义, 此时必存在 $\delta>0$, 使得 $U(x_0,\delta)\cap E=\{x_0\}$.

\begin{example}\label{example:ma-2-4-4}
设
\[
  E=\left\{1,\frac12,\frac13,\cdots,\frac1n,\cdots\right\},
\]
则 $0$ 是它的唯一聚点, 而且 $0\notin E$; $\forall n\in\mathbb{N}$, $\dfrac1n$ 是 $E$ 的孤立点.
\end{example}


\begin{example}\label{example:ma-2-4-5}
设 $E$ 是 $[0,1]$ 中所有有理数组成的集合, 则 $E$ 的聚点全体是 $[0,1]$, 而 $E$ 没有孤立点。注意, 此时 $E$ 是它的聚点集 $[0,1]$ 的真子集.
\end{example}


\begin{theorem}[{聚点原理}]\label{theorem:ma-2-4-3}
$\mathbb{R}$ 中的任何一个有界无穷子集至少有一个聚点.
\end{theorem}


\begin{proof}
设 $E$ 是 $\mathbb{R}$ 中的一个有界无穷子集, 不妨设 $E\subseteq[a,b]$。倘若定理的结论不真, 则 $\forall x\in[a,b]$, $x$ 都不是 $E$ 的聚点, 从而 $\exists\delta_x>0$, 使得 $U(x,\delta_x)\cap E$ 至多只有一个点。显然,
\[
  \mathcal{F}=\{U(x,\delta_x):x\in[a,b]\}
\]
构成 $[a,b]$ 的一个开覆盖, 从而由有限覆盖定理知, 存在有限个开区间 $U(x_j,\delta_{x_j})(j=1,2,\cdots,m)$, 它们覆盖了 $[a,b]$, 即
\[
  [a,b]=\bigcup_{j=1}^m U(x_j,\delta_{x_j}).
\]
注意到 $E\subseteq[a,b]$, 而每个 $U(x_j,\delta_{x_j})$ 至多只有 $E$ 的一个点, 我们便可推知 $E$ 中至多有 $m$ 个点, 这与 $E$ 是 $\mathbb{R}$ 的无穷子集矛盾, 故 $E$ 必有聚点.
\end{proof}


\paragraph*{思考题}
试利用闭区间套定理证明聚点原理.

聚点原理与序列的子序列密切相关。设 $\{x_n\}$ 是一个序列, 则由该序列的一部分元素按原来的顺序构成的序列 $\{x_{n_k}\}$ 称为是 $\{x_n\}$ 的一个\keyterm{子序列}。关于下标 $n_k$ 需要说明如下三点:

(1) $\{n_k\}$ 中的每一项都是正整数, 并且它是一个严格递增序列:
\[
  n_1<n_2<\cdots<n_k<\cdots;
\]

(2) $x_{n_k}$ 表示在子序列中它是第 $k$ 项, 在原序列中它是第 $n_k$ 项。因此, 子序列的序号是 $k$, 而不是 $n_k$;

(3) 必有 $n_k\geqslant k$, 从而 $\lim\limits_{k\to\infty}n_k=+\infty$.

对于子序列, 我们容易证明如下的结果:

\begin{theorem}\label{theorem:ma-2-4-4}
设 $\lim\limits_{n\to\infty}x_n=a$, 则对 $\{x_n\}$ 的任意子序列 $\{x_{n_k}\}$, 都有 $\lim\limits_{k\to\infty}x_{n_k}=a$.
\end{theorem}

此定理给出了判定一个序列发散的一个方法, 即: 若在一个序列 $\{x_n\}$ 中可以找到两个收敛的子序列 $\{x_{n_k}\}$ 和 $\{x_{n'_k}\}$, 使得
\[
  \lim_{k\to\infty}x_{n_k}=a\ne b=\lim_{k\to\infty}x_{n'_k},
\]
则该序列必然发散.

\begin{example}\label{example:ma-2-4-6}
设 $x_n=\sin\dfrac{n\pi}{2},\ n=1,2,\cdots$, 证明 $\{x_n\}$ 发散.
\end{example}


\begin{proof}
由于
\[
  \lim_{k\to\infty}x_{4k}
  =
  \lim_{k\to\infty}\sin2k\pi=0,
\]
而
\[
  \lim_{k\to\infty}x_{4k+1}
  =
  \lim_{k\to\infty}\sin\left(\frac\pi2+2k\pi\right)=1\ne0,
\]
因此 $\lim\limits_{n\to\infty}\sin\dfrac{n\pi}{2}$ 不存在.
\end{proof}


\begin{example}\label{example:ma-2-4-7}
设
\[
  x_n=\left(1+\frac{(-1)^n}{n}\right)^n,\qquad n=1,2,\cdots,
\]
试证明 $\{x_n\}$ 为发散序列.
\end{example}


\begin{proof}
由于
\[
  \lim_{k\to+\infty}\left(1+\frac{(-1)^{2k}}{2k}\right)^{2k}=e,
\]
而
\[
  \lim_{k\to+\infty}\left(1+\frac{(-1)^{2k+1}}{2k+1}\right)^{2k+1}
  =
  \lim_{k\to+\infty}\frac1{\left(1+\dfrac1{2k}\right)^{2k}\left(1+\dfrac1{2k}\right)}
  =
  \frac1e,
\]
因此 $\lim\limits_{n\to\infty}x_n$ 不存在.

下面我们讲述波尔查诺-魏尔斯特拉斯 (Bolzano-Weierstrass) 定理.
\end{proof}


\begin{theorem}[{波尔查诺-魏尔斯特拉斯定理}]\label{theorem:ma-2-4-5}
任何有界序列必有收敛的子序列.
\end{theorem}


\begin{proof}
设 $\{x_n\}$ 是一有界序列。若 $\{x_n\}$ 只由有限多个数组成, 则它必有无穷多项等于同一个数, 此时定理结论自然成立.

现设 $\{x_n\}$ 由无穷多个互不相同的数组成, 则数集
\[
  E=\{x_n:n\in\mathbb{N}\}
\]
就是 $\mathbb{R}$ 中的一个有界无穷子集。由聚点原理知, 它必有一个聚点, 设其为 $a$。由聚点的定义知, 对任意的正整数 $k$, 在 $U(a,1/k)$ 中必有 $\{x_n\}$ 的无穷多项。这样, 我们首先取
\[
  x_{n_1}\in U(a,1)\cap\{x_1,x_2,\cdots\};
\]
然后, 再取
\[
  x_{n_2}\in U\left(a,\frac12\right)\cap\{x_{n_1+1},x_{n_1+2},\cdots\}.
\]
如此进行下去, 就可找到 $\{x_n\}$ 的一个子列 $\{x_{n_k}\}$ 满足 $x_{n_k}\in U(a,1/k)$, 即
\[
  |x_{n_k}-a|<\frac1k,\qquad k=1,2,\cdots,
\]
从而 $\lim\limits_{k\to\infty}x_{n_k}=a$.
\end{proof}


\subsection{柯西收敛准则}
\label{subsec:ma-2-4-4}

一个序列是否收敛是我们研究的中心问题。下面要介绍的柯西收敛准则从序列本身的性质出发给出了该序列收敛的充分必要条件。在理论和应用上, 它都是十分重要的.

我们首先引入如下的定义.

\begin{definition}\label{definition:ma-2-4-2}
设 $\{x_n\}$ 是一个序列, 若 $\forall\varepsilon>0$, $\exists N$, 当 $n,m>N$ 时, 有 $|x_n-x_m|<\varepsilon$, 则称 $\{x_n\}$ 是一个\keyterm{柯西序列}.
\end{definition}


\begin{theorem}[{柯西收敛准则}]\label{theorem:ma-2-4-6}
序列 $\{x_n\}$ 收敛的充分必要条件是它是一个柯西序列.
\end{theorem}


\begin{proof}
\keyterm{必要性}\quad 设序列 $\{x_n\}$ 收敛, 并设 $a=\lim\limits_{n\to\infty}x_n$, 则由极限的定义知, $\forall\varepsilon>0$, $\exists N$, 当 $n>N$ 时, 有
\[
  |x_n-a|<\frac\varepsilon2,
\]
从而当 $n,m>N$ 时, 有
\[
  |x_n-x_m|\leqslant |x_n-a|+|a-x_m|<\varepsilon.
\]

\keyterm{充分性}\quad 设 $\{x_n\}$ 是一个柯西序列, 对 $\varepsilon_0=1$, 则 $\exists N$, 当 $n\geqslant N+1>N$ 时, 有
\[
  |x_n-x_{N+1}|<1,
\]
从而当 $n>N+1$ 时, 有 $|x_n|<1+|x_{N+1}|$。令
\[
  M=\max\{|x_1|,|x_2|,\cdots,|x_N|,1+|x_{N+1}|\},
\]
则 $\forall n\in\mathbb{N}$, 有 $|x_n|\leqslant M$。这表明 $\{x_n\}$ 是有界的。于是由波尔查诺-魏尔斯特拉斯定理知, 存在 $\{x_n\}$ 的一个收敛子列 $\{x_{n_k}\}$, 记 $a=\lim\limits_{k\to\infty}x_{n_k}$.

下证
\[
  \lim_{n\to\infty}x_n=a.
\]
$\forall\varepsilon>0$, 由 $a=\lim\limits_{k\to\infty}x_{n_k}$ 知, $\exists K$, 当 $k\geqslant K$ 时, 有
\[
  |x_{n_k}-a|<\frac\varepsilon2.
\]
再由 $\{x_n\}$ 是柯西序列知, $\exists N$, 当 $n,m>N$ 时, 有
\[
  |x_n-x_m|<\frac\varepsilon2.
\]
这样, 取 $k>K$, 使得 $n_k>N$, 则当 $n>N$ 时, 有
\[
  |x_n-a|\leqslant |x_n-x_{n_k}|+|x_{n_k}-a|<\varepsilon,
\]
即 $\lim\limits_{n\to\infty}x_n=a$.
\end{proof}


\begin{remark}\label{remark:ma-2-source-95}
我们有时候将一些数构成的集合 $K$ 称为是一个空间。如果一个空间 $K$ 中的任何柯西序列 $\{x_n\}$ 都在 $K$ 中存在极限, 即存在 $A\in K$, 使得 $\lim\limits_{n\to\infty}x_n=A$, 则称 $K$ 是\keyterm{完备的}。例如, 上面证明的柯西收敛准则说明了实数域 $\mathbb{R}$ 是完备的。显然, 有理数域 $\mathbb{Q}$ 是不完备的, 这是因为 $\mathbb{Q}$ 中的柯西序列的极限不一定是有理数。另外, 任何开区间 $(a,b)$ 也是不完备的, 这是因为存在柯西序列 $\left\{b-\dfrac{b-a}{n}\right\}\subset(a,b)$ 但
\[
  \lim_{n\to\infty}\left(b-\frac{b-a}{n}\right)=b\notin(a,b).
\]
\end{remark}


\begin{example}\label{example:ma-2-4-8}
证明序列 $\{x_n\}$ 是发散的, 其中
\[
  x_n=1+\frac12+\frac13+\cdots+\frac1n,\qquad n=1,2,\cdots.
\]
\end{example}


\begin{proof}
在本章的第一节我们已经用定义证明了该序列发散, 这里我们再用柯西收敛准则证明它发散。取 $\varepsilon_0=\dfrac12$, 则对任意的正整数 $n$, 有
\[
  |x_{2n}-x_n|
  =
  \frac1{n+1}+\cdots+\frac1{2n}
  \geqslant \frac n{2n}=\frac12,
\]
所以它不是柯西序列, 从而它必是发散的.
\end{proof}


\begin{example}\label{example:ma-2-4-9}
证明序列 $\{y_n\}$ 是收敛的, 其中
\[
  y_n=1+\frac1{2^2}+\frac1{3^2}+\cdots+\frac1{n^2},\qquad n=1,2,\cdots.
\]
\end{example}


\begin{proof}
由于对任意的 $n,m\in\mathbb{N}$, 我们不妨设 $n>m$, 有
\[
\begin{aligned}
  |y_n-y_m|
  &=
  \frac1{(m+1)^2}+\cdots+\frac1{n^2}\\
  &\leqslant
  \frac1{m(m+1)}+\cdots+\frac1{(n-1)n}
  \leqslant \frac1m.
\end{aligned}
\]
因此, $\forall\varepsilon>0$, 取 $N=\left[\dfrac1\varepsilon\right]$, 则当 $n,m>N$ 时, 就有
\[
  |y_n-y_m|<\varepsilon,
\]
即 $\{y_n\}$ 是柯西序列, 所以它是收敛的.
\end{proof}


\begin{example}[{压缩映像原理}]\label{example:ma-2-4-10}
设 $f(x)$ 在 $[a,b]$ 上有定义, $f([a,b])\subseteq[a,b]$, 且满足
\[
  |f(x)-f(y)|\leqslant q|x-y|,\qquad \forall x,y\in[a,b],
\]
其中 $0<q<1$。证明: 存在唯一的 $c\in[a,b]$, 使得 $f(c)=c$.
\end{example}


\begin{proof}
任取 $x_0\in[a,b]$, 由条件 $f([a,b])\subseteq[a,b]$ 知, 我们可递推地定义
\[
  x_n=f(x_{n-1}),\qquad n=1,2,\cdots.
\]
由函数所满足的条件, 我们有
\[
  |x_{n+1}-x_n|
  =
  |f(x_n)-f(x_{n-1})|
  \leqslant q|x_n-x_{n-1}|,\qquad \forall n\in\mathbb{N}.
\]
反复应用上述不等式, 可得
\[
  |x_{n+1}-x_n|\leqslant q^n|x_1-x_0|,\qquad \forall n\in\mathbb{N}.
\]
因此, 对任意的正整数 $n$ 和 $p$, 有
\[
\begin{aligned}
  |x_{n+p}-x_n|
  &\leqslant \sum_{k=1}^p |x_{n+k}-x_{n+k-1}|\\
  &\leqslant \sum_{k=1}^p q^{n+k-1}|x_1-x_0|\\
  &=
  \frac{1-q^p}{1-q}q^n|x_1-x_0|\\
  &<
  \frac{q^n}{1-q}|x_1-x_0|=Lq^n,
\end{aligned}
\]
其中 $L=\dfrac1{1-q}|x_1-x_0|$。由此立即可知 $\{x_n\}$ 是柯西序列, 从而它收敛, 记 $\lim\limits_{n\to\infty}x_n=c$。显然, $c\in[a,b]$。又由于
\[
  |f(x_n)-f(c)|\leqslant q|x_{n-1}-c|,
\]
故有 $\lim\limits_{n\to\infty}f(x_n)=f(c)$。这样, 在等式 $x_{n+1}=f(x_n)$ 两边取极限, 即得 $c=f(c)$, 从而 $c$ 的存在性得证.

下证唯一性。若存在 $c'\in[a,b]$ 也满足 $f(c')=c'$, 而且 $c'\ne c$, 则由
\[
  |c-c'|=|f(c)-f(c')|\leqslant q|c-c'|
\]
即得矛盾, 从而 $c$ 的唯一性得证.
\end{proof}


\section{序列的上、下极限}
\label{sec:ma-2-5}

上极限和下极限的概念是极限概念的拓广, 它为研究序列的敛散性开拓了一条全新的思路.

设 $\{x_n\}$ 是有界序列。令
\[
  \ell_n=\inf\{x_n,x_{n+1},x_{n+2},\cdots\},
  \qquad
  h_n=\sup\{x_n,x_{n+1},x_{n+2},\cdots\},
\]
则有
\[
  \ell_1\leqslant\ell_2\leqslant\cdots\leqslant\ell_n\leqslant\cdots\leqslant h_n\leqslant\cdots\leqslant h_2\leqslant h_1,
\]

即 $\{\ell_n\}$ 和 $\{h_n\}$ 是单调有界序列。令
\[
  \ell=\sup_n\{\ell_n\}=\sup_n\inf_k\{x_{n+k}\},
  \qquad
  h=\inf_n\{h_n\}=\inf_n\sup_k\{x_{n+k}\},
\]
则 $\ell$ 与 $h$ 分别称为 $\{x_n\}$ 的\keyterm{下极限}和\keyterm{上极限}, 记为
\[
  \ell=\underline{\lim_{n\to\infty}}\,x_n,\qquad
  h=\overline{\lim_{n\to\infty}}\,x_n;
\]
有时下极限和上极限也记为 $\liminf\limits_{n\to\infty}x_n$ 和 $\limsup\limits_{n\to\infty}x_n$.

此外, 对于无界序列, 我们规定:

(1) 若 $\{x_n\}$ 无上界, 则 $\overline{\lim\limits_{n\to\infty}}\,x_n=+\infty$。注意此时必有: $h_n=+\infty,\ \forall n\in\mathbb{N}$.

(2) 若 $\{x_n\}$ 无下界, 则 $\underline{\lim\limits_{n\to\infty}}\,x_n=-\infty$。注意此时必有: $\ell_n=-\infty,\ \forall n\in\mathbb{N}$.

(3) 若 $\{x_n\}$ 有下界但无上界, 则 $\{\ell_n\}$ 有定义, 故可定义
\[
  \underline{\lim_{n\to\infty}}\,x_n=\sup_n\{\ell_n\},
\]
但须注意的是, 此时 $\sup_n\{\ell_n\}$ 可能为 $+\infty$.

(4) 若 $\{x_n\}$ 有上界但无下界, 则 $\{h_n\}$ 有定义, 故可定义
\[
  \overline{\lim_{n\to\infty}}\,x_n=\inf_n\{h_n\},
\]
但须注意的是, 此时 $\inf_n\{h_n\}$ 可能为 $-\infty$.

这样, 对任意的序列 $\{x_n\}$, $\overline{\lim\limits_{n\to\infty}}\,x_n$ 和 $\underline{\lim\limits_{n\to\infty}}\,x_n$ 都有明确的定义, 而且恒有
\[
  \underline{\lim_{n\to\infty}}\,x_n
  \leqslant
  \overline{\lim_{n\to\infty}}\,x_n.
\]

\begin{example}\label{example:ma-2-5-1}
试求序列 $\{x_n\}$ 的上极限与下极限, 其中
\[
  (1)\ x_n=\sin\frac n2\pi;\qquad
  (2)\ x_n=n^{(-1)^n};
\]
\[
  (3)\ x_n=(-1)^n n;\qquad
  (4)\ x_n=n.
\]
\end{example}


\begin{solve}
(1) 由 $\forall n\in\mathbb{N}$, $\ell_n=-1,\ h_n=1$, 得
\[
  \underline{\lim_{n\to\infty}}\,\sin\frac n2\pi=-1,
  \qquad
  \overline{\lim_{n\to\infty}}\,\sin\frac n2\pi=1.
\]

(2) 由 $\forall n\in\mathbb{N}$, $\ell_n=0,\ h_n=+\infty$, 即 $\{x_n\}$ 无上界, 得
\[
  \underline{\lim_{n\to\infty}}\, n^{(-1)^n}=0,
  \qquad
  \overline{\lim_{n\to\infty}}\, n^{(-1)^n}=+\infty.
\]

(3) 由 $\forall n\in\mathbb{N}$, $\ell_n=-\infty,\ h_n=+\infty$, 即 $\{x_n\}$ 既无上界又无下界, 得
\[
  \underline{\lim_{n\to\infty}}\,(-1)^n n=-\infty,
  \qquad
  \overline{\lim_{n\to\infty}}\,(-1)^n n=+\infty.
\]

(4) 由 $\forall n\in\mathbb{N}$, $\ell_n=n,\ h_n=+\infty$, 得
\[
  \underline{\lim_{n\to\infty}}\, n=+\infty,
  \qquad
  \overline{\lim_{n\to\infty}}\, n=+\infty.
\]

从上极限的定义, 容易证明如下关于上极限的几个等价的描述.
\end{solve}


\begin{theorem}\label{theorem:ma-2-5-1}
设 $\{x_n\}$ 是一有界序列, $h$ 是一实数, 则下列三个命题等价:

(1) $h$ 是 $\{x_n\}$ 的上极限;

(2) $\forall\varepsilon>0,\ \exists N$, 当 $n>N$ 时, 有 $x_n<h+\varepsilon$, 而且对 $\forall K,\ \exists n_k>K$, 使得 $x_{n_k}>h-\varepsilon$;

(3) 存在子列 $\{x_{n_k}\}$, 使得 $\lim\limits_{k\to\infty}x_{n_k}=h$, 而且对任何其他收敛子列 $\{x_{n'_k}\}$, 有 $\lim\limits_{k\to\infty}x_{n'_k}\leqslant h$.
\end{theorem}


\begin{proof}
(1) $\Rightarrow$ (2)。$\forall n\in\mathbb{N}$, 记 $h_n=\sup\{x_n,x_{n+1},\cdots\}$。由 $\overline{\lim\limits_{n\to\infty}}\,x_n=h,\ \forall\varepsilon>0,\ \exists N$, 当 $n>N$ 时, 有
\[
  |h_n-h|<\varepsilon,
  \qquad\text{即}\qquad h-\varepsilon<h_n<h+\varepsilon.
\]
由于对任意的 $n\in\mathbb{N}$, 有 $x_n\leqslant h_n$, 因此, 当 $n>N$ 时, 有
\[
  x_n<h+\varepsilon.
\]
对 $\forall K>0$, 当 $n'_k>\max\{N,K\}$ 时, 有
\[
  h_{n'_k}>h-\varepsilon.
\]
由 $h_{n'_k}=\sup\{x_{n'_k},x_{n'_k+1},x_{n'_k+2},\cdots\}$ 知, 存在 $n_k\geqslant n'_k>K$ 使得有
\[
  x_{n_k}>h-\varepsilon.
\]
综上所述, (2) 成立.

(2) $\Rightarrow$ (3)。由 (2), 对 $\varepsilon=1$, 显然存在正整数 $n_1\geqslant1$, 使得
\[
  h-1<x_{n_1}<h+1.
\]
现假定对 $\varepsilon_k=\dfrac1k(k\geqslant1,\ k\in\mathbb{N})$, 存在 $n_k\in\mathbb{N}$ 满足
\[
  h-\frac1k<x_{n_k}<h+\frac1k.
\]
在 (2) 中取 $\varepsilon=\dfrac1{k+1}$, 从而存在 $N_{k+1}$, 当 $n>N_{k+1}$ 时, 有
\[
  x_n<h+\frac1{k+1}.
\]
令 $K=\max\{N_{k+1},n_k\}$。由 (2) 知存在 $n_{k+1}>K$, 使得
\[
  h-\frac1{k+1}<x_{n_{k+1}}<h+\frac1{k+1}.
\]
由上述归纳法找到的 $\{x_{n_k}\}$ 显然满足 $\lim\limits_{k\to\infty}x_{n_k}=h$.

现设 $\{x_{n'_k}\}$ 是 $\{x_n\}$ 的一个收敛子列, 并设
\[
  \lim_{k\to\infty}x_{n'_k}=h'.
\]
由 (2) 知, $\forall\varepsilon>0,\ \exists N$, 当 $n>N$ 时, 有
\[
  x_n<h+\varepsilon.
\]
注意到 $n'_k\geqslant k$, 从而当 $k>N$ 时, 有
\[
  x_{n'_k}<h+\varepsilon.
\]
这说明
\[
  h'=\lim_{k\to\infty}x_{n'_k}\leqslant h+\varepsilon.
\]
由 $\varepsilon$ 的任意性, 知 $h'\leqslant h$。因此 (3) 的结论成立.

(3) $\Rightarrow$ (1)。设 $\overline{\lim\limits_{n\to\infty}}\,x_n=h'$, 我们将证明 $h'=h$, 其中 $h$ 满足 (3) 中的条件.

由 (3) 知, 存在 $\{x_n\}$ 的子列 $\{x_{n_k}\}$ 满足:
\[
  \lim_{k\to\infty}x_{n_k}=h.
\]
注意到 $\forall k\in\mathbb{N}$, 有 $h_{n_k}\geqslant x_{n_k}$, 从而有 $h'\geqslant h$.

假若 $h'>h$, 我们将推出与 (3) 相矛盾的结果。事实上, 令 $\varepsilon_0=\dfrac{h'-h}{2}>0$。由 (2) 知存在 $\{x_n\}$ 的子列 $\{x_{n'_k}\}$, 使得对 $\forall k\in\mathbb{N}$, 有
\[
  x_{n'_k}>h'-\varepsilon_0=h+\varepsilon_0.
\]
由于 $\{x_n\}$ 为有界序列, 从而 $\{x_{n'_k}\}$ 存在收敛子列 $\{x_{n''_k}\}$, 显然有 $\lim\limits_{k\to\infty}x_{n''_k}\geqslant h+\varepsilon_0>h$。这与 (3) 矛盾, 从而 $h'=h$, 即 (1) 成立.
\end{proof}


\begin{remark}\label{remark:ma-2-source-106}
对于下极限, 我们有相应于定理~\ref{theorem:ma-2-5-1} 的结果, 请读者自己给出这些结果的叙述与说明.
\end{remark}


\begin{theorem}\label{theorem:ma-2-5-2}
(1) 若有界序列 $\{x_n\}$ 由互不相同的数组成, 则上极限 $\overline{\lim\limits_{n\to\infty}}\,x_n$ 是 $\{x_n\}$ 的最大聚点, 而下极限 $\underline{\lim\limits_{n\to\infty}}\,x_n$ 是 $\{x_n\}$ 的最小聚点;

(2) $\{x_{n_k}\}$ 是 $\{x_n\}$ 的任一子列, 则
\[
  \underline{\lim_{n\to\infty}}\,x_n
  \leqslant
  \underline{\lim_{k\to\infty}}\,x_{n_k}
  \leqslant
  \overline{\lim_{k\to\infty}}\,x_{n_k}
  \leqslant
  \overline{\lim_{n\to\infty}}\,x_n;
\]

(3) $\lim\limits_{n\to\infty}x_n=a$ 的充分必要条件是 $\underline{\lim\limits_{n\to\infty}}\,x_n=\overline{\lim\limits_{n\to\infty}}\,x_n=a$, 其中 $a$ 可以是有限数、$-\infty$ 或 $+\infty$.
\end{theorem}


\begin{proof}
(1) 设 $\overline{\lim\limits_{n\to\infty}}\,x_n=h$。由定理~\ref{theorem:ma-2-5-1}(3) 知, 存在 $\{x_n\}$ 的子列 $\{x_{n_k}\}$, 使得 $\lim\limits_{k\to\infty}x_{n_k}=h$。由于 $\{x_{n_k}\}$ 是由互不相同的数组成, $h$ 是 $\{x_{n_k}\}$ 的聚点, 因此 $h$ 更是 $\{x_n\}$ 的聚点。对于 $\{x_n\}$ 的任何聚点 $h'$, 都存在它的子列 $\{x_{n'_k}\}$, 使得 $\lim\limits_{k\to\infty}x_{n'_k}=h'$。由定理~\ref{theorem:ma-2-5-1}(3) 知必有 $h'\leqslant h$。这就证明了 $h$ 是 $\{x_n\}$ 的最大聚点。同样的方法可证 $\underline{\lim\limits_{n\to\infty}}\,x_n$ 是 $\{x_n\}$ 的最小聚点.

(2) 对 $\forall k\in\mathbb{N}$, 有
\[
  \ell_{n_k}=\inf\{x_{n_k},x_{n_k+1},\cdots\}
  \leqslant
  \inf\{x_{n_k},x_{n_{k+1}},\cdots\}.
\]
因此
\[
  \underline{\lim_{n\to\infty}}\,x_n
  =
  \lim_{n\to\infty}\ell_n
  =
  \lim_{k\to\infty}\ell_{n_k}
  \leqslant
  \lim_{k\to\infty}\inf\{x_{n_k},x_{n_{k+1}},\cdots\}
  =
  \underline{\lim_{k\to\infty}}\,x_{n_k}.
\]
对 $\forall k\in\mathbb{N}$, 有
\[
  h'_k=\sup\{x_{n_k},x_{n_{k+1}},\cdots\}
  \leqslant
  \sup\{x_{n_k},x_{n_k+1},\cdots\}=h_{n_k},
\]
因此
\[
  \overline{\lim_{k\to\infty}}\,x_{n_k}
  =
  \lim_{k\to\infty}h'_k
  \leqslant
  \lim_{k\to\infty}h_{n_k}
  =
  \lim_{n\to\infty}h_n
  =
  \overline{\lim_{n\to\infty}}\,x_n.
\]
(3) 我们只证 $a$ 为有限数的情形, 其余的证明留给读者作为练习.

充分性\quad 假定 $\underline{\lim\limits_{n\to\infty}}\,x_n=\overline{\lim\limits_{n\to\infty}}\,x_n=a$, 由于 $\forall n\in\mathbb{N}$, 我们有
\[
  \ell_n=\inf_k\{x_{n+k}\}\leqslant x_n\leqslant\sup_k\{x_{n+k}\}=h_n,
\]
由夹逼收敛原理知, $\lim\limits_{n\to\infty}x_n=a$ 成立.

必要性\quad 设 $\lim\limits_{n\to\infty}x_n=a$, 则 $\forall\varepsilon>0,\ \exists N$, 当 $n>N$ 时, 有
\[
  a-\varepsilon<x_n<a+\varepsilon,
\]
从而有
\[
  a-\varepsilon\leqslant\ell_n\leqslant x_n\leqslant h_n\leqslant a+\varepsilon,\qquad \forall n>N.
\]
这说明
\[
  \underline{\lim_{n\to\infty}}\,x_n
  =
  \lim_{n\to\infty}\ell_n
  =
  a
  =
  \lim_{n\to\infty}h_n
  =
  \overline{\lim_{n\to\infty}}\,x_n.
\qedhere
\]
\end{proof}


\begin{example}\label{example:ma-2-5-2}
设 $\{x_n\}$ 是 $[0,1]$ 中全体有理数构成的序列, 试求它的上极限与下极限.
\end{example}


\begin{solve}[{解法 1}]
由于 $\{x_n\}$ 在 $[0,1]$ 稠密, 因此, 对任意的 $n\in\mathbb{N}$, $\{x_n,x_{n+1},x_{n+2},\cdots\}$ 仍在 $[0,1]$ 稠密。由此, $\forall n\in\mathbb{N}$, 有 $\ell_n=0,\ h_n=1$, 从而有
\[
  \underline{\lim_{n\to\infty}}\,x_n=0,
  \qquad
  \overline{\lim_{n\to\infty}}\,x_n=1.
\qedhere
\]
\end{solve}


\begin{solve}[{解法 2}]
由于 $\lim\limits_{k\to\infty}\dfrac1k=0,\ \lim\limits_{k\to\infty}\dfrac{k-1}{k}=1$, 因此 $0$ 与 $1$ 都是 $\{x_n\}$ 的聚点。显然 $0$ 是它的最小聚点, $1$ 是它的最大聚点。这说明: $\underline{\lim\limits_{n\to\infty}}\,x_n=0,\ \overline{\lim\limits_{n\to\infty}}\,x_n=1$.

下面的定理给出了上、下极限的保序性和运算性质.
\end{solve}


\begin{theorem}\label{theorem:ma-2-5-3}
设 $\{x_n\}$ 和 $\{y_n\}$ 是任意给定的两个有界序列, 则有

(1) $x_n\leqslant y_n(n=1,2,\cdots)$ 蕴涵着
\[
  \underline{\lim_{n\to\infty}}\,x_n
  \leqslant
  \underline{\lim_{n\to\infty}}\,y_n,
  \qquad
  \overline{\lim_{n\to\infty}}\,x_n
  \leqslant
  \overline{\lim_{n\to\infty}}\,y_n;
\]

(2)
\[
  \underline{\lim_{n\to\infty}}\,(-x_n)
  =
  -\overline{\lim_{n\to\infty}}\,x_n,
  \qquad
  \overline{\lim_{n\to\infty}}\,(-x_n)
  =
  -\underline{\lim_{n\to\infty}}\,x_n;
\]

(3)
\[
  \underline{\lim_{n\to\infty}}\,x_n
  +
  \underline{\lim_{n\to\infty}}\,y_n
  \leqslant
  \underline{\lim_{n\to\infty}}\,(x_n+y_n)
  \leqslant
  \underline{\lim_{n\to\infty}}\,x_n
  +
  \overline{\lim_{n\to\infty}}\,y_n,
\]
\[
  \overline{\lim_{n\to\infty}}\,x_n
  +
  \underline{\lim_{n\to\infty}}\,y_n
  \leqslant
  \overline{\lim_{n\to\infty}}\,(x_n+y_n)
  \leqslant
  \overline{\lim_{n\to\infty}}\,x_n
  +
  \overline{\lim_{n\to\infty}}\,y_n;
\]

(4) 若 $x_n\geqslant0,\ y_n\geqslant0,\ n=1,2,\cdots$, 则
\[
  \underline{\lim_{n\to\infty}}\,x_n\cdot
  \underline{\lim_{n\to\infty}}\,y_n
  \leqslant
  \underline{\lim_{n\to\infty}}\,(x_n\cdot y_n)
  \leqslant
  \underline{\lim_{n\to\infty}}\,x_n\cdot
  \overline{\lim_{n\to\infty}}\,y_n,
\]
\[
  \overline{\lim_{n\to\infty}}\,x_n\cdot
  \underline{\lim_{n\to\infty}}\,y_n
  \leqslant
  \overline{\lim_{n\to\infty}}\,(x_n\cdot y_n)
  \leqslant
  \overline{\lim_{n\to\infty}}\,x_n\cdot
  \overline{\lim_{n\to\infty}}\,y_n.
\]
\end{theorem}


\begin{proof}
我们只证明 (3) 中的前两个不等式, 其余的不等式留给读者作为练习.

取 $\{x_n+y_n\}$ 的一个收敛子列 $\{x_{n_k}+y_{n_k}\}$, 使得
\[
  \lim_{k\to\infty}(x_{n_k}+y_{n_k})
  =
  \underline{\lim_{n\to\infty}}\,(x_n+y_n).
\]
记 $\ell'_n=\inf\{x_n,x_{n+1},\cdots\},\ \ell''_n=\inf\{y_n,y_{n+1},\cdots\}$, 则
\[
  \ell'_{n_k}+\ell''_{n_k}\leqslant x_{n_k}+y_{n_k}.
\]
上式两边令 $k\to\infty$, 取极限即得
\[
  \underline{\lim_{n\to\infty}}\,x_n
  +
  \underline{\lim_{n\to\infty}}\,y_n
  =
  \lim_{k\to\infty}(\ell'_{n_k}+\ell''_{n_k})
  \leqslant
  \lim_{k\to\infty}(x_{n_k}+y_{n_k})
  =
  \underline{\lim_{n\to\infty}}\,(x_n+y_n).
\]
取 $\{x_n\}$ 的子列 $\{x_{n_k}\}$, 使得 $\lim\limits_{k\to\infty}x_{n_k}=\underline{\lim\limits_{n\to\infty}}\,x_n$。由定理~\ref{theorem:ma-2-5-2}(2) 知
\[
  \underline{\lim_{n\to\infty}}\,(x_n+y_n)
  \leqslant
  \underline{\lim_{k\to\infty}}\,(x_{n_k}+y_{n_k}).
\]
再取 $\{y_{n_k}\}$ 的收敛子列 $\{y_{n'_k}\}$。又由定理~\ref{theorem:ma-2-5-2}(2) 及 (3) 得
\[
\begin{aligned}
  \underline{\lim_{k\to\infty}}\,(x_{n_k}+y_{n_k})
  &\leqslant
  \underline{\lim_{k\to\infty}}\,(x_{n'_k}+y_{n'_k})\\
  &=
  \lim_{k\to\infty}x_{n'_k}+\lim_{k\to\infty}y_{n'_k}\\
  &=
  \underline{\lim_{n\to\infty}}\,x_n
  +
  \overline{\lim_{n\to\infty}}\,y_n.
\end{aligned}
\]

作为本节的结束, 我们给出两个例子来说明如何用上、下极限的理论来研究序列的敛散性.
\end{proof}


\begin{example}\label{example:ma-2-5-3}
给定序列 $\{x_n\}$ 及常数 $\alpha\geqslant2$, 令
\[
  y_n=x_n+\alpha x_{n+1},\qquad n=1,2,\cdots,
\]
证明: 序列 $\{y_n\}$ 收敛的充分必要条件是序列 $\{x_n\}$ 收敛.
\end{example}


\begin{proof}
充分性是显然的, 下证必要性.

首先, $\{y_n\}$ 收敛蕴涵着它有界, 从而存在 $M>0$, 使得 $|y_n|\leqslant M$ 和 $|x_1|\leqslant M$。现假定 $|x_n|\leqslant M$, 则有
\[
  |x_{n+1}|\leqslant\frac1\alpha(|y_n|+|x_n|)\leqslant\frac{2M}{\alpha}\leqslant M.
\]
由归纳法原理知, 对一切的正整数 $n$, 有 $|x_n|\leqslant M$, 即 $\{x_n\}$ 有界。记 $\ell=\underline{\lim\limits_{n\to\infty}}\,x_n,\ h=\overline{\lim\limits_{n\to\infty}}\,x_n$, 则 $\ell$ 和 $h$ 均为有限数.

令 $\lim\limits_{n\to\infty}y_n=b$, 在等式 $x_n=y_n-\alpha x_{n+1}$ 两边分别取上、下极限, 并且利用本章习题第 32 题的结论和定理~\ref{theorem:ma-2-5-3} 的 (2), 可得
\[
  \begin{cases}
    h=b-\alpha\ell,\\
    \ell=b-\alpha h.
  \end{cases}
\]
由此可得 $\ell+\alpha h=h+\alpha\ell$, 从而有 $\ell=h$。这样, 由定理~\ref{theorem:ma-2-5-2} 的 (3) 知, 序列 $\{x_n\}$ 收敛.
\end{proof}


\begin{example}\label{example:ma-2-5-4}
证明序列 $\{\sin n\}$ 发散.
\end{example}


\begin{proof}
对 $\forall k\in\mathbb{N}$, 令
\[
  n_k=\left[2k\pi+\frac{3\pi}{4}\right],
  \qquad
  m_k=\left[(2k+1)\pi+\frac{3\pi}{4}\right],
\]
则有
\[
  2k\pi+\frac\pi4<n_k<2k\pi+\frac{3\pi}{4},
  \qquad
  (2k+1)\pi+\frac\pi4<m_k<(2k+1)\pi+\frac{3\pi}{4},
\]
从而
\[
  \sin n_k>\frac{\sqrt2}{2},
  \qquad
  \sin m_k<-\frac{\sqrt2}{2}.
\]
这样, 我们有
\[
  \underline{\lim_{n\to\infty}}\,\sin n
  \leqslant
  \underline{\lim_{k\to\infty}}\,\sin m_k
  \leqslant
  -\frac{\sqrt2}{2}
  <
  \frac{\sqrt2}{2}
  \leqslant
  \overline{\lim_{k\to\infty}}\,\sin n_k
  \leqslant
  \overline{\lim_{n\to\infty}}\,\sin n.
\]
因此, 序列 $\{\sin n\}$ 发散.
\end{proof}


\section*{习题二}
\phantomsection
\addcontentsline{toc}{section}{习题二}
\label{exercises:ma-2}

1。试将有理数集 $\mathbb{Q}$ 排成一个序列.

2。试问以下陈述是否可作为序列极限的定义, 并说明理由:

(1) 对 $\varepsilon=\dfrac1{10^{100}}$, $\exists N>0$, 当 $n>N$ 时, 有 $|x_n-a|<\varepsilon$;

(2) 对 $\forall\varepsilon_k=\dfrac1k(k\in\mathbb{N})$, $\exists N>0$, 当 $n>N$ 时, 有 $|x_n-a|<\varepsilon_k$;

(3) $\exists M>0,\ \forall\varepsilon>0,\ \exists N>0$, 当 $n>N$ 时, 有 $|x_n-a|<M\varepsilon$.

3。用序列极限的定义证明:

(1) $\lim\limits_{n\to\infty}\dfrac{\cos n}{n}=0$;\qquad
(2) $\lim\limits_{n\to\infty}\dfrac{n}{n^3+1}=0$;

% 整合来源：math-3.tex
% \chapter{函数的极限与连续性}
% \label{chap:ma-3}
% 
% \section{函数的极限}
% \label{sec:ma-3-1}
% 
% 上一章我们讨论了序列的极限，本节我们将进一步论述函数的极限。若将序列看成是定义在 $\mathbb{N}$ 中的函数，则在序列极限中，自变量只有一种变化状态。而在函数极限中的自变量却有六种不同的变化状态，这就导致函数极限的定义有各种不同的形式。为了避免繁琐，本节我们主要以自变量趋向于某一给定的点为规范来论述函数极限的定义和性质.
% 
% \subsection{函数极限的定义}
% \label{subsec:ma-3-1-1}
% 
% 在很多实际应用中，我们常常要分析函数的变化情况。特别地，当函数 $f(x)$ 在某点 $x_0$ 没有意义，但在 $U_0(x_0,\delta)$ 内有意义时，我们要研究当 $x$ 充分接近 $x_0$ 时函数的变化情况。例如，考查函数 $y=x\sin\dfrac1x$ 在点 $x=0$ 的附近的变化情况。我们发现 $\sin\dfrac1x$ 在点 $x=0$ 的邻域内是起伏不定的，但当 $x$ 趋于 $0$ 时，由于
% \[
%   |y-0|=\left|x\sin\frac1x\right|\leqslant |x|,
% \]
% $y$ 也随之与 $0$ 无限接近。因此，我们引入函数极限的定义.
% 
% \begin{definition}\label{definition:ma-3-1-1}
% 设函数 $f(x)$ 在 $U_0(x_0,\delta_0)(\delta_0>0)$ 内有定义。若存在实数 $A$，使得 $\forall\varepsilon>0$，$\exists\delta>0$，当 $x\in U_0(x_0,\delta)$ (即 $0<|x-x_0|<\delta$) 时，有 $|f(x)-A|<\varepsilon$，则称当 $x$ 趋于 $x_0$ 时，函数 $f(x)$ 以 $A$ 为极限 (又称函数 $f(x)$ 在点 $x_0$ 的极限存在，其极限为 $A$)，记为
% \[
%   \lim_{x\to x_0} f(x)=A
% \]
% 或者 $f(x)\to A\ (x\to x_0)$.
% \end{definition}
% 
% 从以上定义中可以看出，$f(x)$ 在 $x_0$ 存在极限 $A$ 与 $f(x)$ 在 $x_0$ 是否有定义无关，即使 $f(x)$ 在 $x_0$ 有定义，$A$ 也不一定等于 $f(x_0)$。换句话说，在考虑函数极限时，我们只考虑 $x_0$ 附近 $f(x)$ 的变化趋势，而并不关心 $f(x)$ 在 $x_0$ 的取值情况.
% 
% 再来看 $\lim\limits_{x\to x_0} f(x)=A$ 的几何意义。任给 $\varepsilon>0$，用平行于 $x$ 轴的直线 $y=A-\varepsilon$ 和 $y=A+\varepsilon$ 作一长条带域。由定义，存在 $\delta>0$，使得当 $0<|x-x_0|<\delta$ 时，$|f(x)-A|<\varepsilon$。这就是说，在 $x$ 轴上可以找到一个以 $x_0$ 为中心的开区间 $(x_0-\delta,x_0+\delta)$，使当 $x\in(x_0-\delta,x_0+\delta)$ 且 $x\ne x_0$ 时，点 $P(x,f(x))$ 就必然落在所述的长条阴影带域内 (参见图 3.1.1).
% 
% \begin{example}\label{example:ma-3-1-1}
% 证明极限
% \[
%   \lim_{x\to1}\frac{\sqrt{x}-1}{x-1}=\frac12.
% \]
% \end{example}
% 
% 
% \begin{analyze}
% 当 $x\ne1$ 时，有
% \[
%   \left|\frac{\sqrt{x}-1}{x-1}-\frac12\right|
%   =\left|\frac{x-1}{2(\sqrt{x}+1)^2}\right|.
% \]
% 由于我们只考虑函数在 $x=1$ 附近的变化情况，因此我们无妨限定 $|x-1|<1$，此时 $(\sqrt{x}+1)^2>1$，从而
% \[
%   \left|\frac{\sqrt{x}-1}{x-1}-\frac12\right|
%   =\frac{|x-1|}{2(\sqrt{x}+1)^2}<\frac12|x-1|.
% \]
% 这样，$\forall\varepsilon>0$，欲使
% \[
%   \left|\frac{\sqrt{x}-1}{x-1}-\frac12\right|<\varepsilon,
% \]
% 只需取 $\delta=2\varepsilon$.
% \end{analyze}
% 
% 
% \begin{proof}
% $\forall\varepsilon>0$，取 $\delta=\min\{1,2\varepsilon\}$，则当 $|x-1|<\delta$ 时，有
% \[
%   \left|\frac{\sqrt{x}-1}{x-1}-\frac12\right|
%   <\frac12|x-1|<\varepsilon.
% \]
% 按函数极限的定义知 $\lim\limits_{x\to1}\dfrac{\sqrt{x}-1}{x-1}=\dfrac12$ 成立.
% \end{proof}
% 
% 
% \begin{example}\label{example:ma-3-1-2}
% 证明极限 $\lim\limits_{x\to2}x^4=16$.
% \end{example}
% 
% 
% \begin{proof}
% 由于
% \[
%   |x^4-16|=|x-2||x^3+2x^2+4x+8|,
% \]
% 限定 $|x-2|<1$ 时，有 $|x|<3$，从而有
% \[
%   |x^4-16|<65|x-2|.
% \]
% 因此，$\forall\varepsilon>0$，取 $\delta=\min\left\{\dfrac{\varepsilon}{65},1\right\}$，则当 $0<|x-2|<\delta$ 时，有
% \[
%   |x^4-16|<\varepsilon.
% \]
% 按函数极限的定义知 $\lim\limits_{x\to2}x^4=16$ 成立.
% 
% 用完全类似的方法可证：对任意的实数 $a$ 和任意的正整数 $n$，有
% \[
%   \lim_{x\to a}x^n=a^n.
% \]
% 请读者作为练习给出这一结果的证明.
% 
% 从以上例子的证明可以看出：要证明 $\lim\limits_{x\to x_0}f(x)=A$，我们应该先限定 $x$ 在 $x_0$ 的某个邻域内，这样容易从 $|f(x)-A|$ 中得到所要的关于 $|x-x_0|$ 的不等式。而预先限定 $x$ 的变化范围，可以通过给定某个 $\delta_0>0$ 得到.
% \end{proof}
% 
% 
% \begin{example}\label{example:ma-3-1-3}
% 证明极限 $\lim\limits_{x\to0}x\sin\dfrac1x=0$.
% \end{example}
% 
% 
% \begin{proof}
% $\forall\varepsilon>0$，取 $\delta=\varepsilon$，当 $x\in U_0(0,\delta)$ 时，有
% \[
%   \left|x\sin\frac1x-0\right|\leqslant |x|<\varepsilon,
% \]
% 按函数极限的定义知 $\lim\limits_{x\to0}x\sin\dfrac1x=0$ 成立.
% \end{proof}
% 
% 
% \subsection{函数极限的性质}
% \label{subsec:ma-3-1-2}
% 
% 与序列极限的性质一样，函数极限也具有相应的性质。但是由于函数极限的存在与否只与 $x_0$ 附近的函数值有关，因此一些相应的结论也只能在 $x_0$ 的某个邻域内成立。例如，对于``改变序列的有限多项不改变序列的敛散性''这一性质而言，我们有以下相应的结论：设 $f(x)$ 在 $x_0$ 的某个邻域 $U_0(x_0,\delta_0)$ ($\delta_0>0$) 有定义，则对任何 $\delta_0>\delta_1>0$，改变 $f(x)$ 在 $U_0(x_0,\delta_1)$ 外面的函数值，不影响 $f(x)$ 在 $x_0$ 处的敛散性.
% 
% \begin{theorem}[{唯一性和有界性}]\label{theorem:ma-3-1-1}
% 设极限 $\lim\limits_{x\to x_0}f(x)$ 存在，则
% \begin{enumerate}
%   \item 该极限值是唯一的；
%   \item $\exists\delta_0>0$，使得 $f(x)$ 在 $U_0(x_0,\delta_0)$ 有界.
% \end{enumerate}
% \end{theorem}
% 
% 
% \begin{proof}
% (1) 设 $\lim\limits_{x\to x_0}f(x)=A$，$\lim\limits_{x\to x_0}f(x)=B$，则由函数极限的定义，$\forall\varepsilon>0$，$\exists\delta>0$，当 $0<|x-x_0|<\delta$ 时，有
% \[
%   |f(x)-A|<\frac{\varepsilon}{2},\qquad |f(x)-B|<\frac{\varepsilon}{2},
% \]
% 从而有
% \[
%   |A-B|<|f(x)-A|+|f(x)-B|<\varepsilon.
% \]
% 由于 $A,B$ 为两个常数，它们的差又能小于任给的 $\varepsilon>0$，因此 $A$ 必须等于 $B$.
% 
% (2) 设 $\lim\limits_{x\to x_0}f(x)=A$。对 $\varepsilon_0=1$，$\exists\delta_0>0$，当 $0<|x-x_0|<\delta_0$ 时，有 $|f(x)-A|<1$，从而有
% \[
%   |f(x)|\leqslant |A|+|f(x)-A|<1+|A|,\qquad \forall x\in U_0(x_0,\delta_0),
% \]
% 按有界函数的定义知 $f(x)$ 在 $U_0(x_0,\delta_0)$ 有界.
% \end{proof}
% 
% 
% \begin{theorem}[{保序性}]\label{theorem:ma-3-1-2}
% 设 $\lim\limits_{x\to x_0}f(x)=A$，$\lim\limits_{x\to x_0}g(x)=B$，则
% \begin{enumerate}
%   \item 若 $\exists\delta_0>0$，使得当 $x\in U_0(x_0,\delta_0)$ 时，有 $f(x)\leqslant g(x)$，则 $A\leqslant B$；
%   \item 若 $A<B$，则对任何 $C\in(A,B)$，$\exists\delta_C>0$，使得当 $x\in U_0(x_0,\delta_C)$ 时，有 $f(x)<C<g(x)$.
% \end{enumerate}
% 定理的证明留给读者作为练习.
% \end{theorem}
% 
% 
% \begin{theorem}[{四则运算}]\label{theorem:ma-3-1-3}
% 设 $\lim\limits_{x\to x_0}f(x)=A$，$\lim\limits_{x\to x_0}g(x)=B$，则有
% \begin{enumerate}
%   \item $\lim\limits_{x\to x_0}(f(x)\pm g(x))=A\pm B$；
%   \item $\lim\limits_{x\to x_0}f(x)\cdot g(x)=A\cdot B$；
%   \item $\lim\limits_{x\to x_0}\dfrac{f(x)}{g(x)}=\dfrac{A}{B}$，这里假定 $B\ne0$.
% \end{enumerate}
% \end{theorem}
% 
% 
% \begin{proof}
% 我们只证 (3)，其余的请读者自己证明.
% 
% 在 $x_0$ 的附近，我们有
% \[
%   \left|\frac{f(x)}{g(x)}-\frac{A}{B}\right|
%   =\frac{|Bf(x)-Ag(x)|}{|Bg(x)|}
%   =\frac{|B(f(x)-A)-A(g(x)-B)|}{|B||g(x)|}
% \]
% \[
%   \leqslant \frac{1}{|B||g(x)|}\bigl(|B||f(x)-A|+|A||g(x)-B|\bigr).
% \]
% 由 $\lim\limits_{x\to x_0}g(x)=B\ne0$ 知，$\exists\delta_0>0$，使得当 $x\in U_0(x_0,\delta_0)$ 时，有
% \[
%   |g(x)-B|<\frac{|B|}{2},
% \]
% 从而
% \[
%   |g(x)|\geqslant |B|-|g(x)-B|\geqslant |B|-\frac{|B|}{2}=\frac{|B|}{2}.
% \]
% 这样，与前面的不等式相结合，我们有
% \[
%   \left|\frac{f(x)}{g(x)}-\frac{A}{B}\right|
%   \leqslant \frac{2}{|B|}|f(x)-A|+2\frac{|A|}{|B|^2}|g(x)-B|
% \]
% 对一切的 $x\in U_0(x_0,\delta_0)$ 成立.
% 
% $\forall\varepsilon>0$，由 $\lim\limits_{x\to x_0}g(x)=B$ 知，$\exists\delta_1>0$，使得当 $x\in U_0(x_0,\delta_1)$ 时，有
% \[
%   |g(x)-B|<\frac{B^2}{4(|A|+1)}\varepsilon;
% \]
% 再由 $\lim\limits_{x\to x_0}f(x)=A$ 知，$\exists\delta_2>0$，使得当 $x\in U_0(x_0,\delta_2)$ 时，有
% \[
%   |f(x)-A|<\frac{|B|}{4}\varepsilon.
% \]
% 
% 现在取 $\delta=\min\{\delta_0,\delta_1,\delta_2\}$，则当 $x\in U_0(x_0,\delta)$ 时，有
% \[
%   \left|\frac{f(x)}{g(x)}-\frac{A}{B}\right|
%   <\frac{\varepsilon}{2}+\frac{\varepsilon}{2}=\varepsilon.
% \]
% 按函数极限的定义知 (3) 成立.
% \end{proof}
% 
% 
% \begin{example}\label{example:ma-3-1-4}
% 求极限
% \[
%   \lim_{x\to2}\frac{2x^3+3x-1}{x^3+3}.
% \]
% \end{example}
% 
% 
% \begin{solve}
% 因 $x\to2$ 时，$2x^3+3x-1$ 和 $x^3+3$ 都有极限，而且
% \[
%   \lim_{x\to2}(2x^3+3x-1)=2\times2^3+3\times2-1=21,
% \]
% \[
%   \lim_{x\to2}(x^3+3)=2^3+3=11\ne0,
% \]
% 所以
% \[
%   \lim_{x\to2}\frac{2x^3+3x-1}{x^3+3}
%   =\frac{\lim\limits_{x\to2}(2x^3+3x-1)}
%   {\lim\limits_{x\to2}(x^3+3)}
%   =\frac{21}{11}.
% \qedhere
% \]
% \end{solve}
% 
% 
% \begin{example}\label{example:ma-3-1-5}
% 求
% \[
%   \lim_{x\to1}\frac{x+x^2+\cdots+x^n-n}{x-1},
% \]
% 其中 $n$ 是一个给定的正整数.
% \end{example}
% 
% 
% \begin{solve}
% 将所求极限式恒等变形，有
% \[
% \begin{aligned}
%   \lim_{x\to1}\frac{x+x^2+\cdots+x^n-n}{x-1}
%   &=\lim_{x\to1}\left(\frac{x-1}{x-1}+\frac{x^2-1}{x-1}+\cdots+\frac{x^n-1}{x-1}\right)\\
%   &=\lim_{x\to1}\left[1+(1+x)+\cdots+(x^{n-1}+x^{n-2}+\cdots+1)\right]\\
%   &=1+2+\cdots+n=\frac12 n(n+1).
% \end{aligned}
% \]
% 
% 下面我们来讨论复合函数求极限的问题，有下述结论：
% \end{solve}
% 
% 
% \begin{theorem}\label{theorem:ma-3-1-4}
% 设函数 $f(u)$ 在 $U_0(u_0,\delta_1)(\delta_1>0)$ 有定义且 $\lim\limits_{u\to u_0}f(u)=A$；$u=g(x)$ 在 $U_0(x_0,\delta_0)(\delta_0>0)$ 有定义，当 $x\in U_0(x_0,\delta_0)$ 时有 $g(x)\in U_0(u_0,\delta_1)$ 且 $\lim\limits_{x\to x_0}g(x)=u_0$，则有
% \[
%   \lim_{x\to x_0}f(g(x))=A.
% \]
% \end{theorem}
% 
% 
% \begin{proof}
% 由 $\lim\limits_{u\to u_0}f(u)=A$ 知，$\forall\varepsilon>0$，$\exists 0<\eta<\delta_1$，使得当 $u\in U_0(u_0,\eta)$ 时，有
% \[
%   |f(u)-A|<\varepsilon.
% \]
% 对于 $\eta>0$，由 $\lim\limits_{x\to x_0}g(x)=u_0$，$\exists 0<\delta<\delta_0$，使得当 $x\in U_0(x_0,\delta)$ 时，有
% \[
%   |g(x)-u_0|<\eta.
% \]
% 注意到 $\forall x\in U_0(x_0,\delta_0)$ 有 $g(x)\ne u_0$，因此当 $x\in U_0(x_0,\delta)$ 时，有 $g(x)\in U_0(u_0,\eta)$，因此有
% \[
%   |f(g(x))-A|<\varepsilon.
% \]
% 这就证明了 $\lim\limits_{x\to x_0}f(g(x))=A$。定理
% \end{proof}
% 
% 
% \begin{example}\label{example:ma-3-1-6}
% 求极限
% \[
%   \lim_{x\to0}\frac{\sqrt{x^2+x+1}-1}{x^2+x}.
% \]
% \end{example}
% 
% 
% \begin{solve}
% 取 $f(u)=\dfrac{\sqrt{u}-1}{u-1}$，其中 $u=x^2+x+1$，则 $f(u)$ 在 $u_0=1$ 的某个邻域，$g(x)$ 在 $x_0=0$ 的某个邻域内满足定理~\ref{theorem:ma-3-1-4} 的所有条件。由例~\ref{example:ma-3-1-1} 知 $\lim\limits_{u\to1}f(u)=\dfrac12$。显然有 $\lim\limits_{x\to0}g(x)=\lim\limits_{x\to0}(x^2+x+1)=1$。由定理~\ref{theorem:ma-3-1-4} 知
% \[
%   \lim_{x\to0}\frac{\sqrt{x^2+x+1}-1}{x^2+x}
%   =\lim_{x\to0}f(g(x))=\frac12.
% \]
% 
% 对于定理~\ref{theorem:ma-3-1-4} 中的条件，我们要求当 $x\in U_0(x_0,\delta_0)$ 时，有 $g(x)\in U_0(u_0,\delta_1)$。该条件在复合函数求极限中是必要的。换句话说，若假定 $g(x)\in U(u_0,\delta_1)$，则定理~\ref{theorem:ma-3-1-4} 的结论可能不真。事实上，我们取
% \[
%   f(u)=
%   \begin{cases}
%     0, & u=0,\\
%     1, & u\ne0,
%   \end{cases}
%   \qquad
%   g(x)=
%   \begin{cases}
%     1, & x=\dfrac1n,\\
%     0, & \text{其他},
%   \end{cases}
% \]
% 则
% \[
%   f(g(x))=
%   \begin{cases}
%     1, & x=\dfrac1n,\\
%     0, & \text{其他}.
%   \end{cases}
% \]
% 现取 $x_0=0$，则 $\lim\limits_{x\to0}g(x)=0$。注意到 $\lim\limits_{u\to0}f(u)=1$，但是 $\lim\limits_{x\to0}f(g(x))$ 不存在。这是因为在 $x_0=0$ 的任意小的邻域内 $f(g(x))$ 既有取值为 $1$ 的点，又有取值为 $0$ 的点，故其极限不存在.
% \end{solve}
% 
% 
% \begin{theorem}[{夹逼收敛原理}]\label{theorem:ma-3-1-5}
% 若 $\lim\limits_{x\to x_0}f(x)=\lim\limits_{x\to x_0}h(x)=A$，而且存在 $\delta_0>0$，使得 $f(x)\leqslant g(x)\leqslant h(x)$ 对一切 $x\in U_0(x_0,\delta)$ 成立，则有
% \[
%   \lim_{x\to x_0}g(x)=A.
% \]
% \end{theorem}
% 
% 此定理的证明与序列的逼收敛原理类似，故从略.
% 
% 最后，我们要重复指出的是，函数极限存在时，它所具有的性质只是在局部成立。例如，设 $f(x)=\dfrac1x$，$x\in(0,+\infty)$。对任意的 $x_0\in(0,+\infty)$，有
% \[
%   \lim_{x\to x_0}f(x)=\lim_{x\to x_0}\frac1x=\frac1{x_0}.
% \]
% 由此可以导出 $f(x)$ 在 $x_0$ 附近有界，事实上，我们有
% \[
%   |f(x)|<\frac{2}{x_0},\qquad \forall x\in\left(x_0-\frac{x_0}{2},x_0+\frac{x_0}{2}\right),
% \]
% 但 $f(x)$ 显然在 $(0,+\infty)$ 上无界.
% 
% \subsection{函数极限概念的推广}
% \label{subsec:ma-3-1-3}
% 
% 在实际应用中，前面所引入的函数极限的概念是不够的，例如对闭区间上有定义的函数，就无法讨论其在端点处的极限问题。因此我们就必须拓广极限的概念，以便处理其他形式的极限问题.
% 
% \paragraph*{单侧极限}
% 
% 
% 所谓\keyterm{单侧极限}是考虑函数在某一点的一侧的变化情况。记
% \[
%   U^+(x_0,\delta)=\{x\in\mathbb{R}:x_0\leqslant x<x_0+\delta\}\qquad (x_0\text{ 的右邻域});
% \]
% \[
%   U^-(x_0,\delta)=\{x\in\mathbb{R}:x_0-\delta<x\leqslant x_0\}\qquad (x_0\text{ 的左邻域});
% \]
% \[
%   U_0^+(x_0,\delta)=U^+(x_0,\delta)\setminus\{x_0\}\qquad (x_0\text{ 的右空心邻域});
% \]
% \[
%   U_0^-(x_0,\delta)=U^-(x_0,\delta)\setminus\{x_0\}\qquad (x_0\text{ 的左空心邻域}).
% \]
% 
% \begin{definition}\label{definition:ma-3-1-2}
% 设 $f(x)$ 在 $U_0^+(x_0,\delta_0)(\delta_0>0)$ 上有定义。如果存在实数 $A$，对 $\forall\varepsilon>0$，$\exists\delta>0$，使得当 $x\in U_0^+(x_0,\delta)$ (即 $0<x-x_0<\delta$) 时，有 $|f(x)-A|<\varepsilon$，则称 $f(x)$ 在点 $x_0$ 的\keyterm{右极限存在}，而称 $A$ 为 $f(x)$ 在点 $x_0$ 的\keyterm{右极限}，记为
% \[
%   \lim_{x\to x_0+0}f(x)=A
% \]
% 或者 $f(x_0+0)=A$.
% \end{definition}
% 
% 类似地可定义 $f(x)$ 在点 $x_0$ 的\keyterm{左极限} $\lim\limits_{x\to x_0-0}f(x)$ 或者 $f(x_0-0)$。有关右极限的几何解释参见图 3.1.2.
% 
% \begin{theorem}\label{theorem:ma-3-1-6}
% 函数 $f(x)$ 在点 $x_0$ 极限存在的充分必要条件是 $f(x)$ 在点 $x_0$ 的左、右极限都存在且相等.
% \end{theorem}
% 
% 
% \begin{proof}
% 必要性\quad 设 $\lim\limits_{x\to x_0}f(x)=A$ 存在，则 $\forall\varepsilon>0$，$\exists\delta>0$，使得当 $0<|x-x_0|<\delta$ 时，有
% \[
%   |f(x)-A|<\varepsilon.
% \]
% 特别地，当 $0<x-x_0<\delta$ 或 $0<x_0-x<\delta$ 时，亦有 $|f(x)-A|<\varepsilon$，故 $\lim\limits_{x\to x_0+0}f(x)$ 和 $\lim\limits_{x\to x_0-0}f(x)$ 都存在，而且
% \[
%   \lim_{x\to x_0+0}f(x)=\lim_{x\to x_0-0}f(x)=A.
% \]
% 
% 充分性\quad 设 $\lim\limits_{x\to x_0+0}f(x)$ 和 $\lim\limits_{x\to x_0-0}f(x)$ 都存在，而且
% \[
%   \lim_{x\to x_0+0}f(x)=\lim_{x\to x_0-0}f(x)=A,
% \]
% 则 $\forall\varepsilon>0$，由 $\lim\limits_{x\to x_0+0}f(x)=A$ 知，$\exists\delta_1>0$，使得当 $0<x-x_0<\delta_1$ 时，有 $|f(x)-A|<\varepsilon$；又由 $\lim\limits_{x\to x_0-0}f(x)=A$ 知，$\exists\delta_2>0$，使得当 $0<x_0-x<\delta_2$ 时，有 $|f(x)-A|<\varepsilon$.
% 
% 令 $\delta=\min\{\delta_1,\delta_2\}$，则当 $0<|x-x_0|<\delta$ 时，有 $|f(x)-A|<\varepsilon$。故 $\lim\limits_{x\to x_0}f(x)$ 存在，而且 $\lim\limits_{x\to x_0}f(x)=A$.
% \end{proof}
% 
% 
% \begin{example}\label{example:ma-3-1-7}
% 证明取整函数 $f(x)=[x]$ 在 $x=3$ 的极限不存在.
% \end{example}
% 
% 
% \begin{proof}
% 由取整函数的定义可证
% \[
%   f(3+0)=3>2=f(3-0),
% \]
% 因此，由定理~\ref{theorem:ma-3-1-5} 知，取整函数 $f(x)=[x]$ 在 $x=3$ 的极限不存在.
% \end{proof}
% 
% 
% \paragraph*{自变量趋向无穷大时的极限}
% 
% 
% 如果将序列看成是定义在 $\mathbb{N}$ 上的函数的话，则序列极限就是自变量 $n\to+\infty$ 时函数的极限，类似地也可以考虑在实数集 $\mathbb{R}$ 上定义的函数当 $x$ 趋于 $\infty$ 时的极限问题.
% 
% 称集合 $\{x:|x|>h\}$ ($h>0$) 为 $\infty$ 的邻域，记为 $U(\infty,h)$ 或 $U(\infty)$；称 $\{x:h<x<+\infty\}$ 与 $\{x:-\infty<x<-h\}$ 为 $\infty$ 的单侧邻域，分别记为 $U^+(\infty,h)(U^+(\infty)\text{ 或 }U(+\infty))$ 和 $U^-(\infty,h)(U^-(\infty)\text{ 或 }U(-\infty))$.
% 
% \begin{definition}\label{definition:ma-3-1-3}
% 设函数 $f(x)$ 在 $U^+(\infty)$ 上有定义。若存在实数 $A$，使得 $\forall\varepsilon>0$，$\exists X\in U^+(\infty)$，当 $x>X$ 时，有 $|f(x)-A|<\varepsilon$，则称当 $x$ 趋于 $+\infty$ 时 $f(x)$ 的极限存在，其极限为 $A$，记为 $\lim\limits_{x\to+\infty}f(x)=A$ 或者 $f(x)\to A\ (x\to+\infty)$.
% \end{definition}
% 
% 类似地可以定义 $\lim\limits_{x\to-\infty}f(x)=A$ 和 $\lim\limits_{x\to\infty}f(x)=A$。有关 $x\to+\infty$ 时 $f(x)$ 的极限的几何解释参见图 3.1.3.
% 
% 对这种类型的极限，也有类似定理~\ref{theorem:ma-3-1-6} 结论，另外，这类极限也有四则运算及复合函数相应的结果。请读者作为练习给出它们的表述和证明.
% 
% \begin{example}\label{example:ma-3-1-8}
% 证明极限
% \[
%   \lim_{x\to\infty}\frac{3x^2+2x-1}{x^2-2}=3.
% \]
% \end{example}
% 
% 
% \begin{proof}
% 由于
% \[
%   \left|\frac{3x^2+2x-1}{x^2-2}-3\right|
%   =\left|\frac{2x+5}{x^2-2}\right|
%   \leqslant\frac{3|x|}{3x^2/4}=\frac{4}{|x|}
% \]
% 对一切的 $|x|>5$ 成立，所以 $\forall\varepsilon>0$，取 $X=\max\{5,4/\varepsilon\}$，则当 $|x|>X$ 时，有
% \[
%   \left|\frac{3x^2+2x-1}{x^2-2}-3\right|\leqslant\frac4{|x|}<\varepsilon.
% \]
% 按极限定义知 $\lim\limits_{x\to\infty}\dfrac{3x^2+2x-1}{x^2-2}=3$ 成立.
% \end{proof}
% 
% 
% \begin{example}\label{example:ma-3-1-9}
% 设 $f(x)=\sqrt{4x^2+2x+3}$，试求下列极限：
% \[
%   (1)\ \lim_{x\to+\infty}\frac{f(x)}{x};\qquad
%   (2)\ \lim_{x\to-\infty}\frac{f(x)}{x};
% \]
% \[
%   (3)\ \lim_{x\to+\infty}(f(x)-2x);\qquad
%   (4)\ \lim_{x\to-\infty}(f(x)+2x).
% \]
% \end{example}
% 
% 
% \begin{solve}
% (1)
% \[
%   \lim_{x\to+\infty}\frac{f(x)}{x}
%   =\lim_{x\to+\infty}\sqrt{4+\frac2x+\frac3{x^2}}=2;
% \]
% (2)
% \[
%   \lim_{x\to-\infty}\frac{f(x)}{x}
%   =\lim_{t\to+\infty}\frac{\sqrt{4t^2-2t+3}}{-t}=-2;
% \]
% (3)
% \[
%   \lim_{x\to+\infty}(f(x)-2x)
%   =\lim_{x\to+\infty}\frac{4x^2+2x+3-4x^2}{\sqrt{4x^2+2x+3}+2x}
%   =\frac12;
% \]
% (4)
% \[
% \begin{aligned}
%   \lim_{x\to-\infty}(f(x)+2x)
%   &=\lim_{x\to-\infty}\frac{4x^2+2x+3-4x^2}{\sqrt{4x^2+2x+3}-2x}\\
%   &=\lim_{t\to+\infty}\frac{-2t+3}{\sqrt{4t^2-2t+3}+2t}
%   =-\frac12.
% \end{aligned}
% \qedhere
% \]
% \end{solve}
% 
% 
% \paragraph*{广义极限}
% 
% 
% 当函数 $f(x)$ 在点 $x_0$ 的极限不存在时，这时 $f(x)$ 在 $x_0$ 附近的变化情况非常复杂。但所有在点 $x_0$ 不存在极限的函数中，有一类变化还是具有规律的。如函数 $f(x)=\dfrac1x$，当 $x\to0$ 时，函数的绝对值将目标一致地无限变大.
% 
% \begin{definition}\label{definition:ma-3-1-4}
% 设 $f(x)$ 在 $U_0(x_0,\delta_0)(\delta_0>0)$ 上有定义。若 $\forall M>0$，$\exists\delta>0$，使得当 $0<|x-x_0|<\delta$ 时，有 $f(x)>M$，则称当 $x$ 趋于 $x_0$ 时，$f(x)$ 趋于 $+\infty$，或称当 $x$ 趋于 $x_0$ 时，$f(x)$ 的广义极限为 $+\infty$，并记为 $\lim\limits_{x\to x_0}f(x)=+\infty$ 或 $f(x)\to+\infty\ (x\to x_0)$。此时，也称 $f(x)$ 为当 $x$ 趋于 $x_0$ 时的\keyterm{正无穷大量}.
% \end{definition}
% 
% 类似地可以定义 $\lim\limits_{x\to x_0}f(x)=-\infty$ (称 $f(x)$ 为当 $x$ 趋于 $x_0$ 时的\keyterm{负无穷大量})，$\lim\limits_{x\to x_0}f(x)=\infty$ (称 $f(x)$ 为当 $x$ 趋于 $x_0$ 时的\keyterm{无穷大量})，$\lim\limits_{x\to x_0+0}f(x)=+\infty$ (称 $f(x)$ 为当 $x$ 从右侧趋于 $x_0$ 时的\keyterm{正无穷大量})，$\lim\limits_{x\to+\infty}f(x)=+\infty$ (称 $f(x)$ 为当 $x$ 趋于 $+\infty$ 时的\keyterm{正无穷大量}) 等.
% 
% 注意无穷大量并不是函数值可以取到任意大的值的量。例如，函数 $f(x)=x^2\sin x$ 当 $x\to\infty$ 时可以取到任意大的值，但是它并不是当 $x\to\infty$ 时的无穷大量.
% 
% 对于函数极限而言，自变量 $x$ 共有六种可能的变化情况，即
% \[
%   x\to x_0;\quad x\to x_0+0;\quad x\to x_0-0;\quad x\to\infty;\quad x\to\pm\infty.
% \]
% 对于自变量 $x$ 的每一种趋向，函数又有四种不同的趋向，即
% \[
%   f(x)\to A;\quad f(x)\to\pm\infty;\quad f(x)\to\infty.
% \]
% 因此在研究函数极限时，就有 24 种可能的情形要研究。对于
% \[
%   \lim_{x\to x_0\pm0}f(x)=A,\qquad
%   \lim_{x\to\infty}f(x)=A,\qquad
%   \lim_{x\to\pm\infty}f(x)=A
% \]
% 这五种极限，也相应地有上一小节中所介绍的所有性质；而对于其他类型的极限，在不发生 $(\pm\infty)+(\mp\infty)$，$0\cdot\infty$，$\dfrac00$ 和 $\dfrac{\infty}{\infty}$ 这些类型 (称为不定式) 情况下，也有相应的四则运算公式.
% 
% \begin{example}\label{example:ma-3-1-10}
% 求极限
% \[
%   \lim_{x\to\infty}\frac{5-2x^2}{3x^2+x-2}.
% \]
% \end{example}
% 
% 
% \begin{solve}
% 利用极限的四则运算，可得
% \[
% \begin{aligned}
%   \lim_{x\to\infty}\frac{5-2x^2}{3x^2+x-2}
%   &=\lim_{x\to\infty}\frac{\dfrac5{x^2}-2}{3+\dfrac1x-\dfrac2{x^2}}\\
%   &=\frac{\lim\limits_{x\to\infty}\left(\dfrac5{x^2}-2\right)}
%   {\lim\limits_{x\to\infty}\left(3+\dfrac1x-\dfrac2{x^2}\right)}
%   =-\frac23.
% \end{aligned}
% \qedhere
% \]
% \end{solve}
% 
% 
% \begin{example}\label{example:ma-3-1-11}
% 设 $f(x)=\dfrac1{x-1}$，求 $f(1-0)$，$f(1+0)$ 和 $\lim\limits_{x\to1}f(x)$.
% \end{example}
% 
% 
% \begin{solve}
% $f(1-0)=-\infty$，$f(1+0)=+\infty$ 和 $\lim\limits_{x\to1}f(x)=\infty$.
% \end{solve}
% 
% 
% \subsection{序列极限与函数极限的关系}
% \label{subsec:ma-3-1-4}
% 
% 以 $x$ 趋于 $x_0$ 的情形为例，我们先来给出 $f(x)$ 不趋于 $A$ 的精确描述。如同序列情形，$f(x)$ 不趋于 $A$，则必定 $\exists\varepsilon_0>0$，在 $x_0$ 的附近有充分接近 $x_0$ 的点 $x'$，使得 $|f(x')-A|\geqslant\varepsilon_0$。读者应该充分理解``有充分接近 $x_0$ 的点 $x'$''的精确意义，此即是说，在 $x_0$ 的任何邻域内都存在异于 $x_0$ 的点 $x'$，使得 $|f(x')-A|\geqslant\varepsilon_0$。因此，用肯定的语气来叙述 $\lim\limits_{x\to x_0}f(x)\ne A$，即为
% \[
%   \exists\varepsilon_0>0,\ \forall\delta>0,\ \exists x'\in U_0(x_0,\delta),\ \text{使得 } |f(x')-A|\geqslant\varepsilon_0.
% \]
% 我们这里须特别指出的是，初学者容易犯的一个错误是，将其叙述成：
% \[
%   \exists\varepsilon_0>0,\ \forall\delta>0,\ \text{当 }x\in U_0(x_0,\delta)\text{ 时，有 } |f(x)-A|\geqslant\varepsilon_0.
% \]
% 当然，满足以上条件一定有 $x$ 趋于 $x_0$ 时 $f(x)$ 不趋于 $A$，但并不是所有的 $x$ 趋于 $x_0$ 时 $f(x)$ 不以 $A$ 为极限的都有以上情形发生.
% 
% 下面我们来讨论序列极限与函数极限的关系问题。事实上，我们有
% 
% \begin{theorem}\label{theorem:ma-3-1-7}
% 设 $f(x)$ 在 $U_0(x_0,\delta_0)(\delta_0>0)$ 上有定义，则 $\lim\limits_{x\to x_0}f(x)=A$ 成立的充分必要条件是：对于 $U_0(x_0,\delta_0)$ 内任意收敛于 $x_0$ 的序列 $\{x_n\}$，都有
% \[
%   \lim_{n\to\infty}f(x_n)=A.
% \]
% \end{theorem}
% 
% 
% \begin{proof}
% 必要性\quad 在 $U_0(x_0,\delta_0)$ 中任取一收敛于 $x_0$ 的序列 $\{x_n\}$，即 $\lim\limits_{n\to\infty}x_n=x_0$，要证 $\lim\limits_{n\to\infty}f(x_n)=A$.
% 
% $\forall\varepsilon>0$，由 $\lim\limits_{x\to x_0}f(x)=A$ 知，$\exists\delta>0$，使得当 $0<|x-x_0|<\delta$ 时，有
% \[
%   |f(x)-A|<\varepsilon.
% \]
% 对于上述 $\delta>0$，由 $x_n\to x_0\ (n\to\infty)$ 知，$\exists N$，使得当 $n>N$ 时，有
% \[
%   0<|x_n-x_0|<\delta.
% \]
% 于是，当 $n>N$ 时，有 $|f(x_n)-A|<\varepsilon$，即 $\lim\limits_{n\to\infty}f(x_n)=A$.
% 
% 充分性\quad 如果不然，即当 $x\to x_0$ 时 $f(x)$ 不以 $A$ 为极限，则 $\exists\varepsilon_0>0$，对 $\forall\delta>0$，$\exists x_\delta\in U_0(x_0,\delta)$，使得
% \[
%   |f(x_\delta)-A|\geqslant\varepsilon_0.
% \]
% 这样，分别取 $\delta=\dfrac{\delta_0}{n}\ (n=1,2,\cdots)$，$\exists x_n\in U_0\left(x_0,\dfrac{\delta_0}{n}\right)$，使得
% \[
%   |f(x_n)-A|\geqslant\varepsilon_0.
% \]
% 对于这样得到的序列 $\{x_n\}$，显然有 $x_n\to x_0\ (n\to\infty)$，且 $x_n\in U_0(x_0,\delta_0)$，从而 $\lim\limits_{n\to\infty}f(x_n)=A$，但这与 $|f(x_n)-A|\geqslant\varepsilon_0$ 对一切正整数 $n$ 成立相矛盾.
% \end{proof}
% 
% 
% \begin{example}\label{example:ma-3-1-12}
% 试证明当 $x\to0$ 时函数 $f(x)=\sin\dfrac1x\ (x\ne0)$ 的极限不存在.
% \end{example}
% 
% 
% \begin{proof}
% 取
% \[
%   x_n=\frac1{n\pi},\qquad n=1,2,\cdots,
% \]
% 则有 $x_n\to0\ (n\to\infty)$，且
% \[
%   f(x_n)=\sin n\pi\equiv0\to0\qquad (n\to\infty);
% \]
% 再取
% \[
%   \widetilde{x}_n=\frac1{2n\pi+\dfrac{\pi}{2}},\qquad n=1,2,\cdots,
% \]
% 则有 $\widetilde{x}_n\to0\ (n\to\infty)$，且
% \[
%   f(\widetilde{x}_n)=\sin\left(2n\pi+\frac{\pi}{2}\right)\equiv1\to1\qquad (n\to\infty).
% \]
% 这样，由定理~\ref{theorem:ma-3-1-7} 知 $\lim\limits_{x\to0}\sin\dfrac1x$ 不存在.
% 
% 值得注意的是，对于此例，虽然当 $x\to0$ 时 $\sin\dfrac1x$ 不趋向于 $0$，但找不到 $\varepsilon_0>0$，使得对任何 $\delta>0$，都有 $\left|\sin\dfrac1x-0\right|\geqslant\varepsilon_0$ 在 $U_0(x_0,\delta)$ 上成立.
% \end{proof}
% 
% 
% \subsection{极限存在性定理和两个重要极限}
% \label{subsec:ma-3-1-5}
% 
% \paragraph*{极限存在性定理}
% 
% 
% 如同序列极限一样，如果函数单调，则它必存在广义极限.
% 
% \begin{theorem}\label{theorem:ma-3-1-8}
% 设函数 $f(x)$ 在 $U_0^+(x_0,\delta_0)$ 内有定义，若 $f(x)$ 在 $U_0^+(x_0,\delta_0)$ 上内单调上升，则
% \[
%   \lim_{x\to x_0+0}f(x)=\inf\{f(x):x\in U_0^+(x_0,\delta_0)\};
% \]
% 若 $f(x)$ 在 $U_0^+(x_0,\delta_0)$ 内单调下降，则
% \[
%   \lim_{x\to x_0+0}f(x)=\sup\{f(x):x\in U_0^+(x_0,\delta_0)\}.
% \]
% \end{theorem}
% 
% 
% \begin{proof}
% 我们只给出定理第一部分的证明。记
% \[
%   A=\inf\{f(x):x\in U_0^+(x_0,\delta_0)\},
% \]
% 则 $A$ 为有限数或 $-\infty$.
% 
% 先考虑 $A$ 为有限数的情形。对 $\forall\varepsilon>0$，由下确界的性质知，$\exists x_1\in U_0^+(x_0,\delta_0)$，使得 $A\leqslant f(x_1)<A+\varepsilon$。记 $\delta=x_1-x_0$，则 $\delta>0$，而且当 $x\in U_0^+(x_0,\delta)$ 时，由于 $f(x)$ 在 $U_0^+(x_0,\delta_0)$ 内单调上升，故有 $A\leqslant f(x)\leqslant f(x_1)<A+\varepsilon$，即 $|f(x)-A|<\varepsilon$。这就证明了，此时 $\lim\limits_{x\to x_0+0}f(x)=A$ 成立.
% 
% 再考虑 $A=-\infty$ 的情形。对 $\forall M>0$，由下确界的定义知，$\exists x_1\in U_0^+(x_0,\delta_0)$ 使得 $f(x_1)<-M$。记 $\delta=x_1-x_0$，则 $\delta>0$，而且当 $x\in U_0^+(x_0,\delta)$ 时，由于 $f(x)$ 在 $U_0^+(x_0,\delta_0)$ 内单调上升，故有 $f(x)\leqslant f(x_1)<-M$。这就证明了，此时亦有 $\lim\limits_{x\to x_0+0}f(x)=A$ 成立.
% 
% 对于函数极限，有下面的柯西收敛准则：
% \end{proof}
% 
% 
% \begin{theorem}\label{theorem:ma-3-1-9}
% 设 $f(x)$ 在 $U_0(x_0,\delta_0)$ 内有定义，则 $\lim\limits_{x\to x_0}f(x)$ 存在的充分必要条件是：$\forall\varepsilon>0$，$\exists\delta>0$，当 $x',x''\in U_0(x_0,\delta)$ 时，有
% \[
%   |f(x')-f(x'')|<\varepsilon.
% \]
% \end{theorem}
% 
% 
% \begin{proof}
% 必要性\quad 设 $\lim\limits_{x\to x_0}f(x)=A$ 存在，则 $\forall\varepsilon>0$，$\exists\delta>0$，当 $x\in U_0(x_0,\delta)$ 时，有 $|f(x)-A|<\varepsilon/2$。于是，当 $x',x''\in U_0(x_0,\delta)$ 时，有
% \[
%   |f(x')-f(x'')|\leqslant |f(x')-A|+|A-f(x'')|<\frac{\varepsilon}{2}+\frac{\varepsilon}{2}=\varepsilon.
% \]
% 
% 充分性\quad 在 $U_0(x_0,\delta_0)$ 内任意选取一个收敛于 $x_0$ 的序列 $\{x_n\}$。$\forall\varepsilon>0$，由充分性假定知，$\exists\delta>0$，使得
% \[
%   |f(x')-f(x'')|<\varepsilon/2,\qquad \forall x',x''\in U_0(x_0,\delta).
% \]
% 而对这一 $\delta>0$，由 $\lim\limits_{n\to\infty}x_n=x_0$ 知，存在 $N$，使得当 $n>N$ 时，有 $|x_n-x_0|<\delta$。于是，对 $\forall m,n>N$，有 $x_n,x_m\in U_0(x_0,\delta)$，从而有
% \[
%   |f(x_n)-f(x_m)|<\varepsilon/2.
% \]
% 这表明序列 $\{f(x_n)\}$ 是一个柯西序列，因此由序列的柯西收敛准则知，存在实数 $A$，使得 $\lim\limits_{n\to\infty}f(x_n)=A$.
% 
% 下证 $\lim\limits_{x\to x_0}f(x)=A$。由前面所证知，$\forall\varepsilon>0$，存在 $\delta>0$ 和正整数 $N$，使得
% \[
%   |f(x)-f(x_n)|<\varepsilon/2,\qquad \forall x\in U_0(x_0,\delta),\ \forall n>N.
% \]
% 于是，有
% \[
%   |f(x)-A|=\lim_{n\to\infty}|f(x)-f(x_n)|\leqslant\varepsilon/2<\varepsilon,\qquad \forall x\in U_0(x_0,\delta).
% \]
% 这就证明了 $\lim\limits_{x\to x_0}f(x)=A$ 成立.
% \end{proof}
% 
% 
% \begin{remark}\label{remark:ma-3-source-46}
% 对于其他类型极限也有相应的单调收敛性定理和柯西收敛准则，请读者作为练习给出这些结果的详细表述.
% \end{remark}
% 
% 
% \begin{example}\label{example:ma-3-1-13}
% 证明极限
% \[
%   \lim_{x\to+\infty}\frac{x-1}{x+1}\cos\frac{2\pi x}{3}
% \]
% 不存在.
% \end{example}
% 
% 
% \begin{proof}
% 令 $\varepsilon_0=\dfrac12$，对任意的 $X>0$ (不妨设 $X>3$)，取正整数 $n_0$，使得
% \[
%   x'=3n_0>X,\qquad x''=3n_0+\frac34>X,
% \]
% 则有
% \[
%   |f(x')-f(x'')|=\left|\frac{3n_0-1}{3n_0+1}\right|>\frac12=\varepsilon_0.
% \]
% 故由柯西收敛准则知，$\lim\limits_{x\to+\infty}\dfrac{x-1}{x+1}\cos\dfrac{2\pi x}{3}$ 不存在.
% \end{proof}
% 
% 
% \paragraph*{两个重要极限}
% 
% 
% \begin{theorem}[{第一个重要极限}]\label{theorem:ma-3-source-49}
% \[
%   \lim_{x\to0}\frac{\sin x}{x}=1.
% \]
% \end{theorem}
% 
% 
% \begin{analyze}
% 如图 3.1.4 所示，在单位圆盘 $D=\{(\xi,\eta):\xi^2+\eta^2\leqslant1\}$ 上，$x$ 是圆心角 $\angle AOB$，以弧度计，即它恰好等于 $\overset{\frown}{AB}$ 的弧长，而 $\sin x=\overline{BC}$ 是弦长 $\overline{BB'}$ 之半。因此
% \[
%   \frac{\overline{BB'}}{\overset{\frown}{BB'}}
%   =\frac{2\sin x}{2x}=\frac{\sin x}{x}\to1\qquad (x\to0),
% \]
% 即圆心角趋于 $0$ 时，对应的弦长与弧长之比趋于 $1$.
% \end{analyze}
% 
% 
% \begin{proof}
% 为了给出这一重要极限的证明，我们先证明一个重要的不等式：
% \[
%   \sin x<x<\tan x,\qquad \forall x\in\left(0,\frac{\pi}{2}\right).
% \]
% 事实上，如图 3.1.4 所示，我们有
% \[
%   \triangle AOB\text{ 的面积}<\text{扇形 }AOB\text{ 的面积}<\triangle AOD\text{ 的面积},
% \]
% 即
% \[
%   \frac12\sin x<\frac12x<\frac12\tan x.
% \]
% 这样，我们就证明了所述的不等式.
% 
% 利用这一不等式，容易导出，对一切的 $0\ne x\in\mathbb{R}$，有
% \[
%   |\sin x|<|x|.
% \]
% 由此可以导出，对 $\forall a\in\mathbb{R}$，有
% \[
%   |\cos x-\cos a|
%   =\left|2\sin\frac{x+a}{2}\sin\frac{x-a}{2}\right|
%   \leqslant 2\left|\sin\frac{x-a}{2}\right|\leqslant |x-a|,
% \]
% 从而有
% \[
%   \lim_{x\to a}\cos x=\cos a.
% \]
% 同理可证
% \[
%   \lim_{x\to a}\sin x=\sin a.
% \]
% 利用前面所证不等式，我们有
% \[
%   \cos x<\frac{\sin x}{x}<1,\qquad \forall x\in\left(0,\frac{\pi}{2}\right).
% \]
% 利用偶函数性质，这一不等式在 $-\dfrac{\pi}{2}<x<0$ 时也成立。再注意到，前面我们已经证明了 $\lim\limits_{x\to0}\cos x=1$，由夹逼收敛原理知 $\lim\limits_{x\to0}\dfrac{\sin x}{x}=1$ 成立.
% \end{proof}
% 
% 
% \begin{example}\label{example:ma-3-1-14}
% 求极限
% \[
%   \lim_{x\to0}\frac{\sin(\sin x)}{x}.
% \]
% \end{example}
% 
% 
% \begin{solve}
% 由于 $|\sin x|<|x|$，所以 $\sin x\to0\ (x\to0)$。这样，我们有
% \[
% \begin{aligned}
%   \lim_{x\to0}\frac{\sin(\sin x)}{x}
%   &=\lim_{x\to0}\frac{\sin(\sin x)}{\sin x}\frac{\sin x}{x}\\
%   &=\lim_{x\to0}\frac{\sin(\sin x)}{\sin x}\lim_{x\to0}\frac{\sin x}{x}\\
%   &=\lim_{t\to0}\frac{\sin t}{t}\lim_{x\to0}\frac{\sin x}{x}=1\qquad (t=\sin x).
% \end{aligned}
% \qedhere
% \]
% \end{solve}
% 
% 
% \begin{theorem}[{第二个重要极限}]\label{theorem:ma-3-source-54}
% \[
%   \lim_{x\to\infty}\left(1+\frac1x\right)^x=e.
% \]
% \end{theorem}
% 
% 
% \begin{proof}
% 先考虑 $x\to+\infty$ 的情形。由于当 $x>1$ 时，有
% \[
%   1+\frac1{[x]+1}\leqslant 1+\frac1x\leqslant 1+\frac1{[x]},
% \]
% 所以，有
% \[
%   \left(1+\frac1{[x]+1}\right)^x
%   \leqslant \left(1+\frac1x\right)^x
%   \leqslant \left(1+\frac1{[x]}\right)^x,
% \]
% 从而有
% \[
%   \left(1+\frac1{[x]+1}\right)^{[x]}
%   \leqslant \left(1+\frac1x\right)^x
%   \leqslant \left(1+\frac1{[x]}\right)^{[x]+1}.
% \]
% 这样，应用序列极限所得结果和夹逼收敛原理，知
% \[
%   \lim_{x\to+\infty}\left(1+\frac1x\right)^x=e.
% \]
% 
% 再考虑 $x\to-\infty$ 的情形。令 $x=-y$，则 $y\to+\infty$，从而有
% \[
% \begin{aligned}
%   \lim_{x\to-\infty}\left(1+\frac1x\right)^x
%   &=\lim_{y\to+\infty}\left(1-\frac1y\right)^{-y}\\
%   &=\lim_{y\to+\infty}\left(1+\frac1{y-1}\right)^{y-1}\left(1+\frac1{y-1}\right)=e.
% \end{aligned}
% \]
% 由极限与单侧极限的关系，得
% \[
%   \lim_{x\to\infty}\left(1+\frac1x\right)^x=e.
% \qedhere
% \]
% \end{proof}
% 
% 
% \begin{example}\label{example:ma-3-1-15}
% 求极限
% \[
%   \lim_{x\to0}(1-2x)^{\frac1x}.
% \]
% \end{example}
% 
% 
% \begin{solve}
% 令 $x=\dfrac1t$，则 $x\to0$ 蕴涵着 $t\to\infty$。于是
% \[
% \begin{aligned}
%   \lim_{x\to0}(1-2x)^{\frac1x}
%   &=\lim_{t\to\infty}\left(1-\frac2t\right)^t
%   =\lim_{t\to\infty}\left(\frac{t-2}{t}\right)^t\\
%   &=\lim_{t\to\infty}\left(1+\frac2{t-2}\right)^{-t}
%   =\lim_{s\to\infty}\left(1+\frac1s\right)^{-(2s+2)}\\
%   &=\frac1{\left(\lim\limits_{s\to\infty}\left(1+\dfrac1s\right)^s\right)^2}
%   \lim_{s\to\infty}\left(1+\frac1s\right)^{-2}=e^{-2},
% \end{aligned}
% \]
% 这里 $s=\dfrac{t-2}{2}$.
% \end{solve}
% 
% 
% \paragraph*{函数的上、下极限}
% 
% 
% 作为本节的结束，我们简要地介绍一下函数的上、下极限。为了避免繁琐，我们这里仅考虑当 $x\to x_0$ 时的情形，其他情形请读者自己补出.
% 
% 设函数 $y=f(x)$ 在 $U_0(x_0,\delta_0)$ 内有定义。对任意的 $0<\delta<\delta_0$，设
% \[
%   \ell(\delta)=\inf\{f(x):x\in U_0(x_0,\delta)\},
% \]
% \[
%   h(\delta)=\sup\{f(x):x\in U_0(x_0,\delta)\},
% \]
% 则当 $x\in U_0(x_0,\delta)$ 时，有 $\ell(\delta)\leqslant f(x)\leqslant h(\delta)$。容易看出，作为 $(0,\delta_0)$ 上的函数，$\ell(\delta)$ 关于 $\delta$ 单调递减，而 $h(\delta)$ 关于 $\delta$ 单调递增。因此它们的广义极限都存在。令
% \[
%   \ell=\lim_{\delta\to0}\ell(\delta),\qquad h=\lim_{\delta\to0}h(\delta),
% \]
% 则分别称之为当 $x$ 趋于 $x_0$ 时 $f(x)$ 的下极限和上极限，记为
% \[
%   \ell=\varliminf_{x\to x_0}f(x),\qquad h=\varlimsup_{x\to x_0}f(x).
% \]
% 
% 函数的上、下极限与序列的上、下极限具有相似的性质，请读者作为练习自己补出。特别地，我们有
% 
% \begin{theorem}\label{theorem:ma-3-1-10}
% 设函数 $f(x)$ 在 $U_0(x_0,\delta_0)(\delta_0>0)$ 内有定义，则 $\lim\limits_{x\to x_0}f(x)$ 存在的充分必要条件是
% \[
%   \varliminf_{x\to x_0}f(x)=\varlimsup_{x\to x_0}f(x).
% \]
% \end{theorem}
% 
% 
% \section{函数的连续与间断}
% \label{sec:ma-3-2}
% 
% \subsection{函数的连续与间断}
% \label{subsec:ma-3-2-1}
% 
% 所谓函数在某点的连续与间断，是对该点函数值与其附近点的函数值的变化趋势是否衔接的一种刻画，是函数的一种局部属性，其定义如下：
% 
% \begin{definition}\label{definition:ma-3-2-1}
% 设函数 $f(x)$ 在 $U(x_0,\delta_0)(\delta_0>0)$ 内有定义。若有 $\lim\limits_{x\to x_0}f(x)=f(x_0)$，则称 $f(x)$ 在点 $x_0$ 连续，并称 $x_0$ 为 $f(x)$ 的一个连续点；否则称 $f(x)$ 在点 $x_0$ 间断（或不连续），并称 $x_0$ 为 $f(x)$ 的一个间断点（或不连续点）.
% \end{definition}
% 
% 我们也可以用 $\varepsilon$-$\delta$ 语言来表述 $f(x)$ 在点 $x_0$ 处连续：若 $\forall\varepsilon>0$，$\exists\delta>0$，当 $x\in U(x_0,\delta)$ 时，有 $|f(x)-f(x_0)|<\varepsilon$，则称 $f(x)$ 在点 $x_0$ 连续.
% 
% 设函数 $f(x)$ 在点 $x_0$ 处连续，则有
% \[
%   \lim_{x\to x_0}f(x)=f(x_0)=f\left(\lim_{x\to x_0}x\right).
% \]
% 因此，函数在一点处连续可以看成在该点附近求函数值与求极限的顺序可以互换.
% 
% \begin{example}\label{example:ma-3-2-1}
% 在上一节例~\ref{example:ma-3-1-2} 之后，我们曾经指出，对任意的实数 $x_0$ 和任意的正整数 $k$，有 $\lim\limits_{x\to x_0}x^k=x_0^k$。由函数极限的性质知，对任意的多项式函数
% \[
%   P(x)=a_nx^n+a_{n-1}x^{n-1}+\cdots+a_1x+a_0,
% \]
% 有 $\lim\limits_{x\to x_0}P(x)=P(x_0)$。这表明，多项式函数在 $\mathbb{R}$ 上的任意点处都连续.
% \end{example}
% 
% 
% \begin{example}\label{example:ma-3-2-2}
% 在上一节我们已经证明了
% \[
%   \lim_{x\to x_0}\sin x=\sin x_0,\qquad \lim_{x\to x_0}\cos x=\cos x_0.
% \]
% 这里 $x_0$ 是任意给定的实数。这表明，三角函数 $\sin x$ 和 $\cos x$ 在 $\mathbb{R}$ 上的任意点处都连续.
% \end{example}
% 
% 为了能够对闭区间上定义的函数讨论其端点处的连续性问题，我们引入：
% 
% \begin{definition}\label{definition:ma-3-2-2}
% 若函数 $f(x)$ 在 $U^+(x_0,\delta_0)$ 上有定义，且 $f(x_0+0)=f(x_0)$，则称 $f(x)$ 在点 $x_0$ 右连续；若 $f(x)$ 在 $U^-(x_0,\delta_0)$ 上有定义，且 $f(x_0-0)=f(x_0)$，则称 $f(x)$ 在点 $x_0$ 左连续.
% \end{definition}
% 
% 显然，$f(x)$ 在点 $x_0$ 处连续的充分必要条件是它在该点左、右连续.
% 
% \begin{definition}\label{definition:ma-3-2-3}
% 设函数 $f(x)$ 在 $[a,b]$ 上有定义。若对 $x\in(a,b)$，$f(x)$ 在点 $x$ 处连续，则称 $f(x)$ 在 $(a,b)$ 内连续，此时记为 $f(x)\in C(a,b)$；若 $f(x)\in C(a,b)$，而且在左端点 $a$ 右连续、在右端点 $b$ 左连续，则称 $f(x)$ 在 $[a,b]$ 上连续，此时记为 $f(x)\in C[a,b]$.
% \end{definition}
% 
% 
% \paragraph*{注 1}
% 设 $f(x)\in C(a,b)$，若 $\lim\limits_{x\to a+0}f(x)=A$ 及 $\lim\limits_{x\to b-0}f(x)=B$ 存在，则 $f(x)$ 可连续延拓到 $[a,b]$，即以下函数
% \[
%   \tilde f(x)=
%   \begin{cases}
%     A, & x=a,\\
%     f(x), & x\in(a,b),\\
%     B, & x=b
%   \end{cases}
% \]
% 是 $[a,b]$ 上的连续函数.
% 
% \paragraph*{注 2}
% 若 $f(x)\in C[a,b]$，则 $f(x)$ 可连续延拓到 $(-\infty,+\infty)$。事实上，以下函数就是 $f(x)$ 的一个连续延拓：
% \[
%   \tilde f(x)=
%   \begin{cases}
%     f(a), & x<a,\\
%     f(x), & a\leqslant x\leqslant b,\\
%     f(b), & x>b.
%   \end{cases}
% \]
% 请读者自己给出 $\tilde f(x)$ 的连续性证明.
% 
% 如果 $f(x)$ 在 $x_0$ 处是间断的，则情况比较复杂。一般地，我们可以将间断点分为以下两类：
% 
% (1) 若 $f(x_0+0)$ 与 $f(x_0-0)$ 都存在，则称 $x_0$ 为 $f(x)$ 的第一类间断点。此时，若 $f(x_0+0)=f(x_0-0)\ne f(x_0)$，则称此间断点为可去间断点；否则称其为跳跃间断点.
% 
% (2) 若 $f(x_0+0)$ 与 $f(x_0-0)$ 至少有一个不存在时，则称 $x_0$ 为 $f(x)$ 的第二类间断点.
% 
% \begin{example}\label{example:ma-3-2-3}
% 考查函数
% \[
%   f(x)=
%   \begin{cases}
%     x\sin\dfrac1x, & x\ne0,\\
%     1, & x=0
%   \end{cases}
% \]
% 在 $x=0$ 处的连续性.
% \end{example}
% 
% 
% \begin{solve}
% 由 $\left|x\sin\dfrac1x\right|\leqslant |x|$ 知，$\lim\limits_{x\to0}f(x)=0\ne f(0)$。因此，$x=0$ 是 $f(x)$ 的可去间断点（参见图 3.2.1）.
% \end{solve}
% 
% 
% \begin{example}\label{example:ma-3-2-4}
% 考查函数 $f(x)=\dfrac1x-\left[\dfrac1x\right]$ 在 $x=1$ 处的连续性.
% \end{example}
% 
% 
% \begin{solve}
% 当 $x>1$ 时，有 $0<\dfrac1x<1$，于是，此时 $f(x)=\dfrac1x$，从而 $f(1+0)=1$；而当 $\dfrac12<x\leqslant1$ 时，有 $1\leqslant\dfrac1x<2$，于是，此时 $f(x)=\dfrac1x-1$，从而 $f(1-0)=0=f(1)$。因此，$f(x)$ 在 $x=1$ 处左连续，右不连续，从而 $x=1$ 为跳跃间断点（参见图 3.2.2）.
% \end{solve}
% 
% 
% \begin{example}\label{example:ma-3-2-5}
% 考查函数
% \[
%   f(x)=
%   \begin{cases}
%     e^{\frac1x}, & x\ne0,\\
%     0, & x=0
%   \end{cases}
% \]
% 在 $x=0$ 处的连续性.
% \end{example}
% 
% 
% \begin{solve}
% 由于
% \[
%   f(0+0)=+\infty,
% \]
% \[
%   f(0-0)=0=f(0),
% \]
% 所以 $x=0$ 是 $f(x)$ 的第二类间断点（参见图 3.2.3）.
% \end{solve}
% 
% 
% \begin{example}\label{example:ma-3-2-6}
% 讨论函数
% \[
%   f(x)=
%   \begin{cases}
%     \sin\dfrac1x, & x\ne0,\\
%     0, & x=0
%   \end{cases}
% \]
% 在 $x=0$ 处的连续性.
% \end{example}
% 
% 
% \begin{solve}
% 在上一节我们已经证明了 $\lim\limits_{x\to0+0}f(x)$ 与 $\lim\limits_{x\to0-0}f(x)$ 都不存在。因此，此时 $x=0$ 为 $f(x)$ 的第二类间断点（参见图 3.2.4）.
% \end{solve}
% 
% 
% \begin{example}\label{example:ma-3-2-7}
% 讨论黎曼函数 $R(x)$（参考例~\ref{example:ma-1-2-4}）在 $(0,1)$ 内的连续性.
% \end{example}
% 
% 
% \begin{solve}
% 任意取定 $x_0\in(0,1)$，我们来考查 $R(x)$ 在 $x_0$ 处的连续性.
% 
% $\forall\varepsilon>0$，取定正整数 $p$，使得 $1/p<\varepsilon$。由于在 $(0,1)$ 内分母小于或等于 $p$ 的分数仅为有限个，在这有限个数中，必存在 $r\ne x_0$，使得 $r$ 到 $x_0$ 的距离是在所有和 $x_0$ 不相等的这些数中与 $x_0$ 的距离达到最小。如果我们令 $\delta=\min\{|x_0-r|,x_0,1-x_0\}$，则对 $\forall x\in U_0(x_0,\delta)$，有
% \[
%   |R(x)-0|<1/p<\varepsilon,
% \]
% 这表明 $\lim\limits_{x\to x_0}R(x)=0$。由此可知：当 $x_0$ 为有理数时，$R(x)$ 在 $x_0$ 处不连续，是可去间断点；当 $x_0$ 为无理数时，$R(x)$ 在 $x_0$ 处连续.
% \end{solve}
% 
% 
% \begin{example}\label{example:ma-3-2-8}
% 考查有理函数
% \[
%   f(x)=\frac{p(x)}{q(x)}
%   =
%   \frac{a_nx^n+a_{n-1}x^{n-1}+\cdots+a_1x+a_0}
%        {b_mx^m+b_{m-1}x^{m-1}+\cdots+b_1x+b_0}
% \]
% 的连续性以及 $x\to\infty$ 时的极限，其中 $a_nb_m\ne0$，而且 $p(x),q(x)$ 互素.
% \end{example}
% 
% 
% \begin{solve}
% 设 $x_0\in(-\infty,+\infty)$。当 $q(x_0)\ne0$ 时，有
% \[
%   \lim_{x\to x_0}\frac{p(x)}{q(x)}=\frac{p(x_0)}{q(x_0)},
% \]
% 从而 $f(x)$ 在 $x_0$ 处连续；当 $q(x_0)=0$ 时，$f(x)$ 在 $x_0$ 处没有定义。此外，由于 $p(x),q(x)$ 互素，所以 $p(x_0)\ne0$，从而
% \[
%   \lim_{x\to x_0}\frac{p(x)}{q(x)}=\infty.
% \]
% 现考虑 $x\to\infty$ 的情况。此时，我们有
% \[
%   \lim_{x\to\infty}f(x)
%   =
%   \lim_{x\to\infty}
%   \left(
%     x^{n-m}\frac{a_n+\dfrac{a_{n-1}}x+\cdots+\dfrac{a_0}{x^n}}
%     {b_m+\dfrac{b_{m-1}}x+\cdots+\dfrac{b_0}{x^m}}
%   \right)
%   =
%   \begin{cases}
%     0, & n<m,\\[4pt]
%     \dfrac{a_n}{b_m}, & n=m,\\[6pt]
%     \infty, & n>m.
%   \end{cases}
% \qedhere
% \]
% \end{solve}
% 
% 
% \begin{example}\label{example:ma-3-2-9}
% 设函数 $f(x)$ 在 $(a,b)$ 上单调，则 $f(x)$ 在 $(a,b)$ 内只可能有第一类间断点.
% \end{example}
% 
% 
% \begin{proof}
% 不妨设 $f(x)$ 在 $(a,b)$ 上单调递减。对 $\forall\xi\in(a,b)$，在 $(a,\xi)$ 中，$f(x)$ 单调递减且有下界 $f(\xi)$，因此 $\lim\limits_{x\to\xi-0}f(x)$ 存在；而在 $(\xi,b)$ 中，$f(x)$ 单调递减且有上界 $f(\xi)$，因此 $\lim\limits_{x\to\xi+0}f(x)$ 存在。这样，若 $f(x)$ 在点 $\xi$ 处不连续，则 $\xi$ 只能是第一类间断点.
% \end{proof}
% 
% 
% \begin{remark}\label{remark:ma-3-source-78}
% 由此例可以看出：若 $f(x)$ 在 $(a,b)$ 单调，则 $f(x)$ 在每一个间断点的取值具有跳跃性。例如，假定 $f(x)$ 在 $(a,b)$ 单调递增，并且假设 $\xi\in(a,b)$ 为 $f(x)$ 的一个间断点，则必有 $f(\xi-0)<f(\xi+0)$。因此 $f(x)$ 在 $(a,b)$ 中取不到 $(f(\xi-0),f(\xi+0))\setminus\{f(\xi)\}$ 中的值。换句话说，$f(x)$ 在点 $\xi$ 处有一跳跃。作为练习，请读者试构造一个单调函数，使其具有无穷多个间断点.
% \end{remark}
% 
% 
% \subsection{连续函数的性质}
% \label{subsec:ma-3-2-2}
% 
% 首先，由函数在一点连续的定义和函数极限的性质，读者不难得到：
% 
% (1) 连续函数是局部有界的，即若 $f(x)$ 在点 $x_0$ 处连续，则必存在 $\delta>0$，使得 $f(x)$ 在 $U(x_0,\delta)$ 上有界.
% 
% (2) 连续函数具有局部保号性，即：若 $f(x)$ 在点 $x_0$ 处连续，而且 $f(x_0)>0$，则必存在 $\delta>0$，使得 $f(x)>0$ 对一切 $x\in U(x_0,\delta)$ 成立。以后我们要经常用到的结论是：对任意的 $0<A<f(x_0)$（如 $A=f(x_0)/2$），总存在 $\delta_0>0$，使得当 $x\in U(x_0,\delta_0)$ 时，有 $f(x)>A$.
% 
% (3) 连续函数经四则运算后仍然连续，即：若 $f(x)$ 和 $g(x)$ 在点 $x_0$ 处连续，则函数 $f(x)\pm g(x)$，$f(x)g(x)$ 和 $\dfrac{f(x)}{g(x)}$（$g(x_0)\ne0$）仍然在点 $x_0$ 处连续.
% 
% 此外，连续函数的复合函数和反函数也是连续的.
% 
% \begin{theorem}[{复合函数的连续性}]\label{theorem:ma-3-2-1}
% 设 $u=g(x)$ 在点 $x_0$ 连续，$y=f(u)$ 在点 $u_0=g(x_0)$ 连续，则复合函数 $f(g(x))$ 在点 $x_0$ 连续.
% \end{theorem}
% 
% 
% \begin{proof}
% $\forall\varepsilon>0$，由 $f(u)$ 在点 $u_0$ 连续知，$\exists\eta>0$，当 $|u-u_0|<\eta$ 时，有
% \[
%   |f(u)-f(u_0)|<\varepsilon.
% \]
% 对 $\eta>0$，再由 $u=g(x)$ 在点 $x_0$ 连续知，$\exists\delta>0$，当 $|x-x_0|<\delta$ 时，有
% \[
%   |u-u_0|=|g(x)-g(x_0)|<\eta,
% \]
% 从而当 $|x-x_0|<\delta$ 时，有
% \[
%   |f(g(x))-f(g(x_0))|<\varepsilon.
% \]
% 因此，$f(g(x))$ 在点 $x_0$ 连续.
% \end{proof}
% 
% 
% \begin{theorem}[{反函数的连续性}]\label{theorem:ma-3-2-2}
% 设 $f(x)$ 是区间 $I$ 上严格单调的连续函数，则其反函数 $x=f^{-1}(y)$ 在 $f(I)$ 上连续.
% \end{theorem}
% 
% 
% \begin{proof}
% 由 $f(x)$ 的严格单调性，知 $f^{-1}(y)$ 存在。现不妨设 $f(x)$ 严格递增。显然，当 $f(x)$ 严格递增时，$f^{-1}(x)$ 必定严格递增。另外，由 $f(x)$ 的连续性，我们可以推出 $f(I)$ 必定是一个区间（在下节中，我们将证明任何区间上连续函数的值域都是一个区间）。事实上，倘若 $f(I)$ 不是一个区间，则由 $f(x)$ 的连续性，必存在一个开区间 $(\alpha,\beta)$ 使得 $f(I)\cap(\alpha,\beta)=\varnothing$，并且 $\alpha,\beta\in f(I)$。设 $f(x_0)=\alpha$。由 $f(x)$ 的单调性有
% \[
%   f(x_0-0)\leqslant f(x_0)<\beta\leqslant f(x_0+0).
% \]
% 这与 $f(x)$ 在 $x_0$ 的连续性矛盾.
% 
% 现任取 $y_0\in f(I)$，则存在 $x_0\in I$，使得 $x_0=f^{-1}(y_0)$。不妨假定 $x_0$ 不是 $I$ 的端点，否则考查左、右连续即可.
% 
% $\forall\varepsilon>0$，存在 $x_1,x_2\in I$，使得 $x_0-\varepsilon<x_1<x_0<x_2<x_0+\varepsilon$。记 $y_1=f(x_1),y_2=f(x_2)$，则有
% \[
%   y_1=f(x_1)<f(x_0)=y_0<f(x_2)=y_2.
% \]
% 令 $\delta=\min\{y_0-y_1,y_2-y_0\}>0$，则当 $y\in(y_0-\delta,y_0+\delta)$ 时，有 $y_1<y<y_2$，从而由 $f(I)$ 是一个区间推知，$\exists x\in I$，使得 $y=f(x)$，即 $y\in f(I)$.
% 于是
% \[
%   x_0-\varepsilon<x_1=f^{-1}(y_1)<f^{-1}(y)<f^{-1}(y_2)=x_2<x_0+\varepsilon
% \]
% 即
% \[
%   |f^{-1}(y)-f^{-1}(y_0)|=|f^{-1}(y)-x_0|<\varepsilon.
% \]
% 这就证明了 $f^{-1}(y)$ 在点 $y_0$ 处连续。注意到 $y_0\in f(I)$ 的任意性，即知 $f^{-1}(y)$ 在 $f(I)$ 上连续.
% \end{proof}
% 
% 
% \subsection{初等函数的连续性}
% \label{subsec:ma-3-2-3}
% 
% 首先我们来证明，对 $a>1$，指数函数 $f(x)=a^x$ 在 $(-\infty,+\infty)$ 上连续。我们回忆一下，实际上，在中学数学的学习中，我们并不知道当 $x$ 是无理数时 $a^x$ 是如何定义的。下面我们给出它的严格定义：
% 
% 对任何正有理数 $x=\dfrac qp$（其中 $p,q$ 是互素的正整数），定义 $a^x=(\sqrt[p]{a})^q$；对于任何负有理数 $x$，定义 $a^x=\dfrac1{a^{-x}}$，并定义 $a^0=1$。容易证明，$\forall r,s\in\mathbb Q$，有
% \[
%   \text{(1) } a^ra^s=a^{r+s};
% \]
% \[
%   \text{(2) 若 } r<s,\text{ 则 } a^r<a^s.
% \]
% 而当 $x$ 为一无理数时，定义
% \[
%   a^x=\sup\{a^r:r\text{ 为小于 }x\text{ 的有理数}\}.
% \]
% 这样，对一切 $x\in(-\infty,+\infty)$，$a^x$ 均有定义.
% 
% 现在我们来证明 $a^x$ 在 $(-\infty,+\infty)$ 上严格递增。设 $x_1<x_2$，则存在有理数 $q_1,q_2$，使得
% \[
%   x_1\leqslant q_1<q_2\leqslant x_2,
% \]
% 从而
% \[
%   a^{x_1}=\sup\{a^q:q\text{ 为小于 }x_1\text{ 的有理数}\}
%   \leqslant a^{q_1}<a^{q_2}
%   \leqslant \sup\{a^q:q\text{ 为小于 }x_2\text{ 的有理数}\}=a^{x_2}.
% \]
% 
% 再来证明 $a^x\in C(-\infty,+\infty)$。$\forall x_0\in(-\infty,+\infty)$，$\forall\varepsilon>0$，先取正整数 $N$，使得 $a<(1+\varepsilon)^{N/2}$；再取正整数 $q$，使得 $q\leqslant Nx_0<q+1$；然后令
% \[
%   q_1=\frac{q-1}{N},\quad q_2=\frac{q+1}{N}.
% \]
% 这样，由 $a^x$ 的严格递增性知
% \[
%   a^{q_1}\leqslant f(x_0-0)\leqslant f(x_0)\leqslant f(x_0+0)\leqslant a^{q_2},
% \]
% 从而有
% \[
%   1\leqslant \frac{f(x_0+0)}{f(x_0-0)}
%   \leqslant \frac{a^{q_2}}{a^{q_1}}
%   =a^{q_2-q_1}=a^{2/N}<1+\varepsilon.
% \]
% 由 $\varepsilon>0$ 的任意性，即知
% \[
%   f(x_0+0)=f(x_0-0)=f(x_0)=a^{x_0}.
% \]
% 这样，我们就证明了 $a^x\in C(-\infty,+\infty)$.
% 
% 对于 $0<a<1$，我们定义 $a^x=\dfrac1{(1/a)^x}$，则由连续函数的性质知，此时亦有 $a^x\in C(-\infty,+\infty)$；至于 $a=1$ 的情形，我们规定 $a^x=1$。这样一来，对 $\forall a\in(0,+\infty)$，$y=a^x$ 均有定义，而且 $y=a^x\in C(-\infty,+\infty)$.
% 
% 由反函数的连续性知，对任意的 $a>0$，对数函数 $y=\log_a x\in C(0,+\infty)$；再由复合函数连续性知，幂函数 $y=x^\alpha=e^{\alpha\ln x}\in C(0,+\infty)$.
% 
% 上一节我们已经证明了，对任意的实数 $a$，有 $\lim\limits_{x\to a}\cos x=\cos a$，从而 $\cos x\in C(-\infty,+\infty)$。同理可证 $\sin x\in C(-\infty,+\infty)$。再利用连续函数的性质知，所有的三角函数及反三角函数在其定义域内都是连续的.
% 
% 这样一来，我们已经证明了所有的基本初等函数在其定义域内都是连续的。由于初等函数是由基本初等函数经过有限次四则运算和复合运算所得，利用连续函数的性质，我们有
% 
% \begin{theorem}\label{theorem:ma-3-2-3}
% 初等函数在其定义域内是连续的.
% \end{theorem}
% 
% 
% \begin{example}\label{example:ma-3-2-10}
% 设 $\lim\limits_{x\to x_0}u(x)=a>0$，$\lim\limits_{x\to x_0}v(x)=b$，则
% \[
%   \lim_{x\to x_0}u(x)^{v(x)}=a^b.
% \]
% \end{example}
% 
% 
% \begin{proof}
% 定义 $u(x_0)=a$，$v(x_0)=b$，则 $u(x),v(x)$ 在点 $x_0$ 处都连续，从而
% \[
%   \lim_{x\to x_0}u(x)^{v(x)}
%   =
%   \lim_{x\to x_0}e^{v(x)\ln u(x)}
%   =
%   e^{v(x_0)\ln u(x_0)}
%   =
%   u(x_0)^{v(x_0)}
%   =
%   a^b.
% \qedhere
% \]
% \end{proof}
% 
% 
% \begin{example}\label{example:ma-3-2-11}
% 求极限
% \[
%   \lim_{x\to+\infty}\left(1+\sin\frac1x\right)^x.
% \]
% \end{example}
% 
% 
% \begin{solve}
% 由
% \[
%   \left(1+\sin\frac1x\right)^x
%   =
%   \left(1+\sin\frac1x\right)^{\frac1{\sin(1/x)}\cdot\frac{\sin(1/x)}{1/x}},
% \]
% \[
%   \lim_{x\to+\infty}\left(1+\sin\frac1x\right)^{\frac1{\sin(1/x)}}=e,
%   \qquad
%   \lim_{x\to+\infty}\frac{\sin(1/x)}{1/x}=1,
% \]
% 得
% \[
%   \lim_{x\to+\infty}\left(1+\sin\frac1x\right)^x=e^1=e.
% \qedhere
% \]
% \end{solve}
% 
% 
% \section{闭区间上连续函数的基本性质}
% \label{sec:ma-3-3}
% 
% 上一节中，我们讨论了函数 $f(x)$ 在其连续点 $x_0$ 邻近的局部性质。例如，若一个函数 $f(x)$ 在点 $x_0$ 连续且满足 $f(x_0)<\eta$，则我们能断定在点 $x_0$ 邻近有 $f(x)<\eta$，即存在 $\delta>0$，使得 $f(x)<\eta$ 对一切的 $x\in(x_0-\delta,x_0+\delta)$ 成立。如果一个函数 $f(x)\in C[a,b]$，那么 $f(x)$ 在整个闭区间上能具有什么样的整体性质呢？本节将讨论这一重要问题.
% 
% \begin{theorem}[{有界性}]\label{theorem:ma-3-3-1}
% 设函数 $f(x)\in C[a,b]$，则 $f(x)$ 在 $[a,b]$ 上有界.
% \end{theorem}
% 
% 
% \begin{proof}
% 倘若 $f(x)$ 在 $[a,b]$ 上无界，则将 $[a,b]$ 等分成两个闭区间，$f(x)$ 必在其中一个子区间上无界，取定这样的一个区间并将其设为 $[a_1,b_1]$。同理，再将 $[a_1,b_1]$ 二等分，$f(x)$ 必在其中一个子区间上无界，取定这样的一个区间并将其记为 $[a_2,b_2]$。如此进行下去，我们就得到一个闭区间列 $\{[a_n,b_n]\}$，满足：
% \[
%   \text{(1) } [a_{n+1},b_{n+1}]\subset [a_n,b_n];
% \]
% \[
%   \text{(2) } b_{n+1}-a_{n+1}=\frac12(b_n-a_n);
% \]
% \[
%   \text{(3) } f(x)\text{ 在每个 }[a_n,b_n]\text{ 上都无界.}
% \]
% 于是，由闭区间套定理知，存在唯一的 $\xi\in\bigcap\limits_{n=1}^{+\infty}[a_n,b_n]$。再由 $f(x)\in C[a,b]$ 和 $\xi\in[a_1,b_1]\subset[a,b]$ 知，存在 $\delta>0$，使得 $f(x)$ 在 $U(\xi,\delta)\cap[a,b]$ 上有界；而条件 (2) 又蕴涵着，当 $n$ 充分大时，有 $[a_n,b_n]\subset U(\xi,\delta)$，从而 $f(x)$ 在 $[a_n,b_n]$ 上有界。这与条件 (3) 矛盾.
% \end{proof}
% 
% 
% \begin{example}\label{example:ma-3-3-1}
% 函数 $y=\dfrac1x$ 在开区间 $(0,1)$ 上连续，但它在此区间上无界.
% \end{example}
% 
% 
% \paragraph*{思考题}
% 若不假定 $f(x)$ 在 $[a,b]$ 上连续，仅仅假定 $\forall x_0\in[a,b]$，$\lim\limits_{x\to x_0}f(x)$ 存在（在端点处单侧极限存在），能否推出 $f(x)$ 在 $[a,b]$ 上有界.
% 
% 设 $f(x)\in C[a,b]$。由于 $f(x)$ 在 $[a,b]$ 上有界，故 $\sup\limits_{x\in[a,b]}\{f(x)\}$ 与 $\inf\limits_{x\in[a,b]}\{f(x)\}$ 都为有限数。下面的定理告诉我们它们将属于 $f(x)$ 的值域.
% 
% \begin{theorem}[{最值定理}]\label{theorem:ma-3-3-2}
% 设 $f(x)\in C[a,b]$，则 $f(x)$ 在 $[a,b]$ 上必有最小值和最大值，即存在 $\xi,\zeta\in[a,b]$，使得 $f(\xi)\leqslant f(x)\leqslant f(\zeta)$ 对一切的 $x\in[a,b]$ 成立.
% \end{theorem}
% 
% 
% \begin{proof}
% 令 $M=\sup\limits_{x\in[a,b]}\{f(x)\}$，则由上确界的定义知，存在一个序列 $\{x_n\}\subset[a,b]$，使得 $f(x_n)\to M$（$n\to\infty$）。由于 $\{x_n\}$ 是有界序列，所以必有一收敛子列 $\{x_{n_k}\}$。设 $\lim\limits_{k\to\infty}x_{n_k}=\zeta$，则 $\zeta\in[a,b]$。再由 $f(x)\in C[a,b]$ 知
% \[
%   f(\zeta)=\lim_{k\to\infty}f(x_{n_k})=M.
% \]
% 同理可证存在 $\xi\in[a,b]$，使得
% \[
%   f(\xi)=m=\inf_{x\in[a,b]}\{f(x)\}.
% \]
% 此定理表明：若 $f(x)\in C[a,b]$，则存在 $\xi,\zeta\in[a,b]$，使得
% \[
%   f(\xi)=\min_{x\in[a,b]}\{f(x)\},\qquad
%   f(\zeta)=\max_{x\in[a,b]}\{f(x)\}.
% \qedhere
% \]
% \end{proof}
% 
% 
% \begin{example}\label{example:ma-3-3-2}
% 设 $f(x)\in C[a,b]$，且 $\forall x\in[a,b]$，有 $f(x)>0$。证明存在正数 $\eta$，使得 $f(x)\geqslant\eta$ 对一切的 $x\in[a,b]$ 成立.
% \end{example}
% 
% 
% \begin{proof}
% 由 $f(x)\in C[a,b]$，根据定理~\ref{theorem:ma-3-3-2} 知，$\exists\xi\in[a,b]$，使
% \[
%   f(\xi)=\min_{x\in[a,b]}\{f(x)\}>0,
% \]
% 从而取 $\eta=f(\xi)$ 即可.
% \end{proof}
% 
% 
% \begin{theorem}[{介值定理}]\label{theorem:ma-3-3-3}
% 设 $f(x)\in C[a,b]$，记 $m=\min\limits_{x\in[a,b]}\{f(x)\}$，$M=\max\limits_{x\in[a,b]}\{f(x)\}$，则 $f([a,b])=[m,M]$，即对 $\forall\eta\in(m,M)$，$\exists\xi\in[a,b]$，使得 $f(\xi)=\eta$.
% \end{theorem}
% 
% 
% \begin{proof}
% 由定理~\ref{theorem:ma-3-3-2} 知，$\exists x_1,x_2\in[a,b]$，使得 $f(x_1)=m$，$f(x_2)=M$。下面证明对 $\forall\eta\in(m,M)$，$\exists\xi\in[a,b]$，使得 $f(\xi)=\eta$.
% 
% 不妨设 $x_1<x_2$，并定义
% \[
%   E=\{x\in[x_1,x_2]:f(x)>\eta\}.
% \]
% 显然有 $x_1\notin E$，$x_2\in E$。现令 $E$ 的下确界为 $\xi$，则有 $f(\xi)=\eta$。事实上，若不然，则 $f(\xi)>\eta$。由连续函数的局部保号性知，存在 $\xi$ 的一个邻域 $U(\xi,\delta)$，使得 $f(x)>\eta$ 对一切的 $x\in U(\xi,\delta)$ 成立。由于 $\xi>x_1$，所以 $[x_1,x_2]\cap(\xi-\delta,\xi)\ne\varnothing$，现取 $\xi'\in[x_1,x_2]\cap(\xi-\delta,\xi)$，则有 $f(\xi')>\eta$，从而有 $\xi'<\xi$ 且 $\xi'\in E$。这与 $\xi$ 是 $E$ 的下确界矛盾.
% \end{proof}
% 
% 
% \begin{example}\label{example:ma-3-3-3}
% 设 $f(x)\in C[a,b]$，$x_j\in[a,b]$，$j=1,2,\cdots,n$。证明 $\exists\xi\in[a,b]$，使得
% \[
%   f(\xi)=\frac1n\sum_{j=1}^{n}f(x_j).
% \]
% \end{example}
% 
% 
% \begin{proof}
% 记
% \[
%   m=\min_{x\in[a,b]}\{f(x)\},\qquad
%   M=\max_{x\in[a,b]}\{f(x)\},
% \]
% 则有
% \[
%   m\leqslant\frac1n\sum_{j=1}^{n}f(x_j)\leqslant M.
% \]
% 于是，由介值定理知，$\exists\xi\in[a,b]$，使得
% \[
%   f(\xi)=\frac1n\sum_{j=1}^{n}f(x_j).
% \qedhere
% \]
% \end{proof}
% 
% 
% \begin{remark}\label{remark:ma-3-source-99}
% 将例~\ref{example:ma-3-3-3} 中的 $[a,b]$ 换成开区间 $(a,b)$，其结论仍真.
% \end{remark}
% 
% 
% \begin{theorem}[{零点存在定理}]\label{theorem:ma-3-3-4}
% 设 $f(x)$ 在区间 $I$ 上连续。若 $\alpha,\beta\in I$，$\alpha<\beta$，满足 $f(\alpha)f(\beta)<0$，则存在 $\xi\in(\alpha,\beta)$，使得 $f(\xi)=0$.
% \end{theorem}
% 
% 
% \begin{proof}
% 不妨设 $f(\alpha)<0$，$f(\beta)>0$，则有
% \[
%   m=\min_{x\in[\alpha,\beta]}\{f(x)\}<0<\max_{x\in[\alpha,\beta]}\{f(x)\}=M.
% \]
% 在闭区间 $[\alpha,\beta]$ 上应用介值定理，知存在 $\xi\in[\alpha,\beta]$，使得 $f(\xi)=0$。但 $f(\alpha)<0$，而 $f(\beta)>0$，所以必有 $\xi\in(\alpha,\beta)$ 成立.
% \end{proof}
% 
% 
% \begin{remark}\label{remark:ma-3-source-102}
% 这一定理表明：任何使连续函数 $f(x)$ 的函数值异号的两点之间必有方程 $f(x)=0$ 的一个根。我们常可以用它来证明非线性方程之解的存在性，或者用它来求方程的近似解。显然，定理~\ref{theorem:ma-3-3-4} 是定理~\ref{theorem:ma-3-3-3} 的直接推论。但由于零点存在定理有特别的重要性，我们还是将它叙述成一个单独的定理.
% \end{remark}
% 
% 
% \begin{example}\label{example:ma-3-3-4}
% 证明三次方程 $ax^3+bx^2+cx+d=0$ 至少有一实根.
% \end{example}
% 
% 
% \begin{proof}
% 无妨设 $a>0$。由于
% \[
%   f(x)=ax^3+bx^2+cx+d
%   =
%   ax^3\left(1+\frac b{ax}+\frac c{ax^2}+\frac d{ax^3}\right),
% \]
% 我们有 $\lim\limits_{x\to-\infty}f(x)=-\infty$ 和 $\lim\limits_{x\to+\infty}f(x)=+\infty$。因此易知，必有 $\alpha<0$ 和 $\beta>0$，使 $f(\alpha)<0$，$f(\beta)>0$。在闭区间 $[\alpha,\beta]$ 上应用定理~\ref{theorem:ma-3-3-4} 即知在 $(\alpha,\beta)$ 内必有 $f(x)$ 的一个零点，即该三次方程至少有一个实根.
% \end{proof}
% 
% 
% \begin{example}\label{example:ma-3-3-5}
% 证明方程 $x^4+2x-1=0$ 在 $(0,1)$ 中必有一根，并求它的一个近似解.
% \end{example}
% 
% 
% \begin{proof}
% 令 $f(x)=x^4+2x-1$，则有 $f(0)=-1<0$，$f(1)=2>0$。于是，在 $(0,1)$ 之内必有该方程的一个根。取 $[0,1]$ 的中点 $\dfrac12$，计算可得 $f\left(\dfrac12\right)=0.0625>0$。于是，在 $\left(0,\dfrac12\right)$ 之内必有该方程的一个根。再取 $\left[0,\dfrac12\right]$ 的中点 $\dfrac14$，计算可得 $f\left(\dfrac14\right)=-0.4961<0$。于是，在 $\left(\dfrac14,\dfrac12\right)$ 之内必有该方程的一个根。如此进行，我们可以断定，在 $\left(\dfrac{121}{256},\dfrac{61}{128}\right)$ 之内必有该方程的一个根，因此我们可取其中点
% \[
%   \hat x=\frac{243}{512}\approx0.4746
% \]
% 作为该方程的一个近似解，其误差为
% \[
%   |x-\hat x|<\frac1{512}\approx0.002.
% \qedhere
% \]
% \end{proof}
% 
% 
% \begin{example}\label{example:ma-3-3-6}
% 设函数 $f(x)\in C[a,b]$，且 $f^{-1}(x)$ 存在。证明 $f(x)$ 必在 $[a,b]$ 上严格单调.
% \end{example}
% 
% 
% \begin{proof}
% 倘若 $f(x)$ 在 $[a,b]$ 不严格单调，必存在
% \[
%   a\leqslant x_1<x_2<x_3\leqslant b,
% \]
% 使得以下两种情形之一发生：
% \[
%   \text{(1) } f(x_1)<f(x_2)\text{ 且 }f(x_2)>f(x_3);
% \]
% \[
%   \text{(2) } f(x_1)>f(x_2)\text{ 且 }f(x_2)<f(x_3).
% \]
% 不妨设情形 (1) 发生，并且取 $\eta$ 满足
% \[
%   \max\{f(x_1),f(x_3)\}<\eta<f(x_2),
% \]
% 则由介值定理推知，存在 $\xi_1\in(x_1,x_2)$ 和 $\xi_2\in(x_2,x_3)$，使得 $f(\xi_1)=\eta=f(\xi_2)$。这与 $f^{-1}(x)$ 存在相矛盾.
% \end{proof}
% 
% 
% \begin{remark}\label{remark:ma-3-source-109}
% 前面我们已经指出（见例~\ref{example:ma-3-2-9} 的注）：若单调函数 $f(x)$ 在 $I$ 内不连续，则它的取值不能是一个区间。将这一事实与连续函数的介值性定理相结合便有：若 $f(x)$ 在区间 $I$ 上单调，则 $f(x)$ 在 $I$ 上连续的充分必要条件是 $f(I)$ 是一个区间.
% \end{remark}
% 
% 以下我们来讨论函数的一致连续性问题。为此先来考查函数 $f(x)=\dfrac1x$，$x\in(0,1)$。我们知道 $f(x)$ 在 $(0,1)$ 是连续的。换句话说，对每个 $x_0\in(0,1)$，对任意给定的 $\varepsilon>0$，都存在 $\delta>0$，使得当 $|x-x_0|<\delta$ 时，有 $\left|\dfrac1x-\dfrac1{x_0}\right|<\varepsilon$。这里有个值得注意的现象是，对于同一个 $\varepsilon$，当 $x_0$ 与 $0$ 越近时，相应得到的 $\delta$ 也就越小。换句话说，我们不可能找到一个仅依赖 $\varepsilon$，而不依赖 $x_0$ 的 $\delta$，使得对所有的 $x_0\in(0,1)$，有 $\left|\dfrac1x-\dfrac1{x_0}\right|<\varepsilon$.
% 
% 基于这一现象，我们引入如下的概念：
% 
% \begin{definition}\label{definition:ma-3-3-1}
% 设函数 $f(x)$ 在区间 $I$ 上有定义。若 $\forall\varepsilon>0$，$\exists\delta>0$，当 $x_1,x_2\in I$ 且 $|x_1-x_2|<\delta$ 时，有 $|f(x_1)-f(x_2)|<\varepsilon$，则称 $f(x)$ 在 $I$ 上一致连续.
% \end{definition}
% 
% 显然，$f(x)$ 在 $I$ 上一致连续，它必在 $I$ 上连续.
% 
% \begin{example}\label{example:ma-3-3-7}
% 证明函数 $y=\sin x$ 在 $(-\infty,+\infty)$ 上一致连续.
% \end{example}
% 
% 
% \begin{proof}
% 由于对任意的 $x_1,x_2\in(-\infty,+\infty)$ 有
% \[
%   |\sin x_1-\sin x_2|
%   =
%   2\left|\cos\frac{x_1+x_2}{2}\right|
%   \left|\sin\frac{x_1-x_2}{2}\right|
%   \leqslant |x_1-x_2|,
% \]
% 所以 $\forall\varepsilon>0$，取 $\delta=\varepsilon$，当 $x_1,x_2\in(-\infty,+\infty)$ 且 $|x_1-x_2|<\delta$ 时，有 $|\sin x_1-\sin x_2|<\varepsilon$。这就证明了 $y=\sin x$ 在 $(-\infty,+\infty)$ 上一致连续.
% 
% 对于判别一个函数是否一致连续，下面的定理有时是非常方便的.
% \end{proof}
% 
% 
% \begin{theorem}\label{theorem:ma-3-3-5}
% 设函数 $f(x)$ 在区间 $I$ 上有定义，则 $f(x)$ 在 $I$ 一致连续的充分必要条件是：对任意的两个序列 $\{x_n'\}\subset I,\{x_n''\}\subset I$，若它们满足 $\lim\limits_{n\to\infty}(x_n'-x_n'')=0$，必有
% \[
%   \lim_{n\to\infty}[f(x_n')-f(x_n'')]=0.
% \]
% 此定理的证明留给读者自己完成.
% \end{theorem}
% 
% 
% \begin{example}\label{example:ma-3-3-8}
% 证明函数 $f(x)=\dfrac1x$ 在 $(0,+\infty)$ 上不一致连续；但对任意的 $x_0>0$，它在 $[x_0,+\infty)$ 上一致连续.
% \end{example}
% 
% 
% \begin{proof}
% 由定理~\ref{theorem:ma-3-3-5} 我们注意到，“$f(x)$ 在 $I$ 上不一致连续”与以下断言等价：存在 $\varepsilon_0>0$ 及序列 $\{x_n'\},\{x_n''\}\subset I$ 满足 $x_n'-x_n''\to0$（$n\to\infty$）及 $|f(x_n')-f(x_n'')|\geqslant\varepsilon_0$.
% 
% 对 $f(x)=\dfrac1x$ 而言，我们取 $\varepsilon_0=1$，$x_n'=\dfrac1n$，$x_n''=\dfrac1{n+1}$（$n=1,2,\cdots$），则有
% \[
%   x_n'-x_n''\to0\quad(n\to\infty),
% \]
% 且
% \[
%   \left|\frac1{x_n'}-\frac1{x_n''}\right|=1.
% \]
% 此即证明了 $\dfrac1x$ 在 $(0,+\infty)$ 上不一致连续.
% 
% 对任意给定的 $x_0>0$，由于对 $x_1,x_2\in[x_0,+\infty)$，有
% \[
%   \left|\frac1{x_1}-\frac1{x_2}\right|
%   =
%   \frac{|x_1-x_2|}{x_1x_2}
%   \leqslant \frac{|x_1-x_2|}{x_0^2},
% \]
% 所以 $\forall\varepsilon>0$，取 $\delta=x_0^2\varepsilon>0$，则当 $x_1,x_2\in[x_0,+\infty)$ 且 $|x_1-x_2|<\delta$ 时，就有
% \[
%   \left|\frac1{x_1}-\frac1{x_2}\right|<\varepsilon.
% \]
% 此时的 $\delta$ 只依赖于 $\varepsilon$，所以 $f(x)$ 在 $[x_0,+\infty)$ 上一致连续.
% \end{proof}
% 
% 
% \begin{theorem}[{康托尔定理}]\label{theorem:ma-3-3-6}
% 设函数 $f(x)\in C[a,b]$，则 $f(x)$ 在闭区间 $[a,b]$ 上一致连续.
% \end{theorem}
% 
% 
% \begin{proof}
% 定义
% \[
%   F(x)=
%   \begin{cases}
%     f(a), & x\in(-\infty,a),\\
%     f(x), & x\in[a,b],\\
%     f(b), & x\in(b,+\infty),
%   \end{cases}
% \]
% 则 $F(x)\in C(-\infty,+\infty)$。$\forall\varepsilon>0$ 和 $x\in[a,b]$，由 $F(x)$ 在点 $x$ 处的连续性知，$\exists\delta_x>0$，当 $x',x''\in U(x,2\delta_x)$ 时，有
% \[
%   |F(x')-F(x'')|<\varepsilon.
% \]
% 显然，开区间簇 $\{U(x,\delta_x):x\in[a,b]\}$ 是闭区间 $[a,b]$ 的一个开覆盖。由有限覆盖定理知，必有有限多个这样的开区间
% \[
%   U(x_1,\delta_{x_1}),\quad U(x_2,\delta_{x_2}),\quad \cdots,\quad U(x_m,\delta_{x_m}),
% \]
% 它们也完全覆盖了 $[a,b]$。令 $\delta=\min\{\delta_{x_j}:j=1,2,\cdots,m\}$，则 $\delta>0$，而且当 $x',x''\in[a,b]$ 且 $|x'-x''|<\delta$ 时，必存在某一 $U(x_j,\delta_{x_j})$，使得 $x'\in U(x_j,\delta_{x_j})$，即 $|x'-x_j|<\delta_{x_j}$，从而
% \[
%   |x''-x_j|\leqslant |x''-x'|+|x'-x_j|<\delta+\delta_{x_j}\leqslant2\delta_{x_j},
% \]
% 于是有 $x',x''\in U(x_j,2\delta_{x_j})$。这样，由 $\delta_{x_j}$ 的定义，有
% \[
%   |f(x')-f(x'')|=|F(x')-F(x'')|<\varepsilon.
% \]
% 注意到 $\delta$ 只与 $\varepsilon$ 有关而与 $x',x''$ 无关，这就证明了 $f(x)$ 在 $[a,b]$ 上一致连续.
% \end{proof}
% 
% 
% \begin{example}\label{example:ma-3-3-9}
% 设函数 $f(x)$ 在 $(a,b)$ 上连续。证明：$f(x)$ 在 $(a,b)$ 上一致连续的充分必要条件是 $\lim\limits_{x\to a+0}f(x)$ 与 $\lim\limits_{x\to b-0}f(x)$ 都存在.
% \end{example}
% 
% 
% \begin{proof}
% 必要性\quad 若 $f(x)$ 在 $(a,b)$ 上一致连续，由一致连续的定义及函数极限的柯西收敛准则知，$\lim\limits_{x\to a+0}f(x)$ 与 $\lim\limits_{x\to b-0}f(x)$ 都存在.
% 
% 充分性\quad 令
% \[
%   f(a)=\lim_{x\to a+0}f(x),\qquad f(b)=\lim_{x\to b-0}f(x),
% \]
% 并且定义
% \[
%   \tilde f(x)=
%   \begin{cases}
%     f(a), & x=a,\\
%     f(x), & x\in(a,b),\\
%     f(b), & x=b,
%   \end{cases}
% \]
% 则 $\tilde f\in C[a,b]$。于是由定理~\ref{theorem:ma-3-3-6} 知，它在 $[a,b]$ 上一致连续，从而 $f(x)$ 在 $(a,b)$ 上一致连续.
% \end{proof}
% 
% 
% \begin{remark}\label{remark:ma-3-source-120}
% 从以上例子可知，要考查一个连续函数在有穷区间上是否一致连续，只需考查它在端点的单侧极限是否存在即可。从这个意义上来讲，连续函数在有穷区间的一致连续性问题已经圆满解决。特别地，若 $f(x)$ 在有穷区间 $I$ 上无界，则它在 $I$ 上必定不一致连续.
% \end{remark}
% 
% 当函数 $f(x),g(x)$ 在区间 $I$ 上一致连续时，显然 $f(x)\pm g(x)$ 在区间 $I$ 上也一致连续.
% 
% \paragraph*{思考题}
% 设 $f(x),g(x)$ 在区间 $I$ 上一致连续，是否 $f(x)g(x)$ 也在 $I$ 上一致连续？
% 
% \begin{example}\label{example:ma-3-3-10}
% 证明 $f(x)=\sqrt{x}$ 在 $[0,+\infty)$ 上一致连续.
% \end{example}
% 
% 
% \begin{proof}
% 显然 $f(x)$ 在 $[0,+\infty)$ 连续。下面我们证明 $f(x)=\sqrt{x}$ 在 $[0,+\infty)$ 一致连续.
% 
% 事实上，任取 $\varepsilon>0$，取 $\delta_1=\varepsilon$，则当 $x',x''\in[1,+\infty)$ 且 $|x'-x''|<\delta_1$ 时，有
% \[
%   |f(x')-f(x'')|
%   =
%   |\sqrt{x'}-\sqrt{x''}|
%   =
%   \frac{|x'-x''|}{|\sqrt{x'}+\sqrt{x''}|}
%   \leqslant |x'-x''|<\varepsilon.
% \]
% 其次，由康托尔定理，$f(x)$ 在 $[0,2]$ 上连续，从而一致连续。因此对上述 $\varepsilon>0$，$\exists\delta_2>0$，使得当 $x',x''\in[0,2]$ 且 $|x'-x''|<\delta_2$ 时，有
% \[
%   |f(x')-f(x'')|<\varepsilon.
% \]
% 因此取 $\delta=\min\left\{\delta_1,\delta_2,\dfrac12\right\}$，对任意的 $x',x''\in[0,+\infty)$，当 $|x'-x''|<\delta$ 时，必有 $x',x''$ 同时落在 $[0,2]$ 或 $[1,+\infty)$ 上。因此，由上所证，有
% \[
%   |f(x')-f(x'')|<\varepsilon.
% \]
% 这就证明了 $f(x)=\sqrt{x}$ 在 $[0,+\infty)$ 上一致连续.
% \end{proof}
% 
% 
% \section{无穷小量与无穷大量的阶}
% \label{sec:ma-3-4}
% 
% 首先，我们以 $x\to x_0$ 为规范，给出与无穷小量和无穷大量有关的一些定义和记号.
% 
% \begin{definition}\label{definition:ma-3-4-1}
% 设函数 $f(x)$ 在 $U_0(x_0,\delta_0)(\delta_0>0)$ 上有定义。若 $\lim\limits_{x\to x_0}f(x)=0$，则称 $f(x)$ 为 $x\to x_0$ 时的一个无穷小量；若 $\lim\limits_{x\to x_0}f(x)=\infty$，则称 $f(x)$ 为 $x\to x_0$ 时的一个无穷大量.
% \end{definition}
% 
% 无穷小量有特殊的意义，在数学分析产生与发展的历史上，它起过特别重要的作用。这是因为无穷小量不仅是一种特殊的有极限的变量，而且由于：$x\to x_0$ 时 $f(x)$ 以 $A$ 为极限的充分必要条件是 $f(x)-A$ 为 $x\to x_0$ 时的一个无穷小量，从而无穷小量可以表征一般有极限的变量.
% 
% 关于无穷小量与无穷大量之间有如下的关系：
% 
% \begin{theorem}\label{theorem:ma-3-4-1}
% 设函数 $f(x)$ 在 $U_0(x_0,\delta_0)(\delta_0>0)$ 上有定义且恒不为零，则 $f(x)$ 为 $x\to x_0$ 时的一个无穷小量的充分必要条件是 $\dfrac1{f(x)}$ 为 $x\to x_0$ 时的一个无穷大量.
% \end{theorem}
% 
% 
% 证明略.
% 
% 在前面我们已经知道 $\left\{\dfrac1n\right\}$ 和 $\left\{\dfrac1{n^n}\right\}$ 都是无穷小量，而且后者趋于零的速度比前者要“快”。如何来刻画它们趋于零的速度是一个值得研究的问题。为了便于比较无穷小量（无穷大量），我们引入：
% 
% \begin{definition}\label{definition:ma-3-4-2}
% 设 $f(x)$ 和 $g(x)$ 都是当 $x\to x_0$ 时的无穷小量（无穷大量），且 $g(x)\ne0$，$x\in U_0(x_0,\delta_0)(\delta_0>0)$.
% 
% (1) 若 $\lim\limits_{x\to x_0}\dfrac{f(x)}{g(x)}=0$，则称 $f(x)$ 是比 $g(x)$ 更高阶的无穷小量（更低阶的无穷大量），记为 $f(x)=o(g(x))$（$x\to x_0$），并读做 $f(x)$ 等于小欧 $g(x)$；
% 
% (2) 若 $\lim\limits_{x\to x_0}\dfrac{f(x)}{g(x)}=l\ne0$，则称 $f(x)$ 与 $g(x)$ 是同阶无穷小量（同阶无穷大量）；
% 
% (3) 若 $\lim\limits_{x\to x_0}\dfrac{f(x)}{g(x)}=1$，则称 $f(x)$ 与 $g(x)$ 是等价无穷小量（等价无穷大量），记为 $f(x)\sim g(x)$（$x\to x_0$）；
% 
% (4) 若存在 $M>0$ 和 $\delta>0$，使得 $\forall x\in U_0(x_0,\delta)$，有 $|f(x)|\leqslant M|g(x)|$ 成立，则记为 $f(x)=O(g(x))$（$x\to x_0$）（读做 $f(x)$ 等于大欧 $g(x)$）.
% \end{definition}
% 
% 有的时候我们可以在上述定义中不假定 $g(x)$ 是一个无穷小（大）量。如令 $g(x)\equiv1$，此时我们有以下记号
% \[
%   f(x)=O(1)\quad(x\to x_0).
% \]
% 由 (4) 我们知道此记号表示 $f(x)$ 在 $x_0$ 的某个去心邻域上有界；而记号
% \[
%   f(x)=o(1)\quad(x\to x_0)
% \]
% 表示 $\lim\limits_{x\to x_0}f(x)=0$，即 $f(x)$ 为 $x\to x_0$ 时的无穷小量.
% 
% 在前面所给出的定义和记号中将“$x_0$”换成“$x_0+0$”，“$x_0-0$”，“$-\infty$”，“$+\infty$”或“$\infty$”，就得到相应的极限过程的定义和记号。当然，此时对凡涉及“$x_0$ 的邻域”的地方都要做相应的修改.
% 
% 特别地，我们可以对序列引进上述记号，如
% \[
%   \frac1{n^n}=o\left(\frac1n\right)\quad(n\to\infty).
% \]
% 
% \begin{example}\label{example:ma-3-4-1}
% 函数 $y=x^n\sin\dfrac1x=O(x^n)$（$x\to0$）.
% \end{example}
% 
% 在学习无穷小量的比较时我们要注意以下的一个重要问题，即有时候不能简单地在两个无穷小量之间加上等号。例如：从 $x^2=o(x)$（$x\to0$）我们可以推出 $x^2=O(x)$（$x\to0$）。但我们不能简单地推出以下的等式：$o(x)=O(x)$（$x\to0$）。读者很容易举出这种等式不成立的例子.
% 
% \begin{example}\label{example:ma-3-4-2}
% 设 $f(x)=(x-x_0)^m$，$g(x)=(x-x_0)^n$，则有
% \[
%   \lim_{x\to x_0}\frac{f(x)}{g(x)}
%   =
%   \begin{cases}
%     0, & m>n,\\
%     1, & m=n,\\
%     \infty, & m<n.
%   \end{cases}
% \]
% 这说明：对于极限过程 $x\to x_0$，正整数 $k$ 越大时，无穷小量 $(x-x_0)^k$ 的阶就越高.
% \end{example}
% 
% 
% \begin{example}\label{example:ma-3-4-3}
% 设 $f(x)=\dfrac1{(x-x_0)^m}$，$g(x)=\dfrac1{(x-x_0)^n}$，则有
% \[
%   \lim_{x\to x_0}\frac{f(x)}{g(x)}
%   =
%   \begin{cases}
%     0, & m<n,\\
%     1, & m=n,\\
%     \infty, & m>n.
%   \end{cases}
% \]
% 这说明：对于极限过程 $x\to x_0$，正整数 $k$ 越大时，无穷大量 $\dfrac1{(x-x_0)^k}$ 的阶就越高.
% \end{example}
% 
% 
% \begin{example}\label{example:ma-3-4-4}
% 设 $f(x)=x^\mu$（$\mu>0$），$g(x)=a^x$（$a>1$），证明
% \[
%   \lim_{x\to+\infty}\frac{f(x)}{g(x)}
%   =
%   \lim_{x\to+\infty}\frac{x^\mu}{a^x}=0.
% \]
% \end{example}
% 
% 
% \begin{proof}
% 先考虑 $\mu=k$ 为正整数的情形。令 $a=1+b$（$b>0$），则对任意的正整数 $n$，有
% \[
%   a^n=(1+b)^n
%   =
%   1+nb+\frac{n(n-1)}2b^2+\cdots+b^n
%   >
%   \frac{n(n-1)}2b^2,
% \]
% 从而有
% \[
%   0<\frac n{a^n}<\frac2{(n-1)b^2},
% \]
% 所以
% \[
%   \lim_{n\to\infty}\frac n{a^n}=0.
% \]
% 由此立即可得
% \[
%   \lim_{n\to\infty}\frac{n^k}{a^n}
%   =
%   \lim_{n\to\infty}
%   \left(\frac{n}{(\sqrt[k]{a})^n}\right)^k
%   =
%   0.
% \]
% 这样，由于
% \[
%   0<\frac{x^k}{a^x}\leqslant \frac{(1+[x])^k}{a^{[x]}},
% \]
% 所以我们有
% \[
%   \lim_{x\to+\infty}\frac{x^k}{a^x}=0.
% \]
% 对于一般的 $\mu>0$，我们取一个正整数 $k>\mu$，使得
% \[
%   0<\frac{x^\mu}{a^x}\leqslant\frac{x^k}{a^x}\quad(x>1),
% \]
% 从而有
% \[
%   \lim_{x\to+\infty}\frac{x^\mu}{a^x}=0.
% \]
% 本例说明：对于极限过程 $x\to+\infty$，指数函数 $a^x$（$a>1$）是比任何幂函数 $x^\mu$ 都高阶的无穷大量.
% \end{proof}
% 
% 
% \begin{example}\label{example:ma-3-4-5}
% 设 $f(x)=\log_a x$（$a>1$），$g(x)=x^\mu$（$\mu>0$），证明
% \[
%   \lim_{x\to+\infty}\frac{f(x)}{g(x)}
%   =
%   \lim_{x\to+\infty}\frac{\log_a x}{x^\mu}=0.
% \]
% \end{example}
% 
% 
% \begin{proof}
% 令 $y=\log_a x$，则 $x=a^y$。由此当 $x\to+\infty$ 时，$y\to+\infty$，有
% \[
%   \lim_{x\to+\infty}\frac{\log_a x}{x^\mu}
%   =
%   \lim_{y\to+\infty}\frac{y}{(a^\mu)^y}=0.
% \]
% 本例说明：对于极限过程 $x\to+\infty$，对数函数 $\log_a x$（$a>1$）是比任何幂函数 $x^\mu$ 都低阶的无穷大量.
% 
% 当 $x\to0$ 时，我们有以下重要的例子。这些例子在今后的学习中会经常用到，读者应该熟练掌握它们.
% \end{proof}
% 
% 
% \begin{example}\label{example:ma-3-4-6}
% 证明：当 $x\to0$ 时，下列关系式成立：
% \[
%   \text{(1) }\sin x\sim x;\qquad \sin x=x+o(x).
% \]
% \[
%   \text{(2) }1-\cos x\sim\frac12x^2;\qquad
%   \cos x=1-\frac12x^2+o(x^2).
% \]
% \[
%   \text{(3) }\tan x\sim x;\qquad \tan x=x+o(x).
% \]
% \[
%   \text{(4) }\arcsin x\sim x;\qquad \arcsin x=x+o(x).
% \]
% \[
%   \text{(5) }\ln(1+x)\sim x;\qquad \ln(1+x)=x+o(x).
% \]
% \[
%   \text{(6) }e^x-1\sim x;\qquad e^x=1+x+o(x).
% \]
% \[
%   \text{(7) }(1+x)^\mu-1\sim\mu x;\qquad
%   (1+x)^\mu=1+\mu x+o(x).
% \]
% \end{example}
% 
% 
% \begin{proof}
% (1) $\lim\limits_{x\to0}\dfrac{\sin x}{x}=1$.
% 
% (2)
% \[
%   \lim_{x\to0}\frac{1-\cos x}{\frac12x^2}
%   =
%   \lim_{x\to0}\frac{2\sin^2\frac x2}{\frac12x^2}
%   =
%   \lim_{x\to0}\left(\frac{\sin\frac x2}{\frac x2}\right)^2
%   =
%   1.
% \]
% 
% (3)
% \[
%   \lim_{x\to0}\frac{\tan x}{x}
%   =
%   \lim_{x\to0}\frac{\sin x}{x}\frac1{\cos x}
%   =
%   1.
% \]
% 
% (4) 令 $t=\arcsin x$，则 $x=\sin t$，且当 $x\to0$ 时，$t\to0$。于是，有
% \[
%   \lim_{x\to0}\frac{\arcsin x}{x}
%   =
%   \lim_{t\to0}\frac t{\sin t}
%   =
%   1.
% \]
% 
% (5)
% \[
%   \lim_{x\to0}\frac{\ln(1+x)}x
%   =
%   \lim_{x\to0}\ln(1+x)^{1/x}
%   =
%   \ln e
%   =
%   1.
% \]
% 
% (6) 令 $e^x-1=t$，则 $x=\ln(1+t)$，且当 $x\to0$ 时，$t\to0$。因此
% \[
%   \lim_{x\to0}\frac{e^x-1}{x}
%   =
%   \lim_{t\to0}\frac t{\ln(1+t)}
%   =
%   1.
% \]
% 
% (7)
% \[
%   \lim_{x\to0}\frac{(1+x)^\mu-1}{\mu x}
%   =
%   \lim_{x\to0}
%   \left(\frac{e^{\mu\ln(1+x)}-1}{\mu\ln(1+x)}\right)
%   \left(\frac{\ln(1+x)}x\right)
%   =
%   1.
% \]
% 
% 设 $f(x)$ 与 $g(x)$ 是当 $x\to x_0$ 时的等价无穷小量，则 $\dfrac{f(x)}{g(x)}-1=o(1)$（$x\to x_0$）。由此推出 $f(x)=g(x)+o(g(x))$（$x\to x_0$），从而上述各结论的相应等式成立.
% 
% 下面的定理说明，在求乘积或者商的极限时，可以将任何一个无穷小（大）的因式用它的等价因式来替换.
% \end{proof}
% 
% 
% \begin{theorem}\label{theorem:ma-3-4-2}
% 若无穷小（大）量 $h(x)\sim \ell(x)$（$x\to x_0$），则有
% \[
%   \text{(1) }\lim_{x\to x_0}h(x)f(x)=\lim_{x\to x_0}\ell(x)f(x);
% \]
% \[
%   \text{(2) }\lim_{x\to x_0}\frac{h(x)f(x)}{g(x)}
%   =
%   \lim_{x\to x_0}\frac{\ell(x)f(x)}{g(x)};
% \]
% \[
%   \text{(3) }\lim_{x\to x_0}\frac{f(x)}{h(x)g(x)}
%   =
%   \lim_{x\to x_0}\frac{f(x)}{\ell(x)g(x)}.
% \]
% 这里，我们假定所有的函数都在点 $x_0$ 的某个去心邻域上有定义，作为分母的函数在这个去心领域上不为 $0$，并且假定各式右端的极限存在.
% \end{theorem}
% 
% 
% \begin{proof}
% 我们这里只给出结论 (1) 的证明，其余的请读者自己补出。事实上，我们有
% \[
%   \lim_{x\to x_0}h(x)f(x)
%   =
%   \lim_{x\to x_0}\frac{h(x)}{\ell(x)}\ell(x)f(x)
%   =
%   \lim_{x\to x_0}\ell(x)f(x).
% \qedhere
% \]
% \end{proof}
% 
% 
% \begin{example}\label{example:ma-3-4-7}
% 求当 $x\to0$ 时无穷小量 $\ln\cos 2x$ 的形如 $ax^\alpha$ 的等价无穷小量.
% \end{example}
% 
% 
% \begin{solve}
% 由于
% \[
%   \ln\cos 2x
%   =
%   \ln(1-2\sin^2x)
%   \sim
%   -2\sin^2x
%   \sim
%   -2x^2,
% \]
% 所以
% \[
%   \ln\cos 2x\sim -2x^2.
% \qedhere
% \]
% \end{solve}
% 
% 
% \begin{remark}\label{remark:ma-3-source-139}
% 如果不是乘积和商的极限，则一般不能用等价无穷小（大）量来代替。例如：
% \[
%   \lim_{x\to0}\frac{\tan x-\sin x}{x^3}
%   =
%   \frac12
%   \ne
%   \lim_{x\to0}\frac{x-x}{x^3}
%   =
%   0.
% \]
% \end{remark}
% 
% 
% \section*{习题三}
% \phantomsection
% \addcontentsline{toc}{section}{习题三}
% \label{exercises:ma-3}
% 
% \noindent 1。证明函数极限的定义与下述的 (1) 或 (2) 均等价：
% 
% (1) 对任意给定的正整数 $n$，存在正数 $\delta$，使当 $0<|x-x_0|<\delta$ 时，
% \[
%   |f(x)-A|<\frac1n;
% \]
% 
% (2) 对任意给定的正数 $\varepsilon$，存在正整数 $n_0$，使当 $0<|x-x_0|<\dfrac1{n_0}$ 时，$|f(x)-A|<\varepsilon$.
% 
% \noindent 2。用极限的定义证明下列极限：
% \[
%   \text{(1) }\lim_{x\to a}|x|=|a|;\qquad
%   \text{(2) }\lim_{x\to3}x^3=27;\qquad
%   \text{(3) }\lim_{x\to4}\sqrt{x}=2.
% \]
% 
% \noindent 3。叙述下列极限的定义：
% 
% \begin{remark}
% 原稿在本题处截断，未提供完整题面。
% \end{remark}
% 
% % 整合来源：math-4.tex
% \chapter{导数与微分}
% \label{chap:ma-4}
% 
% 人们所说的 ``微积分'' 实际上包括 ``微分学'' 和 ``积分学'' 两部分。这一章我们就来讲述一元函数的微分学.
% 
% \section{导数}
% \label{sec:ma-4-1}
% 
% \subsection{导数概念的引入}
% \label{subsec:ma-4-1-1}
% 
% 导数是微分学的核心, 历史上这一概念是由对如下两个问题的研究而产生的:
% \begin{enumerate}
%   \item 求变速直线运动的瞬时速度;
%   \item 求曲线上一点处的切线.
% \end{enumerate}
% 这两个问题的解决最终都归结为求一个函数的变化率的问题。牛顿 (Newton) 从第一个问题出发, 莱布尼茨 (Leibniz) 从第二个问题出发, 分别给出了导数的概念.
% 
% \paragraph*{求变速直线运动的瞬时速度}
% 
% 
% 通常人们所说的物体的运动速度, 一般是指物体在某段时间内的平均速度。事实上, 任何物体的运动都很难保持匀速状态, 其速度每时每刻都在变化着。像火箭升空, 火车行驶, 都是从静止慢慢加速而达到某种速度的。此外, 即使以一定速度运动的物体, 也会根据运动过程中遇到的各种情况, 随时调整它的运行速度。因此, 一般来说平均速度只是对运动物体快慢程度的粗略描述, 并不能真正反映物体每一时刻运动的快慢程度。随着科学技术的发展, 仅仅知道运动物体的平均速度是不够的, 而是需要精确地给出其每一时刻的速度 --- 瞬时速度。那么, 如何来计算瞬时速度呢? 一般来说, 人们易于测量的是沿直线运动的物体离开初始位置的距离, 因此一个自然的想法就是利用物体所走过的路程来求瞬时速度.
% 
% 设物体沿直线朝一个方向作非匀速运动, 其运动规律由函数 $s=s(t)\ (t\in[0,T])$ 给出, 其中 $t$ 是时间, $s$ 是物体在时间 $t$ 内所走过的路程。现在要求它在时刻 $t_0\in(0,T)$ 的瞬时速度 $v(t_0)$.
% 
% 显然, 用 $[0,T]$ 上的平均速度来代替 $v(t_0)$, 一般来说误差应很大。一个朴素的想法是, $v(t_0)$ 的大小应该只与 $t_0$ 附近物体移动的路程的改变量有关。为此, 给 $t_0$ 一个微小的改变量 $\Delta t\ne0$, 物体在 $\Delta t$ 这段时间内所经过的路程为
% \[
%   \Delta s=s(t_0+\Delta t)-s(t_0).
% \]
% 物体在 $\Delta t$ 这段时间内的平均速度为
% \[
%   v=\frac{s(t_0+\Delta t)-s(t_0)}{\Delta t}.
% \]
% 容易看出, $\Delta t$ 越接近于 $0$, 该平均速度就应该越接近瞬时速度 $v(t_0)$。于是, 在时刻 $t_0\in(0,T)$ 的瞬时速度 $v(t_0)$ 就应该是当 $\Delta t\to0$ 时, 平均速度 $v$ 的极限, 即
% \[
%   v(t_0)
%   =
%   \lim_{\Delta t\to0}\frac{\Delta s}{\Delta t}
%   =
%   \lim_{\Delta t\to0}
%   \frac{s(t_0+\Delta t)-s(t_0)}{\Delta t}.
% \]
% 上式既给出了瞬时速度的定义, 也给出了计算瞬时速度的方法.
% 
% \paragraph*{求曲线上一点处切线}
% 
% 
% 莱布尼茨在研究如何求曲线在一点的切线这一几何问题时, 同样遇到了类似的极限问题。现假定有一条平面曲线, 它是函数 $y=f(x)\ (x\in(a,b))$ 的图像, 并给定 $x_0\in(a,b)$。那么, 应该怎样来求此曲线在点 $P(x_0,f(x_0))$ 的切线呢 (参见图 4.1.1)?
% 
% 由于所求的切线经过点 $(x_0,f(x_0))$, 因此我们只要将其斜率确定即可。如同计算瞬时速度一样, 一个显然的事实是过该点的切线的斜率应该只与 $x_0$ 附近的 $x$ 所对应的函数值有关。为此, 我们在 $x_0$ 附近取一点 $x$, 并且观察过曲线上两点 $P(x_0,f(x_0))$ 和 $Q(x,f(x))$ 的割线。显然, 该割线的斜率为
% \[
%   \tilde{k}(x)=\frac{f(x)-f(x_0)}{x-x_0}.
% \]
% 当 $x$ 在 $x_0$ 附近变化时, 点 $Q$ 在曲线上变动, 割线的斜率也在变化。当 $x$ 趋于 $x_0$ 时, 点 $Q$ 沿曲线趋近于点 $P$, 割线 $PQ$ 就应该趋近于过点 $P$ 的切线 $PT$.
% 
% 因此, 如果令 $\Delta x=x-x_0,\Delta y=f(x)-f(x_0)$, 当 $\Delta x$ 趋于 $0$ 时, 我们自然定义割线 $PQ$ 的极限位置为该曲线过点 $P$ 的切线 $PT$。这样, 割线 $PQ$ 的斜率 $\tilde{k}(x)=\dfrac{\Delta y}{\Delta x}$ 的极限就是所求切线 $PT$ 的斜率, 即
% \[
%   k=
%   \lim_{\Delta x\to0}\tilde{k}(x)
%   =
%   \lim_{\Delta x\to0}\frac{\Delta y}{\Delta x}
%   =
%   \lim_{x\to x_0}\frac{f(x)-f(x_0)}{x-x_0}.
% \]
% 上面的讨论既给出了切线的定义, 也给出了计算切线斜率的方法。上述两个问题的实际意义完全不同: 一个是运动学中沿直线运动的物体的瞬时速度, 一个是几何学中曲线在一点处的切线。但从数学上来看, 它们是完全一样的, 即都是归结为求函数的改变量 $\Delta y$ 与自变量的改变量 $\Delta x$ 之比的极限.
% 
% \subsection{导数的定义}
% \label{subsec:ma-4-1-2}
% 
% 除了上一小节的例子外, 还有许多实际问题也要遇到类似的求函数的改变量与自变量的改变量之比的极限问题。例如, 由质量函数求密度, 由人口增长函数求增长率等。略去各种背景, 便有以下纯数学的定义:
% 
% \begin{definition}\label{definition:ma-4-1-1}
% 设函数 $y=f(x)$ 在 $U(x_0,\delta_0)$ 内有定义。若极限
% \[
%   \lim_{x\to x_0}\frac{f(x)-f(x_0)}{x-x_0}
% \]
% 存在, 则称 $f(x)$ 在点 $x_0$ 处可导, 并且称此极限值为 $f(x)$ 在 $x_0$ 处的\keyterm{导数}, 记为 $f'(x_0)$ 或者 $\dfrac{df(x_0)}{dx}$, 有时也记为 $\left.\dfrac{df}{dx}\right|_{x=x_0}$ 或者 $\left.y'\right|_{x=x_0}$.
% \end{definition}
% 
% 由定义即知, 函数在一点的导数是它在该点附近各点平均变化率的极限, 是函数变化率的一种定量刻画。从几何上来看, 导数 $f'(x_0)$ 就是曲线 $y=f(x)$ 在点 $P(x_0,f(x_0))$ 的切线的斜率, 即切线与 $x$ 轴正向夹角的正切值 (参见图 4.1.1).
% 
% \begin{example}\label{example:ma-4-1-1}
% 求常数函数 $f(x)=C$ 在点 $x_0\in(-\infty,+\infty)$ 的导数.
% \end{example}
% 
% 
% \begin{solve}
% 对 $\forall x_0\in(-\infty,+\infty)$, 由于
% \[
%   \lim_{x\to x_0}\frac{f(x)-f(x_0)}{x-x_0}
%   =
%   \lim_{x\to x_0}\frac{C-C}{x-x_0}=0,
% \]
% 所以 $f'(x_0)=0$.
% \end{solve}
% 
% 
% \begin{example}\label{example:ma-4-1-2}
% 求 $f(x)=\sqrt{x}$ 在 $x=x_0>0$ 的导数.
% \end{example}
% 
% 
% \begin{solve}
% 由于
% \[
%   \lim_{x\to x_0}\frac{\sqrt{x}-\sqrt{x_0}}{x-x_0}
%   =
%   \lim_{x\to x_0}\frac1{\sqrt{x}+\sqrt{x_0}}
%   =
%   \frac1{2\sqrt{x_0}},
% \]
% 所以
% \[
%   f'(x_0)=\frac1{2\sqrt{x_0}}.
% \]
% 
% 若 $f(x)$ 在 $(a,b)$ 内每点都有导数, 则它的导数 $f'(x)$ 也是 $(a,b)$ 内的一个函数, 称为 $f(x)$ 的\keyterm{导函数}, 简称导数。此时, 如果我们要求 $f(x)$ 的导函数, 则我们先要设 $x_0$ 是任意固定的一点, 然后按定义求出极限, 再将 $x_0$ 换成 $x$。为了避免这一过程, 我们可以引入 $\Delta x=x-x_0$, 从而 $x=x_0+\Delta x$, $f(x)=f(x_0+\Delta x)$。因此, 我们有
% \[
%   f'(x)=
%   \lim_{\Delta x\to0}\frac{f(x+\Delta x)-f(x)}{\Delta x}.
% \]
% 这里 $\Delta x$ 是自变量的一个改变量, 而 $f(x+\Delta x)-f(x)$ 是函数的一个相应改变量, 记之为 $\Delta y$, 便有
% \[
%   f'(x)=\lim_{\Delta x\to0}\frac{\Delta y}{\Delta x}.
% \]
% 初学者应注意的是, 这里的 $\Delta x$ 是可正可负的, 而不是恒大于 $0$ 的.
% 
% 在推导基本初等函数的导数时, 我们需要如下的三个重要极限:
% \[
%   \text{(1) }\lim_{x\to0}\frac{\log_a(1+x)}{x}
%   =
%   \log_a e
%   =
%   \frac1{\ln a}\quad(a>0\text{且}a\ne1);
% \]
% \[
%   \text{(2) }\lim_{x\to0}\frac{a^x-1}{x}
%   =
%   \ln a\quad(a>0);
% \]
% \[
%   \text{(3) }\lim_{x\to0}\frac{(1+x)^\mu-1}{x}
%   =
%   \mu.
% \]
% 在上一章中我们已经给出这些结果的详细推导过程.
% \end{solve}
% 
% 
% \begin{example}\label{example:ma-4-1-3}
% 求 $f(x)=x^\alpha$ 的导数 (定义域与 $\alpha$ 有关, 对任何 $\alpha$, $x>0$ 总有定义).
% \end{example}
% 
% 
% \begin{solve}
% 设 $x\ne0$, 则有
% \[
%   \frac{f(x+\Delta x)-f(x)}{\Delta x}
%   =
%   \frac{(x+\Delta x)^\alpha-x^\alpha}{\Delta x}
%   =
%   x^{\alpha-1}
%   \frac{\left(1+\dfrac{\Delta x}{x}\right)^\alpha-1}
%        {\dfrac{\Delta x}{x}}.
% \]
% 对固定的 $x\ne0$, 由于当 $\Delta x\to0$ 时, $\dfrac{\Delta x}{x}\to0$, 从而推出
% \[
%   f'(x)=\alpha x^{\alpha-1}.
% \]
% 
% 特别地, 我们有
% \[
%   (x^{-1})'=-\frac1{x^2}\quad(x\ne0),
% \]
% \[
%   (\sqrt{x})'=\frac1{2\sqrt{x}}\quad(x>0).
% \qedhere
% \]
% \end{solve}
% 
% 
% \begin{example}\label{example:ma-4-1-4}
% 求 $f(x)=a^x\ (a>0,-\infty<x<+\infty)$ 的导数.
% \end{example}
% 
% 
% \begin{solve}
% 由于 $\forall x\in(-\infty,+\infty)$, 有
% \[
%   \frac{a^{x+\Delta x}-a^x}{\Delta x}
%   =
%   a^x\frac{a^{\Delta x}-1}{\Delta x}
%   \to a^x\ln a\quad(\Delta x\to0),
% \]
% 所以我们有 $f'(x)=a^x\ln a$。特别地, 我们有 $(e^x)'=e^x$.
% \end{solve}
% 
% 
% \begin{example}\label{example:ma-4-1-5}
% 求 $f(x)=\log_a x\ (0<a\ne1,0<x<+\infty)$ 的导数.
% \end{example}
% 
% 
% \begin{solve}
% 由于 $\forall x\in(0,+\infty)$, 有
% \[
%   \frac{\log_a(x+\Delta x)-\log_a x}{\Delta x}
%   =
%   \frac1x
%   \frac{\log_a\left(1+\dfrac{\Delta x}{x}\right)}
%        {\dfrac{\Delta x}{x}}
%   \to
%   \frac{\log_a e}{x}\quad(\Delta x\to0),
% \]
% 所以我们有
% \[
%   f'(x)=\frac{\log_a e}{x}=\frac1{x\ln a}.
% \]
% 特别地, 我们有 $(\ln x)'=\dfrac1x$.
% \end{solve}
% 
% 
% \begin{example}\label{example:ma-4-1-6}
% 求 $f(x)=\sin x\ (-\infty<x<+\infty)$ 的导数.
% \end{example}
% 
% 
% \begin{solve}
% 由于 $\forall x\in(-\infty,+\infty)$, 有
% \[
%   \frac{\sin(x+\Delta x)-\sin x}{\Delta x}
%   =
%   \frac{2\cos\left(x+\dfrac{\Delta x}{2}\right)\sin\dfrac{\Delta x}{2}}
%        {\Delta x}
%   \to
%   \cos x\quad(\Delta x\to0),
% \]
% 所以我们有 $f'(x)=\cos x$.
% \end{solve}
% 
% 
% \begin{example}\label{example:ma-4-1-7}
% 求 $f(x)=\cos x\ (-\infty<x<+\infty)$ 的导数.
% \end{example}
% 
% 
% \begin{solve}
% 由于 $\forall x\in(-\infty,+\infty)$, 有
% \[
%   \frac{\cos(x+\Delta x)-\cos x}{\Delta x}
%   =
%   \frac{-2\sin\left(x+\dfrac{\Delta x}{2}\right)\sin\dfrac{\Delta x}{2}}
%        {\Delta x}
%   \to
%   -\sin x\quad(\Delta x\to0),
% \]
% 所以我们有 $f'(x)=-\sin x$.
% \end{solve}
% 
% 
% \subsection{单侧导数}
% \label{subsec:ma-4-1-3}
% 
% 类似于单侧极限, 我们有以下单侧导数的概念.
% 
% \begin{definition}\label{definition:ma-4-1-2}
% 设函数 $y=f(x)$ 在 $U^+(x_0,\delta)$ 内有定义。若极限
% \[
%   \lim_{\Delta x\to0+0}
%   \frac{f(x_0+\Delta x)-f(x_0)}{\Delta x}
% \]
% 存在, 则称 $f(x)$ 在 $x_0$ 处右可导, 并且称此极限值为 $f(x)$ 在 $x_0$ 处的\keyterm{右导数}, 记为 $f'_+(x_0)$。同理可定义函数 $f(x)$ 在 $x_0$ 处的\keyterm{左可导}和\keyterm{左导数}, 即
% \[
%   f'_-(x_0)=
%   \lim_{\Delta x\to0-0}
%   \frac{f(x_0+\Delta x)-f(x_0)}{\Delta x}.
% \]
% \end{definition}
% 
% 
% \begin{remark}\label{remark:ma-4-source-17}
% 从导数的定义我们可以看出, $f'(x_0)$ 存在的充分必要条件是 $f'_+(x_0)$ 与 $f'_-(x_0)$ 都存在且相等.
% \end{remark}
% 
% 当 $f(x)$ 在 $x_0$ 处可导时, 有
% \[
%   f'(x_0)=
%   \lim_{\Delta x\to0}
%   \frac{f(x_0+\Delta x)-f(x_0)}{\Delta x}.
% \]
% 因此, $f'(x_0)$ 存在的一个必要条件是, 当 $\Delta x\to0$ 时, $\Delta y=f(x+\Delta x)-f(x)$ 必须趋于零。而这正是 $f(x)$ 在 $x_0$ 处连续的定义, 从而我们有
% 
% \begin{theorem}\label{theorem:ma-4-1-1}
% 若函数 $y=f(x)$ 在 $x_0$ 处可导, 则它在 $x_0$ 处连续.
% \end{theorem}
% 
% 这样, 一个自然的问题就是, 是否连续函数一定可导呢? 其答案是否定的。例如, 函数 $y=f(x)=|x|$ 在 $x=0$ 处是连续的, 但 $f(x)$ 在 $x=0$ 处不可导。这是由于当 $\Delta x\to0$ 时,
% \[
%   \frac{f(0+\Delta x)-f(0)}{\Delta x}
%   =
%   \frac{|\Delta x|}{\Delta x}
%   =
%   \begin{cases}
%     1, & \Delta x>0,\\
%     -1, & \Delta x<0
%   \end{cases}
% \]
% 的极限不存在 (参见图 4.1.2).
% 
% \begin{remark}\label{remark:ma-4-source-19}
% 从单侧导数的定义可以看出, $f(x)$ 在 $x_0$ 处左可导, 则 $f(x)$ 在 $x_0$ 处左连续; 右可导则右连续。于是 $f(x)$ 在 $x_0$ 左、右可导 (不要求左导数与右导数相等) 就可以推出 $f(x)$ 在 $x_0$ 连续.
% \end{remark}
% 
% 
% \begin{example}\label{example:ma-4-1-8}
% 令函数
% \[
%   S(x)=
%   \begin{cases}
%     \sin\dfrac1x, & x\ne0,\\
%     0, & x=0,
%   \end{cases}
% \]
% 则有
% \begin{enumerate}
%   \item $S(x)$ 在 $x=0$ 处间断, 是第二类间断点;
%   \item $xS(x)$ 在 $x=0$ 处连续, 但不可导 (其实, 其单侧导数也不存在);
%   \item $x^2S(x)$ 在 $x=0$ 处连续而且可导, 其导数为 $0$.
% \end{enumerate}
% 此外, 按照导数的定义, 曲线 $y=x^2\sin\dfrac1x$ 在 $(0,0)$ 处的切线方程为 $y=0$, 此时, 该切线在 $x=0$ 附近总是与曲线有交点, 这与中学里曲线的切线定义是有很大区别的.
% \end{example}
% 
% 
% \section{求导数的方法}
% \label{sec:ma-4-2}
% 
% 按照导数定义, 导数表示为当 $\Delta x\to0$ 时, $\dfrac{\Delta y}{\Delta x}$ 的极限。但直接用定义对各种函数求导数, 并不总是一件容易的事, 有时甚至很难办到。因此, 这一节我们来介绍各种行之有效的求导方法。利用这些求导方法就可以轻而易举地求出许多函数的导数.
% 
% \subsection{函数四则运算的导数}
% \label{subsec:ma-4-2-1}
% 
% \begin{theorem}\label{theorem:ma-4-2-1}
% 设函数 $f(x)$ 和 $g(x)$ 都在点 $x$ 处可导, 则下列各式在点 $x$ 处成立:
% \[
%   \text{(1) }\quad
%   (f(x)\pm g(x))'=f'(x)\pm g'(x);
% \]
% \[
%   \text{(2) }\quad
%   (f(x)g(x))'=f'(x)g(x)+f(x)g'(x);
% \]
% \[
%   \text{(3) }\quad
%   \left(\frac{f(x)}{g(x)}\right)'
%   =
%   \frac{f'(x)g(x)-f(x)g'(x)}{g^2(x)}
%   \quad(g(x)\ne0).
% \]
% \end{theorem}
% 
% 
% \begin{proof}
% (1) 令 $y=f(x)\pm g(x)$, 则有
% \[
%   (f(x)\pm g(x))'
%   =
%   \lim_{\Delta x\to0}\frac{\Delta y}{\Delta x}
% \]
% \[
%   =
%   \lim_{\Delta x\to0}
%   \frac{[f(x+\Delta x)\pm g(x+\Delta x)]-[f(x)\pm g(x)]}
%        {\Delta x}
% \]
% \[
%   =
%   \lim_{\Delta x\to0}\frac{f(x+\Delta x)-f(x)}{\Delta x}
%   \pm
%   \lim_{\Delta x\to0}\frac{g(x+\Delta x)-g(x)}{\Delta x}
%   =
%   f'(x)\pm g'(x).
% \]
% 
% (2) 令 $y=f(x)g(x)$, 则有
% \[
%   (f(x)g(x))'
%   =
%   \lim_{\Delta x\to0}\frac{\Delta y}{\Delta x}
% \]
% \[
%   =
%   \lim_{\Delta x\to0}
%   \frac{f(x+\Delta x)g(x+\Delta x)-f(x)g(x)}
%        {\Delta x}
% \]
% \[
%   =
%   \lim_{\Delta x\to0}
%   \frac{f(x+\Delta x)-f(x)}{\Delta x}g(x+\Delta x)
%   +
%   \lim_{\Delta x\to0}
%   f(x)\frac{g(x+\Delta x)-g(x)}{\Delta x}
% \]
% \[
%   =
%   f'(x)g(x)+f(x)g'(x).
% \]
% 
% (3) 令 $y=\dfrac1{g(x)}$, 则有
% \[
%   y'=
%   \lim_{\Delta x\to0}\frac{\Delta y}{\Delta x}
% \]
% \[
%   =
%   \lim_{\Delta x\to0}
%   \frac1{\Delta x}
%   \left[
%     \frac1{g(x+\Delta x)}-\frac1{g(x)}
%   \right]
% \]
% \[
%   =
%   -
%   \lim_{\Delta x\to0}
%   \frac{g(x+\Delta x)-g(x)}{\Delta x}
%   \frac1{g(x+\Delta x)g(x)}
%   =
%   -\frac{g'(x)}{g^2(x)}.
% \]
% 再由 $\dfrac{f(x)}{g(x)}=f(x)\dfrac1{g(x)}$ 及 (2), 我们有
% \[
%   \left(\frac{f(x)}{g(x)}\right)'
%   =
%   \frac{f'(x)}{g(x)}
%   +
%   f(x)\left(\frac1{g(x)}\right)'
%   =
%   \frac{f'(x)g(x)-f(x)g'(x)}{g^2(x)}.
% \qedhere
% \]
% \end{proof}
% 
% 
% \begin{corollary}\label{corollary:ma-4-source-23}
% 设 $f_i(x)\ (i=1,2,\cdots,n)$ 都在点 $x$ 处可导, 则下列各式在点 $x$ 处成立:
% \[
%   \text{(1) }\quad
%   \left(\sum_{i=1}^{n}c_if_i(x)\right)'
%   =
%   \sum_{i=1}^{n}c_if_i'(x),
%   \quad\text{其中 }c_i\text{ 是常数};
% \]
% \[
%   \text{(2) }\quad
%   \left(\prod_{i=1}^{n}f_i(x)\right)'
%   =
%   \sum_{k=1}^{n}
%   \left(
%     f_k'(x)\prod_{\substack{i=1\\i\ne k}}^{n}f_i(x)
%   \right).
% \]
% \end{corollary}
% 
% 
% \begin{example}\label{example:ma-4-2-1}
% 设 $f(x)=\tan x$, 求 $f'(x)$.
% \end{example}
% 
% 
% \begin{solve}
% \[
%   f'(x)=(\tan x)'=\left(\frac{\sin x}{\cos x}\right)'
% \]
% \[
%   =
%   \frac{(\sin x)'\cos x-\sin x(\cos x)'}{\cos^2 x}
%   =
%   \frac{\cos^2 x+\sin^2 x}{\cos^2 x}
%   =
%   \frac1{\cos^2 x}
%   =
%   \operatorname{sec}^2x.
% \qedhere
% \]
% \end{solve}
% 
% 
% \begin{example}\label{example:ma-4-2-2}
% 设 $f(x)=\cot x$, 求 $f'(x)$.
% \end{example}
% 
% 
% \begin{solve}
% \[
%   (\cot x)'=\left(\frac1{\tan x}\right)'
%   =
%   -\frac{(\tan x)'}{\tan^2 x}
% \]
% \[
%   =
%   -\frac1{\cos^2 x}\frac{\cos^2 x}{\sin^2 x}
%   =
%   -\frac1{\sin^2 x}
%   =
%   -\operatorname{csc}^2x.
% \qedhere
% \]
% \end{solve}
% 
% 
% \subsection{反函数的求导法则}
% \label{subsec:ma-4-2-2}
% 
% \begin{theorem}\label{theorem:ma-4-2-2}
% 设函数 $y=f(x)$ 在 $(a,b)$ 内严格单调, 并且令
% \[
%   \alpha=\min\{f(a+0),f(b-0)\},
%   \quad
%   \beta=\max\{f(a+0),f(b-0)\}.
% \]
% 如果 $f(x)$ 在 $(a,b)$ 内可导且导数 $f'(x)\ne0$, 则它的反函数 $x=f^{-1}(y)$ 在 $(\alpha,\beta)$ 内可导, 而且有
% \[
%   \frac{df^{-1}(y)}{dy}=\frac1{f'(x)}.
% \]
% \end{theorem}
% 
% 
% \begin{proof}
% 由于 $y=f(x)$ 在 $(a,b)$ 内可导, 从而它在 $(a,b)$ 上连续。再由其在 $(a,b)$ 内严格单调知, 它的反函数在 $(\alpha,\beta)$ 中连续且严格单调。任取 $y_0\in(\alpha,\beta)$, 则存在 $x_0\in(a,b)$, 使得 $y_0=f(x_0)$, 而且
% \[
%   \left.(f^{-1}(y))'\right|_{y=y_0}
%   =
%   \lim_{y\to y_0}\frac{f^{-1}(y)-f^{-1}(y_0)}{y-y_0}
% \]
% \[
%   =
%   \lim_{y\to y_0}\frac{x-x_0}{y-y_0}
%   =
%   \lim_{x\to x_0}
%   \frac1{\dfrac{f(x)-f(x_0)}{x-x_0}}
%   =
%   \frac1{f'(x_0)}.
% \qedhere
% \]
% \end{proof}
% 
% 
% \begin{example}\label{example:ma-4-2-3}
% 求 $y=\arctan x$ 的导数.
% \end{example}
% 
% 
% \begin{solve}
% 由 $y=\arctan x$, 得 $x=\tan y$。因此, 有
% \[
%   y'=(\arctan x)'=\frac1{(\tan y)'}
%   =\frac1{\sec^2 y}
%   =\frac1{1+\tan^2 y}
%   =\frac1{1+x^2}.
% \]
% 类似地, 有
% \[
%   (\operatorname{arccot}x)'=-\frac1{1+x^2}.
% \qedhere
% \]
% \end{solve}
% 
% 
% \begin{example}\label{example:ma-4-2-4}
% 求 $y=\arcsin x$ 的导数.
% \end{example}
% 
% 
% \begin{solve}
% 由 $x=\sin y$, 得
% \[
%   y'=(\arcsin x)'=\frac1{(\sin y)'}
%   =\frac1{\cos y}
%   =\frac1{\sqrt{1-\sin^2y}}
%   =\frac1{\sqrt{1-x^2}}.
% \]
% 类似地, 有
% \[
%   (\arccos x)'=-\frac1{\sqrt{1-x^2}}.
% \]
% 
% 在利用反函数求导法时, 读者需注意两点: 一是要注意每次求导是关于什么变量进行的; 二是如何将求导数后关于 $y$ 的函数转化为 $x$ 的函数.
% 
% 通过一系列例题, 我们已经求出了所有基本初等函数的导数。现将这些结果总结如下, 以便以后查阅:
% \[
% \begin{array}{ll}
%   \text{(1) }(C)'=0; &
%   \text{(2) }(x^\alpha)'=\alpha x^{\alpha-1};\\[4pt]
%   \text{(3) }(\sin x)'=\cos x; &
%   \text{(4) }(\cos x)'=-\sin x;\\[4pt]
%   \text{(5) }(\tan x)'=\dfrac1{\cos^2x}; &
%   \text{(6) }(\cot x)'=-\dfrac1{\sin^2x};\\[8pt]
%   \text{(7) }(\arcsin x)'=\dfrac1{\sqrt{1-x^2}}; &
%   \text{(8) }(\arccos x)'=-\dfrac1{\sqrt{1-x^2}};\\[8pt]
%   \text{(9) }(\arctan x)'=\dfrac1{1+x^2}; &
%   \text{(10) }(\operatorname{arccot}x)'=-\dfrac1{1+x^2};\\[8pt]
%   \text{(11) }(e^x)'=e^x; &
%   \text{(12) }(a^x)'=a^x\ln a\ (a>0);\\[4pt]
%   \text{(13) }(\ln x)'=\dfrac1x; &
%   \text{(14) }(\log_a x)'=\dfrac1{x\ln a}\quad\text{若 }(a>0\text{ 且 }a\ne1).
% \end{array}
% \qedhere
% \]
% \end{solve}
% 
% 
% \subsection{复合函数的求导法则}
% \label{subsec:ma-4-2-3}
% 
% \begin{theorem}\label{theorem:ma-4-2-3}
% 设函数 $y=f(u)$ 在 $U(u_0,\delta_0)$ 内有定义, 函数 $u=g(x)$ 在 $U(x_0,\eta_0)$ 内有定义, 且 $u_0=g(x_0)$。若 $f'(u_0)$ 与 $g'(x_0)$ 都存在, 则复合函数 $F(x)=f(g(x))$ 在点 $x_0$ 可导, 且 $F'(x_0)=f'(g(x_0))g'(x_0)$.
% \end{theorem}
% 
% 
% \begin{proof}
% 定义函数
% \[
%   A(u)=
%   \begin{cases}
%     \dfrac{f(u)-f(u_0)}{u-u_0}, & u\ne u_0,\\[8pt]
%     f'(u_0), & u=u_0,
%   \end{cases}
% \]
% 则 $A(u)$ 在 $u_0$ 点连续, 即有
% \[
%   \lim_{u\to u_0}A(u)=A(u_0)=f'(u_0).
% \]
% 由恒等式 $f(u)-f(u_0)=A(u)(u-u_0)$, 我们有
% \[
%   \frac{F(x)-F(x_0)}{x-x_0}
%   =
%   \frac{f(g(x))-f(g(x_0))}{x-x_0}
%   =
%   A(g(x))\frac{g(x)-g(x_0)}{x-x_0}.
% \]
% 在上式两边令 $x\to x_0$, 得
% \[
%   F'(x_0)=f'(g(x_0))g'(x_0).
% \qedhere
% \]
% \end{proof}
% 
% 
% \begin{remark}\label{remark:ma-4-source-36}
% 若 $f(u)$ 的定义域包含 $u=g(x)$ 的值域, 两函数在各自的定义域上可导, 则复合函数 $F(x)=f(g(x))$ 在 $g(x)$ 的定义域上可导, 且
% \[
%   F'(x)=f'(g(x))g'(x).
% \]
% 上述等式常常简记为
% \[
%   y_x'=y_u'u_x'
%   \quad\text{或者}\quad
%   \frac{dy}{dx}=\frac{dy}{du}\frac{du}{dx}.
% \]
% 我们一般称它为链锁法则.
% \end{remark}
% 
% 
% \begin{example}\label{example:ma-4-2-5}
% 分别求 $y=\sin x^2$ 和 $y=\sin^2x$ 的导数.
% \end{example}
% 
% 
% \begin{solve}
% $y=\sin x^2$ 可看成 $y=\sin u$ 和 $u=x^2$ 的复合, 因此
% \[
%   (\sin x^2)'=2x\cos x^2;
% \]
% 而 $y=\sin^2x$ 可看成 $y=u^2$ 和 $u=\sin x$ 的复合, 因此
% \[
%   (\sin^2x)'=2\sin x(\sin x)'=2\sin x\cos x=\sin2x.
% \qedhere
% \]
% \end{solve}
% 
% 
% \begin{example}\label{example:ma-4-2-6}
% 求 $y=\ln|x|\ (x\ne0)$ 的导数.
% \end{example}
% 
% 
% \begin{solve}
% 当 $x>0$ 时, $y=\ln x$, 从而
% \[
%   y'=(\ln x)'=\frac1x;
% \]
% 当 $x<0$ 时,
% \[
%   y=\ln|x|=\ln(-x),
% \]
% 于是, 它可看成由 $y=\ln u$ 与 $u=-x$ 复合而成, 从而有
% \[
%   (\ln|x|)'=(\ln(-x))'=\frac1{-x}(-x)'=\frac1x.
% \qedhere
% \]
% \end{solve}
% 
% 
% \begin{example}\label{example:ma-4-2-7}
% 求 $y=x^\alpha\ (x>0,\alpha\text{ 为任意实数})$ 的导数.
% \end{example}
% 
% 
% \begin{solve}
% 前面我们用定义已求出 $y'=\alpha x^{\alpha-1}$。现应用复合函数求导法于 $y=e^{\alpha\ln x}$, 则有
% \[
%   y'=(e^{\alpha\ln x})'
%   =
%   e^{\alpha\ln x}(\alpha\ln x)'
%   =
%   \frac{\alpha e^{\alpha\ln x}}x
%   =
%   \alpha\frac{x^\alpha}x
%   =
%   \alpha x^{\alpha-1}.
% \qedhere
% \]
% \end{solve}
% 
% 
% \begin{example}\label{example:ma-4-2-8}
% 证明: 当 $n\geqslant2$ 时, 有
% \[
%   \sum_{k=1}^{n}C_n^kk^2=n(n+1)2^{n-2}.
% \]
% \end{example}
% 
% 
% \begin{proof}
% 在恒等式
% \[
%   (1+x)^n=\sum_{k=0}^{n}C_n^kx^k
% \]
% 两边对 $x$ 求导数得
% \[
%   n(1+x)^{n-1}=\sum_{k=1}^{n}C_n^kkx^{k-1}.
% \]
% 在上述恒等式的两边乘以 $x$ 并对 $x$ 求导数得
% \[
%   n(1+x)^{n-1}+n(n-1)x(1+x)^{n-2}
%   =
%   \sum_{k=1}^{n}C_n^kk^2x^{k-1}.
% \]
% 在上述等式中令 $x=1$ 整理得
% \[
%   n(n+1)2^{n-2}=\sum_{k=1}^{n}C_n^kk^2.
% \qedhere
% \]
% \end{proof}
% 
% 
% \subsection{隐函数的求导法}
% \label{subsec:ma-4-2-4}
% 
% 前面我们所遇到的函数都是由显式给出的。但在解决实际问题时, 常常得到的只是变量 $x$ 和 $y$ 所满足的方程, 如:
% \[
%   x^2+y^2=1,\quad \tan x+\tan y=xy,\quad \ln y+\sin y=x^2y,
% \]
% 等等。通常在一定条件下, 对指定数集 $X$ 中的每一个数 $x$, 由变量 $x$ 和 $y$ 所满足的方程可以唯一地确定一个数 $y$ 与之对应。由这种对应法则所确定的函数 $y=y(x)$ 就称做由该方程所确定的隐函数。例如, 方程 $x^2+y^2=1$ 就隐含了两个定义在 $[-1,1]$ 上的函数 $y=f_1(x)=\sqrt{1-x^2}$ 和 $y=f_2(x)=-\sqrt{1-x^2}$。至于在什么条件下才能保证由所给出的二元方程唯一地确定一个连续可导的隐函数呢? 这要等到学习了多元函数的微分学才能给出解答。在本节的后继部分, 我们总是假定给定的二元方程所确定函数是可导的, 然后来讨论该函数的求导问题.
% 
% \begin{example}\label{example:ma-4-2-9}
% 证明开普勒 (Kepler) 方程
% \[
%   y-x-\varepsilon\sin y=0
%   \quad(0<\varepsilon<1)
% \]
% 可确定一个隐函数 $y=y(x)$, 并且求 $y'(x)$.
% \end{example}
% 
% 
% \begin{proof}
% 令 $x=\varphi(y)=y-\varepsilon\sin y$, 则有
% \[
%   \lim_{y\to-\infty}\varphi(y)=-\infty,
%   \quad
%   \lim_{y\to+\infty}\varphi(y)=+\infty;
% \]
% 而且对任意的 $y_1<y_2$, 有
% \[
%   \varphi(y_2)-\varphi(y_1)
%   =
%   y_2-y_1-\varepsilon(\sin y_2-\sin y_1)
% \]
% \[
%   =
%   y_2-y_1-\varepsilon\left(2\cos\frac{y_2+y_1}{2}\sin\frac{y_2-y_1}{2}\right)
%   \geqslant(1-\varepsilon)(y_2-y_1)>0.
% \]
% 因此, $\varphi(y)$ 是一严格单调递增的连续函数, 从而 $\forall x_0\in(-\infty,+\infty)$, 总存在唯一 $y_0$, 使得 $x_0=y_0-\varepsilon\sin y_0$。这表明, 该方程确定了 $(-\infty,+\infty)$ 上的一个函数 $y=y(x)$。以后我们还将知道, 开普勒方程确定的函数是可导的。下面我们来求出此函数的导数.
% 
% 既然 $y=y(x)$ 是由开普勒方程确定的函数, 则该函数满足恒等式
% \[
%   y(x)-x-\varepsilon\sin y(x)\equiv0.
% \]
% 两边对 $x$ 求导数后得恒等式
% \[
%   y'(x)-1-\varepsilon\cos y(x)y'(x)\equiv0,
% \]
% 从而有
% \[
%   y'(x)=\frac1{1-\varepsilon\cos y(x)}.
% \]
% 
% 这一例子表明, 虽然我们并没有解出函数的显式表达式, 但我们可以由它所满足的方程求出它的导数。当然, 导数的表达式中仍含有隐函数 $y(x)$。同时, 这一例子也告诉我们求这类函数导数的方法, 即只要将 $y$ 看成 $x$ 的函数, 在方程的两边对 $x$ 求导数, 就可得到 $y'(x)$ 满足的恒等式, 然后再从中将 $y'(x)$ 解出即可.
% \end{proof}
% 
% 
% \begin{example}\label{example:ma-4-2-10}
% 给定圆的方程 $x^2+y^2=1$, 求其斜率为 $-\dfrac12$ 的切线方程.
% \end{example}
% 
% 
% \begin{solve}
% 设所求的切线经过圆上一点 $(x_0,y_0)$。由隐函数求导法, 得 $2x+2yy'=0$, 从而有 $y'(x)=-\dfrac{x}{y}$。所以在切点 $(x_0,y_0)$ 处应有
% \[
%   y'(x_0)=-\frac{x_0}{y_0}=-\frac12,
% \]
% 即 $y_0=2x_0$。再由 $(x_0,y_0)$ 在圆上可知, $x_0=\pm\dfrac1{\sqrt5},y_0=\pm\dfrac2{\sqrt5}$。于是, 所求的切线有两条, 其分别过 $\left(\dfrac1{\sqrt5},\dfrac2{\sqrt5}\right)$ 和 $\left(-\dfrac1{\sqrt5},-\dfrac2{\sqrt5}\right)$ 的方程为
% \[
%   y=\frac2{\sqrt5}-\frac12\left(x-\frac1{\sqrt5}\right)
%   =
%   \frac{\sqrt5}2-\frac12x
% \]
% 和
% \[
%   y=-\frac2{\sqrt5}-\frac12\left(x+\frac1{\sqrt5}\right)
%   =
%   -\frac{\sqrt5}2-\frac12x.
% \]
% 
% 对某些函数求导数时, 若从显函数出发来求比较复杂或不好求, 这时可转化为隐函数来求导数, 常用的方法是对函数的两边取对数.
% \end{solve}
% 
% 
% \begin{example}\label{example:ma-4-2-11}
% 设 $y=x^{e^x}\ (x>0)$, 求 $y'$.
% \end{example}
% 
% 
% \begin{solve}
% 对 $y=x^{e^x}$ 两边取对数, 得 $\ln y=e^x\ln x$, 再两边对 $x$ 求导数, 得
% \[
%   \frac{y'}y=e^x\ln x+\frac{e^x}x
%   =
%   e^x\left(\frac{x\ln x+1}{x}\right),
% \]
% 所以
% \[
%   y'=x^{e^x}e^x\left(\frac{x\ln x+1}{x}\right).
% \qedhere
% \]
% \end{solve}
% 
% 
% \begin{example}\label{example:ma-4-2-12}
% 求
% \[
%   y=\left(1+\frac1{x^2}\right)^{x^2}
% \]
% 的导数.
% \end{example}
% 
% 
% \begin{solve}
% 在 $y=\left(1+\dfrac1{x^2}\right)^{x^2}$ 的两边取对数得
% \[
%   \ln y=x^2\ln\left(1+\frac1{x^2}\right)
%   =
%   x^2[\ln(x^2+1)-\ln x^2].
% \]
% 两边对 $x$ 求导数得
% \[
%   \frac{y'}y
%   =
%   2x(\ln(x^2+1)-\ln x^2)
%   +
%   x^2\left(\frac{2x}{1+x^2}-\frac{2x}{x^2}\right)
% \]
% \[
%   =
%   2x\left[\ln\left(1+\frac1{x^2}\right)-\frac1{1+x^2}\right],
% \]
% 所以
% \[
%   y'=2x\left(1+\frac1{x^2}\right)^{x^2}
%   \left[\ln\left(1+\frac1{x^2}\right)-\frac1{1+x^2}\right].
% \qedhere
% \]
% \end{solve}
% 
% 
% \begin{remark}\label{remark:ma-4-source-53}
% 以上例子中的求导方法也称为对数求导法, 它适合对一些比较复杂的函数求导.
% \end{remark}
% 
% 
% \subsection{参数式函数的求导法}
% \label{subsec:ma-4-2-5}
% 
% 下面我们来讨论由参数方程
% \[
%   \begin{cases}
%     x=x(t),\\
%     y=y(t),
%   \end{cases}
%   \quad t\in(\alpha,\beta)
% \]
% 所确定的函数 $y=f(x)$ 的求导问题.
% 
% 假定函数 $x(t)$ 和 $y(t)$ 都在 $(\alpha,\beta)$ 内可导, 且 $x'(t)\ne0$。在这样的条件下, 下一章我们将证明 $x(t)$ 在 $(\alpha,\beta)$ 严格单调。令
% \[
%   a=\min(x(\alpha+0),x(\beta-0)),
% \]
% \[
%   b=\max(x(\alpha+0),x(\beta-0)),
% \]
% 则 $x=x(t)$ 的反函数 $t=t(x)$ 在 $(a,b)$ 上有定义。因此, 由该参数方程可以确定一个函数
% \[
%   y=f(x)=y(t(x)),\quad x\in(a,b).
% \]
% 由复合函数及反函数的求导法则, 得
% \[
%   y_x'=y'(t)t'(x)=\frac{y'(t)}{x'(t)}.
% \]
% 这表明, 由参数方程所确定的函数的导数就等于 $y$ 与 $x$ 关于参数 $t$ 的导数的商。在以后的计算中, 我们只要记住此公式即可, 但理解公式的推导过程是十分有益的.
% 
% \begin{example}\label{example:ma-4-2-13}
% 求椭圆 $\dfrac{x^2}{a^2}+\dfrac{y^2}{b^2}=1\ (a,b>0)$ 上点 $(x_0,y_0)$ 处的切线方程.
% \end{example}
% 
% 
% \begin{solve}[{解法 1}]
% 利用隐函数求导法, 得
% \[
%   \frac{2x}{a^2}+\frac{2yy'}{b^2}=0,
% \]
% 即
% \[
%   y_x'=-\frac{b^2x}{a^2y}.
% \]
% 由此可知, 当 $y_0\ne0$ 时, 在 $(x_0,y_0)$ 处的切线方程为
% \[
%   y-y_0=-\frac{b^2x_0}{a^2y_0}(x-x_0).
% \]
% 注意到 $(x_0,y_0)$ 在椭圆上, 最后得切线方程为
% \[
%   \frac{x_0x}{a^2}+\frac{y_0y}{b^2}=1.
% \]
% 为了求曲线在点 $(a,0)$ 处切线, 在此点的附近我们可以将 $x$ 看成 $y$ 的函数, 再利用隐函数求导法, 得
% \[
%   \left.x_y'\right|_{y=0}
%   =
%   -\left.\frac{a^2y}{b^2x}\right|_{(a,0)}
%   =
%   0,
% \]
% 从而在点 $(a,0)$ 处的切线方程为 $x=a$。同理, 可得在点 $(-a,0)$ 处的切线方程为 $x=-a$.
% \end{solve}
% 
% 
% \begin{solve}[{解法 2}]
% 椭圆的参数方程为
% \[
%   \begin{cases}
%     x=a\cos t,\\
%     y=b\sin t,
%   \end{cases}
%   \quad t\in[0,2\pi).
% \]
% 由参数式函数的求导法, 得
% \[
%   y_x'=\frac{(b\sin t)'}{(a\cos t)'}
%   =
%   -\frac{b\cos t}{a\sin t}
%   \quad(t\ne0,\pi),
% \]
% 从而在 $t_0\ne0,\pi$ 处的切线方程为
% \[
%   y-b\sin t_0=
%   \frac{-b\cos t_0}{a\sin t_0}(x-a\cos t_0).
% \]
% 令 $x_0=a\cos t_0,y_0=b\sin t_0$, 整理即得
% \[
%   \frac{x_0x}{a^2}+\frac{y_0y}{b^2}=1.
% \]
% 完全类似于解法 1, 可求出在 $(\pm a,0)$ 处的切线方程.
% \end{solve}
% 
% 
% \subsection{极坐标式函数的求导法}
% \label{subsec:ma-4-2-6}
% 
% 设曲线由极坐标形式
% \[
%   r=r(\theta),\quad \theta\in(\alpha,\beta)
% \]
% 给出, 其中 $r$ 是极距, $\theta$ 表示极角。现对曲线上一点 $(\theta,r(\theta))$ 来求切线的斜率.
% 
% 由极坐标方程即可得曲线的参数方程为
% \[
%   \begin{cases}
%     x=r(\theta)\cos\theta,\\
%     y=r(\theta)\sin\theta.
%   \end{cases}
% \]
% 这样, 假定 $r$ 作为 $\theta$ 的函数其导数 $r'(\theta)$ 存在, 而且在所考虑的极角 $\theta$ 附近有
% \[
%   x'(\theta)=r'(\theta)\cos\theta-r(\theta)\sin\theta\ne0,
% \]
% 则由参数式函数的求导法知, 极坐标表示下的曲线在点 $(\theta,r(\theta))$ 处的切线的斜率为
% \[
%   \frac{dy}{dx}
%   =
%   \frac{r'(\theta)\sin\theta+r(\theta)\cos\theta}
%        {r'(\theta)\cos\theta-r(\theta)\sin\theta}
%   =
%   \frac{\tan\theta+\dfrac{r(\theta)}{r'(\theta)}}
%        {1-\tan\theta\dfrac{r(\theta)}{r'(\theta)}}.
% \]
% 上式给出了极坐标式函数的求导公式。由导数几何意义知, $\dfrac{dy}{dx}$ 是切线与 $x$ 轴正向的夹角 $\alpha$ 的正切, 即
% \[
%   \tan\alpha=
%   \frac{\tan\theta+\dfrac{r(\theta)}{r'(\theta)}}
%        {1-\tan\theta\dfrac{r(\theta)}{r'(\theta)}}.
% \]
% 由此可解出
% \[
%   \frac{r(\theta)}{r'(\theta)}
%   =
%   \frac{\tan\alpha-\tan\theta}{1+\tan\alpha\tan\theta}
%   =
%   \tan(\alpha-\theta).
% \]
% 它具有鲜明的几何意义。在极坐标系中同时建立直角坐标系, 记向径沿逆时针方向转到切线位置的角度为 $\beta$（通常称做向径与切线的夹角）, 则有
% \[
%   \tan\beta=\tan(\alpha-\theta)=\frac{r(\theta)}{r'(\theta)}.
% \]
% 这是因为这二个角之间有关系式: $\beta=k\pi+(\alpha-\theta)$, 这里 $k$ 是非负整数, 且满足 $0\leqslant k\pi+(\alpha-\theta)\leqslant\pi$。例如, $k=0$ 的情形如图 4.2.1 所示.
% 
% \begin{example}\label{example:ma-4-2-14}
% 求证: 双纽线 $r^2=a^2\cos2\theta\ (0<\theta<\pi/4)$ 的向径与切线的夹角等于极角的两倍加 $\pi/2$.
% \end{example}
% 
% 
% \begin{proof}
% 设向径与切线的夹角为 $\beta$, 则
% \[
%   \tan\beta
%   =
%   \frac{r(\theta)}{r'(\theta)}
%   =
%   \frac{r^2(\theta)}{r(\theta)r'(\theta)}
%   =
%   -\frac{\cos2\theta}{\sin2\theta}
%   =
%   -\cot2\theta
%   =
%   \tan\left(2\theta+\frac{\pi}{2}\right).
% \]
% 注意到 $0<\theta<\dfrac{\pi}{4}$, 便有 $\beta=2\theta+\dfrac{\pi}{2}$.
% 
% 最后指出, 显式表示、隐式表示、参数表示和极坐标表示是表达函数的四种重要形式, 各有不同的适用场合, 有很多函数只能用其中的某一种来表达。即使有些函数能够用它们中的任何一种来表达, 但运算的难易程度也可能是大相径庭的, 因此必须全面地掌握相应的求导数公式, 根据实际情况选择使用.
% \end{proof}
% 
% 
% \section{微分}
% \label{sec:ma-4-3}
% 
% \subsection{微分的定义}
% \label{subsec:ma-4-3-1}
% 
% 与函数在一点的可导性紧密联系着的一个概念是可微性, 本节就来讨论函数的可微性。我们先来看一个具体的例子。若用 $S$ 表示半径为 $r$ 的圆的面积, 则 $S$ 可以表示为 $r$ 的函数: $S=\pi r^2$。如果给半径 $r$ 一个增量 $\Delta r$, 则面积 $S$ 相应地有增量
% \[
%   \Delta S
%   =
%   \pi(r+\Delta r)^2-\pi r^2
%   =
%   2\pi r\Delta r+\pi(\Delta r)^2.
% \]
% 它是半径为 $r$ 与 $r+\Delta r$ 的两个同心圆之间的圆环的面积。从上式可以看出, $\Delta S$ 由两部分组成: 第一部分 $2\pi r\Delta r$ 是 $\Delta r$ 的线性函数, 而第二部分 $\pi(\Delta r)^2$ 是较 $\Delta r$ 高阶的无穷小量, 即 $\pi(\Delta r)^2=o(|\Delta r|)\ (\Delta r\to0)$。由此可知, 当给半径 $r$ 一个微小的增量 $\Delta r$ 时, 所引起的圆面积的增量 $\Delta S$ 就可以近似地用第一部分 $2\pi r\Delta r$ 来代替, 由此所产生的相对误差为
% \[
%   \left|\frac{\pi(\Delta r)^2}{2\pi r\Delta r}\right|
%   =
%   \frac{|\Delta r|}{2r}.
% \]
% 显然, $\Delta r$ 越小, 所产生的相对误差也越小.
% 
% 对于一般的情形, 我们引入如下的定义:
% 
% \begin{definition}\label{definition:ma-4-3-1}
% 设函数 $y=f(x)$ 在 $U(x_0,\delta_0)$ 内有定义。如果存在常数 $A$, 使得
% \[
%   \Delta y=f(x_0+\Delta x)-f(x_0)=A\Delta x+o(\Delta x)
%   \quad(\Delta x\to0),
% \]
% 则称 $f(x)$ 在点 $x_0$ 处可微, 并称 $A\Delta x$ 为 $f(x)$ 在 $x_0$ 处的微分, 记做
% \[
%   dy=A\Delta x\quad\text{或者}\quad df(x_0)=A\Delta x.
% \]
% \end{definition}
% 
% 从定义看到, 一个函数的微分具有两个特点:
% \begin{enumerate}
%   \item 它是自变量的增量 $\Delta x$ 的线性函数.
%   \item 它与函数的增量 $\Delta y$ 之差是较 $\Delta x$ 高阶的无穷小量 $(\Delta x\to0)$.
% \end{enumerate}
% 根据这一特点, 我们就可以用 $dy$ 来近似地表达 $\Delta y$, 所产生的相对误差是一个无穷小量 $(\Delta x\to0)$, 而且 $|\Delta x|$ 越小, 近似程度就越好。因此, 我们称 $dy=A\Delta x$ 为 $\Delta y$ 的\keyterm{线性主部}。此外, 由于 $dy$ 是 $\Delta x$ 的线性函数, 所以它特别容易计算.
% 
% 微分与导数有着密切的关系, 表现为如下定理:
% 
% \begin{theorem}\label{theorem:ma-4-3-1}
% 函数 $f(x)$ 在点 $x_0$ 处可微的充分必要条件是 $f(x)$ 在 $x_0$ 处可导.
% \end{theorem}
% 
% 
% \begin{proof}
% 必要性\quad 设 $f(x)$ 在 $x_0$ 处可微, 即
% \[
%   \Delta y=f(x_0+\Delta x)-f(x_0)=A\Delta x+o(\Delta x)
%   \quad(\Delta x\to0)
% \]
% 成立, 则有
% \[
%   \frac{\Delta y}{\Delta x}
%   =
%   A+o(1)
%   \quad(\Delta x\to0),
% \]
% 从而
% \[
%   f'(x_0)
%   =
%   \lim_{\Delta x\to0}\frac{\Delta y}{\Delta x}
%   =
%   A.
% \]
% 这表明, $f(x)$ 在 $x_0$ 处可导, 且 $f'(x_0)=A$.
% 
% 充分性\quad 设 $f(x)$ 在 $x_0$ 处可导, 即
% \[
%   \lim_{\Delta x\to0}
%   \frac{f(x_0+\Delta x)-f(x_0)}{\Delta x}
%   =
%   f'(x_0)
% \]
% 成立, 于是有
% \[
%   \frac{f(x_0+\Delta x)-f(x_0)}{\Delta x}
%   =
%   f'(x_0)+o(1)
%   \quad(\Delta x\to0),
% \]
% 即
% \[
%   \Delta y=f(x_0+\Delta x)-f(x_0)
%   =
%   f'(x_0)\Delta x+o(\Delta x)
%   \quad(\Delta x\to0).
% \]
% 这表明, $f(x)$ 在点 $x_0$ 处可微, 且 $df=f'(x_0)\Delta x$.
% \end{proof}
% 
% 
% \begin{remark}\label{remark:ma-4-source-62}
% 由于这一定理的缘故, 对于一元函数而言, 人们常把 ``可导'' 和 ``可微'' 等同起来使用, 求导方法又称为微分法。此外, 从定理的证明过程知, 微分定义式中的系数 $A$ 唯一地确定, 而且必有 $A=f'(x_0)$ 成立.
% \end{remark}
% 
% 现考虑一个具体函数 $y=x$ 的微分。由定义, 得 $dy=dx=1\Delta x$, 即 $dx=\Delta x$。因此, 我们就可在微分记号中将 $\Delta x$ 改为 $dx$, 即若 $y=f(x)$ 在 $x_0$ 处可微, 则 $dy=f'(x_0)dx$。因此, 导数的记号
% \[
%   \left.\frac{dy}{dx}\right|_{x=x_0}=f'(x_0)
% \]
% 就可以看成 $dy$ 与 $dx$ 的比值了, 即导数乃是微分之商, 故又称微商.
% 
% 下面我们来看微分的几何意义。设函数 $y=f(x)$ 的图像如图 4.3.1 所示, 令 $N=(x_0+\Delta x,f(x_0))$, $P=(x_0+\Delta x,f(x_0+\Delta x))$, $K=(x_0+\Delta x,f(x_0)+f'(x_0)\Delta x)$, 并且过 $M=(x_0,f(x_0))$ 作曲线的切线, 则有
% \[
% \begin{gathered}
%   dx=\Delta x=MN\text{ 为自变量的增量;}\\
%   \Delta y=PN\text{ 为函数的增量;}\\
%   dy=f'(x_0)dx=KN\text{ 为切线的增量;}\\
%   \Delta y-dy=o(\Delta x)=PK\text{ 为函数的增量与切线的增量之差.}
% \end{gathered}
% \]
% 这表明, 微分就是在点 $x_0$ 的附近以切线的增量拟合曲线的增量。因此, 所谓微分的基本原理就是: 一个可微函数, 在每一点的充分小局部, 均可近似视为增量按比例变化的线性函数.
% 
% 若函数 $y=f(x)$ 在区间 $(a,b)$ 内每一点 $x$ 处可微 (即 $dy=f'(x)dx$), 则称函数 $f(x)$ 在 $(a,b)$ 内可微.
% 
% 利用微分可作误差估计和近似计算, 下面举两个简单的例子.
% 
% \begin{example}\label{example:ma-4-3-1}
% 求 $\sqrt[4]{80}$ 的近似值.
% \end{example}
% 
% 
% \begin{solve}
% 取 $f(x)=\sqrt[4]{x}$, 令 $x_0=81,\Delta x=-1$, 则有
% \[
%   f'(x)=\frac{1}{4\sqrt[4]{x^3}},
%   \quad
%   f(x_0)=3,
%   \quad
%   f'(x_0)=\frac{1}{108}.
% \]
% 于是
% \[
%   \sqrt[4]{80}\approx 3+\frac{1}{108}\times(-1)\approx2.9907.
% \]
% 其实, 其精确值为 $\sqrt[4]{80}=2.990697\cdots$.
% \end{solve}
% 
% 
% \begin{example}\label{example:ma-4-3-2}
% 为了使计算的正方形的面积的相对误差不超过 $1\%$, 度量边长所允许的最大相对误差为多少?
% \end{example}
% 
% 
% \begin{solve}
% 设正方形的边长为 $a$, 面积为 $S$, 则有
% \[
%   S=a^2,\quad |\Delta S|\approx |dS|=|2a\Delta a|,
% \]
% \[
%   \left|\frac{\Delta S}{S}\right|
%   \approx
%   \left|\frac{2a\Delta a}{a^2}\right|
%   =
%   2\left|\frac{\Delta a}{a}\right|.
% \]
% 于是
% \[
%   \left|\frac{\Delta a}{a}\right|
%   \approx
%   \frac{1}{2}\left|\frac{\Delta S}{S}\right|
%   =
%   0.5\%.
% \]
% 这表明度量边长所允许的最大相对误差为 $0.5\%$.
% \end{solve}
% 
% 
% \subsection{一阶微分的形式不变性}
% \label{subsec:ma-4-3-2}
% 
% 设函数 $y=f(u),u=g(x)$。根据复合函数的求导法则, 可得复合函数 $y=f(g(x))$ 的微分公式为
% \[
%   dy=(f(g(x)))'dx=f'(g(x))g'(x)dx.
% \]
% 由于 $du=g'(x)dx$, 代入上式就得到了它的等价表示形式
% \[
%   dy=f'(u)du.
% \]
% 这里 $u=g(x)$ 是 $x$ 的函数。但我们发现, 它与 $u$ 为自变量的函数 $f(u)$ 的微分形式
% \[
%   dy=f'(u)du
% \]
% 一模一样。换句话说, 对 $f(u)$ 进行微分时, 不管 $u$ 是因变量还是自变量, 所得结果具有相同的形式, 这也就是所谓的\keyterm{一阶微分的形式不变性}。当然它们的意义是不一样的, $u$ 是自变量时, $du=\Delta u$, 而当 $u$ 是函数时, 一般来说 $du$ 与 $\Delta u$ 是不相同的.
% 
% 这是一阶微分的一个非常重要的性质, 有了这个 ``形式不变性'' 作保证, 我们拿到一个函数 $y=f(u)$ 之后, 就可以按 $u$ 是自变量去求它的微分 $f'(u)du$, 而无须顾忌 $u$ 是自变量, 还是因变量。这使得微分运算在很多场合比求导数运算优越得多.
% 
% 例如, 设 $y=f(x)$ 在 $x_0$ 可微, 且存在反函数 $x=f^{-1}(y)$, 再假定 $f'(x_0)\ne0$。我们先对 $y=f(x)$ 两边微分, 得 $dy=f'(x_0)dx$。然后将 $x$ 看成因变量, $y$ 看成自变量, 从前面的等式便可导出
% \[
%   \frac{dx}{dy}=\frac{1}{f'(x_0)},
% \]
% 这正是反函数的求导数公式.
% 
% 再如, 求参数方程
% \[
%   \begin{cases}
%     x=x(t),\\
%     y=y(t)
%   \end{cases}
% \]
% 所确定的函数 $y=f(x)$ 的导数。我们可以先对参数方程分别微分得
% \[
%   \begin{cases}
%     dx=x'(t)dt,\\
%     dy=y'(t)dt.
%   \end{cases}
% \]
% 然后利用这两个等式, 便可得到
% \[
%   \frac{dy}{dx}
%   =
%   \frac{y'(t)dt}{x'(t)dt}
%   =
%   \frac{y'(t)}{x'(t)},
% \]
% 这正是参数式函数的求导公式.
% 
% 利用前面得到的基本初等函数的导数公式, 立即可得它们的微分公式为:
% \[
% \begin{array}{ll}
%   (1)\ d(C)=0dx; & (2)\ d(x^\alpha)=\alpha x^{\alpha-1}dx;\\[4pt]
%   (3)\ d(a^x)=a^x\ln a\,dx\ (a>0); &
%   (4)\ d(\ln|x|)=\dfrac{dx}{x};\\[8pt]
%   (5)\ d(\sin x)=\cos x\,dx; &
%   (6)\ d(\cos x)=-\sin x\,dx;\\[4pt]
%   (7)\ d(\tan x)=\dfrac{dx}{\cos^2 x}; &
%   (8)\ d(\arcsin x)=\dfrac{dx}{\sqrt{1-x^2}};\\[8pt]
%   (9)\ d(\arctan x)=\dfrac{dx}{1+x^2}。&
% \end{array}
% \]
% 同样, 由导数的四则运算法则, 可得微分运算的法则。设 $u,v$ 为 $x$ 的可微函数, 则
% \[
% \begin{aligned}
%   (1)\quad &d(u\pm v)=du\pm dv;\\
%   (2)\quad &d(u\cdot v)=u\,dv+v\,du;\\
%   (3)\quad &d\left(\frac{u}{v}\right)=\frac{v\,du-u\,dv}{v^2}.
% \end{aligned}
% \]
% 利用这些公式和运算法则, 我们就可以求任何初等函数的微分了.
% 
% \begin{example}\label{example:ma-4-3-3}
% 设 $y=\arctan e^x$, 求 $dy$.
% \end{example}
% 
% 
% \begin{solve}
% \[
%   dy
%   =
%   \frac{1}{1+e^{2x}}de^x
%   =
%   \frac{e^x}{1+e^{2x}}dx
%   =
%   \frac{dx}{e^x+e^{-x}}.
% \qedhere
% \]
% \end{solve}
% 
% 
% \begin{example}\label{example:ma-4-3-4}
% 设 $y=\dfrac{x-a}{1-ax}\ (|a|<1)$, 求证: 当 $|x|<1$ 时, 有
% \[
%   \frac{dy}{1-y^2}=\frac{dx}{1-x^2}.
% \]
% \end{example}
% 
% 
% \begin{proof}
% 由于
% \[
%   dy=\frac{1-a^2}{(1-ax)^2}dx,
% \]
% \[
%   1-y^2
%   =
%   1-\left(\frac{x-a}{1-ax}\right)^2
%   =
%   \frac{(1-a^2)(1-x^2)}{(1-ax)^2},
% \]
% 故有
% \[
%   \frac{dy}{1-y^2}=\frac{dx}{1-x^2}.
% \]
% 此例的结果在双曲几何中起着重要的作用.
% \end{proof}
% 
% 
% \begin{example}\label{example:ma-4-3-5}
% 求由方程 $\sin y^2=\cos\sqrt{x}$ 所确定的隐函数 $y=y(x)$ 的导数.
% \end{example}
% 
% 
% \begin{solve}
% 对方程
% \[
%   \sin y^2=\cos\sqrt{x}
% \]
% 两边微分, 得
% \[
%   2y\cos y^2\,dy
%   =
%   -\frac{\sin\sqrt{x}}{2\sqrt{x}}\,dx,
% \]
% 由此即得
% \[
%   \frac{dy}{dx}
%   =
%   -\frac{\sin\sqrt{x}}{4y\sqrt{x}\cos y^2}.
% \qedhere
% \]
% \end{solve}
% 
% 
% \begin{remark}\label{remark:ma-4-source-73}
% 读者可以对此例中的隐函数的显式表达式
% \[
%   y=\pm\sqrt{\arcsin(\cos\sqrt{x})}
% \]
% 直接求导数, 并且将运算过程加以比较.
% \end{remark}
% 
% 
% \section{高阶导数与高阶微分}
% \label{sec:ma-4-4}
% 
% \subsection{高阶导数}
% \label{subsec:ma-4-4-1}
% 
% 当知道物体沿直线运动的路程 $s=s(t)$ 需要求加速度时, 我们就会遇到求速度函数的导数, 即求 $s(t)$ 的二阶导数的问题.
% 
% 若一个函数 $y=f(x)$ 的一阶导数 $f'(x)$ 仍是可导函数, 则我们可求 $(f'(x))'$, 记其为 $f''(x)$ 或 $\dfrac{d^2y}{dx^2}$, 并称之为 $f$ 的\keyterm{二阶导数}。类似地, 可定义\keyterm{三阶导数}为 $f'''(x)=(f''(x))'$。一般地, 当 $n\geqslant4$ 时, $f(x)$ 的 $n$ 阶导数定义为 $f(x)$ 的 $n-1$ 阶导数的导数, 并记为 $f^{(n)}(x),y^{(n)}$, 或 $\dfrac{d^ny}{dx^n}$.
% 
% 对于一些较简单的函数, 我们可以导出它们的高阶导数的公式。基本方法是, 先求出前几阶导数, 发现规律, 总结出一般的公式, 然后再加以证明.
% 
% \begin{example}\label{example:ma-4-4-1}
% 设 $y=e^{ax}$, 求 $y^{(n)}$.
% \end{example}
% 
% 
% \begin{solve}
% 由于
% \[
%   y'=ae^{ax},\quad y''=a^2e^{ax},\quad y'''=a^3e^{ax},
% \]
% 因此, 由高阶导数的定义, 容易用数学归纳法证明
% \[
%   y^{(n)}=a^ne^{ax}\quad(n=1,2,\cdots).
% \qedhere
% \]
% \end{solve}
% 
% 
% \begin{example}\label{example:ma-4-4-2}
% 设 $y=x^\alpha$, 求 $y^{(n)}$.
% \end{example}
% 
% 
% \begin{solve}
% 由于
% \[
%   y'=\alpha x^{\alpha-1},\quad y''=\alpha(\alpha-1)x^{\alpha-2},
% \]
% 因此容易由数学归纳法证明
% \[
%   y^{(n)}
%   =
%   \alpha(\alpha-1)\cdots(\alpha-n+1)x^{\alpha-n}
%   \quad(n\geqslant1).
% \qedhere
% \]
% \end{solve}
% 
% 
% \begin{remark}\label{remark:ma-4-source-78}
% 当 $\alpha$ 为正整数时, 若 $\alpha=n$, 则 $y^{(n)}=n!$; 若 $\alpha<n$, 则 $y^{(n)}=0$。由此立即导出: 如果 $P_m(x)$ 是 $m$ 次多项式, 且 $n>m$, 则
% \[
%   P_m^{(n)}(x)=0.
% \]
% \end{remark}
% 
% 
% \begin{example}\label{example:ma-4-4-3}
% 设 $y=\ln(1+x)$, 求 $y^{(n)}$.
% \end{example}
% 
% 
% \begin{solve}
% 由
% \[
%   y'=\frac{1}{1+x},\quad
%   y''=\frac{-1}{(1+x)^2},\quad
%   y'''=\frac{2!}{(1+x)^3},\quad
%   y^{(4)}=\frac{-3!}{(1+x)^4},
% \]
% 可以看出
% \[
%   y^{(n)}
%   =
%   (-1)^{n-1}\frac{(n-1)!}{(1+x)^n}
%   \quad(n\geqslant1),
% \]
% 其中规定 $0!=1$.
% 
% 事实上, 设
% \[
%   y^{(k)}
%   =
%   (-1)^{k-1}\frac{(k-1)!}{(1+x)^k},
% \]
% 则
% \[
%   y^{(k+1)}
%   =
%   (-1)^{k-1}(k-1)!\frac{-k(1+x)^{k-1}}{(1+x)^{2k}}
%   =
%   \frac{(-1)^kk!}{(1+x)^{k+1}}.
% \]
% 这样, 由数学归纳法原理知, 对一切的正整数 $n$ 所得公式都成立.
% \end{solve}
% 
% 
% \begin{example}\label{example:ma-4-4-4}
% 设 $y=\sin x$, 求 $y^{(n)}$.
% \end{example}
% 
% 
% \begin{solve}
% 直接计算, 得
% \[
%   y'=\cos x=\sin\left(x+\frac{\pi}{2}\right),
%   \quad
%   y''=-\sin x=\sin(x+\pi).
% \]
% 现在假定
% \[
%   y^{(n-1)}
%   =
%   \sin\left(x+\frac{(n-1)\pi}{2}\right),
% \]
% 则有
% \[
%   y^{(n)}
%   =
%   \left[\sin\left(x+\frac{(n-1)\pi}{2}\right)\right]'
%   =
%   \cos\left(x+\frac{(n-1)\pi}{2}\right)
%   =
%   \sin\left(x+\frac{n\pi}{2}\right).
% \]
% 同理可得
% \[
%   (\cos x)^{(n)}
%   =
%   \cos\left(x+\frac{n\pi}{2}\right).
% \qedhere
% \]
% \end{solve}
% 
% 
% \begin{remark}\label{remark:ma-4-source-83}
% 上例中的两个公式也可以借助欧拉公式形式地导出。在欧拉公式
% \[
%   e^{ix}=\cos x+i\sin x
% \]
% 的两边形式地求导, 可得
% \[
%   (e^{ix})^{(n)}
%   =
%   (\cos x)^{(n)}+i(\sin x)^{(n)},
% \]
% 而
% \[
%   (e^{ix})^{(n)}
%   =
%   i^ne^{ix}
%   =
%   e^{i\frac{n\pi}{2}}e^{ix}
%   =
%   e^{i\left(x+\frac{n\pi}{2}\right)}
%   =
%   \cos\left(x+\frac{n\pi}{2}\right)
%   +
%   i\sin\left(x+\frac{n\pi}{2}\right),
% \]
% 所以
% \[
%   (\cos x)^{(n)}
%   =
%   \cos\left(x+\frac{n\pi}{2}\right),\quad
%   (\sin x)^{(n)}
%   =
%   \sin\left(x+\frac{n\pi}{2}\right).
% \]
% \end{remark}
% 
% 
% \begin{example}\label{example:ma-4-4-5}
% 设 $y=\arctan x$, 求 $y^{(n)}$.
% \end{example}
% 
% 
% \begin{solve}
% 由于 $x=\tan y$, 直接计算, 可得
% \[
%   y'=\frac{1}{1+x^2}
%   =
%   \frac{1}{1+\tan^2y}
%   =
%   \cos^2y,
% \]
% \[
%   y''
%   =
%   -2\cos y\sin y\cdot y'
%   =
%   \cos^2y\sin2\left(y+\frac{\pi}{2}\right),
% \]
% \[
% \begin{aligned}
%   y'''
%   &=
%   2\cos^3y(-\sin y)\sin2\left(y+\frac{\pi}{2}\right)
%   +
%   2\cos^4y\cos2\left(y+\frac{\pi}{2}\right)  \\
%   &=
%   2\cos^3y\cos\left(2\left(y+\frac{\pi}{2}\right)+y\right)  \\
%   &=
%   2\cos^3y\sin3\left(y+\frac{\pi}{2}\right).
% \end{aligned}
% \]
% 一般地, 由数学归纳法可证
% \[
%   y^{(n)}
%   =
%   (n-1)!\cos^ny\sin n\left(y+\frac{\pi}{2}\right).
% \]
% 此外, 由这一公式容易导出
% \[
%   \left.y^{(2n-1)}\right|_{x=0}
%   =
%   (-1)^{n-1}(2n-2)!,
%   \quad
%   \left.y^{(2n)}\right|_{x=0}=0.
% \qedhere
% \]
% \end{solve}
% 
% 
% \subsection{莱布尼茨公式}
% \label{subsec:ma-4-4-2}
% 
% 对于两个函数 $u,v$ 的和、差的高阶导数, 我们有
% \[
%   (u\pm v)^{(n)}=u^{(n)}\pm v^{(n)};
% \]
% 对常数 $c$ 和函数 $u$ 的积, 有
% \[
%   (cu)^{(n)}=cu^{(n)};
% \]
% 而对于两个函数 $u,v$ 的积, 我们有以下莱布尼茨公式:
% 
% \begin{theorem}[{莱布尼茨公式}]\label{theorem:ma-4-4-1}
% 若函数 $u$ 和 $v$ 有任意阶导数, 则
% \[
%   (u\cdot v)^{(n)}
%   =
%   \sum_{k=0}^{n}C_n^ku^{(n-k)}v^{(k)},
% \]
% 其中
% \[
%   C_n^k=\frac{n!}{k!(n-k)!},\quad u^{(0)}=u,\quad v^{(0)}=v.
% \]
% \end{theorem}
% 
% 
% \begin{proof}
% 用数学归纳法证之。当 $n=1$ 时, 公式显然成立。假设公式对正整数 $n$ 成立, 我们来证公式对正整数 $n+1$ 也成立。事实上, 应用归纳法假定, 我们有
% \[
% \begin{aligned}
%   (u\cdot v)^{(n+1)}
%   &=
%   \left[(u\cdot v)^{(n)}\right]'  \\
%   &=
%   \sum_{k=0}^{n}C_n^k\left[u^{(n-k)}v^{(k)}\right]'  \\
%   &=
%   \sum_{k=0}^{n}C_n^k
%   \left[u^{(n-k+1)}v^{(k)}+u^{(n-k)}v^{(k+1)}\right]  \\
%   &=
%   \sum_{k=0}^{n}C_n^ku^{(n-k+1)}v^{(k)}
%   +
%   \sum_{k=1}^{n+1}C_n^{k-1}u^{(n-k+1)}v^{(k)}  \\
%   &=
%   C_n^0u^{(n+1)}v^{(0)}
%   +
%   \sum_{k=1}^{n}(C_n^{k-1}+C_n^k)u^{(n-k+1)}v^{(k)}
%   +
%   C_n^nu^{(0)}v^{(n+1)}  \\
%   &=
%   \sum_{k=0}^{n+1}C_{n+1}^ku^{(n+1-k)}v^{(k)},
% \end{aligned}
% \]
% 这与牛顿二项式展开公式证明的格式是一致的, 其中最后一步用到了恒等式
% \[
%   C_n^k+C_n^{k-1}=C_{n+1}^k.
% \qedhere
% \]
% \end{proof}
% 
% 
% \begin{example}\label{example:ma-4-4-6}
% 设 $y=\arcsin x$, 求 $y^{(n)}(0)$.
% \end{example}
% 
% 
% \begin{solve}
% 由 $y'=\dfrac{1}{\sqrt{1-x^2}}$, 可得
% \[
%   \sqrt{1-x^2}\cdot y'=1,
% \]
% 再两边平方得
% \begin{equation}
%   (1-x^2)\cdot y'^2=1。\label{eq:ma-4-4-1}
% \end{equation}
% 若对上式用莱布尼茨公式, 则 $y'^2$ 的高阶导数仍不好算。所以对上式再求导数一次, 得
% \[
%   (1-x^2)\cdot 2y'y''-2xy'^2=0.
% \]
% 由 \eqref{eq:ma-4-4-1} 式看出, 当 $|x|<1$ 时, $y'\ne0$, 化简得
% \[
%   (1-x^2)y''-xy'=0.
% \]
% 对上式两边求 $(n-2)$ 阶导数, 并应用莱布尼茨公式, 得
% \[
% \begin{aligned}
%   &(1-x^2)y^{(n)}
%   +(n-2)(-2x)y^{(n-1)}
%   +\frac{(n-2)(n-3)}{2}\cdot(-2)y^{(n-2)}  \\
%   &\qquad
%   -xy^{(n-1)}-(n-2)y^{(n-2)}=0.
% \end{aligned}
% \]
% 将 $x=0$ 代入上式, 就有
% \[
%   y^{(n)}(0)
%   -(n-2)(n-3)y^{(n-2)}(0)
%   -(n-2)y^{(n-2)}(0)=0,
% \]
% 由此即得递推公式
% \[
%   y^{(n)}(0)=(n-2)^2y^{(n-2)}(0).
% \]
% 又因 $y^{(0)}(0)=0$, $y^{(1)}(0)=1$, 所以我们有 $y^{(2n)}(0)=0$, 且
% \[
% \begin{aligned}
%   y^{(2n+1)}(0)
%   &=
%   (2n-1)^2y^{(2n-1)}(0)  \\
%   &=
%   (2n-1)^2(2n-3)^2y^{(2n-3)}(0)  \\
%   &=\cdots  \\
%   &=
%   (2n-1)^2(2n-3)^2\cdots3^2\cdot1^2\cdot y^{(1)}(0)  \\
%   &=
%   [(2n-1)!!]^2.
% \end{aligned}
% \qedhere
% \]
% \end{solve}
% 
% 
% \begin{example}\label{example:ma-4-4-7}
% 设 $y=\tan x$, 求 $y^{(n)}(0)$, $n=1,2,\cdots,8$.
% \end{example}
% 
% 
% \begin{solve}
% 因奇函数的导数为偶函数, 偶函数的导数为奇函数（参考习题四第 4 题）, 以及奇函数在原点的值为零, 所以有
% \[
%   y(0)=y^{(2)}(0)=y^{(4)}(0)=y^{(6)}(0)=y^{(8)}(0)=0.
% \]
% 为了求奇数阶导数在原点的值, 我们设法建立递推公式。等式
% \[
%   y\cdot\cos x=\sin x
% \]
% 两边求 $(2n+1)$ 阶导数, 并利用莱布尼茨公式, 得
% \[
%   \sum_{k=0}^{2n+1}
%   C_{2n+1}^ky^{(2n+1-k)}
%   \cos\left(x+\frac{k\pi}{2}\right)
%   =
%   \sin\left(x+\frac{2n+1}{2}\pi\right),
% \]
% 将 $x=0$ 代入上式, 即得递推公式
% \[
%   \sum_{k=0}^{2n+1}
%   C_{2n+1}^ky^{(2n+1-k)}(0)\cos\frac{k\pi}{2}
%   =
%   (-1)^n,
% \]
% 即
% \[
% \begin{aligned}
%   &y^{(2n+1)}(0)
%   -C_{2n+1}^{2}y^{(2n-1)}(0)
%   +C_{2n+1}^{4}y^{(2n-3)}(0)+\cdots \\
%   &\quad
%   +(-1)^{n-1}C_{2n+1}^{2n-2}y^{(3)}(0)
%   +(-1)^nC_{2n+1}^{2n}y'(0)
%   =
%   (-1)^n.
% \end{aligned}
% \]
% 将 $n=1$ 及 $y'(0)=1$ 代入上式, 得
% \[
%   y^{(3)}(0)-\frac{3\cdot2}{2!}=-1,
% \]
% 解出 $y^{(3)}(0)=2$.
% 
% 在递推公式中令 $n=2$, 得
% \[
%   y^{(5)}(0)-\frac{5\cdot4}{2!}\cdot2
%   +
%   \frac{5\cdot4\cdot3\cdot2}{4!}\cdot1
%   =
%   1.
% \]
% 解出 $y^{(5)}(0)=16$.
% 
% 再在递推公式中令 $n=3$, 得
% \[
%   y^{(7)}(0)-\frac{7\cdot6}{2!}\cdot16
%   +
%   \frac{7\cdot6\cdot5\cdot4}{4!}\cdot2
%   -
%   \frac{7!}{6!}\cdot1
%   =
%   -1.
% \]
% 解出 $y^{(7)}(0)=272$.
% 
% 附带有
% \[
%   \frac{y'(0)}{1!}=1,\quad
%   \frac{y^{(3)}(0)}{3!}=\frac13,\quad
%   \frac{y^{(5)}(0)}{5!}=\frac{2}{15},\quad
%   \frac{y^{(7)}(0)}{7!}=\frac{17}{315}.
% \qedhere
% \]
% \end{solve}
% 
% 
% \subsection{一般函数的高阶导数}
% \label{subsec:ma-4-4-3}
% 
% 对于一般没有什么特性的函数, 通常很难归纳出其 $n$ 阶导数的公式, 只能逐次地求导。特别对于复合函数、反函数、隐函数和参数式函数求高阶导数更是如此.
% 
% \begin{example}\label{example:ma-4-4-8}
% 设函数 $y=y(x)$ 是由方程 $\dfrac{x^2}{a^2}+\dfrac{y^2}{b^2}=1\ (y>0)$ 给出, 试求 $y''$.
% \end{example}
% 
% 
% \begin{solve}[{解法 1}]
% 先从方程解出
% \[
%   y=\frac{b}{a}\sqrt{a^2-x^2},
% \]
% 然后直接计算, 可得
% \[
%   y'
%   =
%   -\frac{b}{a}(a^2-x^2)^{-\frac12}x
%   =
%   -\frac{bx}{a\sqrt{a^2-x^2}},
% \]
% \[
%   y''
%   =
%   -\frac{b}{a}(a^2-x^2)^{-\frac32}x^2
%   -
%   \frac{b}{a}(a^2-x^2)^{-\frac12}
%   =
%   -\frac{b^4}{a^2y^3}.
% \qedhere
% \]
% \end{solve}
% 
% 
% \begin{solve}[{解法 2}]
% 先将方程表示为参数式
% \[
%   \begin{cases}
%     x=a\cos t,\\
%     y=b\sin t,
%   \end{cases}
% \]
% 再直接计算, 有
% \[
%   y'_x=-\frac{b\cos t}{a\sin t}=-\frac{b}{a}\cot t,
% \]
% \[
% \begin{aligned}
%   y''_{xx}
%   &=
%   -\frac{d(y'_x)}{dx}
%   =
%   -\frac{b}{a}\frac{(\cot t)'}{x'_t}  \\
%   &=
%   \frac{b}{a\sin^2t}\cdot\frac{-1}{a\sin t}
%   =
%   \frac{-b}{a^2\sin^3t}
%   =
%   -\frac{b^4}{a^2y^3}.
% \end{aligned}
% \qedhere
% \]
% \end{solve}
% 
% 
% \begin{solve}[{解法 3}]
% 直接对方程两边关于 $x$ 求导数, 得
% \begin{equation}
%   \frac{2x}{a^2}+\frac{2yy'}{b^2}=0。\label{eq:ma-4-4-2}
% \end{equation}
% 由此解得
% \[
%   y'=-\frac{b^2x}{a^2y}.
% \]
% 再对方程 \eqref{eq:ma-4-4-2} 两边关于 $x$ 求导数, 得
% \[
%   \frac1{a^2}+\frac{y'^2+yy''}{b^2}=0.
% \]
% 由此解出 $y''$ 并且将 $y'$ 的表达式代入, 得
% \[
%   y''=-\frac{b^4}{a^2y^3}.
% \]
% 比较上面的三种解法, 不难发现, 似乎隐函数求导法要简单一些.
% \end{solve}
% 
% 
% \subsection{高阶微分}
% \label{subsec:ma-4-4-4}
% 
% 若函数 $f(x)$ 在区间 $(a,b)$ 内可微, 则有 $\mathrm df=f'(x)\mathrm dx$, 其中 $f'(x)$ 是 $x$ 的函数, 而 $\mathrm dx$ 是与 $x$ 无关的一个量。这样, 把一阶微分 $\mathrm df$ 看成 $x$ 的函数, 再求一次微分 $\mathrm d(\mathrm df)$, 就称之为 $f(x)$ 的二阶微分, 记为 $\mathrm d^2f$, 即
% \[
%   \mathrm d^2f
%   =
%   \mathrm d(\mathrm df)
%   =
%   (f'(x)\mathrm dx)'\mathrm dx
%   =
%   f''(x)\mathrm dx\mathrm dx
%   =
%   f''(x)\mathrm dx^2,
% \]
% 其中 $\mathrm dx^2=\mathrm dx\mathrm dx$（读者应注意 $\mathrm dx^2\ne \mathrm d(x^2)=2x\mathrm dx$）.
% 
% 一般地, 我们定义 $n$ 阶微分为
% \[
% \begin{aligned}
%   \mathrm d^nf
%   &=
%   \mathrm d(\mathrm d^{n-1}f)
%   =
%   (f^{(n-1)}(x)\mathrm dx^{n-1})'\mathrm dx  \\
%   &=
%   f^{(n)}(x)\mathrm dx^{n-1}\mathrm dx
%   =
%   f^{(n)}(x)\mathrm dx^n.
% \end{aligned}
% \]
% 由此即可知道, 为什么将 $f^{(n)}(x)$ 记为 $\dfrac{\mathrm d^nf}{\mathrm dx^n}$ 的缘故了.
% 
% 由定义可知, 求高阶微分本质上就是求高阶导数, 因此我们不再详细讨论, 但值得指出的是高阶微分是不具有形式不变性的.
% 
% 事实上, 设 $y=f(u)$, 当 $u$ 为自变量时, 我们有 $\mathrm d^2y=f''(u)\mathrm du^2$; 但当 $u=g(x)$ 是 $x$ 的函数时, 则
% \[
%   \mathrm df=f'(u)g'(x)\mathrm dx,
% \]
% \[
%   \mathrm d^2f
%   =
%   f''(u)(g'(x))^2\mathrm dx^2+f'(u)g''(x)\mathrm dx^2
%   =
%   f''(u)\mathrm du^2+f'(u)\mathrm d^2u,
% \]
% 这比 $u$ 是自变量时多了一项 $f'(u)\mathrm d^2u$! 因此, 当 $\mathrm d^2u\ne0$, 即 $g''(x)\ne0$ 时, 二阶微分就不具有形式不变性了.
% 
% \section*{习题四}
% \phantomsection
% \addcontentsline{toc}{section}{习题四}
% \label{exercises:ma-4}
% 
% \noindent 1。讨论函数 $f(x)$ 和 $g(x)$ 在 $x=0$ 处的可导性, 其中
% \[
%   f(x)=
%   \begin{cases}
%     -x, & x\in\mathbb R\setminus\mathbb Q,\\
%     x, & x\in\mathbb Q;
%   \end{cases}
% \]
% \[
%   g(x)=
%   \begin{cases}
%     -x^2, & x\in\mathbb R\setminus\mathbb Q,\\
%     x^2, & x\in\mathbb Q.
%   \end{cases}
% \]
% 
% \noindent 2。设函数 $f(x)$ 在 $(a,b)$ 上有定义, 且在点 $x_0\in(a,b)$ 处可导, 并假定序列 $\{x_n\}$ 和 $\{y_n\}$ 满足
% \[
%   a<x_n<x_0<y_n<b,\quad n=1,2,\cdots,
% \]
% 且有
% \[
%   \lim_{n\to\infty}x_n=x_0=\lim_{n\to\infty}y_n.
% \]
% 证明:
% \[
%   \lim_{n\to\infty}\frac{f(y_n)-f(x_n)}{y_n-x_n}=f'(x_0).
% \]
% 
% \noindent 3。设函数 $f(x)$ 在 $(-1,1)$ 上可导, 且满足 $|f(x)|\leqslant|\sin x|$, 求证 $|f'(0)|\leqslant1$.
% 
% \noindent 4。求证:
% \begin{enumerate}
%   \item 可导的偶函数, 其导数为奇函数;
%   \item 可导的奇函数, 其导数为偶函数;
%   \item 可导的周期函数, 其导数为相同周期的周期函数.
% \end{enumerate}
% 
% \noindent 5。讨论下列函数在指定点 $x_0$ 处的可导性:
% \begin{enumerate}
%   \item $y=|(x-1)(x-2)^2(x-3)^3|,\quad x_0=1,2,3$;
%   \item
%   $
%   y=
%   \begin{cases}
%     \dfrac14(x-1)(x+1)^2, & |x|\leqslant1,\\
%     |x|-1, & |x|>1,
%   \end{cases}
%   \quad x_0=-1,0,1.
%   $
% \end{enumerate}
% 
% \noindent 6。设函数 $f(x)$ 在 $[-1,1]$ 上有定义, $f'(0)$ 存在, 求
% \[
%   \lim_{n\to\infty}
%   \left[
%     f\left(\frac1{n^2}\right)
%     +
%     f\left(\frac2{n^2}\right)
%     +\cdots+
%     f\left(\frac n{n^2}\right)
%     -
%     nf(0)
%   \right].
% \]
% 
% \noindent 7。求下列序列的极限:
% \begin{enumerate}
%   \item
%   $\displaystyle
%   \lim_{n\to\infty}
%   \left[
%     \sin\frac1{n^2}+\sin\frac2{n^2}+\cdots+\sin\frac n{n^2}
%   \right].
%   $
%   \item
%   $\displaystyle
%   \lim_{n\to\infty}
%   \left(1+\frac1{n^2}\right)
%   \left(1+\frac2{n^2}\right)\cdots
%   \left(1+\frac n{n^2}\right).
%   $
% \end{enumerate}
% 
% \noindent 8。设 $P_n(x)$ 为最高次系数为 1 的 $n(n\geqslant1)$ 次多项式, $M$ 为 $P_n(x)=0$ 的最大实根, 求证 $P'_n(M)\geqslant0$.
% 
% \noindent 9。设函数 $\varphi$ 在点 $a$ 连续, $f(x)=|x-a|\varphi(x)$。求 $f$ 在点 $a$ 的左右导数, 问在什么条件下 $f(x)$ 在点 $a$ 可导?
% 
% \noindent 10。给定曲线 $y=x^2+5x+4$.
% \begin{enumerate}
%   \item 确定 $b$, 使直线 $y=3x+b$ 为该曲线的切线;
%   \item 确定 $m$, 使直线 $y=mx$ 为该曲线的切线.
% \end{enumerate}
% 
% \noindent 11。给定曲线 $y=x^3$ 和直线 $y=px-q$, 其中 $p,q$ 为实数, 且 $p>0$.
% \begin{enumerate}
%   \item $p$ 给定后, 试确定 $q$, 使 $y=px-q$ 是 $y=x^3$ 的切线;
%   \item 利用图形求出方程 $x^3-px+q=0$ 有三个不同实根的条件.
% \end{enumerate}
% 
% \noindent 12.
% \begin{enumerate}
%   \item 确定 $m$, 使 $y=mx$ 为曲线 $y=\ln x$ 的切线;
%   \item 利用图形求出方程 $\ln x-mx=0$ 有两个实根的条件.
% \end{enumerate}
% 
% \noindent 13。问抛物线 $y=x^2-2x-1$ 在哪一点的切线垂直于直线 $x+2y-1=0$?
% 
% \noindent 14。求下列函数的导数:
% \begin{enumerate}
%   \item $y=\sqrt{x}+\dfrac1x-2x^3$;
%   \item $y=x^22^x$;
%   \item $y=\dfrac{2x-x^2}{1+x+x^2}$;
%   \item $y=\dfrac{x\sin x+\cos x}{x\sin x-\cos x}$;
%   \item $y=\dfrac1{\sqrt[3]{x}}+\sqrt[3]{x}$;
%   \item $y=\dfrac1{1+\sqrt{x}}-\dfrac1{1-\sqrt{x}}$;
%   \item $y=x^3\ln x-\dfrac1n x^n$;
%   \item $y=\left(x+\dfrac1x\right)\ln x$;
%   \item $y=\sec x$;
%   \item $y=\dfrac{\cos x}{x^4}\ln\dfrac1x$.
% \end{enumerate}
% 
% \noindent 15。确定常数 $a,b$, 使函数
% \[
%   f(x)=
%   \begin{cases}
%     ax+b, & x>1,\\
%     x^2, & x\leqslant1
%   \end{cases}
% \]
% 有连续的导函数.
% 
% \noindent 16。求下列函数的导数:
% \begin{enumerate}
%   \item $y=x(a^2+x^2)\sqrt{a^2-x^2}$;
%   \item $y=\sqrt[3]{\dfrac{1+x^3}{1-x^3}}$;
%   \item $y=\ln(\ln x)$;
%   \item $y=\dfrac1{2a}\ln\left|\dfrac{a+x}{a-x}\right|$;
%   \item $y=\ln(x+\sqrt{a+x^2})$;
%   \item $y=\ln\tan\dfrac x2$;
%   \item $y=\cos^3x-\cos3x$;
%   \item $y=\sin^nx\cos nx$;
%   \item $y=\cos(\cos\sqrt{x})$;
%   \item $y=\dfrac{\sin^2x}{\sin x^2}$;
%   \item $y=(e^x+e^{-x})^2$;
%   \item $y=\ln\sqrt{\dfrac{1+\cos x}{1-\cos x}}$.
% \end{enumerate}
% 
% \noindent 17。设在 $x=1$ 处有
% \[
%   \frac{d}{dx}f(x^2)=\frac{d}{dx}f^2(x),
% \]
% 求证 $f'(1)=0$ 或 $f(1)=1$.
% 
% \noindent 18。给定可导函数 $f(x)>0$, 证明函数 $y=f(x)\sin ax\ (a>0)$ 永远夹在两条振幅曲线 $y=\pm f(x)$ 之间, 并与之相切.
% 
% \noindent 19。求下列函数的导数:
% \begin{enumerate}
%   \item $y=x^x\ (x>0)$;
%   \item $y=x^{\tan x}\ (x>0)$;
%   \item $y=x^{\ln x}\ (x>0)$;
%   \item $y=(1+x)^{\frac1x}\ (x>0)$;
%   \item $y=x^{x^x}\ (x>0)$;
%   \item $y=e^{-\frac{\sin x}{x^2}}\ (x\ne0)$.
% \end{enumerate}
% 
% \noindent 20。求下列函数的导数:
% \begin{enumerate}
%   \item $y=\arcsin\sqrt{1-x^2}$;
%   \item $y=\arccos(\sin x)$;
%   \item $y=x\arctan x-\dfrac12\ln(1+x^2)$;
%   \item $y=\arctan\dfrac{\sqrt{1-x^2}-1}{x}+\arctan\dfrac{2x}{1-x^2}$;
%   \item $y=\arctan(\tan^2x)$;
%   \item $y=\left(\dfrac ab\right)^x\left(\dfrac bx\right)^a\left(\dfrac xa\right)^b\ (a,b>0)$;
%   \item $y=\ln\left(\arccos\dfrac1{\sqrt{x}}\right)$;
%   \item $y=\ln(e^x+\sqrt{1+e^{2x}})$;
%   \item $y=\dfrac x2\sqrt{a^2-x^2}+\dfrac{a^2}{2}\arctan\dfrac xa\ (a>0)$;
%   \item $y=\dfrac1{4\sqrt2}\ln\dfrac{x^2+\sqrt2x+1}{x^2-\sqrt2x+1}-\dfrac1{2\sqrt2}\arctan\dfrac{\sqrt2x}{x^2-1}$.
% \end{enumerate}
% 
% \noindent 21。求证:
% \[
%   \lim_{x\to a}\frac{f(x)-b}{x-a}=A
% \]
% 的充分必要条件是
% \[
%   \lim_{x\to a}\frac{e^{f(x)}-e^b}{x-a}=e^bA.
% \]
% 
% \noindent 22。设曲线的参数表示: $x=a\cos^3t,\ y=a\sin^3t\ (a>0)$.
% \begin{enumerate}
%   \item 求 $y'(x)$;
%   \item 证明曲线的切线被坐标轴所截的长度为一常数.
% \end{enumerate}
% 
% \noindent 23。求椭圆 $\dfrac{x^2}{a^2}+\dfrac{y^2}{b^2}=1$ 在第一象限中的切线, 使它平行于过 $(0,b)$ 和 $(a,0)$ 的直线.
% 
% \noindent 24。证明曲线
% \[
%   \begin{cases}
%     x=a(\cos t+t\sin t),\\
%     y=a(\sin t-t\cos t)
%   \end{cases}
%   \quad(a>0)
% \]
% 
% \begin{remark}
% 原稿在本题处截断，未提供完整题面。
% \end{remark}
% 
% % 整合来源：math-5.tex
% \chapter{导数的应用}
% \label{chap:ma-5}
% 
% 上一章中, 我们着重介绍了导数的概念及计算导数的各种方法。这一章, 我们将主要利用导数来研究函数的性质.
% 
% \section{微分中值定理}
% \label{sec:ma-5-1}
% 
% 微分中值定理是反映函数和导数之间联系的重要定理, 是利用导数研究函数的桥梁, 因此它是一个强有力的工具。灵活地运用它可使许许多多问题的解决变得轻而易举。在建立微分中值定理时需要用到如下的费马 (Fermat) 定理, 因此我们先来介绍这一定理.
% 
% \subsection{费马定理}
% \label{subsec:ma-5-1-1}
% 
% 极值问题是数学中的重要研究对象, 在本章中我们将利用导数来讨论此类问题。为此我们引入
% 
% \begin{definition}\label{definition:ma-5-1-1}
% 若函数 $f(x)$ 在 $x_0$ 的某个去心邻域 $U_0(x_0,\delta)$ 内恒有
% \[
%   f(x)\leqslant f(x_0)\quad (f(x)\geqslant f(x_0)),
% \]
% 则称 $x_0$ 为 $f(x)$ 的极大 (小) 值点, $f(x_0)$ 称为它的极大 (小) 值; 若在上述定义中, 不等号严格成立, 则称 $f(x_0)$ 为严格极大 (小) 值。极大值点和极小值点统称为极值点, 极大值和极小值统称为极值.
% \end{definition}
% 
% 从图像上来看, 一个函数 $y=f(x)$ 若在 $x_0$ 处取极值, 并且该函数的图像在点 $(x_0,f(x_0))$ 处存在切线, 则该切线一定是水平的。事实上, 我们有
% 
% \begin{theorem}[{费马定理}]\label{theorem:ma-5-1-1}
% 设函数 $f(x)$ 在 $U(x_0,\delta)$ 中有定义。若 $x_0$ 是 $f(x)$ 的极值点, 且 $f'(x_0)$ 存在, 则 $f'(x_0)=0$.
% \end{theorem}
% 
% 
% \begin{proof}
% 无妨设 $f(x_0)$ 为极大值, 则当 $\Delta x>0$ 且 $x_0+\Delta x\in U(x_0,\delta)$ 时, 有
% \[
%   \frac{f(x_0+\Delta x)-f(x_0)}{\Delta x}\leqslant0.
% \]
% 令 $\Delta x\to0+0$, 得 $f'(x_0)\leqslant0$.
% 
% 当 $\Delta x<0$ 且 $x_0+\Delta x\in U(x_0,\delta)$ 时, 有
% \[
%   \frac{f(x_0+\Delta x)-f(x_0)}{\Delta x}\geqslant0.
% \]
% 令 $\Delta x\to0-0$, 得 $f'(x_0)\geqslant0$.
% 
% 综上所证, 推得 $f'(x_0)=0$.
% \end{proof}
% 
% 
% \begin{remark}\label{remark:ma-5-source-4}
% 使得 $f'(x)=0$ 的点 $x_0$ 也称为 $f(x)$ 的驻点。费马定理告诉我们, 若极值点可导, 则它必是驻点。例子 $y=x^3$ 在 $x=0$ 处的性质告诉我们, 驻点未必是极值点.
% \end{remark}
% 
% 
% \subsection{罗尔微分中值定理}
% \label{subsec:ma-5-1-2}
% 
% \begin{theorem}[{罗尔 (Rolle) 微分中值定理}]\label{theorem:ma-5-1-2}
% 设函数 $f(x)$ 在 $[a,b]$ 上连续, 在 $(a,b)$ 内可导, 且 $f(a)=f(b)$, 则在 $(a,b)$ 内至少存在一点 $\xi$, 使得
% \[
%   f'(\xi)=0.
% \]
% \end{theorem}
% 
% 
% \begin{proof}
% 因为 $f(x)$ 在 $[a,b]$ 连续, 所以 $f(x)$ 在 $[a,b]$ 上有最大值 $M$ 与最小值 $m$。如果 $M=m$, 则 $f(x)$ 在 $[a,b]$ 为一常值函数, 从而 $f'(x)=0,\ \forall x\in[a,b]$。此时, 我们可取 $(a,b)$ 中任意一点作为 $\xi$。如果 $M>m$, 则 $M$ 与 $m$ 至少有一个不等于 $f(a)=f(b)$, 不妨设 $M>f(a)$。这时, 必存在 $\xi\in(a,b)$, 使得 $f(\xi)=M$, 即 $\xi$ 为 $f(x)$ 的一个极值点。再由 $f'(\xi)$ 存在, 并且应用费马定理, 得 $f'(\xi)=0$.
% 
% 罗尔微分中值定理的几何意义是, 在一段每点都有切线的曲线上, 若两端点的高度相同, 则在此曲线上至少存在一条水平切线 (参见图 5.1.1).
% 
% 此外, 罗尔微分中值定理只是告诉我们在区间 $(a,b)$ 内存在一点 $\xi$ 使得 $f'(\xi)=0$, 但并没有告诉我们怎样去求这一 $\xi$。事实上, 一般来讲, 要将这一 $\xi$ 求出来是十分困难。其次, 在很多时候我们会发现, 满足这样要求的 $\xi$ 可能有很多, 甚至有可能是无穷多个.
% \end{proof}
% 
% 
% \begin{example}\label{example:ma-5-1-1}
% 设函数 $f(x)$ 在 $(a,b)$ 上可微, 证明 $f(x)$ 在 $(a,b)$ 内的两个零点之间必有 $f(x)+f'(x)$ 的零点.
% \end{example}
% 
% 
% \begin{proof}
% 令 $F(x)=e^x f(x)$, 由罗尔微分中值定理知, 在 $f(x)$ 的两个零点之间必有
% \[
%   F'(x)=e^x(f(x)+f'(x))
% \]
% 的零点, 也即有 $f(x)+f'(x)$ 的零点.
% \end{proof}
% 
% 
% \begin{example}\label{example:ma-5-1-2}
% 设 $P_n(x)$ 为一个 $n$ 次多项式, 证明 $e^x-P_n(x)=0$ 至多有 $n+1$ 个不同的根.
% \end{example}
% 
% 
% \begin{proof}
% 令 $f(x)=e^x-P_n(x)$。假若 $f(x)$ 有 $n+2$ 个不同的零点, 则由罗尔微分中值定理知, $f'(x)=e^x-(P_n(x))'$ 至少有 $n+1$ 个不同的零点; 然后再应用罗尔微分中值定理, 又有 $f''(x)=e^x-(P_n(x))''$ 至少有 $n$ 个不同的零点。如此下去, 即知 $f^{(n+1)}=e^x-(P_n(x))^{(n+1)}$ 至少有一个零点。由于 $P_n(x)$ 是一个 $n$ 次多项式, 所以这就蕴涵着 $e^x$ 必有一个零点, 这显然是不可能的。这一矛盾说明 $e^x-P_n(x)=0$ 至多有 $n+1$ 个不同的根.
% \end{proof}
% 
% 
% \begin{example}\label{example:ma-5-1-3}
% 设函数 $f(x)$ 在 $[0,1]$ 上连续, 在 $(0,1)$ 内可导, 且 $f(0)=f(1)=0$。再假设 $f(t_0)=\alpha$, 其中 $t_0\in(0,1)$。证明: 存在 $\xi\in(0,1)$, 使得 $f'(\xi)=\alpha$.
% \end{example}
% 
% 
% \begin{proof}
% 若 $\alpha=0$, 利用罗尔微分中值定理即得.
% 
% 现设 $\alpha\ne0$。构造辅助函数
% \[
%   F(x)=f(x)-\alpha x,
% \]
% 则 $F(x)$ 在 $[0,1]$ 上连续, 在 $(0,1)$ 内可导, 并且有
% \[
%   F(0)=0,\quad F(1)=-\alpha,\quad F(t_0)=(1-t_0)\alpha.
% \]
% 由于 $F(1)F(t_0)<0$, 因此在 $(t_0,1)$ 中必存在 $\eta$, 使得 $F(\eta)=0$。因此我们在 $[0,\eta]$ 上对 $F(x)$ 应用罗尔微分中值定理, 可知在 $(0,\eta)\subset(0,1)$ 中存在 $\xi$, 使得 $F'(\xi)=0$, 即 $f'(\xi)=\alpha$.
% \end{proof}
% 
% 
% \subsection{拉格朗日微分中值定理}
% \label{subsec:ma-5-1-3}
% 
% 罗尔微分中值定理要求所考虑的函数在所讨论区间的两端点有相同的函数值, 这给它的应用范围带来很大的限制。因此, 一个自然的问题就是: 如果将罗尔微分中值定理中的条件 ``$f(a)=f(b)$'' 去掉, 将有什么结论呢? 其实, 如果将坐标轴去掉, 纯粹从几何上来看, 罗尔微分中值定理是说, 在一段每点都有切线的曲线上至少有一点的切线平行于两端点的连线。而这正是一个可微函数的图像所具有的几何特征。因此, 如果将罗尔微分中值定理中的条件 ``$f(a)=f(b)$'' 去掉, 则有如下的定理:
% 
% \begin{theorem}[{拉格朗日 (Lagrange) 微分中值定理}]\label{theorem:ma-5-1-3}
% 若函数 $f(x)$ 在 $[a,b]$ 上连续, 在 $(a,b)$ 内可导, 则在 $(a,b)$ 内至少存在一点 $\xi$, 使得
% \begin{equation}
%   f'(\xi)=\frac{f(b)-f(a)}{b-a}.
%   \label{eq:ma-5-1-1}
% \end{equation}
% \end{theorem}
% 
% 
% \begin{proof}
% 考查 $f(x)$ 减去过 $(a,f(a))$ 和 $(b,f(b))$ 这两点的直线所确定的线性函数而得到的函数:
% \[
%   F(x)=f(x)-\left[f(a)+\frac{f(b)-f(a)}{b-a}(x-a)\right].
% \]
% 容易验证, 这样定义的函数 $F(x)$ 满足罗尔微分中值定理的所有条件。因此, 存在 $\xi\in(a,b)$, 使得
% \[
%   F'(\xi)=f'(\xi)-\frac{f(b)-f(a)}{b-a}=0,
% \]
% 即
% \[
%   f'(\xi)=\frac{f(b)-f(a)}{b-a}.
% \qedhere
% \]
% \end{proof}
% 
% 
% \begin{remark}\label{remark:ma-5-source-15}
% (1) 在上述定理中, 若我们只假定 $f(x)$ 在以 $a$ 和 $b$ 为端点的闭区间满足所有的条件, 不难看出, 对 $a>b$ 的情形, 拉格朗日微分中值定理仍具有同样的形式.
% 
% (2) 公式 \eqref{eq:ma-5-1-1} 有时也称为拉格朗日中值公式。在实际应用时, 常将拉格朗日中值公式写成如下形式
% \[
%   f(a+h)-f(a)=f'(\xi)h,
% \]
% 这里 $h=b-a$ 可以是正数, 也可以是负数, 而 $\xi$ 则位于 $a$ 与 $a+h$ 之间。有时, 为了避免 ``$\xi$ 位于 $a$ 与 $a+h$ 之间'' 可能出现的两种情形, 常将上式改写成
% \[
%   f(a+h)=f(a)+f'(a+\theta h)h,
% \]
% 其中 $\theta=\dfrac{\xi-a}{h}$。无论 $h$ 是正数还是负数, 总有 $0<\theta<1$, 这是因为 $\xi$ 位于 $a$ 与 $a+h$ 之间的缘故.
% \end{remark}
% 
% 此外, 拉格朗日微分中值定理的几何意义如图 5.1.2 所示, 即在一段每点都有切线的曲线上至少有一点的切线平行于两个端点的连线.
% 
% 由拉格朗日微分中值定理立即可推出如下两个十分有用的结论:
% 
% \begin{corollary}\label{corollary:ma-local-5.1-1}
% 设函数 $f(x)\in C[a,b]$, 且在 $(a,b)$ 内可导。若 $f'(x)=0$, 则 $f(x)\equiv C$, 其中 $C$ 是一常数.
% \end{corollary}
% 
% 
% \begin{proof}
% 对 $\forall x\in[a,b]$, 在 $[a,x]$ 上应用拉格朗日微分中值定理知, 存在 $\xi\in(a,x)$, 使得
% \[
%   0=f'(\xi)=\frac{f(x)-f(a)}{x-a},
% \]
% 由此即得
% \[
%   f(x)\equiv f(a)=C.
% \]
% 推论~\ref{corollary:ma-local-5.1-1} 的结论似乎是显然的, 但如果不利用微分中值定理, 似乎很难证明它.
% \end{proof}
% 
% 
% \begin{corollary}\label{corollary:ma-local-5.1-2}
% 设函数 $f(x),g(x)\in C[a,b]$, 且在 $(a,b)$ 内可导。若 $f'(x)=g'(x)$, 则 $f(x)=g(x)+C$, 其中 $C$ 是一常数.
% \end{corollary}
% 
% 
% \begin{proof}
% 对 $f(x)-g(x)$ 应用推论~\ref{corollary:ma-local-5.1-1} 即得.
% \end{proof}
% 
% 
% \begin{example}\label{example:ma-5-1-4}
% 证明当 $x>-1$ 且 $x\ne0$ 时, 有
% \[
%   \frac{x}{1+x}<\ln(1+x)<x.
% \]
% \end{example}
% 
% 
% \begin{proof}
% 对 $\ln(1+x)$, 在 $0$ 和 $x$ 之间应用拉格朗日微分中值定理, 得
% \[
%   \frac{\ln(1+x)}{x}
%   =
%   \frac{\ln(1+x)-\ln1}{x-0}
%   =
%   \left.(\ln(1+x))'\right|_{x=\xi}
%   =
%   \frac1{1+\xi},
% \]
% 其中 $\xi$ 在 $0$ 和 $x$ 之间。再注意到, 当 $x>0$ 时, 有
% \[
%   \frac1{1+x}<\frac1{1+\xi}<1;
% \]
% 而当 $-1<x<0$ 时, 有
% \[
%   1<\frac1{1+\xi}<\frac1{1+x},
% \]
% 便知所证不等式成立.
% \end{proof}
% 
% 
% \begin{example}\label{example:ma-5-1-5}
% 证明 $f(x)=x^\alpha\ (0<\alpha<1)$ 在 $[0,+\infty)$ 上一致连续.
% \end{example}
% 
% 
% \begin{proof}
% 由于 $f(x)$ 在 $[0,+\infty)$ 上连续, 因此它在 $[0,2]$ 上一致连续。于是, $\forall\varepsilon>0,\exists\delta_1>0$, 使得当 $x_1,x_2\in[0,2]$ 且 $|x_1-x_2|<\delta_1$ 时, 有
% \[
%   |f(x_1)-f(x_2)|<\varepsilon.
% \]
% 另一方面, 因为 $f'(x)=\alpha x^{\alpha-1}$ 在 $[1,+\infty)$ 上严格单调递减, 所以在 $[1,+\infty)$ 上恒有 $|f'(x)|\leqslant\alpha<1$。于是, 对 $\forall x_1,x_2\in[1,+\infty)$, 应用拉格朗日微分中值定理, 得
% \[
%   |f(x_1)-f(x_2)|=|f'(\xi)(x_1-x_2)|\leqslant |x_1-x_2|.
% \]
% 这样, 对给定的 $\varepsilon$, 取 $\delta_2=\varepsilon$, 则当 $x_1,x_2\in[1,+\infty)$ 且 $|x_1-x_2|<\delta_2$ 时, 有
% \[
%   |f(x_1)-f(x_2)|<\varepsilon.
% \]
% 现取 $\delta=\min\{\delta_1,\delta_2,1\}$, 则 $\forall x_1,x_2\in[0,+\infty)$, 当 $|x_1-x_2|<\delta$ 时, 一定有 $x_1,x_2\in[0,2]$ 或 $x_1,x_2\in[1,+\infty)$, 从而必有 $|f(x_1)-f(x_2)|<\varepsilon$。这表明 $f(x)$ 在 $[0,+\infty)$ 上一致连续.
% \end{proof}
% 
% 
% \begin{remark}\label{remark:ma-5-source-24}
% 从上例的证明中可以看出: 若一个函数在某区间上的导数有界, 则它必在该区间上一致连续。此外, 上例也说明了, 一致连续的函数即使导数存在, 其导数也未必在所论区间上有界.
% \end{remark}
% 
% 
% \paragraph*{思考题}
% 当 $\alpha>1$ 时, 函数 $f(x)=x^\alpha$ 是否在 $[0,+\infty)$ 上一致连续?
% 
% \begin{example}\label{example:ma-5-1-6}
% 证明: 对任意的 $x\in\mathbb{R}$, 有
% \[
%   \arctan x=\arcsin\left(\frac{x}{\sqrt{1+x^2}}\right).
% \]
% \end{example}
% 
% 
% \begin{proof}
% 令
% \[
%   f(x)=\arctan x-\arcsin\left(\frac{x}{\sqrt{1+x^2}}\right),
% \]
% 则有
% \[
%   f'(x)
%   =
%   \frac1{1+x^2}
%   -
%   \frac1{\sqrt{1-\frac{x^2}{1+x^2}}}
%   \cdot
%   \frac{\sqrt{1+x^2}-\frac{x^2}{\sqrt{1+x^2}}}{1+x^2}
%   =
%   0,\quad \forall x\in\mathbb{R}.
% \]
% 于是由推论~\ref{corollary:ma-local-5.1-1} 知 $f(x)\equiv C$, 而 $f(0)=0$, 故所证等式成立.
% \end{proof}
% 
% 
% \subsection{柯西微分中值定理}
% \label{subsec:ma-5-1-4}
% 
% 为了将拉格朗日微分中值定理作进一步的推广, 我们现在换个观点来看这个定理。在该定理的条件下, 我们有 $f'(\xi_1)=\dfrac{f(b)-f(a)}{b-a}$, 其分子是函数 $f(x)$ 在区间端点之值的差, 而其分母可以看成函数 $y=x$ 在端点之值的差。从这个观点来看, 一个自然的问题是: 能否将 $y=x$ 换成一个任意的函数呢? 下面的柯西微分中值定理回答了这个问题.
% 
% \begin{theorem}[{柯西微分中值定理}]\label{theorem:ma-5-1-4}
% 若函数 $f(x)$ 和 $g(x)$ 在 $[a,b]$ 上连续, 在 $(a,b)$ 内可导, 而且 $g'(x)\ne0$, 则在 $(a,b)$ 内至少存在一点 $\xi$, 使得
% \begin{equation}
%   \frac{f(b)-f(a)}{g(b)-g(a)}=\frac{f'(\xi)}{g'(\xi)}.
%   \label{eq:ma-5-1-2}
% \end{equation}
% \end{theorem}
% 
% 
% \begin{remark}\label{remark:ma-5-source-28}
% 分别应用拉格朗日微分中值定理于 $f(x)$ 和 $g(x)$, 我们可得
% \[
%   \frac{f(b)-f(a)}{g(b)-g(a)}=\frac{f'(\xi_1)}{g'(\xi_2)}.
% \]
% 注意, 这里的 $\xi_1$ 与 $\xi_2$ 可能是不一样的! 因此, 我们不能简单地应用拉格朗日微分中值定理来证明本定理.
% \end{remark}
% 
% 
% \begin{proof}
% 对函数 $g(x)$ 应用拉格朗日微分中值定理 (或罗尔微分中值定理), 可知 $g'(x)\ne0$ 蕴涵着 $g(a)\ne g(b)$。仿照拉格朗日中值定理的证明方法, 令
% \[
%   G(x)=f(x)-\left[f(a)+\frac{f(b)-f(a)}{g(b)-g(a)}(g(x)-g(a))\right],
% \]
% 则 $G(x)$ 满足罗尔定理的所有条件。于是, 由罗尔定理知, 存在 $\xi\in(a,b)$, 使得 $G'(\xi)=0$。整理即得公式 \eqref{eq:ma-5-1-2}.
% 
% 柯西微分中值定理的几何意义是, 若 $uv$ 坐标平面的曲线由参数方程
% \[
%   \begin{cases}
%     u=g(x),\\
%     v=f(x),
%   \end{cases}
%   \quad x\in[a,b]
% \]
% 给出, 其中 $f(x)$ 和 $g(x)$ 满足定理的条件, 则存在 $\xi\in(a,b)$, 使得过点 $(g(\xi),f(\xi))$ 的切线平行于两端点的连线.
% 
% 此外, 类似于拉格朗日微分中值定理, 柯西微分中值定理也常写成如下形式:
% \[
%   \frac{f(b)-f(a)}{g(b)-g(a)}
%   =
%   \frac{f'(a+\theta(b-a))}{g'(a+\theta(b-a))},
% \]
% 其中 $0<\theta<1$, 且不要求 $b>a$.
% \end{proof}
% 
% 
% \begin{example}\label{example:ma-5-1-7}
% 设函数 $f(x)$ 在区间 $(0,1]$ 上连续, 在区间 $(0,1)$ 内可导, 而且
% \[
%   \lim_{x\to0+0}\sqrt{x}f'(x)=\ell.
% \]
% 证明 $f(x)$ 在区间 $(0,1]$ 上一致连续.
% \end{example}
% 
% 
% \begin{proof}
% 由 $\lim\limits_{x\to0+0}\sqrt{x}f'(x)=\ell$ 可知, 存在 $M>0$ 和 $0<c<1$, 使得
% \[
%   |\sqrt{x}f'(x)|\leqslant M,\quad \forall x\in(0,c].
% \]
% 现在取 $x_1,x_2\in(0,c]$, $x_1\ne x_2$, 应用柯西微分中值定理, 可得
% \[
%   \frac{f(x_1)-f(x_2)}{\sqrt{x_1}-\sqrt{x_2}}
%   =
%   \left.\frac{f'(x)}{(\sqrt{x})'}\right|_{x=\xi}
%   =
%   2\sqrt{\xi}f'(\xi),
% \]
% 其中 $\xi$ 位于 $x_1$ 和 $x_2$ 之间。于是, 有
% \[
%   |f(x_1)-f(x_2)|\leqslant 2M|\sqrt{x_1}-\sqrt{x_2}|.
% \]
% 当然, 上式对 $x_1=x_2$ 亦成立。再注意到
% \[
%   |\sqrt{x_1}-\sqrt{x_2}|^2
%   \leqslant |\sqrt{x_1}-\sqrt{x_2}|\,|\sqrt{x_1}+\sqrt{x_2}|
%   =
%   |x_1-x_2|,
% \]
% 对 $\forall\varepsilon>0$, 只需取 $\delta=\dfrac{\varepsilon^2}{(2M)^2}$, 则当 $x_1,x_2\in(0,c]$ 且 $|x_1-x_2|<\delta$ 时, 便有
% \[
%   |f(x_1)-f(x_2)|
%   \leqslant 2M|\sqrt{x_1}-\sqrt{x_2}|
%   \leqslant 2M|x_1-x_2|^{\frac12}
%   <\varepsilon.
% \]
% 这表明 $f(x)$ 在 $(0,c]$ 上一致连续。又 $f(x)$ 在区间 $(0,1]$ 上连续蕴涵着 $f(x)$ 在区间 $[c,1]$ 上一致连续, 从而便有 $f(x)$ 在区间 $(0,1]$ 上一致连续.
% 
% 从本节内容可以看出, 罗尔微分中值定理是一个基本的结果, 拉格朗日微分中值定理是罗尔微分中值定理的直接推广, 而柯西微分中值定理则是拉格朗日微分中值定理的一个重要推广。它们各自适合不同的场合, 是我们利用导数研究函数的强有力工具.
% \end{proof}
% 
% 
% \section{洛必达法则}
% \label{sec:ma-5-2}
% 
% 在建立求导法则和求导公式的过程中, 极限的理论和一些具体的极限起着决定性的作用。而有了导数理论和求导公式之后, 又可利用它解决极限理论中某些不定式的极限问题.
% 
% 大家已经知道, 如果我们已知
% \[
%   \lim_{x\to a}f(x)=A,\quad \lim_{x\to a}g(x)=B,
% \]
% 这里 $A,B,a$ 可以是有限数或无穷大, 那么除去某些例外情形, 我们可以利用下面的极限运算法则去求它们的和、差、积、商及幂--指数函数式的极限:
% \begin{enumerate}
%   \item $\lim\limits_{x\to a} f(x)\pm g(x)=A\pm B$;
%   \item $\lim\limits_{x\to a} f(x)\cdot g(x)=AB$;
%   \item $\lim\limits_{x\to a}\dfrac{f(x)}{g(x)}=\dfrac AB\ (B\ne0)$;
%   \item $\lim\limits_{x\to a} f(x)^{g(x)}=A^B$.
% \end{enumerate}
% 对于以上各式, 所说的例外情况分别为:
% \begin{enumerate}
%   \item[(a)] 在 (1) 中出现两个同号无穷大相减;
%   \item[(b)] 在 (2) 中 $A$ 与 $B$ 之一为 $0$, 另一为无穷大;
%   \item[(c)] 在 (3) 中 $A$ 与 $B$ 同时为 $0$, 或者同为无穷大;
%   \item[(d)] 在 (4) 中 $A=1,B=\infty$, 或者 $A=B=0$, 或者 $A=\infty,B=0$.
% \end{enumerate}
% 这些例外情形所涉及的极限类型统称为\keyterm{不定式}, 分别记为:
% \[
%   \infty-\infty,\quad 0\cdot\infty,\quad \frac00,\quad \frac\infty\infty,\quad 1^\infty,\quad 0^0,\quad \infty^0.
% \]
% 所有这些不定式本质上只有一种: $\dfrac00$。其他的不定式经过简单的变换可以化为这种不定式, 如 $0\cdot\infty=0\cdot\dfrac10$ 就成为 $\dfrac00$ 型, 或者 $0\cdot\infty=\dfrac1\infty\cdot\infty$ 变成 $\dfrac\infty\infty$ 型, 而 $\dfrac\infty\infty=\dfrac10/\dfrac10=\dfrac00$, 等等。下面我们利用柯西微分中值定理来给出这种不定式的定值方法 --- 洛必达 (L'Hospital) 法则。以后我们还会发现, 虽然 $\dfrac\infty\infty$ 型可以化为 $\dfrac00$ 型, 但要利用洛必达法则, 我们还必须对 $\dfrac\infty\infty$ 的情况进行研究.
% 
% \subsection{\texorpdfstring{$\dfrac00$ 型不定式}{0/0 型不定式}}
% \label{subsec:ma-5-2-1}
% 
% \begin{theorem}\label{theorem:ma-5-2-1}
% 设函数 $f(x)$ 和 $g(x)$ 在 $a$ 点的某一去心邻域 $U_0(a,\delta)$ 上可导, 而且满足:
% \begin{enumerate}
%   \item $\lim\limits_{x\to a}f(x)=\lim\limits_{x\to a}g(x)=0$;
%   \item $g'(x)\ne0,\ \forall x\in U_0(a,\delta)$;
%   \item $\lim\limits_{x\to a}\dfrac{f'(x)}{g'(x)}=\ell\ (\ell\text{ 为有限数或 }\pm\infty,\infty)$,
% \end{enumerate}
% 则有
% \[
%   \lim_{x\to a}\frac{f(x)}{g(x)}
%   =
%   \lim_{x\to a}\frac{f'(x)}{g'(x)}
%   =
%   \ell.
% \]
% \end{theorem}
% 
% 
% \begin{proof}
% 先考虑 $\ell$ 为有限数的情形, 并且我们先证 $\lim\limits_{x\to a-0}\dfrac{f(x)}{g(x)}=\ell$.
% 由条件 (1) 知, 若补充定义 $f(a)=g(a)=0$, 则 $f(x),g(x)$ 在 $a$ 点连续。于是, 对 $\forall x\in(a-\delta,a)$, 在区间 $[x,a]$ 上应用柯西微分中值定理, 有
% \[
%   \frac{f(x)}{g(x)}
%   =
%   \frac{f(x)-f(a)}{g(x)-g(a)}
%   =
%   \frac{f'(\xi)}{g'(\xi)},\quad x<\xi<a.
% \]
% 又由条件 (3), 可得
% \[
%   \lim_{x\to a-0}\frac{f'(\xi)}{g'(\xi)}
%   =
%   \lim_{x\to a-0}\frac{f'(x)}{g'(x)}
%   =
%   \ell.
% \]
% 所以
% \[
%   \lim_{x\to a-0}\frac{f(x)}{g(x)}=\ell.
% \]
% 同理可证 $\lim\limits_{x\to a+0}\dfrac{f(x)}{g(x)}=\ell$。综合起来就有 $\lim\limits_{x\to a}\dfrac{f(x)}{g(x)}=\ell$.
% 当 $\ell=+\infty,-\infty$ 或 $\infty$ 时, 证明类似, 请读者作为练习自己给出.
% \end{proof}
% 
% 
% \begin{remark}\label{remark:ma-5-source-34}
% 从定理的证明可以看出, 把 $x\to a$ 改为 $x\to a-0$ 或 $x\to a+0$, 上述定理的结论仍然成立.
% \end{remark}
% 
% 当 $x\to\infty$ 时, 我们有
% 
% \begin{theorem}\label{theorem:ma-5-2-2}
% 设函数 $f(x),g(x)$ 在 $U=\{x: |x|>a>0\}$ 上可导, 而且满足:
% \begin{enumerate}
%   \item $\lim\limits_{x\to\infty}f(x)=\lim\limits_{x\to\infty}g(x)=0$;
%   \item $g'(x)\ne0,\ \forall x\in U$;
%   \item $\lim\limits_{x\to\infty}\dfrac{f'(x)}{g'(x)}=\ell\ (\ell\text{ 为有限数或 }\pm\infty,\infty)$,
% \end{enumerate}
% 则有
% \[
%   \lim_{x\to\infty}\frac{f(x)}{g(x)}=\ell.
% \]
% \end{theorem}
% 
% 
% \begin{proof}
% 我们先证 $x\to+\infty$ 的情形。作自变量变换 $x=\dfrac1t$, 则 $x\to+\infty$ 相应于 $t\to0+0$。于是有
% \[
%   \lim_{t\to0+0}
%   \frac{f'\left(\frac1t\right)}
%        {g'\left(\frac1t\right)}
%   =
%   \lim_{x\to+\infty}\frac{f'(x)}{g'(x)},
% \]
% 并且由条件 (1), 有
% \[
%   \lim_{t\to0+0}f\left(\frac1t\right)=0,\quad
%   \lim_{t\to0+0}g\left(\frac1t\right)=0.
% \]
% 应用定理~\ref{theorem:ma-5-2-1} 于开区间 $\left(0,\dfrac1a\right)$ 上新变量 $t$ 的函数 $f\left(\dfrac1t\right)$ 和 $g\left(\dfrac1t\right)$, 并注意到它们关于 $t$ 的导数为
% \[
%   f'\left(\frac1t\right)\left(-\frac1{t^2}\right),\quad
%   g'\left(\frac1t\right)\left(-\frac1{t^2}\right),
% \]
% 可得
% \[
%   \lim_{t\to0+0}
%   \frac{f\left(\frac1t\right)}
%        {g\left(\frac1t\right)}
%   =
%   \lim_{t\to0+0}
%   \frac{f'\left(\frac1t\right)\left(-\frac1{t^2}\right)}
%        {g'\left(\frac1t\right)\left(-\frac1{t^2}\right)}
%   =
%   \lim_{x\to+\infty}\frac{f'(x)}{g'(x)}
%   =
%   \ell.
% \]
% 因此, 我们有
% \[
%   \lim_{x\to+\infty}\frac{f(x)}{g(x)}=\ell.
% \]
% 同理可证 $x\to-\infty$ 的情形。这样我们就证明了 $\lim\limits_{x\to\infty}\dfrac{f(x)}{g(x)}=\ell$ 成立.
% \end{proof}
% 
% 
% \begin{remark}\label{remark:ma-5-source-37}
% 从定理的证明可以看出, 把 $x\to\infty$ 改为 $x\to-\infty$ 或 $x\to+\infty$, 上述定理的结论仍然成立.
% \end{remark}
% 
% 
% \begin{example}\label{example:ma-5-2-1}
% 求极限
% \[
%   \lim_{x\to0}\frac{\sinh x-x}{\sin x-x}.
% \]
% \end{example}
% 
% 
% \begin{solve}
% 这是一个 $\dfrac00$ 型不定式。应用洛必达法则, 可得
% \[
%   \lim_{x\to0}\frac{\sinh x-x}{\sin x-x}
%   =
%   \lim_{x\to0}\frac{\cosh x-1}{\cos x-1}
%   =
%   \lim_{x\to0}\frac{\sinh x}{-\sin x}
%   =
%   \lim_{x\to0}\frac{\cosh x}{-\cos x}
%   =
%   -1.
% \qedhere
% \]
% \end{solve}
% 
% 
% \begin{example}\label{example:ma-5-2-2}
% 求极限
% \[
%   \lim_{x\to0+0}\frac{\sqrt{x}}{1-e^{2\sqrt{x}}}.
% \]
% \end{example}
% 
% 
% \begin{solve}
% 这也是一个 $\dfrac00$ 型不定式。令 $t=\sqrt{x}$, 则有
% \[
%   \lim_{x\to0+0}\frac{\sqrt{x}}{1-e^{2\sqrt{x}}}
%   =
%   \lim_{t\to0+0}\frac{t}{1-e^{2t}}
%   =
%   \lim_{t\to0+0}\frac1{-2e^{2t}}
%   =
%   -\frac12.
% \qedhere
% \]
% \end{solve}
% 
% 
% \begin{example}\label{example:ma-5-2-3}
% 求极限
% \[
%   \lim_{x\to+\infty}x\left(\frac\pi2-\arctan x\right).
% \]
% \end{example}
% 
% 
% \begin{solve}
% 这是一个 $0\cdot\infty$ 型不定式。先将其转化为 $\dfrac00$ 型不定式, 然后应用洛必达法则, 可得
% \[
%   \lim_{x\to+\infty}x\left(\frac\pi2-\arctan x\right)
%   =
%   \lim_{x\to+\infty}
%   \frac{\frac\pi2-\arctan x}{\frac1x}
%   =
%   \lim_{x\to+\infty}
%   \frac{-\frac1{1+x^2}}{-\frac1{x^2}}
%   =
%   1.
% \qedhere
% \]
% \end{solve}
% 
% 
% \subsection{\texorpdfstring{$\dfrac\infty\infty$ 型不定式}{∞/∞ 型不定式}}
% \label{subsec:ma-5-2-2}
% 
% 从理论上来讲, $\dfrac\infty\infty$ 型不定式可化为 $\dfrac00$ 型。但是有时候这样做会使问题更为复杂, 致使问题的解决失败。例如, 求 $\lim\limits_{x\to+\infty}\dfrac{\ln x}{x^\alpha}\ (\alpha>0)$。如果将它化成
% \[
%   \frac{\ln x}{x^\alpha}
%   =
%   \frac{\frac1{x^\alpha}}{\frac1{\ln x}},
% \]
% 则分子、分母分别求导后所得表达式比原来的函数还要复杂。因此, 我们有必要专门研究求 $\dfrac\infty\infty$ 型不定式极限的方法.
% 
% \begin{theorem}\label{theorem:ma-5-2-3}
% 设函数 $f(x),g(x)$ 在 $a$ 点的某一去心邻域 $U_0(a,\delta_0)\ (\delta_0>0)$ 上可导, 且满足:
% \begin{enumerate}
%   \item $\lim\limits_{x\to a}g(x)=\infty$;
%   \item $g'(x)\ne0,\ \forall x\in U_0(a,\delta_0)$;
%   \item $\lim\limits_{x\to a}\dfrac{f'(x)}{g'(x)}=\ell\ (\ell\text{ 有限数或 }\pm\infty,\infty)$,
% \end{enumerate}
% 则有
% \[
%   \lim_{x\to a}\frac{f(x)}{g(x)}=\ell.
% \]
% \end{theorem}
% 
% 
% \begin{proof}
% 只对 $\ell\in\mathbb{R}$ 和 $x\to a-0$ 的情形证明, 其他的情形请读者自己给出证明.
% 
% $\forall\varepsilon>0$, 由条件 (2) 和 (3) 知, $\exists\delta_1>0,0<\delta_1<\delta_0$, 当 $a-\delta_1<\xi<a$ 时, 有
% \[
%   \left|\frac{f'(\xi)}{g'(\xi)}-\ell\right|<\frac\varepsilon3.
% \]
% 对已经取定的 $\delta_1$ 及 $x\in(a-\delta_1,a)$, 在 $[a-\delta_1,x]$ 上应用柯西微分中值定理, 存在 $\xi\in[a-\delta_1,x]$, 使
% \[
%   \frac{f(x)-f(x_1)}{g(x)-g(x_1)}-\ell
%   =
%   \frac{f'(\xi)}{g'(\xi)}-\ell,
% \]
% 这里 $x_1=a-\delta_1,\ \xi\in(a-\delta_1,a)$。上式可化为
% \[
%   f(x)-f(x_1)-\ell[g(x)-g(x_1)]
%   =
%   \left[\frac{f'(\xi)}{g'(\xi)}-\ell\right][g(x)-g(x_1)],
% \]
% 整理得
% \[
%   f(x)-\ell g(x)
%   =
%   [f(x_1)-\ell g(x_1)]
%   +
%   \left[\frac{f'(\xi)}{g'(\xi)}-\ell\right][g(x)-g(x_1)].
% \]
% 在上式两边同除以 $g(x)$, 可得
% \[
%   \frac{f(x)}{g(x)}-\ell
%   =
%   \left[\frac{f'(\xi)}{g'(\xi)}-\ell\right]\left[1-\frac{g(x_1)}{g(x)}\right]
%   +
%   \frac{f(x_1)-\ell g(x_1)}{g(x)}.
% \]
% 又由于 $\lim\limits_{x\to a-0}g(x)=\infty$, 故对固定的 $x_1$, 有
% \[
%   \lim_{x\to a-0}\frac{f(x_1)-\ell g(x_1)}{g(x)}=0,\quad
%   \lim_{x\to a-0}\frac{g(x_1)}{g(x)}=0.
% \]
% 所以 $\exists\delta_2(0<\delta_2<\delta_1)$, 使得当 $a-\delta_2<x<a$ 时, 有
% \[
%   \left|\frac{f(x_1)-\ell g(x_1)}{g(x)}\right|<\frac\varepsilon2,\quad
%   \left|\frac{g(x_1)}{g(x)}\right|<\frac12.
% \]
% 令 $\delta=\min\{\delta_1,\delta_2\}$, 则当 $a-\delta<x<a$ 时, 有
% \[
%   \left|\frac{f(x)}{g(x)}-\ell\right|
%   \leqslant
%   \left|\frac{f'(\xi)}{g'(\xi)}-\ell\right|
%   \left|1-\frac{g(x_1)}{g(x)}\right|
%   +
%   \left|\frac{f(x_1)-\ell g(x_1)}{g(x)}\right|
%   \leqslant
%   \frac32\cdot\frac\varepsilon3+\frac\varepsilon2
%   =
%   \varepsilon.
% \]
% 因此 $\lim\limits_{x\to a-0}\dfrac{f(x)}{g(x)}=\ell$.
% \end{proof}
% 
% 
% \begin{remark}\label{remark:ma-5-source-46}
% 把 $x\to a$ 改为 $x\to a-0$ 或 $x\to a+0$, 上述定理的结论仍然成立.
% \end{remark}
% 
% 
% \begin{theorem}\label{theorem:ma-5-2-4}
% 设函数 $f(x),g(x)$ 在 $U=\{x: |x|>h>0\}$ 上可导, 而且满足:
% \[
%   (1)\quad \lim_{x\to\infty}g(x)=\infty;
% \]
% \[
%   (2)\quad g'(x)\ne0,\ \forall x\in U;
% \]
% \[
%   (3)\quad \lim_{x\to\infty}\frac{f'(x)}{g'(x)}=\ell\ (\ell\text{ 有限数或 }\pm\infty,\infty),
% \]
% 则有
% \[
%   \lim_{x\to\infty}\frac{f(x)}{g(x)}
%   =
%   \lim_{x\to\infty}\frac{f'(x)}{g'(x)}
%   =
%   \ell.
% \]
% 本定理的证明类似于定理~\ref{theorem:ma-5-2-2} 的证明, 请读者作为练习自己补证.
% \end{theorem}
% 
% 
% \begin{remark}\label{remark:ma-5-source-48}
% 把 $x\to\infty$ 改为 $x\to-\infty$ 或 $x\to+\infty$, 上述定理的结论仍然成立.
% \end{remark}
% 
% 
% \begin{example}\label{example:ma-5-2-4}
% 求极限
% \[
%   \lim_{x\to+\infty}\frac{\ln x}{x^\varepsilon},\quad \text{其中 }\varepsilon>0.
% \]
% \end{example}
% 
% 
% \begin{solve}
% 这是一个 $\dfrac\infty\infty$ 型不定式, 应用定理~\ref{theorem:ma-5-2-4}, 有
% \[
%   \lim_{x\to+\infty}\frac{\ln x}{x^\varepsilon}
%   =
%   \lim_{x\to+\infty}\frac{\frac1x}{\varepsilon x^{\varepsilon-1}}
%   =
%   \lim_{x\to+\infty}\frac1{\varepsilon x^\varepsilon}=0.
% \qedhere
% \]
% \end{solve}
% 
% 
% \begin{example}\label{example:ma-5-2-5}
% 求极限
% \[
%   \lim_{x\to+\infty}\frac{x^\alpha}{e^x},\quad \text{其中 }\alpha>0.
% \]
% \end{example}
% 
% 
% \begin{solve}
% 这也是一个 $\dfrac\infty\infty$ 型不定式, 应用定理~\ref{theorem:ma-5-2-4}, 有
% \[
%   \lim_{x\to+\infty}\frac{x^\alpha}{e^x}
%   =
%   \lim_{x\to+\infty}\frac{\alpha x^{\alpha-1}}{e^x}
%   =
%   \cdots
%   =
%   \lim_{x\to+\infty}
%   \frac{\alpha(\alpha-1)\cdots(\alpha-[\alpha])x^{\alpha-[\alpha]-1}}{e^x}
%   =
%   0.
% \]
% 
% 设 $I$ 是一个区间, $n$ 是一个非负整数, 在今后我们用记号 $C^nI$ 表示 $I$ 上所有具有 $n$ 阶连续导数的函数的集合; 用记号 $C^\infty I$ 表示 $I$ 上具有任意阶导数的函数的集合。特别地, $C^0I$ 表示 $I$ 上所有连续函数的集合.
% \end{solve}
% 
% 
% \begin{example}\label{example:ma-5-2-6}
% 设函数
% \[
%   f(x)=
%   \begin{cases}
%     e^{-\frac1{x^2}}, & x\ne0,\\
%     0, & x=0.
%   \end{cases}
% \]
% 证明 $f(x)\in C^\infty(-\infty,+\infty)$.
% \end{example}
% 
% 
% \begin{proof}
% 当 $x\ne0$ 时,
% \[
%   f'(x)=\frac2{x^3}e^{-\frac1{x^2}}.
% \]
% 当 $x=0$ 时,
% \[
%   f'(0)=\lim_{x\to0}\frac{f(x)-f(0)}x
%   =
%   \lim_{t\to\infty}\frac{t}{e^{t^2}}
%   =
%   0\quad \left(x=\frac1t\right),
% \]
% 且
% \[
%   \lim_{x\to0}f'(x)
%   =
%   \lim_{t\to\infty}\frac{2t^3}{e^{t^2}}
%   =
%   0,
% \]
% 所以 $f(x)\in C^1(-\infty,+\infty)$.
% 
% 继续求导, 我们发现 $f^{(n)}(x)$ 具有如下形式:
% \[
%   f^{(n)}(x)=
%   \begin{cases}
%     P_{3n}\left(\dfrac1x\right)e^{-\frac1{x^2}}, & x\ne0,\\
%     0, & x=0,
%   \end{cases}
% \]
% 其中 $P_{3n}(u)$ 为变元 $u$ 的 $3n$ 次多项式。用数学归纳法可证上述断言为真。事实上, 若它对 $n$ 成立, 则当 $x\ne0$ 时,
% \[
%   f^{(n+1)}(x)
%   =
%   \left[
%     \frac2{x^3}P_{3n}\left(\frac1x\right)
%     -
%     \frac1{x^2}P'_{3n}\left(\frac1x\right)
%   \right]e^{-\frac1{x^2}}
%   =
%   P_{3(n+1)}\left(\frac1x\right)e^{-\frac1{x^2}},
% \]
% 而
% \[
%   f^{(n+1)}(0)
%   =
%   \lim_{x\to0}
%   \frac{P_{3n}\left(\frac1x\right)e^{-\frac1{x^2}}}{x}
%   =
%   \lim_{t\to\infty}\frac{tP_{3n}(t)}{e^{t^2}}
%   =
%   0.
% \]
% 由此推知 $f(x)$ 任意次可微, 即 $f(x)\in C^\infty(-\infty,+\infty)$.
% \end{proof}
% 
% 
% \begin{remark}\label{remark:ma-5-source-55}
% 上例中的函数是一个重要的函数, 在今后的学习中还会多次遇到它。这个函数可以用来说明很多问题.
% \end{remark}
% 
% 
% \subsection{其他类型不定式}
% \label{subsec:ma-5-2-3}
% 
% 通过变量替换, 我们总可把其他形式的不定式转化为 $\dfrac00$ 型和 $\dfrac\infty\infty$ 型, 具体采用什么样替换, 这要对具体问题做具体分析, 灵活对待.
% 
% \begin{example}\label{example:ma-5-2-7}
% 求极限
% \[
%   \lim_{x\to0+0}x^\alpha\ln x\quad(\alpha>0).
% \]
% \end{example}
% 
% 
% \begin{solve}
% 这是一个 $0\cdot\infty$ 型不定式。我们有
% \[
%   \lim_{x\to0+0}x^\alpha\ln x
%   =
%   \lim_{x\to0+0}\frac{\ln x}{\frac1{x^\alpha}}
%   =
%   \lim_{x\to0+0}\frac{\frac1x}{-\frac{\alpha}{x^{\alpha+1}}}
%   =
%   0.
% \qedhere
% \]
% \end{solve}
% 
% 
% \begin{example}\label{example:ma-5-2-8}
% 求极限
% \[
%   \lim_{x\to a}\left(\frac{a_1^x+a_2^x+\cdots+a_n^x}{n}\right)^{\frac1x}
%   \quad (a_k>0,\ k=1,2,\cdots,n),
% \]
% 这里 $a=0$ 或 $\pm\infty$.
% \end{example}
% 
% 
% \begin{solve}
% 令
% \[
%   y=\left(\frac{a_1^x+a_2^x+\cdots+a_n^x}{n}\right)^{\frac1x},
%   \quad
%   \ln y=\frac1x\ln\left(\frac{a_1^x+a_2^x+\cdots+a_n^x}{n}\right).
% \]
% 当 $x\to0$ 时, 应用洛必达法则得
% \[
%   \lim_{x\to0}\ln y
%   =
%   \frac1n\ln(a_1a_2\cdots a_n)
%   =
%   \ln\sqrt[n]{a_1a_2\cdots a_n},
% \]
% 故 $\lim\limits_{x\to0}y=\sqrt[n]{a_1a_2\cdots a_n}$.
% 
% 当 $x\to+\infty$ 时, 记 $M=\max\{a_1,a_2,\cdots,a_n\}$, 可得
% \[
%   \lim_{x\to+\infty}\ln y=\ln M,\quad
%   \lim_{x\to+\infty}y=M.
% \]
% 当 $x\to-\infty$ 时, 令 $m=\min\{a_1,a_2,\cdots,a_n\}$, 可得
% \[
%   \lim_{x\to-\infty}\ln y=\ln m,\quad
%   \lim_{x\to-\infty}y=m.
% \qedhere
% \]
% \end{solve}
% 
% 
% \begin{example}\label{example:ma-5-2-9}
% 设 $x_0\in(0,1)$, 定义 $x_n=\sin x_{n-1},\ n=1,2,\cdots$, 证明
% \[
%   \lim_{n\to\infty}\sqrt{n}x_n=\sqrt3.
% \]
% \end{example}
% 
% 
% \begin{proof}
% 首先有
% \[
%   \lim_{x\to0}\left(\frac1{\sin^2x}-\frac1{x^2}\right)=\frac13.
% \]
% 事实上,
% \[
% \begin{aligned}
%   \lim_{x\to0}\left(\frac1{\sin^2x}-\frac1{x^2}\right)
%   &=
%   \lim_{x\to0}\frac{x^2-\sin^2x}{x^2\sin^2x}
%   =
%   \lim_{x\to0}\frac{x^2-\sin^2x}{x^4} \\
%   &=
%   \lim_{x\to0}\frac{2x-\sin2x}{4x^3}
%   =
%   \lim_{x\to0}\frac{1-\cos2x}{6x^2}
%   =
%   \frac13.
% \end{aligned}
% \]
% 又
% \[
%   0<x_{n+1}=\sin x_n<x_n,\quad n=1,2,\cdots,
% \]
% 故 $x_n\to0$。于是
% \[
%   \lim_{n\to\infty}\left(\frac1{x_{n+1}^2}-\frac1{x_n^2}\right)=\frac13,
% \]
% 从而由第二章 第~\ref{sec:ma-2-1}~节 中例~\ref{example:ma-2-1-13} 的结果,
% \[
%   \lim_{n\to\infty}\frac1{nx_n^2}
%   =
%   \lim_{n\to\infty}\frac1n\sum_{k=1}^{n-1}
%   \left(\frac1{x_{k+1}^2}-\frac1{x_k^2}\right)
%   =
%   \frac13.
% \]
% 于是
% \[
%   \lim_{n\to\infty}\sqrt{n}x_n=\sqrt3.
% \]
% 
% 最后我们须指出的是:
% 
% (1) 并非所有的 $\dfrac00$ 或 $\dfrac\infty\infty$ 型不定式都可用洛必达法则求其极限。例如, 容易证明
% \[
%   \lim_{x\to\infty}\frac{x+\sin x}{x}=1,
% \]
% 但是
% \[
%   \lim_{x\to\infty}\frac{(x+\sin x)'}{(x)'}
%   =
%   \lim_{x\to\infty}(1+\cos x)
% \]
% 不存在.
% 
% (2) 应用洛必达法则时, 每步必须验证 $\lim\dfrac{f'(x)}{g'(x)}$ 是否存在的条件。例如, 显然有
% \[
%   \lim_{x\to\infty}\frac{x-\sin x}{x+\sin x}=1,
% \]
% 但若盲目地使用洛必达法则, 就会得到
% \[
%   \lim_{x\to\infty}\frac{x-\sin x}{x+\sin x}
%   =
%   \lim_{x\to\infty}\frac{1-\cos x}{1+\cos x}
%   =
%   \lim_{x\to\infty}\frac{\sin x}{-\sin x}
%   =
%   -1
% \]
% 的错误结论.
% \end{proof}
% 
% 
% \section{泰勒公式}
% \label{sec:ma-5-3}
% 
% 用简单函数来逼近复杂函数, 这是一个重要的研究课题。一般来说, 我们最熟悉而又十分简单的函数就是多项式函数。在本节中, 我们将考虑两类多项式的逼近问题: 一类是在一点附近的逼近问题; 另一类则是在一个区间上的逼近问题.
% 
% 设 $f(x)$ 是区间 $I$ 上的一个函数。若我们用 $P_n(x)$ 去近似它, 则
% \[
%   R_n(x)=f(x)-P_n(x)
% \]
% 就反映了近似的误差。一个好的逼近就应该能够根据函数的性质求出所需多项式以及相应的控制误差, 并且还要使得误差尽可能地小.
% 
% \subsection{带佩亚诺余项的泰勒公式}
% \label{subsec:ma-5-3-1}
% 
% 设函数 $f(x)$ 在 $U(x_0,\delta)$ 内有定义。若 $f(x)$ 在点 $x_0$ 处连续, 则
% \[
%   f(x)-f(x_0)\to0\quad (x\to x_0),
% \]
% 即
% \[
%   f(x)=f(x_0)+o(1)\quad (x\to x_0);
% \]
% 若 $f(x)$ 在点 $x_0$ 处可导, 则
% \[
%   f(x)=f(x_0)+f'(x_0)(x-x_0)+o(x-x_0)\quad (x\to x_0).
% \]
% 由此可知, 若用线性函数 $f(x_0)+f'(x_0)(x-x_0)$ 来近似 $f(x)$, 则在 $x_0$ 附近, 比用 $f(x_0)$ 来近似它误差要小得多。若 $f(x)$ 在 $x_0$ 处具有高阶的导数, 则我们有
% 
% \begin{theorem}\label{theorem:ma-5-3-1}
% 设函数 $f(x)$ 在 $x_0$ 处具有 $n(n\geqslant1)$ 阶导数, 则有
% \[
%   f(x)=f(x_0)+f'(x_0)(x-x_0)+\frac{f''(x_0)}{2!}(x-x_0)^2+\cdots
%   +\frac{f^{(n)}(x_0)}{n!}(x-x_0)^n+o((x-x_0)^n)
%   \quad (x\to x_0).
% \]
% \end{theorem}
% 
% 
% \begin{proof}
% 要证上式成立, 只要证
% \[
% \begin{aligned}
% &\lim_{x\to x_0}\frac1{(x-x_0)^n}
% \left\{
% f(x)-\left[
% f(x_0)+f'(x_0)(x-x_0)+\frac{f''(x_0)}{2!}(x-x_0)^2 \right.\right.\\
% &\left.\left.\hspace{8em}
% \cdots+\frac{f^{(n)}(x_0)}{n!}(x-x_0)^n
% \right]
% \right\}=0
% \end{aligned}
% \]
% 即可。对上式左边应用洛必达法则 $(n-1)$ 次, 可知它等于
% \[
%   \lim_{x\to x_0}\frac1{n!}
%   \left[
%     \frac{f^{(n-1)}(x)-f^{(n-1)}(x_0)}{x-x_0}
%     -
%     f^{(n)}(x_0)
%   \right]
%   =
%   0.
% \]
% 最后的等式成立, 是因为 $f^{(n)}(x_0)$ 存在。定理
% \end{proof}
% 
% 
% \begin{remark}\label{remark:ma-5-source-64}
% 当 $f(x)$ 在 $x_0$ 处具有 $n(n\geqslant1)$ 阶导数时, 我们把多项式
% \[
%   P_n(x)=f(x_0)+f'(x_0)(x-x_0)+\frac{f''(x_0)}{2!}(x-x_0)^2
%   +\cdots+\frac{f^{(n)}(x_0)}{n!}(x-x_0)^n
% \]
% 称为 $f(x)$ 在 $x_0$ 处的泰勒 (Taylor) 多项式, 而将
% \[
%   f(x)=f(x_0)+f'(x_0)(x-x_0)+\frac{f''(x_0)}{2!}(x-x_0)^2
%   +\cdots+\frac{f^{(n)}(x_0)}{n!}(x-x_0)^n+R_n(x)
% \]
% 称为 $f(x)$ 在 $x_0$ 处的泰勒公式, 其中 $R_n(x)=f(x)-P_n(x)$ 称为泰勒公式的余项。特别地, 当 $x_0=0$ 时, 我们称此时的泰勒公式为麦克劳林 (Maclaurin) 公式.
% \end{remark}
% 
% 定理~\ref{theorem:ma-5-3-1} 告诉我们, 当 $f(x)$ 在 $x_0$ 处存在 $n$ 阶导数时, 泰勒公式的余项为
% \[
%   R_n(x)=o((x-x_0)^n)\quad (x\to x_0).
% \]
% 这种余项称为佩亚诺 (Peano) 余项.
% 
% 值得指出的是, 若 $f(x)$ 在点 $x_0$ 存在 $n$ 阶导数, 且在点 $x_0$ 附近成立
% \begin{equation}
%   f(x)=a_0+a_1(x-x_0)+\cdots+a_n(x-x_0)^n+o((x-x_0)^n)\quad (x\to x_0),
%   \label{eq:ma-5-3-1}
% \end{equation}
% 则必有
% \[
%   a_k=\frac{f^{(k)}(x_0)}{k!}\quad (k=0,1,2,\cdots,n).
% \]
% 事实上, 在 \eqref{eq:ma-5-3-1} 式中令 $x\to x_0$, 即得 $a_0=f(x_0)$。现将 \eqref{eq:ma-5-3-1} 式改写并逐次除以 $(x-x_0),(x-x_0)^2,\cdots$, 再令 $x\to x_0$, 即可依次推出
% \[
%   a_1=f'(x_0),\quad a_2=\frac{f''(x_0)}{2!},\quad \cdots,\quad
%   a_k=\frac{f^{(k)}(x_0)}{k!}.
% \]
% 以上论述说明了 $f(x)$ 在 $x_0$ 处的多项式逼近的某种唯一性, 它将给我们寻找 $f(x)$ 的泰勒公式提供许多方便.
% 
% \begin{example}\label{example:ma-5-3-1}
% 求 $f(x)=e^x$ 的带佩亚诺余项的麦克劳林公式.
% \end{example}
% 
% 
% \begin{solve}
% 由 $f^{(k)}(x)=e^x$, 且 $e^0=1$, 得
% \[
%   a_k=\frac{f^{(k)}(0)}{k!}=\frac1{k!},\quad k=0,1,\cdots,n,
% \]
% 所以
% \[
%   e^x=1+\frac{x}{1!}+\frac{x^2}{2!}+\cdots+\frac{x^n}{n!}+o(x^n)\quad (x\to0).
% \qedhere
% \]
% \end{solve}
% 
% 
% \begin{example}\label{example:ma-5-3-2}
% 求 $f(x)=\sin x$ 的带佩亚诺余项的麦克劳林公式.
% \end{example}
% 
% 
% \begin{solve}
% 由 $f^{(k)}(x)=\sin\left(x+\dfrac{k\pi}{2}\right)$, 得
% \[
%   f^{(2k)}(0)=0,\quad f^{(2k+1)}(0)=(-1)^k,\quad k=0,1,\cdots,n,
% \]
% 所以
% \[
%   \sin x=x-\frac{x^3}{3!}+\frac{x^5}{5!}+\cdots+(-1)^{n-1}\frac{x^{2n-1}}{(2n-1)!}+o(x^{2n})\quad (x\to0).
% \qedhere
% \]
% \end{solve}
% 
% 
% \begin{example}\label{example:ma-5-3-3}
% 求 $f(x)=\cos x$ 的带佩亚诺余项的麦克劳林公式.
% \end{example}
% 
% 
% \begin{solve}
% 由 $f^{(k)}(x)=\cos\left(x+\dfrac{k\pi}{2}\right)$, 得
% \[
%   f^{(2k)}(0)=(-1)^k,\quad f^{(2k+1)}(0)=0,\quad k=0,1,\cdots,n,
% \]
% 所以
% \[
%   \cos x=1-\frac{x^2}{2!}+\frac{x^4}{4!}+\cdots+(-1)^n\frac{x^{2n}}{(2n)!}+o(x^{2n+1})\quad (x\to0).
% \qedhere
% \]
% \end{solve}
% 
% 
% \begin{example}\label{example:ma-5-3-4}
% 求 $f(x)=\ln(1+x)$ 的带佩亚诺余项的麦克劳林公式.
% \end{example}
% 
% 
% \begin{solve}
% 由
% \[
%   f^{(k)}(x)=(-1)^{k-1}\frac{(k-1)!}{(1+x)^k},
% \]
% 得
% \[
%   f(0)=0,\quad f^{(k)}(0)=(-1)^{k-1}(k-1)!,\quad k=1,2,\cdots,n,
% \]
% 所以
% \[
%   \ln(1+x)=x-\frac{x^2}{2}+\frac{x^3}{3}+\cdots+(-1)^{n-1}\frac{x^n}{n}+o(x^n)\quad (x\to0).
% \qedhere
% \]
% \end{solve}
% 
% 
% \begin{example}\label{example:ma-5-3-5}
% 求 $f(x)=(1+x)^\alpha\ (\alpha\in\mathbb{R})$ 的带佩亚诺余项的麦克劳林公式.
% \end{example}
% 
% 
% \begin{solve}
% 由
% \[
%   f^{(k)}(x)=\alpha(\alpha-1)\cdots(\alpha-k+1)(1+x)^{\alpha-k},
% \]
% 得
% \[
%   a_0=1,\quad
%   a_k=\frac{\alpha(\alpha-1)\cdots(\alpha-k+1)}{k!},\quad k=1,2,\cdots,n,
% \]
% 所以
% \[
%   (1+x)^\alpha
%   =
%   1+\alpha x+\frac{\alpha(\alpha-1)}{2!}x^2+\cdots
%   +\frac{\alpha(\alpha-1)\cdots(\alpha-n+1)}{n!}x^n+o(x^n)
%   \quad (x\to0).
% \]
% 
% 例~\ref{example:ma-5-3-1}--例~\ref{example:ma-5-3-5} 中函数的麦克劳林公式在今后的各章节及后续课程中经常要用到, 读者应牢牢地记住它们.
% \end{solve}
% 
% 
% \begin{example}\label{example:ma-5-3-6}
% 求 $f(x)=\sin^2(1+x^2)$ 的带佩亚诺余项的麦克劳林公式.
% \end{example}
% 
% 
% \begin{solve}
% 由于
% \[
% \begin{aligned}
%   \sin^2(1+x^2)
%   &=
%   \frac12-\frac12\cos(2+2x^2) \\
%   &=
%   \frac12-\frac12\bigl(\cos2\cos(2x^2)-\sin2\sin(2x^2)\bigr),
% \end{aligned}
% \]
% 而
% \[
%   \cos(2x^2)=1-\frac{(2x^2)^2}{2!}+\frac{(2x^2)^4}{4!}+\cdots+(-1)^n\frac{(2x^2)^{2n}}{(2n)!}+o(x^{4n+2}),
% \]
% \[
%   \sin(2x^2)=2x^2-\frac{(2x^2)^3}{3!}+\cdots+(-1)^{n-1}\frac{(2x^2)^{2n-1}}{(2n-1)!}+o(x^{4n}),
% \]
% 因此
% \[
% \begin{aligned}
%   \sin^2(1+x^2)
%   &=
%   \frac12-\frac12\cos2+(\sin2)x^2+(\cos2)x^4-\frac{2^2\sin2}{3!}x^6
%   -\frac{2^3\cos2}{4!}x^8 \\
%   &\quad+\cdots+(-1)^{n-1}\frac{2^{2n-2}\sin2}{(2n-1)!}x^{4n-2}
%   +(-1)^{n+1}\frac{2^{2n-1}\cos2}{(2n)!}x^{4n}
%   +o(x^{4n})
%   \quad (x\to0).
% \end{aligned}
% \]
% 
% 由于 $f'(x)$ 的 $n$ 阶导数是 $f(x)$ 的 $n+1$ 阶导数, 因此若
% \[
%   f'(x)=a_1+a_2x+a_3x^2+\cdots+a_nx^{n-1}+o(x^{n-1})\quad (x\to0),
% \]
% 则
% \[
%   f(x)=f(0)+a_1x+\frac{a_2}{2}x^2+\frac{a_3}{3}x^3+\cdots+\frac{a_n}{n}x^n+o(x^n)\quad (x\to0).
% \]
% 下面我们利用这一结果, 导出 $\arctan x$ 和 $\arcsin x$ 的麦克劳林公式.
% \end{solve}
% 
% 
% \begin{example}\label{example:ma-5-3-7}
% 求 $\arctan x$ 的带佩亚诺余项的麦克劳林公式.
% \end{example}
% 
% 
% \begin{solve}
% 由于
% \[
%   (\arctan x)'=\frac1{1+x^2}=1-x^2+x^4+\cdots+(-1)^nx^{2n}+o(x^{2n+1})\quad (x\to0),
% \]
% 因此
% \[
%   \arctan x=x-\frac13x^3+\frac15x^5+\cdots+\frac{(-1)^n}{2n+1}x^{2n+1}+o(x^{2n+2})\quad (x\to0).
% \qedhere
% \]
% \end{solve}
% 
% 
% \begin{example}\label{example:ma-5-3-8}
% 求 $\arcsin x$ 的带佩亚诺余项的麦克劳林公式.
% \end{example}
% 
% 
% \begin{solve}
% 由于
% \[
%   (\arcsin x)'=(1-x^2)^{-\frac12}
%   =
%   1+\frac1{2!!}x^2+\frac{3!!}{4!!}x^4+\cdots+
%   \frac{(2n-1)!!}{(2n)!!}x^{2n}+o(x^{2n+1})
%   \quad (x\to0),
% \]
% 因此
% \[
%   \arcsin x
%   =
%   x+\frac13\frac1{2!!}x^3+\frac15\frac{3!!}{4!!}x^5+\cdots
%   +\frac1{2n+1}\frac{(2n-1)!!}{(2n)!!}x^{2n+1}
%   +o(x^{2n+2})
%   \quad (x\to0).
% \]
% 
% 下面我们再举几个例子来说明如何求一个函数的泰勒公式.
% \end{solve}
% 
% 
% \begin{example}\label{example:ma-5-3-9}
% 求 $e^{\cos x}$ 的带佩亚诺余项的麦克劳林公式 (展到 $x^4$).
% \end{example}
% 
% 
% \begin{solve}
% 利用 $e^u$ 在 $u=0$ 及 $\cos x$ 在 $x=0$ 的泰勒公式,
% \[
% \begin{aligned}
%   e^{\cos x}
%   &=
%   e\cdot e^{\cos x-1} \\
%   &=
%   e\left[1+\cos x-1+\frac{(\cos x-1)^2}{2!}+o((\cos x-1)^2)\right] \\
%   &=
%   e\left[1+\left(-\frac{x^2}{2!}+\frac{x^4}{4!}+o(x^4)\right)
%   +\frac1{2!}\left(-\frac{x^2}{2!}+o(x^2)\right)^2+o(x^4)\right] \\
%   &=
%   e-\frac e2x^2+\frac e6x^4+o(x^4)\quad (x\to0).
% \end{aligned}
% \qedhere
% \]
% \end{solve}
% 
% 
% \begin{example}\label{example:ma-5-3-10}
% 求 $e^x\ln(1+x)$ 的带佩亚诺余项的麦克劳林公式 (展到 $x^4$).
% \end{example}
% 
% 
% \begin{solve}
% 我们利用 $e^x$ 和 $\ln(1+x)$ 的麦克劳林公式得
% \[
% \begin{aligned}
%   e^x\ln(1+x)
%   &=
%   \left(1+x+\frac{x^2}{2!}+\frac{x^3}{3!}+o(x^3)\right)
%   \left(x-\frac{x^2}{2}+\frac{x^3}{3}-\frac{x^4}{4}+o(x^4)\right) \\
%   &=
%   x+\frac{x^2}{2}+\frac{x^3}{3}+o(x^4)\quad (x\to0).
% \end{aligned}
% \]
% 
% 下面我们来给出泰勒公式的一些简单应用.
% \end{solve}
% 
% 
% \begin{example}\label{example:ma-5-3-11}
% 求极限
% \[
%   \lim_{x\to+\infty}\left(\sqrt[3]{x^3+3x}-\sqrt{x^2-2x}\right).
% \]
% \end{example}
% 
% 
% \begin{solve}
% 这是一个 $\infty-\infty$ 型不定式。下面我们用泰勒展开来求极限.
% \[
% \begin{aligned}
%   &\lim_{x\to+\infty}\left(\sqrt[3]{x^3+3x}-\sqrt{x^2-2x}\right) \\
%   &=
%   \lim_{x\to+\infty}x
%   \left\{
%     \left[1+\frac3{x^2}\right]^{\frac13}
%     -
%     \left[1-\frac2x\right]^{\frac12}
%   \right\} \\
%   &=
%   \lim_{x\to+\infty}x
%   \left[
%     \frac1{x^2}+\frac1x+o\left(\frac1x\right)
%   \right]
%   =
%   1.
% \end{aligned}
% \qedhere
% \]
% \end{solve}
% 
% 
% \begin{example}\label{example:ma-5-3-12}
% 求极限
% \[
%   \lim_{n\to\infty}n^2\left(1-n\sin\frac1n\right).
% \]
% \end{example}
% 
% 
% \begin{solve}
% 将原式恒等变形, 利用泰勒公式得
% \[
% \begin{aligned}
%   \lim_{n\to\infty}n^2\left(1-n\sin\frac1n\right)
%   &=
%   \lim_{n\to\infty}n^3\left(\frac1n-\sin\frac1n\right) \\
%   &=
%   \lim_{n\to\infty}n^3
%   \left\{
%     \frac1n-
%     \left[\frac1n-\frac1{3!}\frac1{n^3}+o\left(\frac1{n^4}\right)\right]
%   \right\}
%   =
%   \frac16.
% \end{aligned}
% \qedhere
% \]
% \end{solve}
% 
% 
% \subsection{带拉格朗日余项的泰勒公式}
% \label{subsec:ma-5-3-2}
% 
% 上面我们讨论了函数在一点附近用多项式逼近的问题, 相应的误差只给出定性描述, 不能具体估计误差的大小, 所以它只适用于求无穷小量的阶或求极限等问题。若要具体计算函数值并且达到预先指定的误差, 就需要给出误差的定量描述。带拉格朗日余项的泰勒公式就是为了解决这一问题而产生的.
% 
% \begin{theorem}\label{theorem:ma-5-3-2}
% 设 $f(x)\in C^n[a,b]$, 而且在 $(a,b)$ 内存在 $n+1$ 阶导数, 则对任意 $x,x_0\in[a,b]$, 有
% \begin{equation}
% \begin{aligned}
%   f(x)
%   &=
%   f(x_0)+\frac{f'(x_0)}{1!}(x-x_0)+\frac{f''(x_0)}{2!}(x-x_0)^2
%   +\cdots+\frac{f^{(n)}(x_0)}{n!}(x-x_0)^n \\
%   &\quad+
%   \frac{f^{(n+1)}(\xi)}{(n+1)!}(x-x_0)^{n+1},
% \end{aligned}
%   \label{eq:ma-5-3-2}
% \end{equation}
% 其中 $\xi$ 介于 $x$ 与 $x_0$ 之间.
% \end{theorem}
% 
% 
% \begin{proof}
% 作两个辅助函数:
% \[
%   F(t)=f(x)-\left[
%   f(t)+\frac{f'(t)}{1!}(x-t)+\frac{f''(t)}{2!}(x-t)^2
%   +\cdots+\frac{f^{(n)}(t)}{n!}(x-t)^n
%   \right]
% \]
% 和
% \[
%   G(t)=(x-t)^{n+1}.
% \]
% 容易验证它们在相应闭区间上连续, 在相应开区间内可导, 且 $F(x)=G(x)=0$。应用柯西微分中值定理, 并且注意到
% \[
%   F'(t)=-\frac{f^{(n+1)}(t)}{n!}(x-t)^n,\quad
%   G'(t)=-(n+1)(x-t)^n,
% \]
% 便有
% \[
%   \frac{F(x_0)}{G(x_0)}
%   =
%   \frac{F'(\xi)}{G'(\xi)}
%   =
%   \frac{f^{(n+1)}(\xi)}{(n+1)!},
% \]
% 这里 $\xi$ 介于 $x$ 与 $x_0$ 之间。由此即得结论.
% \end{proof}
% 
% 
% \begin{remark}\label{remark:ma-5-source-91}
% \eqref{eq:ma-5-3-2} 中的余项
% \[
%   \frac{f^{(n+1)}(\xi)}{(n+1)!}(x-x_0)^{n+1}
% \]
% 称为拉格朗日余项。另外, 我们可将上述定理中的 $\xi$ 表为
% \[
%   \xi=x_0+\theta(x-x_0),\quad \text{其中 }0<\theta<1.
% \]
% \end{remark}
% 
% 下面我们列出几个常见函数的带拉格朗日余项的泰勒公式:
% \[
%   e^x=1+x+\frac{x^2}{2!}+\cdots+\frac{x^n}{n!}+\frac{e^{\theta x}}{(n+1)!}x^{n+1}
%   \quad (-\infty<x<+\infty,\ 0<\theta<1);
% \]
% \[
%   \sin x=x-\frac{x^3}{3!}+\cdots+(-1)^{n-1}\frac{x^{2n-1}}{(2n-1)!}
%   +(-1)^n\frac{\cos\theta x}{(2n+1)!}x^{2n+1}
%   \quad (-\infty<x<+\infty,\ 0<\theta<1);
% \]
% \[
%   \cos x=1-\frac{x^2}{2!}+\frac{x^4}{4!}+\cdots+(-1)^n\frac{x^{2n}}{(2n)!}
%   +(-1)^{n+1}\frac{\cos\theta x}{(2n+2)!}x^{2n+2}
%   \quad (-\infty<x<+\infty,\ 0<\theta<1);
% \]
% \[
%   \ln(1+x)=x-\frac{x^2}{2}+\cdots+(-1)^{n-1}\frac{x^n}{n}
%   +(-1)^n\frac{x^{n+1}}{(n+1)(1+\theta x)^{n+1}}
%   \quad (|x|<1,\ 0<\theta<1);
% \]
% \[
% \begin{aligned}
%   (1+x)^\alpha
%   &=
%   1+\alpha x+\frac{\alpha(\alpha-1)}{2!}x^2+\cdots
%   +\frac{\alpha(\alpha-1)\cdots(\alpha-n+1)}{n!}x^n \\
%   &\quad+
%   \frac{\alpha(\alpha-1)\cdots(\alpha-n)}{(n+1)!}(1+\theta x)^{\alpha-n-1}x^{n+1}
%   \quad (|x|<1,\ 0<\theta<1).
% \end{aligned}
% \]
% 利用拉格朗日余项, 我们可以估计在整个区间上 $f(x)$ 用泰勒公式逼近时的误差。这在实际应用中是十分重要的.
% 
% \begin{example}\label{example:ma-5-3-13}
% 求 $e$ 的近似值, 使得其误差不超过 $10^{-5}$.
% \end{example}
% 
% 
% \begin{solve}
% 在 $e^x$ 的带拉格朗日余项的泰勒公式中令 $x=1$, 得
% \[
%   e=1+1+\frac1{2!}+\cdots+\frac1{n!}+\frac{e^\theta}{(n+1)!},\quad 0<\theta<1.
% \]
% 欲使
% \[
%   \frac{e^\theta}{(n+1)!}<\frac3{(n+1)!}<10^{-5},
% \]
% 只需取 $n=8$ 即可。于是, 有
% \[
%   e\approx1+1+\frac1{2!}+\cdots+\frac1{8!}\approx2.71828.
% \qedhere
% \]
% \end{solve}
% 
% 
% \begin{example}\label{example:ma-5-3-14}
% 试证 $e$ 是无理数.
% \end{example}
% 
% 
% \begin{proof}
% 用反证法。假设 $e$ 是有理数, 它表示成分数时的分母为 $N$。现取正整数 $n>\max\{N,3\}$。根据 $e^x$ 的带拉格朗日余项的泰勒公式, 得
% \[
%   e-\left(1+1+\frac1{2!}+\cdots+\frac1{n!}\right)=\frac{e^\theta}{(n+1)!}.
% \]
% 这里 $0<\theta<1$, 从而 $1<e^\theta<3$。上式两边乘以 $n!$, 得
% \[
%   n!\left[e-\left(1+1+\frac1{2!}+\cdots+\frac1{n!}\right)\right]
%   =
%   \frac{e^\theta}{n+1}.
% \]
% 上式左边是整数, 而右边满足 $0<\dfrac{e^\theta}{n+1}<1$。这一矛盾说明 $e$ 是无理数.
% \end{proof}
% 
% 
% \begin{example}\label{example:ma-5-3-15}
% 求极限
% \[
%   \lim_{n\to\infty}n\sin(2\pi en!).
% \]
% \end{example}
% 
% 
% \begin{solve}
% 根据 $e^x$ 的带拉格朗日余项的泰勒公式, 得
% \[
%   n!e
%   =
%   n!\left(1+1+\frac1{2!}+\cdots+\frac1{n!}+\frac1{(n+1)!}+\frac{e^\theta}{(n+2)!}\right)
%   =
%   k+\frac1{n+1}+\frac{e^\theta}{(n+2)(n+1)},
% \]
% 其中 $0<\theta<1$, 而
% \[
%   k=n!\left(1+1+\frac1{2!}+\cdots+\frac1{n!}\right)
% \]
% 是正整数。于是
% \[
% \begin{aligned}
%   \lim_{n\to\infty}n\sin(2\pi en!)
%   &=
%   \lim_{n\to\infty}
%   n\sin\left(2k\pi+\frac{2\pi}{n+1}
%   +\frac{2\pi e^\theta}{(n+2)(n+1)}\right) \\
%   &=
%   \lim_{n\to\infty}
%   n\left[\frac{2\pi}{n+1}+o\left(\frac1n\right)\right]
%   =
%   2\pi.
% \end{aligned}
% \qedhere
% \]
% \end{solve}
% 
% 
% \subsection{拉格朗日插值多项式}
% \label{subsec:ma-5-3-3}
% 
% 在实际应用中大量的函数都是用表格给出的, 为了进行理论研究和工程设计, 需要寻找与已知函数值相符而形式简单的函数去近似它。解决这类问题的最基本方法之一就是多项式插值法。设 $y=f(x)$ 是在区间 $[a,b]$ 上定义的某个函数, 已知它在该区间上 $n+1$ 个不同点 $x_0,x_1,x_2,\cdots,x_n$ 处的函数值为
% \[
%   y_i=f(x_i)\quad (i=0,1,2,\cdots,n),
% \]
% 寻找一个多项式 $P(x)$, 使得
% \begin{equation}
%   P(x_i)=y_i\quad (i=0,1,2,\cdots,n).
%   \label{eq:ma-5-3-interpolation}
% \end{equation}
% 对于给定的插值数据 $\{(x_i,y_i)\}_{i=0}^n$, 满足条件 \eqref{eq:ma-5-3-interpolation} 的多项式 $P(x)$ 称为 $f(x)$ 的插值多项式.
% 
% 先来看 $n=1$ 的情形。通过 $(x_0,y_0)$ 和 $(x_1,y_1)$ 两点的直线可表示为
% \[
%   L_1(x)=y_0\cdot\frac{x-x_1}{x_0-x_1}+y_1\cdot\frac{x-x_0}{x_1-x_0}.
% \]
% 其中
% \[
%   \ell_0(x)=\frac{x-x_1}{x_0-x_1},\quad
%   \ell_1(x)=\frac{x-x_0}{x_1-x_0},
% \]
% 满足
% \[
%   \ell_0(x_0)=1,\quad \ell_0(x_1)=0,\quad \ell_1(x_0)=0,\quad \ell_1(x_1)=1.
% \]
% 若能构造出具有如下性质的特殊插值多项式 $\ell_i(x)$:
% \[
%   \ell_i(x_j)=\delta_{ij}=
%   \begin{cases}
%     1, & i=j,\\
%     0, & i\ne j,
%   \end{cases}
% \]
% 则多项式
% \[
%   P(x)=\sum_{i=0}^n y_i\cdot\ell_i(x)
% \]
% 就是满足插值条件的插值多项式。注意到 $x_0,\cdots,x_n$ 中除 $x_i$ 外, 均为多项式 $\ell_i(x)$ 的零点, 即知
% \[
%   \ell_i(x)
%   =
%   \frac{(x-x_0)\cdots(x-x_{i-1})(x-x_{i+1})\cdots(x-x_n)}
%        {(x_i-x_0)\cdots(x_i-x_{i-1})(x_i-x_{i+1})\cdots(x_i-x_n)}.
% \]
% 我们把上面所确定的插值多项式记为
% \[
%   L_n(x)=\sum_{i=0}^n y_i\frac{\omega(x)}{(x-x_i)\omega'(x_i)},
% \]
% 其中
% \begin{equation}
%   \omega(x)=(x-x_0)(x-x_1)\cdots(x-x_n),
%   \label{eq:ma-5-3-3}
% \end{equation}
% 并且称之为 $n$ 阶拉格朗日插值多项式, 而称 $\{\ell_i(x)\}_{i=0}^n$ 为拉格朗日插值基函数。另外, 我们称 $x_0,x_1,\cdots,x_n$ 为插值节点.
% 
% \begin{example}\label{example:ma-5-3-16}
% 已知函数 $f(x)$ 在三点的函数值为
% \[
% \begin{array}{c|ccc}
%   x & 1.0 & 1.5 & 2.0\\
%   \hline
%   f(x) & 0.0000 & 0.4055 & 0.6931
% \end{array}
% \]
% 求它的二阶插值多项式 $P_2(x)$, 并用 $P_2(x)$ 估算 $f(1.2)$.
% \end{example}
% 
% 
% \begin{solve}
% 记 $x_0=1.0,x_1=1.5,x_2=2.0$, 则有
% \[
% \begin{aligned}
%   L_2(x)
%   &=
%   0.0000\times\frac{(x-1.5)(x-2)}{(1-1.5)(1-2)}
%   +0.4055\times\frac{(x-1)(x-2)}{(1.5-1)(1.5-2)} \\
%   &\quad
%   +0.6931\times\frac{(x-1)(x-1.5)}{(2-1)(2-1.5)} \\
%   &=
%   -1.6220(x-1)(x-2)+1.3862(x-1)(x-1.5).
% \end{aligned}
% \]
% 取 $x=1.2$, 得
% \[
% \begin{aligned}
%   f(1.2)\approx L_2(1.2)
%   &=
%   -1.6220(1.2-1)(1.2-2)
%   +1.3862(1.2-1)(1.2-1.5) \\
%   &=
%   0.176348.
% \end{aligned}
% \]
% 由于原始数据只有 $4$ 位小数, 我们可以取 $f(1.2)=0.1763$.
% 
% 用插值多项式 $L_n(x)$ 作为被插函数 $f(x)$ 的逼近, 除了在插值节点处以外, $L_n(x)$ 与被插函数 $f(x)$ 是有差别的, 下面的定理给出了插值余项 (误差) 的一个表达式.
% \end{solve}
% 
% 
% \begin{theorem}\label{theorem:ma-5-3-3}
% 设被插函数 $f(x)\in C^{(n+1)}[a,b]$, 且插值节点互不相同, 则对任意的 $x\in[a,b]$, 都存在 $\xi\in[a,b]$, 使得
% \[
%   R_n(x)=f(x)-L_n(x)=\frac{f^{(n+1)}(\xi)}{(n+1)!}\omega(x),
% \]
% 其中 $\omega(x)$ 如 \eqref{eq:ma-5-3-3} 所定义.
% \end{theorem}
% 
% 
% \begin{proof}
% 由于
% \[
%   f(x_j)=L_n(x_j),\quad j=0,1,2,\cdots,n,
% \]
% 所以 $x_0,x_1,\cdots,x_n$ 为误差函数 $R_n(x)$ 的零点。于是
% \[
%   R_n(x)=K(x)\omega(x).
% \]
% 为给出 $K(x)$ 的表达式, 引入变量为 $t$ 的函数
% \[
%   F(t)=f(t)-L_n(t)-K(x)\omega(t).
% \]
% 任取 $x\in[a,b]$, 且 $x\ne x_j$, 则有 $F(x)=0,\ F(x_j)=0$。反复应用罗尔微分中值定理, 可得 $F^{(n+1)}(t)$ 至少有一个零点 $\xi$, 即
% \[
%   0=F^{(n+1)}(\xi)=f^{(n+1)}(\xi)-K(x)(n+1)!.
% \]
% 因此
% \[
%   K(x)=\frac{f^{(n+1)}(\xi)}{(n+1)!}.
% \qedhere
% \]
% \end{proof}
% 
% 
% \section{利用导数研究函数}
% \label{sec:ma-5-4}
% 
% 这一节我们将利用导数来研究函数的单调性、极值、凹凸性和拐点等。最后再利用函数的这些特性, 给出描绘函数图像的方法.
% 
% \subsection{函数的单调性}
% \label{subsec:ma-5-4-1}
% 
% 这一小节, 作为拉格朗日中值微分定理的应用, 我们来讨论如何利用导数的符号来判断函数的单调性。事实上, 从几何上来看, 函数的升降, 与切线的方向密切相关。曲线单调上升时, 切线与 $x$ 轴正向的夹角为锐角; 曲线单调下降时, 切线与 $x$ 轴正向的夹角为钝角。这样, 我们便有
% 
% \begin{theorem}\label{theorem:ma-5-4-1}
% 设函数 $f(x)$ 在区间 $I$ 内可导, 则
% 
% (1) $f(x)$ 在 $I$ 内单调上升的充分必要条件是 $f'(x)\geqslant0,\ \forall x\in I$;
% 
% (2) $f(x)$ 在 $I$ 内单调下降的充分必要条件是 $f'(x)\leqslant0,\ \forall x\in I$.
% \end{theorem}
% 
% 
% \begin{proof}
% 我们只证 (1)。(2) 的证明是类似的.
% 
% 当 $f(x)$ 在 $I$ 内单调上升时, 对任意的 $x\in I$, 取充分小的 $\Delta x>0$ 使得 $x+\Delta x\in I$。由于 $f(x+\Delta x)-f(x)\geqslant0$, 从而有
% \[
%   \frac{f(x+\Delta x)-f(x)}{\Delta x}\geqslant0.
% \]
% 由于 $f(x)$ 在 $I$ 内可导, 因此有
% \[
%   f'(x)=f'(x+0)=\lim_{\Delta x\to0+0}\frac{f(x+\Delta x)-f(x)}{\Delta x}\geqslant0.
% \]
% 这就证明了必要性.
% 
% 现证充分性。任取 $x_1,x_2\in I$, 并且 $x_1<x_2$, 在 $[x_1,x_2]$ 上应用拉格朗日微分中值定理, 知 $\exists\xi\in(x_1,x_2)$, 使得
% \[
%   f(x_2)-f(x_1)=f'(\xi)(x_2-x_1)\geqslant0.
% \]
% 因此 $f(x_1)\leqslant f(x_2)$, 即 $f(x)$ 在 $I$ 内单调上升.
% \end{proof}
% 
% 
% \begin{remark}\label{remark:ma-5-source-104}
% 大家已经知道, 一个可微函数在一个区间上为常数的充分必要条件是该函数的导数在该区间上恒等于零。因此我们立即可得: 一个可导函数 $f(x)$ 在 $I$ 内严格单调上升 (下降) 的充分必要条件是在区间 $I$ 上 $f'(x)\geqslant0\ (f'(x)\leqslant0)$, 而且在该区间的任何子区间上 $f'(x)$ 都不恒等于 $0$.
% \end{remark}
% 
% 
% \begin{example}\label{example:ma-5-4-1}
% 设函数
% \[
%   f(x)=
%   \begin{cases}
%     e^{\sin\frac1x-\frac1x}, & x>0,\\
%     0, & x=0.
%   \end{cases}
% \]
% 对于任意的 $x>0$, 有
% \[
%   f'(x)=e^{\sin\frac1x-\frac1x}\left(1-\cos\frac1x\right)\frac1{x^2}\geqslant0,
% \]
% 等号成立当且仅当 $x=\dfrac1{2k\pi}\ (k=1,2,\cdots)$。此外, 由于
% \[
%   \lim_{x\to0+0}f(x)=0=f(0),
% \]
% 所以 $f(x)$ 在 $[0,+\infty)$ 上连续。这样, 应用定理~\ref{theorem:ma-5-4-1} 的注即知, $f(x)$ 在 $[0,+\infty)$ 上严格单调上升.
% \end{example}
% 
% 
% \begin{example}\label{example:ma-5-4-2}
% 证明不等式
% \[
%   \tan x>x+\frac13x^3\quad \left(0<x<\frac\pi2\right).
% \]
% \end{example}
% 
% 
% \begin{proof}
% 令
% \[
%   f(x)=\tan x-x-\frac13x^3,
% \]
% 则有
% \[
%   f'(x)=\frac1{\cos^2x}-1-x^2=\tan^2x-x^2>0,\quad \forall x\in\left(0,\frac\pi2\right),
% \]
% 这是因为在该区间上 $\tan x>x$ 之故。因此, $f(x)$ 在 $[0,\pi/2)$ 上严格单调上升。又 $f(0)=0$, 故有
% \[
%   0=f(0)<f(x)=\tan x-x-\frac13x^3,\quad \forall x\in\left(0,\frac\pi2\right),
% \]
% 从而不等式得证.
% \end{proof}
% 
% 
% \subsection{函数的极值}
% \label{subsec:ma-5-4-2}
% 
% 在本章的第一节我们已经介绍了极值的定义和费马定理。我们已经知道, 一个在区间 $I$ 内可导的函数 $f(x)$, 若它在 $I$ 内的一点 $x_0$ 上取极值, 则必有 $f'(x_0)=0$ (即 $x_0$ 是 $f(x)$ 的一个驻点), 但反之不真。这样, 假如我们已经知道 $x_0$ 是 $f(x)$ 的一个驻点, 那么在什么条件下它是 $f(x)$ 的一个极值点呢? 下面我们利用导数的符号来回答这一问题.
% 
% \begin{theorem}\label{theorem:ma-5-4-2}
% 设函数 $f(x)$ 在 $U_0(x_0,\delta)$ 内可导且在点 $x_0$ 处连续.
% 
% (1) 若当 $x_0-\delta<x<x_0$ 时, $f'(x)>0$, 而当 $x_0<x<x_0+\delta$ 时, $f'(x)<0$, 则 $f(x)$ 在点 $x_0$ 取得严格极大值;
% 
% (2) 若当 $x_0-\delta<x<x_0$ 时, $f'(x)<0$, 而当 $x_0<x<x_0+\delta$ 时, $f'(x)>0$, 则 $f(x)$ 在点 $x_0$ 取得严格极小值;
% 
% (3) 若当 $x_0-\delta<x<x_0$ 及 $x_0<x<x_0+\delta$ 时, 都有 $f'(x)>0$, 或者 $f'(x)<0$, 则 $x_0$ 不是 $f(x)$ 的极值点.
% \end{theorem}
% 
% 请读者作为练习给出定理的证明.
% 
% 定理~\ref{theorem:ma-5-4-2} 就是说, 当 $x$ 从点 $x_0$ 的左侧经过 $x_0$ 而变到其右侧时, 若 $f'(x)$ 改变符号, 则 $f(x)$ 在点 $x_0$ 取得极值: 由正变负, 取极大值; 由负变正, 取极小值。若 $f'(x)$ 不改变符号, 则 $f(x)$ 在点 $x_0$ 不取极值。对于上述定理中的条件与结论, 读者应该根据函数图像的几何特征及导数的几何意义来加以理解与记忆, 而不是简单地记住该定理。如图 5.4.1 所示, $x_0$ 是 $f(x)$ 的极大值点, $x_1$ 是 $f(x)$ 的极小值点, $x_2$ 不是 $f(x)$ 的极值点.
% 
% \begin{example}\label{example:ma-5-4-3}
% 求函数 $f(x)=(x+2)^2(x-1)^3$ 在 $(-\infty,+\infty)$ 上的极值.
% \end{example}
% 
% 
% \begin{solve}
% 直接计算, 得
% \[
%   f'(x)=(x+2)(x-1)^2(5x+4).
% \]
% 由此可知, 该函数有 $3$ 个驻点:
% \[
%   x_1=-2,\qquad x_2=-\frac45,\qquad x_3=1.
% \]
% 这 $3$ 个点将实数轴分成 $4$ 个区间:
% \[
%   (-\infty,-2),\qquad \left(-2,-\frac45\right),\qquad
%   \left(-\frac45,1\right),\qquad (1,+\infty).
% \]
% 容易看出, 导数 $f'(x)$ 在这 $4$ 个区间上的符号分别为:
% \[
%   +,\quad -,\quad +,\quad +.
% \]
% 因此, 函数 $f(x)$ 在点 $x_1=-2$ 取得极大值 $0$; 在点 $x_2=-\dfrac45$ 取得极小值
% \[
%   f\left(-\frac45\right)=\left(\frac65\right)^2\left(-\frac95\right)^3\approx -8.4;
% \]
% 而在点 $x_3=1$ 不取极值.
% 
% 我们也可以根据函数在一点的高阶导数的符号来判定极值点。我们有
% \end{solve}
% 
% 
% \begin{theorem}\label{theorem:ma-5-4-3}
% 设函数 $f(x)$ 在 $U(x_0,\delta)$ 内 $n$ 阶可导, $f'(x_0)=f''(x_0)=\cdots=f^{(n-1)}(x_0)=0$, 且 $f^{(n)}(x_0)\ne0$, 则
% 
% (1) 当 $n$ 为奇数时, $f(x)$ 在点 $x_0$ 不取极值;
% 
% (2) 当 $n$ 为偶数且 $f^{(n)}(x_0)>0$ 时, $f(x)$ 在 $x_0$ 取严格极小值;
% 
% (3) 当 $n$ 为偶数且 $f^{(n)}(x_0)<0$ 时, $f(x)$ 在 $x_0$ 取严格极大值.
% \end{theorem}
% 
% 
% \begin{proof}
% 由于函数 $f(x)$ 在 $U(x_0,\delta)$ 内 $n$ 阶可导, 而且
% \[
%   f'(x_0)=f''(x_0)=\cdots=f^{(n-1)}(x_0)=0,
% \]
% 利用带佩亚诺余项的泰勒公式, 可得
% \begin{equation}
% \begin{aligned}
%   f(x)-f(x_0)
%   &=\frac{f^{(n)}(x_0)}{n!}(x-x_0)^n+o((x-x_0)^n)\\
%   &=\left(\frac{f^{(n)}(x_0)}{n!}+o(1)\right)(x-x_0)^n\quad (x\to x_0).
% \end{aligned}
% \label{eq:ma-5-4-1}
% \end{equation}
% 而
% \[
%   \lim_{x\to x_0}\left(\frac{f^{(n)}(x_0)}{n!}+o(1)\right)
%   =\frac{f^{(n)}(x_0)}{n!},
% \]
% 因此存在 $\delta_1\in(0,\delta)$, 使得当 $x\in U(x_0,\delta_1)$ 时, 有 $\dfrac{f^{(n)}(x_0)}{n!}+o(1)$ 与 $f^{(n)}(x_0)$ 同号。于是, 我们有:
% 
% (1) 当 $n$ 为奇数时, 若 $x$ 从点 $x_0$ 的左侧经过 $x_0$ 而变到其右侧, 等式 \eqref{eq:ma-5-4-1} 的右边将改变符号, 从而 $f(x)-f(x_0)$ 也改变符号。这表明在 $x_0$ 的任何邻域内总有两点 $x_1$ 和 $x_2$, 使得
% \[
%   (f(x_0)-f(x_1))(f(x_0)-f(x_2))<0.
% \]
% 因此 $f(x)$ 在 $x_0$ 不取极值;
% 
% (2) 当 $n$ 为偶数且 $f^{(n)}(x_0)>0$ 时, 对 $\forall x\in U_0(x_0,\delta_1)$, \eqref{eq:ma-5-4-1} 式的右边值大于零, 从而有 $f(x)>f(x_0)$, 即 $f(x)$ 在 $x_0$ 取严格极小值; 而当 $n$ 为偶数且 $f^{(n)}(x_0)<0$ 时, $\forall x\in U_0(x_0,\delta_1)$, \eqref{eq:ma-5-4-1} 式的右边值小于零, 从而有 $f(x)<f(x_0)$, 即 $f(x)$ 在 $x_0$ 取严格极大值。定理
% \end{proof}
% 
% 
% \begin{remark}\label{remark:ma-5-source-113}
% 对于定理~\ref{theorem:ma-5-4-3}, 最常见的是 $n=2$ 的情形, 即若 $f'(x_0)=0$, 但 $f''(x_0)\ne0$, 则
% 
% (1) 当 $f''(x_0)>0$ 时, $f(x)$ 在 $x_0$ 取严格极小值;
% 
% (2) 当 $f''(x_0)<0$ 时, $f(x)$ 在 $x_0$ 取严格极大值.
% \end{remark}
% 
% 若对函数 $y=x^n\ (n\in\mathbb{N})$ 来讨论其在点 $x=0$ 处的极值情况, 则定理~\ref{theorem:ma-5-4-3} 的结论是显然的。事实上, 定理~\ref{theorem:ma-5-4-3} 告诉我们, 利用泰勒公式, 任何满足该定理条件的函数 $f(x)$ 在点 $x=x_0$ 处取极值情况和函数 $f^{(n)}(x_0)(x-x_0)^n$ 完全一样。换句话说, 读者可利用 $ax^n(a\ne0)$ 在 $x=0$ 的取值情况来记忆定理~\ref{theorem:ma-5-4-3}.
% 
% 如果我们利用定理~\ref{theorem:ma-5-4-3} 来解例~\ref{example:ma-5-4-3}, 则由 $f''(-2)=-54<0$ 立刻可以得出 $f(x)$ 在 $x_1=-2$ 取极大值; 由 $f''(1)=0$ 及 $f'''(1)=54>0$ 知 $f(x)$ 在 $x=1$ 不取极值。再注意到 $f''\left(-\dfrac45\right)>0$ 即可推出 $f(x)$ 在 $x_2=-\dfrac45$ 也取极小值.
% 
% \begin{example}[{达布 (Darboux) 定理}]\label{example:ma-5-4-4}
% 设函数 $f(x)$ 在 $[a,b]$ 上可导, 则对于任意介于 $f'(a+0)$ 与 $f'(b-0)$ 之间的数 $\eta$, 都存在 $\xi\in(a,b)$, 使得 $f'(\xi)=\eta$.
% \end{example}
% 
% 
% \begin{proof}
% 先假设 $f'(a+0)>0,\ f'(b-0)<0,\ \eta=0$。根据导数的定义, 由 $f'(a+0)>0$ 可导出, 存在 $\delta_1>0$, 使得 $\forall x\in[a,a+\delta_1]$, 有
% \[
%   f(x)>f(a);
% \]
% 再由 $f'(b-0)<0$ 可导出, 存在 $\delta_2>0$, 使得 $\forall x\in[b-\delta_2,b]$, 有
% \[
%   f(x)>f(b).
% \]
% 此外, 由于 $f(x)$ 在 $[a,b]$ 上可导, 从而 $f(x)$ 在 $[a,b]$ 上连续。现设
% \[
%   M=\max_{x\in[a,b]}\{f(x)\},
% \]
% 则存在 $\xi\in[a,b]$, 使得 $f(\xi)=M$。由以上分析可知 $\xi\ne a,b$, 因此费马定理告诉我们 $f'(\xi)=0$.
% 
% 同理可证 $f'(a+0)<0,\ f'(b-0)>0$ 和 $\eta=0$ 的情形.
% 
% 对一般情形, 我们构造辅助函数 $F(x)=f(x)-\eta x$。由 $F'(a+0)=f'(a+0)-\eta$ 和 $F'(b-0)=f'(b-0)-\eta$ 异号知, 存在 $\xi\in(a,b)$, 使得
% \[
%   0=F'(\xi)=f'(\xi)-\eta.
% \qedhere
% \]
% \end{proof}
% 
% 
% \begin{remark}\label{remark:ma-5-source-116}
% 值得注意的是, 一个函数即使在一个区间上可导, 其导函数也不一定是连续的 (参见下例)。而上述例子告诉我们, 即使导函数不连续, 仍然具有类似于连续函数的介值定理成立.
% \end{remark}
% 
% 
% \begin{example}\label{example:ma-5-4-5}
% 设函数
% \[
%   f(x)=
%   \begin{cases}
%     x^2\sin\dfrac1x, & x\ne0,\\
%     0, & x=0,
%   \end{cases}
% \]
% 则
% \[
%   f'(x)=
%   \begin{cases}
%     2x\sin\dfrac1x-\cos\dfrac1x, & x\ne0,\\
%     0, & x=0.
%   \end{cases}
% \]
% 显然, $f'(x)$ 在 $x=0$ 处不连续.
% \end{example}
% 
% 
% \paragraph*{思考题}
% 设函数 $f(x)$ 在区间 $I$ 上可导, 且 $f'(x)\ne0,\ \forall x\in I$。试证明 $f(x)$ 是 $I$ 上的严格单调函数.
% 
% 此外, 我们需要特别指出的是:
% 
% (1) 函数的极值只是函数的局部特性, 一个函数的某些极小值完全可以远远大于它的某些极大值;
% 
% (2) 函数的极值点必须位于函数所定义区间的内部, 区间的端点不涉及函数的极值问题.
% 
% 在实际应用中我们常常遇到的是求一个函数在指定范围内的最大值或者最小值问题。大家已经知道, 若函数 $f(x)$ 在闭区间 $[a,b]$ 上连续, 则 $f(x)$ 在 $[a,b]$ 上取得它的最大值和最小值。现在我们可以通过比较函数 $f(x)$ 在 $[a,b]$ 上所有的驻点、不可导点和端点的函数值的大小来确定它的最大值和最小值.
% 
% \begin{example}\label{example:ma-5-4-6}
% 求函数 $f(x)=\sin^3x+\cos^3x$ 在 $\left[-\dfrac\pi4,\dfrac{3\pi}4\right]$ 上的最大值和最小值.
% \end{example}
% 
% 
% \begin{solve}
% 直接计算, 得
% \[
%   f'(x)=3\sin x\cos x(\sin x-\cos x).
% \]
% 于是, 在所考虑区间内有函数的 $3$ 个驻点:
% \[
%   x_1=0,\qquad x_2=\frac\pi4,\qquad x_3=\frac\pi2.
% \]
% 易知
% \[
%   f\left(\frac\pi2\right)=f(0)=1,\quad
%   f\left(\frac\pi4\right)=\frac{\sqrt2}{2},\quad
%   f\left(-\frac\pi4\right)=f\left(\frac{3\pi}4\right)=0.
% \]
% 因此, 该函数的最大值和最小值分别为 $1$ 和 $0$.
% \end{solve}
% 
% 
% \begin{example}\label{example:ma-5-4-7}
% 作一个无盖圆柱形茶缸, 使得它的底的厚度是侧面厚度的三倍。试问: 容积一定时如何确定它的底面半径和高才能使得用料最省?
% \end{example}
% 
% 
% \begin{solve}
% 设茶缸的侧面的厚度为 $\delta$, 则它的底的厚度为 $3\delta$。记茶缸的容积为 $V=$ 常数, 底面半径为 $r$, 则高为 $h=\dfrac{V}{\pi r^2}$。于是, 若茶缸的用料量用 $\rho S(r)$ 表示, 其中 $\rho$ 为所用材料的密度, 则有
% \[
%   S(r)=\delta(3\pi r^2+2\pi rh)
%   =\delta\left(3\pi r^2+\frac{2V}{r}\right).
% \]
% 这样, 所考虑的问题就等价于求函数 $S(r)$ 的极小值.
% 
% 对 $S(r)$ 求导, 得
% \[
%   S'(r)=\delta\left(6\pi r-\frac{2V}{r^2}\right).
% \]
% 因此 $S'(r)$ 在 $[0,+\infty)$ 内有唯一的零点 $r_0=\sqrt[3]{\dfrac{V}{3\pi}}$。又
% \[
%   S''(r)=\delta\left(6\pi+\frac{4V}{r^3}\right)>0,\quad \forall r\in(0,+\infty),
% \]
% 所以 $r_0$ 是 $S(r)$ 的唯一极小值点。这时, 对应的高为
% \[
%   h_0=\frac{V}{\pi r_0^2}
%   =\frac{3\pi r_0^3}{\pi r_0^2}=3r_0.
% \]
% 这也就是说, 当高为底面半径的 $3$ 倍时用料最省.
% \end{solve}
% 
% 
% \subsection{函数的凹凸性}
% \label{subsec:ma-5-4-3}
% 
% 函数的凹凸性是函数一个重要的几何特征。设函数 $f(x)$ 在区间 $I$ 内有定义。从几何上来看, 若 $y=f(x)$ 的图像上任意两点 $(x_1,f(x_1))$ 和 $(x_2,f(x_2))$ 之间的曲线段总位于连接这两点的线段之下 (上), 则称该函数是凸 (凹) 的 (参见图 5.4.2)。这个概念用解析的语言可以表述成定义~\ref{definition:ma-5-4-1}.
% 
% \begin{definition}\label{definition:ma-5-4-1}
% 设函数 $f(x)$ 在区间 $I$ 内有定义。若 $\forall x_1,x_2\in I,\ \forall t\in(0,1)$, 有
% \[
%   f(tx_1+(1-t)x_2)\leqslant tf(x_1)+(1-t)f(x_2),
% \]
% 则称 $f(x)$ 为 $I$ 上的凸函数; 若当 $x_1\ne x_2$ 时, 总有严格不等式成立, 则称 $f(x)$ 为 $I$ 上的严格凸函数.
% \end{definition}
% 
% 
% \begin{remark}\label{remark:ma-5-source-123}
% 在上述定义中将 ``$\leqslant(<)$'' 改为 ``$\geqslant(>)$'', 就得到了凹 (严格凹) 函数的定义。所以有一个凸函数的定理, 就有一个相对应的凹函数的定理。下面我们只讨论凸函数.
% \end{remark}
% 
% 设 $f(x)$ 在 $\mathbb{R}$ 中有定义, 容易看出 $f(x)$ 既是凸函数又是凹函数的充分必要条件是 $f(x)$ 是一个线性函数; 若 $f(x)=ax^2+bx+c\ (a,b,c\in\mathbb{R})$, 则当 $a>0$ 时, $f(x)$ 是凸函数; 而当 $a<0$ 时, $f(x)$ 是凹函数.
% 
% 从定义不难证明: $f(x)(x\in I)$ 为凸函数的充分必要条件是对任意的 $x_1,x_2,x_3\in I$, 只要 $x_1<x_3<x_2$, 便有
% \begin{equation}
%   \frac{f(x_3)-f(x_1)}{x_3-x_1}
%   \leqslant
%   \frac{f(x_2)-f(x_3)}{x_2-x_3}.
% \label{eq:ma-5-4-2}
% \end{equation}
% 事实上, 取 $t=\dfrac{x_2-x_3}{x_2-x_1}$, 则 $x_3=tx_1+(1-t)x_2$, 再利用凸函数定义即得 \eqref{eq:ma-5-4-2} 式。反之, 若 \eqref{eq:ma-5-4-2} 成立, 可将其化简为
% \[
%   f(x_3)\leqslant
%   \frac{x_2-x_3}{x_2-x_1}f(x_1)
%   +\frac{x_3-x_1}{x_2-x_1}f(x_2).
% \]
% 令 $t=\dfrac{x_2-x_3}{x_2-x_1}$, 则 $t\in(0,1),\ x_1<x_3=tx_1+(1-t)x_2\leqslant x_2$, 则有
% \[
%   f(tx_1+(1-t)x_2)\leqslant tf(x_1)+(1-t)f(x_2),
% \]
% 所以 $f(x)$ 为凸函数.
% 
% 当 $f(x)$ 为严格凸函数时, 不等式 \eqref{eq:ma-5-4-2} 成立严格不等式。不等式 \eqref{eq:ma-5-4-2} 的几何解释是: 对一个凸函数而言, 过其图像上任意一点向两边作割线, 则左边割线的斜率总是小于或等于右边割线的斜率 (参见图 5.4.3).
% 
% 下面我们主要研究可导函数的凸凹性。我们有:
% 
% \begin{theorem}\label{theorem:ma-5-4-4}
% 设函数 $f(x)\in C[a,b]$ 且在 $(a,b)$ 上可导, 则 $f(x)$ 为凸函数的充分必要条件是 $f'(x)$ 在 $(a,b)$ 内单调上升; 而 $f(x)$ 为严格凸函数的充分必要条件是 $f'(x)$ 在 $(a,b)$ 内严格单调上升.
% \end{theorem}
% 
% 
% \begin{proof}
% 必要性\quad $\forall x_1,x_2\in(a,b),\ x_1<x_2$, 要证 $f'(x_1)\leqslant f'(x_2)$。令 $h>0$, 使得 $x_1-h,\ x_2+h\in(a,b)$, 则由凸函数的等价条件 \eqref{eq:ma-5-4-2}, 有
% \[
%   \frac{f(x_1)-f(x_1-h)}{h}
%   \leqslant
%   \frac{f(x_2)-f(x_1)}{x_2-x_1}
%   \leqslant
%   \frac{f(x_2+h)-f(x_2)}{h}.
% \]
% 令 $h\to0$, 得
% \[
%   f'(x_1)\leqslant\frac{f(x_2)-f(x_1)}{x_2-x_1}\leqslant f'(x_2),
% \]
% 即 $f'(x)$ 在 $(a,b)$ 内单调上升.
% 
% 若 $f(x)$ 为严格凸函数, 在 $(x_1,x_2)$ 中任取一点 $x^*$, 这时有
% \[
%   \frac{f(x_1)-f(x_1-h)}{h}
%   <
%   \frac{f(x^*)-f(x_1)}{x^*-x_1}
%   <
%   \frac{f(x_2)-f(x^*)}{x_2-x^*}
%   <
%   \frac{f(x_2+h)-f(x_2)}{h}.
% \]
% 令 $h\to0$, 得
% \[
%   f'(x_1)\leqslant
%   \frac{f(x^*)-f(x_1)}{x^*-x_1}
%   <
%   \frac{f(x_2)-f(x^*)}{x_2-x^*}
%   \leqslant f'(x_2),
% \]
% 即 $f'(x)$ 在 $(a,b)$ 内严格单调上升.
% 
% 充分性\quad 要证 $f(x)$ 为凸函数, 由等价条件 \eqref{eq:ma-5-4-2} 知, 只需证对任意的 $x_1,x_2,x_3\in[a,b]$, 当 $x_1<x_2<x_3$ 时, 有
% \[
%   \frac{f(x_2)-f(x_1)}{x_2-x_1}
%   \leqslant
%   \frac{f(x_3)-f(x_2)}{x_3-x_2}.
% \]
% 由拉格朗日微分中值定理, 可得
% \[
%   \frac{f(x_2)-f(x_1)}{x_2-x_1}=f'(\xi_1),\qquad
%   \frac{f(x_3)-f(x_2)}{x_3-x_2}=f'(\xi_2),
% \]
% 其中 $x_1<\xi_1<x_2<\xi_2<x_3$。由 $f'(x)$ 在 $(a,b)$ 内单调上升知, $f'(\xi_1)\leqslant f'(\xi_2)$, 从而
% \[
%   \frac{f(x_2)-f(x_1)}{x_2-x_1}
%   \leqslant
%   \frac{f(x_3)-f(x_2)}{x_3-x_2},
% \]
% 即 $f(x)$ 为凸函数。若 $f'(x)$ 严格单调上升, 则可导出上述不等式严格成立, 从而 $f(x)$ 为严格凸函数.
% \end{proof}
% 
% 
% \begin{theorem}\label{theorem:ma-5-4-5}
% 设函数 $f(x)\in C[a,b]$ 且在 $(a,b)$ 上二阶可导, 则 $f(x)$ 为凸函数的充分必要条件为 $f''(x)\geqslant0$; 而 $f(x)$ 为严格凸函数的充分必要条件为:
% 
% (1) $f''(x)\geqslant0$;
% 
% (2) $f''(x)$ 在 $(a,b)$ 的任一子区间上都不恒等于 $0$.
% \end{theorem}
% 
% 证明留作练习.
% 
% \begin{example}\label{example:ma-5-4-8}
% 设 $f(x)$ 是 $[a,b]$ 上的凸函数。证明: 对任意的
% \[
%   x_1,x_2,\cdots,x_n\in[a,b],\qquad
%   \sum_{i=1}^n t_i=1\quad (t_i>0),
% \]
% 有
% \[
%   f(t_1x_1+\cdots+t_nx_n)\leqslant t_1f(x_1)+\cdots+t_nf(x_n),
% \]
% 而且当 $f(x)$ 为严格凸函数且 $x_i$ 不全相等时, 上述不等式严格成立.
% \end{example}
% 
% 
% \begin{proof}
% 用数学归纳法。当 $n=2$ 时, 就是凸函数的定义。现假定该命题对 $n=k$ 成立, 下证 $n=k+1$ 也成立.
% 
% 对任意的 $x_i\in[a,b],\ t_i>0\ (i=1,2,\cdots,k+1),\ \sum_{i=1}^{k+1}t_i=1$, 令
% \[
%   \lambda_i=\frac{t_i}{1-t_{k+1}},\qquad i=1,2,\cdots,k,
% \]
% 则有 $\lambda_i>0\ (i=1,2,\cdots,k)$, 而且 $\sum_{i=1}^k\lambda_i=1$。于是, 由凸函数的定义和归纳法假设, 可得
% \[
% \begin{aligned}
%   &f(t_1x_1+\cdots+t_kx_k+t_{k+1}x_{k+1})\\
%   &\quad=f[(1-t_{k+1})(\lambda_1x_1+\cdots+\lambda_kx_k)+t_{k+1}x_{k+1}]\\
%   &\quad\leqslant(1-t_{k+1})f(\lambda_1x_1+\cdots+\lambda_kx_k)+t_{k+1}f(x_{k+1})\\
%   &\quad\leqslant(1-t_{k+1})(\lambda_1f(x_1)+\cdots+\lambda_kf(x_k))+t_{k+1}f(x_{k+1})\\
%   &\quad=t_1f(x_1)+\cdots+t_kf(x_k)+t_{k+1}f(x_{k+1}).
% \end{aligned}
% \]
% 当 $f(x)$ 为严格凸函数且 $x_i$ 不全相等时, 分两种情况: 一种情况是 $x_1,x_2,\cdots,x_k$ 不全相等, 此时由归纳法假设, 可知严格不等式成立; 另一种情况是 $x_1,x_2,\cdots,x_k$ 相等, 但不等于 $x_{k+1}$, 此时由 $\lambda_1x_1+\cdots+\lambda_kx_k\ne x_{k+1}$, 可导出严格不等式也成立。这样, 由归纳法原理知这一命题对一切正整数 $n$ 成立.
% \end{proof}
% 
% 
% \begin{example}\label{example:ma-5-4-9}
% 设 $a_i>0\ (i=1,2,\cdots,n)$ 不全相等, 证明: 当 $x\ne0$ 时, 有
% \[
%   x\frac{a_1^x\ln a_1+\cdots+a_n^x\ln a_n}{a_1^x+\cdots+a_n^x}
%   -\ln\frac{a_1^x+\cdots+a_n^x}{n}>0.
% \]
% \end{example}
% 
% 
% \begin{proof}
% 简单的运算可知, 所要证的不等式等价于
% \[
%   \frac1n a_1^x\ln a_1^x+\cdots+\frac1n a_n^x\ln a_n^x
%   >\frac{a_1^x+\cdots+a_n^x}{n}\ln\frac{a_1^x+\cdots+a_n^x}{n}.
% \]
% 现在在区间 $(0,+\infty)$ 上考虑函数 $f(u)=u\ln u$, 则有
% \[
%   f'(u)=\ln u+1,\qquad f''(u)=\frac1u>0\quad (u>0).
% \]
% 所以 $f(u)$ 是 $(0,+\infty)$ 上的严格凸函数。又 $a_i^x\ (i=1,2,\cdots,n)$ 不全相等, 应用上例的结论, 得
% \[
%   f\left(\frac{a_1^x+\cdots+a_n^x}{n}\right)
%   <\frac1n f(a_1^x)+\cdots+\frac1n f(a_n^x).
% \]
% 这正是所要证明的不等式.
% \end{proof}
% 
% 
% \begin{example}\label{example:ma-5-4-10}
% 设 $a_i>0\ (i=1,2,\cdots,n)$ 不全相等, 证明
% \[
%   f(x)=
%   \begin{cases}
%     \left(\dfrac{a_1^x+\cdots+a_n^x}{n}\right)^{\frac1x}, & x\ne0,\\[6pt]
%     \sqrt[n]{a_1a_2\cdots a_n}, & x=0,
%   \end{cases}
% \]
% 是 $(-\infty,+\infty)$ 上的严格单调递增函数.
% \end{example}
% 
% 
% \begin{proof}
% 由例~\ref{example:ma-5-2-8} 可知, $f(x)\in C(-\infty,+\infty)$。当 $x\ne0$ 时, 对 $\ln f(x)$ 求导, 得
% \[
%   \frac{f'(x)}{f(x)}
%   =\frac1{x^2}\left(
%   x\frac{a_1^x\ln a_1+\cdots+a_n^x\ln a_n}{a_1^x+\cdots+a_n^x}
%   -\ln\frac{a_1^x+\cdots+a_n^x}{n}\right).
% \]
% 因为 $f(x)>0$, 利用上例所证不等式可知, 当 $x\ne0$ 时, $f'(x)>0$。因此, $f(x)$ 在实轴上是严格单调递增的.
% \end{proof}
% 
% 
% \begin{remark}\label{remark:ma-5-source-133}
% 对上例的 $f(x)$, 我们有:
% \[
%   f(-1)=\frac{n}{\dfrac1{a_1}+\dfrac1{a_2}+\cdots+\dfrac1{a_n}},
% \]
% 称之为 $a_1,a_2,\cdots,a_n$ 的调和平均;
% \[
%   f(0)=\sqrt[n]{a_1a_2\cdots a_n},
% \]
% 称之为 $a_1,a_2,\cdots,a_n$ 的几何平均;
% \[
%   f(1)=\frac{a_1+a_2+\cdots+a_n}{n},
% \]
% 称之为 $a_1,a_2,\cdots,a_n$ 的算术平均。另外, 例~\ref{example:ma-5-2-8} 表明
% \[
%   f(-\infty)=\lim_{x\to-\infty}f(x)=\min\{a_1,a_2,\cdots,a_n\},
% \]
% \[
%   f(+\infty)=\lim_{x\to+\infty}f(x)=\max\{a_1,a_2,\cdots,a_n\}.
% \]
% 这样, 由 $f(x)$ 的严格单调递增性, 我们有
% \[
%   f(-\infty)<f(-1)<f(0)<f(1)<f(+\infty),
% \]
% 即
% \[
%   \min\{a_1,a_2,\cdots,a_n\}
%   <\frac{n}{\dfrac1{a_1}+\dfrac1{a_2}+\cdots+\dfrac1{a_n}}
%   <\sqrt[n]{a_1a_2\cdots a_n}
% \]
% \[
%   <\frac{a_1+a_2+\cdots+a_n}{n}
%   <\max\{a_1,a_2,\cdots,a_n\}.
% \]
% 这就是著名的调和--几何--算术平均不等式.
% \end{remark}
% 
% 
% \begin{example}\label{example:ma-5-4-11}
% 设 $a_i>0,\ b_i>0\ (i=1,2,\cdots,n)$, 证明:
% \[
%   \sum_{i=1}^n a_i b_i
%   \leqslant
%   \left(\sum_{i=1}^n a_i^p\right)^{\frac1p}
%   \left(\sum_{i=1}^n b_i^q\right)^{\frac1q},
% \]
% 其中 $0<p,q<+\infty,\ \dfrac1p+\dfrac1q=1$。此不等式称为赫尔德 (H\"older) 不等式, 当 $p=q=2$ 时, 又称为施瓦茨 (Schwarz) 不等式或柯西不等式, 它表明 $n$ 维空间中的两个向量夹角余弦之绝对值不超过 1.
% \end{example}
% 
% 
% \begin{proof}
% 令 $f(x)=x^{\frac1q}$, 则
% \[
%   f''(x)=\frac1q\left(\frac1q-1\right)x^{\frac1q-2}<0,\qquad \forall x>0.
% \]
% 因此, $f(x)$ 为 $(0,+\infty)$ 上的严格凹函数。于是, 若
% \[
%   x_i>0,\qquad t_i>0,\qquad \sum_{i=1}^n t_i=1,
% \]
% 则有
% \[
%   t_1x_1^{\frac1q}+t_2x_2^{\frac1q}+\cdots+t_nx_n^{\frac1q}
%   \leqslant (t_1x_1+t_2x_2+\cdots+t_nx_n)^{\frac1q}.
% \]
% 现取
% \[
%   t_i=\frac{a_i^p}{\sum_{i=1}^n a_i^p},\qquad x_i=\frac{b_i^q}{a_i^p},
% \]
% 并且代入上述不等式, 得
% \[
%   \frac{a_1b_1+\cdots+a_nb_n}{\sum_{i=1}^n a_i^p}
%   \leqslant
%   \frac{(b_1^q+\cdots+b_n^q)^{\frac1q}}{\left(\sum_{i=1}^n a_i^p\right)^{\frac1q}},
% \]
% 整理即得
% \[
%   \sum_{i=1}^n a_i b_i
%   \leqslant
%   \left(\sum_{i=1}^n a_i^p\right)^{\frac1p}
%   \left(\sum_{i=1}^n b_i^q\right)^{\frac1q}.
% \qedhere
% \]
% \end{proof}
% 
% 
% \subsection{拐点}
% \label{subsec:ma-5-4-4}
% 
% \begin{definition}\label{definition:ma-5-4-2}
% 设函数 $f(x)$ 在 $U(x_0,\delta)$ 上连续。如果存在 $\delta_0>0$, 使得 $f(x)$ 在 $U_0^-(x_0,\delta_0)$ 是凹 (凸) 的, 而 $U_0^+(x_0,\delta_0)$ 是凸 (凹) 的, 则称 $x_0$ 为 $f(x)$ 的一个拐点.
% \end{definition}
% 
% 
% \begin{theorem}\label{theorem:ma-5-4-6}
% 如果 $x_0$ 是 $f(x)$ 的拐点, 且 $f''(x_0)$ 存在, 则
% \[
%   f''(x_0)=0.
% \]
% \end{theorem}
% 
% 
% \begin{proof}
% $f''(x_0)$ 存在表明 $f'(x)$ 在点 $x_0$ 附近存在。$f(x)$ 在点 $x_0$ 的左右凹凸性相反, 表明 $f'(x)$ 在点 $x_0$ 的左右升降性相反, 即 $x_0$ 是 $f'(x)$ 一个极值点。由费马定理知, $f''(x_0)=0$.
% \end{proof}
% 
% 
% \begin{theorem}\label{theorem:ma-5-4-7}
% 如果函数 $f(x)$ 在 $U(x_0,\delta)$ 二阶可导, $f''(x_0)=0$, $f'''(x_0)$ 存在而且不为零, 则 $x_0$ 是 $f(x)$ 的拐点.
% \end{theorem}
% 
% 
% \begin{proof}
% $f''(x)$ 的带佩亚诺余项的泰勒公式为
% \[
%   f''(x)=f''(x_0)+f'''(x_0)(x-x_0)+o(x-x_0)
%   =(f'''(x_0)+o(1))(x-x_0).
% \]
% 所以 $\exists\delta_1<\delta$, 使得当 $x\in U(x_0,\delta_1)$ 时, $f'''(x_0)+o(1)$ 与 $f'''(x_0)$ 有相同符号, 从而 $f''(x)$ 在 $x_0$ 左右符号相反, 即 $f(x)$ 在 $x_0$ 左右凸凹性相反, 故 $x_0$ 是拐点.
% 
% 从拐点定义容易看出, 若函数 $f(x)$ 具有二阶导数, 则 $f(x)$ 的拐点即为 $f'(x)$ 的极值点。因此将前面关于 $f(x)$ 的极值点的结果应用于 $f'(x)$, 即得到了关于拐点的相应结果.
% \end{proof}
% 
% 
% \subsection{渐近线}
% \label{subsec:ma-5-4-5}
% 
% 设函数 $y=f(x)$ 的图像有一向无穷远无限伸展的分支, 如果当一个点沿着这一分支趋向无穷远时, 该点与某一条直线 $\ell$ 的距离趋向于零, 则称直线 $\ell$ 为 $f(x)$ 的一条渐近线.
% 
% 首先, 如果有 $\lim_{x\to x_0}f(x)=\infty$, 则 $x=x_0$ 就是 $f(x)$ 的一条渐近线。此时, 我们称它为 $f(x)$ 的一条垂直渐近线。当然也可以考虑 $x\to x_0+0$ 和 $x\to x_0-0$ 的情形.
% 
% 下面考虑非垂直的渐近线。设 $Y=kx+b$ 为 $y=f(x)$ 的一条渐近线 (参见图 5.4.4)。由解析几何的知识知, 函数 $y=f(x)$ 的图像上一点 $M=(x,f(x))$ 到直线 $Y=kx+b$ 的距离为
% \[
%   \delta=|y-Y|\cos\alpha=\frac{|f(x)-kx-b|}{\sqrt{1+k^2}}.
% \]
% 因此, 当这点沿着曲线 $y=f(x)$ 趋于无穷时, 它与直线 $Y=kx+b$ 的距离 $\delta$ 趋于零的充要条件是
% \[
%   \lim_{x\to+\infty}[f(x)-(kx+b)]=0.
% \]
% 于是
% \[
%   \lim_{x\to+\infty}\frac{f(x)}x
%   =\lim_{x\to+\infty}\left(\frac{f(x)-kx-b}{x}+\frac{kx+b}{x}\right)=k.
% \]
% 这样, 所求渐近线的斜率 $k$ 和截距 $b$ 为
% \[
%   k=\lim_{x\to+\infty}\frac{f(x)}x,\qquad
%   b=\lim_{x\to+\infty}(f(x)-kx).
% \]
% 应当注意, 非垂直渐近线包括了斜渐近线和水平渐近线。当由上面的式子确定的 $k=0$ 时, 则有
% \[
%   b=\lim_{x\to+\infty}f(x).
% \]
% 此时, $y=b$ 就是一条水平渐近线.
% 
% 同理可以考虑 $x\to-\infty$ 和 $x\to\infty$ 的情形.
% 
% \subsection{函数的作图}
% \label{subsec:ma-5-4-6}
% 
% 解决实际问题时, 遇到的总是具体函数, 作出函数的图像, 对我们分析问题和解决问题会有很大的帮助.
% 
% 最基本的作图方法是描点画图法, 计算机就是利用这个方法来作图的。一般来说, 我们先把 $[a,b]$ 分得充分细, 在每个分点 $x_i$ 上计算 $f(x_i)$; 然后描出点 $(x_i,f(x_i))$, 并且用光滑的曲线将各点连起来, 便得到了 $y=f(x)$ 的图像。本书中的图全部是这样做出来的.
% 
% 手工作图不能这样, 因为计算量太大。但利用导数作为工具, 先把函数各种性质尽可能地搞清楚, 就可很快作出函数的大致图像。其具体步骤如下:
% 
% (1) 求出函数的定义域;
% 
% (2) 研究函数的有界性、奇偶性和周期性;
% 
% (3) 解方程 $f'(x)=0$, 列表确定函数升降区间和极值点;
% 
% (4) 解方程 $f''(x)=0$, 列表确定函数的凸凹区间和拐点;
% 
% (5) 求出函数的斜渐近线与垂直渐近线;
% 
% (6) 计算一些重要点上 (如点 $x=0$) 的函数值;
% 
% (7) 根据以上数据及函数的变化趋势描出其草图.
% 
% \begin{example}\label{example:ma-5-4-12}
% 作出函数 $y=e^{-x^2}$ 的草图.
% \end{example}
% 
% 
% \begin{solve}
% 显然, 函数 $y=e^{-x^2}$ 定义域为 $(-\infty,+\infty)$, 且为偶函数。因此它的图像关于 $y$ 轴对称。当 $x=0$ 时, $y=1$, 并且总有 $y>0$, 可见它的图像仅与 $y$ 轴相交而与 $x$ 轴不相交。此外, 由于
% \[
%   \lim_{x\to\pm\infty}e^{-x^2}=0,
% \]
% 所以 $x$ 轴为它的一条水平渐近线。求此函数的导数并且解方程 $y'=-2xe^{-x^2}=0$, 知其有唯一的驻点 $x=0$。再求此函数的二阶导数并且解方程
% \[
%   y''=4e^{-x^2}\left(x^2-\frac12\right)=0,
% \]
% 得 $x=\pm1/\sqrt2$。关于这一函数的升降性、极值、凹凸性及拐点的情况, 如表 5.4.1 所示.
% 
% 根据上面的分析, 大致描出 $y=e^{-x^2}$ 的图像如图 5.4.5 所示.
% \end{solve}
% 
% 
% \begin{example}\label{example:ma-5-4-13}
% 作出函数
% \[
%   f(x)=\frac{(x-1)^3}{(x+1)^2}
% \]
% 的草图.
% \end{example}
% 
% 
% \begin{solve}
% 除 $x=-1$ 外, $f(x)$ 在实轴上都有意义。显然有 $\lim_{x\to-1}f(x)=-\infty$, 因此 $x=-1$ 是该函数的一条垂直渐近线。又
% \[
%   \lim_{x\to\infty}\frac{f(x)}x=1,\qquad
%   \lim_{x\to\infty}(f(x)-x)=-5,
% \]
% 所以 $y=x-5$ 是它的一条斜渐近线。此外, 直接计算有 $f(1)=0$ 及
% \[
%   f'(x)=\frac{(x-1)^2(x+5)}{(x+1)^3},\qquad
%   f'(1)=0,\qquad f'(-5)=0,
% \]
% \[
%   f''(x)=\frac{24(x-1)}{(x+1)^4},\qquad f''(1)=0.
% \]
% 根据上面的结果, 可得这一函数的升降性、极值、凹凸性及拐点的情况, 如表 5.4.2 所示.
% 
% 根据上面的分析, 大致描出该函数的图像如图 5.4.6 所示.
% \end{solve}
% 
% 
% \section*{习题五}
% \phantomsection
% \addcontentsline{toc}{section}{习题五}
% \label{exercises:ma-5}
% 
% 1。证明广义罗尔微分中值定理: 设函数 $f(x)$ 在 $(a,b)$ 上可导, $f(a+0)=f(b-0)=A$, 则存在 $\xi\in(a,b)$, 使得 $f'(\xi)=0$, 其中 $a$ 可为 $-\infty$, $b$ 可为 $+\infty$, $A$ 可为 $+\infty$ 或 $-\infty$.
% 
% 2。证明: 函数 $f(x)=[(x-a)^n(x-b)^n]^{(n)}$ 在 $(a,b)$ 内有 $n$ 个互不相同的零点.
% 
% 3。证明: 若函数 $f(x)=x^m(1-x)^n$, 其中 $m,n$ 为正整数, 则存在 $\xi\in(0,1)$, 使得
% \[
%   \frac mn=\frac{\xi}{1-\xi}.
% \]
% 
% 4。求证: $4ax^3+3bx^2+2cx=a+b+c$ 在 $(0,1)$ 上至少有一个根.
% 
% 5。求证: 切比雪夫--拉盖尔 (Chebyshev-Laguerre) 多项式
% \[
%   L_n(x)=e^x\frac{\mathrm d^n}{\mathrm dx^n}(x^ne^{-x})
% \]
% 有 $n$ 个不同的零点.
% 
% 6。求证: 切比雪夫--埃尔米特 (Chebyshev-Hermite) 多项式
% \[
%   H_n(x)=(-1)^n\frac1{n!}e^{\frac{x^2}2}\frac{\mathrm d^n}{\mathrm dx^n}\left(e^{-\frac{x^2}2}\right)
% \]
% 有 $n$ 个不同的零点.
% 
% 7。证明:
% 
% (1) 若函数 $f(x)$ 在 $(a,b)$ 上的导数有界, 则 $f(x)$ 在 $(a,b)$ 上一致连续;
% 
% (2) 对任意的两个实数 $x_1,x_2$, 有
% \[
%   |\sin x_1-\sin x_2|\leqslant |x_1-x_2|;
% \]
% 
% (3) 对任意的两个实数 $x_1,x_2$, 有
% \[
%   |\arctan x_1-\arctan x_2|\leqslant |x_1-x_2|.
% \]
% 
% 8。设 $\lim_{x\to+\infty}f'(x)=a$, 求证: 对任意的 $T>0$, 有
% \[
%   \lim_{x\to+\infty}(f(x+T)-f(x))=Ta.
% \]
% 
% 9。证明下列不等式:
% 
% (1) $x-\dfrac{x^3}6<\sin x<x\quad (x>0)$;
% 
% (2) $x-\dfrac{x^2}2<\ln(1+x)<x\quad (x>0)$;
% 
% (3) $2\sqrt x>3-\dfrac1x\quad (x>1)$;
% 
% (4) $e^x-1>x+\dfrac{x^2}2\quad (x>0)$;
% 
% (5) $\ln(1+x)>\dfrac{\arctan x}{1+x}\quad (x>0)$.
% 
% 10。设函数 $f(x)$ 和 $g(x)$ 在 $(-\infty,+\infty)$ 上可导, 且有
% \[
%   f'(x)>g'(x),\qquad f(a)=g(a).
% \]
% 证明: 当 $x>a$ 时, 有 $f(x)>g(x)$; 而当 $x<a$ 时, 有 $f(x)<g(x)$.
% 
% 11。设函数 $f(x)$ 在区间 $[a,b]\ (0<a<b)$ 上连续, 在 $(a,b)$ 内可导。求证: 存在 $\xi_1,\xi_2,\xi_3\in(a,b)$, 使得
% \[
%   f'(\xi_1)=(b+a)\frac{f'(\xi_2)}{2\xi_2}
%   =(b^2+ab+a^2)\frac{f'(\xi_3)}{3\xi_3^2}.
% \]
% 
% 12。设函数 $f(x)$ 在 $(a,b)$ 上可导, 且 $f'(x)$ 单调。证明 $f'(x)$ 在 $(a,b)$ 上连续.
% 
% 13。设函数 $f(x)$ 在点 $a$ 二阶可导, 求证:
% \[
%   \lim_{h\to0}\frac{f(a+h)+f(a-h)-2f(a)}{h^2}=f''(a).
% \]
% 
% 14。设函数 $f(x)$ 在 $[a,b]$ 上二阶可导, $f(a)=f(b)=0$, 且存在 $c\in(a,b)$, 使得 $f(c)>0$。求证: 存在 $\xi\in(a,b)$, 使得 $f''(\xi)<0$.
% 
% 15。证明: 设 $x\geqslant0$, 则
% 
% (1)
% \[
%   \sqrt{x+1}-\sqrt x=\frac1{2\sqrt{x+\theta(x)}},\qquad
%   \frac14\leqslant\theta(x)\leqslant\frac12;
% \]
% 
% (2)
% \[
%   \lim_{x\to+0}\theta(x)=\frac14,\qquad
%   \lim_{x\to+\infty}\theta(x)=\frac12.
% \]
% 
% 16。设函数 $f(x)$ 在 $[a,b]$ 上三阶可导, 证明: 存在一点 $\xi\in(a,b)$, 使得
% \[
%   f(b)=f(a)+\frac12(b-a)[f'(a)+f'(b)]-\frac1{12}(b-a)^3f^{(3)}(\xi).
% \]
% 
% 17。设函数
% \[
%   f(x)=
%   \begin{cases}
%     x^2\sin\dfrac1x, & x\ne0,\\[4pt]
%     0, & x=0,
%   \end{cases}
% \]
% 在 $[0,x]$ 上应用中值定理, 得
% \[
%   x^2\sin\frac1x=x\left(2\xi\sin\frac1\xi-\cos\frac1\xi\right),\qquad 0<\xi<x.
% \]
% 由此得
% \[
%   \cos\frac1\xi=2\xi\sin\frac1\xi-x\sin\frac1x.
% \]
% 当 $x\to0$ 时, 有 $\xi\to0$, 所以由上式得到
% \[
%   \lim_{\xi\to0}\cos\frac1\xi=0.
% \]
% 但已知 $\lim_{\xi\to0}\cos\dfrac1\xi$ 不存在。试说明矛盾何在.
% 
% 18。设函数 $f(x)$ 在 $(0,+\infty)$ 上可导, 且 $\lim_{x\to+\infty}f'(x)=+\infty$。求证 $f(x)$ 在 $(0,+\infty)$ 上不一致连续.
% 
% 19。证明函数 $f(x)=x\ln x$ 在 $(0,+\infty)$ 上不一致连续.
% 
% 20。设函数 $f(x),g(x),h(x)$ 都在 $[a,b]$ 上连续, 在 $(a,b)$ 内可导, 并定义
% \[
%   F(x)=
%   \begin{vmatrix}
%     f(a) & g(a) & h(a)\\
%     f(b) & g(b) & h(b)\\
%     f(x) & g(x) & h(x)
%   \end{vmatrix}.
% \]
% 证明: 存在 $\xi\in(a,b)$, 使得 $F'(\xi)=0$。问若 $h(x)\equiv1$, 则上述结论与柯西中值公式有什么关系? 若令 $g(x)=x,\ h(x)\equiv1$, 则会有什么结论?
% 
% 21。设函数 $f(x)$ 在 $[a,b]$ 上连续, 在 $(a,b)$ 内可导。求证: 存在 $\xi\in(a,b)$, 使得
% \[
%   2\xi(f(b)-f(a))=(b^2-a^2)f'(\xi).
% \]
% 
% 22。设函数 $f(x)$ 在 $[a,b]$ 上连续, 在 $(a,b)$ 内可导。求证: 存在 $\xi\in(a,b)$, 使得
% \[
%   f(b)-f(a)=\xi\ln\frac ba f'(\xi).
% \]
% 
% 23。设函数 $f(x)$ 在 $[a,b]$ 上连续, 在 $(a,b)$ 内可导, 且 $f(a)=f(b)=1$。求证: 存在 $\xi,\eta\in(a,b)$, 使得
% \[
%   e^{\eta-\xi}[f(\eta)+f'(\eta)]=1.
% \]
% 
% 24。证明下列恒等式:
% 
% (1)
% \[
%   2\arctan x+\arcsin\frac{2x}{1+x^2}=\pi\operatorname{sgn}x\quad (|x|\geqslant1);
% \]
% 
% (2)
% \[
%   3\arccos x-\arccos(3x-4x^3)=\pi\quad (|x|\leqslant1/2).
% \]
% 
% 25。求下列极限:
% 
% (1) $\displaystyle\lim_{x\to1}\frac{\ln\cos(x-1)}{1-\sin\frac{\pi x}2}$;\qquad
% (2) $\displaystyle\lim_{x\to+\infty}(\pi-2\arctan x)\ln x$;
% 
% (3) $\displaystyle\lim_{x\to1}\left(\frac1{\ln x}-\frac1{x-1}\right)$;\qquad
% (4) $\displaystyle\lim_{x\to1}\frac{x-1}{\ln x}$;
% 
% (5) $\displaystyle\lim_{x\to+0}x^{\sin x}$;\qquad
% (6) $\displaystyle\lim_{x\to0}\left(\frac{(\ln(1+x))^{1+x}}{x^2}-\frac1x\right)$;
% 
% (7) $\displaystyle\lim_{x\to0}\left(\cot x-\frac1x\right)$;\qquad
% (8) $\displaystyle\lim_{x\to0}\frac{(1+x)^{\frac1x}-e}{x}$;
% 
% (9) $\displaystyle\lim_{x\to+\infty}\left(\frac\pi2-\arctan x\right)^{\frac1{\ln x}}$;\qquad
% (10) $\displaystyle\lim_{x\to+0}\sin x\ln x$;
% 
% (11) $\displaystyle\lim_{x\to0}(1-\cos x)^{1-\cos x}$;\qquad
% (12) $\displaystyle\lim_{x\to0}\frac{\tan x-\sin x}{\sin x-x\cos x}$;
% 
% (13) $\displaystyle\lim_{x\to1}x^{\frac1{1-x}}$;\qquad
% (14) $\displaystyle\lim_{x\to0}\frac{xe^x-\ln(1+x)}{x^2}$;
% 
% (15) $\displaystyle\lim_{x\to0}\frac{x^2\sin\dfrac1x}{\sin x}$;\qquad
% (16) $\displaystyle\lim_{x\to+0}\left(\frac{1+x^a}{1+x^b}\right)^{\frac1{\ln x}}\quad (a,b\text{ 为正实数})$;
% 
% (17) $\displaystyle\lim_{x\to\frac\pi4-0}(\tan x)^{\tan 2x}$.
% 
% 26。设函数 $f(x)$ 在 $[a,+\infty)$ 上有界, $f'(x)$ 存在且 $\lim_{x\to+\infty}f'(x)=b$。求证: $b=0$.
% 
% 27。由拉格朗日微分中值定理, 有
% \[
%   \ln(1+x)=\frac{x}{1+\theta x}\qquad (0<\theta<1).
% \]
% 求证 $\lim_{x\to0}\theta=\dfrac12$.
% 
% 28。由拉格朗日微分中值定理, 有
% \[
%   \arcsin x=\frac{x}{\sqrt{1-\theta^2x^2}}\qquad (0<\theta<1).
% \]
% 求证 $\lim_{x\to0}\theta=\dfrac1{\sqrt3}$.
% 
% 29。设函数 $f(x)$ 在 $(a,+\infty)$ 上可导, 且 $\lim_{x\to+\infty}[f(x)+f'(x)]=k$ ($k$ 有限或 $\pm\infty$)。求证: $\lim_{x\to+\infty}f(x)=k$.
% 
% 30。设函数 $f(x)$ 在实轴上任意阶可导, 令 $F(x)=f(x^2)$, 求证:
% \[
%   F^{(2n+1)}(0)=0,\qquad
%   \frac{F^{(2n)}(0)}{(2n)!}=\frac{f^{(n)}(0)}{n!}.
% \]
% 
% 31。写出下列函数在 $x=0$ 的带佩亚诺余项的泰勒公式:
% 
% (1) $\cos x^2$;\qquad
% (2) $\sin^3x$;
% 
% (3) $\dfrac1{(1+x)^2}$;\qquad
% (4) $\dfrac{x^3+2x+1}{x-1}$.
% 
% 32。写出下列函数在 $x=0$ 的带有佩亚诺余项的泰勒公式, 要求至所指定的阶数.
% 
% (1) $\dfrac{x}{\sin x}\quad (x^4)$;\qquad
% (2) $\ln(\cos x+\sin x)\quad (x^4)$;
% 
% (3) $\dfrac{x}{2x^2+x-1}\quad (x^3)$;\qquad
% (4) $\ln\dfrac{1+x}{1-2x}\quad (x^n)$;
% 
% (5) $\ln(1+x+x^2+x^3)\quad (x^6)$;\qquad
% (6) $\dfrac{1+x+x^2}{1-x+x^2}\quad (x^4)$.
% 
% 33。确定常数 $a,b$, 使当 $x\to0$ 时,
% 
% (1) $f(x)=(a+b\cos x)\sin x-x$ 为 $x$ 的 5 阶无穷小量;
% 
% (2) $f(x)=e^x-\dfrac{1+ax}{1-bx}$ 是 $x$ 的 3 阶无穷小量.
% 
% 34。求下列极限:
% 
% (1) $\displaystyle\lim_{x\to0}\left(\frac1x-\frac1{\sin x}\right)$;\qquad
% (2) $\displaystyle\lim_{x\to0}\frac{e^{x^3}-1-x^3}{\sin^6 2x}$;
% 
% (3) $\displaystyle\lim_{x\to+\infty}\left(\sqrt[3]{x^3-3x}-\sqrt{x^2-2x}\right)$;\qquad
% (4) $\displaystyle\lim_{x\to\infty}\left(x+\frac12\right)\ln\left(1+\frac1x\right)$;
% 
% (5) $\displaystyle\lim_{x\to0}\left(\frac{\tan x}{x}\right)^{\frac1{x^2}}$;\qquad
% (6) $\displaystyle\lim_{x\to\infty}x^2\ln\left(x\sin\frac1x\right)$.
% 
% 35。设函数 $f(x)$ 在 $x=0$ 的某个邻域内二阶可导, 且
% \[
%   \lim_{x\to0}\left(\frac{\sin3x}{x^3}+\frac{f(x)}{x^2}\right)=0.
% \]
% 
% (1) 求 $f(0),f'(0),f''(0)$;
% 
% (2) 求 $\displaystyle\lim_{x\to0}\left(\frac3{x^2}+\frac{f(x)}{x^2}\right)$.
% 
% 36。设函数 $f(x)$ 在 $x=0$ 的某个邻域内二阶可导, 且
% \[
%   \lim_{x\to0}\left(1+x+\frac{f(x)}x\right)^{\frac1x}=e^3.
% \]
% 
% (1) 求 $f(0),f'(0),f''(0)$;
% 
% (2) 求 $\displaystyle\lim_{x\to0}\left(1+\frac{f(x)}x\right)^{\frac1x}$.
% 
% 37。设函数 $f(x)$ 在 $(a,+\infty)$ 上有直到 $n$ 阶导数, 且有
% \[
%   \lim_{x\to+\infty}f(x)=A,\qquad
%   \lim_{x\to+\infty}f^{(n)}(x)=B.
% \]
% 求证 $B=0$.
% 
% 38。设函数 $f(x)$ 在 $x=a$ 的某个邻域内二阶导数连续, 且 $f''(a)\ne0$, 由拉格朗日微分中值定理有
% \[
%   f(a+h)-f(a)=f'(a+\theta h)h\qquad (0<\theta<1).
% \]
% 求证:
% \[
%   \lim_{h\to0}\theta=\frac12.
% \]
% 
% 39。设 $P(x)$ 为 $n$ 次多项式, 证明:
% 
% (1) 若 $P(a),P'(a),\cdots,P^{(n)}(a)$ 皆为正数, 则 $P(x)$ 在 $(a,+\infty)$ 上无零点;
% 
% (2) 若 $P(a),P'(a),\cdots,P^{(n)}(a)$ 正负号相间, 则 $P(x)$ 在 $(-\infty,a)$ 上无零点;
% 
% 40。设 $P(x)$ 为 $n$ 次多项式, 证明 $x=a$ 是 $P(x)=0$ 的 $k$ 重根的充分必要条件是
% \[
%   P^{(j)}(a)=0\ (j=0,1,\cdots,k-1),\qquad P^{(k)}(a)\ne0.
% \]
% 
% 41。设
% \[
%   P_n(x)=\frac1{2^n n!}\frac{\mathrm d^n}{\mathrm dx^n}(x^2-1)^n,
% \]
% 求证:
% 
% (1) $P_n(1)=1$;
% 
% (2) $P_n(-1)=(-1)^n$;
% 
% (3) $P_{2n-1}(0)=0$;
% 
% (4)
% \[
%   P_{2n}(0)=\frac{(-1)^n(2n)!}{2^{2n}(n!)^2}.
% \]
% 
% 42。设函数 $f(x)=e^{x^2}$.
% 
% (1) 求证: $f^{(n)}(x)=P_n(x)e^{x^2}$, 其中 $P_n(x)$ 为 $n$ 次多项式, 且满足 $P_0(x)=1,\ P_1(x)=2x,\ P_{n+1}(x)=2xP_n(x)+2nP_{n-1}(x)$;
% 
% (2) 求 $f^{(n)}(0)$ 的值.
% 
% 43。设函数 $f(x)$ 在 $(0,+\infty)$ 上三阶可导, 而且有
% \[
%   |f(x)|\leqslant M_0,\qquad |f'''(x)|\leqslant M_3,\qquad \forall x\in(0,+\infty).
% \]
% 求证 $f'(x)$ 和 $f''(x)$ 在 $(0,+\infty)$ 上有界.
% 
% 44。求证:
% 
% (1) $0<x-\ln(1+x)<\dfrac12x^2\quad (0<x\leqslant1)$;
% 
% (2) $\displaystyle\lim_{n\to\infty}\sum_{k=1}^n\left[\frac1k-\ln\left(1+\frac1k\right)\right]$ 存在.
% 
% 45。设函数 $f(x)$ 在 $[a,b]$ 上有二阶导数, 且 $f'(a)=f'(b)=0$。证明: 存在 $c\in(a,b)$, 使得
% \[
%   |f''(c)|\geqslant\frac4{(b-a)^2}|f(b)-f(a)|.
% \]
% 
% 46。若函数 $f(x)$ 在 $[-1,1]$ 上三阶可导, 且有
% \[
%   f(0)=f'(0)=0,\qquad f(1)=1,\qquad f(-1)=0.
% \]
% 求证: 存在 $\xi\in(-1,1)$, 使得 $f'''(\xi)\geqslant3$.
% 
% 47。设函数 $f(x)$ 在 $(-\infty,+\infty)$ 上二阶可导, 且有
% \[
%   |f(x)|\leqslant M_0,\qquad |f''(x)|\leqslant M_2,\qquad \forall x\in(-\infty,+\infty).
% \]
% 
% (1) 写出 $f(x+h)$ 和 $f(x-h)$ 的泰勒公式;
% 
% (2) 求证: $|f'(x)|\leqslant\dfrac{M_0}h+\dfrac h2M_2,\ \forall h>0$;
% 
% (3) 求 $\dfrac{M_0}h+\dfrac h2M_2$ 在 $(0,+\infty)$ 上的最小值;
% 
% % 整合来源：math-6.tex
% \chapter{不定积分}
% \label{chap:ma-6}
% 
% \section{原函数与不定积分}
% \label{sec:ma-6-1}
% 
% \subsection{原函数与不定积分的概念}
% \label{subsec:ma-6-1-1}
% 
% 大家已经知道, 寻求一个已知函数的导数是一个十分有意义的问题, 这是因为各种各样的物理问题和几何问题的解决, 常常归结为求某一函数的导数。但是, 我们也必须看到另一方面, 那就是很多的实际问题的解决, 不是归结为求某一已知函数的导数, 而恰恰相反, 是要求一个未知函数, 使之以某一已知函数为其导数。例如, 考虑质点沿直线运动, 已知质点所走过的路程 $s=s(t)$, 如果要求瞬时速度, 则我们有 $v(t)=s'(t)$。反过来, 如果知道每个时刻的瞬时速度 $v(t)$, 要求质点所走过的路程 $s(t)$, 则是求导数的逆运算。此时我们要找一个函数 $s(t)$, 使得 $s'(t)=v(t)$。这个 $s(t)$ 称为 $v(t)$ 的一个原函数。原函数的定义如下:
% 
% \begin{definition}\label{definition:ma-6-1-1}
% 在区间 $I$ 上给定函数 $f(x)$, 若存在定义在 $I$ 上的函数 $F(x)$, 使得
% \[
%   F'(x)=f(x),\quad \forall x\in I
%   \quad\text{或}\quad
%   \mathrm dF(x)=f(x)\mathrm dx,\quad \forall x\in I,
% \]
% 则称 $F(x)$ 是 $f(x)$ 的一个原函数.
% \end{definition}
% 
% 由定义不难看出, 若 $f(x)$ 有原函数, 则它就有无穷多个原函数。这是因为若 $F(x)$ 是 $f(x)$ 的一个原函数, 则对任意的常数 $C$, 函数 $F(x)+C$ 也是 $f(x)$ 的一个原函数。几何上看这是明显的, 曲线 $F(x)$ 和 $F(x)+C$ 在点 $x$ 有相同切线斜率 (参见图 6.1.1)。另一方面, 由拉格朗日微分中值定理的推论~\ref{corollary:ma-local-5.1-2} 知, 若 $F(x)$ 和 $G(x)$ 是 $f(x)$ 的两个原函数, 则它们之间只能是相差一个常数, 即 $G(x)=F(x)+C$。原函数的这一特点, 也可以从实际问题中反映出来。例如, 两个粒子如果以相同的速度 $v(t)$ 运动, 则它们在任何相同时间走过的路程是相同的。因此, 它们的路程函数只差一个常数, 这是由它们的初始位置的不同而引起的.
% 
% 这样, 设 $f(x)$ 在区间 $I$ 上有定义, 函数 $F(x)$ 是 $f(x)$ 的一个原函数, 则函数族 $F(x)+C$ ($C$ 为任意常数) 包含了 $f(x)$ 的所有原函数.
% 
% \begin{definition}\label{definition:ma-6-1-2}
% 若函数 $f(x)$ 在区间 $I$ 上存在原函数, 则称 $f(x)$ 的全体原函数为 $f(x)$ 的不定积分, 记为
% \[
%   \int f(x)\mathrm dx,
% \]
% 这里符号 ``$\int$'' 称为积分号, $f(x)$ 称为被积函数, $x$ 称为积分变量.
% \end{definition}
% 
% 根据上述定义以及上面的分析我们有, 若 $F(x)$ 是 $f(x)$ 的一个原函数, 则
% \[
%   \int f(x)\mathrm dx=F(x)+C,
% \]
% 其中 $C$ 为任意常数。另外我们有
% \[
%   \left(\int f(x)\mathrm dx\right)'=f(x),
%   \qquad
%   \mathrm d\left(\int f(x)\mathrm dx\right)=f(x)\mathrm dx
% \]
% 和
% \[
%   \int F'(x)\mathrm dx=F(x)+C,
%   \qquad
%   \int \mathrm dF(x)=F(x)+C,
% \]
% 其中 $C$ 为任意常数。因此, 在允许相差一个常数的意义下, 求不定积分的这一运算恰好是求导或求微分的逆运算.
% 
% \subsection{基本不定积分表和不定积分的线性性质}
% \label{subsec:ma-6-1-2}
% 
% 由于微分和积分互为逆运算, 因此对应于微分的每一种方法相应地就有一种积分的方法。首先, 从求导数的公式中我们可以导出如下不定积分公式, 它是我们计算不定积分的基础, 务必牢记.
% \[
% \begin{aligned}
%   &(1)\quad \int x^\alpha\mathrm dx=\frac1{\alpha+1}x^{\alpha+1}+C
%   \quad(\alpha\ne-1);\\
%   &(2)\quad \int\frac{\mathrm dx}{x}=\ln|x|+C;\\
%   &(3)\quad \int e^x\mathrm dx=e^x+C;\\
%   &(4)\quad \int\cos x\mathrm dx=\sin x+C;\\
%   &(5)\quad \int\sin x\mathrm dx=-\cos x+C;\\
%   &(6)\quad \int\frac{\mathrm dx}{\cos^2x}=\tan x+C;\\
%   &(7)\quad \int\frac{\mathrm dx}{1+x^2}=\arctan x+C;\\
%   &(8)\quad \int\frac{\mathrm dx}{\sqrt{1-x^2}}=\arcsin x+C;\\
%   &(9)\quad \int\sinh x\mathrm dx=\cosh x+C;\\
%   &(10)\quad \int\cosh x\mathrm dx=\sinh x+C;\\
%   &(11)\quad \int\frac{\mathrm dx}{\cosh^2x}=\tanh x+C.
% \end{aligned}
% \]
% 其次, 对应于求导法则
% \[
%   (f(x)\pm g(x))'=f'(x)\pm g'(x),
% \]
% 就有如下的求不定积分性质:
% 
% \begin{property}\label{property:ma-6-1-1}
% 设函数 $f(x),g(x)$ 存在原函数, 则 $f(x)\pm g(x)$ 也存在原函数, 而且有
% \[
%   \int[f(x)\pm g(x)]\mathrm dx
%   =
%   \int f(x)\mathrm dx\pm\int g(x)\mathrm dx.
% \]
% \end{property}
% 
% 
% \begin{proof}
% 令
% \[
%   \int f(x)\mathrm dx=F(x)+C
%   \quad\text{或}\quad
%   F'(x)=f(x),
% \]
% \[
%   \int g(x)\mathrm dx=G(x)+C
%   \quad\text{或}\quad
%   G'(x)=g(x),
% \]
% 则
% \[
%   [F(x)\pm G(x)]'=f(x)\pm g(x),
% \]
% 所以
% \[
%   \int[f(x)\pm g(x)]\mathrm dx=F(x)\pm G(x)+C.
% \]
% 
% 对应于求导法则 $(kf(x))'=kf'(x)$, 其中 $k$ 是常数, 就有如下的求不定积分性质:
% \end{proof}
% 
% 
% \begin{property}\label{property:ma-6-1-2}
% 设函数 $f(x)$ 存在原函数, 则 $kf(x)$ 存在原函数, 且
% \[
%   \int kf(x)\mathrm dx=k\int f(x)\mathrm dx,
% \]
% 其中 $k$ 是常数, 且 $k\ne0$.
% \end{property}
% 
% 
% 证明留作练习.
% 
% 性质~\ref{property:ma-6-1-1} 和~\ref{property:ma-6-1-2} 说明求不定积分运算是一种线性运算.
% 
% \begin{example}\label{example:ma-6-1-1}
% 求不定积分
% \[
%   \int\frac{x^4}{1+x^2}\mathrm dx.
% \]
% \end{example}
% 
% 
% \begin{solve}
% 将被积函数恒等变形, 有
% \[
% \begin{aligned}
%   \int\frac{x^4}{1+x^2}\mathrm dx
%   &=
%   \int\frac{x^4-1+1}{1+x^2}\mathrm dx\\
%   &=
%   \int\left(x^2-1+\frac1{1+x^2}\right)\mathrm dx\\
%   &=
%   \frac{x^3}{3}-x+\arctan x+C.
% \end{aligned}
% \qedhere
% \]
% \end{solve}
% 
% 
% \begin{example}\label{example:ma-6-1-2}
% 求不定积分
% \[
%   \int\tan^2x\mathrm dx.
% \]
% \end{example}
% 
% 
% \begin{solve}
% 将被积函数恒等变形, 有
% \[
% \begin{aligned}
%   \int\tan^2x\mathrm dx
%   &=
%   \int\frac{\sin^2x}{\cos^2x}\mathrm dx\\
%   &=
%   \int\left(\frac1{\cos^2x}-1\right)\mathrm dx
%   =
%   \tan x-x+C.
% \end{aligned}
% \qedhere
% \]
% \end{solve}
% 
% 
% \begin{example}\label{example:ma-6-1-3}
% 求不定积分
% \[
%   \int\frac1{\sin^22x}\mathrm dx.
% \]
% \end{example}
% 
% 
% \begin{solve}
% \[
% \begin{aligned}
%   \int\frac1{\sin^22x}\mathrm dx
%   &=
%   \int\frac1{4\sin^2x\cos^2x}\mathrm dx\\
%   &=
%   \int\frac{\sin^2x+\cos^2x}{4\sin^2x\cos^2x}\mathrm dx\\
%   &=
%   \frac14\left(\int\frac1{\cos^2x}\mathrm dx+\int\frac1{\sin^2x}\mathrm dx\right)\\
%   &=
%   \frac14(\tan x-\cot x)+C.
% \end{aligned}
% \qedhere
% \]
% \end{solve}
% 
% 
% \section{换元法与分部积分法}
% \label{sec:ma-6-2}
% 
% 不定积分换元法是对应于复合函数的求导法则的。大家已经知道, 对于复合函数而言, 我们有如下的一阶微分形式不变性
% \[
%   f(u(x))u'(x)\mathrm dx=f(u)\mathrm du,
% \]
% 其中 $u=u(x)$。由此立得
% \begin{equation}
%   \int f(u(x))u'(x)\mathrm dx=\int f(u)\mathrm du.
%   \label{eq:ma-6-2-1}
% \end{equation}
% 
% 利用这一等式可导出两种求不定积分的方法: 如果所求的不定积分有 \eqref{eq:ma-6-2-1} 的左边所示的形式, 则我们可以利用变量替换 $u=u(x)$ 将其转化为 \eqref{eq:ma-6-2-1} 右边的形式, 然后通过求右边关于变量 $u$ 的不定积分而得到所要求的不定积分, 这就是所谓的\keyterm{第一换元法}; 如果所求的不定积分有 \eqref{eq:ma-6-2-1} 的右边所示的形式, 则我们可以利用变量替换 $u=u(x)$ 将其转化为 \eqref{eq:ma-6-2-1} 左边的形式, 然后通过求左边关于变量 $x$ 的不定积分而得到所要求的不定积分, 这就是所谓的\keyterm{第二换元法}.
% 
% \subsection{第一换元法}
% \label{subsec:ma-6-2-1}
% 
% \begin{theorem}\label{theorem:ma-6-2-1}
% 如果
% \[
%   \int f(u)\mathrm du=F(u)+C,
% \]
% 而 $u=u(x)$ 是关于 $x$ 的可微函数, 则有
% \[
%   \int f(u(x))u'(x)\mathrm dx=F(u(x))+C.
% \]
% \end{theorem}
% 
% 
% \begin{proof}
% 由定理的条件, 我们有
% \[
%   \mathrm dF(u)=f(u)\mathrm du.
% \]
% 再根据一阶微分有形式不变性, 可得
% \[
%   \mathrm dF(u(x))=f(u(x))\mathrm du(x)=f(u(x))u'(x)\mathrm dx,
% \]
% 所以
% \[
%   \int f(u(x))u'(x)\mathrm dx=F(u(x))+C.
% \qedhere
% \]
% \end{proof}
% 
% 
% \begin{example}\label{example:ma-6-2-1}
% 求不定积分
% \[
%   \int\frac1{x-1}\mathrm dx.
% \]
% \end{example}
% 
% 
% \begin{solve}
% 我们已经知道
% \[
%   \int\frac{\mathrm du}{u}=\ln|u|+C.
% \]
% 如果令 $u=x-1$, 则 $u$ 是关于 $x$ 的可微函数, 而且 $\mathrm du=\mathrm dx$, 因此我们有
% \[
%   \int\frac1{x-1}\mathrm dx
%   =
%   \int\frac{\mathrm du}{u}
%   =
%   \ln|u|+C
%   =
%   \ln|x-1|+C.
% \qedhere
% \]
% \end{solve}
% 
% 
% \begin{example}\label{example:ma-6-2-2}
% 求不定积分
% \[
%   \int\frac{\mathrm dx}{ax+b}\quad (a\ne0).
% \]
% \end{example}
% 
% 
% \begin{solve}
% 令 $u=ax+b$, 则有 $\mathrm du=a\mathrm dx$, 即 $\mathrm dx=\dfrac{\mathrm du}{a}$, 于是有
% \[
%   \int\frac{\mathrm dx}{ax+b}
%   =
%   \int\frac1a\frac{\mathrm du}{u}
%   =
%   \frac1a\ln|u|+C
%   =
%   \frac1a\ln|ax+b|+C.
% \qedhere
% \]
% \end{solve}
% 
% 
% \begin{example}\label{example:ma-6-2-3}
% 求不定积分
% \[
%   \int(ax+b)^n\mathrm dx\quad (a\ne0,\ n\ne-1).
% \]
% \end{example}
% 
% 
% \begin{solve}
% 令 $u=ax+b$, 则有
% \[
%   \int(ax+b)^n\mathrm dx
%   =
%   \frac1a\int u^n\mathrm du
%   =
%   \frac{u^{n+1}}{a(n+1)}+C
%   =
%   \frac{(ax+b)^{n+1}}{a(n+1)}+C.
% \qedhere
% \]
% \end{solve}
% 
% 
% \begin{example}\label{example:ma-6-2-4}
% 求不定积分
% \[
%   \int\frac{\mathrm dx}{a^2+x^2}\quad (a\ne0).
% \]
% \end{example}
% 
% 
% \begin{solve}
% 由于
% \[
%   \int\frac{\mathrm dx}{a^2+x^2}
%   =
%   \frac1{a^2}\int\frac{\mathrm dx}{1+\left(\dfrac xa\right)^2},
% \]
% 因此令 $u=\dfrac xa$, 则有 $\mathrm du=\dfrac{\mathrm dx}{a}$, 即 $\mathrm dx=a\mathrm du$, 从而
% \[
%   \int\frac{\mathrm dx}{a^2+x^2}
%   =
%   \frac1a\int\frac{\mathrm du}{1+u^2}
%   =
%   \frac1a\arctan\frac xa+C.
% \qedhere
% \]
% \end{solve}
% 
% 
% \begin{example}\label{example:ma-6-2-5}
% 求不定积分
% \[
%   \int\frac{\mathrm dx}{x^4-1}.
% \]
% \end{example}
% 
% 
% \begin{solve}
% 将被积函数恒等变形, 有
% \[
% \begin{aligned}
%   \int\frac{\mathrm dx}{x^4-1}
%   &=
%   \int\frac{\mathrm dx}{(x^2-1)(x^2+1)}\\
%   &=
%   \frac12\left(\int\frac{\mathrm dx}{x^2-1}-\int\frac{\mathrm dx}{x^2+1}\right)\\
%   &=
%   \frac12\left[\frac12\left(\int\frac{\mathrm dx}{x-1}-\int\frac{\mathrm dx}{x+1}\right)
%   -\int\frac{\mathrm dx}{x^2+1}\right]\\
%   &=
%   \frac14\ln\left|\frac{x-1}{x+1}\right|-\frac12\arctan x+C.
% \end{aligned}
% \]
% 
% 从以上例子可以看出, 采用第一换元法求不定积分的关键是, 要将所求不定积分的被积表达式 $f(x)\mathrm dx$ 写成两个因式的乘积, 前一因式形如 $g(u(x))$, 后一因式形如 $\mathrm du(x)$。对于与上面例子类似的一些简单的不定积分, 我们很容易做到这一点。但对于较为复杂的不定积分, 要做到这一点就比较困难。要对具体问题做具体分析, 灵活运用。下面看一些较为复杂的例子.
% \end{solve}
% 
% 
% \begin{example}\label{example:ma-6-2-6}
% 求不定积分
% \[
%   \int\frac{x\mathrm dx}{x^2+1}.
% \]
% \end{example}
% 
% 
% \begin{solve}
% 令 $u=x^2+1$, 则有 $\mathrm du=2x\mathrm dx$, 即 $x\mathrm dx=\dfrac12\mathrm du$, 所以
% \[
%   \int\frac{x\mathrm dx}{x^2+1}
%   =
%   \frac12\int\frac{\mathrm d(x^2+1)}{x^2+1}
%   =
%   \frac12\ln(x^2+1)+C.
% \]
% 在此例中, 我们在计算过程中没有写出关于 $u$ 的积分。在熟练后, 我们不必每次写出代表中间变量的符号 $u$, 而只要在心目中将变量 $u$ 所代表的函数 $u(x)$ 看做一个整体就可以了。此外, 从这个例子也可以看出, 确定 $u=u(x)$ 的时候, 要仔细观察在除去因子 $g(u(x))$ 之后, 是否剩下的函数恰好为 $u'(x)$ (或相差一个常数倍)。这些技巧读者应该在不断的练习过程中加以总结和掌握.
% \end{solve}
% 
% 
% \begin{example}\label{example:ma-6-2-7}
% 求不定积分
% \[
%   \int\frac{\mathrm dx}{a^2-x^2}\quad (a\ne0).
% \]
% \end{example}
% 
% 
% \begin{solve}[{解法 1}]
% \[
% \begin{aligned}
%   \int\frac{\mathrm dx}{a^2-x^2}
%   &=
%   \frac1a\int\frac{\mathrm d\left(\dfrac xa\right)}{1-\left(\dfrac xa\right)^2}\\
%   &=
%   \frac1{2a}\ln\left|\frac{1+\dfrac xa}{1-\dfrac xa}\right|+C
%   =
%   \frac1{2a}\ln\left|\frac{a+x}{a-x}\right|+C.
% \end{aligned}
% \qedhere
% \]
% \end{solve}
% 
% 
% \begin{solve}[{解法 2}]
% \[
% \begin{aligned}
%   \int\frac{\mathrm dx}{a^2-x^2}
%   &=
%   \frac1{2a}\int\left(\frac1{a+x}+\frac1{a-x}\right)\mathrm dx\\
%   &=
%   \frac1{2a}\left[\int\frac{\mathrm d(a+x)}{a+x}
%   -\int\frac{\mathrm d(a-x)}{a-x}\right]\\
%   &=
%   \frac1{2a}\ln\left|\frac{a+x}{a-x}\right|+C.
% \end{aligned}
% \qedhere
% \]
% \end{solve}
% 
% 
% \begin{example}\label{example:ma-6-2-8}
% 求不定积分
% \[
%   I=\int\frac{\mathrm dx}{\sqrt{a^2-x^2}}\quad (a>0).
% \]
% \end{example}
% 
% 
% \begin{solve}
% 将被积函数变形得
% \[
%   I=
%   \int\frac{\mathrm d\left(\dfrac xa\right)}
%   {\sqrt{1-\left(\dfrac xa\right)^2}}
%   =
%   \arcsin\frac xa+C.
% \qedhere
% \]
% \end{solve}
% 
% 
% \begin{example}\label{example:ma-6-2-9}
% 求不定积分
% \[
%   I=\int\frac{\mathrm dx}{\sqrt{x(1-x)}}.
% \]
% \end{example}
% 
% 
% \begin{solve}[{解法 1}]
% 将被积函数变形得
% \[
%   I=
%   \int\frac{\mathrm dx}{\sqrt{\dfrac14-\left(x-\dfrac12\right)^2}}
%   =
%   \arcsin(2x-1)+C.
% \qedhere
% \]
% \end{solve}
% 
% 
% \begin{solve}[{解法 2}]
% 由 $\mathrm dx/\sqrt x=2\mathrm d\sqrt x$, 有
% \[
%   I=
%   \int\frac{\mathrm dx}{\sqrt x\sqrt{1-x}}
%   =
%   2\int\frac{\mathrm d\sqrt x}{\sqrt{1-(\sqrt x)^2}}
%   =
%   2\arcsin(\sqrt x)+C.
% \qedhere
% \]
% \end{solve}
% 
% 
% \begin{remark}\label{remark:ma-6-source-34}
% 虽然两种方法得到的答案形式不同, 但实质上它们只相差一个常数。这一点请读者自行验证之.
% \end{remark}
% 
% 
% \begin{example}\label{example:ma-6-2-10}
% 求不定积分
% \[
%   I=\int\tan x\mathrm dx.
% \]
% \end{example}
% 
% 
% \begin{solve}
% 利用 $\sin x\mathrm dx=-\mathrm d\cos x$, 有
% \[
%   I=-\int\frac{\mathrm d\cos x}{\cos x}
%   =
%   -\ln|\cos x|+C.
% \qedhere
% \]
% \end{solve}
% 
% 
% \begin{example}\label{example:ma-6-2-11}
% 求不定积分
% \[
%   I_n=\int\tan^nx\mathrm dx.
% \]
% \end{example}
% 
% 
% \begin{solve}
% 利用 $\tan^2x=\sec^2x-1$, 有
% \[
% \begin{aligned}
%   I_n
%   &=
%   \int\tan^{n-2}x\frac{1-\cos^2x}{\cos^2x}\mathrm dx\\
%   &=
%   \int\tan^{n-2}x\mathrm d\tan x-\int\tan^{n-2}x\mathrm dx\\
%   &=
%   \frac1{n-1}\tan^{n-1}x-I_{n-2}.
% \end{aligned}
% \]
% 这是一个递推公式。利用这一公式, 我们可以归纳地求出 $I_n$。例如
% \[
%   I_3=\frac12\tan^2x-\int\tan x\mathrm dx
%   =
%   \frac12\tan^2x+\ln|\cos x|+C;
% \]
% \[
% \begin{aligned}
%   I_4
%   &=
%   \frac13\tan^3x-\int\tan^2x\mathrm dx\\
%   &=
%   \frac13\tan^3x-\tan x+\int1\mathrm dx\\
%   &=
%   \frac13\tan^3x-\tan x+x+C.
% \end{aligned}
% \qedhere
% \]
% \end{solve}
% 
% 
% \begin{example}\label{example:ma-6-2-12}
% 求不定积分
% \[
%   \int\frac{\mathrm dx}{\cos x}.
% \]
% \end{example}
% 
% 
% \begin{solve}[{解法 1}]
% 将被积函数变形, 并利用 $\cos x\mathrm dx=\mathrm d\sin x$ 得
% \[
% \begin{aligned}
%   I
%   &=
%   \int\frac{\cos x\mathrm dx}{\cos^2x}
%   =
%   \int\frac{\mathrm d(\sin x)}{1-\sin^2x}\\
%   &=
%   \frac12\ln\left|\frac{1+\sin x}{1-\sin x}\right|+C
%   =
%   \frac12\ln\left|\frac{1+\sin x}{\cos x}\right|^2+C\\
%   &=
%   \ln|\sec x+\tan x|+C.
% \end{aligned}
% \qedhere
% \]
% \end{solve}
% 
% 
% \begin{solve}[{解法 2}]
% 将被积函数变形得
% \[
% \begin{aligned}
%   I
%   &=
%   \int\frac{\mathrm dx}{\cos^2\dfrac x2-\sin^2\dfrac x2}
%   =
%   2\int\frac{\mathrm d\left(\dfrac x2\right)}
%   {\cos^2\dfrac x2\left(1-\tan^2\dfrac x2\right)}\\
%   &=
%   2\int\frac{\mathrm d\tan\left(\dfrac x2\right)}
%   {1-\tan^2\dfrac x2}
%   =
%   \ln\left|\frac{1+\tan\dfrac x2}{1-\tan\dfrac x2}\right|+C\\
%   &=
%   \ln\left|\tan\left(\frac x2+\frac\pi4\right)\right|+C.
% \end{aligned}
% \qedhere
% \]
% \end{solve}
% 
% 
% \begin{example}\label{example:ma-6-2-13}
% 求不定积分
% \[
%   I=\int\frac{\mathrm dx}{1+x^3}
%   \quad\text{和}\quad
%   J=\int\frac{x\mathrm dx}{1+x^3}.
% \]
% \end{example}
% 
% 
% \begin{solve}
% 由于
% \[
% \begin{aligned}
%   I+J
%   &=
%   \int\frac{1+x}{1+x^3}\mathrm dx
%   =
%   \int\frac{\mathrm dx}{1-x+x^2}
%   =
%   \int\frac{\mathrm dx}{\left(x-\dfrac12\right)^2+\dfrac34}\\
%   &=
%   \frac2{\sqrt3}\arctan\frac{2x-1}{\sqrt3}+C,
% \end{aligned}
% \]
% \[
% \begin{aligned}
%   I-J
%   &=
%   \int\frac{1-x}{1+x^3}\mathrm dx
%   =
%   \int\frac{\mathrm dx}{1+x}
%   -
%   \int\frac{x^2\mathrm dx}{1+x^3}\\
%   &=
%   \ln|1+x|-\frac13\ln|1+x^3|+C,
% \end{aligned}
% \]
% 所以
% \[
%   I=
%   \frac1{\sqrt3}\arctan\frac{2x-1}{\sqrt3}
%   +
%   \frac12\ln|1+x|-\frac16\ln|1+x^3|+C;
% \]
% \[
%   J=
%   \frac1{\sqrt3}\arctan\frac{2x-1}{\sqrt3}
%   -
%   \frac12\ln|1+x|+\frac16\ln|1+x^3|+C.
% \qedhere
% \]
% \end{solve}
% 
% 
% \subsection{第二换元法}
% \label{subsec:ma-6-2-2}
% 
% \begin{theorem}\label{theorem:ma-6-2-2}
% 设函数 $x=x(t)$ 在某一开区间上可导, 且 $x'(t)\ne0$。如果
% \[
%   \int f(x(t))x'(t)\mathrm dt=G(t)+C,
% \]
% 则有
% \[
%   \int f(x)\mathrm dx=G(t(x))+C,
% \]
% 其中 $t=t(x)$ 为 $x=x(t)$ 的反函数.
% \end{theorem}
% 
% 
% \begin{proof}
% 由不定积分的定义, 可得 $G'(t)=f(x(t))x'(t)$。再由 $x'(t)\ne0$ 和达布定理知, $x'(t)$ 恒大于 $0$ 或恒小于 $0$。所以 $x(t)$ 是严格单调连续函数, 从而其反函数 $t=t(x)$ 存在、连续、且严格单调。注意到
% \[
%   t'(x)=\frac1{x'(t(x))},
% \]
% 便有
% \[
%   (G(t(x)))'=G'(t(x))t'(x)=f(x)x'(t(x))t'(x)=f(x),
% \]
% 所以
% \[
%   \int f(x)\mathrm dx=G(t(x))+C.
% \]
% 
% 第二换元法主要用来求含有根式的不定积分.
% \end{proof}
% 
% 
% \begin{example}\label{example:ma-6-2-14}
% 求不定积分
% \[
%   I=\int\sqrt{a^2-x^2}\mathrm dx.
% \]
% \end{example}
% 
% 
% \begin{solve}
% 令 $x=a\sin t,\ |t|<\pi/2$, 则
% \[
%   I=\int a^2\cos^2t\mathrm dt
%   =
%   \frac{a^2}2t+\frac{a^2}4\sin2t+C
%   =
%   \frac{a^2}2\arcsin\frac xa+\frac12x\sqrt{a^2-x^2}+C.
% \]
% 在将上面积分结果中的新变量 $t$ 返回到变量 $x$ 时, 我们根据变换 $x=a\sin t$ 作直角三角形 (见图 6.2.1), 可得出
% \[
%   \cos t=\frac{\sqrt{a^2-x^2}}{a}.
% \qedhere
% \]
% \end{solve}
% 
% 
% \begin{example}\label{example:ma-6-2-15}
% 求不定积分
% \[
%   I=\int\frac1{\sqrt{x^2-a^2}}\mathrm dx\quad (a>0,\ |x|>a).
% \]
% \end{example}
% 
% 
% \begin{solve}[{解法 1}]
% 令 $x=a\sec t\ (0<t<\pi/2)$, 则
% \[
% \begin{aligned}
%   I
%   &=
%   \int\frac{a\sec t\tan t}{a\tan t}\mathrm dt
%   =
%   \int\frac{\mathrm dt}{\cos t}
%   =
%   \ln|\sec t+\tan t|+C\\
%   &=
%   \ln\left|\frac xa+\frac{\sqrt{x^2-a^2}}a\right|+C
%   =
%   \ln|x+\sqrt{x^2-a^2}|+C.
% \end{aligned}
% \]
% 在将上面积分结果中的新变量 $t$ 返回到变量 $x$ 时, 我们利用了图 6.2.2, 其中 $\sec t=x/a,\ \tan t=\sqrt{x^2-a^2}/a$.
% \end{solve}
% 
% 
% \begin{solve}[{解法 2}]
% 令 $x=a\cosh t,\ 0<t<+\infty$, 则 $t=\ln|x+\sqrt{x^2-a^2}|-\ln a$, 而且
% \[
%   I=\int\frac{a\sinh t}{a\sinh t}\mathrm dt
%   =
%   t+C
%   =
%   \ln|x+\sqrt{x^2-a^2}|+C.
% \qedhere
% \]
% \end{solve}
% 
% 
% \begin{remark}\label{remark:ma-6-source-51}
% 上面是对 $x>a$ 进行积分的, 对于 $x<-a$ 情形, 可以用同样方法处理.
% \end{remark}
% 
% 
% \begin{example}\label{example:ma-6-2-16}
% 求不定积分
% \[
%   I=\int\frac1{\sqrt{x^2+a^2}}\mathrm dx\quad (a>0).
% \]
% \end{example}
% 
% 
% \begin{solve}[{解法 1}]
% 令 $x=a\sinh t$, 类似于上面的解法 2, 我们有
% \[
%   I=t+C=\operatorname{arcsinh}\frac xa+C
%   =
%   \ln(x+\sqrt{x^2+a^2})+C.
% \qedhere
% \]
% \end{solve}
% 
% 
% \begin{solve}[{解法 2}]
% 令 $x=a\tan t$, 则 $\mathrm dx=a\sec^2t\mathrm dt$,
% \[
% \begin{aligned}
%   I
%   &=
%   \int\frac1{a\sec t}\cdot a\sec^2t\mathrm dt
%   =
%   \int\frac1{\cos t}\mathrm dt\\
%   &=
%   \ln|\sec t+\tan t|+C
%   =
%   \ln\left|\frac{\sqrt{a^2+x^2}}a+\frac xa\right|+C\\
%   &=
%   \ln|x+\sqrt{a^2+x^2}|+C.
% \end{aligned}
% \]
% 将 $t$ 还原 $x$ 的过程中, 用到图 6.2.3 中 $\tan t=x/a$ 的关系.
% \end{solve}
% 
% 
% \begin{example}\label{example:ma-6-2-17}
% 求不定积分
% \[
%   I=\int\frac1{\sqrt x+\sqrt[3]x}\mathrm dx.
% \]
% \end{example}
% 
% 
% \begin{solve}
% 令 $x=t^6$, 则有
% \[
% \begin{aligned}
%   I
%   &=
%   \int\frac{6t^5\mathrm dt}{t^3+t^2}
%   =
%   6\int\frac{t^3}{1+t}\mathrm dt\\
%   &=
%   6\int\left(t^2-t+1-\frac1{1+t}\right)\mathrm dt\\
%   &=
%   2t^3-3t^2+6t-6\ln|1+t|+C\\
%   &=
%   2\sqrt x-3\sqrt[3]x+6\sqrt[6]x-6\ln|1+\sqrt[6]x|+C.
% \end{aligned}
% \qedhere
% \]
% \end{solve}
% 
% 
% \subsection{分部积分法}
% \label{subsec:ma-6-2-3}
% 
% 对应于求导法则
% \[
%   (u(x)v(x))'=u'(x)v(x)+u(x)v'(x),
% \]
% 有如下的分部积分法:
% 
% \begin{theorem}\label{theorem:ma-6-2-3}
% 设 $u(x),v(x)$ 可导, 若 $\int u'(x)v(x)\mathrm dx$ 存在, 则
% \[
%   \int u(x)v'(x)\mathrm dx
%   =
%   u(x)v(x)-\int u'(x)v(x)\mathrm dx.
% \]
% \end{theorem}
% 
% 
% \begin{proof}
% 由 $(uv)'=u'v+uv'$, 我们有 $uv'=(uv)'-u'v$。而该等式的右端两项的原函数都存在, 故其左端项的原函数也存在, 且有
% \[
%   \int uv'\mathrm dx=uv-\int u'v\mathrm dx.
% \qedhere
% \]
% \end{proof}
% 
% 
% \begin{remark}\label{remark:ma-6-source-59}
% 为了便于记忆, 常将分部积分公式写成
% \[
%   \int u\mathrm dv=uv-\int v\mathrm du.
% \]
% 用分部积分法求不定积分的一般步骤为:
% 
% (1) 把被积函数拆成 $uv'$, 并且将 $v'\mathrm dx$ 改写成 $\mathrm dv$, 通常 $v'$ 取
% \[
%   e^{\pm x},\ \sin x,\ \cos x,\ x^n,\ \sinh x,\ \cosh x
% \]
% 等函数;
% 
% (2) 使用分部积分公式;
% 
% (3) 求不定积分 $\int u'v\mathrm dx$, 或者设法建立函数方程来求解.
% \end{remark}
% 
% 
% \begin{example}\label{example:ma-6-2-18}
% 求不定积分
% \[
%   I=\int x^3\ln x\mathrm dx.
% \]
% \end{example}
% 
% 
% \begin{solve}
% 令 $u=\ln x,\ \mathrm dv=x^3\mathrm dx=\mathrm d\dfrac{x^4}4$, 即 $v=\dfrac{x^4}4$, 则
% \[
%   I=\frac14x^4\ln x-\frac14\int x^3\mathrm dx
%   =
%   \frac14x^4\ln x-\frac1{16}x^4+C.
% \qedhere
% \]
% \end{solve}
% 
% 
% \begin{example}\label{example:ma-6-2-19}
% 求不定积分
% \[
%   I=\int\arctan x\mathrm dx.
% \]
% \end{example}
% 
% 
% \begin{solve}
% 设 $v=x,\ u=\arctan x$, 则
% \[
%   I
%   =
%   x\arctan x-\int x\mathrm d(\arctan x)
%   =
%   x\arctan x-\int\frac{x\mathrm dx}{1+x^2}
%   =
%   x\arctan x-\frac12\ln(1+x^2)+C.
% \qedhere
% \]
% \end{solve}
% 
% 
% \begin{example}\label{example:ma-6-2-20}
% 求不定积分
% \[
%   I=\int x^2\sin x\mathrm dx.
% \]
% \end{example}
% 
% 
% \begin{solve}
% 设 $v=\cos x,\ u=x^2$, 即 $\sin x\mathrm dx=-\mathrm dv$, 则
% \[
% \begin{aligned}
%   I
%   &=
%   -\int x^2\mathrm d(\cos x)
%   =
%   -x^2\cos x+\int\cos x\mathrm d x^2\\
%   &=
%   -x^2\cos x+2\int x\mathrm d(\sin x)\\
%   &=
%   -x^2\cos x+2x\sin x-2\int\sin x\mathrm dx\\
%   &=
%   -x^2\cos x+2x\sin x+2\cos x+C.
% \end{aligned}
% \qedhere
% \]
% \end{solve}
% 
% 
% \begin{example}\label{example:ma-6-2-21}
% 求不定积分
% \[
%   I=\int\sqrt{x^2-1}\mathrm dx.
% \]
% \end{example}
% 
% 
% \begin{solve}[{解法 1}]
% 由于
% \[
% \begin{aligned}
%   I
%   &=
%   x\sqrt{x^2-1}-\int x\mathrm d(\sqrt{x^2-1})\\
%   &=
%   x\sqrt{x^2-1}-\int\frac{x^2}{\sqrt{x^2-1}}\mathrm dx\\
%   &=
%   x\sqrt{x^2-1}-\int\sqrt{x^2-1}\mathrm dx-\int\frac{\mathrm dx}{\sqrt{x^2-1}}\\
%   &=
%   x\sqrt{x^2-1}-\ln|x+\sqrt{x^2-1}|-I,
% \end{aligned}
% \]
% 所以
% \[
%   I=\frac12x\sqrt{x^2-1}-\frac12\ln|x+\sqrt{x^2-1}|+C.
% \qedhere
% \]
% \end{solve}
% 
% 
% \begin{solve}[{解法 2}]
% 令 $x=\cosh t$, 则有
% \[
% \begin{aligned}
%   I
%   &=
%   \int\sinh^2t\mathrm dt
%   =
%   \int\frac{\cosh2t-1}{2}\mathrm dt\\
%   &=
%   \frac14\sinh2t-\frac12t+C\\
%   &=
%   \frac12x\sqrt{x^2-1}-\frac12\ln|x+\sqrt{x^2-1}|+C.
% \end{aligned}
% \qedhere
% \]
% \end{solve}
% 
% 
% \begin{example}\label{example:ma-6-2-22}
% 求不定积分
% \[
%   I=\int e^{ax}\cos bx\mathrm dx
%   \quad\text{和}\quad
%   J=\int e^{ax}\sin bx\mathrm dx.
% \]
% \end{example}
% 
% 
% \begin{solve}
% 利用分部积分公式, 可得
% \[
%   I=\frac1b\int e^{ax}\mathrm d(\sin bx)
%   =
%   \frac1be^{ax}\sin bx-\frac abJ;
% \]
% \[
%   J=-\frac1b\int e^{ax}\mathrm d(\cos bx)
%   =
%   -\frac1be^{ax}\cos bx+\frac abI.
% \]
% 所以
% \[
%   \begin{cases}
%     bI+aJ=e^{ax}\sin bx,\\
%     aI-bJ=e^{ax}\cos bx.
%   \end{cases}
% \]
% 解此方程组, 可得
% \[
%   I=\frac{e^{ax}(a\cos bx+b\sin bx)}{a^2+b^2}+C,
% \]
% \[
%   J=\frac{e^{ax}(a\sin bx-b\cos bx)}{a^2+b^2}+C.
% \qedhere
% \]
% \end{solve}
% 
% 
% \begin{example}\label{example:ma-6-2-23}
% 求不定积分
% \[
%   I=\int\frac{\sin x}{a\cos x+b\sin x}\mathrm dx
%   \quad\text{和}\quad
%   J=\int\frac{\cos x}{a\cos x+b\sin x}\mathrm dx.
% \]
% \end{example}
% 
% 
% \begin{solve}
% 由于
% \[
%   bI+aJ=
%   \int\frac{a\cos x+b\sin x}{a\cos x+b\sin x}\mathrm dx
%   =
%   x+C,
% \]
% \[
%   bJ-aI=
%   \int\frac{\mathrm d(a\cos x+b\sin x)}{a\cos x+b\sin x}
%   =
%   \ln|a\cos x+b\sin x|+C,
% \]
% 所以
% \[
%   I=\frac1{a^2+b^2}(bx-a\ln|a\cos x+b\sin x|)+C,
% \]
% \[
%   J=\frac1{a^2+b^2}(ax+b\ln|a\cos x+b\sin x|)+C.
% \qedhere
% \]
% \end{solve}
% 
% 
% \begin{example}\label{example:ma-6-2-24}
% 求不定积分
% \[
%   K_n=\int\cos^nx\mathrm dx.
% \]
% \end{example}
% 
% 
% \begin{solve}
% 由于
% \[
% \begin{aligned}
%   K_n
%   &=
%   \int\cos^{n-1}x\mathrm d(\sin x)
%   =
%   \sin x\cos^{n-1}x-\int\sin x\mathrm d(\cos^{n-1}x)\\
%   &=
%   \sin x\cos^{n-1}x+(n-1)\int\sin^2x\cos^{n-2}x\mathrm dx\\
%   &=
%   \sin x\cos^{n-1}x+(n-1)\int\cos^{n-2}x\mathrm dx
%   -(n-1)\int\cos^nx\mathrm dx\\
%   &=
%   \sin x\cos^{n-1}x+(n-1)K_{n-2}-(n-1)K_n,
% \end{aligned}
% \]
% 所以
% \[
%   K_n=\frac1n\sin x\cos^{n-1}x+\frac{n-1}nK_{n-2}.
% \]
% 由此可推出:
% \[
%   K_0=x+C;\quad
%   K_1=\sin x+C;\quad
%   K_2=\frac14\sin2x+\frac12x+C;\quad
%   \cdots\cdots\cdots
% \qedhere
% \]
% \end{solve}
% 
% 
% \begin{example}\label{example:ma-6-2-25}
% 求不定积分
% \[
%   I_n=\int\sin^nx\mathrm dx.
% \]
% \end{example}
% 
% 
% \begin{solve}
% 由于
% \[
% \begin{aligned}
%   I_n
%   &=
%   -\int\sin^{n-1}x\mathrm d\cos x\\
%   &=
%   -\sin^{n-1}x\cos x+(n-1)\int\cos^2x\sin^{n-2}x\mathrm dx\\
%   &=
%   -\sin^{n-1}x\cos x+(n-1)I_{n-2}-(n-1)I_n,
% \end{aligned}
% \]
% 所以
% \[
%   I_n=-\frac{\sin^{n-1}x\cos x}{n}+\frac{n-1}{n}I_{n-2}.
% \]
% 再注意到
% \[
%   I_0=x+C,\qquad I_1=-\cos x+C,
% \]
% 便有
% \[
%   I_2=-\frac12\sin x\cos x+\frac12x+C;
% \]
% \[
%   I_3=-\frac13\sin^2x\cos x-\frac23\cos x+C;
% \]
% \[
%   I_4=-\frac14\sin^3x\cos x-\frac38\sin x\cos x+\frac38x+C;
% \]
% \[
%   I_5=-\frac15\sin^4x\cos x-\frac4{15}\sin^2x\cos x-\frac8{15}\cos x+C;
% \]
% \[
%   \cdots\cdots\cdots
% \qedhere
% \]
% \end{solve}
% 
% 
% \begin{example}\label{example:ma-6-2-26}
% 求不定积分
% \[
%   I_n=\int\frac{\mathrm dx}{(x^2+a^2)^n}.
% \]
% \end{example}
% 
% 
% \begin{solve}
% 由于
% \[
% \begin{aligned}
%   I_n
%   &=
%   \frac{x}{(x^2+a^2)^n}
%   +
%   2n\int\frac{x^2}{(x^2+a^2)^{n+1}}\mathrm dx\\
%   &=
%   \frac{x}{(x^2+a^2)^n}
%   +
%   2nI_n-2na^2I_{n+1},
% \end{aligned}
% \]
% 所以
% \[
%   I_{n+1}=
%   \frac{x}{2na^2(x^2+a^2)^n}+\frac{2n-1}{2na^2}I_n.
% \]
% 特别地, 我们有
% \[
%   I_2=
%   \int\frac{\mathrm dx}{(x^2+a^2)^2}
%   =
%   \frac{x}{2a^2(x^2+a^2)}+\frac1{2a^3}\arctan\frac xa+C.
% \]
% 
% 从理论上来讲, 任何初等函数都有原函数 (见下章的定积分理论)。上面我们所举的例子中遇到的初等函数的原函数均为初等函数, 但并非所有的初等函数的原函数都是初等函数。例如, 已经证明如下的积分就都不能表成初等函数:
% \[
%   \int e^{-x^2}\mathrm dx,\qquad
%   \int\frac{\sin x}{x}\mathrm dx,\qquad
%   \int\frac{\cos x}{x}\mathrm dx,
% \]
% \[
%   \int\frac1{\ln x}\mathrm dx,\qquad
%   \int\sin x^2\mathrm dx,\qquad
%   \int\cos x^2\mathrm dx.
% \]
% 再如, 积分 $\int(a+bx)^px^q\mathrm dx$ 和椭圆积分
% \[
%   \int R(x,\sqrt{ax^3+bx^2+cx+d})\mathrm dx
% \]
% 以及
% \[
%   \int R(x,\sqrt{ax^4+bx^3+cx^2+dx+e})\mathrm dx
% \]
% 也都不能用初等函数表示, 其中 $p,q$ 和 $p+q$ 非整数, $R$ 为两个变元的有理函数.
% \end{solve}
% 
% 
% \section{其他类型函数的不定积分}
% \label{sec:ma-6-3}
% 
% \subsection{有理函数的不定积分}
% \label{subsec:ma-6-3-1}
% 
% \keyterm{有理函数}是指两个多项式之比, 即有理函数有如下形式
% \[
%   R(x)=\frac{P(x)}{Q(x)},
% \]
% 其中 $P(x),Q(x)$ 为互素多项式。理论上讲, 它的不定积分一定能够表成初等函数.
% 
% 有理函数可分为真分式和假分式: \keyterm{真分式}是指分子次数小于分母次数的有理函数; \keyterm{假分式}是指分子次数大于或等于分母次数的有理函数。显然, 通过作多项式除法, 任何一个假分式都可以写成一个多项式与一个真分式之和。因此我们只需讨论真分式的不定积分即可.
% 
% 从理论上来讲, 任何一个真分式都可以写成最简真分式之和。\keyterm{最简真分式}是指分母为素多项式或素多项式之幂的真分式。对于系数均为实数的多项式来讲, 素多项式只有两种, 即 $x-a$ 和 $x^2+px+q$, 其中 $p^2-4q<0$。因此, 最简真分式只有如下四种:
% \[
%   \frac A{x-a},\qquad
%   \frac A{(x-a)^m}\ (m>1),
% \]
% \[
%   \frac{Ax+B}{x^2+px+q},\qquad
%   \frac{Ax+B}{(x^2+px+q)^n}\ (n>1),
% \]
% 其中 $p^2-4q<0$。而前面所介绍的不定积分方法和例子表明, 我们已经有办法把这些最简真分式的原函数求出来, 而且它们的原函数均为初等函数.
% 
% 事实上, 我们有
% \[
%   (1)\quad
%   \int\frac A{x-a}\mathrm dx
%   =
%   A\int\frac1{x-a}\mathrm d(x-a)
%   =
%   A\ln|x-a|+C.
% \]
% \[
%   (2)\quad
%   \int\frac A{(x-a)^m}\mathrm dx
%   =
%   A\int\frac1{(x-a)^m}\mathrm d(x-a)
%   =
%   -\frac A{m-1}(x-a)^{1-m}+C.
% \]
% (3) 将 $x^2+px+q$ 进行配方得
% \[
%   x^2+px+q=\left(x+\frac p2\right)^2+\left(q-\frac{p^2}4\right),
% \]
% 利用替换
% \[
%   z=x+\frac p2,\qquad x=z-\frac p2,\qquad \mathrm dx=\mathrm dz,
% \]
% 又由于 $p^2-4q<0$, 有 $q-\dfrac{p^2}4>0$, 记 $a^2=q-\dfrac{p^2}4$, 得
% \[
% \begin{aligned}
%   \int\frac{Ax+B}{x^2+px+q}\mathrm dx
%   &=
%   \int\frac{Ax+B}{\left(x+\dfrac p2\right)^2+\left(q-\dfrac{p^2}4\right)}\mathrm dx\\
%   &=
%   \int\frac{A(z-p/2)+B}{z^2+a^2}\mathrm dz\\
%   &=
%   \int\frac{Az}{z^2+a^2}\mathrm dz+
%   \int\frac{B-Ap/2}{z^2+a^2}\mathrm dz\\
%   &=
%   \frac A2\ln(z^2+a^2)+\left(B-\frac{Ap}2\right)\frac1a\arctan\frac za+C\\
%   &=
%   \frac A2\ln(x^2+px+q)+
%   \frac{2B-Ap}{\sqrt{4q-p^2}}\arctan\frac{2x+p}{\sqrt{4q-p^2}}+C.
% \end{aligned}
% \]
% 
% (4) 同样用变量替换 $z=x+\dfrac p2,\ x=z-\dfrac p2,\ \mathrm dx=\mathrm dz$, 且记 $a^2=q-\dfrac{p^2}4$, 得
% \[
% \begin{aligned}
%   \int\frac{Ax+B}{(x^2+px+q)^n}\mathrm dx
%   &=
%   \int\frac{Az}{(z^2+a^2)^n}\mathrm dz+
%   \left(B-\frac{Ap}2\right)\int\frac1{(z^2+a^2)^n}\mathrm dz.
% \end{aligned}
% \]
% 上式右边的第一项不定积分容易求出, 用换元积分法即可; 第二项积分要利用 第~\ref{sec:ma-6-2}~节 中例~\ref{example:ma-6-2-26} 的递推公式。至此上述四种情况的不定积分全部得以解决.
% 
% 现在我们讨论有理真分式的分解。因为多项式 $Q(x)$ 在实数范围内能分解成一次因式和二次因式的乘积, 因此我们有
% \[
%   Q(x)=(x-a)^k\cdots(x-b)^t(x^2+px+q)^l\cdots(x^2+rx+s)^h,
% \]
% 其中 $a,\ldots,b,p,q,\ldots,r,s$ 为常数; 且 $p^2-4q<0,\ldots,r^2-4s<0$; $k,\ldots,t,l,\ldots,h$ 为正整数。可以证明真分式 $\dfrac{P(x)}{Q(x)}$ 必可分解成如下部分分式之和:
% \[
% \begin{aligned}
%   \frac{P(x)}{Q(x)}
%   &=
%   \frac{A_1}{x-a}+\frac{A_2}{(x-a)^2}+\cdots+\frac{A_k}{(x-a)^k}+\cdots\\
%   &\quad+
%   \frac{B_1}{x-b}+\frac{B_2}{(x-b)^2}+\cdots+\frac{B_t}{(x-b)^t}\\
%   &\quad+
%   \frac{C_1x+D_1}{x^2+px+q}+\frac{C_2x+D_2}{(x^2+px+q)^2}
%   +\cdots+\frac{C_lx+D_l}{(x^2+px+q)^l}+\cdots\\
%   &\quad+
%   \frac{E_1x+F_1}{x^2+rx+s}+\frac{E_2x+F_2}{(x^2+rx+s)^2}
%   +\cdots+\frac{E_hx+F_h}{(x^2+rx+s)^h},
% \end{aligned}
% \]
% 其中 $A_1,A_2,\ldots,A_k; B_1,\ldots,B_t; C_1,\ldots,C_l; D_1,\ldots,D_l; E_1,\ldots,E_h; F_1,\ldots,F_h$ 都是常数。因此有理函数的原函数都可由初等函数表示.
% 
% \begin{example}\label{example:ma-6-3-1}
% 求不定积分
% \[
%   I=\int\frac{x\mathrm dx}{(x^2+1)(x-1)}.
% \]
% \end{example}
% 
% \begin{solve}
% 设
% \[
%   \frac{x}{(x^2+1)(x-1)}
%   =
%   \frac{Ax+B}{x^2+1}+\frac C{x-1}.
% \]
% 将右端通分, 并比较等式两端分子同次幂的系数, 得
% \[
%   \begin{cases}
%     A+C=0,\\
%     -A+B=1,\\
%     -B+C=0.
%   \end{cases}
% \]
% 解之得 $A=-\dfrac12,\ B=\dfrac12,\ C=\dfrac12$。因此
% \[
% \begin{aligned}
%   I
%   &=
%   -\frac12\int\frac{x-1}{x^2+1}\mathrm dx+\frac12\int\frac{\mathrm dx}{x-1}\\
%   &=
%   -\frac14\ln(x^2+1)+\frac12\arctan x+\frac12\ln|x-1|+C.
% \end{aligned}
% \qedhere
% \]
% \end{solve}
% 
% 
% \begin{example}\label{example:ma-6-3-2}
% 求不定积分
% \[
%   I=\int\frac{x^3+1}{x^4-3x^3+3x^2-x}\mathrm dx.
% \]
% \end{example}
% 
% \begin{solve}
% 将分母作因式分解, 得
% \[
%   x^4-3x^3+3x^2-x=x(x-1)^3.
% \]
% 于是, 我们可设
% \[
%   \frac{x^3+1}{x(x-1)^3}
%   =
%   \frac Ax+\frac B{(x-1)^3}+\frac C{(x-1)^2}+\frac D{x-1}.
% \]
% 将两边乘 $x$, 然后令 $x=0$, 得 $A=-1$; 两边乘 $(x-1)^3$, 然后令 $x=1$, 得 $B=2$; 两边乘 $x-1$, 然后令 $x\to-\infty$, 得 $D=2$; 最后令 $x=-1$, 得 $C=1$。这样, 我们便有
% \[
% \begin{aligned}
%   I
%   &=
%   -\int\frac{\mathrm dx}x+2\int\frac{\mathrm dx}{(x-1)^3}
%   +\int\frac{\mathrm dx}{(x-1)^2}+2\int\frac{\mathrm dx}{x-1}\\
%   &=
%   -\frac{x}{(x-1)^2}+\ln\frac{(x-1)^2}{|x|}+C.
% \end{aligned}
% \qedhere
% \]
% \end{solve}
% 
% 
% \subsection{三角函数有理式的不定积分}
% \label{subsec:ma-6-3-2}
% 
% 二元有理函数是形如
% \[
%   R(u,v)=\frac{P(u,v)}{Q(u,v)}
% \]
% 的两个变元的函数, 其中 $P(u,v)$ 和 $Q(u,v)$ 是二元多项式(即 $u^i,v^j,u^pv^q$ 的线性组合)。三角函数有理式是指形如 $R(\sin x,\cos x)$ 的函数, 其中 $R(u,v)$ 为二元有理函数, 它是由基本三角函数 $\sin x,\cos x,\tan x,\cot x$ 经有限次四则运算所得的函数。但像如下形式的函数
% \[
%   \frac1{\sqrt{\sin^2x+5\cos^2x}},\qquad
%   \sin x^2,\qquad
%   \frac{\sin x}{x}
% \]
% 则不属此列.
% 
% 下面我们来说明三角函数有理式的不定积分
% \[
%   I=\int R(\sin x,\cos x)\mathrm dx
% \]
% 一定可以表成初等函数。令
% \[
%   t=\tan\frac x2\quad\text{或}\quad x=2\arctan t
% \]
% (通常称之为万能变换), 并且注意到
% \[
% \begin{aligned}
%   \sin x
%   &=
%   2\sin\frac x2\cos\frac x2
%   =
%   \frac{2\tan\dfrac x2}{1+\tan^2\dfrac x2}
%   =
%   \frac{2t}{1+t^2},\\
%   \cos x
%   &=
%   \cos^2\frac x2-\sin^2\frac x2
%   =
%   \frac{1-\tan^2\dfrac x2}{1+\tan^2\dfrac x2}
%   =
%   \frac{1-t^2}{1+t^2},\\
%   \mathrm dx&=\frac{2\mathrm dt}{1+t^2},
% \end{aligned}
% \]
% 即有
% \[
%   I=
%   \int R\left(\frac{2t}{1+t^2},\frac{1-t^2}{1+t^2}\right)
%   \frac{2\mathrm dt}{1+t^2}.
% \]
% 这样, 我们就将求三角函数有理式的不定积分变成了有理函数的不定积分。再由上一小节所介绍的有理函数的积分方法就可解决求三角函数有理式的积分问题。这也表明, 三角函数有理式的原函数一定可以表成初等函数.
% 
% \begin{example}\label{example:ma-6-3-3}
% 求不定积分
% \[
%   I=\int\frac{\mathrm dx}{5+4\sin x}.
% \]
% \end{example}
% 
% \begin{solve}
% 令 $t=\tan\dfrac x2$, 则
% \[
% \begin{aligned}
%   I
%   &=
%   \int\frac{\dfrac2{1+t^2}}{5+\dfrac{8t}{1+t^2}}\mathrm dt
%   =
%   \int\frac{2\mathrm dt}{5t^2+8t+5}\\
%   &=
%   \frac25\int\frac{\mathrm dt}{\left(t+\dfrac45\right)^2+\dfrac9{25}}
%   =
%   \frac23\arctan\frac{5t+4}3+C\\
%   &=
%   \frac23\arctan\frac{5\tan\dfrac x2+4}3+C.
% \end{aligned}
% \]
% 
% 这里需指出的是, 具体解题时千万不要死搬硬套, 要具体问题作具体分析, 灵活运用。事实上, 这种``万能变换''往往导致所要求的有理函数的积分比较复杂, 而对某些三角函数有理式的不定积分来说, 采用其他变换会更加方便。例如, 当 $R(u,-v)=-R(u,v)$ 时, 我们可以证明, 这时必有 $R(u,v)=vR_1(u,v^2)$ 成立, 其中 $R_1$ 是一个二变元的有理函数。于是若令 $t=\sin x$, 则有
% \[
% \begin{aligned}
%   \int R(\sin x,\cos x)\mathrm dx
%   &=
%   \int\cos xR_1(\sin x,\cos^2x)\mathrm dx\\
%   &=
%   \int R_1(t,1-t^2)\mathrm dt.
% \end{aligned}
% \]
% 这表明, 若被积函数 $R(\sin x,\cos x)$ 关于 $\cos x$ 是奇函数, 则我们可用变换 $t=\sin x$ 将其转化为关于 $t$ 的有理函数的不定积分.
% 
% 同理, 若 $R(-u,v)=-R(u,v)$, 则有 $R(u,v)=uR_1(u^2,v)$, 其中 $R_1$ 是一个二变元的有理函数。于是若令 $t=\cos x$, 则有
% \[
% \begin{aligned}
%   \int R(\sin x,\cos x)\mathrm dx
%   &=
%   \int\sin xR_1(\sin^2x,\cos x)\mathrm dx\\
%   &=
%   -\int R_1(1-t^2,t)\mathrm dt.
% \end{aligned}
% \]
% 这表明, 若被积函数 $R(\sin x,\cos x)$ 关于 $\sin x$ 是奇函数, 则我们可用变换 $t=\cos x$ 将其转化为关于 $t$ 的有理函数的不定积分.
% 
% 再有若 $R(-u,-v)=R(u,v)$, 则可证此时有 $R(u,v)=R_1(u/v,v^2)$, 其中 $R_1$ 是一个二变元的有理函数。于是, 若我们用变换 $t=\tan x$, 则有
% \[
% \begin{aligned}
%   \int R(\sin x,\cos x)\mathrm dx
%   &=
%   \int R_1(\tan x,\cos^2x)\mathrm dx\\
%   &=
%   \int R_1\left(t,\frac1{1+t^2}\right)\frac{\mathrm dt}{1+t^2}.
% \end{aligned}
% \]
% 这表明, 若被积函数 $R(\sin x,\cos x)$ 关于 $\sin x$ 和 $\cos x$ 是偶函数, 则我们可用变换 $t=\tan x$ 将其转化为关于 $t$ 的有理函数的不定积分.
% 
% 上述三种变换一般要比``万能变换''简单得多.
% \end{solve}
% 
% 
% \begin{example}\label{example:ma-6-3-4}
% 求不定积分
% \[
%   I=\int\frac{\cos^3x}{1+\sin^2x}\mathrm dx.
% \]
% \end{example}
% 
% \begin{solve}
% 由于 $R(\sin x,-\cos x)=-R(\sin x,\cos x)$, 故令 $t=\sin x$, 则
% \[
% \begin{aligned}
%   I
%   &=
%   \int\frac{1-t^2}{1+t^2}\mathrm dt
%   =
%   \int\frac{2\mathrm dt}{1+t^2}-\int\mathrm dt\\
%   &=
%   2\arctan t-t+C
%   =
%   2\arctan\sin x-\sin x+C.
% \end{aligned}
% \qedhere
% \]
% \end{solve}
% 
% 
% \begin{example}\label{example:ma-6-3-5}
% 求不定积分
% \[
%   I=\int\frac{\sin x\cos x}{\sin^2x+\cos x}\mathrm dx.
% \]
% \end{example}
% 
% \begin{solve}
% 由于 $R(-\sin x,\cos x)=-R(\sin x,\cos x)$, 故令 $t=\cos x$, 则
% \[
% \begin{aligned}
%   I
%   &=
%   -\int\frac{\sin x\,\mathrm d\cos x}{\sin^2x+\cos x}
%   =
%   -\int\frac{t\mathrm dt}{1-t^2+t}\\
%   &=
%   \frac12\int\frac{\mathrm d(1+t-t^2)}{1+t-t^2}
%   -\frac12\int\frac{\mathrm dt}{1+t-t^2}\\
%   &=
%   \frac12\ln|1+t-t^2|-\frac12\int\frac{\mathrm dt}{\dfrac54-\left(\dfrac12-t\right)^2}\\
%   &=
%   \frac12\ln|1+t-t^2|+
%   \frac1{2\sqrt5}\ln\left|\frac{\sqrt5+1-2t}{\sqrt5-1+2t}\right|+C\\
%   &=
%   \frac12\left[
%   \ln|1+\cos x-\cos^2x|+
%   \frac1{\sqrt5}\ln\left|\frac{\sqrt5+1-2\cos x}{\sqrt5-1+2\cos x}\right|
%   \right]+C.
% \end{aligned}
% \qedhere
% \]
% \end{solve}
% 
% 
% \begin{example}\label{example:ma-6-3-6}
% 求不定积分
% \[
%   I=\int\frac{\cos^2x}{(a^2\sin^2x+b^2\cos^2x)^2}\mathrm dx.
% \]
% \end{example}
% 
% \begin{solve}
% 由于 $R(-\sin x,-\cos x)=R(\sin x,\cos x)$, 故令 $t=\tan x$, 则
% \[
% \begin{aligned}
%   I
%   &=
%   \int\frac{\cos^4x}{(a^2\sin^2x+b^2\cos^2x)^2\cos^2x}\mathrm dx\\
%   &=
%   \int\frac{1}{(a^2\tan^2x+b^2)^2}\mathrm d\tan x\\
%   &=
%   \frac1{b^4}\int\frac{\mathrm dt}{\left(1+\dfrac{a^2}{b^2}t^2\right)^2}
%   =
%   \frac1{ab^3}\int\frac{\mathrm du}{(1+u^2)^2}\quad\left(u=\frac abt\right)\\
%   &=
%   \frac1{ab^3}\left\{\frac12\frac u{1+u^2}+\frac12\arctan u\right\}+C\\
%   &=
%   \frac{\tan x}{2b^2}\frac1{b^2+a^2\tan^2x}
%   +\frac1{2ab^3}\arctan\left(\frac ab\tan x\right)+C.
% \end{aligned}
% \qedhere
% \]
% \end{solve}
% 
% 
% \begin{example}\label{example:ma-6-3-7}
% 求不定积分
% \[
%   I=\frac12\int\frac{1-r^2}{1-2r\cos x+r^2}\mathrm dx
%   \qquad(0<r<1,\ |x|<\pi).
% \]
% \end{example}
% 
% \begin{solve}
% 令 $\tan\dfrac x2=t$, 则有
% \[
%   \cos x=\frac{1-t^2}{1+t^2},\qquad
%   \mathrm dx=\frac{2\mathrm dt}{1+t^2}.
% \]
% 由此得
% \[
% \begin{aligned}
%   I
%   &=
%   \int\frac{(1-r^2)\mathrm dt}{(1-r)^2+(1+r)^2t^2}\\
%   &=
%   \arctan\left(\frac{1+r}{1-r}t\right)+C\\
%   &=
%   \arctan\left(\frac{1+r}{1-r}\tan\frac x2\right)+C.
% \end{aligned}
% \qedhere
% \]
% \end{solve}
% 
% 
% \subsection{无理函数的不定积分}
% \label{subsec:ma-6-3-3}
% 
% 这里我们考虑两类无理函数的不定积分。第一类是如下的积分:
% \[
%   I=\int R\left(x,\sqrt[m]{\frac{ax+b}{cx+d}}\right)\mathrm dx,
% \]
% 其中 $R(u,v)$ 是变元 $u,v$ 的二元函数, $m$ 为正整数, $a,b,c,d$ 为常数, 且 $ad-bc\ne0$。令
% \[
%   \frac{ax+b}{cx+d}=t^m,
% \]
% 则有
% \[
%   x=\frac{dt^m-b}{a-ct^m},\qquad
%   \mathrm dx=\frac{m(ad-bc)t^{m-1}}{(a-ct^m)^2}\mathrm dt.
% \]
% 于是, 我们有
% \[
%   I=
%   \int R\left(\frac{dt^m-b}{a-ct^m},t\right)
%   \frac{m(ad-bc)t^{m-1}}{(a-ct^m)^2}\mathrm dt.
% \]
% 这样, 原不定积分就变成了有理函数的积分, 从而这类无理函数一定有初等函数作为其原函数.
% 
% \begin{example}\label{example:ma-6-3-8}
% 求不定积分
% \[
%   I=\int\frac{1-\sqrt{x+1}}{1+\sqrt[3]{x+1}}\mathrm dx.
% \]
% \end{example}
% 
% \begin{solve}
% 令 $\sqrt[6]{x+1}=t$, 则有
% \[
%   I=6\int\frac{t^5-t^8}{1+t^2}\mathrm dt.
% \]
% 设
% \[
%   \frac{t^5-t^8}{1+t^2}=P_6(t)+\frac{Bt+C}{1+t^2},
% \]
% 其中 $P_6(t)$ 为 $t$ 的多项式, $B,C$ 是实数。在上式两边乘以 $1+t^2$, 并且将 $t=\mathrm i=\sqrt{-1}$ 代入, 得
% \[
%   \mathrm i-1=B\mathrm i+C,
% \]
% 于是有 $B=1,C=-1$, 而且
% \[
% \begin{aligned}
%   P_6(t)
%   &=
%   \frac{t^5-t^8}{1+t^2}-\frac{t-1}{1+t^2}\\
%   &=
%   \frac{t^5-t^8-t+1}{1+t^2}
%   =
%   1-t-t^2+t^3+t^4-t^6.
% \end{aligned}
% \]
% 所以
% \[
% \begin{aligned}
%   I
%   &=
%   6\int\left(1-t-t^2+t^3+t^4-t^6+\frac{t-1}{1+t^2}\right)\mathrm dt\\
%   &=
%   -\frac67t^7+\frac65t^5+\frac32t^4-2t^3-3t^2+6t
%   +3\ln(1+t^2)-6\arctan t+C\\
%   &=
%   -\frac67(x+1)^{7/6}+\frac65(x+1)^{5/6}
%   +\frac32(x+1)^{2/3}-2(x+1)^{1/2}-3(x+1)^{1/3}\\
%   &\quad+
%   6(x+1)^{1/6}+3\ln\left(1+(x+1)^{1/3}\right)
%   -6\arctan(x+1)^{1/6}+C.
% \end{aligned}
% \]
% 
% 我们要考虑的第二类无理函数的不定积分是如下的二项式微分式的积分:
% \[
%   I=\int x^m(a+bx^n)^p\mathrm dx,
% \]
% 其中 $a,b$ 为常数, $m,n,p$ 为有理数。令 $x^n=t$, 则 $\mathrm dx=\dfrac1n t^{1/n-1}\mathrm dt$, 从而有
% \[
% \begin{aligned}
%   I
%   &=
%   \int t^{m/n}(a+bt)^p\frac1n t^{1/n-1}\mathrm dt\\
%   &=
%   \frac1n\int t^{(m+1)/n-1}(a+bt)^p\mathrm dt.
% \end{aligned}
% \]
% 为简明起见, 令 $q=\dfrac{m+1}n-1$, 则
% \begin{equation}
%   I=\int x^m(a+bx^n)^p\mathrm dx
%   =
%   \frac1n\int t^q(a+bt)^p\mathrm dt.
%   \label{eq:ma-6-3-1}
% \end{equation}
% 当 $p$ 是整数而 $q=r/s(r,s$ 为整数) 时, 被积函数最多含有一个根式, 根式内的线性分式为 $t$, 因此只要作变换 $z=\sqrt[s]{t}=\sqrt[s]{x^n}$, 就可将原积分化为有理函数的积分.
% 
% 当 $q$ 为整数而 $p=\mu/\lambda(\mu,\lambda$ 为整数) 时, 被积函数也最多含有一个根式, 根式内的线性分式为 $a+bt$, 因此只要作变换 $z=\sqrt[\lambda]{a+bt}=\sqrt[\lambda]{a+bx^n}$, 就可将原积分化为有理函数的积分.
% 
% 最后, 因为积分 \eqref{eq:ma-6-3-1} 还可以改写成
% \[
%   I=\int x^m(a+bx^n)^p\mathrm dx
%   =
%   \frac1n\int t^{p+q}\left(\frac{a+bt}{t}\right)^p\mathrm dt,
% \]
% 所以当 $p+q$ 是整数而 $p=\mu/\lambda(\mu,\lambda$ 为整数) 时, 被积函数最多含有一个根式, 根式内的线性分式为 $\dfrac{a+bt}{t}$, 因此只要作变换
% \[
%   z=\sqrt[\lambda]{\frac{a+bt}{t}}
%   =
%   \sqrt[\lambda]{\frac{a+bx^n}{x^n}},
% \]
% 就可将原积分化为有理函数的积分.
% 
% 总之, 如果在 $p,q$ 和 $p+q$ 之中至少有一个是整数, 则二项式微分式的积分就可以化成有理函数的积分, 从而也就可以表成初等函数.
% 
% 19 世纪中叶, 切比雪夫曾经证明: 二项式微分式的积分, 除上述三种情形外, 都不能由初等函数表示。比如, 积分
% \[
%   \int\frac{\mathrm dx}{\sqrt{1-x^4}}
% \]
% 就不能用初等函数表示.
% \end{solve}
% 
% 
% \begin{example}\label{example:ma-6-3-9}
% 求不定积分
% \[
%   I=\int\sqrt{x^2+\frac1{x^2}}\,\mathrm dx\qquad(x>0).
% \]
% \end{example}
% 
% \begin{solve}
% 由于原积分可以写成
% \[
%   I=\int x^{-1}(1+x^4)^{1/2}\mathrm dx,
% \]
% 因此
% \[
%   q=\frac{m+1}n-1=\frac{-1+1}4-1=-1
% \]
% 为整数, 从而可用初等函数表示出。具体计算如下:
% \[
% \begin{aligned}
%   I
%   &=
%   \frac12\int\frac{\sqrt{1+x^4}}{x^2}\mathrm d(x^2)
%   =
%   \frac12\int\frac{\sqrt{1+t^2}}{t}\mathrm dt
%   \qquad(x^2=t>0)\\
%   &=
%   \frac12\int\frac{1+t^2}{t\sqrt{1+t^2}}\mathrm dt\\
%   &=
%   \frac12\int\frac{t\mathrm dt}{\sqrt{1+t^2}}
%   +\frac12\int\frac{\mathrm dt}{t^2\sqrt{1+1/t^2}}\\
%   &=
%   \frac12\sqrt{1+t^2}-\frac12\ln\left(\frac1t+\sqrt{1+\frac1{t^2}}\right)+C\\
%   &=
%   \frac12\sqrt{1+x^4}
%   -\frac12\ln\left|\frac{1+\sqrt{1+x^4}}{x^2}\right|+C.
% \end{aligned}
% \qedhere
% \]
% \end{solve}
% 
% 
% \section*{习题六}
% \phantomsection
% \addcontentsline{toc}{section}{习题六}
% \label{exercises:ma-6}
% 
% 1。证明: 若 $\displaystyle \int f(x)\mathrm dx=F(x)+C$, 则
% \[
%   \int f(ax+b)\mathrm dx=\frac1aF(ax+b)+C\qquad(a\ne0).
% \]
% 
% 2。求下列不定积分:
% 
% (1) $\displaystyle \int\left(\sqrt{2x}+\sqrt[3]{3x}\right)\mathrm dx$;
% 
% (2) $\displaystyle \int\left(1+\frac1{x^2}\sqrt{x\sqrt x}\right)\mathrm dx$;
% 
% (3) $\displaystyle \int\left(3^{x+1}+3^{-x}+\frac13e^x\right)\mathrm dx$;
% 
% (4) $\displaystyle \int\frac{2-\sin^2x}{\cos^2x}\mathrm dx$;
% 
% (5) $\displaystyle \int\left(\sqrt{\frac{1+x}{1-x}}+\sqrt{\frac{1-x}{1+x}}\right)\mathrm dx$;
% 
% (6) $\displaystyle \int\frac{\cos2x}{\cos^2x\sin^2x}\mathrm dx$;
% 
% (7) $\displaystyle \int\sqrt{1-\sin2x}\,\mathrm dx$;
% 
% (8) $\displaystyle \int(x+|x|)\mathrm dx$;
% 
% (9) $\displaystyle \int\operatorname{sgn}(\sin x)\mathrm dx$;
% 
% (10) $\displaystyle \int[x]\mathrm dx$.
% 
% 3。用换元法求下列不定积分:
% 
% (1) $\displaystyle \int\frac{2\mathrm dx}{3-5x}$;
% 
% (2) $\displaystyle \int\frac{\mathrm dx}{\sqrt{2-x^2}}$;
% 
% (3) $\displaystyle \int x\sqrt[3]{1-3x}\,\mathrm dx$;
% 
% (4) $\displaystyle \int\frac{\mathrm dx}{1+\sin x}$;
% 
% (5) $\displaystyle \int x(1+x^2)^5\mathrm dx$;
% 
% (6) $\displaystyle \int\frac{x\mathrm dx}{\sqrt{1-x^2}}$;
% 
% (7) $\displaystyle \int\frac{x\mathrm dx}{4+x^4}$;
% 
% (8) $\displaystyle \int\frac{e^x}{1+e^x}\mathrm dx$;
% 
% (9) $\displaystyle \int\frac{\mathrm dx}{\sqrt x(1+x)}$;
% 
% (10) $\displaystyle \int\frac{\mathrm dx}{x\sqrt{x^2-1}}$;
% 
% (11) $\displaystyle \int\sin x\sin3x\,\mathrm dx$;
% 
% (12) $\displaystyle \int\sin^2x\cos^2x\,\mathrm dx$;
% 
% (13) $\displaystyle \int\frac{\sin2x}{\sqrt{1+\sin^2x}}\mathrm dx$;
% 
% (14) $\displaystyle \int\frac{\sin x\cos x}{\sqrt{a^2\sin^2x+b^2\cos^2x}}\mathrm dx\quad(a^2>b^2)$;
% 
% (15) $\displaystyle \int\frac{\sqrt{a^2-x^2}}x\mathrm dx$;
% 
% (16) $\displaystyle \int\frac{\mathrm dx}{x^2\sqrt{1+x^2}}$;
% 
% (17) $\displaystyle \int\frac{\mathrm dx}{\sqrt{(a^2+x^2)^3}}$;
% 
% (18) $\displaystyle \int\frac{x}{1+\sqrt x}\mathrm dx$;
% 
% (19) $\displaystyle \int\sqrt{\frac{2+x}{2-x}}\,\mathrm dx$;
% 
% (20) $\displaystyle \int\frac{\mathrm dx}{x\sqrt{4-x^2}}$.
% 
% 4。用分部积分法求下列不定积分:
% 
% (1) $\displaystyle \int\ln(1+x^2)\mathrm dx$;
% 
% (2) $\displaystyle \int xe^{-3x}\mathrm dx$;
% 
% (3) $\displaystyle \int x\cos nx\,\mathrm dx$;
% 
% (4) $\displaystyle \int\arcsin x\,\mathrm dx$;
% 
% (5) $\displaystyle \int\frac{x}{\sin^2x}\mathrm dx$;
% 
% (6) $\displaystyle \int\frac{\arctan x}{x^2(1+x^2)}\mathrm dx$;
% 
% (7) $\displaystyle \int\frac{e^{\arctan x}}{(1+x^2)^{3/2}}\mathrm dx$;
% 
% (8) $\displaystyle \int\frac{x^2}{(1+x^2)^2}\mathrm dx$;
% 
% (9) $\displaystyle \int\sin(\ln x)\mathrm dx$;
% 
% (10) $\displaystyle \int e^x\cos^2x\,\mathrm dx$.
% 
% 5。对下列不定积分建立递推公式:
% 
% (1) $\displaystyle I_n=\int x^m\ln^nx\,\mathrm dx$;
% 
% (2) $\displaystyle I_n=\int x^ne^{-x}\mathrm dx$;
% 
% (3) $\displaystyle I_n=\int\frac{x^n}{\sqrt{1-x^2}}\mathrm dx$;
% 
% (4) $\displaystyle I_n=\int\frac{\mathrm dx}{x^n\sqrt{1+x^2}}$.
% 
% 6。求下列有理函数的不定积分:
% 
% (1) $\displaystyle \int\frac{2x}{(x-1)(x+1)^2}\mathrm dx$;
% 
% (2) $\displaystyle \int\frac{2x^2+2x+13}{(x-2)(1+x^2)^2}\mathrm dx$;
% 
% (3) $\displaystyle \int\frac{x^2-x+1}{2x+1}\mathrm dx$;
% 
% (4) $\displaystyle \int\frac{3x+5}{(x^2+2x+2)^2}\mathrm dx$;
% 
% (5) $\displaystyle \int\frac{x+1}{x^3+2x^2-x-2}\mathrm dx$;
% 
% (6) $\displaystyle \int\frac{\mathrm dx}{x(x^3+2)}$;
% 
% (7) $\displaystyle \int\frac{\mathrm dx}{(x-2)^2(x+3)^3}$;
% 
% (8) $\displaystyle \int\frac{x\mathrm dx}{x^3-1}$.
% 
% 7。求下列三角函数有理式的不定积分:
% 
% (1) $\displaystyle \int\frac{\mathrm dx}{\sin^2x\cos x}$;
% 
% (2) $\displaystyle \int\frac{\mathrm dx}{\cos^2x\sin x}$;
% 
% (3) $\displaystyle \int\frac{\sin2x}{1+\cos^2x}\mathrm dx$;
% 
% (4) $\displaystyle \int\sin5x\sin3x\,\mathrm dx$;
% 
% (5) $\displaystyle \int\frac{\sin^2x}{1+\cos^2x}\mathrm dx$;
% 
% (6) $\displaystyle \int\frac{\sin^5x}{\cos^4x}\mathrm dx$;
% 
% (7) $\displaystyle \int\frac1{1+\cos x}\mathrm dx$;
% 
% (8) $\displaystyle \int\frac{1+\sin x}{1-\sin x}\mathrm dx$;
% 
% (9) $\displaystyle \int\frac{\sin x}{a\sin x+b\cos x}\mathrm dx$;
% 
% (10) $\displaystyle \int\frac{\sin x\mathrm dx}{\sin x-\cos x+2}$.
% 
% 8。求下列无理函数的不定积分:
% 
% (1) $\displaystyle \int\frac{\sqrt x}{\sqrt[4]{x^3}+1}\mathrm dx$;
% 
% (2) $\displaystyle \int\frac{\mathrm dx}{1+\sqrt x+\sqrt{x+1}}$;
% 
% (3) $\displaystyle \int\frac{\mathrm dx}{x^4\sqrt{1-x^2}}$;
% 
% (4) $\displaystyle \int\frac{\mathrm dx}{x\sqrt{x^2+x+1}}$;
% 
% (5) $\displaystyle \int\frac{\mathrm dx}{(x+1)\sqrt{x^2+4x+5}}$;
% 
% (6) $\displaystyle \int\frac{\mathrm dx}{x+\sqrt{x^2+x+1}}$;
% 
% (7) $\displaystyle \int\sqrt{\frac{1-x}{1+x}}\,\mathrm dx$;
% 
% (8) $\displaystyle \int\frac{\mathrm dx}{x\sqrt{4-x^2}}$;
% 
% (9) $\displaystyle \int\sqrt{\tan x}\,\mathrm dx$;
% 
% (10) $\displaystyle \int\sqrt[3]{\frac{2-x}{2+x}}\frac{\mathrm dx}{(2-x)^2}$.
% 
% % 整合来源：math-7.tex
% \chapter{定积分}
% \label{chap:ma-7}
% 
% \section{定积分的概念与微积分基本定理}
% \label{sec:ma-7-1}
% 
% \subsection{曲边梯形的面积}
% \label{subsec:ma-7-1-1}
% 
% 定积分最朴素的思想可以追溯到阿基米德时代, 阿基米德的 ``由直线组成平面图形, 由平面组成立体图形'' 的思想可以认为是定积分的萌芽。他曾借助圆柱和圆锥的体积并利用细分的方法求出了球的体积。但由于用细分的方法无法解决一般函数的定积分计算问题, 因此定积分的研究在漫长的岁月里一直没有进展。直到 17 世纪, 在天文学和物理学等学科的促进下, 数学研究才取得了很大的进展。伽利略和开普勒的一系列发现, 如开普勒的行星运动三定律等, 导致了导数的诞生。导数对现代数学理论的建立起了巨大的促进作用。到了 17 世纪下半叶, 牛顿和莱布尼茨两人在综合、发展前人工作的基础上, 几乎同时建立了微积分理论。正是这种理论, 使得定积分有了简便的计算方法, 从而也使得它具有极大的理论和实用价值.
% 
% 可以引入定积分概念的问题很多, 如开普勒研究的椭圆面积问题; 已知速度求路程的问题; 已知密度求质量的问题; 等等。下面我们以求由区间 $[a,b]$ 上非负函数 $f(x)$ 的图像(曲线)与 $x$ 轴及直线 $x=a,x=b$ 所围图形 $S$(如图 7.1.1, 我们称该图形为曲边梯形)的面积为例来详细阐述定积分的基本思想.
% 
% 回忆一下, 我们目前所能精确求出面积的图形只能是由有限条线段所围成的平面图形。如果 $y=f(x)$ 不是一个线性函数, 则不可能简单地用某个我们已知的计算公式算出该曲边梯形的面积。因此我们需要探究如何来求这种图形面积的方法.
% 
% 为了使问题简化, 我们假定 $f(x)$ 在区间 $[a,b]$ 上连续, 并记 $M,m$ 分别为它在区间 $[a,b]$ 上的最大值和最小值。在重积分理论中, 我们将讨论一般的平面图形的面积问题并证明曲边梯形 $S$ 的面积必定存在。现在我们假定该曲边梯形的面积存在并将其面积仍记为 $S$, 易于看出
% \[
%   m(b-a)\leq S\leq M(b-a).
% \]
% 
% 设函数 $y=f(x)(x\in[a,b])$ 的图像为 $T$。由连续函数的性质可以知道 $T$ 上必有一点 $(c,f(c))$, 使得 $S=f(c)(b-a)$。我们事先无法知道 $c$ 在何处, 能做的是在区间 $[a,b]$ 上取一点 $\xi$, 计算 $f(\xi)(b-a)$。$f(\xi)(b-a)$ 的几何意义很明显, 它是以区间 $[a,b]$ 为底, 以 $f(\xi)$ 为高的矩形的面积。它与 $S$ 的误差可以由下式估计:
% \[
%   |S-f(\xi)(b-a)|=|f(c)-f(\xi)|(b-a)\leq (M-m)(b-a),
% \]
% 其中 $M-m$ 也称为 $f(x)$ 在 $[a,b]$ 上的振幅.
% 
% 如果我们等分区间 $[a,b]$ 成两个小区间, 并令 $x_1=\dfrac{a+b}{2}$(见图 7.1.1), 则直线 $x=x_1$ 将 $S$ 分成两个曲边梯形 $S_1$ 及 $S_2$。现将它们的面积仍然分别记为 $S_1$ 和 $S_2$, 显然有
% \[
%   S=S_1+S_2.
% \]
% 在 $S_1$ 和 $S_2$ 中重复我们刚才在区间 $[a,b]$ 上所做的过程, 记 $M_1,M_2$ 分别为 $f(x)$ 在区间 $[a,x_1]$ 和 $[x_1,b]$ 上的最大值, $m_1,m_2$ 分别为 $f(x)$ 在区间 $[a,x_1]$ 和 $[x_1,b]$ 上的最小值。在区间 $[a,x_1]$ 上任意取 $\xi_1$, 在区间 $[x_1,b]$ 上任意取 $\xi_2$, 则有以下的估计式:
% \[
% \begin{aligned}
%   |f(\xi_1)(x_1-a)-S_1|&\leq (M_1-m_1)\frac{b-a}{2},\\
%   |f(\xi_2)(b-x_1)-S_2|&\leq (M_2-m_2)\frac{b-a}{2}.
% \end{aligned}
% \]
% 因此
% \[
% \begin{aligned}
%   &|f(\xi_1)(x_1-a)+f(\xi_2)(b-x_1)-S|\\
%   &\quad\leq [(M_1-m_1)+(M_2-m_2)]\frac{b-a}{2}\\
%   &\quad\leq (M-m)(b-a).
% \end{aligned}
% \]
% 
% 上述成立是因为分割后的误差估计由小区间上的振幅给出, 而一个函数在小区间上的振幅不会比在包含该区间的大区间上的振幅大的缘故。上式还告诉我们: 对于分成两块所得到的和式 $f(\xi_1)(x_1-a)+f(\xi_2)(b-x_1)$ 与 $S$ 的误差和原来 $f(\xi)(b-a)$ 与 $S$ 的误差, 分割后的误差的估计式不会变大。在许多情况下, 比如 $f(x)$ 在区间 $[a,b]$ 单调, 分割后任意取值得到的和式与 $S$ 的误差估计比分割前的误差估计会变得更小.
% 
% 当区间 $[a,b]$ 分割后的区间长度很小时, 由 $f(x)$ 的一致连续性, $f(x)$ 在每个小区间上的振幅都可以很小。基于这种思想, 我们将区间 $[a,b]$ 分割成更小的区间, 如对充分大的 $n$, 将区间 $[a,b]$ 等分为 $n$ 个小区间(见图 7.1.2)。设其分点为
% \[
%   a=x_0<x_1<x_2<\cdots<x_n=b,
% \]
% 记 $f(x)$ 在区间 $[x_{i-1},x_i](i=1,2,\cdots,n)$ 上的最大值为 $M_i$, 最小值为 $m_i$。在每个区间上任取 $\xi_i$, 作和式
% \[
%   \sum_{i=1}^n f(\xi_i)(x_i-x_{i-1});
% \]
% 则此和式与 $S$ 的误差可以有下面的估计:
% \begin{equation}
%   \left|\sum_{i=1}^n f(\xi_i)(x_i-x_{i-1})-S\right|
%   \leq \sum_{i=1}^n (M_i-m_i)(x_i-x_{i-1}).
%   \label{eq:ma-7-1-1}
% \end{equation}
% 上述不等式右边即为图 7.1.2 中与曲线相交的小矩形面积之和。由于 $f(x)$ 在区间 $[a,b]$ 上连续, 从而一致连续, 因此对于 $\forall\varepsilon>0$, $\exists\delta>0$, 当每个 $x_i-x_{i-1}<\delta$ 时(这只要 $n$ 充分大即可), 有 $M_i-m_i<\varepsilon$。于是我们有
% \[
%   \left|\sum_{i=1}^n f(\xi_i)(x_i-x_{i-1})-S\right|
%   \leq \sum_{i=1}^n (M_i-m_i)(x_i-x_{i-1})<\varepsilon(b-a).
% \]
% 
% 以上分析告诉我们, 只要 $f(x)$ 在区间 $[a,b]$ 上连续, 当区间 $[a,b]$ 分割得充分 ``细'' 时, 对每个小曲边梯形的曲边 ``以直代曲'', 作出的和式的值可以任意地接近该曲边梯形的面积。因此有
% \[
%   \lim_{n\to\infty}\sum_{i=1}^n f(\xi_i)(x_i-x_{i-1})=S.
% \]
% 
% 下面我们以由曲线 $f(x)=x^2(x\in[0,1])$ 与 $x$ 轴及直线 $x=1$ 所围的图形为例来求其面积 $S$(见图 7.1.3).
% 
% 将区间 $[0,1]$ 进行 $n$ 等分, 其分点为 $0,\dfrac1n,\dfrac2n,\cdots,\dfrac{n-1}{n},1$。在区间 $\left[\dfrac{i-1}{n},\dfrac{i}{n}\right](i=1,2,\cdots,n)$ 上任取 $\xi_i$, 作和式
% \[
%   \sum_{i=1}^n f(\xi_i)\left(\frac{i}{n}-\frac{i-1}{n}\right)
%   =\frac1n\sum_{i=1}^n \xi_i^2.
% \]
% 由于 $\dfrac{i-1}{n}\leq \xi_i\leq \dfrac{i}{n}(i=1,2,\cdots,n)$, 且 $f(x)$ 在区间 $[0,1]$ 单调上升, 因此有
% \[
%   m_i=\left(\frac{i-1}{n}\right)^2\leq \xi_i^2\leq \left(\frac{i}{n}\right)^2=M_i.
% \]
% 注意到
% \[
%   \lim_{n\to\infty}\frac1{n^3}\sum_{i=1}^n (i-1)^2
%   =\lim_{n\to\infty}\frac1{n^3}\sum_{i=1}^n i^2
%   =\frac13,
% \]
% 我们有
% \[
%   \left|\frac1n\sum_{i=1}^n \xi_i^2-S\right|
%   \leq \left|\frac1{n^3}\sum_{i=1}^n i^2-\frac1{n^3}\sum_{i=1}^n (i-1)^2\right|
%   \to 0\quad(n\to\infty).
% \]
% 因此有
% \[
%   S=\lim_{n\to\infty}\frac1n\sum_{i=1}^n \xi_i^2=\frac13.
% \]
% 
% 我们再来举一个质点做直线运动的例子。假定一个质点在直线上朝着一个方向运动, 其速度 $v=v(t)$ 是时间 $t$ 的函数。为简便计, 假定 $v(t)\geq 0$。现在我们来求在时间段 $[a,b]$ 内质点走过的路程 $s$.
% 
% 若 $v(t)=c$ 是常数函数, 则 $s=c(b-a)$。如果 $v(t)$ 不是常数函数, 即质点做非匀速运动时, 我们如何来求 $s$ 呢? 为此我们来分析一下, 假定 $v(t)$ 是连续变化的(在许多实际问题中大都是如此), 则对任意 $t_0\in[a,b]$, $v(t)$ 在 $t_0$ 附近的值与 $v(t_0)$ 很接近。换句话说, 对很小的 $\delta>0$, 在时间段 $[t_0-\delta,t_0+\delta]$ 内, 质点走过的路程 $\tilde s\approx 2\delta v(t_0)$(简称 ``以不变代变'')。这就启发我们将区间 $[a,b]$ 作分割:
% \[
%   a=t_0<t_1<\cdots<t_n=b,
% \]
% 当 $\lambda=\max\limits_{1\leq i\leq n}\{t_i-t_{i-1}\}$ 很小时, 在每个时间段 $[t_{i-1},t_i](i=1,2,\cdots,n)$ 内质点走过的路程 $s_i\approx v(\xi_i)(t_i-t_{i-1})$, 其中 $\xi_i\in[t_{i-1},t_i]$。因此, 总的路程
% \begin{equation}
%   s=\sum_{i=1}^n s_i\approx \sum_{i=1}^n v(\xi_i)(t_i-t_{i-1}).
%   \label{eq:ma-7-1-2}
% \end{equation}
% 当时间段 $[a,b]$ 的分割对应的小时间段的长度趋于零时, 上面路程的近似值与路程的准确值无限地接近。因此若和式 \eqref{eq:ma-7-1-2} 有极限, 此极限便是我们要求的路程。容易看出以上已知速度函数求路程与我们在前面求曲边梯形的面积具有相同的 ``纯数学'' 极限过程.
% 
% \subsection{定积分的定义}
% \label{subsec:ma-7-1-2}
% 
% 上面我们对若干问题进行了讨论, 在自然界还有许多类似的问题, 在解决它们的过程中可以抽象出一个相同的 ``纯数学'' 极限过程。为此我们给出下面的定义.
% 
% \begin{definition}\label{definition:ma-7-1-1}
% 设函数 $f(x)$ 在区间 $[a,b]$ 上有定义, 对于区间 $[a,b]$ 的一个分割
% \[
%   \Delta:a=x_0<x_1<\cdots<x_n=b,
% \]
% 记 $\Delta x_i=x_i-x_{i-1}(i=1,2,\cdots,n)$, $\lambda(\Delta)=\max\limits_{1\leq i\leq n}\{\Delta x_i\}$。在每个小区间 $[x_{i-1},x_i](i=1,2,\cdots,n)$ 上任取 $\xi_i$, 作和式
% \[
%   \sum_{i=1}^n f(\xi_i)\Delta x_i.
% \]
% 如果当 $\lambda(\Delta)\to0$ 时, 上述和式存在极限 $I$, 且 $I$ 不依赖分割 $\Delta$ 的选取及 $\xi_i(i=1,2,\cdots,n)$ 在 $[x_{i-1},x_i]$ 上的选取, 则称 $f(x)$ 在区间 $[a,b]$ 上是黎曼可积(简称可积)的, 同时称 $I$ 为 $f(x)$ 在区间 $[a,b]$ 上的定积分, 记为
% \[
%   I=\int_a^b f(x)\mathrm dx,
% \]
% 其中 $a$ 与 $b$ 分别称为定积分的下限和上限, $f(x)$ 称为被积函数, $x$ 称为积分变量.
% \end{definition}
% 
% 在定积分理论的讨论中, 我们经常要遇到分割求和。为了使行文简洁, 在今后对每个给定的区间 $[a,b]$, 我们总是用
% \[
%   \Delta:a=x_0<x_1<\cdots<x_n=b
% \]
% 来表示它的一个分割, 对于这个分割总是记 $\Delta x_i=x_i-x_{i-1}(i=1,2,\cdots,n)$, 并且记 $\lambda(\Delta)=\max\limits_{1\leq i\leq n}\{\Delta x_i\}$。在 $[x_{i-1},x_i](i=1,2,\cdots,n)$ 上任取一点 $\xi_i$, 我们称和式 $\sum\limits_{i=1}^n f(\xi_i)\Delta x_i$ 为函数 $f(x)$ 关于 $\Delta$ 的黎曼和.
% 
% 用 ``$\varepsilon$-$\delta$'' 语言叙述定积分的定义则更为精确.
% 
% \begin{definition}\label{definition:ma-7-1-1-prime}
% 设函数 $f(x)$ 在区间 $[a,b]$ 上有定义, 若存在常数 $I\in\mathbb R$, 使得对于 $\forall\varepsilon>0$, $\exists\delta>0$, 对区间 $[a,b]$ 的任何一个分割 $\Delta$, 当 $\lambda(\Delta)<\delta$ 时, 在每个 $[x_{i-1},x_i](i=1,2,\cdots,n)$ 上任取 $\xi_i$, 都有
% \[
%   \left|\sum_{i=1}^n f(\xi_i)\Delta x_i-I\right|<\varepsilon,
% \]
% 则称函数 $f(x)$ 在区间 $[a,b]$ 是黎曼可积的, 并称 $I$ 为函数 $f(x)$ 在区间 $[a,b]$ 上的定积分.
% \end{definition}
% 
% 今后, 我们用记号 $f(x)\in R[a,b]$ 表示函数 $f(x)$ 在区间 $[a,b]$ 上可积.
% 
% 从定义~\ref{definition:ma-7-1-1}(或定义~\ref{definition:ma-7-1-1-prime}) 可以看出: 若函数 $f(x)$ 在区间 $[a,b]$ 上可积, 则它在区间 $[a,b]$ 上的定积分 $\int_a^b f(x)\mathrm dx$ 是一种特殊和式的极限。因此, 和式写成 $\sum\limits_{i=1}^n f(\xi_i)\Delta x_i$ 与 $\sum\limits_{i=1}^n f(\tau_i)\Delta t_i$ 并不影响和式的极限, 此和式是由函数关系 $f$ 决定的, 与自变量取 $x$ 还是取 $t$ 无关, 即
% \[
%   \int_a^b f(x)\mathrm dx=\int_a^b f(t)\mathrm dt.
% \]
% 
% 值得注意的是, 上述和式的极限与我们所学过的序列极限和函数极限是不同的。在 $\lambda(\Delta)\to0$ 的过程中, 这个极限过程具有两个任意性, 即分割的任意性及在每个小区间上 $\xi_i$ 选取的任意性。尽管这些极限过程有所不同, 但定积分定义中和式的极限与序列极限或函数极限仍具有一些相似的性质, 如一个函数的定积分存在, 则它必定唯一。这些性质的证明可以由定义直接推出, 读者可自证之.
% 
% \subsection{定积分的几何意义}
% \label{subsec:ma-7-1-3}
% 
% 设函数 $f(x)$ 在区间 $[a,b]$ 上连续。我们先看 $f(x)\geq0$ 的情形。根据我们前面的讨论可知: 定积分 $\int_a^b f(x)\mathrm dx$ 是由直线 $x=a,x=b,y=0$ 与曲线 $y=f(x)$ 所围成的曲边梯形的面积 $S$.
% 
% 当 $f(x)\leq0$ 时, 同样地设由直线 $x=a,x=b,y=0$ 与曲线 $y=f(x)$ 所围成的曲边梯形的面积为 $S$, 则此时
% \[
%   S=-\int_a^b f(x)\mathrm dx.
% \]
% 当 $f(x)$ 有正有负时, 则 $\int_a^b f(x)\mathrm dx$ 是由直线 $x=a,x=b,y=0$ 与 $y=f(x)$ 所围成的几个曲边梯形中, 位于 $x$ 轴上方的各个曲边梯形面积之和减去位于 $x$ 轴下方的各个曲边梯形面积之和。如图 7.1.4, 则有
% \[
%   \int_a^b f(x)\mathrm dx=S_1-S_2+S_3-S_4.
% \]
% 
% \subsection{连续函数的可积性}
% \label{subsec:ma-7-1-4}
% 
% 在对曲边梯形面积的讨论中, 如果假定该面积存在, 我们找到了一种求出该面积的方法。有了这一基础, 对一个连续函数 $y=f(x),x\in[a,b]$, 我们比较容易证明它在区间 $[a,b]$ 上必定是可积的.
% 
% \begin{theorem}\label{theorem:ma-7-1-1}
% 设函数 $f(x)\in C[a,b]$, 则 $f(x)\in R[a,b]$.
% \end{theorem}
% 
% 
% \begin{proof}
% 记 $m,M$ 分别为函数 $f(x)$ 在区间 $[a,b]$ 上的最小和最大值。我们注意到, 对区间 $[a,b]$ 的任意分割 $\Delta$, $f(x)$ 关于 $\Delta$ 的任一黎曼和满足
% \begin{equation}
%   m(b-a)\leq \sum_{i=1}^n m_i\Delta x_i
%   \leq \sum_{i=1}^n f(\xi_i)\Delta x_i
%   \leq \sum_{i=1}^n M_i\Delta x_i
%   \leq M(b-a),
%   \label{eq:ma-7-1-3}
% \end{equation}
% 其中 $m_i$ 与 $M_i$ 分别为 $f(x)$ 在区间 $[x_{i-1},x_i](i=1,2,\cdots,n)$ 上的最小值和最大值。特别地, 对每一个 $n\in\mathbb N$, 将区间 $[a,b]$ 的 $n$ 等分的分割记为 $\Delta_n$, 在此分割的每个小区间 $[x_{j-1},x_j](j=1,2,\cdots,n)$ 上取定 $\xi_j=x_j$, 我们可得一个特殊的黎曼和
% \[
%   S_n=\sum_{j=1}^n f(x_j)\frac{b-a}{n}.
% \]
% 这样一来, 我们便得到了一个序列 $\{S_n\}$, 由 \eqref{eq:ma-7-1-3} 知它是一个有界序列, 因此存在收敛子列 $\{S_{n_k}\}\subset\{S_n\}$。设 $\lim\limits_{k\to\infty}S_{n_k}=S$, 下面我们证明 $S$ 即为函数 $f(x)$ 在区间 $[a,b]$ 上的定积分.
% 
% 对于 $\forall\varepsilon>0$, 由函数 $f(x)$ 在区间 $[a,b]$ 上一致连续知, $\exists\delta>0$, 当 $x',x''\in[a,b]$ 且 $|x'-x''|<\delta$ 时, 有
% \[
%   |f(x')-f(x'')|<\frac{\varepsilon}{4(b-a)},
% \]
% 从而对任意的分割 $\Delta$, 当 $\lambda(\Delta)<\delta$ 时, 有
% \[
%   \sum_{i=1}^n (M_i-m_i)\Delta x_i<\frac{\varepsilon}{4}.
% \]
% 由于 $\lim\limits_{k\to\infty}S_{n_k}=S$, 因此存在 $K>0$, 使得当 $n_k>K$ 时, 有 $\lambda(\Delta_{n_k})<\delta$ 和 $|S_{n_k}-S|<\dfrac{\varepsilon}{2}$ 同时成立。现取定一个 $n_{k_0}>K$, 使得它同时满足这两个条件, 即 $\lambda(\Delta_{n_{k_0}})<\delta$ 和 $|S_{n_{k_0}}-S|<\dfrac{\varepsilon}{2}$.
% 
% 对于任意的分割 $\Delta$, 当 $\lambda(\Delta)<\delta$ 时, 我们有以下的估计:
% \begin{equation}
% \begin{aligned}
%   \left|\sum_{i=1}^n f(\xi_i)\Delta x_i-S\right|
%   &\leq \left|\sum_{i=1}^n f(\xi_i)\Delta x_i-S_{n_{k_0}}\right|+|S_{n_{k_0}}-S|\\
%   &< \left|\sum_{i=1}^n f(\xi_i)\Delta x_i-\sum_{j=1}^{n_{k_0}} f(x_j)\frac{b-a}{n_{k_0}}\right|+\frac{\varepsilon}{2}.
% \end{aligned}
% \label{eq:ma-7-1-4}
% \end{equation}
% 为了估计上述不等式右边两个黎曼和的差, 记 $\Delta'$ 为 $\Delta$ 与 $\Delta_{n_{k_0}}$ 的分点集的并所构成的区间 $[a,b]$ 的分割, 显然有 $\lambda(\Delta')<\delta$。容易看出对关于 $\Delta',\Delta,\Delta_{n_{k_0}}$ 的任意黎曼和, 有下述估计:
% \[
%   \left|\sum_{i=1}^n f(\xi_i)\Delta x_i-\sum_{k=1}^{n'}f(\xi_k')\Delta x_k'\right|
%   \leq \sum_{i=1}^n(M_i-m_i)\Delta x_i<\frac{\varepsilon}{4}
% \]
% 和
% \[
%   \left|\sum_{j=1}^{n_{k_0}}f(x_j)\frac{b-a}{n_{k_0}}-\sum_{k=1}^{n'}f(\xi_k')\Delta x_k'\right|
%   \leq \sum_{j=1}^{n_{k_0}}(M_j'-m_j')\frac{b-a}{n_{k_0}}<\frac{\varepsilon}{4},
% \]
% 其中 $\sum\limits_{k=1}^{n'}f(\xi_k')\Delta x_k'$ 是关于 $\Delta'$ 的任意一个黎曼和, $\sum\limits_{i=1}^n(M_i-m_i)\Delta x_i$ 是关于分割 $\Delta$ 求和, $\sum\limits_{j=1}^{n_{k_0}}(M_j'-m_j')\dfrac{b-a}{n_{k_0}}$ 是关于分割 $\Delta_{n_{k_0}}$ 求和。因此
% \[
% \begin{aligned}
%   \left|\sum_{i=1}^n f(\xi_i)\Delta x_i-\sum_{j=1}^{n_{k_0}}f(x_j)\frac{b-a}{n_{k_0}}\right|
%   &\leq \left|\sum_{i=1}^n f(\xi_i)\Delta x_i-\sum_{k=1}^{n'}f(\xi_k')\Delta x_k'\right|\\
%   &\quad+\left|\sum_{j=1}^{n_{k_0}}f(x_j)\frac{b-a}{n_{k_0}}-\sum_{k=1}^{n'}f(\xi_k')\Delta x_k'\right|\\
%   &<\frac{\varepsilon}{4}+\frac{\varepsilon}{4}=\frac{\varepsilon}{2}.
% \end{aligned}
% \]
% 将此估计代入 \eqref{eq:ma-7-1-4} 式即得
% \[
%   \left|\sum_{i=1}^n f(\xi_i)\Delta x_i-S\right|<\varepsilon.
% \]
% 这就证明了函数 $f(x)$ 在区间 $[a,b]$ 上可积, 并且 $\int_a^b f(x)\mathrm dx=S$.
% \end{proof}
% 
% 
% \subsection{微积分基本定理}
% \label{subsec:ma-7-1-5}
% 
% 从定积分的定义可知, 定积分应在许多问题中有重要应用, 因此如何计算一个函数的定积分是一个十分关键的问题。如果对每个函数都要从定积分的定义出发, 对区间分割, 在分割后的每个小区间任取一点作和式, 再求极限, 当函数稍微复杂一点, 求定积分的值就会面临巨大的困难。如果没有一种相对简单的定积分计算方法, 定积分或许不会成为一种有用的理论和工具.
% 
% 如何发现定积分简便的计算方法呢? 我们可以先考查一下以下具体问题。假定我们已知一个质点做直线运动, 其速度函数为 $v(t)>0(t\in[t_0,t_1])$, 现在来求它在该时间段内走过的路程。前面我们已经知道, 若速度函数 $v(t)$ 的性质比较好(如 $v(t)$ 是连续的), $v(t)$ 在 $[t_0,t_1]$ 上的定积分
% \[
%   \int_{t_0}^{t_1}v(t)\mathrm dt
% \]
% 刚好是质点在该时间段内走过的路程。从另外一方面来看, 该段路程刚好是路程函数 $s(t)$ 在 $t=t_1$ 和 $t=t_0$ 两点的值之差。换句话说, 我们有
% \[
%   \int_{t_0}^{t_1}v(t)\mathrm dt=s(t_1)-s(t_0).
% \]
% 在微分学理论中, 我们已经知道了 $s'(t)=v(t)$.
% 
% 那么, 若一个函数 $f(x)$ 在区间 $[a,b]$ 上可积, 且存在原函数 $F(x)$, 是否总是成立
% \[
%   \int_a^b f(x)\mathrm dx=F(b)-F(a)
% \]
% 呢? 这个问题已由牛顿和莱布尼茨解决。因此人们称下面定理给出的计算定积分的公式为牛顿-莱布尼茨公式。由于它是微积分理论中最重要的结果, 人们也称它为微积分基本定理.
% 
% \begin{theorem}[{微积分基本定理}]\label{theorem:ma-7-1-2}
% 设函数 $f(x)$ 在区间 $[a,b]$ 上有定义, 并且满足以下两个条件:
% 
% (1) 在区间 $[a,b]$ 上可积;
% 
% (2) 在区间 $[a,b]$ 上存在原函数 $F(x)$,
% 
% 则有
% \begin{equation}
%   \int_a^b f(x)\mathrm dx=F(b)-F(a)\triangleq F(x)\big|_a^b.
%   \label{eq:ma-7-1-5}
% \end{equation}
% \end{theorem}
% 
% 
% \begin{proof}
% 首先, 由于函数 $f(x)$ 在区间 $[a,b]$ 上可积, 任取区间 $[a,b]$ 的分割
% \[
%   \Delta:a=x_0<x_1<\cdots<x_n=b,
% \]
% 在 $[x_{i-1},x_i](i=1,2,\cdots,n)$ 上任取一点 $\xi_i$, 则有
% \[
%   \lim_{\lambda(\Delta)\to0}\sum_{i=1}^n f(\xi_i)\Delta x_i=\int_a^b f(x)\mathrm dx.
% \]
% 其次, 对于分割 $\Delta$, 我们有
% \[
%   F(b)-F(a)=\sum_{i=1}^n[F(x_i)-F(x_{i-1})].
% \]
% 在 $[x_{i-1},x_i](i=1,2,\cdots,n)$ 上对 $F(x)$ 应用拉格朗日微分中值定理得
% \[
%   F(x_i)-F(x_{i-1})=f(\xi_i')\Delta x_i,
% \]
% 其中 $\xi_i'\in(x_{i-1},x_i)$。因此我们有
% \[
% \begin{aligned}
%   F(b)-F(a)
%   &=\lim_{\lambda(\Delta)\to0}\sum_{i=1}^n[F(x_i)-F(x_{i-1})]\\
%   &=\lim_{\lambda(\Delta)\to0}\sum_{i=1}^n f(\xi_i')\Delta x_i\\
%   &=\int_a^b f(x)\mathrm dx.
% \end{aligned}
% \qedhere
% \]
% \end{proof}
% 
% 
% 利用牛顿-莱布尼茨公式可以使许多定积分的计算变得十分简单.
% 
% 例如前面我们利用分割、求和以及计算极限求出了 $\int_0^1x^2\mathrm dx=\dfrac13$。由于 $f(x)=x^2\in C[0,1]$, 由定理~\ref{theorem:ma-7-1-1}, 我们知 $f(x)=x^2\in R[0,1]$。在不定积分中, 我们知道 $F(x)=\dfrac{x^3}{3}$ 是函数 $f(x)$ 的一个原函数, 现利用牛顿-莱布尼茨公式即有
% \[
%   \int_0^1 x^2\mathrm dx=\left.\frac{x^3}{3}\right|_0^1=\frac13.
% \]
% 
% \begin{example}\label{example:ma-7-1-1}
% 设
% \[
%   I_n=\frac1{n^{\alpha+1}}(1+2^\alpha+3^\alpha+\cdots+n^\alpha)\quad(\alpha>0),
% \]
% 求 $\lim\limits_{n\to\infty}I_n$.
% \end{example}
% 
% 
% \begin{solve}
% 由于
% \[
%   I_n=\frac1n\left(\frac1{n^\alpha}+\frac{2^\alpha}{n^\alpha}+\frac{3^\alpha}{n^\alpha}+\cdots+\frac{n^\alpha}{n^\alpha}\right)
%   =\frac1n\sum_{i=1}^n\frac{i^\alpha}{n^\alpha},
% \]
% 因此 $I_n$ 可以看成是函数 $x^\alpha$ 在区间 $[0,1]$ 上关于分割
% \[
%   0<\frac1n<\frac2n<\cdots<\frac nn=1
% \]
% 的一个黎曼和。由 $x^\alpha$ 在 $[0,1]$ 上连续, 知它在 $[0,1]$ 上可积。在不定积分中, 我们已知 $\dfrac1{\alpha+1}x^{\alpha+1}$ 是 $x^\alpha$ 的一个原函数, 因此由牛顿-莱布尼茨公式得
% \[
%   \lim_{n\to\infty}I_n=\int_0^1 x^\alpha\mathrm dx
%   =\left.\frac1{\alpha+1}x^{\alpha+1}\right|_0^1
%   =\frac1{\alpha+1}.
% \qedhere
% \]
% \end{solve}
% 
% 
% \section{可积性问题}
% \label{sec:ma-7-2}
% 
% 我们已经指出, 在微积分的理论中, 微积分基本定理是最重要的定理。从该定理可以看出, 在定积分理论中我们要解决的两个基本问题是:
% 
% (1) 函数的可积性;
% 
% (2) 原函数的存在性.
% 
% 在以下各节中, 我们将逐步展开来讨论这两个问题。在本节中, 我们将讨论函数的可积性问题。尽管我们已经证明了连续函数的可积性, 但在许多数学理论中我们还需要讨论更为广泛的可积函数类。事实上, 正是在研究更为广泛的函数的可积性时, 人们创造了新的积分理论——勒贝格 (Lebesgue) 积分理论。关于这类积分, 我们将在实变函数论这门课程中加以仔细讨论.
% 
% 要解决满足什么样条件的函数是可积的这个基本的问题, 我们必须从定积分的定义着手。在定积分的定义中, 分割是任意的, 对于给定的分割, 在每个小区间上的取值也是任意的。因此, 在函数的可积性的研究中, 我们自然有以下两个问题: 在定积分的定义中,
% 
% (a) 是否可取特殊的分割来代替任意的分割?
% 
% (b) 能否通过特殊的取值来作和式?
% 
% 通过本节的讨论, 我们将回答这两个问题.
% 
% \subsection{可积的必要条件}
% \label{subsec:ma-7-2-1}
% 
% 首先我们观察到, 如果一个函数 $f(x)$ 在区间 $[a,b]$ 上无界, 则对于任意的分割
% \[
%   \Delta:a=x_0<x_1<\cdots<x_n=b,
% \]
% 必存在 $1\leq i_0\leq n$, 使得 $f(x)$ 在 $[x_{i_0-1},x_{i_0}]$ 上无界。对每个 $i(1\leq i\leq n$ 且 $i\ne i_0)$, 在 $[x_{i-1},x_i]$ 上取定 $\xi_i$, 作和式
% \[
%   \sum_{i=1,i\ne i_0}^n f(\xi_i)\Delta x_i,
% \]
% 则上述和式在 $\xi_i$ 取定后是一个定数。记
% \[
%   M_1=\sum_{i=1,i\ne i_0}^n f(\xi_i)\Delta x_i.
% \]
% 对任意的 $M>0$, 由于 $f(x)$ 在 $[x_{i_0-1},x_{i_0}]$ 上无界, 因此在 $[x_{i_0-1},x_{i_0}]$ 上可取 $\xi_{i_0}$, 使得
% \[
%   |f(\xi_{i_0})|>\frac{M+|M_1|}{\Delta x_{i_0}},
% \]
% 从而有
% \[
% \begin{aligned}
%   \left|\sum_{i=1}^n f(\xi_i)\Delta x_i\right|
%   &=\left|\sum_{i=1,i\ne i_0}^n f(\xi_i)\Delta x_i+f(\xi_{i_0})\Delta x_{i_0}\right|\\
%   &\geq |f(\xi_{i_0})|\Delta x_{i_0}-|M_1|\\
%   &>M.
% \end{aligned}
% \]
% 由于 $M$ 是任意给定的正数, 因此当 $\lambda(\Delta)\to0$ 时, $\sum\limits_{i=1}^n f(\xi_i)\Delta x_i$ 不可能有极限。所以我们证明了下面的定理.
% 
% \begin{theorem}\label{theorem:ma-7-2-1}
% 若函数 $f(x)\in R[a,b]$, 则 $f(x)$ 在区间 $[a,b]$ 上有界.
% \end{theorem}
% 
% 在下面的可积性讨论中, 我们总是假定函数 $f(x)$ 是区间 $[a,b]$ 上的有界函数。一个自然的问题是: 有界函数是否一定可积呢? 此答案是否定的, 如下例所示.
% 
% \begin{example}\label{example:ma-7-2-1}
% 证明狄利克雷函数
% \[
%   y=D(x)=
%   \begin{cases}
%     1, & x\in\mathbb Q,\\
%     0, & x\in\mathbb R\setminus\mathbb Q
%   \end{cases}
% \]
% 在任何区间 $[a,b]$ 上不可积.
% \end{example}
% 
% 
% \begin{proof}
% 对区间 $[a,b]$ 的任意分割 $\Delta:a=x_0<x_1<\cdots<x_n=b$, 如果在 $[x_{i-1},x_i](i=1,2,\cdots,n)$ 上取一点 $\xi_i(i=1,2,\cdots,n)$, 且 $\xi_i$ 为有理数, 则
% \[
%   \sum_{i=1}^n D(\xi_i)\Delta x_i=1;
% \]
% 若取 $\xi_i$ 为无理数, 则
% \[
%   \sum_{i=1}^n D(\xi_i)\Delta x_i=0.
% \]
% 因此, 当 $\lambda(\Delta)\to0$ 时, $\sum\limits_{i=1}^n D(\xi_i)\Delta x_i$ 不可能存在极限。这说明 $D(x)\notin R[a,b]$.
% 
% 例~\ref{example:ma-7-2-1} 实际上回答了我们在本节开始时提出的关于可积性的第二个问题(问题 (b)), 即在定积分的定义中, 我们不能通过取特殊值的黎曼和来确定函数的可积性.
% \end{proof}
% 
% 
% \subsection{达布理论}
% \label{subsec:ma-7-2-2}
% 
% 关于可积性的讨论有点类似于序列或函数的上、下极限的讨论。对区间 $[a,b]$ 上的有界函数 $f(x)$ 和 $[a,b]$ 的一个分割 $\Delta:a=x_0<x_1<\cdots<x_n=b$, 由于 $f(x)$ 在 $[x_{i-1},x_i](i=1,2,\cdots,n)$ 上取值的任意性, 使得其黎曼和具有一定的任意性。我们注意到: 如果 $f(x)$ 在 $[a,b]$ 上可积, 由定积分的定义可知, 在小区间上任意取值的黎曼和当 $\lambda(\Delta)\to0$ 时有极限, 那么, 对 $\lambda(\Delta)$ 很小时的每个分割 $\Delta$, 所有可能的黎曼和的上确界和下确界必相差不大。显然 $f(x)$ 在每个小区间上的上确界和下确界只与分割 $\Delta$ 有关, 而与 $\xi_i$ 的取法无关。为此, 对于分割 $\Delta$ 及 $i=1,2,\cdots,n$, 令
% \[
%   M_i=\sup_{x\in[x_{i-1},x_i]}\{f(x)\},\qquad
%   m_i=\inf_{x\in[x_{i-1},x_i]}\{f(x)\}.
% \]
% 我们作如下的和式:
% \[
%   \overline S(\Delta)=\sum_{i=1}^n M_i\Delta x_i,\qquad
%   \underline S(\Delta)=\sum_{i=1}^n m_i\Delta x_i.
% \]
% $\overline S(\Delta),\underline S(\Delta)$ 分别称为 $f(x)$ 关于分割 $\Delta$ 的达布 (Darboux) 大和与达布小和.
% 
% 由于达布大和与达布小和仅仅依赖分割的选取, 对达布大和与达布小和的研究可以避免因 $f(x)$ 在 $[x_{i-1},x_i](i=1,2,\cdots,n)$ 上取值的任意性而使得其黎曼和产生任意性。如果记 $S(\Delta)=\sum\limits_{i=1}^n f(\xi_i)\Delta x_i$ 为 $f(x)$ 关于分割 $\Delta$ 的任意一个黎曼和, 则显然有
% \[
%   \underline S(\Delta)\leq S(\Delta)\leq \overline S(\Delta).
% \]
% 另外不难看出, 当 $f(x)\in C[a,b]$ 时, $\underline S(\Delta),\overline S(\Delta)$ 均是黎曼和.
% 
% 下面我们来研究达布大和与达布小和的一些性质。若 $\Delta'$ 和 $\Delta''$ 都是区间 $[a,b]$ 的分割, 并且 $\Delta'$ 中的分点均是 $\Delta''$ 中的分点, 我们则称 $\Delta''$ 是 $\Delta'$ 的细分, 记为 $\Delta'\subset\Delta''$.
% 
% 以下我们总假定函数 $f(x)$ 在区间 $[a,b]$ 上有界, 涉及的分割均是 $[a,b]$ 的分割, 并且所有的黎曼和都是关于 $f(x)$ 的。对一个给定的分割 $\Delta$, 我们记
% \[
%   \omega_i=M_i-m_i(i=1,2,\cdots,n).
% \]
% $\omega_i$ 通常称为 $f(x)$ 在 $[x_{i-1},x_i](i=1,2,\cdots,n)$ 上的振幅。不难证明
% \[
%   \omega_i=\sup_{x',x''\in[x_{i-1},x_i]}\{|f(x')-f(x'')|\}.
% \]
% 
% 达布大和与达布小和具有以下性质.
% 
% \begin{property}\label{property:ma-7-2-2}
% 设 $\Delta'$ 与 $\Delta''$ 为区间 $[a,b]$ 的两个分割, 若 $\Delta'\subset\Delta''$, 则有 $\overline S(\Delta')\geq\overline S(\Delta'')$ 和 $\underline S(\Delta')\leq\underline S(\Delta'')$, 即对于分割的细分, 达布大和不增, 达布小和不减.
% \end{property}
% 
% 
% \begin{proof}
% 设
% \[
%   \Delta':a=x_0<x_1<\cdots<x_n=b.
% \]
% 我们先假定
% \[
%   \Delta'':a=x_0<x_1<\cdots<x_{i-1}<x'<x_i<\cdots<x_n=b,
% \]
% 即 $\Delta''$ 只是在 $\Delta'$ 中加一个分点 $x'$ 所成。记
% \[
%   M_i=\sup_{x\in[x_{i-1},x_i]}\{f(x)\},\quad
%   M_i'=\sup_{x\in[x_{i-1},x']}\{f(x)\},\quad
%   M_i''=\sup_{x\in[x',x_i]}\{f(x)\},
% \]
% 则有 $M_i'\leq M_i$ 和 $M_i''\leq M_i$, 因此有
% \[
% \begin{aligned}
%   \overline S(\Delta')-\overline S(\Delta'')
%   &=M_i(x_i-x_{i-1})-[M_i'(x'-x_{i-1})+M_i''(x_i-x')]\\
%   &\geq0.
% \end{aligned}
% \]
% 对于 $\Delta'$ 中加入多个分点的情形可以比较逐个每次加入一个分点的分割的达布大和并利用上述证明的结论得到。对于达布小和的情形可类似讨论.
% \end{proof}
% 
% 
% \begin{property}\label{property:ma-7-2-3}
% 设 $\Delta'$ 和 $\Delta''$ 是区间 $[a,b]$ 的任意两个分割, 则有 $\underline S(\Delta')\leq\overline S(\Delta'')$, 即任意分割的达布小和不大于任意分割的达布大和.
% \end{property}
% 
% 
% \begin{proof}
% 将 $\Delta',\Delta''$ 中的分点合并在一起构成区间 $[a,b]$ 的一个新的分割, 记其为 $\Delta=\Delta'\cup\Delta''$。由性质~\ref{property:ma-7-2-2} 得
% \[
%   \underline S(\Delta')\leq \underline S(\Delta)\leq \overline S(\Delta)\leq \overline S(\Delta'').
% \qedhere
% \]
% \end{proof}
% 
% 
% 现记
% \[
%   m=\inf_{x\in[a,b]}\{f(x)\},\qquad M=\sup_{x\in[a,b]}\{f(x)\},
% \]
% 则对任意的分割 $\Delta$, 有
% \[
%   m(b-a)\leq \underline S(\Delta)\leq \overline S(\Delta)\leq M(b-a).
% \]
% 当分割越来越细时, 性质~\ref{property:ma-7-2-2} 告诉我们达布大和与达布小和在某种意义下具有单调性。因此类似于单调序列极限的性质, 定义
% \[
%   \underline{\int_a^b} f(x)\mathrm dx=\sup_\Delta\{\underline S(\Delta)\}
%   \quad\text{和}\quad
%   \overline{\int_a^b} f(x)\mathrm dx=\inf_\Delta\{\overline S(\Delta)\}.
% \]
% 我们分别称 $\underline{\int_a^b} f(x)\mathrm dx$ 和 $\overline{\int_a^b} f(x)\mathrm dx$ 为函数 $f(x)$ 在区间 $[a,b]$ 上的下积分和上积分.
% 
% \begin{theorem}[{达布定理}]\label{theorem:ma-7-2-4}
% 设函数 $f(x)$ 在区间 $[a,b]$ 上有界, 则
% \[
%   \lim_{\lambda(\Delta)\to0}\overline S(\Delta)=\overline{\int_a^b} f(x)\mathrm dx,\qquad
%   \lim_{\lambda(\Delta)\to0}\underline S(\Delta)=\underline{\int_a^b} f(x)\mathrm dx.
% \]
% \end{theorem}
% 
% 
% \begin{proof}
% 只证上积分情形, 关于下积分的证明留给读者.
% 
% 对于 $\forall\varepsilon>0$, 由上积分的定义, 存在一个分割 $\Delta_1:a=x_0'<x_1'<\cdots<x_N'=b(N>1)$, 使得下式成立:
% \[
%   0\leq \overline S(\Delta_1)-\overline{\int_a^b} f(x)\mathrm dx<\frac{\varepsilon}{2}.
% \]
% 对任意分割 $\Delta:a=x_0<x_1<\cdots<x_n=b$, 由上积分的定义有
% \[
%   0\leq \overline S(\Delta)-\overline{\int_a^b} f(x)\mathrm dx.
% \]
% 我们记 $\Delta^*=\Delta\cup\Delta_1$, 并考查分割 $\Delta$ 中的每一个小区间 $[x_{i-1},x_i](i=1,2,\cdots,n)$。如果 $(x_{i-1},x_i)$ 中无 $\Delta_1$ 中的分点, 则在 $\overline S(\Delta)$ 与 $\overline S(\Delta^*)$ 中的相应项同为 $M_i\Delta x_i$。现在设 $(x_{i-1},x_i)$ 中含有 $\Delta_1$ 中的分点。记 $\delta_1$ 为 $\Delta_1$ 中 $N$ 个小区间的最小长度, 若预先取 $\lambda(\Delta)<\delta_1$, 则在 $(x_{i-1},x_i)$ 中只有 $\Delta_1$ 中的一个分点 $x'$。因此 $\overline S(\Delta)$ 与 $\overline S(\Delta^*)$ 中相应项之差为
% \[
%   M_i(x_i-x_{i-1})-[M_i'(x'-x_{i-1})+M_i''(x_i-x')]
%   \leq (M-m)(x_i-x_{i-1})\leq (M-m)\lambda(\Delta),
% \]
% 其中
% \[
%   M_i=\sup_{x\in[x_{i-1},x_i]}\{f(x)\},
% \]
% \[
%   M_i'=\sup_{x\in[x_{i-1},x']}\{f(x)\},\qquad
%   M_i''=\sup_{x\in[x',x_i]}\{f(x)\};
% \]
% \[
%   m=\inf_{x\in[a,b]}\{f(x)\},\qquad M=\sup_{x\in[a,b]}\{f(x)\}.
% \]
% 因此我们得到以下的估计:
% \[
%   0\leq \overline S(\Delta)-\overline S(\Delta^*)\leq (N-1)(M-m)\lambda(\Delta).
% \]
% 我们不妨设 $m<M$, 否则 $f(x)$ 为常数函数, 此时结论显然成立。取
% \[
%   \delta=\min\left\{\delta_1,\frac{\varepsilon}{2(N-1)(M-m)}\right\},
% \]
% 当 $\lambda(\Delta)<\delta$ 时, 有
% \[
% \begin{aligned}
%   0\leq \overline S(\Delta)-\overline{\int_a^b} f(x)\mathrm dx
%   &=[\overline S(\Delta)-\overline S(\Delta^*)]+[\overline S(\Delta^*)-\overline S(\Delta_1)]\\
%   &\quad+\left[\overline S(\Delta_1)-\overline{\int_a^b} f(x)\mathrm dx\right]\\
%   &<\frac{\varepsilon}{2}+0+\frac{\varepsilon}{2}=\varepsilon.
% \end{aligned}
% \qedhere
% \]
% \end{proof}
% 
% 
% \begin{remark}\label{remark:ma-7-source-18}
% 如果用 ``$\varepsilon$-$\delta$'' 语言来叙述定理~\ref{theorem:ma-7-2-4}, 则它的形式是: 对于 $\forall\varepsilon>0$, $\exists\delta>0$, 对区间 $[a,b]$ 的任意分割 $\Delta$, 当 $\lambda(\Delta)<\delta$ 时, 就有
% \[
%   0\leq \overline S(\Delta)-\overline{\int_a^b} f(x)\mathrm dx<\varepsilon,\qquad
%   0\leq \underline{\int_a^b} f(x)\mathrm dx-\underline S(\Delta)<\varepsilon.
% \]
% \end{remark}
% 
% 在研究什么样的函数可积这一问题之前, 我们先来考查一下定理~\ref{theorem:ma-7-2-4} 的证明过程。从该定理的证明中读者不难发现以下事实: 对于 $\forall\varepsilon>0$, 若我们能找到一个特别的分割 $\Delta_1$, 使得
% \[
%   0\leq \overline S(\Delta_1)-\overline{\int_a^b} f(x)\mathrm dx<\frac{\varepsilon}{2},
% \]
% 则必定存在 $\delta>0$, 使得对任意分割 $\Delta$, 当 $\lambda(\Delta)<\delta$ 时, 有
% \[
%   0\leq \overline S(\Delta)-\overline{\int_a^b} f(x)\mathrm dx<\varepsilon.
% \]
% 上述事实告诉我们如下结论: $\lim\limits_{\lambda(\Delta)\to0}\overline S(\Delta)=\overline{\int_a^b} f(x)\mathrm dx$ 的充分必要条件是: 对于 $\forall\varepsilon>0$, 存在分割 $\Delta$, 使得
% \[
%   0\leq \overline S(\Delta)-\overline{\int_a^b} f(x)\mathrm dx<\varepsilon.
% \]
% 
% 同理, 对于下积分, 我们也有相应的性质。上积分与下积分的这个性质将应用到可积性问题的讨论中.
% 
% 现在我们来证明下面的定理.
% 
% \begin{theorem}\label{theorem:ma-7-2-5}
% 设函数 $f(x)$ 在区间 $[a,b]$ 上有界, 则 $f(x)\in R[a,b]$ 的充分必要条件是
% \[
%   \underline{\int_a^b} f(x)\mathrm dx=\overline{\int_a^b} f(x)\mathrm dx.
% \]
% \end{theorem}
% 
% 
% \begin{proof}
% 必要性\quad 我们观察到, 对任意的分割, 通过在每个小区间上的特殊取值, 可以使黎曼和任意地接近达布大和或达布小和。设 $f(x)\in R[a,b]$, 则对于 $\forall\varepsilon>0$, $\exists\delta>0$, 对 $[a,b]$ 的任意分割
% \[
%   \Delta:a=x_0<x_1<x_2<\cdots<x_n=b,
% \]
% 只要 $\lambda(\Delta)<\delta$, 就有
% \[
%   \left|\sum_{i=1}^n f(\xi_i)\Delta x_i-\int_a^b f(x)\mathrm dx\right|<\varepsilon,
% \]
% 其中 $\xi_i\in[x_{i-1},x_i](i=1,2,\cdots,n)$.
% 由上、下确界的定义, 存在 $\xi_i',\xi_i''\in[x_{i-1},x_i](i=1,2,\cdots,n)$, 使得
% \[
%   f(\xi_i')\geq M_i-\frac{\varepsilon}{b-a}\quad\text{及}\quad f(\xi_i'')\leq m_i+\frac{\varepsilon}{b-a},
% \]
% 其中
% \[
%   M_i=\sup_{x\in[x_{i-1},x_i]}\{f(x)\},\qquad
%   m_i=\inf_{x\in[x_{i-1},x_i]}\{f(x)\}\quad(i=1,2,\cdots,n).
% \]
% 因此有
% \[
%   \overline{\int_a^b} f(x)\mathrm dx-\varepsilon
%   \leq \sum_{i=1}^n M_i\Delta x_i-\varepsilon
%   \leq \sum_{i=1}^n f(\xi_i')\Delta x_i
%   < \int_a^b f(x)\mathrm dx+\varepsilon
% \]
% 和
% \[
%   \underline{\int_a^b} f(x)\mathrm dx+\varepsilon
%   \geq \sum_{i=1}^n m_i\Delta x_i+\varepsilon
%   \geq \sum_{i=1}^n f(\xi_i'')\Delta x_i
%   > \int_a^b f(x)\mathrm dx-\varepsilon.
% \]
% 由此得
% \[
%   0\leq \overline{\int_a^b} f(x)\mathrm dx-\underline{\int_a^b} f(x)\mathrm dx<4\varepsilon.
% \]
% 由于 $\overline{\int_a^b} f(x)\mathrm dx-\underline{\int_a^b} f(x)\mathrm dx$ 是固定常数, 由 $\varepsilon$ 的任意性, 必有
% \[
%   \underline{\int_a^b} f(x)\mathrm dx=\overline{\int_a^b} f(x)\mathrm dx.
% \]
% 必要性得证。
% 
% 下面证明充分性。
% 由定理~\ref{theorem:ma-7-2-4} 可知, 对于 $\forall\varepsilon>0$, $\exists\delta>0$, 对于任意的分割 $\Delta$, 当 $\lambda(\Delta)<\delta$ 时, 对于 $\Delta$ 的任意黎曼和 $S(\Delta)$, 有
% \[
%   \underline{\int_a^b} f(x)\mathrm dx-\varepsilon
%   \leq \underline S(\Delta)\leq S(\Delta)\leq \overline S(\Delta)
%   \leq \overline{\int_a^b} f(x)\mathrm dx+\varepsilon.
% \]
% 记 $I=\underline{\int_a^b} f(x)\mathrm dx=\overline{\int_a^b} f(x)\mathrm dx$, 则由上式得
% \[
%   |S(\Delta)-I|<\varepsilon.
% \]
% 由定积分的定义知 $f(x)\in R[a,b]$。充分性得证。
% \end{proof}
% 
% 为了叙述简便, 下面对于 $[a,b]$ 的一个分割 $\Delta:a=x_0<x_1<\cdots<x_n=b$, 我们总记 $\omega_i$ 为 $f(x)$ 在区间 $[x_{i-1},x_i]$ 上的振幅.
% 由定理~\ref{theorem:ma-7-2-4} 和后面的注及定理~\ref{theorem:ma-7-2-5}, 我们有下面的结果.
% 
% \begin{theorem}\label{theorem:ma-7-2-6}
% 设函数 $f(x)$ 在区间 $[a,b]$ 上有界, 则下面的三个结论相互等价:
% 
% (1) $f(x)$ 在区间 $[a,b]$ 上可积;
% 
% (2) 对于 $\forall\varepsilon>0$, 存在区间 $[a,b]$ 的一个分割 $\Delta$, 使得
% \[
%   \sum_{i=1}^n \omega_i\Delta x_i<\varepsilon;
% \]
% 
% (3) 对于 $\forall\varepsilon>0$, $\forall\sigma>0$, 存在区间 $[a,b]$ 的一个分割 $\Delta$, 使得 $\omega_i\geq\varepsilon$ 的小区间 $[x_{i-1},x_i]$ 的长度总和小于 $\sigma$.
% \end{theorem}
% 
% 
% \begin{proof}
% 我们先证 (1)$\Rightarrow$(2)。设 $f(x)\in R[a,b]$, 由定理~\ref{theorem:ma-7-2-5} 我们有
% \[
%   \underline{\int_a^b} f(x)\mathrm dx=\overline{\int_a^b} f(x)\mathrm dx.
% \]
% 再由上、下积分的定义易知, 对于 $\forall\varepsilon>0$, 存在 $[a,b]$ 的分割 $\Delta$, 使得
% \[
%   \underline{\int_a^b} f(x)\mathrm dx-\frac{\varepsilon}{2}<\underline S(\Delta)\leq \overline S(\Delta)<\overline{\int_a^b} f(x)\mathrm dx+\frac{\varepsilon}{2}.
% \]
% 因此有
% \[
%   \sum_{i=1}^n \omega_i\Delta x_i=\overline S(\Delta)-\underline S(\Delta)<\varepsilon.
% \]
% (2) 得证.
% 
% 现证 (2)$\Rightarrow$(1)。由于对于 $\forall\varepsilon>0$, 存在 $[a,b]$ 的一个分割 $\Delta:a=x_0<x_1<\cdots<x_n=b$, 使得
% \[
%   \sum_{i=1}^n \omega_i\Delta x_i<\varepsilon,
% \]
% 因此有
% \[
%   0\leq \overline{\int_a^b} f(x)\mathrm dx-\underline{\int_a^b} f(x)\mathrm dx
%   \leq \overline S(\Delta)-\underline S(\Delta)
%   =\sum_{i=1}^n\omega_i\Delta x_i<\varepsilon.
% \]
% 由 $\varepsilon$ 的任意性, 我们有
% \[
%   \underline{\int_a^b} f(x)\mathrm dx=\overline{\int_a^b} f(x)\mathrm dx.
% \]
% 再由定理~\ref{theorem:ma-7-2-5} 我们知 $f(x)\in R[a,b]$.
% 
% 下面我们来证明 (2)$\Longleftrightarrow$(3)。首先注意到, 对于 $\forall\varepsilon>0$, 对 $[a,b]$ 的分割
% \[
%   \Delta:a=x_0<x_1<\cdots<x_n=b,
% \]
% 总有
% \[
%   \sum_{i=1}^n\omega_i\Delta x_i=\sum_{(1)}\omega_i\Delta x_i+\sum_{(2)}\omega_i\Delta x_i,
% \]
% 其中 $\sum_{(1)}$ 表示对 $\omega_i<\varepsilon$ 的求和, $\sum_{(2)}$ 表示对 $\omega_i\geq\varepsilon$ 的求和.
% 先证 (2)$\Rightarrow$(3)。对于 $\forall\varepsilon>0$, $\forall\sigma>0$, 由 (2), 存在分割
% \[
%   \Delta:a=x_0<x_1<\cdots<x_n=b,
% \]
% 使得
% \[
%   \sum_{i=1}^n\omega_i\Delta x_i<\sigma\varepsilon.
% \]
% 因此有
% \[
%   \varepsilon\sum_{(2)}\Delta x_i\leq \sum_{(2)}\omega_i\Delta x_i\leq \sum_{i=1}^n\omega_i\Delta x_i<\sigma\varepsilon,
% \]
% 即
% \[
%   \sum_{(2)}\Delta x_i<\sigma.
% \]
% 
% 再证 (3)$\Rightarrow$(2)。记 $\omega$ 为 $f(x)$ 在 $[a,b]$ 上的振幅, 不妨设 $f(x)$ 不是常数函数, 因此 $\omega>0$。由 (3), 对于 $\dfrac{\varepsilon}{2(b-a)}$ 及 $\sigma=\dfrac{\varepsilon}{2\omega}$, 存在分割
% \[
%   \Delta:a=x_0<x_1<\cdots<x_n=b,
% \]
% 使得 $\omega_i\geq\dfrac{\varepsilon}{2(b-a)}$ 的小区间 $[x_{i-1},x_i]$ 长度的总和小于 $\sigma$。现在记 $\sum_{(1)'}\omega_i\Delta x_i$ 是对 $\omega_i<\dfrac{\varepsilon}{2(b-a)}$ 的求和, $\sum_{(2)'}\omega_i\Delta x_i$ 是对 $\omega_i\geq\dfrac{\varepsilon}{2(b-a)}$ 的求和, 则有
% \[
% \begin{aligned}
%   \sum_{i=1}^n\omega_i\Delta x_i
%   &=\sum_{(1)'}\omega_i\Delta x_i+\sum_{(2)'}\omega_i\Delta x_i\\
%   &<\frac{\varepsilon}{2(b-a)}\sum_{(1)'}\Delta x_i+\omega\sum_{(2)'}\Delta x_i\\
%   &<\frac{\varepsilon}{2(b-a)}(b-a)+\omega\frac{\varepsilon}{2\omega}
%   =\varepsilon.
% \end{aligned}
% \qedhere
% \]
% \end{proof}
% 
% 
% \begin{remark}\label{remark:ma-7-source-23}
% 定理~\ref{theorem:ma-7-2-6} 实际上回答了我们在本节开始时提出的第一个关于可积性的问题(问题 (a)), 即在定积分的定义中, 我们可以取一列特殊的分割 $\Delta_n$, 只要当 $\lambda(\Delta_n)\to0$ 时, $f(x)$ 关于 $\Delta_n$ 的任意黎曼和可以任意接近一个常数 $I$, 则 $f(x)\in R[a,b]$ 且 $I=\int_a^b f(x)\mathrm dx$.
% \end{remark}
% 
% 
% \subsection{可积函数类}
% \label{subsec:ma-7-2-3}
% 
% 由上一节的结果, 我们可以找到一些常见的可积函数类.
% 
% \begin{theorem}\label{theorem:ma-7-2-7}
% 设函数 $f(x)$ 在区间 $[a,b]$ 上有界, 且只有有限个间断点, 则 $f(x)\in R[a,b]$.
% \end{theorem}
% 
% 
% \begin{proof}
% 记 $f(x)$ 在 $[a,b]$ 上的振幅为 $\omega$, 我们不妨设 $\omega>0$(否则 $f(x)$ 为一个常数函数)。再设 $f(x)$ 在 $[a,b]$ 上的间断点的个数为 $N$。对于 $\forall\varepsilon>0$, 在 $[a,b]$ 上取 $N$ 个互不相交的闭区间, 使得满足: (1) 每个区间的长度大于零小于 $\dfrac{\varepsilon}{2N\omega}$; (2) 它们包含了 $f(x)$ 在 $[a,b]$ 上的所有的间断点; (3) 它们的内部(设 $[c,d]$ 为一闭区间, 称 $(c,d)$ 为该区间的内部)包含了 $f(x)$ 在 $(a,b)$ 内所有的间断点。$[a,b]$ 除去这组区间后得到的每个小区间加上端点则为有限个闭区间: $I_1,I_2,\cdots,I_K$。由于 $f(x)$ 在每个小闭区间 $I_j(j=1,2,\cdots,K)$ 上连续, 从而一致连续, 因此存在 $\delta>0$, 使得当 $x',x''$ 属于同一个闭区间 $I_j$, 且 $|x'-x''|<\delta$ 时, 有
% \[
%   |f(x')-f(x'')|<\frac{\varepsilon}{2(b-a)}.
% \]
% 将每个小区间 $I_j$ 适当等分, 使得等分后的小区间长度小于 $\delta$, 则所有这些小区间的等分点与原来的 $N$ 个区间的端点构成 $[a,b]$ 的一个分割 $\Delta:a=x_0<x_1<\cdots<x_n=b$。对此分割有
% \[
%   \sum_{i=1}^n\omega_i\Delta x_i< \frac{\varepsilon}{2N\omega}\omega N+\frac{\varepsilon(b-a)}{2(b-a)}=\varepsilon.
% \qedhere
% \]
% \end{proof}
% 
% 
% 另外, 由定积分的定义以及以上我们得到的可积性的判别法容易看出, 若函数 $f(x)\in R[a,b]$, 函数 $g(x)$ 在 $[a,b]$ 上有定义且与 $f(x)$ 在 $[a,b]$ 上只有有限多个点上的取值不同, 则 $g(x)\in R[a,b]$, 并且有
% \[
%   \int_a^b f(x)\mathrm dx=\int_a^b g(x)\mathrm dx.
% \]
% 
% \begin{theorem}\label{theorem:ma-7-2-8}
% 设函数 $f(x)$ 在区间 $[a,b]$ 上单调, 则 $f(x)\in R[a,b]$.
% \end{theorem}
% 
% 
% \begin{proof}
% 不妨设 $f(x)$ 在 $[a,b]$ 上单调上升。对于 $\forall\varepsilon>0$, 将 $[a,b]$ 等分成 $n$ 个小区间, 使得
% \[
%   \frac{(b-a)[f(b)-f(a)]}{n}<\varepsilon.
% \]
% 设其对应的分割为 $a=x_0<x_1<\cdots<x_n=b$, 则有
% \[
% \begin{aligned}
%   \sum_{i=1}^n \omega_i\Delta x_i
%   &=\sum_{i=1}^n [f(x_i)-f(x_{i-1})]\frac{b-a}{n}\\
%   &=\frac{b-a}{n}[f(b)-f(a)]<\varepsilon.
% \end{aligned}
% \qedhere
% \]
% \end{proof}
% 
% 
% \begin{example}\label{example:ma-7-2-2}
% 证明黎曼函数
% \[
% R(x)=
% \begin{cases}
%   \dfrac{1}{p}, & x\in[0,1],\ x=\dfrac{q}{p}(p,q>0,\text{为既约正整数}),\\
%   0, & x\in(0,1)\setminus\mathbb Q,\\
%   1, & x=0
% \end{cases}
% \]
% 在区间 $[0,1]$ 上可积, 且 $\int_0^1 R(x)\mathrm dx=0$.
% \end{example}
% 
% 
% \begin{proof}
% 对于 $\forall\varepsilon>0$, $\forall\sigma>0$, 在 $[0,1]$ 中仅存在有限多个点使得 $R(x)\geq\varepsilon$。设这有限多个点为 $x_1',\cdots,x_N'$。在 $[0,1]$ 中取 $N$ 个互不相交的小闭区间, 使得每一个区间的长度大于零小于 $\dfrac{\sigma}{N}$, 并且它们包含这 $N$ 个点。进一步我们要求当 $x_i'\neq0,1$ 时, $x_i'$ 落在某个小区间的内部。对于这 $N$ 个小区间的端点构成 $[0,1]$ 的分割 $\Delta$, 那些振幅大于或等于 $\varepsilon$ 的小区间的总和显然小于 $\dfrac{\sigma}{N}\cdot N=\sigma$。由定理~\ref{theorem:ma-7-2-6} 的 (3) 知 $R(x)\in R[0,1]$.
% 
% 由于 $R(x)$ 在 $[0,1]$ 上可积, 并且对任意的分割
% \[
%   \Delta:a=x_0<x_1<\cdots<x_n=b,
% \]
% 在每个 $[x_{i-1},x_i]$ 中取 $\xi_i$ 为无理数, 则其黎曼和为
% \[
%   \sum_{i=1}^n R(\xi_i)\Delta x_i=0,
% \]
% 因此
% \[
%   \int_0^1 R(x)\mathrm dx=0.
% \qedhere
% \]
% \end{proof}
% 
% 
% \begin{example}\label{example:ma-7-2-3}
% 证明函数
% \[
% f(x)=
% \begin{cases}
%   \sin\dfrac{1}{x}, & x\neq0,\\
%   0, & x=0
% \end{cases}
% \]
% 在任何区间 $[a,b]$ 上可积.
% \end{example}
% 
% 
% \begin{proof}
% 当 $0\notin[a,b]$ 时, $f(x)$ 在 $[a,b]$ 上连续, 从而 $f(x)\in R[a,b]$; 当 $0\in[a,b]$ 时, $f(x)$ 在 $[a,b]$ 上有界, 且只有一个间断点 $x_0=0$, 因此由定理~\ref{theorem:ma-7-2-7} 知 $f(x)\in R[a,b]$.
% 
% 从本例中可以看出, 尽管函数 $f(x)=\sin\dfrac{1}{x}$ 在 $x_0=0$ 的邻域中变化非常剧烈, 但它仍然在包含 $x_0$ 的区间上可积.
% 
% 函数的可积性问题将在实变函数论课程中得到完全解决。在实变函数论课程中将证明以下结论.
% \end{proof}
% 
% 
% \begin{theorem}[{勒贝格定理}]\label{theorem:ma-7-2-9}
% 设函数 $f(x)$ 在区间 $[a,b]$ 上有界, 记 $E$ 为 $f(x)$ 在 $[a,b]$ 上的间断点集, 则 $f(x)\in R[a,b]$ 的充分必要条件是: 对于 $\forall\varepsilon>0$, 存在一列开区间 $(x_{i-1},x_i)(i=1,2,\cdots)$, 使得
% \[
%   E\subset \bigcup_{i=1}^{\infty}(x_{i-1},x_i),
%   \qquad
%   \text{并且对于 } \forall n\in\mathbb N,\ \text{有 } \sum_{i=1}^n(x_i-x_{i-1})<\varepsilon.
% \]
% \end{theorem}
% 
% 
% \section{定积分的性质}
% \label{sec:ma-7-3}
% 
% 前面我们已经对可积性问题进行了讨论, 在本节中将讨论定积分的一些基本性质.
% 
% 首先, 我们作两条自然的规定:
% \[
%   \text{(1) }\int_a^a f(x)\mathrm dx=0;
% \]
% \[
%   \text{(2) 当函数 } f(x) \text{ 在区间 } [a,b] \text{ 上可积时, }
%   \int_b^a f(x)\mathrm dx=-\int_a^b f(x)\mathrm dx.
% \]
% 定积分具有以下基本性质.
% 
% \begin{theorem}[{线性性质}]\label{theorem:ma-7-3-1}
% 设函数 $f(x),g(x)\in R[a,b]$, $\alpha,\beta\in\mathbb R$, 则 $\alpha f(x)+\beta g(x)\in R[a,b]$, 并且有
% \[
%   \int_a^b[\alpha f(x)+\beta g(x)]\mathrm dx
%   =\alpha\int_a^b f(x)\mathrm dx+\beta\int_a^b g(x)\mathrm dx.
% \]
% 此性质可由定积分的定义直接推出, 请读者自证.
% \end{theorem}
% 
% 
% \begin{theorem}\label{theorem:ma-7-3-2}
% 设函数 $f(x)\in R[a,b]$, 则 $|f(x)|\in R[a,b]$, 并且有
% \[
%   \left|\int_a^b f(x)\mathrm dx\right|\leq \int_a^b |f(x)|\mathrm dx.
% \]
% \end{theorem}
% 
% 
% \begin{proof}
% 对于 $\forall\varepsilon>0$, 由 $f(x)\in R[a,b]$, 存在 $[a,b]$ 的分割
% \[
%   \Delta:a=x_0<x_1<\cdots<x_n=b,
% \]
% 记 $\omega_i(f)$ 为 $f(x)$ 在 $[x_{i-1},x_i](i=1,2,\cdots,n)$ 上的振幅, 则有
% \[
%   \sum_{i=1}^n \omega_i(f)\Delta x_i<\varepsilon.
% \]
% 现记 $\omega_i(|f|)$ 为 $|f(x)|$ 在 $[x_{i-1},x_i](i=1,2,\cdots,n)$ 上的振幅, 由于当 $x',x''\in[x_{i-1},x_i]$ 时, 有
% \[
%   \bigl||f(x')|-|f(x'')|\bigr|\leq |f(x')-f(x'')|,
% \]
% 因此有
% \[
%   \omega_i(|f|)\leq \omega_i(f)\quad(i=1,2,\cdots,n),
% \]
% 从而有
% \[
%   \sum_{i=1}^n \omega_i(|f|)\Delta x_i<\varepsilon.
% \]
% 所以 $|f(x)|\in R[a,b]$.
% 
% 对 $[a,b]$ 的任意分割 $\Delta$ 及 $[x_{i-1},x_i](i=1,2,\cdots,n)$ 上任取的 $\xi_i$, 我们有
% \[
%   \left|\sum_{i=1}^n f(\xi_i)\Delta x_i\right|
%   \leq \sum_{i=1}^n |f(\xi_i)|\Delta x_i.
% \]
% 因此, 当 $\lambda(\Delta)\to0$ 时, 有
% \[
%   \left|\int_a^b f(x)\mathrm dx\right|\leq \int_a^b |f(x)|\mathrm dx.
% \qedhere
% \]
% \end{proof}
% 
% 
% \begin{theorem}\label{theorem:ma-7-3-3}
% 设 $a<c<b$, 则函数 $f(x)\in R[a,b]$ 的充分必要条件是: $f(x)\in R[a,c]$ 和 $f(x)\in R[c,b]$。当 $f(x)\in R[a,b]$ 时, 下式成立:
% \begin{equation}
%   \int_a^b f(x)\mathrm dx=\int_a^c f(x)\mathrm dx+\int_c^b f(x)\mathrm dx.
%   \label{eq:ma-7-3-1}
% \end{equation}
% \end{theorem}
% 
% 
% \begin{proof}
% 必要性\quad 设函数 $f(x)\in R[a,b]$, 我们证明 $f(x)$ 在 $[a,b]$ 的任何子区间 $[a_1,b_1]\subset[a,b]$ 上可积。事实上, 由于 $f(x)$ 在 $[a,b]$ 上可积, 对于 $\forall\varepsilon>0$, 存在 $[a,b]$ 的分割
% \[
%   \Delta:a=x_0<x_1<\cdots<x_n=b,
% \]
% 记 $\omega_i$ 为 $f(x)$ 在 $[x_{i-1},x_i](i=1,2,\cdots,n)$ 上的振幅, 则有
% \[
%   \sum_{i=1}^n \omega_i\Delta x_i<\varepsilon.
% \]
% 由于 $\Delta$ 落在 $[a_1,b_1]$ 中的分割点与 $a_1,b_1$ 构成了 $[a_1,b_1]$ 的一个分割
% \[
%   \Delta':a_1=x_0'<x_1'<\cdots<x_m'=b_1,
% \]
% 对于此分割显然有
% \[
%   \sum_{i=1}^m \omega_i'\Delta x_i'<\varepsilon,
% \]
% 其中 $\omega_i'$ 为 $f(x)$ 在 $[x_{i-1}',x_i']$ 上的振幅, 而 $\Delta x_i'=x_i'-x_{i-1}'(i=1,2,\cdots,m)$, 因此 $f(x)$ 在 $[a_1,b_1]$ 上可积。特别地, 有 $f(x)\in R[a,c]$, $f(x)\in R[c,b]$.
% 
% 充分性\quad 对于 $\forall\varepsilon>0$, 由 $f(x)\in R[a,c]$, 存在 $[a,c]$ 的分割
% \[
%   \Delta_1:a=x_0<x_1<\cdots<x_n=c,
% \]
% 记 $\omega_i$ 为 $f(x)$ 在 $[x_{i-1},x_i](i=1,2,\cdots,n)$ 上的振幅, 则有
% \[
%   \sum_{i=1}^n\omega_i\Delta x_i<\frac{\varepsilon}{2}.
% \]
% 由 $f(x)\in R[c,b]$, 存在 $[c,b]$ 上的分割
% \[
%   \Delta_2:c=x_0'<x_1'<\cdots<x_m'=b,
% \]
% 记 $\omega_i'$ 为 $f(x)$ 在 $[x_{i-1}',x_i'](i=1,2,\cdots,m)$ 上的振幅及 $\Delta x_i'=x_i'-x_{i-1}'$, 则有
% \[
%   \sum_{i=1}^m\omega_i'\Delta x_i'<\frac{\varepsilon}{2}.
% \]
% 显然, $\Delta_1\cup\Delta_2$ 是 $[a,b]$ 的一个分割 $\Delta$, 且有
% \[
%   \sum_{i=1}^n\omega_i\Delta x_i+\sum_{i=1}^m\omega_i'\Delta x_i'<\varepsilon.
% \]
% 而上式左边刚好是 $f(x)$ 关于 $\Delta$ 的每个小区间的振幅乘以对应小区间的长度之和, 因此 $f(x)\in R[a,b]$.
% 
% 现证定积分等式 \eqref{eq:ma-7-3-1} 成立。事实上, 对 $[a,c]$ 的任意分割
% \[
%   \Delta_1:a=x_0<x_1<\cdots<x_n=c
% \]
% 作黎曼和 $\sum_{i=1}^n f(\xi_i)\Delta x_i$, 对 $[c,b]$ 的任意分割
% \[
%   \Delta_2:c=x_0'<x_1'<\cdots<x_m'=b
% \]
% 作黎曼和 $\sum_{i=1}^m f(\xi_i')\Delta x_i'$, 则
% \[
%   \sum_{i=1}^n f(\xi_i)\Delta x_i+\sum_{i=1}^m f(\xi_i')\Delta x_i'
% \]
% 是 $f(x)$ 在 $[a,b]$ 上关于分割 $\Delta_1\cup\Delta_2$ 的黎曼和。由于以上的各个黎曼和当小区间最大长度趋于零时都存在极限, 因此定积分等式 \eqref{eq:ma-7-3-1} 成立.
% \end{proof}
% 
% 
% 关于本性质我们要指出的是: 如果不假定 $a<c<b$, 但假定函数 $f(x)$ 在所涉及的区间是可积的, 利用在本节开始时的两个规定, 则等式 \eqref{eq:ma-7-3-1} 总是成立的。例如, 设 $a<b<c$, 假定 $f(x)\in R[a,c]$, 则仍然成立
% \[
%   \int_a^b f(x)\mathrm dx=\int_a^c f(x)\mathrm dx+\int_c^b f(x)\mathrm dx.
% \]
% 事实上, 我们有
% \[
% \begin{aligned}
%   \int_a^b f(x)\mathrm dx
%   &=\int_a^b f(x)\mathrm dx+\int_b^c f(x)\mathrm dx-\int_b^c f(x)\mathrm dx\\
%   &=\int_a^c f(x)\mathrm dx+\int_c^b f(x)\mathrm dx.
% \end{aligned}
% \]
% 
% \begin{theorem}\label{theorem:ma-7-3-4}
% 设函数 $f(x),g(x)\in R[a,b]$, 并且对于 $\forall x\in[a,b]$, 有 $f(x)\geq g(x)$, 则
% \[
%   \int_a^b f(x)\mathrm dx\geq \int_a^b g(x)\mathrm dx.
% \]
% 定理~\ref{theorem:ma-7-3-4} 的证明是简单的。事实上对函数 $f(x)-g(x)$ 用定积分的定义及定积分的线性性质, 即可推出该定理的结论.
% \end{theorem}
% 
% 设函数 $g(x)$ 在区间 $[a,b]$ 上有定义, 如果 $[a,b]$ 是有限个两两互不相交的小区间的并, 并且 $g(x)$ 在每个小区间上均为常数函数, 则称 $g(x)$ 为 $[a,b]$ 上的一个阶梯函数。显然闭区间 $[a,b]$ 上的任何阶梯函数是有界的, 而且至多只有有限个间断点, 因此阶梯函数是可积的.
% 
% \begin{example}\label{example:ma-7-3-1}
% 设函数 $f(x)\in R[a,b]$, 证明: 对于 $\forall\varepsilon>0$,
% 
% (1) 存在 $[a,b]$ 上的阶梯函数 $h(x)$, 使得 $\int_a^b |f(x)-h(x)|\mathrm dx<\varepsilon$;
% 
% (2) 存在 $[a,b]$ 上的连续函数 $g(x)$, 使得 $\int_a^b |f(x)-g(x)|\mathrm dx<\varepsilon$.
% \end{example}
% 
% 
% \begin{proof}
% 由于 $f(x)\in R[a,b]$, 对于 $\forall\varepsilon>0$, 存在 $[a,b]$ 的分割
% \[
%   \Delta:a=x_0<x_1<\cdots<x_n=b,
% \]
% 记 $M_i,m_i$ 和 $\omega_i$ 分别为 $f(x)$ 在 $[x_{i-1},x_i](i=1,2,\cdots,n)$ 上的上、下确界和振幅, 则有
% \[
%   \sum_{i=1}^n\omega_i\Delta x_i<\varepsilon.
% \]
% (1) 定义阶梯函数
% \[
%   \overline f(x)=
%   \begin{cases}
%     M_i, & x\in[x_{i-1},x_i),\ i=1,2,\cdots,n-1,\\
%     M_n, & x\in[x_{n-1},b]
%   \end{cases}
% \]
% 和
% \[
%   \underline f(x)=
%   \begin{cases}
%     m_i, & x\in[x_{i-1},x_i),\ i=1,2,\cdots,n-1,\\
%     m_n, & x\in[x_{n-1},b],
%   \end{cases}
% \]
% 则 $\overline f(x)$ 与 $\underline f(x)$ 均为区间 $[a,b]$ 上的阶梯函数。为了证明 $\overline f(x)$ 与 $\underline f(x)$ 满足 (1) 中的结论, 我们令
% \[
%   \overline f_1(x)=
%   \begin{cases}
%     \overline f(x), & x\in[a,b]\text{ 且 }x\neq x_i,\\
%     f(x_i), & x=x_i
%   \end{cases}
%   \quad(i=0,1,\cdots,n).
% \]
% \[
%   \underline f_1(x)=
%   \begin{cases}
%     \underline f(x), & x\in[a,b]\text{ 且 }x\neq x_i,\\
%     f(x_i), & x=x_i
%   \end{cases}
%   \quad(i=0,1,\cdots,n).
% \]
% 容易看出 $\overline f(x)$, $\overline f_1(x)$, $\underline f(x)$ 与 $\underline f_1(x)$ 均为 $[a,b]$ 上的可积函数, 且有
% \[
%   \int_a^b |f(x)-\overline f(x)|\mathrm dx
%   =\int_a^b |f(x)-\overline f_1(x)|\mathrm dx
% \]
% 和
% \[
%   \int_a^b |f(x)-\underline f(x)|\mathrm dx
%   =\int_a^b |f(x)-\underline f_1(x)|\mathrm dx.
% \]
% 由振幅的定义, 对于 $\forall x\in[x_{i-1},x_i](i=1,2,\cdots,n)$, 有
% \begin{equation}
%   |f(x)-\overline f_1(x)|\leq\omega_i
%   \label{eq:ma-7-3-2}
% \end{equation}
% 和
% \begin{equation}
%   |f(x)-\underline f_1(x)|\leq\omega_i,
%   \label{eq:ma-7-3-3}
% \end{equation}
% 从而有
% \[
% \begin{aligned}
%   \int_a^b |f(x)-\underline f(x)|\mathrm dx
%   &=\int_a^b |f(x)-\underline f_1(x)|\mathrm dx\\
%   &=\sum_{i=1}^n\int_{x_{i-1}}^{x_i}|f(x)-\underline f_1(x)|\mathrm dx\\
%   &\leq \sum_{i=1}^n\omega_i\Delta x_i<\varepsilon
% \end{aligned}
% \]
% 和
% \[
%   \int_a^b |f(x)-\overline f(x)|\mathrm dx<\varepsilon.
% \]
% 取 $h(x)=\underline f(x)$(或 $\overline f(x)$), 即得 (1).
% 
% (2) 作折线依次连接 $(x_{i-1},f(x_{i-1}))$ 和 $(x_i,f(x_i))(i=1,2,\cdots,n)$。设 $g(x)$ 为 $[a,b]$ 上的函数, 其图像恰为该折线, 则 $g(x)$ 在 $[a,b]$ 上连续($g(x)$ 也称为是分段线性函数).
% 
% 同理由振幅的定义, 对于 $\forall x\in[x_{i-1},x_i](i=1,2,\cdots,n)$, 有
% \begin{equation}
%   |f(x)-g(x)|\leq\omega_i,
%   \label{eq:ma-7-3-4}
% \end{equation}
% 因此有
% \[
%   \int_a^b |f(x)-g(x)|\mathrm dx<\varepsilon,
% \]
% 从而 (2) 得证.
% \end{proof}
% 
% 
% \begin{remark}\label{remark:ma-7-source-41}
% 在第十二章傅里叶(Fourier)级数的研究中, 我们还要用到以下结论: 设函数 $f(x)\in R[a,b]$, 则对于 $\forall\varepsilon>0$, 存在区间 $[a,b]$ 上的连续函数 $g(x)$, 使得
% \[
%   \int_a^b [f(x)-g(x)]^2\mathrm dx<\varepsilon.
% \]
% 利用例~\ref{example:ma-7-3-1}, 我们很容易证明此结论。事实上, 由 $f(x)\in R[a,b]$, 从而存在 $M>0$, 使得对于 $x\in[a,b]$, 有 $|f(x)|\leq M$。对于 $\forall\varepsilon>0$, 由例~\ref{example:ma-7-3-1}, 我们可以找到 $[a,b]$ 上的连续函数 $g(x)$, 使得
% \[
%   \int_a^b |f(x)-g(x)|\mathrm dx<\frac{\varepsilon}{2M}.
% \]
% 显然, 从例~\ref{example:ma-7-3-1} 中 $g(x)$ 的构造不难看出, 对于 $\forall x\in[a,b]$, 有 $|g(x)|\leq M$, 因此有
% \[
% \begin{aligned}
%   \int_a^b [f(x)-g(x)]^2\mathrm dx
%   &=\int_a^b |f(x)-g(x)|^2\mathrm dx\\
%   &\leq \int_a^b (|f(x)|+|g(x)|)|f(x)-g(x)|\mathrm dx\\
%   &\leq \int_a^b 2M|f(x)-g(x)|\mathrm dx\\
%   &<2M\cdot\frac{\varepsilon}{2M}=\varepsilon.
% \end{aligned}
% \]
% \end{remark}
% 
% 我们下面继续讨论定积分的性质.
% 
% \begin{theorem}\label{theorem:ma-7-3-5}
% 设函数 $f(x),g(x)\in R[a,b]$, 则 $f(x)g(x)\in R[a,b]$.
% \end{theorem}
% 
% 
% \begin{proof}
% 由 $f(x),g(x)$ 的可积性知, $f(x)$ 与 $g(x)$ 均是有界函数。我们取 $M>0$, 使得当 $x\in[a,b]$ 时, 有 $|f(x)|\leq M$ 及 $|g(x)|\leq M$。对于 $\forall\varepsilon>0$, 再由 $f(x),g(x)$ 的可积性, 存在 $[a,b]$ 的分割
% \[
%   \Delta:a=x_0<x_1<\cdots<x_n=b,
% \]
% 记 $\omega_i(h)$ 为函数 $h(x)$ 在 $[x_{i-1},x_i](i=1,2,\cdots,n)$ 上的振幅, 则有
% \begin{equation}
%   \sum_{i=1}^n \omega_i(f)\Delta x_i<\frac{\varepsilon}{2M}
%   \label{eq:ma-7-3-5}
% \end{equation}
% 及
% \begin{equation}
%   \sum_{i=1}^n \omega_i(g)\Delta x_i<\frac{\varepsilon}{2M}.
%   \label{eq:ma-7-3-6}
% \end{equation}
% 对任意的 $x',x''\in[x_{i-1},x_i]$, 由
% \[
% \begin{aligned}
%   |f(x'')g(x'')-f(x')g(x')|
%   &\leq |f(x'')g(x'')-f(x'')g(x')|\\
%   &\quad+|f(x'')g(x')-f(x')g(x')|\\
%   &\leq |f(x'')||g(x'')-g(x')|+|g(x')||f(x'')-f(x')|
% \end{aligned}
% \]
% 推出
% \[
%   \omega_i(fg)\leq M[\omega_i(f)+\omega_i(g)]
%   \quad(i=1,2,\cdots,n).
% \]
% 因此
% \[
%   \sum_{i=1}^n\omega_i(fg)\Delta x_i
%   \leq M\left[\sum_{i=1}^n\omega_i(f)\Delta x_i+\sum_{i=1}^n\omega_i(g)\Delta x_i\right]<\varepsilon.
% \qedhere
% \]
% \end{proof}
% 
% 
% 容易看出, 当 $f(x),g(x)\in R[a,b]$ 时, $\dfrac{f(x)}{g(x)}$ 未必在 $[a,b]$ 上可积(如此时 $\dfrac{f(x)}{g(x)}$ 可以是无界函数); 若 $f(x)$ 与 $g(x)$ 可以复合, $f(g(x))$ 在 $[a,b]$ 上也不一定可积。例如, 取 $g(x)=R(x)$, 其中 $R(x)$ 是区间 $[0,1]$ 上的黎曼函数, 而取
% \[
%   f(x)=
%   \begin{cases}
%     1, & x\neq0,\\
%     0, & x=0,
%   \end{cases}
% \]
% 则 $f(g(x))$ 为区间 $[0,1]$ 上的狄利克雷函数。因此 $f(g(x))$ 在 $[0,1]$ 上不可积。但当 $f(x)$ 连续时, 若 $g(x)$ 是 $[a,b]$ 上的可积函数并且 $f(g(x))$ 在 $[a,b]$ 上有定义, 则必定有 $f(g(x))\in R[a,b]$(见本章习题).
% 
% \begin{example}\label{example:ma-7-3-2}
% 设函数 $f(x)$ 在区间 $[a,b]$ 上连续, 证明:
% 
% (1) 如果对 $[a,b]$ 的任何子区间 $[a_1,b_1]$ 有 $\int_{a_1}^{b_1}f(x)\mathrm dx=0$, 则 $f(x)\equiv0,\ x\in[a,b]$;
% 
% (2) 若 $f(x)\geq0$, 并且 $\int_a^b f(x)\mathrm dx=0$, 则 $f(x)\equiv0,\ x\in[a,b]$.
% \end{example}
% 
% 
% \begin{proof}
% (1) 用反证法。倘若 $f(x)$ 不恒为 $0$, 则存在 $x_0\in[a,b]$, 使得 $f(x_0)\neq0$。不妨设 $f(x_0)>0$, 且 $x_0\in(a,b)$。由 $f(x)$ 的连续性, 存在 $\delta_0>0$, 使得 $[x_0-\delta_0,x_0+\delta_0]\subset[a,b]$, 且当 $x\in[x_0-\delta_0,x_0+\delta_0]$ 时, 有
% \[
%   f(x)\geq\frac{f(x_0)}{2}.
% \]
% 因此
% \[
%   \int_{x_0-\delta_0}^{x_0+\delta_0} f(x)\mathrm dx
%   \geq \frac{f(x_0)}{2}\int_{x_0-\delta_0}^{x_0+\delta_0}\mathrm dx
%   =\delta_0f(x_0)>0.
% \]
% 这与已知矛盾。(1) 得证.
% 
% (2) 用反证法。当 $f(x)$ 不恒为 $0$ 时, 必存在 $x_0\in(a,b)$, 使得 $f(x_0)>0$。由 (1) 的证明过程知, 存在 $\delta_0>0$, 使得
% \[
%   \int_{x_0-\delta_0}^{x_0+\delta_0} f(x)\mathrm dx>0.
% \]
% 但此时有
% \[
%   \int_a^b f(x)\mathrm dx\geq \int_{x_0-\delta_0}^{x_0+\delta_0} f(x)\mathrm dx>0.
% \]
% 这与已知矛盾。(2) 得证.
% \end{proof}
% 
% 
% \begin{example}\label{example:ma-7-3-3}
% 设 $0<q\leq p$, 证明 $\ln\dfrac{p}{q}\leq\dfrac{p-q}{q}$.
% \end{example}
% 
% 
% \begin{proof}
% 当 $p=q$ 时, 结论显然成立。现在设 $q<p$。对于 $x\in[q,p]$, 显然有 $\dfrac{1}{x}\leq\dfrac{1}{q}$。注意到 $\dfrac{1}{x}\in C[q,p]$, 因此 $\dfrac{1}{x}\in R[q,p]$。再注意到 $\ln x$ 是 $\dfrac{1}{x}$ 的一个原函数, 因此对 $\dfrac{1}{x}\leq\dfrac{1}{q}$ 两边从 $q$ 到 $p$ 积分得
% \[
%   \ln\frac{p}{q}=\ln x\bigg|_q^p=\int_q^p\frac{1}{x}\mathrm dx
%   \leq \int_q^p\frac{1}{q}\mathrm dx=\frac{p-q}{q}.
% \qedhere
% \]
% \end{proof}
% 
% 
% \section{原函数的存在性与定积分的计算}
% \label{sec:ma-7-4}
% 
% \subsection{变限定积分}
% \label{subsec:ma-7-4-1}
% 
% 在求函数的不定积分时, 我们知道有些初等函数的不定积分不能由初等函数表示。另外, 当一个函数比较复杂时, 求不定积分也不是一件容易的事。因此我们很难知道某类函数是否存在原函数。利用定积分理论, 我们则可以知道对于非常广泛的函数类, 它们总是存在原函数的.
% 
% 设函数 $f(t)\in R[a,b]$, 则对于 $\forall x\in[a,b]$, $f(t)\in R[a,x]$。因此
% \[
%   \Phi(x)=\int_a^x f(t)\mathrm dt
% \]
% 是定义在区间 $[a,b]$ 上的一个函数。我们称 $\Phi(x)$ 为 $f(t)$ 的变上限定积分。同样我们也可以在 $[a,b]$ 上定义 $f(t)$ 的变下限定积分
% \[
%   \Psi(x)=\int_x^b f(t)\mathrm dt\quad(x\in[a,b]).
% \]
% 下面我们将要看到 $\Phi(x),\Psi(x)$ 总比 $f(t)$ 具有更好的性质。在这里我们只讨论 $\Phi(x)$ 的情况, 对于 $\Psi(x)$ 的情况请读者自己讨论.
% 
% \begin{theorem}\label{theorem:ma-7-4-1}
% 设 $\Phi(x)$ 是函数 $f(t)\in R[a,b]$ 的变上限定积分.
% 
% (1) 若 $f(t)\in R[a,b]$, 则 $\Phi(x)\in C[a,b]$;
% 
% (2) 若 $f(t)\in C[a,b]$, 则 $\Phi(x)$ 在 $[a,b]$ 上可导, 并且 $\Phi'(x)=f(x)$(在端点处为单侧导数).
% \end{theorem}
% 
% 
% \begin{proof}
% (1) 由 $f(t)\in R[a,b]$, 存在常数 $M>0$, 使得对一切 $t\in[a,b]$, 有 $|f(t)|\leq M$。任取 $x_0\in[a,b]$, 则当 $x\in[a,b]$ 时, 有
% \[
% \begin{aligned}
%   |\Phi(x)-\Phi(x_0)|
%   &=\left|\int_a^x f(t)\mathrm dt-\int_a^{x_0} f(t)\mathrm dt\right|\\
%   &=\left|\int_{x_0}^x f(t)\mathrm dt\right|
%   \leq \left|\int_{x_0}^x |f(t)|\mathrm dt\right|\\
%   &\leq M|x-x_0|.
% \end{aligned}
% \]
% 由此推知 $\Phi(x)\in C[a,b]$.
% 
% (2) 现设 $f(t)\in C[a,b]$。下面我们证明 $\Phi(x)$ 在 $x_0\in(a,b)$ 处满足 $\Phi'(x_0)=f(x_0)$(对于 $x_0=a$ 或 $b$ 的情形请读者自证)。由 $f(t)$ 的连续性, 对于 $\forall\varepsilon>0$, $\exists\delta>0$(不妨设 $\delta<\min\{b-x_0,x_0-a\}$), 当 $|t-x_0|<\delta$ 时, 有
% \[
%   |f(t)-f(x_0)|<\varepsilon.
% \]
% 因此, 当 $0<|x-x_0|<\delta$ 时, 有
% \[
% \begin{aligned}
%   \left|\frac{\Phi(x)-\Phi(x_0)}{x-x_0}-f(x_0)\right|
%   &=\left|\frac{\int_a^x f(t)\mathrm dt-\int_a^{x_0} f(t)\mathrm dt}{x-x_0}
%   -\frac{\int_{x_0}^x f(x_0)\mathrm dt}{x-x_0}\right|\\
%   &=\left|\frac{\int_{x_0}^x [f(t)-f(x_0)]\mathrm dt}{x-x_0}\right|\\
%   &\leq \frac{\left|\int_{x_0}^x |f(t)-f(x_0)|\mathrm dt\right|}{|x-x_0|}\\
%   &<\varepsilon.
% \end{aligned}
% \]
% 这就证明了
% \[
%   \Phi'(x_0)=\lim_{x\to x_0}\frac{\Phi(x)-\Phi(x_0)}{x-x_0}=f(x_0).
% \qedhere
% \]
% \end{proof}
% 
% 
% 由上述定理我们知道, 区间 $[a,b]$ 上的连续函数 $f(x)$ 总存在原函数, 且其变上限定积分 $\Phi(x)=\int_a^x f(t)\mathrm dt(x\in[a,b])$ 即是它的一个原函数。在以上定理的证明中, 我们实际上证明了更强的结论: 若函数 $f(t)\in R[a,b]$, 并且在 $x_0\in[a,b]$ 处连续, 则其变上限定积分 $\Phi(x)$ 在 $x_0$ 处可导, 并且有 $\Phi'(x_0)=f(x_0)$.
% 
% \paragraph*{注 1}
% 若假定函数 $f(x)$ 在区间 $[a,b]$ 上连续, 从定理~\ref{theorem:ma-7-4-1} 可以简单地推导出牛顿--莱布尼茨公式。事实上, 设 $F(x)$ 是 $f(x)$ 的任一原函数, 则 $F(x)$ 与 $\int_a^x f(t)\mathrm dt$ 同为 $f(x)$ 的原函数, 因此存在常数 $C$, 使得
% \[
%   F(x)=\int_a^x f(t)\mathrm dt+C.
% \]
% 在上式中, 令 $x=a$, 得 $F(a)=C$; 令 $x=b$, 得
% \[
%   \int_a^b f(t)\mathrm dt=F(b)-F(a).
% \]
% 再将积分变量换回 $x$ 即得牛顿--莱布尼茨公式.
% 
% \paragraph*{注 2}
% 牛顿--莱布尼茨公式是定积分计算的基本方法。但值得指出的是: 即使一个函数 $f(x)$ 存在原函数 $F(x)$, 有时候 $f(x)$ 也不可积。例如
% \[
%   F(x)=
%   \begin{cases}
%     x^2\sin\dfrac{1}{x^2}, & x\neq0,\\
%     0, & x=0,
%   \end{cases}
% \]
% 它是
% \[
%   f(x)=
%   \begin{cases}
%     2x\sin\dfrac{1}{x^2}-\dfrac{2}{x}\cos\dfrac{1}{x^2}, & x\neq0,\\
%     0, & x=0
%   \end{cases}
% \]
% 的原函数。显然 $f(x)$ 在 $x=0$ 的邻域内无界, 因此 $f(x)$ 在包含 $x=0$ 的任何闭区间 $[a,b]$ 上不可积。对于这类存在原函数的无界函数积分问题, 在下一章中我们还将仔细研究.
% 
% 下面我们进一步考虑积分上(下)限是函数时变限定积分的求导问题.
% 
% 若函数 $f(t)\in C[a,b]$, $u=\varphi(x)$ 在区间 $[\alpha,\beta]$ 上可导, 且对于 $\forall x\in[\alpha,\beta]$, $\varphi(x)\in[a,b]$, 则对于 $\forall x\in[\alpha,\beta]$, $\int_a^{\varphi(x)} f(t)\mathrm dt$ 在 $[\alpha,\beta]$ 上有定义, 因此它是 $[\alpha,\beta]$ 上的一个函数。若令
% \[
%   G(x)=\int_a^{\varphi(x)} f(t)\mathrm dt,
% \]
% 则 $G(x)$ 可以看成 $\Phi(u)=\int_a^u f(t)\mathrm dt$ 与 $u=\varphi(x)$ 的复合函数。因此
% \[
%   G'(x)=\Phi'(\varphi(x))\varphi'(x)=f(\varphi(x))\varphi'(x).
% \]
% 类似地, 我们也可考虑积分下限是函数的情况.
% 
% \begin{example}\label{example:ma-7-4-1}
% 求导数 $\dfrac{\mathrm d}{\mathrm dx}\left(\int_x^{x^2} e^{-t^2}\mathrm dt\right)$.
% \end{example}
% 
% 
% \begin{solve}
% 由于 $e^{-x^2}\in C(-\infty,+\infty)$, 因此它在任何有限闭区间上均可积。显然有
% \[
%   \int_x^{x^2} e^{-t^2}\mathrm dt=\int_0^{x^2} e^{-t^2}\mathrm dt-\int_0^x e^{-t^2}\mathrm dt.
% \]
% 令 $\Phi(u)=\int_0^u e^{-t^2}\mathrm dt$, $u=x^2$, 则应用复合函数求导法则得
% \[
% \begin{aligned}
%   \frac{\mathrm d}{\mathrm dx}\left(\int_x^{x^2} e^{-t^2}\mathrm dt\right)
%   &=\frac{\mathrm d}{\mathrm dx}\left(\int_0^{x^2} e^{-t^2}\mathrm dt-\int_0^x e^{-t^2}\mathrm dt\right)\\
%   &=\left.\frac{\mathrm d(\Phi(u))}{\mathrm du}\right|_{u=x^2}\frac{\mathrm du}{\mathrm dx}-e^{-x^2}\\
%   &=2xe^{-x^4}-e^{-x^2}.
% \end{aligned}
% \qedhere
% \]
% \end{solve}
% 
% 
% \begin{example}\label{example:ma-7-4-2}
% 求极限 $\lim\limits_{x\to0}\dfrac{\int_0^{x^2}\sin t\mathrm dt}{x^4}$.
% \end{example}
% 
% 
% \begin{solve}
% 此极限为 $\dfrac{0}{0}$ 型不定式。利用洛必达法则得
% \[
%   \lim_{x\to0}\frac{\int_0^{x^2}\sin t\mathrm dt}{x^4}
%   =\lim_{x\to0}\frac{2x\sin x^2}{4x^3}=\frac{1}{2}.
% \qedhere
% \]
% \end{solve}
% 
% 
% \subsection{定积分的计算}
% \label{subsec:ma-7-4-2}
% 
% 从牛顿--莱布尼茨公式知, 对于连续函数的定积分, 只要能找到它的一个原函数, 则定积分的问题即可解决。因此计算不定积分的公式以及求不定积分的方法均可直接应用到定积分的计算中.
% 
% \begin{example}\label{example:ma-7-4-3}
% 求 $I=\int_{-\frac{\pi}{2}}^{\frac{\pi}{2}} f(x)\mathrm dx$, 其中
% \[
%   f(x)=
%   \begin{cases}
%     -\sin x, & x\geq0,\\
%     x, & x<0.
%   \end{cases}
% \]
% \end{example}
% 
% 
% \begin{solve}
% 注意到 $f(x)$ 是分段函数, 我们有
% \[
% \begin{aligned}
%   I&=\int_{-\frac{\pi}{2}}^0 x\mathrm dx-\int_0^{\frac{\pi}{2}}\sin x\mathrm dx\\
%   &=\left.\frac{x^2}{2}\right|_{-\frac{\pi}{2}}^0+\left.\cos x\right|_0^{\frac{\pi}{2}}
%   =-\left(\frac{\pi^2}{8}+1\right).
% \end{aligned}
% \]
% 尽管定积分的计算可以先由不定积分找出原函数, 然后计算原函数端点值的差而得到, 但在定积分的计算中, 有时要比不定积分来得简单。如在不定积分的换元法中必须将所换变量换回原来的变量, 但在定积分的换元法中, 则无此必要.
% \end{solve}
% 
% 
% \begin{theorem}[{换元法}]\label{theorem:ma-7-4-2}
% 设函数 $f(x)\in C[a,b]$, $\varphi(t)$ 在区间 $[\alpha,\beta]$ 上具有连续导数, 且 $\varphi(\alpha)=a,\varphi(\beta)=b$, $a\leq\varphi(t)\leq b(\alpha\leq t\leq\beta)$, 则有
% \[
%   \int_a^b f(x)\mathrm dx=\int_\alpha^\beta f(\varphi(t))\varphi'(t)\mathrm dt.
% \]
% \end{theorem}
% 
% 
% \begin{proof}
% 一方面, 由于 $f(x)\in C[a,b]$, 因此它在区间 $[a,b]$ 上存在原函数 $F(x)$。由牛顿--莱布尼茨公式, 我们有
% \[
%   \int_a^b f(x)\mathrm dx=F(b)-F(a).
% \]
% 另一方面, $F(\varphi(t))$ 是连续函数 $f(\varphi(t))\varphi'(t)$ 在 $[\alpha,\beta]$ 上的一个原函数, 因此也有
% \[
% \begin{aligned}
%   \int_\alpha^\beta f(\varphi(t))\varphi'(t)\mathrm dt
%   &=\left.F(\varphi(t))\right|_\alpha^\beta
%   =F(\varphi(\beta))-F(\varphi(\alpha))\\
%   &=F(b)-F(a).
% \end{aligned}
% \qedhere
% \]
% \end{proof}
% 
% 
% \begin{remark}\label{remark:ma-7-source-58}
% 若在定理~\ref{theorem:ma-7-4-2} 中, 假定函数 $\varphi(t)$ 在区间 $[\alpha,\beta]$ 上具有连续导数, 且 $\varphi(\alpha)=b,\varphi(\beta)=a,a\leq\varphi(t)\leq b(\alpha\leq t\leq\beta)$, 则有
% \[
%   \int_a^b f(x)\mathrm dx=\int_\beta^\alpha f(\varphi(t))\varphi'(t)\mathrm dt.
% \]
% 另外, 函数 $\varphi'(t)$ 在 $[\alpha,\beta]$ 上连续的条件可以减弱为 $\varphi'(t)\in R[\alpha,\beta]$。关于此结果的证明请读者自己给出.
% \end{remark}
% 
% 
% \begin{example}\label{example:ma-7-4-4}
% 设函数 $f(x)\in C[-a,a](a>0)$, 证明:
% 
% (1) 若 $f(x)$ 为偶函数, 则 $\int_{-a}^a f(x)\mathrm dx=2\int_0^a f(x)\mathrm dx$;
% 
% (2) 若 $f(x)$ 为奇函数, 则 $\int_{-a}^a f(x)\mathrm dx=0$.
% \end{example}
% 
% 
% \begin{proof}
% 在 $\int_{-a}^0 f(x)\mathrm dx$ 中令 $x=-t$, 则有
% \[
%   \int_{-a}^0 f(x)\mathrm dx=-\int_a^0 f(-t)\mathrm dt=\int_0^a f(-x)\mathrm dx.
% \]
% 当 $f(x)$ 是偶函数时, 我们有 $f(-x)=f(x)(x\in[-a,a])$。因此由定积分的性质有
% \[
% \begin{aligned}
%   \int_{-a}^a f(x)\mathrm dx
%   &=\int_{-a}^0 f(x)\mathrm dx+\int_0^a f(x)\mathrm dx\\
%   &=\int_0^a f(-x)\mathrm dx+\int_0^a f(x)\mathrm dx
%   =2\int_0^a f(x)\mathrm dx.
% \end{aligned}
% \]
% (1) 得证.
% 
% 当 $f(x)$ 是奇函数时, 我们有 $f(-x)=-f(x)(x\in[-a,a])$。因此由定积分的性质有
% \[
% \begin{aligned}
%   \int_{-a}^a f(x)\mathrm dx
%   &=\int_{-a}^0 f(x)\mathrm dx+\int_0^a f(x)\mathrm dx\\
%   &=\int_0^a f(-x)\mathrm dx+\int_0^a f(x)\mathrm dx\\
%   &=-\int_0^a f(x)\mathrm dx+\int_0^a f(x)\mathrm dx=0.
% \end{aligned}
% \]
% (2) 得证.
% \end{proof}
% 
% 
% \begin{example}\label{example:ma-7-4-5}
% 求定积分 $\int_0^a \sqrt{a^2-x^2}\mathrm dx$, $a>0$.
% \end{example}
% 
% 
% \begin{solve}
% 令 $x=a\sin t\left(t\in\left[0,\dfrac{\pi}{2}\right]\right)$, 则 $x=0\Rightarrow t=0$, $x=a\Rightarrow t=\dfrac{\pi}{2}$。由此得
% \[
% \begin{aligned}
%   \int_0^a \sqrt{a^2-x^2}\mathrm dx
%   &=\int_0^{\frac{\pi}{2}} a\sqrt{a^2\cos^2 t}\cos t\mathrm dt\\
%   &=a^2\int_0^{\frac{\pi}{2}}\cos^2 t\mathrm dt
%   =a^2\int_0^{\frac{\pi}{2}}\frac{1+\cos 2t}{2}\mathrm dt\\
%   &=\frac{a^2}{2}\left.\left(t+\frac{1}{2}\sin 2t\right)\right|_0^{\frac{\pi}{2}}
%   =\frac{\pi}{4}a^2.
% \end{aligned}
% \qedhere
% \]
% \end{solve}
% 
% 
% \begin{example}\label{example:ma-7-4-6}
% 设 $f(x)$ 是以 $T(T>0)$ 为周期的连续函数, 证明: 对于任意的 $a$, 有
% \[
%   \int_a^{a+T} f(x)\mathrm dx=\int_0^T f(x)\mathrm dx.
% \]
% \end{example}
% 
% 
% \begin{proof}
% 由定积分的性质, 我们有
% \begin{equation}
%   \int_a^{a+T} f(x)\mathrm dx
%   =\int_a^0 f(x)\mathrm dx+\int_0^T f(x)\mathrm dx+\int_T^{a+T} f(x)\mathrm dx.
%   \label{eq:ma-7-3-7}
% \end{equation}
% 令 $x=t+T$, 则有
% \[
%   \int_T^{a+T} f(x)\mathrm dx
%   =\int_0^a f(t+T)\mathrm dt
%   =\int_0^a f(t)\mathrm dt
%   =-\int_a^0 f(t)\mathrm dt.
% \]
% 将此等式中的积分变量 $t$ 换回 $x$ 后, 我们有
% \[
%   \int_T^{a+T} f(x)\mathrm dx=-\int_a^0 f(x)\mathrm dx.
% \]
% 将此等式代入 \eqref{eq:ma-7-3-7} 即得所证.
% \end{proof}
% 
% 
% \begin{example}\label{example:ma-7-4-7}
% 求定积分 $I=\int_0^\pi \dfrac{x\sin x}{1+\cos^2 x}\mathrm dx$.
% \end{example}
% 
% 
% \begin{solve}
% 由定积分的性质, 我们有
% \[
%   I=\int_0^{\frac{\pi}{2}}\frac{x\sin x}{1+\cos^2 x}\mathrm dx
%   +\int_{\frac{\pi}{2}}^\pi \frac{x\sin x}{1+\cos^2 x}\mathrm dx.
% \]
% 在 $\int_{\frac{\pi}{2}}^\pi \dfrac{x\sin x}{1+\cos^2 x}\mathrm dx$ 中令 $x=\pi-t$, 则 $x=\dfrac{\pi}{2}\Rightarrow t=\dfrac{\pi}{2},x=\pi\Rightarrow t=0$, 从而有
% \[
% \begin{aligned}
%   \int_{\frac{\pi}{2}}^\pi \frac{x\sin x}{1+\cos^2 x}\mathrm dx
%   &=-\int_{\frac{\pi}{2}}^0
%   \frac{(\pi-t)\sin(\pi-t)}{1+\cos^2(\pi-t)}\mathrm dt\\
%   &=\pi\int_0^{\frac{\pi}{2}}\frac{\sin x}{1+\cos^2 x}\mathrm dx
%   -\int_0^{\frac{\pi}{2}}\frac{x\sin x}{1+\cos^2 x}\mathrm dx.
% \end{aligned}
% \]
% 最后我们得到
% \[
%   I=\pi\int_0^{\frac{\pi}{2}}\frac{\sin x}{1+\cos^2 x}\mathrm dx
%   =-\left.\pi\arctan(\cos x)\right|_0^{\frac{\pi}{2}}
%   =\frac{\pi^2}{4}.
% \qedhere
% \]
% \end{solve}
% 
% 
% \begin{example}\label{example:ma-7-4-8}
% 设 $x_1,x_2\in(-1,1)$, 定义
% \[
%   d(x_1,x_2)=\left|\int_{x_1}^{x_2}\frac{\mathrm dx}{1-x^2}\right|.
% \]
% 试证明以下结论:
% 
% (1) 对任何 $x_1,x_2\in(-1,1)$, 有 $d(x_1,x_2)\geq0$, 且
% \[
%   d(x_1,x_2)=0\Longleftrightarrow x_1=x_2;
% \]
% 
% (2) $d(x_1,x_2)=d(x_2,x_1)$;
% 
% (3) 对任何 $x_1,x_2,x_3\in(-1,1)$, 成立
% \[
%   d(x_1,x_3)\leq d(x_1,x_2)+d(x_2,x_3);
% \]
% 
% (4) 设 $x_0,x_1,x_2\in(-1,1)$, $w_j=\dfrac{x_j-x_0}{1-x_0x_j}(j=1,2)$, 则 $w_1,w_2\in(-1,1)$, 并且 $d(w_1,w_2)=d(x_1,x_2)$;
% 
% (5) 设 $x_0\in(-1,1)$, 则 $\lim\limits_{x\to\pm1}d(x_0,x)=+\infty$.
% \end{example}
% 
% 
% \begin{proof}
% (1) 当 $x_1=x_2$ 时, 显然 $d(x_1,x_2)=0$。当 $x_1\neq x_2$ 时, 有
% \[
%   d(x_1,x_2)
%   =\left|\int_{x_1}^{x_2}\frac{\mathrm dx}{1-x^2}\right|
%   =\left|\frac{1}{2}\left(\ln\frac{1+x_2}{1-x_2}-\ln\frac{1+x_1}{1-x_1}\right)\right|>0.
% \]
% 因此若 $d(x_1,x_2)=0$, 必有 $x_1=x_2$, 从而 (1) 成立.
% 
% (2) 由于 $\int_{x_1}^{x_2}\dfrac{\mathrm dx}{1-x^2}=-\int_{x_2}^{x_1}\dfrac{\mathrm dx}{1-x^2}$, 两边取绝对值即得 (2).
% 
% (3) 当以 $x_1,x_3$ 为端点的开区间包含在以 $x_1,x_2$ 或 $x_2,x_3$ 为端点的开区间时, 显然结论成立。现不妨设 $x_1<x_2<x_3$, 则有
% \[
% \begin{aligned}
%   d(x_1,x_3)
%   &=\int_{x_1}^{x_3}\frac{\mathrm dx}{1-x^2}
%   =\int_{x_1}^{x_2}\frac{\mathrm dx}{1-x^2}
%   +\int_{x_2}^{x_3}\frac{\mathrm dx}{1-x^2}\\
%   &=d(x_1,x_2)+d(x_2,x_3).
% \end{aligned}
% \]
% 
% (4) 当 $x_0,x_1\in(-1,1)$ 时, 我们有
% \[
%   0<(1-x_0^2)(1-x_1^2)=1-(x_0^2+x_1^2)+x_0^2x_1^2,
% \]
% 从而有
% \[
%   (x_1-x_0)^2=x_1^2+x_0^2-2x_0x_1<1+x_0^2x_1^2-2x_0x_1=(1-x_0x_1)^2.
% \]
% 因此有 $w_1\in(-1,1)$。同理有 $w_2\in(-1,1)$.
% 
% 当 $x_0\in(-1,1)$ 时, $w(x)=\dfrac{x-x_0}{1-x_0x}$ 在 $x\in(-1,1)$ 时有定义且有
% \[
%   w'(x)=\frac{1-x_0^2}{(1-x_0x)^2}>0,
% \]
% 因此 $w(x)$ 在 $(-1,1)$ 内严格单调上升。注意到我们曾经证明过等式
% \[
%   \frac{\mathrm dw}{1-w^2}=\frac{\mathrm dx}{1-x^2},
% \]
% 因此由定积分的换元法有
% \[
%   d(x_1,x_2)=\left|\int_{x_1}^{x_2}\frac{\mathrm dx}{1-x^2}\right|
%   =\left|\int_{w_1}^{w_2}\frac{\mathrm dw}{1-w^2}\right|
%   =d(w_1,w_2).
% \]
% 
% (5) 直接计算可得
% \[
% \begin{aligned}
%   \lim_{x\to\pm1}d(x_0,x)
%   &=\lim_{x\to\pm1}\left|\int_{x_0}^x\frac{\mathrm dx}{1-x^2}\right|\\
%   &=\lim_{x\to\pm1}\frac{1}{2}
%   \left|\ln\frac{1+x}{1-x}-\ln\frac{1+x_0}{1-x_0}\right|\\
%   &=+\infty.
% \end{aligned}
% \qedhere
% \]
% \end{proof}
% 
% 
% \begin{remark}\label{remark:ma-7-source-69}
% 若在区间 $(-1,1)$ 内如上定义两点间距离 $d$, 则称该距离为庞加莱 (Poincaré) 距离或双曲距离。容易看出, 在该距离下, $(-1,1)$ 内的柯西序列都是收敛的, 并且其极限是 $(-1,1)$ 内的点。因此在该距离下, $(-1,1)$ 是完备的.
% \end{remark}
% 
% 
% \begin{theorem}[{分部积分法}]\label{theorem:ma-7-4-3}
% 设函数 $u(x),v(x)$ 在区间 $[a,b]$ 上可导, 并且 $u'(x),v'(x)\in R[a,b]$, 则
% \[
%   \int_a^b u(x)v'(x)\mathrm dx
%   =\left.u(x)v(x)\right|_a^b-\int_a^b u'(x)v(x)\mathrm dx.
% \]
% \end{theorem}
% 
% 
% \begin{proof}
% 在恒等式 $[u(x)v(x)]'=u'(x)v(x)+u(x)v'(x)$ 中, 我们知等式两边的函数在 $[a,b]$ 上均可积, 因此两边从 $a$ 到 $b$ 对 $x$ 积分并利用牛顿--莱布尼茨公式得
% \[
%   \left.u(x)v(x)\right|_a^b
%   =\int_a^b [u(x)v(x)]'\mathrm dx
%   =\int_a^b u'(x)v(x)\mathrm dx+\int_a^b u(x)v'(x)\mathrm dx.
% \]
% 将上述等式整理即得所证.
% \end{proof}
% 
% 
% \begin{example}\label{example:ma-7-4-9}
% 计算定积分 $I=\int_0^1 x\ln(1+x)\mathrm dx$.
% \end{example}
% 
% 
% \begin{solve}
% \[
% \begin{aligned}
%   I&=\int_0^1 \ln(1+x)\mathrm d\left(\frac{x^2}{2}\right)
%   =\left.\frac{x^2}{2}\ln(1+x)\right|_0^1
%   -\int_0^1 \frac{x^2}{2}\mathrm d(\ln(1+x))\\
%   &=\frac{1}{2}\ln2-\frac{1}{2}\int_0^1\left(x-1+\frac{1}{1+x}\right)\mathrm dx\\
%   &=\frac{1}{2}\ln2-\frac{1}{2}\left.\left[\frac{x^2}{2}-x+\ln(1+x)\right]\right|_0^1
%   =\frac{1}{4}.
% \end{aligned}
% \qedhere
% \]
% \end{solve}
% 
% 
% \begin{example}\label{example:ma-7-4-10}
% 证明: 对任意的 $n\in\mathbb N$, 有
% \[
%   \int_0^{\frac{\pi}{2}}\sin^n x\mathrm dx
%   =\int_0^{\frac{\pi}{2}}\cos^n x\mathrm dx\triangleq I_n,
% \]
% 并求 $I_n$ 的值.
% \end{example}
% 
% 
% \begin{proof}
% 对 $\int_0^{\frac{\pi}{2}}\sin^n x\mathrm dx$ 作变换 $x=\dfrac{\pi}{2}-t$, 得
% \[
%   \int_0^{\frac{\pi}{2}}\sin^n x\mathrm dx
%   =\int_{\frac{\pi}{2}}^0 \cos^n t\,\mathrm d\left(\frac{\pi}{2}-t\right)
%   =-\int_{\frac{\pi}{2}}^0 \cos^n t\mathrm dt
%   =\int_0^{\frac{\pi}{2}}\cos^n x\mathrm dx.
% \]
% 下面我们来求 $I_n$ 的递推公式。对于 $n=0,1$, 我们有
% \[
%   I_0=\int_0^{\frac{\pi}{2}}\mathrm dx=\frac{\pi}{2},\qquad
%   I_1=\int_0^{\frac{\pi}{2}}\sin x\mathrm dx=1.
% \]
% 当 $n>1$ 时, 有
% \[
% \begin{aligned}
%   I_n&=\int_0^{\frac{\pi}{2}}\sin^{n-1}x\,\mathrm d(-\cos x)\\
%   &=-\left.\sin^{n-1}x\cos x\right|_0^{\frac{\pi}{2}}
%   +\int_0^{\frac{\pi}{2}}\cos x\,\mathrm d(\sin^{n-1}x)\\
%   &=(n-1)\int_0^{\frac{\pi}{2}}\sin^{n-2}x(1-\sin^2x)\mathrm dx\\
%   &=(n-1)I_{n-2}-(n-1)I_n,
% \end{aligned}
% \]
% 即 $I_n=\dfrac{n-1}{n}I_{n-2}$。因此得
% \[
%   I_{2k}=\frac{(2k-1)!!}{(2k)!!}\cdot\frac{\pi}{2},\qquad
%   I_{2k+1}=\frac{(2k)!!}{(2k+1)!!}\quad(k=1,2,\cdots).
% \]
% 作为上述定积分的应用, 可以推出下面的瓦利斯 (Wallis) 公式:
% \[
%   \lim_{n\to\infty}\left[\frac{(2n)!!}{(2n-1)!!}\right]^2\frac{1}{2n+1}=\frac{\pi}{2}.
% \]
% 事实上, 当 $x\in\left(0,\dfrac{\pi}{2}\right)$ 时, 对任意 $n\in\mathbb N$, 有
% \[
%   \sin^{2n+1}x<\sin^{2n}x<\sin^{2n-1}x,
% \]
% 从而有
% \[
%   \int_0^{\frac{\pi}{2}}\sin^{2n+1}x\mathrm dx
%   <\int_0^{\frac{\pi}{2}}\sin^{2n}x\mathrm dx
%   <\int_0^{\frac{\pi}{2}}\sin^{2n-1}x\mathrm dx.
% \]
% 利用例~\ref{example:ma-7-4-10} 有
% \[
%   \frac{(2n)!!}{(2n+1)!!}
%   <\frac{(2n-1)!!}{(2n)!!}\cdot\frac{\pi}{2}
%   <\frac{(2n-2)!!}{(2n-1)!!}.
% \]
% 从上式容易推出
% \[
%   J_n\triangleq\left[\frac{(2n)!!}{(2n-1)!!}\right]^2\frac{1}{2n+1}
%   <\frac{\pi}{2}
%   <\left[\frac{(2n)!!}{(2n-1)!!}\right]^2\frac{1}{2n}\triangleq J'_n.
% \]
% 由于
% \[
% \begin{aligned}
%   0&\leq\lim_{n\to\infty}(J'_n-J_n)\\
%   &=\lim_{n\to\infty}
%   \left[\frac{(2n)!!}{(2n-1)!!}\right]^2\frac{1}{2n(2n+1)}\\
%   &\leq\lim_{n\to\infty}\frac{1}{2n}\cdot\frac{\pi}{2}=0,
% \end{aligned}
% \]
% 因此瓦利斯公式成立.
% \end{proof}
% 
% 
% \section{定积分中值定理}
% \label{sec:ma-7-5}
% 
% 在定积分的理论中, 我们最后介绍定积分的两个重要性质.
% 
% \subsection{定积分第一中值定理}
% \label{subsec:ma-7-5-1}
% 
% 设函数 $f(x)$ 在区间 $[a,b]$ 上连续, 由定积分的几何意义知, 必存在一点 $\xi\in[a,b]$, 使得
% \[
%   \int_a^b f(x)\mathrm dx=f(\xi)\int_a^b\mathrm dx=f(\xi)(b-a).
% \]
% 这里 $\int_a^b f(x)\mathrm dx$ 可以看成函数 $f(x)g(x)=f(x)\cdot1$ 在 $[a,b]$ 上的定积分, 其中 $g(x)\equiv1$。一个自然的问题是: 函数 $g(x)\equiv1$ 能否推广成一般的函数? 定积分第一中值定理将回答这个问题.
% 
% \begin{theorem}[{定积分第一中值定理}]\label{theorem:ma-7-5-1}
% 设函数 $f(x)\in C[a,b],g(x)\in R[a,b]$ 且在区间 $[a,b]$ 上不变号, 则存在 $\xi\in[a,b]$, 使得
% \[
%   \int_a^b f(x)g(x)\mathrm dx=f(\xi)\int_a^b g(x)\mathrm dx.
% \]
% 特别地, 当 $g(x)\equiv1$ 时, 有 $\int_a^b f(x)\mathrm dx=f(\xi)(b-a)$.
% \end{theorem}
% 
% 
% \begin{proof}
% 不妨设 $g(x)\geq0$, 并令 $m,M$ 分别是函数 $f(x)$ 在区间 $[a,b]$ 上的最小值和最大值, 则对于 $\forall x\in[a,b]$, 有
% \[
%   mg(x)\leq f(x)g(x)\leq Mg(x).
% \]
% 因此有
% \[
%   m\int_a^b g(x)\mathrm dx
%   \leq\int_a^b f(x)g(x)\mathrm dx
%   \leq M\int_a^b g(x)\mathrm dx.
% \]
% 若 $\int_a^b g(x)\mathrm dx=0$, 由上式得 $\int_a^b f(x)g(x)\mathrm dx=0$, 此时任取 $\xi\in[a,b]$ 即可。当 $\int_a^b g(x)\mathrm dx\neq0$ 时, 则有
% \[
%   m\leq
%   \frac{\int_a^b f(x)g(x)\mathrm dx}{\int_a^b g(x)\mathrm dx}
%   \leq M.
% \]
% 由连续函数的介值定理, 必存在 $\xi\in[a,b]$, 使得
% \[
%   f(\xi)=
%   \frac{\int_a^b f(x)g(x)\mathrm dx}{\int_a^b g(x)\mathrm dx}.
% \qedhere
% \]
% \end{proof}
% 
% 
% 当 $g(x)\equiv1$ 且 $f(x)\geq0$ 时, 定积分第一中值定理具有鲜明的几何意义, 即由直线 $x=a,x=b,y=0$ 和曲线 $y=f(x)$ 所围的曲边梯形的面积必与由直线 $x=a,x=b,y=0,y=f(\xi)(\xi\in[a,b])$ 围成的矩形面积相等 (见图 7.5.1).
% 
% \paragraph*{注 1}
% 在定积分第一中值定理中, 若只假定 $f(x)\in R[a,b]$, 令
% \[
%   m=\inf_{x\in[a,b]}\{f(x)\},\qquad
%   M=\sup_{x\in[a,b]}\{f(x)\},
% \]
% 则用同样的方法可以证明存在 $\mu\in[m,M]$, 使得
% \[
%   \int_a^b f(x)g(x)\mathrm dx=\mu\int_a^b g(x)\mathrm dx.
% \]
% 
% \paragraph*{注 2}
% 若在定积分第一中值定理中假定 $f(x)$ 与 $g(x)$ 均是连续函数, 且 $g(x)$ 在 $[a,b]$ 上不变号, 则可以证明必存在 $\xi\in(a,b)$, 使得
% \[
%   \int_a^b f(x)g(x)\mathrm dx=f(\xi)\int_a^b g(x)\mathrm dx
% \]
% (见本章习题).
% 
% \begin{example}\label{example:ma-7-5-1}
% 设函数 $f(x)\in C[0,1]$, 证明
% \[
%   \lim_{n\to\infty}\int_0^1\frac{nf(x)}{1+n^2x^2}\mathrm dx=\frac{\pi}{2}f(0).
% \]
% \end{example}
% 
% 
% \begin{proof}
% 由于函数 $f(x)\in C[0,1]$, 从而存在 $M>0$, 使得对于 $\forall x\in[0,1]$, 有 $|f(x)|\leq M$。由定积分的性质, 对于 $\forall n\in\mathbb N$, 我们有
% \[
%   \int_0^1\frac{nf(x)}{1+n^2x^2}\mathrm dx
%   =\int_0^{n^{-\frac13}}\frac{nf(x)}{1+n^2x^2}\mathrm dx
%   +\int_{n^{-\frac13}}^1\frac{nf(x)}{1+n^2x^2}\mathrm dx.
% \]
% 下面我们分别来估计上面等式右边中的两个定积分。对于第二个定积分我们有
% \[
%   \left|\int_{n^{-\frac13}}^1\frac{nf(x)}{1+n^2x^2}\mathrm dx\right|
%   \leq\int_{n^{-\frac13}}^1\left|\frac{nf(x)}{1+n^2x^2}\right|\mathrm dx
%   \leq\frac{nM}{1+n^{1+\frac13}}\to0\quad(n\to\infty).
% \]
% 而在第一个定积分 $\int_0^{n^{-\frac13}}\dfrac{nf(x)}{1+n^2x^2}\mathrm dx$ 中, $\dfrac{n}{1+n^2x^2}$ 在 $[0,1]$ 上不变号, 因此存在 $\xi\in[0,n^{-\frac13}]$, 使得
% \[
% \begin{aligned}
%   \int_0^{n^{-\frac13}}\frac{nf(x)}{1+n^2x^2}\mathrm dx
%   &=f(\xi)\int_0^{n^{-\frac13}}\frac{n}{1+n^2x^2}\mathrm dx\\
%   &=\left.f(\xi)\arctan nx\right|_0^{n^{-\frac13}}
%   =f(\xi)\arctan n^{\frac23}.
% \end{aligned}
% \]
% 当 $n\to\infty$ 时, 由于 $\xi\in[0,n^{-\frac13}]$, 有 $\xi\to0$, 而 $\arctan n^{\frac23}\to\dfrac{\pi}{2}$, 因此
% \[
%   \int_0^{n^{-\frac13}}\frac{nf(x)}{1+n^2x^2}\mathrm dx\to\frac{\pi}{2}f(0).
% \]
% 最终我们有
% \[
% \begin{aligned}
%   \lim_{n\to\infty}\int_0^1\frac{nf(x)}{1+n^2x^2}\mathrm dx
%   &=\lim_{n\to\infty}\int_0^{n^{-\frac13}}\frac{nf(x)}{1+n^2x^2}\mathrm dx
%   +\lim_{n\to\infty}\int_{n^{-\frac13}}^1\frac{nf(x)}{1+n^2x^2}\mathrm dx\\
%   &=\frac{\pi}{2}f(0).
% \end{aligned}
% \]
% 
% 利用定积分的分部积分法我们可以得到带积分余项的泰勒公式; 再利用定积分第一中值定理, 我们还可以从带积分余项的泰勒公式推出带拉格朗日余项和带柯西余项的泰勒公式.
% \end{proof}
% 
% 
% \begin{theorem}\label{theorem:ma-7-5-2}
% 设函数 $f(x)$ 在 $x_0$ 的邻域 $U(x_0,h)(h>0)$ 内具有 $n+1$ 阶连续导数 $f^{(n+1)}(x)$, 则对 $\forall x\in U(x_0,h)$, 有
% \begin{equation}
% \begin{aligned}
%   f(x)&=f(x_0)+f'(x_0)(x-x_0)+\frac{f''(x_0)}{2!}(x-x_0)^2\\
%   &\quad+\cdots+\frac{f^{(n)}(x_0)}{n!}(x-x_0)^n+R_n(x),
% \end{aligned}
% \label{eq:ma-7-5-1}
% \end{equation}
% 其中
% \[
%   R_n(x)=\frac{1}{n!}\int_{x_0}^x f^{(n+1)}(t)(x-t)^n\mathrm dt.
% \]
% \end{theorem}
% 
% 
% \begin{proof}
% 由牛顿--莱布尼茨公式有
% \[
%   f(x)=f(x_0)+\int_{x_0}^x f'(t)\mathrm dt
%   =f(x_0)+\int_{x_0}^x (x-t)^0f'(t)\mathrm dt.
% \]
% 利用分部积分我们有
% \[
% \begin{aligned}
%   \int_{x_0}^x (x-t)^0f'(t)\mathrm dt
%   &=-\int_{x_0}^x f'(t)\mathrm d(x-t)\\
%   &=-\left.f'(t)(x-t)\right|_{x_0}^x+\int_{x_0}^x (x-t)f''(t)\mathrm dt\\
%   &=f'(x_0)(x-x_0)+\int_{x_0}^x (x-t)f''(t)\mathrm dt\\
%   &=f'(x_0)(x-x_0)-\frac{1}{2}\int_{x_0}^x f''(t)\mathrm d(x-t)^2\\
%   &=f'(x_0)(x-x_0)-\left.\frac{f''(t)}{2}(x-t)^2\right|_{x_0}^x
%   +\frac{1}{2}\int_{x_0}^x (x-t)^2f'''(t)\mathrm dt\\
%   &=f'(x_0)(x-x_0)+\frac{f''(x_0)}{2}(x-x_0)^2
%   +\frac{1}{2}\int_{x_0}^x (x-t)^2f'''(t)\mathrm dt.
% \end{aligned}
% \]
% 再将上式中的积分继续分部积分, 逐次类推即可得到所要结论.
% \end{proof}
% 
% 
% \paragraph*{注 1}
% 定理~\ref{theorem:ma-7-5-2} 中的带余项
% \[
%   R_n(x)=\frac{1}{n!}\int_{x_0}^x f^{(n+1)}(t)(x-t)^n\mathrm dt
% \]
% 的公式 \eqref{eq:ma-7-5-1} 称为带积分余项的泰勒公式, 该余项也称为积分余项。在定理~\ref{theorem:ma-7-5-2} 的条件下, 我们可以推出带拉格朗日余项的泰勒公式。事实上, 在 $x_0$ 与 $x$ 之间, $g(t)=(x-t)^n$ 是关于 $t$ 的连续函数并且是不变号的, 而 $f^{(n+1)}(t)$ 是连续的, 因此, 由定积分第一中值定理的注 2, 在 $x$ 与 $x_0$ 之间存在 $\xi$, 使得下式成立:
% \[
% \begin{aligned}
%   R_n(x)&=\frac{1}{n!}\int_{x_0}^x f^{(n+1)}(t)(x-t)^n\mathrm dt\\
%   &=\frac{f^{(n+1)}(\xi)}{n!}\int_{x_0}^x (x-t)^n\mathrm dt\\
%   &=\frac{f^{(n+1)}(\xi)}{(n+1)!}(x-x_0)^{n+1}.
% \end{aligned}
% \]
% 这就是在泰勒公式中的拉格朗日余项.
% 
% \paragraph*{注 2}
% 从定理~\ref{theorem:ma-7-5-2} 中我们还可以推出带以下柯西余项的泰勒公式, 即在定理~\ref{theorem:ma-7-5-2} 的条件下, 我们有
% \begin{equation}
%   R_n(x)=\frac{1}{n!}f^{(n+1)}[x_0+\theta(x-x_0)](1-\theta)^n(x-x_0)^{n+1},
%   \label{eq:ma-7-5-2}
% \end{equation}
% 其中 $0<\theta<1$。事实上, 对积分余项应用定积分第一中值定理的注 2 知, 在 $x$ 与 $x_0$ 之间存在 $\xi$, 使得下式成立:
% \[
%   R_n(x)=\frac{1}{n!}\int_{x_0}^x f^{(n+1)}(t)(x-t)^n\mathrm dt
%   =\frac{1}{n!}f^{(n+1)}(\xi)(x-\xi)^n(x-x_0).
% \]
% 注意到 $x-\xi=(x-x_0)-\theta(x-x_0)=(1-\theta)(x-x_0)$, 其中 $0<\theta<1$。将此代入上式整理即得 \eqref{eq:ma-7-5-2}.
% 
% 形如 \eqref{eq:ma-7-5-2} 的余项称为柯西余项, 带此余项的泰勒公式称为带柯西余项的泰勒公式.
% 
% \subsection{定积分第二中值定理}
% \label{subsec:ma-7-5-2}
% 
% 设 $f(x)$ 是区间 $[a,b]$ 上单调上升的函数且 $f(x)\geq0$, 则 $\int_a^b f(x)\mathrm dx\geq0$, 从而有
% \[
%   f(b)(b-b)\leq\int_a^b f(x)\mathrm dx\leq f(b)(b-a).
% \]
% 于是存在 $\xi\in[a,b]$, 使得
% \[
%   \int_a^b f(x)\mathrm dx=f(b)(b-\xi)=f(b)\int_\xi^b\mathrm dx.
% \]
% 因此人们不禁要问: 对区间 $[a,b]$ 上的可积函数 $g(x)$, 是否也存在 $\xi\in[a,b]$, 使得 $\int_a^b f(x)g(x)\mathrm dx=f(b)\int_\xi^b g(x)\mathrm dx$? 定积分第二中值定理将回答这个问题, 该定理在后面章节中将多次用到.
% 
% \begin{theorem}[{定积分第二中值定理}]\label{theorem:ma-7-5-3}
% 设函数 $g(x)\in R[a,b]$.
% 
% (1) 若函数 $f(x)$ 在区间 $[a,b]$ 上单调上升, 且对于 $\forall x\in[a,b]$, 有 $f(x)\geq0$, 则存在 $\xi_1\in[a,b]$, 使得
% \[
%   \int_a^b f(x)g(x)\mathrm dx=f(b)\int_{\xi_1}^b g(x)\mathrm dx;
% \]
% 
% (2) 若函数 $f(x)$ 在区间 $[a,b]$ 上单调下降, 且对于 $\forall x\in[a,b]$, 有 $f(x)\geq0$, 则存在 $\xi_2\in[a,b]$, 使得
% \[
%   \int_a^b f(x)g(x)\mathrm dx=f(a)\int_a^{\xi_2} g(x)\mathrm dx;
% \]
% 
% (3) 若函数 $f(x)$ 在区间 $[a,b]$ 上单调, 则存在 $\xi\in[a,b]$, 使得
% \[
%   \int_a^b f(x)g(x)\mathrm dx
%   =f(a)\int_a^\xi g(x)\mathrm dx+f(b)\int_\xi^b g(x)\mathrm dx.
% \]
% \end{theorem}
% 
% 由于定理~\ref{theorem:ma-7-5-3} 的条件比较弱, 从而它的证明也比较复杂。为了看出定理~\ref{theorem:ma-7-5-3} 的证明思路, 我们的方法是先尝试证明该定理在一些较强条件下的相同的结果; 然后对其证明过程加以考察, 分析出证明的关键所在, 从而找到在较弱条件下的证明思路。以 (1) 为例, 若在原来假定下, 再假设 $g(x)\in C[a,b],f(x)\in C^1[a,b]$, 则有 $f'(x)\geq0,x\in[a,b]$, 且 $f(a)\geq0$。在这种情况下, 我们可以用下面简单方法来证明 (1).
% 
% 令 $G(x)=\int_x^b g(t)\mathrm dt$, 则 $G(x)$ 在 $[a,b]$ 上连续可导。设 $m,M$ 分别为 $G(x)$ 在 $[a,b]$ 上的最小值和最大值。由 $G'(x)=-g(x),G(b)=0$, 得
% \[
% \begin{aligned}
%   \int_a^b f(x)g(x)\mathrm dx
%   &=-\int_a^b f(x)\mathrm dG(x)\\
%   &=-\left.f(x)G(x)\right|_a^b+\int_a^b G(x)f'(x)\mathrm dx\\
%   &=f(a)G(a)+\int_a^b G(x)f'(x)\mathrm dx.
% \end{aligned}
% \]
% 再由
% \[
%   m[f(b)-f(a)]\leq\int_a^b G(x)f'(x)\mathrm dx\leq M[f(b)-f(a)],
% \]
% 我们推出
% \[
%   f(a)G(a)+\int_a^b G(x)f'(x)\mathrm dx
%   \leq f(a)M+M[f(b)-f(a)]=Mf(b)
% \]
% 和
% \[
%   f(a)G(a)+\int_a^b G(x)f'(x)\mathrm dx
%   \geq f(a)m+m[f(b)-f(a)]=mf(b).
% \]
% 因此有
% \[
%   mf(b)\leq\int_a^b f(x)g(x)\mathrm dx\leq Mf(b).
% \]
% 若 $f(b)=0$, 则 $f(x)\equiv0$, 结论显然成立。现设 $f(b)>0$, 则有
% \[
%   m\leq\frac{\int_a^b f(x)g(x)\mathrm dx}{f(b)}\leq M.
% \]
% 由连续函数的介值定理, 存在 $\xi_1\in[a,b]$, 使得
% \[
%   G(\xi_1)=\int_{\xi_1}^b g(x)\mathrm dx
%   =\frac{1}{f(b)}\int_a^b f(x)g(x)\mathrm dx.
% \]
% 因此 (1) 成立.
% 
% 从上述讨论可以看出, 若要证明定理~\ref{theorem:ma-7-5-3} 中的 (1), 我们仍然可以设 $G(x)=\int_x^b g(t)\mathrm dt$。由 $g(x)$ 的可积性知 $G(x)$ 在区间 $[a,b]$ 上连续。如果在这种情形下还是能证明不等式
% \[
%   mf(b)\leq\int_a^b f(x)g(x)\mathrm dx\leq Mf(b)
% \]
% 成立, 其中 $m,M$ 分别为 $G(x)$ 在 $[a,b]$ 上的最小值和最大值, 我们就可以证明 (1) 成立。但现在我们的假定只是 $f(x)$ 与 $g(x)$ 在 $[a,b]$ 上可积, 因此我们只能从定积分的定义出发来证明该结论。在定理~\ref{theorem:ma-7-5-3} 的证明过程中, 还需要用到下面的阿贝尔变换。为此, 我们将其以引理的形式给出.
% 
% \begin{lemma}[{阿贝尔变换}]\label{lemma:ma-7-5-4}
% 设 $\alpha_1,\alpha_2,\cdots,\alpha_n$ 和 $\beta_1,\beta_2,\cdots,\beta_n$ 为两组数, 记 $B_k=\sum_{i=1}^k\beta_i(k=1,2,\cdots,n)$, 则有
% \[
%   \sum_{i=1}^n\alpha_i\beta_i
%   =\sum_{i=1}^{n-1}(\alpha_i-\alpha_{i+1})B_i+\alpha_nB_n.
% \]
% \end{lemma}
% 
% 
% \begin{proof}
% 由 $\beta_i=B_i-B_{i-1}(i=2,3,\cdots,n)$, 得
% \[
% \begin{aligned}
%   \sum_{i=1}^n\alpha_i\beta_i
%   &=\sum_{i=2}^n\alpha_i(B_i-B_{i-1})+\alpha_1B_1\\
%   &=B_1(\alpha_1-\alpha_2)+B_2(\alpha_2-\alpha_3)+\cdots\\
%   &\quad+B_{n-1}(\alpha_{n-1}-\alpha_n)+B_n\alpha_n.
% \end{aligned}
% \qedhere
% \]
% \end{proof}
% 
% 
% \begin{proof}[{定理~\ref{theorem:ma-7-5-2} 的证明}]
% 我们只证 (2) 和 (3), (1) 的证明留给读者。为了证明 (2), 我们令 $h(x)=\int_a^x g(t)\mathrm dt$。由变限定积分的性质知 $h(x)\in C[a,b]$。记
% \[
%   m=\min_{x\in[a,b]}\{h(x)\},\qquad
%   M=\max_{x\in[a,b]}\{h(x)\}.
% \]
% 由于函数 $g(x)\in R[a,b]$, 存在 $0<M_1<+\infty$, 使得对于 $\forall x\in[a,b]$, 有 $|g(x)|\leq M_1$.
% 
% 设 $\Delta:a=x_0<x_1<\cdots<x_n=b$ 为 $[a,b]$ 的任一个分割, 则有
% \[
%   \int_a^b f(x)g(x)\mathrm dx
%   =\lim_{\lambda(\Delta)\to0}\sum_{i=1}^n\int_{x_{i-1}}^{x_i} f(x)g(x)\mathrm dx.
% \]
% 由于
% \[
%   \left|\sum_{i=1}^n\int_{x_{i-1}}^{x_i}[f(x)-f(x_{i-1})]g(x)\mathrm dx\right|
%   \leq M_1\sum_{i=1}^n\omega_i\Delta x_i\to0\quad(\lambda(\Delta)\to0),
% \]
% 其中 $\omega_i$ 为 $f(x)$ 在 $[x_{i-1},x_i](i=1,2,\cdots,n)$ 上的振幅, 我们有
% \[
%   \int_a^b f(x)g(x)\mathrm dx
%   =\lim_{\lambda(\Delta)\to0}\sum_{i=1}^n f(x_{i-1})\int_{x_{i-1}}^{x_i}g(x)\mathrm dx.
% \]
% 因此我们只要对
% \[
%   \sum_{i=1}^n f(x_{i-1})\int_{x_{i-1}}^{x_i}g(x)\mathrm dx
%   =\sum_{i=1}^n f(x_{i-1})[h(x_i)-h(x_{i-1})]
% \]
% 进行估计即可.
% 
% 利用阿贝尔变换 (引理~\ref{lemma:ma-7-5-4}), 我们有
% \[
% \sum_{i=1}^n f(x_{i-1})[h(x_i)-h(x_{i-1})]
% =\sum_{i=1}^{n-1}[f(x_{i-1})-f(x_i)]h(x_i)+f(x_{n-1})h(b).
% \]
% 由于 $f(x)$ 在 $[a,b]$ 上是单调下降的, 我们有 $f(x_{i-1})-f(x_i)\geq0(i=1,2,\cdots,n-1)$, 再注意到 $f(b)\geq0$ 和 $m\leq h(x)\leq M$, 可以推出
% \[
% \begin{aligned}
%   mf(a)&=\sum_{i=1}^{n-1}[f(x_{i-1})-f(x_i)]m+f(x_{n-1})m\\
%   &\leq\sum_{i=1}^{n-1}[f(x_{i-1})-f(x_i)]h(x_i)+f(x_{n-1})h(b)\\
%   &\leq\sum_{i=1}^{n-1}[f(x_{i-1})-f(x_i)]M+f(x_{n-1})M\\
%   &=Mf(a).
% \end{aligned}
% \]
% 由此得到
% \[
%   mf(a)\leq\int_a^b f(x)g(x)\mathrm dx
%   =\lim_{\lambda(\Delta)\to0}\sum_{i=1}^n f(x_{i-1})[h(x_i)-h(x_{i-1})]\leq Mf(a).
% \]
% 若 $f(a)=0$, 则 $f(x)\equiv0$, 结论显然成立。若 $f(a)\neq0$, 则有
% \[
%   m\leq\frac{\int_a^b f(x)g(x)\mathrm dx}{f(a)}\leq M.
% \]
% 因此存在 $\xi_2\in[a,b]$, 使得
% \[
%   h(\xi_2)=\int_a^{\xi_2}g(x)\mathrm dx
%   =\frac{\int_a^b f(x)g(x)\mathrm dx}{f(a)}.
% \]
% (2)
% \end{proof}
% 
% 
% 现在证 (3)。不妨设 $f(x)$ 单调上升, 则 $F(x)=f(x)-f(a)$ 满足 (1)。由 (1) 知, 存在 $\xi\in[a,b]$, 使得
% \[
%   \int_a^b F(x)g(x)\mathrm dx=F(b)\int_\xi^b g(x)\mathrm dx=[f(b)-f(a)]\int_\xi^b g(x)\mathrm dx.
% \]
% 因此
% \[
% \begin{aligned}
%   \int_a^b f(x)g(x)\mathrm dx
%   &=f(a)\int_a^b g(x)\mathrm dx+\int_a^b F(x)g(x)\mathrm dx\\
%   &=f(a)\int_a^b g(x)\mathrm dx+f(b)\int_\xi^b g(x)\mathrm dx
%   -f(a)\int_\xi^b g(x)\mathrm dx\\
%   &=f(a)\int_a^\xi g(x)\mathrm dx+f(b)\int_\xi^b g(x)\mathrm dx.
% \end{aligned}
% \]
% 证毕.
% 
% \begin{example}\label{example:ma-7-5-2}
% 设函数 $f(x)$ 在区间 $[a,b]$ 上单调上升, 证明
% \[
%   \int_a^b xf(x)\mathrm dx\geq\frac{a+b}{2}\int_a^b f(x)\mathrm dx.
% \]
% \end{example}
% 
% 
% \begin{proof}
% 对函数 $f(x)$ 及 $g(x)=x-\dfrac{a+b}{2}$ 在区间 $[a,b]$ 上应用定积分第二中值定理知, 存在 $\xi\in[a,b]$, 使得下式成立:
% \[
% \begin{aligned}
%   \int_a^b f(x)\left(x-\frac{a+b}{2}\right)\mathrm dx
%   &=f(a)\int_a^\xi\left(x-\frac{a+b}{2}\right)\mathrm dx
%   +f(b)\int_\xi^b\left(x-\frac{a+b}{2}\right)\mathrm dx\\
%   &=[f(b)-f(a)]\frac{b-\xi}{2}(\xi-a)\geq0.
% \end{aligned}
% \]
% 整理即得所要结论.
% \end{proof}
% 
% 
% \begin{example}\label{example:ma-7-5-3}
% 设函数
% \[
%   f(x)=
%   \begin{cases}
%     \displaystyle\int_0^x \sin\frac{1}{t}\mathrm dt, & x\neq0,\\
%     0, & x=0,
%   \end{cases}
% \]
% 证明 $f(x)$ 在 $x=0$ 处可导, 并且 $f'(0)=0$.
% \end{example}
% 
% 
% \begin{proof}
% 对任给的 $x>0$, 由于 $\sin\dfrac{1}{t}$ 在区间 $[0,x]$ 上可积, 从而有
% \[
%   \int_0^x \sin\frac{1}{t}\mathrm dt=\lim_{\delta\to0+0}\int_\delta^x \sin\frac{1}{t}\mathrm dt.
% \]
% 对于 $\forall 0<\delta<x$, 在上面等式右端中的定积分作积分变换 $t=\dfrac{1}{u}$, 得
% \[
%   \int_\delta^x \sin\frac{1}{t}\mathrm dt
%   =\int_{\frac{1}{x}}^{\frac{1}{\delta}}\frac{\sin u}{u^2}\mathrm du.
% \]
% 对上面等式右端中的定积分应用定积分第二中值定理知, 存在 $\xi_\delta\in\left[\dfrac{1}{x},\dfrac{1}{\delta}\right]$, 使得
% \[
%   \int_{\frac{1}{x}}^{\frac{1}{\delta}}\frac{\sin u}{u^2}\mathrm du
%   =x^2\int_{\frac{1}{x}}^{\xi_\delta}\sin u\mathrm du.
% \]
% 由于 $\left|\int_{\frac{1}{x}}^{\xi_\delta}\sin u\mathrm du\right|\leq2$, 因此, 对任意的 $\delta>0$, 有
% \[
%   \left|\int_\delta^x \sin\frac{1}{t}\mathrm dt\right|\leq2x^2.
% \]
% 于是对于任意给定的 $x>0$, 令 $\delta\to0+0$, 得
% \[
%   \left|\int_0^x \sin\frac{1}{t}\mathrm dt\right|\leq2x^2.
% \]
% 注意到 $\int_0^x\sin\dfrac{1}{t}\mathrm dt$ 是偶函数, 从而上式对 $x<0$ 也成立。因此有
% \[
%   0\leq\lim_{x\to0}\left|\frac{\int_0^x\sin\frac{1}{t}\mathrm dt}{x}\right|
%   \leq\lim_{x\to0}\frac{2x^2}{|x|}=0.
% \]
% 从上式我们知, $f(x)$ 在 $x=0$ 处可导, 并且
% \[
%   f'(0)=\lim_{x\to0}\frac{\int_0^x\sin\frac{1}{t}\mathrm dt}{x}=0.
% \qedhere
% \]
% \end{proof}
% 
% 
% \section{定积分在几何学中的应用}
% \label{sec:ma-7-6}
% 
% \subsection{直角坐标系下平面图形的面积}
% \label{subsec:ma-7-6-1}
% 
% 根据定积分的几何意义可以直接求出一些简单平面图形的面积。设平面图形 $S$ 的边界是由直线 $x=a,x=b$ 及连续函数 $y=f(x),y=g(x)(x\in[a,b])$ 的图像所组成 (其中 $g(x)\leq f(x),x\in[a,b]$, 见图 7.6.1), 我们称这种形状的平面图形 $S$ 为 X 型区域。对于这种区域的面积 $A$, 由定积分的几何意义立即得到
% \[
%   A=\int_a^b [f(x)-g(x)]\mathrm dx.
% \]
% 
% 若平面图形 $S$ 的边界是由直线 $y=c,y=d$ 及连续函数 $x=\varphi(y),x=\psi(y),y\in[c,d]$ 的图像所组成 (其中 $\varphi(y)\leq\psi(y),y\in[c,d]$, 见图 7.6.2), 则我们称这种平面图形 $S$ 为 Y 型区域。对于这种区域的面积 $A$, 由定积分的几何意义, 我们有以下的计算公式
% \[
%   A=\int_c^d [\psi(y)-\varphi(y)]\mathrm dy.
% \]
% 
% 因此当一个平面图形 $S$ 可以分成有限个互不相交的 X 型或 Y 型区域时, 我们就可以求出 $S$ 的面积.
% 
% \begin{example}\label{example:ma-7-6-1}
% 求由直线 $y=2,y=x$ 及曲线 $xy=1$ 所围平面图形 $S$ 的面积 $A$.
% \end{example}
% 
% 
% \begin{solve}
% 从图 7.6.3 可以看出, $S$ 是 Y 型区域, 因此有
% \[
%   A=\int_1^2\left(y-\frac{1}{y}\right)\mathrm dy
%   =\left.\left(\frac{y^2}{2}-\ln y\right)\right|_1^2
%   =\frac{3}{2}-\ln2.
% \]
% 
% 我们也可通过添加辅助线 $x=1$ 将图形 $S$ 分割成两个 X 型区域来进行计算:
% \[
% \begin{aligned}
%   A&=\int_{\frac12}^1\left(2-\frac{1}{x}\right)\mathrm dx+\int_1^2(2-x)\mathrm dx\\
%   &=\left.(2x-\ln x)\right|_{\frac12}^1
%   +\left.\left(2x-\frac{x^2}{2}\right)\right|_1^2\\
%   &=\frac{3}{2}-\ln2.
% \end{aligned}
% \qedhere
% \]
% \end{solve}
% 
% 
% \begin{example}[{Young 不等式}]\label{example:ma-7-6-2}
% 设 $y=f(x)$ 是区间 $[0,+\infty)$ 上严格上升的连续函数, 且满足 $f(0)=0$, 证明: 对任何的 $a>0,b>0$, 有
% \begin{equation}
%   ab\leq\int_0^a f(x)\mathrm dx+\int_0^b f^{-1}(y)\mathrm dy,
%   \label{eq:ma-7-6-1}
% \end{equation}
% 并且上述不等式中等号成立的充分必要条件是 $b=f(a)$.
% \end{example}
% 
% 
% \begin{proof}
% 首先我们注意到, 若 $y=\varphi(x)$ 是区间 $[\alpha,\beta]$ 上严格上升的非负连续函数, 则由定积分的几何意义有 (见图 7.6.4)
% \begin{equation}
%   \beta\varphi(\beta)-\alpha\varphi(\alpha)
%   =\int_\alpha^\beta \varphi(x)\mathrm dx
%   +\int_{\varphi(\alpha)}^{\varphi(\beta)}\varphi^{-1}(y)\mathrm dy.
%   \label{eq:ma-7-6-2}
% \end{equation}
% 现在我们回到 $f(x)$。先设 $0<b\leq f(a)$, 由于 $f(x)$ 的严格上升性, 我们有以下不等式
% \[
%   [a-f^{-1}(b)]b\leq\int_{f^{-1}(b)}^a f(x)\mathrm dx.
% \]
% 整理得
% \begin{equation}
%   ab-\int_0^a f(x)\mathrm dx
%   \leq bf^{-1}(b)-\int_0^{f^{-1}(b)}f(x)\mathrm dx.
%   \label{eq:ma-7-6-3}
% \end{equation}
% 在式 \eqref{eq:ma-7-6-2} 中令 $\alpha=0,\beta=b,\varphi(x)=f^{-1}(x)$, 则有
% \[
%   bf^{-1}(b)-\int_0^{f^{-1}(b)}f(x)\mathrm dx
%   =\int_0^b f^{-1}(y)\mathrm dy,
% \]
% 将上述等式代入 \eqref{eq:ma-7-6-3} 即得所证.
% 
% 当 $0<f(a)\leq b$ 时, 我们可以从 $y=f^{-1}(x)$ 出发, 重复以上论证即可证明所要的结论.
% 
% 显然, 当 $b=f(a)$ 时, 上面的每一不等式均成立等式, 从而 \eqref{eq:ma-7-6-1} 中成立等式。反之, 若 \eqref{eq:ma-7-6-1} 式的等号成立, 则由 $f(x)$ 的严格上升性知必有 $b=f(a)$.
% \end{proof}
% 
% 
% \subsection{参数方程表示的曲线所围平面图形的面积}
% \label{subsec:ma-7-6-2}
% 
% 设平面图形 $S$ 的边界曲线 $\gamma$ 由参数方程
% \[
%   \begin{cases}
%     x=x(t),\\
%     y=y(t)
%   \end{cases}
%   \quad(\alpha\leq t\leq\beta)
% \]
% 给出, 其中 $x(\alpha)=x(\beta),y(\alpha)=y(\beta)$, 曲线 $\gamma$ 除了端点重合外再无自交点 (如图 7.6.5)。假定 $x'(t),y'(t)$ 在区间 $[\alpha,\beta]$ 上存在且连续。为了叙述简便起见, 我们称这种曲线为 $C^1$ 曲线。再假定 $t$ 从 $\alpha$ 连续变化到 $\beta$ 时, 如果一个人沿着 $\gamma$ 行走, 图形 $S$ 总在他的左边。这种规定给出了 $\gamma$ 的一个正定向。以后我们总假定所求面积的图形是由满足上述条件的曲线所围.
% 
% 现在假设曲线 $\gamma$:
% \[
%   \begin{cases}
%     x=x(t),\\
%     y=y(t)
%   \end{cases}
%   \quad(\alpha\leq t\leq\beta)
% \]
% 围成一个 X 型区域, 如图 7.6.5 所示, 我们希望利用直角坐标求面积的公式来导出参数方程求面积的公式。注意到当 $t$ 从 $\alpha$ 严格上升到 $t_1$ 时, $x$ 从 $x(\alpha)$ 严格下降到 $x(t_1)$, 因此 $x(t)$ 在区间 $[\alpha,t_1]$ 上存在反函数 $t=t^{-1}(x)$。将其代入 $y=y(t)$ 得 $y=y(t^{-1}(x))\triangleq\varphi(x)(x\in[x(t_1),x(\alpha)])$。同理, $x(t)$ 在区间 $[t_1,t_2]$ 上也存在反函数 $t^{-1}(x)$, 因此当 $t\in[t_1,t_2]$ 时, $y=y(t^{-1}(x))\triangleq\psi(x)(x\in[x(t_1),x(t_2)])$。于是由定积分的几何意义及定积分的变量替换公式, 该 X 型区域的面积 $A$ 由以下等式表示:
% \[
% \begin{aligned}
%   A&=\int_{x(t_1)}^{x(\alpha)}\varphi(x)\mathrm dx
%   -\int_{x(t_1)}^{x(t_2)}\psi(x)\mathrm dx\\
%   &=-\int_\alpha^{t_1}y(t)x'(t)\mathrm dt
%   -\int_{t_1}^{t_2}y(t)x'(t)\mathrm dt
%   -\int_{t_2}^{\beta}y(t)\cdot0\,\mathrm dt\\
%   &=-\int_\alpha^\beta y(t)x'(t)\mathrm dt.
% \end{aligned}
% \]
% 同理, 若上述图形是一个 Y 型区域。读者可以推出 $\gamma$ 所围区域的面积为
% \[
%   A=\int_\alpha^\beta x(t)y'(t)\mathrm dt.
% \]
% 如果 $\gamma$ 所围图形可以添加几条 $C^1$ 曲线后成为若干块 X 型或 Y 型区域 (见图 7.6.6), 则 $\gamma$ 所围区域的面积仍具有形式
% \begin{equation}
% \begin{aligned}
%   A&=\int_\alpha^\beta x(t)y'(t)\mathrm dt
%   =-\int_\alpha^\beta y(t)x'(t)\mathrm dt\\
%   &=\frac{1}{2}\int_\alpha^\beta [x(t)y'(t)-y(t)x'(t)]\mathrm dt.
% \end{aligned}
% \label{eq:ma-7-6-4}
% \end{equation}
% 这是因为 $\gamma$ 所围图形的面积是各小块图形的面积之和, 而添加的曲线必定是两块小图形的共同边界。沿着添加的曲线的积分必须要积两次, 但由于曲线相对它所围的图形的定向相反, 因此在这些添加曲线上的积分互相抵消.
% 
% 为了简便起见, 可以将 \eqref{eq:ma-7-6-4} 简记为
% \[
%   A=\int_\gamma x\mathrm dy=-\int_\gamma y\mathrm dx
%   =\frac{1}{2}\int_\gamma x\mathrm dy-y\mathrm dx.
% \]
% 若不指明 $\gamma$ 的定向, 则积分总是取 $\gamma$ 的正向.
% 
% \begin{example}\label{example:ma-7-6-3}
% 计算椭圆 $\dfrac{x^2}{a^2}+\dfrac{y^2}{b^2}=1$ 所围图形 $S$ 的面积 $A$(见图 7.6.7).
% \end{example}
% 
% 
% \begin{solve}[{解法 1}]
% 因为 $S$ 是一个 X 型区域, 其中上、下曲线分别为 $y=\pm\dfrac{b}{a}\sqrt{a^2-x^2}$, 因此
% \[
% \begin{aligned}
%   A&=2\int_{-a}^a\frac{b}{a}\sqrt{a^2-x^2}\mathrm dx\\
%   &=\frac{2b}{a}\left.\left(\frac{x}{2}\sqrt{a^2-x^2}
%   +\frac{a^2}{2}\arcsin\frac{x}{a}\right)\right|_{-a}^a
%   =\pi ab.
% \end{aligned}
% \qedhere
% \]
% \end{solve}
% 
% 
% \begin{solve}[{解法 2}]
% 将椭圆方程写成参数方程
% \[
%   \gamma:
%   \begin{cases}
%     x=a\cos t,\\
%     y=b\sin t,
%   \end{cases}
%   \quad t\in[0,2\pi],
% \]
% 则当 $t$ 从 $0$ 连续变为 $2\pi$ 时, 该曲线为正定向。因此
% \[
% \begin{aligned}
%   A&=\int_\gamma x\mathrm dy
%   =\int_0^{2\pi}a\cos t\cdot b\cos t\mathrm dt\\
%   &=ab\int_0^{2\pi}\frac{1+\cos2t}{2}\mathrm dt
%   =\pi ab.
% \end{aligned}
% \qedhere
% \]
% \end{solve}
% 
% 
% \subsection{微元法}
% \label{subsec:ma-7-6-3}
% 
% 在以下各个小节中我们将推导一些平面或空间几何体的面积公式或体积公式。在这里我们利用定积分的思想归纳出推导这些公式的一个普遍的方法, 即所谓的微元法.
% 
% 我们考查正连续函数 $y=f(x)(x\in[a,b])$ 的定积分 $S=\int_a^b f(x)\mathrm dx$。我们已经知道它的几何意义是由曲线 $y=f(x)(x\in[a,b])$, $x$ 轴及直线 $x=a,x=b$ 所围的面积.
% 
% 对区间 $[0,S]$ 作分割 $\Delta_S:0=s_0<s_1<\cdots<s_n=S$, 则有
% \begin{equation}
%   \sum_{i=1}^n(s_i-s_{i-1})=\sum_{i=1}^n\Delta s_i=S=\int_a^b f(x)\mathrm dx.
%   \label{eq:ma-7-6-5}
% \end{equation}
% 由于 $f(x)>0$, 因此对于 $[0,S]$ 的分割相应地有 $[a,b]$ 的分割 $\Delta:a=x_0<x_1<\cdots<x_n=b$, 使得 $\Delta s_i$ 恰好是曲线 $y=f(x)(x\in[x_{i-1},x_i])$ 与直线 $x=x_{i-1},x=x_i$ 及 $x$ 轴所围的面积。容易看出当 $\lambda_S=\max_{1\leq i\leq n}\{\Delta s_i\}\to0$ 时, 有 $\lambda(\Delta)=\max_{1\leq i\leq n}\{\Delta x_i\}\to0$。因此, $f(x)$ 的黎曼和
% \begin{equation}
%   \sum_{i=1}^n f(\xi_i)\Delta x_i\to\sum_{i=1}^n\Delta s_i=S=\int_a^b f(x)\mathrm dx.
%   \label{eq:ma-7-6-6}
% \end{equation}
% 区间 $[0,S]$ 的长度和将它分成一些小区间再加起来是一样的。现在我们换一种比较高的观点来看此问题。任取 $[s,s+\Delta s]\subset[0,S]$, 当 $s$ 是自变量时, $\Delta s=\mathrm ds$。因此将小区间的长度加起来, 也就是将 $\mathrm ds$ 在 $[0,S]$ 上积分起来, 即 $S=\int_0^S\mathrm ds$.
% 
% 另外, 我们从 \eqref{eq:ma-7-6-6} 可以知道, $[s,s+\Delta s]$ 对应有一个 $[x,x+\Delta x]$, 使得 $\Delta s=f(\xi)\Delta x$, 其中 $\xi\in[x,x+\Delta x]$。当 $\Delta s$ 很小时, $\Delta x$ 也很小。现在我们将 $x$ 看成自变量, 则有 $\Delta x=\mathrm dx$。若将 $\xi$ 取成 $x$, 则由于 $\Delta x\to0$ 时有 $\Delta s=f(x)\Delta x+o(\Delta x)$。由此我们有 $\mathrm ds=f(x)\mathrm dx$。因此得到
% \begin{equation}
%   S=\int_0^S\mathrm ds=\int_a^b f(x)\mathrm dx.
%   \label{eq:ma-7-6-7}
% \end{equation}
% 
% 将上面的过程总结一下。我们希望求 $S$, 即曲线 $y=f(x)(x\in[a,b])$ 与直线 $x=a,x=b$ 及 $y=0$ 所围的曲边梯形的面积。假定我们不知道求该面积的公式, 现在来找出该公式。为此我们任取一小曲边梯形: 由过 $x,x+\Delta x$ 的两条平行于 $y$ 轴的直线、直线 $y=0$ 与曲线 $y=f(x)$ 所构成的小曲边梯形 (如图 7.6.8)。同时得到 $[0,S]$ 上的一个小区间 $[s,s+\Delta s]$, 使得 $\Delta s$ 等于该小曲边梯形的面积, 即有
% \[
%   \Delta s=f(\xi)\Delta x\approx f(x)\Delta x=f(x)\mathrm dx,
% \]
% 其中 $\xi$ 是 $[x,x+\Delta x]$ 中的某一点。如果 $f(x)$ 连续, 则有 $\mathrm ds=f(x)\mathrm dx$。因此我们得到了 \eqref{eq:ma-7-6-7}.
% 
% 去掉上述例子的具体背景, 我们可以归纳出以下求某些量的一般方法。设所求的量是 $S$, 它与区间 $[a,b]$ 有关, 当区间 $[a,b]$ 给定后, $S$ 就是一个确定的量, 而且该量对该区间具有可加性, 即对于 $[a,b]$ 的任意分割 $\Delta:a=x_0<x_1<\cdots<x_n=b$, 若记 $[x_{i-1},x_i]$ 对应的部分量为 $\Delta s_i(i=1,2,\cdots,n)$, 则 $S=\sum_{i=1}^n\Delta s_i$。为求 $S$, 我们可以任取小区间 $[x,x+\mathrm dx]\subset[a,b]$, 若 $[x,x+\mathrm dx]$ 所对应的部分量 $\Delta s=f(\xi)\mathrm dx(\xi\in[x,x+\mathrm dx])$, 其中 $f(x)$ 是 $[a,b]$ 上的一个连续函数, 则我们即有 $\mathrm ds=f(x)\mathrm dx$, 并且 $S=\int_a^b f(x)\mathrm dx$。这个求 $S$ 的过程也称微元法, 其中 $\mathrm ds=f(x)\mathrm dx$ 称为量 $S$ 的微元。微元法在后续章节中也多次用到.
% 
% \subsection{极坐标方程表示的曲线所围平面图形的面积}
% \label{subsec:ma-7-6-4}
% 
% 在本小节中, 我们设曲线 $\gamma$ 由 $r=r(\theta)(\alpha\leq\theta\leq\beta)$ 所表示, 其中 $r(\theta)$ 是 $\theta$ 的连续函数。现在我们来求曲线 $\gamma$ 与射线 $\theta=\alpha,\theta=\beta$ 所围平面图形 (见图 7.6.9) 的面积 $A$.
% 
% 我们用微元法来推导其面积公式。任取 $[\theta,\theta+\mathrm d\theta]\subset[\alpha,\beta]$。因此由曲线 $r=r(\theta)$ 及向径 $\theta_1=\theta,\theta_2=\theta+\mathrm d\theta$ 所围的图形近似于一个三角形, 其高近似为 $r(\theta)$, 而底边长近似为 $r(\theta)\mathrm d\theta$, 因此其面积微元为
% \[
%   \mathrm dA=\frac{1}{2}r^2(\theta)\mathrm d\theta.
% \]
% 从 $\alpha$ 到 $\beta$ 将其面积微元积分得
% \[
%   A=\int_\alpha^\beta\frac{1}{2}r^2(\theta)\mathrm d\theta.
% \]
% 
% \begin{example}\label{example:ma-7-6-4}
% 求心脏线 $r=a(1+\cos\theta)(a>0,0\leq\theta\leq2\pi)$ 所围平面图形 $S$ 的面积 $A$.
% \end{example}
% 
% 
% \begin{solve}
% \[
%   A=\frac{1}{2}\int_0^{2\pi}a^2(1+\cos\theta)^2\mathrm d\theta=\frac{3}{2}\pi a^2.
% \qedhere
% \]
% \end{solve}
% 
% 
% \begin{example}\label{example:ma-7-6-5}
% 求由圆弧
% \[
%   \begin{cases}
%     x=\cos t,\\
%     y=\sin t
%   \end{cases}
%   \quad\left(0\leq t\leq\frac{3}{4}\pi\right)
% \]
% 及 $x$ 轴和直线 $x=-\dfrac{\sqrt2}{2}$ 所围平面图形的面积 (见图 7.6.10).
% \end{example}
% 
% 
% \begin{solve}[{解法 1}]
% 我们用参数方程所围平面图形的面积公式来求, 注意到在线段 $BC$ 上有 $\mathrm dx=0$, 而在线段 $CD$ 上有 $y=0$, 因此该图形的面积为
% \[
% \begin{aligned}
%   A&=-\int_0^{\frac34\pi}y\mathrm dx
%   =-\int_0^{\frac34\pi}\sin t(-\sin t)\mathrm dt\\
%   &=\int_0^{\frac34\pi}\frac{1-\cos2t}{2}\mathrm dt
%   =\left.\left(\frac{t}{2}-\frac{\sin2t}{4}\right)\right|_0^{\frac34\pi}\\
%   &=\frac{3}{8}\pi+\frac{1}{4}.
% \end{aligned}
% \qedhere
% \]
% \end{solve}
% 
% 
% \begin{solve}[{解法 2}]
% 现在我们用极坐标表示的曲线所围平面图形的面积公式。参数方程表示的圆弧方程为 $r=1\left(0\leq\theta\leq\dfrac{3}{4}\pi\right)$。所围图形面积 $A$ 应为扇形 $BOD$ 的面积加上三角形 $OBC$ 的面积:
% \[
%   A=\int_0^{\frac34\pi}\frac{1}{2}\cdot1^2\mathrm d\theta
%   +\frac{1}{2}\left(\frac{\sqrt2}{2}\right)^2
%   =\frac{3}{8}\pi+\frac{1}{4}.
% \qedhere
% \]
% \end{solve}
% 
% 
% \subsection{平行截面面积为已知的立体的体积}
% \label{subsec:ma-7-6-5}
% 
% 利用定积分, 我们还可以求出空间一些规则立体图形的体积。设空间某立体 $\Omega$ 介于平面 $x=a$ 和 $x=b$ 之间, 假定过 $x\in[a,b]$ 并且与 $x$ 轴垂直的平面截该立体 $\Omega$ 的截面的面积为已知函数 $A(x)$(见图 7.6.11), 进一步假定 $A(x)$ 是区间 $[a,b]$ 上的连续函数。现在我们来求 $\Omega$ 的体积 $V$.
% 
% 利用微元法, 任取 $[x,x+\mathrm dx]\subset[a,b]$, $\Omega$ 夹在过点 $x=x$ 和 $x=x+\mathrm dx$ 并且与 $x$ 轴垂直的两个平面之间的部分近似于一个柱体, 因此体积微元 $\mathrm dV=A(x)\mathrm dx$。将体积微元从 $a$ 到 $b$ 积分得
% \[
%   V=\int_a^b A(x)\mathrm dx.
% \]
% 特别地, 若该立体 $\Omega$ 的边界曲面是由平面 $x=a,x=b$ 以及 $Oxy$ 平面内一条曲线 $y=f(x)\geq0(x\in[a,b])$ 绕 $x$ 轴旋转而成, 则称这个立体为旋转体。这种旋转体 $\Omega$ 与过 $x\in[a,b]$ 并且与 $x$ 轴垂直的平面相截的面积为 $A(x)=\pi f^2(x)$。因此 $\Omega$ 的体积 $V$ 由下面公式给出:
% \[
%   V=\pi\int_a^b f^2(x)\mathrm dx.
% \]
% 
% \begin{example}\label{example:ma-7-6-6}
% 求椭球体 $\dfrac{x^2}{a^2}+\dfrac{y^2}{b^2}+\dfrac{z^2}{c^2}\leq1(a,b,c>0)$ 的体积 $V$.
% \end{example}
% 
% 
% \begin{solve}
% 对于每个 $x\in[-a,a]$, 过点 $x$ 且与 $x$ 轴垂直的平面截该椭球体的截面为椭圆
% \[
%   \frac{y^2}{b^2\left(1-\frac{x^2}{a^2}\right)}
%   +\frac{z^2}{c^2\left(1-\frac{x^2}{a^2}\right)}=1
% \]
% 所围的平面图形, 由例~\ref{example:ma-7-6-2} 知该图形的面积为 $\pi bc\left(1-\dfrac{x^2}{a^2}\right)$, 因此
% \[
%   V=\int_{-a}^a \pi bc\left(1-\frac{x^2}{a^2}\right)\mathrm dx
%   =\left.\pi bc\left(x-\frac{x^3}{3a^2}\right)\right|_{-a}^a
%   =\frac{4}{3}\pi abc.
% \qedhere
% \]
% \end{solve}
% 
% 
% \begin{example}\label{example:ma-7-6-7}
% 求星形线 $x^{\frac23}+y^{\frac23}=a^{\frac23}(a>0)$ 绕 $x$ 轴旋转一周所成的立体体积 $V$.
% \end{example}
% 
% 
% \begin{solve}
% 由旋转体的体积公式, 我们有
% \[
%   V=\pi\int_{-a}^a y^2\mathrm dx
%   =\pi\int_{-a}^a (a^{\frac23}-x^{\frac23})^3\mathrm dx
%   =\frac{32}{105}\pi a^3.
% \qedhere
% \]
% \end{solve}
% 
% 
% \subsection{曲线的弧长}
% \label{subsec:ma-7-6-6}
% 
% 设平面曲线 $\Gamma$ 由参数方程
% \[
%   \begin{cases}
%     x=x(t),\\
%     y=y(t)
%   \end{cases}
%   \quad(\alpha\leq t\leq\beta)
% \]
% 给出。若 $x(t)$ 与 $y(t)$ 是区间 $[\alpha,\beta]$ 上的连续函数, 则我们称 $\Gamma$ 是一条连续曲线 (简称曲线)。我们下面涉及的曲线, 除特别说明外, 均指连续曲线.
% 
% 下面我们来讨论曲线的弧长问题。为此我们将面临以下三个问题:
% 
% (1) 如何来定义曲线的弧长?
% 
% (2) 是不是曲线都有弧长?
% 
% (3) 对有弧长的曲线怎样来求出它的弧长?
% 
% 首先来讨论问题 (1)。设平面曲线 $\Gamma$ 的起点为 $A$, 终点为 $B$。由于我们目前能够精确求出线段长度, 因此希望用线段的长度去逼近曲线的长度。为此我们在 $\Gamma$ 上从 $A$ 到 $B$ 依次取 $n+1$ 个点 $A=M_0,M_1,M_2,\cdots,M_{n-1},M_n=B$。这 $n+1$ 个点将 $\Gamma$ 分成了 $n$ 段小弧段。记 $L_i(i=1,2,\cdots,n)$ 为以 $M_{i-1},M_i$ 两点为端点的线段的长度。如果分点无限增加, 且当 $\lambda=\max_{1\leq i\leq n}\{L_i\}\to0$ 时, $\sum_{i=1}^n L_i$ 有极限 $L$, 则称 $\Gamma$ 为可求长曲线, 并且定义 $\Gamma$ 的长度为 $L$.
% 
% 现在我们来讨论问题 (2): 是不是曲线都可以求长度呢? 为此我们考虑由下述参数方程表示的曲线 $\Gamma$: 当 $t\in(0,1]$ 时, $x=t,y=t\sin\dfrac{1}{t}$; 当 $t=0$ 时, $x=y=0$.
% 
% 我们取 $t'_n=\dfrac{1}{2n\pi+\frac12\pi}$, $t''_n=\dfrac{1}{2n\pi-\frac12\pi}(n=1,2,\cdots,N)$, 设 $\Gamma$ 上对应 $t'_n$ 的点为 $M'_n$, 而对应 $t''_n$ 的点为 $M''_n$, 则我们有
% \[
%   \sum_{n=1}^N L_n
%   >\sum_{n=1}^N\left|t'_n\sin\frac{1}{t'_n}-t''_n\sin\frac{1}{t''_n}\right|
%   =\sum_{n=1}^N(t'_n+t''_n)>\frac{1}{2\pi}\sum_{n=1}^N\frac{1}{n},
% \]
% 其中 $L_n(n=1,2,\cdots,N)$ 是连接 $M'_n$ 和 $M''_n$ 的线段的长度。因为当分点个数 $N\to\infty$ 时, 不等式右端将趋于 $+\infty$, 从而 $\Gamma$ 是不可求长的曲线.
% 
% 从上面的例子可知, 并非所有的曲线都是可求长的, 我们必须对曲线加上适当条件才能使得它是可求长的.
% 
% 下面我们将讨论问题 (3)。如果我们要对一般的可求长曲线来求出它的长度, 除了利用弧长的定义外, 似乎很难找到其他的求弧长公式。为了利用微积分的工具来讨论曲线的求长问题, 我们必须对可求长曲线加上适当的光滑条件, 因此在下面主要讨论光滑曲线的求长问题。设 $\Gamma$:
% \[
%   \begin{cases}
%     x=x(t),\\
%     y=y(t)
%   \end{cases}
%   \quad(\alpha\leq t\leq\beta)
% \]
% 是一条曲线, 如果它满足 $x'(t),y'(t)$ 在区间 $[\alpha,\beta]$ 上存在且连续, 并且 $[x'(t)]^2+[y'(t)]^2\neq0$, 则称 $\Gamma$ 是一条光滑曲线.
% 
% 我们下面还是用微元法来推导这类曲线的弧长公式。任取 $[t,t+\mathrm dt]\subset[\alpha,\beta]$, 得到曲线上两个点 $M_t(x(t),y(t)),M_{t+\mathrm dt}(x(t+\mathrm dt),y(t+\mathrm dt))$, 则弦长 $\overline{M_tM_{t+\mathrm dt}}$ 有以下计算公式:
% \[
% \begin{aligned}
%   \overline{M_tM_{t+\mathrm dt}}
%   &=\sqrt{[x(t+\mathrm dt)-x(t)]^2+[y(t+\mathrm dt)-y(t)]^2}\\
%   &=\sqrt{[x'(\xi_1)]^2+[y'(\xi_2)]^2}\mathrm dt\\
%   &\approx\sqrt{[x'(t)]^2+[y'(t)]^2}\mathrm dt,
% \end{aligned}
% \]
% 其中 $\xi_1,\xi_2$ 在 $t$ 与 $t+\mathrm dt$ 之间。因此, 当 $\mathrm dt$ 很小时, 两点之间的弧长 $\Delta L$ 可以由弦长近似表示:
% \[
%   \Delta L\approx\sqrt{[x'(t)]^2+[y'(t)]^2}\mathrm dt.
% \]
% 这样我们得到弧长微元如下:
% \[
%   \mathrm dL=\sqrt{[x'(t)]^2+[y'(t)]^2}\mathrm dt.
% \]
% 因此 $\Gamma$ 的弧长公式为
% \[
%   L=\int_\alpha^\beta\sqrt{[x'(t)]^2+[y'(t)]^2}\mathrm dt.
% \]
% 特别地, 当曲线 $\Gamma$ 为函数 $y=f(x)(x\in[a,b])$ 的图像时, 若 $f(x)$ 在 $[a,b]$ 上具有连续导数, 则有
% \[
%   L=\int_a^b\sqrt{1+[f'(x)]^2}\mathrm dx.
% \]
% 如果曲线 $\Gamma$ 由极坐标 $r=r(\theta)(\alpha\leq\theta\leq\beta)$ 方程给出, 其中 $r'(\theta)$ 在 $\alpha\leq\theta\leq\beta$ 上存在并且连续, 则我们可将其写成参数式
% \[
%   \begin{cases}
%     x=r(\theta)\cos\theta,\\
%     y=r(\theta)\sin\theta
%   \end{cases}
%   \quad(\alpha\leq\theta\leq\beta),
% \]
% 从而弧长微元为
% \[
% \begin{aligned}
%   \mathrm dL
%   &=\sqrt{[r'(\theta)\cos\theta-r(\theta)\sin\theta]^2
%   +[r'(\theta)\sin\theta+r(\theta)\cos\theta]^2}\mathrm d\theta\\
%   &=\sqrt{[r(\theta)]^2+[r'(\theta)]^2}\mathrm d\theta.
% \end{aligned}
% \]
% 因此 $\Gamma$ 的长度 $L$ 有下面的计算公式:
% \[
%   L=\int_\alpha^\beta\sqrt{[r(\theta)]^2+[r'(\theta)]^2}\mathrm d\theta.
% \]
% 如果 $\Gamma$ 是空间曲线, $\Gamma$ 由参数方程
% \[
%   \begin{cases}
%     x=x(t),\\
%     y=y(t),\\
%     z=z(t)
%   \end{cases}
%   \quad(\alpha\leq t\leq\beta)
% \]
% 给出, 并且 $x'(t),y'(t),z'(t)$ 在区间 $[\alpha,\beta]$ 上连续且不同时为零, 则 $\Gamma$ 的弧长 $L$ 有下面的计算公式:
% \[
%   L=\int_\alpha^\beta\sqrt{[x'(t)]^2+[y'(t)]^2+[z'(t)]^2}\mathrm dt.
% \]
% 
% \begin{example}\label{example:ma-7-6-8}
% 求曲线 $\Gamma:y=x^2(x\in[0,1])$ 的弧长 $L$.
% \end{example}
% 
% 
% \begin{solve}
% \[
% \begin{aligned}
%   L&=\int_0^1\sqrt{1+[(x^2)']^2}\mathrm dx
%   =\int_0^1\sqrt{1+4x^2}\mathrm dx\\
%   &=\left.\left[x\sqrt{x^2+\frac{1}{4}}
%   +\frac{1}{4}\ln\left(x+\sqrt{x^2+\frac{1}{4}}\right)\right]\right|_0^1\\
%   &=\frac{\sqrt5}{2}+\frac{1}{4}\ln\left(1+\frac{\sqrt5}{2}\right)-\frac{1}{4}\ln\frac{1}{2}.
% \end{aligned}
% \qedhere
% \]
% \end{solve}
% 
% 
% \begin{example}\label{example:ma-7-6-9}
% 求双纽线 $r^2=2a^2\cos2\theta(a>0)$ 从 $\theta=0$ 到 $\theta=\dfrac{\pi}{6}$ 之间的弧的长度 $L$.
% \end{example}
% 
% 
% \begin{solve}
% 对 $r^2=2a^2\cos2\theta$ 两边关于 $\theta$ 求导得 $2rr'=-4a^2\sin2\theta$, 因此
% \[
%   r'(\theta)=\frac{-2a^2\sin2\theta}{r},
% \]
% 且
% \[
% \begin{aligned}
%   \sqrt{r^2(\theta)+[r'(\theta)]^2}
%   &=\sqrt{r^2+\frac{(-2a^2\sin2\theta)^2}{r^2}}
%   =\sqrt{\frac{r^4+4a^4\sin^22\theta}{r^2}}\\
%   &=\frac{\sqrt2a}{\sqrt{\cos2\theta}}
%   =\frac{\sqrt2a}{\sqrt{1-2\sin^2\theta}}.
% \end{aligned}
% \]
% 最后得到
% \[
%   L=\int_0^{\frac{\pi}{6}}\frac{\sqrt2a}{\sqrt{1-2\sin^2\theta}}\mathrm d\theta.
% \]
% 此积分称为椭圆积分, 由于被积函数没有初等原函数, 因此需要用其他的方法算出其值.
% \end{solve}
% 
% 
% \subsection{旋转体的侧面积}
% \label{subsec:ma-7-6-7}
% 
% 下面我们来讨论旋转体的侧面积问题。设 $\Gamma$ 为一光滑曲线, 其参数方程为
% \[
%   \begin{cases}
%     x=x(t),\\
%     y=y(t)
%   \end{cases}
%   \quad(\alpha\leq t\leq\beta).
% \]
% 再设当 $t\in[\alpha,\beta]$ 时, $y=y(t)\geq0$, 则曲线 $\Gamma$ 绕 $x$ 轴旋转一周后与 $x=x(\alpha)$ 和 $x=x(\beta)$ 两个平面围成一个旋转体 $\Omega$。现在我们用微元法来推导 $\Omega$ 的侧面积 $S$ 的计算公式.
% 
% 任取 $[t,t+\mathrm dt]\subset[\alpha,\beta]$, 曲线 $\Gamma$ 在 $[t,t+\mathrm dt]$ 的部分绕 $x$ 轴旋转所形成的立体可近似看做一个圆台 (图 7.6.12), 其侧面积为
% \[
%   \Delta S\approx\pi[y(t)+y(t+\mathrm dt)]\overline{M(t)M(t+\mathrm dt)},
% \]
% 其中 $M(t)=(x(t),y(t)),M(t+\mathrm dt)=(x(t+\mathrm dt),y(t+\mathrm dt))$。当 $\mathrm dt$ 很小时, $\overline{M(t)M(t+\mathrm dt)}\approx\mathrm dL$, 其中 $\mathrm dL=\sqrt{[x'(t)]^2+[y'(t)]^2}\mathrm dt$ 是 $\Gamma$ 的弧长微元, 而 $y(t+\mathrm dt)\approx y(t)$, 因此我们得到面积微元为
% \[
%   \mathrm dS=2\pi y(t)\sqrt{[x'(t)]^2+[y'(t)]^2}\mathrm dt.
% \]
% 对 $t$ 从 $\alpha$ 到 $\beta$ 积分即得所求侧面积:
% \[
%   S=2\pi\int_\alpha^\beta y(t)\sqrt{[x'(t)]^2+[y'(t)]^2}\mathrm dt.
% \]
% 类似地, 我们可以给出曲线绕 $y$ 轴旋转一周所得的旋转体的侧面积公式.
% 
% 对于在其他坐标下的曲线绕坐标轴旋转一周所得立体的侧面积公式, 请读者自己给出.
% 
% \begin{example}\label{example:ma-7-6-10}
% 求圆周 $x^2+y^2=R^2$ 上 $x$ 介于 $[a,a+h]$(其中 $-R\leq a<a+h\leq R$) 的部分绕 $x$ 轴旋转一周所形成的球台的体积及侧面积.
% \end{example}
% 
% 
% \begin{solve}
% 该球台可由曲线
% \[
%   y=f(x)=\sqrt{R^2-x^2},\quad x\in[a,a+h]
% \]
% 绕 $x$ 轴旋转一周而成, 因此其体积
% \[
% \begin{aligned}
%   V&=\pi\int_a^{a+h}(R^2-x^2)\mathrm dx
%   =\left.\pi\left(R^2x-\frac{x^3}{3}\right)\right|_a^{a+h}\\
%   &=\pi\left[R^2h-\left(a^2h+ah^2+\frac{h^3}{3}\right)\right].
% \end{aligned}
% \]
% 而该球台的侧面积为
% \[
% \begin{aligned}
%   S&=2\pi\int_a^{a+h}\sqrt{R^2-x^2}\sqrt{1+[f'(x)]^2}\mathrm dx\\
%   &=2\pi R\int_a^{a+h}\frac{\sqrt{R^2-x^2}}{\sqrt{R^2-x^2}}\mathrm dx
%   =2\pi hR.
% \end{aligned}
% \qedhere
% \]
% \end{solve}
% 
% 
% \begin{remark}\label{remark:ma-7-source-112}
% 在上述结果中, 令 $a=-R,h=2R$, 则我们得到半径为 $R$ 的球的体积为 $\dfrac{4\pi R^3}{3}$, 球面面积为 $4\pi R^2$。另外, 我们还可看出球台的体积不仅依赖于球台的高, 而且还依赖于 $a$ 的选取。但它的侧面积则为球面上大圆周长与球台的高的乘积, 而与 $a$ 的选取无关.
% \end{remark}
% 
% 
% \section{定积分在物理学中的应用}
% \label{sec:ma-7-7}
% 
% 定积分在物理学中有着广泛的应用。在本节中, 我们主要讨论一些规则物体的质心、转动惯量的计算问题.
% 
% 关于这些物理量的计算, 读者首先应该知道在有限个质点的情形下相应物理量的计算公式, 然后再根据定积分的定义来推导出其他形状物体的相应物理量的计算公式.
% 
% 先来考虑质心坐标问题。首先来复习一下物理学的有关概念。设平面上有 $n$ 个质点, 它们在平面内的位置为 $A_1(x_1,y_1),A_2(x_2,y_2),\cdots,A_n(x_n,y_n)$, 在点 $A_i$ 处的质点的质量为 $m_i(i=1,2,\cdots,n)$。由物理学中的定义, 这 $n$ 个质点关于 $x$ 轴和 $y$ 轴的静力矩分别为
% \[
%   M_x=\sum_{i=1}^n m_i y_i\quad\text{和}\quad M_y=\sum_{i=1}^n m_i x_i.
% \]
% 记 $M=\sum_{i=1}^n m_i$ 为这 $n$ 个质点的总质量, 则我们有以下的静力矩定律:
% 
% 在上述条件下, 存在点 $A(\overline x,\overline y)$, 满足
% \[
%   M_x=\sum_{i=1}^n m_i y_i=M\overline y,\qquad
%   M_y=\sum_{i=1}^n m_i x_i=M\overline x.
% \]
% 因此我们可得到以下结论: 这 $n$ 个质点关于 $x$ 轴 ($y$ 轴) 的静力矩之和相当于一个位于在点 $(\overline x,\overline y)$ 处质量为 $M$ 的质点关于 $x$ 轴 ($y$ 轴) 的静力矩。因此 $A(\overline x,\overline y)$ 可以认为是这 $n$ 个质点的质心, 其坐标由以下公式给出:
% \[
%   \overline x=\frac{M_y}{M}
%   =\frac{\sum_{i=1}^n m_i x_i}{\sum_{i=1}^n m_i},\qquad
%   \overline y=\frac{M_x}{M}
%   =\frac{\sum_{i=1}^n m_i y_i}{\sum_{i=1}^n m_i}.
% \]
% 在物理学中, 我们还分别称
% \[
%   I_x=\sum_{i=1}^n m_i y_i^2\quad\text{与}\quad I_y=\sum_{i=1}^n m_i x_i^2
% \]
% 为这 $n$ 个质点关于 $x$ 轴与 $y$ 轴的转动惯量.
% 
% 现在我们来求平面内一物质曲线的质心和关于坐标轴的转动惯量。所谓物质曲线是一个理想的模型, 即平面 (或空间) 有一条光滑曲线 $\Gamma$, 它由参数方程 $x=x(t),y=y(t)(\alpha\leq t\leq\beta)$ 给出, 同时在 $\Gamma$ 上有一个质量分布函数, 它由线密度 $\rho(t)(\alpha\leq t\leq\beta)$ 给出, 其中 $\rho(t)\in C[\alpha,\beta]$.
% 
% 由微元法, 任取 $[t,t+\mathrm dt]\subset[a,b]$, 则曲线弧长微元
% \[
%   \mathrm dL=\sqrt{[x'(t)]^2+[y'(t)]^2}\mathrm dt.
% \]
% 因此得到质量微元
% \[
%   \mathrm dM=\rho(t)\mathrm dL=\rho(t)\sqrt{[x'(t)]^2+[y'(t)]^2}\mathrm dt,
% \]
% 从而整个曲线的质量为
% \[
%   M=\int_\alpha^\beta \rho(t)\sqrt{[x'(t)]^2+[y'(t)]^2}\mathrm dt.
% \]
% 另外, 由质点的静力矩的定义易知, 该物质曲线关于 $x$ 轴及 $y$ 轴的静力矩微元
% \[
%   \mathrm dM_x=\mathrm dM\cdot y(t),\qquad
%   \mathrm dM_y=\mathrm dM\cdot x(t),
% \]
% 从而 $\Gamma$ 关于 $x$ 轴和 $y$ 轴的静力矩分别为
% \[
%   M_x=\int_\alpha^\beta \rho(t)y(t)\sqrt{[x'(t)]^2+[y'(t)]^2}\mathrm dt,
% \]
% \[
%   M_y=\int_\alpha^\beta \rho(t)x(t)\sqrt{[x'(t)]^2+[y'(t)]^2}\mathrm dt.
% \]
% 由此我们可以求出 $\Gamma$ 的质心坐标 $(\overline x,\overline y)$:
% \[
%   \overline x=\frac{M_y}{M}
%   =\frac{\int_\alpha^\beta \rho(t)x(t)\sqrt{[x'(t)]^2+[y'(t)]^2}\mathrm dt}
%   {\int_\alpha^\beta \rho(t)\sqrt{[x'(t)]^2+[y'(t)]^2}\mathrm dt},
% \]
% \[
%   \overline y=\frac{M_x}{M}
%   =\frac{\int_\alpha^\beta \rho(t)y(t)\sqrt{[x'(t)]^2+[y'(t)]^2}\mathrm dt}
%   {\int_\alpha^\beta \rho(t)\sqrt{[x'(t)]^2+[y'(t)]^2}\mathrm dt}.
% \]
% 若在上面公式中, 特别地令 $\rho(t)\equiv1$, 即曲线 $\Gamma$ 具有均匀的质量分布, 则得到质心坐标
% \[
%   \overline x=\frac{\int_\alpha^\beta x(t)\sqrt{[x'(t)]^2+[y'(t)]^2}\mathrm dt}{L},
% \]
% \[
%   \overline y=\frac{\int_\alpha^\beta y(t)\sqrt{[x'(t)]^2+[y'(t)]^2}\mathrm dt}{L},
% \]
% 其中 $L$ 是曲线 $\Gamma$ 的长度.
% 
% 特别地, 我们有
% \[
%   2\pi\overline yL
%   =2\pi\int_\alpha^\beta y(t)\sqrt{[x'(t)]^2+[y'(t)]^2}\mathrm dt=S,
% \]
% 其中 $S$ 是曲线 $\Gamma$ 绕 $x$ 轴旋转一周得到的旋转体的侧面积。因此上面的等式告诉我们, 一条曲线绕 $x$ 轴旋转一周所成旋转体的侧面积等于该曲线长度与质心绕 $x$ 轴旋转一周的周长的乘积。这个结果也称为是古鲁金 (Guldin) 第一定理.
% 
% 对于物质曲线 $\Gamma$ 关于 $x$ 轴和 $y$ 轴的转动惯量, 我们用同样的方法可得到如下公式:
% \[
%   I_x=\int_\alpha^\beta y^2(t)\rho(t)\sqrt{[x'(t)]^2+[y'(t)]^2}\mathrm dt,
% \]
% \[
%   I_y=\int_\alpha^\beta x^2(t)\rho(t)\sqrt{[x'(t)]^2+[y'(t)]^2}\mathrm dt.
% \]
% 
% \begin{example}\label{example:ma-7-7-1}
% 求圆周
% \[
%   x^2+(y-a)^2=R^2\quad(a>R>0)
% \]
% 绕 $x$ 轴旋转所成的轮胎体的表面积 (见图 7.7.1).
% \end{example}
% 
% 
% \begin{solve}
% 将该圆周看成线密度函数 $\rho\equiv1$ 的质量曲线。显然, 它的质心坐标为 $(0,a)$。注意到圆的周长为 $L=2\pi R$, 因此由古鲁金第一定理得其表面积为
% \[
%   S=L\cdot2\pi\overline y=2\pi R\cdot2\pi a=4\pi^2Ra.
% \]
% 
% 现在我们来考虑平面图形的质心问题。设平面图形 $S$ 由曲线 $y=y_1(x),y=y_2(x)(y_1(x)\leq y_2(x),a\leq x\leq b)$ 及直线 $x=a,x=b$ 所围成, 并设该平面图形 $S$ 的质量为均匀分布, 且假定其面密度函数为 $1$.
% 
% 下面我们仍由微元法来求其质心坐标。任取 $[x,x+\mathrm dx]\subset[a,b]$, 则该图形 $S$ 落在过点 $x$ 和 $x+\mathrm dx$ 并且与 $x$ 轴垂直的直线之间的部分可近似看成一个小矩形, 从而质量微元为 $\mathrm dM=[y_2(x)-y_1(x)]\mathrm dx$, 质心坐标近似为
% \[
%   \left(x,\frac{1}{2}[y_1(x)+y_2(x)]\right)
% \]
% (见图 7.7.2)。因此得到 $S$ 关于 $x$ 轴和 $y$ 轴的静力矩微元分别为
% \[
% \begin{aligned}
%   \mathrm dM_x
%   &=\frac{1}{2}[y_2(x)+y_1(x)][y_2(x)-y_1(x)]\mathrm dx\\
%   &=\frac{1}{2}[y_2^2(x)-y_1^2(x)]\mathrm dx,
% \end{aligned}
% \]
% \[
%   \mathrm dM_y=x[y_2(x)-y_1(x)]\mathrm dx.
% \]
% 对 $x$ 从 $a$ 到 $b$ 积分即得到整个图形 $S$ 关于 $x$ 轴和 $y$ 轴的静力矩分别为
% \[
%   M_x=\frac{1}{2}\int_a^b [y_2^2(x)-y_1^2(x)]\mathrm dx,
% \]
% \[
%   M_y=\int_a^b x[y_2(x)-y_1(x)]\mathrm dx.
% \]
% 注意到 $S$ 的总质量为
% \[
%   M=\int_a^b [y_2(x)-y_1(x)]\mathrm dx,
% \]
% 因此该图形 $S$ 的质心坐标为
% \[
%   \overline x=\frac{M_y}{M}
%   =\frac{\int_a^b x[y_2(x)-y_1(x)]\mathrm dx}
%   {\int_a^b [y_2(x)-y_1(x)]\mathrm dx},
% \]
% \[
%   \overline y=\frac{M_x}{M}
%   =\frac{\dfrac{1}{2}\int_a^b [y_2^2(x)-y_1^2(x)]\mathrm dx}
%   {\int_a^b [y_2(x)-y_1(x)]\mathrm dx}.
% \]
% 
% 由上面关于 $\overline y$ 的计算公式可得
% \[
% \begin{aligned}
%   2\pi\overline yS
%   &=2\pi\overline y\int_a^b [y_2(x)-y_1(x)]\mathrm dx\\
%   &=\pi\int_a^b [y_2^2(x)-y_1^2(x)]\mathrm dx=V.
% \end{aligned}
% \]
% 其中 $S$ 是该图形的面积, $V$ 是该图形绕 $x$ 轴旋转一周所得旋转体的体积。上述公式表明一个平面图形绕 $x$ 轴旋转一周得到的旋转体的体积等于该平面图形的面积与以质心到 $x$ 轴距离为半径的圆周长的乘积。这个结果也称为古鲁金第二定理.
% \end{solve}
% 
% 
% \begin{example}\label{example:ma-7-7-2}
% 求半径为 $R>0$ 的半圆盘的质心.
% \end{example}
% 
% 
% \begin{solve}
% 如图 7.7.3 建立坐标系。设半圆盘的质心坐标为 $(\overline x,\overline y)$, 由对称性显然有 $\overline x=0$。下面求 $\overline y$。注意到该图形绕 $x$ 轴旋转一周得到一个球体, 利用古鲁金第二定理有
% \[
%   2\pi\overline y\cdot\frac{\pi}{2}R^2=\frac{4}{3}\pi R^3.
% \]
% 解上述方程得 $\overline y=\dfrac{4}{3\pi}R$, 因此圆盘质心为 $\left(0,\dfrac{4}{3\pi}R\right)$.
% 
% 下面我们举一个例子来讨论极坐标方程表示的物质曲线转动惯量的计算.
% \end{solve}
% 
% 
% \begin{example}\label{example:ma-7-7-3}
% 设平面质量曲线为 $r=r(\theta)(\alpha\leq\theta\leq\beta)$, 线密度为 $\rho\equiv1$, 求该曲线关于极轴的转动惯量.
% \end{example}
% 
% 
% \begin{solve}
% 任取 $[\theta,\theta+\mathrm d\theta]\subset(\alpha,\beta)$, 则由转动惯量的定义, 曲线上位于 $[\theta,\theta+\mathrm d\theta]$ 上的弧段质量微元为 $\mathrm dM=\sqrt{r^2(\theta)+r'^2(\theta)}\mathrm d\theta$。这小段质量弧到极轴的距离近似为 $r(\theta)\sin\theta$ (见图 7.7.4), 它关于极轴的转动惯量, 即转动惯量微元为
% \[
%   \mathrm dI=r^2(\theta)\sin^2\theta\,\mathrm dM,
% \]
% 因此所求曲线关于极轴的转动惯量为
% \[
%   I=\int_\alpha^\beta r^2(\theta)\sin^2\theta
%   \sqrt{r^2(\theta)+r'^2(\theta)}\mathrm d\theta.
% \qedhere
% \]
% \end{solve}
% 
% 
% \section*{习题七}
% \phantomsection
% \addcontentsline{toc}{section}{习题七}
% \label{exercises:ma-7}
% 
% 1。从定义出发计算定积分 $\int_0^T(\alpha+\beta t)\mathrm dt$, 这里 $\alpha,\beta$ 是常数.
% 
% % 整合来源：math-8.tex
% \chapter{广义积分}
% \label{chap:ma-8}
% 
% 在定积分中讨论的函数必须是定义在有限区间上的有界函数, 在本章中我们将定积分作两个方面的推广: 一是将有限区间推广到无穷区间, 二是将有界函数推广到无界函数。对于前者的积分, 我们称之为无穷积分, 而将后者的积分称为瑕积分。通常将无穷积分和瑕积分统称为广义积分.
% 
% \section{无穷积分的基本概念与性质}
% \label{sec:ma-8-1}
% 
% 在本节中我们研究无穷积分。首先我们来考查一个火箭升空的例子。假设人们在地球上某个地方向天空发射一个质量为 $m$ 的火箭, 其目的是探索太阳系中某个遥远的星球。设地球半径为 $R$, 质量为 $M$, 则由万有引力定律知, 火箭从地球表面发射升至高度 $X$, 它克服地球引力所做的功为
% \[
%   W(X)=\int_R^{R+X}G\frac{Mm}{r^2}\mathrm dr,
% \]
% 其中 $G=\dfrac{R^2g}{M}$, 而 $g$ 为大家熟悉的重力加速度, 即 $g\approx9.8\mathrm{m}/\mathrm{s}^2$。对任意固定的 $X$, 我们有
% \[
%   W(X)=mgR^2\left(\frac1R-\frac1{R+X}\right).
% \]
% 火箭逐步摆脱地球引力的过程可以看成是 $X\to+\infty$ 的变化过程。因此, 从理论上来看, 火箭要摆脱地球引力所做的功应该为
% \[
% \begin{aligned}
%   \lim_{X\to+\infty}W(X)
%   &=\lim_{X\to+\infty}\int_R^{R+X}G\frac{Mm}{r^2}\mathrm dr\\
%   &=\lim_{X\to+\infty}mgR^2\left(\frac1R-\frac1{R+X}\right)=mgR.
% \end{aligned}
% \]
% 利用上述分析, 我们可以求出所谓的第二宇宙速度, 即火箭具有的最小初速度, 使得它能够飞离地球的引力范围。事实上, 设火箭向上发射时的初速度为 $v_0$, 它得到的动能为 $\dfrac12mv_0^2$。当它超过 $\lim\limits_{h\to+\infty}W(h)=mgR$ 时, 则火箭可以脱离地球的引力。因此由能量守恒定律, 有
% \[
%   \frac12mv_0^2=mgR,
% \]
% 取 $g=9.81\mathrm{m}/\mathrm{s}^2,R=6371\ \mathrm{km}$, 代入得
% \[
%   v_0\approx11.2\ \mathrm{km}/\mathrm{s}.
% \]
% 
% 现在我们只关注火箭在升空的过程中克服地球引力做功所涉及的积分问题。在火箭达到的每一个高度 $X$, 我们都面临一个有限区间的积分 $\int_R^X G\dfrac{Mm}{r^2}\mathrm dr$, 该积分对固定的 $X$ 是一个定积分, 但由于 $X$ 不断变大, 实际我们将面临的是考查当 $X\to+\infty$ 时, 积分 $\int_R^X G\dfrac{Mm}{r^2}\mathrm dr$ 的变化趋势。我们把这个考查过程记为 $\int_R^{+\infty}G\dfrac{Mm}{r^2}\mathrm dr$, 并称之为无穷积分。为此给出下面的定义.
% 
% \begin{definition}\label{definition:ma-8-1-1}
% 设函数 $f(x)$ 在 $[a,+\infty)$ 上有定义, 并且对于 $\forall X\in(a,+\infty)$, 在 $[a,X]$ 上可积。如果极限
% \[
%   \lim_{X\to+\infty}\int_a^X f(x)\mathrm dx
% \]
% 存在, 则称无穷积分 $\int_a^{+\infty}f(x)\mathrm dx$ 收敛, 或称函数 $f(x)$ 在 $[a,+\infty)$ 上可积, 并记
% \[
%   \int_a^{+\infty}f(x)\mathrm dx=\lim_{X\to+\infty}\int_a^X f(x)\mathrm dx;
% \]
% 如果极限 $\lim\limits_{X\to+\infty}\int_a^X f(x)\mathrm dx$ 不存在, 则称无穷积分 $\int_a^{+\infty}f(x)\mathrm dx$ 发散.
% \end{definition}
% 
% 从上述定义可以看出: 记号 $\int_a^{+\infty}f(x)\mathrm dx$ 仅代表了我们考虑 $f(x)$ 在 $[a,+\infty)$ 上的积分这件事情, 只有当无穷积分 $\int_a^{+\infty}f(x)\mathrm dx$ 收敛时, 它才是一个有限实数.
% 
% 类似地我们也可以讨论一个函数在 $(-\infty,b]$ 上的积分问题.
% 
% \begin{definition}\label{definition:ma-8-1-2}
% 设函数 $f(x)$ 在 $(-\infty,b]$ 上有定义, 并且对于 $\forall X\in(-\infty,b)$, 在区间 $[X,b]$ 上可积。如果极限
% \[
%   \lim_{X\to-\infty}\int_X^b f(x)\mathrm dx
% \]
% 存在, 则称无穷积分 $\int_{-\infty}^b f(x)\mathrm dx$ 收敛, 或称函数 $f(x)$ 在 $(-\infty,b]$ 上可积, 并记
% \[
%   \int_{-\infty}^b f(x)\mathrm dx=\lim_{X\to-\infty}\int_X^b f(x)\mathrm dx;
% \]
% 如果极限 $\lim\limits_{X\to-\infty}\int_X^b f(x)\mathrm dx$ 不存在, 则称无穷积分 $\int_{-\infty}^b f(x)\mathrm dx$ 发散.
% \end{definition}
% 
% 
% \begin{definition}\label{definition:ma-8-1-3}
% 设函数 $f(x)$ 在 $(-\infty,+\infty)$ 上有定义, 且在任何的闭区间 $[a,b]$ 上可积。任取 $c\in\mathbb R$, 若无穷积分 $\int_{-\infty}^c f(x)\mathrm dx$ 与 $\int_c^{+\infty}f(x)\mathrm dx$ 都收敛, 则称无穷积分 $\int_{-\infty}^{+\infty}f(x)\mathrm dx$ 收敛, 并记
% \[
%   \int_{-\infty}^{+\infty}f(x)\mathrm dx
%   =
%   \int_{-\infty}^c f(x)\mathrm dx
%   +
%   \int_c^{+\infty}f(x)\mathrm dx;
% \]
% 若无穷积分 $\int_{-\infty}^c f(x)\mathrm dx$ 和 $\int_c^{+\infty}f(x)\mathrm dx$ 中至少有一个发散, 则称无穷积分 $\int_{-\infty}^{+\infty}f(x)\mathrm dx$ 发散.
% \end{definition}
% 
% 由于我们在定义~\ref{definition:ma-8-1-3} 中假定了 $f(x)$ 在任何有限闭区间均可积, 读者易于验证, $\int_{-\infty}^{+\infty}f(x)\mathrm dx$ 收敛与否与 $c$ 的选取无关; 当 $\int_{-\infty}^{+\infty}f(x)\mathrm dx$ 收敛时, 其值也与 $c$ 的选取无关.
% 
% 由无穷积分敛散性的定义我们可以看出: 若函数 $f(x)$ 在 $[a,+\infty)$ 上有原函数 $F(x)$, 并形式地记 $F(+\infty)=\lim\limits_{x\to+\infty}F(x)$, 则有
% \[
%   \int_a^{+\infty}f(x)\mathrm dx=F(x)\big|_a^{+\infty}=F(+\infty)-F(a).
% \]
% 同样, 若 $f(x)$ 在 $(-\infty,b]$ 上有原函数 $G(x)$, 记 $G(-\infty)=\lim\limits_{x\to-\infty}G(x)$, 则
% \[
%   \int_{-\infty}^b f(x)\mathrm dx=G(x)\big|_{-\infty}^b=G(b)-G(-\infty).
% \]
% 类似地, 若 $f(x)$ 在 $(-\infty,+\infty)$ 上有原函数 $H(x)$, 则
% \[
%   \int_{-\infty}^{+\infty}f(x)\mathrm dx=H(x)\big|_{-\infty}^{+\infty}=H(+\infty)-H(-\infty).
% \]
% 
% 上述的公式给出了判定无穷积分是否收敛以及计算收敛无穷积分的值的方法.
% 
% 由定积分几何意义, 不难理解无穷积分 $\int_a^{+\infty}f(x)\mathrm dx$, $\int_{-\infty}^b f(x)\mathrm dx$ 和 $\int_{-\infty}^{+\infty}f(x)\mathrm dx$ 的几何意义。例如, 当 $f(x)$ 是 $[a,+\infty)$ 上的连续函数时, $\int_a^{+\infty}f(x)\mathrm dx$ 的几何意义是图 8.1.1 中阴影部分面积的代数和.
% 
% \begin{example}\label{example:ma-8-1-1}
% 证明无穷积分 $\int_{-\infty}^{+\infty}\dfrac{\mathrm dx}{1+x^2}$ 收敛并求其值.
% \end{example}
% 
% 
% \begin{solve}
% 由于 $\arctan x$ 是 $\dfrac1{1+x^2}$ 在 $(-\infty,+\infty)$ 上的一个原函数, 因此我们有
% \[
%   \lim_{X\to-\infty}\int_X^0\frac{\mathrm dx}{1+x^2}
%   =
%   0-\lim_{X\to-\infty}\arctan X
%   =
%   \frac{\pi}{2},
% \]
% \[
%   \lim_{X\to+\infty}\int_0^X\frac{\mathrm dx}{1+x^2}
%   =
%   \lim_{X\to+\infty}\arctan X
%   =
%   \frac{\pi}{2}.
% \]
% 由无穷积分的定义知 $\int_{-\infty}^{+\infty}\dfrac{\mathrm dx}{1+x^2}$ 收敛, 并且有
% \[
%   \int_{-\infty}^{+\infty}\frac{\mathrm dx}{1+x^2}
%   =
%   \arctan x\big|_{-\infty}^{+\infty}
%   =
%   \frac{\pi}{2}-\left(-\frac{\pi}{2}\right)=\pi.
% \qedhere
% \]
% \end{solve}
% 
% 
% \begin{example}\label{example:ma-8-1-2}
% 讨论无穷积分 $\int_1^{+\infty}\dfrac{\mathrm dx}{x^p}$ 的敛散性, 其中 $p\in\mathbb R$.
% \end{example}
% 
% 
% \begin{solve}
% 当 $p=1$ 时,
% \[
% \begin{aligned}
%   \int_1^{+\infty}\frac{\mathrm dx}{x}
%   &=
%   \lim_{X\to+\infty}\int_1^X\frac{\mathrm dx}{x}
%   =
%   \lim_{X\to+\infty}\ln x\big|_1^X\\
%   &=
%   \lim_{X\to+\infty}\ln X=+\infty;
% \end{aligned}
% \]
% 当 $p\ne1$ 时,
% \[
% \begin{aligned}
%   \int_1^{+\infty}\frac{\mathrm dx}{x^p}
%   &=
%   \lim_{X\to+\infty}\int_1^X\frac{\mathrm dx}{x^p}
%   =
%   \lim_{X\to+\infty}\frac{x^{-p+1}}{1-p}\bigg|_1^X\\
%   &=
%   \lim_{X\to+\infty}\frac{X^{-p+1}-1}{1-p}
%   =
%   \begin{cases}
%     \dfrac1{p-1}, & p>1,\\
%     +\infty, & p<1.
%   \end{cases}
% \end{aligned}
% \]
% 因此无穷积分 $\int_1^{+\infty}\dfrac{\mathrm dx}{x^p}$ 当 $p>1$ 时收敛, 当 $p\leqslant1$ 时发散.
% \end{solve}
% 
% 
% \begin{example}\label{example:ma-8-1-3}
% 讨论无穷积分 $\int_2^{+\infty}\dfrac{\mathrm dx}{x\ln^p x}$ 的敛散性, 其中 $p\in\mathbb R$.
% \end{example}
% 
% 
% \begin{solve}
% 当 $p=1$ 时,
% \[
% \begin{aligned}
%   \int_2^{+\infty}\frac{\mathrm dx}{x\ln x}
%   &=
%   \lim_{X\to+\infty}\int_2^X\frac{\mathrm dx}{x\ln x}\\
%   &=
%   \lim_{X\to+\infty}(\ln\ln X-\ln\ln2)=+\infty.
% \end{aligned}
% \]
% 当 $p\ne1$ 时,
% \[
% \begin{aligned}
%   \int_2^{+\infty}\frac{\mathrm dx}{x\ln^p x}
%   &=
%   \lim_{X\to+\infty}\int_2^X\frac{\mathrm dx}{x\ln^p x}
%   =
%   \lim_{X\to+\infty}\frac{\ln^{-p+1}x}{1-p}\bigg|_2^X\\
%   &=
%   \lim_{X\to+\infty}\frac{\ln^{1-p}X-\ln^{1-p}2}{1-p}
%   =
%   \begin{cases}
%     \dfrac1{(p-1)\ln^{p-1}2}, & p>1,\\
%     +\infty, & p<1.
%   \end{cases}
% \end{aligned}
% \]
% 因此无穷积分 $\int_2^{+\infty}\dfrac{\mathrm dx}{x\ln^p x}$ 当 $p>1$ 时收敛, 当 $p\leqslant1$ 时发散.
% 
% 类似地, 我们可以证明: 无穷积分
% \[
%   \int_{e^2}^{+\infty}
%   \frac{\mathrm dx}{x(\ln x)(\ln\ln x)^p}
% \]
% 当 $p>1$ 时收敛, 当 $p\leqslant1$ 时发散.
% 
% 从无穷积分的定义出发, 我们可以容易地证明下面的关于无穷积分的换元公式与分部积分公式.
% 
% 设函数 $f(x)$ 在 $[a,+\infty)$ 上有定义, 且对于 $\forall X>a$, 在区间 $[a,X]$ 上可积, 再设函数 $\varphi(t)$ 在区间 $[\alpha,\beta)(\beta$ 可以是 $+\infty)$ 上连续可微, 严格单调上升, 并且满足
% \[
%   a=\varphi(\alpha)\leq \varphi(t)\leq \lim_{t\to\beta-0}\varphi(t)=+\infty,
% \]
% 则我们有以下的换元公式:
% \[
%   \int_a^{+\infty}f(x)\mathrm dx
%   =
%   \int_\alpha^\beta f(\varphi(t))\varphi'(t)\mathrm dt.
% \qedhere
% \]
% \end{solve}
% 
% 
% \begin{remark}\label{remark:ma-8-source-10}
% 上述等式是在以下意义下成立的: 当等式两端的积分有一个收敛时, 则另一个积分也收敛并且上述等式必成立; 若上面的两个积分中有一个发散, 则另外一个积分也发散。以下若有两个无穷积分相等的等式, 则读者均可按上述意义来加以理解。另外, 对 $\varphi(t)$ 是单调下降的情形, 请读者给出相应的换元公式.
% \end{remark}
% 
% 现设函数 $u(x),v(x)$ 在 $[a,+\infty)$ 上连续可微, 且极限
% \[
%   u(+\infty)v(+\infty)\triangleq \lim_{x\to+\infty}u(x)v(x)
% \]
% 存在, 则有以下分部积分公式
% \[
%   \int_a^{+\infty}u(x)v'(x)\mathrm dx
%   =
%   u(x)v(x)\big|_a^{+\infty}
%   -
%   \int_a^{+\infty}u'(x)v(x)\mathrm dx.
% \]
% 
% 读者可类似地给出无穷积分 $\int_{-\infty}^b f(x)\mathrm dx$ 及 $\int_{-\infty}^{+\infty}f(x)\mathrm dx$ 的换元公式与分部积分公式.
% 
% \begin{example}\label{example:ma-8-1-4}
% 计算无穷积分 $\int_1^{+\infty}\dfrac{\mathrm dx}{x\sqrt{1+x^2}}$.
% \end{example}
% 
% 
% \begin{solve}
% 令 $x=\dfrac1t$, 则原积分可化为
% \[
%   \int_1^{+\infty}\frac{\mathrm dx}{x\sqrt{1+x^2}}
%   =
%   \int_0^1\frac{\mathrm dt}{\sqrt{1+t^2}}
%   =
%   \ln(t+\sqrt{1+t^2})\big|_0^1
%   =
%   \ln(1+\sqrt2).
% \]
% 从上例可知, 一个无穷积分经过换元后有时可变成一个定积分.
% \end{solve}
% 
% 
% \begin{example}\label{example:ma-8-1-5}
% 计算无穷积分 $\int_1^{+\infty}\dfrac{\arctan x}{x^3}\mathrm dx$.
% \end{example}
% 
% 
% \begin{solve}
% 取 $u(x)=\arctan x,v(x)=-\dfrac1{2x^2}$, 则利用分部积分有
% \[
% \begin{aligned}
%   \int_1^{+\infty}\frac{\arctan x}{x^3}\mathrm dx
%   &=
%   -\frac{\arctan x}{2x^2}\bigg|_1^{+\infty}
%   +
%   \int_1^{+\infty}\frac{\mathrm dx}{2x^2(1+x^2)}\\
%   &=
%   \frac{\pi}{8}
%   +
%   \frac12\int_1^{+\infty}\left(\frac1{x^2}-\frac1{1+x^2}\right)\mathrm dx\\
%   &=
%   \frac{\pi}{8}
%   +
%   \frac12\left(-\frac1x-\arctan x\right)\bigg|_1^{+\infty}
%   =
%   \frac12.
% \end{aligned}
% \qedhere
% \]
% \end{solve}
% 
% 
% \begin{example}\label{example:ma-8-1-6}
% 计算无穷积分
% \[
%   I_1=\int_0^{+\infty}e^{-ax}\sin bx\mathrm dx,\quad
%   I_2=\int_0^{+\infty}e^{-ax}\cos bx\mathrm dx,\quad
%   \text{其中 }a>0.
% \]
% \end{example}
% 
% 
% \begin{solve}
% 在下一节中, 我们将知道以上无穷积分是收敛的。现在我们来求其值:
% \[
% \begin{aligned}
%   I_1
%   &=
%   -\lim_{X\to+\infty}\frac1a\int_0^X\sin bx\,\mathrm d(e^{-ax})\\
%   &=
%   -\lim_{X\to+\infty}\frac1a\left(e^{-ax}\sin bx\big|_0^X
%   -
%   b\int_0^X e^{-ax}\cos bx\mathrm dx\right)\\
%   &=
%   \frac ba\int_0^{+\infty}e^{-ax}\cos bx\mathrm dx
%   =
%   \frac ba I_2,
% \end{aligned}
% \]
% \[
% \begin{aligned}
%   I_2
%   &=
%   -\lim_{X\to+\infty}\frac1a\int_0^X\cos bx\,\mathrm d(e^{-ax})\\
%   &=
%   -\lim_{X\to+\infty}\frac1a\left(e^{-ax}\cos bx\big|_0^X
%   +
%   b\int_0^X e^{-ax}\sin bx\mathrm dx\right)\\
%   &=
%   \frac1a-\frac ba I_1.
% \end{aligned}
% \]
% 由以上两式可得
% \[
%   I_1=\frac b{a^2+b^2},\quad I_2=\frac a{a^2+b^2}.
% \]
% 
% 回忆一下, 我们说无穷积分 $\int_{-\infty}^{+\infty}f(x)\mathrm dx$ 收敛是指关于 $X_1$ 和 $X_2$ 的极限 $\lim\limits_{\substack{X_1\to-\infty\\ X_2\to+\infty}}\int_{X_1}^{X_2}f(x)\mathrm dx$ 存在, 其中关于 $X_1$ 和 $X_2$ 的极限过程是两个独立的过程。如果仅考虑 $\lim\limits_{X\to+\infty}\int_{-X}^X f(x)\mathrm dx$, 则当它收敛时, 我们称 $\int_{-\infty}^{+\infty}f(x)\mathrm dx$ 在柯西主值意义下是收敛的, 此时我们用
% \[
%   \mathrm{V.P.}\int_{-\infty}^{+\infty}f(x)\mathrm dx
% \]
% 来表示极限 $\lim\limits_{X\to+\infty}\int_{-X}^X f(x)\mathrm dx$, 同时称该极限值为积分 $\int_{-\infty}^{+\infty}f(x)\mathrm dx$ 的柯西主值。显然, 对任何 $(-\infty,+\infty)$ 上可积的奇函数 $f(x)$, 有
% \[
%   \mathrm{V.P.}\int_{-\infty}^{+\infty}f(x)\mathrm dx=0.
% \]
% 
% 从某种意义上来说, 当讨论无穷积分 $\int_{-\infty}^{+\infty}f(x)\mathrm dx$ 的敛散性问题时, 我们必须讨论两个独立的极限过程。而讨论该无穷积分的柯西主值时, 我们只要讨论一个极限过程。无穷积分的柯西主值在一些数学分支中有应用.
% \end{solve}
% 
% 
% \section{无穷积分敛散性的判别法}
% \label{sec:ma-8-2}
% 
% 在对无穷积分 $\int_a^{+\infty}f(x)\mathrm dx$ 的讨论中, 我们假定了函数 $f(x)$ 在 $[a,+\infty)$ 的有限子区间上都是可积的, 而要研究的 $\int_a^{+\infty}f(x)\mathrm dx$ 敛散性问题, 实际上就是研究当 $X\to+\infty$ 时连续函数 $F(X)=\int_a^X f(x)\mathrm dx$ 的极限是否存在的问题。因此从理论上讲, 无穷积分的敛散性可以认为是已经解决的问题。但是人们很快就发现, 对于许多的函数 $f(x)$, 相应的 $F(X)$ 并不能容易求出, 因此我们要寻找直接从被积函数的性质来判定广义积分是否收敛的方法.
% 
% 下面我们主要讨论 $\int_a^{+\infty}f(x)\mathrm dx$ 的敛散性问题, 关于 $\int_{-\infty}^b f(x)\mathrm dx$ 及 $\int_{-\infty}^{+\infty}f(x)\mathrm dx$ 的敛散性, 请读者自己给出相应的结论.
% 
% 首先由无穷积分 $\int_a^{+\infty}f(x)\mathrm dx$ 收敛的定义, 我们容易得到下面的柯西准则.
% 
% \begin{theorem}[{柯西准则}]\label{theorem:ma-8-2-1}
% 设函数 $f(x)$ 在 $[a,+\infty)$ 上有定义, 对于 $\forall X>a$, 在区间 $[a,X]$ 上可积, 则无穷积分 $\int_a^{+\infty}f(x)\mathrm dx$ 收敛的充分必要条件是: 对于 $\forall\varepsilon>0$, $\exists M>a$, 当 $X''>X'>M$ 时, 有
% \[
%   \left|\int_{X'}^{X''}f(x)\mathrm dx\right|<\varepsilon.
% \]
% \end{theorem}
% 
% 此定理的证明可以从无穷积分敛散性的定义直接推出, 故省略.
% 
% \begin{definition}\label{definition:ma-8-2-1}
% 设函数 $f(x)$ 在 $[a,+\infty)$ 上有定义, 对于 $\forall X>a$, 在区间 $[a,X]$ 上可积。若无穷积分 $\int_a^{+\infty}|f(x)|\mathrm dx$ 收敛, 则称无穷积分 $\int_a^{+\infty}f(x)\mathrm dx$ 绝对收敛; 若无穷积分 $\int_a^{+\infty}f(x)\mathrm dx$ 收敛, 但无穷积分 $\int_a^{+\infty}|f(x)|\mathrm dx$ 发散, 则称无穷积分 $\int_a^{+\infty}f(x)\mathrm dx$ 条件收敛.
% \end{definition}
% 
% 
% \begin{theorem}\label{theorem:ma-8-2-2}
% 设函数 $f(x)$ 在 $[a,+\infty)$ 上有定义, 对于 $\forall X>a$, 在区间 $[a,X]$ 上可积。若无穷积分 $\int_a^{+\infty}f(x)\mathrm dx$ 绝对收敛, 则无穷积分 $\int_a^{+\infty}f(x)\mathrm dx$ 必收敛.
% \end{theorem}
% 
% 
% \begin{proof}
% 假设 $\int_a^{+\infty}|f(x)|\mathrm dx$ 收敛, 由柯西准则, 对于 $\forall\varepsilon>0$, $\exists M>a$, 当 $X''>X'>M$ 时, 有
% \[
%   \left|\int_{X'}^{X''}|f(x)|\mathrm dx\right|
%   =
%   \int_{X'}^{X''}|f(x)|\mathrm dx<\varepsilon.
% \]
% 因此, 当 $X''>X'>M$ 时, 有
% \[
%   \left|\int_{X'}^{X''}f(x)\mathrm dx\right|
%   \leq
%   \int_{X'}^{X''}|f(x)|\mathrm dx<\varepsilon.
% \]
% 再由柯西准则知 $\int_a^{+\infty}f(x)\mathrm dx$ 收敛.
% \end{proof}
% 
% 
% 从上述定理知, 无穷积分是否绝对收敛的问题可以转化为非负函数的无穷积分的敛散性问题.
% 
% \begin{theorem}\label{theorem:ma-8-2-3}
% 设非负函数 $f(x)$ 在 $[a,+\infty)$ 上有定义, 对于 $\forall X>a$, 在 $[a,X]$ 上可积, 则无穷积分 $\int_a^{+\infty}f(x)\mathrm dx$ 收敛的充分必要条件是: 存在 $A>0$, 使得对一切 $X\geq a$, 有
% \[
%   \int_a^X f(x)\mathrm dx\leq A.
% \]
% \end{theorem}
% 
% 
% \begin{proof}
% 事实上, 由定积分关于积分区间的可加性及 $f(x)\geq0$ 知, $F(X)=\int_a^X f(x)\mathrm dx$ 在 $[a,+\infty)$ 上是 $X$ 的单调上升函数。因此由单调有界定理即可推出定理~\ref{theorem:ma-8-2-3}.
% \end{proof}
% 
% 
% 对于非负函数的无穷积分, 我们有下面的比较定理.
% 
% \begin{theorem}\label{theorem:ma-8-2-4}
% 设非负函数 $f(x),g(x)$ 在 $[a,+\infty)$ 上有定义, 且对于 $\forall X>a$, 在 $[a,X]$ 上可积。若存在常数 $c_1>0,c_2>0$ 及 $M_0\geq a$, 使得当 $x\geq M_0$ 时, 成立不等式
% \[
%   c_1f(x)\leq c_2g(x),
% \]
% 则我们有下述结论:
% 
% (1) 若 $\int_a^{+\infty}g(x)\mathrm dx$ 收敛, 则 $\int_a^{+\infty}f(x)\mathrm dx$ 也收敛;
% 
% (2) 若 $\int_a^{+\infty}f(x)\mathrm dx$ 发散, 则 $\int_a^{+\infty}g(x)\mathrm dx$ 也发散.
% \end{theorem}
% 
% 
% \begin{proof}
% 我们只证 (1), (2) 的证明留给读者.
% 
% 由假设 $\int_a^{+\infty}g(x)\mathrm dx$ 收敛, 根据柯西准则, 对于 $\forall\varepsilon>0$, $\exists M_1>0$, 当 $X''>X'>M_1$ 时, 有
% \[
%   \int_{X'}^{X''}g(x)\mathrm dx<\frac{c_1}{c_2}\varepsilon.
% \]
% 由于当 $x\geq M_0$ 时, 有
% \[
%   f(x)\leq \frac{c_2}{c_1}g(x),
% \]
% 因此我们取 $M=\max\{M_0,M_1\}$, 当 $X''>X'>M$ 时, 有
% \[
%   0\leq \int_{X'}^{X''}f(x)\mathrm dx
%   \leq
%   \frac{c_2}{c_1}\int_{X'}^{X''}g(x)\mathrm dx<\varepsilon.
% \]
% (1) 得证.
% \end{proof}
% 
% 
% \begin{corollary}\label{corollary:ma-8-source-25}
% 设非负函数 $f(x),g(x)(g(x)\ne0)$ 在 $[a,+\infty)$ 上有定义, 且对于 $\forall X>a$, 在区间 $[a,X]$ 上可积。若
% \[
%   \lim_{x\to+\infty}\frac{f(x)}{g(x)}=l,
% \]
% 则
% 
% (1) 当 $0<l<+\infty$ 时, $\int_a^{+\infty}g(x)\mathrm dx$ 与 $\int_a^{+\infty}f(x)\mathrm dx$ 同时收敛或同时发散;
% 
% (2) 当 $l=0$ 时, 若 $\int_a^{+\infty}g(x)\mathrm dx$ 收敛, 则 $\int_a^{+\infty}f(x)\mathrm dx$ 收敛;
% 
% (3) 当 $l=+\infty$ 时, 若 $\int_a^{+\infty}g(x)\mathrm dx$ 发散, 则 $\int_a^{+\infty}f(x)\mathrm dx$ 发散.
% \end{corollary}
% 
% 上述推论可由极限的保号性和定理~\ref{theorem:ma-8-2-4} 推出.
% 
% 显然, 对两个非负函数的无穷积分, 以上定理和推论告诉我们, 取值大的函数的无穷积分收敛, 则取值小的函数的无穷积分也收敛; 反之, 取值小的函数的无穷积分发散, 则取值大的函数的无穷积分也发散.
% 
% 在判断一些无穷积分的敛散性时, 我们有一些 ``标准函数'', 其他的一些函数可与它们作比较。这些 ``标准函数'' 有 $\dfrac1{x^p},\dfrac1{x^p\ln^q x},e^{-ax}$ 等。读者应该熟悉这些函数中参数取哪些值是使得无穷积分收敛的, 而哪些值是使得无穷积分发散的。例如, 在 $\int_2^{+\infty}\dfrac{\mathrm dx}{x^p\ln^q x}$ 中, 当 $p>1$ 时, 对任意的 $q\in\mathbb R$, 无穷积分收敛; 而当 $p=1$ 时, 对于 $q>1$ 收敛, 对于其他的情形该无穷积分均发散。再如, 在 $\int_0^{+\infty}P(x)e^{-ax}\mathrm dx$ 中(其中 $P(x)$ 是一个多项式), 当 $a>0$ 时, 无穷积分是收敛的.
% 
% 在上节例~\ref{example:ma-8-1-6} 中, 由于被积函数在 $[0,+\infty)$ 上满足
% \[|e^{-ax}\sin bx|\leq e^{-ax},\qquad |e^{-ax}\cos bx|\leq e^{-ax},\]
%  因此 $I_1=\int_0^{+\infty}e^{-ax}\sin bx\mathrm dx$ 与 $I_2=\int_0^{+\infty}e^{-ax}\cos bx\mathrm dx$ 均是绝对收敛的.
% 
% \begin{example}\label{example:ma-8-2-1}
% 讨论无穷积分 $\int_0^{+\infty}\dfrac{x^2}{x^n+x^2+1}\mathrm dx$ 的敛散性, 其中 $n$ 为一正整数.
% \end{example}
% 
% 
% \begin{solve}
% 由于
% \[
%   \lim_{x\to+\infty}
%   \frac{\dfrac{x^2}{x^n+x^2+1}}{\dfrac1{x^{n-2}}}
%   =
%   1,
% \]
% 而当 $n-2>1$, 即 $n>3$ 时, $\int_0^{+\infty}\dfrac1{x^{n-2}}\mathrm dx$ 收敛, 因此当 $n>3$ 时, 原无穷积分收敛; 而当 $n\leq3$ 时, 无穷积分发散.
% 
% 若一个有理函数 $R(x)=\dfrac{P(x)}{Q(x)}$ 在 $[a,+\infty)$ 上有定义, 其中 $P(x)$ 与 $Q(x)$ 为两个既约多项式, 它们的次数分别为 $m,n$, 则 $\int_a^{+\infty}R(x)\mathrm dx$ 收敛的充分必要条件是 $n-m\geq2$。对于有理函数的无穷积分, 虽然我们很容易判断其收敛性, 但计算其值并不容易。在复变函数课程中我们有系统的办法来求其值.
% \end{solve}
% 
% 
% \begin{example}\label{example:ma-8-2-2}
% 讨论无穷积分 $\int_1^{+\infty}\dfrac{\cos x\sin\dfrac1x}{x}\mathrm dx$ 的敛散性.
% \end{example}
% 
% 
% \begin{solve}
% 由于
% \[
%   \left|\frac{\cos x\sin\dfrac1x}{x}\right|
%   \leq
%   \frac1x\sin\frac1x,\quad x\geq1,
% \]
% 且
% \[
%   \lim_{x\to+\infty}
%   \frac{\dfrac1x\sin\dfrac1x}{\dfrac1{x^2}}=1,
% \]
% 而 $\int_1^{+\infty}\dfrac1{x^2}\mathrm dx$ 收敛, 所以 $\int_1^{+\infty}\dfrac{\cos x\sin\dfrac1x}{x}\mathrm dx$ 绝对收敛.
% \end{solve}
% 
% 
% \begin{example}\label{example:ma-8-2-3}
% 讨论无穷积分 $\int_1^{+\infty}\left[\dfrac1{x^p}-\ln\left(1+\dfrac1{x^p}\right)\right]\mathrm dx\ (p>0)$ 的敛散性.
% \end{example}
% 
% 
% \begin{solve}
% 在 $\ln(1+t)$ 的麦克劳林展开式中令 $t=\dfrac1{x^p}$, 我们有
% \[
%   \ln\left(1+\frac1{x^p}\right)
%   =
%   \frac1{x^p}-\frac1{2x^{2p}}+o\left(\frac1{x^{2p}}\right)
%   \quad(x\to+\infty),
% \]
% 从而有
% \[
%   0\leq
%   \frac1{x^p}-\ln\left(1+\frac1{x^p}\right)
%   =
%   \frac1{2x^{2p}}+o\left(\frac1{x^{2p}}\right)
%   \sim
%   \frac1{2x^{2p}}
%   \quad(x\to+\infty).
% \]
% 因此, $\int_1^{+\infty}\left[\dfrac1{x^p}-\ln\left(1+\dfrac1{x^p}\right)\right]\mathrm dx$ 当 $p>\dfrac12$ 时收敛, 当 $p\leq\dfrac12$ 时发散.
% 
% 以上我们讨论了非负函数无穷积分的敛散性问题。正如我们所指出的那样, 这些判别法则可以用来判定无穷积分是否绝对收敛。值得注意的是, 当 $f(x)\geq0$ 时, 由 $\int_a^{+\infty}f(x)\mathrm dx$ 的几何意义来看, 似乎当 $x$ 越大时, $f(x)$ 的值要很小时才能保证无穷积分 $\int_a^{+\infty}f(x)\mathrm dx$ 的收敛性。当 $f(x)$ 不连续时, 我们容易看出这个想法是不对的, 这可以从以下事实得到验证。取严格单调上升序列 $\{x_n\}\subset[a,+\infty)$, 使得 $\lim\limits_{n\to\infty}x_n=+\infty$, 由于改变 $f(x)$ 在 $\{x_n\}$ 上的取值并不改变 $f(x)$ 在任何区间的可积性, 因此我们可以使得 $f(x)$ 在该序列的值改变后的函数在这一序列上的值趋于 $\infty$。即使 $f(x)$ 是连续的, 我们也容易构造出一个在 $[a,+\infty)$ 上无界的函数 $f(x)$, 使得广义积分 $\int_a^{+\infty}f(x)\mathrm dx$ 收敛.
% \end{solve}
% 
% 
% \begin{example}\label{example:ma-8-2-4}
% 设函数 $f(x)$ 在 $[1,+\infty)$ 上的定义如下:
% \[
%   f(x)=
%   \begin{cases}
%     \dfrac1{x^2}, & x\in[1,x_2'],\\
%     \dfrac1{x^2}, & x\in[x_n',x_{n+1}],\\
%     f(x_{n+1})+3^{n+1}\bigl[2^{n+1}-f(x_{n+1})\bigr](x-x_{n+1}), & x\in(x_{n+1},n+1],\\
%     2^{n+1}-3^{n+1}\bigl[2^{n+1}-f(x_{n+1}')\bigr](x-n-1), & x\in(n+1,x_{n+1}'],
%   \end{cases}
% \]
% 其中 $x_n=n-\dfrac1{3^n},x_n'=n+\dfrac1{3^n},n=2,3,\cdots$ ($f(x)$ 的大致图像如图 8.2.1 所示), 证明无穷积分 $\int_1^{+\infty}f(x)\mathrm dx$ 收敛.
% \end{example}
% 
% 
% \begin{proof}
% 从 $f(x)$ 的构造可以看出 $f(x)>0(x\in[1,+\infty)), f(x)\in C[1,+\infty)$, 并且 $\lim\limits_{n\to\infty}f(n)=+\infty$。对于 $\forall X>1$, 取正整数 $N>X$, 则有
% \[
% \begin{aligned}
%   \int_1^X f(x)\mathrm dx
%   &\leq
%   \int_1^N f(x)\mathrm dx
%   \leq
%   \int_1^N\frac{\mathrm dx}{x^2}
%   +
%   \sum_{n=2}^N\int_{x_n}^{x_n'}f(x)\mathrm dx\\
%   &\leq
%   1+\sum_{n=2}^N\int_{x_n}^{x_n'}2^n\mathrm dx
%   <
%   1+2\sum_{n=2}^N\frac{2^n}{3^n}<4.
% \end{aligned}
% \]
% 因此无穷积分 $\int_1^{+\infty}f(x)\mathrm dx$ 收敛.
% 
% 作为练习, 读者可以证明: 当 $f(x)$ 在 $[a,+\infty)$ 上单调时, 若无穷积分 $\int_a^{+\infty}f(x)\mathrm dx$ 收敛, 则必有 $\lim\limits_{x\to+\infty}f(x)=0$.
% 
% 现在我们转向讨论条件收敛的无穷积分。对于这类无穷积分, 被积函数 $f(x)$ 的绝对值的无穷积分趋于 $+\infty$。若 $f(x)$ 在 $[a,+\infty)$ 上连续, 当 $x\to+\infty$ 时, $f(x)$ 取值必须在 $x$ 轴的两侧波动, 使得其图像与 $x$ 轴所围的图形在 $x$ 轴的上面与下面的面积不断抵消, 才能使得 $\int_a^{+\infty}f(x)\mathrm dx$ 收敛。在此方面我们有以下两个有用的判别法.
% \end{proof}
% 
% 
% \begin{theorem}[{狄利克雷判别法}]\label{theorem:ma-8-2-5}
% 设函数 $f(x),g(x)$ 在 $[a,+\infty)$ 上有定义, 且满足下面两个条件:
% 
% (1) 对于 $\forall X>a$, $f(x)$ 在区间 $[a,X]$ 上可积, 并且 $\exists M>0$, 使得对于 $\forall X>a$, 有
% \[
%   \left|\int_a^X f(x)\mathrm dx\right|\leq M;
% \]
% 
% (2) $g(x)$ 在 $[a,+\infty)$ 单调, 并且 $\lim\limits_{x\to+\infty}g(x)=0$,
% 
% 则无穷积分 $\int_a^{+\infty}f(x)g(x)\mathrm dx$ 收敛.
% \end{theorem}
% 
% 
% \begin{proof}
% 对于 $\forall\varepsilon>0$, 由 $\lim\limits_{x\to+\infty}g(x)=0$, $\exists X>0$, 当 $x>X$ 时, 有
% \[
%   |g(x)|<\frac{\varepsilon}{4M}.
% \]
% 对于任意的 $X''>X'>X$, 由条件 (1) 和 (2), 我们可以对定积分 $\int_{X'}^{X''}f(x)g(x)\mathrm dx$ 应用定积分第二中值定理, 从而在 $(X',X'')$ 内存在 $\xi$, 使得
% \[
% \begin{aligned}
%   \left|\int_{X'}^{X''}f(x)g(x)\mathrm dx\right|
%   &=
%   \left|g(X')\int_{X'}^\xi f(x)\mathrm dx
%   +
%   g(X'')\int_\xi^{X''}f(x)\mathrm dx\right|\\
%   &<
%   \frac{\varepsilon}{4M}
%   \left\{
%   \left|\int_a^\xi f(x)\mathrm dx-\int_a^{X'}f(x)\mathrm dx\right|
%   +
%   \left|\int_a^{X''}f(x)\mathrm dx-\int_a^\xi f(x)\mathrm dx\right|
%   \right\}\\
%   &\leq \frac{\varepsilon}{4M}\cdot4M=\varepsilon.
% \end{aligned}
% \]
% 由柯西准则我们知无穷积分 $\int_a^{+\infty}f(x)g(x)\mathrm dx$ 收敛.
% \end{proof}
% 
% 
% \begin{theorem}[{阿贝尔判别法}]\label{theorem:ma-8-2-6}
% 设函数 $f(x),g(x)$ 在 $[a,+\infty)$ 上有定义, 并且满足下面两个条件:
% 
% (1) 对于 $\forall X>a$, $f(x)$ 在 $[a,X]$ 上可积, 并且 $\int_a^{+\infty}f(x)\mathrm dx$ 收敛;
% 
% (2) $g(x)$ 在 $[a,+\infty)$ 单调有界,
% 
% 则无穷积分 $\int_a^{+\infty}f(x)g(x)\mathrm dx$ 收敛.
% \end{theorem}
% 
% 
% \begin{proof}
% 由条件 (2) 我们知 $\lim\limits_{x\to+\infty}g(x)$ 存在, 设该极限值为 $B$。容易验证 $f(x)$ 和 $G(x)=g(x)-B$ 满足狄利克雷判别法(定理~\ref{theorem:ma-8-2-5}) 的所有条件, 从而 $\int_a^{+\infty}f(x)G(x)\mathrm dx$ 收敛。由
% \[
%   \int_a^{+\infty}f(x)g(x)\mathrm dx
%   =
%   \int_a^{+\infty}f(x)G(x)\mathrm dx
%   +
%   B\int_a^{+\infty}f(x)\mathrm dx
% \]
% 即知定理结论成立.
% \end{proof}
% 
% 
% \begin{example}\label{example:ma-8-2-5}
% 证明无穷积分 $\int_{-\infty}^{+\infty}\dfrac{\sin x}{x}\mathrm dx$ 条件收敛.
% \end{example}
% 
% 
% \begin{solve}
% 注意到被积函数为偶函数以及
% \[
%   f(x)=
%   \begin{cases}
%     \dfrac{\sin x}{x}, & x\ne0,\\
%     1, & x=0
%   \end{cases}
% \]
% 在 $(-\infty,+\infty)$ 上连续, 因此我们只要证明无穷积分 $\int_{2\pi}^{+\infty}\dfrac{\sin x}{x}\mathrm dx$ 条件收敛即可.
% 
% 对于 $\forall X>2\pi$, 我们有
% \[
%   \left|\int_{2\pi}^X\sin x\mathrm dx\right|
%   =
%   |\cos2\pi-\cos X|<2,
% \]
% 又由于 $\dfrac1x$ 在 $[2\pi,+\infty)$ 上单调下降且 $\lim\limits_{x\to+\infty}\dfrac1x=0$, 由狄利克雷判别法知 $\int_{2\pi}^{+\infty}\dfrac{\sin x}{x}\mathrm dx$ 收敛.
% 
% 由于
% \[
%   \frac{|\sin x|}{x}\geq \frac{\sin^2x}{x}=\frac1{2x}-\frac{\cos x}{2x},
% \]
% 类似上面讨论, 我们可证 $\int_{2\pi}^{+\infty}\dfrac{\cos x}{2x}\mathrm dx$ 收敛。由于 $\int_{2\pi}^{+\infty}\dfrac{\mathrm dx}{2x}$ 发散, 因此 $\int_{2\pi}^{+\infty}\dfrac{\sin^2x}{x}\mathrm dx$ 发散。再由比较判别法知 $\int_{2\pi}^{+\infty}\dfrac{|\sin x|}{x}\mathrm dx$ 发散。因此无穷积分 $\int_{2\pi}^{+\infty}\dfrac{\sin x}{x}\mathrm dx$ 条件收敛.
% 
% 从此例子可以清楚地看出, 当 $x$ 充分大时, 由 $\sin x$ 的周期性, 对曲线 $y=\dfrac{\sin x}{x}$ 与 $x$ 轴所围的每一小块的面积 $A_i(i=1,2,\cdots)$(见图 8.2.2), 有
% \[
%   A_1>A_2>A_3>\cdots,
% \]
% 且 $\lim\limits_{n\to+\infty}A_n=0$。由于 $x$ 轴上面小块与下面小块的面积不断抵消, 所以 $\int_{2\pi}^{+\infty}\dfrac{\sin x}{x}\mathrm dx$ 收敛.
% 
% 另外, 在不定积分中我们知道 $\int\dfrac{\sin x}{x}\mathrm dx$ 不能用初等函数表示, 因此对 $-\infty<a<b<+\infty$, 定积分 $\int_a^b\dfrac{\sin x}{x}\mathrm dx$ 就很难计算出来。但以后我们可以用多种方法求出 $\int_0^{+\infty}\dfrac{\sin x}{x}\mathrm dx$ 的值。类似的例子还有 $\int_0^{+\infty}e^{-x^2}\mathrm dx$ 等.
% \end{solve}
% 
% 
% \begin{example}\label{example:ma-8-2-6}
% 试讨论 $I=\int_1^{+\infty}\ln\left(1+\dfrac{\sin x}{x^p}\right)\mathrm dx(p>0)$ 的敛散性.
% \end{example}
% 
% 
% \begin{solve}
% 由
% \[
%   \ln(1+t)=t-\frac{t^2}{2}+o(t^2)\quad(t\to0)
% \]
% 知
% \[
% \begin{aligned}
%   \ln\left(1+\frac{\sin x}{x^p}\right)
%   &=
%   \frac{\sin x}{x^p}
%   -
%   \frac{\sin^2x}{2x^{2p}}
%   +
%   o\left(\frac1{x^{2p}}\right)\\
%   &=
%   \frac{\sin x}{x^p}
%   +
%   \frac{\cos2x}{4x^{2p}}
%   -
%   [1-o(1)]\frac1{2x^{2p}}
%   \quad(x\to+\infty).
% \end{aligned}
% \]
% 因此我们有
% \[
% \begin{aligned}
%   I
%   &=
%   \int_1^{+\infty}\ln\left(1+\frac{\sin x}{x^p}\right)\mathrm dx\\
%   &=
%   \int_1^{+\infty}\frac{\sin x}{x^p}\mathrm dx
%   +
%   \int_1^{+\infty}\frac{\cos2x}{4x^{2p}}\mathrm dx
%   -
%   \int_1^{+\infty}[1-o(1)]\frac{\mathrm dx}{2x^{2p}}\\
%   &\triangleq I_1+I_2-I_3.
% \end{aligned}
% \]
% 当 $0<p\leq\dfrac12$ 时, 由狄利克雷判别法知 $I_1$ 和 $I_2$ 收敛, 而注意到此时 $I_3$ 发散, 因此 $I$ 发散; 当 $1\geq p>\dfrac12$ 时, $I_1$ 条件收敛, 而 $I_2$ 和 $I_3$ 绝对收敛, 因此 $I$ 条件收敛; 当 $p>1$ 时, $I_1,I_2$ 和 $I_3$ 均绝对收敛, 从而 $I$ 绝对收敛.
% \end{solve}
% 
% 
% \begin{remark}\label{remark:ma-8-source-42}
% 注意到对所有的 $p>0$, 我们都有
% \[
%   \lim_{x\to+\infty}
%   \frac{\ln\left(1+\dfrac{\sin x}{x^p}\right)}{\dfrac{\sin x}{x^p}}
%   =
%   1.
% \]
% 此例说明若被积函数不保号时, 不能用比较判别法来确定它是否收敛.
% \end{remark}
% 
% 
% \section{瑕积分}
% \label{sec:ma-8-3}
% 
% \subsection{瑕积分的概念}
% \label{subsec:ma-8-3-1}
% 
% 正如我们在本章开始所指出的, 定积分的第二个推广是研究在有限区间上无界函数的积分问题。回忆一下, 在研究函数可积性时, 我们知道
% \[
%   F(x)=
%   \begin{cases}
%     x^2\sin\dfrac1{x^2}, & x\ne0,\\
%     0, & x=0
%   \end{cases}
% \]
% 是函数
% \[
%   f(x)=
%   \begin{cases}
%     2x\sin\dfrac1{x^2}-\dfrac2x\cos\dfrac1{x^2}, & x\ne0,\\
%     0, & x=0
%   \end{cases}
% \]
% 的一个原函数。函数 $f(x)$ 是无界的, 因此它在包含 $x=0$ 的任何闭区间 $[a,b]$ 上不可积, 但此时函数 $f(x)$ 却有一个有界的原函数。在本节中, 我们希望研究此类函数的积分问题。我们称这类函数的积分为瑕积分。因为 $f(x)$ 仅仅在 $x=0$ 附近无界, 所以 $x=0$ 也称为 $f(x)$ 的一个瑕点。下面我们称 $x_0$ 是 $f(x)$ 的一个瑕点即是指 $f(x)$ 在 $x_0$ 的某个去心(左或右)邻域内有定义, 但在该去心(左或右)邻域内无界.
% 
% 首先我们讨论函数 $f(x)$ 在区间 $[a,b]$ 上只有一个瑕点的情形.
% 
% \begin{definition}\label{definition:ma-8-3-1}
% 设函数 $f(x)$ 在区间 $(a,b]$ 上有定义, $a$ 是 $f(x)$ 的一个瑕点。若对于 $\forall 0<\delta<b-a$, $f(x)$ 在区间 $[a+\delta,b]$ 上可积, 且极限
% \begin{equation}
%   \lim_{\delta\to+0}\int_{a+\delta}^b f(x)\mathrm dx
%   \label{eq:ma-8-3-1}
% \end{equation}
% 存在, 则称瑕积分 $\int_a^b f(x)\mathrm dx$ 收敛, 并记
% \[
%   \int_a^b f(x)\mathrm dx
%   =
%   \lim_{\delta\to+0}\int_{a+\delta}^b f(x)\mathrm dx;
% \]
% \end{definition}
% 
% 若极限 \eqref{eq:ma-8-3-1} 不存在, 则称瑕积分 $\int_a^b f(x)\mathrm dx$ 发散.
% 
% 如果 $b$ 为函数 $f(x)$ 的瑕点, 我们可以类似定义
% \[
%   \int_a^b f(x)\mathrm dx
%   =
%   \lim_{\delta\to+0}\int_a^{b-\delta}f(x)\mathrm dx.
% \]
% 
% 当 $c\in(a,b)$ 为 $f(x)$ 在 $[a,b]$ 上的唯一瑕点时, 我们称 $\int_a^b f(x)\mathrm dx$ 收敛是指瑕积分 $\int_a^c f(x)\mathrm dx$ 与 $\int_c^b f(x)\mathrm dx$ 同时收敛。此时, 我们还可类似地定义柯西主值意义下的收敛(请读者给出).
% 
% 当 $a$ 是函数 $f(x)$ 在区间 $[a,b]$ 上的唯一瑕点时, 若 $F(x)$ 是 $f(x)$ 在 $(a,b]$ 上的一个原函数, 则瑕积分 $\int_a^b f(x)\mathrm dx$ 可由下式计算出:
% \[
%   \int_a^b f(x)\mathrm dx
%   =
%   \lim_{\delta\to+0}\int_{a+\delta}^b f(x)\mathrm dx
%   =
%   F(b)-\lim_{\delta\to+0}F(a+\delta).
% \]
% 例如, 对于前面的例子我们有
% \[
%   \int_0^1\left(2x\sin\frac1{x^2}-\frac2x\cos\frac1{x^2}\right)\mathrm dx
%   =
%   \lim_{\delta\to+0}x^2\sin\frac1{x^2}\bigg|_\delta^1
%   =
%   \sin1.
% \]
% 以上的分析实际上给出了我们计算瑕积分的一个方法。对于瑕积分, 我们也有相应的换元公式与分部积分公式, 这些公式请读者自己给出.
% 
% \begin{example}\label{example:ma-8-3-1}
% 计算瑕积分 $\int_{-1}^1\dfrac{\arccos x}{\sqrt{1-x^2}}\mathrm dx$.
% \end{example}
% 
% 
% \begin{solve}
% 由洛必达法则我们有
% \[
%   \lim_{x\to1-0}\frac{\arccos x}{\sqrt{1-x^2}}
%   =
%   \lim_{x\to1-0}
%   \frac{-\dfrac1{\sqrt{1-x^2}}}{-\dfrac{2x}{2\sqrt{1-x^2}}}
%   =
%   1,
% \]
% 因此被积函数有唯一的瑕点 $x=-1$。所以有
% \[
%   \int_{-1}^1\frac{\arccos x}{\sqrt{1-x^2}}\mathrm dx
%   =
%   \lim_{\delta\to+0}\left(-\frac12\arccos^2x\right)\bigg|_{-1+\delta}^1
%   =
%   \frac{\pi^2}{2}.
% \qedhere
% \]
% \end{solve}
% 
% 
% \begin{example}\label{example:ma-8-3-2}
% 计算瑕积分 $\int_0^1\ln x\mathrm dx$.
% \end{example}
% 
% 
% \begin{solve}
% 容易看出 $x=0$ 是瑕点。利用分部积分法得
% \[
%   \int_0^1\ln x\mathrm dx
%   =
%   \lim_{\delta\to+0}x\ln x\big|_\delta^1
%   -
%   \int_0^1\mathrm dx
%   =
%   -1.
% \qedhere
% \]
% \end{solve}
% 
% 
% \begin{example}\label{example:ma-8-3-3}
% 讨论瑕积分 $\int_a^b\dfrac{\mathrm dx}{(b-x)^p}$ 的敛散性.
% \end{example}
% 
% 
% \begin{solve}[{解法 1}]
% 当 $p=1$ 时,
% \[
% \begin{aligned}
%   \int_a^b\frac{\mathrm dx}{b-x}
%   &=
%   \lim_{\delta\to+0}\int_a^{b-\delta}\frac{\mathrm dx}{b-x}
%   =
%   \lim_{\delta\to+0}[-\ln(b-x)]\big|_a^{b-\delta}\\
%   &=
%   \lim_{\delta\to+0}[\ln(b-a)-\ln\delta]=+\infty.
% \end{aligned}
% \]
% 当 $p\ne1$ 时,
% \[
% \begin{aligned}
%   \int_a^b\frac{\mathrm dx}{(b-x)^p}
%   &=
%   \lim_{\delta\to+0}\int_a^{b-\delta}\frac{\mathrm dx}{(b-x)^p}\\
%   &=
%   -\lim_{\delta\to+0}\frac1{1-p}(b-x)^{1-p}\bigg|_a^{b-\delta}\\
%   &=
%   \begin{cases}
%     \dfrac{(b-a)^{1-p}}{1-p}, & p<1,\\
%     +\infty, & p>1.
%   \end{cases}
% \end{aligned}
% \]
% 因此, $\int_a^b\dfrac{\mathrm dx}{(b-x)^p}$ 当 $p<1$ 时收敛, 当 $p\geq1$ 时发散.
% \end{solve}
% 
% 
% \begin{solve}[{解法 2}]
% 我们也可以用以下办法来讨论 $\int_a^b\dfrac{\mathrm dx}{(b-x)^p}$ 的敛散性。作变换 $t=\dfrac1{b-x}$, 则当 $x=a$ 时, 有 $t=\dfrac1{b-a}$, 当 $x\to b-0$ 时, 有 $t\to+\infty$, 并且
% \[
%   \mathrm dt=\frac{\mathrm dx}{(b-x)^2}=t^2\mathrm dx,\quad
%   \mathrm dx=t^{-2}\mathrm dt.
% \]
% 将它们代入 $\int_a^b\dfrac{\mathrm dx}{(b-x)^p}$, 得
% \[
%   \int_a^b\frac{\mathrm dx}{(b-x)^p}
%   =
%   \int_{\frac1{b-a}}^{+\infty}t^{-2}t^p\mathrm dt
%   =
%   \int_{\frac1{b-a}}^{+\infty}\frac{\mathrm dt}{t^{2-p}}.
% \]
% 由无穷积分的结果知, 当 $2-p>1$, 即 $p<1$ 时, $\int_a^b\dfrac{\mathrm dx}{(b-x)^p}$ 收敛; 当 $2-p\leq1$, 即 $p\geq1$ 时, $\int_a^b\dfrac{\mathrm dx}{(b-x)^p}$ 发散.
% 
% 一般来说, 通过类似于上述解法 2 中的变换, 一个瑕积分总可以转化成一个无穷积分。另外, 一个无穷积分 $\int_a^{+\infty}f(x)\mathrm dx(a>0)$ 通过变换 $x=\dfrac1t$ 可以化为 $\int_0^{\frac1a}\dfrac1{t^2}f\left(\dfrac1t\right)\mathrm dt$, 因此一个无穷积分就有可能化成一个定积分或者瑕积分.
% \end{solve}
% 
% 
% \subsection{瑕积分敛散性的判别法}
% \label{subsec:ma-8-3-2}
% 
% 对于瑕积分的敛散性的判定, 一个方法是将其化为无穷积分来讨论。但有时由于化成无穷积分后被积函数会更加复杂, 因此我们需要讨论直接从瑕积分本身来判定的方法。以下我们列出函数 $f(x)$ 在区间 $[a,b]$ 中仅有瑕点 $b$ 时瑕积分的有关判定定理, 其证明则留给读者。因此, 在本小节中, 定义在 $[a,b)$ 上的函数 $f(x),g(x)$ 等, 总假定 $b$ 是它们的瑕点, 并且它们在 $[a,b)$ 的任何子区间上可积。另外, 我们要指出的是, 对于瑕积分, 同样有绝对收敛、条件收敛等概念和相应的结论.
% 
% \begin{theorem}[{柯西准则}]\label{theorem:ma-8-3-1}
% 瑕积分 $\int_a^b f(x)\mathrm dx$ 收敛的充分必要条件是: 对于 $\forall\varepsilon>0$, $\exists\delta>0$, 当 $0<\delta''<\delta'<\delta$ 时, 有
% \[
%   \left|\int_{b-\delta'}^{b-\delta''}f(x)\mathrm dx\right|<\varepsilon.
% \]
% \end{theorem}
% 
% 
% \begin{theorem}\label{theorem:ma-8-3-2}
% 设非负函数 $f(x),g(x)$ 在区间 $[a,b)$ 上满足: 存在正常数 $c_1,c_2$, 使得当 $x\in[b-\delta_0,b)(0<\delta_0<b-a)$ 时, 有
% \[
%   c_1f(x)\leq c_2g(x),
% \]
% 则
% 
% (1) 若 $\int_a^b g(x)\mathrm dx$ 收敛, 则 $\int_a^b f(x)\mathrm dx$ 也收敛;
% 
% (2) 若 $\int_a^b f(x)\mathrm dx$ 发散, 则 $\int_a^b g(x)\mathrm dx$ 也发散.
% \end{theorem}
% 
% 
% \begin{corollary}\label{corollary:ma-8-source-53}
% 设非负函数 $f(x),g(x)(g(x)\ne0)$ 满足
% \[
%   \lim_{x\to b-0}\frac{f(x)}{g(x)}=l,
% \]
% 则
% 
% (1) 当 $0<l<+\infty$ 时, $\int_a^b f(x)\mathrm dx$ 与 $\int_a^b g(x)\mathrm dx$ 同时收敛或同时发散;
% 
% (2) 当 $l=0$ 时, 若 $\int_a^b g(x)\mathrm dx$ 收敛, 则 $\int_a^b f(x)\mathrm dx$ 收敛;
% 
% (3) 当 $l=+\infty$ 时, 若 $\int_a^b g(x)\mathrm dx$ 发散, 则 $\int_a^b f(x)\mathrm dx$ 发散.
% \end{corollary}
% 
% 同样, 对于瑕积分, 我们也有一些 ``标准函数'' 供其他函数来比较, 如 $\dfrac1{(b-x)^p}$.
% 
% \begin{example}\label{example:ma-8-3-4}
% 讨论瑕积分 $\int_{-1}^1\dfrac{\mathrm dx}{(1-x^2)^p}$ 的敛散性.
% \end{example}
% 
% 
% \begin{solve}
% 非负被积函数有两个瑕点 $x=\pm1$。由于
% \[
%   \lim_{x\to1-0}
%   \frac{\dfrac1{(1-x^2)^p}}{\dfrac1{(1-x)^p}}
%   =
%   \frac1{2^p},\quad
%   \lim_{x\to-1+0}
%   \frac{\dfrac1{(1-x^2)^p}}{\dfrac1{(1+x)^p}}
%   =
%   \frac1{2^p},
% \]
% 因此可知, 当 $p<1$ 时, 瑕积分收敛; 当 $p\geq1$ 时, 瑕积分发散.
% 
% 对于瑕积分, 我们还有以下的判别法.
% \end{solve}
% 
% 
% \begin{theorem}[{狄利克雷判别法}]\label{theorem:ma-8-3-3}
% 设函数 $f(x),g(x)$ 在区间 $[a,b)$ 上满足下述条件:
% 
% (1) $\exists M>0$, 使得对于 $\forall\delta>0$, 有
% \[
%   \left|\int_a^{b-\delta}f(x)\mathrm dx\right|\leq M;
% \]
% 
% (2) $g(x)$ 在 $[a,b)$ 上单调, 且 $\lim\limits_{x\to b-0}g(x)=0$,
% 
% 则瑕积分 $\int_a^b f(x)g(x)\mathrm dx$ 收敛.
% \end{theorem}
% 
% 
% \begin{theorem}[{阿贝尔判别法}]\label{theorem:ma-8-3-4}
% 设函数 $f(x),g(x)$ 在区间 $[a,b)$ 上满足下述条件:
% 
% (1) $\int_a^b f(x)\mathrm dx$ 收敛;
% 
% (2) $g(x)$ 在 $[a,b)$ 上单调有界,
% 
% 则瑕积分 $\int_a^b f(x)g(x)\mathrm dx$ 收敛.
% \end{theorem}
% 
% 
% \begin{example}\label{example:ma-8-3-5}
% 讨论广义积分 $\int_0^{+\infty}\dfrac{\ln(1+x)}{x^\alpha}\mathrm dx(\alpha\in\mathbb R)$ 的敛散性.
% \end{example}
% 
% 
% \begin{solve}
% 这个广义积分不仅是一个无穷积分, 而且当 $\alpha>1$ 时, 还有一瑕点 $x=0$。为了讨论该广义积分的敛散性, 我们必须分别讨论
% \[
%   \int_0^1\frac{\ln(1+x)}{x^\alpha}\mathrm dx,\quad
%   \int_1^{+\infty}\frac{\ln(1+x)}{x^\alpha}\mathrm dx
% \]
% 的敛散性.
% 
% 当 $x\to0+0$ 时, 我们有
% \[
%   \frac{\ln(1+x)}{x^\alpha}\sim\frac1{x^{\alpha-1}}.
% \]
% 因此, 当 $\alpha<2$ 时, $\int_0^1\dfrac{\ln(1+x)}{x^\alpha}\mathrm dx$ 收敛; 当 $\alpha\geq2$ 时, $\int_0^1\dfrac{\ln(1+x)}{x^\alpha}\mathrm dx$ 发散.
% 
% 当 $x\to+\infty$ 时, 如果 $\alpha\leq1$, 则有
% \[
%   \lim_{x\to+\infty}
%   \frac{\dfrac{\ln(1+x)}{x^\alpha}}{\dfrac1{x^\alpha}}
%   =
%   \lim_{x\to+\infty}\ln(1+x)=+\infty.
% \]
% 由于 $\int_1^{+\infty}\dfrac{\mathrm dx}{x^\alpha}$ 发散, 因此有 $\int_1^{+\infty}\dfrac{\ln(1+x)}{x^\alpha}\mathrm dx$ 发散.
% 
% 当 $\alpha>1$ 时, 取 $\alpha'$, 使得 $1<\alpha'<\alpha$, 则有
% \[
%   \lim_{x\to+\infty}
%   \frac{\dfrac{\ln(1+x)}{x^\alpha}}{\dfrac1{x^{\alpha'}}}
%   =
%   \lim_{x\to+\infty}\frac{\ln(1+x)}{x^{\alpha-\alpha'}}=0.
% \]
% 由于 $\int_1^{+\infty}\dfrac{\mathrm dx}{x^{\alpha'}}$ 收敛, 从而 $\int_1^{+\infty}\dfrac{\ln(1+x)}{x^\alpha}\mathrm dx$ 收敛.
% 
% 因此, 当 $1<\alpha<2$ 时, $\int_0^{+\infty}\dfrac{\ln(1+x)}{x^\alpha}\mathrm dx$ 收敛; 当 $\alpha\leq1$ 或 $\alpha\geq2$ 时, $\int_0^{+\infty}\dfrac{\ln(1+x)}{x^\alpha}\mathrm dx$ 发散.
% \end{solve}
% 
% 
% \section*{习题八}
% \phantomsection
% \addcontentsline{toc}{section}{习题八}
% \label{exercises:ma-8}
% 
% 1。计算下列无穷积分:
% 
% (1) $\int_1^{+\infty}\dfrac{\ln x}{(1+x)^2}\mathrm dx$;
% 
% (2) $\int_0^{+\infty}\dfrac{x\mathrm dx}{(x^2+a^2)^{\frac32}}\ (a\ne0)$;
% 
% (3) $\int_1^{+\infty}\dfrac{\mathrm dx}{x\sqrt{x^2+x+1}}$;
% 
% (4) $\int_1^{+\infty}\dfrac{\arctan x}{x^2}\mathrm dx$;
% 
% (5) $\int_0^{+\infty}\dfrac{\mathrm dx}{(1+x^2)^n}$.
% 
% 2。讨论下列无穷积分的敛散性:
% 
% (1) $\int_0^{+\infty}\dfrac{\sin x}{1+e^{-x}}\mathrm dx$;
% 
% (2) $\int_0^{+\infty}\dfrac{x^2+100x+1000}{x^4-x^3+1}\mathrm dx$;
% 
% % 整合来源：math-9.tex
% \chapter{数项级数}
% \label{chap:ma-9}
% 
% 无穷级数的理论具有悠久的历史, 许多有趣的数学问题与它有关系, 在日常生活中我们也常常会遇到与它有关的问题。如在古代人们为了计算圆的面积, 用圆内接正多边形的面积来逼近圆的面积, 该过程实际上就是通过一个无穷级数的部分和序列来求得圆面积具有任何精度的近似值。另外, 无穷级数中有许多研究技巧与方法, 在许多学科中也有着广泛的用途。当然, 若仅仅限于以上这些因素, 无穷级数未必能成为数学分析课程中专门讨论的一个对象。我们在数学分析中研究无穷级数, 最主要的原因还是它对函数的研究起着重要的作用。在本书的后续章节中可以看出这一点.
% 
% \section{数项级数的基本概念}
% \label{sec:ma-9-1}
% 
% 人们为了拓广函数类, 很原始的想法是将无穷多个函数相加起来。我们注意到无穷多个函数相加时, 对定义域中固定的每一点来说, 也就是无穷多个数相加, 因此必须先研究形如
% \begin{equation}
%   a_1+\cdots+a_n+\cdots\triangleq \sum_{n=1}^{+\infty}a_n
%   \label{eq:ma-9-1-1}
% \end{equation}
% 的和式。另外, 在许多实际问题中, 我们也会遇到类似的和式, 如大家在中学学过的等比数列的和等。在学习序列极限时, 我们也曾遇到过形如
% \[
%   x_n=1+\frac12+\frac13+\cdots+\frac1n,\quad n=1,2,\cdots,
% \]
% \[
%   y_n=1+\frac1{2^2}+\frac1{3^2}+\cdots+\frac1{n^2},\quad n=1,2,\cdots
% \]
% 的序列。读者不难发现, 对 $n\geq1$, $x_n$ 其实就是和式 $\sum\limits_{k=1}^{+\infty}\dfrac1k$ 的前 $n$ 项之和; 而对 $n\geq1$, $y_n$ 则是 $\sum\limits_{k=1}^{+\infty}\dfrac1{k^2}$ 的前 $n$ 项之和。通过前面的学习我们已经知道 $\lim\limits_{n\to+\infty}x_n=+\infty$, 而 $\lim\limits_{n\to+\infty}y_n$ 是存在的。因此我们很自然地就说 $\sum\limits_{n=1}^{+\infty}\dfrac1n$ 是发散的, 而 $\sum\limits_{n=1}^{+\infty}\dfrac1{n^2}$ 是收敛的, 且 $\sum\limits_{n=1}^{+\infty}\dfrac1{n^2}=\lim\limits_{n\to+\infty}y_n$.
% 
% \subsection{数项级数的基本概念}
% \label{subsec:ma-9-1-1}
% 
% 设 $a_1,a_2,\cdots,a_n,\cdots$ 为一个序列, 我们称
% \[
%   a_1+a_2+\cdots+a_n+\cdots
% \]
% 为一个数项级数(简称级数), 记为 $\sum\limits_{n=1}^{+\infty}a_n$, 其中 $a_n(n=1,2,\cdots)$ 称为级数的通项, 并称 $S_n=\sum\limits_{k=1}^n a_k(n=1,2,\cdots)$ 为级数的前 $n$ 项部分和, 而称 $\{S_n\}$ 为部分和序列.
% 
% \begin{definition}\label{definition:ma-9-1-1}
% 设 $\sum\limits_{n=1}^{+\infty}a_n$ 是一个数项级数。如果它的部分和序列 $\{S_n\}$ 是收敛的, 即 $\lim\limits_{n\to\infty}S_n$ 存在, 则称级数 $\sum\limits_{n=1}^{+\infty}a_n$ 收敛, 并且记 $\sum\limits_{n=1}^{+\infty}a_n=\lim\limits_{n\to\infty}S_n$。这时也称极限 $\lim\limits_{n\to\infty}S_n$ 为级数 $\sum\limits_{n=1}^{+\infty}a_n$ 的和。如果 $\{S_n\}$ 是发散序列, 则称级数 $\sum\limits_{n=1}^{+\infty}a_n$ 发散。特别地, 若 $\lim\limits_{n\to\infty}S_n=\infty(\pm\infty)$, 我们也称 $\sum\limits_{n=1}^{+\infty}a_n$ 发散到 $\infty(\pm\infty)$.
% \end{definition}
% 
% 从以上定义知, 级数的敛散性本质上可以归结到序列的敛散性。但我们以后会发现, 对级数的研究有许多特别的工具与方法, 似乎这些工具与方法不太容易直接从序列的研究中得到.
% 
% \begin{example}\label{example:ma-9-1-1}
% 对 $q\in\mathbb R$, 形如
% \begin{equation}
%   \sum_{n=1}^{+\infty}q^{n-1}=1+q+q^2+\cdots
%   \label{eq:ma-9-1-2}
% \end{equation}
% 的级数称为等比级数(或几何级数)。试讨论该级数的敛散性.
% \end{example}
% 
% 
% \begin{solve}
% 对于 $\forall n\in\mathbb N$, 该级数的前 $n$ 项部分和为
% \[
%   S_n(q)=\sum_{k=1}^n q^{k-1}=\frac{1-q^n}{1-q}\quad(q\neq1).
% \]
% 当 $|q|<1$ 时, 我们有 $\lim\limits_{n\to\infty}S_n(q)=\dfrac1{1-q}$。因此此时级数 \eqref{eq:ma-9-1-2} 收敛, 并且有
% \[
%   \sum_{n=1}^{+\infty}q^{n-1}=\frac1{1-q}.
% \]
% 当 $q=-1$ 时, $S_n(-1)=\dfrac{1-(-1)^n}{2}$, 显然 $\lim\limits_{n\to\infty}S_n(q)$ 不存在。因此此时级数 \eqref{eq:ma-9-1-2} 发散, 即 $\sum\limits_{n=1}^{+\infty}(-1)^{n-1}$ 发散.
% 
% 当 $q=1$ 时, 有 $S_n(1)=n(n=1,2,\cdots)$, 因此 $\lim\limits_{n\to\infty}S_n(1)=+\infty$, 从而此时级数 \eqref{eq:ma-9-1-2} 发散.
% 
% 当 $|q|>1$ 时, $\lim\limits_{n\to\infty}S_n(q)=\lim\limits_{n\to\infty}\dfrac{q^n-1}{q-1}=\infty$, 因此此时级数 \eqref{eq:ma-9-1-2} 发散.
% 
% 所以当 $|q|<1$ 时, 等比级数 \eqref{eq:ma-9-1-2} 收敛; 当 $|q|\geq1$ 时, 级数 \eqref{eq:ma-9-1-2} 发散。特别地, 当 $q<-1$ 时, 级数 \eqref{eq:ma-9-1-2} 发散到 $\infty$; 当 $q>1$ 时, 级数 \eqref{eq:ma-9-1-2} 发散到 $+\infty$.
% \end{solve}
% 
% 
% \begin{example}\label{example:ma-9-1-2}
% 对固定的 $n_0\in\mathbb N$, 证明数项级数 $\sum\limits_{n=1}^{+\infty}\dfrac1{n(n+n_0)}$ 收敛, 并求其值.
% \end{example}
% 
% 
% \begin{solve}
% 对于 $\forall n\in\mathbb N$, 该级数的前 $n$ 项部分和 $S_n$ 可以表示成
% \[
% \begin{aligned}
%   S_n
%   &=\sum_{k=1}^n\frac1{k(k+n_0)}
%     =\sum_{k=1}^n\frac1{n_0}\left(\frac1k-\frac1{k+n_0}\right)\\
%   &=\frac1{n_0}\left(\sum_{k=1}^{n_0}\frac1k+\sum_{k=n_0+1}^n\frac1k-\sum_{k=n_0+1}^{n+n_0}\frac1k\right)\\
%   &=\frac1{n_0}\sum_{k=1}^{n_0}\frac1k+\frac1{n_0}\sum_{k=n+1}^{n+n_0}\frac1k.
% \end{aligned}
% \]
% 由于
% \[
%   0<\frac1{n_0}\sum_{k=n+1}^{n+n_0}\frac1k\leq \frac1{n_0}\cdot\frac{n_0}{n+1}\to0\quad(n\to\infty),
% \]
% 因此 $\sum\limits_{n=1}^{+\infty}\dfrac1{n(n+n_0)}$ 收敛, 并且
% \[
%   \sum_{n=1}^{+\infty}\frac1{n(n+n_0)}
%   =\lim_{n\to\infty}S_n=\frac1{n_0}\sum_{k=1}^{n_0}\frac1k.
% \qedhere
% \]
% \end{solve}
% 
% 
% \begin{example}\label{example:ma-9-1-3}
% 证明 $\sum\limits_{n=0}^{+\infty}\dfrac1{n!}=e$, 并证明 $e$ 是无理数.
% \end{example}
% 
% 
% \begin{proof}
% 记 $S_n=\sum\limits_{k=0}^{n-1}\dfrac1{k!}(n=1,2,\cdots)$。对任意固定的 $n>1$, 任取 $k>n$, 有
% \[
%   e>\left(1+\frac1k\right)^k
%   >1+1+\frac1{2!}\left(1-\frac1k\right)+\cdots+\frac1{(n-1)!}\left(1-\frac1k\right)\cdots\left(1-\frac{n-2}{k}\right).
% \]
% 在上式中令 $k\to\infty$, 得
% \[
%   1+1+\frac1{2!}+\frac1{3!}+\cdots+\frac1{(n-1)!}\leq e.
% \]
% 再由
% \[
%   \left(1+\frac1{n-1}\right)^{n-1}<1+1+\frac1{2!}+\frac1{3!}+\cdots+\frac1{(n-1)!}\leq e,
% \]
% 得 $\lim\limits_{n\to\infty}S_n=e$。由级数收敛的定义知 $e=\sum\limits_{n=0}^{+\infty}\dfrac1{n!}$.
% 
% 对任意正整数 $n,m$, 有
% \[
% \begin{aligned}
%   0<S_{n+m}-S_n
%   &=\frac1{n!}+\frac1{(n+1)!}+\cdots+\frac1{(n+m-1)!}\\
%   &\leq\frac1{n!}\left[1+\frac1{n+1}+\frac1{(n+1)^2}+\cdots+\frac1{(n+1)^{m-1}}\right]\\
%   &<\frac1{n!}\cdot\frac{n+1}{n}.
% \end{aligned}
% \]
% 在上式中固定 $n$, 让 $m\to+\infty$, 得
% \begin{equation}
%   0<e-S_n\leq\frac1{n!}\cdot\frac{n+1}{n}.
%   \label{eq:ma-9-1-3}
% \end{equation}
% 从上式我们容易推出 $e$ 是无理数。事实上, 若 $e$ 是有理数, 则存在既约正整数 $p,q$, 使得 $e=\dfrac qp$。取定 $n\geq\max\{2,p\}+1$, 在 \eqref{eq:ma-9-1-3} 后一不等式的两边乘上 $(n-1)!$, 则其左边是一个正整数, 而右边是一个真分数。此矛盾便证明了 $e$ 是一个无理数.
% 
% 在第一册里, 我们曾利用指数函数 $e^x$ 的带拉格朗日余项的泰勒公式证明过 $e$ 是无理数。在这里, 我们仅仅利用数项级数收敛的定义就证明了该结论.
% 
% 关于数项级数我们有以下简单的性质, 其证明作为练习留给读者.
% \end{proof}
% 
% 
% \begin{property}\label{property:ma-9-1-1}
% 改变数项级数有限项的值后得到的新级数与原级数敛散性相同.
% \end{property}
% 
% 
% \begin{property}\label{property:ma-9-1-2}
% 对于任意常数 $k\neq0$, 数项级数 $\sum\limits_{n=1}^{+\infty}a_n$ 与 $\sum\limits_{n=1}^{+\infty}ka_n$ 的敛散性相同.
% \end{property}
% 
% 
% \begin{property}\label{property:ma-9-1-3}
% 设 $\alpha,\beta\in\mathbb R$, 数项级数 $\sum\limits_{n=1}^{+\infty}a_n$ 和 $\sum\limits_{n=1}^{+\infty}b_n$ 都收敛, 则数项级数 $\sum\limits_{n=1}^{+\infty}(\alpha a_n+\beta b_n)$ 也收敛, 并且成立下述等式:
% \[
%   \sum_{n=1}^{+\infty}(\alpha a_n+\beta b_n)
%   =\alpha\sum_{n=1}^{+\infty}a_n+\beta\sum_{n=1}^{+\infty}b_n.
% \]
% \end{property}
% 
% 在性质~\ref{property:ma-9-1-3} 中, 特别地, 当 $\beta=0$ 时, 对于收敛的数项级数 $\sum\limits_{n=1}^{+\infty}a_n$, 有
% \[
%   \sum_{n=1}^{+\infty}\alpha a_n=\alpha\sum_{n=1}^{+\infty}a_n\quad(\alpha\in\mathbb R).
% \]
% 
% \subsection{柯西准则}
% \label{subsec:ma-9-1-2}
% 
% 对于很多数项级数, 并不能简单地计算出它的部分和, 因此很难从定义来判别其敛散性, 从而我们必须直接从数项级数的本身来研究它是否收敛。在下面的两节中我们将详细地研究这个问题。在这里, 我们根据数项级数敛散性的定义, 给出以下判别级数收敛的柯西准则.
% 
% \begin{theorem}[{柯西准则}]\label{theorem:ma-9-1-4}
% 设 $\sum\limits_{n=1}^{+\infty}a_n$ 是一个数项级数, 则它收敛的充分必要条件是: 对于 $\forall\varepsilon>0$, $\exists N>0$, 当 $n>m>N$ 时, 有
% \[
%   \left|\sum_{k=m+1}^n a_k\right|
%   =|a_{m+1}+a_{m+2}+\cdots+a_n|<\varepsilon.
% \]
% 此定理的证明可以从数项级数收敛的定义直接推出, 故省略.
% \end{theorem}
% 
% 若数项级数 $\sum\limits_{n=1}^{+\infty}a_n$ 收敛, 由柯西准则, 对于 $\forall\varepsilon>0$, $\exists N>0$, 特别地取 $n,n-1\geq N$ 时, 我们有 $|a_n|<\varepsilon$, 从而我们有下面的重要结论.
% 
% \begin{corollary}\label{corollary:ma-9-source-12}
% 设 $\sum\limits_{n=1}^{+\infty}a_n$ 是一个数项级数, 若它收敛, 则必有
% \[
%   \lim_{n\to\infty}a_n=0.
% \]
% \end{corollary}
% 
% 显然, 推论的逆命题是不成立的, 如级数 $\sum\limits_{n=1}^{+\infty}\dfrac1n$(称为调和级数)是发散的, 但是有 $\lim\limits_{n\to\infty}\dfrac1n=0$.
% 
% \begin{example}\label{example:ma-9-1-4}
% 证明数项级数 $\sum\limits_{n=1}^{+\infty}n\ln\left(1+\dfrac1n\right)$ 发散.
% \end{example}
% 
% 
% \begin{proof}
% 由于
% \[
%   \lim_{n\to\infty}n\ln\left(1+\frac1n\right)
%   =\lim_{n\to\infty}\ln\left(1+\frac1n\right)^n=1,
% \]
% 因此由定理~\ref{theorem:ma-9-1-4} 的推论即知, 该数项级数是发散的.
% \end{proof}
% 
% 
% \begin{example}\label{example:ma-9-1-5}
% 证明: 当 $\alpha\neq k\pi(k\in\mathbb Z)$ 时, 数项级数 $\sum\limits_{n=1}^{+\infty}\sin n\alpha$ 发散.
% \end{example}
% 
% 
% \begin{proof}
% 倘若 $\sum\limits_{n=1}^{+\infty}\sin n\alpha$ 收敛, 由推论~\ref{corollary:ma-9-source-12} 知 $\lim\limits_{n\to\infty}\sin n\alpha=0$。由于对 $\forall n\in\mathbb N$, 我们有
% \[
%   \sin(n+1)\alpha=\cos\alpha\sin n\alpha+\cos n\alpha\sin\alpha,
% \]
% 因此我们推出 $\lim\limits_{n\to\infty}\cos n\alpha=0$。显然这与恒等式 $\sin^2n\alpha+\cos^2n\alpha=1$ 矛盾。因此数项级数 $\sum\limits_{n=1}^{+\infty}\sin n\alpha$ 发散.
% \end{proof}
% 
% 
% \section{正项级数}
% \label{sec:ma-9-2}
% 
% \subsection{比较判别法}
% \label{subsec:ma-9-2-1}
% 
% 对于数项级数 $\sum\limits_{n=1}^{+\infty}a_n$ 而言, 与无穷积分一样, 也有绝对收敛和条件收敛的概念.
% 
% \begin{definition}\label{definition:ma-9-2-1}
% 设 $\sum\limits_{n=1}^{+\infty}a_n$ 为一个数项级数。如果 $\sum\limits_{n=1}^{+\infty}|a_n|$ 收敛, 则称 $\sum\limits_{n=1}^{+\infty}a_n$ 绝对收敛。如果 $\sum\limits_{n=1}^{+\infty}a_n$ 收敛, 但 $\sum\limits_{n=1}^{+\infty}|a_n|$ 发散, 则称 $\sum\limits_{n=1}^{+\infty}a_n$ 条件收敛.
% \end{definition}
% 
% 由柯西准则及三角不等式易于推出: 若 $\sum\limits_{n=1}^{+\infty}a_n$ 绝对收敛, 则 $\sum\limits_{n=1}^{+\infty}a_n$ 必定收敛.
% 
% 由定义~\ref{definition:ma-9-2-1} 可看出, 一个级数是否为绝对收敛的问题可以转化为所谓的正项级数的敛散性问题.
% 
% \begin{definition}\label{definition:ma-9-2-2}
% 设 $\sum\limits_{n=1}^{+\infty}a_n$ 是一个数项级数。若对于 $\forall n\in\mathbb N$, 有 $a_n\geq0$, 则称 $\sum\limits_{n=1}^{+\infty}a_n$ 为一个正项级数.
% \end{definition}
% 
% 对于正项级数, 以下结论是显然的.
% 
% \begin{theorem}\label{theorem:ma-9-2-1}
% 设 $\sum\limits_{n=1}^{+\infty}a_n$ 为一个正项级数, 则它收敛的充分必要条件是它的部分和序列是有界的。如果 $\sum\limits_{n=1}^{+\infty}a_n$ 发散, 则它必发散到 $+\infty$.
% \end{theorem}
% 
% 判别正项级数敛散性最基本的方法是如下的比较判别法.
% 
% \begin{theorem}[{比较判别法}]\label{theorem:ma-9-2-2}
% 设 $\sum\limits_{n=1}^{+\infty}a_n$ 与 $\sum\limits_{n=1}^{+\infty}b_n$ 为两个正项级数, $c_1,c_2$ 是两个正数。若存在 $N\in\mathbb N$, 当 $n>N$ 时, 有
% \begin{equation}
%   c_1a_n\leq c_2b_n,
%   \label{eq:ma-9-2-1}
% \end{equation}
% 则
% 
% (1) 当 $\sum\limits_{n=1}^{+\infty}b_n$ 收敛时, $\sum\limits_{n=1}^{+\infty}a_n$ 也收敛;
% 
% (2) 当 $\sum\limits_{n=1}^{+\infty}a_n$ 发散时, $\sum\limits_{n=1}^{+\infty}b_n$ 也发散.
% \end{theorem}
% 
% 
% \begin{proof}
% 因为一个级数改变有限多项的值后得到的级数与原级数具有相同的敛散性, 因此不妨设 \eqref{eq:ma-9-2-1} 式对一切 $n\geq1$ 都成立.
% 
% 对于 $n\in\mathbb N$, 分别记 $S_n$ 和 $S'_n$ 为 $\sum\limits_{n=1}^{+\infty}a_n$ 和 $\sum\limits_{n=1}^{+\infty}b_n$ 的前 $n$ 项部分和, 则由 \eqref{eq:ma-9-2-1} 式有
% \[
%   c_1S_n\leq c_2S'_n\quad(n=1,2,\cdots).
% \]
% (1) 当 $\sum\limits_{n=1}^{+\infty}b_n$ 收敛时, 由定理~\ref{theorem:ma-9-2-1} 知, $\exists M>0$, 使得对于 $\forall n\in\mathbb N$, 有
% \[
%   S'_n\leq M.
% \]
% 因此, 对一切 $n\in\mathbb N$, 有
% \[
%   S_n\leq\frac{c_2}{c_1}M.
% \]
% 由定理~\ref{theorem:ma-9-2-1} 知, $\sum\limits_{n=1}^{+\infty}a_n$ 也收敛.
% 
% (2) 当 $\sum\limits_{n=1}^{+\infty}a_n$ 发散时, 有 $S_n\to+\infty$, 从而有
% \[
%   \lim_{n\to\infty}S'_n\geq\lim_{n\to\infty}\frac{c_1}{c_2}S_n=+\infty.
% \]
% 因此 $\sum\limits_{n=1}^{+\infty}b_n$ 也发散.
% 
% 由极限的保号性及定理~\ref{theorem:ma-9-2-2}, 我们有以下的推论.
% \end{proof}
% 
% 
% \begin{corollary}\label{corollary:ma-9-source-22}
% 设正项级数 $\sum\limits_{n=1}^{+\infty}a_n$ 和 $\sum\limits_{n=1}^{+\infty}b_n(b_n\neq0,\forall n\in\mathbb N)$ 满足
% \[
%   \lim_{n\to\infty}\frac{a_n}{b_n}=l,
% \]
% 则
% 
% (1) 当 $0<l<+\infty$ 时, $\sum\limits_{n=1}^{+\infty}a_n$ 和 $\sum\limits_{n=1}^{+\infty}b_n$ 具有相同的敛散性;
% 
% (2) 当 $l=0$ 时, 若 $\sum\limits_{n=1}^{+\infty}b_n$ 收敛, 则 $\sum\limits_{n=1}^{+\infty}a_n$ 也收敛;
% 
% (3) 当 $l=+\infty$ 时, 若 $\sum\limits_{n=1}^{+\infty}b_n$ 发散, 则 $\sum\limits_{n=1}^{+\infty}a_n$ 也发散.
% \end{corollary}
% 
% 
% \begin{example}\label{example:ma-9-2-1}
% 证明: 正项级数 $\sum\limits_{n=1}^{+\infty}\dfrac1{n^p}$ 当 $p\leq1$ 时发散, 当 $p>1$ 时收敛.
% \end{example}
% 
% 
% \begin{proof}
% 当 $p=1$ 时, 我们已经知道 $\sum\limits_{n=1}^{+\infty}\dfrac1n$ 是发散的。当 $p<1$ 时, 对于 $\forall n>1$, 有 $\dfrac1n<\dfrac1{n^p}$, 因此 $\sum\limits_{n=1}^{+\infty}\dfrac1{n^p}$ 发散.
% 
% 为了证明 $\sum\limits_{n=1}^{+\infty}\dfrac1{n^p}$ 当 $p>1$ 时收敛, 我们首先注意到以下简单事实: 因为一个正项级数 $\sum\limits_{n=1}^{+\infty}a_n$ 的部分和序列 $\{S_n\}$ 是单调上升的, 因此若它有一个子列是有界的, 则 $\sum\limits_{n=1}^{+\infty}a_n$ 必收敛.
% 
% 注意到 $\left\{\dfrac1{n^p}\right\}$ 是一个单调下降序列, 对于 $\forall k\in\mathbb N$, 我们有
% \[
% \begin{aligned}
%   S_{2^{k+1}-1}
%   &=1+\left(\frac1{2^p}+\frac1{3^p}\right)
%     +\left(\frac1{4^p}+\cdots+\frac1{7^p}\right)+\cdots\\
%   &\quad+\left[\frac1{(2^k)^p}+\frac1{(2^k+1)^p}+\cdots+\frac1{(2^{k+1}-1)^p}\right]\\
%   &\leq1+2\cdot\frac1{2^p}+2^2\cdot\frac1{2^{2p}}+\cdots+2^k\cdot\frac1{2^{kp}}\\
%   &=\sum_{j=0}^k2^{(1-p)j}<\frac1{1-2^{1-p}}.
% \end{aligned}
% \]
% 上式说明, 当 $p>1$ 时, $\sum\limits_{n=1}^{+\infty}\dfrac1{n^p}$ 的部分和序列 $\{S_n\}$ 的子列 $\{S_{2^{k+1}-1}\}$ 有界, 因此该级数是收敛的.
% 
% 在正项级数敛散性的讨论中, 几何级数 $\sum\limits_{n=1}^{+\infty}q^n$ 与 $\sum\limits_{n=1}^{+\infty}\dfrac1{n^p}$ 是两个常用的参考级数。不仅许多其他的级数常用它们来作比较, 而且很多判别法也是建立在这两个级数的基础上的.
% \end{proof}
% 
% 
% \begin{example}\label{example:ma-9-2-2}
% 设 $F_1=1,F_2=2,F_n=F_{n-1}+F_{n-2}(n=3,4,\cdots)$, 证明正项级数 $\sum\limits_{n=1}^{+\infty}\dfrac1{F_n}$ 收敛.
% \end{example}
% 
% 
% \begin{proof}
% 对于 $n=1,2,\cdots$, 记 $S_n=\sum\limits_{k=1}^n\dfrac1{F_k}$。对任意 $n\geq3$, 我们有
% \[
%   F_{n-2}<F_{n-1}<2F_{n-2},
% \]
% 因此有
% \[
%   F_n>F_{n-1}+\frac12F_{n-1}=\frac32F_{n-1}.
% \]
% 利用上式容易看出
% \[
%   \lim_{n\to\infty}F_n=+\infty.
% \]
% 同时我们还可以推出
% \[
%   \sum_{k=3}^n\frac1{F_k}\leq\frac23\sum_{k=3}^n\frac1{F_{k-1}},
% \]
% 即
% \[
%   S_n-1-\frac12<\frac23S_n-\frac23-\frac23F_n^{-1}.
% \]
% 将上式化简得
% \[
%   S_n<\frac52-2F_n^{-1}<3.
% \]
% 因此 $\sum\limits_{n=1}^{+\infty}\dfrac1{F_n}$ 收敛.
% \end{proof}
% 
% 
% \begin{example}\label{example:ma-9-2-3}
% 试讨论数项级数 $\sum\limits_{n=3}^{+\infty}\ln\cos\dfrac{\pi}{n}$ 的敛散性.
% \end{example}
% 
% 
% \begin{solve}
% 记 $a_n=\ln\cos\dfrac{\pi}{n}$, 则 $a_n<0$。由
% \[
%   \ln\cos\frac{\pi}{n}
%   =\ln\left(1+\cos\frac{\pi}{n}-1\right)
%   \sim \cos\frac{\pi}{n}-1\sim-\frac{\pi^2}{2n^2}\quad(n\to\infty)
% \]
% 及 $\sum\limits_{n=1}^{+\infty}\dfrac{\pi^2}{2n^2}$ 收敛, 可知正项级数 $\sum\limits_{n=1}^{+\infty}\left(-\ln\cos\dfrac{\pi}{n}\right)$ 收敛, 因此原级数收敛.
% \end{solve}
% 
% 
% \begin{example}\label{example:ma-9-2-4}
% 试讨论正项级数 $\sum\limits_{n=1}^{+\infty}\left[1-\left(\dfrac{n-1}{n+1}\right)^k\right]^p(k>0,p>0)$ 的敛散性.
% \end{example}
% 
% 
% \begin{solve}
% 记 $a_n=\left[1-\left(\dfrac{n-1}{n+1}\right)^k\right]^p(n=1,2,\cdots)$。由
% \[
% \begin{aligned}
%   \left(\frac{n-1}{n+1}\right)^k
%   &=\left(1-\frac1n\right)^k\left(1+\frac1n\right)^{-k}\\
%   &=\left[1-\frac{k}{n}+o\left(\frac1n\right)\right]\left[1-\frac{k}{n}+o\left(\frac1n\right)\right]\\
%   &=1-\frac{2k}{n}+o\left(\frac1n\right)\quad(n\to\infty),
% \end{aligned}
% \]
% 得
% \[
%   a_n=\left[\frac{2k}{n}+o\left(\frac1n\right)\right]^p\sim\frac{(2k)^p}{n^p}\quad(n\to\infty).
% \]
% 由比较判别法知, 当 $p>1$ 时, $\sum\limits_{n=1}^{+\infty}a_n$ 收敛; 当 $p\leq1$ 时, $\sum\limits_{n=1}^{+\infty}a_n$ 发散.
% \end{solve}
% 
% 
% \begin{example}\label{example:ma-9-2-5}
% 讨论正项级数 $\sum\limits_{n=1}^{+\infty}\dfrac{n^{n-2}}{e^n n!}$ 的敛散性.
% \end{example}
% 
% 
% \begin{solve}
% 记 $a_n=\dfrac{n^{n-2}}{e^n n!}(n=1,2,\cdots)$。对任意 $n\geq2$, 我们有
% \[
%   \frac{a_{n+1}}{a_n}
%   =\frac{\left(1+\frac1n\right)^{n-2}}{e}
%   <\frac{\left(1+\frac1n\right)^{n-2}}{\left(1+\frac1n\right)^n}
%   =\frac{\dfrac1{(n+1)^2}}{\dfrac1{n^2}}
%  =\frac{b_{n+1}}{b_n},
% \]
% 其中 $b_n=\dfrac1{n^2}$.
% 
% 由上面的不等式我们推出
% \[
% \begin{aligned}
%   \frac{a_n}{a_2}
%   &=\frac{a_n}{a_{n-1}}\frac{a_{n-1}}{a_{n-2}}\cdots\frac{a_3}{a_2}\\
%   &<\frac{b_n}{b_{n-1}}\frac{b_{n-1}}{b_{n-2}}\cdots\frac{b_3}{b_2}
%   =\frac{b_n}{b_2}=\frac4{n^2}.
% \end{aligned}
% \]
% 因此当 $n\geq2$ 时, 有 $a_n<\dfrac2{e^2}\cdot\dfrac1{n^2}$。由定理~\ref{theorem:ma-9-2-2} 知 $\sum\limits_{n=1}^{+\infty}\dfrac{n^{n-2}}{e^n n!}$ 收敛.
% \end{solve}
% 
% 
% \subsection{达朗贝尔判别法与柯西判别法}
% \label{subsec:ma-9-2-2}
% 
% 以下两个判别法给出了从正项级数 $\sum a_n$ 本身的性质来判断其敛散性的方法, 但这些判别法的本质还是基于将正项级数与几何级数进行比较而得到的.
% 
% \begin{theorem}[{达朗贝尔(D'Alembert)判别法}]\label{theorem:ma-9-2-3}
% 设 $\sum a_n(a_n\neq0,n=1,2,\cdots)$ 为正项级数, 则
% 
% (1) 当
% \[
%   \overline{\lim}_{n\to\infty}\frac{a_{n+1}}{a_n}=\bar r<1
% \]
% 时, $\sum a_n$ 收敛;
% 
% (2) 当
% \[
%   \underline{\lim}_{n\to\infty}\frac{a_{n+1}}{a_n}=\underline r>1
% \]
% 时, $\sum a_n$ 发散.
% \end{theorem}
% 
% 
% \begin{proof}
% 我们先证 (1)。取 $r_1$, 使得其满足 $\bar r<r_1<1$, 则存在 $N_1\in\mathbb{N}$, 当 $n>N_1$ 时, 有 $\dfrac{a_n}{a_{n-1}}<r_1$。因此有
% \[
%   a_n=\frac{a_n}{a_{n-1}}\frac{a_{n-1}}{a_{n-2}}\cdots\frac{a_{N_1+1}}{a_{N_1}}a_{N_1}
%   <r_1^{n-N_1}a_{N_1}=C_1r_1^n,
% \]
% 其中 $C_1=a_{N_1}r_1^{-N_1}$。由于 $\sum r_1^n$ 是收敛的, 因此由定理~\ref{theorem:ma-9-2-2} 知 $\sum a_n$ 也收敛.
% 
% 现证 (2)。取 $r_2$, 使得其满足 $\underline r>r_2>1$, 则存在 $N_2\in\mathbb{N}$, 当 $n>N_2$ 时, 有 $\dfrac{a_n}{a_{n-1}}>r_2$。因此有
% \[
%   a_n=\frac{a_n}{a_{n-1}}\frac{a_{n-1}}{a_{n-2}}\cdots\frac{a_{N_2+1}}{a_{N_2}}a_{N_2}
%   >r_2^{n-N_2}a_{N_2}=C_2r_2^n,
% \]
% 其中 $C_2=a_{N_2}r_2^{-N_2}$。由此推出 $\lim\limits_{n\to\infty}a_n=+\infty$, 因此 $\sum a_n$ 发散.
% \end{proof}
% 
% 
% 对于 $\forall p>0$, 我们有
% \[
%   \lim_{n\to\infty}\frac{1/(n+1)^p}{1/n^p}
%   =\lim_{n\to\infty}\left(\frac{n}{n+1}\right)^p=1,
% \]
% 因此用达朗贝尔判别法我们无法判别正项级数 $\sum\dfrac1{n^p}$ 的敛散性。换句话说, 当定理~\ref{theorem:ma-9-2-3} 中的 $\underline r$ 和 $\bar r$ 等于 1 时, 达朗贝尔判别法是失效的.
% 
% \begin{example}\label{example:ma-9-2-6}
% 试讨论正项级数 $\sum\limits_{n=1}^{+\infty}\dfrac{x^n n!}{n^n}$ 的敛散性, 其中 $x\geq0$.
% \end{example}
% 
% 
% \begin{solve}
% 对于 $\forall n\in\mathbb{N}$, 记 $a_n=\dfrac{x^n n!}{n^n}$, 则有
% \[
%   \frac{a_{n+1}}{a_n}
%   =\frac{\dfrac{x^{n+1}(n+1)!}{(n+1)^{n+1}}}{\dfrac{x^n n!}{n^n}}
%   =\frac{x(n+1)}{\left(1+\frac1n\right)^n(n+1)}
%   =\frac{x}{\left(1+\frac1n\right)^n}\to\frac{x}{e}\quad(n\to\infty).
% \]
% 因此由达朗贝尔判别法知, $\sum\limits_{n=1}^{+\infty}\dfrac{x^n n!}{n^n}$ 当 $0\leq x<e$ 时收敛, 当 $x>e$ 时发散.
% 
% 当 $x=e$ 时, 由于
% \[
%   (n+1)^n<e^n n!\quad(n>1),
% \]
% 从而 $\lim\limits_{n\to\infty}a_n\neq0$, 因此该级数发散.
% \end{solve}
% 
% 
% \begin{theorem}[{柯西判别法}]\label{theorem:ma-9-2-4}
% 设 $\sum a_n$ 为一个正项级数, 记
% \[
%   \overline{\lim}_{n\to\infty}\sqrt[n]{a_n}=r,
% \]
% 则
% 
% (1) $r<1$ 时, $\sum a_n$ 收敛;
% 
% (2) $r>1$ 时, $\sum a_n$ 发散.
% \end{theorem}
% 
% 
% \begin{proof}
% (1) 当 $r<1$ 时, 取 $r_1$ 满足 $r<r_1<1$, 则 $\exists N_1$, 当 $n>N_1$ 时, 有 $\sqrt[n]{a_n}<r_1$, 即 $a_n<r_1^n$。而 $\sum r_1^n$ 收敛, 因此 $\sum a_n$ 收敛.
% 
% (2) 当 $r>1$ 时, 取 $r_2$ 满足 $r>r_2>1$, 则对于 $\forall N\in\mathbb{N}$, $\exists n'>N$, 使得 $\sqrt[n']{a_{n'}}>r_2$, 即 $a_{n'}>r_2^{n'}$。这说明 $\lim\limits_{n\to\infty}a_n\neq0$, 因此 $\sum a_n$ 发散.
% \end{proof}
% 
% 
% 对于 $\forall p>0$, 我们有 $\lim\limits_{n\to\infty}\sqrt[n]{\dfrac1{n^p}}=1$, 因此用柯西判别法我们也无法判别正项级数 $\sum\dfrac1{n^p}$ 的敛散性。换句话说, 当定理~\ref{theorem:ma-9-2-4} 中的 $r$ 为 1 时, 柯西判别法也是失效的.
% 
% 从理论上来说, 凡是能用达朗贝尔判别法来判别敛散性的正项级数均可由柯西判别法来加以判别。事实上, 对一般正项级数 $\sum a_n(a_n>0(n=1,2,\cdots))$, 下面的不等式总是成立的:
% \[
%   \underline{\lim}_{n\to\infty}\frac{a_{n+1}}{a_n}
%   \leq \underline{\lim}_{n\to\infty}\sqrt[n]{a_n}
%   \leq \overline{\lim}_{n\to\infty}\sqrt[n]{a_n}
%   \leq \overline{\lim}_{n\to\infty}\frac{a_{n+1}}{a_n}.
% \]
% 上述不等式的证明留给读者.
% 
% \begin{example}\label{example:ma-9-2-7}
% 试讨论正项级数 $\sum\limits_{n=1}^{+\infty}\left(\sqrt[n]{n}-\sin\dfrac1n\right)^{n^2}$ 的敛散性.
% \end{example}
% 
% 
% \begin{solve}
% 记 $a_n=\left(\sqrt[n]{n}-\sin\dfrac1n\right)^{n^2}$。由
% \[
% \begin{aligned}
%   \sqrt[n]{n}-\sin\frac1n
%   &=e^{\frac{\ln n}{n}}-\sin\frac1n\\
%   &=1+\frac{\ln n}{n}+o\left(\frac{\ln n}{n}\right)-\left[\frac1n+o\left(\frac1n\right)\right]\\
%   &=1+\frac{\ln n-1}{n}+o\left(\frac{\ln n}{n}\right)\quad(n\to\infty)
% \end{aligned}
% \]
% 得
% \[
%   \sqrt[n]{a_n}=\left[1+\frac{\ln n-1}{n}+o\left(\frac{\ln n}{n}\right)\right]^n\to+\infty\quad(n\to\infty).
% \]
% 由柯西判别法知 $\sum a_n$ 发散.
% \end{solve}
% 
% 
% \begin{example}\label{example:ma-9-2-8}
% 试讨论正项级数 $\sum C_n$ 的敛散性, 其中
% \[
%   C_n=\begin{cases}
%     a^{(n+1)/2}, & n\text{ 为奇数},\\
%     b^{n/2}, & n\text{ 为偶数}
%   \end{cases}\quad(0<a<b<1).
% \]
% \end{example}
% 
% 
% \begin{solve}
% 由于
% \[
%   \overline{\lim}_{n\to\infty}\sqrt[n]{C_n}\leq\lim_{n\to\infty}\sqrt[n]{b^{n/2}}=b^{1/2}<1,
% \]
% 由柯西判别法知级数 $\sum C_n$ 是收敛的.
% \end{solve}
% 
% 
% \begin{remark}\label{remark:ma-9-source-43}
% 注意到
% \[
%   \overline{\lim}_{n\to\infty}\frac{C_{n+1}}{C_n}
%   =\lim_{n\to\infty}\frac{b^{(n+1)/2}}{a^{(n+1)/2}}=+\infty,
% \]
% 而
% \[
%   \underline{\lim}_{n\to\infty}\frac{C_{n+1}}{C_n}
%   =\lim_{n\to\infty}\frac{a^{(n+2)/2}}{b^{n/2}}=0,
% \]
% 因此达朗贝尔判别法无法判别 $\sum C_n$ 的敛散性。值得注意的是, 在某些情况下, 用达朗贝尔判别法则较为简便, 如前面的例~\ref{example:ma-9-2-5}.
% \end{remark}
% 
% 
% \subsection{拉贝判别法}
% \label{subsec:ma-9-2-3}
% 
% 以下我们考虑以 $\sum a_n=\sum\dfrac1{n^p}(p>0)$ 作为比较标准时的判别法。考查比值 $\dfrac{a_n}{a_{n+1}}$, 读者不难找到规律并得到以下的判别法.
% 
% \begin{theorem}[{拉贝(Raabe)判别法}]\label{theorem:ma-9-2-5}
% 设 $\sum a_n$ 为一正项级数.
% 
% (1) 若 $\lim\limits_{n\to\infty}n\left(\dfrac{a_n}{a_{n+1}}-1\right)=r>1$, 则 $\sum a_n$ 收敛.
% 
% (2) 若 $\overline{\lim}_{n\to\infty}n\left(\dfrac{a_n}{a_{n+1}}-1\right)=r'<1$, 则 $\sum a_n$ 发散.
% \end{theorem}
% 
% 
% \begin{proof}
% (1) 取 $r_1(r>r_1>1)$, 则存在 $N_1\in\mathbb{N}$, 当 $n\geq N_1$ 时, 有
% \[
%   \frac{a_n}{a_{n+1}}>1+\frac{r_1}{n}.
% \]
% 再取 $r_2(r_1>r_2>1)$, 由于 $f(x)=1+r_1x-(1+x)^{r_2}$ 满足 $f(0)=0$, 且 $f'(x)=r_1-r_2(1+x)^{r_2-1}$ 在 $x=0$ 的某个邻域内恒大于 0, 因此存在 $N_2\geq N_1$, 当 $n\geq N_2$ 时, 有
% \[
%   \frac{a_n}{a_{n+1}}>1+\frac{r_1}{n}>\left(1+\frac1n\right)^{r_2}
%   =\frac{(n+1)^{r_2}}{n^{r_2}}.
% \]
% 上式等价于
% \[
%   (n+1)^{r_2}a_{n+1}<n^{r_2}a_n.
% \]
% 因此当 $n>N_2$ 时, 我们有
% \[
%   a_n<\frac{N_2^{r_2}a_{N_2}}{n^{r_2}}=\frac{C}{n^{r_2}}\quad(C=N_2^{r_2}a_{N_2}).
% \]
% 由比较判别法, 我们推知 $\sum a_n$ 是收敛的.
% 
% (2) $r'<1$, 则存在 $N\in\mathbb{N}$, 当 $n\geq N$ 时, 有
% \[
%   \frac{a_n}{a_{n+1}}<1+\frac1n=\frac{n+1}{n}.
% \]
% 因此当 $n\geq N$ 时, 有
% \[
%   na_n<(n+1)a_{n+1},
% \]
% 进一步我们有 $na_n>Na_N$, 即
% \[
%   a_n>\frac{Na_N}{n}=\frac{C}{n}\quad(C=Na_N).
% \]
% 由比较判别法, 我们推知 $\sum a_n$ 是发散的.
% \end{proof}
% 
% 
% \begin{example}\label{example:ma-9-2-9}
% 讨论正项级数 $\sum\limits_{n=1}^{+\infty}\dfrac{(2n-1)!!}{(2n)!!}$ 的敛散性.
% \end{example}
% 
% 
% \begin{solve}
% 记 $a_n=\dfrac{(2n-1)!!}{(2n)!!}$, 则有
% \[
% \begin{aligned}
%   n\left(\frac{a_n}{a_{n+1}}-1\right)
%   &=n\left[\frac{\dfrac{(2n-1)!!}{(2n)!!}}{\dfrac{(2n+1)!!}{(2n+2)!!}}-1\right]\\
%   &=n\left(\frac{2n+2}{2n+1}-1\right)\\
%   &=\frac{n}{2n+1}\to\frac12.
% \end{aligned}
% \]
% 由拉贝判别法知, 级数 $\sum\dfrac{(2n-1)!!}{(2n)!!}$ 发散.
% \end{solve}
% 
% 
% \begin{remark}\label{remark:ma-9-source-48}
% 若对例~\ref{example:ma-9-2-9} 中的级数用达朗贝尔与柯西判别法, 则无法确定该级数的敛散性.
% \end{remark}
% 
% 从定理~\ref{theorem:ma-9-2-5} 的证明中可以看出, 若 (2) 中的条件用以下条件替换其结论也成立: 存在 $N\in\mathbb{N}$, 当 $n>N$ 时, 有
% \[
%   n\left(\frac{a_n}{a_{n+1}}-1\right)\leq1.
% \]
% 
% 细心的读者不难发现, 以上的一些正项级数判别法实际上可以通过几何级数或级数 $\sum\dfrac1{n^p}$ 中的前后两项的比较找出规律而得到。按照这种思想, 读者可以从研究其他的``标准级数''来发现一些规律从而找到一些新的判别法.
% 
% \subsection{柯西积分判别法}
% \label{subsec:ma-9-2-4}
% 
% 在许多正项级数 $\sum a_n$ 中, 序列 $\{a_n\}$ 是单调下降趋于零的。$\{a_n\}$ 作为定义在正整数集 $\mathbb{N}$ 上的一个函数, 它的图像是 $Oxy$ 平面内的一个点集。任取一个 $[1,+\infty)$ 上的单调下降连续函数 $f(x)$, 使得 $f(n)=a_n(n=1,2,\cdots)$(见图 9.2.1).
% 
% 我们可以建立如下级数敛散性与无穷积分敛散性的关系.
% 
% \begin{theorem}\label{theorem:ma-9-2-6}
% 设函数 $f(x)$ 在 $[1,+\infty)$ 上单调下降趋于零, 记 $a_n=f(n)(n=1,2,\cdots)$, 则正项级数 $\sum a_n$ 收敛的充分必要条件是无穷积分 $\int_1^{+\infty}f(x)\,dx$ 收敛.
% \end{theorem}
% 
% 
% \begin{proof}
% 对任何正整数 $n$, 当 $n\leq x\leq n+1$ 时, 有 $f(n+1)\leq f(x)\leq f(n)$, 因此有
% \[
%   a_{n+1}=\int_n^{n+1}f(n+1)\,dx\leq\int_n^{n+1}f(x)\,dx\leq\int_n^{n+1}f(n)\,dx=a_n.
% \]
% 对上面不等式关于 $n$ 求和得
% \[
%   \sum_{n=2}^{+\infty}a_n\leq\int_1^{+\infty}f(x)\,dx
%   =\sum_{n=1}^{+\infty}\int_n^{n+1}f(x)\,dx
%   \leq\sum_{n=1}^{+\infty}a_n.
% \]
% 因此当 $\int_1^{+\infty}f(x)\,dx$ 收敛时, $\sum\limits_{n=2}^{+\infty}a_n$ 有上界 $\int_1^{+\infty}f(x)\,dx$, 从而 $\sum\limits_{n=1}^{+\infty}a_n$ 收敛。反之, 当 $\sum a_n$ 收敛时, 无穷积分 $\int_1^{+\infty}f(x)\,dx$ 有上界 $\sum a_n$, 从而积分收敛.
% \end{proof}
% 
% 
% 在广义积分中, 我们讨论过以下积分的敛散性:
% \[
%   \int_1^{+\infty}\frac{dx}{x^p},\qquad
%   \int_2^{+\infty}\frac{dx}{x^p\ln^q x}.
% \]
% 由定理~\ref{theorem:ma-9-2-6} 知, $\sum\dfrac1{n^p}$ 当 $p>1$ 时收敛, 当 $p\leq1$ 时发散。对于级数 $\sum\limits_{n=2}^{+\infty}\dfrac1{n^p\ln^q n}$, 当 $p>1$ 时, 对任何 $q$, 该级数都收敛; 当 $p=1$ 时, 对 $q>1$ 该级数收敛, 而对其他情形该级数发散.
% 
% 值得注意的是, 对于正项级数
% \[
%   \sum_{n=2}^{+\infty}a_n=\sum_{n=2}^{+\infty}\frac1{n^p\ln^q n},
% \]
% 用柯西判别法及拉贝判别法都无法判别它的敛散性.
% 
% \begin{example}\label{example:ma-9-2-10}
% 设正项级数 $\sum a_n$ 收敛, 记 $r_n=\sum\limits_{k=n}^{+\infty}a_k(n=1,2,\cdots)$, 证明: 当 $p<1$ 时, 正项级数 $\sum\dfrac{a_n}{r_n^p}$ 收敛.
% \end{example}
% 
% 
% \begin{proof}
% 设 $a_0=0$, 并记 $S=\sum a_n$。由
% \[
%   r_n=S-\sum_{k=0}^{n-1}a_k\quad(n=1,2,\cdots)
% \]
% 及 $\lim\limits_{n\to\infty}\sum\limits_{k=0}^{n-1}a_k=S$ 知, $\{r_n\}$ 单调下降趋于零, 因此我们有
% \[
%   S=r_1\geq r_2\geq r_3\geq r_4\geq\cdots,
% \]
% 且 $r_n>0$.
% 
% 注意到对于 $\forall n\in\mathbb{N}$, 有
% \[
%   \frac{a_n}{r_n^p}\leq\int_{r_{n+1}}^{r_n}\frac{dx}{x^p},
% \]
% 从而当 $p<1$ 时, 对于 $\forall n\in\mathbb{N}$, 我们有
% \[
%   \sum_{k=1}^n\frac{a_k}{r_k^p}
%   \leq\sum_{k=1}^n\int_{r_{k+1}}^{r_k}\frac{dx}{x^p}
%   <\int_0^S\frac{dx}{x^p}<+\infty.
% \]
% 因此 $\sum\dfrac{a_n}{r_n^p}$ 收敛.
% 
% 上例告诉我们, 当正项级数 $\sum a_n(a_n>0)$ 收敛时, 我们总能找到另一个收敛的正项级数 $\sum b_n$, 使得 $\lim\limits_{n\to\infty}\dfrac{b_n}{a_n}=+\infty$.
% \end{proof}
% 
% 
% \section{任意项级数}
% \label{sec:ma-9-3}
% 
% \subsection{交错级数的敛散性}
% \label{subsec:ma-9-3-1}
% 
% 对任意项级数 $\sum\limits_{n=1}^{+\infty}a_n$ 而言, 若要讨论它是否绝对收敛时, 我们可以转化为讨论正项级数的敛散性问题。因此对于任意项级数, 我们现在要讨论的问题是如何判定它们的条件收敛性.
% 
% 数项级数 $\sum\limits_{n=1}^{+\infty}a_n$ 收敛的一个必要条件是 $\lim\limits_{n\to\infty}a_n=0$。由于该条件不是充分的, 因此我们还需附加一定的条件才能保证 $\sum\limits_{n=1}^{+\infty}a_n$ 的收敛性。我们有以下的定理.
% 
% \begin{theorem}\label{theorem:ma-9-3-1}
% 设数项级数 $\sum\limits_{n=1}^{+\infty}a_n$ 满足条件:
% 
% (1) $\lim\limits_{n\to\infty}a_n=0$;
% 
% (2) 存在 $N\in\mathbb{N}$, 使得在 $\sum\limits_{n=1}^{+\infty}a_n$ 中加上一些括号, 并且在每个括号中的加数个数 $\leq N$, 得到的级数 $\sum\limits_{k=1}^{+\infty}b_k$ 是收敛的,
% 
% 则级数 $\sum\limits_{n=1}^{+\infty}a_n$ 收敛.
% \end{theorem}
% 
% 
% \begin{proof}
% 记 $S_n=\sum\limits_{k=1}^n a_k(n=1,2,\cdots)$。由于加括号后的级数收敛, 因此存在 $\{S_n\}$ 的一个子列 $\{S_{n_k}\}$ 是收敛的。对任意的正整数 $n>n_1$, 必存在 $n_k$, 使得 $n_k\leq n<n_{k+1}$, 因此有
% \[
%   |S_n-S_{n_k}|\leq\sum_{j=1}^N|a_{n_k+j}|.
% \]
% 注意到 $n\to\infty$ 时必有 $n_k\to\infty$, 由 $\lim\limits_{j\to\infty}\sum\limits_{k=j}^{j+N}|a_k|=0$, 我们推出 $\lim\limits_{n\to\infty}S_n$ 存在, 即 $\sum\limits_{n=1}^{+\infty}a_n$ 收敛.
% \end{proof}
% 
% 
% 下面我们来研究一类特殊的任意项级数——交错级数。设 $\{a_n\}$ 是一个序列, 且对 $\forall n\in\mathbb{N}$, $a_n\geq0$, 我们称级数 $\sum\limits_{n=1}^{+\infty}(-1)^{n-1}a_n$ 或 $\sum\limits_{n=1}^{+\infty}(-1)^na_n$ 为一个交错级数。由定理~\ref{theorem:ma-9-3-1}, 可推出下面重要的结论.
% 
% \begin{corollary}[{莱布尼茨交错级数判别法}]\label{corollary:ma-9-source-55}
% 设 $\{a_n\}$ 是一个单调序列, 并且 $\lim\limits_{n\to\infty}a_n=0$, 则交错级数 $\sum\limits_{n=1}^{+\infty}(-1)^{n-1}a_n$ 收敛.
% \end{corollary}
% 
% 
% \begin{proof}
% 不妨设序列 $\{a_n\}$ 单调下降趋于零, 从而必有 $a_n\geq0(n=1,2,\cdots)$。显然 $\sum\limits_{k=1}^{+\infty}(a_{2k-1}-a_{2k})$ 是在 $\sum\limits_{n=1}^{+\infty}(-1)^{n-1}a_n$ 中加上括号后得到的正项级数。记该正项级数的前 $n$ 项部分和为 $S_n$, 我们有
% \[
% \begin{aligned}
%   0\leq S_n
%   &=(a_1-a_2)+(a_3-a_4)+\cdots+(a_{2n-1}-a_{2n})\\
%   &\leq a_1-(a_2-a_3)-(a_4-a_5)-\cdots-(a_{2n-2}-a_{2n-1})-a_{2n}\\
%   &\leq a_1,
% \end{aligned}
% \]
% 因此 $\sum\limits_{k=1}^{+\infty}(a_{2k-1}-a_{2k})$ 收敛。由定理~\ref{theorem:ma-9-3-1} 知 $\sum\limits_{n=1}^{+\infty}(-1)^{n-1}a_n$ 也收敛.
% \end{proof}
% 
% 
% \begin{remark}\label{remark:ma-9-source-57}
% 对上述定理的证明略加分析, 我们可以得到以下结论: 设 $\{a_n\}$ 是单调下降趋于零的序列, 则对任意的正整数 $N$, 以下不等式成立:
% \[
%   0\leq\left|\sum_{n=N}^{+\infty}(-1)^{n-1}a_n\right|\leq a_N.
% \]
% 另外, 我们今后将收敛的交错级数称为莱布尼茨交错级数.
% \end{remark}
% 
% 
% \begin{example}\label{example:ma-9-3-1}
% 试讨论交错级数 $\sum\limits_{n=2}^{+\infty}\dfrac{(-1)^n}{n^p\ln^q n}(p,q\in\mathbb{R})$ 的敛散性.
% \end{example}
% 
% 
% \begin{solve}
% 从上节知, 当 $p>1$ 时, 对于 $\forall q\in\mathbb{R}$, 该级数绝对收敛; 当 $p=1,q>1$ 时, 该级数也绝对收敛.
% 
% 当 $0<p<1$ 时, 对于 $\forall q\in\mathbb{R}$, 我们有 $\lim\limits_{n\to\infty}\dfrac1{n^p\ln^q n}=0$。注意到当 $x$ 充分大时, 有
% \[
% \begin{aligned}
%   (x^{-p}\ln^{-q}x)'
%   &=-px^{-p-1}\ln^{-q}x-x^{-p}(-q)(\ln^{-q-1}x)\frac1x\\
%   &=x^{-p-1}\ln^{-q}x\left(-p+\frac{q}{\ln x}\right)<0,
% \end{aligned}
% \]
% 我们推出, 存在 $N\in\mathbb{N}$, 当 $n>N$ 时, $\left\{\dfrac1{n^p\ln^q n}\right\}$ 单调下降。因此当 $0<p<1$ 时, 对于 $\forall q\in\mathbb{R}$, $\sum\limits_{n=2}^{+\infty}\dfrac{(-1)^n}{n^p\ln^q n}$ 收敛。显然此时级数 $\sum\limits_{n=2}^{+\infty}\dfrac1{n^p\ln^q n}$ 发散, 从而原级数条件收敛.
% 
% 当 $p=0$ 时,
% \[
%   \sum_{n=2}^{+\infty}\frac{(-1)^n}{n^p\ln^q n}
%   =\sum_{n=1}^{+\infty}\frac{(-1)^n}{\ln^q n},
% \]
% 显然该级数当 $q>0$ 时条件收敛, 当 $q\leq0$ 发散.
% 
% 当 $p<0$ 时, 由于对于 $\forall q\in\mathbb{R}$, 有 $\lim\limits_{n\to\infty}\dfrac1{n^p\ln^q n}=\infty$, 因此 $\sum\limits_{n=2}^{+\infty}\dfrac{(-1)^n}{n^p\ln^q n}$ 发散.
% \end{solve}
% 
% 
% \subsection{狄利克雷判别法和阿贝尔判别法}
% \label{subsec:ma-9-3-2}
% 
% 如同广义积分一样, 当数项级数具有形式 $\sum\limits_{n=1}^{+\infty}a_nb_n$ 时, 在一些条件下我们也有相应的狄利克雷判别法和阿贝尔判别法.
% 
% \begin{theorem}[{狄利克雷判别法}]\label{theorem:ma-9-3-2}
% 设数项级数 $\sum\limits_{n=1}^{+\infty}a_n$ 的部分和序列 $\left\{S_n=\sum\limits_{k=1}^n a_k\right\}$ 是有界的, $\{b_n\}$ 是单调序列且 $\lim\limits_{n\to\infty}b_n=0$, 则数项级数 $\sum\limits_{n=1}^{+\infty}a_nb_n$ 收敛.
% \end{theorem}
% 
% 
% \begin{proof}
% 由 $\{S_n\}$ 有界, 存在 $M>0$, 使得
% \[
%   |S_n|\leq\frac{M}{2}\quad(n=1,2,\cdots),
% \]
% 从而对任意的正整数 $n,p$, 有
% \[
%   |S_{n+p}-S_n|\leq M.
% \]
% 对于 $\forall\varepsilon>0$, 由 $\lim\limits_{n\to\infty}b_n=0$, 存在 $N\in\mathbb{N}$, 当 $n>N$ 时, 有
% \[
%   |b_n|<\frac{\varepsilon}{6M}.
% \]
% 因此当 $n>N$ 时, 对任意正整数 $p$, 利用阿贝尔变换(引理~\ref{lemma:ma-7-5-4})得
% \[
% \begin{aligned}
%   \left|\sum_{k=n+1}^{n+p}a_kb_k\right|
%   &=\left|b_{n+p}(S_{n+p}-S_n)-\sum_{k=n+1}^{n+p-1}(S_k-S_n)(b_{k+1}-b_k)\right|\\
%   &\leq M\left\{|b_{n+p}|+\sum_{k=n+1}^{n+p-1}|b_{k+1}-b_k|\right\}\\
%   &=M\left\{|b_{n+p}|+\left|\sum_{k=n+1}^{n+p-1}(b_{k+1}-b_k)\right|\right\}\\
%   &\leq M\{|b_{n+1}|+2|b_{n+p}|\}<\varepsilon.
% \end{aligned}
% \]
% 由柯西准则知 $\sum\limits_{n=1}^{+\infty}a_nb_n$ 收敛.
% \end{proof}
% 
% 
% \begin{theorem}[{阿贝尔判别法}]\label{theorem:ma-9-3-3}
% 设数项级数 $\sum\limits_{n=1}^{+\infty}a_n$ 收敛, 序列 $\{b_n\}$ 单调有界, 则数项级数 $\sum\limits_{n=1}^{+\infty}a_nb_n$ 收敛.
% \end{theorem}
% 
% 
% \begin{proof}
% 由于 $\{b_n\}$ 是单调有界序列, 从而收敛。设 $\lim\limits_{n\to\infty}b_n=b$, 我们有
% \[
%   \sum_{n=1}^{+\infty}a_nb_n
%   =\sum_{n=1}^{+\infty}(b_n-b)a_n+\sum_{n=1}^{+\infty}ba_n.
% \]
% 显然级数 $\sum\limits_{n=1}^{+\infty}ba_n$ 是收敛的, 而在级数 $\sum\limits_{n=1}^{+\infty}(b_n-b)a_n$ 中令 $b'_n=b_n-b(n=1,2,\cdots)$, 则 $\{b'_n\}$ 及 $\sum\limits_{n=1}^{+\infty}a_n$ 满足狄利克雷判别法(定理~\ref{theorem:ma-9-3-2})的条件, 因此 $\sum\limits_{n=1}^{+\infty}(b_n-b)a_n$ 收敛, 所以 $\sum\limits_{n=1}^{+\infty}a_nb_n$ 收敛.
% \end{proof}
% 
% 
% \begin{remark}\label{remark:ma-9-source-64}
% 容易看出来布尼茨交错级数判别法可以由狄利克雷判别法直接推出.
% \end{remark}
% 
% 
% \begin{example}\label{example:ma-9-3-2}
% 试讨论交错级数 $\sum\limits_{n=1}^{+\infty}\dfrac{(-1)^n}{n^{\alpha+\frac1n}}(\alpha>0)$ 的敛散性.
% \end{example}
% 
% 
% \begin{solve}
% 显然, 当 $\alpha>1$ 时, 该级数绝对收敛。当 $0<\alpha\leq1$ 时, 对 $n\in\mathbb{N}$, 记 $a_n=\dfrac{(-1)^n}{n^\alpha}$ 和 $b_n=\dfrac1{n^{\frac1n}}$。由莱布尼茨交错级数判别法知 $\sum\limits_{n=1}^{+\infty}a_n$ 收敛。又记 $n^{\frac1n}=e^{\frac{\ln n}{n}}$, 则容易证明: 当 $n$ 充分大时, 序列 $\left\{\dfrac{\ln n}{n}\right\}$ 是单调下降的。因此序列 $\left\{\dfrac1{n^{\frac1n}}\right\}$ 当 $n$ 充分大时单调上升, 且有上界 1。于是由阿贝尔判别法知 $\sum\limits_{n=1}^{+\infty}\dfrac{(-1)^n}{n^{\alpha+\frac1n}}$ 收敛。由于
% \[
%   \lim_{n\to\infty}\frac{\dfrac1{n^{\alpha+\frac1n}}}{\dfrac1{n^\alpha}}=1,
% \]
% 而当 $0<\alpha\leq1$ 时, $\sum\limits_{n=1}^{+\infty}\dfrac1{n^\alpha}$ 发散, 从而 $\sum\limits_{n=1}^{+\infty}\dfrac1{n^{\alpha+\frac1n}}$ 发散, 故此时原级数是条件收敛的.
% \end{solve}
% 
% 
% \begin{example}\label{example:ma-9-3-3}
% 设序列 $\{a_n\}$ 单调且 $\lim\limits_{n\to\infty}a_n=0$, 证明:
% 
% (1) 当 $\sum\limits_{n=1}^{+\infty}a_n$ 收敛时, $\sum\limits_{n=1}^{+\infty}a_n\sin nx$ 与 $\sum\limits_{n=1}^{+\infty}a_n\cos nx$ 对一切 $x\in\mathbb{R}$ 都绝对收敛;
% 
% (2) 当 $\sum\limits_{n=1}^{+\infty}a_n$ 发散时, $\sum\limits_{n=1}^{+\infty}a_n\sin nx$ 对一切 $x\neq k\pi(k\in\mathbb{Z})$ 条件收敛, 而 $\sum\limits_{n=1}^{+\infty}a_n\cos nx$ 对一切 $x\neq2k\pi(k\in\mathbb{Z})$ 条件收敛.
% \end{example}
% 
% 
% \begin{proof}
% (1) 对于 $\forall n\in\mathbb{N}$, 我们不妨设 $a_n>0$, 此时显然有
% \[
%   |a_n\sin nx|\leq a_n,\qquad |a_n\cos nx|\leq a_n.
% \]
% 因此当 $\sum\limits_{n=1}^{+\infty}a_n$ 收敛时, $\sum\limits_{n=1}^{+\infty}a_n\sin nx$ 与 $\sum\limits_{n=1}^{+\infty}a_n\cos nx$ 对一切 $x\in\mathbb{R}$ 都绝对收敛.
% 
% (2) 我们先证两个级数在给定的条件下是收敛的。由狄利克雷判别法, 只要证在所给条件下, 对于 $\forall n\in\mathbb{N}$,
% \[
%   S_n=\sum_{k=1}^n\sin kx,\qquad S'_n=\sum_{k=1}^n\cos kx,
% \]
% 有界即可。事实上, 由
% \[
% \begin{aligned}
%   2\sin\frac{x}{2}\cdot S_n
%   &=\sum_{k=1}^n\left[\cos\left(k-\frac12\right)x-\cos\left(k+\frac12\right)x\right]\\
%   &=\cos\frac{x}{2}-\cos\left(n+\frac12\right)x,
% \end{aligned}
% \]
% 当 $x\neq k\pi(k\in\mathbb{Z})$ 时(当 $x=k\pi$ 时, 该级数的每一项均为 0), 对任意的 $n\in\mathbb{N}$, 有
% \[
%   |S_n|\leq\frac1{\left|\sin\dfrac{x}{2}\right|}.
% \]
% 而对于 $\forall n\in\mathbb{N}$, 由
% \[
% \begin{aligned}
%   2\sin\frac{x}{2}\cdot S'_n
%   &=-\sum_{k=1}^n\left[\sin\left(k-\frac12\right)x-\sin\left(k+\frac12\right)x\right]\\
%   &=\sin\left(n+\frac12\right)x-\sin\frac{x}{2},
% \end{aligned}
% \]
% 当 $x\neq2k\pi(k\in\mathbb{Z})$ 时, 有
% \[
%   |S'_n|\leq\frac1{\left|\sin\dfrac{x}{2}\right|}.
% \]
% 因此由狄利克雷判别法知, $\sum\limits_{n=1}^{+\infty}a_n\sin nx$ 对一切 $x\in\mathbb{R}$ 都收敛, 而 $\sum\limits_{n=1}^{+\infty}a_n\cos nx$ 对 $x\neq2k\pi(k\in\mathbb{Z})$ 收敛.
% 
% 当 $\sum\limits_{n=1}^{+\infty}a_n$ 发散时, 注意到当 $x=k\pi(k\in\mathbb{Z})$ 时, $\sum\limits_{n=1}^{+\infty}a_n\sin nx=\sum\limits_{n=1}^{+\infty}0$, 下面假定 $x\neq k\pi(k\in\mathbb{Z})$, 我们证明级数 $\sum\limits_{n=1}^{+\infty}a_n\sin nx$ 条件收敛。事实上, 由于对于任意的 $n\in\mathbb{N}$, 有
% \[
%   |a_n\sin nx|\geq |a_n\sin^2 nx|=\frac{a_n}{2}(1-\cos2nx),
% \]
% 而 $\sum\limits_{n=1}^{+\infty}a_n\cos2nx$ 收敛, $\sum\limits_{n=1}^{+\infty}a_n=+\infty$, 因此 $\sum\limits_{n=1}^{+\infty}|a_n\sin nx|=+\infty$, 从而 $\sum\limits_{n=1}^{+\infty}a_n\sin nx$ 条件收敛.
% 
% 同理, 当 $x\neq2k\pi(k\in\mathbb{Z})$ 且 $\sum\limits_{n=1}^{+\infty}a_n$ 发散时, $\sum\limits_{n=1}^{+\infty}a_n\cos nx$ 条件收敛.
% \end{proof}
% 
% 
% \section{数项级数的性质}
% \label{sec:ma-9-4}
% 
% 由于数项级数是由一列数相加而成, 人们自然要问: 对级数而言加法结合律、交换律是否还成立? 当讨论两个数项级数相乘的问题时, 我们还会遇到分配律的问题。在本节中我们将仔细研究这些问题.
% 
% \subsection{结合律}
% \label{subsec:ma-9-4-1}
% 
% 设 $\sum\limits_{n=1}^{+\infty}a_n$ 是一个数项级数, 对该级数加上一些括号并将每个括号中的加数求和后作为一项将得到一个新的级数 $\sum\limits_{k=1}^{+\infty}b_k$。关于数项级数的结合律, 我们有以下的结论.
% 
% \begin{theorem}\label{theorem:ma-9-4-1}
% 设 $\sum\limits_{n=1}^{+\infty}a_n$ 是收敛级数, 则在 $\sum\limits_{n=1}^{+\infty}a_n$ 中任意加括号后得到的级数 $\sum\limits_{k=1}^{+\infty}b_k$ 也收敛, 并且 $\sum\limits_{n=1}^{+\infty}a_n=\sum\limits_{k=1}^{+\infty}b_k$.
% \end{theorem}
% 
% 
% \begin{proof}
% 对于 $\forall n\in\mathbb{N}$, 设 $S_n$ 是 $\sum\limits_{n=1}^{+\infty}a_n$ 前 $n$ 项部分和; 对于 $\forall k\in\mathbb{N}$, 设 $S'_k$ 是 $\sum\limits_{k=1}^{+\infty}b_k$ 的前 $k$ 项的部分和。容易看出, $\{S'_k\}$ 是 $\{S_n\}$ 的一个子列, 因此当 $\lim\limits_{n\to\infty}S_n$ 存在时, $\lim\limits_{k\to\infty}S'_k$ 也存在, 并且 $\lim\limits_{k\to\infty}S'_k=\lim\limits_{n\to\infty}S_n$.
% \end{proof}
% 
% 
% 数项级数 $\sum\limits_{n=1}^{+\infty}(-1)^n$ 发散告诉我们, 级数加上括号后得到的级数收敛并不能推出原来的级数收敛.
% 
% 对于正项级数 $\sum\limits_{n=1}^{+\infty}a_n$, 我们容易看出, 如果在 $\sum\limits_{n=1}^{+\infty}a_n$ 中加上括号后得到的级数是收敛的, 则 $\sum\limits_{n=1}^{+\infty}a_n$ 也收敛。对于一般的级数来说, 如果在 $\sum\limits_{n=1}^{+\infty}a_n$ 中加上括号, 并且在每个括号中的项都具有相同的符号, 将这样得到的级数记为 $\sum\limits_{k=1}^{+\infty}b_k$, 则 $\sum\limits_{n=1}^{+\infty}a_n$ 收敛的充分必要条件是 $\sum\limits_{k=1}^{+\infty}b_k$ 收敛, 并且当它们收敛时它们的和也相同(见本章习题).
% 
% \subsection{交换律}
% \label{subsec:ma-9-4-2}
% 
% 交换律对无穷多个数相加是否还成立呢? 为此我们要讨论数项级数 $\sum\limits_{n=1}^{+\infty}a_n$ 与它的各项重排后所得的级数 $\sum\limits_{n=1}^{+\infty}a'_n$ 的关系。为了精确描述级数的重排, 我们首先来定义正整数集 $\mathbb{N}$ 的重排: 若函数 $f(n)$ 是 $\mathbb{N}$ 到 $\mathbb{N}$ 的一个一一对应, 则称 $f(n)(n\in\mathbb{N})$ 是 $\mathbb{N}$ 的一个重排.
% 
% 现设 $\sum\limits_{n=1}^{+\infty}a_n$ 是一个数项级数。我们称 $\sum\limits_{n=1}^{+\infty}a'_n$ 是它的一个重排, 如果存在 $\mathbb{N}$ 的一个重排 $f(n)$, 使得 $a'_n=a_{f(n)}(n=1,2,\cdots)$.
% 
% 对于一般的数项级数来说, 我们有以下的结果.
% 
% \begin{theorem}\label{theorem:ma-9-4-2}
% 设 $\sum\limits_{n=1}^{+\infty}a_n$ 为一个数项级数, $f(n):\mathbb{N}\to\mathbb{N}$ 为一个重排, 再设 $\exists M>0$, 使得对于 $\forall n\in\mathbb{N}$, 有 $|f(n)-n|\leq M$, 则级数 $\sum\limits_{n=1}^{+\infty}a_n$ 收敛的充分必要条件是级数 $\sum\limits_{n=1}^{+\infty}a_{f(n)}$ 收敛, 并且当它们收敛时, 有 $\sum\limits_{n=1}^{+\infty}a_n=\sum\limits_{n=1}^{+\infty}a_{f(n)}$.
% \end{theorem}
% 
% 
% \begin{proof}
% 必要性\quad 设 $\sum\limits_{n=1}^{+\infty}a_n$ 收敛, 下证 $\sum\limits_{n=1}^{+\infty}a_{f(n)}$ 也收敛。由于 $\sum\limits_{n=1}^{+\infty}a_n$ 收敛, 我们有 $\lim\limits_{n\to\infty}a_n=0$。又因对于 $\forall n\in\mathbb{N}$, 有 $|f(n)-n|\leq M$, 从而有
% \[
%   \lim_{n\to\infty}\sum_{j=-M}^{M}|a_{n+j}|=0.
% \]
% 设 $\{S_n\}$ 和 $\{S'_n\}$ 分别为 $\sum\limits_{n=1}^{+\infty}a_n$ 和 $\sum\limits_{n=1}^{+\infty}a_{f(n)}$ 的部分和序列, 容易看出
% \[
%   |S_n-S'_n|\leq\sum_{j=-M}^{M}|a_{n+j}|\to0\quad(n\to\infty).
% \]
% 因此 $\lim\limits_{n\to\infty}S'_n=\lim\limits_{n\to\infty}S_n$, 即 $\sum\limits_{n=1}^{+\infty}a_{f(n)}$ 收敛, 且 $\sum\limits_{n=1}^{+\infty}a_n=\sum\limits_{n=1}^{+\infty}a_{f(n)}$.
% 
% 充分性\quad 设 $\sum\limits_{n=1}^{+\infty}a_{f(n)}$ 收敛, 显然 $\sum\limits_{n=1}^{+\infty}a_n$ 是 $\sum\limits_{n=1}^{+\infty}a_{f(n)}$ 的一个重排, 其重排函数为 $f^{-1}(n):\mathbb{N}\to\mathbb{N}$。容易看出, 对于 $\forall n\in\mathbb{N}$, 有
% \[
%   |f^{-1}(f(n))-f(n)|=|n-f(n)|\leq M.
% \]
% 由必要性的证明知 $\sum\limits_{n=1}^{+\infty}a_n$ 收敛, 且 $\sum\limits_{n=1}^{+\infty}a_n=\sum\limits_{n=1}^{+\infty}a_{f(n)}$.
% \end{proof}
% 
% 
% \begin{remark}\label{remark:ma-9-source-73}
% 定理~\ref{theorem:ma-9-4-2} 告诉我们: 若一个级数交换各项的位置, 只要每项离原来所在的位置不要太远, 则不会改变级数的敛散性.
% \end{remark}
% 
% 对于绝对收敛的级数, 我们有下面满意的结果:
% 
% \begin{theorem}\label{theorem:ma-9-4-3}
% 设 $\sum\limits_{n=1}^{+\infty}a_n$ 是绝对收敛的, 则它的任意一个重排 $\sum\limits_{n=1}^{+\infty}a_{f(n)}$ 也绝对收敛, 并且 $\sum\limits_{n=1}^{+\infty}a_n=\sum\limits_{n=1}^{+\infty}a_{f(n)}$.
% \end{theorem}
% 
% 
% \begin{proof}
% 先证 $\sum\limits_{n=1}^{+\infty}|a_{f(n)}|$ 收敛并且 $\sum\limits_{n=1}^{+\infty}|a_{f(n)}|=\sum\limits_{n=1}^{+\infty}|a_n|$。记 $\{S'_n\}$ 为正项级数 $\sum\limits_{n=1}^{+\infty}|a_{f(n)}|$ 的部分和序列, 显然, 对于 $\forall n\in\mathbb{N}$, 有
% \[
%   S'_n\leq\sum_{n=1}^{+\infty}|a_n|.
% \]
% 这说明了 $\sum\limits_{n=1}^{+\infty}|a_{f(n)}|$ 是收敛的, 即 $\sum\limits_{n=1}^{+\infty}a_{f(n)}$ 是绝对收敛的。再令 $n\to\infty$, 得
% \[
%   \sum_{n=1}^{+\infty}|a_{f(n)}|\leq\sum_{n=1}^{+\infty}|a_n|.
% \]
% 由于 $\sum\limits_{n=1}^{+\infty}a_n$ 是 $\sum\limits_{n=1}^{+\infty}a_{f(n)}$ 的一个重排, 我们同样有
% \[
%   \sum_{n=1}^{+\infty}|a_n|\leq\sum_{n=1}^{+\infty}|a_{f(n)}|.
% \]
% 因此有
% \[
%   \sum_{n=1}^{+\infty}|a_n|=\sum_{n=1}^{+\infty}|a_{f(n)}|.
% \]
% 为了证明 $\sum\limits_{n=1}^{+\infty}a_n=\sum\limits_{n=1}^{+\infty}a_{f(n)}$, 对于任意的 $n\in\mathbb{N}$, 我们定义
% \begin{equation}
%   a_n^+=
%   \begin{cases}
%     a_n, & a_n>0,\\
%     0, & a_n\leq0,
%   \end{cases}
%   \qquad
%   a_n^-=
%   \begin{cases}
%     -a_n, & a_n<0,\\
%     0, & a_n\geq0.
%   \end{cases}
%   \label{eq:ma-9-4-1}
% \end{equation}
% 显然对于 $\forall n\in\mathbb{N}$, 有 $a_n^+\leq|a_n|$, $a_n^-\leq|a_n|$。因此 $\sum\limits_{n=1}^{+\infty}a_n^+$ 与 $\sum\limits_{n=1}^{+\infty}a_n^-$ 均为收敛的正项级数, 且有 $\sum\limits_{n=1}^{+\infty}a_n^+=\sum\limits_{n=1}^{+\infty}a_{f(n)}^+$ 和 $\sum\limits_{n=1}^{+\infty}a_n^-=\sum\limits_{n=1}^{+\infty}a_{f(n)}^-$。由此得
% \[
%   \sum_{n=1}^{+\infty}a_n
%   =\sum_{n=1}^{+\infty}a_n^+-\sum_{n=1}^{+\infty}a_n^-
%   =\sum_{n=1}^{+\infty}a_{f(n)}^+-\sum_{n=1}^{+\infty}a_{f(n)}^-
%   =\sum_{n=1}^{+\infty}a_{f(n)}.
% \qedhere
% \]
% \end{proof}
% 
% 
% 对于条件收敛的级数, 则上述定理可能不成立。我们来考查以下例子。我们知道数项级数
% \[
%   \sum_{n=1}^{+\infty}\frac{(-1)^{n-1}}{n}=1-\frac12+\frac13-\frac14+\cdots
% \]
% 是条件收敛的级数。利用 $1+\dfrac12+\cdots+\dfrac1n=\ln n+c+\varepsilon_n$, 其中 $c$ 为欧拉常数, $\varepsilon_n\to0(n\to\infty)$, 我们得到级数 $\sum\limits_{n=1}^{+\infty}\dfrac{(-1)^{n-1}}{n}$ 前 $2n$ 项的部分和
% \[
% \begin{aligned}
%   S_{2n}
%   &=1-\frac12+\frac13-\cdots-\frac1{2n}\\
%   &=\left(1+\frac12+\frac13+\cdots+\frac1{2n}\right)-2\left(\frac12+\frac14+\cdots+\frac1{2n}\right)\\
%   &=\ln2n+\varepsilon_{2n}+c-(\ln n+c+\varepsilon_n)\\
%   &=\ln2+\varepsilon_{2n}-\varepsilon_n\to\ln2\quad(n\to\infty).
% \end{aligned}
% \]
% 再由级数 $\sum\limits_{n=1}^{+\infty}\dfrac{(-1)^{n-1}}{n}$ 前 $2n+1$ 项的部分和
% \[
%   S_{2n+1}=S_{2n}+\frac1{2n+1}\to\ln2\quad(n\to\infty),
% \]
% 知 $\sum\limits_{n=1}^{+\infty}\dfrac{(-1)^{n-1}}{n}$ 收敛到 $\ln2$.
% 
% 我们将此级数重排, 使得每两个正项后面跟一个负项, 即
% \[
%   1+\frac13-\frac12+\frac15+\frac17-\frac14+\cdots+\frac1{4n-3}+\frac1{4n-1}-\frac1{2n}+\cdots.
% \]
% 对于 $n\in\mathbb{N}$, 设它的前 $n$ 项部分和 $S'_n$, 则
% \begin{equation}
% \begin{aligned}
%   S'_{3n}
%   &=\sum_{k=1}^n\left(\frac1{4k-3}+\frac1{4k-1}-\frac1{2k}\right)\\
%   &=\sum_{k=1}^{4n}\frac1k-\sum_{k=1}^n\left(\frac1{4k-2}+\frac1{4k}+\frac1{2k}\right).
% \end{aligned}
%   \label{eq:ma-9-4-2}
% \end{equation}
% 而
% \[
%   \sum_{k=1}^n\left(\frac1{4k-2}+\frac1{4k}+\frac1{2k}\right)
%   =\frac12\sum_{k=1}^{2n}\frac1k+\frac12\sum_{k=1}^n\frac1k,
% \]
% 将它与 $1+\dfrac12+\cdots+\dfrac1n=\ln n+c+\varepsilon_n$ 代入 \eqref{eq:ma-9-4-2}, 得
% \[
%   S'_{3n}
%   =\ln4n+c+\varepsilon_{4n}
%   -\frac12(\ln2n+c+\varepsilon_{2n})
%   -\frac12(\ln n+c+\varepsilon_n)\to\frac32\ln2\quad(n\to\infty).
% \]
% 再由
% \[
%   S'_{3n+1}=S'_{3n}+\frac1{4n+1}\to\frac32\ln2\quad(n\to\infty),
% \]
% \[
%   S'_{3n+2}=S'_{3n}+\frac1{4n+1}+\frac1{4n+3}\to\frac32\ln2\quad(n\to\infty),
% \]
% 我们推知重排后的级数收敛到 $\dfrac32\ln2$。这说明了对于条件收敛的级数, 重排有可能改变级数的和.
% 
% 对于条件收敛的级数, 我们有下面的一般性结果.
% 
% \begin{theorem}[{黎曼定理}]\label{theorem:ma-9-4-4}
% 设数项级数 $\sum\limits_{n=1}^{+\infty}a_n$ 条件收敛, 则对任意的 $S(S$ 为有限实数或者 $\pm\infty)$, 都存在 $\sum\limits_{n=1}^{+\infty}a_n$ 的重排 $\sum\limits_{n=1}^{+\infty}a_{f(n)}$, 使得
% \[
%   \sum_{n=1}^{+\infty}a_{f(n)}=S.
% \]
% \end{theorem}
% 
% 
% \begin{proof}
% 对于条件收敛级数 $\sum\limits_{n=1}^{+\infty}a_n$, 我们有 $\lim\limits_{n\to\infty}a_n=0$ 和
% \[
%   \sum_{n=1}^{+\infty}a_n^+=+\infty,\qquad \sum_{n=1}^{+\infty}a_n^-=+\infty,
% \]
% 其中 $a_n^+$, $a_n^-(n=1,2,\cdots)$ 由 \eqref{eq:ma-9-4-1} 定义.
% 
% 先设 $0\leq S<+\infty$, 因此存在正整数 $p_1$, 使得
% \[
%   S<\sum_{k=1}^{p_1}a_k^+\leq S+a_{p_1}^+,
% \]
% 同理存在正整数 $q_1$, 使得
% \[
%   S-a_{q_1}^-\leq\sum_{k=1}^{p_1}a_k^+-\sum_{k=1}^{q_1}a_k^-<S.
% \]
% 由归纳法, 我们可以得到 $\sum\limits_{n=1}^{+\infty}a_n$ 的一个重排
% \[
% \begin{aligned}
%   &(a_1^++\cdots+a_{p_1}^+)-(a_1^-+\cdots+a_{q_1}^-)\\
%   &\quad +(a_{p_1+1}^++\cdots+a_{p_2}^+)-(a_{q_1+1}^-+\cdots+a_{q_2}^-)+\cdots\\
%   &\quad +(a_{p_{k-1}+1}^++\cdots+a_{p_k}^+)-(a_{q_{k-1}+1}^-+\cdots+a_{q_k}^-)+\cdots,
% \end{aligned}
% \]
% 使得它满足
% \[
% \begin{aligned}
%   S
%   &< (a_1^++\cdots+a_{p_1}^+)-(a_1^-+\cdots+a_{q_1}^-)\\
%   &\quad +(a_{p_1+1}^++\cdots+a_{p_2}^+)-(a_{q_1+1}^-+\cdots+a_{q_2}^-)\\
%   &\quad +\cdots+(a_{p_{k-1}+1}^++\cdots+a_{p_k}^+)\\
%   &\leq S+a_{p_k}^+
% \end{aligned}
% \]
% 和
% \[
% \begin{aligned}
%   S-a_{q_k}^-
%   &\leq (a_1^++\cdots+a_{p_1}^+)-(a_1^-+\cdots+a_{q_1}^-)\\
%   &\quad +(a_{p_1+1}^++\cdots+a_{p_2}^+)-(a_{q_1+1}^-+\cdots+a_{q_2}^-)\\
%   &\quad +\cdots+(a_{p_{k-1}+1}^++\cdots+a_{p_k}^+)-(a_{q_{k-1}+1}^-+\cdots+a_{q_k}^-)\\
%   &<S.
% \end{aligned}
% \]
% 由 $q_k\geq k$, $p_k\geq k$, 知
% \[
%   \lim_{k\to\infty}a_{p_k}^+=\lim_{k\to\infty}a_{q_k}^-=0.
% \]
% 因此去掉括号得到重排后的级数也收敛到 $S$(见本章习题 13).
% 
% 当 $S<0$ 时, 由于 $\sum\limits_{n=1}^{+\infty}(-a_n)$ 也条件收敛, 因此存在重排 $\sum\limits_{n=1}^{+\infty}(-a_{f(n)})$ 收敛于 $-S$, 从而 $\sum\limits_{n=1}^{+\infty}a_{f(n)}$ 收敛于 $S$.
% 
% 当 $S=+\infty$ 时, 存在正整数 $p_1$, 使得
% \[
%   \sum_{k=1}^{p_1}a_k^+>1+a_1^-,
% \]
% 同样存在正整数 $p_2$, 使得
% \[
%   \sum_{k=1}^{p_1}a_k^++\sum_{k=p_1+1}^{p_2}a_k^+>2+a_1^-+a_2^-.
% \]
% 继续下去, 由于 $\sum\limits_{n=1}^{+\infty}a_n^+=+\infty$, 我们可取正整数 $p_3,p_4,\cdots,p_k$, 使得
% \[
%   \sum_{k=1}^{p_1}a_k^++\sum_{k=p_1+1}^{p_2}a_k^++\cdots+\sum_{k=p_{k-1}+1}^{p_k}a_k^+
%   > k+a_1^-+a_2^-+\cdots+a_k^-,
% \]
% 因此我们得到级数的重排
% \[
%   a_1^++\cdots+a_{p_1}^+-a_1^-+a_{p_1+1}^++\cdots+a_{p_2}^+-a_2^-+\cdots
%   +a_{p_{k-1}+1}^++\cdots+a_{p_k}^+-a_k^-+\cdots.
% \]
% 不难看出该级数必发散到 $+\infty$.
% 
% 当 $S=-\infty$ 时, 利用上述结果知, $\sum\limits_{n=1}^{+\infty}(-a_n)$ 存在重排, 使得它发散到 $+\infty$, 因此 $\sum\limits_{n=1}^{+\infty}a_n$ 存在重排, 使得它发散到 $-\infty$.
% \end{proof}
% 
% 
% \subsection{级数的乘法（分配律）}
% \label{subsec:ma-9-4-3}
% 
% 现在我们来讨论级数的乘法问题。设 $\sum\limits_{n=1}^{+\infty}a_n$ 与 $\sum\limits_{n=1}^{+\infty}b_n$ 为两个数项级数, 类似于有限项求和的乘法规则, 作出两级数各项所有可能的乘积 $a_kb_j(k=1,2,\cdots,j=1,2,\cdots)$, 将它们排成以下的无穷矩阵(称为这两个级数的乘积矩阵):
% \[
%   \begin{pmatrix}
%     a_1b_1 & a_1b_2 & \cdots & a_1b_n & \cdots\\
%     a_2b_1 & a_2b_2 & \cdots & a_2b_n & \cdots\\
%     \cdots & \cdots & \cdots & \cdots & \cdots\\
%     a_nb_1 & a_nb_2 & \cdots & a_nb_n & \cdots\\
%     \cdots & \cdots & \cdots & \cdots & \cdots
%   \end{pmatrix}
% \]
% 任何将该矩阵中的元素排成一个序列的方法都将得到 $\sum\limits_{n=1}^{+\infty}a_n$ 与 $\sum\limits_{n=1}^{+\infty}b_n$ 乘积的一个表示。显然, 这个矩阵中元素的排序方式是多种多样的, 我们常用的有对角线及正方形方法.
% 
% 在用对角线方法排序时, $\sum\limits_{n=1}^{+\infty}a_n$ 与 $\sum\limits_{n=1}^{+\infty}b_n$ 的乘积为 $\sum\limits_{n=1}^{+\infty}c_n$, 其中
% \[
%   c_n=\sum_{k+j=n+1}a_kb_j=a_1b_n+a_2b_{n-1}+\cdots+a_nb_1.
% \]
% 用对角线方法得到的级数的乘积也称为 $\sum\limits_{n=1}^{+\infty}a_n$ 与 $\sum\limits_{n=1}^{+\infty}b_n$ 的柯西乘积.
% 
% 在用正方形方法排序时, $\sum\limits_{n=1}^{+\infty}a_n$ 与 $\sum\limits_{n=1}^{+\infty}b_n$ 的乘积为 $\sum\limits_{n=1}^{+\infty}d_n$, 其中
% \[
% \begin{aligned}
%   d_1&=a_1b_1,\\
%   d_2&=a_1b_2+a_2b_2+a_2b_1,\\
%   &\cdots\cdots\cdots\\
%   d_n&=a_1b_n+a_2b_n+\cdots+a_nb_n+a_nb_{n-1}+\cdots+a_nb_1.
% \end{aligned}
% \]
% 分别记 $\sum\limits_{n=1}^{+\infty}a_n$, $\sum\limits_{n=1}^{+\infty}b_n$, $\sum\limits_{n=1}^{+\infty}d_n$ 的前 $n$ 项部分和为 $S_n,S'_n,S''_n$, 则我们有 $S''_n=S_nS'_n$。因此当 $\sum\limits_{n=1}^{+\infty}a_n$ 与 $\sum\limits_{n=1}^{+\infty}b_n$ 收敛时, $\sum\limits_{n=1}^{+\infty}d_n$ 也收敛, 并且 $\sum\limits_{n=1}^{+\infty}d_n=\sum\limits_{n=1}^{+\infty}a_n\cdot\sum\limits_{n=1}^{+\infty}b_n$.
% 
% 对于两个绝对收敛的级数的乘积, 我们有以下定理.
% 
% \begin{theorem}\label{theorem:ma-9-4-5}
% 设级数 $\sum\limits_{n=1}^{+\infty}a_n$ 与 $\sum\limits_{n=1}^{+\infty}b_n$ 都绝对收敛, 则其乘积矩阵中所有元素的任何一个排列构成的级数也绝对收敛, 并且它的和为
% \[
%   \sum_{n=1}^{+\infty}a_n\cdot\sum_{n=1}^{+\infty}b_n.
% \]
% \end{theorem}
% 
% 
% \begin{proof}
% 设 $\sum\limits_{n=1}^{+\infty}a_{k_n}b_{j_n}$ 为 $\sum\limits_{n=1}^{+\infty}a_n$ 与 $\sum\limits_{n=1}^{+\infty}b_n$ 的乘积矩阵中所有元素的一个排列构成的级数, 其前 $n$ 项部分和记为 $S_n(n=1,2,\cdots)$。对于 $\forall n\in\mathbb{N}$, 记
% \[
%   M_n=\max_{1\leq i\leq n}\{k_i,j_i\},
% \]
% 则
% \[
%   \sum_{i=1}^n|a_{k_i}b_{j_i}|
%   \leq\sum_{k=1}^{M_n}|a_k|\cdot\sum_{j=1}^{M_n}|b_j|
%   \leq\sum_{k=1}^{+\infty}|a_k|\cdot\sum_{j=1}^{+\infty}|b_j|.
% \]
% 因此 $\sum\limits_{n=1}^{+\infty}a_{k_n}b_{j_n}$ 绝对收敛。由定理~\ref{theorem:ma-9-4-3} 知, 它的任何重排也绝对收敛并且和不变.
% 
% 由于正方形方法排列得到的乘积 $\sum\limits_{n=1}^{+\infty}d_n$ 收敛到 $\sum\limits_{n=1}^{+\infty}a_n\cdot\sum\limits_{n=1}^{+\infty}b_n$, 而 $\sum\limits_{n=1}^{+\infty}d_n$ 是在乘积矩阵中所有元素的一个排列构成的级数加上一些括号所得, 因此乘积矩阵中所有元素的任何一个排列构成的级数都收敛到 $\sum\limits_{n=1}^{+\infty}a_n\cdot\sum\limits_{n=1}^{+\infty}b_n$, 即 $\sum\limits_{n=1}^{+\infty}a_{k_n}b_{j_n}$ 收敛到 $\sum\limits_{n=1}^{+\infty}a_n\cdot\sum\limits_{n=1}^{+\infty}b_n$.
% \end{proof}
% 
% 
% \begin{example}\label{example:ma-9-4-1}
% 设级数 $\sum\limits_{n=1}^{+\infty}a_n=\sum\limits_{n=1}^{+\infty}\dfrac{(-1)^{n+1}}{n^\alpha}$, 证明: 当 $\alpha\leq\dfrac12$ 时, 级数 $\sum\limits_{n=1}^{+\infty}a_n$ 与自身的柯西乘积发散.
% \end{example}
% 
% 
% \begin{proof}
% 设 $\sum\limits_{n=1}^{+\infty}a_n$ 与自身的柯西乘积为 $\sum\limits_{n=1}^{+\infty}c_n$, 则
% \[
%   c_n=\sum_{k=1}^n\frac{(-1)^{k+n+1-k}}{k^\alpha(n+1-k)^\alpha}
%   =(-1)^{n+1}\sum_{k=1}^n\frac1{k^\alpha(n+1-k)^\alpha}.
% \]
% 由于
% \[
%   \frac1{k^\alpha(n+1-k)^\alpha}
%   \geq\frac1{k^{\frac12}(n+1-k)^{\frac12}}
%   \geq\frac2{n+1},
% \]
% 因此对于 $\forall n\in\mathbb{N}$, 有
% \[
%   |c_n|\geq\frac{2n}{n+1}\geq1\quad(n=1,2,\cdots).
% \]
% 所以 $\sum\limits_{n=1}^{+\infty}c_n$ 发散.
% 
% 值得注意的是, 该级数自乘的正方形方法排序得到的级数是收敛的。请读者证明下列级数与自身的柯西乘积收敛：
% \[\sum_{n=1}^{+\infty}\frac{(-1)^{n+1}}{n}.\qedhere
% \]
% \end{proof}
% 
% 
% \section{无穷乘积}
% \label{sec:ma-9-5}
% 
% 在本节中我们讨论无穷个数相乘的问题。设 $\{a_n\}$ 是一个序列, 我们用 $\prod\limits_{n=1}^{+\infty}a_n$ 表示该序列中所有项的乘积, 并称它为该序列的无穷乘积。对于 $\forall n\in\mathbb{N}$, 记
% \[
%   T_n=\prod_{k=1}^n a_k,
% \]
% 我们称它为该无穷乘积的前 $n$ 项部分积, 而称 $\{T_n\}$ 为该无穷乘积的部分积序列.
% 
% 对于无穷乘积, 我们容易看出它的一个简单性质: 序列 $\{a_n\}$ 中若有一个元素为零, 则对所有充分大的 $n$ 有 $T_n=0$。此时部分积序列不能给我们提供任何信息。在 $\lim\limits_{n\to\infty}T_n=0$ 的情形下, 部分积序列也对无穷乘积的研究很难提供有用的信息。因此我们给出以下的定义.
% 
% \begin{definition}\label{definition:ma-9-5-1}
% 设 $\{a_n\}$ 是一个序列, $\{T_n\}$ 为无穷乘积 $\prod\limits_{n=1}^{+\infty}a_n$ 的部分积序列。若 $\lim\limits_{n\to\infty}T_n=a(a$ 为非零常数), 则称无穷乘积 $\prod\limits_{n=1}^{+\infty}a_n$ 收敛, 并记 $a=\prod\limits_{n=1}^{+\infty}a_n$。此时也称 $a$ 为该无穷乘积的积。若 $\{T_n\}$ 不收敛到一个非零常数, 则称无穷乘积 $\prod\limits_{n=1}^{+\infty}a_n$ 发散.
% \end{definition}
% 
% 由上述定义, 我们容易推出下面的定理.
% 
% \begin{theorem}\label{theorem:ma-9-5-1}
% 若无穷乘积 $\prod\limits_{n=1}^{+\infty}a_n$ 收敛, 则 $\lim\limits_{n\to\infty}a_n=1$.
% \end{theorem}
% 
% 
% \begin{proof}
% 对于 $n=1,2,\cdots$, 记 $T_n=\prod\limits_{k=1}^n a_k$。由 $\prod\limits_{n=1}^{+\infty}a_n$ 收敛, 设其积为 $a$, 则
% \[
%   \lim_{n\to\infty}a_n=\lim_{n\to\infty}\frac{T_n}{T_{n-1}}=\frac{a}{a}=1.
% \qedhere
% \]
% \end{proof}
% 
% 
% 根据定理~\ref{theorem:ma-9-5-1}, 若无穷乘积 $\prod\limits_{n=1}^{+\infty}a_n$ 收敛, 则对于 $\forall\varepsilon>0$, 当 $n$ 充分大时, 有
% \[
%   1-\varepsilon<a_n<1+\varepsilon.
% \]
% 由于改变该乘积前面的有限多项(不能变为 0)不改变无穷乘积的收敛性。因此收敛的无穷乘积的理论中基本上是对正数序列的无穷乘积的研究.
% 
% 由定理~\ref{theorem:ma-9-5-1}, 我们可将无穷乘积表成 $\prod\limits_{n=1}^{+\infty}(1+a_n)$ 的形式, 不失一般性, 我们还可以假定 $|a_n|<1(n=1,2,\cdots)$。在上述假定下, 我们容易建立以下定理.
% 
% \begin{theorem}\label{theorem:ma-9-5-2}
% 无穷乘积 $\prod\limits_{n=1}^{+\infty}(1+a_n)$ 收敛的充分必要条件是数项级数 $\sum\limits_{n=1}^{+\infty}\ln(1+a_n)$ 收敛.
% \end{theorem}
% 
% 
% \begin{proof}
% 对于 $\forall n\in\mathbb{N}$, 记 $T_n=\prod\limits_{k=1}^n(1+a_k)$, $S_n=\sum\limits_{k=1}^n\ln(1+a_k)$, 则有 $T_n=e^{S_n}(n=1,2,\cdots)$。因此由 $e^x$ 的连续性, 当 $\lim\limits_{n\to\infty}S_n=S$ 时, 有 $\lim\limits_{n\to\infty}T_n=e^S$; 反之, 当 $\lim\limits_{n\to\infty}T_n=T>0$ 时, 有 $\lim\limits_{n\to\infty}S_n=\ln T$.
% \end{proof}
% 
% 
% \begin{corollary}\label{corollary:ma-9-source-87}
% 若 $a_n>0(n=1,2,\cdots)$, 则无穷乘积 $\prod\limits_{n=1}^{+\infty}(1+a_n)$ 收敛的充分必要条件是数项级数 $\sum\limits_{n=1}^{+\infty}a_n$ 收敛.
% \end{corollary}
% 
% 
% \begin{proof}
% 对于正项级数 $\sum\limits_{n=1}^{+\infty}a_n$, 我们容易看出它与 $\sum\limits_{n=1}^{+\infty}\ln(1+a_n)$ 同时收敛与发散。由定理~\ref{theorem:ma-9-5-2} 立得推论.
% \end{proof}
% 
% 
% \begin{example}\label{example:ma-9-5-1}
% 讨论无穷乘积 $\prod\limits_{n=1}^{+\infty}\left(1+\dfrac1{n^x}\right)$ 的敛散性.
% \end{example}
% 
% 
% \begin{solve}
% 当 $x\leq1$ 时, 由 $\sum\limits_{n=1}^{+\infty}\dfrac1{n^x}$ 发散, 知 $\prod\limits_{n=1}^{+\infty}\left(1+\dfrac1{n^x}\right)$ 发散。而当 $x>1$ 时, $\sum\limits_{n=1}^{+\infty}\dfrac1{n^x}$ 收敛, 所以 $\prod\limits_{n=1}^{+\infty}\left(1+\dfrac1{n^x}\right)$ 也收敛。因此, $\prod\limits_{n=1}^{+\infty}\left(1+\dfrac1{n^x}\right)$ 当 $x>1$ 时收敛, 当 $x\leq1$ 时发散.
% \end{solve}
% 
% 
% \begin{example}\label{example:ma-9-5-2}
% 求无穷乘积 $\prod\limits_{n=1}^{+\infty}\left[1-\dfrac1{(2n)^2}\right]$ 的积.
% \end{example}
% 
% 
% \begin{solve}
% 对于 $\forall n\in\mathbb{N}$, 无穷乘积 $\prod\limits_{n=1}^{+\infty}\left[1-\dfrac1{(2n)^2}\right]$ 的前 $n$ 项部分积为
% \[
%   T_n=\prod_{k=1}^n\left[1-\frac1{(2k)^2}\right]
%   =\frac{[(2n-1)!!]^2}{[(2n)!!]^2}(2n+1).
% \]
% 在定积分中我们利用 $\sin^n x$ 在 $\left[0,\dfrac{\pi}{2}\right]$ 的积分求出了 $\lim\limits_{n\to\infty}T_n=\dfrac2\pi$, 因此
% \[
%   \prod_{n=1}^{+\infty}\left[1-\frac1{(2n)^2}\right]=\frac2\pi.
% \]
% 最后我们举一个有趣的例子结束本章.
% 
% 我们已经知道, 当 $s>1$ 时, 正项级数 $\sum\limits_{n=1}^{+\infty}\dfrac1{n^s}$ 是收敛的。现在我们来建立该级数与无穷乘积的联系。对每个固定的 $s>1$, 令
% \[
%   \zeta(s)=\sum_{n=1}^{+\infty}\frac1{n^s}.
% \]
% 首先我们观察
% \[
%   \left(1-\frac1{2^s}\right)\zeta(s)
%   =\sum_{n=1}^{+\infty}\frac1{n^s}-\sum_{n=1}^{+\infty}\frac1{(2n)^s}
%   =\sum\frac1{m^s},
% \]
% 在上式右边的求和中 $m$ 将取遍所有奇数, 即所有不是 2 的整数倍的正整数。然后再考虑
% \[
%   \left(1-\frac1{2^s}\right)\left(1-\frac1{3^s}\right)\zeta(s)
%   =\sum\frac1{m^s}-\sum\frac1{(3m)^s}
%   =\sum\frac1{k^s},
% \]
% 在上式右边求和中 $k$ 取遍所有不是 2,3 的整数倍的正整数。依此类推, 取所有素数 $2,3,5,7,\cdots,p_n,\cdots$, 则有
% \[
%   \lim_{n\to\infty}\prod_{k=1}^n\left(1-\frac1{p_k^s}\right)\zeta(s)=1,
% \]
% 即
% \[
%   \prod_{n=1}^{+\infty}\left(1-\frac1{p_n^s}\right)=\frac1{\zeta(s)}.
% \]
% 由上述等式我们易知素数的个数必定是无穷的。事实上, 倘若结论不真, 设 $p_N$ 为最大的素数, 则对任意的 $s>1$, 成立
% \[
%   \prod_{n=1}^N\left(1-\frac1{p_n^s}\right)=\frac1{\zeta(s)}.
% \]
% 令 $s\to1+0$。由 $\lim\limits_{s\to1+0}\zeta(s)=+\infty$，等式右边趋于零；左边是有限乘积，其极限为 $\prod_{n=1}^N\left(1-\frac1{p_n}\right)>0$，矛盾。因此素数有无穷多个。
% \end{solve}
% 
% % 整合来源：math-10.tex
% \chapter{函数序列与函数项级数}
% \label{chap:ma-10}
% 
% 我们知道数学分析主要是研究函数的一门学科。到目前为止, 我们经常遇到的函数, 除了一些很特别的函数外(如黎曼函数, 狄利克雷函数等), 大都是一些初等函数。一个重要的问题是: 是否存在更为广泛的函数类呢? 例如, 是否存在处处连续但处处不可导的函数呢? 再如, 设 $\{n_k\}$ 是一个严格单调上升的正整数序列, 是否存在 $C^\infty(-\infty,+\infty)$ 的函数 $f(x)$, 满足当 $i\in\{n_k\}$ 时, $f^{(i)}(0)=1$, 而当 $i\notin\{n_k\}$ 时, $f^{(i)}(0)=0$ 呢? 仅用前面所学的知识, 似乎很难解决这类问题。另外, 人们还希望利用简单函数来逼近复杂的函数, 以达到对复杂函数的深入研究。若要研究上述这些问题, 我们不可避免地要遇到函数序列或函数项级数.
% 
% 在本章中, 我们将对函数序列及函数项级数进行系统地学习。本章内容在今后几章和许多后续课程中均起着重要作用.
% 
% \section{函数序列与函数项级数的基本问题}
% \label{sec:ma-10-1}
% 
% 设有一个序列 $\{f_n(x)\}$, 其中序列的第 $n$ 项是一个函数 $f_n(x)$, 而所有 $f_n(x)(n=1,2,\cdots)$ 都有一个共同的定义域 $I_0\subset\mathbb R$。对于这样的序列, 我们今后称之为定义在 $I_0$ 上的函数序列或称该函数序列在 $I_0$ 上是有定义的。对于固定的点 $x_0\in I_0$, 该函数序列在点 $x_0$ 的取值组成了一个序列 $\{f_n(x_0)\}$。若该序列收敛, 我们则称函数序列在 $x_0$ 处收敛, $x_0$ 也称为该函数序列的收敛点。函数序列 $\{f_n(x)\}$ 的所有收敛点的集合称为它的收敛域。设 $\{f_n(x)\}$ 的收敛域为 $I\ne\varnothing$, 则对于 $x\in I$, 我们得到 $I$ 上的一个函数 $f(x)$, 它的定义为: $\forall x\in I$, $f(x)=\lim\limits_{n\to\infty}f_n(x)$。$f(x)$ 也称为 $\{f_n(x)\}$ 在 $I$ 上的极限函数.
% 
% 同样我们可以考虑和式
% \[
%   \sum_{n=1}^{+\infty}u_n(x),
% \]
% 其中 $\{u_n(x)\}$ 是定义在集合 $I_0$ 上的一个函数序列。今后我们称该和式为定义在 $I_0$ 上的函数项级数(简称级数) 或称该函数项级数在 $I_0$ 有定义。记
% \[
%   S_n(x)=\sum_{k=1}^n u_k(x)\quad(n=1,2,\cdots),
% \]
% 称之为函数项级数 $\sum\limits_{n=1}^{+\infty}u_n(x)$ 的前 $n$ 项部分和, 并称 $\{S_n(x)\}$ 为该函数项级数的部分和序列。若 $\{S_n(x)\}$ 在 $x_0\in I_0$ 处收敛, 即
% \[
%   \lim_{n\to\infty}S_n(x_0)=S(x_0),
% \]
% 则称 $\sum\limits_{n=1}^{+\infty}u_n(x)$ 在 $x_0$ 处收敛, $x_0$ 称为 $\sum\limits_{n=1}^{+\infty}u_n(x)$ 的一个收敛点, 其收敛点的全体称为 $\sum\limits_{n=1}^{+\infty}u_n(x)$ 的收敛域。设 $\sum\limits_{n=1}^{+\infty}u_n(x)$ 的收敛域为 $I$, 则对于 $\forall x\in I$, 记
% \[
%   \sum_{n=1}^{+\infty}u_n(x)=\lim_{n\to\infty}\sum_{k=1}^n u_k(x)=\lim_{n\to\infty}S_n(x)=S(x),
% \]
% 并称 $S(x)$ 为该函数项级数在 $I$ 上的和函数.
% 
% 从以上的定义来看, 求函数序列的极限函数 $f(x)$ 与函数项级数的和函数 $S(x)$ 实质上是分别求序列的极限和数项级数的和.
% 
% 由函数项级数 $\sum\limits_{n=1}^{+\infty}u_n(x)$ 收敛性的定义, 它的敛散性可以转化为函数序列
% \[
%   \left\{S_n(x)=\sum_{k=1}^n u_k(x)\right\}
% \]
% 的敛散性问题。反过来, 对于函数序列 $\{f_n(x)\}$, 则它的敛散性也可转化为函数项级数 $\sum\limits_{n=1}^{+\infty}[f_n(x)-f_{n-1}(x)]$ (令 $f_0(x)\equiv0$) 的敛散性问题.
% 
% 在本章中, 我们所关心的是极限(和)函数与函数序列(函数项级数)中的函数之间具有什么联系? 为此我们先来看以下例子.
% 
% \begin{example}\label{example:ma-10-1-1}
% 函数序列 $f_n(x)=x^n(n=1,2,\cdots)$ 的收敛域是 $(-1,1]$, 其极限函数是
% \[
%   f(x)=
%   \begin{cases}
%     0, & |x|<1,\\
%     1, & x=1.
%   \end{cases}
% \]
% 注意到函数序列 $\{x^n\}$ 中的每一个函数都是连续函数, 其极限函数在 $(-1,1)$ 内连续, 但在 $(-1,1]$ 上却不连续。同时我们还注意到对于 $\forall n\in\mathbb N$, $f_n(x)$ 均为可导函数, 但其极限函数却在 $x=1$ 处左导数不存在.
% \end{example}
% 
% 
% \begin{example}\label{example:ma-10-1-2}
% 设 $f_n(x)=nx(1-x^2)^n(x\in[0,1],n=1,2,\cdots)$, 则
% \[
%   \lim_{n\to\infty}f_n(x)\equiv0,\quad x\in[0,1],
% \]
% 但
% \[
%   \lim_{n\to\infty}\int_0^1 f_n(x)\mathrm dx
%   =\lim_{n\to\infty}\int_0^1 nx(1-x^2)^n\mathrm dx
%   =\frac12
%   \ne \int_0^1\lim_{n\to\infty}f_n(x)\mathrm dx=0.
% \]
% \end{example}
% 
% 
% \begin{example}\label{example:ma-10-1-3}
% 设 $f_n(x)=\dfrac{\sin nx}{\sqrt n}(n=1,2,\cdots)$, 显然
% \[
%   \lim_{n\to\infty}\frac{\sin nx}{\sqrt n}\equiv0=f(x),\quad x\in(-\infty,+\infty),
% \]
% 但对 $n=1,2,\cdots$, 有
% \[
%   f_n'(x)=\sqrt n\cos nx.
% \]
% 容易看出 $\{f_n'(x)\}$ 在许多点是不收敛的, 因此该函数序列的导函数序列不收敛于 $f(x)$ 的导数.
% \end{example}
% 
% 
% \begin{example}\label{example:ma-10-1-4}
% 设 $f_n(x)=\dfrac{2x}{\pi}\arctan nx(n=1,2,\cdots)$, 显然
% \[
%   \lim_{n\to\infty}f_n(x)=|x|=f(x),\quad x\in(-\infty,+\infty),
% \]
% 而对 $\forall n\in\mathbb N$, 有
% \[
%   f_n'(x)=\frac2\pi\arctan nx+\frac{2nx}{\pi(1+n^2x^2)}.
% \]
% 容易看出 $\{f_n'(x)\}$ 处处收敛, 但 $f(x)$ 在 $x=0$ 处不可导.
% \end{example}
% 
% 在讨论函数序列(函数项级数)时, 一个十分重要的问题是研究其极限函数(和函数)能从该函数序列(函数项级数)继承到什么样的性质, 其中主要的有以下三方面的问题(我们主要对函数序列叙述, 读者不难对函数项级数提出相应的问题):
% 
% (1) 若 $f_n(x)(n=1,2,\cdots)$ 是区间 $I$ 上的连续函数, 其极限函数 $f(x)$ 是否为 $I$ 上的连续函数?
% 
% (2) 若 $f_n(x)(n=1,2,\cdots)$ 是区间 $[a,b]$ 上的可积函数, 其极限函数 $f(x)$ 是否为 $[a,b]$ 上的可积函数? 若 $f(x)\in R[a,b]$, 是否成立
% \[
%   \lim_{n\to\infty}\int_a^b f_n(x)\mathrm dx=\int_a^b f(x)\mathrm dx?
% \]
% 
% (3) 若 $f_n(x)(n=1,2,\cdots)$ 是区间 $I$ 上的可导函数, 其极限函数 $f(x)$ 是否为可导函数? 若可导, 是否成立 $\lim\limits_{n\to\infty}f_n'(x)=f'(x)$?
% 
% 在上面的这些例子中, 我们可以发现: 一般来说, 上述的回答都是否定的。因此我们必须对所给的函数序列(函数项级数)加上一些条件, 才有可能得到肯定的回答。在本章以下几节中, 我们将主要研究这些问题.
% 
% \section{一致收敛的概念}
% \label{sec:ma-10-2}
% 
% 在上节中, 我们发现了很多极限函数(和函数)没有保留函数序列(函数项级数)中每个函数的相应性质。因此一个重要的问题是: 在什么条件下才能保证极限函数能保留相应函数序列中的函数的性质呢? 为此让我们再来考查一下函数序列 $\{f_n(x)=x^n\}$。在上节中我们已知其收敛域为 $(-1,1]$, 极限函数为
% \[
%   f(x)=
%   \begin{cases}
%     0, & |x|<1,\\
%     1, & x=1.
%   \end{cases}
% \]
% 我们发现 $f(x)$ 在 $(-1,1)$ 内还是保持了函数的连续性, 不仅如此, 当 $x\in(-1,1)$ 时, 还可容易地证明
% \[
%   \lim_{n\to\infty}(x^n)'=\lim_{n\to\infty}nx^{n-1}=0=\left(\lim_{n\to\infty}x^n\right)'.
% \]
% 为什么极限函数 $f(x)$ 在 $x=1$ 处不是左连续的呢? 我们可以粗略地先来分析一下。极限函数 $f(x)$ 它在 $x=1$ 左连续与它在 $x=1$ 左边邻域的函数值有关, 而现在 $f(x)$ 的函数值是 $\{x^n\}$ 的极限。尽管当 $n$ 趋于 $\infty$ 时, $\{x^n\}$ 在 $(1-\delta,1)$ 中都趋于零, 但当 $x$ 靠近 $1$ 时, 使得 $x^n$ 接近零所需要的 $n$ 也越大(见图 10.2.1)。换句话说, 对任给 $\varepsilon>0$, 不管 $n$ 多么大, 对每个固定的 $n$, $|x^n-f(x)|<\varepsilon$ 在 $(1-\delta,1)$ 中不能对所有的 $x$ 同时成立。因此极限函数就很难保持函数序列中的每一个固定函数所具有的特性了.
% 
% 另一方面, 在 $(-1,1)$ 中极限函数 $f(x)\equiv0$ 与函数序列 $\{x^n\}$ 收敛到零的过程具有何种联系呢?
% 
% 我们知道 $f(x)$ 在 $x_0$ 处连续是一个局部的性质。换句话说, $f(x)$ 在 $x_0$ 处是否连续只与 $f(x)$ 在 $x_0$ 的某个邻域内的函数值有关。设 $x_0\in(-1,1)$, 可先取定 $\delta_0$ 充分小, 使得 $[x_0-\delta_0,x_0+\delta_0]\subset(-1,1)$。由 $f(x)$ 在 $x_0$ 处连续的定义, 对于 $\forall\varepsilon>0,\exists\delta>0(\delta<\delta_0)$, 当 $|x-x_0|<\delta$ 时有
% \[
%   |f(x)-f(x_0)|<\varepsilon.
% \]
% 换句话说, $f(x)$ 在 $x_0$ 处连续本质上是 $f(x)$ 在 $x_0$ 的附近变化不大。而在 $[x_0-\delta_0,x_0+\delta_0]$ 中, 由于 $(x_0-\delta_0)^n\to0$ 与 $(x_0+\delta_0)^n\to0$, 从而 $\exists N$, 当 $n>N$ 时, 有
% \[
%   |x^n-f(x)|<\varepsilon
% \]
% 对一切 $x\in[x_0-\delta_0,x_0+\delta_0]$ 成立。特别地, 对一切 $x\in[x_0-\delta_0,x_0+\delta_0]$, 有
% \[
%   |x^{N+1}-f(x)|<\varepsilon.
% \]
% 上式说明, 在整个区间 $[x_0-\delta_0,x_0+\delta_0]$ 上, $f(x)$ 与一个连续函数 $x^{N+1}$ 的差的绝对值可以小于任意给定的正数。由于 $x^{N+1}$ 在 $x_0$ 附近的变化不大, 从而 $f(x)$ 在 $x_0$ 附近的变化也不大, 因此我们从 $n$ 充分大时 $x^n$ 的连续性就可以推出 $f(x)$ 在 $x_0$ 处的连续性.
% 
% 从上面的分析可以看出 $\{x^n\}$ 在 $(-1,1)$ 中的收敛性与在 $[a,b]\subset(-1,1)$ 中的收敛性具有本质的区别。对于后者的收敛性, 我们引入以下的重要概念.
% 
% \begin{definition}\label{definition:ma-10-2-1}
% 设 $f(x),f_n(x)(n=1,2,\cdots)$ 为定义在 $I\subset\mathbb R$ 上的函数。若对于 $\forall\varepsilon>0$, 存在 $N\in\mathbb N$, 当 $n>N$ 时, 对一切 $x\in I$, 有
% \[
%   |f_n(x)-f(x)|<\varepsilon,
% \]
% 则称函数序列 $\{f_n(x)\}$ 在 $I$ 上一致收敛于 $f(x)$, 记为 $f_n(x)\rightrightarrows f(x)(x\in I)$.
% \end{definition}
% 
% 显然, 当 $f_n(x)\rightrightarrows f(x)(x\in I)$ 时, $\{f_n(x)\}$ 在 $I$ 上收敛于 $f(x)$.
% 
% 对于函数项级数, 我们有
% 
% \begin{definition}\label{definition:ma-10-2-2}
% 设 $\sum\limits_{n=1}^{+\infty}u_n(x)$ 为定义在 $I\subset\mathbb R$ 上的函数项级数。若存在 $I$ 上的函数 $S(x)$, 使得 $\sum\limits_{n=1}^{+\infty}u_n(x)$ 的部分和序列 $S_n(x)=\sum\limits_{k=1}^n u_k(x)(n=1,2,\cdots)$ 在 $I$ 上一致收敛到 $S(x)$, 则称 $\sum\limits_{n=1}^{+\infty}u_n(x)$ 在 $I$ 上一致收敛于 $S(x)$.
% \end{definition}
% 
% 用 ``$\varepsilon-N$'' 语言来叙述函数项级数的一致收敛性, 我们有: 若对于 $\forall\varepsilon>0$, $\exists N\in\mathbb N$, 当 $n>N$ 时, 对一切 $x\in I$, 有
% \[
%   |S_n(x)-S(x)|=\left|\sum_{k=n+1}^{+\infty}u_k(x)\right|<\varepsilon,
% \]
% 则称 $\sum\limits_{n=1}^{+\infty}u_n(x)$ 在 $I$ 上一致收敛于 $S(x)$.
% 
% 下面我们来研究一下一致收敛的几何意义。设 $\{f_n(x)\}$ 为定义在 $I=[a,b]$ 上的函数序列。从几何上来看, 对于 $\forall\varepsilon>0$, 则 $y=f(x)\pm\varepsilon(x\in I)$ 的图像与 $x=a,x=b$ 围成了一个带形区域。当 $\{f_n(x)\}$ 在 $I$ 上一致收敛于 $f(x)$ 时, 则存在 $N\in\mathbb N$, 当 $n>N$ 时, $f_n(x)$ 在 $I$ 上的图像整个地落在该带形区域内(见图 10.2.2).
% 
% \begin{example}\label{example:ma-10-2-1}
% 试讨论函数序列 $f_n(x)=x^n(n=1,2,\cdots)$ 在
% 
% (1) $(-r,r)(0<r<1)$ 内的一致收敛性;
% 
% (2) $(-1,1)$ 内的一致收敛性.
% \end{example}
% 
% 
% \begin{solve}
% 我们已经知道 $\lim\limits_{n\to\infty}f_n(x)\equiv0,x\in(-1,1)$.
% 
% (1) 对于 $\forall\varepsilon>0$, 由 $r^n\to0(n\to\infty)$, 存在 $N\in\mathbb N$, 当 $n>N$ 时, 有 $|r^n|<\varepsilon$。因此当 $n>N$ 时, 对于 $\forall x\in(-r,r)$, 都有
% \[
%   |x^n-0|\leq |r^n|<\varepsilon.
% \]
% 由一致收敛的定义我们有 $x^n\rightrightarrows0(x\in(-r,r))$.
% 
% (2) 为了讨论 $f_n(x)$ 在 $(-1,1)$ 内的一致收敛性, 我们首先用肯定的语气来叙述一个函数序列 $\{f_n(x)\}$ 在 $I$ 上不一致收敛于 $f(x)$:
% 
% 若存在 $\varepsilon_0>0$, 对任意 $N>0$, 存在 $n'>N$ 以及 $x'\in I$, 使得
% \[
%   |f_{n'}(x')-f(x')|\geq\varepsilon_0,
% \]
% 则 $\{f_n(x)\}$ 在 $I$ 上不一致收敛于 $f(x)$.
% 
% 现在我们来证明 $\{x^n\}$ 在 $(-1,1)$ 内不一致收敛于零.
% 
% 事实上, 存在 $\varepsilon_0=\dfrac12$, 对任意 $N>0$, 存在 $n'=N+1>N$, 由于 $\lim\limits_{x\to1}x^{n'}=1$(对固定的 $n'$), 因此存在 $0<x'<1$, 使得
% \[
%   |(x')^{n'}-0|=(x')^{n'}>1-\frac12=\frac12.
% \]
% 这就证明了 $\{x^n\}$ 在 $(-1,1)$ 内不一致收敛于零.
% \end{solve}
% 
% 
% \begin{example}\label{example:ma-10-2-2}
% 设函数 $f(x)\in C[0,1]$ 且 $f(1)=0$, 证明:
% \[
%   x^nf(x)\rightrightarrows0\quad(x\in[0,1]).
% \]
% \end{example}
% 
% 
% \begin{proof}
% 对于 $\forall x\in[0,1]$, 显然有 $\lim\limits_{n\to\infty}x^nf(x)=0$.
% 
% 现证一致收敛性。对于 $\forall\varepsilon>0$, 由 $f(x)$ 在 $x=1$ 处左连续, 存在 $\delta>0$, 当 $x\in(1-\delta,1]$ 时, 有
% \[
%   |f(x)-0|<\varepsilon,
% \]
% 从而对于 $\forall x\in(1-\delta,1]$ 及 $\forall n\in\mathbb N$, 有
% \begin{equation}
%   |x^nf(x)|\leq |f(x)|<\varepsilon.
%   \label{eq:ma-10-2-1}
% \end{equation}
% 设 $\max\limits_{x\in[0,1]}\{|f(x)|\}=M$, 对于 $\forall x\in[0,1-\delta]$, 则有
% \[
%   |x^nf(x)|\leq M(1-\delta)^n.
% \]
% 因为 $\lim\limits_{n\to\infty}(1-\delta)^n=0$, 所以 $\exists N\in\mathbb N$, 当 $n>N$ 时, 对于 $\forall x\in[0,1-\delta]$, 有
% \[
%   |x^nf(x)|<\varepsilon.
% \]
% 注意到当 $x\in(1-\delta,1]$ 时, 对任意的 $n$, \eqref{eq:ma-10-2-1} 成立, 从而当 $n>N$ 时, 对于 $\forall x\in[0,1]$, 有 $|x^nf(x)-0|<\varepsilon$。因此 $x^nf(x)\rightrightarrows0(x\in[0,1])$.
% \end{proof}
% 
% 
% \begin{example}\label{example:ma-10-2-3}
% 试讨论函数项级数
% \[
%   \sum_{n=1}^{+\infty}\frac{(2n+1)x}{[1+(n+1)^2x](1+n^2x)}
% \]
% 在
% 
% (1) $[a,+\infty)(a>0)$ 上的一致收敛性;
% 
% (2) $(0,+\infty)$ 上的一致收敛性.
% \end{example}
% 
% 
% \begin{solve}
% 对于 $\forall n\in\mathbb N$ 和 $\forall x\in(0,+\infty)$, 记
% \[
%   S_n(x)=\sum_{k=1}^n\frac{(2k+1)x}{[1+(k+1)^2x](1+k^2x)}
%   =\frac1{1+x}-\frac1{1+(n+1)^2x},
% \]
% 则有
% \[
%   \lim_{n\to\infty}S_n(x)=S(x)=\frac1{1+x}.
% \]
% 因此我们有
% \[
%   \frac1{1+x}=\sum_{n=1}^{+\infty}\frac{(2n+1)x}{[1+(n+1)^2x](1+n^2x)},\quad x\in(0,+\infty).
% \]
% 
% (1) 当 $x\in[a,+\infty)$ 时, 我们有
% \[
%   |S_n(x)-S(x)|=\frac1{1+(n+1)^2x}\leq\frac1{1+(n+1)^2a}.
% \]
% 因此对于 $\forall\varepsilon>0$, $\exists N\in\mathbb N$, 当 $n>N$ 时, 有
% \[
%   0<\frac1{1+(n+1)^2a}<\varepsilon,
% \]
% 从而当 $n>N$ 时, 对于 $\forall x\in[a,+\infty)$, 有
% \[
%   |S_n(x)-S(x)|<\varepsilon.
% \]
% 这就证明了 $\sum\limits_{n=1}^{+\infty}\dfrac{(2n+1)x}{[1+(n+1)^2x](1+n^2x)}$ 在 $[a,+\infty)(a>0)$ 上一致收敛.
% 
% (2) 当 $x\in(0,+\infty)$ 时, 我们观察到: $\exists\varepsilon_0=\dfrac12$, 对于 $\forall N\in\mathbb N$, $\exists n'=N+1>N$, $\exists x'=\dfrac1{(N+1)^2}$, 使得
% \[
%   |S_{N+1}(x')-S(x')|=\frac1{1+(N+1)^2\dfrac1{(N+1)^2}}=\frac12.
% \]
% 这说明 $\sum\limits_{n=1}^{+\infty}\dfrac{(2n+1)x}{[1+(n+1)^2x](1+n^2x)}$ 在 $(0,+\infty)$ 上不一致收敛.
% 
% 从定义可以看出, 函数序列的一致收敛具有以下简单性质:
% 
% (1) 当函数序列 $\{f_n(x)\}$ 在 $I\subset\mathbb R$ 一致收敛时, 若 $I'\subset I$, 则 $\{f_n(x)\}$ 在 $I'$ 也一致收敛。另外若 $\{f_n(x)\}$ 分别在 $I_1,I_2\subset\mathbb R$ 上一致收敛, 则它在 $I_1\cup I_2$ 上也一致收敛.
% 
% (2) 若函数序列 $\{f_n(x)\}$ 与 $\{g_n(x)\}$ 在 $I\subset\mathbb R$ 上均一致收敛, 则对 $\forall\alpha,\beta\in\mathbb R$, $\{\alpha f_n(x)+\beta g_n(x)\}$ 在 $I$ 上也一致收敛.
% 
% 容易看出, 对函数项级数而言, 上述两个相应的性质也成立.
% 
% 人们不禁要问: 如果两个函数序列 $\{f_n(x)\},\{g_n(x)\}$ 都在 $I\subset\mathbb R$ 上一致收敛, 是否 $\{f_n(x)g_n(x)\}$ 在 $I$ 上还一致收敛呢? 一般来说, 答案是否定的。例如, 取 $I=(0,1)$, $f_n(x)=(1-x)x^n(n=1,2,\cdots)$, 由前面例~\ref{example:ma-10-2-2} 知, $\{f_n(x)\}$ 在 $(0,1)$ 内一致收敛。另外, 令 $g_n(x)=\dfrac1{1-x}(n=1,2,\cdots)$, 则 $\{g_n(x)\}$ 在 $(0,1)$ 内一致收敛。由于 $f_n(x)g_n(x)=x^n(n=1,2,\cdots)$, 从而 $\{f_n(x)g_n(x)\}$ 在 $(0,1)$ 内不一致收敛.
% \end{solve}
% 
% 
% \begin{definition}\label{definition:ma-10-2-3}
% 设 $\{f_n(x)\}$ 为定义在 $I\subset\mathbb R$ 上的函数序列。若存在 $M>0$, 使得对于 $\forall n\in\mathbb N$ 和 $\forall x\in I$, 有 $|f_n(x)|\leq M$, 则称 $\{f_n(x)\}$ 在 $I$ 上一致有界.
% \end{definition}
% 
% 若函数序列 $\{f_n(x)\}$ 与 $\{g_n(x)\}$ 在 $I\subset\mathbb R$ 都一致收敛, 并且都是一致有界的, 则可以证明 $\{f_n(x)g_n(x)\}$ 在 $I$ 上一致收敛(见本章习题).
% 
% \section{函数序列与函数项级数一致收敛的判别法}
% \label{sec:ma-10-3}
% 
% 如果一个函数序列 $\{f_n(x)\}$ 一致收敛于 $f(x)$, 此时函数序列中 $f_n(x)$ 是点点收敛于 $f(x)$。要判别函数序列(函数项级数)是否一致收敛, 一个基本的原则是: 将已有的点点收敛的判别法中的条件换成关于 $x\in I$ 中一致成立, 这样便能得到相应的关于一致收敛的判别法则.
% 
% \subsection{柯西准则}
% \label{subsec:ma-10-3-1}
% 
% 对于函数序列的一致收敛, 我们首先有以下的柯西准则.
% 
% \begin{theorem}[{柯西准则}]\label{theorem:ma-10-3-1}
% 设 $\{f_n(x)\}$ 是定义在 $I\subset\mathbb R$ 上的函数序列, 则 $\{f_n(x)\}$ 在 $I$ 上一致收敛的充分必要条件是: 对于 $\forall\varepsilon>0$, 存在 $N\in\mathbb N$, 当 $n,m>N$ 时, 对一切 $x\in I$, 有
% \[
%   |f_n(x)-f_m(x)|<\varepsilon.
% \]
% \end{theorem}
% 
% 
% \begin{proof}
% 必要性\quad 设 $f_n(x)\rightrightarrows f(x)(x\in I)$, 则对于 $\forall\varepsilon>0$, 存在 $N\in\mathbb N$, 当 $n>N$ 时, 对于 $\forall x\in I$, 有
% \[
%   |f_n(x)-f(x)|<\frac\varepsilon2.
% \]
% 因此当 $n,m>N$ 时, 对于 $\forall x\in I$, 有
% \[
%   |f_n(x)-f_m(x)|\leq |f_n(x)-f(x)|+|f_m(x)-f(x)|<\varepsilon.
% \]
% 必要性得证.
% 
% 充分性\quad 设 $\{f_n(x)\}$ 满足: 对于 $\forall\varepsilon>0$, 存在 $N\in\mathbb N$, 当 $n,m>N$ 时, 对一切 $x\in I$, 有
% \[
%   |f_n(x)-f_m(x)|<\frac\varepsilon2.
% \]
% 特别地, 上式对每个固定的点 $x\in I$ 也成立, 从而序列 $\{f_n(x)\}$ 收敛, 设其极限为 $f(x)$, 因此我们在区间 $I$ 上可得到 $\{f_n(x)\}$ 的极限函数 $f(x)$。在
% \[
%   |f_n(x)-f_m(x)|<\frac\varepsilon2
% \]
% 中令 $m\to\infty$, 我们有
% \[
%   |f_n(x)-f(x)|\leq\frac\varepsilon2<\varepsilon
% \]
% 对一切 $n>N$ 及 $x\in I$ 成立。这就证明了 $f_n(x)\rightrightarrows f(x)(x\in I)$.
% \end{proof}
% 
% 
% 对于函数项级数, 我们有
% 
% \begin{theorem}[{柯西准则}]\label{theorem:ma-10-3-1-prime}
% 设 $\sum\limits_{n=1}^{+\infty}u_n(x)$ 是定义在 $I\subset\mathbb R$ 上的函数项级数, 则 $\sum\limits_{n=1}^{+\infty}u_n(x)$ 在 $I$ 上一致收敛的充分必要条件是: 对于 $\forall\varepsilon>0$, 存在 $N\in\mathbb N$, 当 $n>m>N$ 时, 对一切 $x\in I$, 有
% \[
%   \left|\sum_{k=m+1}^n u_k(x)\right|<\varepsilon.
% \]
% 定理~\ref{theorem:ma-10-3-1-prime} 的证明留给读者.
% \end{theorem}
% 
% 从定理~\ref{theorem:ma-10-3-1-prime} 可以推出下面有用的推论.
% 
% \begin{corollary}\label{corollary:ma-10-source-17}
% 设 $\sum\limits_{n=1}^{+\infty}u_n(x)$ 是定义在 $I\subset\mathbb R$ 上的函数项级数, 并且它在 $I$ 上一致收敛, 则 $u_n(x)\rightrightarrows0(x\in I)$.
% \end{corollary}
% 
% 
% \begin{example}\label{example:ma-10-3-1}
% 设函数 $f_n(x)(n=1,2,\cdots)$ 在闭区间 $[a,b]$ 上连续, 并且 $\{f_n(x)\}$ 在开区间 $(a,b)$ 内一致收敛。证明 $\{f_n(x)\}$ 在 $[a,b]$ 上一致收敛.
% \end{example}
% 
% 
% \begin{proof}
% 由于 $\{f_n(x)\}$ 在 $(a,b)$ 内一致收敛, 由柯西准则, 对于 $\forall\varepsilon>0$, 存在 $N\in\mathbb N$, 当 $n,m>N$ 时, 对于 $\forall x\in(a,b)$, 有
% \[
%   |f_n(x)-f_m(x)|<\frac\varepsilon2.
% \]
% 对每个 $n,m>N$, 由于 $f_n(x),f_m(x)$ 在 $[a,b]$ 上连续, 因此在上述不等式中令 $x\to a+0$ 和 $x\to b-0$ 得
% \[
%   |f_n(a)-f_m(a)|\leq\frac\varepsilon2<\varepsilon
% \]
% 和
% \[
%   |f_n(b)-f_m(b)|\leq\frac\varepsilon2<\varepsilon.
% \]
% 这说明当 $n,m>N$ 时, 对一切 $x\in[a,b]$, 有
% \[
%   |f_n(x)-f_m(x)|<\varepsilon.
% \]
% 再由柯西准则知 $\{f_n(x)\}$ 在 $[a,b]$ 上一致收敛.
% 
% 读者同样可以证明以下结果:
% \end{proof}
% 
% 
% \begin{example}\label{example:ma-10-3-1-prime}
% 设函数 $u_n(x)(n=1,2,\cdots)$ 在闭区间 $[a,b]$ 上连续, 且 $\sum\limits_{n=1}^{+\infty}u_n(x)$ 在开区间 $(a,b)$ 内一致收敛, 则 $\sum\limits_{n=1}^{+\infty}u_n(x)$ 在 $[a,b]$ 上一致收敛.
% \end{example}
% 
% 
% \begin{example}\label{example:ma-10-3-2}
% 试讨论函数项级数 $\sum\limits_{n=1}^{+\infty}\dfrac{(-1)^n}{1+nx}$ 在
% 
% (1) $[a,+\infty)(a>0)$ 上的一致收敛性;
% 
% (2) $(0,+\infty)$ 上的一致收敛性.
% \end{example}
% 
% 
% \begin{proof}
% 显然, 对任意固定的 $x\in(0,+\infty)$, 该级数为一个莱布尼茨交错级数, 因此是收敛的。对于 $\forall n\in\mathbb N$, 记
% \[
%   S_n(x)=\sum_{k=1}^n\frac{(-1)^k}{1+kx}\quad\text{及}\quad S(x)=\sum_{n=1}^{+\infty}\frac{(-1)^n}{1+nx}.
% \]
% 
% (1) 当 $x\in[a,+\infty)$ 时, 由交错级数余项的性质得(见定理~\ref{theorem:ma-9-3-1} 推论后面的注)
% \[
%   |S_n(x)-S(x)|=\left|\sum_{k=n+1}^{+\infty}\frac{(-1)^k}{1+kx}\right|\leq\frac1{1+(n+1)x}\leq\frac1{1+(n+1)a}.
% \]
% 由于对于 $\forall\varepsilon>0$, $\exists N\in\mathbb N$, 当 $n>N$ 时, 有
% \[
%   \frac1{1+(n+1)a}<\varepsilon,
% \]
% 因此当 $n>N$ 时, 对于 $\forall x\in[a,+\infty)$, 有
% \[
%   |S_n(x)-S(x)|<\varepsilon.
% \]
% 这就证明了 $\sum\limits_{n=1}^{+\infty}\dfrac{(-1)^n}{1+nx}$ 在 $[a,+\infty)$ 上一致收敛.
% 
% (2) 注意到 $u_n(x)=\dfrac{(-1)^n}{1+nx}$ 在 $[0,+\infty)$ 上连续, 若 $\sum\limits_{n=1}^{+\infty}u_n(x)$ 在 $(0,+\infty)$ 上一致收敛, 由例~\ref{example:ma-10-3-1-prime} 知它在 $x=0$ 必收敛。由于 $\sum\limits_{n=1}^{+\infty}u_n(0)=\sum\limits_{n=1}^{+\infty}(-1)^n$ 是发散的, 因此 $\sum\limits_{n=1}^{+\infty}u_n(x)$ 在 $(0,+\infty)$ 上不一致收敛.
% \end{proof}
% 
% 
% \begin{remark}\label{remark:ma-10-source-23}
% (2) 也可以用以下的方法证明。由定理~\ref{theorem:ma-10-3-1-prime} 的推论可知, 我们只要证明 $u_n(x)=\dfrac{(-1)^n}{1+nx}$ 在 $(0,+\infty)$ 上不一致收敛到零即可。事实上, $\exists\varepsilon_0=\dfrac12$, 对于 $\forall N\in\mathbb N$, $\exists N+1>N$ 及 $x'=\dfrac1{N+1}\in(0,+\infty)$, 使得
% \[
%   \left|\frac{(-1)^{N+1}}{1+(N+1)x'}\right|=\frac12=\varepsilon_0.
% \]
% 这就证明了 $\sum\limits_{n=1}^{+\infty}\dfrac{(-1)^n}{1+nx}$ 在 $(0,+\infty)$ 上不一致收敛.
% \end{remark}
% 
% 
% \subsection{一致收敛的判别法}
% \label{subsec:ma-10-3-2}
% 
% 柯西准则是一个在理论上非常重要的定理, 但有时在具体应用上不是很方便。以下的判别法在实用上是比较方便的.
% 
% \begin{theorem}\label{theorem:ma-10-3-2}
% 设函数序列 $\{f_n(x)\}$ 在集合 $I\subset\mathbb R$ 上有定义, 则 $f_n(x)\rightrightarrows f(x)(x\in I)$ 的充分必要条件是 $\lim\limits_{n\to\infty}\sup\limits_{x\in I}\{|f_n(x)-f(x)|\}=0$.
% \end{theorem}
% 
% 
% \begin{proof}
% 必要性\quad 设 $f_n(x)\rightrightarrows f(x)(x\in I)$, 则对于 $\forall\varepsilon>0$, 存在 $N\in\mathbb N$, 当 $n>N$ 时, 对一切 $x\in I$, 有 $|f_n(x)-f(x)|<\dfrac\varepsilon2$。因此我们有
% \[
%   0\leq \sup_{x\in I}\{|f_n(x)-f(x)|\}\leq\frac\varepsilon2<\varepsilon,
% \]
% 此即 $\lim\limits_{n\to\infty}\sup\limits_{x\in I}\{|f_n(x)-f(x)|\}=0$。必要性得证.
% 
% 充分性\quad 由于对于 $\forall\varepsilon>0$, 存在 $N\in\mathbb N$, 当 $n>N$ 时, 有
% \[
%   \sup_{x\in I}\{|f_n(x)-f(x)|\}<\varepsilon,
% \]
% 因此, 当 $n>N$ 时, 对一切 $x\in I$, 有
% \[
%   |f_n(x)-f(x)|\leq \sup_{x\in I}\{|f_n(x)-f(x)|\}<\varepsilon,
% \]
% 此即 $f_n(x)\rightrightarrows f(x)(x\in I)$.
% \end{proof}
% 
% 
% 以上定理告诉我们, 当我们知道一个函数序列的极限函数时, 判断它的一致收敛性的问题可转化为求上确界以及讨论序列的收敛性问题, 以上的判别法被人们称为最值判别法.
% 
% \begin{example}\label{example:ma-10-3-3}
% 试讨论 $f_n(x)=\dfrac{x}{1+n^2x^2}(x\in(-\infty,+\infty), n=1,2,\cdots)$ 的一致收敛性.
% \end{example}
% 
% 
% \begin{solve}
% 对于 $\forall x\in(-\infty,+\infty)$, 不难看出 $f_n(x)\to0(n\to\infty)$。对每个固定的 $n$, 我们来求 $|f_n(x)|$ 在 $(-\infty,+\infty)$ 上的上确界。由于 $f_n(x)$ 是奇函数且 $f_n(x)\geq0(x\geq0)$, 因此只要计算 $f_n(x)$ 在 $x\geq0$ 时的上确界即可.
% 
% 由
% \[
%   f_n'(x)=\frac{(1+n^2x^2)-x\cdot2n^2x}{(1+n^2x^2)^2}
%   =\frac{1-n^2x^2}{(1+n^2x^2)^2}=0
% \]
% 得 $x=\pm\dfrac1n$, 注意到 $f(0)=0$ 和 $\lim\limits_{x\to+\infty}f_n(x)=0$, 知 $f_n(x)$ 在 $\dfrac1n$ 取到最大值 $\dfrac1{2n}$.
% 
% 因为
% \[
%   \lim_{n\to\infty}\sup_{x\in(-\infty,+\infty)}|f_n(x)-0|
%   =\lim_{n\to\infty}\frac1{2n}=0,
% \]
% 所以
% \[
%   \frac{x}{1+n^2x^2}\rightrightarrows0\quad(x\in(-\infty,+\infty)).
% \qedhere
% \]
% \end{solve}
% 
% 
% \begin{example}\label{example:ma-10-3-4}
% 证明函数序列
% \[
%   f_n(x)=n^2\left(e^{\frac1{nx}}-1\right)\sin\frac1{nx}\quad(n=1,2,\cdots)
% \]
% 在 $[a,+\infty)(a>0)$ 上一致收敛.
% \end{example}
% 
% 
% \begin{proof}
% 对每个 $x\in[a,+\infty)$, 我们有
% \[
%   \lim_{n\to\infty}f_n(x)
%   =\lim_{n\to\infty}\frac{\sin\frac1{nx}}{\frac1{nx}}\cdot
%   \frac{e^{\frac1{nx}}-1}{\frac1{nx}}\cdot\frac1{x^2}
%   =\frac1{x^2}.
% \]
% 注意到当 $t\to0$ 时, 有
% \[
%   e^t-1\sim t+o(t),\qquad \sin t\sim t+o(t),
% \]
% 从而当 $x\in[a,+\infty)$ 且 $n\to\infty$ 时, 有
% \[
% \begin{aligned}
%   \left|n^2\left(e^{\frac1{nx}}-1\right)\sin\frac1{nx}-\frac1{x^2}\right|
%   &=n^2\left|\left(e^{\frac1{nx}}-1\right)\sin\frac1{nx}-\frac1{n^2x^2}\right|\\
%   &=n^2\left|\left[\frac1{nx}+o\left(\frac1{nx}\right)\right]
%   \left[\frac1{nx}+o\left(\frac1{nx}\right)\right]-\frac1{n^2x^2}\right|\\
%   &=n^2o\left(\frac1{n^2x^2}\right).
% \end{aligned}
% \]
% 因此我们有
% \[
%   \sup_{x\in[a,+\infty)}
%   \left|n^2\left(e^{\frac1{nx}}-1\right)\sin\frac1{nx}-\frac1{x^2}\right|
%   \leq n^2o\left(\frac1{n^2a^2}\right)\to0\quad(n\to\infty).
% \]
% 这就证明了 $f_n(x)\rightrightarrows\dfrac1{x^2}(x\in[a,+\infty))$.
% 
% 对于函数项级数来说, 我们也可以对其部分和序列利用定理~\ref{theorem:ma-10-3-2} 来判别其收敛性。如我们可以用定理~\ref{theorem:ma-10-3-2} 来讨论 $\left\{\dfrac{1-x^{n+1}}{1-x}\right\}$, 从而可以解决 $\sum\limits_{n=1}^{+\infty}x^n$ 在 $(-1,1)$ 内的一致收敛性问题。但在一般情形, 要求出一个函数项级数部分和序列的通项表达式不是一件容易的事情, 因此就很难用定理~\ref{theorem:ma-10-3-2} 来判别其一致收敛性.
% 
% 对于函数项级数, 我们有以下的概念.
% \end{proof}
% 
% 
% \begin{definition}\label{definition:ma-10-3-1}
% 设 $\sum\limits_{n=1}^{+\infty}u_n(x)$ 是定义在 $I$ 上的函数项级数, 并且 $\sum\limits_{n=1}^{+\infty}|u_n(x)|$ 在 $I$ 上一致收敛, 则称 $\sum\limits_{n=1}^{+\infty}u_n(x)$ 在 $I$ 上绝对一致收敛.
% \end{definition}
% 
% 利用柯西准则, 读者易于证明, 若 $\sum\limits_{n=1}^{+\infty}u_n(x)$ 在 $I$ 上绝对一致收敛, 则它在 $I$ 上一致收敛.
% 
% 对于函数项级数, 以下的判别法是经常用到的.
% 
% \begin{theorem}[{魏尔斯特拉斯 M 判别法}]\label{theorem:ma-10-3-3}
% 设函数项级数 $\sum\limits_{n=1}^{+\infty}u_n(x)$ 在 $I\subset\mathbb R$ 有定义。若存在正数序列 $\{M_n\}$, 使得对每个 $n\in\mathbb N$, 对一切 $x\in I$, 有
% \[
%   |u_n(x)|\leq M_n,
% \]
% 并且 $\sum\limits_{n=1}^{+\infty}M_n$ 收敛, 则 $\sum\limits_{n=1}^{+\infty}u_n(x)$ 在 $I$ 上绝对一致收敛.
% \end{theorem}
% 
% 
% \begin{proof}
% 由于 $\sum\limits_{n=1}^{+\infty}M_n$ 收敛, 根据数项级数收敛的柯西准则, 对于 $\forall\varepsilon>0$, 存在 $N\in\mathbb N$, 当 $n>m>N$ 时, 有
% \[
%   \sum_{k=m+1}^n M_k<\varepsilon.
% \]
% 从而对一切 $x\in I$, 有
% \[
%   \sum_{k=m+1}^n |u_k(x)|\leq\sum_{k=m+1}^n M_k<\varepsilon,
% \]
% 再由函数项级数一致收敛的柯西准则知 $\sum\limits_{n=1}^{+\infty}u_n(x)$ 在 $I$ 上绝对一致收敛.
% \end{proof}
% 
% 
% \begin{example}\label{example:ma-10-3-5}
% 证明 $\sum\limits_{n=1}^{+\infty}x^ke^{-nx}(k>1)$ 在 $[0,+\infty)$ 上一致收敛.
% \end{example}
% 
% 
% \begin{proof}
% 对每个正整数 $n$, 当 $x\in(0,+\infty)$ 时, $u_n(x)=x^ke^{-nx}>0$。注意到 $u_n(0)=0$ 和 $\lim\limits_{x\to+\infty}x^ke^{-nx}=0(n=1,2,\cdots)$, 因此 $u_n(x)$ 在 $(0,+\infty)$ 上取到最大值。由
% \[
%   u_n'(x)=kx^{k-1}e^{-nx}-nx^ke^{-nx}=0
% \]
% 得 $x=0$ 和 $x=\dfrac{k}{n}$。因此 $u_n(x)$ 在 $x=\dfrac{k}{n}$ 取最大值 $e^{-k}\dfrac{k^k}{n^k}$.
% 
% 由于 $\sum\limits_{n=1}^{+\infty}e^{-k}\dfrac{k^k}{n^k}$ 当 $k>1$ 时收敛, 由魏尔斯特拉斯 M 判别法知 $\sum\limits_{n=1}^{+\infty}x^ke^{-nx}$ 在 $[0,+\infty)$ 上一致收敛.
% 
% 对于不是绝对一致收敛的函数项级数, 我们有下面的狄利克雷判别法与阿贝尔判别法.
% \end{proof}
% 
% 
% \begin{theorem}[{狄利克雷判别法}]\label{theorem:ma-10-3-4}
% 设函数序列 $u_n(x),v_n(x)(n=1,2,\cdots)$ 在 $I\subset\mathbb R$ 上有定义, 并且满足以下条件:
% 
% (1) $\sum\limits_{n=1}^{+\infty}u_n(x)$ 的部分和序列在 $I$ 上一致有界, 即存在 $M>0$, 使得对于 $\forall n\geq1$ 及 $\forall x\in I$, 有
% \[
%   |S_n(x)|=\left|\sum_{k=1}^n u_k(x)\right|\leq M;
% \]
% 
% (2) 对每个 $x\in I$, $\{v_n(x)\}$ 关于 $n$ 是单调的, 且 $v_n(x)\rightrightarrows0(x\in I)$,
% 
% 则 $\sum\limits_{n=1}^{+\infty}u_n(x)v_n(x)$ 在 $I$ 上一致收敛.
% \end{theorem}
% 
% 
% \begin{proof}
% 由条件 (1), 当 $n>m\geq1$ 时, 对一切 $x\in I$, 我们有
% \[
%   \left|\sum_{k=m+1}^n u_k(x)\right|=|S_n(x)-S_m(x)|\leq2M.
% \]
% 由 $v_n(x)\rightrightarrows0(x\in I)$, 对于 $\forall\varepsilon>0$, $\exists N\in\mathbb N$, 当 $n>N$ 时, 对于 $\forall x\in I$ 有
% \[
%   |v_n(x)|<\frac{\varepsilon}{4M}.
% \]
% 注意到对每个 $x\in I$, $\{v_k(x)\}$ 关于 $n$ 是单调的, 如同数项级数的狄利克雷判别法的证明, 利用阿贝尔变换 (引理~\ref{lemma:ma-7-5-4}) 知, 当 $n>m>N$ 时, 对 $\forall x\in I$, 有
% \[
%   \left|\sum_{k=m+1}^n u_k(x)v_k(x)\right|
%   \leq M(|v_{m+1}(x)|+2|v_n(x)|)<\varepsilon.
% \]
% 由柯西准则知 $\sum\limits_{n=1}^{+\infty}u_n(x)v_n(x)$ 在 $I$ 上一致收敛.
% \end{proof}
% 
% 
% \begin{theorem}[{阿贝尔判别法}]\label{theorem:ma-10-3-5}
% 设函数序列 $u_n(x),v_n(x)(n=1,2,\cdots)$ 在 $I\subset\mathbb R$ 上有定义, 并且满足以下条件:
% 
% (1) $\sum\limits_{n=1}^{+\infty}u_n(x)$ 在 $I$ 上一致收敛;
% 
% (2) 对固定的 $x\in I$, $\{v_n(x)\}$ 关于 $n$ 是单调函数, 且 $\{v_n(x)\}$ 在 $I$ 上一致有界,
% 
% 则 $\sum\limits_{n=1}^{+\infty}u_n(x)v_n(x)$ 在 $I$ 上一致收敛.
% \end{theorem}
% 
% 
% \begin{proof}
% 由条件 (2), 存在 $M>0$, 使得对于 $\forall n\geq1$ 及 $\forall x\in I$, 有
% \[
%   |v_n(x)|\leq M.
% \]
% 再由 (1) 及柯西准则, 对于 $\forall\varepsilon>0$, $\exists N\in\mathbb N$, 当 $n>m>N$ 时, 对一切 $x\in I$, 有
% \[
%   \left|\sum_{k=m+1}^n u_k(x)\right|\leq\frac{\varepsilon}{3M}.
% \]
% 由于 $\{v_n(x)\}$ 关于 $n$ 是单调的, 由阿贝尔变换 (引理~\ref{lemma:ma-7-5-4}), 对任意的 $n>m>N$ 及 $\forall x\in I$, 有
% \[
%   \left|\sum_{k=m+1}^n u_k(x)v_n(x)\right|
%   \leq\frac{\varepsilon}{3M}(|v_{m+1}(x)|+2|v_n(x)|)<\varepsilon.
% \]
% 再由柯西准则知 $\sum\limits_{n=1}^{+\infty}u_n(x)v_n(x)$ 在 $I$ 上一致收敛.
% \end{proof}
% 
% 
% 回忆一下, 在数项级数 $\sum\limits_{n=1}^{+\infty}a_nb_n$ 的阿贝尔判别法的证明中, 我们直接设单调序列 $\{b_n\}$ 的极限为 $b$, 然后对 $\sum\limits_{n=1}^{+\infty}a_n(b_n-b)$ 应用狄利克雷判别法即可得到证明。对函数项级数 $\sum\limits_{n=1}^{+\infty}u_n(x)v_n(x)$ 而言, 对每点 $x\in I$, $\lim\limits_{n\to\infty}v_n(x)$ 存在, 因此我们也有极限函数 $v(x)$。为了要应用函数项级数的狄利克雷判别法, 我们必须要证明 $v_n(x)\rightrightarrows v(x)(x\in I)$。这个结论是否正确呢? 一般来说此结论是不成立, 例如 $\{x^n\}$ 在 $[0,1]$ 上关于 $n$ 单调并且一致有界, 但 $\{x^n\}$ 在该区间上并不一致收敛.
% 
% \begin{example}\label{example:ma-10-3-6}
% 证明函数项级数 $\sum\limits_{n=1}^{+\infty}\dfrac{(-1)^n}{n+x}$ 在
% 
% (1) $[0,+\infty)$ 上一致收敛;
% 
% (2) 任何不含负整数点的闭区间 $I$ 上一致收敛;
% 
% (3) 任何区间 $I$ 上不绝对收敛.
% \end{example}
% 
% 
% \begin{proof}
% (1) 对任意的正整数 $n$, 令 $u_n(x)=(-1)^n$, 则对于 $x\in[0,+\infty)$, 有
% \[
%   \left|\sum_{k=1}^n u_k(x)\right|\leq1.
% \]
% 令 $v_n(x)=\dfrac1{n+x}(n=1,2,\cdots)$, 显然 $v_n(x)$ 对一切 $x\in[0,+\infty)$ 均是关于 $n$ 单调下降的函数, 且当 $x\in[0,+\infty)$ 时, 有
% \[
%   |v_n(x)|\leq\frac1n\to0\quad(n\to\infty).
% \]
% 因此对于 $x\in[0,+\infty)$, 有
% \[
%   v_n(x)\rightrightarrows0\quad(x\in[0,+\infty)).
% \]
% 根据狄利克雷判别法, $\sum\limits_{n=1}^{+\infty}\dfrac{(-1)^n}{n+x}$ 在 $[0,+\infty)$ 上一致收敛.
% 
% (2) 若闭区间 $I\subset[0,+\infty)$, 显然结论成立。不妨设 $I\subset(-K,-K+1), K\in\mathbb N$。如同 (1), 我们容易验证: 对于 $\forall n\in\mathbb N$ 及 $\forall x\in I$, $\left|\sum\limits_{k=1}^n u_k(x)\right|\leq1$。当 $n>K$ 时, $v_n(x)$ 对于 $x\in I$ 均是关于 $n$ 单调上升的且 $v_n(x)\rightrightarrows0(x\in I)$。因此由狄利克雷判别法, $\sum\limits_{n=1}^{+\infty}\dfrac{(-1)^n}{n+x}$ 在 $I$ 上一致收敛.
% 
% (3) $\sum\limits_{n=1}^{+\infty}\dfrac1{n+x}$ 对所有非负整数的点 $x$, 当 $n$ 充分大时, 均是正项级数, 且 $\dfrac1{n+x}\sim\dfrac1n(n\to\infty)$。由于 $\sum\limits_{n=1}^{+\infty}\dfrac1n$ 发散, 因此 $\sum\limits_{n=1}^{+\infty}\dfrac1{n+x}$ 也发散。因此 $\sum\limits_{n=1}^{+\infty}\dfrac{(-1)^n}{n+x}$ 在任何区间都不是绝对收敛.
% \end{proof}
% 
% 
% \begin{example}\label{example:ma-10-3-7}
% 证明函数项级数 $\sum\limits_{n=1}^{+\infty}\dfrac{a_n}{n^x}$ 在 $[1,+\infty)$ 上一致收敛的充分必要条件是 $\sum\limits_{n=1}^{+\infty}\dfrac{a_n}{n}$ 收敛.
% \end{example}
% 
% 
% \begin{proof}
% 必要性是显然的, 现证充分性.
% 
% 设 $\sum\limits_{n=1}^{+\infty}\dfrac{a_n}{n}$ 收敛, 在
% \[
%   \sum_{n=1}^{+\infty}\frac{a_n}{n^x}
%   =\sum_{n=1}^{+\infty}\frac{a_n}{nn^{x-1}}
% \]
% 中, 对于 $n=1,2,\cdots$, 我们令
% \[
%   u_n(x)=\frac{a_n}{n},\qquad v_n(x)=\frac1{n^{x-1}}.
% \]
% 因为 $\sum\limits_{n=1}^{+\infty}u_n(x)=\sum\limits_{n=1}^{+\infty}\dfrac{a_n}{n}$ 收敛, 故它在 $[1,+\infty)$ 上一致收敛。由于 $\{v_n(x)\}$ 对固定的 $x$ 是 $n$ 的单调函数, 且对于 $\forall x\in[1,+\infty)$ 及 $\forall n\in\mathbb N$, 有
% \[
%   |v_n(x)|=\left|\frac1{n^{x-1}}\right|\leq1,
% \]
% 根据阿贝尔判别法, $\sum\limits_{n=1}^{+\infty}\dfrac{a_n}{n^x}$ 在 $[1,+\infty)$ 上一致收敛.
% \end{proof}
% 
% 
% \begin{example}\label{example:ma-10-3-8}
% 试讨论函数项级数 $\sum\limits_{n=1}^{+\infty}(-1)^n(1-x)x^n$ 在区间 $[0,1]$ 上的
% 
% (1) 一致收敛性;
% 
% (2) 绝对收敛性;
% 
% (3) 绝对一致收敛性.
% \end{example}
% 
% 
% \begin{solve}
% (1) 取 $u_n(x)=(-1)^n$, $v_n(x)=(1-x)x^n(n=1,2,\cdots)$, 则 $\left\{\sum\limits_{k=1}^n u_k(x)\right\}$ 在 $[0,1]$ 上一致有界, 而 $v_n(x)$ 对固定的 $x$ 是 $n$ 的单调下降函数, 而且 $v_n(x)\rightrightarrows0(x\in[0,1])$ (见例~\ref{example:ma-10-2-2})。根据狄利克雷判别法, $\sum\limits_{n=1}^{+\infty}(-1)^n(1-x)x^n$ 在 $[0,1]$ 上一致收敛.
% 
% (2) 对于 $n=1,2,\cdots$, 记
% \[
%   S_n(x)=\sum_{k=1}^n(1-x)x^n=
%   \begin{cases}
%     x-x^{n+1}, & x\in[0,1),\\
%     0, & x=1.
%   \end{cases}
% \]
% 注意当 $x\in[0,1)$ 时,
% \[
%   \lim_{n\to\infty}S_n(x)=S(x)=x.
% \]
% 因此当 $x\in[0,1]$ 时, $\sum\limits_{n=1}^{+\infty}|(-1)^n(1-x)x^n|$ 收敛, 从而原级数在 $[0,1]$ 上绝对收敛.
% 
% (3) 继续使用 (2) 中的记号。由于在 $[0,1)$ 中 $|S_n(x)-S(x)|=x^{n+1}$, 而我们已经知道 $\{x^n\}$ 在 $[0,1)$ 上不一致收敛, 因此 $\sum\limits_{n=1}^{+\infty}(-1)^n(1-x)x^n$ 在 $[0,1]$ 不是绝对一致收敛的.
% 
% 此例说明了即使一个函数项级数在一个区间上一致收敛并且绝对收敛, 但该级数在该区间上不一定绝对一致收敛.
% \end{solve}
% 
% 
% \section{一致收敛的函数序列和函数项级数}
% \label{sec:ma-10-4}
% 
% 在本节中, 我们来讨论一致收敛的函数序列(函数项级数)的极限函数(和函数)的性质, 并且回答 第~\ref{sec:ma-10-1}~节 中提出的问题.
% 
% \subsection{极限函数的连续性}
% \label{subsec:ma-10-4-1}
% 
% \begin{theorem}\label{theorem:ma-10-4-1}
% 设函数 $f_n(x)\in C[a,b](n=1,2,\cdots)$, 且 $f_n(x)\rightrightarrows f(x)(x\in[a,b])$, 则 $f(x)\in C[a,b]$.
% \end{theorem}
% 
% 
% \begin{proof}
% 任取 $x_0\in[a,b]$, 我们证 $f(x)$ 在 $x_0$ 处连续。对于 $\forall\varepsilon>0$, 由于 $f_n(x)\rightrightarrows f(x)(x\in[a,b])$, 所以 $\exists N\in\mathbb N$, 当 $n>N$ 时, 对一切 $x\in[a,b]$, 有
% \[
%   |f_n(x)-f(x)|<\frac\varepsilon3.
% \]
% 取定 $n_0>N$(例如取 $n_0=N+1$), 由于 $f_{n_0}(x)$ 在 $x_0\in[a,b]$ 处连续, 从而存在 $\delta>0$, 使得当 $x\in U(x_0,\delta)\cap[a,b]$ 时, 有
% \[
%   |f_{n_0}(x)-f_{n_0}(x_0)|<\frac\varepsilon3.
% \]
% 因此当 $x\in U(x_0,\delta)\cap[a,b]$ 时, 有
% \[
% \begin{aligned}
%   |f(x)-f(x_0)|
%   &\leq |f(x)-f_{n_0}(x)|+|f_{n_0}(x)-f_{n_0}(x_0)|+|f_{n_0}(x_0)-f(x_0)|\\
%   &<\varepsilon.
% \end{aligned}
% \qedhere
% \]
% \end{proof}
% 
% 
% \begin{remark}\label{remark:ma-10-source-47}
% 仿照定理~\ref{theorem:ma-10-4-1} 的证明方法我们可以证明以下结果: 设 $\{f_n(x)\}$ 是定义在 $[a,b]\setminus\{x_0\}$ 上的函数序列, 其中 $x_0\in[a,b]$。如果 $f_n(x)\rightrightarrows f(x), (x\in[a,b]\setminus\{x_0\})$, 并且对每个 $n\geq1$, 有 $\lim\limits_{x\to x_0}f_n(x)=c_n$, 则 $\lim\limits_{n\to+\infty}c_n$ 存在, 并且有
% \[
%   \lim_{x\to x_0}f(x)=\lim_{n\to+\infty}c_n.
% \]
% \end{remark}
% 
% 对函数项级数, 用完全类似于定理~\ref{theorem:ma-10-4-1} 的证明方法, 可以证明
% 
% \begin{theorem}\label{theorem:ma-10-4-1-prime}
% 设函数 $u_n(x)\in C[a,b](n=1,2,\cdots)$, 且 $\sum\limits_{n=1}^{+\infty}u_n(x)$ 在 $[a,b]$ 上一致收敛, 则 $\sum\limits_{n=1}^{+\infty}u_n(x)$ 的和函数在 $[a,b]$ 上连续.
% \end{theorem}
% 
% 
% \begin{remark}\label{remark:ma-10-source-49}
% 用类似的方法我们还可以证明以下结论: 若 $\sum\limits_{n=1}^{+\infty}u_n(x)$ 在 $I\setminus\{x_0\}$ 上一致收敛, 且对每个 $n\geq1$, 有 $\lim\limits_{x\to x_0}u_n(x)=a_n$, 则 $\sum\limits_{n=1}^{+\infty}a_n$ 收敛, 且
% \[
%   \lim_{x\to x_0}\sum_{n=1}^{+\infty}u_n(x)
%   =\sum_{n=1}^{+\infty}\left[\lim_{x\to x_0}u_n(x)\right]
%   =\sum_{n=1}^{+\infty}a_n.
% \]
% \end{remark}
% 
% 从定理~\ref{theorem:ma-10-4-1} 我们可知, 当连续的函数序列 $\{f_n(x)\}$(即对每个 $n$, $f_n(x)$ 连续) 一致收敛时, 有以下极限顺序的交换性, 即
% \[
%   \lim_{x\to x_0}\lim_{n\to\infty}f_n(x)
%   =\lim_{n\to\infty}\lim_{x\to x_0}f_n(x).
% \]
% 从定理~\ref{theorem:ma-10-4-1-prime} 可知, 当连续的函数项级数 $\sum\limits_{n=1}^{+\infty}u_n(x)$(即对每个 $n$, $u_n(x)$ 连续) 一致收敛时, 有以下极限号与求和号的可交换性
% \[
%   \lim_{x\to x_0}\sum_{n=1}^{+\infty}u_n(x)
%   =\sum_{n=1}^{+\infty}\lim_{x\to x_0}u_n(x).
% \]
% 我们曾经多次指出, 函数连续是一个局部性的概念。换句话说, 函数 $f(x)$ 在 $x=x_0$ 处连续仅与 $f(x)$ 在 $x_0$ 附近的函数值有关。因此若一个连续函数序列在 $x_0$ 的一个邻域内一致收敛到 $f(x)$(此时也称它在 $x_0$ 处是局部一致收敛的), 则 $f(x)$ 在 $x_0$ 处连续。在数学的一些分支中, 以下内闭一致收敛的概念经常用到.
% 
% \begin{definition}\label{definition:ma-10-4-1}
% 设函数序列 $\{f_n(x)\}$ (函数项级数 $\sum\limits_{n=1}^{+\infty}u_n(x)$) 在区间 $I$ 上有定义, 若对任何的闭区间 $[a,b]\subset I$, $\{f_n(x)\}$ ($\sum\limits_{n=1}^{+\infty}u_n(x)$) 在 $[a,b]$ 上一致收敛, 则称 $\{f_n(x)\}$ ($\sum\limits_{n=1}^{+\infty}u_n(x)$) 在 $I$ 内闭一致收敛.
% \end{definition}
% 
% 读者不难证明, 对于开区间 $I$ 来说, 内闭一致收敛等价于局部一致收敛。因此我们有
% 
% \begin{corollary}\label{corollary:ma-10-source-51}
% (1) 设函数 $f_n(x)\in C(a,b)(n=1,2,\cdots)$ 且 $\{f_n(x)\}$ 在 $(a,b)$ 内闭一致收敛于 $f(x)$, 则 $f(x)\in C(a,b)$;
% 
% (2) 设函数 $u_n(x)\in C(a,b)$, 且 $\sum\limits_{n=1}^{+\infty}u_n(x)$ 在 $(a,b)$ 内闭一致收敛, 则 $\sum\limits_{n=1}^{+\infty}u_n(x)$ 的和函数在 $(a,b)$ 内连续.
% \end{corollary}
% 
% 我们前面已经知道 $\{x^n\}$ 在 $(-1,1)$ 内闭一致收敛, 从而它的极限函数 $0$ 在 $(-1,1)$ 中连续.
% 
% \begin{example}\label{example:ma-10-4-1}
% 设 $\{a_n\}$ 是一个单调序列且 $\lim\limits_{n\to\infty}a_n=0$, 证明级数 $\sum\limits_{n=1}^{+\infty}a_n\cos nx$ 与 $\sum\limits_{n=1}^{+\infty}a_n\sin nx$ 的和函数在 $(0,2\pi)$ 内连续.
% \end{example}
% 
% 
% \begin{proof}
% 对于 $\forall n\in\mathbb N$, 令
% \[
%   u_n(x)=a_n,\qquad v_n(x)=\cos nx.
% \]
% 下面证明 $\sum\limits_{n=1}^{+\infty}u_n(x)v_n(x)$ 在 $(0,2\pi)$ 内闭一致收敛。事实上, 任取闭区间 $[\delta_0,2\pi-\delta_0](\delta_0>0)$, 对于 $\forall n\in\mathbb N$, 我们有
% \[
%   \left|\sum_{k=1}^n\cos kx\right|
%   =\frac{\left|\sin\left(n+\frac12\right)x-\sin\frac{x}{2}\right|}{2\sin\frac{x}{2}}
%   \leq\frac1{\sin\frac{\delta_0}{2}}.
% \]
% 由于 $u_n(x)$ 与 $x$ 无关, $\{u_n(x)\}$ 关于 $n$ 单调且在 $[\delta_0,2\pi-\delta_0]$ 上一致趋于零, 根据狄利克雷判别法, $\sum\limits_{n=1}^{+\infty}a_n\cos nx$ 在 $[\delta_0,2\pi-\delta_0]$ 上一致收敛.
% 
% 由于 $a_n\cos nx\in C(-\infty,+\infty)(n=1,2,\cdots)$, 因此
% \[
%   \sum_{n=1}^{+\infty}a_n\cos nx\in C[\delta_0,2\pi-\delta_0].
% \]
% 由 $\delta_0$ 的任意性, 我们有
% \[
%   \sum_{n=1}^{+\infty}a_n\cos nx\in C(0,2\pi).
% \]
% 利用
% \[
%   \left|\sum_{k=1}^n\sin kx\right|
%   =\frac{\left|\cos\left(n+\frac12\right)x-\cos\frac{x}{2}\right|}{2\sin\frac{x}{2}}
% \]
% 同理可证 $\sum\limits_{n=1}^{+\infty}a_n\sin nx\in C(0,2\pi)$.
% 
% 值得指出的是, 定理~\ref{theorem:ma-10-4-1} 及定理~\ref{theorem:ma-10-4-1-prime} 都是给出了极限函数连续的充分条件。容易举出例子来说明这些条件不是必要的, 例如 $f_n(x)=nxe^{-nx^2}(n=1,2,\cdots)$ 在 $[0,1]$ 点点收敛到连续函数 $0$, 但它在 $[0,1]$ 上不一致收敛; 又如函数项级数 $\sum\limits_{n=0}^{+\infty}(1-x^2)x^n=1+x, x\in(-1,1)$, 其和函数在 $(-1,1)$ 内连续, 但该函数项级数在 $(-1,1)$ 内不一致收敛.
% 
% 下面的狄尼定理告诉我们, 如果附加一定的条件, 我们可以从点收敛推出一致收敛性.
% \end{proof}
% 
% 
% \begin{theorem}[{狄尼(Dini)定理}]\label{theorem:ma-10-4-2}
% 设函数 $f_n(x)(n=1,2,\cdots)$ 在区间 $[a,b]$ 上连续, 且对于 $\forall x\in[a,b]$ 及 $\forall n\in\mathbb N$ 成立 $f_n(x)\leq f_{n+1}(x)$, 再设对于 $\forall x\in[a,b]$, 有 $\lim\limits_{n\to\infty}f_n(x)=f(x)$(即 $\{f_n(x)\}$ 是点点收敛的单调函数序列), 则 $f(x)$ 在区间 $[a,b]$ 上连续的充分必要条件是
% \[
%   f_n(x)\rightrightarrows f(x)\quad(x\in[a,b]).
% \]
% \end{theorem}
% 
% 
% \begin{proof}
% 充分性是显然的, 现证必要性。我们用反证法, 倘若 $\{f_n(x)\}$ 在 $[a,b]$ 上不一致收敛于 $f(x)$, 则 $\exists\varepsilon_0>0$, 对于 $\forall N\in\mathbb N$, 总存在 $n'>N$ 及 $x'\in[a,b]$, 使得
% \[
%   |f_{n'}(x')-f(x')|\geq\varepsilon_0.
% \]
% 因此存在趋于无穷的正整数序列 $\{n_k\}$ 及序列 $\{x_{n_k}\}\subset[a,b]$, 使得
% \[
%   |f_{n_k}(x_{n_k})-f(x_{n_k})|\geq\varepsilon_0.
% \]
% 我们不妨设 $x_{n_k}\to x_0\in[a,b](k\to\infty)$。由 $f_n(x_0)\to f(x_0)(n\to\infty)$, 对于 $\varepsilon=\dfrac{\varepsilon_0}{4}$, 存在 $N\in\mathbb N$, 当 $n>N$ 时, 有
% \[
%   |f_n(x_0)-f(x_0)|<\varepsilon.
% \]
% 取 $n_0=N+1$, 再由 $f_{n_0}(x)$ 及 $f(x)$ 的连续性, $\exists\delta>0$, 当 $x\in U(x_0,\delta)\cap[a,b]$ 时, 有
% \[
%   |f(x)-f(x_0)|<\varepsilon\quad\text{及}\quad |f_{n_0}(x)-f_{n_0}(x_0)|<\varepsilon.
% \]
% 现在取一个 $n_k>n_0$, 使得 $|x_{n_k}-x_0|<\delta$, 则有
% \[
%   f_{n_0}(x_{n_k})\leq f_{n_k}(x_{n_k})\leq f(x_{n_k}).
% \]
% 因此有
% \[
% \begin{aligned}
%   \varepsilon_0
%   &\leq |f_{n_k}(x_{n_k})-f(x_{n_k})|
%   \leq |f_{n_0}(x_{n_k})-f(x_{n_k})|\\
%   &\leq |f_{n_0}(x_{n_k})-f_{n_0}(x_0)|+|f_{n_0}(x_0)-f(x_0)|+|f(x_0)-f(x_{n_k})|\\
%   &<3\varepsilon=\frac34\varepsilon_0.
% \end{aligned}
% \]
% 此矛盾便证明了必要性.
% \end{proof}
% 
% 
% 由函数序列的狄尼定理, 我们容易推出下述关于函数项级数的狄尼定理.
% 
% \begin{theorem}[{狄尼定理}]\label{theorem:ma-10-4-2-prime}
% 设函数项级数 $\sum\limits_{n=1}^{+\infty}u_n(x)$ 在区间 $[a,b]$ 上收敛, 并且对于 $\forall n\in\mathbb N$, $u_n(x)$ 在 $[a,b]$ 上连续且非负, 则 $\sum\limits_{n=1}^{+\infty}u_n(x)$ 的和函数在 $[a,b]$ 上连续的充分必要条件是 $\sum\limits_{n=1}^{+\infty}u_n(x)$ 在 $[a,b]$ 上一致收敛.
% \end{theorem}
% 
% 
% \begin{remark}\label{remark:ma-10-source-57}
% 在狄尼定理中, 闭区间不能改为开区间或无穷区间, 请读者自己举例说明之.
% \end{remark}
% 
% 
% \begin{example}\label{example:ma-10-4-2}
% 证明 $\sum\limits_{n=1}^{+\infty}x^n\dfrac{\ln x}{1+\left|\ln\ln\dfrac1x\right|}$ 在 $(0,1)$ 内一致收敛.
% \end{example}
% 
% 
% \begin{proof}
% 对于 $n=1,2,\cdots$, 令
% \[
%   u_n(x)=x^n\frac{\ln x}{1+\left|\ln\ln\dfrac1x\right|}\quad(x\in(0,1)).
% \]
% 我们有 $\lim\limits_{x\to0+0}u_n(x)=0$ 和 $\lim\limits_{x\to1-0}u_n(x)=0$。记
% \[
%   \widetilde u_n(x)=
%   \begin{cases}
%     0, & x=0,1,\\
%     -u_n(x), & x\in(0,1),
%   \end{cases}
% \]
% 则 $\widetilde u_n(x)(n=1,2,\cdots)$ 在 $[0,1]$ 上连续、非负, 且
% \[
%   \sum_{k=1}^n\widetilde u_k(x)=
%   \begin{cases}
%     0, & x=0,1,\\
%     -\dfrac{x-x^{n+1}}{1-x}\dfrac{\ln x}{1+\left|\ln\ln\dfrac1x\right|}, & x\in(0,1).
%   \end{cases}
% \]
% 因此我们有
% \[
%   \sum_{n=1}^{+\infty}\widetilde u_n(x)=
%   \begin{cases}
%     0, & x=0,1,\\
%     -\dfrac{x}{1-x}\dfrac{\ln x}{1+\left|\ln\ln\dfrac1x\right|}, & x\in(0,1).
%   \end{cases}
% \]
% 容易看出它在 $[0,1]$ 上连续。由定理~\ref{theorem:ma-10-4-2-prime} 知 $\sum\limits_{n=1}^{+\infty}\widetilde u_n(x)$ 在 $[0,1]$ 一致收敛, 从而 $\sum\limits_{n=1}^{+\infty}u_n(x)$ 在 $(0,1)$ 内一致收敛, 即 $\sum\limits_{n=1}^{+\infty}x^n\dfrac{\ln x}{1+\left|\ln\ln\dfrac1x\right|}$ 在 $(0,1)$ 内一致收敛.
% 
% 读者可以验证, 例~\ref{example:ma-10-4-2} 中的函数项级数不能用魏尔斯特拉斯 M 判别法来证明它是一致收敛的.
% \end{proof}
% 
% 
% \subsection{极限函数的积分}
% \label{subsec:ma-10-4-2}
% 
% 对于 $[a,b]$ 上可积函数序列 $\{f_n(x)\}$, 设它在 $[a,b]$ 上一致收敛到其极限函数 $f(x)$, 我们来研究 $f(x)$ 在 $[a,b]$ 上的可积性。我们有以下的定理.
% 
% \begin{theorem}\label{theorem:ma-10-4-3}
% 设函数 $f_n(x)(n=1,2,\cdots)$ 在 $[a,b]$ 上可积, 且 $f_n(x)\rightrightarrows f(x)(x\in[a,b])$, 则 $f(x)$ 在 $[a,b]$ 上可积, 并且成立
% \[
%   \lim_{n\to\infty}\int_a^b f_n(x)\mathrm dx
%   =\int_a^b f(x)\mathrm dx
%   =\int_a^b\lim_{n\to\infty}f_n(x)\mathrm dx.
% \]
% \end{theorem}
% 
% 
% \begin{proof}
% 对于 $\forall\varepsilon>0$, 由 $f_n(x)\rightrightarrows f(x)(x\in[a,b])$, 存在 $N\in\mathbb N$, 当 $n>N$ 时, 对于 $\forall x\in[a,b]$, 有
% \begin{equation}
%   |f_n(x)-f(x)|<\frac{\varepsilon}{3(b-a)}.
%   \label{eq:ma-10-4-1}
% \end{equation}
% 特别地, 取 $n_0=N+1$, 则 $f_{n_0}(x)$ 满足 \eqref{eq:ma-10-4-1} 式。由 $f_{n_0}(x)$ 在 $[a,b]$ 上可积, 从而存在 $[a,b]$ 的分割
% \[
%   a=x_0<x_1<\cdots<x_n=b,
% \]
% 若记 $\Delta x_i=x_i-x_{i-1}(i=1,2,\cdots,n)$, 则有
% \[
%   \sum_{i=1}^n \omega_i(f_{n_0})\Delta x_i<\frac\varepsilon3,
% \]
% 其中我们用 $\omega_i(g)$ 表示函数 $g(x)$ 在 $[x_{i-1},x_i](i=1,2,\cdots,n)$ 上的振幅.
% 
% 对 $x',x''\in[x_{i-1},x_i]$, 由 \eqref{eq:ma-10-4-1}, 我们有
% \[
% \begin{aligned}
%   |f(x')-f(x'')|
%   &\leq |f(x')-f_{n_0}(x')|+|f_{n_0}(x')-f_{n_0}(x'')|+|f_{n_0}(x'')-f(x'')|.
% \end{aligned}
% \]
% 由上式我们推出
% \[
%   \omega_i(f)\leq\omega_i(f_{n_0})+\frac{2\varepsilon}{3(b-a)}.
% \]
% 因此我们有
% \[
% \begin{aligned}
%   \sum_{i=1}^n \omega_i(f)\Delta x_i
%   &\leq \sum_{i=1}^n \omega_i(f_{n_0})\Delta x_i+\frac{2\varepsilon(b-a)}{3(b-a)}\\
%   &\leq \frac\varepsilon3+\frac{2\varepsilon}{3}=\varepsilon.
% \end{aligned}
% \]
% 这就证明了 $f(x)$ 在 $[a,b]$ 上可积。再由 \eqref{eq:ma-10-4-1}, 当 $n>N$ 时, 有
% \[
%   \left|\int_a^b f_n(x)\mathrm dx-\int_a^b f(x)\mathrm dx\right|
%   \leq \int_a^b |f_n(x)-f(x)|\mathrm dx
%   \leq \frac{\varepsilon}{3(b-a)}(b-a)<\varepsilon,
% \]
% 因此有
% \[
%   \lim_{n\to\infty}\int_a^b f_n(x)\mathrm dx=\int_a^b f(x)\mathrm dx.
% \qedhere
% \]
% \end{proof}
% 
% 
% 同理可证下面的定理.
% 
% \begin{theorem}\label{theorem:ma-10-4-3-prime}
% 设函数 $u_n(x)(n=1,2,\cdots)$ 在区间 $[a,b]$ 上可积, 且 $\sum\limits_{n=1}^{+\infty}u_n(x)$ 在 $[a,b]$ 上一致收敛, 则 $\sum\limits_{n=1}^{+\infty}u_n(x)$ 的和函数在 $[a,b]$ 上可积, 并且成立
% \begin{equation}
%   \int_a^b\left(\sum_{n=1}^{+\infty}u_n(x)\right)\mathrm dx
%   =\sum_{n=1}^{+\infty}\int_a^b u_n(x)\mathrm dx.
%   \label{eq:ma-10-4-2}
% \end{equation}
% \end{theorem}
% 
% 为了方便起见, 当 \eqref{eq:ma-10-4-2} 成立时, 我们也称函数项级数 $\sum\limits_{n=1}^{+\infty}u_n(x)$ 可在 $[a,b]$ 上逐项积分.
% 
% \begin{example}\label{example:ma-10-4-3}
% 证明: 当 $x\in(-1,1)$ 时, 成立以下等式:
% \[
%   \frac12\ln\frac{1+x}{1-x}
%   =\sum_{n=0}^{+\infty}\frac{x^{2n+1}}{2n+1}.
% \]
% \end{example}
% 
% 
% \begin{proof}
% 设 $f(x)=\dfrac1{1-x^2}$, 则当 $|x|<1$ 时, $f(x)$ 是以下几何级数的和函数:
% \[
%   f(x)=\frac1{1-x^2}=\sum_{n=0}^{+\infty}x^{2n}.
% \]
% 由于其部分和序列
% \[
%   S_n(x)=\sum_{k=0}^n x^{2k}=\frac{1-x^{2n+2}}{1-x^2}\quad(n=1,2,\cdots)
% \]
% 在 $(-1,1)$ 内闭一致收敛到 $\dfrac1{1-x^2}$, 因此对固定的 $x\in(-1,1)$, 该函数项级数可在 $[0,x](x>0)$ 或 $[x,0](x<0)$ 上逐项积分。因此有
% \[
%   \frac12\ln\frac{1+x}{1-x}
%   =\int_0^x\frac{\mathrm dt}{1-t^2}
%   =\sum_{n=0}^{+\infty}\int_0^x t^{2n}\mathrm dt
%   =\sum_{n=0}^{+\infty}\frac{x^{2n+1}}{2n+1}.
% \qedhere
% \]
% \end{proof}
% 
% 
% \begin{example}\label{example:ma-10-4-4}
% 求 $I=\int_0^1\sum\limits_{n=1}^{+\infty}\dfrac{x}{n(n+x)}\mathrm dx$ 的值.
% \end{example}
% 
% 
% \begin{solve}
% 对于 $\forall x\in[0,1]$ 及 $n=1,2,\cdots$, 有
% \[
%   0\leq\frac{x}{n(n+x)}\leq\frac1{n^2},
% \]
% 因此 $\sum\limits_{n=1}^{+\infty}\dfrac{x}{n(n+x)}$ 在 $[0,1]$ 上一致收敛。由定理~\ref{theorem:ma-10-4-3-prime}, 我们可以对该函数项级数逐项积分
% \[
% \begin{aligned}
%   I
%   &=\sum_{n=1}^{+\infty}\int_0^1\frac{x}{n(n+x)}\mathrm dx
%    =\sum_{n=1}^{+\infty}\int_0^1\left(\frac1n-\frac1{n+x}\right)\mathrm dx\\
%   &=\sum_{n=1}^{+\infty}\left[\frac1n-\ln(1+n)+\ln n\right]
%    =\lim_{n\to\infty}\left[\sum_{k=1}^n\frac1k-\ln(1+n)\right]\\
%   &=c\text{(欧拉常数)}.
% \end{aligned}
% \qedhere
% \]
% \end{solve}
% 
% 
% \subsection{极限函数的导数}
% \label{subsec:ma-10-4-3}
% 
% 在本小节中, 我们研究函数序列（函数项级数）逐项求导组成的函数序列（函数项级数）的敛散性问题。在这里我们注意到以下的一个事实: 即使一个函数序列 $\{f_n(x)\}$ 中每个函数在区间 $I$ 内都有连续的导函数 $f_n'(x)$, 且 $\{f_n'(x)\}$ 在区间 $I$ 上一致收敛, 也不能保证 $\{f_n(x)\}$ 在区间 $I$ 上一致收敛, 甚至不能保证 $\{f_n(x)\}$ 在区间 $I$ 上的点收敛性。这是因为一个函数的原函数可相差一个常数的缘故。要使得相应的 $\{f_n(x)\}$ 收敛, 则必须要求 $\{f_n(x)\}$ 至少在 $I$ 上的一个点上是收敛的。我们有以下定理.
% 
% \begin{theorem}\label{theorem:ma-10-4-4}
% 设函数 $f_n(x)(n=1,2,\cdots)$ 在区间 $[a,b]$ 上可微, 且满足
% 
% (a) 存在 $x_0\in[a,b]$, 使得 $\lim\limits_{n\to\infty}f_n(x_0)$ 存在;
% 
% (b) $f_n'(x)\rightrightarrows g(x)(x\in[a,b])$,
% 
% 则我们有以下结论:
% 
% (1) 存在 $[a,b]$ 上的函数 $f(x)$, 使得 $f_n(x)\rightrightarrows f(x)(x\in[a,b])$;
% 
% (2) $f(x)$ 在 $[a,b]$ 上可微（端点单侧可微）, 并且 $f'(x)=g(x)$, 即
% \[
%   \lim_{n\to\infty}f_n'(x)=\left[\lim_{n\to\infty}f_n(x)\right]'.
% \]
% \end{theorem}
% 
% 
% \begin{proof}
% 我们主要利用柯西准则来证明上述结论。对于 $\forall\varepsilon>0$, 由 $\lim\limits_{n\to\infty}f_n(x_0)$ 存在和 $f_n'(x)\rightrightarrows g(x)(x\in[a,b])$ 知, $\exists N\in\mathbb N$, 当 $n>m\geq N$ 时, 有
% \begin{equation}
%   |f_n(x_0)-f_m(x_0)|<\frac\varepsilon2;
%   \label{eq:ma-10-4-3}
% \end{equation}
% 并且对 $\forall x\in[a,b]$, 有
% \begin{equation}
%   |f_n'(x)-f_m'(x)|<\frac{\varepsilon}{2(b-a)}.
%   \label{eq:ma-10-4-4}
% \end{equation}
% 对任意 $n>m\geq N$, 对 $f_n(x)-f_m(x)$ 在 $x_0$ 与 $x$ 之间应用拉格朗日微分中值定理知: 在 $x_0$ 与 $x$ 之间存在 $\xi$, 使得
% \begin{equation}
% \begin{aligned}
%   &\left|[f_n(x)-f_m(x)]-[f_n(x_0)-f_m(x_0)]\right|\\
%   &\qquad =|f_n'(\xi)-f_m'(\xi)||x-x_0|
%    <\frac{|x-x_0|\varepsilon}{2(b-a)}\leq\frac\varepsilon2.
% \end{aligned}
%   \label{eq:ma-10-4-5}
% \end{equation}
% 因此对于 $\forall x\in[a,b]$, 由 \eqref{eq:ma-10-4-3} 和 \eqref{eq:ma-10-4-5} 得
% \[
% \begin{aligned}
%   |f_n(x)-f_m(x)|
%   &\leq \left|[f_n(x)-f_m(x)]-[f_n(x_0)-f_m(x_0)]\right|
%       +|f_n(x_0)-f_m(x_0)|\\
%   &\leq\frac\varepsilon2+\frac\varepsilon2=\varepsilon.
% \end{aligned}
% \]
% 由柯西准则知 $\{f_n(x)\}$ 在 $[a,b]$ 上一致收敛, 从而存在 $[a,b]$ 上的函数 $f(x)$, 使得 $f_n(x)\rightrightarrows f(x)(x\in[a,b])$, 即 (1) 得证.
% 
% 为了证明 (2), 对于 $\forall x^*\in[a,b]$, 对 $n=1,2,\cdots$, 令
% \[
%   h_n(x)=\frac{f_n(x)-f_n(x^*)}{x-x^*}\ (x\ne x^*),\qquad h_n(x^*)=f_n'(x^*),
% \]
% \[
%   h(x)=\frac{f(x)-f(x^*)}{x-x^*}\quad (x\ne x^*),
% \]
% 则 $h_n(x)$ 在 $[a,b]$ 上连续。现在我们证明 $\{h_n(x)\}$ 在 $[a,b]\setminus\{x^*\}$ 上一致收敛到 $h(x)$.
% 
% 事实上, 从 \eqref{eq:ma-10-4-5} 的证明过程中可以推出, 对于 $\forall\varepsilon>0$, $\exists N\in\mathbb N$, 当 $n>m\geq N$ 时, 对于 $\forall x\in[a,b]\setminus\{x^*\}$, 有
% \[
% \begin{aligned}
%   |h_n(x)-h_m(x)|
%   &\leq \frac1{|x-x^*|}\left|[f_n(x)-f_m(x)]-[f_n(x^*)-f_m(x^*)]\right|\\
%   &\leq\frac{\varepsilon}{2(b-a)}.
% \end{aligned}
% \]
% 因此 $h_n(x)$ 在 $[a,b]\setminus\{x^*\}$ 上一致收敛到 $h(x)$。注意到 $\lim\limits_{x\to x^*}h_n(x)=f_n'(x^*)$, 由定理~\ref{theorem:ma-10-4-1} 后的注, 我们有 $f'(x^*)=\lim\limits_{x\to x^*}h(x)$ 存在, 并且 $f'(x^*)=\lim\limits_{n\to\infty}f_n'(x^*)$。由 $x^*\in[a,b]$ 的任意性, 因此 (2) 成立.
% \end{proof}
% 
% 
% 对于函数项级数, 我们有
% 
% \begin{theorem}\label{theorem:ma-10-4-4-prime}
% 设函数 $u_n(x)(n=1,2,\cdots)$ 在区间 $[a,b]$ 上可微, 且满足:
% 
% (a) 存在 $x_0\in[a,b]$, 使得 $\sum\limits_{n=1}^{+\infty}u_n(x_0)$ 收敛;
% 
% (b) $\sum\limits_{n=1}^{+\infty}u_n'(x)$ 在 $[a,b]$ 上一致收敛,
% 
% 则
% 
% (1) $\sum\limits_{n=1}^{+\infty}u_n(x)$ 在 $[a,b]$ 上一致收敛;
% 
% (2) $\sum\limits_{n=1}^{+\infty}u_n(x)$ 的和函数在 $[a,b]$ 上可导, 并且
% \begin{equation}
%   \left(\sum_{n=1}^{+\infty}u_n(x)\right)'
%   =\sum_{n=1}^{+\infty}u_n'(x).
%   \label{eq:ma-10-4-6}
% \end{equation}
% \end{theorem}
% 
% 同样为了便于叙述, 当 \eqref{eq:ma-10-4-6} 成立时, 我们称函数项级数 $\sum\limits_{n=1}^{+\infty}u_n(x)$ 可在 $[a,b]$ 上逐项求导数（或逐项可微）.
% 
% \begin{remark}\label{remark:ma-10-source-70}
% 由于函数可微也是一种局部性质, 因此要讨论函数序列的极限函数的可微性只需要求该函数序列点点收敛, 其导函数序列内闭一致收敛即可。显然, 对于函数项级数也有类似的结论.
% \end{remark}
% 
% 
% \begin{example}\label{example:ma-10-4-5}
% 证明 $\sum\limits_{n=1}^{+\infty}\dfrac1{n^x}$ 在 $(1,+\infty)$ 上无穷次可微.
% \end{example}
% 
% 
% \begin{proof}
% 任取 $\delta_0>0$, 我们已经知道 $\sum\limits_{n=1}^{+\infty}\dfrac1{n^x}$ 在 $(1+\delta_0,+\infty)$ 上一致收敛。由于
% \[
%   \left(\frac1{n^x}\right)'=\frac{-\ln n}{n^x}\in C(0,+\infty),
% \]
% 并且对于 $\forall x\in(1+\delta_0,+\infty)$, 有
% \[
%   \left|\frac{-\ln n}{n^x}\right|<\frac{\ln n}{n^{1+\delta_0}},
% \]
% 注意到 $\sum\limits_{n=1}^{+\infty}\dfrac{\ln n}{n^{1+\delta_0}}$ 收敛, 因此 $\sum\limits_{n=1}^{+\infty}\dfrac{-\ln n}{n^x}$ 在 $(1+\delta_0,+\infty)$ 上一致收敛.
% 
% 由定理~\ref{theorem:ma-10-4-4-prime} 知 $\sum\limits_{n=1}^{+\infty}\dfrac1{n^x}$ 的和函数在 $(1+\delta_0,+\infty)$ 上有连续的导函数, 且
% \[
%   \left(\sum_{n=1}^{+\infty}\frac1{n^x}\right)'
%   =\sum_{n=1}^{+\infty}\frac{-\ln n}{n^x}.
% \]
% 由于 $\sum\limits_{n=1}^{+\infty}\dfrac{-\ln n}{n^x}$ 逐项求导后得到的函数项级数 $\sum\limits_{n=1}^{+\infty}\dfrac{\ln^2 n}{n^x}$ 仍然在 $(1+\delta_0,+\infty)$ 上一致收敛, 所以 $\sum\limits_{n=1}^{+\infty}\dfrac1{n^x}$ 的和函数在 $(1+\delta_0,+\infty)$ 上有连续的二阶导数。依次下去知, $\sum\limits_{n=1}^{+\infty}\dfrac1{n^x}$ 的和函数在 $(1+\delta_0,+\infty)$ 上具有任意阶导数。再由 $\delta_0$ 的任意性, 知 $\sum\limits_{n=1}^{+\infty}\dfrac1{n^x}$ 的和函数在 $(1,+\infty)$ 上无穷次可微.
% 
% 利用函数项级数, 我们可以构造许多具有特殊性质的函数.
% \end{proof}
% 
% 
% \begin{example}\label{example:ma-10-4-6}
% 设 $\{x_n\}\subset[0,1]$ 是一个各项互不相同的序列, 证明 $\sum\limits_{n=1}^{+\infty}\dfrac{\operatorname{sgn}(x-x_n)}{2^n}$ 的和函数当且仅当 $x=x_n(n=1,2,\cdots)$ 时不连续.
% \end{example}
% 
% 
% \begin{proof}
% 首先注意到, 对任意的 $n$ 及 $x\in[0,1]$, 有
% \[
%   \left|\frac{\operatorname{sgn}(x-x_n)}{2^n}\right|\leq\frac1{2^n}.
% \]
% 由于 $\sum\limits_{n=1}^{+\infty}\dfrac1{2^n}$ 收敛, 因此 $\sum\limits_{n=1}^{+\infty}\dfrac{\operatorname{sgn}(x-x_n)}{2^n}$ 在 $[0,1]$ 上一致收敛。对于 $\forall k\in\mathbb N$, 该函数项级数中仅有一项 $\dfrac{\operatorname{sgn}(x-x_k)}{2^k}$ 在 $x=x_k$ 处不连续, 而其他项均在 $x=x_k$ 处连续。因此 $\sum\limits_{n=1}^{+\infty}\dfrac{\operatorname{sgn}(x-x_n)}{2^n}$ 的和函数在 $x=x_k$ 处不连续.
% 
% 当 $x_0\notin\{x_n\}$ 时, 由于该函数项级数每一项都在 $x=x_0$ 处连续, 从而 $\sum\limits_{n=1}^{+\infty}\dfrac{\operatorname{sgn}(x-x_n)}{2^n}$ 的和函数在 $x=x_0$ 处连续.
% \end{proof}
% 
% 
% \begin{remark}\label{remark:ma-10-source-75}
% 若在上例中, 将符号函数 $\operatorname{sgn}(x)$ 换为
% \[
%   f(x)=
%   \begin{cases}
%     x^2\sin\dfrac1x, & x\ne0,\\
%     0, & x=0,
%   \end{cases}
% \]
% 作
% \[
%   g(x)=\sum_{n=1}^{+\infty}\frac{f(x-x_n)}{2^n},
% \]
% 则可证明 $g(x)$ 在 $[0,1]$ 上可导, 但 $g'(x)$ 在 $x=x_n\in[0,1](n=1,2,\cdots)$ 处不连续, 而在 $[0,1]$ 中的其他点连续.
% \end{remark}
% 
% 历史上有许多人曾经以为一个连续函数仅有可能在一列点上是不可导的。魏尔斯特拉斯利用函数项级数首先构造出了一个处处不可导的连续函数。后来, 人们借助函数项级数构造了一些更为简单的处处不可导的连续函数的例子, 如 V。D。Waerden 构造了以下例子:
% 
% 记函数 $\widetilde f(x)=|x|, x\in\left[-\dfrac12,\dfrac12\right]$, 然后记将 $\widetilde f(x)$ 以 $1$ 为周期延拓到 $(-\infty,+\infty)$ 上的函数为 $f(x)$, 则 $\sum\limits_{n=1}^{+\infty}\dfrac{f(4^n x)}{4^n}$ 的和函数是 $(-\infty,+\infty)$ 上的一个连续函数, 但它在 $(-\infty,+\infty)$ 上处处不可导.
% 
% 由于这个事实的证明较为复杂, 我们在这里省略之.
% 
% \section*{习题十}
% \phantomsection
% \addcontentsline{toc}{section}{习题十}
% \label{exercises:ma-10}
% 
% 1。求出下列函数项级数的收敛域:
% \[
%   (1)\ \sum_{n=1}^{+\infty}\frac{n}{n+1}\left(\frac{x}{2x+1}\right)^n;\qquad
%   (2)\ \sum_{n=1}^{+\infty}\frac{x^n}{1-x^n}.
% \]
% 
% 2。求下列函数序列 $\{f_n(x)\}$ 的极限函数, 并讨论在给定的区间上是否一致收敛:
% \[
% \begin{gathered}
%   (1)\ f_n(x)=x(1-x)^n,\ x\in[0,1].\\
%   (2)\ f_n(x)=nxe^{-nx^2},\ x\in(-\infty,+\infty).\\
%   (3)\ f_n(x)=n^2e^{-nx^2},\ x\in[a,+\infty),\ \text{这里 }a>0.\\
%   (4)\ f_n(x)=\frac{x^2}{x^2+(1-nx)^2},\ x\in[0,1].\\
%   (5)\ f_n(x)=n^2x(1-x^2)^n,\ x\in[0,1].\\
%   (6)\ f_n(x)=\sin\frac{x}{n^n};\ \text{(a) }x\in[a,b];\ \text{(b) }x\in(-\infty,+\infty).\\
%   (7)\ f_n(x)=\frac{\sin n^n x}{n^\alpha},\ \alpha>0,\ x\in(-\infty,+\infty).\\
%   (8)\ f_n(x)=(\sin x)^{n^\alpha},\ \alpha>0,\ x\in[0,\pi].
% \end{gathered}
% \]
% 
% 3。讨论下列函数序列或函数项级数在指定区间的一致收敛性.
% \[
% \begin{gathered}
%   (1)\ \left\{f_n(x)=x^{\sum_{k=1}^n\frac1{2^k}}\right\},\ x\in(0,1);\\
%   (2)\ \sum_{n=1}^{+\infty}\frac1{n^2+x^2},\ x\in(-\infty,+\infty);\\
%   (3)\ \sum_{n=1}^{+\infty}\frac{\sin x\sin nx}{\sqrt{n+x^2}},\ x\in[0,+\infty);\\
%   (4)\ \sum_{n=1}^{+\infty}\frac{(-1)^n}{n+x^2},\ x\in(-\infty,+\infty);
% \end{gathered}
% \]
% 
% % 整合来源：math-11.tex
% \chapter{幂级数}
% \label{chap:ma-11}
% 
% 形如
% \[
%   \sum_{n=0}^{+\infty}a_n(x-x_0)^n
%   \quad(\text{其中 }x_0,a_n\in\mathbb R)
% \]
% 的函数项级数称为\keyterm{幂级数}.
% 
% 对任意的正整数 $n$, 这类级数的前 $n$ 项部分和
% \[
%   S_n=\sum_{k=0}^{n}a_k(x-x_0)^k
% \]
% 是一个至多为 $n$ 次的多项式。因此从某种意义上说, 幂级数是一种最简单的函数项级数。在函数项级数中, 无论在理论上还是实际应用中, 幂级数都是最重要的函数项级数之一。在许多后续课程学习过程中, 我们还将常常遇到幂级数。一个函数 $f(x)$ 若能展开成幂级数, 该函数也称为实解析的。我们将会看到实解析函数比 $C^\infty$ 函数（无穷次可微的函数）具有更好的性质。在本章中, 我们将系统介绍幂级数的理论.
% 
% \section{幂级数的收敛半径与收敛域}
% \label{sec:ma-11-1}
% 
% \subsection{幂级数的收敛半径与收敛域}
% \label{subsec:ma-11-1-1}
% 
% 对幂级数
% \begin{equation}
%   \sum_{n=0}^{+\infty}a_n(x-x_0)^n
%   =a_0+a_1(x-x_0)+a_2(x-x_0)^2+\cdots ,
%   \label{eq:ma-11-1-1}
% \end{equation}
% 作变量替换 $t=x-x_0$, 则该幂级数可化为
% \[
%   \sum_{n=0}^{+\infty}a_nt^n=a_0+a_1t+a_2t^2+\cdots ,
% \]
% 因此我们不妨在 \eqref{eq:ma-11-1-1} 中假定 $x_0=0$。所以我们将主要讨论形如
% \begin{equation}
%   \sum_{n=0}^{+\infty}a_nx^n=a_0+a_1x+a_2x^2+\cdots
%   \label{eq:ma-11-1-2}
% \end{equation}
% 的幂级数。显然当 $x=0$ 时, \eqref{eq:ma-11-1-2} 是收敛的.
% 
% \begin{example}\label{example:ma-11-1-1}
% 求幂级数
% \[
%   \sum_{n=0}^{+\infty}n^nx^n
% \]
% 的收敛域.
% \end{example}
% 
% 
% \begin{solve}
% 当 $x=0$ 时, 由于
% \[
%   \sum_{n=0}^{+\infty}n^nx^n=0^0=1,
% \]
% 因此该幂级数是收敛的。当 $x\ne0$ 时, 由于
% \[
%   \lim_{n\to\infty}|n^nx^n|=+\infty,
% \]
% 从而 $\sum_{n=0}^{+\infty}n^nx^n$ 发散。因此 $\sum_{n=0}^{+\infty}n^nx^n$ 的收敛域为 $\{0\}$.
% \end{solve}
% 
% 
% \begin{example}\label{example:ma-11-1-2}
% 求幂级数
% \[
%   \sum_{n=0}^{+\infty}x^n
% \]
% 的收敛域.
% \end{example}
% 
% 
% \begin{solve}
% 由几何级数的性质知, 当 $|x|<1$ 时, 该幂级数是收敛的。当 $|x|\geqslant1$ 时, 由
% \[
%   \lim_{n\to\infty}x^n\ne0,
% \]
% 推知该幂级数发散。因此 $\sum_{n=0}^{+\infty}x^n$ 的收敛域为 $(-1,1)$.
% \end{solve}
% 
% 
% \begin{example}\label{example:ma-11-1-3}
% 求幂级数
% \[
%   \sum_{n=1}^{+\infty}\frac{x^n}{n}
% \]
% 的收敛域.
% \end{example}
% 
% 
% \begin{solve}
% 当 $|x|<1$ 时, 由
% \[
%   \left|\frac{x^n}{n}\right|\leqslant |x|^n,
% \]
% 知该幂级数收敛。当 $x=-1$ 时, 它是一个交错级数, 并且通项趋于零, 由此可知此时该幂级数是收敛的。当 $x=1$ 时, 它是调和级数, 因此该幂级数发散。当 $|x|>1$ 时, 由于
% \[
%   \lim_{n\to\infty}\frac{x^n}{n}=\infty,
% \]
% 从而该幂级数发散。因此 $\sum_{n=1}^{+\infty}\frac{x^n}{n}$ 的收敛域为 $[-1,1)$.
% \end{solve}
% 
% 
% \begin{example}\label{example:ma-11-1-4}
% 求幂级数
% \[
%   \sum_{n=1}^{+\infty}\frac{x^n}{n^2}
% \]
% 的收敛域.
% \end{example}
% 
% 
% \begin{solve}
% 如同上题, 易知当 $|x|<1$ 时, 该幂级数收敛; 当 $|x|>1$ 时, 该幂级数发散; 而当 $x=\pm1$ 时, 显然该幂级数收敛。因此 $\sum_{n=1}^{+\infty}\frac{x^n}{n^2}$ 的收敛域为 $[-1,1]$.
% \end{solve}
% 
% 
% \begin{example}\label{example:ma-11-1-5}
% 求幂级数
% \[
%   \sum_{n=0}^{+\infty}\frac{x^n}{n!}
% \]
% 的收敛域.
% \end{example}
% 
% 
% \begin{solve}
% 对于 $\forall x\ne0$, 我们有
% \[
%   \lim_{n\to\infty}
%   \frac{|x|^{n+1}}{(n+1)!}\Big/\frac{|x|^n}{n!}
%   =
%   \lim_{n\to\infty}\frac{|x|}{n+1}=0.
% \]
% 由此推知, 该幂级数在 $x$ 处收敛。由于 $x$ 是 $(-\infty,+\infty)$ 上任意一点, 因此
% \[
%   \sum_{n=0}^{+\infty}\frac{x^n}{n!}
% \]
% 的收敛域为 $(-\infty,+\infty)$.
% 
% 以上例子说明了幂级数具有各种不同的收敛域, 是否所有幂级数的收敛域一定具有上面列举的收敛域的形式之一呢? 为了证实这一点, 我们先来证明以下关于幂级数收敛域的一个基本结果.
% \end{solve}
% 
% 
% \begin{theorem}\label{theorem:ma-11-1-1}
% 设幂级数
% \[
%   \sum_{n=0}^{+\infty}a_nx^n
% \]
% 在 $x_0\ne0$ 处收敛, 则该级数在 $(-|x_0|,|x_0|)$ 内闭绝对一致收敛.
% \end{theorem}
% 
% 
% \begin{proof}
% 任取闭区间 $[a,b]\subset(-|x_0|,|x_0|)$, 则存在常数 $0<\delta_0<1$（依赖 $[a,b]$）, 使得
% \[
%   [a,b]\subset[-\delta_0|x_0|,\delta_0|x_0|]\subset(-|x_0|,|x_0|).
% \]
% 下面我们证明 $\sum_{n=0}^{+\infty}a_nx^n$ 在 $[-\delta_0|x_0|,\delta_0|x_0|]$ 上绝对一致收敛.
% 
% 由于 $\sum_{n=0}^{+\infty}a_nx_0^n$ 收敛, 我们有 $|a_nx_0^n|\to0\ (n\to\infty)$, 因此存在 $M>0$, 使得对于 $\forall n\geqslant0$, 有
% \[
%   |a_nx_0^n|\leqslant M.
% \]
% 所以对于 $\forall x\in[-\delta_0|x_0|,\delta_0|x_0|]$, 有
% \[
%   |a_nx^n|
%   =
%   \left|a_nx_0^n\frac{x^n}{x_0^n}\right|
%   \leqslant M\delta_0^n.
% \]
% 由于 $\sum_{n=0}^{+\infty}M\delta_0^n$ 收敛, 因此 $\sum_{n=0}^{+\infty}|a_nx^n|$ 在 $[-\delta_0|x_0|,\delta_0|x_0|]$ 上一致收敛.
% \end{proof}
% 
% 
% 由定理~\ref{theorem:ma-11-1-1} 可知, 若 $\sum_{n=0}^{+\infty}a_nx^n$ 在点 $x_0\ne0$ 处收敛而在点 $x_1\ne0$ 处发散时, 则必存在开区间 $(-R,R)$, $|x_0|\leqslant R\leqslant |x_1|$, 使得它在 $(-R,R)$ 内收敛, 而在 $[-R,R]$ 外面发散。事实上, 令
% \[
%   R=\sup\left\{|x|:\sum_{n=0}^{+\infty}a_nx^n\text{收敛}\right\},
% \]
% 则由定理~\ref{theorem:ma-11-1-1} 有 $R\geqslant |x_0|$, 但 $R\leqslant |x_1|$。这是因为倘若 $R>|x_1|$, 则存在 $|x_2|>|x_1|$, 使得 $\sum_{n=0}^{+\infty}a_nx_2^n$ 收敛。由定理~\ref{theorem:ma-11-1-1} 知 $\sum_{n=0}^{+\infty}a_nx_1^n$ 也必收敛, 这与假设矛盾.
% 
% 容易看出 $\sum_{n=0}^{+\infty}a_nx^n$ 在 $(-R,R)$ 内闭绝对一致收敛, 而在 $(-\infty,-R)$ 和 $(R,+\infty)$ 上发散, 因此我们称 $R$ 是幂级数 $\sum_{n=0}^{+\infty}a_nx^n$ 的收敛半径.
% 
% 从上面讨论可知, 此时 $\sum_{n=0}^{+\infty}a_nx^n$ 的收敛域必包含 $(-R,R)$。该幂级数在 $x=\pm R$ 是否收敛, 则没有一般性的结论, 需对每个幂级数分别进行讨论。从前面的例子可知各种情形均可发生, 读者应当注意到这一点.
% 
% 另外, 当 $\sum_{n=0}^{+\infty}a_nx^n$ 仅在 $x=0$ 处收敛时, 规定它的收敛半径为 $0$。而当 $\sum_{n=0}^{+\infty}a_nx^n$ 处处收敛时, 则自然规定它的收敛半径为 $+\infty$, 这时它的收敛域为 $(-\infty,+\infty)$.
% 
% \subsection{收敛半径的求法}
% \label{subsec:ma-11-1-2}
% 
% 在某种意义下, 幂级数可以看成是几何级数的一种推广。因此有关几何级数的一些敛散性判别法则对幂级数也是适用的.
% 
% \begin{theorem}[{柯西--哈达玛定理}]\label{theorem:ma-11-1-2}
% 设幂级数
% \[
%   \sum_{n=0}^{+\infty}a_nx^n
% \]
% 的收敛半径为 $R$, 记
% \[
%   \rho=\varlimsup_{n\to\infty}\sqrt[n]{|a_n|},
% \]
% 则
% \begin{enumerate}
%   \item 当 $\rho=+\infty$ 时, $R=0$;
%   \item 当 $\rho=0$ 时, $R=+\infty$;
%   \item 当 $0<\rho<+\infty$ 时, $R=\dfrac1\rho$.
% \end{enumerate}
% \end{theorem}
% 
% 
% \begin{proof}
% 在定理的条件下, 对于固定的 $x\in\mathbb R$, 有
% \[
%   \varlimsup_{n\to\infty}\sqrt[n]{|a_nx^n|}
%   =
%   \rho|x|.
% \]
% 因此由数项级数的柯西判别法知, 当 $\rho|x|<1$ 时, $\sum_{n=0}^{+\infty}a_nx^n$ 收敛; 而当 $\rho|x|>1$ 时, $\sum_{n=0}^{+\infty}a_nx^n$ 发散。所以我们有
% \begin{enumerate}
%   \item 当 $\rho=+\infty$ 时, $\sum_{n=0}^{+\infty}a_nx^n$ 仅当 $x=0$ 时收敛, 因此 $R=0$;
%   \item 当 $\rho=0$ 时, $\sum_{n=0}^{+\infty}a_nx^n$ 对一切 $x\in(-\infty,+\infty)$ 收敛, 因此 $R=+\infty$;
%   \item 当 $0<\rho<+\infty$ 时, $\sum_{n=0}^{+\infty}a_nx^n$ 在 $|x|<\dfrac1\rho$ 时收敛, 而在 $|x|>\dfrac1\rho$ 时发散, 因此 $R=\dfrac1\rho$。
% \end{enumerate}
% \end{proof}
% 
% 
% 有时候, 用以下定理计算收敛半径要方便些.
% 
% \begin{theorem}\label{theorem:ma-11-1-3}
% 设幂级数
% \[
%   \sum_{n=0}^{+\infty}a_nx^n
% \]
% 的收敛半径为 $R$。若
% \[
%   \lim_{n\to\infty}\left|\frac{a_{n+1}}{a_n}\right|=\rho,
% \]
% 则
% \begin{enumerate}
%   \item 当 $\rho=+\infty$ 时, $R=0$;
%   \item 当 $\rho=0$ 时, $R=+\infty$;
%   \item 当 $0<\rho<+\infty$ 时, $R=\dfrac1\rho$.
% \end{enumerate}
% \end{theorem}
% 
% 
% \begin{proof}
% 由于
% \[
%   \underline{\lim_{n\to\infty}}\left|\frac{a_{n+1}}{a_n}\right|
%   \leqslant
%   \underline{\lim_{n\to\infty}}\sqrt[n]{|a_n|}
%   \leqslant
%   \overline{\lim_{n\to\infty}}\sqrt[n]{|a_n|}
%   \leqslant
%   \overline{\lim_{n\to\infty}}\left|\frac{a_{n+1}}{a_n}\right|,
% \]
% 因此在定理的条件下, 必有
% \[
%   \lim_{n\to\infty}\sqrt[n]{|a_n|}=\rho.
% \]
% 所以, 由定理~\ref{theorem:ma-11-1-2} 即可推出定理~\ref{theorem:ma-11-1-3}.
% \end{proof}
% 
% 
% \paragraph*{注 1}
% 对于幂级数
% \[
%   \sum_{n=0}^{+\infty}a_n(x-x_0)^n,
% \]
% 若
% \[
%   \overline{\lim_{n\to\infty}}\sqrt[n]{|a_n|}=\rho,
% \]
% 则它的收敛半径仍然为
% \[
%   R=
%   \begin{cases}
%     0, & \rho=+\infty,\\
%     +\infty, & \rho=0,\\
%     \dfrac1\rho, & 0<\rho<+\infty.
%   \end{cases}
% \]
% 值得指出的是, 此时幂级数 $\sum_{n=0}^{+\infty}a_n(x-x_0)^n$ 的收敛域则是以 $x_0$ 为中心且包含 $(x_0-R,x_0+R)$ 的区间（端点的情况视不同的幂级数而定）.
% 
% \begin{example}\label{example:ma-11-1-6}
% 求幂级数
% \[
%   \sum_{n=0}^{+\infty}n!x^{n^n}
% \]
% 的收敛半径与收敛域.
% \end{example}
% 
% 
% \begin{solve}
% 该幂级数有许多项的系数为 $0$, 这样的幂级数我们称之为缺项幂级数。由
% \[
%   1<(n!)^{\frac1{n^n}}<(n^n)^{\frac1{n^n}}
% \]
% 及
% \[
%   \lim_{n\to\infty}(n^n)^{\frac1{n^n}}=1,
% \]
% 得
% \[
%   \overline{\lim_{n\to\infty}}\sqrt[n]{a_n}
%   =
%   \lim_{n\to\infty}\sqrt[n^n]{n!}
%   =
%   \lim_{n\to\infty}(n!)^{\frac1{n^n}}
%   =1.
% \]
% 因此该幂级数的收敛半径 $R=1$。显然在 $x=\pm1$ 时, 该幂级数是发散的, 因此
% \[
%   \sum_{n=0}^{+\infty}n!x^{n^n}
% \]
% 的收敛域为 $(-1,1)$.
% \end{solve}
% 
% 
% \begin{example}\label{example:ma-11-1-7}
% 求幂级数
% \[
%   \sum_{n=1}^{+\infty}\frac{\ln(n+1)}{n^2}(x-3)^n
% \]
% 的收敛半径和收敛域.
% \end{example}
% 
% 
% \begin{solve}
% 由
% \[
%   \lim_{n\to\infty}\left|\frac{a_{n+1}}{a_n}\right|
%   =
%   \lim_{n\to\infty}
%   \left[
%     \frac{\ln(n+2)}{(n+1)^2}\Big/\frac{\ln(n+1)}{n^2}
%   \right]
%   =1,
% \]
% 知该幂级数的收敛半径为 $1$.
% 
% 在端点处, 即 $x=2$ 或 $x=4$ 时, 存在 $N\in\mathbb N$, 当 $n>N$ 时, 成立以下不等式
% \[
%   \left|\frac{\ln(n+1)(x-3)^n}{n^2}\right|
%   \leqslant
%   \frac{\ln(n+1)}{n^2}
%   <
%   \frac1{n^{3/2}},
% \]
% 因此该幂级数在两个端点处绝对收敛。由此可知它的收敛域为 $[2,4]$.
% \end{solve}
% 
% 
% \begin{example}\label{example:ma-11-1-8}
% 求幂级数
% \[
%   \sum_{n=0}^{+\infty}\frac{(x+1)^n}{2^{n^\alpha}}
%   \quad(\text{其中 }\alpha>0)
% \]
% 的收敛半径和收敛域.
% \end{example}
% 
% 
% \begin{solve}
% 由于
% \[
%   \overline{\lim_{n\to\infty}}\sqrt[n]{\frac1{2^{n^\alpha}}}
%   =
%   \lim_{n\to\infty}\frac1{2^{n^{\alpha-1}}}
%   =
%   \begin{cases}
%     1, &0<\alpha<1,\\
%     \dfrac12, &\alpha=1,\\
%     0, &\alpha>1,
%   \end{cases}
% \]
% 我们分别有以下各种情况:
% \begin{enumerate}
%   \item 当 $0<\alpha<1$ 时, $\sum_{n=0}^{+\infty}\dfrac{(x+1)^n}{2^{n^\alpha}}$ 的收敛半径为 $1$。当 $x=-2$ 或 $x=0$ 时, 存在 $N\in\mathbb N$, 使得当 $n>N$ 时,
%   \[
%     n^\alpha>\frac2{\ln2}\ln n,
%   \]
%   从而有
%   \[
%     \left|\frac{(x+1)^n}{2^{n^\alpha}}\right|
%     =
%     \frac1{2^{n^\alpha}}
%     \leqslant
%     \frac1{2^{\frac2{\ln2}\ln n}}
%     =
%     \frac1{n^{\frac2{\ln2}\ln2}}
%     =
%     \frac1{n^2}.
%   \]
%   因此该级数在两个端点处绝对收敛, 此时它的收敛域为 $[-2,0]$.
% 
%   \item 当 $\alpha=1$ 时, 该幂级数的收敛半径 $R=2$。当 $x=-3,1$ 时, 我们有
%   \[
%     \left|\frac{(x+1)^n}{2^n}\right|\equiv1\nrightarrow0\quad(n\to\infty).
%   \]
%   因此该级数在两个端点处都不收敛, 此时它的收敛域为 $(-3,1)$.
% 
%   \item 当 $\alpha>1$ 时, 该幂级数的收敛半径 $R=+\infty$, 此时收敛域为 $(-\infty,+\infty)$.
% \end{enumerate}
% \end{solve}
% 
% 
% \section{幂级数的性质}
% \label{sec:ma-11-2}
% 
% 设幂级数
% \[
%   \sum_{n=0}^{+\infty}a_nx^n
% \]
% 的收敛半径为 $R>0$, 则该幂级数的收敛域是一个区间。人们自然要问: 该幂级数的和函数在收敛域内具有什么样的性质? 由于幂级数的每一项都具有很好的性质, 为了回答这个问题, 我们只要研究幂级数在其收敛域的一致收敛性即可.
% 
% 以下是这方面的主要结果.
% 
% \begin{theorem}[{阿贝尔定理}]\label{theorem:ma-11-2-1}
% 设幂级数
% \[
%   \sum_{n=0}^{+\infty}a_nx^n
% \]
% 的收敛半径为 $R>0$, 则
% \begin{enumerate}
%   \item $\sum_{n=0}^{+\infty}a_nx^n$ 在 $(-R,R)$ 内闭一致收敛;
%   \item 若 $\sum_{n=0}^{+\infty}a_nR^n$ 收敛, 则 $\sum_{n=0}^{+\infty}a_nx^n$ 在 $(-R,R]$ 的任何闭子区间一致收敛;
%   \item 若 $\sum_{n=0}^{+\infty}(-1)^na_nR^n$ 收敛, 则 $\sum_{n=0}^{+\infty}a_nx^n$ 在 $[-R,R)$ 的任何闭子区间一致收敛.
% \end{enumerate}
% \end{theorem}
% 
% 
% \begin{proof}
% 我们容易看出 (1) 可由上节中的定理~\ref{theorem:ma-11-1-1} 直接推出, 因此请读者自证。由于 (2) 和 (3) 的证明类似, 我们只证 (3), (2) 的证明也留给读者.
% 
% 首先我们注意到由 (1) 可知, 要证明 (3) 只要证该幂级数在 $[-R,0]$ 上一致收敛即可。在
% \[
%   \sum_{n=0}^{+\infty}a_nx^n
%   =
%   \sum_{n=0}^{+\infty}a_n(-R)^n\left(\frac{x}{-R}\right)^n
% \]
% 中设
% \[
%   u_n(x)=a_n(-R)^n,\quad
%   v_n(x)=\left(\frac{x}{-R}\right)^n
%   \quad(n=0,1,2,\cdots).
% \]
% 由于函数序列 $\{u_n(x)\}$ 与 $x$ 无关且 $\sum_{n=0}^{+\infty}u_n(x)$ 收敛, 因此它关于 $x\in[-R,0]$ 是一致收敛。而 $v_n(x)$ 关于固定的 $x\in[-R,0]$ 是 $n$ 的单调下降函数, 且对所有 $x\in[-R,0]$ 及所有 $n\in\mathbb N$, 有 $|v_n(x)|\leqslant1$。因此由函数项级数的阿贝尔判别法知在 $[-R,0]$ 上一致收敛.
% \end{proof}
% 
% 
% 人们习惯上将定理~\ref{theorem:ma-11-2-1} 中的结论 (1) 称为\keyterm{阿贝尔第一定理}, 而将结论 (2), (3) 称为\keyterm{阿贝尔第二定理}.
% 
% 由定理~\ref{theorem:ma-11-2-1}, 我们推出
% 
% \begin{corollary}\label{corollary:ma-local-11.2-1}
% 设幂级数
% \[
%   \sum_{n=0}^{+\infty}a_n(x-x_0)^n
% \]
% 的收敛半径为 $R(0<R<+\infty)$, 则
% \begin{enumerate}
%   \item 当该幂级数在 $x=x_0+R$ 收敛时, 有
%   \[
%     \lim_{x\to x_0+R-0}
%     \sum_{n=0}^{+\infty}a_n(x-x_0)^n
%     =
%     \sum_{n=0}^{+\infty}a_nR^n;
%   \]
%   \item 当该幂级数在 $x=x_0-R$ 收敛时, 有
%   \[
%     \lim_{x\to x_0-R+0}
%     \sum_{n=0}^{+\infty}a_n(x-x_0)^n
%     =
%     \sum_{n=0}^{+\infty}a_n(-R)^n.
%   \]
% \end{enumerate}
% 特别地, 我们有
% \end{corollary}
% 
% 
% \begin{corollary}\label{corollary:ma-local-11.2-2}
% 设幂级数
% \[
%   \sum_{n=0}^{+\infty}a_n(x-x_0)^n
% \]
% 的收敛半径 $R>0$, 则和函数
% \[
%   \sum_{n=0}^{+\infty}a_n(x-x_0)^n
% \]
% 在其收敛域上连续.
% \end{corollary}
% 
% 下面我们讨论幂级数在其收敛域上逐项积分的问题。由定理~\ref{theorem:ma-11-2-1} 及其推论, 我们有以下结论.
% 
% \begin{theorem}\label{theorem:ma-11-2-2}
% 设幂级数
% \[
%   \sum_{n=0}^{+\infty}a_n(x-x_0)^n
% \]
% 的收敛半径为 $R>0$（$R$ 可为 $+\infty$）, 则对于其收敛域内任意两点的 $t_1,t_2$, 有
% \[
%   \int_{t_1}^{t_2}\sum_{n=0}^{+\infty}a_n(x-x_0)^n\,dx
%   =
%   \sum_{n=0}^{+\infty}a_n\int_{t_1}^{t_2}(x-x_0)^n\,dx
% \]
% \[
%   =
%   \sum_{n=0}^{+\infty}\frac{a_n}{n+1}
%   \bigl[(t_2-x_0)^{n+1}-(t_1-x_0)^{n+1}\bigr].
% \]
% \end{theorem}
% 
% 
% \begin{remark}\label{remark:ma-11-source-28}
% 在定理~\ref{theorem:ma-11-2-2} 中, 取 $t_1=x_0,t_2=x$, 则幂级数
% \[
%   \sum_{n=0}^{+\infty}a_n(x-x_0)^n
% \]
% 逐项积分得到幂级数
% \[
%   \sum_{n=0}^{+\infty}\frac{a_n}{n+1}(x-x_0)^{n+1},
% \]
% 容易看出这两个幂级数具有相同的收敛半径 $R>0$。但在 $(x_0-R,x_0+R)$ 端点处这两个级数可能具有不同的敛散性。一般来说, 若
% \[
%   \sum_{n=0}^{+\infty}a_n(x-x_0)^n
% \]
% 在端点处收敛, 则
% \[
%   \sum_{n=0}^{+\infty}\frac{a_n}{n+1}(x-x_0)^{n+1}
% \]
% 在端点处也收敛。反之一般不真。例如幂级数
% \[
%   \sum_{n=0}^{+\infty}x^n
% \]
% 的收敛域为 $(-1,1)$, 其逐项积分后得到的幂级数
% \[
%   \sum_{n=0}^{+\infty}\frac{x^{n+1}}{n+1}
% \]
% 的收敛域为 $[-1,1)$, 再对它逐项积分后得到的幂级数
% \[
%   \sum_{n=0}^{+\infty}\frac{x^{n+2}}{(n+1)(n+2)}
% \]
% 的收敛域则为 $[-1,1]$.
% \end{remark}
% 
% 下面我们来讨论幂级数的和函数的可微性问题.
% 
% \begin{theorem}\label{theorem:ma-11-2-3}
% 设幂级数
% \[
%   f(x)=\sum_{n=0}^{+\infty}a_n(x-x_0)^n
% \]
% 的收敛半径为 $R>0$（$R$ 可为 $+\infty$）, 则对于 $\forall x\in(x_0-R,x_0+R)$, $f(x)$ 在 $x$ 处具有任意阶导数, 并且对 $k=1,2,\cdots$, 有
% \[
%   f^{(k)}(x)
%   =
%   \sum_{n=k}^{+\infty}n(n-1)\cdots(n-k+1)a_n(x-x_0)^{n-k}.
% \]
% \end{theorem}
% 
% 
% \begin{proof}
% 设
% \[
%   \varlimsup_{n\to\infty}\sqrt[n]{|a_n|}=\rho,
% \]
% 则对固定的正整数 $k$, 有
% \[
%   \varlimsup_{n\to\infty}
%   \sqrt[n]{n(n-1)\cdots(n-k+1)|a_n|}
%   =\rho.
% \]
% 这说明幂级数
% \[
%   \sum_{n=0}^{+\infty}a_n(x-x_0)^n
% \]
% 逐项 $k$ 次求导后得到幂级数与原幂级数具有相同的收敛半径, 因此它们在 $(x_0-R,x_0+R)$ 均内闭一致收敛。分别对 $k=1,2,\cdots$ 应用函数项级数逐项求导定理即可证明所要结论.
% \end{proof}
% 
% 
% 对于一个幂级数
% \[
%   f(x)=\sum_{n=0}^{+\infty}a_n(x-x_0)^n,
% \]
% 逐项求导后的幂级数与原级数具有相同的收敛半径 $R$。但在 $(x_0-R,x_0+R)$ 的端点处, 它们可能具有不同的敛散性。一般来说, 若逐项求导后的级数在端点处收敛, 则原来的级数在端点处必收敛, 但反之不真.
% 
% \begin{example}\label{example:ma-11-2-1}
% 设
% \[
%   f(x)=\sum_{n=1}^{+\infty}\frac{x^n}{n^p\ln^q(1+n)}
%   \quad(p>0,\ q>0),
% \]
% 试求:
% \begin{enumerate}
%   \item 该幂级数的收敛域;
%   \item 该幂级数逐项积分后所得幂级数的收敛域;
%   \item 该幂级数逐项求导后所得幂级数的收敛域.
% \end{enumerate}
% \end{example}
% 
% 
% \begin{solve}
% (1) 由
% \[
%   \lim_{n\to\infty}
%   \sqrt[n]{\frac1{n^p\ln^q(1+n)}}=1,
% \]
% 知
% \[
%   \sum_{n=1}^{+\infty}\frac{x^n}{n^p\ln^q(1+n)}
% \]
% 的收敛半径为 $1$。当 $x=-1$ 时, 它是一个交错级数, 并且
% \[
%   \left\{\frac1{n^p\ln^q(1+n)}\right\}
% \]
% 单调下降趋于零（可对 $x^{-p}\ln^{-q}x$ 求导验证之）, 因此该级数在 $x=-1$ 时收敛。当 $x=1$ 时, 由数项级数的判别法知, 当 $p>1$ 或者 $p=1$ 且 $q>1$ 时原幂级数收敛, 否则发散。因此
% \[
%   \sum_{n=1}^{+\infty}\frac{x^n}{n^p\ln^q(1+n)}
% \]
% 的收敛域为
% \[
%   \begin{cases}
%     [-1,1], & p>1\text{ 或 }p=1\text{ 且 }q>1,\\
%     [-1,1), & p<1\text{ 或 }p=1\text{ 且 }q\leqslant1.
%   \end{cases}
% \]
% 
% (2) 逐项积分后所得幂级数为
% \[
%   \sum_{n=1}^{+\infty}
%   \frac{x^{n+1}}{(n+1)n^p\ln^q(1+n)},
% \]
% 它的收敛半径为 $1$, 且易看出其收敛域为 $[-1,1]$.
% 
% (3) 逐项求导后所得幂级数为
% \[
%   \sum_{n=1}^{+\infty}\frac{x^{n-1}}{n^{p-1}\ln^q(1+n)},
% \]
% 它的收敛半径为 $1$。当 $p\geqslant1$ 时, 在 $x=-1$ 处, 它是一个莱布尼茨交错级数, 因此收敛。当 $p<1$ 时, 由于
% \[
%   \lim_{n\to\infty}\frac1{n^{p-1}\ln^q(1+n)}=+\infty,
% \]
% 因此该级数在 $x=-1$ 处发散。当 $x=1$ 时, 由数项级数的判别法知, 当 $p>2$, 或者 $p=2$ 且 $q>1$ 时, 它是收敛的, 否则发散。因此
% \[
%   \sum_{n=1}^{+\infty}\frac{x^n}{n^{p-1}\ln^q(1+n)}
% \]
% 的收敛域为
% \[
%   \begin{cases}
%     [-1,1], & p>2\text{ 或 }p=2\text{ 且 }q>1,\\
%     [-1,1), & 1\leqslant p<2\text{ 或 }p=2\text{ 且 }q\leqslant1,\\
%     (-1,1), & p<1.
%   \end{cases}
% \qedhere
% \]
% \end{solve}
% 
% 
% \begin{example}\label{example:ma-11-2-2}
% 求数项级数
% \[
%   \sum_{n=1}^{+\infty}\frac1{n(2n+1)}
% \]
% 的和.
% \end{example}
% 
% 
% \begin{solve}
% 设
% \[
%   f(x)=\sum_{n=1}^{+\infty}\frac{x^{2n+1}}{n(2n+1)},
% \]
% 则我们要求的值为 $f(1)$。由
% \[
%   f'(x)=\sum_{n=1}^{+\infty}\frac{x^{2n}}n
% \]
% 及
% \[
%   f''(x)=2\sum_{n=1}^{+\infty}x^{2n-1}
%   =2x\sum_{n=0}^{+\infty}x^{2n}
%   =\frac{2x}{1-x^2},\quad x\in(-1,1),
% \]
% 我们可求出
% \[
%   f'(x)-f'(0)=\int_0^x f''(t)\,dt
%   =
%   \int_0^x\frac{2t}{1-t^2}\,dt
%   =
%   -\ln(1-x^2).
% \]
% 注意到 $f'(0)=0$, 我们有
% \[
%   f(x)-f(0)
%   =
%   \int_0^x f'(t)\,dt
%   =
%   -\int_0^x\ln(1-t^2)\,dt
% \]
% \[
%   =
%   -t\ln(1-t^2)\Big|_0^x
%   -
%   \int_0^x\frac{2t^2}{1-t^2}\,dt
% \]
% \[
%   =
%   -x\ln(1-x^2)+\int_0^x2\,dt-\int_0^x\frac2{1-t^2}\,dt
% \]
% \[
%   =
%   -x\ln(1-x^2)+2x+\ln\frac{1-x}{1+x}
% \]
% \[
%   =
%   (1-x)\ln(1-x)+2x-(1+x)\ln(1+x).
% \]
% 注意到 $f(0)=0$, 并且
% \[
%   \sum_{n=1}^{+\infty}\frac{x^{2n+1}}{n(2n+1)}
% \]
% 的收敛域为 $[-1,1]$, 因此 $f(x)$ 在 $[-1,1]$ 上连续。由此推出
% \[
%   \sum_{n=1}^{+\infty}\frac1{n(2n+1)}
%   =
%   \lim_{x\to1-0}f(x)
%   =
%   2-2\ln2.
% \qedhere
% \]
% \end{solve}
% 
% 
% \section{初等函数的幂级数展开}
% \label{sec:ma-11-3}
% 
% \subsection{泰勒级数}
% \label{subsec:ma-11-3-1}
% 
% 设 $f(x)$ 在 $(x_0-\delta,x_0+\delta)$（$\delta>0$）内成立
% \[
%   f(x)=\sum_{n=0}^{+\infty}a_n(x-x_0)^n,
% \]
% 则称 $f(x)$ 在 $x_0$ 处可展成幂级数。一个函数 $f(x)$ 若能在某个区间 $I$ 上展开成幂级数, 我们则称 $f(x)$ 在 $I$ 上是实解析的。由上节的定理知道, 此时 $f(x)$ 在 $I$ 内必具有任意阶导数。那么什么样的函数能展开成幂级数呢? 这个问题可以在复变函数课程中得到完满的解决。在数学分析课程中, 我们则主要是通过考查函数的泰勒公式的余项的敛散性来研究这一问题.
% 
% 首先我们来考查以下问题: 若一个函数 $f(x)$ 在 $(-R,R)$（$R>0$）中能展开成幂级数
% \[
%   \sum_{n=0}^{+\infty}a_nx^n,
% \]
% 该幂级数的系数与 $f(x)$ 具有什么样的联系呢?
% 
% 为此设 $f(x)$ 在 $x\in(-R,R)$（$R>0$）时成立等式
% \begin{equation}
%   f(x)=\sum_{n=0}^{+\infty}a_nx^n,
%   \label{eq:ma-11-3-1}
% \end{equation}
% 则首先我们有
% \[
%   f(0)=a_0.
% \]
% 对于 $\forall n\in\mathbb N$, 在 \eqref{eq:ma-11-3-1} 两边求 $n$ 阶导数后令 $x=0$, 我们有
% \[
%   a_n=\frac{f^{(n)}(0)}{n!}.
% \]
% 上述讨论告诉我们, 若 $f(x)$ 是 $(-R,R)$ 中一个幂级数
% \[
%   \sum_{n=0}^{+\infty}a_nx^n
% \]
% 的和函数, 则成立如下等式:
% \[
%   f(x)=\sum_{n=0}^{+\infty}\frac{f^{(n)}(0)}{n!}x^n.
% \]
% 
% \begin{definition}\label{definition:ma-11-3-1}
% 设函数 $f(x)$ 在 $x=x_0$ 处具有任意阶导数, 则称幂级数
% \[
%   \sum_{n=0}^{+\infty}\frac{f^{(n)}(x_0)}{n!}(x-x_0)^n
% \]
% 为 $f(x)$ 在 $x_0$ 处的\keyterm{泰勒级数}; 当 $x_0=0$ 时, 该级数也称为 $f(x)$ 的\keyterm{麦克劳林级数}。如果在 $x_0$ 的某个邻域内成立
% \[
%   f(x)=\sum_{n=0}^{+\infty}\frac{f^{(n)}(x_0)}{n!}(x-x_0)^n,
% \]
% 则
% \[
%   \sum_{n=0}^{+\infty}\frac{f^{(n)}(x_0)}{n!}(x-x_0)^n
% \]
% 称为 $f(x)$ 在该邻域内的\keyterm{泰勒展开式}。特别地, 当 $x_0=0$ 时, 该泰勒展开式也称为 $f(x)$ 的\keyterm{麦克劳林展开式}.
% \end{definition}
% 
% 由上面的讨论, 我们有以下结论.
% 
% \begin{theorem}\label{theorem:ma-11-3-1}
% 若
% \[
%   f(x)=\sum_{n=0}^{+\infty}a_n(x-x_0)^n
% \]
% 在 $(x_0-R,x_0+R)$（$R>0$）成立, 则必有
% \[
%   a_n=\frac{f^{(n)}(x_0)}{n!}\quad(n=0,1,\cdots),
% \]
% 即
% \[
%   \sum_{n=0}^{+\infty}a_n(x-x_0)^n
% \]
% 必是 $f(x)$ 的泰勒展开式（即 $f(x)$ 的幂级数展开式是唯一的）.
% \end{theorem}
% 
% 下面的例子说明 $C^\infty$ 的函数未必是实解析的.
% 
% \begin{example}\label{example:ma-11-3-1}
% 设函数
% \[
%   f(x):=
%   \begin{cases}
%     e^{-\frac1{x^2}}, & x\ne0,\\
%     0, & x=0.
%   \end{cases}
% \]
% 证明: $f(x)\in C^\infty(-\infty,+\infty)$, 但在 $x=0$ 的邻域内不是实解析的（即它不能展成幂级数）.
% \end{example}
% 
% 
% \begin{proof}
% 在第一册中, 我们用洛必达法则证明了 $f(x)$ 在 $(-\infty,+\infty)$ 上具有任意阶导数, 并且对于 $\forall n\in\mathbb N$, 有 $f^{(n)}(0)=0$。但 $f(x)$ 在 $x=0$ 的任何邻域内都不能展成幂级数。否则的话, 在 $x=0$ 的某个邻域内有
% \[
%   f(x)=\sum_{n=0}^{+\infty}\frac{f^{(n)}(0)}{n!}x^n\equiv0.
% \]
% 这显然是不成立的。因此 $f(x)$ 在 $x=0$ 的邻域内不是实解析的.
% \end{proof}
% 
% 
% \subsection{初等函数的泰勒展开式}
% \label{subsec:ma-11-3-2}
% 
% 在本节中我们主要讨论如何来求初等函数的泰勒展开式。设函数
% $f(x)\in C^\infty(-\delta_0,\delta_0)$, 则对 $\forall n\in\mathbb N$, 我们可以求得
% \[
%   \sum_{k=0}^{n}\frac{f^{(k)}(0)}{k!}x^k.
% \]
% 因此 $f(x)$ 在 $(-\delta_0,\delta_0)$ 内有泰勒展开式的充分必要条件是: 对于 $\forall x\in(-\delta_0,\delta_0)$, 有
% \[
%   \lim_{n\to\infty}R_n(x)
%   =
%   \lim_{n\to\infty}
%   \left[
%     f(x)-\sum_{k=0}^{n}\frac{f^{(k)}(0)}{k!}x^k
%   \right]
%   =0.
% \]
% 换句话说, $f(x)$ 在 $(-\delta_0,\delta_0)$ 内有泰勒展开式的充分必要条件是 $f(x)$ 的泰勒公式中的余项趋于零.
% 
% \begin{example}\label{example:ma-11-3-2}
% 证明:
% \[
%   e^x=\sum_{n=0}^{+\infty}\frac{x^n}{n!},
%   \quad x\in(-\infty,+\infty).
% \]
% \end{example}
% 
% 
% \begin{proof}
% 对于 $\forall n\in\mathbb N$, 我们有 $(e^x)^{(n)}=e^x$。因此对于 $\forall x\in(-\infty,+\infty)$, 我们有带拉格朗日余项的泰勒公式
% \[
%   e^x=\sum_{k=0}^{n}\frac{x^k}{k!}
%   +\frac{e^{\theta x}x^{n+1}}{(n+1)!},
% \]
% 其中 $0<\theta<1$。由
% \[
%   \left|\frac{e^{\theta x}x^{n+1}}{(n+1)!}\right|
%   \leqslant
%   e^{|x|}\frac{|x|^{n+1}}{(n+1)!}
%   \to0\quad(n\to\infty),
% \]
% 及 $x$ 的任意性得
% \[
%   e^x=\sum_{n=0}^{+\infty}\frac{x^n}{n!},
%   \quad x\in(-\infty,+\infty).
% \qedhere
% \]
% \end{proof}
% 
% 
% \begin{example}\label{example:ma-11-3-3}
% 证明:
% \[
%   \text{(1) }\quad
%   \sin x=\sum_{n=0}^{+\infty}\frac{(-1)^nx^{2n+1}}{(2n+1)!},
%   \quad x\in(-\infty,+\infty);
% \]
% \[
%   \text{(2) }\quad
%   \cos x=\sum_{n=0}^{+\infty}\frac{(-1)^nx^{2n}}{(2n)!},
%   \quad x\in(-\infty,+\infty).
% \]
% \end{example}
% 
% 
% \begin{proof}
% 对于 $\forall n\in\mathbb N$, 我们有
% \[
%   (\sin x)^{(n)}=\sin\left(x+\frac{n\pi}{2}\right).
% \]
% 对于 $\forall x\in(-\infty,+\infty)$, 其泰勒公式中的拉格朗日余项满足
% \[
%   \left|
%     \frac{\sin\left[\theta x+(2n+1)\frac{\pi}{2}\right]}{(2n+1)!}
%     x^{2n+1}
%   \right|
%   \leqslant
%   \frac{|x|^{2n+1}}{(2n+1)!}
%   \to0\quad(n\to\infty),
% \]
% 其中 $0<\theta<1$。因此我们有
% \[
%   \sin x=\sum_{n=0}^{+\infty}\frac{(-1)^nx^{2n+1}}{(2n+1)!},
%   \quad x\in(-\infty,+\infty).
% \]
% 
% 对于 $\forall n\in\mathbb N$, 我们有
% \[
%   (\cos x)^{(n)}=\cos\left(x+\frac{n\pi}{2}\right).
% \]
% 对于 $\forall x\in(-\infty,+\infty)$, 其泰勒公式中的拉格朗日余项满足
% \[
%   \left|
%     \frac{\cos[\theta x+(n+1)\pi]}{(2n+2)!}
%     x^{2n+2}
%   \right|
%   \leqslant
%   \frac{|x|^{2n+2}}{(2n+2)!}
%   \to0\quad(n\to\infty),
% \]
% 其中 $0<\theta<1$。因此我们有
% \[
%   \cos x=\sum_{n=0}^{+\infty}\frac{(-1)^nx^{2n}}{(2n)!},
%   \quad x\in(-\infty,+\infty).
% \qedhere
% \]
% \end{proof}
% 
% 
% \begin{example}\label{example:ma-11-3-4}
% 设 $\alpha\in\mathbb R$, 证明:
% \begin{equation}
%   (1+x)^\alpha
%   =
%   1+\sum_{n=1}^{+\infty}
%   \frac{\alpha(\alpha-1)\cdots(\alpha-n+1)}{n!}x^n
%   =
%   \sum_{n=0}^{+\infty}\binom{\alpha}{n}x^n,
%   \quad x\in(-1,1),
%   \label{eq:ma-11-3-2}
% \end{equation}
% 其中我们规定 $\binom{\alpha}{0}=1$, 并讨论上述幂级数在端点的收敛情况.
% \end{example}
% 
% 
% \begin{proof}
% 当 $\alpha$ 为一正整数时, $(1+x)^\alpha$ 为一多项式, 由二项式定理即知 \eqref{eq:ma-11-3-2} 成立。因此下面假定 $\alpha$ 不是正整数。由于
% \[
%   \lim_{n\to\infty}
%   \left|
%     \frac{\binom{\alpha}{n+1}}{\binom{\alpha}{n}}
%   \right|
%   =
%   \lim_{n\to\infty}\left|\frac{\alpha-n}{n+1}\right|
%   =1,
% \]
% 因此 \eqref{eq:ma-11-3-2} 式右边的幂级数的收敛半径为 $1$, 从而该级数在 $(-1,1)$ 内收敛。下面证明 \eqref{eq:ma-11-3-2} 式成立.
% 
% 对于 $\forall n\in\mathbb N$, 我们有
% \[
%   [(1+x)^\alpha]^{(n)}
%   =
%   \alpha(\alpha-1)\cdots(\alpha-n+1)(1+x)^{\alpha-n}.
% \]
% 对于 $(1+x)^\alpha$ 的泰勒公式中的积分余项, 我们有
% \[
%   |R_n(x)|
%   =
%   \left|
%     \frac{(\alpha-1)\cdots(\alpha-n)}{n!}\alpha
%     \int_0^x(1+t)^{\alpha-n-1}(x-t)^n\,dt
%   \right|
% \]
% \[
%   =
%   \left|
%     \left[\frac{(\alpha-1)\cdots(\alpha-n)}{n!}x^n\right]
%     \alpha\int_0^x(1+t)^{\alpha-1}
%     \frac{(x-t)^n}{x^n(1+t)^n}\,dt
%   \right|
% \]
% \[
%   \leqslant
%   \left|
%     \left[\frac{(\alpha-1)\cdots(\alpha-n)}{n!}x^n\right]
%     \alpha\int_0^x(1+t)^{\alpha-1}\,dt
%   \right|
% \]
% \begin{equation}
%   =
%   \left|
%     \left[\frac{(\alpha-1)\cdots(\alpha-n)}{n!}x^n\right]
%     [(1+x)^\alpha-1]
%   \right|.
%   \label{eq:ma-11-3-3}
% \end{equation}
% 在 \eqref{eq:ma-11-3-3} 的证明中, 我们利用了以下不等式:
% \[
%   \left|\frac{x-t}{x(1+t)}\right|\leqslant1,
% \]
% 其中 $|x|<1$, 而 $t$ 在 $0$ 和 $x$ 之间。事实上, 当 $x>0$ 时, 我们有
% \[
%   \left|\frac{x-t}{x(1+t)}\right|
%   =
%   \frac{x-t}{x(1+t)}
%   =
%   \frac1{1+t}\left(1-\frac{t}{x}\right)
%   \leqslant
%   \frac1{1+t}
%   \leqslant1;
% \]
% 当 $x<0$ 时, 有
% \[
%   \left|\frac{x-t}{x(1+t)}\right|
%   =
%   \left|\frac{1-\left|\frac{t}{x}\right|}{1-|t|}\right|
%   \leqslant1.
% \]
% 因此该不等式成立, 从而 \eqref{eq:ma-11-3-3} 成立.
% 
% 由于当 $|x|<1$ 时,
% \[
%   \sum_{n=1}^{+\infty}
%   \frac{\alpha(\alpha-1)\cdots(\alpha-n)}{n!}x^n
% \]
% 是收敛的, 因此它的一般项
% \[
%   \frac{\alpha(\alpha-1)\cdots(\alpha-n)}{n!}x^n
%   \to0\quad(n\to\infty).
% \]
% 此时从 \eqref{eq:ma-11-3-3} 可以看出, 对任意的 $|x|<1$, 有
% \[
%   \lim_{n\to\infty}R_n(x)=0.
% \]
% 这就证明了当 $|x|<1$ 时成立
% \[
%   (1+x)^\alpha=\sum_{n=0}^{+\infty}\binom{\alpha}{n}x^n.
% \]
% 
% 为了讨论在端点处 \eqref{eq:ma-11-3-2} 式是否成立, 我们分几种情况进行考虑.
% 
% (1) 考虑 $\alpha>0$ 的情形.
% 
% 由于 $(1+x)^\alpha$ 在 $x=-1$ 处连续, 此时 \eqref{eq:ma-11-3-2} 右边的级数为
% \[
%   1+\sum_{n=1}^{+\infty}
%   \frac{\alpha(\alpha-1)\cdots(\alpha-n+1)}{n!}(-1)^n
% \]
% \[
%   =
%   1+\sum_{n=1}^{+\infty}
%   \frac{-\alpha(-\alpha+1)\cdots[-\alpha+(n-1)]}{n!}
%   \triangleq
%   \sum_{n=0}^{+\infty}a_n.
% \]
% 由于
% \[
%   n\left(\frac{|a_n|}{|a_{n+1}|}-1\right)
%   =
%   n\left(\frac{n+1}{n-\alpha}-1\right)
%   =
%   \frac{n(1+\alpha)}{n-\alpha}
%   \to1+\alpha\quad(n\to\infty),
% \]
% 由拉贝判别法知
% \[
%   \sum_{n=0}^{+\infty}\binom{\alpha}{n}(-1)^n
% \]
% 绝对收敛.
% 
% 同理当 $x=1$ 时,
% \[
%   \sum_{n=0}^{+\infty}\binom{\alpha}{n}
% \]
% 也绝对收敛。因此当 $\alpha>0$ 时, \eqref{eq:ma-11-3-2} 式在 $[-1,1]$ 上成立.
% 
% (2) 考虑 $\alpha\leqslant-1$ 的情形.
% 
% 此时由于
% \[
%   \left|
%     \frac{\binom{\alpha}{n}}{\binom{\alpha}{n+1}}
%   \right|
%   =
%   \left|\frac{n+1}{\alpha-n}\right|
%   \leqslant1,
% \]
% 因此
% \[
%   \lim_{n\to\infty}\binom{\alpha}{n}\ne0,
% \]
% \eqref{eq:ma-11-3-2} 式在 $x=1$ 时不成立.
% 
% (3) 考虑 $-1<\alpha<0$ 的情形.
% 
% 此时
% \[
%   \sum_{n=1}^{+\infty}
%   \frac{\alpha(\alpha-1)\cdots(\alpha-n+1)}{n!}
% \]
% 是一交错级数, 注意到
% \[
%   \left|
%     \frac{\alpha(\alpha-1)\cdots(\alpha-n+1)}{n!}
%   \right|
%   \geqslant
%   \left|
%     \frac{\alpha(\alpha-1)\cdots(\alpha-n+1)}{n!}
%     \cdot\frac{n-\alpha}{n+1}
%   \right|
% \]
% \[
%   =
%   \left|
%     \frac{\alpha(\alpha-1)\cdots(\alpha-(n+1)+1)}{(n+1)!}
%   \right|
% \]
% 且
% \[
%   \left|
%     \frac{\alpha(\alpha-1)\cdots(\alpha-n+1)}{n!}
%   \right|
%   =
%   \left|
%     \left(1-\frac{\alpha+1}{1}\right)
%     \left(1-\frac{\alpha+1}{2}\right)
%     \cdots
%     \left(1-\frac{\alpha+1}{n}\right)
%   \right|
% \]
% \[
%   =
%   e^{\sum_{k=1}^{n}\ln\left(1-\frac{\alpha+1}{k}\right)}
%   \to0\quad(n\to\infty).
% \]
% 由莱布尼茨判别法知, \eqref{eq:ma-11-3-2} 式在 $x=1$ 处成立.
% 
% 当 $\alpha<0$ 时, 由于 $f(x)$ 在 $x=-1$ 处没有定义, 从而此时 \eqref{eq:ma-11-3-2} 式在 $x=-1$ 处不成立.
% 
% 因此, 我们有以下结论:
% 
% 当 $\alpha\leqslant-1$ 时, \eqref{eq:ma-11-3-2} 式在 $(-1,1)$ 内成立;
% 
% 当 $-1<\alpha<0$ 时, \eqref{eq:ma-11-3-2} 式在 $(-1,1]$ 上成立;
% 
% 当 $\alpha>0$ 时, \eqref{eq:ma-11-3-2} 式在 $[-1,1]$ 上成立.
% 
% 下面我们继续讨论一些初等函数的泰勒展开式。有了前面几个初等函数的泰勒展开式, 很多初等函数的泰勒展开式都可以借助它们的泰勒展开式求得。正如我们已经指出的, 一个函数的幂级数展开式如果存在, 则必唯一。因此我们可以对已知函数的泰勒展开式通过自变量的变换或在其收敛域中进行微分或积分得到新的函数的泰勒展开式.
% \end{proof}
% 
% 
% \begin{example}\label{example:ma-11-3-5}
% 证明
% \[
%   \arcsin x
%   =
%   x+\sum_{n=1}^{+\infty}
%   \frac{(2n-1)!!}{(2n)!!}\frac{x^{2n+1}}{2n+1},
%   \quad x\in[-1,1].
% \]
% \end{example}
% 
% 
% \begin{proof}
% 由
% \[
%   \frac1{\sqrt{1-x^2}}=(1-x^2)^{-\frac12},
% \]
% 在 $(1+t)^{-\frac12}$ 的泰勒展开式中令 $t=-x^2$, 得
% \[
%   (1-x^2)^{-\frac12}
%   =
%   \sum_{n=0}^{+\infty}\binom{-1/2}{n}(-x^2)^n
% \]
% \[
%   =
%   1+\frac12x^2+\frac38x^4+\cdots
%   +\frac{(2n-1)!!}{(2n)!!}x^{2n}+\cdots .
% \]
% 对 $x\in(0,1)$, 两边从 $0$ 到 $x$ 积分得
% \[
%   \arcsin x
%   =
%   \int_0^x\frac{dt}{\sqrt{1-t^2}}
%   =
%   \int_0^x
%   \sum_{n=0}^{+\infty}\binom{-1/2}{n}(-t^2)^n\,dt
% \]
% \[
%   =
%   \sum_{n=0}^{+\infty}
%   \int_0^x\binom{-1/2}{n}(-t^2)^n\,dt
%   =
%   x+\sum_{n=1}^{+\infty}
%   \frac{(2n-1)!!}{(2n)!!}\frac{x^{2n+1}}{2n+1}.
% \]
% 当 $x=\pm1$ 时, 由拉贝判别法知上式也成立.
% \end{proof}
% 
% 
% \begin{remark}\label{remark:ma-11-source-47}
% 在 $\arcsin x$ 的泰勒展开式中令 $x=1$, 我们得到
% \[
%   \frac{\pi}{2}
%   =
%   1+\sum_{n=1}^{+\infty}
%   \frac{(2n-1)!!}{(2n)!!}\frac1{2n+1}.
% \]
% \end{remark}
% 
% 
% \begin{example}\label{example:ma-11-3-6}
% 证明
% \[
%   \ln(1+x)
%   =
%   x-\frac{x^2}{2}+\frac{x^3}{3}-\cdots
%   +(-1)^{n+1}\frac{x^n}{n}+\cdots
% \]
% \[
%   =
%   \sum_{n=1}^{+\infty}\frac{(-1)^{n+1}x^n}{n},
%   \quad x\in(-1,1].
% \]
% \end{example}
% 
% 
% \begin{proof}
% 因为 $\dfrac1{1+x}$ 是一个几何级数的和函数, 我们有
% \[
%   \frac1{1+x}
%   =
%   \frac1{1-(-x)}
%   =
%   \sum_{n=1}^{+\infty}(-1)^{n+1}x^{n-1},
%   \quad x\in(-1,1).
% \]
% 对上述等式两边从 $0$ 到 $x\in(-1,1)$ 积分得
% \[
%   \ln(1+x)
%   =
%   \int_0^x\frac1{1+t}\,dt
%   =
%   \sum_{n=1}^{+\infty}(-1)^{n+1}\int_0^x t^{n-1}\,dt
% \]
% \[
%   =
%   \sum_{n=1}^{+\infty}(-1)^{n+1}\frac{x^n}{n},
%   \quad x\in(-1,1).
% \]
% 由于当 $x=1$ 时, 上式右边的级数为莱布尼茨交错级数, 而 $\ln(1+x)$ 在 $x=1$ 处连续, 上述等式在 $x=1$ 时成立。当 $x=-1$ 时, 该级数为
% \[
%   -\sum_{n=1}^{+\infty}\frac1n,
% \]
% 从而是发散的。故
% \[
%   \ln(1+x)
%   =
%   \sum_{n=1}^{+\infty}\frac{(-1)^{n+1}x^n}{n},
%   \quad x\in(-1,1].
% \qedhere
% \]
% \end{proof}
% 
% 
% \begin{example}\label{example:ma-11-3-7}
% 求 $f(x)=\sin(x^2-2x)$ 在 $x=1$ 处的泰勒展开式.
% \end{example}
% 
% 
% \begin{solve}
% 由
% \[
%   \sin(x^2-2x)
%   =
%   \sin[(x-1)^2-1]
% \]
% \[
%   =
%   \sin(x-1)^2\cos1-\sin1\cos(x-1)^2,
% \]
% 利用 $\sin x$ 及 $\cos x$ 的泰勒展开式得
% \[
%   \sin(x-1)^2\cos1-\sin1\cos(x-1)^2
% \]
% \[
%   =
%   \cos1\sum_{n=0}^{+\infty}
%   \frac{(-1)^n(x-1)^{4n+2}}{(2n+1)!}
%   -
%   \sin1\sum_{n=0}^{+\infty}
%   \frac{(-1)^n(x-1)^{4n}}{(2n)!},
%   \quad x\in(-\infty,+\infty).
% \qedhere
% \]
% \end{solve}
% 
% 
% \begin{remark}\label{remark:ma-11-source-52}
% 对于像例~\ref{example:ma-11-3-7} 中一类函数的泰勒展开式, 一般来说不应该直接求 $n$ 阶导数, 而是要应用已知函数的展开式来求之.
% \end{remark}
% 
% 
% \begin{example}\label{example:ma-11-3-8}
% 将
% \[
%   \frac1{1+x+x^2+\cdots+x^{k-1}}\quad(k>1)
% \]
% 展成麦克劳林级数.
% \end{example}
% 
% 
% \begin{solve}
% 直接计算得
% \[
%   \frac1{1+x+x^2+\cdots+x^{k-1}}
%   =
%   \frac{1-x}{1-x^k}
%   =
%   (1-x)\sum_{n=0}^{+\infty}x^{kn}
% \]
% \[
%   =
%   \sum_{n=0}^{+\infty}x^{kn}
%   -
%   \sum_{n=0}^{+\infty}x^{kn+1},
%   \quad x\in(-1,1).
% \]
% 
% 设级数
% \[
%   f(x)=\sum_{n=0}^{+\infty}a_nx^n
%   \quad\text{与}\quad
%   g(x)=\sum_{n=0}^{+\infty}b_nx^n
% \]
% 的收敛半径分别为 $R_1>0$, $R_2>0$ 且 $R_1\leqslant R_2$, 则这两个级数在 $(-R_1,R_1)$ 均内闭绝对一致收敛。因此 $f(x)g(x)$ 的麦克劳林展开式可以通过它们对应的幂级数的乘积得到.
% \end{solve}
% 
% 
% \begin{example}\label{example:ma-11-3-9}
% 求 $\ln^2(1+x)$ 的麦克劳林展开式.
% \end{example}
% 
% 
% \begin{solve}
% 由于
% \[
%   \ln(1+x)=\sum_{n=0}^{+\infty}\frac{(-1)^nx^{n+1}}{n+1},
%   \quad x\in(-1,1],
% \]
% 我们有
% \[
%   \ln^2(1+x)
%   =
%   \left[\sum_{n=0}^{+\infty}\frac{(-1)^nx^{n+1}}{n+1}\right]^2
%   =
%   x^2\left[\sum_{n=0}^{+\infty}\frac{(-1)^nx^n}{n+1}\right]^2
% \]
% \[
%   =
%   x^2\sum_{n=0}^{+\infty}c_nx^n,
%   \quad x\in(-1,1),
% \]
% 其中
% \[
%   c_n
%   =
%   (-1)^n\sum_{k=0}^{n}\frac1{(k+1)(n-k+1)}
% \]
% \[
%   =
%   \frac{(-1)^n}{n+2}
%   \sum_{k=0}^{n}
%   \frac{(k+1)+(n-k+1)}{(k+1)(n-k+1)}
% \]
% \[
%   =
%   \frac{(-1)^n}{n+2}
%   \sum_{k=0}^{n}
%   \left(\frac1{k+1}+\frac1{n-k+1}\right)
%   =
%   \frac{(-1)^n2}{n+2}\sum_{k=0}^{n}\frac1{k+1}.
% \]
% 因此
% \[
%   \ln^2(1+x)
%   =
%   2\sum_{n=1}^{+\infty}
%   \frac{(-1)^{n+1}}{n+1}
%   \left(1+\frac12+\cdots+\frac1n\right)x^{n+1},
%   \quad x\in(-1,1).
% \]
% 我们可以证明当 $x=1$ 时, 该级数收敛; 而当 $x=-1$ 时, 该级数发散。因此我们有
% \[
%   \ln^2(1+x)
%   =
%   2\sum_{n=1}^{+\infty}
%   \frac{(-1)^{n+1}}{n+1}
%   \left(1+\frac12+\cdots+\frac1n\right)x^{n+1},
%   \quad x\in(-1,1].
% \qedhere
% \]
% \end{solve}
% 
% 
% \section{连续函数的多项式逼近}
% \label{sec:ma-11-4}
% 
% 从函数的复杂程度来看, 多项式可以认为是最简单的函数类。用简单函数去逼近复杂函数无论从理论上还是在实际应用中都是十分重要的。但如何去逼近一个函数, 则有许多方法。如果一个函数 $y=f(x)$ 在 $(x_0-R,x_0+R)$（$R>0$）内能展开成幂级数
% \[
%   \sum_{n=0}^{+\infty}a_n(x-x_0)^n,
% \]
% 则容易看出对 $(x_0-R,x_0+R)$ 内的闭区间 $[a,b]$ 和任给 $\varepsilon>0$, 必存在多项式 $P(x)$（事实上取级数的前 $n$ 项（$n$ 充分大）部分和即可）, 使得 $|P(x)-f(x)|<\varepsilon$ 对一切 $x\in[a,b]$ 成立.
% 
% 显然, 由幂级数来逼近一个函数具有很大的限制性: 它要求 $f(x)$ 是实解析的。因此一个自然的问题是: 对一些性质不那么好的函数（如连续函数）是否存在多项式逼近呢?
% 
% 为了讨论这个问题, 我们先给出以下定义.
% 
% \begin{definition}\label{definition:ma-11-4-1}
% 设函数 $f(x)$ 在区间 $I$ 上有定义。若对于 $\forall\varepsilon>0$, 存在多项式 $P(x)$, 使得对一切 $x\in I$, 有 $|f(x)-P(x)|<\varepsilon$, 则称 $f(x)$ 在 $I$ 上\keyterm{可被多项式逼近}.
% \end{definition}
% 
% 显然, $f(x)$ 在 $I$ 上可被多项式逼近的充分必要条件是存在多项式序列 $\{P_n(x)\}$, 使得 $P_n(x)\rightrightarrows f(x)$（$x\in I$）。因此, 若 $f(x)$ 在 $I$ 上可被多项式逼近, $f(x)$ 在 $I$ 上必定连续.
% 
% 可以证明, 若 $f(x)$ 在有限开区间 $(a,b)$ 内可被多项式逼近, 则 $f(x)$ 必可以连续延拓到 $[a,b]$（见本章习题）。另外, 若 $I$ 是无界区间, 当 $f(x)$ 不是多项式时, $f(x)$ 在 $I$ 上则一定不能被多项式逼近（见本章习题）.
% 
% 因此, 一个自然的问题是: 任何闭区间上的连续函数 $f(x)$ 是否一定可被多项式逼近呢? 以下魏尔斯特拉斯定理对此问题给了一个圆满的回答.
% 
% \begin{theorem}[{魏尔斯特拉斯定理}]\label{theorem:ma-11-4-1}
% 设函数 $f(x)\in C[a,b]$, 则 $f(x)$ 可被多项式逼近.
% \end{theorem}
% 
% 为了证明该定理, 我们先证明以下几个引理.
% 
% \begin{lemma}\label{lemma:ma-11-4-2}
% 设函数 $f_j(x)$（$j=1,2,\cdots,J$）在区间 $[a,b]$ 上可被多项式逼近, 并且设 $c_j\in\mathbb R$（$j=1,2,\cdots,J$）, 则
% \[
%   \sum_{j=1}^{J}c_jf_j(x)
% \]
% 在 $[a,b]$ 上可被多项式逼近.
% \end{lemma}
% 
% 
% \begin{lemma}\label{lemma:ma-11-4-3}
% 设函数序列 $\{f_n(x)\}$ 在区间 $[a,b]$ 上有定义, 并且 $f_n(x)\rightrightarrows f(x)$（$x\in[a,b]$）, 再设对于 $\forall n\in\mathbb N$, $f_n(x)$ 可被多项式逼近, 则 $f(x)$ 在 $[a,b]$ 上可被多项式逼近.
% \end{lemma}
% 
% 引理~\ref{lemma:ma-11-4-2} 和引理~\ref{lemma:ma-11-4-3} 的证明是简单的, 请读者自证.
% 
% \begin{lemma}\label{lemma:ma-11-4-4}
% 设 $c\in\mathbb R$, 则存在多项式序列 $\{P_n(x)\}$, 使得 $\{P_n(x)\}$ 在任何有限闭区间上一致收敛到 $|x-c|$.
% \end{lemma}
% 
% 
% \begin{proof}
% 我们首先证明在 $[-1,1]$ 上, $|x|$ 可被多项式逼近.
% 
% 在幂级数理论中我们已知: 当 $\alpha>0$ 时, $(1+t)^\alpha$ 可在 $[-1,1]$ 展成幂级数:
% \[
%   (1+t)^\alpha=\sum_{n=0}^{+\infty}\binom{\alpha}{n}t^n,
%   \quad t\in[-1,1].
% \]
% 因此该幂级数的部分和序列 $\{S_n(t)\}$ 为关于 $t$ 的多项式序列, 它在 $[-1,1]$ 上一致收敛到 $(1+t)^\alpha$。所以, 对于 $\alpha=\dfrac12$,
% \[
%   |x|
%   =
%   [1+(x^2-1)]^{\frac12}
%   =
%   \sum_{n=0}^{+\infty}\binom{1/2}{n}(x^2-1)^n
% \]
% 也在 $|x|\leqslant1$ 上成立。由于它的部分和序列 $\{S_n(x^2-1)\}$ 是关于 $x$ 的多项式序列, 根据函数项级数一致收敛的定义, 对于 $\forall\varepsilon>0$, 存在 $N\in\mathbb N$, 当 $n>N$ 时, 对一切 $x\in[-1,1]$, 有
% \[
%   |S_n(x^2-1)-|x||<\varepsilon.
% \]
% 这就是说, 在 $[-1,1]$ 上, $|x|$ 可被多项式逼近.
% 
% 对于 $\forall n\geqslant1$, 由刚才所证明事实知, 存在多项式 $Q_n(x)$（不一定是 $n$ 次多项式）, 使得对一切 $x\in[-1,1]$, 有
% \[
%   |Q_n(x)-|x||<\frac1{n^2}.
% \]
% 我们设
% \[
%   P_n(x)=nQ_n\left(\frac{x-c}{n}\right),
% \]
% 且 $[A,B]$ 为任意一个区间, 则当
% \[
%   n>\max\{|A|+|c|,|B|+|c|\}
% \]
% 时,
% \[
%   |P_n(x)-|x-c||<\frac1n
% \]
% 对一切 $x\in[A,B]$ 成立.
% \end{proof}
% 
% 
% 为了证明定理~\ref{theorem:ma-11-4-1}, 我们引进符号
% \[
%   H(x)=\frac{x+|x|}{2}=\max\{x,0\},\quad x\in\mathbb R.
% \]
% 再设 $g(x)\in C[a,b]$, 我们称 $g(x)$ 在 $[a,b]$ 上分段线性, 若存在 $[a,b]$ 的分割
% \[
%   \Delta:\quad a=x_0<x_1<\cdots<x_n=b,
% \]
% 使得 $g(x)$ 在 $[x_{i-1},x_i]$（$i=1,2,\cdots,n$）均为线性函数.
% 
% \begin{lemma}\label{lemma:ma-11-4-5}
% 设函数 $g(x)$ 在区间 $[a,b]$ 上连续, 且分段线性, 则 $g(x)$ 可以表示成
% \[
%   g(x)=g(a)+\sum_{i=1}^{n}c_iH(x-x_{i-1})
%   \quad(x\in[a,b]),
% \]
% 其中 $c_i\in\mathbb R$（$i=1,2,\cdots,n$）.
% \end{lemma}
% 
% 此引理的证明是初等的, 可以对分点作归纳法证之, 请读者自己给出.
% 
% \begin{proof}[{定理~\ref{theorem:ma-11-4-1} 的证明}]
% 首先注意到我们可以找到 $[a,b]$ 上的一列分段线性函数 $\{g_m(x)\}$, 使得
% \[
%   g_m(x)\rightrightarrows f(x),\quad x\in[a,b]
% \]
% （例如用例~\ref{example:ma-7-3-1} 的方法）.
% 
% 由引理~\ref{lemma:ma-11-4-5}, 对于分段线性函数 $g_m(x)$, 存在 $[a,b]$ 的分割
% \[
%   \Delta:\quad a=x_0<x_1<\cdots<x_{n_m}=b,
% \]
% 使得
% \[
%   g_m(x)=g_m(a)+\sum_{i=1}^{n_m}c_i^{(m)}H(x-x_{i-1}^{(m)}).
% \]
% 由于
% \[
%   H(x-x_{i-1}^{(m)})
%   =
%   \frac{x-x_{i-1}^{(m)}+|x-x_{i-1}^{(m)}|}{2},
%   \quad i=1,2,\cdots,n_m.
% \]
% 由引理~\ref{lemma:ma-11-4-4} 知, $H(x-x_{i-1}^{(m)})$（$i=1,2,\cdots,n_m$）在 $[a,b]$ 上可被多项式逼近。根据引理~\ref{lemma:ma-11-4-2}, $g_m(x)$ 可被多项式逼近。注意到
% \[
%   g_m(x)\rightrightarrows f(x),\quad x\in[a,b],
% \]
% 再由引理~\ref{lemma:ma-11-4-3} 知 $f(x)$ 在 $[a,b]$ 上可被多项式逼近.
% \end{proof}
% 
% 
% \section*{习题十一}
% \phantomsection
% \addcontentsline{toc}{section}{习题十一}
% \label{exercises:ma-11}
% 
% \noindent 1。求下列幂级数的收敛半径和收敛域:
% 
% \begin{remark}
% 原稿在本题处截断，未提供完整题面。
% \end{remark}
% 
% % 整合来源：math-12.tex
% \chapter{傅里叶级数}
% \label{chap:ma-12}
% 
% 在 18 世纪中叶, 法国数学家傅里叶 (Fourier) 在研究热传导问题时首先引进了三角级数, 后来人们将其称为傅里叶级数。傅里叶级数理论 (人们后来称与之有关的理论为傅里叶分析) 不仅是数学研究中的一门重要的分支, 而且在物理学、工程等方面都有重要的应用。在数学分析课程中, 我们主要介绍傅里叶级数的一些最基本性质.
% 
% 我们可以从下面纯数学的观点来引进傅里叶级数.
% 
% 在上一章中, 我们已经较系统地学习了幂级数。现在换一种观点来看幂级数, 我们可以对幂级数的形式作如下的描述.
% 
% 取定一组函数系, 即
% \[
%   1,\ x,\ x^2,\ \cdots,\ x^n,\ \cdots,
% \]
% 则任何一个多项式 $P_n(x)$ 可以认为是从该函数系中取有限个元素的线性组合, 而一个幂级数 $\sum\limits_{n=0}^{+\infty}a_nx^n$ 则是该函数系的一个无穷的线性组合.
% 
% 从这种观点来看, 我们可以取另一组函数系来考虑其线性组合。要得到系统的理论, 必须要求这个函数系具有很好的性质以及其线性组合也具有很好的性质.
% 
% 我们所熟悉的函数中, 除了多项式外, $\sin nx,\ \cos nx(n=1,2,\cdots)$ 则是遇到比较多的一类初等函数。而三角函数系的无穷线性组合就是我们前面提到的傅里叶级数。在本章中, 我们将考虑三角函数系的无穷线性组合——傅里叶级数, 并研究其基本性质.
% 
% \section{函数的傅里叶级数}
% \label{sec:ma-12-1}
% 
% \subsection{基本三角函数系}
% \label{subsec:ma-12-1-1}
% 
% 形如
% \[
%   \frac{a_0}{2}+\sum_{n=1}^{+\infty}(a_n\cos nx+b_n\sin nx)
% \]
% 的级数称为傅里叶级数, 其中 $a_0,a_n,b_n(n=1,2,\cdots)$ 为实数。我们可以将傅里叶级数看成函数系
% \begin{equation}
%   1,\ \cos x,\ \sin x,\ \cos 2x,\ \sin 2x,\ \cdots,\ \cos nx,\ \sin nx,\ \cdots
%   \label{eq:ma-12-1-1}
% \end{equation}
% 的一个线性组合, 这个函数系我们以后称之为基本三角函数系。为了以后叙述方便, 我们称基本三角函数系中有限个元素的线性组合为一个三角多项式。特别地, 称
% \[
%   \frac{a_0}{2}+\sum_{k=1}^n(a_k\cos kx+b_k\sin kx)=T_n(x)
% \]
% 为一个 $n$ 阶三角多项式.
% 
% 为了讨论傅里叶级数, 首先我们来研究函数系 \eqref{eq:ma-12-1-1}。基本三角函数系 \eqref{eq:ma-12-1-1} 的第一个特性是它的周期性, 即基本函数系 \eqref{eq:ma-12-1-1} 中的函数具有公共的周期 $2\pi$.
% 
% 我们对闭区间 $[a,b]$ 上的可积函数 $f(x),g(x)$, 定义它们之间的内积为
% \[
%   \int_a^b f(x)g(x)\mathrm dx.
% \]
% 如果 $[a,b]$ 上的可积函数 $f(x),g(x)$ 满足 $\int_a^b f(x)g(x)\mathrm dx=0$, 则称 $f(x)$ 和 $g(x)$ 在 $[a,b]$ 上正交.
% 
% 基本三角函数系 \eqref{eq:ma-12-1-1} 的第二个特性是它的正交性, 即基本函数系 \eqref{eq:ma-12-1-1} 中任意两个不同的函数在长度为 $2\pi$ 的任意区间上正交。由于该函数系的周期性, 上述断言等价于下述等式成立:
% \[
%   (1)\quad \int_{-\pi}^{\pi}\sin mx\sin nx\mathrm dx=0
%   \quad (m\ne n,\ \text{并且 }n,m\in\mathbb N);
% \]
% \[
%   (2)\quad \int_{-\pi}^{\pi}\cos mx\cos nx\mathrm dx=0
%   \quad (m\ne n\ \text{并且 }n,m\in\mathbb N\cup\{0\});
% \]
% \[
%   (3)\quad \int_{-\pi}^{\pi}\cos mx\sin nx\mathrm dx=0
%   \quad (m\in\mathbb N\cup\{0\},\ n\in\mathbb N).
% \]
% 事实上, 对于 $\forall n,m\in\mathbb N$, 我们有
% \[
% \begin{aligned}
%   \int_{-\pi}^{\pi}\sin mx\sin nx\mathrm dx
%   &=\frac12\int_{-\pi}^{\pi}[\cos(m-n)x-\cos(m+n)x]\mathrm dx \\
%   &=
%   \begin{cases}
%     \dfrac12\left[\dfrac{\sin(m-n)x}{m-n}-\dfrac{\sin(m+n)x}{m+n}\right]\bigg|_{-\pi}^{\pi}=0,
%       & m\ne n,\\[1.2em]
%     \dfrac12\left[x-\dfrac{\sin(m+n)x}{m+n}\right]\bigg|_{-\pi}^{\pi}=\pi,
%       & m=n;
%   \end{cases}
% \end{aligned}
% \]
% \[
% \begin{aligned}
%   \int_{-\pi}^{\pi}\cos mx\cos nx\mathrm dx
%   &=\frac12\int_{-\pi}^{\pi}[\cos(m-n)x+\cos(m+n)x]\mathrm dx \\
%   &=
%   \begin{cases}
%     0, & m\ne n,\\
%     \pi, & m=n;
%   \end{cases}
% \end{aligned}
% \]
% \[
%   \int_{-\pi}^{\pi}\sin mx\cos nx\mathrm dx
%   =\frac12\int_{-\pi}^{\pi}[\sin(m+n)x+\sin(m-n)x]\mathrm dx=0.
% \]
% 另外, 我们有
% \[
%   \int_{-\pi}^{\pi}1^2\mathrm dx=2\pi
% \]
% 和
% \[
%   \int_{-\pi}^{\pi}\sin nx\cdot1\mathrm dx
%   =\int_{-\pi}^{\pi}\cos nx\cdot1\mathrm dx=0.
% \]
% 因此 (1), (2) 和 (3) 成立.
% 
% \subsection{\texorpdfstring{周期为 $2\pi$ 的函数的傅里叶级数}{周期为 2π 的函数的傅里叶级数}}
% \label{subsec:ma-12-1-2}
% 
% 与幂级数理论一样, 一个函数能否展成傅里叶级数是傅里叶级数理论中的一个重要问题.
% 
% 现在设
% \[
%   f(x)=\frac{a_0}{2}+\sum_{n=1}^{+\infty}(a_n\cos nx+b_n\sin nx)
% \]
% 在 $[-\pi,\pi]$ 上成立, 且右边的级数一致收敛到 $f(x)$。显然, 此时 $f(x)$ 是一个 $[-\pi,\pi]$ 上的连续函数。由函数项级数的一般性理论知 $f(x)$ 在 $[-\pi,\pi]$ 上的定积分可通过该级数逐项积分得到。对于 $\forall k\geqslant0$, 利用基本三角函数系的正交性得
% \[
% \begin{aligned}
%   \int_{-\pi}^{\pi}f(x)\cos kx\mathrm dx
%   &=\int_{-\pi}^{\pi}\frac{a_0}{2}\cos kx\mathrm dx
%     +\sum_{n=1}^{+\infty}a_n\int_{-\pi}^{\pi}\cos nx\cos kx\mathrm dx \\
%   &\quad+\sum_{n=1}^{+\infty}b_n\int_{-\pi}^{\pi}\sin nx\cos kx\mathrm dx \\
%   &=\pi a_k,
% \end{aligned}
% \]
% 因此有
% \[
%   a_k=\frac1\pi\int_{-\pi}^{\pi}f(x)\cos kx\mathrm dx,\quad k=0,1,2,\cdots.
% \]
% 由于对于 $\forall k\geqslant1$, 有
% \[
% \begin{aligned}
%   \int_{-\pi}^{\pi}f(x)\sin kx\mathrm dx
%   &=\int_{-\pi}^{\pi}\frac{a_0}{2}\sin kx\mathrm dx
%     +\sum_{n=1}^{+\infty}a_n\int_{-\pi}^{\pi}\cos nx\sin kx\mathrm dx \\
%   &\quad+\sum_{n=1}^{+\infty}b_n\int_{-\pi}^{\pi}\sin nx\sin kx\mathrm dx \\
%   &=\pi b_k,
% \end{aligned}
% \]
% 因此有
% \[
%   b_k=\frac1\pi\int_{-\pi}^{\pi}f(x)\sin kx\mathrm dx,\quad k=1,2,\cdots.
% \]
% 以上的分析给我们提供了两点信息:
% 
% (1) 如果一个函数 $f(x)$ 能展成傅里叶级数并且其傅里叶级数能逐项积分, 则该 $f(x)$ 的傅里叶级数具有唯一性.
% 
% (2) 如果一个函数 $f(x)$ 在 $[-\pi,\pi]$ 上可积, 我们就可计算出
% \begin{equation}
%   a_n=\frac1\pi\int_{-\pi}^{\pi}f(x)\cos nx\mathrm dx,\quad n=0,1,2,\cdots
%   \label{eq:ma-12-1-2}
% \end{equation}
% 及
% \begin{equation}
%   b_n=\frac1\pi\int_{-\pi}^{\pi}f(x)\sin nx\mathrm dx,\quad n=1,2,\cdots,
%   \label{eq:ma-12-1-3}
% \end{equation}
% 并形式地得到一个与 $f(x)$ 有关系的傅里叶级数。为此我们形式地将它记为
% \begin{equation}
%   f(x)\sim\frac{a_0}{2}+\sum_{n=1}^{+\infty}(a_n\cos nx+b_n\sin nx),
%   \label{eq:ma-12-1-4}
% \end{equation}
% 其中 $a_n(n=0,1,2,\cdots), b_n(n=1,2,\cdots)$ 分别由 \eqref{eq:ma-12-1-2} 和 \eqref{eq:ma-12-1-3} 式给出。今后我们将称 $a_n,b_n(n=1,2,\cdots)$ 为 $f(x)$ 的傅里叶系数, 并称
% \[
%   \frac{a_0}{2}+\sum_{n=1}^{+\infty}(a_n\cos nx+b_n\sin nx)
% \]
% 为 $f(x)$ 的傅里叶级数.
% 
% 在这里读者应注意到, 我们只是说 $f(x)$ 与该傅里叶级数有关 (在 \eqref{eq:ma-12-1-4} 式中, 我们用记号 $\sim$, 而不用等号 $=$), 这是因为我们并不知道该傅里叶级数是否能收敛于 $f(x)$。我们很容易举出例子来说明该傅里叶级数很可能不点点收敛于 $f(x)$。事实上, 改变 $f(x)$ 的有限个点的值, 并不改变 \eqref{eq:ma-12-1-2} 和 \eqref{eq:ma-12-1-3} 中的每个积分的值。设 $f(x)$ 改变有限个点的值后得到的函数为 $g(x)$, 则 $f(x)$ 与 $g(x)$ 具有相同的傅里叶级数。但是在这有限个点, 该傅里叶级数显然不可能同时收敛到 $f(x)$ 与 $g(x)$.
% 
% 但对于连续函数 $f(x)$, 我们有以下结论.
% 
% \begin{theorem}\label{theorem:ma-12-1-1}
% 设函数 $f(x)$ 在 $(-\infty,+\infty)$ 上连续, 以 $2\pi$ 为周期, 且 $f(x)$ 的傅里叶系数全为 $0$, 则 $f(x)$ 在 $(-\infty,+\infty)$ 上恒为 $0$.
% \end{theorem}
% 
% 
% \begin{proof}
% 倘若 $f(x)$ 在 $[-\pi,\pi]$ 上不恒为 $0$, 则在 $(-\pi,\pi)$ 内存在 $x_0$, 使得 $f(x_0)\ne0$。我们不妨设 $f(x_0)>0$。由 $f(x)$ 的连续性, 存在 $\delta>0$, 当 $x\in(x_0-\delta,x_0+\delta)$ 时, 有
% \[
%   f(x)>\frac{f(x_0)}{2}=M>0.
% \]
% 由于 $f(x)$ 的傅里叶系数全为 $0$, 容易看出对于任何的三角多项式 $T(x)$, 有
% \[
%   \int_{-\pi}^{\pi}f(x)T(x)\mathrm dx=0.
% \]
% 我们现在取定以下三角多项式
% \[
% \begin{aligned}
%   T_0(x)&=1+\cos(x-x_0)-\cos\delta\\
%   &=1+\cos x\cos x_0+\sin x\sin x_0-\cos\delta.
% \end{aligned}
% \]
% 显然存在 $r>1$, 对于 $x\in\left(x_0-\dfrac{\delta}{2},x_0+\dfrac{\delta}{2}\right)$, 有 $T_0(x)\geqslant r>1$。注意到当 $x\in[x_0-\pi,x_0+\pi]\setminus(x_0-\delta,x_0+\delta)$ 时, 有 $|T_0(x)|\leqslant1$.
% 
% 对于 $\forall n\in\mathbb N$, 由于 $T_0^n(x)$ 是一个三角多项式, 因此有
% \begin{equation}
%   \int_{x_0-\pi}^{x_0+\pi}f(x)T_0^n(x)\mathrm dx
%   =\int_{-\pi}^{\pi}f(x)T_0^n(x)\mathrm dx=0.
%   \label{eq:ma-12-1-5}
% \end{equation}
% 另一方面, 对于 $\forall n\in\mathbb N$, 我们有
% \[
%   \left|
%     \int_{x_0-\pi}^{x_0-\delta}f(x)T_0^n(x)\mathrm dx
%     +\int_{x_0+\delta}^{x_0+\pi}f(x)T_0^n(x)\mathrm dx
%   \right|\leqslant 2\pi M\cdot1^n=2\pi M.
% \]
% 而当 $n\to+\infty$ 时, 我们有
% \[
%   \int_{x_0-\delta}^{x_0+\delta}f(x)T_0^n(x)\mathrm dx
%   \geqslant \int_{x_0-\frac{\delta}{2}}^{x_0+\frac{\delta}{2}}f(x)T_0^n(x)\mathrm dx
%   \geqslant Mr^n\delta\to+\infty.
% \]
% 因此有
% \[
%   \int_{x_0-\pi}^{x_0+\pi}f(x)T_0^n(x)\mathrm dx\to+\infty
%   \quad (n\to+\infty).
% \]
% 这与 \eqref{eq:ma-12-1-5} 矛盾.
% \end{proof}
% 
% 
% 对于定理~\ref{theorem:ma-12-1-1}, 我们以后有更加简单的方法来证明之.
% 
% 在举例来熟悉函数的傅里叶级数之前, 我们必须注意到以下事实: 对于在长度为 $2\pi$ 的半开半闭区间上定义的一个函数 $f(x)$, 我们总是可以以 $2\pi$ 为周期将它延拓到 $(-\infty,+\infty)$。因此在下面的例子中的函数均是在一个长度为 $2\pi$ 的区间给出, 但它们可以看成是一个周期为 $2\pi$ 的函数的限制.
% 
% \begin{example}\label{example:ma-12-1-1}
% 设函数 $f(x)=x(-\pi\leqslant x<\pi)$, 求 $f(x)$ 的傅里叶级数.
% \end{example}
% 
% 
% \begin{solve}
% $f(x)$ 周期延拓后的图像见图 12.1.1。注意到 $\sin kx(k\in\mathbb N)$ 是奇函数, 而 $\cos kx(k\in\mathbb N\cup\{0\})$ 是偶函数, 因此, 当 $f(x)=x$ 为奇函数时, 有 $a_n=0(n=0,1,2,\cdots)$。对于 $\forall n\in\mathbb N$, 我们有
% \[
% \begin{aligned}
%   b_n
%   &=\frac1\pi\int_{-\pi}^{\pi}x\sin nx\mathrm dx
%   =\frac2\pi\int_0^\pi x\sin nx\mathrm dx\\
%   &=-\frac2{n\pi}x\cos nx\bigg|_0^\pi+\frac2{n\pi}\int_0^\pi\cos nx\mathrm dx\\
%   &=(-1)^{n-1}\frac2n.
% \end{aligned}
% \]
% 因此
% \[
%   f(x)\sim\sum_{n=1}^{+\infty}(-1)^{n-1}\frac2n\sin nx.
% \]
% 
% 在例~\ref{example:ma-12-1-1} 中, 严格说来 $f(x)$ 只在 $(-\pi,\pi)$ 中是奇函数。但由于改变 $x=-\pi$ 一点处的函数值不会改变 $f(x)$ 的傅里叶系数, 因此我们不妨设 $f(x)$ 在 $[-\pi,\pi]$ 上是奇函数。在下面的例子中, 我们说一个函数 $f(x)$ 是奇 (偶) 函数, 在大部分情形下仅是对 $f(x)$ 在一对称区间中除去若干个点后是奇 (偶) 函数而言的.
% \end{solve}
% 
% 
% \paragraph*{思考题}
% 设函数 $f(x)=x(x\in[0,2\pi))$, 是否 $f(x)$ 的傅里叶级数也是
% \[
%   \sum_{n=1}^{+\infty}(-1)^n\frac2n\sin nx?
% \]
% 
% \begin{example}\label{example:ma-12-1-2}
% 设函数 $f(x)=\dfrac{\pi-x}{2}(x\in[0,2\pi))$, 求 $f(x)$ 的傅里叶级数.
% \end{example}
% 
% 
% \begin{solve}
% 由图 12.1.2 可见, $f(x)$ 是奇函数, 因此 $a_n=0(n=0,1,2,\cdots)$。对于 $\forall n\in\mathbb N$, 我们有
% \[
% \begin{aligned}
%   b_n
%   &=\frac2\pi\int_0^\pi\frac{\pi-x}{2}\sin nx\mathrm dx\\
%   &=-\frac1{n\pi}(\pi-x)\cos nx\bigg|_0^\pi
%     -\frac1{n\pi}\int_0^\pi\cos nx\mathrm dx\\
%   &=\frac1n.
% \end{aligned}
% \]
% 因此我们有
% \[
%   f(x)\sim\sum_{n=1}^{+\infty}\frac{\sin nx}{n}.
% \qedhere
% \]
% \end{solve}
% 
% 
% \begin{remark}\label{remark:ma-12-source-7}
% 从上例可以看出, 当考查 $f(x)$ 的奇偶性时, 可结合其延拓后的图像来考查, 不要仅仅从所给函数的解析式来下结论.
% \end{remark}
% 
% 
% \begin{example}\label{example:ma-12-1-3}
% 设函数 $f(x)=x^2(x\in[0,2\pi))$, 求 $f(x)$ 的傅里叶级数.
% \end{example}
% 
% 
% \begin{solve}
% $f(x)$ 周期延拓后的图像如图 12.1.3 所示.
% 
% 从图 12.1.3 可以看出 $f(x)$ 不具有奇偶性。由于 $f(x)$ 的傅里叶系数的积分中的被积函数均为以 $2\pi$ 为周期的函数, 因此我们只需在一个长度为 $2\pi$ 的区间来求之即可。因此有
% \[
%   a_0=\frac1\pi\int_{-\pi}^{\pi}f(x)\mathrm dx
%   =\frac1\pi\int_0^{2\pi}x^2\mathrm dx
%   =\frac83\pi^2.
% \]
% \[
% \begin{aligned}
%   a_n
%   &=\frac1\pi\int_0^{2\pi}x^2\cos nx\mathrm dx\\
%   &=\frac1{n\pi}x^2\sin nx\bigg|_0^{2\pi}
%     -\frac2{n\pi}\int_0^{2\pi}x\sin nx\mathrm dx\\
%   &=\frac2{n^2\pi}x\cos nx\bigg|_0^{2\pi}
%     -\frac2{n^2\pi}\int_0^{2\pi}\cos nx\mathrm dx\\
%   &=\frac4{n^2},\quad n=1,2,\cdots,
% \end{aligned}
% \]
% \[
% \begin{aligned}
%   b_n
%   &=\frac1\pi\int_0^{2\pi}x^2\sin nx\mathrm dx\\
%   &=-\frac1{n\pi}x^2\cos nx\bigg|_0^{2\pi}
%     +\frac2{n\pi}\int_0^{2\pi}x\cos nx\mathrm dx\\
%   &=-\frac{4\pi}{n}-\frac2{n^2\pi}\int_0^{2\pi}\sin nx\mathrm dx\\
%   &=-\frac{4\pi}{n},\quad n=1,2,\cdots,
% \end{aligned}
% \]
% 因此
% \[
%   f(x)\sim\frac{4\pi^2}{3}
%   +4\sum_{n=1}^{+\infty}\frac{\cos nx}{n^2}
%   -4\pi\sum_{n=1}^{+\infty}\frac{\sin nx}{n}.
% \qedhere
% \]
% \end{solve}
% 
% 
% \subsection{正弦级数与余弦级数}
% \label{subsec:ma-12-1-3}
% 
% 若函数 $f(x)$ 在 $(a,b)$ 内有定义, 而且它是一个在 $[a,b]$ 上可积函数的限制, 我们将称 $f(x)$ 在 $(a,b)$ 内可积。设 $f(x)$ 在 $(0,\pi)$ 内可积, 则我们可以将 $f(x)$ 延拓成 $(-\pi,\pi)$ 中的奇函数 $\widetilde f(x)$ (称之为奇延拓), 然后再将它周期地延拓到 $(-\infty,+\infty)$。在这种情形下, $\widetilde f(x)$ 的傅里叶系数中
% \[
%   a_n=0,\quad n=0,1,2,\cdots,
% \]
% 而
% \[
%   b_n=\frac1\pi\int_{-\pi}^{\pi}\widetilde f(x)\sin nx\mathrm dx
%   =\frac2\pi\int_0^\pi f(x)\sin nx\mathrm dx,\quad n=1,2,\cdots,
% \]
% 因此
% \[
%   f(x)\sim\sum_{n=1}^{+\infty}b_n\sin nx.
% \]
% 此时我们称 $\sum\limits_{n=1}^{+\infty}b_n\sin nx$ 为 $f(x)$ 在 $(0,\pi)$ 上的正弦级数.
% 
% 设 $f(x)$ 在 $(0,\pi)$ 内可积, 则我们也可以将 $f(x)$ 延拓成 $(-\pi,\pi)$ 中的偶函数 $\widetilde f(x)$ (称之为偶延拓), 然后再将它周期地延拓到 $(-\infty,+\infty)$。在这种情形下, $\widetilde f(x)$ 的傅里叶系数中
% \[
%   b_n=0,\quad n=1,2,\cdots,
% \]
% 而
% \[
%   a_n=\frac2\pi\int_0^\pi f(x)\cos nx\mathrm dx,\quad n=0,1,2,\cdots,
% \]
% 因此
% \[
%   f(x)\sim\frac{a_0}{2}+\sum_{n=1}^{+\infty}a_n\cos nx.
% \]
% 此时我们称 $\dfrac{a_0}{2}+\sum\limits_{n=1}^{+\infty}a_n\cos nx$ 为 $f(x)$ 在 $(0,\pi)$ 上的余弦级数。由于在上述的奇延拓与偶延拓中, 延拓后的函数虽然不是唯一的, 但是仅仅在 $x=0,\pm\pi$ 处的值可能不同, 因此它们的傅里叶系数是相同的。这说明了上面定义的正弦级数与余弦级数是由 $f(x)$ 在 $(0,\pi)$ 内的值唯一确定的.
% 
% \begin{example}\label{example:ma-12-1-4}
% 分别求函数 $f(x)=x^2(x\in(0,\pi))$ 的正弦级数与余弦级数.
% \end{example}
% 
% 
% \begin{solve}
% 先求正弦级数。由
% \[
% \begin{aligned}
%   b_n
%   &=\frac2\pi\int_0^\pi x^2\sin nx\mathrm dx\\
%   &=-\frac2{n\pi}x^2\cos nx\bigg|_0^\pi
%     +\frac4{n\pi}\int_0^\pi x\cos nx\mathrm dx\\
%   &=(-1)^{n-1}\frac{2\pi}{n}
%     +\frac4{n^2\pi}\left(x\sin nx\bigg|_0^\pi-\int_0^\pi\sin nx\mathrm dx\right)\\
%   &=(-1)^{n-1}\frac{2\pi}{n}
%     +\frac4{n^3\pi}[(-1)^n-1],\quad n=1,2,\cdots,
% \end{aligned}
% \]
% 我们求得 $f(x)$ 的正弦级数为
% \[
%   x^2\sim
%   2\pi\sum_{n=1}^{+\infty}\frac{(-1)^n}{n}\sin nx
%   -\frac8\pi\sum_{n=1}^{+\infty}\frac{\sin(2n-1)x}{(2n-1)^3}.
% \]
% 
% 再求余弦级数。由
% \[
%   a_0=\frac2\pi\int_0^\pi x^2\mathrm dx=\frac23\pi^2,
% \]
% \[
% \begin{aligned}
%   a_n
%   &=\frac2\pi\int_0^\pi x^2\cos nx\mathrm dx\\
%   &=\frac2{n\pi}x^2\sin nx\bigg|_0^\pi
%     -\frac4{n\pi}\int_0^\pi x\sin nx\mathrm dx\\
%   &=\frac4{n^2\pi}\left(x\cos nx\bigg|_0^\pi-\int_0^\pi\cos nx\mathrm dx\right)\\
%   &=(-1)^n\frac4{n^2},\quad n=1,2,\cdots,
% \end{aligned}
% \]
% 我们求得 $f(x)$ 的余弦级数为
% \[
%   x^2\sim\frac{\pi^2}{3}
%   +4\sum_{n=1}^{+\infty}\frac{(-1)^n}{n^2}\cos nx.
% \qedhere
% \]
% \end{solve}
% 
% 
% \subsection{\texorpdfstring{周期为 $2T$ 的函数的傅里叶级数}{周期为 2T 的函数的傅里叶级数}}
% \label{subsec:ma-12-1-4}
% 
% 我们现在来讨论一下周期为 $2T>0$ 的函数的傅里叶级数问题.
% 
% 设函数 $f(x)$ 以 $2T$ 为周期且在任何有限闭区间上可积, 作自变量变换 $x=\dfrac{T}{\pi}t$, 则函数 $F(t)=f\left(\dfrac{T}{\pi}t\right)$ 以 $2\pi$ 为周期且在任何有限闭区间上可积。设
% \[
%   F(t)\sim\frac{a_0}{2}+\sum_{n=1}^{+\infty}(a_n\cos nt+b_n\sin nt),
% \]
% 则将变量 $t$ 变回变量 $x$ 后, 我们便得到 $f(x)$ 的傅里叶级数:
% \[
%   f(x)\sim\frac{a_0}{2}
%   +\sum_{n=1}^{+\infty}\left(a_n\cos\frac{n\pi}{T}x+b_n\sin\frac{n\pi}{T}x\right),
% \]
% 其中
% \[
%   a_n=\frac1\pi\int_{-\pi}^{\pi}F(t)\cos nt\mathrm dt
%   =\frac1T\int_{-T}^{T}f(x)\cos\frac{n\pi}{T}x\mathrm dx,\quad n=0,1,2,\cdots,
% \]
% \[
%   b_n=\frac1\pi\int_{-\pi}^{\pi}F(t)\sin nt\mathrm dt
%   =\frac1T\int_{-T}^{T}f(x)\sin\frac{n\pi}{T}x\mathrm dx,\quad n=1,2,\cdots.
% \]
% 
% \begin{example}\label{example:ma-12-1-5}
% 求 $f(x)=x\cos x\left(-\dfrac{\pi}{2}\leqslant x<\dfrac{\pi}{2}\right)$ 的傅里叶级数.
% \end{example}
% 
% 
% \begin{solve}
% $f(x)$ 是以 $\pi$ 为周期的函数, 且为奇函数, 因此 $a_n=0(n=0,1,2,\cdots)$。记 $T=\dfrac{\pi}{2}$, 则
% \[
% \begin{aligned}
%   b_n
%   &=\frac2T\int_0^T x\cos x\sin 2nx\mathrm dx\\
%   &=\frac1T\int_0^T x[\sin(2n-1)x+\sin(2n+1)x]\mathrm dx\\
%   &=\frac1T\left[-\frac1{2n-1}x\cos(2n-1)x\bigg|_0^T
%     +\frac1{2n-1}\int_0^T\cos(2n-1)x\mathrm dx\right]\\
%   &\quad+\frac1T\left[-\frac1{2n+1}x\cos(2n+1)x\bigg|_0^T
%     +\frac1{2n+1}\int_0^T\cos(2n+1)x\mathrm dx\right]\\
%   &=\frac1T\left[\frac1{(2n-1)^2}\sin(2n-1)T
%     +\frac1{(2n+1)^2}\sin(2n+1)T\right]\\
%   &=\frac{(-1)^{n-1}}{T}\left[\frac1{(2n-1)^2}-\frac1{(2n+1)^2}\right]\\
%   &=\frac{(-1)^{n-1}}{\pi}\frac{16n}{(4n^2-1)^2},\quad n=1,2,\cdots.
% \end{aligned}
% \]
% 因此
% \[
%   x\cos x\sim
%   \sum_{n=1}^{+\infty}\frac{(-1)^{n-1}16n}{\pi(4n^2-1)^2}\sin 2nx.
% \]
% 
% 对于周期为 $2T$ 的函数 $f(x)$ 的傅里叶级数, 我们也可以用如下的办法来加以讨论。取基本三角函数系
% \[
%   1,\ \cos\frac{\pi}{T}x,\ \sin\frac{\pi}{T}x,\ \cos\frac{2\pi}{T}x,\ \sin\frac{2\pi}{T}x,\ \cdots,\ 
%   \cos\frac{n\pi}{T}x,\ \sin\frac{n\pi}{T}x,\ \cdots,
% \]
% 容易看出它们是长度为 $2T$ 的区间上的正交函数系, 平行于周期为 $2\pi$ 的函数的傅里叶级数的讨论, 即可得到周期为 $2T$ 的 $f(x)$ 的傅里叶级数公式.
% \end{solve}
% 
% 
% \section{傅里叶级数的敛散性}
% \label{sec:ma-12-2}
% 
% 从函数项级数 $\sum\limits_{n=1}^{+\infty}u_n(x)$ 的一般研究方法知, 若其前 $n$ 项部分和
% \[
%   S_n(x)=\sum_{k=1}^n u_k(x)\quad(n=1,2,\cdots)
% \]
% 具有某种有规律的表达式, 则该级数可以容易通过对其部分和序列的研究来讨论其敛散性。傅里叶级数正是具有这种特性, 由于它的部分和序列具有很有规律的表达式, 因此我们可以从它的研究中得到傅里叶级数敛散性的许多信息.
% 
% \subsection{狄利克雷积分}
% \label{subsec:ma-12-2-1}
% 
% 从上节中我们已经知道, 若周期函数 $f(x)$ 在 $[-\pi,\pi]$ 上可积, 则可形式地得到其傅里叶级数
% \[
%   f(x)\sim\frac{a_0}{2}+\sum_{n=1}^{+\infty}(a_n\cos nx+b_n\sin nx),
% \]
% 其中
% \[
%   a_n=\frac1\pi\int_{-\pi}^{\pi}f(x)\cos nx\mathrm dx,\quad n=0,1,2,\cdots
% \]
% \[
%   b_n=\frac1\pi\int_{-\pi}^{\pi}f(x)\sin nx\mathrm dx,\quad n=1,2,\cdots.
% \]
% 我们现在来考查 $f(x)$ 的傅里叶级数部分和序列。对于 $\forall n\in\mathbb N$, 我们有
% \[
% \begin{aligned}
%   S_n(x)
%   &=\frac{a_0}{2}+\sum_{k=1}^n(a_k\cos kx+b_k\sin kx)\\
%   &=\frac1{2\pi}\int_{-\pi}^{\pi}f(u)\mathrm du
%     +\sum_{k=1}^n\frac1\pi\int_{-\pi}^{\pi}f(u)[\cos kx\cos ku+\sin kx\sin ku]\mathrm du\\
%   &=\frac1\pi\int_{-\pi}^{\pi}f(u)\left[\frac12+\sum_{k=1}^n\cos k(u-x)\right]\mathrm du\\
%   &=\frac1\pi\int_{-\pi}^{\pi}f(u)\frac{\sin\left(n+\frac12\right)(u-x)}{2\sin\frac{u-x}{2}}\mathrm du.
% \end{aligned}
% \]
% 为了更加便于分析, 我们作变换 $u=x+t$, 并利用 $f(x)$ 的周期性得
% \[
% \begin{aligned}
%   S_n(x)
%   &=\frac1\pi\int_{-\pi+x}^{\pi+x}
%     f(u)\frac{\sin\left(n+\frac12\right)(u-x)}{2\sin\frac{u-x}{2}}\mathrm du\\
%   &=\frac1\pi\int_{-\pi}^{\pi}
%     f(x+t)\frac{\sin\left(n+\frac12\right)t}{2\sin\frac t2}\mathrm dt\\
%   &=\frac1\pi\int_{-\pi}^{0}
%     f(x+t)\frac{\sin\left(n+\frac12\right)t}{2\sin\frac t2}\mathrm dt
%     +\frac1\pi\int_0^\pi
%     f(x+t)\frac{\sin\left(n+\frac12\right)t}{2\sin\frac t2}\mathrm dt\\
%   &=\frac1\pi\int_0^\pi
%     (f(x+t)+f(x-t))\frac{\sin\left(n+\frac12\right)t}{2\sin\frac t2}\mathrm dt.
% \end{aligned}
% \]
% 上述 $f(x)$ 的傅里叶级数部分和的积分形式将给我们提供许多有用信息。下面我们先来对它们作一些粗略的分析。若取 $f(x)\equiv1$, 则其傅里叶级数的前 $n$ 项的部分和 $S_n(x)\equiv1$。这说明了对于 $\forall n\in\mathbb N$, 有
% \[
%   1=\frac2\pi\int_0^\pi\frac{\sin\left(n+\frac12\right)t}{2\sin\frac t2}\mathrm dt
%   =2\int_0^\pi D_n(t)\mathrm dt,
% \]
% 其中
% \[
%   D_n(t)=\frac{\sin\left(n+\frac12\right)t}{2\pi\sin\frac t2}.
% \]
% 今后我们将称
% \[
%   \int_0^\pi [f(x+t)+f(x-t)]D_n(t)\mathrm dt
% \]
% 为 $f(x)$ 的狄利克雷积分, 并且称 $D_n(t)$ 为狄利克雷核.
% 
% 现取定 $x_0\in[-\pi,\pi]$, 我们来研究 $f(x)$ 的傅里叶级数部分和序列 $\{S_n(x_0)\}$ 是否以 $S_0$ 为极限。由于
% \[
% \begin{aligned}
%   S_n(x_0)-S_0
%   &=\int_0^\pi (f(x_0+t)+f(x_0-t)-2S_0)D_n(t)\mathrm dt\\
%   &=\int_0^\pi\frac{f(x_0+t)+f(x_0-t)-2S_0}{2\pi\sin\frac t2}
%     \sin\left(n+\frac12\right)t\mathrm dt\\
%   &=\int_0^\pi G(t)\sin\left(n+\frac12\right)t\mathrm dt,
% \end{aligned}
% \]
% 其中
% \begin{equation}
%   G(t)=\frac{f(x_0+t)+f(x_0-t)-2S_0}{2\pi\sin\frac t2}.
%   \label{eq:ma-12-2-1}
% \end{equation}
% 对上述积分我们粗略地观察一下可以看出, 当 $n$ 很大时, $\sin\left(n+\dfrac12\right)t$ 的值随着 $t$ 在 $[0,\pi]$ 上的变化, 不断地在 $x$ 轴的上下波动。若 $G(t)$ 是个常数的话, 显然有
% \[
%   \int_0^\pi G(t)\sin\left(n+\frac12\right)t\mathrm dt\to0
%   \quad (n\to\infty).
% \]
% 若 $G(t)$ 不是常数, 但 $G(t)$ 可积, 当 $n\to\infty$ 时, 则由 $\sin\left(n+\dfrac12\right)t$ 在 $x$ 轴上下不断波动, 使得 $G(t)\sin\left(n+\dfrac12\right)t$ 的积分不断相抵, 因此上面的积分也应该趋于 $0$.
% 
% 在以下行文中, 我们经常要用到以下假定: $f(x)$ 在 $[a,b]$ 上可积或有瑕点时绝对可积。对此假定的确切意义是: $f(x)$ 在 $[a,b]$ 上或者是黎曼可积的, 或者具有有限个瑕点, 在 $[a,b]$ 中不含瑕点的任何闭子区间 $f(x)$ 均是黎曼可积的, 并且 $|f(x)|$ 在 $[a,b]$ 上的瑕积分是收敛的.
% 
% 我们有下面的黎曼-勒贝格引理.
% 
% \begin{lemma}[{黎曼-勒贝格引理}]\label{lemma:ma-12-2-1}
% 设函数 $f(x)$ 在区间 $[a,b]$ 上可积或有瑕点时绝对可积, 则
% \[
%   \lim_{\lambda\to+\infty}\int_a^b f(x)\sin\lambda x\mathrm dx
%   =\lim_{\lambda\to+\infty}\int_a^b f(x)\cos\lambda x\mathrm dx=0.
% \]
% \end{lemma}
% 
% 
% \begin{proof}
% 我们分两种情形加以证明.
% 
% (1) $f(x)$ 在 $[a,b]$ 上可积.
% 
% 由 $f(x)$ 在 $[a,b]$ 上可积, 从而它是有界的, 因此存在 $M>0$, 使得 $\forall x\in[a,b]$, 有 $|f(x)|\leqslant M$。再由 $f(x)$ 的可积性, 对于 $\forall\varepsilon>0$, 存在 $[a,b]$ 的分割
% \[
%   a=x_0<x_1<\cdots<x_N=b,
% \]
% 对于 $i=1,2,\cdots,N$, 记 $\Delta x_i=x_i-x_{i-1}$ 和
% \[
%   M_i=\sup_{x\in[x_{i-1},x_i]}\{f(x)\},\quad
%   m_i=\inf_{x\in[x_{i-1},x_i]}\{f(x)\},
% \]
% 则有
% \[
%   \sum_{i=1}^N(M_i-m_i)\Delta x_i<\frac{\varepsilon}{2}.
% \]
% 现在我们取定一个满足上述条件的固定分割, 注意到对每个 $1\leqslant i\leqslant N$, 当 $x\in[x_{i-1},x_i]$ 时, 有
% \[
%   0\leqslant f(x)-m_i\leqslant M_i-m_i,
% \]
% 因此有
% \[
% \begin{aligned}
%   \left|\int_a^b f(x)\sin\lambda x\mathrm dx\right|
%   &=\left|\sum_{i=1}^N\int_{x_{i-1}}^{x_i}f(x)\sin\lambda x\mathrm dx\right|\\
%   &\leqslant\left|\sum_{i=1}^N\int_{x_{i-1}}^{x_i}(f(x)-m_i)\sin\lambda x\mathrm dx\right|
%     +\left|\sum_{i=1}^N m_i\int_{x_{i-1}}^{x_i}\sin\lambda x\mathrm dx\right|\\
%   &\leqslant\sum_{i=1}^N(M_i-m_i)\Delta x_i+\frac2\lambda\sum_{i=1}^N|m_i|\\
%   &<\frac{\varepsilon}{2}+\frac{2NM}{\lambda}.
% \end{aligned}
% \]
% 由于这个分割是固定的, 因此存在 $X=\dfrac{4NM}{\varepsilon}$, 当 $\lambda>X$ 时, 有 $\dfrac{2NM}{\lambda}<\dfrac{\varepsilon}{2}$。所以, 当 $\lambda>X$ 时, 有
% \[
%   \left|\int_a^b f(x)\sin\lambda x\mathrm dx\right|<\varepsilon.
% \]
% 这就证明了当 $f(x)$ 在 $[a,b]$ 上可积时, 有
% \[
%   \lim_{\lambda\to+\infty}\int_a^b f(x)\sin\lambda x\mathrm dx=0.
% \]
% 
% (2) $f(x)$ 在 $[a,b]$ 上有瑕点.
% 
% 不妨设 $f(x)$ 只有一个瑕点 $c\in(a,b)$。由于 $f(x)$ 绝对可积, 因此存在 $\delta>0$, 使得 $a<c-\delta<c+\delta<b$ 和 $\int_{c-\delta}^{c+\delta}|f(x)|\mathrm dx<\dfrac{\varepsilon}{3}$, 从而对任意的 $\lambda$, 有
% \[
%   \left|\int_{c-\delta}^{c+\delta}f(x)\sin\lambda x\mathrm dx\right|
%   \leqslant\int_{c-\delta}^{c+\delta}|f(x)|\mathrm dx<\frac{\varepsilon}{3}.
% \]
% 再由于 $f(x)$ 在 $[a,c-\delta]$ 及 $[c+\delta,b]$ 上可积, 因此存在 $X>0$, 当 $\lambda>X$ 时, 有
% \[
%   \left|\int_a^{c-\delta}f(x)\sin\lambda x\mathrm dx\right|<\frac{\varepsilon}{3}
% \]
% 及
% \[
%   \left|\int_{c+\delta}^b f(x)\sin\lambda x\mathrm dx\right|<\frac{\varepsilon}{3},
% \]
% 从而当 $\lambda>X$ 时, 有
% \[
% \begin{aligned}
%   \left|\int_a^b f(x)\sin\lambda x\mathrm dx\right|
%   &\leqslant\left|\int_a^{c-\delta}f(x)\sin\lambda x\mathrm dx\right|
%     +\left|\int_{c-\delta}^{c+\delta}f(x)\sin\lambda x\mathrm dx\right|\\
%   &\quad+\left|\int_{c+\delta}^b f(x)\sin\lambda x\mathrm dx\right|\\
%   &<\frac{\varepsilon}{3}+\frac{\varepsilon}{3}+\frac{\varepsilon}{3}=\varepsilon.
% \end{aligned}
% \]
% 因此由 (1) 和 (2) 推出在定理的假设下有
% \[
%   \lim_{\lambda\to+\infty}\int_a^b f(x)\sin\lambda x\mathrm dx=0.
% \]
% 同理可证
% \[
%   \lim_{\lambda\to+\infty}\int_a^b f(x)\cos\lambda x\mathrm dx=0.
% \qedhere
% \]
% \end{proof}
% 
% 
% 现在我们回过头来考查以下等式:
% \[
%   S_n(x_0)-S_0
%   =\int_0^\pi\frac{f(x_0+t)+f(x_0-t)-2S_0}{2\pi\sin\frac t2}
%     \sin\left(n+\frac12\right)t\mathrm dt.
% \]
% 设 $f(x)$ 在 $[-\pi,\pi]$ 上可积或有瑕点时绝对可积, 我们来看该式当 $n\to\infty$ 时的变化情况.
% 
% 在 $t=0$ 的某个邻域外, 若 $G(t)$ ($G(t)$ 由 \eqref{eq:ma-12-2-1} 定义) 没有瑕点显然是可积的, 并且有瑕点时也绝对可积, 因此由黎曼-勒贝格引理 (引理~\ref{lemma:ma-12-2-1}) 有
% \[
%   \int_\delta^\pi G(t)\sin\left(n+\frac12\right)t\mathrm dt\to0
%   \quad (n\to\infty).
% \]
% 所以当 $n$ 趋于无穷时, $S_n(x_0)-S_0$ 是否趋于 $0$ 仅与 $f(x_0\pm t)$ 在 $t=0$ 附近的值有关, 换句话说仅与 $f(x)$ 在 $x_0$ 附近的值有关。因此我们有
% 
% \begin{theorem}[{黎曼局部化定理}]\label{theorem:ma-12-2-2}
% 设周期为 $2\pi$ 的函数 $f(x)$ 在 $[-\pi,\pi]$ 上可积或有瑕点时绝对可积, 则 $f(x)$ 的傅里叶级数在 $x_0\in[-\pi,\pi]$ 处的敛散性只与 $f(x)$ 在 $(x_0-\delta,x_0+\delta)$ 的值有关, 其中 $\delta>0$ 是一任意给定的正常数.
% \end{theorem}
% 
% 值得指出的是, 一个函数 $f(x)$ 的傅里叶系数是由 $f(x)$ 与基本三角函数系中每个函数的乘积在 $[-\pi,\pi]$ 上的积分得到。换句话说, 它们依赖于 $[-\pi,\pi]$ 上 $f(x)$ 的取值。而以上定理却告诉我们: 傅里叶级数在 $x_0$ 处的敛散性只依赖于 $f(x)$ 在 $x_0$ 附近的局部性质。这个结果有点出乎人们的意料.
% 
% 现在我们继续来研究 $f(x)$ 的傅里叶级数在 $x_0$ 处的敛散性。为了更加容易估计 $S_n(x_0)-S_0$, 我们注意到在 $f(x)$ 的狄利克雷积分中被积函数的 $2\sin\dfrac t2$ 可以用 $t$ 代替。换句话说, 若 $f(x)$ 在 $[0,\delta]$ 上可积或具有瑕点时绝对可积, 则当 $n\to\infty$ 时,
% \[
%   \int_0^\delta\frac{\sin\left(n+\frac12\right)t}{2\sin\frac t2}f(t)\mathrm dt
%   \quad\text{与}\quad
%   \int_0^\delta\frac{\sin\left(n+\frac12\right)t}{t}f(t)\mathrm dt
% \]
% 具有相同的敛散性, 并且它们收敛时具有相同的极限.
% 
% 事实上, 由
% \[
%   \lim_{t\to0}\left(\frac1{2\sin\frac t2}-\frac1t\right)
%   =\lim_{t\to0}\frac{t-2\sin\frac t2}{2t\sin\frac t2}
%   =\lim_{t\to0}\frac{t-2\left(\frac t2+o(t^2)\right)}{2t\sin\frac t2}=0,
% \]
% 定义
% \[
%   \varphi(t)=
%   \begin{cases}
%     \dfrac1{2\sin\frac t2}-\dfrac1t, & t\ne0,\\
%     0, & t=0,
%   \end{cases}
% \]
% 则 $\varphi(t)$ 在 $[0,\pi]$ 上连续, 因此有界, 从而 $f(t)\varphi(t)$ 在 $[0,\delta]$ 上可积或有瑕点时绝对可积。由黎曼-勒贝格引理有
% \[
%   \lim_{n\to\infty}\int_0^\delta f(t)\varphi(t)\sin\left(n+\frac12\right)t\mathrm dt=0.
% \]
% 由此即得
% \[
%   \lim_{n\to\infty}\int_0^\delta f(t)\frac{\sin\left(n+\frac12\right)t}{2\sin\frac t2}\mathrm dt
%   =
%   \lim_{n\to\infty}\int_0^\delta\frac{f(t)\sin\left(n+\frac12\right)t}{t}\mathrm dt.
% \]
% 
% \begin{remark}\label{remark:ma-12-source-17}
% 上述等式的确切意义是, 等式两边的极限具有相同的敛散性, 并且它们收敛时具有相同的极限。另外, 由黎曼-勒贝格引理及上面的分析, 若 $f(x)$ 在 $[0,\pi]$ 上可积或有瑕点时绝对可积, 则有
% \[
%   \lim_{n\to\infty}\int_0^\pi f(t)\frac{\sin\left(n+\frac12\right)t}{2\sin\frac t2}\mathrm dt
%   =
%   \lim_{n\to\infty}\int_0^\pi f(t)\frac{\sin\left(n+\frac12\right)t}{t}\mathrm dt.
% \]
% 利用上述等式, 我们可以计算出积分 $\int_0^{+\infty}\dfrac{\sin x}{x}\mathrm dx$.
% \end{remark}
% 
% 事实上, 我们已经知道该积分是收敛的, 取 $f(t)\equiv1$, 则有
% \[
% \begin{aligned}
%   1
%   &\equiv\lim_{n\to\infty}\frac2\pi\int_0^\pi
%     \frac{\sin\left(n+\frac12\right)t}{2\sin\frac t2}\mathrm dt\\
%   &=\lim_{n\to\infty}\frac2\pi\int_0^\pi
%     \frac{\sin\left(n+\frac12\right)t}{t}\mathrm dt\\
%   &=\lim_{n\to\infty}\frac2\pi\int_0^{(n+\frac12)\pi}\frac{\sin x}{x}\mathrm dx
%   =\frac2\pi\int_0^{+\infty}\frac{\sin x}{x}\mathrm dx.
% \end{aligned}
% \]
% 因此
% \[
%   \int_0^{+\infty}\frac{\sin x}{x}\mathrm dx=\frac{\pi}{2}.
% \]
% 
% \subsection{傅里叶级数的收敛判别法}
% \label{subsec:ma-12-2-2}
% 
% 对于给定的 $x_0\in[-\pi,\pi]$, 我们继续进行周期为 $2\pi$ 的函数 $f(x)$ 的傅里叶级数在 $x_0$ 处敛散性的讨论。根据上一小节推导, 我们知道 $f(x)$ 的傅里叶级数在 $x_0$ 处收敛到 $S_0$ 的充分必要条件是: 对充分小的正数 $\delta>0$, 有
% \begin{equation}
%   \lim_{n\to\infty}\int_0^\delta
%   \frac{f(x_0+t)+f(x_0-t)-2S_0}{t}
%   \sin\left(n+\frac12\right)t\mathrm dt=0.
%   \label{eq:ma-12-2-2}
% \end{equation}
% 该极限的收敛性还可转化为如何选择 $S_1$ 及 $S_2$, 使得
% \begin{equation}
%   \lim_{n\to\infty}\int_0^\delta
%   \frac{f(x_0+t)-S_1}{t}\sin\left(n+\frac12\right)t\mathrm dt=0
%   \label{eq:ma-12-2-3}
% \end{equation}
% 及
% \begin{equation}
%   \lim_{n\to\infty}\int_0^\delta
%   \frac{f(x_0-t)-S_2}{t}\sin\left(n+\frac12\right)t\mathrm dt=0.
%   \label{eq:ma-12-2-4}
% \end{equation}
% 显然当 \eqref{eq:ma-12-2-3} 及 \eqref{eq:ma-12-2-4} 成立时, 令 $S_0=\dfrac{S_1+S_2}{2}$, 则 \eqref{eq:ma-12-2-2} 必成立.
% 
% 我们首先来看一个简单的情形: 设 $f(x)$ 在 $x_0$ 处可导, 则
% \[
%   \lim_{t\to0}\frac{f(x_0+t)-f(x_0)}{t}=f'(x_0).
% \]
% 特别地，上述差商在 $t=0$ 的某个去心邻域内有界，所以 $t=0$ 不是该差商的无界瑕点。 因此取 $S_1=f(x_0)$, 则 \eqref{eq:ma-12-2-3} 成立。同理, 当 $f(x)$ 在 $x_0$ 处可导时, 取 $S_2=f(x_0)$, 则 \eqref{eq:ma-12-2-4} 成立。因此, 当 $f(x)$ 在 $x_0$ 处可导时, 取 $S_0=f(x_0)$, 则 \eqref{eq:ma-12-2-2} 成立.
% 
% 基于上述的结果, 我们进一步讨论傅里叶级数在一个区间上的敛散性。为此我们引入一个函数分段可微的概念.
% 
% 设函数 $f(x)$ 在区间 $[a,b]$ 上有定义。若存在 $[a,b]$ 的分割
% \[
%   \Delta: a=x_0<x_1<\cdots<x_n=b,
% \]
% 使得 $f(x)$ 仅以 $x_i$ 为第一类间断点, 并且在 $x_i(i=1,2,\cdots,n)$ 处, 存在推广的单侧导数, 即
% \[
%   \lim_{h\to0+0}\frac{f(x_i+h)-f(x_i+0)}{h}\triangleq f'_+(x_0)
% \]
% 与
% \[
%   \lim_{h\to0-0}\frac{f(x_i+h)-f(x_i-0)}{h}\triangleq f'_-(x_0)
% \]
% 存在 (在 $[a,b]$ 的端点处存在推广的左右导数); 而当 $x\in(x_{i-1},x_i)(i=1,2,\cdots,n)$ 时, $f'(x)$ 存在。若 $f(x)$ 满足以上条件, 则称 $f(x)$ 在 $[a,b]$ 上是分段可微的.
% 
% 对于分段可微函数, 若在 \eqref{eq:ma-12-2-3} 式左边取 $S_1=f(x_0+0)$, 在 \eqref{eq:ma-12-2-4} 式左边取 $S_2=f(x_0-0)$, 则 \eqref{eq:ma-12-2-3}, \eqref{eq:ma-12-2-4} 式同时成立。因此我们有
% 
% \begin{theorem}\label{theorem:ma-12-2-3}
% 设 $f(x)$ 是周期为 $2\pi$ 的函数, 且在 $[-\pi,\pi]$ 内分段可微, 则 $f(x)$ 的傅里叶级数处处收敛到
% \[
%   \frac{f(x-0)+f(x+0)}{2},
% \]
% 即
% \[
%   \frac{a_0}{2}+\sum_{n=1}^{+\infty}(a_n\cos nx+b_n\sin nx)
%   =
%   \frac{f(x-0)+f(x+0)}{2},
%   \quad x\in[-\pi,\pi].
% \]
% \end{theorem}
% 
% 下面我们来看一些更为精细的结果。对于 $x_0\in[-\pi,\pi]$, 我们令
% \begin{equation}
%   \varphi(t)=f(x_0+t)+f(x_0-t)-2S_0,
%   \label{eq:ma-12-2-5}
% \end{equation}
% 则 $\varphi(t)$ 在 $t=0$(即 $f(x)$ 在 $x_0$) 附近的性质对傅里叶级数的收敛起着至关重要的作用.
% 
% 由黎曼-勒贝格引理, 若要 \eqref{eq:ma-12-2-2} 成立, 只要 $\dfrac{\varphi(t)}{t}$ 在 $t=0$ 的邻域绝对可积即可。因此我们有
% 
% \begin{theorem}[{狄尼定理}]\label{theorem:ma-12-2-4}
% 设 $f(x)$ 是周期为 $2\pi$ 的函数, 在 $[-\pi,\pi]$ 上可积或有瑕点时绝对可积, 并且对于 $x_0\in[-\pi,\pi]$, 存在 $\delta>0$, 使得
% \[
%   \int_0^\delta\frac{|\varphi(t)|}{t}\mathrm dt<+\infty,
% \]
% 则 $f(x)$ 的傅里叶级数在 $x_0$ 处收敛到 $S_0$(其中 $\varphi(t)$ 见 \eqref{eq:ma-12-2-5}).
% \end{theorem}
% 
% 
% \begin{proof}
% 由定理所给条件知 $\dfrac{\varphi(t)}{(t)}$ 在 $[0,\delta]$ 上绝对可积。又由 $f(x)$ 的条件知 $\dfrac{\varphi(t)}{(t)}$ 在 $[\delta,\pi]$ 上可积或有瑕点时绝对可积, 因此 $\dfrac{\varphi(t)}{t}$ 在 $[0,\pi]$ 上绝对可积。由黎曼-勒贝格引理推出
% \[
%   \lim_{n\to\infty}[S_n(x_0)-S_0]
%   =\lim_{n\to\infty}\int_0^\pi
%   \frac{\varphi(t)}{t}\sin\left(n+\frac12\right)t\mathrm dt=0.
% \qedhere
% \]
% \end{proof}
% 
% 
% 从定理~\ref{theorem:ma-12-2-4} 立即可推出下面的李普西茨 (Lipschitz) 定理.
% 
% \begin{theorem}[{李普西茨定理}]\label{theorem:ma-12-2-5}
% 设 $f(x)$ 是周期为 $2\pi$ 的函数, 在 $[-\pi,\pi]$ 上瑕可积或有瑕点时绝对可积, 再设 $f(x)$ 在 $x_0\in[-\pi,\pi]$ 处满足 $\alpha$-赫尔德 ($\alpha>0$) 连续性, 即存在 $L>0,\delta>0$, 使得对于 $x\in U(x_0,\delta)$, 有
% \[
%   |f(x_0+t)-f(x_0)|\leqslant L|t|^\alpha,
% \]
% 则 $f(x)$ 的傅里叶级数在 $x_0$ 处收敛到 $f(x_0)$.
% \end{theorem}
% 
% 
% \begin{proof}
% 由 $f(x)$ 在 $x_0$ 处满足 $\alpha$-赫尔德 ($\alpha>0$) 连续性, 存在 $\delta>0$, 使得当 $t\in(-\delta,\delta)$ 时, 我们有
% \[
%   \frac{|f(x_0+t)-f(x_0)|}{t}\leqslant\frac{L}{t^{1-\alpha}}.
% \]
% 由于 $\alpha>0$, 我们推出
% \[
%   \int_0^\delta\frac{|\varphi(t)|}{t}\mathrm dt<+\infty,
% \]
% 其中 $\varphi(t)$ 见 \eqref{eq:ma-12-2-5}。由狄尼定理知 $f(x)$ 的傅里叶级数在 $x_0$ 处收敛到 $f(x_0)$.
% \end{proof}
% 
% 
% 最后我们来证明一个在理论上及应用上都十分有用的结果, 即单调函数的傅里叶级数必定收敛.
% 
% \begin{theorem}[{狄利克雷定理}]\label{theorem:ma-12-2-6}
% 设 $f(x)$ 是周期为 $2\pi$ 的函数, 在 $[-\pi,\pi]$ 上可积或有瑕点时绝对可积, 再设 $x_0\in[-\pi,\pi]$ 不是 $f(x)$ 的瑕点, 并且存在 $\delta_0>0$, 使得 $f(x)$ 在 $(x_0-\delta_0,x_0)$ 及 $(x_0,x_0+\delta_0)$ 内分别单调, 则 $f(x)$ 的傅里叶级数在 $x_0$ 处收敛到
% \[
%   \frac{f(x_0-0)+f(x_0+0)}{2}.
% \]
% \end{theorem}
% 
% 
% \begin{proof}
% 先设 $f(x)$ 在 $(x_0,x_0+\delta_0)$ 上单调。我们证明
% \begin{equation}
%   \lim_{\lambda\to+\infty}\int_0^{\delta_0}
%   \frac{f(x_0+t)-f(x_0+0)}{t}\sin\lambda t\mathrm dt=0.
%   \label{eq:ma-12-2-6}
% \end{equation}
% 不妨设 $f(x)$ 是单调上升的。由于
% \[
%   \int_0^{+\infty}\frac{\sin t}{t}\mathrm dt=\frac{\pi}{2},
% \]
% 故连续函数 $G(x)=\displaystyle\int_0^x\dfrac{\sin t}{t}\mathrm dt$ 在 $[0,+\infty)$ 上有界。由此存在常数 $M>1$, 使得对任意 $0\leqslant t_1<t_2<+\infty$, 有
% \[
%   \left|\int_{t_1}^{t_2}\frac{\sin t}{t}\mathrm dt\right|\leqslant M.
% \]
% 由 $f(x)$ 的单调性, 对于 $\forall\varepsilon>0$, $\exists 0<\delta<\delta_0$, 使得当 $0<t<\delta$ 时, 有
% \[
%   0\leqslant f(x_0+t)-f(x_0+0)<\frac{\varepsilon}{2M}.
% \]
% 为了估计积分, 我们记
% \[
% \begin{aligned}
%   \int_0^{\delta_0}\frac{f(x_0+t)-f(x_0+0)}{t}\sin\lambda t\mathrm dt
%   &=
%   \int_0^\delta\frac{f(x_0+t)-f(x_0+0)}{t}\sin\lambda t\mathrm dt\\
%   &\quad+
%   \int_\delta^{\delta_0}\frac{f(x_0+t)-f(x_0+0)}{t}\sin\lambda t\mathrm dt\\
%   &\triangleq I+J.
% \end{aligned}
% \]
% 
% 对于 $J$, 注意到 $\dfrac{f(x_0+t)-f(x_0+0)}{t}$ 在 $[\delta,\delta_0]$ 上可积或有瑕点时绝对可积, 因此由黎曼-勒贝格引理知, $\exists\lambda_0>0$, 当 $\lambda>\lambda_0$ 时, 有
% \[
%   |J|<\frac{\varepsilon}{2}.
% \]
% 
% 现在来估计积分 $I$。由定积分第二中值定理, 存在 $\xi\in(0,\delta)$, 使得
% \[
% \begin{aligned}
%   |I|
%   &=
%   \left|\int_0^\delta\frac{f(x_0+t)-f(x_0+0)}{t}\sin\lambda t\mathrm dt\right|\\
%   &=
%   |f(x_0+\delta)-f(x_0+0)|
%   \left|\int_\xi^\delta\frac1t\sin\lambda t\mathrm dt\right|\\
%   &=
%   |f(x_0+\delta)-f(x_0+0)|
%   \left|\int_{\lambda\xi}^{\lambda\delta}\frac{\sin t}{t}\mathrm dt\right|\\
%   &<
%   \frac{\varepsilon}{2M}M=\frac{\varepsilon}{2}.
% \end{aligned}
% \]
% 因此
% \[
%   \left|\int_0^{\delta_0}\frac{f(x_0+t)-f(x_0+0)}{t}\sin\lambda t\mathrm dt\right|
%   \leqslant |I|+|J|<\frac{\varepsilon}{2}+\frac{\varepsilon}{2}=\varepsilon.
% \]
% \eqref{eq:ma-12-2-6} 得证。同理可证当 $f(x)$ 在 $[0,\delta_0]$ 上单调下降时, \eqref{eq:ma-12-2-6} 也成立.
% 
% 由于 $f(x_0+t)$ 和 $f(x_0-t)$ 在 $[0,\delta_0]$ 上分别单调, 利用上面已证结果我们有
% \[
%   \lim_{n\to\infty}\int_0^{\delta_0}
%   \left[
%     f(x_0+t)+f(x_0-t)-2\frac{f(x_0+0)+f(x_0-0)}{2}
%   \right]
%   \frac{\sin\left(n+\frac12\right)t}{t}\mathrm dt=0.
% \qedhere
% \]
% \end{proof}
% 
% 
% \begin{remark}\label{remark:ma-12-source-25}
% 在定理的假设下, 由黎曼-勒贝格引理我们有
% \[
%   \lim_{n\to\infty}\int_{\delta_0}^\pi
%   \left[
%     f(x_0+t)+f(x_0-t)-2\frac{f(x_0+0)+f(x_0-0)}{2}
%   \right]
%   \frac{\sin\left(n+\frac12\right)t}{t}\mathrm dt=0,
% \]
% 因此
% \[
%   \lim_{n\to\infty}\int_0^\pi
%   \frac{f(x_0+t)+f(x_0-t)}{2}
%   \frac{\sin\left(n+\frac12\right)t}{t}\mathrm dt
%   =
%   \frac{f(x_0+0)+f(x_0-0)}{2}.
% \]
% \end{remark}
% 
% 另外, 从定理~\ref{theorem:ma-12-2-6} 的证明中我们可以看出, 当 $f(x)$ 在 $x_0$ 的 $\delta_0$ 邻域内单调时, 下述等式成立:
% \[
%   \lim_{\lambda\to+\infty}\int_0^\pi f(x_0\pm t)\frac{\sin\lambda t}{t}\mathrm dt
%   =\frac{\pi}{2}f(x_0\pm0).
% \]
% 
% 最后我们来小结一下本节的内容。在实际应用中, 下述的结论是经常用到的.
% 
% 设 $f(x)$ 是周期为 $2\pi$ 的函数, 若 $f(x)$ 满足以下条件之一:
% 
% (1) $f(x)$ 在 $[-\pi,\pi]$ 分段单调;
% 
% (2) $f(x)$ 在 $[-\pi,\pi]$ 分段可微,
% 
% 则对于 $\forall x\in[-\pi,\pi]$, $f(x)$ 的傅里叶级数收敛到
% \[
%   \frac{f(x-0)+f(x+0)}{2}.
% \]
% 
% 因此, 对 第~\ref{sec:ma-12-1}~节 中的例子, 我们有
% 
% \begin{example}\label{example:ma-12-2-1}
% 函数 $f(x)=x(-\pi\leqslant x<\pi)$ 的傅里叶级数为
% \[
%   \sum_{n=1}^{+\infty}(-1)^{n-1}\frac2n\sin nx
%   =
%   \begin{cases}
%     x, & x\in(-\pi,\pi),\\
%     0, & x=-\pi,\pi.
%   \end{cases}
% \]
% \end{example}
% 
% 
% \begin{example}\label{example:ma-12-2-2}
% 函数 $f(x)=\dfrac{\pi-x}{2}(x\in[0,2\pi))$ 的傅里叶级数为
% \[
%   \sum_{n=1}^{+\infty}\frac{\sin nx}{n}
%   =
%   \begin{cases}
%     \dfrac{\pi-x}{2}, & x\in(0,2\pi),\\
%     0, & x=0,2\pi.
%   \end{cases}
% \]
% \end{example}
% 
% 
% \begin{example}\label{example:ma-12-2-3}
% 函数 $f(x)=x^2(x\in[0,2\pi))$ 的傅里叶级数为
% \[
%   \frac{4\pi^2}{3}+4\sum_{n=1}^{+\infty}\frac{\cos nx}{n^2}
%   -4\pi\sum_{n=1}^{+\infty}\frac{\sin nx}{n}
%   =
%   \begin{cases}
%     x^2, & x\in(0,2\pi),\\
%     2\pi^2, & x=0,2\pi.
%   \end{cases}
% \]
% \end{example}
% 
% 
% \begin{example}\label{example:ma-12-2-4}
% 函数 $f(x)=x^2(x\in[0,\pi])$ 的正弦级数与余弦级数分别为
% \[
%   2\pi\sum_{n=1}^{+\infty}\frac{(-1)^n}{n}\sin nx
%   -\frac8\pi\sum_{n=1}^{+\infty}\frac{\sin(2n-1)x}{(2n-1)^3}
%   =
%   \begin{cases}
%     x^2, & x\in[0,\pi),\\
%     0, & x=\pi
%   \end{cases}
% \]
% 和
% \[
%   \frac{\pi^2}{3}+4\sum_{n=1}^{+\infty}\frac{(-1)^n}{n^2}\cos nx=x^2,\quad x\in[0,\pi].
% \]
% \end{example}
% 
% 
% \begin{example}\label{example:ma-12-2-5}
% 函数 $f(x)=x\cos x\left(-\dfrac{\pi}{2}\leqslant x<\dfrac{\pi}{2}\right)$ 的傅里叶级数为
% \[
%   \sum_{n=1}^{+\infty}\frac{(-1)^{n-1}}{\pi}\frac{16n}{(4n^2-1)^2}\sin2nx
%   =x\cos x,\quad x\in\left[-\frac{\pi}{2},\frac{\pi}{2}\right].
% \]
% \end{example}
% 
% 
% \begin{remark}\label{remark:ma-12-source-31}
% 在例~\ref{example:ma-12-2-5} 中函数的周期为 $\pi$, 但相应的傅里叶级数收敛性的结论仍真, 请读者自己证明该结论.
% \end{remark}
% 
% 
% \begin{example}\label{example:ma-12-2-6}
% 设 $0<\alpha<1$, 试求函数 $f(x)=\cos\alpha x(x\in(-\pi,\pi))$ 的傅里叶级数, 并求出该傅里叶级数的和函数.
% \end{example}
% 
% 
% \begin{solve}
% 显然, $f(x)$ 为偶函数, 因此
% \[
%   b_n=0,\quad n=1,2,\cdots;
% \]
% \[
%   a_0=\frac2\pi\int_0^\pi\cos\alpha x\mathrm dx=\frac{2}{\pi\alpha}\sin\alpha\pi,
% \]
% \[
% \begin{aligned}
%   a_n
%   &=\frac2\pi\int_0^\pi\cos\alpha x\cos nx\mathrm dx\\
%   &=\frac1\pi\int_0^\pi[\cos(\alpha-n)x+\cos(\alpha+n)x]\mathrm dx\\
%   &=\frac1\pi\left[
%     \frac{\sin(\alpha-n)x}{\alpha-n}
%     +\frac{\sin(\alpha+n)x}{\alpha+n}
%   \right]_0^\pi\\
%   &=(-1)^n\frac{2\alpha\sin\alpha\pi}{\pi(\alpha^2-n^2)},\quad n=1,2,\cdots.
% \end{aligned}
% \]
% 注意到 $\cos\alpha x$ 分段可微且连续, 因此当 $x\in[-\pi,\pi]$ 时, 有
% \[
%   \frac{\sin\alpha\pi}{\pi}\left(
%     \frac1\alpha+\sum_{n=1}^{+\infty}\frac{(-1)^n2\alpha}{\alpha^2-n^2}\cos nx
%   \right)=\cos\alpha x.
% \]
% 
% 利用以上展开式, 我们可以得到一些有趣的推论 (见本章习题)。最后我们简单地介绍一下关于傅里叶级数敛散性的研究历史。傅里叶在研究热传导问题时引进了傅里叶级数后, 当时他曾认为基本上所有函数的傅里叶级数都收敛。后来狄利克雷给出了傅里叶级数敛散性判别法。在狄利克雷的研究工作之后, 人们以为任何连续函数的傅里叶级数都收敛到该函数, 但人们在 1873 年给出了例子, 从这些例子可以看出一个连续函数的傅里叶级数可以在一个无穷点集上发散, 而且后来人们还举出了可积函数的傅里叶级数处处发散的例子。在傅里叶级数的敛散性这个问题上直到近代还有一些重要的研究工作, 由于这些工作的介绍要用到实变函数课程中的知识, 我们在这里不再赘述.
% \end{solve}
% 
% 
% \section{傅里叶级数的其他收敛性}
% \label{sec:ma-12-3}
% 
% 在上节中我们主要讨论了函数 $f(x)$ 的傅里叶级数的点收敛性。在本节中, 我们将讨论傅里叶级数有关的其他一些收敛性。不失一般性, 在本节中我们假定 $f(x)$ 是周期为 $2\pi$ 的函数。读者不难将本节的结果推广到 $f(x)$ 是以 $2T$ 为周期的函数的情形.
% 
% \subsection{连续函数的三角多项式一致逼近}
% \label{subsec:ma-12-3-1}
% 
% 我们在幂级数理论中已经证明了 $[a,b]$ 上连续函数必可由多项式一致逼近。在本小节中, 我们证明以 $2\pi$ 为周期的连续函数必可由三角多项式一致逼近, 即证明
% 
% \begin{theorem}[{魏尔斯特拉斯第二逼近定理}]\label{theorem:ma-12-3-1}
% 设 $f(x)$ 是以 $2\pi$ 为周期的连续函数, 则存在
% \[
%   T_n(x)=\frac{\alpha_0}{2}+\sum_{k=1}^n(\alpha_k\cos kx+\beta_k\sin kx)\quad(n=1,2,\cdots),
% \]
% 使得 $\forall\varepsilon>0$, 存在 $N\in\mathbb N$, 当 $n>N$ 时, 对一切 $x\in(-\infty,+\infty)$, 有
% \[
%   |T_n(x)-f(x)|<\varepsilon.
% \]
% \end{theorem}
% 
% 为了证明定理~\ref{theorem:ma-12-3-1}, 我们需要做些准备工作。从上节我们介绍的傅里叶级数敛散性研究的历史中可知, 一般来说, 定理~\ref{theorem:ma-12-3-1} 中的三角多项式序列不能取 $f(x)$ 的傅里叶级数的部分和序列。因此我们必须另找其他的三角多项式序列。我们回忆一下, 对一个序列 $\{x_n\}$ 来说, 若 $\lim\limits_{n\to\infty}x_n=a$, 则必有
% \[
%   \lim_{n\to\infty}\frac{x_1+\cdots+x_n}{n}=a,
% \]
% 而反之一般不真。因此 $\left\{y_n=\dfrac{x_1+\cdots+x_n}{n}\right\}$ 比 $\{x_n\}$ 来说有可能具有更好的收敛性.
% 
% 设 $f(x)$ 的傅里叶级数为
% \[
%   \frac{a_0}{2}+\sum_{n=1}^{+\infty}(a_n\cos nx+b_n\sin nx),
% \]
% $S_n(x)(n=0,1,2,\cdots)$ 为其部分和序列, 对于 $n=0,1,2,\cdots$, 我们记
% \[
%   S_n^*(x)=\frac{S_0(x)+S_1(x)+\cdots+S_n(x)}{n+1}.
% \]
% 我们将看到, 对于 $n\in\mathbb N$, $S_n^*(x)$ 也有很好的积分表达式。事实上, 我们有
% \[
% \begin{aligned}
%   S_n^*(x)
%   &=\frac{1}{(n+1)\pi}\int_{-\pi}^{\pi}f(x+t)
%     \left[
%       \frac{\displaystyle\sum_{k=0}^n\sin\left(k+\frac12\right)t}{2\sin\dfrac t2}
%     \right]\mathrm dt\\
%   &=\frac{1}{(n+1)\pi}\int_{-\pi}^{\pi}f(x+t)
%     \left[
%       \frac{\displaystyle\sum_{k=0}^n2\sin\dfrac t2\sin\left(k+\frac12\right)t}{4\sin^2\dfrac t2}
%     \right]\mathrm dt\\
%   &=\frac{1}{(n+1)\pi}\int_{-\pi}^{\pi}f(x+t)
%     \left[
%       \frac{\displaystyle\sum_{k=0}^n(\cos kt-\cos(k+1)t)}{4\sin^2\dfrac t2}
%     \right]\mathrm dt\\
%   &=\frac1\pi\int_{-\pi}^{\pi}f(x+t)
%     \frac{\sin^2\dfrac{n+1}{2}t}{2(n+1)\sin^2\dfrac t2}\mathrm dt\\
%   &=\frac1\pi\int_{-\pi}^{\pi}f(x+t)\Phi_n(t)\mathrm dt,
% \end{aligned}
% \]
% 其中
% \[
%   \Phi_n(t)=\frac{\sin^2\dfrac{n+1}{2}t}{2(n+1)\sin^2\dfrac t2}
% \]
% 称为费叶 (Fejer) 核。显然对任意的 $t$, 有 $\Phi_n(t)\geqslant0$。在上述积分表达式中令 $f(x)\equiv1$, 我们推出
% \[
%   \frac1\pi\int_{-\pi}^{\pi}\Phi_n(t)\mathrm dt=1.
% \]
% 
% 下面我们来证明定理~\ref{theorem:ma-12-3-1}.
% 
% \begin{proof}[{定理~\ref{theorem:ma-12-3-1} 的证明}]
% 我们将证明 $f(x)$ 的傅里叶级数部分和序列 $\{S_n(x)\}$ 的算术平均构成的序列 $\{S_n^*(x)\}$ 在 $[-\pi,\pi]$ 上一致收敛于 $f(x)$。对于任意的 $x\in[-\pi,\pi]$, 我们有下述不等式:
% \begin{equation}
% \begin{aligned}
%   |f(x)-S_n^*(x)|
%   &=\left|\frac1\pi\int_{-\pi}^{\pi}\Phi_n(t)[f(x)-f(x+t)]\mathrm dt\right|\\
%   &\leqslant\frac1\pi\int_{-\pi}^{\pi}\Phi_n(t)|f(x)-f(x+t)|\mathrm dt.
% \end{aligned}
% \label{eq:ma-12-3-1}
% \end{equation}
% 设 $|f(x)|$ 在 $[-\pi,\pi]$ 上的最大值为 $M$。注意到 $f(x)$ 以 $2\pi$ 为周期, 由 $f(x)$ 的连续性我们推出它在 $(-\infty,+\infty)$ 上一致连续, 从而对于 $\forall\varepsilon>0$, $\exists\delta>0$, 当 $x_1,x_2$ 满足 $|x_1-x_2|<\delta$ 时, 有
% \[
%   |f(x_1)-f(x_2)|<\frac{\varepsilon}{3}.
% \]
% 将 \eqref{eq:ma-12-3-1} 式右边的积分分解成 $I_1+I_2+I_3$, 其中
% \[
%   I_1=\frac1\pi\int_{-\pi}^{-\delta}\Phi_n(t)|f(x)-f(x+t)|\mathrm dt,
% \]
% \[
%   I_2=\frac1\pi\int_{-\delta}^{\delta}\Phi_n(t)|f(x)-f(x+t)|\mathrm dt,
% \]
% \[
%   I_3=\frac1\pi\int_{\delta}^{\pi}\Phi_n(t)|f(x)-f(x+t)|\mathrm dt.
% \]
% 现在我们来估计上面的每一个积分.
% 
% 对于 $I_2$, 我们有
% \[
% \begin{aligned}
%   I_2
%   &=\frac1\pi\int_{-\delta}^{\delta}\Phi_n(t)|f(x)-f(x+t)|\mathrm dt\\
%   &\leqslant\frac{\varepsilon}{3\pi}\int_{-\delta}^{\delta}\Phi_n(t)\mathrm dt\\
%   &\leqslant\frac{\varepsilon}{3\pi}\int_{-\pi}^{\pi}\Phi_n(t)\mathrm dt
%   =\frac{\varepsilon}{3}.
% \end{aligned}
% \]
% 
% 对于 $I_1$, 我们有
% \[
% \begin{aligned}
%   I_1
%   &=\frac1\pi\int_{-\pi}^{-\delta}\Phi_n(t)|f(x)-f(x+t)|\mathrm dt\\
%   &\leqslant 2M\frac1\pi\int_{-\pi}^{-\delta}
%   \frac{\sin^2\dfrac{n+1}{2}t}{2(n+1)\sin^2\dfrac t2}\mathrm dt\\
%   &\leqslant 2M\frac1\pi\int_{-\pi}^{-\delta}
%   \frac{1}{2(n+1)\sin^2\dfrac\delta2}\mathrm dt\\
%   &<\frac{M\pi}{\sin^2\dfrac\delta2}\cdot\frac1{n+1}.
% \end{aligned}
% \]
% 因此存在 $N_1>0$, 当 $n>N_1$ 时, 有
% \[
%   I_1<\frac{\varepsilon}{3}.
% \]
% 同理可证, 存在 $N_2>0$, 当 $n>N_2$ 时, 有
% \[
%   I_3<\frac{\varepsilon}{3}.
% \]
% 因此, 当 $n>N=\max\{N_1,N_2\}$ 时, 对于 $\forall x\in[-\pi,\pi]$, 有
% \[
%   |f(x)-S_n^*(x)|<\varepsilon.
% \]
% 注意到对于 $\forall n\in\mathbb N$, $S_n^*(x)$ 仍为一个至多 $n$ 阶的三角多项式 (且与 $S_n(x)$ 的阶相同), 定理~\ref{theorem:ma-12-3-1} 得证.
% \end{proof}
% 
% 
% \begin{example}\label{example:ma-12-3-1}
% 设函数 $f(x)$ 在 $(-\infty,+\infty)$ 上连续且以 $2\pi$ 为周期, 再设 $f(x)$ 的傅里叶级数处处收敛。证明 $f(x)$ 的傅里叶级数必处处收敛于 $f(x)$.
% \end{example}
% 
% 
% \begin{proof}
% 设
% \[
%   f(x)\sim\frac{a_0}{2}+\sum_{n=1}^{+\infty}(a_n\cos nx+b_n\sin nx).
% \]
% 对于 $\forall x\in(-\infty,+\infty)$ 记
% \[
%   S_0(x)=\frac{a_0}{2},
% \]
% \[
%   S_n(x)=\frac{a_0}{2}+\sum_{k=1}^n(a_k\cos kx+b_k\sin kx)\quad(n=1,2,\cdots),
% \]
% \[
%   S_n^*(x)=\frac1{n+1}\sum_{k=0}^nS_k(x)\quad(n=1,2,\cdots).
% \]
% 由于 $f(x)$ 的傅里叶级数处处收敛, 因此存在 $(-\infty,+\infty)$ 上定义的函数 $g(x)$, 使得对于 $\forall x\in(-\infty,+\infty)$, 有
% \[
%   \lim_{n\to\infty}S_n(x)=g(x).
% \]
% 因此, 对于 $\forall x\in(-\infty,+\infty)$, 有
% \[
%   \lim_{n\to\infty}S_n^*(x)=g(x).
% \]
% 再根据定理~\ref{theorem:ma-12-3-1}, 我们有 $S_n^*(x)\to f(x)(x\in(-\infty,+\infty))$, 从而有 $f(x)\equiv g(x)(x\in(-\infty,+\infty))$, 即
% \[
%   f(x)=\frac{a_0}{2}+\sum_{n=1}^{+\infty}(a_n\cos nx+b_n\sin nx),\quad x\in(-\infty,+\infty).
% \qedhere
% \]
% \end{proof}
% 
% 
% \subsection{傅里叶级数的均方收敛}
% \label{subsec:ma-12-3-2}
% 
% 我们已经对函数序列 (函数项级数) 的点收敛及一致收敛有了许多研究。对于傅里叶级数来说, 由于它与积分密切相关, 我们可以引入另一种意义下的收敛---均方收敛.
% 
% 为了讨论均方收敛, 我们必须对所涉及的函数附加一个稍强的条件。设函数 $f(x)$ 在区间 $[a,b]$ 上除了有限个瑕点外有定义, 在 $[a,b]$ 上的任何不包含瑕点的子区间 $[c,d]$ 上可积, 并且 $f^2(x)$ 在 $[a,b]$ 上的瑕积分是收敛的, 则我们称 $f(x)$ 在 $[a,b]$ 上平方可积。特别地, 当 $f(x)$ 在 $[a,b]$ 上可积时, $f(x)$ 在 $[a,b]$ 上必平方可积.
% 
% \begin{definition}\label{definition:ma-12-3-1}
% 设函数 $f(x),f_n(x)(n=1,2,\cdots)$ 在区间 $[a,b]$ 上平方可积, 并且满足
% \[
%   \lim_{n\to\infty}\int_a^b[f_n(x)-f(x)]^2\mathrm dx=0,
% \]
% 则称函数序列 $\{f_n(x)\}$ 在 $[a,b]$ 上均方收敛于 $f(x)$.
% \end{definition}
% 
% 函数序列的这种均方收敛性在一些数学分支的研究中起着重要的作用, 在傅里叶级数理论中也具有重要的理论和应用价值。下面我们先来研究一下平方可积函数类的一些初等性质.
% 
% 由于对区间 $[a,b]$ 上任意的函数 $f(x)$ 总成立
% \[
%   |f(x)|\leqslant\frac{f^2(x)+1}{2}.
% \]
% 因此, 当 $f(x)$ 在 $[a,b]$ 上平方可积时, $f(x)$ 在 $[a,b]$ 上必定绝对可积。这说明 $f(x)$ 平方可积的条件要比 $f(x)$ 绝对可积的条件强一些。对于区间 $[a,b]$ 上平方可积的函数 $f(x),g(x)$, 我们有以下的结论:
% 
% (1) $f(x)g(x)$ 在 $[a,b]$ 上绝对可积;
% 
% (2) $f(x)+g(x)$ 在 $[a,b]$ 上平方可积.
% 
% 事实上, 对于 $\forall x\in[a,b]$, 我们有
% \[
%   |f(x)g(x)|\leqslant\frac{f^2(x)+g^2(x)}{2},
% \]
% 因此 (1) 成立.
% 
% 记
% \[
% \begin{aligned}
%   \varphi(t)
%   &=\int_a^b(t|f|+|g|)^2\mathrm dx\\
%   &=t^2\int_a^bf^2(x)\mathrm dx
%     +2t\int_a^b|f(x)g(x)|\mathrm dx
%     +\int_a^bg^2(x)\mathrm dx.
% \end{aligned}
% \]
% 则对于 $\forall t\in\mathbb R$, 有 $\varphi(t)\geqslant0$。由此我们推知上式右边关于 $t$ 的一元二次方程的判别式 $\Delta\leqslant0$, 即
% \[
%   \int_a^b|f(x)g(x)|\mathrm dx
%   \leqslant
%   \left[\int_a^bf^2(x)\mathrm dx\right]^{\frac12}
%   \left[\int_a^bg^2(x)\mathrm dx\right]^{\frac12}.
% \]
% 由此不等式我们也可推出 $f(x)g(x)$ 在 $[a,b]$ 上的绝对可积性。利用这个不等式, 我们还有
% \[
% \begin{aligned}
%   \int_a^b[f(x)+g(x)]^2\mathrm dx
%   &=\int_a^bf^2(x)\mathrm dx+2\int_a^bf(x)g(x)\mathrm dx+\int_a^bg^2(x)\mathrm dx\\
%   &\leqslant\int_a^bf^2(x)\mathrm dx
%     +2\sqrt{\int_a^bf^2(x)\mathrm dx}\sqrt{\int_a^bg^2(x)\mathrm dx}
%     +\int_a^bg^2(x)\mathrm dx\\
%   &=\left[
%     \sqrt{\int_a^bf^2(x)\mathrm dx}
%     +\sqrt{\int_a^bg^2(x)\mathrm dx}
%   \right]^2.
% \end{aligned}
% \]
% 由此得到
% \[
%   \left\{\int_a^b[f(x)+g(x)]^2\mathrm dx\right\}^{\frac12}
%   \leqslant
%   \left[\int_a^bf^2(x)\mathrm dx\right]^{\frac12}
%   +
%   \left[\int_a^bg^2(x)\mathrm dx\right]^{\frac12}.
% \]
% 上面的不等式也称为柯西不等式, 从它可以直接看出 $f(x)+g(x)$ 是平方可积的。(2) 得证.
% 
% 现在我们来考虑函数 $f(x)$ 的傅里叶级数的均方收敛性问题。显然, 当 $f(x)$ 在 $[-\pi,\pi]$ 上可积或有瑕点时平方可积时, 对任意的三角多项式
% \[
%   T_n(x)=\frac{\alpha_0}{2}+\sum_{k=1}^n(\alpha_k\cos kx+\beta_k\sin kx),
% \]
% $f(x)-T_n(x)$ 也在 $[-\pi,\pi]$ 上可积或有瑕点时平方可积。设 $f(x)$ 的傅里叶级数为
% \[
%   \frac{a_0}{2}+\sum_{n=1}^{+\infty}(a_n\cos nx+b_n\sin nx),
% \]
% 令
% \[
%   \Delta_n=\left\{\frac1{2\pi}\int_{-\pi}^{\pi}[f(x)-T_n(x)]^2\mathrm dx\right\}^{\frac12},
% \]
% 则我们称 $\Delta_n$ 为 $f(x)-T_n(x)$ 的平均偏差。下面我们来计算 $\Delta_n^2$:
% \[
% \begin{aligned}
%   \Delta_n^2
%   &=\frac1{2\pi}\int_{-\pi}^{\pi}[f(x)-T_n(x)]^2\mathrm dx\\
%   &=\frac1{2\pi}\left[
%     \int_{-\pi}^{\pi}f^2(x)\mathrm dx
%     -2\int_{-\pi}^{\pi}f(x)T_n(x)\mathrm dx
%     +\int_{-\pi}^{\pi}T_n^2(x)\mathrm dx
%   \right]\\
%   &=\frac1{2\pi}\int_{-\pi}^{\pi}f^2(x)\mathrm dx
%     -\frac1\pi\int_{-\pi}^{\pi}f(x)
%     \left[\frac{\alpha_0}{2}+\sum_{k=1}^n(\alpha_k\cos kx+\beta_k\sin kx)\right]\mathrm dx\\
%   &\quad+\frac1{2\pi}\int_{-\pi}^{\pi}
%     \left[\frac{\alpha_0}{2}+\sum_{k=1}^n(\alpha_k\cos kx+\beta_k\sin kx)\right]^2\mathrm dx\\
%   &=\frac1{2\pi}\int_{-\pi}^{\pi}f^2(x)\mathrm dx
%     -\frac{\alpha_0a_0}{2}
%     -\sum_{k=1}^n(\alpha_ka_k+\beta_kb_k)
%     +\frac{\alpha_0^2}{4}
%     +\frac12\sum_{k=1}^n(\alpha_k^2+\beta_k^2)\\
%   &=\frac1{2\pi}\int_{-\pi}^{\pi}f^2(x)\mathrm dx
%     +\frac12\left[
%       -\alpha_0a_0+\frac{\alpha_0^2}{2}
%       -2\sum_{k=1}^n(\alpha_ka_k+\beta_kb_k)
%       +\sum_{k=1}^n(\alpha_k^2+\beta_k^2)
%     \right]\\
%   &=\frac1{2\pi}\int_{-\pi}^{\pi}f^2(x)\mathrm dx
%     +\frac12\left[
%       \frac{(\alpha_0-a_0)^2}{2}
%       +\sum_{k=1}^n\{(\alpha_k-a_k)^2+(\beta_k-b_k)^2\}
%     \right]\\
%   &\quad-\frac{a_0^2}{4}
%     -\frac12\sum_{k=1}^n(a_k^2+b_k^2).
% \end{aligned}
% \]
% 从上面的等式我们可以得出以下一系列的结论.
% 
% \begin{theorem}[{傅里叶级数最佳逼近}]\label{theorem:ma-12-3-2}
% 设函数 $f(x)$ 在 $[-\pi,\pi]$ 上可积或有瑕点时平方可积, 并设其傅里叶级数的部分和序列为 $\{S_n(x)\}$, 则对任何 $n\in\mathbb N$ 阶三角多项式 $T_n(x)$, 成立
% \[
%   \frac1{2\pi}\int_{-\pi}^{\pi}[f(x)-S_n(x)]^2\mathrm dx
%   \leqslant
%   \frac1{2\pi}\int_{-\pi}^{\pi}[f(x)-T_n(x)]^2\mathrm dx,
% \]
% 并且当且仅当 $S_n(x)\equiv T_n(x)$ 时上面不等式中等号成立.
% \end{theorem}
% 
% 
% \begin{remark}\label{remark:ma-12-source-40}
% 从上述分析我们还知道, 在定理~\ref{theorem:ma-12-3-2} 的条件下, 对于 $f(x)$ 的傅里叶级数部分和序列有以下性质: 当 $m<n$ 时, 有
% \[
%   \int_{-\pi}^{\pi}[f(x)-S_n(x)]^2\mathrm dx
%   \leqslant
%   \int_{-\pi}^{\pi}[f(x)-S_m(x)]^2\mathrm dx.
% \]
% \end{remark}
% 
% 由定理~\ref{theorem:ma-12-3-2} 可以解决我们在本小节的中心问题:
% 
% \begin{theorem}\label{theorem:ma-12-3-3}
% 设函数 $f(x)$ 在 $[-\pi,\pi]$ 上可积或有瑕点时平方可积, 则对 $f(x)$ 的傅里叶级数部分和序列 $\{S_n(x)\}$, 有
% \[
%   \lim_{n\to\infty}\int_{-\pi}^{\pi}[f(x)-S_n(x)]^2\mathrm dx=0.
% \]
% \end{theorem}
% 
% 
% \begin{proof}
% 先设 $f(x)$ 在 $[-\pi,\pi]$ 上连续且以 $2\pi$ 为周期, 由魏尔斯特拉斯第二逼近定理, 对于 $\forall\varepsilon>0$, 存在三角多项式 $T_{n_0}(x)$, 使得对 $\forall x\in[-\pi,\pi]$, 有
% \[
%   |f(x)-T_{n_0}(x)|<\sqrt{\frac{\varepsilon}{2\pi}}.
% \]
% 因此我们有
% \[
%   \int_{-\pi}^{\pi}[f(x)-T_{n_0}(x)]^2\mathrm dx<\varepsilon.
% \]
% 现设 $f(x)$ 在 $[-\pi,\pi]$ 上可积。我们知道存在连续函数 (可以是分段线性函数 (见例~\ref{example:ma-7-3-1} 的注)) $g(x)$, 使得
% \[
%   \int_{-\pi}^{\pi}[f(x)-g(x)]^2\mathrm dx<\frac{\varepsilon}{4}.
% \]
% 在 $x=\pm\pi$ 的充分小的邻域内对 $g(x)$ 的值进行修改, 我们可以假定满足上述不等式的 $g(x)$ 同时满足 $g(-\pi)=g(\pi)$.
% 
% 由上面所证, 存在 $n_0\in\mathbb N$, 使得
% \[
%   \int_{-\pi}^{\pi}[g(x)-T_{n_0}(x)]^2\mathrm dx<\frac{\varepsilon}{4}.
% \]
% 因此
% \[
% \begin{aligned}
%   \int_{-\pi}^{\pi}[f(x)-T_{n_0}(x)]^2\mathrm dx
%   &\leqslant
%   2\left\{
%     \int_{-\pi}^{\pi}[f(x)-g(x)]^2\mathrm dx
%     +\int_{-\pi}^{\pi}[g(x)-T_{n_0}(x)]^2\mathrm dx
%   \right\}\\
%   &<\varepsilon.
% \end{aligned}
% \]
% 最后设 $f(x)$ 有瑕点且在 $[-\pi,\pi]$ 上平方可积。此时不妨设 $f(x)$ 仅以 $x=\pi$ 为瑕点。由假设, $\exists\delta>0$, 使得
% \[
%   \int_{\pi-\delta}^{\pi}f^2(x)\mathrm dx<\frac{\varepsilon}{8}
%   \quad\text{及}\quad
%   \int_{-\pi}^{-\pi+\delta}f^2(x)\mathrm dx<\frac{\varepsilon}{8}.
% \]
% 令
% \[
%   \widetilde f(x)=
%   \begin{cases}
%     0, & x\in(\pi-\delta,\pi],\\
%     f(x), & x\in[-\pi+\delta,\pi-\delta],\\
%     0, & x\in[-\pi,-\pi+\delta).
%   \end{cases}
% \]
% 显然 $\widetilde f(x)$ 在 $[-\pi,\pi]$ 上可积, 且 $\widetilde f(-\pi)=\widetilde f(\pi)$, 从而存在三角多项式 $T_{n_0}(x)$, 使得
% \[
%   \int_{-\pi}^{\pi}[\widetilde f(x)-T_{n_0}(x)]^2\mathrm dx<\frac{\varepsilon}{4}.
% \]
% 因此
% \[
% \begin{aligned}
%   \int_{-\pi}^{\pi}[f(x)-T_{n_0}(x)]^2\mathrm dx
%   &\leqslant
%   2\int_{-\pi}^{\pi}[\widetilde f(x)-f(x)]^2\mathrm dx
%   +2\int_{-\pi}^{\pi}[\widetilde f(x)-T_{n_0}(x)]^2\mathrm dx\\
%   &\leqslant
%   2\int_{\pi-\delta}^{\pi}f^2(x)\mathrm dx
%   +2\int_{-\pi}^{-\pi+\delta}f^2(x)\mathrm dx
%   +\frac{\varepsilon}{2}\\
%   &<\frac{\varepsilon}{4}+\frac{\varepsilon}{4}+\frac{\varepsilon}{2}
%   =\varepsilon.
% \end{aligned}
% \]
% 综合上述的证明我们知, 若 $f(x)$ 在 $[-\pi,\pi]$ 上可积或有瑕点时平方可积, 则存在三角多项式 $T_{n_0}(x)$, 使得
% \[
%   \int_{-\pi}^{\pi}[f(x)-T_{n_0}(x)]^2\mathrm dx<\varepsilon.
% \]
% 由定理~\ref{theorem:ma-12-3-2}, 对于 $f(x)$ 的傅里叶级数部分和 $S_{n_0}(x)$, 有
% \[
%   \int_{-\pi}^{\pi}[f(x)-S_{n_0}(x)]^2\mathrm dx
%   \leqslant
%   \int_{-\pi}^{\pi}[f(x)-T_{n_0}(x)]^2\mathrm dx<\varepsilon.
% \]
% 因此当 $n>N(=n_0)$ 时, 有
% \[
%   \int_{-\pi}^{\pi}[f(x)-S_n(x)]^2\mathrm dx
%   \leqslant
%   \int_{-\pi}^{\pi}[f(x)-S_{n_0}(x)]^2\mathrm dx<\varepsilon.
% \qedhere
% \]
% \end{proof}
% 
% 
% 由定理~\ref{theorem:ma-12-3-3} 我们可以推出下面的帕塞瓦尔 (Parseval) 等式.
% 
% \begin{theorem}[{帕塞瓦尔等式}]\label{theorem:ma-12-3-4}
% 设函数 $f(x)$ 在 $[-\pi,\pi]$ 上可积或有瑕点时平方可积, 则有
% \[
%   \frac{a_0^2}{2}+\sum_{n=1}^{+\infty}(a_n^2+b_n^2)
%   =
%   \frac1\pi\int_{-\pi}^{\pi}f^2(x)\mathrm dx.
% \]
% \end{theorem}
% 
% 
% \begin{proof}
% 由定理~\ref{theorem:ma-12-3-3}, 对于 $\forall\varepsilon>0$, $\exists N\in\mathbb N$, 当 $n>N$ 时, 有
% \[
%   \frac1\pi\int_{-\pi}^{\pi}[f(x)-S_n(x)]^2\mathrm dx<\varepsilon,
% \]
% 其中 $\{S_n(x)\}$ 是 $f(x)$ 的傅里叶级数部分和序列。因此我们有
% \[
% \begin{aligned}
%   0
%   &\leqslant\frac1\pi\int_{-\pi}^{\pi}f^2(x)\mathrm dx
%   -\left[\frac{a_0^2}{2}+\sum_{k=1}^n(a_k^2+b_k^2)\right]\\
%   &=\frac1\pi\int_{-\pi}^{\pi}[f(x)-S_n(x)]^2\mathrm dx
%   <\varepsilon.
% \end{aligned}
% \qedhere
% \]
% \end{proof}
% 
% 
% 从帕塞瓦尔等式我们可以知道, 并不是任意一个三角级数都是某个平方可积函数的傅里叶级数。例如, 对于三角级数
% \[
%   \sum_{n=1}^{+\infty}\frac{\cos nx}{\sqrt n},
% \]
% 由于
% \[
%   \sum_{n=1}^{+\infty}\frac1{(\sqrt n)^2}=+\infty,
% \]
% 由帕塞瓦尔等式, 它不可能是某个平方可积函数 $f(x)$ 的傅里叶级数.
% 
% \begin{remark}\label{remark:ma-12-source-45}
% 设 $f(x)$ 与 $g(x)$ 均是以 $2\pi$ 为周期的函数且在 $[-\pi,\pi]$ 上平方可积, 并设 $f(x)$ 与 $g(x)$ 的傅里叶级数分别为
% \[
%   f(x)\sim\frac{a_0}{2}+\sum_{n=1}^{+\infty}(a_n\cos nx+b_n\sin nx),
% \]
% \[
%   g(x)\sim\frac{\alpha_0}{2}+\sum_{n=1}^{+\infty}(\alpha_n\cos nx+\beta_n\sin nx),
% \]
% 则成立
% \begin{equation}
%   \frac1\pi\int_{-\pi}^{\pi}f(x)g(x)\mathrm dx
%   =
%   \frac{a_0\alpha_0}{2}+\sum_{n=1}^{+\infty}(a_n\alpha_n+b_n\beta_n),
%   \label{eq:ma-12-3-2}
% \end{equation}
% 事实上, 从 $f+g$ 与 $f-g$ 的帕塞瓦尔等式相减即可得到 \eqref{eq:ma-12-3-2}。\eqref{eq:ma-12-3-2} 也称为广义帕塞瓦尔等式.
% \end{remark}
% 
% 
% \subsection{傅里叶级数的一致收敛性}
% \label{subsec:ma-12-3-3}
% 
% 利用帕塞瓦尔等式, 我们可以容易解决傅里叶级数的一致收敛、逐项求导以及逐项积分的问题.
% 
% \begin{theorem}[{傅里叶级数的一致收敛性}]\label{theorem:ma-12-3-5}
% 设函数 $f(x)$ 以 $2\pi$ 为周期, 在 $[-\pi,\pi]$ 上可导, 且 $f'(x)$ 在 $[-\pi,\pi]$ 上可积, 则 $f(x)$ 的傅里叶级数在 $(-\infty,+\infty)$ 上一致收敛到 $f(x)$.
% \end{theorem}
% 
% 
% \begin{proof}
% 设 $f(x)$ 的傅里叶级数为
% \[
%   f(x)\sim\frac{a_0}{2}+\sum_{n=1}^{+\infty}(a_n\cos nx+b_n\sin nx),
% \]
% $f'(x)$ 的傅里叶级数为
% \[
%   f'(x)\sim\frac{a_0'}{2}+\sum_{n=1}^{+\infty}(a_n'\cos nx+b_n'\sin nx).
% \]
% 我们有
% \[
%   a_0'=\frac1\pi\int_{-\pi}^{\pi}f'(x)\mathrm dx
%   =\frac1\pi f(x)\bigg|_{-\pi}^{\pi}=0,
% \]
% \[
% \begin{aligned}
%   a_n'
%   &=\frac1\pi\int_{-\pi}^{\pi}f'(x)\cos nx\mathrm dx\\
%   &=\frac1\pi f(x)\cos nx\bigg|_{-\pi}^{\pi}
%     +\frac n\pi\int_{-\pi}^{\pi}f(x)\sin nx\mathrm dx\\
%   &=nb_n,\quad n=1,2,\cdots.
% \end{aligned}
% \]
% 同理我们有
% \[
%   b_n'=-na_n,\quad n=1,2,\cdots.
% \]
% 因此对于 $\forall N\in\mathbb N$, 有
% \[
% \begin{aligned}
%   \sum_{n=1}^N(|a_n|+|b_n|)
%   &=\sum_{n=1}^N\frac{|a_n'|+|b_n'|}{n}\\
%   &\leqslant
%   \left[\sum_{n=1}^N(a_n'^2+b_n'^2)\right]^{\frac12}
%   \left(\sum_{n=1}^N\frac2{n^2}\right)^{\frac12}\\
%   &<\left(\frac{\pi^2}{3}\right)^{\frac12}
%   \left[\frac1\pi\int_{-\pi}^{\pi}(f'(x))^2\mathrm dx\right]^{\frac12}.
% \end{aligned}
% \]
% 因此 $f(x)$ 的傅里叶级数在 $(-\infty,+\infty)$ 上绝对一致收敛。由于 $f(x)$ 处处可导, $f(x)$ 的傅里叶级数处处收敛到 $f(x)$, 因此 $f(x)$ 的傅里叶级数在 $(-\infty,+\infty)$ 上一致收敛到 $f(x)$.
% \end{proof}
% 
% 
% \begin{theorem}[{傅里叶级数逐项微分定理}]\label{theorem:ma-12-3-6}
% 设函数 $f(x)$ 以 $2\pi$ 为周期, 对于 $x\in[-\pi,\pi]$, $f''(x)$ 存在且在 $[-\pi,\pi]$ 上可积, 再设 $f(x)$ 的傅里叶级数为
% \[
%   f(x)=\frac{a_0}{2}+\sum_{n=1}^{+\infty}(a_n\cos nx+b_n\sin nx),\quad x\in(-\infty,+\infty),
% \]
% 则
% \[
%   f'(x)=\sum_{n=1}^{+\infty}(nb_n\cos nx-na_n\sin nx),\quad x\in(-\infty,+\infty).
% \]
% \end{theorem}
% 
% 
% \begin{proof}
% 直接应用定理~\ref{theorem:ma-12-3-5} 及函数项级数逐项微分定理即得.
% \end{proof}
% 
% 
% 对于 $f(x)$ 的傅里叶级数的逐项积分问题, 我们在相当弱的条件下即能得到相应的结果.
% 
% \begin{theorem}[{傅里叶级数逐项积分定理}]\label{theorem:ma-12-3-7}
% 设函数 $f(x)$ 在 $[0,2\pi]$ 上可积, 且以 $2\pi$ 为周期, 再设
% \[
%   f(x)\sim\frac{a_0}{2}+\sum_{n=1}^{+\infty}(a_n\cos nx+b_n\sin nx),
% \]
% 则
% \[
%   \int_0^xf(t)\mathrm dt
%   =
%   \frac{a_0}{2}x+\sum_{n=1}^{+\infty}
%   \left[
%     \frac{a_n}{n}\sin nx+\frac{b_n(1-\cos nx)}{n}
%   \right]
% \]
% 在 $[0,2\pi]$ 成立.
% \end{theorem}
% 
% 
% \begin{proof}
% 对于 $x\in[0,2\pi]$, 作
% \[
%   g(t)=
%   \begin{cases}
%     \dfrac{\pi}{2}, & t=0,t=x,\\
%     \pi, & t\in(0,x),\\
%     0, & t\in(x,2\pi).
%   \end{cases}
% \]
% 设 $g(t)$ 的傅里叶级数为
% \[
%   \frac{\alpha_0}{2}+\sum_{n=1}^{+\infty}(\alpha_n\cos nt+\beta_n\sin nt),
% \]
% 则由计算直接得出
% \[
%   \alpha_0=x,\quad
%   \alpha_n=\frac{\sin nx}{n}\quad(n=1,2,\cdots),\quad
%   \beta_n=\frac{1-\cos nx}{n}\quad(n=1,2,\cdots).
% \]
% 由关于 $f(t)$ 与 $g(t)$ 的帕塞瓦尔等式即得
% \[
%   \frac1\pi\int_0^x\pi f(t)\mathrm dt
%   =
%   \frac{a_0}{2}x
%   +\sum_{n=1}^{+\infty}
%   \left[
%     \frac{a_n}{n}\sin nx+\frac{b_n(1-\cos nx)}{n}
%   \right],
%   \quad x\in[0,2\pi].
% \qedhere
% \]
% \end{proof}
% 
% 
% \begin{example}\label{example:ma-12-3-2}
% 求函数 $f(x)=\dfrac{x^2}{2}-\pi x(x\in[0,2\pi])$ 的傅里叶级数, 并求其和函数.
% \end{example}
% 
% 
% \begin{solve}
% 由例~\ref{example:ma-12-2-2} 有
% \[
%   \frac{\pi-x}{2}\sim\sum_{n=1}^{+\infty}\frac{\sin nx}{n},
% \]
% 因此由定理~\ref{theorem:ma-12-3-7}, 对于 $\forall x\in[0,2\pi]$, 有
% \[
%   \int_0^x\frac{\pi-t}{2}\mathrm dt
%   =
%   \sum_{n=1}^{+\infty}\int_0^x\frac{\sin nt}{n}\mathrm dt,
% \]
% 即
% \[
%   \frac{x^2}{4}-\frac{\pi}{2}x
%   =
%   -\sum_{n=1}^{+\infty}\frac1{n^2}
%   +\sum_{n=1}^{+\infty}\frac{\cos nx}{n^2},\quad x\in[0,2\pi].
% \]
% 所以
% \[
%   \frac{x^2}{2}-\pi x
%   =
%   -\frac{\pi^2}{3}+\sum_{n=1}^{+\infty}\frac{2\cos nx}{n^2},\quad x\in[0,2\pi].
% \qedhere
% \]
% \end{solve}
% 
% 
% \section*{习题十二}
% \phantomsection
% \addcontentsline{toc}{section}{习题十二}
% \label{exercises:ma-12}
% 
% 1。求下列周期为 $2\pi$ 的函数的傅里叶级数:
% 
% (1) $f(x)=|x|,\ -\pi\leqslant x<\pi$;
% 
% (2) $f(x)=e^{\alpha x},\ -\pi\leqslant x<\pi$;
% 
% (3) $f(x)=x\cos x,\ -\pi\leqslant x<\pi$;
% 
% (4) $f(x)=\cos^2x,\ -\pi\leqslant x<\pi$;
% 
% (5)
% \[
%   f(x)=
%   \begin{cases}
%     \dfrac{\pi}{4}, & 0<x<\pi,\\
%     -\dfrac{\pi}{4}, & -\pi\leqslant x\leqslant0;
%   \end{cases}
% \]
% 
% (6) $f(x)=\sin^4x$;
% 
% (7)
% \[
%   f(x)=\frac{1-r^2}{1-2r\cos x+r^2}\quad(|r|<1);
% \]
% 
% (8) $f(x)=\ln\left|\sin\dfrac x2\right|$.
% 
% % 整合来源：math-13.tex
% \chapter{多元函数的极限和连续}
% \label{chap:ma-13}
% 
% 在本套教材的第一册与第二册中, 我们已经系统地学习了一元微积分与级数理论。但在理论与实践中, 仅仅一元函数远远不能满足需要。这是因为, 在许多事物的变化过程中, 一个变量的变化过程往往依赖于多个变量。就拿我们每天生活的空间来说, 它是一个三维的立体空间, 因此几乎所有跟空间位置有关的变量一般都要用空间点的坐标来描述, 从而它们就不太可能用一元函数来刻画.
% 
% 另外, 即使在数学研究中, 由于一元函数的研究仅仅是局限于数轴 $\mathbb R$ 的子集上定义的函数, 它们基本上已不再是现代数学研究的主要对象。在当今的数学研究中, 大部分的研究对象都是关于高维空间 ($n(n\geqslant2)$ 维空间) 的一些问题。因此, 我们对多元函数 (映射) 的学习是十分必要的.
% 
% 多元微积分的主要内容是将一元函数的微积分理论推广到高维空间上的多元函数。大家会发现, 我们将平行于一元微积分的基本理论来研究多元微积分。值得指出的是, 由于多元函数的微积分理论是建立在一元微积分的基础之上的, 读者如果具备一元微积分的坚实基础, 并且有较好的空间想象能力, 就能学好多元微积分.
% 
% \section{\texorpdfstring{欧氏空间 $\mathbb R^n$}{欧氏空间 Rⁿ}}
% \label{sec:ma-13-1}
% 
% \subsection{\texorpdfstring{欧氏空间 $\mathbb R^n$}{欧氏空间 Rⁿ}}
% \label{subsec:ma-13-1-1}
% 
% 在本节中, 我们先来讨论多元函数定义域的问题。多元函数的定义域是高维空间的子集, 这些子集相对于 $\mathbb R$ 中的子集将更为复杂。为了研究高维空间的子集, 我们必须先研究一下它们所在的空间
% \[
%   \mathbb R^n=\{(x_1,x_2,\cdots,x_n):x_i\in\mathbb R,\ i=1,2,\cdots,n\}
%   =\underbrace{\mathbb R\times\mathbb R\times\cdots\times\mathbb R}_{n\text{ 个}}.
% \]
% 
% 读者应特别注意当 $n=2$ 时的情形。多元微积分与一元微积分的许多本质区别将在 $n=1$ 和 $n=2$ 两者之间发生, 对于 $n>2$ 的情形则与 $n=2$ 的情形具有很大的相似性.
% 
% 在以下讨论中我们将假定 $n\geqslant2$。记 $x=(x_1,x_2,\cdots,x_n)\in\mathbb R^n$, 我们称 $x$ 为 $\mathbb R^n$ 中的一个点或向量, $x_i(i=1,2,\cdots,n)$ 称为 $x$ 的第 $i$ 个坐标或分量。今后, 我们也常常用列向量 $(x_1,x_2,\cdots,x_n)^{\mathrm T}$ 来表示同一个 $x$, 这里 $(x_1,x_2,\cdots,x_n)^{\mathrm T}$ 是 $(x_1,x_2,\cdots,x_n)$ 的转置。另外, 记 $\mathbf 0=(0,0,\cdots,0)$ 为 $\mathbb R^n$ 中的原点或零向量。在代数课程中, 我们已经在 $\mathbb R^n$ 中引进了加法与数乘运算。由于在今后要经常用到它们, 在这里我们做一下简单介绍.
% 
% $\mathbb R^n$ 中的加法运算定义如下: 设 $x=(x_1,x_2,\cdots,x_n), y=(y_1,y_2,\cdots,y_n)\in\mathbb R^n$, 定义
% \[
%   x+y=(x_1+y_1,x_2+y_2,\cdots,x_n+y_n)\in\mathbb R^n,
% \]
% 并称 $x+y$ 为 $x$ 与 $y$ 的和.
% 
% $\mathbb R^n$ 中的数乘运算定义如下: 设 $\alpha\in\mathbb R, x=(x_1,x_2,\cdots,x_n)\in\mathbb R^n$, 定义
% \[
%   \alpha x=(\alpha x_1,\alpha x_2,\cdots,\alpha x_n)\in\mathbb R^n,
% \]
% 并称 $\alpha x$ 为 $\alpha$ 与 $x$ 的数乘.
% 
% 这两种运算称为 $\mathbb R^n$ 中的线性运算。在 $n=2,3$ 时, 它们具有鲜明的几何意义。对于 $\forall x,y,z\in\mathbb R^n,\alpha,\beta\in\mathbb R$, 容易验证它们满足:
% \[
% \begin{gathered}
%   (1)\quad \text{\keyterm{交换律}}\quad x+y=y+x;\\
%   (2)\quad \text{\keyterm{结合律}}\quad (x+y)+z=x+(y+z),\quad (\alpha\beta)x=\alpha(\beta x);\\
%   (3)\quad \text{\keyterm{分配律}}\quad \alpha(x+y)=\alpha x+\alpha y,\quad
%   (\alpha+\beta)x=\alpha x+\beta x.
% \end{gathered}
% \]
% 另外, 在加法运算中, 存在零元素 $\mathbf 0\in\mathbb R^n$, 它满足: 对于 $\forall x\in\mathbb R^n$, 有
% \[
%   x+\mathbf 0=x;
% \]
% 在数乘运算中, 存在单位元 $1\in\mathbb R$, 它满足: 对于 $\forall x\in\mathbb R^n$, 有
% \[
%   1x=x.
% \]
% 
% 对 $\mathbb R^n$ 赋予上述线性运算后, 我们称 $\mathbb R^n$ 为一个 $n$ 维向量空间 (简称空间)。在这个空间中, 它还有一个重要的运算——内积运算, 它的定义如下:
% 
% 设 $x=(x_1,x_2,\cdots,x_n), y=(y_1,y_2,\cdots,y_n)\in\mathbb R^n$, 则 $x$ 与 $y$ 的内积定义为
% \[
%   xy=\sum_{i=1}^n x_iy_i\in\mathbb R.
% \]
% 从内积的定义容易看出它具有以下一些基本性质:
% \[
% \begin{aligned}
%   &(1)\quad \text{\keyterm{正定性}}\quad \text{对于 } \forall x\in\mathbb R^n,\ \text{有 } xx\geqslant0,\ \text{并且上述等号成立当且仅当 } x=\mathbf 0;\\
%   &(2)\quad \text{\keyterm{对称性}}\quad \text{对于 } \forall x,y\in\mathbb R^n,\ \text{有 } xy=yx;\\
%   &(3)\quad \text{对于 } \forall x,y,z\in\mathbb R^n,\ \text{有 } x(y+z)=xy+xz;\\
%   &(4)\quad \text{对于 } \forall\alpha\in\mathbb R \text{ 和 } \forall x,y\in\mathbb R^n,\ \text{有 }(\alpha x)y=\alpha(xy).
% \end{aligned}
% \]
% 向量空间 $\mathbb R^n$ 有了内积运算后, 我们称 $\mathbb R^n$ 为欧几里得 (Euclid) 空间或欧氏空间。利用内积运算, 我们定义向量 $x\in\mathbb R^n$ 的模如下:
% \[
%   \|x\|=\sqrt{xx}=\sqrt{\sum_{i=1}^n x_i^2}.
% \]
% 在代数学中, 我们已经知道, $\mathbb R^n(n\geqslant2)$ 中两个非零向量 $x$ 与 $y$ 的内积
% \[
%   xy=\|x\|\cdot\|y\|\cos(x,y),
% \]
% 其中 $(x,y)$ 是向量 $x$ 与 $y$ 的夹角。由此我们可以清楚地知道内积的几何意义.
% 
% 利用向量的模, 我们可以给出 $\mathbb R^n$ 中两个点之间的距离的定义.
% 
% \begin{definition}\label{definition:ma-13-1-1}
% 设 $x=(x_1,x_2,\cdots,x_n)$ 与 $y=(y_1,y_2,\cdots,y_n)$ 为 $\mathbb R^n$ 中任意两个点, 则 $x$ 与 $y$ 的距离定义为
% \[
%   |x-y|=\|x-y\|=\sqrt{\sum_{i=1}^n(x_i-y_i)^2}.
% \]
% 显然, 在 $\mathbb R,\mathbb R^2$ 或 $\mathbb R^3$ 中, 两个点 $x$ 与 $y$ 的距离即是连接 $x$ 与 $y$ 的线段的长度。请读者注意, 我们用 $|x-y|$ (而不是用 $\|x-y\|$) 来表示 $x,y$ 之间的距离主要是为了记号简便, 且这个记号与 $\mathbb R$ 中两点之间距离的相应记号保持一致。由于对于 $\forall x\in\mathbb R^n$, 我们有 $\|x\|=|x-\mathbf 0|=|x|$, 因此, 今后我们也用 $|x|$ 来记 $x$ 的模.
% \end{definition}
% 
% 从距离的定义容易推出它满足以下的性质:
% \[
% \begin{aligned}
%   &(1)\quad \text{\keyterm{正定性}}\quad
%   \text{对于 } \forall x,y\in\mathbb R^n,\ \text{有 } |x-y|\geqslant0,\ \text{并且 } |x-y|=0 \text{ 的充分必要条件是 } x=y;\\
%   &(2)\quad \text{\keyterm{对称性}}\quad
%   \text{对于 } \forall x,y\in\mathbb R^n,\ \text{有 } |x-y|=|y-x|;\\
%   &(3)\quad \text{\keyterm{三角不等式}}\quad
%   \text{对于 } \forall x,y,z\in\mathbb R^n,\ \text{有 } |x-z|\leqslant |x-y|+|y-z|.
% \end{aligned}
% \]
% 距离与向量的模两个概念是可以相互转化的, 因此我们也称定义了距离后的空间 $\mathbb R^n$ 为欧氏空间。当 $n=2$ 时, 我们常用 $(x,y)$ 来表示平面 $\mathbb R^2$ 中的点; 当 $n=3$ 时, 用 $(x,y,z)$ 来表示空间 $\mathbb R^3$ 中的点。今后在 $\mathbb R^2$ 中, 为了记号方便, 我们也用 $\boldsymbol{i}$ 来表示单位向量 $(1,0)$, 用 $\boldsymbol{j}$ 来表示单位向量 $(0,1)$, 并称它们为单位坐标向量。因此, 对于 $\forall(x,y)\in\mathbb R^2$, 我们有
% \[
%   (x,y)=x\boldsymbol{i}+y\boldsymbol{j}.
% \]
% 相应地, 在 $\mathbb R^3$ 中, 我们则记 $\boldsymbol{i}=(1,0,0), \boldsymbol{j}=(0,1,0), \boldsymbol{k}=(0,0,1)$。于是对于 $\forall(x,y,z)\in\mathbb R^3$, 有
% \[
%   (x,y,z)=x\boldsymbol{i}+y\boldsymbol{j}+z\boldsymbol{k}.
% \]
% 
% \begin{example}\label{example:ma-13-1-1}
% 设
% \[
%   A=
%   \begin{pmatrix}
%     a_{11} & a_{12} & \cdots & a_{1n}\\
%     a_{21} & a_{22} & \cdots & a_{2n}\\
%     \vdots & \vdots &        & \vdots\\
%     a_{m1} & a_{m2} & \cdots & a_{mn}
%   \end{pmatrix}
% \]
% 是一个 $m\times n$ 矩阵, 其中 $a_{ji}\in\mathbb R(j=1,2,\cdots,m;\ i=1,2,\cdots,n)$。定义
% \[
%   \|A\|=\left(\sum_{j=1}^m\sum_{i=1}^n a_{ji}^2\right)^{\frac12}
% \]
% 对于 $\forall x=(x_1,x_2,\cdots,x_n)^{\mathrm T}\in\mathbb R^n$, 证明 $|Ax|\leqslant \|A\||x|$.
% \end{example}
% 
% 
% \begin{proof}
% 利用柯西-施瓦茨不等式我们有
% \[
% \begin{aligned}
%   |Ax|^2
%   &=\sum_{j=1}^m\left(\sum_{i=1}^n a_{ji}x_i\right)^2
%     \leqslant \sum_{j=1}^m\left(\sum_{i=1}^n a_{ji}^2\sum_{i=1}^n x_i^2\right)\\
%   &=\left(\sum_{j=1}^m\sum_{i=1}^n a_{ji}^2\right)
%     \left(\sum_{i=1}^n x_i^2\right)
%     =(\|A\||x|)^2.
% \end{aligned}
% \]
% 对上面不等式的两边分别开方即得所证.
% \end{proof}
% 
% 
% \begin{remark}\label{remark:ma-13-source-4}
% $\|A\|=\left(\sum\limits_{j=1}^m\sum\limits_{i=1}^n a_{ji}^2\right)^{\frac12}$ 称为矩阵 $A$ 的范数, 在本书后面章节我们还会遇到它.
% \end{remark}
% 
% 
% \subsection{点列极限}
% \label{subsec:ma-13-1-2}
% 
% 下面我们给出欧氏空间 $\mathbb R^n$ 中邻域的概念.
% 
% \begin{definition}\label{definition:ma-13-1-2}
% 设 $x_0=(x_1^0,x_2^0,\cdots,x_n^0)\in\mathbb R^n,\delta>0$, 称集合
% \[
%   U(x_0,\delta)=\{x=(x_1,x_2,\cdots,x_n)\in\mathbb R^n:|x-x_0|<\delta\}
% \]
% 为以 $x_0$ 为心的 $\delta$ 邻域; 称集合 $U_0(x_0,\delta)=U(x_0,\delta)\setminus\{x_0\}$ 为 $x_0$ 的 $\delta$ 去心邻域.
% \end{definition}
% 
% 上述定义的邻域通常称为球形邻域。我们经常要用到的另外一种邻域是方形邻域。设 $x_0=(x_1^0,x_2^0,\cdots,x_n^0)\in\mathbb R^n,\delta>0$, 定义
% \[
%   N(x_0,\delta)=\{x=(x_1,x_2,\cdots,x_n): |x_i-x_i^0|<\delta,\ i=1,2,\cdots,n\};
% \]
% 并称它为 $x_0$ 的方形邻域。容易证明这两种邻域有下面的包含关系:
% \[
%   U(x_0,\delta)\subset N(x_0,\delta)\subset U(x_0,\sqrt n\delta).
% \]
% 另外, 我们称 $N_0(x_0,\delta)=N(x_0,\delta)\setminus\{x_0\}$ 为方形去心邻域.
% 
% 以上所定义的邻域中, 当 $n=1$ 时, 球形邻域与方形邻域是一致的, 它们就是实数轴上的开区间。当 $n=2$ 时, 球形邻域就是开圆盘, 方形邻域就是开正方形。当 $n=3$ 时, 球形邻域就是开球, 方形邻域就是开立方体.
% 
% 下面我们利用距离来描述 $\mathbb R^n$ 中点列的极限。如果点列 $\{x_k\}$ 中的点 $x_k$ 到某个点 $x_0$ 的距离趋于零, 即
% \[
%   \lim_{k\to\infty}|x_k-x_0|=0,
% \]
% 那么我们就说点列 $\{x_k\}$ 收敛于 $x_0$。用邻域的语言可以给出如下定义.
% 
% \begin{definition}\label{definition:ma-13-1-3}
% 设 $\{x_k\}$ 是 $\mathbb R^n$ 中的一个点列, 若存在 $x_0\in\mathbb R^n$, 使得对于 $\forall\varepsilon>0,\exists K\in\mathbb N$, 当 $k>K$ 时, 有 $|x_k-x_0|<\varepsilon$, 即 $x_k\in U(x_0,\varepsilon)$, 则称 $\{x_k\}$ 是收敛点列, 并称 $\{x_k\}$ 收敛于 $x_0$, 记做
% \[
%   \lim_{k\to\infty}x_k=x_0.
% \]
% 这时也称 $x_0$ 为 $\{x_k\}$ 的极限。若不存在 $x_0\in\mathbb R^n$, 使得 $\lim\limits_{k\to\infty}|x_k-x_0|=0$, 则称 $\{x_k\}$ 发散.
% \end{definition}
% 
% 在本章中, 为了避免与 $\mathbb R^n$ 中的 $n$ 混淆, 我们常用 $\{x_k\}$ 来表示 $\mathbb R^n$ 中的点列, 而不用 $\{x_n\}$。设 $x_0=(x_1^0,x_2^0,\cdots,x_n^0)$, $x_k=(x_1^k,x_2^k,\cdots,x_n^k)$, 则有下面的定理.
% 
% \begin{theorem}\label{theorem:ma-13-1-1}
% 设 $\{x_k\}$ 是 $\mathbb R^n$ 中的一个点列, $x_0=(x_1^0,x_2^0,\cdots,x_n^0)\in\mathbb R^n$, 则 $\lim\limits_{k\to\infty}x_k=x_0$ 的充分必要条件是: 对于 $\forall i(1\leqslant i\leqslant n)$, 有
% \[
%   \lim_{k\to\infty}x_i^k=x_i^0.
% \]
% \end{theorem}
% 
% 
% \begin{proof}
% 对于 $\forall i=1,2,\cdots,n$ 和 $\forall k\in\mathbb N$, 有
% \[
%   |x_i^k-x_i^0|\leqslant |x_k-x_0|\leqslant \sum_{j=1}^n|x_j^k-x_j^0|.
% \]
% 由此不等式容易推出定理~\ref{theorem:ma-13-1-1}.
% \end{proof}
% 
% 
% 若 $\mathbb R^n$ 中的集合 $E$ 包含于某个球形邻域 $U(0,M)$ 中, 则称 $E$ 是有界集。若点列 $\{x_k\}$ 的所有点组成的集合是有界集, 则称 $\{x_k\}$ 是有界点列。与数列的情形类似, $\mathbb R^n$ 中的收敛点列具有下面的性质.
% 
% \begin{property}\label{property:ma-13-1-2}
% 收敛点列的极限是唯一的.
% \end{property}
% 
% 
% \begin{property}\label{property:ma-13-1-3}
% 收敛点列一定是有界的.
% \end{property}
% 
% 
% \begin{property}\label{property:ma-13-1-4}
% 若 $\mathbb R^n$ 中的点列 $\{x_k\},\{y_k\}$ 满足
% \[
%   \lim_{k\to\infty}x_k=x_0,\qquad \lim_{k\to\infty}y_k=y_0,
% \]
% 且 $\alpha,\beta\in\mathbb R$, 则
% \[
%   (1)\quad \lim_{k\to\infty}(\alpha x_k+\beta y_k)=\alpha x_0+\beta y_0;
% \]
% \[
%   (2)\quad \lim_{k\to\infty}x_ky_k=x_0y_0.
% \]
% 对于以上性质, 我们只证性质~\ref{property:ma-13-1-4} 的 (2), 读者应该注意这个性质说的是两个收敛点列的内积收敛到点列极限的内积.
% \end{property}
% 
% 
% \begin{proof}
% 设 $x_k=(x_1^k,x_2^k,\cdots,x_n^k)$, $y_k=(y_1^k,y_2^k,\cdots,y_n^k)$, $x_0=(x_1^0,x_2^0,\cdots,x_n^0)$, $y_0=(y_1^0,y_2^0,\cdots,y_n^0)$。由定理~\ref{theorem:ma-13-1-1} 知, 对于 $\forall i=1,2,\cdots,n$, 有
% \[
%   \lim_{k\to\infty}x_i^k=x_i^0,\qquad \lim_{k\to\infty}y_i^k=y_i^0.
% \]
% 由数列极限的性质知
% \[
%   \lim_{k\to\infty}x_i^ky_i^k=x_i^0y_i^0.
% \]
% 于是
% \[
%   \lim_{k\to\infty}x_ky_k
%   =\lim_{k\to\infty}\sum_{i=1}^n x_i^ky_i^k
%   =\sum_{i=1}^n x_i^0y_i^0=x_0y_0.
% \qedhere
% \]
% \end{proof}
% 
% 
% 性质~\ref{property:ma-13-1-4} 的 (2) 是实数乘积极限性质的推广, 但当 $n\geqslant2$ 时, 这个性质与实数乘积极限性质有明显的差别。在 $\mathbb R^n$ 中, 两个非零向量的内积可能等于零, 例如两个非零向量正交时.
% 
% \begin{example}\label{example:ma-13-1-2}
% 设
% \[
%   x_k=\left(\frac{\sqrt k\sin k}{k},\,k^4e^{-k^2},\,\cos k\ln\left(1+\frac1k\right),\,k\tan\frac1k\right)\in\mathbb R^4
% \]
% $(k=1,2,\cdots)$, 求 $\lim\limits_{k\to\infty}x_k$.
% \end{example}
% 
% 
% \begin{solve}
% 由于
% \[
%   \lim_{k\to\infty}\frac{\sqrt k\sin k}{k}=0,\qquad
%   \lim_{k\to\infty}k^4e^{-k^2}=0,
% \]
% \[
%   \lim_{k\to\infty}\cos k\ln\left(1+\frac1k\right)=0,\qquad
%   \lim_{k\to\infty}k\tan\frac1k=1,
% \]
% 故
% \[
%   \lim_{k\to\infty}x_k=(0,0,0,1).
% \qedhere
% \]
% \end{solve}
% 
% 
% \subsection{聚点}
% \label{subsec:ma-13-1-3}
% 
% 聚点是极限理论中重要的概念。我们曾在 $\mathbb R$ 中讨论过它(见第一册)。对于 $\mathbb R^n$ 的一般情形, 聚点的定义如下:
% 
% \begin{definition}\label{definition:ma-13-1-4}
% 设 $E\subset\mathbb R^n$。若 $x\in\mathbb R^n$ 的每个 $\delta$ 邻域 $U(x,\delta)(\delta>0)$ 中都含有异于 $x$ 的 $E$ 中的点, 则称 $x$ 是 $E$ 的聚点或极限点.
% \end{definition}
% 
% 在上述定义中, 我们也可以用方形邻域来代替球形邻域。容易看出, 若 $x$ 是 $E$ 的聚点, 则 $x$ 的每个邻域中都含有无穷多个 $E$ 中的点。另外, $x$ 可以属于 $E$, 也可以不属于 $E$.
% 
% 若 $x$ 不是 $E$ 的聚点, 则存在 $\delta_0>0$, 使得 $U_0(x,\delta_0)$ 中不含有 $E$ 中的点。如果 $x\in E$ 且 $x$ 不是 $E$ 的聚点, 则称 $x$ 是 $E$ 的孤立点。这时存在 $\delta_0>0$, 使得
% \[
%   U(x,\delta_0)\cap E=\{x\}.
% \]
% 
% 利用点列的极限, 我们可以描述一个集合的聚点.
% 
% \begin{theorem}\label{theorem:ma-13-1-5}
% 设 $\varnothing\ne E\subset\mathbb R^n$, 则 $x$ 是 $E$ 的聚点的充分必要条件是: 存在 $E$ 中的互异点列 $\{x_k\}$, 使得
% \[
%   \lim_{k\to\infty}x_k=x.
% \]
% \end{theorem}
% 
% 定理~\ref{theorem:ma-13-1-5} 的证明与 $\mathbb R$ 中的情形类似, 请读者自己给出.
% 
% \begin{example}\label{example:ma-13-1-3}
% 设
% \[
%   E=\{(x_1,x_2,\cdots,x_n)\in\mathbb R^n:\text{对于 }\forall i(1\leqslant i\leqslant n),\ x_i\text{ 是有理数}\}.
% \]
% 求 $E$ 的聚点集.
% \end{example}
% 
% 
% \begin{solve}
% 设 $x_0=(x_1^0,x_2^0,\cdots,x_n^0)\in\mathbb R^n$。对于 $\forall\delta>0$, 考虑 $x_0$ 的方形邻域 $N(x_0,\delta)$。对于每一个 $i$, 开区间 $(x_i^0-\delta,x_i^0+\delta)$ 中都含有有理数 $x_i'$ 且 $x_i'\ne x_i^0$。于是
% \[
%   (x_1',x_2',\cdots,x_n')\in N(x_0,\delta)\cap E,
% \]
% 且 $(x_1',x_2',\cdots,x_n')\ne x_0$。这说明 $x_0$ 是 $E$ 的聚点。由于 $x_0$ 是 $\mathbb R^n$ 中任意一点, 故 $E$ 的聚点集是 $\mathbb R^n$.
% 
% 如果集合 $A$ 中的每个点的任意邻域中都含有集合 $B$ 中的点, 则称 $B$ 在 $A$ 中稠密。上面的例子表明, $E$ 在 $\mathbb R^n$ 中稠密。请读者利用例~\ref{example:ma-13-1-3} 的结论, 将 $E$ 中的点排成一个点列.
% \end{solve}
% 
% 
% \subsection{开集与闭集}
% \label{subsec:ma-13-1-4}
% 
% 设 $E\subset\mathbb R^n$, 记 $E$ 的余集为 $E^c=\mathbb R^n\setminus E$。下面我们利用邻域来对 $\mathbb R^n$ 中的点进行分类.
% 
% \begin{definition}\label{definition:ma-13-1-5}
% 设 $E\subset\mathbb R^n,\ x\in\mathbb R^n$.
% 
% (1) 若存在 $\delta>0$, 使得 $U(x,\delta)\subset E$, 则称 $x$ 是 $E$ 的内点。$E$ 的全体内点组成的集合记为 $E^\circ$, 称为 $E$ 的内部.
% 
% (2) 若存在 $\delta>0$, 使得 $U(x,\delta)\cap E=\varnothing$, 则称 $x$ 是 $E$ 的外点。$E$ 的全体外点组成的集合称为 $E$ 的外部.
% 
% (3) 若对于 $\forall\delta>0$, 都有 $U(x,\delta)\cap E\ne\varnothing$ 且 $U(x,\delta)\cap E^c\ne\varnothing$, 则称 $x$ 是 $E$ 的边界点。$E$ 的全体边界点组成的集合记为 $\partial E$, 称为 $E$ 的边界.
% \end{definition}
% 
% 由定义可知, $E$ 的内点一定属于 $E$; $E$ 的外点一定不属于 $E$, 且 $E$ 的外部就是 $(E^c)^\circ$; $E$ 的边界点可能属于 $E$, 也可能不属于 $E$。一个点 $x$ 是 $E$ 的边界点的充分必要条件是: $x$ 既不是 $E$ 的内点, 也不是 $E$ 的外点.
% 
% \begin{example}\label{example:ma-13-1-4}
% 设
% \[
%   E=\{(x,y):x^2+y^2<1,\ y\geqslant x\},
% \]
% 求 $E^\circ,(E^c)^\circ$ 和 $\partial E$.
% \end{example}
% 
% 
% \begin{solve}
% 参见图 13.1.1。容易看出
% \[
%   E^\circ=\{(x,y):x^2+y^2<1,\ y>x\}=E\setminus\{(x,y):y=x\},
% \]
% \[
%   (E^c)^\circ=\mathbb R^2\setminus\{(x,y):x^2+y^2\leqslant1,\ y\geqslant x\},
% \]
% \[
%   \partial E=\{(x,y):x^2+y^2=1,\ y\geqslant x\}\cup\{(x,y):x^2+y^2<1,\ y=x\}.
% \]
% 
% 在一元微积分中, 开区间与闭区间起着重要的作用。在 $\mathbb R^n$ 中与它们密切相关的概念是开集与闭集。开集与闭集是 $\mathbb R^n$ 中两类重要的子集, 它们在多元函数的极限与连续理论中有重要作用.
% \end{solve}
% 
% 
% \begin{definition}\label{definition:ma-13-1-6}
% 设 $E\subset\mathbb R^n$。若 $E=E^\circ$, 则称 $E$ 是开集。规定空集 $\varnothing$ 是开集.
% \end{definition}
% 
% 显然, 每个球形邻域和方形邻域都是开集。在 $\mathbb R$ 中, 开集可以表示为至多可数个互不相交的开区间的并。在 $\mathbb R^2$ 中, 开集可以看成是不含边界的平面点集。例如
% \[
%   E=\{(x,y): |(x,y)-(0,k)|<1/2,\ k\in\mathbb Z\}
% \]
% 是一个开集, 它由无穷多个两两不相交的开圆盘组成。即使一个开集只有一块, 它的边界也可能很复杂。例如
% \[
%   E=\{(x,y):0\leqslant y\leqslant 1/2,\ x=1/k,\ k=1,2,\cdots\},
% \]
% 则
% \[
%   D=((0,1)\times(0,1))\setminus E
% \]
% 的边界为 $E\cup\partial([0,1]\times[0,1])$, 且 $D$ 是开集.
% 
% \begin{theorem}\label{theorem:ma-13-1-6}
% $\mathbb R^n$ 中的开集具有以下性质:
% 
% (1) $\mathbb R^n$ 和 $\varnothing$ 都是开集;
% 
% (2) 任意多个开集的并是开集;
% 
% (3) 有限多个开集的交是开集.
% \end{theorem}
% 
% 定理~\ref{theorem:ma-13-1-6} 的证明很容易, 请读者自己完成.
% 
% 需要注意的是, 任意多个开集的交不一定是开集。事实上, 对于 $x\in\mathbb R^n$, 有
% \[
%   \{x\}=\bigcap_{k=1}^{+\infty}U\left(x,\frac1k\right).
% \]
% 在 $\mathbb R$ 中, 闭区间的余集是开集。由此我们可以引入闭集的概念.
% 
% \begin{definition}\label{definition:ma-13-1-7}
% 设 $E\subset\mathbb R^n$。若 $E^c$ 是开集, 则称 $E$ 是闭集.
% \end{definition}
% 
% 在 $\mathbb R$ 中, 闭集不一定能表示为至多可数个闭区间的并。由定义可知, 对于任意开集 $E$, 集合 $F=E\cup\partial E$ 是闭集.
% 
% \begin{theorem}\label{theorem:ma-13-1-7}
% $\mathbb R^n$ 中的闭集具有以下性质:
% 
% (1) $\mathbb R^n$ 和 $\varnothing$ 都是闭集;
% 
% (2) 有限多个闭集的并是闭集;
% 
% (3) 任意多个闭集的交是闭集.
% \end{theorem}
% 
% 定理~\ref{theorem:ma-13-1-7} 可以直接由闭集的定义和定理~\ref{theorem:ma-13-1-6} 得到。事实上, 对任意一族集合 $\{E_\lambda\}_{\lambda\in\Lambda}$, 我们有德摩根公式
% \[
%   \left(\bigcup_{\lambda\in\Lambda}E_\lambda\right)^c
%   =\bigcap_{\lambda\in\Lambda}E_\lambda^c,\qquad
%   \left(\bigcap_{\lambda\in\Lambda}E_\lambda\right)^c
%   =\bigcup_{\lambda\in\Lambda}E_\lambda^c.
% \]
% 
% 设 $E\subset\mathbb R^n$, 记 $E$ 的全体聚点组成的集合为 $E'$, 称为 $E$ 的导集。集合
% \[
%   \overline E=E\cup E'
% \]
% 称为 $E$ 的闭包.
% 
% \begin{theorem}\label{theorem:ma-13-1-8}
% 设 $E\subset\mathbb R^n$, 则 $E$ 是闭集的充分必要条件是 $E=\overline E$.
% \end{theorem}
% 
% 
% \begin{proof}
% 由闭包的定义可知, 只需证明 $E$ 是闭集的充分必要条件是 $E'\subset E$.
% 
% 充分性。若 $E$ 的所有聚点都属于 $E$, 则 $E^c$ 中没有 $E$ 的聚点。于是对于 $\forall x\in E^c$, 存在 $\delta_x>0$, 使得
% \[
%   U(x,\delta_x)\cap E=\varnothing.
% \]
% 这说明 $U(x,\delta_x)\subset E^c$, 因而 $E^c$ 是开集, 即 $E$ 是闭集.
% 
% 必要性。若 $E$ 是闭集, 则 $E^c$ 是开集。对于 $\forall x\in E^c$, 存在 $\delta_x>0$, 使得 $U(x,\delta_x)\subset E^c$, 因而 $x$ 不是 $E$ 的聚点。这说明 $E^c$ 中没有 $E$ 的聚点, 即 $E'\subset E$, 从而 $E=E\cup E'=\overline E$.
% \end{proof}
% 
% 
% \subsection{\texorpdfstring{欧氏空间 $\mathbb R^n$ 中的基本定理}{欧氏空间 Rⁿ 中的基本定理}}
% \label{subsec:ma-13-1-5}
% 
% 本小节我们将实直线中的一些基本定理推广到欧氏空间 $\mathbb R^n$ 中, 包括完备性定理、闭集套定理、聚点原理和有限覆盖定理.
% 
% \paragraph*{完备性}
% 
% 
% 设 $\{x_k\}$ 是 $\mathbb R^n$ 中的点列。若对于 $\forall\varepsilon>0,\exists K\in\mathbb N$, 当 $k',k''>K$ 时, 有
% \[
%   |x_{k'}-x_{k''}|<\varepsilon,
% \]
% 则称 $\{x_k\}$ 是柯西点列。若空间 $X\subset\mathbb R^n$ 中的每个柯西点列都收敛, 且极限属于 $X$, 则称 $X$ 是完备的.
% 
% \begin{theorem}\label{theorem:ma-13-1-9}
% 欧氏空间 $\mathbb R^n$ 是完备的.
% \end{theorem}
% 
% 
% \begin{proof}
% 设 $x_k=(x_1^k,x_2^k,\cdots,x_n^k)\in\mathbb R^n(k=1,2,\cdots)$ 是一个柯西点列。则对于 $\forall i=1,2,\cdots,n$, 数列 $\{x_i^k\}$ 是 $\mathbb R$ 中的柯西数列。由 $\mathbb R$ 的完备性知, 每个数列 $\{x_i^k\}$ 都收敛。由定理~\ref{theorem:ma-13-1-1} 可知 $\{x_k\}$ 收敛.
% \end{proof}
% 
% 
% 完备性是欧氏空间 $\mathbb R^n$ 的一个重要性质, 在数学分析中有广泛应用.
% 
% \paragraph*{闭集套定理}
% 
% 
% 设 $E$ 是 $\mathbb R^n$ 中的有界集, 定义
% \[
%   \operatorname{diam}(E)=\sup_{x,y\in E}\{|x-y|\},
% \]
% 称为 $E$ 的直径.
% 
% \begin{theorem}\label{theorem:ma-13-1-10}
% 设 $F_k\subset\mathbb R^n(k=1,2,\cdots)$ 是一列非空闭集, 且满足
% 
% (1) 对于 $\forall k\in\mathbb N$, 有 $F_k\supset F_{k+1}$;
% 
% (2) $\lim\limits_{k\to\infty}\operatorname{diam}(F_k)=0$.
% 
% 则存在唯一的 $x_0\in\mathbb R^n$, 使得 $x_0\in F_k(k=1,2,\cdots)$, 即
% \[
%   \{x_0\}=\bigcap_{k=1}^{+\infty}F_k.
% \]
% \end{theorem}
% 
% 
% \begin{proof}
% 对于每个 $k$, 任取 $x_k\in F_k$。下面证明 $\{x_k\}$ 是柯西点列。对于 $\forall\varepsilon>0$, 存在 $K\in\mathbb N$, 当 $k>K$ 时, 有 $\operatorname{diam}(F_k)<\varepsilon$。若 $k''>k'>K$, 则 $x_{k'}\in F_{k'}$, $x_{k''}\in F_{k''}\subset F_{k'}$, 因而
% \[
%   |x_{k'}-x_{k''}|\leqslant \operatorname{diam}(F_{k'})<\varepsilon.
% \]
% 这说明 $\{x_k\}$ 是柯西点列。由定理~\ref{theorem:ma-13-1-9} 知 $\{x_k\}$ 收敛, 记 $x_0=\lim\limits_{k\to\infty}x_k$。对于任意给定的 $k_0$, 当 $k>k_0$ 时, 有 $x_k\in F_{k_0}$。由于 $F_{k_0}$ 是闭集, 故 $x_0\in F_{k_0}$。由 $k_0$ 的任意性知 $x_0\in\bigcap\limits_{k=1}^{+\infty}F_k$.
% 
% 若另有 $x'\in\bigcap\limits_{k=1}^{+\infty}F_k$, 则
% \[
%   |x'-x_0|\leqslant \lim_{k\to\infty}\operatorname{diam}(F_k)=0.
% \]
% 于是 $x'=x_0$.
% \end{proof}
% 
% 
% 实直线上的闭区间套定理是定理~\ref{theorem:ma-13-1-10} 的特殊情形。定理~\ref{theorem:ma-13-1-10} 中的闭集不能换成开集; 如果 $F_k$ 不是闭集, 结论可能不成立.
% 
% \paragraph*{聚点原理}
% 
% 
% 在 $\mathbb R^n$ 中, 我们同样有以下的聚点原理.
% 
% \begin{theorem}\label{theorem:ma-13-1-11}
% $\mathbb R^n$ 中的有界点列必有收敛子列.
% \end{theorem}
% 
% 
% \begin{proof}
% 设 $\{x_k\}$ 是 $\mathbb R^n$ 中的有界点列, 其中 $x_k=(x_1^k,x_2^k,\cdots,x_n^k)$。则对于每个 $i=1,2,\cdots,n$, 数列 $\{x_i^k\}$ 都有界。由数列的聚点原理知, $\{x_1^k\}$ 有收敛子列 $\{x_1^{k_j'}\}$。再从 $\{x_2^{k_j'}\}$ 中取收敛子列 $\{x_2^{k_j''}\}$。依此进行, 经过 $n$ 次后得到 $\{x_k\}$ 的一个子列 $\{x_{k_j}\}$, 使得它的每个分量数列都收敛。因而由定理~\ref{theorem:ma-13-1-1} 知 $\{x_{k_j}\}$ 收敛.
% \end{proof}
% 
% 
% \begin{theorem}\label{theorem:ma-13-1-12}
% $\mathbb R^n$ 中的任一有界无限集 $E$ 至少有一个聚点.
% \end{theorem}
% 
% 
% \begin{proof}
% 由于 $E$ 是无限集, 可从 $E$ 中取出互异点列 $\{x_k\}$。又由于 $E$ 有界, 故 $\{x_k\}$ 是有界点列。由定理~\ref{theorem:ma-13-1-11} 知 $\{x_k\}$ 有收敛子列。设该收敛子列的极限为 $x_0$。由于 $\{x_k\}$ 中的点互异, 故 $x_0$ 是集合 $\{x_k\}$ 的聚点, 从而也是 $E$ 的聚点.
% \end{proof}
% 
% 
% \paragraph*{紧集}
% 
% 
% 在实直线上, 有限覆盖定理是一个重要定理。在 $\mathbb R^n$ 中, 与有限覆盖定理密切相关的概念是紧集.
% 
% 设 $E\subset\mathbb R^n$, $\{O_\lambda\}_{\lambda\in\Lambda}$ 是 $\mathbb R^n$ 中的一族开集。若
% \[
%   E\subset \bigcup_{\lambda\in\Lambda}O_\lambda,
% \]
% 则称 $\{O_\lambda\}_{\lambda\in\Lambda}$ 是 $E$ 的一个开覆盖。若 $\Lambda$ 是有限集, 则称它为有限开覆盖.
% 
% \begin{definition}\label{definition:ma-13-1-8}
% 设 $E\subset\mathbb R^n$。若 $E$ 的每一个开覆盖都存在有限子覆盖, 即对于 $E$ 的任意开覆盖 $\{O_\lambda\}_{\lambda\in\Lambda}$, 都存在其中有限个开集 $O_1,O_2,\cdots,O_K(K<\infty)$, 使得
% \[
%   E\subset\bigcup_{k=1}^K O_k,
% \]
% 则称 $E$ 是紧集.
% \end{definition}
% 
% 紧集的定义强调的是: $E$ 的每一个开覆盖都存在有限子覆盖, 而不是只要找到某一个有限开覆盖即可.
% 
% 在 $\mathbb R$ 中, 有限个有界闭区间的并是紧集。若 $E$ 是 $\mathbb R^n$ 中的无界集, 则 $E$ 不是紧集。事实上, $\{U(0,k)\}_{k\in\mathbb N}$ 是 $E$ 的一个开覆盖, 但它不存在有限子覆盖。例如 $E=\{1/k:k\in\mathbb N\}$ 不是紧集; 但 $E\cup E'=\{1/k:k\in\mathbb N\}\cup\{0\}$ 是紧集.
% 
% \begin{theorem}\label{theorem:ma-13-1-13}
% 设 $E\subset\mathbb R^n$, 则 $E$ 是紧集的充分必要条件是 $E$ 是有界闭集.
% \end{theorem}
% 
% 
% \begin{proof}
% 必要性。对于任意集合 $E\subset\mathbb R^n$, $\{U(0,k)\}_{k\in\mathbb N}$ 是 $E$ 的一个开覆盖。由于 $E$ 是紧集, 故它存在有限子覆盖 $U(0,k_1),U(0,k_2),\cdots,U(0,k_J)$, 不妨设 $0<k_1<k_2<\cdots<k_J<\infty$。于是
% \[
%   E\subset\bigcup_{j=1}^J U(0,k_j)=U(0,k_J),
% \]
% 故 $E$ 有界.
% 
% 下面证明 $E$ 是闭集。反设 $E$ 不是闭集, 则存在 $E$ 的聚点 $x_0\notin E$。对于 $\forall x\in E$, 有 $x\ne x_0$, 令
% \[
%   r_x=\frac12|x-x_0|>0,
% \]
% 则
% \[
%   U(x,r_x)\cap U(x_0,r_x)=\varnothing.
% \]
% 由于 $\bigcup\limits_{x\in E}U(x,r_x)$ 是 $E$ 的一个开覆盖, 故存在有限子覆盖
% \[
%   U(x_1,r_{x_1}),U(x_2,r_{x_2}),\cdots,U(x_K,r_{x_K}).
% \]
% 令
% \[
%   r_K=\min_{1\leqslant k\leqslant K}r_{x_k}>0,
% \]
% 则
% \[
%   U(x_0,r_K)\cap\left\{\bigcup_{k=1}^K U(x_k,r_{x_k})\right\}=\varnothing.
% \]
% 于是 $U(x_0,r_K)\cap E=\varnothing$, 这与 $x_0$ 是 $E$ 的聚点矛盾。因此 $E$ 是闭集.
% 
% 充分性。设 $E$ 是有界闭集但不是紧集, 则存在 $E$ 的一个开覆盖 $\{O_\lambda\}_{\lambda\in\Lambda}$, 使得它不存在有限子覆盖。为简单起见, 我们只证明 $E\subset\mathbb R^2$ 的情形, $\mathbb R^n$ 中的情形类似。由于 $E$ 有界, 存在 $a>0$, 使得 $E\subset[-a,a]\times[-a,a]=I_0$。将 $I_0$ 等分为四个闭正方形, 则其中至少有一个闭正方形与 $E$ 的交不能被 $\{O_\lambda\}_{\lambda\in\Lambda}$ 中的有限个开集覆盖, 记这个闭正方形为 $I_1$。再将 $I_1$ 等分为四个闭正方形, 继续上述过程, 得到一列闭正方形 $I_k(k=1,2,\cdots)$, 满足:
% 
% (1) $I_k\supset I_{k+1}$;
% 
% (2) $\lim\limits_{k\to\infty}\operatorname{diam}(I_k)=0$;
% 
% (3) 对于每个 $k$, $I_k\cap E$ 都不能被 $\{O_\lambda\}_{\lambda\in\Lambda}$ 中的有限个开集覆盖.
% 
% 令 $E_k=I_k\cap E$。则 $\{E_k\}$ 是满足定理~\ref{theorem:ma-13-1-10} 条件的一列非空闭集。于是存在唯一的点 $x_0\in\bigcap\limits_{k=1}^{+\infty}E_k$。因为 $x_0\in E$, 所以存在 $O_{\lambda_0}$, 使得 $x_0\in O_{\lambda_0}$。又由于 $O_{\lambda_0}$ 是开集, 当 $k$ 充分大时有 $E_k\subset O_{\lambda_0}$, 这与 (3) 矛盾.
% \end{proof}
% 
% 
% \section{多元函数与向量函数的极限}
% \label{sec:ma-13-2}
% 
% \subsection{多元函数的概念}
% \label{subsec:ma-13-2-1}
% 
% 在前面的学习中, 我们已经接触到许多一个量依赖于多个变量的情形。例如, 空间中某点的温度不仅依赖于该点的位置, 还依赖于时间; 火箭在某一时刻达到的高度依赖于推力、重力、燃料质量和火箭自身质量等多个因素。有了前面的准备, 现在我们可以很自然地引入 $n$ 元函数的概念.
% 
% \begin{definition}\label{definition:ma-13-2-1}
% 设 $E\subset\mathbb R^n$。如果对应法则 $f$ 对于 $E$ 中的每一个点 $x=(x_1,x_2,\cdots,x_n)$, 都有唯一的实数 $u$ 与之对应, 则称 $f$ 是定义在 $E$ 上的 $n$ 元函数, 记为
% \[
%   f:E\to\mathbb R,\qquad x\mapsto u=f(x)=f(x_1,x_2,\cdots,x_n).
% \]
% 称 $u$ 为函数 $f$ 在点 $x$ 处的函数值, 称 $E$ 为 $f$ 的定义域, 称集合 $\{f(x):x\in E\}\subset\mathbb R$ 为 $f$ 的值域.
% \end{definition}
% 
% 当 $n\geqslant2$ 时, $n$ 元函数也称为多元函数。一般地, 我们用 $u=f(x)$ 表示多元函数; 必要时也用 $g(x),h(x)$ 等表示。当 $n=1$ 时, 这与一元函数的记号是一致的, 常写作 $y=f(x)$。当 $n=2$ 时, 常写作 $z=f(x,y)$; 当 $n=3$ 时, 常写作 $u=f(x,y,z)$.
% 
% 对于二元函数 $z=f(x,y)$, 如果它具有较好的性质, 那么它的图形是空间中的一个曲面
% \[
%   \{(x,y,f(x,y)):(x,y)\in E\}.
% \]
% 参见图 13.2.1.
% 
% \begin{example}\label{example:ma-13-2-1}
% 单位球面在 $\mathbb R^3$ 中的方程为
% \[
%   x^2+y^2+z^2=1.
% \]
% 于是
% \[
%   z=\sqrt{1-x^2-y^2}
% \]
% 是定义在闭单位圆盘
% \[
%   \overline\Delta=\{(x,y):x^2+y^2\leqslant1\}
% \]
% 上的二元函数, 它的图形是单位球面的上半球面。类似地,
% \[
%   z=-\sqrt{1-x^2-y^2}
% \]
% 的图形是单位球面的下半球面.
% \end{example}
% 
% 设 $u=f(x)(x\in E\subset\mathbb R^n)$。若 $f(E)$ 是 $\mathbb R$ 中的有界集, 则称 $f$ 在 $E$ 上有界。类似地, 可以定义函数在 $E$ 中的上界、下界以及上、下确界等。请读者对一个二元函数来描述上述概念的几何意义.
% 
% \subsection{多元函数的极限}
% \label{subsec:ma-13-2-2}
% 
% \begin{definition}\label{definition:ma-13-2-2}
% 设函数 $f(x)$ 定义在 $E\subset\mathbb R^n$ 上, $x_0=(x_1^0,x_2^0,\cdots,x_n^0)\in E'$。若存在 $A\in\mathbb R$, 使得对于 $\forall\varepsilon>0$, 存在 $\delta>0$, 当 $x=(x_1,x_2,\cdots,x_n)\in U_0(x_0,\delta)\cap E$, 即 $0<|x-x_0|<\delta$ 且 $x\in E$ 时, 有
% \[
%   |f(x)-A|<\varepsilon,
% \]
% 则称 $A$ 为 $x$ 在 $E$ 中趋于 $x_0$ 时函数 $f(x)$ 的极限, 记做
% \[
%   \lim_{E\ni x\to x_0}f(x)=A,
% \]
% 或
% \[
%   \lim_{E\ni(x_1,x_2,\cdots,x_n)\to(x_1^0,x_2^0,\cdots,x_n^0)}
%   f(x_1,x_2,\cdots,x_n)=A.
% \]
% 如果存在 $x_0$ 的某个去心邻域 $U_0(x_0,\delta_0)\subset E$, 则在极限记号中可以省略 $x\in E$.
% \end{definition}
% 
% 上述定义对于一元函数同样适用。例如, 若 $y=f(x)$ 定义在 $[a,b]$ 上, 则它实际上给出了自变量 $x$ 趋于 $x_0\in(a,b)$ 时的极限以及趋于 $a$ 或 $b$ 时的单侧极限的统一定义。不仅如此, 上述极限也不要求函数在一个区间上处处有定义, 它的定义域可以是很广泛的集合。对于多元函数, 上述定义则更为广泛, 因为这时函数的定义域有可能更加复杂.
% 
% 多元函数极限同一元函数极限一样, 具有唯一性、局部有界性、四则运算法则等性质, 也可以类似地定义无穷小量和无穷大量。这里不再一一列出.
% 
% \begin{theorem}\label{theorem:ma-13-2-1}
% 设函数 $f(x)$ 定义在 $E\subset\mathbb R^n$ 上, $x_0\in E'$, 则
% \[
%   \lim_{E\ni x\to x_0}f(x)=A
% \]
% ($A$ 可以为 $\infty$) 的充分必要条件是: 对于 $E\setminus\{x_0\}$ 中任意收敛于 $x_0$ 的点列 $\{x_k\}$, 都有
% \[
%   \lim_{k\to\infty}f(x_k)=A.
% \]
% \end{theorem}
% 
% 由函数极限的定义知, 当一个函数 $f(x)$ 在 $x_0\in E'$ 存在极限时, 由于 $\mathbb R^n$ 中点 $x$ 趋向 $x_0$ 的方式非常复杂, 因此该极限过程蕴含着许多信息。反之, 若 $x$ 以某种特殊的方式(如沿曲线、点列等)趋于 $x_0$ 时, $f(x)$ 的极限不存在, 则极限 $\lim\limits_{E\ni x\to x_0}f(x)$ 必不存在.
% 
% \begin{example}\label{example:ma-13-2-2}
% 求极限
% \[
%   \lim_{(x,y)\to(0,0)}\frac{\tan(x^3+y^3)}{x^2+y^2}.
% \]
% \end{example}
% 
% 
% \begin{solve}
% 由于
% \[
%   \tan(x^3+y^3)\sim x^3+y^3,\qquad (x,y)\to(0,0),
% \]
% 而
% \[
%   |x^3+y^3|=|x+y||x^2+y^2-xy|
%   \leqslant \frac32 |x+y|(x^2+y^2),
% \]
% 故当 $(x,y)\ne(0,0)$ 时, 有
% \[
%   \left|\frac{x^3+y^3}{x^2+y^2}\right|\leqslant \frac32|x+y|.
% \]
% 由于
% \[
%   \lim_{(x,y)\to(0,0)}\frac32|x+y|=0,
% \]
% 从而有
% \[
%   \lim_{(x,y)\to(0,0)}
%   \left|\frac{\tan(x^3+y^3)}{x^2+y^2}\right|
%   =\lim_{(x,y)\to(0,0)}
%   \left|\frac{x^3+y^3}{x^2+y^2}\right|=0.
% \]
% 从上式我们推出
% \[
%   \lim_{(x,y)\to(0,0)}\frac{\tan(x^3+y^3)}{x^2+y^2}=0.
% \qedhere
% \]
% \end{solve}
% 
% 
% \begin{example}\label{example:ma-13-2-3}
% 设函数 $f(x,y)=\dfrac{xy}{x^2+y^2}$, 试问: 当 $(x,y)\to(0,0)$ 时, $f(x,y)$ 是否存在极限?
% \end{example}
% 
% 
% \begin{solve}[{解法 1}]
% 取点列 $\{(0,\frac1k)\}$, 则有
% \[
%   \lim_{k\to\infty}f\left(0,\frac1k\right)=\lim_{k\to\infty}0=0.
% \]
% 再取点列 $\{(\frac1k,\frac1k)\}$, 我们有
% \[
%   \lim_{k\to\infty}f\left(\frac1k,\frac1k\right)
%   =\lim_{k\to\infty}\frac{\frac1{k^2}}{2\frac1{k^2}}=\frac12.
% \]
% 因此当 $(x,y)$ 趋于 $(0,0)$ 时, $f(x,y)$ 不存在极限.
% \end{solve}
% 
% 
% \begin{solve}[{解法 2}]
% 取 $y=kx(k\in\mathbb R)$, 则当 $x\to0$ 时, 有 $y\to0$。由于
% \[
%   \lim_{x\to0}f(x,kx)
%   =\lim_{x\to0}\frac{kx^2}{(1+k^2)x^2}
%   =\frac{k}{1+k^2},
% \]
% 这说明自变量 $(x,y)$ 沿射线 $y=kx$ 趋于 $(0,0)$ 时, $f(x,y)$ 的极限依赖于射线的选取, 从而极限 $\lim\limits_{(x,y)\to(0,0)}f(x,y)$ 必不存在.
% \end{solve}
% 
% 
% \begin{solve}[{解法 3}]
% 在 $Oxy$ 平面上取极坐标
% \[
%   \begin{cases}
%     x=r\cos\theta,\\
%     y=r\sin\theta,
%   \end{cases}
% \]
% 其中 $r\in(0,+\infty),\theta\in[0,2\pi)$, 则 $f(x,y)$ 在极坐标下为
% \[
%   f(x,y)=\frac12\sin2\theta.
% \]
% 由于 $(x,y)\to(0,0)$ 时必有 $r\to0+0$, 而 $\frac12\sin2\theta$ 与 $r$ 无关但又不是常数函数, 从而极限 $\lim\limits_{(x,y)\to(0,0)}f(x,y)$ 也不存在.
% \end{solve}
% 
% 
% \begin{example}\label{example:ma-13-2-4}
% 设函数
% \[
%   f(x,y)=\frac{x^2y}{x^4+y^2},
% \]
% 讨论 $\lim\limits_{(x,y)\to(0,0)}f(x,y)$ 是否存在.
% \end{example}
% 
% 
% \begin{solve}
% 如果取 $y=kx(k\in\mathbb R)$, 则当 $(x,y)$ 沿该射线趋于 $(0,0)$ 时, 有
% \[
%   \lim_{x\to0}f(x,kx)=\lim_{x\to0}\frac{kx^3}{(x^2+k^2)x^2}=0;
% \]
% 如果取 $y=x^2$, 则当 $(x,y)$ 沿该抛物线趋于 $(0,0)$ 时, 有
% \[
%   \lim_{x\to0}f(x,x^2)=\lim_{x\to0}\frac{x^4}{x^4+x^4}=\frac12.
% \]
% 这说明 $\lim\limits_{(x,y)\to(0,0)}f(x,y)$ 不存在.
% 
% 此例说明, 一个函数沿任何射线的极限都存在且相等也不足以保证该函数的极限存在.
% \end{solve}
% 
% 
% \begin{example}\label{example:ma-13-2-5}
% 设函数
% \[
%   f(x)=f(x,y,z)=[\sin1-\sin(x^2+y^2+z^2)]\ln[1-(x^2+y^2+z^2)].
% \]
% 记 $E=\{x=(x,y,z):|x|<1\}$, 并设 $x_0\in\partial E$, 求极限 $\lim\limits_{E\ni x\to x_0}f(x)$.
% \end{example}
% 
% 
% \begin{solve}
% 显然 $f$ 的定义域为 $E$。当 $E\ni x\to x_0$ 时, $t=x^2+y^2+z^2\to1-0$, 因此
% \[
%   \lim_{E\ni x\to x_0}f(x)
%   =\lim_{t\to1-0}(\sin1-\sin t)\ln(1-t)=0.
% \]
% 最后一步的极限可由求一元函数极限的方法求得(如用洛必达法则).
% \end{solve}
% 
% 
% \subsection{累次极限}
% \label{subsec:ma-13-2-3}
% 
% 设 $f(x)=f(x_1,x_2,\cdots,x_n)$ 在 $x_0\in\mathbb R^n$ 的邻域内有定义, 现在继续考虑 $x$ 趋于 $x_0$ 时极限行为的方式。当 $x$ 趋于 $x_0$ 时具有许多种方式, 且 $x$ 沿着某些特殊方式趋于 $x_0$ 时具有特别的意义。在以后的章节中我们常常会遇到以下形式的累次极限。我们以一个二元函数的情形为例来讲述这种极限, 对 $n(n\geqslant3)$ 元函数的情形是类似的, 只是形式复杂而已.
% 
% \begin{definition}\label{definition:ma-13-2-3}
% 设函数 $z=f(x,y)$ 在 $E\subset\mathbb R^2$ 上有定义, 且邻域 $N_0((x_0,y_0),\delta_0)\subset E(\delta_0>0)$。若在 $N_0((x_0,y_0),\delta_0)$ 内, 对每个固定的 $y\ne y_0$, $\lim\limits_{x\to x_0}f(x,y)=\varphi(y)$ 存在, 并且 $\lim\limits_{y\to y_0}\varphi(y)=A(A$ 为常数), 则称 $A$ 为 $f(x,y)$ 当 $(x,y)$ 趋于 $(x_0,y_0)$ 时先 $x$ 后 $y$ 的累次极限, 记为
% \[
%   A=\lim_{y\to y_0}\lim_{x\to x_0}f(x,y).
% \]
% 类似地, 我们可以定义先 $y$ 后 $x$ 的累次极限 $\lim\limits_{x\to x_0}\lim\limits_{y\to y_0}f(x,y)$.
% \end{definition}
% 
% 从定义~\ref{definition:ma-13-2-3} 可以看出, 一个二元函数在 $(x_0,y_0)$ 处的累次极限存在与否与函数在直线 $x=x_0$ 与 $y=y_0$ 上的值无关。因此容易举出例子来说明一个函数的两个累次极限都存在, 但有可能极限却不存在。例如, 设函数
% \[
%   f(x,y)=
%   \begin{cases}
%     0, & xy=0,\\
%     1, & xy\ne0,
%   \end{cases}
% \]
% 则有
% \[
%   \lim_{x\to0}\lim_{y\to0}f(x,y)
%   =\lim_{y\to0}\lim_{x\to0}f(x,y)=1,
% \]
% 而 $\lim\limits_{(x,y)\to(0,0)}f(x,y)$ 却不存在.
% 
% 另外, 当函数的极限存在时, 人们很自然地会猜测累次极限必定存在且等于函数的极限。对于该猜测, 我们的答案是否定的。例如, 设
% \[
%   f(x,y)=
%   \begin{cases}
%     (x+y)\sin\dfrac1x\sin\dfrac1y, & xy\ne0,\\
%     0, & xy=0.
%   \end{cases}
% \]
% 由于对于 $\forall(x,y)\in\mathbb R^2$, 有 $|f(x,y)|\leqslant |x|+|y|$, 因此
% \[
%   \lim_{(x,y)\to(0,0)}f(x,y)=0.
% \]
% 注意到对每个 $x\ne0,\frac1{k\pi}(k$ 为整数), 极限 $\lim\limits_{y\to0}f(x,y)$ 不存在; 而对每个 $y\ne0,\frac1{k\pi}(k$ 为整数), 极限 $\lim\limits_{x\to0}f(x,y)$ 也不存在。因此无法讨论 $(x,y)$ 趋于 $(0,0)$ 时的累次极限, 所以 $f(x,y)$ 的两个累次极限都不存在.
% 
% 考查以上例子, 其累次极限不存在的本质是存在任意小的 $x\ne0$, 极限 $\lim\limits_{y\to0}f(x,y)$ 不存在。若假定后者处处存在, 则可以证明人们的猜测是正确的.
% 
% \begin{theorem}\label{theorem:ma-13-2-2}
% 设函数 $f(x,y)$ 在 $N_0((x_0,y_0),\delta_0)(\delta_0>0)$ 上有定义, 并且满足:
% 
% (1) $\exists A$(可以是 $\infty$), 使得
% \[
%   \lim_{(x,y)\to(x_0,y_0)}f(x,y)=A;
% \]
% 
% (2) 对于 $U_0(y_0,\delta_0)$ 内任意固定的 $y\ne y_0$, 极限 $\lim\limits_{x\to x_0}f(x,y)=\varphi(y)$ 存在,
% 
% 则有
% \[
%   \lim_{y\to y_0}\varphi(y)
%   =\lim_{y\to y_0}\lim_{x\to x_0}f(x,y)=A.
% \]
% \end{theorem}
% 
% 
% \begin{proof}
% 我们只证 $A\ne\infty$ 的情形.
% 
% 由 $\lim\limits_{(x,y)\to(x_0,y_0)}f(x,y)=A$ 可知, 对于 $\forall\varepsilon>0,\exists\delta(0<\delta<\delta_0)$, 当 $(x,y)\in N_0((x_0,y_0),\delta)$ 时, 有
% \begin{equation}
%   |f(x,y)-A|<\frac{\varepsilon}{2}.
%   \label{eq:ma-13-2-1}
% \end{equation}
% 对任意的 $y\ne y_0$, 当 $(x,y)\in N((x_0,y_0),\delta)$ 时, 由于 $\lim\limits_{x\to x_0}f(x,y)=\varphi(y)$ 存在, 因此可在式 \eqref{eq:ma-13-2-1} 中令 $x\to x_0$, 从而当 $0<|y-y_0|<\delta$ 时, 有
% \[
%   |\varphi(y)-A|\leqslant\frac{\varepsilon}{2}<\varepsilon,
% \]
% 即 $\lim\limits_{y\to y_0}\varphi(y)=A$.
% \end{proof}
% 
% 
% \subsection{向量函数的定义与极限}
% \label{subsec:ma-13-2-4}
% 
% 设 $A,B$ 为任意两个非空集合(不一定是数集), 若对于 $A$ 中每个元素 $a$, 根据对应法则 $f$, 在 $B$ 中有唯一一个元素 $b$ 与之对应, 则称 $f$ 是 $A$ 到 $B$ 的一个映射, $b$ 称为 $a$ 在映射 $f$ 下的像, 记为
% \[
%   f:A\to B,\qquad a\mapsto b=f(a).
% \]
% 类似于函数的有关定义, 我们可以给出一个映射的定义域、值域等概念。同样, 对于映射我们可以引进单射、满射以及一一对应等概念。利用映射的概念, 容易给出向量函数的定义.
% 
% \begin{definition}\label{definition:ma-13-2-4}
% 设 $E\subset\mathbb R^n$, 若 $f$ 是 $E$ 到 $\mathbb R^m(m\geqslant1)$ 的一个映射, 则称 $f$ 是 $E$ 上的一个 $m$ 维 $(n$ 元)向量函数(简称向量函数), 记为
% \[
%   u=f(x),\qquad x\in E.
% \]
% \end{definition}
% 
% 设 $E\subset\mathbb R^n$, 显然, 一个 $E\to\mathbb R$ 的向量函数即是前面我们讨论的 $n$ 元函数。对于上述定义的 $m$ 维向量函数, 记 $x=(x_1,x_2,\cdots,x_n)\in E$, $u=(u_1,u_2,\cdots,u_m)=f(x_1,x_2,\cdots,x_n)\in\mathbb R^m$, 则 $f(x)$ 的每个分量 $u_j$ 都是 $E$ 上的一个 $n$ 元函数, 即对于 $j=1,2,\cdots,m$, 有 $u_j=f_j(x)(x\in E)$。因此 $u=f(x):E\to\mathbb R^m$ 为一个向量函数的充分必要条件是: 存在 $m$ 个定义在 $E$ 上的 $n$ 元函数 $f_j(x)(j=1,2,\cdots,m)$, 使得
% \[
%   f(x)=(f_1(x),f_2(x),\cdots,f_m(x)).
% \]
% 自然地, 若 $f(E)$ 是 $\mathbb R^m$ 的有界集, 则称 $f(x)$ 在 $E$ 是有界的.
% 
% 例如, 平面 $\mathbb R^2$ 中子集
% \[
%   E=\{(r,\theta):0\leqslant r<+\infty,\ 0\leqslant\theta<2\pi\}
% \]
% 到 $\mathbb R^2$ 上的极坐标变换 $f(r,\theta)=(x,y)=(r\cos\theta,r\sin\theta)$ 是一个二维(二元)向量函数。作为映射, 容易看出它是 $E\setminus\{(0,\theta):0\leqslant\theta<2\pi\}$ 到 $\mathbb R^2\setminus\{(0,0)\}$ 的一一对应。在 $(0,0)$ 处它不是单射。事实上, 对于 $\forall\theta(0\leqslant\theta<2\pi)$, 我们有 $f(0,\theta)=(0,0)$.
% 
% \begin{definition}\label{definition:ma-13-2-5}
% 设 $E\subset\mathbb R^n,x_0\in E'$, $f(x)=(f_1(x),f_2(x),\cdots,f_m(x))$ 是定义在 $E$ 上的一个 $m$ 维向量函数。若对于 $\forall j(1\leqslant j\leqslant m)$, 有
% \[
%   \lim_{E\ni x\to x_0}f_j(x)=A_j,
% \]
% 其中 $A_j\in\mathbb R(j=1,2,\cdots,m)$, 则称 $f(x)$ 当 $x$(在 $E$ 中)趋于 $x_0$ 时的极限为 $(A_1,A_2,\cdots,A_m)$, 记为
% \[
%   \lim_{E\ni x\to x_0}f(x)=(A_1,A_2,\cdots,A_m).
% \]
% 若 $\exists j_0(1\leqslant j_0\leqslant m)$, $\lim\limits_{E\ni x\to x_0}f_{j_0}(x)$ 不存在, 则称 $f(x)$ 当 $x$(在 $E$ 中)趋于 $x_0$ 时极限不存在.
% \end{definition}
% 
% 
% \begin{example}\label{example:ma-13-2-6}
% 设向量函数 $f(r,\theta)=(r\cos\theta,r\sin\theta)$, 其中 $r\in(0,+\infty),\theta\in[0,2\pi)$, 再设 $r_0>0$, 求极限
% \[
%   \lim_{(r,\theta)\to(r_0,0)}f(r,\theta)
%   \quad\text{和}\quad
%   \lim_{(r,\theta)\to(r_0,2\pi)}f(r,\theta).
% \]
% \end{example}
% 
% 
% \begin{solve}
% 由定义~\ref{definition:ma-13-2-5} 我们有
% \[
%   \lim_{(r,\theta)\to(r_0,0)}f(r,\theta)
%   =\lim_{(r,\theta)\to(r_0,0)}(r\cos\theta,r\sin\theta)=(r_0,0)
% \]
% 和
% \[
%   \lim_{(r,\theta)\to(r_0,2\pi)}f(r,\theta)
%   =\lim_{(r,\theta)\to(r_0,2\pi)}(r\cos\theta,r\sin\theta)=(r_0,0).
% \qedhere
% \]
% \end{solve}
% 
% 
% \section{多元连续函数}
% \label{sec:ma-13-3}
% 
% \subsection{多元连续函数}
% \label{subsec:ma-13-3-1}
% 
% 与一元函数类似, 我们可以引进多元函数连续性的相关概念.
% 
% \begin{definition}\label{definition:ma-13-3-1}
% 设函数 $u=f(x)$ 在 $E\subset\mathbb R^n$ 上有定义, $x_0\in E$。若 $x_0$ 是 $E$ 的孤立点或当 $x_0\in E'$ 时有
% \[
%   \lim_{E\ni x\to x_0}f(x)=f(x_0),
% \]
% 则称 $f(x)$ 在 $x_0$ 处连续。若 $f(x)$ 在 $E$ 中的每一点都连续, 则称 $f(x)$ 在 $E$ 上连续。若 $\lim\limits_{E\ni x\to x_0}f(x)$ 不存在, 或存在但不等于 $f(x_0)$, 则称 $f(x)$ 在 $x_0$ 处不连续或间断, 此时 $x_0$ 也称做 $f(x)$ 的间断点.
% \end{definition}
% 
% 首先, 当 $x_0$ 是 $E$ 的孤立点时, 我们自然地认为 $f(x)$ 在 $x_0$ 处连续。当 $x_0$ 是 $E$ 的内点时, $f(x)$ 在 $x_0$ 处的连续性是我们在一元函数里通常见到的情形, 但是以上定义还包括了非常广泛的情形, 例如一元函数在闭区间的连续的定义等。读者可以容易将一元连续函数的一些局部性质(即局部有界、保号等)以及和、差、积、商等结果平行推广到多元连续函数的情形.
% 
% \begin{example}\label{example:ma-13-3-1}
% 讨论函数
% \[
%   f(x,y)=
%   \begin{cases}
%     x\sin\dfrac1y, & xy\ne0,\\
%     0, & xy=0
%   \end{cases}
% \]
% 在 $\mathbb R^2$ 中的连续性.
% \end{example}
% 
% 
% \begin{solve}
% 显然当 $(x_0,y_0)\in\mathbb R^2$ 且 $x_0y_0\ne0$ 时, $f(x,y)$ 在 $(x_0,y_0)$ 处连续。由于对于 $\forall y\in(-\infty,+\infty)\setminus\{0\}$, 有
% \[
%   \left|x\sin\frac1y\right|\leqslant |x|,
% \]
% 注意到当 $y=0$ 时, 有 $f(x,0)=0$, 于是对于 $\forall y_0\in(-\infty,+\infty)$, 有
% \[
%   \lim_{(x,y)\to(0,y_0)}f(x,y)=0=f(0,y_0).
% \]
% 这说明 $f(x,y)$ 在 $y$ 轴上连续.
% 
% 当 $x_0\ne0$ 时, 由于
% \[
%   \lim_{(x_0,y)\to(x_0,0)}f(x_0,y)
%   =\lim_{(x_0,y)\to(x_0,0)}x_0\sin\frac1y
% \]
% 不存在, 所以 $\lim\limits_{(x,y)\to(x_0,0)}f(x,y)$ 也不存在。因此 $f(x,y)$ 在 $\mathbb R^2$ 中的间断点集是正负实轴构成的集合, 即 $\{(x,y):x\ne0,y=0\}$.
% \end{solve}
% 
% 
% \subsection{多元连续向量函数}
% \label{subsec:ma-13-3-2}
% 
% 对于多元向量函数, 我们自然地引进下述定义.
% 
% \begin{definition}\label{definition:ma-13-3-2}
% 设 $f(x)=(f_1(x),f_2(x),\cdots,f_m(x))$ 是 $E\subset\mathbb R^n$ 到 $\mathbb R^m$ 的一个 $m$ 维向量函数, 并设 $x_0\in E$。若 $x_0$ 是 $E$ 的孤立点或当 $x_0\in E'$ 时有
% \[
% \begin{aligned}
%   \lim_{E\ni x\to x_0}f(x)
%   &=
%   \left(\lim_{E\ni x\to x_0}f_1(x),
%   \lim_{E\ni x\to x_0}f_2(x),
%   \cdots,\lim_{E\ni x\to x_0}f_m(x)\right)\\
%   &=(f_1(x_0),f_2(x_0),\cdots,f_m(x_0))=f(x_0),
% \end{aligned}
% \]
% 则称 $f(x)$ 在 $x_0$ 处连续。若 $f(x)$ 在 $E$ 中的每一点都连续, 则称 $f(x)$ 在 $E$ 上连续。若存在 $j_0(1\leqslant j_0\leqslant m)$, $f_{j_0}(x)$ 在 $x_0$ 处不连续, 则称 $f(x)$ 在 $x_0$ 处不连续, 并称 $x_0$ 是 $f(x)$ 的间断点.
% \end{definition}
% 
% 由定义可知, $f(x)$ 在 $x_0$ 处连续的充分必要条件是它的每个分量函数 $f_j(x)(j=1,2,\cdots,m)$ 都在 $x_0$ 处连续。容易验证, 对于定义域相同的两个多元连续向量函数, 它们的和、差以及数乘还是多元连续向量函数.
% 
% 下面我们不加证明地给出经常用到的多元复合向量函数的连续性定理.
% 
% \begin{theorem}\label{theorem:ma-13-3-1}
% 设 $E\subset\mathbb R^n$, 向量函数 $y=f(x)=(f_1(x),f_2(x),\cdots,f_m(x))$ 在 $E$ 上连续; $g(y)=(g_1(y),g_2(y),\cdots,g_l(y))$ 在 $f(E)\subset\mathbb R^m$ 上连续, 则
% \[
% \begin{aligned}
%   g(f(x))={}&(g_1(f_1(x),f_2(x),\cdots,f_m(x)),\\
%   &g_2(f_1(x),f_2(x),\cdots,f_m(x)),\\
%   &\cdots,\ g_l(f_1(x),f_2(x),\cdots,f_m(x)))
% \end{aligned}
% \]
% 在 $E$ 上连续.
% \end{theorem}
% 
% 
% \begin{example}\label{example:ma-13-3-2}
% 设函数
% \[
%   f(x,y)=
%   \begin{cases}
%     \sqrt{1-x^2-y^2}\ln(1-x^2-y^2), & x^2+y^2<1,\\
%     0, & x^2+y^2=1,
%   \end{cases}
% \]
% 讨论 $f(x,y)$ 的连续性.
% \end{example}
% 
% 
% \begin{solve}
% 记 $\Delta=\{(x,y):x^2+y^2<1\}$, 则 $f(x,y)$ 的定义域是 $\overline\Delta$。任取 $(x_0,y_0)\in\Delta$, 容易看出 $1-x^2-y^2$ 在 $(x_0,y_0)$ 处连续且值大于零, 从而由复合函数的连续性知 $\sqrt{1-x^2-y^2}$ 及 $\ln(1-x^2-y^2)$ 在 $(x_0,y_0)$ 处连续。所以 $f(x,y)$ 在 $(x_0,y_0)$ 处连续。当 $(x_0,y_0)\in\partial\Delta$ 且 $\Delta\ni(x,y)\to(x_0,y_0)$ 时, 我们有 $t=x^2+y^2\to1-0$。由于
% \[
%   \lim_{t\to1-0}\sqrt{1-t}\ln(1-t)=0,
% \]
% 所以
% \[
%   \lim_{\Delta\ni(x,y)\to(x_0,y_0)}
%   \sqrt{1-x^2-y^2}\ln(1-x^2-y^2)=0=f(x_0,y_0).
% \]
% 另外, 我们显然有
% \[
%   \lim_{\partial\Delta\ni(x,y)\to(x_0,y_0)}f(x,y)=0=f(x_0,y_0),
% \]
% 从而有
% \[
%   \lim_{\overline\Delta\ni(x,y)\to(x_0,y_0)}f(x,y)=0=f(x_0,y_0).
% \]
% 这说明 $f(x,y)$ 在 $\overline\Delta$ 上连续.
% \end{solve}
% 
% 
% \subsection{集合的连通性}
% \label{subsec:ma-13-3-3}
% 
% 设 $D=[\alpha,\beta]\subset\mathbb R(\alpha<\beta)$, 我们通常称 $D$ 上的一个 $n$ 维连续函数
% \[
%   h(t)=(x_1(t),x_2(t),\cdots,x_n(t))
% \]
% 为 $\mathbb R^n$ 中的一条连续曲线(简称曲线)。当 $n=1,2,3$ 时, 连续曲线具有鲜明的几何意义。一条连续曲线也称为 $\mathbb R^n$ 中的一条道路(或路径).
% 
% \begin{definition}\label{definition:ma-13-3-3}
% 设 $E\subset\mathbb R^n$ 是一个非空集合, 若对于 $\forall x,y\in E$, 都存在一条道路 $h(t):t\in[\alpha,\beta]$, 使得 $h(\alpha)=x,h(\beta)=y$, 且对于 $\forall t\in[\alpha,\beta]$, 有 $h(t)\in E$, 则称 $E$ 是道路连通的.
% \end{definition}
% 
% 
% \begin{remark}\label{remark:ma-13-source-68}
% 对于一般集合的连通性在后续课程中有更为广泛的定义与讨论。但在本书中, 我们提到的连通性则是指上述定义中的道路连通。今后, 我们将省略``道路''两字, 仅仅说一个集合是否是连通的.
% \end{remark}
% 
% 在 $\mathbb R$ 中, 一个集合 $E$ 连通的充分必要条件是 $E$ 是一个区间。在 $\mathbb R^2$ 或 $\mathbb R^3$ 中, 若一个集合 $E$ 是连通的, 则它的形式可以多种多样。特别地, 若 $E$ 是连通的开集, 则直观地看 $E$ 只能由``一块''组成.
% 
% 若 $E\subset\mathbb R^n$ 是一个连通的开集, 则称 $E$ 为一个区域。一般用 $D,\Omega$ 等来表示区域。设 $D$ 是一个区域, 称 $D\cup\partial D$ 是一个闭区域。今后我们也经常用 $D,\Omega$ 等来表示一个闭区域。例如, 以原点为心, 1 为半径的单位圆盘 $\Delta=\{(x,y):x^2+y^2<1\}$ 与 $\mathbb R^2\setminus\overline\Delta$ 都是区域。在今后的学习中, 我们遇到比较多的是 $\mathbb R^2$ 和 $\mathbb R^3$ 中的区域。一般来说, 在微积分课程中, 我们在 $\mathbb R^2$ 中遇到的有界区域一般是由有限多条分段光滑曲线所围的区域。如图 13.3.1 所示的区域 $D$, 它是由 $K+1$ 条封闭曲线所围成的有界区域。而在 $\mathbb R^3$ 中, 我们则主要考虑由若干块曲面所围的区域。例如, 单位球面 $x^2+y^2+z^2=1$ 的内部位于 $Oxy$ 平面上方的部分就是一个区域, 它由单位球面的上半部分与 $Oxy$ 平面所围成.
% 
% 区域在多元微积分中所起的作用与开区间在一元微积分中所起的作用相似, 而闭区域在多元微积分中所起的作用与闭区间在一元微积分中所起的作用相似.
% 
% 今后我们还会遇到凸域的概念。设 $D\subset\mathbb R^n$ 是一个区域, 若对于 $\forall x,y\in D$, 有 $tx+(1-t)y\in D(t\in[0,1])$, 即连接 $x$ 与 $y$ 的线段在 $D$ 内, 则称 $D$ 为凸域.
% 
% \subsection{连续函数的性质}
% \label{subsec:ma-13-3-4}
% 
% 在本小节中, 我们将闭区间上一元连续函数的性质推广到多元(或向量)函数的情形。将一个已知数学定理作推广, 考查该定理的证明过程非常重要。例如, 对于有界闭区间上的一元连续函数必有界并取到最大、最小值这一定理, 如果利用有界序列存在收敛子列这一结论来证的话, 则只要求该闭区间是紧集即可。而用类似的论证方法很容易就可以证明高维情形的相应结果.
% 
% \begin{theorem}\label{theorem:ma-13-3-2}
% 设 $E\subset\mathbb R^n$ 为一个紧集, 向量函数
% \[
%   u=f(x)=(f_1(x),f_2(x),\cdots,f_m(x))
% \]
% 在 $E$ 上连续, 则 $f(E)$ 是 $\mathbb R^m$ 中的紧集.
% \end{theorem}
% 
% 
% \begin{proof}
% 首先我们证明 $f(E)$ 是有界集合。倘若结论不真, 则存在 $E$ 中的点列 $\{x_k\}$, 使得
% \[
%   \lim_{k\to\infty}|f(x_k)|=+\infty.
% \]
% 由 $E$ 的紧性, 存在 $\{x_k\}$ 的子序列 $\{x_{k_j}\}$, 使得
% \[
%   \lim_{j\to\infty}x_{k_j}=x_0\in E.
% \]
% 再由 $f(x)$ 在 $E$ 上的连续性知
% \[
%   \lim_{j\to\infty}f(x_{k_j})=f(x_0).
% \]
% 此矛盾证明了 $f(E)$ 的有界性.
% 
% 再证明 $f(E)$ 是闭集。设 $u_0$ 是 $f(E)$ 的一个聚点, 我们要证 $u_0\in f(E)$。由于 $u_0$ 是 $f(E)$ 的聚点, 从而存在 $\{x_k'\}\subset E$, 使得
% \[
%   \lim_{k\to\infty}f(x_k')=u_0.
% \]
% 再由 $E$ 是紧集知, $\{x_k'\}$ 必存在收敛子列 $\{x_{k_j}'\}$。设
% \[
%   \lim_{j\to\infty}x_{k_j}'=x_0',
% \]
% 则有 $x_0'\in E$。最后由 $f(x)$ 在 $x_0'$ 处连续, 有
% \[
%   u_0=\lim_{j\to\infty}f(x_{k_j}')=f(x_0')\in f(E).
% \qedhere
% \]
% \end{proof}
% 
% 
% 由于对 $\mathbb R$ 中的紧集 $E$ 必有 $\sup E=\max E$ 和 $\inf E=\min E$, 从定理~\ref{theorem:ma-13-3-2} 我们可以推出下面的结论.
% 
% \begin{corollary}\label{corollary:ma-13-source-71}
% 设 $E\subset\mathbb R^n$ 为一个紧集, 函数 $f(x)$ 在 $E$ 上连续, 则
% 
% (1) $f(x)$ 在 $E$ 上有界;
% 
% (2) $f(x)$ 在 $E$ 上取到最大、最小值.
% \end{corollary}
% 
% 一个区间上的一元连续函数具有介值性质。这一性质成立的充分必要条件是该函数的值域是一个区间, 而这正是连通集合的特征。对于多元向量函数, 我们有以下定理.
% 
% \begin{theorem}\label{theorem:ma-13-3-3}
% 设 $E\subset\mathbb R^n$ 是连通集, 向量函数 $f(x)=(f_1(x),f_2(x),\cdots,f_m(x))$ 在 $E$ 上连续, 则 $f(E)$ 是 $\mathbb R^m$ 中的连通集.
% \end{theorem}
% 
% 
% \begin{proof}
% 设 $u_1,u_2\in f(E)\subset\mathbb R^m$, 则 $\exists x_1,x_2\in E$, 使得
% \[
%   u_j=f(x_j)\qquad (j=1,2).
% \]
% 由 $E$ 是连通的, 存在连接 $x_1,x_2$ 的道路 $h(t)=(x_1(t),x_2(t),\cdots,x_n(t))(t\in[0,1])$, 即对于 $\forall t\in[0,1]$, 有 $h(t)\in E$, 且 $h(0)=x_1,h(1)=x_2$。容易看出
% \[
% \begin{aligned}
%   f(h(t))={}&(f_1(x_1(t),x_2(t),\cdots,x_n(t)),\\
%   &f_2(x_1(t),x_2(t),\cdots,x_n(t)),\\
%   &\cdots,\ f_m(x_1(t),x_2(t),\cdots,x_n(t)))
% \end{aligned}
% \]
% 是 $f(E)$ 中连接 $u_1,u_2$ 的一条道路。这说明 $f(E)$ 是 $\mathbb R^m$ 中的连通集.
% \end{proof}
% 
% 
% \begin{corollary}\label{corollary:ma-13-source-74}
% 设 $E\subset\mathbb R^n$ 是连通集, 函数 $u=f(x)$ 在 $E$ 上连续, 再设 $u_1,u_2\in f(E)$, 且 $u_1<u_2$, 则对于 $\forall c\in(u_1,u_2)$, 存在 $\xi\in E$, 使得
% \[
%   f(\xi)=c.
% \]
% \end{corollary}
% 
% 
% \begin{proof}
% 取 $x_1,x_2\in E$, 使得 $u_1=f(x_1),u_2=f(x_2)$, 并任取 $E$ 中的道路 $h(t)(0\leqslant t\leqslant1)$ 连接 $x_1,x_2$, 则 $f(h(t))$ 是 $[0,1]$ 上的连续函数。由定理~\ref{theorem:ma-13-3-1}(或一元连续函数的介值定理)知, 存在 $\eta\in(0,1)$, 使得
% \[
%   f(h(\eta))=c.
% \]
% 记 $\xi=h(\eta)\in E$, 即得 $f(\xi)=c$.
% \end{proof}
% 
% 
% 在上述推论中, 当 $E$ 是 $\mathbb R^n(n\geqslant2)$ 的一个区域时, 满足 $f(\xi)=c$ 的集合 $\{\xi\in E:f(\xi)=c\}$ 一般有无穷多个点。这是因为 $E$ 中连接两点可以有无穷多条道路并且任两条道路仅仅只在端点重合.
% 
% 设 $f(x)(x\in E\subset\mathbb R^n)$ 是一个 $n$ 元函数, 如果对于 $\forall\varepsilon>0,\exists\delta>0$, 当 $x',x''\in E$ 且 $|x'-x''|<\delta$ 时, 有
% \[
%   |f(x')-f(x'')|<\varepsilon,
% \]
% 则称 $f(x)$ 在 $E$ 上一致连续。设 $f(x)=(f_1(x),f_2(x),\cdots,f_m(x))(x\in E\subset\mathbb R^n)$ 是一个 $m$ 维向量函数, 如果对于 $\forall j(1\leqslant j\leqslant m)$, $f_j(x)$ 在 $E$ 上一致连续, 则称 $f(x)$ 在 $E$ 上一致连续.
% 
% \begin{theorem}\label{theorem:ma-13-3-4}
% 设 $E\subset\mathbb R^n$ 是紧集, 若向量函数 $f(x)$ 在 $E$ 上连续, 则 $f(x)$ 在 $E$ 一致连续.
% \end{theorem}
% 
% 定理~\ref{theorem:ma-13-3-4} 的证明与有界闭区间上的一元连续函数的情形类似, 请读者自证.
% 
% \begin{example}\label{example:ma-13-3-3}
% 设函数 $f(x,y)=\sin xy$, 证明: $f(x,y)$ 在 $\mathbb R^2$ 中的任何紧集上是一致连续的, 但它在 $\mathbb R^2$ 中不是一致连续的.
% \end{example}
% 
% 
% \begin{proof}
% 显然 $f(x,y)=\sin xy$ 在 $\mathbb R^2$ 内处处连续, 由定理~\ref{theorem:ma-13-3-4} 知 $f(x,y)$ 在 $\mathbb R^2$ 中的任何紧集上是一致连续的.
% 
% 下证 $f(x,y)$ 在 $\mathbb R^2$ 中不一致连续。事实上, 对 $k=1,2,\cdots$, 令
% \[
%   x_k'=\left(k,\frac1k\right)\quad\text{和}\quad
%   x_k''=\left(k,\frac2k\right),
% \]
% 则
% \[
%   \lim_{k\to\infty}|x_k'-x_k''|=\lim_{k\to\infty}\frac1k=0.
% \]
% 但对于 $\forall k\in\mathbb N$, 我们有
% \[
%   \left|f\left(k,\frac1k\right)-f\left(k,\frac2k\right)\right|
%   =|\sin1-\sin2|\ne0.
% \]
% 这就证明了 $f(x,y)$ 在 $\mathbb R^2$ 中不一致连续.
% \end{proof}
% 
% 
% \subsection{同胚映射}
% \label{subsec:ma-13-3-5}
% 
% 设 $E\subset\mathbb R^n$, $y=f(x)=(f_1(x),f_2(x),\cdots,f_m(x))$ 是定义在 $E$ 上的一个向量函数。作为映射, 若 $f(x):E\to f(E)$ 是一个一一对应, 则存在逆映射
% \[
%   x=f^{-1}(y):f(E)\to E.
% \]
% 如果 $f(x)$ 在 $E$ 上连续以及 $f^{-1}(y)$ 在 $f(E)$ 上连续, 则称 $y=f(x)$ 是 $E\to f(E)$ 的同胚映射(简称同胚)。同胚映射 $f(x)$ 也称为 $E$ 到 $f(E)$ 的变换, 它的逆映射也称为逆变换.
% 
% \begin{example}\label{example:ma-13-3-4}
% 证明: 不存在 $\mathbb R$ 到 $\mathbb R^2$ 的同胚映射.
% \end{example}
% 
% 
% \begin{proof}
% 用反证法。倘若存在 $\mathbb R$ 到 $\mathbb R^2$ 的同胚映射 $y=f(x)$。取 $E=\mathbb R\setminus\{0\}$, 则 $E$ 是 $\mathbb R$ 中的不连通集。令 $\mathbb R^2\setminus\{f(0)\}=f(E)$, 则 $f(E)$ 在 $\mathbb R^2$ 中仍然是连通集。由于 $f^{-1}(y)$ 在 $f(E)$ 上连续, 从而它将 $f(E)$ 映成 $\mathbb R$ 中的连通集, 但 $f^{-1}(y)$ 将 $f(E)$ 映成 $E$。此矛盾便证明了我们的结论.
% 
% 对于一个区间上的一元连续函数, 若它存在反函数, 则其反函数必定连续。这是因为在一元函数的情形, 我们涉及的函数具有单调性的缘故。对于多元函数, 则上述结论可以不真.
% 
% 事实上, 考虑极坐标变换
% \[
%   (x,y)=f(r,\theta)=(r\cos\theta,r\sin\theta),
% \]
% 则 $f(r,\theta)$ 将
% \[
%   D=\{(r,\theta):0<r<+\infty,\ 0\leqslant\theta<2\pi\}
% \]
% 连续地并且一一地映成 $\Omega=\mathbb R^2\setminus\{(0,0)\}$, 因此它的逆映射 $(r,\theta)=f^{-1}(x,y)$ 存在。我们可以清楚地写出它的表达式:
% \[
%   r=\sqrt{x^2+y^2},
% \]
% \[
%   \theta=
%   \begin{cases}
%     \arctan(y/x), & x>0,\ y\geqslant0,\\
%     \pi/2, & x=0,\ y>0,\\
%     \pi+\arctan(y/x), & x<0,\\
%     3\pi/2, & x=0,\ y<0,\\
%     2\pi+\arctan(y/x), & x>0,\ y<0.
%   \end{cases}
% \]
% 从几何上来看, 对于极坐标变换
% \[
%   f:
%   \begin{cases}
%     x=r\cos\theta,\\
%     y=r\sin\theta
%   \end{cases}
% \]
% 来说, 例~\ref{example:ma-13-2-6} 告诉我们, 若取定 $r_0>0$, 则当 $(r,\theta)$ 趋于 $(r_0,0)$ 时, 对应的 $(x,y)$ 从第一象限趋于 $(r_0,0)$; 而当 $(r,\theta)$ 趋于 $(r_0,2\pi)$ 时, 对应的 $(x,y)$ 从第四象限趋于 $(r_0,0)$。这说明当 $(x,y)\to(r_0,0)$ 时, $(r,\theta)=f^{-1}(x,y)$ 不存在极限, 即 $f^{-1}(x,y)$ 在正实轴上处处不连续.
% 
% 从上述的例子我们可以看出, 极坐标变换的定义域 $D$ 不是区域, 而 $f(D)$ 是一个区域。若取 $D=\{(r,\theta):0<r<+\infty,\ 0<\theta<2\pi\}$, 则
% \[
%   f(D)=\mathbb R^2\setminus\{(x,y):x\geqslant0,\ y=0\},
% \]
% 即 $f(D)$ 为 $\mathbb R^2$ 中挖去正 $x$ 轴与原点的区域。此时读者不难看出, $f(r,\theta)$ 是 $D$ 到 $f(D)$ 的变换.
% \end{proof}
% 
% 
% \section*{习题十三}
% \phantomsection
% \addcontentsline{toc}{section}{习题十三}
% \label{exercises:ma-13}
% 
% 1。证明 $\mathbb R^n$ 中两点间的距离满足三角不等式: 对于 $\forall x,y$ 和 $z\in\mathbb R^n$, 成立
% \[
%   |x-z|\leqslant |x-y|+|y-z|.
% \]
% 
% 2。若 $\lim\limits_{k\to\infty}|x_k|=+\infty$, 则称 $\mathbb R^n$ 中的点列 $\{x_k\}$ 趋于 $\infty$。现在设点列 $\{x_k=(x_1^k,x_2^k,\cdots,x_n^k)\}$ 趋于 $\infty$, 试判断下列命题是否正确:
% 
% (1) 对于 $\forall i(1\leqslant i\leqslant n)$, 序列 $\{x_i^k\}$ 趋于 $\infty$;
% 
% (2) $\exists i_0(1\leqslant i_0\leqslant n)$, 序列 $\{x_{i_0}^k\}$ 趋于 $\infty$.
% 
% 3。求下列集合的聚点集:
% \[
%   (1)\quad E=\left\{\left(\frac qp,\frac qp,1\right)\in\mathbb R^3:p,q\in\mathbb N\text{ 互素, 且 }q<p\right\};
% \]
% \[
%   (2)\quad E=\left\{\left(\left(\ln\left(1+\frac1k\right)\right)^k,\sin\frac{k\pi}{2}\right):k=1,2,\cdots\right\}.
% \]
% 
% % 整合来源：math-14.tex
% \chapter{多元微分学}
% \label{chap:ma-14}
% 
% 在本章中, 我们将讨论多元函数的微分学。将一元函数微分学的理论推广到多元函数 (向量) 函数是本章的中心内容。除此之外, 由于多元函数与一元函数具有一些本质的差别, 在本章中我们还要研究一些只有在多元函数的情形才具有的问题, 如隐函数存在定理等.
% 
% \section{偏导数与全微分}
% \label{sec:ma-14-1}
% 
% 对于一个多元函数 $u=f(x_1,x_2,\cdots,x_n)$, 用一元微分学的工具来研究它是一个很自然的问题, 因此引进偏导数以及方向导数等概念是水到渠成的事情。在本节中, 我们主要引进这些概念, 并讨论它们的简单性质.
% 
% \subsection{偏导数}
% \label{subsec:ma-14-1-1}
% 
% \begin{definition}\label{definition:ma-14-1-1}
% 设函数 $u=f(x)=f(x_1,x_2,\cdots,x_n)$ 在区域 $D\subset\mathbb R^n$ 上有定义, $x_0=(x_1^0,x_2^0,\cdots,x_n^0)\in D$。对于 $1\leqslant i\leqslant n$, 若一元函数
% \[
%   f(x_1^0,\cdots,x_{i-1}^0,x_i,x_{i+1}^0,\cdots,x_n^0)
% \]
% 在 $x_i^0$ 处的导数, 即
% \[
%   \lim_{x_i\to x_i^0}
%   \frac{f(x_1^0,\cdots,x_{i-1}^0,x_i,x_{i+1}^0,\cdots,x_n^0)
%   -f(x_1^0,x_2^0,\cdots,x_n^0)}{x_i-x_i^0}
% \]
% 存在, 则称 $f(x)$ 在 $x_0$ 处关于 $x_i$ 可偏导, 并称上述极限为 $f(x)$ 在 $x_0$ 处关于 $x_i$ 的偏导数, 记为 $\dfrac{\partial f(x_0)}{\partial x_i}, f'_{x_i}(x_0), \left.\dfrac{\partial u}{\partial x_i}\right|_{x_0}$ 等.
% \end{definition}
% 
% 容易看出, 当一个 $n$ 元函数 $f(x)$ 关于每个分量都可偏导时, $f(x)$ 具有 $n$ 个偏导数。若 $f(x)$ 在 $D$ 内的每一点关于 $x_i$ 都可偏导, 则将其偏导数记为 $\dfrac{\partial f(x)}{\partial x_i}, f'_{x_i}(x), \dfrac{\partial u}{\partial x_i}$ 等。此时, 对每个固定的分量 $x_i$, $\dfrac{\partial f(x)}{\partial x_i}$ 仍然是 $D$ 中的一个 $n$ 元函数.
% 
% 特别地, 若 $z=f(x,y)$ 是一个二元函数, 则我们用 $\dfrac{\partial f(x,y)}{\partial x}$ (或 $f'_x(x,y), \dfrac{\partial z}{\partial x}$) 与 $\dfrac{\partial f(x,y)}{\partial y}$ (或 $f'_y(x,y), \dfrac{\partial z}{\partial y}$) 来记它的两个偏导数。同理, 对三元函数 $u=f(x,y,z)$, 可以类似地给出相应记号: $\dfrac{\partial f(x,y,z)}{\partial x}, \dfrac{\partial f(x,y,z)}{\partial y}, \dfrac{\partial f(x,y,z)}{\partial z}$, 或 $f'_x(x,y,z), f'_y(x,y,z), f'_z(x,y,z)$, 或 $\dfrac{\partial u}{\partial x}, \dfrac{\partial u}{\partial y}, \dfrac{\partial u}{\partial z}$.
% 
% 对一个二元函数 $z=f(x,y)((x,y)\in D)$, 我们可以讨论其在 $(x_0,y_0)\in D$ 处偏导数的几何意义。当函数具有较好的性质时, 该函数的图像是 $\mathbb R^3$ 中的一块曲面。$y=y_0$ 是一个过点 $(x_0,y_0,0)$ 且平行于 $Ozx$ 平面的平面, 它与曲面 $z=f(x,y)((x,y)\in D)$ 的交线为
% \[
%   L:\begin{cases}
%     x=x,\\
%     y=y_0,\\
%     z=f(x,y_0).
%   \end{cases}
% \]
% 因此, $\dfrac{\partial f(x_0,y_0)}{\partial x}$ 是平面 $y=y_0$ 内的曲线 $L$ 在 $x=x_0$ 处的切线斜率 $k_1$。同样, $\dfrac{\partial f(x_0,y_0)}{\partial y}$ 是平面 $x=x_0$ 与曲面 $z=f(x,y)((x,y)\in D)$ 的交线在 $y=y_0$ 处的切线斜率 $k_2$ (在图 14.1.1 中, $k_1=\tan\alpha, k_2=\tan\beta$).
% 
% 读者应该注意的是, 函数 $z=f(x,y)$ 在 $(x_0,y_0)$ 处的两个偏导数存在与否只与曲面与 $x=x_0$ 及 $y=y_0$ 的交线性质有关。因此, 一个二元函数在点 $(x_0,y_0)$ 处即使两个偏导数都存在, 该函数在点 $(x_0,y_0)$ 处也未必连续。例如, 设函数
% \[
%   f(x,y)=
%   \begin{cases}
%     0, & xy=0,\\
%     1, & xy\ne0,
%   \end{cases}
% \]
% 容易看出 $f'_x(0,0)=f'_y(0,0)=0$, 但 $f(x,y)$ 在 $(0,0)$ 处不连续.
% 
% 下面我们举两个例子:
% 
% \begin{example}\label{example:ma-14-1-1}
% 设函数 $f(x,y)=\tan\dfrac{x^2}{y}$, 求 $\dfrac{\partial f(0,1)}{\partial x}$ 及 $\dfrac{\partial f(0,1)}{\partial y}$.
% \end{example}
% 
% 
% \begin{solve}
% 由偏导数的定义有
% \[
%   \frac{\partial f(0,1)}{\partial x}
%   =(\tan x^2)'\big|_{x=0}
%   =2x\sec^2 x^2\big|_{x=0}=0,
% \]
% \[
%   \frac{\partial f(0,1)}{\partial y}
%   =(0)'\big|_{y=1}=0.
% \qedhere
% \]
% \end{solve}
% 
% 
% \begin{example}\label{example:ma-14-1-2}
% 设函数 $f(x,y)=\arcsin\dfrac{x}{\sqrt{x^2+y^2}}$, 求 $f'_x(x,y)$ 及 $f'_y(x,y)$.
% \end{example}
% 
% 
% \begin{solve}
% 对这种计算题, 我们在求 $f'_x(x,y)$ 时将 $y$ 看成常数即可。因此, 我们有
% \[
% \begin{aligned}
%   f'_x(x,y)
%   &=
%   \frac{1}{\sqrt{1-\dfrac{x^2}{x^2+y^2}}}\cdot
%   \frac{1}{x^2+y^2}
%   \left(\sqrt{x^2+y^2}-\frac{x^2}{\sqrt{x^2+y^2}}\right)\\
%   &=\frac{1}{\sqrt{y^2}}\cdot\frac{y^2}{x^2+y^2}
%   =\frac{y\operatorname{sgn} y}{x^2+y^2};
% \end{aligned}
% \]
% 同理,
% \[
%   f'_y(x,y)
%   =
%   \frac{1}{\sqrt{1-\dfrac{x^2}{x^2+y^2}}}
%   \left(-\frac{xy}{(x^2+y^2)^{3/2}}\right)
%   =\frac{-x\operatorname{sgn} y}{x^2+y^2},
% \]
% 其中 $\operatorname{sgn} y$ 是符号函数.
% \end{solve}
% 
% 
% \subsection{方向导数}
% \label{subsec:ma-14-1-2}
% 
% 一个多元函数的方向导数是用来刻画定义域内从一点出发的射线上函数的变化情况的。从一点出发的一条射线可以看成是一个方向, 而一个方向可由以原点为心的单位球面上的点所确定。设 $n$ 维单位球面为 $x_1^2+x_2^2+\cdots+x_n^2=1$, 则每一个以原点为起点, 终点在该球面上的有向线段表示一个单位向量 $v$, 它可以记为 $v=(\cos\theta_1,\cos\theta_2,\cdots,\cos\theta_n)$, 其中 $\theta_i(1\leqslant i\leqslant n)$ 是该向量与 $x_i$ 轴正向的夹角, $\cos\theta_i(1\leqslant i\leqslant n)$ 也称为是 $v$ 的方向余弦。特别地, 对于 $\mathbb R^2$ 内任一个单位向量 $v$, 记它与正实轴的夹角为 $\theta$, 则有 $v=(\cos\theta,\sin\theta)$。有了这些准备, 我们可以给出以下定义.
% 
% \begin{definition}\label{definition:ma-14-1-2}
% 设函数 $u=f(x)$ 在区域 $D\subset\mathbb R^n$ 上有定义, $x_0\in D$,
% \[
%   v=(\cos\theta_1,\cos\theta_2,\cdots,\cos\theta_n)
% \]
% 为一方向。如果极限
% \[
%   \lim_{t\to0+0}\frac{f(x_0+tv)-f(x_0)}{t}
% \]
% 存在, 则称该极限为 $f(x)$ 在 $x_0$ 处沿方向 $v$ 的方向导数, 记为 $\dfrac{\partial f(x_0)}{\partial v}$ 或 $\left.\dfrac{\partial u}{\partial v}\right|_{x_0}$.
% \end{definition}
% 
% 
% \paragraph*{注 1}
% 在上述定义中, 若记 $x_0+tv=x$, 则
% \[
%   \frac{\partial f(x_0)}{\partial v}
%   =\lim_{x\to x_0}\frac{f(x)-f(x_0)}{(x-x_0)v}.
% \]
% 从上述定义可以看出, 一个函数在某点处的方向导数不依赖于坐标系的选取.
% 
% \paragraph*{注 2}
% 设 $f(x)$ 在 $x_0=(x_1^0,x_2^0,\cdots,x_n^0)$ 处关于 $x_i$ 的偏导数 $\dfrac{\partial f(x_0)}{\partial x_i}(1\leqslant i\leqslant n)$ 存在并等于 $A$, 若记 $v_i$ 是第 $i$ 个分量为 $1$ 的单位向量, 则 $\dfrac{\partial f(x_0)}{\partial v_i}=A$。但读者应该注意的是 $\dfrac{\partial f(x_0)}{\partial(-v_i)}=-A$.
% 
% \begin{example}\label{example:ma-14-1-3}
% 设函数 $f(x,y)=\sqrt{x^2+y^2}$, 试求 $f(x,y)$ 在 $(0,0)$ 处各个方向的方向导数.
% \end{example}
% 
% 
% \begin{solve}
% 设 $v=(\cos\theta,\sin\theta)(0\leqslant\theta<2\pi)$, 则
% \[
% \begin{aligned}
%   \frac{\partial f(0,0)}{\partial v}
%   &=\lim_{t\to0+0}
%   \frac{f(t\cos\theta,t\sin\theta)-f(0,0)}{t}\\
%   &=\lim_{t\to0+0}
%   \frac{\sqrt{t^2\cos^2\theta+t^2\sin^2\theta}}{t}\\
%   &=\lim_{t\to0+0}\frac{t}{t}=1.
% \end{aligned}
% \]
% 
% 从注 2 我们可知, 在上例中, $f(x,y)$ 在 $(0,0)$ 处两个偏导数都不存在。当然, 我们也可从偏导数的定义直接证明这一结果.
% \end{solve}
% 
% 
% \begin{example}\label{example:ma-14-1-4}
% 设函数
% \[
%   f(x,y)=
%   \begin{cases}
%     1, & y=x^2,\ \text{且 } x\ne0,\\
%     0, & \text{其他},
%   \end{cases}
% \]
% 即 $f(x,y)$ 在抛物线 $y=x^2$ 上除去原点外取值为 $1$, 其余各处取值为 $0$。证明: 对任意方向 $v$, 有 $\dfrac{\partial f(0,0)}{\partial v}=0$, 但 $f(x,y)$ 在 $(0,0)$ 处不连续.
% \end{example}
% 
% 
% \begin{proof}
% 记 $\mathbf 0=(0,0)$, 对任意方向 $v$, 当 $t$ 很小时, $\mathbf 0+tv$ 与抛物线不交, 从而 $f(tv)=0=f(0,0)$, 因此 $\dfrac{\partial f(0,0)}{\partial v}=0$。由于
% \[
%   \lim_{k\to\infty}f\left(\frac1k,\frac1{k^2}\right)=1\ne f(0,0),
% \]
% 所以 $f(x,y)$ 在 $(0,0)$ 处不连续.
% \end{proof}
% 
% 
% \subsection{全微分}
% \label{subsec:ma-14-1-3}
% 
% 一元函数 $y=f(x)$ 的微分是函数关于自变量的一阶线性的近似, 由其几何意义, 实际上是用切线来局部近似曲线 $y=f(x)$。对于一个二元函数, 由于它的图像一般是一块曲面, 因此需用切平面来局部近似它, 而这些都与多元函数的全微分有关。对一个 $n$ 元函数, 我们有以下定义.
% 
% \begin{definition}\label{definition:ma-14-1-3}
% 设函数 $f(x)=f(x_1,x_2,\cdots,x_n)$ 在区域 $D\subset\mathbb R^n$ 上有定义, 且 $x_0=(x_1^0,x_2^0,\cdots,x_n^0)\in D$。记 $\Delta x=(\Delta x_1,\Delta x_2,\cdots,\Delta x_n)$, 并称它为自变量的全增量。再设
% \[
%   x=x_0+\Delta x=(x_1+\Delta x_1,x_2+\Delta x_2,\cdots,x_n+\Delta x_n)\in D.
% \]
% 若存在仅依赖于 $x_0$ 的常数 $A_i(i=1,2,\cdots,n)$, 使得
% \[
%   \Delta f(x_0)=f(x_0+\Delta x)-f(x_0)
%   =\sum_{i=1}^n A_i\Delta x_i+o(|\Delta x|),\quad |\Delta x|\to0,
% \]
% 则称 $f(x)$ 在 $x_0$ 处可微, 并称 $\sum\limits_{i=1}^n A_i\Delta x_i$ 为 $f(x)$ 在 $x_0$ 处的全微分, 记为 $\mathrm df(x_0)$, 即
% \[
%   \mathrm df(x_0)=\sum_{i=1}^n A_i\Delta x_i.
% \]
% 若 $f(x)$ 在 $D$ 内每一点处均可微, 则称 $f(x)$ 在 $D$ 内是可微函数.
% \end{definition}
% 
% 当 $x_i(1\leqslant i\leqslant n)$ 是自变量时, 定义 $\mathrm dx_i=\Delta x_i$。因此 $f(x)$ 在 $x_0$ 处的全微分可以记为 $\mathrm df(x_0)=\sum\limits_{i=1}^n A_i\mathrm dx_i$.
% 
% 从全微分的定义可立即推出, 当 $f(x)$ 在 $x_0$ 处可微时, 它的全微分是函数改变量的线性主要部分。本书中涉及一个多元函数的微分时, 一般情形下总是指上述的全微分, 因此今后我们将全微分简称为微分.
% 
% \begin{theorem}\label{theorem:ma-14-1-1}
% 设函数 $f(x)$ 在区域 $D\subset\mathbb R^n$ 上有定义, 在 $x_0=(x_1^0,\cdots,x_n^0)\in D$ 处可微, 记其微分为 $\mathrm df(x_0)=\sum\limits_{i=1}^n A_i\mathrm dx_i$, 则
% \[
% \begin{gathered}
%   (1)\quad f(x)\text{ 在 }x_0\text{ 处连续};\\
%   (2)\quad \text{对于 } i(1\leqslant i\leqslant n), f(x)\text{ 关于 }x_i\text{ 可偏导, 并且有 }
%   \frac{\partial f(x_0)}{\partial x_i}=A_i.
% \end{gathered}
% \]
% \end{theorem}
% 
% 
% \begin{proof}
% (1) 记 $\Delta x=x-x_0$, 则 $x\to x_0$ 等价于 $\Delta x\to0$, 即它等价于: 对于 $\forall i(1\leqslant i\leqslant n)$, 有 $\Delta x_i\to0$。由此我们有
% \[
%   \lim_{\Delta x\to0}[f(x_0+\Delta x)-f(x_0)]
%   =
%   \lim_{\Delta x\to0}\left[\sum_{i=1}^n A_i\Delta x_i+o(|\Delta x|)\right]=0,
% \]
% 即 $f(x)$ 在 $x_0$ 处连续.
% 
% (2) 对任何固定的 $i(1\leqslant i\leqslant n)$, 当 $j=1,2,\cdots,n$ 且 $j\ne i$ 时, 令 $x_j=x_j^0$, 即 $\Delta x_j=0$。此时我们有 $\Delta x=(0,\cdots,\Delta x_i,\cdots,0)$ 和 $|\Delta x|=|\Delta x_i|$, 因此有
% \[
% \begin{aligned}
% &\lim_{\Delta x_i\to0}
%   \frac{f(x_1^0,x_2^0,\cdots,x_{i-1}^0,x_i^0+\Delta x_i,x_{i+1}^0,\cdots,x_n^0)-f(x_0)}{\Delta x_i}\\
% &\quad =
%   \lim_{\Delta x_i\to0}\left(A_i+\frac{o(|\Delta x_i|)}{\Delta x_i}\right)=A_i,
% \end{aligned}
% \]
% 从而 $f(x)$ 关于 $x_i$ 可偏导, 并且有 $\dfrac{\partial f(x_0)}{\partial x_i}=A_i$.
% \end{proof}
% 
% 
% 上述定理说明, 当函数 $f(x)$ 在 $x_0$ 处可微时, 必有
% \[
%   \mathrm df(x_0)=\sum_{i=1}^n\frac{\partial f(x_0)}{\partial x_i}\mathrm dx_i.
% \]
% 特别地, 当 $z=f(x,y)$ 在 $(x_0,y_0)$ 处可微时, 我们有
% \[
%   \mathrm df(x_0,y_0)=\frac{\partial f(x_0,y_0)}{\partial x}\mathrm dx+
%   \frac{\partial f(x_0,y_0)}{\partial y}\mathrm dy;
% \]
% 而当 $u=f(x,y,z)$ 在 $(x_0,y_0,z_0)$ 处可微时, 则有
% \[
%   \mathrm df(x_0,y_0,z_0)=
%   \frac{\partial f(x_0,y_0,z_0)}{\partial x}\mathrm dx+
%   \frac{\partial f(x_0,y_0,z_0)}{\partial y}\mathrm dy+
%   \frac{\partial f(x_0,y_0,z_0)}{\partial z}\mathrm dz.
% \]
% 
% 我们已经有了例子表明存在可偏导但不连续的多元函数, 因此这样的函数不可微。这说明, 可偏导与可微在多元函数的情形不是等价的。但是我们有以下的定理.
% 
% \begin{theorem}\label{theorem:ma-14-1-2}
% 设函数 $f(x)$ 在区域 $D\subset\mathbb R^n$ 上有定义, $x_0=(x_1^0,x_2^0,\cdots,x_n^0)\in D$, 再设 $f(x)$ 在 $x_0$ 的邻域 $U(x_0,\delta_0)(\delta_0>0)$ 内存在各个偏导数, 并且这些偏导数在 $x_0$ 处连续, 则 $f(x)$ 在 $x_0$ 处可微.
% \end{theorem}
% 
% 
% \begin{proof}
% 我们对 $n$ 作归纳法来证明定理.
% 
% 当 $n=1$ 时, 定理结论显然成立.
% 
% 假设当 $n=k$ 时定理成立, 即
% \[
% \begin{aligned}
% &f(x_1^0+\Delta x_1,x_2^0+\Delta x_2,\cdots,x_k^0+\Delta x_k)
%   -f(x_1^0,x_2^0,\cdots,x_k^0)\\
% &\quad =
%   \sum_{i=1}^k
%   \frac{\partial f(x_1^0,x_2^0,\cdots,x_k^0)}{\partial x_i}\Delta x_i
%   +o(1)\left(\sqrt{\sum_{i=1}^k(\Delta x_i)^2}\right)
% \end{aligned}
% \]
% \[
%   \left(\sqrt{\sum_{i=1}^k(\Delta x_i)^2}\to0\right).
% \]
% 当 $n=k+1$ 时,
% \[
% \begin{aligned}
% &\Delta f(x_1^0,\cdots,x_k^0,x_{k+1}^0)\\
% &=f(x_1^0+\Delta x_1,\cdots,x_k^0+\Delta x_k,x_{k+1}^0+\Delta x_{k+1})
%   -f(x_1^0,\cdots,x_k^0,x_{k+1}^0)\\
% &=\big[f(x_1^0+\Delta x_1,\cdots,x_k^0+\Delta x_k,x_{k+1}^0+\Delta x_{k+1})
%   -f(x_1^0+\Delta x_1,\cdots,x_k^0+\Delta x_k,x_{k+1}^0)\big]\\
% &\quad+\big[f(x_1^0+\Delta x_1,\cdots,x_k^0+\Delta x_k,x_{k+1}^0)
%   -f(x_1^0,\cdots,x_k^0,x_{k+1}^0)\big].
% \end{aligned}
% \]
% 由一元函数的拉格朗日微分中值定理及 $\dfrac{\partial f(x)}{\partial x_{k+1}}$ 在 $x_0$ 处的连续性知
% \begin{equation}
% \begin{aligned}
% &f(x_1^0+\Delta x_1,\cdots,x_k^0+\Delta x_k,x_{k+1}^0+\Delta x_{k+1})
%   -f(x_1^0+\Delta x_1,\cdots,x_k^0+\Delta x_k,x_{k+1}^0)\\
% &=f'_{x_{k+1}}(x_1^0+\Delta x_1,\cdots,x_k^0+\Delta x_k,
%   x_{k+1}^0+\theta\Delta x_{k+1})\Delta x_{k+1}\\
% &=f'_{x_{k+1}}(x_1^0,\cdots,x_k^0,x_{k+1}^0)\Delta x_{k+1}
%   +o(1)(|\Delta x_{k+1}|),
% \end{aligned}
% \label{eq:ma-14-1-1}
% \end{equation}
% 其中 $0<\theta<1$。由归纳法假设有
% \begin{equation}
% \begin{aligned}
% &f(x_1^0+\Delta x_1,\cdots,x_k^0+\Delta x_k,x_{k+1}^0)
%   -f(x_1^0,\cdots,x_k^0,x_{k+1}^0)\\
% &=
% \sum_{i=1}^k
% \frac{\partial f(x_1^0,\cdots,x_k^0,x_{k+1}^0)}{\partial x_i}\Delta x_i
% +o(1)\left(\sqrt{\sum_{i=1}^k(\Delta x_i)^2}\right)
% \end{aligned}
% \label{eq:ma-14-1-2}
% \end{equation}
% \[
%   \left(\sqrt{\sum_{i=1}^k(\Delta x_i)^2}\to0\right).
% \]
% 注意到当 $\sqrt{\sum\limits_{i=1}^{k+1}(\Delta x_i)^2}\to0$ 时, 有
% \[
%   \frac{
%   o(1)(|\Delta x_{k+1}|)
%   +o(1)\left(\sqrt{\sum\limits_{i=1}^k(\Delta x_i)^2}\right)}
%   {\sqrt{\sum\limits_{i=1}^{k+1}(\Delta x_i)^2}}\to0,
% \]
% 因此由式 \eqref{eq:ma-14-1-1}, \eqref{eq:ma-14-1-2} 即得
% \[
% \begin{aligned}
% &\Delta f(x_1^0,\cdots,x_k^0,x_{k+1}^0)\\
% &=
%   f(x_1^0+\Delta x_1,\cdots,x_k^0+\Delta x_k,x_{k+1}^0+\Delta x_{k+1})
%   -f(x_1^0,\cdots,x_k^0,x_{k+1}^0)\\
% &=\sum_{i=1}^{k+1}
%   \frac{\partial f(x_1^0,x_2^0,\cdots,x_k^0,x_{k+1}^0)}{\partial x_i}\Delta x_i
%   +o(1)\left(\sqrt{\sum_{i=1}^{k+1}(\Delta x_i)^2}\right),
% \end{aligned}
% \]
% 即 $f(x_1,x_2,\cdots,x_{k+1})$ 在 $(x_1^0,x_2^0,\cdots,x_{k+1}^0)$ 处可微, 且
% \[
%   \mathrm df(x_1^0,x_2^0,\cdots,x_{k+1}^0)
%   =\sum_{i=1}^{k+1}
%   \frac{\partial f(x_1^0,x_2^0,\cdots,x_{k+1}^0)}{\partial x_i}\mathrm dx_i.
% \]
% 这就证明了定理的结论对一切 $n\in\mathbb N$ 成立.
% \end{proof}
% 
% 
% \begin{remark}\label{remark:ma-14-source-16}
% 若函数 $f(x)$ 在区域 $D\subset\mathbb R^n$ 上关于自变量的各个分量都具有连续偏导数, 通常我们称 $f(x)$ 在 $D$ 内是 $C^1$ 的, 记为 $f(x)\in C^1(D)$, 此时我们也称 $f(x)$ 在 $D$ 内连续可微.
% \end{remark}
% 
% 
% \begin{example}\label{example:ma-14-1-5}
% 证明函数
% \[
%   f(x,y)=
%   \begin{cases}
%     (x^2+y^2)\sin\dfrac1{x^2+y^2}, & x^2+y^2\ne0,\\
%     0, & x^2+y^2=0
%   \end{cases}
% \]
% 在 $(0,0)$ 处可微, 但 $f'_x(x,y)$ 及 $f'_y(x,y)$ 在 $(0,0)$ 处不连续.
% \end{example}
% 
% 
% \begin{proof}
% 显然 $f'_x(0,0)=f'_y(0,0)=0$。由于
% \[
%   \Delta f(0,0)=f(\Delta x,\Delta y)
%   =[(\Delta x)^2+(\Delta y)^2]\sin\frac1{(\Delta x)^2+(\Delta y)^2}
%   =o(1)\sqrt{(\Delta x)^2+(\Delta y)^2},
% \]
% \[
%   \left(\sqrt{(\Delta x)^2+(\Delta y)^2}\to0\right),
% \]
% 因此 $f(x,y)$ 在 $(0,0)$ 处可微。当 $(x,y)\ne(0,0)$ 时, 有
% \[
%   f'_x(x,y)=2x\sin\frac1{x^2+y^2}-\frac{2x}{x^2+y^2}\cos\frac1{x^2+y^2}.
% \]
% 由于 $(x_k,y_k)=\left(\dfrac1{\sqrt{2k\pi}},0\right)\to(0,0)$ 时, $f'_x(x_k,y_k)\to\infty$, 因此 $f'_x(x,y)$ 在 $(0,0)$ 处不连续.
% 
% 由函数的对称性, 同理 $f'_y(x,y)$ 在 $(0,0)$ 处也不连续.
% 
% 对于可微函数, 它的方向导数与偏导数密切相关。对此我们有下面的结论.
% \end{proof}
% 
% 
% \begin{theorem}\label{theorem:ma-14-1-3}
% 设函数 $f(x)$ 在区域 $D\subset\mathbb R^n$ 上有定义, 且在 $x_0\in D$ 处可微, 则 $f(x)$ 在 $x_0$ 处沿方向 $v=(\cos\theta_1,\cos\theta_2,\cdots,\cos\theta_n)$ 的方向导数为
% \[
%   \frac{\partial f(x_0)}{\partial v}
%   =\sum_{i=1}^n\frac{\partial f(x_0)}{\partial x_i}\cos\theta_i.
% \]
% \end{theorem}
% 
% 
% \begin{proof}
% 由 $f(x)$ 的可微性及方向导数的定义, 我们有
% \[
% \begin{aligned}
%   \lim_{t\to0+0}\frac{f(x_0+tv)-f(x_0)}{t}
%   &=\lim_{t\to0+0}
%   \frac{\sum\limits_{i=1}^n\dfrac{\partial f(x_0)}{\partial x_i}t\cos\theta_i+o(|t|)}{t}\\
%   &=\sum_{i=1}^n\frac{\partial f(x_0)}{\partial x_i}\cos\theta_i.
% \end{aligned}
% \qedhere
% \]
% \end{proof}
% 
% 
% 在本小节的最后, 我们再举一个例子来说明多元函数微分的在近似计算中的应用.
% 
% \begin{example}\label{example:ma-14-1-6}
% 求 $2.99^2\times1.02^3$ 的近似值.
% \end{example}
% 
% 
% \begin{solve}
% 取函数 $f(x,y)=x^2y^3$, 显然 $f(x,y)$ 的两个偏导数在 $\mathbb R^2$ 上连续, 从而它在 $\mathbb R^2$ 处处可微。取 $(x_0,y_0)=(3,1), \Delta x=-0.01, \Delta y=0.02$, 则
% \[
% \begin{aligned}
%   2.99^2\times1.02^3
%   &=f(2.99,1.02)=f(3-0.01,1+0.02)\\
%   &\approx f(3,1)+f'_x(3,1)\Delta x+f'_y(3,1)\Delta y\\
%   &=3^2\times1^3+f'_x(3,1)\cdot(-0.01)+f'_y(3,1)\cdot(0.02)\\
%   &=9+6\cdot(-0.01)+27\cdot(0.02)=9.48.
% \end{aligned}
% \qedhere
% \]
% \end{solve}
% 
% 
% \subsection{梯度}
% \label{subsec:ma-14-1-4}
% 
% 与方向导数密切相关的概念是梯度。在本节我们先介绍梯度的概念与基本的性质, 在本书后面的章节中我们还会多次讨论它的进一步的性质.
% 
% 在上一小节中, 我们知道当 $f(x)$ 在 $x_0\in\mathbb R^n$ 处可微时, $f(x)$ 沿任意方向 $v=(\cos\theta_1,\cos\theta_2,\cdots,\cos\theta_n)$ 的方向导数为
% \[
%   \frac{\partial f(x_0)}{\partial v}=
%   \sum_{i=1}^n\frac{\partial f(x_0)}{\partial x_i}\cos\theta_i.
% \]
% 当 $f(x)$ 在 $x_0$ 处的 $n$ 个偏导数不全为零时, $f(x)$ 沿方向
% \[
%   v_0=
%   \frac1{\sqrt{\sum\limits_{i=1}^n\left(\dfrac{\partial f(x_0)}{\partial x_i}\right)^2}}
%   \left(\frac{\partial f(x_0)}{\partial x_1},\frac{\partial f(x_0)}{\partial x_2},\cdots,\frac{\partial f(x_0)}{\partial x_n}\right)
% \]
% 的方向导数达到最大值。因此向量
% \[
%   \left(\frac{\partial f(x_0)}{\partial x_1},\frac{\partial f(x_0)}{\partial x_2},\cdots,\frac{\partial f(x_0)}{\partial x_n}\right)
% \]
% 是 $f(x)$ 在 $x_0$ 处方向导数达到最大的方向, 同时它的模就是该方向的方向导数。由此我们引进下面的定义.
% 
% \begin{definition}\label{definition:ma-14-1-4}
% 设函数 $f(x)$ 在 $x_0$ 处可微, 则称向量
% \[
%   \left(\frac{\partial f(x_0)}{\partial x_1},\frac{\partial f(x_0)}{\partial x_2},\cdots,\frac{\partial f(x_0)}{\partial x_n}\right)
% \]
% 为 $f(x)$ 在 $x_0$ 处的梯度, 记为 $\operatorname{grad}f(x_0)$, 即
% \[
%   \operatorname{grad}f(x_0)=
%   \left(\frac{\partial f(x_0)}{\partial x_1},\frac{\partial f(x_0)}{\partial x_2},\cdots,\frac{\partial f(x_0)}{\partial x_n}\right).
% \]
% \end{definition}
% 
% 设 $f,g$ 均是可微函数, 则从定义容易推出梯度有以下简单性质:
% \[
% \begin{aligned}
%   &(1)\quad \operatorname{grad}C=0,\ C\text{ 为常数};\\
%   &(2)\quad \text{对于 } \forall\alpha,\beta\in\mathbb R,\ \text{有 }
%   \operatorname{grad}(\alpha f+\beta g)=\alpha\operatorname{grad}f+\beta\operatorname{grad}g;\\
%   &(3)\quad \operatorname{grad}(f\cdot g)=f\cdot\operatorname{grad}g+g\cdot\operatorname{grad}f;\\
%   &(4)\quad \operatorname{grad}\frac fg=\frac1{g^2}(g\cdot\operatorname{grad}f-f\cdot\operatorname{grad}g)\quad(g\ne0).
% \end{aligned}
% \]
% 从梯度的定义可以看出, 当 $f(x)$ 在 $x_0$ 处可微时, $f(x)$ 沿方向 $v$ 的方向导数可以简单地记成
% \[
%   \frac{\partial f(x_0)}{\partial v}=\operatorname{grad}f(x_0)\cdot v.
% \]
% 读者易于给出方向导数与梯度的关系的几何意义.
% 
% 设函数 $u=f(x,y,z)$ 在 $(x_0,y_0,z_0)$ 处可微, $\operatorname{grad}f(x_0,y_0,z_0)$ 不是零向量。记 $H=|\operatorname{grad}f(x_0,y_0,z_0)|$, $l_v$ 为从 $(x_0,y_0,z_0)$ 出发由方向 $v$ 确定的射线, 则对于 $\forall h\in(-H,H)$, 所有使得 $\dfrac{\partial f}{\partial v}=h$ 的 $l_v$ 的集合构成了顶点在 $(x_0,y_0,z_0)$, 母线是 $l_v$ 的锥面。当 $h$ 从 $-H$ 连续变为 $H$ 时, 其母线 $l_v$ 与 $\operatorname{grad}f(x_0,y_0,z_0)$ 的夹角从 $\pi$ 变到 $0$。特别地, 当它们的夹角为 $\dfrac\pi2$ 时, $l_v$ 的全体构成了过 $(x_0,y_0,z_0)$ 且以 $\operatorname{grad}f(x_0,y_0,z_0)$ 为法线的平面.
% 
% \begin{example}\label{example:ma-14-1-7}
% 设 $k\geqslant2$ 为正整数, 函数 $f(x,y)$ 在极坐标 $(r,\theta)$ 下有表示式
% \[
%   f(x,y)=
%   \begin{cases}
%     r\sin k\theta, &(x,y)\ne(0,0),\\
%     0, &(x,y)=(0,0).
%   \end{cases}
% \]
% (1) 求 $f(x,y)$ 在 $(0,0)$ 处的方向导数; (2) $f(x,y)$ 在 $(0,0)$ 处是否连续? (3) $f(x,y)$ 在 $(0,0)$ 处是否可微?
% \end{example}
% 
% 
% \begin{solve}
% (1) 对固定的 $\theta_0\in[0,2\pi)$, $\theta=\theta_0$ 即确定一个从原点出发的方向 $v_0=(\cos\theta_0,\sin\theta_0)$。我们有
% \[
%   \frac{\partial f(0,0)}{\partial v_0}
%   =\lim_{r\to0}\frac{r\sin k\theta_0-0}{r}
%   =\sin k\theta_0.
% \]
% 因此 $f(x,y)$ 在 $(0,0)$ 处沿方向 $(\cos\theta,\sin\theta)$ 的方向导数即为 $\sin k\theta$.
% 
% (2) 显然 $\lim\limits_{(x,y)\to(0,0)}f(x,y)=0$, 所以 $f(x,y)$ 在 $(0,0)$ 处连续.
% 
% (3) 当 $k\geqslant2$ 时, 由于 $\sin k\theta$ 至少在 $\theta=\dfrac{\pi}{2k}, \dfrac{5\pi}{2k}\in[0,2\pi)$ 处取最大值 $1$, 我们断言 $f(x,y)$ 在 $(0,0)$ 处不可微。如若断言不真, 假设 $f(x,y)$ 在 $(0,0)$ 处可微, 则 $\operatorname{grad}f(0,0)$ 是非零向量, 因此它的方向导数只能在一个方向即梯度的方向达到最大值。此矛盾证明了断言.
% \end{solve}
% 
% 
% \begin{remark}\label{remark:ma-14-source-26}
% 当 $k=2$ 且 $r\ne0$ 时,
% \[
%   r\sin2\theta=\frac{2r^2\sin\theta\cos\theta}{r}
%   =\frac{2xy}{\sqrt{x^2+y^2}}.
% \]
% 由此我们可将 $f(x,y)$ 写成如下形式:
% \[
%   f(x,y)=
%   \begin{cases}
%     \dfrac{2xy}{\sqrt{x^2+y^2}}, &x^2+y^2\ne0,\\
%     0, &x^2+y^2=0.
%   \end{cases}
% \]
% 同理, 当 $k=3$ 且 $r\ne0$ 时,
% \[
%   f(x,y)=
%   \begin{cases}
%     \dfrac{3x^2y-y^3}{x^2+y^2}, &x^2+y^2\ne0,\\
%     0, &x^2+y^2=0.
%   \end{cases}
% \]
% 请读者自己对上述函数在 $(x,y)$ 坐标下重新解例~\ref{example:ma-14-1-7}.
% \end{remark}
% 
% 
% \subsection{向量函数的导数与全微分}
% \label{subsec:ma-14-1-5}
% 
% 在本小节中, 我们将定义向量函数的导数与全微分。由于涉及向量函数, 今后我们有时用行向量, 但有时也用列向量来记同一个向量, 不过从上下文中, 读者可以清楚地知道它何时是行向量, 何时是列向量.
% 
% \begin{definition}\label{definition:ma-14-1-5}
% 设向量函数 $u=f(x)=(f_1(x),f_2(x),\cdots,f_m(x))^\mathrm T$ 在区域 $D\subset\mathbb R^n$ 上有定义, $x_0\in D$, $\Delta x=(\Delta x_1,\cdots,\Delta x_n)^\mathrm T$ 为 $x$ 在 $x_0$ 处的全增量。如果存在 $m\times n$ 矩阵
% \[
%   A=
%   \begin{pmatrix}
%     A_{11}&A_{12}&\cdots&A_{1n}\\
%     A_{21}&A_{22}&\cdots&A_{2n}\\
%     \cdots&\cdots&\cdots&\cdots\\
%     A_{m1}&A_{m2}&\cdots&A_{mn}
%   \end{pmatrix}
% \]
% 使得当 $|\Delta x|\to0$ 时, 下式成立:
% \begin{equation}
% \begin{aligned}
%   \Delta f(x_0)
%   &=(\Delta f_1(x_0),\cdots,\Delta f_m(x_0))^\mathrm T\\
%   &=A(\Delta x_1,\cdots,\Delta x_n)^\mathrm T
%   +(\alpha_1(|\Delta x|),\cdots,\alpha_m(|\Delta x|))^\mathrm T,
% \end{aligned}
% \label{eq:ma-14-1-3}
% \end{equation}
% 其中 $A$ 中的元素仅依赖于 $x_0$ 而不依赖于 $\Delta x$, 对于 $\forall j(1\leqslant j\leqslant m)$, $\alpha_j(|\Delta x|)$ 依赖于 $\Delta x$, 并且满足 $\lim\limits_{|\Delta x|\to0}\dfrac{\alpha_j(|\Delta x|)}{|\Delta x|}=0$, 则称 $f(x)$ 在 $x_0$ 处可微或可导, 矩阵 $A$ 称为 $f(x)$ 在 $x_0$ 处的 Fréchet 导数 (简称导数), 记做 $f'(x_0)$ 或 $Df(x_0)$; $A\Delta x$ 称为 $f(x)$ 在 $x_0$ 处的全微分 (简称微分), 记做 $\mathrm df(x_0)$, 即
% \[
%   \mathrm df(x_0)=A\Delta x=f'(x_0)\Delta x=Df(x_0)\Delta x.
% \]
% 在上述定义中, 若规定 $\mathrm dx=\Delta x$, 则有 $\mathrm df(x_0)=f'(x_0)\mathrm dx$.
% \end{definition}
% 
% 式 \eqref{eq:ma-14-1-3} 用矩阵可以简写成
% \[
%   \Delta f(x_0)=f(x_0+\Delta x)-f(x_0)=A\Delta x+\alpha(|\Delta x|),
% \]
% 其中 $\alpha(|\Delta x|)=(\alpha_1(|\Delta x|),\cdots,\alpha_m(|\Delta x|))^\mathrm T$ 满足
% \[
%   \lim_{|\Delta x|\to0}\frac{\alpha(|\Delta x|)}{|\Delta x|}=0.
% \]
% 
% 下面我们考虑如何判断一个向量函数的可微性, 以及当它可微时, 如何求它的导数。下面的定理对此给出了回答.
% 
% \begin{theorem}\label{theorem:ma-14-1-4}
% 设 $D$ 是 $\mathbb R^n$ 中的区域, $x_0\in D$, 向量函数 $f(x)=(f_1(x),\cdots,f_m(x))^\mathrm T$ 在 $D$ 上有定义, 则 $f(x)$ 在 $x_0$ 处可微的充分必要条件是对于 $\forall j(1\leqslant j\leqslant m)$, $f_j(x)$ 在 $x_0$ 处可微。记
% \[
%   A=
%   \begin{pmatrix}
%     \dfrac{\partial f_1(x_0)}{\partial x_1}&\dfrac{\partial f_1(x_0)}{\partial x_2}&\cdots&\dfrac{\partial f_1(x_0)}{\partial x_n}\\
%     \dfrac{\partial f_2(x_0)}{\partial x_1}&\dfrac{\partial f_2(x_0)}{\partial x_2}&\cdots&\dfrac{\partial f_2(x_0)}{\partial x_n}\\
%     \cdots&\cdots&\cdots&\cdots\\
%     \dfrac{\partial f_m(x_0)}{\partial x_1}&\dfrac{\partial f_m(x_0)}{\partial x_2}&\cdots&\dfrac{\partial f_m(x_0)}{\partial x_n}
%   \end{pmatrix},
% \]
% 则当 $f(x)$ 在 $x_0$ 处可微时, 有
% \[
%   \mathrm df(x_0)=A\,\mathrm dx\quad\text{或}\quad f'(x_0)=A.
% \]
% \end{theorem}
% 
% 
% \begin{proof}
% \keyterm{必要性}\quad 设 $f(x)$ 在 $x_0$ 处可微, 从而存在 $m\times n$ 矩阵 $A=(A_{ji})$ 使得当 $|\Delta x|\to0$ 时, 有
% \[
% \begin{aligned}
%   \Delta f(x_0)
%   &=(\Delta f_1(x_0),\cdots,\Delta f_m(x_0))^\mathrm T\\
%   &=\left(\sum_{i=1}^n A_{1i}\Delta x_i,\cdots,\sum_{i=1}^n A_{mi}\Delta x_i\right)^\mathrm T
%   +(\alpha_1(|\Delta x|),\cdots,\alpha_m(|\Delta x|))^\mathrm T,
% \end{aligned}
% \]
% 其中 $\alpha_j(|\Delta x|)$ 依赖于 $\Delta x$, 且 $\lim\limits_{|\Delta x|\to0}\dfrac{\alpha_j(|\Delta x|)}{|\Delta x|}=0(j=1,\cdots,m)$。比较上式两边向量的分量, 对 $j=1,\cdots,m$, 当 $|\Delta x|\to0$ 时, 有
% \[
%   \Delta f_j(x_0)=\sum_{i=1}^n A_{ji}\Delta x_i+\alpha_j(|\Delta x|).
% \]
% 由多元函数可微的定义知, $f_j(x)$ 在 $x_0$ 处可微, 并且
% \[
%   A_{ji}=\frac{\partial f_j(x_0)}{\partial x_i}
%   \quad (i=1,\cdots,n;\ j=1,\cdots,m).
% \]
% 
% \keyterm{充分性}\quad 设对于 $\forall j(1\leqslant j\leqslant m)$, $f_j(x)$ 在 $x_0$ 处可微, 则有
% \[
%   \Delta f_j(x_0)=\sum_{i=1}^n\frac{\partial f_j(x_0)}{\partial x_i}\Delta x_i+\alpha_j(|\Delta x|),
% \]
% 其中 $\alpha_j(|\Delta x|)$ 依赖于 $\Delta x$, 且 $\lim\limits_{|\Delta x|\to0}\dfrac{\alpha_j(|\Delta x|)}{|\Delta x|}=0$。因此
% \[
% \begin{aligned}
%   \Delta f(x_0)
%   &=(\Delta f_1(x_0),\cdots,\Delta f_m(x_0))^\mathrm T\\
%   &=
%   \begin{pmatrix}
%     \dfrac{\partial f_1(x_0)}{\partial x_1}&\cdots&\dfrac{\partial f_1(x_0)}{\partial x_n}\\
%     \dfrac{\partial f_2(x_0)}{\partial x_1}&\cdots&\dfrac{\partial f_2(x_0)}{\partial x_n}\\
%     \cdots&\cdots&\cdots\\
%     \dfrac{\partial f_m(x_0)}{\partial x_1}&\cdots&\dfrac{\partial f_m(x_0)}{\partial x_n}
%   \end{pmatrix}
%   \begin{pmatrix}
%     \Delta x_1\\
%     \Delta x_2\\
%     \vdots\\
%     \Delta x_n
%   \end{pmatrix}
%   +
%   \begin{pmatrix}
%     \alpha_1(|\Delta x|)\\
%     \alpha_2(|\Delta x|)\\
%     \vdots\\
%     \alpha_m(|\Delta x|)
%   \end{pmatrix}\\
%   &=A\Delta x+\alpha(|\Delta x|),
% \end{aligned}
% \]
% 其中 $\alpha(|\Delta x|)$ 满足 $\lim\limits_{|\Delta x|\to0}\dfrac{\alpha(|\Delta x|)}{|\Delta x|}=0$。由定义知 $f(x)$ 在 $x_0$ 处可微, 且 $f'(x_0)=A$.
% \end{proof}
% 
% 
% 当向量函数 $f(x)=(f_1,\cdots,f_m)^\mathrm T$ 中的每个分量函数在 $x_0$ 处均可微时, 矩阵 $A=f'(x_0)$ 称为 $f(x)$ 在 $x_0$ 处的雅可比 (Jacobi) 矩阵, 记为 $J_f(x_0)$。特别地, 当 $f(x)$ 是 $n$ 维向量函数时, $A$ 是 $n\times n$ 矩阵, 此时 $J_f(x_0)$ 的行列式称为 $f(x)$ 在 $x_0$ 处的雅可比行列式, 记为
% \[
%   |f'(x_0)|\quad\text{或}\quad
%   \left.\frac{\partial(f_1,f_2,\cdots,f_n)}{\partial(x_1,x_2,\cdots,x_n)}\right|_{x_0}.
% \]
% 显然, 对于一个向量函数 $f(x)$, 若其分量函数在 $x_0$ 处具有各个偏导数时, 我们就可以形式地定义 $J_f(x_0)$, 但当 $f(x)$ 在 $x_0$ 处不可微, 仅仅存在所有的偏导数时, $J_f(x_0)$ 不能刻画 $f(x)$ 在 $x_0$ 附近的变化情况.
% 
% 当 $f_j(x)(1\leqslant j\leqslant m)$ 的各个偏导数都在 $x_0$ 处连续时, $f_j(x)$ 在 $x_0$ 处可微, 因此 $f(x)$ 在 $x_0$ 处可微。另外, 若对于 $\forall j(1\leqslant j\leqslant m)$, $f_j(x)$ 的各个偏导数在区域 $D$ 上连续, 我们称 $f(x)$ 在 $D$ 上是 $C^1$ 的, 记为 $f(x)\in C^1(D)$。特别地, 我们称 $\mathbb R^n$ 中区域 $D$ 到 $\Omega$ 的变换 $y=f(x)$ 是 $C^1$ 的, 如果 $f(x)\in C^1(D)$, 并且 $f^{-1}(y)\in C^1(\Omega)$.
% 
% \begin{remark}\label{remark:ma-14-source-30}
% 利用向量函数导数的记号, 对一个多元函数 $f(x)$, 当它可微时, 我们有
% \[
%   f'(x)=\left(\frac{\partial f}{\partial x_1},\cdots,\frac{\partial f}{\partial x_n}\right)=\operatorname{grad}f(x).
% \]
% \end{remark}
% 
% 
% \begin{example}\label{example:ma-14-1-8}
% 设向量 $w=(r,\varphi,\theta)$, 向量函数
% \[
%   f(w)=(x(w),y(w),z(w))^\mathrm T
%   =(r\sin\varphi\cos\theta,r\sin\varphi\sin\theta,r\cos\varphi)^\mathrm T,
% \]
% 其中 $D=\{(r,\varphi,\theta):0<r<+\infty,0<\varphi<\pi,0<\theta<2\pi\}$; 试求 $f'(w)$ 以及 $f(w)$ 在 $w$ 处的雅可比行列式.
% \end{example}
% 
% 
% \begin{solve}
% 由于 $x=r\sin\varphi\cos\theta$, $y=r\sin\varphi\sin\theta$, $z=r\cos\varphi$ 均在 $D$ 上具有连续偏导数, 因此 $f'(w)$ 存在。由计算得
% \[
% \begin{aligned}
%   f'(w)
%   &=
%   \begin{pmatrix}
%     \dfrac{\partial x}{\partial r}&\dfrac{\partial x}{\partial\varphi}&\dfrac{\partial x}{\partial\theta}\\
%     \dfrac{\partial y}{\partial r}&\dfrac{\partial y}{\partial\varphi}&\dfrac{\partial y}{\partial\theta}\\
%     \dfrac{\partial z}{\partial r}&\dfrac{\partial z}{\partial\varphi}&\dfrac{\partial z}{\partial\theta}
%   \end{pmatrix}\\
%   &=
%   \begin{pmatrix}
%     \sin\varphi\cos\theta&r\cos\varphi\cos\theta&-r\sin\varphi\sin\theta\\
%     \sin\varphi\sin\theta&r\cos\varphi\sin\theta&r\sin\varphi\cos\theta\\
%     \cos\varphi&-r\sin\varphi&0
%   \end{pmatrix}.
% \end{aligned}
% \]
% 对上述矩阵取行列式, 即得 $f(w)$ 在 $w=(r,\varphi,\theta)$ 处的雅可比行列式
% \[
%   \frac{\partial(x,y,z)}{\partial(r,\varphi,\theta)}=r^2\sin\varphi.
% \qedhere
% \]
% \end{solve}
% 
% 
% \section{多元函数求导法}
% \label{sec:ma-14-2}
% 
% \subsection{导数的四则运算}
% \label{subsec:ma-14-2-1}
% 
% 在本小节中, 我们总假定涉及的多元 (向量) 函数存在各个偏导数, 然后来讨论多元函数求偏导数的一些公式。读者应该注意的是, 一般情况下一个函数的导数是一个向量 (或矩阵).
% 
% \begin{theorem}\label{theorem:ma-14-2-1}
% 设函数 $f(x),g(x)$ 在区域 $D\subset\mathbb R^n$ 上可导, 则对于 $\forall x\in D$, 有
% \[
% \begin{aligned}
%   &(1)\quad (f(x)\pm g(x))'=f'(x)\pm g'(x);\\
%   &(2)\quad (f(x)g(x))'=f(x)g'(x)+g(x)f'(x);\\
%   &(3)\quad \left(\frac{f(x)}{g(x)}\right)'=
%   \frac{g(x)f'(x)-f(x)g'(x)}{g^2(x)}\quad (g(x)\ne0).
% \end{aligned}
% \]
% 如果 $f(x):\mathbb R^n\to\mathbb R$ 是一个 $n$ 元可微函数, $g(x):\mathbb R^n\to\mathbb R^m$ 是一个可微的 $m$ 维列向量函数, 则有
% \[
%   (4)\quad (f(x)g(x))'=f(x)g'(x)+g(x)f'(x).
% \]
% \end{theorem}
% 
% 
% \begin{proof}
% 我们只证 (3)。(1), (2) 和 (4) 的证明留给读者。由定义有
% \[
%   \left(\frac{f(x)}{g(x)}\right)'=
%   \left(
%     \frac{\partial\left(\dfrac{f(x)}{g(x)}\right)}{\partial x_1},
%     \frac{\partial\left(\dfrac{f(x)}{g(x)}\right)}{\partial x_2},
%     \cdots,
%     \frac{\partial\left(\dfrac{f(x)}{g(x)}\right)}{\partial x_n}
%   \right).
% \]
% 由于
% \[
%   \frac{\partial\left(\dfrac{f(x)}{g(x)}\right)}{\partial x_i}
%   =
%   \frac{g(x)\dfrac{\partial f(x)}{\partial x_i}-f(x)\dfrac{\partial g(x)}{\partial x_i}}{g^2(x)}
%   \quad (i=1,2,\cdots,n),
% \]
% 将其代入上式即知 (3) 成立.
% \end{proof}
% 
% 
% \begin{remark}\label{remark:ma-14-source-35}
% 注意到 (4) 中等式左边 $(f(x)g(x))'$ 是一个 $m\times n$ 矩阵, $g(x)$ 是一个 $m\times1$ 矩阵, 而 $f'(x)$ 是一个 $1\times n$ 矩阵, 因此 $g(x)f'(x)$ 是一个 $m\times n$ 矩阵。显然, $f(x)g'(x)$ 为一个 $m\times n$ 矩阵。至于 (4) 的证明, 读者只要将 (4) 中等式左边矩阵中的每个分量与其右边矩阵中相对应的分量比较即可.
% \end{remark}
% 
% 
% \subsection{复合函数的求导法}
% \label{subsec:ma-14-2-2}
% 
% 下面我们主要来考虑复合函数的求导法.
% 
% \begin{theorem}\label{theorem:ma-14-2-2}
% 设函数 $f(u)=f(u_1,u_2,\cdots,u_m)$ 在区域 $\Omega\subset\mathbb R^m$ 上有定义, 并且在 $u_0=(u_1^0,u_2^0,\cdots,u_m^0)^\mathrm T\in\Omega$ 处可微, 再设
% \[
%   u=u(x)=(u_1(x),u_2(x),\cdots,u_m(x))^\mathrm T
% \]
% 在区域 $D\subset\mathbb R^n$ 上有定义, 在 $x_0=(x_1^0,x_2^0,\cdots,x_n^0)\in D$ 处可微, 并且 $u_0=u(x_0)$, 则 $f(u(x))$ 在 $x_0$ 处可微, 并且
% \[
%   \mathrm df(u(x_0))=f'(u(x_0))u'(x_0)\mathrm dx.
% \]
% \end{theorem}
% 
% 
% \begin{proof}
% 由 $u$ 在 $x_0$ 处可微, 我们有
% \[
%   \Delta u(x_0)=u(x_0+\Delta x)-u(x_0)
%   =u'(x_0)\Delta x+\alpha(|\Delta x|),
% \]
% 其中 $\alpha(|\Delta x|)$ 依赖于 $\Delta x$, 且满足
% \[
%   \frac{\alpha(|\Delta x|)}{|\Delta x|}\to0\quad(|\Delta x|\to0).
% \]
% 再由 $f(u)$ 在 $u_0$ 处可微, 有
% \[
%   \Delta f(u_0)=f(u_0+\Delta u)-f(u_0)=f'(u_0)\Delta u+\beta(|\Delta u|),
% \]
% 其中 $\beta(|\Delta u|)$ 依赖于 $\Delta u$, 且满足
% \[
%   \frac{\beta(|\Delta u|)}{|\Delta u|}\to0\quad(|\Delta u|\to0).
% \]
% 我们规定 $\beta(0)=0$。在这样的规定下, $\beta(|\Delta u|)$ 在 $\Delta u=0$ 时连续.
% 
% 我们有
% \[
% \begin{aligned}
%   \Delta f(u(x_0))
%   &=f(u(x_0+\Delta x))-f(u(x_0))\\
%   &=f'(u(x_0))(u'(x_0)\Delta x+\alpha(|\Delta x|))
%   +\beta(|u'(x_0)\Delta x+\alpha(|\Delta x|)|)\\
%   &=f'(u(x_0))u'(x_0)\Delta x+\gamma(|\Delta x|),
% \end{aligned}
% \]
% 其中
% \[
% \begin{aligned}
%   \gamma(|\Delta x|)
%   &=f'(u(x_0))\alpha(|\Delta x|)
%   +\beta(|\Delta u(x_0)|)\\
%   &=f'(u(x_0))\alpha(|\Delta x|)
%   +\beta(|u'(x_0)\Delta x+\alpha(|\Delta x|)|).
% \end{aligned}
% \]
% 下面我们证明
% \[
%   \frac{\gamma(|\Delta x|)}{|\Delta x|}\to0\quad(|\Delta x|\to0).
% \]
% 首先, 显然有
% \begin{equation}
%   \lim_{|\Delta x|\to0}
%   \frac{|f'(u(x_0))\alpha(|\Delta x|)|}{|\Delta x|}
%   \leqslant
%   \lim_{|\Delta x|\to0}
%   \frac{|f'(u(x_0))|\cdot|\alpha(|\Delta x|)|}{|\Delta x|}=0.
% \label{eq:ma-14-2-1}
% \end{equation}
% 其次, 注意到当 $|\Delta x|$ 很小时, 有
% \[
% \begin{aligned}
%   \frac{|\Delta u(x_0)|}{|\Delta x|}
%   &=\frac{|u'(x_0)\Delta x+\alpha(|\Delta x|)|}{|\Delta x|}\\
%   &\leqslant \frac1{|\Delta x|}\bigl(\|u'(x_0)\|\,|\Delta x|+|\alpha(|\Delta x|)|\bigr)
%   \leqslant \|u'(x_0)\|+1,
% \end{aligned}
% \]
% 其中 $\|u'(x_0)\|$ 是矩阵 $u'(x_0)$ 的范数。因此当 $|\Delta x|\to0$ 时, 必有 $|\Delta u(x_0)|\to0$。我们有
% \[
%   \frac{|\beta(|\Delta u(x_0)|)|}{|\Delta x|}
%   =
%   \begin{cases}
%     \dfrac{|\beta(|\Delta u(x_0)|)|}{|\Delta u(x_0)|}
%     \cdot\dfrac{|\Delta u(x_0)|}{|\Delta x|}, &|\Delta u(x_0)|\ne0,\\[1.2em]
%     0, &|\Delta u(x_0)|=0,
%   \end{cases}
% \]
% 从而有
% \begin{equation}
%   \frac{|\beta(|\Delta u(x_0)|)|}{|\Delta x|}\to0
%   \quad(|\Delta x|\to0).
% \label{eq:ma-14-2-2}
% \end{equation}
% 结合式 \eqref{eq:ma-14-2-1}, \eqref{eq:ma-14-2-2} 得
% \[
%   \frac{\gamma(|\Delta x|)}{|\Delta x|}\to0\quad(|\Delta x|\to0).
% \]
% 我们最后有
% \[
%   \Delta f(u(x_0))=f'(u(x_0))u'(x_0)\Delta x+\gamma(|\Delta x|),
% \]
% 其中 $\gamma(|\Delta x|)$ 依赖于 $\Delta x$ 且 $\dfrac{\gamma(|\Delta x|)}{|\Delta x|}\to0(|\Delta x|\to0)$。由微分的定义知, $f(u(x))$ 在 $x_0$ 处可微, 并且
% \[
%   \mathrm df(u(x_0))=f'(u(x_0))u'(x_0)\mathrm dx.
% \qedhere
% \]
% \end{proof}
% 
% 
% 从定理~\ref{theorem:ma-14-2-2} 我们可以推出以下结论.
% 
% \begin{corollary}\label{corollary:ma-local-14.2-1}
% 设向量函数 $f(u)=(f_1(u),f_2(u),\cdots,f_l(u))^\mathrm T$ 在区域 $\Omega\subset\mathbb R^m$ 上可微, $u(x)=(u_1(x),u_2(x),\cdots,u_m(x))^\mathrm T$ 在区域 $D\subset\mathbb R^n$ 上可微, 且 $u(D)\subset\Omega$, 则 $g=f(u(x))$ 在 $D$ 上可微, 并且
% \[
%   \mathrm df(u(x))=f'(u(x))u'(x)\mathrm dx.
% \]
% 顺便我们有
% \[
%   [f(u(x))]'=f'(u(x))u'(x).
% \]
% \end{corollary}
% 
% 
% \begin{proof}
% 由定理~\ref{theorem:ma-14-2-2} 知 $\mathrm df_j(u(x))=f'_j(u(x))u'(x)\mathrm dx\ (j=1,2,\cdots,l)$, 从而
% \[
%   \mathrm df(u(x))=
%   \begin{pmatrix}
%     f'_1(u(x))u'(x)\mathrm dx\\
%     f'_2(u(x))u'(x)\mathrm dx\\
%     \vdots\\
%     f'_l(u(x))u'(x)\mathrm dx
%   \end{pmatrix}
%   =f'(u(x))u'(x)\mathrm dx.
% \qedhere
% \]
% \end{proof}
% 
% 
% \begin{remark}\label{remark:ma-14-source-40}
% 读者应该注意到 $f'(u(x))u'(x)$ 是一个 $l\times n$ 矩阵.
% \end{remark}
% 
% 
% \begin{corollary}\label{corollary:ma-local-14.2-2}
% 设 $D$ 与 $\Omega$ 为 $\mathbb R^n$ 中的区域, $y=f(x)$ 是 $D$ 到 $\Omega$ 的 $C^1$ 变换, 则对于 $\forall x\in D$, 有
% \begin{equation}
%   (f^{-1})'(y)\cdot f'(x)=E,
% \label{eq:ma-14-2-3}
% \end{equation}
% 其中 $y=f(x)$; 对于 $\forall y\in\Omega$, 有
% \begin{equation}
%   f'(x)\cdot(f^{-1})'(y)=E,
% \label{eq:ma-14-2-4}
% \end{equation}
% 其中 $x=f^{-1}(y)$。特别地, 当 $y=f(x)$ 时, 有
% \begin{equation}
%   (f^{-1})'(y)=[f'(x)]^{-1},
% \label{eq:ma-14-2-5}
% \end{equation}
% 其中 $E$ 是 $n\times n$ 单位矩阵, $[f'(x)]^{-1}$ 为 $f'(x)$ 的逆矩阵.
% \end{corollary}
% 
% 
% \begin{proof}
% 由 $C^1$ 变换的定义知 $f$ 与 $f^{-1}$ 都是可微向量函数, 因此 $(f^{-1}\circ f)(x)=x$ 两边对 $x$ 求导数即得式 \eqref{eq:ma-14-2-3}, 而 $(f\circ f^{-1})(y)=y$ 两边对 $y$ 求导数即得式 \eqref{eq:ma-14-2-4}。在式 \eqref{eq:ma-14-2-3} 的两边同时右乘 $[f'(x)]^{-1}$ 即得式 \eqref{eq:ma-14-2-5}.
% \end{proof}
% 
% 
% \begin{remark}\label{remark:ma-14-source-43}
% 在推论~\ref{corollary:ma-local-14.2-2} 中, 如果我们在式 \eqref{eq:ma-14-2-4} 两边取行列式, 则有
% \begin{equation}
%   \frac{\partial(y_1,\cdots,y_n)}{\partial(x_1,\cdots,x_n)}
%   \cdot
%   \frac{\partial(x_1,\cdots,x_n)}{\partial(y_1,\cdots,y_n)}
%   =1.
% \label{eq:ma-14-2-6}
% \end{equation}
% 下面的推论~\ref{corollary:ma-local-14.2-3} 是多元复合函数求偏导数的基础.
% \end{remark}
% 
% 
% \begin{corollary}\label{corollary:ma-local-14.2-3}
% 设函数 $f(u)=f(u_1,u_2,\cdots,u_m)$ 在区域 $\Omega\subset\mathbb R^m$ 上有定义, 并且在 $u_0=(u_1^0,u_2^0,\cdots,u_m^0)^\mathrm T\in\Omega$ 处可微, 又设向量函数
% \[
%   u=u(x)=(u_1(x),u_2(x),\cdots,u_m(x))^\mathrm T
% \]
% 在区域 $D\subset\mathbb R^n$ 上有定义, 在 $x_0=(x_1^0,x_2^0,\cdots,x_n^0)\in D$ 处可微, 并且 $u_0=u(x_0)$, 则对于 $\forall i(1\leqslant i\leqslant n)$, $f(u(x))$ 在 $x_0$ 处关于 $x_i$ 可偏导, 并且
% \[
%   \frac{\partial f(u(x_0))}{\partial x_i}
%   =\sum_{j=1}^m\left(
%   \frac{\partial f(u_0)}{\partial u_j}\cdot
%   \frac{\partial u_j(x_0)}{\partial x_i}\right).
% \]
% \end{corollary}
% 
% 
% \begin{proof}
% 由定理~\ref{theorem:ma-14-1-4} 的注有
% \[
%   \left.(f(u(x)))'\right|_{x=x_0}
%   =\left(
%   \frac{\partial f(u(x_0))}{\partial x_1},
%   \frac{\partial f(u(x_0))}{\partial x_2},
%   \cdots,
%   \frac{\partial f(u(x_0))}{\partial x_n}
%   \right).
% \]
% 用矩阵形式表示 $f'(u(x_0))u'(x_0)$, 有
% \[
% \begin{aligned}
%   f'(u(x_0))u'(x_0)
%   &=
%   \left(
%   \frac{\partial f(u(x_0))}{\partial u_1},
%   \frac{\partial f(u(x_0))}{\partial u_2},
%   \cdots,
%   \frac{\partial f(u(x_0))}{\partial u_m}
%   \right)\\
%   &\quad\cdot
%   \begin{pmatrix}
%     \dfrac{\partial u_1(x_0)}{\partial x_1}&\dfrac{\partial u_1(x_0)}{\partial x_2}&\cdots&\dfrac{\partial u_1(x_0)}{\partial x_n}\\
%     \dfrac{\partial u_2(x_0)}{\partial x_1}&\dfrac{\partial u_2(x_0)}{\partial x_2}&\cdots&\dfrac{\partial u_2(x_0)}{\partial x_n}\\
%     \vdots&\vdots&\ddots&\vdots\\
%     \dfrac{\partial u_m(x_0)}{\partial x_1}&\dfrac{\partial u_m(x_0)}{\partial x_2}&\cdots&\dfrac{\partial u_m(x_0)}{\partial x_n}
%   \end{pmatrix}\\
%   &=
%   \left(
%   \sum_{j=1}^m\frac{\partial f(u_0)}{\partial u_j}\frac{\partial u_j(x_0)}{\partial x_1},
%   \sum_{j=1}^m\frac{\partial f(u_0)}{\partial u_j}\frac{\partial u_j(x_0)}{\partial x_2},
%   \cdots,
%   \sum_{j=1}^m\frac{\partial f(u_0)}{\partial u_j}\frac{\partial u_j(x_0)}{\partial x_n}
%   \right).
% \end{aligned}
% \]
% 由定理~\ref{theorem:ma-14-2-2} 有 $\left.(f(u(x)))'\right|_{x=x_0}=f'(u(x_0))u'(x_0)$, 再比较它们的分量即得推论~\ref{corollary:ma-local-14.2-3}.
% \end{proof}
% 
% 
% \begin{remark}\label{remark:ma-14-source-46}
% 利用定理~\ref{theorem:ma-14-2-2} 类似的证明方法可以证明, 若将推论~\ref{corollary:ma-local-14.2-3} 中的条件 ``$u(x)$ 在 $x_0$ 处可微'' 减弱成 ``$u(x)$ 在 $x_0$ 处存在各个偏导数'', 则推论~\ref{corollary:ma-local-14.2-3} 的结论仍然成立。请读者自己对此加以证明 (见本章习题)。但是, 在推论~\ref{corollary:ma-local-14.2-3} 中, 条件 ``$f(u)$ 在 $u_0$ 处可微'' 不能减弱成 ``$f(u)$ 在 $u_0$ 处存在各个偏导数''。很容易找到例子来说明这一点, 例如取
% \[
%   f(u,v)=
%   \begin{cases}
%     1, &uv\ne0,\\
%     0, &uv=0,
%   \end{cases}
%   \quad (u,v)\in\mathbb R^2,
% \]
% 而令
% \[
%   \begin{cases}
%     u=t,\\
%     v=t
%   \end{cases}
%   \quad(-\infty<t<+\infty),
% \]
% 则
% \[
%   g(t)=f(t,t)=
%   \begin{cases}
%     1, &t\ne0,\\
%     0, &t=0
%   \end{cases}
% \]
% 在 $t=0$ 处不连续, 从而不可导.
% \end{remark}
% 
% 如同一元函数, 我们将推论~\ref{corollary:ma-local-14.2-3} 中给出的复合函数求导公式称为链锁法则。为了今后应用方便, 下面我们列出常见的二元复合函数的求导公式。设 $z=f(u,v), u=u(x,y), v=v(x,y)$ 都为可微函数, 则
% \[
%   \frac{\partial z}{\partial x}
%   =\frac{\partial f}{\partial u}\cdot\frac{\partial u}{\partial x}
%   +\frac{\partial f}{\partial v}\cdot\frac{\partial v}{\partial x},
% \]
% \[
%   \frac{\partial z}{\partial y}
%   =\frac{\partial f}{\partial u}\cdot\frac{\partial u}{\partial y}
%   +\frac{\partial f}{\partial v}\cdot\frac{\partial v}{\partial y}.
% \]
% 有时, 若给出的函数具有形式 $z=f(u(x,y),v(x,y))$, 我们可以用以下记号:
% \[
%   \frac{\partial z}{\partial x}=f'_1\frac{\partial u}{\partial x}+f'_2\frac{\partial v}{\partial x},
%   \qquad
%   \frac{\partial z}{\partial y}=f'_1\frac{\partial u}{\partial y}+f'_2\frac{\partial v}{\partial y},
% \]
% 其中 $f'_i(i=1,2)$ 指的是 $f$ 对其第 $i$ 个变量求偏导数。对于具有多个中间变量的复合函数, 我们可以引进类似的记号, 这对抽象形式给出的函数的求导运算提供了简洁的记号.
% 
% \begin{example}\label{example:ma-14-2-1}
% 设函数 $u=e^{x^2+y^2+z^2}; z=x^2\sin y$, 求 $\dfrac{\partial u}{\partial x}, \dfrac{\partial u}{\partial y}$.
% \end{example}
% 
% 
% \begin{solve}
% \[
% \begin{aligned}
%   \frac{\partial u}{\partial x}
%   &=\frac{\partial e^{x^2+y^2+z^2}}{\partial x}\cdot1
%   +\frac{\partial e^{x^2+y^2+z^2}}{\partial z}\cdot\frac{\partial z}{\partial x}\\
%   &=2xe^{x^2+y^2+z^2}+e^{x^2+y^2+z^2}(2z)(2x\sin y)\\
%   &=2x(1+2x^2\sin^2 y)e^{x^2+y^2+z^2},
% \end{aligned}
% \]
% \[
% \begin{aligned}
%   \frac{\partial u}{\partial y}
%   &=\frac{\partial e^{x^2+y^2+z^2}}{\partial y}\cdot1
%   +\frac{\partial e^{x^2+y^2+z^2}}{\partial z}\cdot\frac{\partial z}{\partial y}\\
%   &=2ye^{x^2+y^2+z^2}+e^{x^2+y^2+z^2}(2z)(x^2\cos y)\\
%   &=(2y+x^4\sin2y)\cdot e^{x^2+y^2+z^2}.
% \end{aligned}
% \qedhere
% \]
% \end{solve}
% 
% 
% \begin{example}\label{example:ma-14-2-2}
% 设函数 $u=f\left(xy,\dfrac yx,yz\right)$, 并设 $f$ 是可微函数, 试求 $\dfrac{\partial u}{\partial x}, \dfrac{\partial u}{\partial y}$ 和 $\dfrac{\partial u}{\partial z}$.
% \end{example}
% 
% 
% \begin{solve}
% 对于这种抽象形式给出的函数, 我们采用上述约定的记号来求偏导数:
% \[
%   \frac{\partial u}{\partial x}
%   =f'_1\frac{\partial(xy)}{\partial x}+f'_2\frac{\partial(y/x)}{\partial x}+f'_3\cdot0
%   =yf'_1-\frac{y}{x^2}f'_2,
% \]
% \[
%   \frac{\partial u}{\partial y}
%   =f'_1\frac{\partial(xy)}{\partial y}
%   +f'_2\frac{\partial(y/x)}{\partial y}
%   +f'_3\frac{\partial(yz)}{\partial y}
%   =xf'_1+\frac1x f'_2+zf'_3,
% \]
% \[
%   \frac{\partial u}{\partial z}
%   =f'_1\cdot0+f'_2\cdot0+f'_3\frac{\partial(yz)}{\partial z}
%   =yf'_3.
% \qedhere
% \]
% \end{solve}
% 
% 
% \begin{example}\label{example:ma-14-2-3}
% 记向量 $v=(r,\theta_1,\cdots,\theta_{n-1})$, 并设向量函数
% \[
%   f(v)=
%   \begin{pmatrix}
%     x_1(v)\\
%     x_2(v)\\
%     \vdots\\
%     x_{n-1}(v)\\
%     x_n(v)
%   \end{pmatrix}
%   =
%   \begin{pmatrix}
%     r\cos\theta_1\\
%     r\sin\theta_1\cos\theta_2\\
%     \vdots\\
%     r\sin\theta_1\sin\theta_2\cdots\cos\theta_{n-1}\\
%     r\sin\theta_1\sin\theta_2\cdots\sin\theta_{n-1}
%   \end{pmatrix},
% \]
% 试求 $f'(v)$ 及其雅可比行列式.
% \end{example}
% 
% 
% \begin{solve}
% 由定义有
% 
% \[
%   f'(v)=
%   \begin{pmatrix}
%     \cos\theta_1&-r\sin\theta_1&\cdots&0\\
%     \sin\theta_1\cos\theta_2&r\cos\theta_1\cos\theta_2&\cdots&0\\
%     \vdots&\vdots&\ddots&\vdots\\
%     \sin\theta_1\cdots\cos\theta_{n-1}&r\cos\theta_1\cdots\cos\theta_{n-1}&\cdots&-r\sin\theta_1\cdots\sin\theta_{n-1}\\
%     \sin\theta_1\cdots\sin\theta_{n-1}&r\cos\theta_1\cdots\sin\theta_{n-1}&\cdots&r\sin\theta_1\cdots\cos\theta_{n-1}
%   \end{pmatrix}.
% \]
% 下面用归纳法来求 $f'(v)$ 的行列式 $|f'(v)|$。当 $n=2$ 时, 有 $|f'(v)|=r$; 当 $n=3$ 时, 有 $|f'(v)|=r^2\sin\theta_1$。我们猜测对于 $n\geqslant3$, 有
% \[
%   |f'(v)|=r^{n-1}\sin^{n-2}\theta_1\sin^{n-3}\theta_2\cdots\sin\theta_{n-2}.
% \]
% 事实上, 已知 $n=3$ 时上式成立。假设 $n=k-1$ 时上式成立, 现在我们来证明当 $n=k$ 时上式成立。为此将变换 $f$ 拆成下列两个变换的复合:
% \[
%   f_1:\begin{cases}
%     x_1=y_1,\\
%     x_2=y_2\cos y_3,\\
%     x_3=y_2\sin y_3\cos y_4,\\
%     \cdots\\
%     x_{n-1}=y_2\sin y_3\sin y_4\cdots\cos y_n,\\
%     x_n=y_2\sin y_3\sin y_4\cdots\sin y_n,
%   \end{cases}
%   \qquad
%   f_2:\begin{cases}
%     y_1=r\cos\theta_1,\\
%     y_2=r\sin\theta_1,\\
%     y_3=\theta_2,\\
%     \cdots\\
%     y_{n-1}=\theta_{n-2},\\
%     y_n=\theta_{n-1}.
%   \end{cases}
% \]
% 则有 $f=f_1\circ f_2$。当 $n=k$ 时, 由归纳法假设得
% \[
% \begin{aligned}
%   |f'(v)|
%   &=
%   \frac{\partial(x_1,x_2,\cdots,x_k)}
%   {\partial(r,\theta_1,\cdots,\theta_{k-1})}\\
%   &=
%   \left.
%   \frac{\partial(x_1,x_2,\cdots,x_k)}
%   {\partial(y_1,y_2,\cdots,y_k)}
%   \right|_{f_2(v)}
%   \frac{\partial(y_1,y_2,\cdots,y_k)}
%   {\partial(r,\theta_1,\cdots,\theta_{k-1})}\\
%   &=
%   \left.(y_2^{k-2}\sin^{k-3}y_3\sin^{k-4}y_4\cdots\sin y_{k-1})\right|_{f_2(v)}
%   \cdot r\\
%   &=r^{k-1}\sin^{k-2}\theta_1\sin^{k-3}\theta_2\cdots\sin\theta_{k-2}.
% \end{aligned}
% \qedhere
% \]
% \end{solve}
% 
% 
% \begin{example}\label{example:ma-14-2-4}
% 设 $f(x)=(f_1(x),f_2(x),\cdots,f_m(x))$ 和 $g(x)=(g_1(x),g_2(x),\cdots,g_m(x))$ 为 $D\subset\mathbb R^n$ 到 $\mathbb R^m$ 的可微向量函数, 证明:
% \[
%   [f(x)(g(x))^\mathrm T]'=f(x)[(g(x))^\mathrm T]'+g(x)[(f(x))^\mathrm T]'.
% \]
% \end{example}
% 
% 
% \begin{proof}
% 令
% \[
%   F(x)=
%   \begin{pmatrix}
%     (f(x))^\mathrm T\\
%     (g(x))^\mathrm T
%   \end{pmatrix}
%   =(f_1(x),\cdots,f_m(x),g_1(x),\cdots,g_m(x))^\mathrm T,
% \]
% 则 $F(x)$ 是 $D\to\mathbb R^{2m}$ 的可微函数。记 $y=(y_1,y_2,\cdots,y_{2m})^\mathrm T$, 并令
% \[
%   G(y)=\sum_{i=1}^m y_i y_{m+i},
% \]
% 则 $G(y)$ 是 $\mathbb R^{2m}\to\mathbb R$ 的可微函数。我们可以验证
% \[
%   f(x)(g(x))^\mathrm T=G(F(x)).
% \]
% 由链锁法则则有
% \[
% \begin{aligned}
%   [f(x)(g(x))^\mathrm T]'
%   &=[G(F(x))]'\\
%   &=(y_{m+1},\cdots,y_{2m},y_1,\cdots,y_m)\big|_{y=F(x)}
%   \begin{pmatrix}
%     [(f(x))^\mathrm T]'\\
%     [(g(x))^\mathrm T]'
%   \end{pmatrix}\\
%   &=f(x)[(g(x))^\mathrm T]'+g(x)[(f(x))^\mathrm T]'.
% \end{aligned}
% \qedhere
% \]
% \end{proof}
% 
% 
% \begin{example}\label{example:ma-14-2-5}
% 设 $f(x,y)$ 在 $\mathbb R^2$ 上具有连续偏导数, 试用极坐标来表示
% \[
%   |f'(x,y)|^2=
%   \left(\frac{\partial f}{\partial x}\right)^2+
%   \left(\frac{\partial f}{\partial y}\right)^2
%   =|\operatorname{grad}f(x,y)|^2.
% \]
% \end{example}
% 
% 
% \begin{solve}
% 由
% \[
%   \begin{cases}
%     x=r\cos\theta,\\
%     y=r\sin\theta
%   \end{cases}
% \]
% 得
% \[
%   \frac{\partial f}{\partial x}
%   =\frac{\partial f}{\partial r}\cdot\frac{\partial r}{\partial x}
%   +\frac{\partial f}{\partial\theta}\cdot\frac{\partial\theta}{\partial x},
%   \quad
%   \frac{\partial f}{\partial y}
%   =\frac{\partial f}{\partial r}\cdot\frac{\partial r}{\partial y}
%   +\frac{\partial f}{\partial\theta}\cdot\frac{\partial\theta}{\partial y}.
% \]
% 由 $r=\sqrt{x^2+y^2}$ 及 $\theta=\arctan\dfrac yx$ (此处分象限讨论, 我们以第一象限为例求之) 得
% \[
%   \frac{\partial r}{\partial x}=\frac{x}{\sqrt{x^2+y^2}}=\cos\theta,\quad
%   \frac{\partial r}{\partial y}=\frac{y}{\sqrt{x^2+y^2}}=\sin\theta,
% \]
% \[
%   \frac{\partial\theta}{\partial x}=-\frac{y}{x^2+y^2}=-\frac{\sin\theta}{r},\quad
%   \frac{\partial\theta}{\partial y}=\frac{x}{x^2+y^2}=\frac{\cos\theta}{r}.
% \]
% 因此
% \[
%   \frac{\partial f}{\partial x}
%   =\cos\theta\frac{\partial f}{\partial r}
%   -\frac{\sin\theta}{r}\cdot\frac{\partial f}{\partial\theta},
%   \quad
%   \frac{\partial f}{\partial y}
%   =\sin\theta\frac{\partial f}{\partial r}
%   +\frac{\cos\theta}{r}\cdot\frac{\partial f}{\partial\theta};
% \]
% 从而有
% \[
%   |\operatorname{grad}f(x,y)|^2
%   =
%   \left(\frac{\partial f}{\partial x}\right)^2+
%   \left(\frac{\partial f}{\partial y}\right)^2
%   =
%   \left(\frac{\partial f}{\partial r}\right)^2+
%   \frac1{r^2}\left(\frac{\partial f}{\partial\theta}\right)^2.
% \qedhere
% \]
% \end{solve}
% 
% 
% \begin{example}\label{example:ma-14-2-6}
% 设函数 $u=f(x-\lambda t)$ ($\lambda$ 是常数), 其中 $f$ 是 $\mathbb R$ 中的一个可微函数 (即 $f(y)$ 是可微函数), 证明它满足微分方程
% \begin{equation}
%   \frac{\partial u}{\partial t}+\lambda\frac{\partial u}{\partial x}=0,
% \label{eq:ma-14-2-7}
% \end{equation}
% 并证明方程 \eqref{eq:ma-14-2-7} 的解一定具有形式 $u=g(x-\lambda t)$, 其中 $g$ 是任一个 $\mathbb R$ 中的可微函数.
% \end{example}
% 
% 
% \begin{proof}
% 记 $y=x-\lambda t$, 则
% \[
%   \frac{\partial u}{\partial t}
%   =\frac{\mathrm df}{\mathrm dy}\cdot\frac{\partial y}{\partial t}
%   =f'(y)(-\lambda)=-\lambda f'(y),
%   \quad
%   \frac{\partial u}{\partial x}
%   =\frac{\mathrm df}{\mathrm dy}\cdot\frac{\partial y}{\partial x}
%   =\frac{\partial f}{\partial y}=f'(y),
% \]
% 从而
% \[
%   \frac{\partial u}{\partial t}+\lambda\frac{\partial u}{\partial x}
%   =-\lambda f'(y)+\lambda f'(y)=0.
% \]
% 设可微函数 $u=u(x,t)$ 是方程 \eqref{eq:ma-14-2-7} 的解, 作变量替换
% \[
%   \begin{cases}
%     x=y+\lambda z,\\
%     z=t,
%   \end{cases}
% \]
% 则 $u(x,t)=u(y+\lambda z,z)\triangleq u_1(y,z)$。由
% \[
%   \frac{\partial u}{\partial t}
%   =\frac{\partial u_1}{\partial y}\cdot\frac{\partial y}{\partial t}
%   +\frac{\partial u_1}{\partial z}\cdot\frac{\partial z}{\partial t}
%   =-\lambda\frac{\partial u_1}{\partial y}
%   +\frac{\partial u_1}{\partial z},
% \]
% \[
%   \frac{\partial u}{\partial x}
%   =\frac{\partial u_1}{\partial y}\cdot\frac{\partial y}{\partial x}
%   +\frac{\partial u_1}{\partial z}\cdot\frac{\partial z}{\partial x}
%   =\frac{\partial u_1}{\partial y},
% \]
% 有
% \[
%   \frac{\partial u}{\partial t}+\lambda\frac{\partial u}{\partial x}
%   =\frac{\partial u_1}{\partial z}=0.
% \]
% 这说明 $u_1(y,z)$ 与 $z$ 无关, 即 $u_1=g(y)$, 亦即 $u=g(x-\lambda t)$, 其中 $g$ 为 $\mathbb R$ 中的任一个可微函数.
% \end{proof}
% 
% 
% \subsection{高阶偏导数}
% \label{subsec:ma-14-2-3}
% 
% 设 $f(x)=f(x_1,x_2,\cdots,x_n)$ 在区域 $D\subset\mathbb R^n$ 上具有各个偏导数。由定义, 它的每个偏导数 $\dfrac{\partial f(x)}{\partial x_i}(i=1,2,\cdots,n)$ 是 $D$ 上的一个 $n$ 元函数。若它们仍具有各个偏导数, 则称它们的偏导数为 $f(x)$ 的二阶偏导数。类似地可以定义三阶以及更高阶的偏导数。二阶及二阶以上的偏导数统称为高阶偏导数.
% 
% 显然, 若一个函数 $f(x)$ 的各个偏导数都存在 $n$ 个偏导数, 则 $f(x)$ 具有 $n^2$ 个二阶偏导数。当
% \[
%   \frac{\partial\left(\dfrac{\partial f(x)}{\partial x_i}\right)}{\partial x_k}
%   \quad(1\leqslant i,k\leqslant n)
% \]
% 存在时, 我们将其记为 $\dfrac{\partial^2 f(x)}{\partial x_k\partial x_i}$, $f''_{x_kx_i}(x)$ 或 $f''_{ki}(x)$.
% 
% \begin{example}\label{example:ma-14-2-7}
% 设函数
% \[
%   f(x,y)=
%   \begin{cases}
%     xy\dfrac{x^2-y^2}{x^2+y^2}, &x^2+y^2\ne0,\\
%     0, &x^2+y^2=0,
%   \end{cases}
% \]
% 求 $f''_{xy}(0,0)$ 及 $f''_{yx}(0,0)$.
% \end{example}
% 
% 
% \begin{solve}
% 由偏导数的定义得
% \[
%   f'_x(x,y)=
%   \begin{cases}
%     y\left[\dfrac{x^2-y^2}{x^2+y^2}+\dfrac{4x^2y^2}{(x^2+y^2)^2}\right], &x^2+y^2\ne0,\\
%     0, &x^2+y^2=0,
%   \end{cases}
% \]
% \[
%   f'_y(x,y)=
%   \begin{cases}
%     x\left[\dfrac{x^2-y^2}{x^2+y^2}-\dfrac{4x^2y^2}{(x^2+y^2)^2}\right], &x^2+y^2\ne0,\\
%     0, &x^2+y^2=0.
%   \end{cases}
% \]
% 特别地, 我们有
% \[
%   f'_x(0,y)=-y,\qquad f'_y(x,0)=x.
% \]
% 由此得
% \[
%   f''_{yx}(0,0)=-1,\qquad f''_{xy}(0,0)=1.
% \]
% 
% 从上例可知, 一个函数的各个高阶偏导数都存在时, 若改变对自变量分量的求导顺序, 高阶偏导数的值也有可能改变。但如果我们附加一些条件, 可使高阶偏导数的值对求导顺序无关。事实上, 我们有下述结果.
% \end{solve}
% 
% 
% \begin{theorem}\label{theorem:ma-14-2-3}
% 设函数 $f(x)$ 在区域 $D\subset\mathbb R^n$ 上有定义, $x_0\in D$, 且对于 $1\leqslant j<k\leqslant n$, $f''_{kj}(x)$ 与 $f''_{jk}(x)$ 在 $U(x_0,\delta)(\delta>0)$ 内存在, 并且在 $x_0$ 处连续, 则有
% \[
%   f''_{kj}(x_0)=f''_{jk}(x_0).
% \]
% \end{theorem}
% 
% 
% \begin{proof}
% 记 $f(x)=f(x_1,x_2,\cdots,x_n)$。为了使记号简便, 我们不妨设 $j=1,k=2$, 并记 $x=(x,y,x')$, $x_0=(x_0,y_0,x'_0)$ 其中 $x'=(x_3,x_4,\cdots,x_n)$, $x'_0=(x_3^0,x_4^0,\cdots,x_n^0)$。对于充分小的 $\Delta x,\Delta y$, 观察:
% \[
% \begin{aligned}
%   I(\Delta x,\Delta y)
%   &\triangleq
%   \frac{f(x_0+\Delta x,y_0+\Delta y,x'_0)-f(x_0+\Delta x,y_0,x'_0)}{\Delta x\Delta y}\\
%   &\quad-\frac{f(x_0,y_0+\Delta y,x'_0)-f(x_0)}{\Delta x\Delta y}.
% \end{aligned}
% \]
% 记
% \[
%   g(x)=f(x,y_0+\Delta y,x'_0)-f(x,y_0,x'_0),
%   \qquad
%   h(y)=f(x_0+\Delta x,y,x'_0)-f(x_0,y,x'_0).
% \]
% 一方面, 由一元函数的微分中值定理得
% \[
% \begin{aligned}
%   I(\Delta x,\Delta y)
%   &=\frac1{\Delta x\Delta y}
%   \{[f(x_0+\Delta x,y_0+\Delta y,x'_0)-f(x_0+\Delta x,y_0,x'_0)]\\
%   &\quad-[f(x_0,y_0+\Delta y,x'_0)-f(x_0)]\}
% \end{aligned}
% \]
% \[
% \begin{aligned}
%   &=\frac{g(x_0+\Delta x)-g(x_0)}{\Delta x\Delta y}
%   =\frac{g'(x_0+\theta_1\Delta x)}{\Delta y}\\
%   &=\frac{f'_x(x_0+\theta_1\Delta x,y_0+\Delta y,x'_0)
%   -f'_x(x_0+\theta_1\Delta x,y_0,x'_0)}{\Delta y}\\
%   &=f''_{yx}(x_0+\theta_1\Delta x,y_0+\theta_2\Delta y,x'_0),
% \end{aligned}
% \]
% 其中 $0<\theta_1,\theta_2<1$。另一方面, 我们有
% \[
% \begin{aligned}
%   I(\Delta x,\Delta y)
%   &=\frac1{\Delta x\Delta y}
%   \{[f(x_0+\Delta x,y_0+\Delta y,x'_0)-f(x_0,y_0+\Delta y,x'_0)]\\
%   &\quad-[f(x_0+\Delta x,y_0,x'_0)-f(x_0)]\}\\
%   &=\frac{h(y_0+\Delta y)-h(y_0)}{\Delta x\Delta y}
%   =\frac{h'(y_0+\theta_3\Delta y)}{\Delta x}\\
%   &=\frac{f'_y(x_0+\Delta x,y_0+\theta_3\Delta y,x'_0)
%   -f'_y(x_0,y_0+\theta_3\Delta y,x'_0)}{\Delta x}\\
%   &=f''_{xy}(x_0+\theta_4\Delta x,y_0+\theta_3\Delta y,x'_0),
% \end{aligned}
% \]
% 其中 $0<\theta_3,\theta_4<1$。由此得
% \[
%   f''_{yx}(x_0+\theta_1\Delta x,y_0+\theta_2\Delta y,x'_0)
%   =f''_{xy}(x_0+\theta_4\Delta x,y_0+\theta_3\Delta y,x'_0).
% \]
% 在上式中令 $\Delta x\to0,\Delta y\to0$, 并利用 $f''_{yx}(x)$ 与 $f''_{xy}(x)$ 在 $x_0$ 处的连续性即得
% \[
%   f''_{yx}(x_0)=f''_{xy}(x_0).
% \qedhere
% \]
% \end{proof}
% 
% 
% 在今后, 我们碰到的大部分多元函数都具有各阶连续偏导数, 因此, 我们可以不考虑对自变量分量的求导顺序。例如, 当 $y=f(x)(x\in D\subset\mathbb R^n)$ 具有各个二阶连续偏导数时, $f(x)$ 在 $x_0\in D$ 处具有 $\dfrac{n^2+n}{2}$ 个不同的二阶偏导数.
% 
% \subsection{复合函数的高阶偏导数}
% \label{subsec:ma-14-2-4}
% 
% 对于复合函数的高阶偏导数, 很难得出一般性的公式, 需要逐次求之。例如, 设 $u=f(x,y),x=\varphi(s,t),y=\psi(s,t)$, 并假设所有涉及的函数均具有二阶连续偏导数, 则有
% \[
%   \frac{\partial u}{\partial s}
%   =\frac{\partial f}{\partial x}\cdot\frac{\partial\varphi}{\partial s}
%   +\frac{\partial f}{\partial y}\cdot\frac{\partial\psi}{\partial s},
% \]
% \[
% \begin{aligned}
%   \frac{\partial^2u}{\partial s^2}
%   &=\frac{\partial^2 f}{\partial x^2}\left(\frac{\partial\varphi}{\partial s}\right)^2
%   +\frac{\partial^2 f}{\partial x\partial y}\cdot
%   \frac{\partial\varphi}{\partial s}\cdot\frac{\partial\psi}{\partial s}
%   +\frac{\partial f}{\partial x}\cdot\frac{\partial^2\varphi}{\partial s^2}\\
%   &\quad+\frac{\partial^2 f}{\partial x\partial y}\cdot
%   \frac{\partial\varphi}{\partial s}\cdot\frac{\partial\psi}{\partial s}
%   +\frac{\partial^2 f}{\partial y^2}\left(\frac{\partial\psi}{\partial s}\right)^2
%   +\frac{\partial f}{\partial y}\cdot\frac{\partial^2\psi}{\partial s^2}\\
%   &=\frac{\partial^2 f}{\partial x^2}\left(\frac{\partial\varphi}{\partial s}\right)^2
%   +2\frac{\partial^2 f}{\partial x\partial y}\cdot
%   \frac{\partial\varphi}{\partial s}\cdot\frac{\partial\psi}{\partial s}
%   +\frac{\partial^2 f}{\partial y^2}\left(\frac{\partial\psi}{\partial s}\right)^2\\
%   &\quad+\frac{\partial f}{\partial x}\cdot\frac{\partial^2\varphi}{\partial s^2}
%   +\frac{\partial f}{\partial y}\cdot\frac{\partial^2\psi}{\partial s^2}.
% \end{aligned}
% \]
% 
% \begin{example}\label{example:ma-14-2-8}
% 设函数 $u=f\left(xy,\dfrac yx\right)$, 试求 $\dfrac{\partial^2u}{\partial x^2}$ 及 $\dfrac{\partial^2u}{\partial x\partial y}$ (假定其二阶偏导数均连续).
% \end{example}
% 
% 
% \begin{solve}
% 由
% \[
%   \frac{\partial u}{\partial x}=yf'_1-\frac y{x^2}f'_2
% \]
% 得
% \[
% \begin{aligned}
%   \frac{\partial^2u}{\partial x^2}
%   &=y^2f''_{11}-\frac y{x^2}f''_{12}
%   -\frac y{x^2}f''_{12}+\frac{y^2}{x^4}f''_{22}
%   +\frac{2y}{x^3}f'_2\\
%   &=y^2f''_{11}-\frac{2y}{x^2}f''_{12}
%   +\frac{y^2}{x^4}f''_{22}+\frac{2y}{x^3}f'_2.
% \end{aligned}
% \]
% \[
% \begin{aligned}
%   \frac{\partial^2u}{\partial x\partial y}
%   &=f'_1+xyf''_{11}+\frac yx f''_{12}
%   -\frac1{x^2}f'_2-\frac{xy}{x^2}f''_{12}
%   -\frac y{x^3}f''_{22}\\
%   &=xyf''_{11}-\frac y{x^3}f''_{22}+f'_1-\frac1{x^2}f'_2.
% \end{aligned}
% \qedhere
% \]
% \end{solve}
% 
% 
% \begin{example}\label{example:ma-14-2-9}
% 设函数 $u=f(r),r=\sqrt{x^2+y^2+z^2}$, 且假设 $u$ 满足拉普拉斯方程
% \[
%   \frac{\partial^2u}{\partial x^2}
%   +\frac{\partial^2u}{\partial y^2}
%   +\frac{\partial^2u}{\partial z^2}=0,
% \]
% 求 $f(r)$ 的表达式.
% \end{example}
% 
% 
% \begin{solve}
% 由
% \[
%   \frac{\partial u}{\partial x}=f'(r)\frac xr
% \]
% 得
% \[
%   \frac{\partial^2u}{\partial x^2}
%   =f''(r)\frac{x^2}{r^2}+f'(r)\frac{r-\dfrac{x^2}{r}}{r^2}.
% \]
% 同理,
% \[
%   \frac{\partial^2u}{\partial y^2}
%   =f''(r)\frac{y^2}{r^2}+f'(r)\frac{r-\dfrac{y^2}{r}}{r^2},
% \]
% \[
%   \frac{\partial^2u}{\partial z^2}
%   =f''(r)\frac{z^2}{r^2}+f'(r)\frac{r-\dfrac{z^2}{r}}{r^2}.
% \]
% 由于 $u$ 满足拉普拉斯方程, 因此有
% \[
% \begin{aligned}
%   \frac{\partial^2u}{\partial x^2}
%   +\frac{\partial^2u}{\partial y^2}
%   +\frac{\partial^2u}{\partial z^2}
%   &=f''(r)\frac{x^2+y^2+z^2}{r^2}
%   +f'(r)\frac{3r^2-(x^2+y^2+z^2)}{r^3}\\
%   &=f''(r)+\frac2r f'(r)=0.
% \end{aligned}
% \]
% 由此推出
% \[
%   (r^2f'(r))'=r^2f''(r)+2rf'(r)
%   =r^2\left(f''(r)+\frac2r f'(r)\right)=0.
% \]
% 从上式我们进一步推出: 存在常数 $C'$, 使得
% \[
%   f'(r)=\frac{C'}{r^2}.
% \]
% 对上式两边求不定积分得
% \[
%   f(r)=-\frac{C'}r+C_1=\frac C{\sqrt{x^2+y^2+z^2}}+C_1,
% \]
% 其中 $C=-C'$, $C_1$ 为常数.
% \end{solve}
% 
% 
% \subsection{一阶微分的形式不变性与高阶微分}
% \label{subsec:ma-14-2-5}
% 
% 设 $D\subset\mathbb R^n$ 是区域, 函数 $f(x)$ 在 $x\in D$ 处可微, 即有
% \[
%   \mathrm df(x)=\sum_{i=1}^n\frac{\partial f(x)}{\partial x_i}\mathrm dx_i.
% \]
% 若对于 $\forall i(1\leqslant i\leqslant n)$, $\dfrac{\partial f}{\partial x_i}$ 仍是可微函数, 则我们称
% \[
%   \sum_{i=1}^n\left(\sum_{k=1}^n
%   \frac{\partial^2 f(x)}{\partial x_k\partial x_i}\mathrm dx_k\right)\mathrm dx_i
% \]
% 为 $f(x)$ 的二阶微分, 并记为 $\mathrm d^2f(x)$, 即
% \[
%   \mathrm d^2f(x)=\sum_{i=1}^n\sum_{k=1}^n
%   \frac{\partial^2 f(x)}{\partial x_k\partial x_i}\mathrm dx_k\mathrm dx_i.
% \]
% 对于 $k\geqslant2(k\in\mathbb N)$, 我们可以归纳地给出多元函数 $f(x)$ 的 $k$ 阶微分
% \[
%   \mathrm d^kf(x)=\mathrm d(\mathrm d^{k-1}f(x)).
% \]
% 若形式地记 $\mathrm df(x)=\left(\sum_{i=1}^n\mathrm dx_i\frac{\partial}{\partial x_i}\right)f(x)$, 当 $f(x)$ 的每个 $k$ 阶偏导数都连续时, 我们可以将 $\mathrm d^kf(x)$ 记为
% \[
%   \mathrm d^kf(x)=\left(\sum_{i=1}^n\mathrm dx_i\frac{\partial}{\partial x_i}\right)^kf(x).
% \]
% 这样的记号对于高阶微分 (二阶及二阶以上的微分) 来说比较容易记忆。特别地, 对于一个二元函数 $f(x,y)$, 若它各个 $k$ 阶偏导数都存在且连续, 则我们有
% \[
% \begin{aligned}
%   \mathrm d^kf(x,y)
%   &=\left(\mathrm dx\frac{\partial}{\partial x}
%   +\mathrm dy\frac{\partial}{\partial y}\right)^kf(x,y)\\
%   &=\sum_{j=0}^k C_k^j
%   \frac{\partial^k f(x,y)}{\partial x^{k-j}\partial y^j}
%   \mathrm dx^{k-j}\mathrm dy^j,
% \end{aligned}
% \]
% 其中 $C_k^j=\dfrac{k!}{j!(k-j)!}(j=1,2,\cdots,k),C_k^0=1$.
% 
% 现设 $f(u)=f(u_1,u_2,\cdots,u_m)$ 在区域 $D\subset\mathbb R^m$ 上可微, 则 $f(u)$ 的微分 $\mathrm df(u)=f'(u)\mathrm du$。注意到此时 $u$ 是自变量.
% 
% 现在设 $u=(u_1(x),u_2(x),\cdots,u_m(x))^\mathrm T$ 是区域 $\Omega\subset\mathbb R^n$ 上的一个 $n$ 元可微向量函数, 并且 $u(\Omega)\subset D$, 则复合函数 $f(u(x))$ 在 $\Omega$ 处可微, 从而有
% \[
%   \mathrm df(u(x))=f'(u(x))u'(x)\mathrm dx.
% \]
% 由于 $\mathrm du=(u'_1(x)\mathrm dx,u'_2(x)\mathrm dx,\cdots,u'_m(x)\mathrm dx)^\mathrm T=u'(x)\mathrm dx$, 因此, 当 $u$ 是中间变量时, 我们仍然有
% \[
%   \mathrm df(u)=f'(u(x))u'(x)\mathrm dx=f'(u)\mathrm du.
% \]
% 以上讨论说明了对于多元函数的一阶微分仍具有形式不变性。在一元微积分中, 我们已经知道了对于二阶以上的微分不再具有形式不变性, 因此多元函数情形也不可能具有高阶微分的形式不变性.
% 
% \begin{example}\label{example:ma-14-2-10}
% 设函数 $f\left(xy,\dfrac yx\right)$ 具有二阶连续偏导数, 求它的所有二阶偏导数.
% \end{example}
% 
% 
% \begin{solve}
% 在例~\ref{example:ma-14-2-8} 中我们曾经直接对该函数求过二阶偏函数, 现在我们利用 $f$ 的二阶微分来求之。由
% \[
%   \mathrm df=f'_1\mathrm d(xy)+f'_2\mathrm d\left(\frac yx\right)
% \]
% 得
% \[
% \begin{aligned}
%   \mathrm d^2f
%   &=\left(f''_{11}\mathrm d(xy)+f''_{12}\mathrm d\left(\frac yx\right)\right)\mathrm d(xy)
%   +f'_1\mathrm d^2(xy)\\
%   &\quad+\left(f''_{21}\mathrm d(xy)+f''_{22}\mathrm d\left(\frac yx\right)\right)
%   \mathrm d\left(\frac yx\right)+f'_2\mathrm d^2\left(\frac yx\right)\\
%   &=\left(y^2f''_{11}-\frac{2y^2}{x^2}f''_{12}
%   +\frac{y^2}{x^4}f''_{22}+\frac{2y}{x^3}f'_2\right)\mathrm dx^2\\
%   &\quad+2\left(xyf''_{11}-\frac y{x^3}f''_{22}
%   +f'_1-\frac1{x^2}f'_2\right)\mathrm dx\mathrm dy\\
%   &\quad+\left(x^2f''_{11}+2f''_{12}+\frac1{x^2}f''_{22}\right)\mathrm dy^2.
% \end{aligned}
% \]
% 因此
% \[
%   \frac{\partial^2 f}{\partial x^2}
%   =y^2f''_{11}-\frac{2y^2}{x^2}f''_{12}
%   +\frac{y^2}{x^4}f''_{22}+\frac{2y}{x^3}f'_2,
% \]
% \[
%   \frac{\partial^2 f}{\partial x\partial y}
%   =xyf''_{11}-\frac y{x^3}f''_{22}
%   +f'_1-\frac1{x^2}f'_2,
% \]
% \[
%   \frac{\partial^2 f}{\partial y^2}
%   =x^2f''_{11}+2f''_{12}+\frac1{x^2}f''_{22}.
% \qedhere
% \]
% \end{solve}
% 
% 
% \section{泰勒公式}
% \label{sec:ma-14-3}
% 
% 在利用高阶导数研究函数时, 泰勒公式是一个强有力的工具, 对于多元函数我们有以下定理.
% 
% \begin{theorem}[{泰勒公式}]\label{theorem:ma-14-3-1}
% 设函数 $f(x)$ 在 $x_0=(x_1^0,x_2^0,\cdots,x_n^0)\in\mathbb R^n$ 的邻域 $U(x_0,\delta_0)(\delta_0>0)$ 内具有 $K+1$ 阶连续偏导数, 则对于 $\forall x_0+h=(x_1^0+h_1,x_2^0+h_2,\cdots,x_n^0+h_n)\in U(x_0,\delta_0)$, 有
% \begin{equation}
% \begin{aligned}
%   f(x_0+h)
%   &=f(x_0)+\sum_{k=1}^K\frac1{k!}
%   \left(\sum_{i=1}^n h_i\frac{\partial}{\partial x_i}\right)^k f(x_0)\\
%   &\quad+\frac1{(K+1)!}
%   \left(\sum_{i=1}^n h_i\frac{\partial}{\partial x_i}\right)^{K+1}
%   f(x_0+\theta h),
% \end{aligned}
% \label{eq:ma-14-3-1}
% \end{equation}
% 其中 $0<\theta<1$.
% \end{theorem}
% 
% 
% \begin{proof}
% 设 $h=(h_1,h_2,\cdots,h_n)$ 满足 $x_0+h\in U(x_0,\delta_0)$。构造一元函数
% \[
%   \varphi(t)=f(x_0+th),\qquad t\in[0,1].
% \]
% 由复合函数求导公式知, $\varphi(t)$ 具有 $K+1$ 阶连续导数。因此, 由一元函数的泰勒公式有
% \begin{equation}
%   \varphi(t)=\varphi(0)+\sum_{k=1}^K
%   \frac{\varphi^{(k)}(0)}{k!}t^k
%   +\frac{\varphi^{(K+1)}(\theta t)}{(K+1)!}t^{K+1},
%   \label{eq:ma-14-3-2}
% \end{equation}
% 其中 $0<\theta<1$。注意到
% \[
%   \varphi'(t)=\sum_{i=1}^n\frac{\partial f(x_0+th)}{\partial x_i}h_i
%   =\left(\sum_{i=1}^n h_i\frac{\partial}{\partial x_i}\right)f(x_0+th),
% \]
% \[
%   \varphi''(t)=\sum_{j=1}^n\sum_{i=1}^n
%   \frac{\partial^2 f(x_0+th)}{\partial x_i\partial x_j}h_ih_j
%   =\left(\sum_{i=1}^n h_i\frac{\partial}{\partial x_i}\right)^2 f(x_0+th),
% \]
% 一般地, 对于 $k=1,2,\cdots,K+1$, 有
% \[
%   \varphi^{(k)}(t)
%   =\left(\sum_{i=1}^n h_i\frac{\partial}{\partial x_i}\right)^k f(x_0+th).
% \]
% 因此, 将 $\varphi^{(k)}(t)(k=1,2,\cdots,K+1)$ 的表达式代入式 \eqref{eq:ma-14-3-2}, 然后令 $t=1$ 即得所证.
% \end{proof}
% 
% 
% \begin{remark}\label{remark:ma-14-source-71}
% 当 $f(x)$ 在区域 $D$ 内具有 $K+1$ 阶连续偏导数时, 定理~\ref{theorem:ma-14-3-1} 中的余项
% \[
%   R_{K+1}(h)=\frac1{(K+1)!}
%   \left(\sum_{i=1}^n h_i\frac{\partial}{\partial x_i}\right)^{K+1}f(x_0+\theta h)
%   \quad(0<\theta<1)
% \]
% 称为拉格朗日余项, 并称式 \eqref{eq:ma-14-3-1} 为 $f(x)$ 带拉格朗日余项的泰勒公式.
% \end{remark}
% 
% 
% \begin{corollary}\label{corollary:ma-local-14.3-1}
% 设函数 $f(x)\in C^K(U(x_0,\delta_0))(\delta_0>0)$, 即 $f(x)$ 在邻域 $U(x_0,\delta_0)$ 内具有 $K(K\geqslant1)$ 阶连续偏导数, $x_0+h\in U(x_0,\delta_0)$, 则
% \begin{equation}
%   f(x_0+h)=f(x_0)+\sum_{k=1}^K\frac1{k!}
%   \left(\sum_{i=1}^n h_i\frac{\partial}{\partial x_i}\right)^k f(x_0)
%   +o(|h|^K)
%   \quad(|h|\to0).
%   \label{eq:ma-14-3-3}
% \end{equation}
% \end{corollary}
% 
% 
% \begin{proof}
% 由定理~\ref{theorem:ma-14-3-1} 我们有
% \[
% \begin{aligned}
%   f(x_0+h)
%   &=f(x_0)+\sum_{k=1}^K\frac1{k!}
%   \left(\sum_{i=1}^n h_i\frac{\partial}{\partial x_i}\right)^kf(x_0)\\
%   &\quad+\frac1{K!}
%   \left(\sum_{i=1}^n h_i\frac{\partial}{\partial x_i}\right)^K
%   (f(x_0+\theta h)-f(x_0)),
% \end{aligned}
% \]
% 其中 $0<\theta<1$.
% 
% 记
% \[
%   R_K(h)=\frac1{K!}
%   \left(\sum_{i=1}^n h_i\frac{\partial}{\partial x_i}\right)^K
%   (f(x_0+\theta h)-f(x_0)).
% \]
% 由 $f(x)$ 在 $U(x_0,\delta_0)$ 内具有 $K(K\geqslant1)$ 阶连续偏导数, 知
% \[
%   \lim_{|h|\to0}\frac{R_K(h)}{|h|^K}
%   =\lim_{|h|\to0}\frac1{K!}
%   \left(\sum_{i=1}^n\frac{h_i}{|h|}\frac{\partial}{\partial x_i}\right)^K
%   (f(x_0+\theta h)-f(x_0))=0.
% \]
% 上述极限为零是因为当 $|h|\to0$ 时, 和式中的每一项的系数都是有界的, 而 $f(x)$ 在 $x_0+\theta h$ 处的每个 $K$ 阶偏导数由于其连续性都趋于 $f(x)$ 在 $x_0$ 相应的 $K$ 阶偏导数.
% \end{proof}
% 
% 
% \begin{remark}\label{remark:ma-14-source-74}
% 我们将推论~\ref{corollary:ma-local-14.3-1} 中的余项 $o(|h|^K)$ 称为皮亚诺余项, 并称式 \eqref{eq:ma-14-3-3} 为 $f(x)$ 带皮亚诺余项的泰勒公式。由推论~\ref{corollary:ma-local-14.3-1} 我们可知, 若 $f(x)$ 在 $x_0$ 处能展成
% \[
%   f(x_0+h)=P_K(h_1,\cdots,h_n)+o(|h|^K)
%   \quad(|h|\to0),
% \]
% 其中 $P_K$ 是 $n$ 元 $K$ 次多项式, 上式必定是 $f(x)$ 在 $U(x_0,\delta_0)$ 内的泰勒公式.
% \end{remark}
% 
% 现设 $f(x)$ 在区域 $D$ 内具有 $K+1$ 阶连续偏导数, 从定理~\ref{theorem:ma-14-3-1} 的证明中可以看出, 若要求 $f(x)(x\in D)$ 在 $x_0\in D$ 处的泰勒公式, 我们不必要求 $x$ 在 $x_0$ 的某个邻域内, 而只要求连接 $x_0$ 和 $x$ 的线段 $x_0+t(x-x_0)(t\in[0,1])$ 落在 $D$ 内即可.
% 
% 在定理~\ref{theorem:ma-14-3-1} 中取 $K=0$, 我们有以下多元函数的拉格朗日微分中值定理.
% 
% \begin{corollary}\label{corollary:ma-local-14.3-2}
% 设 $f(x)$ 在区域 $D\subset\mathbb R^n$ 内具有连续偏导数, 且对于 $\forall t\in[0,1],x_0+t(x-x_0)\in D$, 则有
% \[
% \begin{aligned}
%   f(x)-f(x_0)
%   &=\sum_{i=1}^n
%   \frac{\partial f(x_0+\theta(x-x_0))}{\partial x_i}(x_i-x_i^0)\\
%   &=f'(x_0+\theta(x-x_0))\cdot(x-x_0)^\mathrm T,
% \end{aligned}
% \]
% 其中 $0<\theta<1$.
% \end{corollary}
% 
% 由上述推论~\ref{corollary:ma-local-14.3-2}, 我们容易推出下述结论.
% 
% \begin{corollary}\label{corollary:ma-local-14.3-3}
% 设函数 $f(x)$ 在区域 $D\subset\mathbb R^n$ 内的各个偏导数均为 0, 则 $f(x)$ 在 $D$ 内为常数函数.
% \end{corollary}
% 
% 请读者给出上述推论的证明.
% 
% 以后我们还常常要用到泰勒公式中 $K=1$ 的情形。在定理~\ref{theorem:ma-14-3-1} 中, 令 $K=1$, 我们有
% \[
%   f(x_0+h)=f(x_0)+f'(x_0)\cdot h^\mathrm T
%   +\frac12 h\cdot H_f(x_0)h^\mathrm T+o(|h|^2)
%   \quad(|h|\to0),
% \]
% 其中
% \[
%   H_f(x_0)=
%   \begin{pmatrix}
%     \dfrac{\partial^2 f(x_0)}{\partial x_1^2} &
%     \dfrac{\partial^2 f(x_0)}{\partial x_1\partial x_2} &
%     \cdots &
%     \dfrac{\partial^2 f(x_0)}{\partial x_1\partial x_n}\\
%     \dfrac{\partial^2 f(x_0)}{\partial x_2\partial x_1} &
%     \dfrac{\partial^2 f(x_0)}{\partial x_2^2} &
%     \cdots &
%     \dfrac{\partial^2 f(x_0)}{\partial x_2\partial x_n}\\
%     \vdots & \vdots & & \vdots\\
%     \dfrac{\partial^2 f(x_0)}{\partial x_n\partial x_1} &
%     \dfrac{\partial^2 f(x_0)}{\partial x_n\partial x_2} &
%     \cdots &
%     \dfrac{\partial^2 f(x_0)}{\partial x_n^2}
%   \end{pmatrix}
% \]
% 称为 $f(x)$ 在 $x_0$ 处的海色 (Hessi) 矩阵.
% 
% \begin{example}\label{example:ma-14-3-1}
% 求函数 $f(x,y)=\dfrac{x-y}{x+y}$ 在 $(1,0)$ 附近带皮亚诺余项的泰勒公式 (直到二次项).
% \end{example}
% 
% 
% \begin{solve}
% 对于这类函数, 一般可直接求偏导, 然后写出泰勒公式, 但我们下面利用一元函数已知的泰勒公式来求之。由于
% \[
% \begin{aligned}
%   f(x,y)
%   &=\frac{(x-1)+1-y}{1+[(x-1)+y]}\\
%   &=[(x-1)+1-y]\{1-(x-1)-y+[(x-1)+y]^2
%   +o((x-1)^2+y^2)\}
% \end{aligned}
% \]
% $((x-1)^2+y^2\to0)$, 将上式整理得
% \[
%   f(x,y)=1-2y+2(x-1)y+2y^2+o((x-1)^2+y^2),
%   \quad ((x-1)^2+y^2\to0).
% \qedhere
% \]
% \end{solve}
% 
% 
% \begin{example}\label{example:ma-14-3-2}
% 求函数 $f(x,y)=xe^{x+y}$ 在原点 $(0,0)$ 处的所有四阶偏导数.
% \end{example}
% 
% 
% \begin{solve}
% 由 $e^{x+y}$ 在 $(0,0)$ 处的泰勒公式有
% \[
%   f(x,y)=xe^{x+y}
%   =x\left[1+\sum_{k=1}^K\frac{(x+y)^k}{k!}
%   +o\left((\sqrt{x^2+y^2})^K\right)\right]
%   \quad(\sqrt{x^2+y^2}\to0).
% \]
% 令 $K=3$, 得
% \[
% \begin{aligned}
%   f(x,y)
%   &=x+x^2+xy+\frac12(x^3+2x^2y+xy^2)\\
%   &\quad+\frac16(x^4+3x^3y+3x^2y^2+xy^3)
%   +o\left((\sqrt{x^2+y^2})^4\right)
%   \quad(\sqrt{x^2+y^2}\to0),
% \end{aligned}
% \]
% 于是有
% \[
%   \frac{C_4^4}{4!}\cdot\frac{\partial^4 f(0,0)}{\partial x^4}=\frac16,\quad
%   \frac{C_4^3}{4!}\cdot\frac{\partial^4 f(0,0)}{\partial x^3\partial y}=\frac12,\quad
%   \frac{C_4^2}{4!}\cdot\frac{\partial^4 f(0,0)}{\partial x^2\partial y^2}=\frac12,
% \]
% \[
%   \frac{C_4^1}{4!}\cdot\frac{\partial^4 f(0,0)}{\partial x\partial y^3}=\frac16,\quad
%   \frac{C_4^0}{4!}\cdot\frac{\partial^4 f(0,0)}{\partial y^4}=0,
% \]
% 从而得
% \[
%   \frac{\partial^4 f(0,0)}{\partial x^4}=4,\quad
%   \frac{\partial^4 f(0,0)}{\partial x^3\partial y}=3,\quad
%   \frac{\partial^4 f(0,0)}{\partial x^2\partial y^2}=2,
% \]
% \[
%   \frac{\partial^4 f(0,0)}{\partial x\partial y^3}=1,\quad
%   \frac{\partial^4 f(0,0)}{\partial y^4}=0.
% \qedhere
% \]
% \end{solve}
% 
% 
% \section{隐函数存在定理}
% \label{sec:ma-14-4}
% 
% \subsection{单个方程的情形}
% \label{subsec:ma-14-4-1}
% 
% 在讨论一般的隐函数存在问题之前, 我们先来讨论一下由一个二元方程确定隐函数的简单情形。设 $F(x,y)=0$ 是区域 $D\subset\mathbb R^2$ 上的一个二元方程, 并且满足 $F(x_0,y_0)=0$, 其中 $(x_0,y_0)\in D$。我们要研究的问题是: 何时方程 $F(x,y)=0$ 在点 $(x_0,y_0)$ 附近能唯一确定一个函数 $y=f(x)$, 使得 $f(x_0)=y_0$, 且 $F(x,f(x))=0$ 对 $x\in(x_0-\delta,x_0+\delta)(\delta>0)$ 成立? 我们称这样定义的函数 $y=f(x)$ 为由方程 $F(x,y)=0$ 所确定的隐函数。显然这个问题等价于: 何时 $z=F(x,y)$ 的零点集合在 $(x_0,y_0)$ 附近是 $Oxy$ 平面内的一条曲线, 而且这条曲线是 $y=f(x)$ 的图像?
% 
% 我们先来考查一个例子。$\mathbb R^2$ 中的单位圆周满足的方程为 $x^2+y^2=1$, 我们可将其写成 $x^2+y^2-1=0$。现考虑函数 $z=F(x,y)=x^2+y^2-1$。先看一个简单情况, 令 $(x_0,y_0)=(0,1)$, 我们来考查一下为什么 $F(x,y)=0$ 在该点的附近能确定一个隐函数 $y=\sqrt{1-x^2}$。首先 $(x_0,y_0)=(0,1)$ 是 $F(x,y)=x^2+y^2-1$ 的零点, 并且平面 $x=0$ 与曲面 $z=F(x,y)$ 的交线是
% \[
%   \begin{cases}
%     x=0,\\
%     z=y^2-1.
%   \end{cases}
% \]
% 对于该曲线容易看出, 对于 $0<\delta_0<1$, 当 $y$ 从 $1-\delta_0$ 连续变化到 $1+\delta_0$ 时, $z$ 严格单调上升并且连续地从负变正。因此, 对于充分小的 $\delta>0$, 对 $(-\delta,\delta)$ 中的每个 $x'$, 曲线
% \[
%   \begin{cases}
%     x=x',\\
%     z=y^2-1+x'^2
%   \end{cases}
% \]
% 都有此性质。显然, 当具有上述性质时, $F(x,y)$ 在点 $(x_0,y_0)$ 附近的零点集合刚好是一个函数 $y=f(x)$ 的图像。而什么时候能保证一个函数 $z=F(x,y)$ 在固定 $x$ 时关于 $y$ 是一个严格单调函数呢? 从多元函数偏导数的几何意义容易看出, 一个充分条件是 $F'_y(x,y)\ne0$.
% 
% 对于满足 $F'_y(x,y)=0$ 的点 $(x_0,y_0)$, 读者容易发现此时未必有隐函数 $y=y(x)$ 存在, 如 $F(x,y)=x^2+y^2-1$ 在 $(1,0)$ 处.
% 
% 下面的定理给出了隐函数存在的充分条件.
% 
% \begin{theorem}[{隐函数存在定理}]\label{theorem:ma-14-4-1}
% 设二元函数 $F(x,y)$ 在 $U((x_0,y_0),\delta)(\delta>0)$ 内满足以下条件:
% 
% (1) $F(x_0,y_0)=0$;
% 
% (2) $F(x,y),F'_y(x,y)$ 在 $U((x_0,y_0),\delta)$ 内连续;
% 
% (3) $F'_y(x_0,y_0)\ne0$,
% 
% 则 $\exists\delta_0>0(0<\delta_0<\delta)$, 使得在 $U(x_0,\delta_0)$ 内存在唯一满足下述条件的连续函数 $y=f(x)$:
% 
% (a) $y_0=f(x_0)$;
% 
% (b) $F(x,f(x))=0,\forall x\in U(x_0,\delta_0)$;
% 
% (c) 如果 $F'_x(x,y)$ 在 $U((x_0,y_0),\delta)$ 内连续, 则 $f(x)$ 在 $U(x_0,\delta_0)$ 存在连续导数, 并且有
% \[
%   f'(x)=-\frac{F'_x(x,f(x))}{F'_y(x,f(x))}.
% \]
% \end{theorem}
% 
% 
% \begin{proof}
% 不妨设 $F'_y(x_0,y_0)>0$。由 $F'_y(x,y)$ 的连续性及极限的保号性可知, 存在 $0<\delta_1,\delta_2<\delta$, 使得对于 $\forall(x,y)\in U(x_0,\delta_1)\times U(y_0,\delta_2)$, 有
% \[
%   F'_y(x,y)>0.
% \]
% 特别地, 若固定 $x=x_0$, 则对于 $\forall y\in U(y_0,\delta_2)$, 有 $F'_y(x_0,y)>0$。因此 $F(x_0,y)$ 在 $U(y_0,\delta_2)$ 内是 $y$ 的严格上升函数。注意到 $F(x_0,y_0)=0$, 我们可知 $F(x_0,y_0-\delta_2)<0,F(x_0,y_0+\delta_2)>0$。再由 $F(x,y)$ 在 $(x_0,y_0-\delta_2)$ 和 $(x_0,y_0+\delta_2)$ 处的连续性, 并利用极限的保号性, 我们可以找到 $\delta_0(0<\delta_0<\delta_1)$, 使得当 $x\in U(x_0,\delta_0)$ 时, 有
% \[
%   F(x,y_0-\delta_2)<0,\qquad F(x,y_0+\delta_2)>0.
% \]
% 对任意给定的 $\tilde x\in U(x_0,\delta_0)$, 由 $F'_y(\tilde x,y)>0$ 推知, 当 $y\in U(y_0,\delta_2)$ 时, $F(\tilde x,y)$ 连续地从负数 $F(\tilde x,y_0-\delta_2)$ 严格上升到正数 $F(\tilde x,y_0+\delta_2)$ (见图 14.4.1)。因此存在唯一的 $\tilde y\in U(y_0,\delta_2)$, 使得
% \[
%   F(\tilde x,\tilde y)=0.
% \]
% 这说明了对于 $\forall\tilde x\in U(x_0,\delta_0)$, 有唯一确定的 $\tilde y$ 与之对应。由函数的定义, 若令 $\tilde y=f(\tilde x)$, 则函数 $\tilde y=f(\tilde x)$ 在 $U(x_0,\delta_0)$ 内有定义, 并且满足 (a),(b)。该函数的唯一性由 $F(x,y)$ 在 $U(x_0,\delta_0)\times U(y_0,\delta_2)$ 内关于 $y$ 严格上升所保证.
% 
% 下面证 $y=f(x)$ 在 $U(x_0,\delta_0)$ 内连续。任意取定 $\bar x\in U(x_0,\delta_0)$, 记 $\bar y=f(\bar x)$。对于充分小的 $\varepsilon>0$, 则有
% \[
%   F(\bar x,\bar y-\varepsilon)<0,\qquad F(\bar x,\bar y+\varepsilon)>0.
% \]
% 由 $F(x,y)$ 在 $(\bar x,\bar y-\varepsilon)$ 及 $(\bar x,\bar y+\varepsilon)$ 处的连续性, 必存在充分小的 $\delta'(0<\delta'<\delta_0)$, 使得当 $x\in U(\bar x,\delta')\subset U(x_0,\delta_0)$ 时, 有
% \[
%   F(x,\bar y-\varepsilon)<0,\qquad F(x,\bar y+\varepsilon)>0.
% \]
% 从上面关于隐函数存在性的证明过程知, 当 $x\in U(\bar x,\delta')$ 时, 必有 $f(x)\in(\bar y-\varepsilon,\bar y+\varepsilon)$, 即
% \[
%   |f(x)-f(\bar x)|<\varepsilon.
% \]
% 这说明 $f(x)$ 在 $U(x_0,\delta_0)$ 内连续.
% 
% 最后, 设 $F'_x(x,y)$ 在 $U((x_0,y_0),\delta)$ 内存在且连续, 我们证明 $y=f(x)$ 在 $U(x_0,\delta_0)$ 内具有连续导数。任取 $\bar x\in U(x_0,\delta_0)$, 再取 $\Delta x$ 充分小, 使得 $\bar x+\Delta x\in U(x_0,\delta_0)$。记
% \[
%   \bar y=f(\bar x),\qquad
%   \Delta y=f(\bar x+\Delta x)-f(\bar x).
% \]
% 由多元函数拉格朗日微分中值定理, $\exists0<\theta<1$, 使得下式成立:
% \[
% \begin{aligned}
%   0&=F(\bar x+\Delta x,\bar y+\Delta y)-F(\bar x,\bar y)\\
%   &=F'_x(\bar x+\theta\Delta x,\bar y+\theta\Delta y)\Delta x
%   +F'_y(\bar x+\theta\Delta x,\bar y+\theta\Delta y)\Delta y.
% \end{aligned}
% \]
% 再由 $F'_y(x,y)\ne0((x,y)\in U((x_0,y_0),\delta))$, 得
% \[
%   \frac{\Delta y}{\Delta x}
%   =-\frac{F'_x(\bar x+\theta\Delta x,\bar y+\theta\Delta y)}
%   {F'_y(\bar x+\theta\Delta x,\bar y+\theta\Delta y)}.
% \]
% 令 $\Delta x\to0$, 由 $F'_x(x,y)$ 及 $F'_y(x,y)$ 的连续性得
% \[
%   f'(\bar x)=-\frac{F'_x(\bar x,f(\bar x))}{F'_y(\bar x,f(\bar x))}.
% \]
% 从上式可以看出 $f'(x)$ 在 $U(x_0,\delta_0)$ 内连续.
% \end{proof}
% 
% 
% 关于隐函数存在定理, 我们要注意以下几点:
% 
% \paragraph*{注 1}
% 隐函数存在定理是一个具有重要理论价值的定理, 这在许多后续课程中读者将会有所体验。值得指出的是, 在定理的条件下, $F(x,y)=0$ 所确定的隐函数是局部存在的, 并且定理的结论只给出了存在性。在很多情况下, 我们未必能找到解析式 $y=f(x)$ 来表示这个隐函数。例如, 对于开普勒 (Kepler) 方程
% \[
%   y-x-\varepsilon\sin y=0\quad(0<\varepsilon<1),
% \]
% 容易验证它在 $(-\infty,+\infty)$ 上能确定函数 $y=f(x)$, 但人们却无法写出 $y=f(x)$ 的解析式.
% 
% \paragraph*{注 2}
% 定理~\ref{theorem:ma-14-4-1} 只是给了隐函数存在的一个充分条件, 若对定理证明过程加以分析, 读者容易举出例子来说明, 当 $F(x,y)$ 不满足定理的条件时, 有时也可以唯一确定一个隐函数 (见本章习题).
% 
% 在隐函数存在定理中, 将 $F(x,y)$ 的 $x\in U(x_0,\delta)$ 换成 $\mathbb R^n$ 中的点 $x\in U(x_0,\delta)(\delta>0)$ 时, 相应的隐函数存在和唯一性的结论不仅成立, 而且其证明与定理~\ref{theorem:ma-14-4-1} 的证明相似。下面不加证明地给出此结论.
% 
% \begin{theorem}[{隐函数存在定理}]\label{theorem:ma-14-4-2}
% 记 $x=(x_1,x_2,\cdots,x_n),x_0=(x_1^0,x_2^0,\cdots,x_n^0)\in\mathbb R^n$, 假设函数 $F(x,y)=F(x_1,x_2,\cdots,x_n,y)$ 在 $U(x_0,\delta)\times U(y_0,\delta)(\delta>0)$ 内有定义, 并且满足
% 
% (1) $F(x_0,y_0)=0$;
% 
% (2) $F(x,y)$ 和 $F'_y(x,y)$ 在 $U(x_0,\delta)\times U(y_0,\delta)$ 内连续;
% 
% (3) $F'_y(x_0,y_0)\ne0$,
% 
% 则 $\exists\delta_0(0<\delta_0<\delta)$, 使得在 $U(x_0,\delta_0)$ 内存在唯一满足下述条件的连续函数 $y=f(x)$:
% 
% (a) $y_0=f(x_0)$;
% 
% (b) 对于 $\forall x\in U(x_0,\delta_0),F(x,f(x))=0$;
% 
% (c) 如果 $F(x,y)$ 在 $U(x_0,\delta)\times U(y_0,\delta)$ 内存在各个连续偏导数, 那么 $y=f(x)$ 在 $U(x_0,\delta_0)$ 具有各个连续偏导数, 并且对于 $i=1,2,\cdots,n$ 及 $\forall x\in U(x_0,\delta_0)$, 有
% \[
%   \frac{\partial f(x)}{\partial x_i}
%   =-\frac{F'_{x_i}(x_1,x_2,\cdots,x_n,y)}
%   {F'_y(x_1,x_2,\cdots,x_n,y)},\qquad
%   \text{其中 }y=f(x).
% \]
% \end{theorem}
% 
% 
% \begin{remark}\label{remark:ma-14-source-84}
% 在定理~\ref{theorem:ma-14-4-1} 和定理~\ref{theorem:ma-14-4-2} 中, 若 $F(x,y)$ (或 $F(x,y)$) 除了满足相应定理的条件外, 再假设它们具有各个高阶偏导数, 从隐函数的一阶导数或偏导数的表达式直接看出, 所确定的隐函数也具有相应的高阶导数或高阶偏导数.
% \end{remark}
% 
% 
% \begin{example}\label{example:ma-14-4-1}
% 设 $x=x(y,z),y=y(x,z),z=z(x,y)$ 都是由方程 $F(x,y,z)=0$ 确定的隐函数, 并且它们具有连续偏导数。证明:
% \[
%   (1)\quad \frac{\partial x}{\partial y}\cdot
%   \frac{\partial y}{\partial z}\cdot
%   \frac{\partial z}{\partial x}=-1;
% \]
% \[
%   (2)\quad \frac{\partial x}{\partial y}\cdot
%   \frac{\partial x}{\partial z}\cdot
%   \frac{\partial y}{\partial x}\cdot
%   \frac{\partial y}{\partial z}\cdot
%   \frac{\partial z}{\partial x}\cdot
%   \frac{\partial z}{\partial y}=1.
% \]
% \end{example}
% 
% 
% \begin{proof}
% (1) 因为
% \[
%   \frac{\partial x}{\partial y}=-\frac{F'_y(x,y,z)}{F'_x(x,y,z)},\quad
%   \frac{\partial y}{\partial z}=-\frac{F'_z(x,y,z)}{F'_y(x,y,z)},\quad
%   \frac{\partial z}{\partial x}=-\frac{F'_x(x,y,z)}{F'_z(x,y,z)},
% \]
% 所以
% \[
%   \frac{\partial x}{\partial y}\cdot
%   \frac{\partial y}{\partial z}\cdot
%   \frac{\partial z}{\partial x}=-1.
% \]
% 
% (2) 由 (1), 又因为
% \[
%   \frac{\partial x}{\partial z}=-\frac{F'_z(x,y,z)}{F'_x(x,y,z)},\quad
%   \frac{\partial y}{\partial x}=-\frac{F'_x(x,y,z)}{F'_y(x,y,z)},\quad
%   \frac{\partial z}{\partial y}=-\frac{F'_y(x,y,z)}{F'_z(x,y,z)},
% \]
% 所以
% \[
%   \frac{\partial x}{\partial y}\cdot
%   \frac{\partial x}{\partial z}\cdot
%   \frac{\partial y}{\partial x}\cdot
%   \frac{\partial y}{\partial z}\cdot
%   \frac{\partial z}{\partial x}\cdot
%   \frac{\partial z}{\partial y}=1.
% \qedhere
% \]
% \end{proof}
% 
% 
% \begin{remark}\label{remark:ma-14-source-87}
% 在一元微分学中, 当 $y$ 是 $x$ 的函数时, $\dfrac{\mathrm dy}{\mathrm dx}$ 可以看成是 $\mathrm dy$ 与 $\mathrm dx$ 的商。上例说明, 在多元微分学中, 当 $z=f(x,y)$ 可偏导时, $\dfrac{\partial z}{\partial x}$ 则是一个整体记号.
% \end{remark}
% 
% 
% \begin{example}\label{example:ma-14-4-2}
% 证明在 $(0,0,0)$ 的邻域内方程
% \begin{equation}
%   -2x+y-x^2+y^2+z+\sin z=0
%   \label{eq:ma-14-4-1}
% \end{equation}
% 确定隐函数 $z=f(x,y)$, 并求 $f(x,y)$ 在 $(0,0)$ 处带皮亚诺余项的泰勒公式 (直到二次项).
% \end{example}
% 
% 
% \begin{solve}
% 令 $F(x,y,z)=-2x+y-x^2+y^2+z+\sin z$, 则 $F(x,y,z)$ 具有各阶连续偏导数且满足 $F(0,0,0)=0$ 和 $F'_z(0,0,0)=1+\cos0=2\ne0$。因此 $F(x,y,z)=0$ 在 $(0,0,0)$ 的某一邻域内唯一确定隐函数 $z=f(x,y)$, 并且 $f(x,y)$ 在该邻域内具有各阶连续偏导数.
% 
% 注意到 $f(0,0)=0$, 因此 $z=f(x,y)$ 的泰勒公式具有下述形式:
% \begin{equation}
%   f(x,y)=a_1x+a_2y+b_{11}x^2+b_{12}xy+b_{22}y^2+o(x^2+y^2)
%   \label{eq:ma-14-4-2}
% \end{equation}
% \[
%   (\sqrt{x^2+y^2}\to0),
% \]
% 其中 $a_i,b_{ij}(i,j=1,2)$ 为常数。由于
% \begin{equation}
%   \sin z=z+o(|z|^2)
%   =a_1x+a_2y+b_{11}x^2+b_{12}xy+b_{22}y^2+o(x^2+y^2)
%   \label{eq:ma-14-4-3}
% \end{equation}
% \[
%   (\sqrt{x^2+y^2}\to0),
% \]
% 将式 \eqref{eq:ma-14-4-2},\eqref{eq:ma-14-4-3} 代入式 \eqref{eq:ma-14-4-1} 得
% \[
%   -2x+y-x^2+y^2+2a_1x+2a_2y+2b_{11}x^2+2b_{12}xy+2b_{22}y^2+o(x^2+y^2)=0.
% \]
% 由此推出
% \[
%   (2a_1-2)x+(2a_2+1)y+(2b_{11}-1)x^2
%   +2b_{12}xy+(2b_{22}+1)y^2+o(x^2+y^2)=0,
% \]
% 从而有
% \[
%   a_1=1,\qquad a_2=-\frac12,\qquad b_{11}=\frac12,\qquad
%   b_{12}=0,\qquad b_{22}=-\frac12.
% \]
% 这样我们就求出了如下 $f(x,y)$ 的带皮亚诺余项的泰勒公式:
% \[
%   f(x,y)=x-\frac12y+\frac12x^2-\frac12y^2+o(x^2+y^2)
%   \quad(\sqrt{x^2+y^2}\to0).
% \qedhere
% \]
% \end{solve}
% 
% 
% \subsection{方程组的情形}
% \label{subsec:ma-14-4-2}
% 
% 我们下面来考查两个方程所组成的方程组
% \[
%   \begin{cases}
%     F(x,u,v)=0,\\
%     G(x,u,v)=0
%   \end{cases}
% \]
% 何时能在 $(x_0,u_0,v_0)$ 的邻域内确定 $u,v$ 是 $x$ 的函数。对一个方程 $F(x,y)=0$ 的情形, $F'_y(x_0,y_0)\ne0$ 起着重要作用。因此, 当我们考虑向量函数
% \[
%   \begin{pmatrix}
%     z_1\\
%     z_2
%   \end{pmatrix}
%   =
%   \begin{pmatrix}
%     F(x,u,v)\\
%     G(x,u,v)
%   \end{pmatrix}
% \]
% 时, 若将 $x$ 看做常量, 它在 $(x_0,u_0,v_0)$ 处关于 $(u,v)$ 的导数将起着隐函数存在定理中 $F'_y(x_0,y_0)$ 的作用。注意到此时该导数是
% \[
%   \left.
%   \begin{pmatrix}
%     F'_u & F'_v\\
%     G'_u & G'_v
%   \end{pmatrix}
%   \right|_{(x_0,u_0,v_0)}.
% \]
% 由于要确定两个隐函数, 因此该雅可比矩阵应该是非退化的, 即
% \[
%   \left.
%   \begin{vmatrix}
%     F_u & F_v\\
%     G_u & G_v
%   \end{vmatrix}
%   \right|_{(x_0,u_0,v_0)}\ne0.
% \]
% 对于上述问题, 我们有以下一般性的定理.
% 
% \begin{theorem}[{隐函数组存在定理}]\label{theorem:ma-14-4-3}
% 设向量函数
% \[
%   F(x,u)=(F_1(x,u),F_2(x,u),\cdots,F_m(x,u))
% \]
% 在 $U(x_0,\delta)\times U(u_0,\delta)(\delta>0)$ 内有定义, 其中 $x=(x_1,x_2,\cdots,x_n)\in U(x_0,\delta)$, $u=(u_1,u_2,\cdots,u_m)\in U(u_0,\delta)$, 并且满足以下条件:
% 
% (1) $F_j(x_0,u_0)=0,\quad j=1,2,\cdots,m$;
% 
% (2) 对于 $j=1,2,\cdots,m$, $F_j(x,u)$ 以及它的各个偏导数在 $U(x_0,\delta)\times U(u_0,\delta)$ 内连续;
% 
% (3)
% \[
%   \left.
%   \frac{\partial(F_1,F_2,\cdots,F_m)}
%   {\partial(u_1,u_2,\cdots,u_m)}
%   \right|_{(x_0,u_0)}
%   \ne0,
% \]
% 
% 则 $\exists\delta_0(0<\delta_0<\delta)$, 使得在 $U(x_0,\delta_0)$ 内存在唯一的 $m$ 维 $n$ 元向量函数
% \[
%   f(x)=(f_1(x),f_2(x),\cdots,f_m(x)),
% \]
% 满足
% 
% (a) $u_0=(f_1(x_0),f_2(x_0),\cdots,f_m(x_0))$;
% 
% (b) 对于 $\forall j(1\leqslant j\leqslant m)$ 及 $\forall x\in U(x_0,\delta_0)$, 有
% \[
%   F_j(x,f_1(x),\cdots,f_m(x))=0;
% \]
% 
% (c) $f(x)$ 的每个分量函数 $f_j(x)(j=1,2,\cdots,m)$ 在 $U(x_0,\delta_0)$ 内存在连续偏导数, 记
% \[
%   A=
%   \begin{pmatrix}
%     \dfrac{\partial F_1(x,u)}{\partial x_1} &
%     \dfrac{\partial F_1(x,u)}{\partial x_2} &
%     \cdots &
%     \dfrac{\partial F_1(x,u)}{\partial x_n}\\
%     \dfrac{\partial F_2(x,u)}{\partial x_1} &
%     \dfrac{\partial F_2(x,u)}{\partial x_2} &
%     \cdots &
%     \dfrac{\partial F_2(x,u)}{\partial x_n}\\
%     \vdots & \vdots & & \vdots\\
%     \dfrac{\partial F_m(x,u)}{\partial x_1} &
%     \dfrac{\partial F_m(x,u)}{\partial x_2} &
%     \cdots &
%     \dfrac{\partial F_m(x,u)}{\partial x_n}
%   \end{pmatrix},
% \]
% \[
%   B=
%   \begin{pmatrix}
%     \dfrac{\partial F_1(x,u)}{\partial u_1} &
%     \dfrac{\partial F_1(x,u)}{\partial u_2} &
%     \cdots &
%     \dfrac{\partial F_1(x,u)}{\partial u_m}\\
%     \dfrac{\partial F_2(x,u)}{\partial u_1} &
%     \dfrac{\partial F_2(x,u)}{\partial u_2} &
%     \cdots &
%     \dfrac{\partial F_2(x,u)}{\partial u_m}\\
%     \vdots & \vdots & & \vdots\\
%     \dfrac{\partial F_m(x,u)}{\partial u_1} &
%     \dfrac{\partial F_m(x,u)}{\partial u_2} &
%     \cdots &
%     \dfrac{\partial F_m(x,u)}{\partial u_m}
%   \end{pmatrix},
% \]
% 那么
% \[
%   f'(x)=-B^{-1}\cdot A,
% \]
% 其中矩阵 $B^{-1}$ 为 $B$ 的逆矩阵.
% \end{theorem}
% 
% 
% \begin{proof}
% 为了证明本定理, 可以对 $m$ 作数学归纳法。对 $m=2$ 的证明是本质的, 而对 $m\geqslant3$ 的证明则只具有形式的复杂性。在此我们只给出 $m=2$ 的证明.
% 
% 由于
% \[
%   \left.
%   \frac{\partial(F_1,F_2)}{\partial(u_1,u_2)}
%   \right|_{(x_0,u_0)}
%   =
%   \left.
%   \begin{vmatrix}
%     \dfrac{\partial F_1}{\partial u_1} &
%     \dfrac{\partial F_1}{\partial u_2}\\
%     \dfrac{\partial F_2}{\partial u_1} &
%     \dfrac{\partial F_2}{\partial u_2}
%   \end{vmatrix}
%   \right|_{(x_0,u_0)}
%   \ne0,
% \]
% 不失一般性, 我们假定 $\dfrac{\partial F_2(x_0,u_0)}{\partial u_2}\ne0$。记 $u_0=(u_1^0,u_2^0)$。由定理~\ref{theorem:ma-14-4-2} 知, $\exists\delta_1(0<\delta_1<\delta)$, 使得方程 $F_2(x,u_1,u_2)=0$ 在 $U(x_0,\delta_1)\times U(u_1^0,\delta_1)$ 内唯一确定一个函数
% \[
%   u_2=g(x,u_1),
% \]
% 满足
% \[
%   F_2(x,u_1,g(x,u_1))=0,\qquad u_2^0=g(x_0,u_1^0),
% \]
% 且 $g(x,u_1)$ 在该邻域内具有各个连续偏导数。将 $u_2=g(x,u_1)$ 代入方程 $F_1(x,u_1,u_2)=0$, 得到一个关于 $(x,u_1)$ 的方程
% \[
%   H(x,u_1)\triangleq F_1(x,u_1,g(x,u_1))=0.
% \]
% 由于
% \[
% \begin{aligned}
%   \frac{\partial H}{\partial u_1}
%   &=\frac{\partial F_1}{\partial u_1}
%   +\frac{\partial F_1}{\partial u_2}\cdot
%   \frac{\partial g(x,u_1)}{\partial u_1}\\
%   &=\frac{\partial F_1}{\partial u_1}
%   +\frac{\partial F_1}{\partial u_2}
%   \left(-\frac{\partial F_2}{\partial u_1}\bigg/
%   \frac{\partial F_2}{\partial u_2}\right)\\
%   &=\frac1{\dfrac{\partial F_2}{\partial u_2}}\cdot
%   \frac{\partial(F_1,F_2)}{\partial(u_1,u_2)},
% \end{aligned}
% \]
% 我们有
% \[
%   \frac{\partial H(x_0,u_1^0)}{\partial u_1}
%   =\frac1{\dfrac{\partial F_2(x_0,u_1^0,u_2^0)}{\partial u_2}}\cdot
%   \left.
%   \frac{\partial(F_1,F_2)}{\partial(u_1,u_2)}
%   \right|_{(x_0,u_1^0,u_2^0)}
%   \ne0.
% \]
% 因此 $\exists\delta_0(0<\delta_0<\delta_1)$, 使得 $H(x,u_1)=0$ 在 $U(x_0,\delta_0)$ 内唯一确定一个函数 $u_1=f_1(x)$, 满足 $u_1^0=f_1(x_0)$, 且
% \[
%   F_1(x,f_1(x),g(x,f_1(x)))=0.
% \]
% 现记
% \[
%   u_1=f_1(x),\qquad u_2=f_2(x)=g(x,f_1(x)),
% \]
% 则在 $U(x_0,\delta_0)$ 内, 有
% \[
%   \begin{cases}
%     F_1(x,f_1(x),f_2(x))=0,\\
%     F_2(x,f_1(x),f_2(x))=0
%   \end{cases}
% \]
% 和 $u_0=(f_1(x_0),f_2(x_0))$。再由 $F(x,u)$ 具有各个连续偏导数, 我们可以推出 $f_1(x),f_2(x)$ 在 $U(x_0,\delta_0)$ 内具有各个连续偏导数 (从而 $f_1(x)$ 与 $f_2(x)$ 在该邻域内连续).
% 
% 对于 $\forall i(1\leqslant i\leqslant n)$, 利用复合函数的求导法则有
% \[
%   \begin{cases}
%     \dfrac{\partial F_1}{\partial x_i}
%     +\dfrac{\partial F_1}{\partial u_1}\cdot\dfrac{\partial f_1}{\partial x_i}
%     +\dfrac{\partial F_1}{\partial u_2}\cdot\dfrac{\partial f_2}{\partial x_i}=0,\\[1.2ex]
%     \dfrac{\partial F_2}{\partial x_i}
%     +\dfrac{\partial F_2}{\partial u_1}\cdot\dfrac{\partial f_1}{\partial x_i}
%     +\dfrac{\partial F_2}{\partial u_2}\cdot\dfrac{\partial f_2}{\partial x_i}=0.
%   \end{cases}
% \]
% 将上式写成向量形式有
% \[
%   \begin{pmatrix}
%     \dfrac{\partial F_1}{\partial u_1} &
%     \dfrac{\partial F_1}{\partial u_2}\\
%     \dfrac{\partial F_2}{\partial u_1} &
%     \dfrac{\partial F_2}{\partial u_2}
%   \end{pmatrix}
%   \begin{pmatrix}
%     \dfrac{\partial f_1}{\partial x_i}\\
%     \dfrac{\partial f_2}{\partial x_i}
%   \end{pmatrix}
%   =-
%   \begin{pmatrix}
%     \dfrac{\partial F_1}{\partial x_i}\\
%     \dfrac{\partial F_2}{\partial x_i}
%   \end{pmatrix},
% \]
% 因此有
% 
% \[
%   \begin{pmatrix}
%     \dfrac{\partial F_1}{\partial u_1} &
%     \dfrac{\partial F_1}{\partial u_2}\\
%     \dfrac{\partial F_2}{\partial u_1} &
%     \dfrac{\partial F_2}{\partial u_2}
%   \end{pmatrix}
%   \begin{pmatrix}
%     \dfrac{\partial f_1}{\partial x_1} &
%     \dfrac{\partial f_1}{\partial x_2} &
%     \cdots &
%     \dfrac{\partial f_1}{\partial x_n}\\
%     \dfrac{\partial f_2}{\partial x_1} &
%     \dfrac{\partial f_2}{\partial x_2} &
%     \cdots &
%     \dfrac{\partial f_2}{\partial x_n}
%   \end{pmatrix}
%   =-
%   \begin{pmatrix}
%     \dfrac{\partial F_1}{\partial x_1} &
%     \dfrac{\partial F_1}{\partial x_2} &
%     \cdots &
%     \dfrac{\partial F_1}{\partial x_n}\\
%     \dfrac{\partial F_2}{\partial x_1} &
%     \dfrac{\partial F_2}{\partial x_2} &
%     \cdots &
%     \dfrac{\partial F_2}{\partial x_n}
%   \end{pmatrix}.
% \]
% 由此得
% \[
%   \begin{pmatrix}
%     \dfrac{\partial f_1}{\partial x_1} &
%     \dfrac{\partial f_1}{\partial x_2} &
%     \cdots &
%     \dfrac{\partial f_1}{\partial x_n}\\
%     \dfrac{\partial f_2}{\partial x_1} &
%     \dfrac{\partial f_2}{\partial x_2} &
%     \cdots &
%     \dfrac{\partial f_2}{\partial x_n}
%   \end{pmatrix}
%   =-
%   \begin{pmatrix}
%     \dfrac{\partial F_1}{\partial u_1} &
%     \dfrac{\partial F_1}{\partial u_2}\\
%     \dfrac{\partial F_2}{\partial u_1} &
%     \dfrac{\partial F_2}{\partial u_2}
%   \end{pmatrix}^{-1}
%   \begin{pmatrix}
%     \dfrac{\partial F_1}{\partial x_1} &
%     \dfrac{\partial F_1}{\partial x_2} &
%     \cdots &
%     \dfrac{\partial F_1}{\partial x_n}\\
%     \dfrac{\partial F_2}{\partial x_1} &
%     \dfrac{\partial F_2}{\partial x_2} &
%     \cdots &
%     \dfrac{\partial F_2}{\partial x_n}
%   \end{pmatrix}.
% \]
% 现证所确定的隐函数组是唯一的。事实上, 倘若由
% \[
%   F(x,u_1,u_2)=0
% \]
% 在 $U(x_0,\delta_0)$ 内确定另一组函数
% \[
%   \begin{cases}
%     u_1=\tilde f_1(x),\\
%     u_2=\tilde f_2(x),
%   \end{cases}
% \]
% 并且它们仍然满足定理结论 (a), (b) 和 (c), 则容易看出 $\tilde f_1$ 与 $f_1$, $\tilde f_2$ 与 $f_2$ 在该邻域内具有相同的各个偏导数, 从而 $f_1(x)-\tilde f_1(x)$, $f_2(x)-\tilde f_2(x)$ 在 $U(x_0,\delta_0)$ 内均为常数函数。但它们在 $x_0$ 处有 $f_1(x_0)=\tilde f_1(x_0)$, $f_2(x_0)=\tilde f_2(x_0)$, 因此在 $U(x_0,\delta_0)$ 内有 $f_1(x)=\tilde f_1(x)$, $f_2(x)=\tilde f_2(x)$.
% \end{proof}
% 
% 
% \begin{example}\label{example:ma-14-4-3}
% 证明方程组
% \begin{equation}
%   \begin{cases}
%     x^2+y^2+u^2-v^2=2,\\
%     u+v+x+y=0
%   \end{cases}
%   \label{eq:ma-14-4-4}
% \end{equation}
% 在 $(x_0,y_0,u_0,v_0)=(1,1,-1,-1)$ 的某个邻域 $U((1,1,-1,-1),\delta_0)(\delta_0>0)$ 内确定隐函数组
% \[
%   \begin{cases}
%     u=u(x,y),\\
%     v=v(x,y),
%   \end{cases}
% \]
% 并在该邻域内求 $\dfrac{\partial u}{\partial x}$, $\dfrac{\partial v}{\partial y}$.
% \end{example}
% 
% 
% \begin{solve}
% 显然点 $x_0=(1,1,-1,-1)$ 满足方程组 \eqref{eq:ma-14-4-4}。记
% \[
%   F(x,y,u,v)=x^2+y^2+u^2-v^2-2,
%   \qquad
%   G(x,y,u,v)=u+v+x+y,
% \]
% 则 $F,G$ 都具有连续偏导数, 且
% \[
%   \left.
%   \frac{\partial(F,G)}{\partial(u,v)}
%   \right|_{x_0}
%   =
%   \left.
%   \begin{vmatrix}
%     2u & -2v\\
%     1 & 1
%   \end{vmatrix}
%   \right|_{x_0}
%   =-2-2=-4\ne0.
% \]
% 因此, $\exists\delta_0>0$, 使得方程组 \eqref{eq:ma-14-4-4} 在 $U((1,1,-1,-1),\delta_0)$ 内唯一确定隐函数组
% \[
%   \begin{cases}
%     u=u(x,y),\\
%     v=v(x,y),
%   \end{cases}
% \]
% 并且 $u(x,y)$ 与 $v(x,y)$ 在 $U((1,1),\delta_0)$ 内存在连续偏导数。下面我们来求偏导数。对方程组 \eqref{eq:ma-14-4-4} 微分得
% \begin{equation}
%   \begin{cases}
%     2x\,dx+2y\,dy+2u\,du-2v\,dv=0,\\
%     du+dv+dx+dy=0.
%   \end{cases}
%   \label{eq:ma-14-4-5}
% \end{equation}
% 令 $dy=0$, 得
% \[
%   \begin{cases}
%     u\,du-v\,dv=-x\,dx,\\
%     du+dv=-dx.
%   \end{cases}
% \]
% 从中解出 $du=-\dfrac{x+v}{u+v}\,dx$, 从而
% \[
%   \frac{\partial u}{\partial x}=-\frac{x+v}{u+v}.
% \]
% 在方程组 \eqref{eq:ma-14-4-5} 中令 $dx=0$, 即得
% \[
%   \begin{cases}
%     u\,du-v\,dv=-y\,dy,\\
%     du+dv=-dy,
%   \end{cases}
% \]
% 从而可求出
% \[
%   \frac{\partial v}{\partial y}=\frac{y-u}{u+v}.
% \qedhere
% \]
% \end{solve}
% 
% 
% \subsection{逆映射存在定理}
% \label{subsec:ma-14-4-3}
% 
% 对于一个可微映射, 如何确定它的逆映射的存在性呢? 隐函数组存在定理可以给出一个连续可微映射局部存在逆映射的充分条件.
% 
% \begin{theorem}[{逆映射存在定理}]\label{theorem:ma-14-4-4}
% 设
% \[
%   y=(y_1,y_2,\cdots,y_n)=(f_1(x),f_2(x),\cdots,f_n(x))=f(x)
% \]
% 是区域 $D\subset\mathbb R^n$ 到区域 $\Omega\subset\mathbb R^n$ 的一个 $C^1$ 映射, 并且在 $x_0\in D$ 处有
% \[
%   \left.
%   \frac{\partial(f_1,f_2,\cdots,f_n)}
%        {\partial(x_1,x_2,\cdots,x_n)}
%   \right|_{x_0}
%   \ne0.
% \]
% 记 $y_0=f(x_0)$, 则存在 $x_0$ 的邻域 $U(x_0,\delta_0)\subset D$, 使得映射 $y=f(x)$ 是 $U(x_0,\delta_0)$ 到 $f(U(x_0,\delta_0))$ 的 $C^1$ 同胚映射 (变换), 其中 $f(U(x_0,\delta_0))$ 是包含 $y_0$ 的一个区域.
% \end{theorem}
% 
% 
% \begin{proof}
% 对于 $j=1,2,\cdots,n$, 记 $F_j(x,y)=y_j-f_j(x)$, 并考虑下面 $n$ 个方程组成的方程组
% \[
%   \begin{cases}
%     F_1(x,y)=0,\\
%     F_2(x,y)=0,\\
%     \cdots\\
%     F_n(x,y)=0.
%   \end{cases}
% \]
% 由定理条件知它们满足
% 
% \noindent (1) 对任意的 $j(1\leqslant j\leqslant n)$, $F_j(x_0,y_0)=0$;
% 
% \noindent (2) $\exists\delta'>0$, 对任意的 $j(1\leqslant j\leqslant n)$, 使得 $F_j(x,y)$ 在 $U((x_0,y_0),\delta')$ 内具有连续偏导数 (因此 $F_j(x,y)$ 在该邻域连续), 并且
% \[
%   \left.
%   \frac{\partial(F_1,F_2,\cdots,F_n)}
%        {\partial(x_1,x_2,\cdots,x_n)}
%   \right|_{(x_0,y_0)}
%   =
%   (-1)^n
%   \left.
%   \frac{\partial(f_1,f_2,\cdots,f_n)}
%        {\partial(x_1,x_2,\cdots,x_n)}
%   \right|_{x_0}
%   \ne0,
% \]
% 因此由定理~\ref{theorem:ma-14-4-3} 知, 在 $y_0$ 的某个邻域 $U(y_0,\delta_1)(0<\delta_1<\delta')$ 内存在一个 $n$ 维向量函数
% \[
%   x=g(y)=(g_1(y),\cdots,g_n(y)),
% \]
% 满足 $x_0=(g_1(y_0),\cdots,g_n(y_0))$ 和
% \[
%   \begin{cases}
%     y_1-f_1(g_1(y),\cdots,g_n(y))=0,\\
%     \cdots\\
%     y_n-f_n(g_1(y),\cdots,g_n(y))=0,
%   \end{cases}
% \]
% 并且 $(g_1(y),\cdots,g_n(y))$ 在 $U(y_0,\delta_1)$ 内具有各个连续偏导数。上面的 $n$ 个等式同时表明, 在 $U(y_0,\delta_1)$ 内 $f(g(y))\equiv y$, 即 $x=g(y)$ 是 $y=f(x)$ 的逆映射。我们从定理~\ref{theorem:ma-14-4-3} 推知 $f'(x)g'(y)=E$ ($n$ 阶单位矩阵), 因此由 $|f'(x_0)|\ne0$ 有 $g'(y_0)=[f'(x_0)]^{-1}$, 从而 $g(y)$ 在 $y_0$ 处的雅可比行列式也不为零。如果将 $f(x)$ 与 $g(y)$ 的地位对换, 则在 $x_0$ 的某个邻域 $U(x_0,\delta_1')(0<\delta_1'<\delta')$ 内存在 $g(y)$ 的逆映射, 显然该逆映射是 $f(x)$.
% 
% 现证存在 $x_0$ 的邻域 $U(x_0,\delta_0)$, 使得 $f(x)$ 是 $U(x_0,\delta_0)$ 到区域 $f(U(x_0,\delta_0))$ 的 $C^1$ 同胚映射。首先取 $\delta_2(0<\delta_2<\delta_1)$ 充分小, 使得对于 $\forall x\in U(x_0,\delta_2)$, $f(x)$ 在 $x$ 处的雅可比行列式不等于零。下面我们证明, 对于 $\forall x'\in U(x_0,\delta_2)$ 及任何包含 $x'$ 的开集 $U_{x'}$, $f(x')$ 是 $f(U_{x'})$ 的内点。由于我们证明该事实只是利用 $f(x)$ 是 $C^1$ 映射, 并且它在 $x'$ 处的雅可比行列式不为零, 容易看出, 我们只要对 $x_0$ 证明上述结论即可。任取包含 $x_0$ 的开集 $U_{x_0}$, 则由 $g(y)$ 在 $y_0$ 处的连续性, 存在 $y_0$ 的邻域 $U(y_0,\delta_3)$ (当然是开集), 使得 $g(U(y_0,\delta_3))\subset U_{x_0}$。因此有
% \[
%   U(y_0,\delta_3)=f(g(U(y_0,\delta_3)))\subset f(U_{x_0}).
% \]
% 再注意到连续映射将连通集映成连通集, 因此 $f(U(x_0,\delta_2))$ 是一个包含 $y_0$ 的区域。同理, $\exists\delta_2'(0<\delta_2'<\delta_1')$, 使得 $g(U(y_0,\delta_2'))$ 是包含 $x_0$ 的一个区域。因此取 $\delta_0(0<\delta_0<\delta_2')$ 充分小, 使得 $f(U(x_0,\delta_0))\subset U(y_0,\delta_2')$, 则 $f(U(x_0,\delta_0))$ 是 $\mathbb R^n$ 中包含 $y_0$ 的区域。由于在 $U(x_0,\delta_0)$ 满足 $g(f(x))=x$, 并且 $f$ 在 $U(x_0,\delta_0)$ 内是 $C^1$ 映射, $g(y)=f^{-1}(y)$ 在 $f(U(x_0,\delta_0))$ 内是 $C^1$ 映射, 因此 $f(x)$ 是 $U(x_0,\delta_0)$ 到 $f(U(x_0,\delta_0))$ 的 $C^1$ 同胚映射 (变换).
% \end{proof}
% 
% 
% 从逆映射存在定理的证明可知, 如果一个 $n$ 维 $n$ 元向量函数
% \[
%   y=f(x)=(f_1(x),f_2(x),\cdots,f_n(x))
% \]
% 在区域 $D\subset\mathbb R^n$ 内具有连续偏导数, 并且对于 $\forall x\in D$, $f'(x)$ 非退化, 即它的雅可比行列式
% \[
%   \left.
%   \frac{\partial(f_1,f_2,\cdots,f_n)}
%        {\partial(x_1,x_2,\cdots,x_n)}
%   \right|_x
%   \ne0,
% \]
% 则局部总存在逆映射, 因此 $f$ 必将开集映成开集。另外, 由于此时 $y=f(x)$ 是连续函数, 因此推出 $f(D)$ 为一个区域.
% 
% \paragraph*{注 1}
% 当一个 $n$ 维 $n$ 元向量函数 $y=f(x)$ 在区域 $D\subset\mathbb R^n$ 内具有连续偏导数, 并且对于 $\forall x\in D$ 有 $|f'(x)|\ne0$ 时, 我们仅仅能推出 $y=f(x)$ 是局部同胚映射, 但不能推出它是 $D$ 到 $f(D)$ 的整体同胚映射。例如:
% \[
%   \begin{pmatrix}
%     u\\
%     v
%   \end{pmatrix}
%   =
%   \begin{pmatrix}
%     e^x\cos y\\
%     e^x\sin y
%   \end{pmatrix},
%   \qquad (x,y)\in\mathbb R^2
% \]
% 是 $\mathbb R^2$ 到 $\mathbb R^2\setminus\{(0,0)\}$ 的映射, 对于 $\forall(x,y)\in\mathbb R^2$, 有
% \[
%   \frac{\partial(u,v)}{\partial(x,y)}
%   =
%   \begin{vmatrix}
%     e^x\cos y & -e^x\sin y\\
%     e^x\sin y & e^x\cos y
%   \end{vmatrix}
%   =e^x\ne0,
% \]
% 因此该映射局部总是同胚映射。由 $\cos y$ 及 $\sin y$ 的周期性, 上述映射显然不是 $\mathbb R^2\to\mathbb R^2\setminus\{(0,0)\}$ 的同胚映射.
% 
% \paragraph*{注 2}
% 对于逆映射存在定理, 我们容易举例说明, 如果仅仅要求存在逆映射, 而不要求逆映射具有连续偏导数, 雅可比行列式不为零的条件不是必要条件。例如,
% \[
%   \begin{pmatrix}
%     u\\
%     v
%   \end{pmatrix}
%   =
%   \begin{pmatrix}
%     x^3\\
%     y
%   \end{pmatrix}
% \]
% 是 $\mathbb R^2\to\mathbb R^2$ 的同胚映射。事实上, 它的逆映射为
% \[
%   \begin{pmatrix}
%     x\\
%     y
%   \end{pmatrix}
%   =
%   \begin{pmatrix}
%     u^{1/3}\\
%     v
%   \end{pmatrix},
% \]
% 但在 $(0,0)$ 处有
% \[
%   \left.
%   \frac{\partial(u,v)}{\partial(x,y)}
%   \right|_{(0,0)}
%   =0.
% \]
% 此例同时说明了当 $\dfrac{\partial x}{\partial u}$ (或 $\dfrac{\partial y}{\partial u}$) 不是处处存在时, 仍然具有连续的逆映射.
% 
% \section{多元函数的极值}
% \label{sec:ma-14-5}
% 
% \subsection{通常极值问题}
% \label{subsec:ma-14-5-1}
% 
% 在一元函数中, 我们利用导数讨论了函数的极值。在本节中, 我们来讨论多元函数的极值问题。对于多元函数来说, 极值的定义与一元函数是相似的.
% 
% \begin{definition}\label{definition:ma-14-5-1}
% 设函数 $u=f(x)$ 在区域 $D$ 内有定义, $x_0\in D$, 若存在 $x_0$ 的邻域 $U(x_0,\delta_0)\subset D(\delta_0>0)$, 使得当 $x\in U(x_0,\delta_0)$ 时, 有 $f(x)\leqslant f(x_0)$ (或 $f(x)\geqslant f(x_0)$), 则称 $f(x)$ 在 $x_0$ 处取极大值 (或极小值), $x_0$ 称为 $f(x)$ 的极大值点 (或极小值点)。如果上述不等式为严格不等式, 则称 $x_0$ 为 $f(x)$ 的严格极大值点 (或严格极小值点).
% \end{definition}
% 
% 对于一个二元可微函数 $z=f(x,y)((x,y)\in D)$, 一般来说它的图像是一块曲面。从几何上来看, 若 $(x_0,y_0)\in D$ 是 $z=f(x,y)$ 的极值点, 则 $(x_0,y_0,f(x_0,y_0))$ 是曲面上局部的最高点或最低点。显然, 对于曲面上的任何一条过 $(x_0,y_0,f(x_0,y_0))$ 的曲线 $\Gamma$, 该点也是 $\Gamma$ 在它附近的最高点或最低点。当 $\Gamma$ 在该点具有切线时, 它的切线应该平行于 $Oxy$ 平面。所以, 当 $z=f(x,y)$ 在 $(x_0,y_0)$ 处可微时, 有
% \[
%   \frac{\partial f(x_0,y_0)}{\partial x}
%   =
%   \frac{\partial f(x_0,y_0)}{\partial y}
%   =0,
% \]
% 即 $f'(x_0,y_0)=\operatorname{grad} f(x_0)=0$.
% 
% \begin{theorem}\label{theorem:ma-14-5-1}
% 设函数 $f(x)$ 在区域 $D\subset\mathbb R^n$ 内有定义, $x_0=(x_1^0,x_2^0,\cdots,x_n^0)\in D$, 再设 $f(x)$ 在 $x_0$ 处取极值并且 $f(x)$ 在该点关于 $x_i(1\leqslant i\leqslant n)$ 可偏导, 则有
% \[
%   \frac{\partial f(x_0)}{\partial x_i}=0.
% \]
% 特别地, 若 $f(x)$ 在该点可微, 则 $f'(x_0)=0$.
% \end{theorem}
% 
% 
% \begin{proof}
% 在定理的条件下, 一元函数
% \[
%   F_i(x_i)=f(x_1^0,\cdots,x_{i-1}^0,x_i,x_{i+1}^0,\cdots,x_n^0)
%   \qquad (1\leqslant i\leqslant n)
% \]
% 在 $x_i^0$ 处取极值, 因此
% \[
%   F_i'(x_i^0)=\frac{\partial f(x_0)}{\partial x_i}=0
%   \qquad (1\leqslant i\leqslant n).
% \qedhere
% \]
% \end{proof}
% 
% 
% 如同一元函数, 对于一个 $n$ 元可微函数 $f(x)$, 若 $f'(x_0)=0$, 则称 $x_0$ 为 $f(x)$ 的一个驻点或临界点。因此, 对可微函数 $f(x)$ 而言, $f(x)$ 的极值点一定是驻点, 但反命题一般不真.
% 
% 例如, $f(x,y)=y^2-x^2$ 在 $(0,0)$ 处显然不取极值。(如图 14.5.1 所示, 该函数的图像在原点附近有点像一个马鞍, 因此人们称之为马鞍面)。当一个可微函数 $f(x)$ 的驻点不是极值点时, 该驻点也称为 $f(x)$ 的一个鞍点.
% 
% 在研究一元函数的极值问题中, 我们曾经利用它的高阶导数来判别驻点是否为极值点。对于多元函数, 我们也可以利用高阶偏导数来讨论其极值问题。由于多元函数高阶偏导数的复杂性, 我们只考虑涉及二阶偏导数的情形.
% 
% \begin{theorem}\label{theorem:ma-14-5-2}
% 设函数 $f(x)=f(x_1,x_2,\cdots,x_n)$ 在区域 $D\subset\mathbb R^n$ 内具有二阶连续偏导数, 且 $f'(x_0)=0(x_0\in D)$, 再设 $f(x)$ 在 $x_0$ 处的海色矩阵 $H_f(x_0)$ 为满秩矩阵, 则
% \end{theorem}
% 
% \noindent (1) 当 $H_f(x_0)$ 正定时, $f(x)$ 在 $x_0$ 处取极小值;
% 
% \noindent (2) 当 $H_f(x_0)$ 负定时, $f(x)$ 在 $x_0$ 处取极大值;
% 
% \noindent (3) 当 $H_f(x_0)$ 不定时, $f(x)$ 在 $x_0$ 处不取极值.
% 
% \begin{proof}
% 在 $x_0$ 的充分小邻域 $U(x_0,\delta_0)(\delta_0>0)$ 内, $f(x)$ 有泰勒公式
% \[
% \begin{aligned}
%   f(x)
%   &=f(x_0)+f'(x_0)(x-x_0)^{\mathrm T}
%     +\frac12(x-x_0)H_f(x_0)(x-x_0)^{\mathrm T}
%     +o(|x-x_0|^2)\\
%   &=f(x_0)+\frac12
%     \frac{x-x_0}{|x-x_0|}
%     H_f(x_0)
%     \frac{(x-x_0)^{\mathrm T}}{|x-x_0|}
%     |x-x_0|^2
%     +o(|x-x_0|^2)
%     \qquad (|x-x_0|\to0).
% \end{aligned}
% \]
% 对于 $\forall x\in U_0(x_0,\delta_0)$, 记 $x'=\dfrac{x-x_0}{|x-x_0|}$, 则 $|x'|=1$。由于 $x'H_f(x_0)x'^{\mathrm T}$ 是 $S(0;1)=\{x':|x'|=1\}$ 上的一个连续函数, 注意到 $S(0;1)$ 是紧集, $x'H_f(x_0)x'^{\mathrm T}$ 在 $S(0;1)$ 上取到最大值 $M$ 和最小值 $m$.
% 
% \noindent (1) 当 $H_f(x_0)$ 正定时, 对任意非零向量 $x-x_0$ 均有
% \[
%   (x-x_0)H_f(x_0)(x-x_0)^{\mathrm T}>0,
% \]
% 因此有 $m>0$, 从而存在 $0<\delta_1<\delta_0$, 使得当 $x\in U_0(x_0,\delta_1)$ 时, 有
% \[
% \begin{aligned}
%   f(x)
%   &=f(x_0)+\frac12
%     \frac{x-x_0}{|x-x_0|}
%     H_f(x_0)
%     \frac{(x-x_0)^{\mathrm T}}{|x-x_0|}
%     |x-x_0|^2+o(|x-x_0|^2)\\
%   &>f(x_0)+\frac m4|x-x_0|^2>f(x_0).
% \end{aligned}
% \]
% 这就证明了 $f(x)$ 在 $x_0$ 处取极小值.
% 
% \noindent (2) 当 $H_f(x_0)$ 负定时, 有 $M<0$。因此, 存在 $0<\delta_2<\delta_0$, 当 $x\in U_0(x_0,\delta_2)$ 时, 有
% \[
% \begin{aligned}
%   f(x)
%   &=f(x_0)+\frac12
%     \frac{x-x_0}{|x-x_0|}
%     H_f(x_0)
%     \frac{(x-x_0)^{\mathrm T}}{|x-x_0|}
%     |x-x_0|^2+o(|x-x_0|^2)\\
%   &<f(x_0)+\frac M4|x-x_0|^2<f(x_0).
% \end{aligned}
% \]
% 这就证明了 $f(x)$ 在 $x_0$ 取极大值.
% 
% \noindent (3) 当 $H_f(x_0)$ 不定时, 有 $m<0$, $M>0$。任取 $x_1\in\mathbb R^n$, 使得
% \[
%   \frac{x_1-x_0}{|x_1-x_0|}
%   H_f(x_0)
%   \frac{(x_1-x_0)^{\mathrm T}}{|x_1-x_0|}
%   =m,
% \]
% 则对充分小的 $t>0$, 有
% \[
%   f\left(x_0+t\frac{x_1-x_0}{|x_1-x_0|}\right)
%   =f(x_0)+\frac12mt^2+o(t^2)
%   \leqslant f(x_0)+\frac14mt^2<f(x_0).
% \]
% 再取 $x_2\in\mathbb R^n$, 使得
% \[
%   \frac{x_2-x_0}{|x_2-x_0|}
%   H_f(x_0)
%   \frac{(x_2-x_0)^{\mathrm T}}{|x_2-x_0|}
%   =M,
% \]
% 则对充分小的 $t>0$, 有
% \[
%   f\left(x_0+t\frac{x_2-x_0}{|x_2-x_0|}\right)
%   =f(x_0)+\frac12Mt^2+o(t^2)
%   \geqslant f(x_0)+\frac14Mt^2>f(x_0).
% \]
% 这说明 $f(x)$ 在 $x_0$ 处不取极值.
% \end{proof}
% 
% 
% 在代数课程中, 有关于 $H_f(x_0)$ 正定、负定和不定的判定法则, 在这里我们就不再赘述了。我们经常用到 $n=2$ 的情况, 下面给出定理~\ref{theorem:ma-14-5-2} 在 $\mathbb R^2$ 情形下的特殊形式.
% 
% \begin{corollary}\label{corollary:ma-14-source-101}
% 设函数 $f(x,y)$ 在区域 $D\subset\mathbb R^2$ 内具有二阶连续偏导数, $(x_0,y_0)\in D$, 再设 $f'(x_0,y_0)=0$, 且记
% \[
%   H_f(x_0,y_0)
%   =
%   \left.
%   \begin{pmatrix}
%     \dfrac{\partial^2 f}{\partial x^2} &
%     \dfrac{\partial^2 f}{\partial x\partial y}\\
%     \dfrac{\partial^2 f}{\partial x\partial y} &
%     \dfrac{\partial^2 f}{\partial y^2}
%   \end{pmatrix}
%   \right|_{(x_0,y_0)}
%   \triangleq
%   \begin{pmatrix}
%     A & B\\
%     B & C
%   \end{pmatrix},
% \]
% 则
% \end{corollary}
% 
% \noindent (1) 当 $A>0$,
% \[
%   \begin{vmatrix}
%     A & B\\
%     B & C
%   \end{vmatrix}
%   =AC-B^2>0
% \]
% 时, $f(x,y)$ 在 $(x_0,y_0)$ 处取极小值;
% 
% \noindent (2) 当 $A<0$,
% \[
%   \begin{vmatrix}
%     A & B\\
%     B & C
%   \end{vmatrix}
%   =AC-B^2>0
% \]
% 时, $f(x,y)$ 在 $(x_0,y_0)$ 处取极大值;
% 
% \noindent (3) 当
% \[
%   \begin{vmatrix}
%     A & B\\
%     B & C
%   \end{vmatrix}
%   =AC-B^2<0
% \]
% 时, $f(x,y)$ 在 $(x_0,y_0)$ 处不取极值.
% 
% \begin{example}\label{example:ma-14-5-1}
% 试求函数 $f(x,y)=x^2+3xy+3y^2-6x-3y$ 的极值.
% \end{example}
% 
% 
% \begin{solve}
% 由
% \[
%   f_x'(x,y)=2x+3y-6=0,\qquad
%   f_y'(x,y)=3x+6y-3=0
% \]
% 得唯一驻点 $(9,-4)$, 再由于
% \[
%   f_{xx}''(x,y)=2,\qquad f_{xy}''(x,y)=3,\qquad f_{yy}''(x,y)=6,
% \]
% 从而
% \[
%   f_{xx}''(9,-4)=2>0,
% \]
% \[
%   \begin{vmatrix}
%     f_{xx}''(9,-4) & f_{xy}''(9,-4)\\
%     f_{xy}''(9,-4) & f_{yy}''(9,-4)
%   \end{vmatrix}
%   =
%   \begin{vmatrix}
%     2 & 3\\
%     3 & 6
%   \end{vmatrix}
%   =3>0,
% \]
% 因此 $f(x,y)$ 在 $(9,-4)$ 处取极小值 $-21$。因为 $f(x,y)$ 在 $\mathbb R^2$ 处处可微, 所以 $(9,-4)$ 是 $f(x,y)$ 的唯一极值点.
% \end{solve}
% 
% 
% \begin{example}\label{example:ma-14-5-2}
% 某工厂生产一批长方体无盖盒子, 要求其体积为 $1\mathrm{m}^3$, 盒子的底的厚度是侧面厚度的三倍。问: 如何设定盒子的长、宽、高才能使得用料最省?
% \end{example}
% 
% 
% \begin{solve}
% 设盒子的长、宽、高分别为 $x,y,z$, 并令
% \[
%   F(x,y,z)=3xy+2yz+2xz.
% \]
% 由于盒子的体积为 $1\mathrm{m}^3$, 即 $xyz=1$, 因此, 若令
% \[
%   S(x,y)=3xy+2\left(\frac1x+\frac1y\right),
% \]
% 则本题的最值问题转化成求 $S(x,y)$ 的最小值问题。由
% \[
%   \frac{\partial S(x,y)}{\partial x}=3y-\frac2{x^2}=0,\qquad
%   \frac{\partial S(x,y)}{\partial y}=3x-\frac2{y^2}=0
% \]
% 解得
% \[
%   x=y=\sqrt[3]{\frac23}.
% \]
% 由于该实际问题总是存在最小值, 且 $\left(\sqrt[3]{\dfrac23},\sqrt[3]{\dfrac23}\right)$ 是 $S(x,y)$ 唯一的驻点, 因此它是 $S(x,y)$ 的极小值点, 且 $S(x,y)$ 必在该点取最小值。所以当
% \[
%   x=y=\sqrt[3]{\frac23},\qquad z=\sqrt[3]{\frac94}
% \]
% 时可使用料最省.
% 
% 我们也可以由 $S(x,y)$ 的海色矩阵的正定性来判断该点为极小值点。事实上, 在 $\left(\sqrt[3]{\dfrac23},\sqrt[3]{\dfrac23}\right)$ 处, $S(x,y)$ 的海色矩阵为
% \[
%   H_S\left(\sqrt[3]{\frac23},\sqrt[3]{\frac23}\right)
%   =
%   \begin{pmatrix}
%     6 & 3\\
%     3 & 6
%   \end{pmatrix}.
% \]
% 容易看出它是正定的, 从而 $\left(\sqrt[3]{\dfrac23},\sqrt[3]{\dfrac23}\right)$ 为 $S(x,y)$ 的极小值点.
% \end{solve}
% 
% 
% \begin{example}\label{example:ma-14-5-3}
% 设 $(x_k,y_k)(k=1,2,\cdots,K)$ 为 $\mathbb R^2$ 内 $K$ 个两两不同的点, 且它们不同在 $Oxy$ 平面内的一条垂直于 $x$ 轴的直线上。证明: 存在唯一直线 $L_0:y=ax+b$, 使得函数
% \[
%   f(s,t)=\sum_{k=1}^K(sx_k+t-y_k)^2
% \]
% 在 $(a,b)$ 处达到最小.
% \end{example}
% 
% 
% \begin{proof}
% 显然 $f(s,t)=\sum\limits_{k=1}^K(sx_k+t-y_k)^2$ 在 $\mathbb R^2$ 中处处可微。由
% \[
%   \begin{cases}
%     \dfrac{\partial f}{\partial s}
%     =2\displaystyle\sum_{k=1}^K x_k(sx_k+t-y_k)=0,\\[1.2ex]
%     \dfrac{\partial f}{\partial t}
%     =2\displaystyle\sum_{k=1}^K (sx_k+t-y_k)=0
%   \end{cases}
% \]
% 得
% \[
%   \begin{cases}
%     s\displaystyle\sum_{k=1}^K x_k^2
%     +t\displaystyle\sum_{k=1}^K x_k
%     =\displaystyle\sum_{k=1}^K x_ky_k,\\[1.2ex]
%     s\displaystyle\sum_{k=1}^K x_k
%     +Kt
%     =\displaystyle\sum_{k=1}^K y_k.
%   \end{cases}
% \]
% 由此求出上述方程组的解为
% \[
%   s=
%   \frac{
%     K\displaystyle\sum_{k=1}^K x_ky_k
%     -\left(\displaystyle\sum_{k=1}^K x_k\right)
%      \left(\displaystyle\sum_{k=1}^K y_k\right)}
%   {
%     K\displaystyle\sum_{k=1}^K x_k^2
%     -\left(\displaystyle\sum_{k=1}^K x_k\right)^2}
%   \triangleq a,
% \]
% \[
%   t=
%   \frac{
%     \left(\displaystyle\sum_{k=1}^K x_k^2\right)
%     \left(\displaystyle\sum_{k=1}^K y_k\right)
%     -\left(\displaystyle\sum_{k=1}^K x_k\right)
%      \left(\displaystyle\sum_{k=1}^K x_ky_k\right)}
%   {
%     K\displaystyle\sum_{k=1}^K x_k^2
%     -\left(\displaystyle\sum_{k=1}^K x_k\right)^2}
%   \triangleq b.
% \]
% 由于该问题中的最小值总是存在的, 且 $f(s,t)$ 只有唯一的驻点 $(a,b)$, 因此 $(a,b)$ 必为 $f(s,t)$ 的最小值点.
% 
% 在上述例子中, 已知平面 $\mathbb R^2$ 内 $K$ 个不落在一条垂直于 $x$ 轴的直线上的点, 求满足定理条件的唯一直线 $y=ax+b$ 的方法称为是最小二乘法.
% \end{proof}
% 
% 
% \subsection{条件极值问题}
% \label{subsec:ma-14-5-2}
% 
% 为了导出条件极值问题, 我们先来看一个例子。设 $\Gamma\subset\mathbb R^3$ 是一条曲线, 现在要求它到原点的最短距离。为了求解此问题, 我们可令 $f(x,y,z)=\sqrt{x^2+y^2+z^2}$, 该函数表示了 $\mathbb R^3$ 中任一点到原点的距离。于是该问题转化为: 当 $(x,y,z)$ 在 $\Gamma$ 上变化时, 求 $f(x,y,z)$ 的最小值。实际上, 这是 $(x,y,z)$ 在 $\Gamma$ 上变化的条件下, 求 $f(x,y,z)$ 的极小值问题。像这样具有外加约束条件的极值问题称为条件极值问题。为了求解这一条件极值问题, 一个自然的想法是: 求出 $\Gamma$ 的参数方程然后代入 $f(x,y,z)$, 从而将此问题转化成上一小节讨论的极值问题。但在很多情形下, 这绝非易事。因此我们需用其他方法来求解条件极值问题.
% 
% 首先, 设 $\Gamma$ 是一条平面曲线, 它由方程 $\varphi(x,y)=0$ 所确定。我们假定 $\varphi(x,y)$ 在区域 $D\subset\mathbb R^2$ 内具有连续偏导数, 且对于 $\forall(x,y)\in D$,
% \[
%   \left.
%   \left(\frac{\partial\varphi}{\partial x},\frac{\partial\varphi}{\partial y}\right)
%   \right|_{(x,y)}
% \]
% 的秩为 $1$, 再设 $(x_0,y_0)\in\Gamma$ 且是函数 $f(x,y)=\sqrt{x^2+y^2}$ 的一个极值点。在理论上 (由隐函数存在定理) 我们可以从 $\varphi(x,y)=0$ 在 $(x_0,y_0)$ 附近解出 $y=y(x)$ 或 $x=x(y)$。现在不妨设
% \[
%   \frac{\partial\varphi(x_0,y_0)}{\partial y}\ne0,
% \]
% 从而 $\exists\delta>0$, 使得在 $U(x_0,\delta)$ 内存在 $y=y(x)$, 满足 $y_0=y(x_0)$ 且 $\varphi(x,y(x))=0$。将 $y=y(x)$ 代入 $f(x,y)$ 得 $f(x,y(x))$, 则 $x_0$ 是 $f(x,y(x))$ 的一个通常极值点, 从而有
% \[
%   \frac{\partial f(x_0,y_0)}{\partial x}
%   +
%   \frac{\partial f(x_0,y_0)}{\partial y}y'(x_0)=0.
% \]
% 由
% \[
%   y'(x_0)=-\frac{\varphi_x'(x_0,y_0)}{\varphi_y'(x_0,y_0)}
% \]
% 知
% \[
%   \frac{\dfrac{\partial f(x_0,y_0)}{\partial x}}
%        {\dfrac{\partial\varphi(x_0,y_0)}{\partial x}}
%   =
%   \frac{\dfrac{\partial f(x_0,y_0)}{\partial y}}
%        {\dfrac{\partial\varphi(x_0,y_0)}{\partial y}},
% \]
% 即 $\exists\lambda\in\mathbb R$, 使得
% \[
%   \begin{cases}
%     \dfrac{\partial f(x_0,y_0)}{\partial x}
%     +\lambda\dfrac{\partial\varphi(x_0,y_0)}{\partial x}=0,\\[1.2ex]
%     \dfrac{\partial f(x_0,y_0)}{\partial y}
%     +\lambda\dfrac{\partial\varphi(x_0,y_0)}{\partial y}=0.
%   \end{cases}
% \]
% 再设 $\Gamma$ 是由下述方程组确定的一条空间曲线:
% \[
%   \begin{cases}
%     \varphi_1(x,y,z)=0,\\
%     \varphi_2(x,y,z)=0,
%   \end{cases}
% \]
% 并假定 $\varphi_1,\varphi_2$ 在区域 $D\subset\mathbb R^3$ 内具有连续偏导数, 且对于 $\forall(x,y,z)\in D$,
% \[
%   \left.
%   \begin{pmatrix}
%     \dfrac{\partial\varphi_1}{\partial x} &
%     \dfrac{\partial\varphi_1}{\partial y} &
%     \dfrac{\partial\varphi_1}{\partial z}\\
%     \dfrac{\partial\varphi_2}{\partial x} &
%     \dfrac{\partial\varphi_2}{\partial y} &
%     \dfrac{\partial\varphi_2}{\partial z}
%   \end{pmatrix}
%   \right|_{(x,y,z)}
% \]
% 的秩为 $2$。由此条件, 理论上从方程组我们仍然可解出曲线的参数方程, 我们将其代入 $f(x,y,z)$ 就可以推出在 $f(x,y,z)$ 的极值点 $(x_0,y_0,z_0)$ 处, 存在常数 $\lambda_1,\lambda_2$ 成立下述方程组:
% \[
%   \begin{cases}
%     \dfrac{\partial f(x_0,y_0,z_0)}{\partial x}
%     +\lambda_1\dfrac{\partial\varphi_1(x_0,y_0,z_0)}{\partial x}
%     +\lambda_2\dfrac{\partial\varphi_2(x_0,y_0,z_0)}{\partial x}=0,\\[1.2ex]
%     \dfrac{\partial f(x_0,y_0,z_0)}{\partial y}
%     +\lambda_1\dfrac{\partial\varphi_1(x_0,y_0,z_0)}{\partial y}
%     +\lambda_2\dfrac{\partial\varphi_2(x_0,y_0,z_0)}{\partial y}=0,\\[1.2ex]
%     \dfrac{\partial f(x_0,y_0,z_0)}{\partial z}
%     +\lambda_1\dfrac{\partial\varphi_1(x_0,y_0,z_0)}{\partial z}
%     +\lambda_2\dfrac{\partial\varphi_2(x_0,y_0,z_0)}{\partial z}=0.
%   \end{cases}
% \]
% 对于一般情形, 我们可以证明下面的定理.
% 
% \begin{theorem}\label{theorem:ma-14-5-3}
% 设函数 $f(x)$, $\varphi(x)=(\varphi_1(x),\varphi_2(x),\cdots,\varphi_m(x))$ 在区域 $D\subset\mathbb R^n(m<n)$ 内具有各个连续偏导数, 再设 $x_0=(x_1^0,x_2^0,\cdots,x_n^0)\in D$ 为 $f(x)$ 在约束条件
% \begin{equation}
%   \begin{cases}
%     \varphi_1(x)=0,\\
%     \varphi_2(x)=0,\\
%     \cdots\\
%     \varphi_m(x)=0
%   \end{cases}
%   \label{eq:ma-14-5-1}
% \end{equation}
% 下的极值点, 并且 $\varphi'(x_0)$ 的秩为 $m$, 则存在常数 $\lambda_1,\lambda_2,\cdots,\lambda_m\in\mathbb R$, 使得在 $x_0$ 处成立下述等式:
% \begin{equation}
%   \begin{cases}
%     \dfrac{\partial f(x_0)}{\partial x_i}
%     +\displaystyle\sum_{j=1}^m\lambda_j
%     \dfrac{\partial\varphi_j(x_0)}{\partial x_i}=0,
%     \qquad i=1,2,\cdots,n,\\[1.2ex]
%     \varphi_j(x_0)=0,\qquad j=1,2,\cdots,m.
%   \end{cases}
%   \label{eq:ma-14-5-2}
% \end{equation}
% \end{theorem}
% 
% 
% \begin{proof}
% 由于 $\varphi'(x_0)$ 的秩为 $m$, 我们不妨设行列式
% \[
%   \left.
%   \frac{\partial(\varphi_1,\varphi_2,\cdots,\varphi_m)}
%        {\partial(x_{n-m+1},x_{n-m+2},\cdots,x_n)}
%   \right|_{x_0}
%   \ne0.
% \]
% 因此, 在 $x_0$ 的某个邻域内唯一确定一组具有各个连续偏导数的隐函数
% \begin{equation}
%   \begin{cases}
%     x_{n-m+1}=g_1(x_1,x_2,\cdots,x_{n-m}),\\
%     x_{n-m+2}=g_2(x_1,x_2,\cdots,x_{n-m}),\\
%     \cdots\\
%     x_n=g_m(x_1,x_2,\cdots,x_{n-m}),
%   \end{cases}
%   \label{eq:ma-14-5-3}
% \end{equation}
% 满足
% \[
%   x_j^0=g_j(x_1^0,x_2^0,\cdots,x_{n-m}^0)
%   \qquad (j=n-m+1,n-m+2,\cdots,n)
% \]
% 及对 $k=1,2,\cdots,m$, 有
% \[
%   \varphi_k(x_1,\cdots,x_{n-m},g_1(x_1,x_2,\cdots,x_{n-m}),\cdots,g_m(x_1,x_2,\cdots,x_{n-m}))=0.
% \]
% 将隐函数组 \eqref{eq:ma-14-5-3} 代入 $f(x)$, 得
% \[
%   f(x_1,\cdots,x_{n-m},g_1(x_1,x_2,\cdots,x_{n-m}),\cdots,g_m(x_1,x_2,\cdots,x_{n-m})).
% \]
% 因此, $x_0$ 是条件极值点转化为 $(x_1^0,x_2^0,\cdots,x_{n-m}^0)$ 为上述函数的通常极值点.
% 
% 令 $x_0'=(x_1^0,x_2^0,\cdots,x_{n-m}^0)$, 则对 $i=1,2,\cdots,n-m$, 有
% \[
%   \frac{\partial f(x_0)}{\partial x_i}
%   +\frac{\partial f(x_0)}{\partial x_{n-m+1}}\cdot
%   \frac{\partial g_1(x_0')}{\partial x_i}
%   +\cdots+
%   \frac{\partial f(x_0)}{\partial x_n}\cdot
%   \frac{\partial g_m(x_0')}{\partial x_i}
%   =0.
% \]
% 令 $g(x')=(g_1(x'),g_2(x'),\cdots,g_m(x'))^{\mathrm T}$, 其中 $x'=(x_1,x_2,\cdots,x_{n-m})$。将上述 $n-m$ 个等式写成向量形式, 有
% \begin{equation}
%   \left(
%     \frac{\partial f(x_0)}{\partial x_1},
%     \cdots,
%     \frac{\partial f(x_0)}{\partial x_{n-m}}
%   \right)
%   +
%   \left(
%     \frac{\partial f(x_0)}{\partial x_{n-m+1}},
%     \cdots,
%     \frac{\partial f(x_0)}{\partial x_n}
%   \right)
%   g'(x_0')=0.
%   \label{eq:ma-14-5-4}
% \end{equation}
% 由于
% \begin{equation}
% \begin{aligned}
%   g'(x_0')
%   &=
%   -
%   \begin{pmatrix}
%     \dfrac{\partial\varphi_1(x_0)}{\partial x_{n-m+1}} &
%     \dfrac{\partial\varphi_1(x_0)}{\partial x_{n-m+2}} &
%     \cdots &
%     \dfrac{\partial\varphi_1(x_0)}{\partial x_n}\\
%     \dfrac{\partial\varphi_2(x_0)}{\partial x_{n-m+1}} &
%     \dfrac{\partial\varphi_2(x_0)}{\partial x_{n-m+2}} &
%     \cdots &
%     \dfrac{\partial\varphi_2(x_0)}{\partial x_n}\\
%     \vdots & \vdots & & \vdots\\
%     \dfrac{\partial\varphi_m(x_0)}{\partial x_{n-m+1}} &
%     \dfrac{\partial\varphi_m(x_0)}{\partial x_{n-m+2}} &
%     \cdots &
%     \dfrac{\partial\varphi_m(x_0)}{\partial x_n}
%   \end{pmatrix}^{-1}\\
%   &\quad\cdot
%   \begin{pmatrix}
%     \dfrac{\partial\varphi_1(x_0)}{\partial x_1} &
%     \dfrac{\partial\varphi_1(x_0)}{\partial x_2} &
%     \cdots &
%     \dfrac{\partial\varphi_1(x_0)}{\partial x_{n-m}}\\
%     \dfrac{\partial\varphi_2(x_0)}{\partial x_1} &
%     \dfrac{\partial\varphi_2(x_0)}{\partial x_2} &
%     \cdots &
%     \dfrac{\partial\varphi_2(x_0)}{\partial x_{n-m}}\\
%     \vdots & \vdots & & \vdots\\
%     \dfrac{\partial\varphi_m(x_0)}{\partial x_1} &
%     \dfrac{\partial\varphi_m(x_0)}{\partial x_2} &
%     \cdots &
%     \dfrac{\partial\varphi_m(x_0)}{\partial x_{n-m}}
%   \end{pmatrix}
%   \triangleq -A^{-1}B.
% \end{aligned}
% \label{eq:ma-14-5-5}
% \end{equation}
% 注意到
% \[
%   -
%   \left(
%     \frac{\partial f(x_0)}{\partial x_{n-m+1}},
%     \frac{\partial f(x_0)}{\partial x_{n-m+2}},
%     \cdots,
%     \frac{\partial f(x_0)}{\partial x_n}
%   \right)\cdot A^{-1}
% \]
% 是一个 $m$ 维行向量, 我们可以将其记为
% \begin{equation}
%   -
%   \left(
%     \frac{\partial f(x_0)}{\partial x_{n-m+1}},
%     \frac{\partial f(x_0)}{\partial x_{n-m+2}},
%     \cdots,
%     \frac{\partial f(x_0)}{\partial x_n}
%   \right)\cdot A^{-1}
%   =(\lambda_1,\lambda_2,\cdots,\lambda_m).
%   \label{eq:ma-14-5-6}
% \end{equation}
% 将式 \eqref{eq:ma-14-5-5} 与 \eqref{eq:ma-14-5-6} 代入式 \eqref{eq:ma-14-5-4} 得
% \begin{equation}
% \begin{aligned}
%   &\left(
%     \frac{\partial f(x_0)}{\partial x_1},
%     \frac{\partial f(x_0)}{\partial x_2},
%     \cdots,
%     \frac{\partial f(x_0)}{\partial x_{n-m}}
%   \right)\\
%   &\quad+
%   (\lambda_1,\lambda_2,\cdots,\lambda_m)
%   \begin{pmatrix}
%     \dfrac{\partial\varphi_1(x_0)}{\partial x_1} &
%     \dfrac{\partial\varphi_1(x_0)}{\partial x_2} &
%     \cdots &
%     \dfrac{\partial\varphi_1(x_0)}{\partial x_{n-m}}\\
%     \dfrac{\partial\varphi_2(x_0)}{\partial x_1} &
%     \dfrac{\partial\varphi_2(x_0)}{\partial x_2} &
%     \cdots &
%     \dfrac{\partial\varphi_2(x_0)}{\partial x_{n-m}}\\
%     \vdots & \vdots & & \vdots\\
%     \dfrac{\partial\varphi_m(x_0)}{\partial x_1} &
%     \dfrac{\partial\varphi_m(x_0)}{\partial x_2} &
%     \cdots &
%     \dfrac{\partial\varphi_m(x_0)}{\partial x_{n-m}}
%   \end{pmatrix}
%   =0.
% \end{aligned}
% \label{eq:ma-14-5-7}
% \end{equation}
% 另外, 我们可以将式 \eqref{eq:ma-14-5-6} 改写成下述形式
% \begin{equation}
%   \left(
%     \frac{\partial f(x_0)}{\partial x_{n-m+1}},
%     \frac{\partial f(x_0)}{\partial x_{n-m+2}},
%     \cdots,
%     \frac{\partial f(x_0)}{\partial x_n}
%   \right)
%   +
%   (\lambda_1,\lambda_2,\cdots,\lambda_m)
%   A=0.
%   \label{eq:ma-14-5-8}
% \end{equation}
% 最后将向量方程 \eqref{eq:ma-14-5-7} 与 \eqref{eq:ma-14-5-8} 写成分量形式的方程, 再加上约束条件 \eqref{eq:ma-14-5-1}, 即得式 \eqref{eq:ma-14-5-2}.
% \end{proof}
% 
% 
% \begin{remark}\label{remark:ma-14-source-110}
% 若构造函数
% \[
%   F(x_1,\cdots,x_n,\lambda_1,\cdots,\lambda_m)
%   =
%   f(x)+\sum_{j=1}^m\lambda_j\varphi_j(x),
% \]
% 则上述条件极值的必要条件形式上化为 $F$ 的通常极值的必要条件
% \begin{equation}
%   \begin{cases}
%     \dfrac{\partial F(x_0)}{\partial x_i}=0
%     \qquad (i=1,2,\cdots,n),\\[1.2ex]
%     \dfrac{\partial F(x_0)}{\partial\lambda_j}=0
%     \qquad (j=1,2,\cdots,m).
%   \end{cases}
%   \label{eq:ma-14-5-9}
% \end{equation}
% 上述求条件极值点的必要条件的方法称为拉格朗日乘数法.
% 
% 由于在拉格朗日乘数法中, 条件极值问题转化为
% \[
%   F(x)=f(x)+\sum_{j=1}^m\lambda_j\varphi_j(x)
% \]
% 的通常极值问题, 而当 $x_0$ 是驻点时, $\lambda_j(j=1,2,\cdots,m)$ 是 $m$ 个常数, 因此我们利用 $F(x)$ 在 $x_0$ 处的海色矩阵 $H_F(x_0)$ 来判定 $x_0$ 是否为极值点。可以证明: 当 $H_F(x_0)$ 正定时, 条件极值点为极小值点; 当 $H_F(x_0)$ 负定时, 条件极值为极大值点.
% 
% 值得提出的是, 当 $H_F(x_0)$ 不定时, 条件极值仍可以取到。例如, 设函数 $f(x,y)=x^2-y^2$, $\varphi(x,y)=y=0$, 则 $\varphi(x,y)=0$ 确定直线方程 $y=0$, 显然 $f(x,0)=x^2$ 在 $(0,0)$ 处取极小值 $0$。容易算出
% \[
%   F(x,y)=x^2-y^2+\lambda y
% \]
% 的海色矩阵在 $(0,0)$ 处为
% \[
%   \begin{pmatrix}
%     2 & 0\\
%     0 & -2
%   \end{pmatrix},
% \]
% 它是一个不定矩阵.
% \end{remark}
% 
% 
% \begin{example}\label{example:ma-14-5-4}
% 求椭球面 $16x^2+4y^2+9z^2=144$ 内接且各面均平行于坐标平面的长方体的最大体积.
% \end{example}
% 
% 
% \begin{solve}
% 设 $(x,y,z)$ 是该椭球面内接长方体在第一卦限内的顶点, 由对称性, 该长方体的体积为
% \[
%   V=f(x,y,z)=8xyz.
% \]
% 因此所求的最大体积转化为求 $f(x,y,z)$ 在约束条件
% \[
%   g(x,y,z)=16x^2+4y^2+9z^2-144=0
% \]
% 下的最大值。由拉格朗日乘数法, 我们作函数
% \[
%   F(x,y,z,\lambda)=f(x,y,z)+\lambda g(x,y,z).
% \]
% 求 $F(x,y,z,\lambda)$ 的各个偏导数得方程组
% \begin{equation}
%   \begin{cases}
%     8yz+32\lambda x=0,\\
%     8xz+8\lambda y=0,\\
%     8xy+18\lambda z=0,\\
%     16x^2+4y^2+9z^2-144=0.
%   \end{cases}
%   \label{eq:ma-14-5-10-to-14-5-13}
% \end{equation}
% 由上面方程组易推出
% \begin{equation}
%   xyz=-12\lambda.
%   \label{eq:ma-14-5-14}
% \end{equation}
% 将式 \eqref{eq:ma-14-5-14} 代入式 \eqref{eq:ma-14-5-10-to-14-5-13}中的第 1 式 得
% \[
%   -32\lambda(3-x^2)=0.
% \]
% 因此得 $\lambda=0$ 或 $x=\sqrt3$。当 $\lambda=0$ 时得 $xyz=0$。显然不合题意。因此必有 $\lambda\ne0$, 从而有 $x=\sqrt3$.
% 
% 同理将式 \eqref{eq:ma-14-5-14} 分别代入式 \eqref{eq:ma-14-5-10-to-14-5-13}中的第 2 式 和 \eqref{eq:ma-14-5-10-to-14-5-13}中的第 3 式 可得 $y=2\sqrt3$ 和 $z=4\sqrt3/3$.
% 
% 由于在该问题中最大值总是存在的, 因此所求最大值为
% \[
%   V_{\max}=8\times\sqrt3\times2\sqrt3\times\frac{4\sqrt3}{3}=64\sqrt3.
% \qedhere
% \]
% \end{solve}
% 
% 
% \begin{example}\label{example:ma-14-5-5}
% 设 $n$ 个正数之和为 $c$, 试求它们乘积的最大值, 并证明对任何正数 $a_1,a_2,\cdots,a_n$, 有
% \[
%   \sqrt[n]{a_1a_2\cdots a_n}
%   \leqslant
%   \frac{a_1+a_2+\cdots+a_n}{n}.
% \]
% \end{example}
% 
% 
% \begin{solve}
% 在 $E=\{(x_1,x_2,\cdots,x_n):x_i\geqslant0,\ i=1,2,\cdots,n\}$ 上定义函数
% \[
%   f(x_1,x_2,\cdots,x_n)=\prod_{i=1}^n x_i.
% \]
% 依题意, 我们要求 $f(x_1,x_2,\cdots,x_n)$ 在约束条件
% \[
%   \sum_{i=1}^n x_i=c
% \]
% 下的最大值.
% 
% 作辅助函数
% \[
%   F(x_1,x_2,\cdots,x_n,\lambda)
%   =
%   \prod_{i=1}^n x_i
%   +\lambda\left(\sum_{i=1}^n x_i-c\right).
% \]
% 求 $F$ 的 $n$ 个偏导数及 $\dfrac{\partial F}{\partial\lambda}$, 并令其为零, 得方程组
% \[
%   \begin{cases}
%     x_2x_3\cdots x_n+\lambda=0,\\
%     x_1x_3\cdots x_n+\lambda=0,\\
%     \cdots\\
%     x_1x_2\cdots x_{n-1}+\lambda=0,\\
%     \displaystyle\sum_{i=1}^n x_i=c,
%   \end{cases}
% \]
% 解出
% \[
%   x_1=x_2=\cdots=x_n=\frac cn,\qquad
%   \lambda=-\left(\frac cn\right)^{n-1}.
% \]
% 我们知道 $f(x_1,x_2,\cdots,x_n)$ 在紧集 $E$ 上取到最大值。另外, 容易看出, 若 $(x_1^0,x_2^0,\cdots,x_n^0)$ 是 $f(x_1,x_2,\cdots,x_n)$ 的最大值点, 必有 $x_i^0\ne0(i=1,2,\cdots,n)$。由于
% \[
%   \left(\frac cn,\frac cn,\cdots,\frac cn,-\left(\frac cn\right)^{n-1}\right)
% \]
% 是 $F(x_1,x_2,\cdots,x_n,\lambda)$ 的唯一驻点, 因此它为 $F(x_1,x_2,\cdots,x_n,\lambda)$ 的极大值点, 从而 $f(x_1,x_2,\cdots,x_n)$ 在该点取最大值
% \[
%   \frac{c^n}{n^n}.
% \]
% 现对 $n$ 个正数 $a_1,a_2,\cdots,a_n$, 令 $c=\displaystyle\sum_{i=1}^n a_i$, 并利用上述函数 $f(x)$ 所证结论, 则对任何满足 $\displaystyle\sum_{i=1}^n x_i=c$ 的正数 $x_i$ 均有
% \[
%   \prod_{i=1}^n x_i\leqslant\frac{c^n}{n^n}.
% \]
% 取 $x_i=a_i(i=1,2,\cdots,n)$, 则有
% \[
%   \sqrt[n]{a_1a_2\cdots a_n}
%   \leqslant
%   \frac{a_1+a_2+\cdots+a_n}{n}.
% \qedhere
% \]
% \end{solve}
% 
% 
% \section{多元微分学的几何应用}
% \label{sec:ma-14-6}
% 
% 在本章的最后一节, 我们来讨论多元微分学在几何上的一些应用.
% 
% \subsection{空间曲线的切线与法平面}
% \label{subsec:ma-14-6-1}
% 
% 我们回忆一下, 在上一章曾经给出的曲线定义: $\mathbb R^n$ 中的一条曲线是 $[\alpha,\beta]\to\mathbb R^n$ 的一个连续映射
% \[
%   h(t)=(x_1(t),x_2(t),\cdots,x_n(t)).
% \]
% 在 $n=1,2,3$ 时, 曲线是我们经常遇到的。但是值得指出的是, 曲线有时可能是非常复杂的。在第二册中, 我们曾经遇到过不可求长的曲线。皮亚诺甚至构造了一条连续曲线 $\to D=[0,1]\times[0,1]$:
% \[
%   \begin{cases}
%     x=x(t),\\
%     y=y(t),
%   \end{cases}
%   \quad (0\leqslant t\leqslant1);
% \]
% 使得它充满了整个 $D$.
% 
% 在以后各章节讨论的曲线中, 我们作如下约定: 设曲线 $\Gamma$ 由连续映射
% \[
%   h(t)=(x_1(t),x_2(t),\cdots,x_n(t)),\quad t\in[\alpha,\beta]
% \]
% 所确定。若对于 $\forall t_1,t_2\in[\alpha,\beta]$, $h(t_1)\ne h(t_2)$, 则称 $\Gamma$ 是一条简单曲线。若对于 $\forall t_1,t_2\in[\alpha,\beta)$, 有 $h(t_1)\ne h(t_2)$, 但 $h(\alpha)=h(\beta)$, 则 $\Gamma$ 是一条封闭曲线, 则称 $\Gamma$ 是一条简单闭曲线。简单闭曲线也称为约当 (Jordan) 曲线.
% 
% 我们首先讨论比较简单的情形, 即空间曲线是由参数方程
% \[
%   \Gamma:
%   \begin{cases}
%     x=x(t),\\
%     y=y(t),\\
%     z=z(t),
%   \end{cases}
%   \quad t\in(\alpha,\beta)
% \]
% 给出, 并假定 $x(t),y(t)$ 与 $z(t)$ 都是 $t$ 的可微函数, 且 $x'^2(t)+y'^2(t)+z'^2(t)\ne0$。对于这类曲线, 它局部总是简单曲线。在上述条件下, 我们现在来求过 $t_0\in(\alpha,\beta)$ 对应的曲线上的点 $x(t_0)=(x(t_0),y(t_0),z(t_0))$ 处的切线方程.
% 
% 由切线的定义, 它是曲线 $\Gamma$ 上过点 $x(t)=(x(t),y(t),z(t))$ 与 $x(t_0)$ 的割线当 $t\to t_0$ 时的极限位置。由于过点 $x(t_0)$ 与 $x(t)$ 的割线方程为
% \[
%   \frac{x-x(t_0)}{x(t)-x(t_0)}
%   =
%   \frac{y-y(t_0)}{y(t)-y(t_0)}
%   =
%   \frac{z-z(t_0)}{z(t)-z(t_0)};
% \]
% 在上式分母中同除以 $t-t_0$, 并令 $t\to t_0$ 即得所求切线方程:
% \[
%   \frac{x-x(t_0)}{x'(t_0)}
%   =
%   \frac{y-y(t_0)}{y'(t_0)}
%   =
%   \frac{z-z(t_0)}{z'(t_0)}.
% \]
% 从上述方程可以看出, 若记曲线 $h(t)=(x(t),y(t),z(t))(t\in(\alpha,\beta))$, 则
% \[
%   h'(t_0)=(x'(t_0),y'(t_0),z'(t_0))
% \]
% 即为曲线 $\Gamma$ 在 $x(t_0)=(x(t_0),y(t_0),z(t_0))$ 处的切向量.
% 
% 我们同时也可得到曲线 $\Gamma$ 在 $x(t_0)=(x(t_0),y(t_0),z(t_0))$ 处的法平面方程, 即过点 $x(t_0)$ 且以切向量 $h'(t_0)$ 为法向的平面方程:
% \[
%   x'(t_0)(x-x(t_0))+y'(t_0)(y-y(t_0))+z'(t_0)(z-z(t_0))=0.
% \]
% 上式用向量内积记之为
% \[
%   h'(t_0)\cdot(x-x(t_0))=0.
% \]
% 
% 其次, 我们来讨论由两个方程所确定的曲线的切线方程。给定两个方程所组成的方程组
% \[
%   \begin{cases}
%     F_1(x,y,z)=0,\\
%     F_2(x,y,z)=0,
%   \end{cases}
% \]
% 其中 $(x,y,z)\in D\subset\mathbb R^3$, 这里 $D$ 是一个区域。假定 $F_1$ 与 $F_2$ 均是 $D$ 上的可微函数, 并设 $x_0\in D$。当
% \[
%   \begin{pmatrix}
%     F_1'(x_0)\\
%     F_2'(x_0)
%   \end{pmatrix}
%   =
%   \left.
%   \begin{pmatrix}
%     \dfrac{\partial F_1}{\partial x} & \dfrac{\partial F_1}{\partial y} & \dfrac{\partial F_1}{\partial z}\\
%     \dfrac{\partial F_2}{\partial x} & \dfrac{\partial F_2}{\partial y} & \dfrac{\partial F_2}{\partial z}
%   \end{pmatrix}
%   \right|_{x_0}
% \]
% 的秩为 $2$ 时, 由定理~\ref{theorem:ma-14-4-3}, 在 $x_0=(x_0,y_0,z_0)$ 的附近, 其中有两个变量可确定为第三个变量的函数, 从而该方程组能确定一条过点 $x_0$ 的曲线 $\Gamma$。从定理~\ref{theorem:ma-14-4-3} 我们还可以知道, 在 $x_0$ 附近 $\Gamma$ 还可以具有参数形式: $(x(t),y(t),z(t)),t\in(\alpha,\beta)$, 其中
% \[
%   (x(t_0),y(t_0),z(t_0))=(x_0,y_0,z_0)=x_0.
% \]
% 由上面的讨论知, 曲线在点 $x_0$ 处的切向量为 $(x'(t_0),y'(t_0),z'(t_0))$。由于曲线 $\Gamma$ 由方程组确定, 从而有
% \[
%   \begin{cases}
%     F_1(x(t),y(t),z(t))=0,\\
%     F_2(x(t),y(t),z(t))=0.
%   \end{cases}
% \]
% 对此方程组的第一个方程两边关于 $t$ 求导数, 得
% \[
%   \frac{\partial F_1(x_0)}{\partial x}x'(t_0)
%   +
%   \frac{\partial F_1(x_0)}{\partial y}y'(t_0)
%   +
%   \frac{\partial F_1(x_0)}{\partial z}z'(t_0)
%   =0.
% \]
% 将上式写成内积形式有
% \[
%   (x'(t_0),y'(t_0),z'(t_0))\cdot
%   \left(
%     \frac{\partial F_1(x_0)}{\partial x},
%     \frac{\partial F_1(x_0)}{\partial y},
%     \frac{\partial F_1(x_0)}{\partial z}
%   \right)
%   =0.
% \]
% 同理,
% \[
%   (x'(t_0),y'(t_0),z'(t_0))\cdot
%   \left(
%     \frac{\partial F_2(x_0)}{\partial x},
%     \frac{\partial F_2(x_0)}{\partial y},
%     \frac{\partial F_2(x_0)}{\partial z}
%   \right)
%   =0.
% \]
% 因此, 曲线 $\Gamma$ 在 $x_0$ 处的切向量与
% \[
%   F_1'(x_0)\times F_2'(x_0)
%   =
%   \left.
%   \begin{vmatrix}
%     \mathbf i & \mathbf j & \mathbf k\\
%     \dfrac{\partial F_1}{\partial x} & \dfrac{\partial F_1}{\partial y} & \dfrac{\partial F_1}{\partial z}\\
%     \dfrac{\partial F_2}{\partial x} & \dfrac{\partial F_2}{\partial y} & \dfrac{\partial F_2}{\partial z}
%   \end{vmatrix}
%   \right|_{x_0}
% \]
% 平行。记
% \[
%   A=\left.\frac{\partial(F_1,F_2)}{\partial(y,z)}\right|_{x_0},\quad
%   B=\left.\frac{\partial(F_1,F_2)}{\partial(z,x)}\right|_{x_0},\quad
%   C=\left.\frac{\partial(F_1,F_2)}{\partial(x,y)}\right|_{x_0},
% \]
% 则 $F_1'(x_0)\times F_2'(x_0)=A\mathbf i+B\mathbf j+C\mathbf k$。因此曲线 $\Gamma$ 在 $x_0$ 处的切线方程为
% \[
%   \frac{x-x_0}{A}=\frac{y-y_0}{B}=\frac{z-z_0}{C},
% \]
% 而法平面方程为
% \[
%   A(x-x_0)+B(y-y_0)+C(z-z_0)=0.
% \]
% 
% \subsection{曲面的切平面与法线}
% \label{subsec:ma-14-6-2}
% 
% 设曲面 $S$ 由方程
% \[
%   F(x,y,z)=0
% \]
% 给出, 其中 $F(x,y,z)\in C^1(D)$, $D\subset\mathbb R^3$ 为一个区域, 再设点 $(x_0,y_0,z_0)\in D$, 使得 $F(x_0,y_0,z_0)=0$。现在我们来求曲面 $S$ 在点 $(x_0,y_0,z_0)$ 处的切平面及法线的方程.
% 
% 在曲面 $S$ 上任取一条过点 $(x_0,y_0,z_0)$ 的光滑曲线 $(x(t),y(t),z(t))$, $t\in(\alpha,\beta)$。因此有
% \[
%   F(x(t),y(t),z(t))=0,\quad \forall t\in(\alpha,\beta).
% \]
% 设 $t_0\in(\alpha,\beta)$, 使得 $(x(t_0),y(t_0),z(t_0))=(x_0,y_0,z_0)$。将上述方程两边在 $t_0$ 处求导数得
% \[
%   F_x'(x_0,y_0,z_0)x'(t_0)+F_y'(x_0,y_0,z_0)y'(t_0)+F_z'(x_0,y_0,z_0)z'(t_0)=0.
% \]
% 注意到 $(x'(t_0),y'(t_0),z'(t_0))$ 是该曲线在 $(x_0,y_0,z_0)$ 处的切向量, 因此上式说明了该切向量与向量 $F'(x_0,y_0,z_0)$ 正交, 从而当 $F'(x_0,y_0,z_0)$ 不是零向量时, 上式表明曲面 $S$ 过点 $(x_0,y_0,z_0)$ 的任何光滑曲线的切向量与固定向量 $F'(x_0,y_0,z_0)$ 正交。由此我们推知曲面 $S$ 过点 $(x_0,y_0,z_0)$ 的任何光滑曲线在 $(x_0,y_0,z_0)$ 处的切线都在平面
% \[
%   F_x'(x_0,y_0,z_0)(x-x_0)+F_y'(x_0,y_0,z_0)(y-y_0)+F_z'(x_0,y_0,z_0)(z-z_0)=0
% \]
% 上。我们自然地称上述平面为曲面 $S$ 在 $(x_0,y_0,z_0)$ 处的切平面.
% 
% 另外, 我们称直线
% \[
%   \frac{x-x_0}{F_x'(x_0,y_0,z_0)}
%   =
%   \frac{y-y_0}{F_y'(x_0,y_0,z_0)}
%   =
%   \frac{z-z_0}{F_z'(x_0,y_0,z_0)}
% \]
% 为曲面 $S$ 在 $(x_0,y_0,z_0)$ 处的法线.
% 
% 当 $F'(x_0,y_0,z_0)\ne0$ 时, 我们可以证明: 曲面 $S$ 在 $(x_0,y_0,z_0)$ 处的切平面上任何过 $(x_0,y_0,z_0)$ 的直线都是该曲面上某光滑曲线的切线 (见本章习题).
% 
% 现在设曲面 $S$ 由参数方程
% \[
%   \begin{cases}
%     x=x(u,v),\\
%     y=y(u,v),\\
%     z=z(u,v),
%   \end{cases}
%   \quad (u,v)\in D
% \]
% 给出, 其中 $D$ 是 $\mathbb R^2$ 中的区域, 而上述三个函数均具有连续偏导数。我们来求曲面 $S$ 在 $(x_0,y_0,z_0)=(x(u_0,v_0),y(u_0,v_0),z(u_0,v_0))$ 处的切平面与法线方程。为此, 我们在曲面 $S$ 上取两条特殊的曲线:
% \[
%   \begin{cases}
%     x=x(u,v_0),\\
%     y=y(u,v_0),\\
%     z=z(u,v_0),
%   \end{cases}
%   \quad\text{和}\quad
%   \begin{cases}
%     x=x(u_0,v),\\
%     y=y(u_0,v),\\
%     z=z(u_0,v).
%   \end{cases}
% \]
% 它们在 $(x_0,y_0,z_0)$ 处的切向量分别为
% \[
%   (x_u'(u_0,v_0),y_u'(u_0,v_0),z_u'(u_0,v_0)),
% \]
% \[
%   (x_v'(u_0,v_0),y_v'(u_0,v_0),z_v'(u_0,v_0)).
% \]
% 若这两个向量不平行, 它们对应的切线将确定曲面 $S$ 在 $(x_0,y_0,z_0)$ 处的法向量:
% \[
%   n=(x_u(u_0,v_0),y_u(u_0,v_0),z_u(u_0,v_0))
%   \times(x_v(u_0,v_0),y_v(u_0,v_0),z_v(u_0,v_0))
% \]
% \[
%   =
%   \left.
%   \begin{vmatrix}
%     \mathbf i & \mathbf j & \mathbf k\\
%     x_u' & y_u' & z_u'\\
%     x_v' & y_v' & z_v'
%   \end{vmatrix}
%   \right|_{(u_0,v_0)}.
% \]
% 因此当
% \[
%   \left.
%   \begin{pmatrix}
%     x_u' & y_u' & z_u'\\
%     x_v' & y_v' & z_v'
%   \end{pmatrix}
%   \right|_{(u_0,v_0)}
% \]
% 的秩为 $2$ 时, 记
% \[
%   A=\left.\frac{\partial(y,z)}{\partial(u,v)}\right|_{(u_0,v_0)},\quad
%   B=\left.\frac{\partial(z,x)}{\partial(u,v)}\right|_{(u_0,v_0)},\quad
%   C=\left.\frac{\partial(x,y)}{\partial(u,v)}\right|_{(u_0,v_0)},
% \]
% 则该曲面在 $(x_0,y_0,z_0)$ 处的切平面方程与法线方程分别为
% \[
%   A(x-x_0)+B(y-y_0)+C(z-z_0)=0,
% \]
% \[
%   \frac{x-x_0}{A}=\frac{y-y_0}{B}=\frac{z-z_0}{C}.
% \]
% 
% 到现在为止, 我们已经研究了 $\mathbb R^3$ 中的曲线切线与曲面的切平面等问题。作为推广, 我们还可以考虑 $\mathbb R^n$ 中的曲线。在 $\mathbb R^n$ 中, 我们曾称 $[0,1]$ 到 $\mathbb R^n$ 的一个连续映射 $h(t)=(x_1(t),x_2(t),\cdots,x_n(t))$ 为一条曲线。当 $h(t)$ 具有连续导数且 $h'(t)\ne0$ 时, 称此曲线为光滑曲线, 并称 $h'(t_0)$ 是曲线在 $h(t_0)$ 处的切向量。类似地, 我们称 $F(x)=F(x_1,x_2,\cdots,x_n)=0$ 为 $\mathbb R^n$ 中的一个曲面, 其中 $F(x)$ 是一个 $C^1$ 函数。当 $F'(x_0)\ne0$ 时, 称 $F'(x_0)$ 为此曲面在 $x_0$ 处的法向量, 并称
% \[
%   F'(x_0)(x-x_0)
%   =
%   \sum_{i=1}^n \frac{\partial F(x)}{\partial x_i}(x_i-x_i^0)=0
% \]
% 为该曲面在 $x_0$ 处的切平面方程。作为练习, 请读者写出 $\mathbb R^n$ 中曲线的切线方程与法平面方程以及曲面的法线方程.
% 
% \begin{example}\label{example:ma-14-6-1}
% 求曲线
% \[
%   \begin{cases}
%     x^2-y^2+2z^2=2,\\
%     x+y+z=3
%   \end{cases}
% \]
% 在 $(1,1,1)$ 处的切线方程与法平面方程.
% \end{example}
% 
% 
% \begin{solve}
% 令 $F(x,y,z)=x^2-y^2+2z^2-2$, $G(x,y,z)=x+y+z-3$, 则
% \[
%   \frac{\partial F}{\partial x}=2x,\quad
%   \frac{\partial F}{\partial y}=-2y,\quad
%   \frac{\partial F}{\partial z}=4z,\quad
%   \frac{\partial G}{\partial x}
%   =
%   \frac{\partial G}{\partial y}
%   =
%   \frac{\partial G}{\partial z}
%   =1,
% \]
% 于是
% \[
%   \left.\frac{\partial(F,G)}{\partial(y,z)}\right|_{(1,1,1)}
%   =
%   \left.
%   \begin{vmatrix}
%     \dfrac{\partial F}{\partial y} & \dfrac{\partial F}{\partial z}\\
%     \dfrac{\partial G}{\partial y} & \dfrac{\partial G}{\partial z}
%   \end{vmatrix}
%   \right|_{(1,1,1)}
%   =
%   \begin{vmatrix}
%     -2 & 4\\
%     1 & 1
%   \end{vmatrix}
%   =-6,
% \]
% \[
%   \left.\frac{\partial(F,G)}{\partial(z,x)}\right|_{(1,1,1)}
%   =
%   \left.
%   \begin{vmatrix}
%     \dfrac{\partial F}{\partial z} & \dfrac{\partial F}{\partial x}\\
%     \dfrac{\partial G}{\partial z} & \dfrac{\partial G}{\partial x}
%   \end{vmatrix}
%   \right|_{(1,1,1)}
%   =
%   \begin{vmatrix}
%     4 & 2\\
%     1 & 1
%   \end{vmatrix}
%   =2,
% \]
% \[
%   \left.\frac{\partial(F,G)}{\partial(x,y)}\right|_{(1,1,1)}
%   =
%   \left.
%   \begin{vmatrix}
%     \dfrac{\partial F}{\partial x} & \dfrac{\partial F}{\partial y}\\
%     \dfrac{\partial G}{\partial x} & \dfrac{\partial G}{\partial y}
%   \end{vmatrix}
%   \right|_{(1,1,1)}
%   =
%   \begin{vmatrix}
%     2 & -2\\
%     1 & 1
%   \end{vmatrix}
%   =4.
% \]
% 因此, 所求切线方程为
% \[
%   \frac{x-1}{-3}=\frac{y-1}{1}=\frac{z-1}{2},
% \]
% 法平面方程为
% \[
%   -3(x-1)+(y-1)+2(z-1)=0,\quad \text{即}\quad 3x-y-2z=0.
% \]
% 
% 对于上述切线方程, 还可用另一方法求解。由于曲面 $F(x,y,z)=0$ 在 $(1,1,1)$ 处的法向量为 $(F_x',F_y',F_z')|_{(1,1,1)}=(2,-2,4)$, 因此在该点处的切平面方程为
% \[
%   2(x-1)-2(y-1)+4(z-1)=0,
% \]
% 即 $x-y+2z=2$。由于切线是两曲面在该点的切平面的交线, 而另一曲面 $G(x,y,z)=0$ 为一平面, 因此所求切线方程为
% \[
%   \begin{cases}
%     x-y+2z=2,\\
%     x+y+z=3.
%   \end{cases}
% \qedhere
% \]
% \end{solve}
% 
% 
% \begin{example}\label{example:ma-14-6-2}
% 证明曲面 $S:x^2+y^2+2xyz=4$ 与曲面族 $S_t:2x+ty^2-(1+t)z^2=1(t\in\mathbb R)$ 在 $(1,1,1)$ 处正交 (即它们的法向量相互垂直), 并求出它们的交线在该点的切线方程.
% \end{example}
% 
% 
% \begin{solve}
% 令
% \[
%   F(x,y,z)=x^2+y^2+2xyz-4,
% \]
% \[
%   G_t(x,y,z)=2x+ty^2-(1+t)z^2-1,
% \]
% 则 $S$ 在 $(1,1,1)$ 处的法向量为
% \[
%   \left(
%     \frac{\partial F(1,1,1)}{\partial x},
%     \frac{\partial F(1,1,1)}{\partial y},
%     \frac{\partial F(1,1,1)}{\partial z}
%   \right)
%   =(2,2,2),
% \]
% 而 $S_t$ 在 $(1,1,1)$ 处的法向量为
% \[
%   \left(
%     \frac{\partial G_t(1,1,1)}{\partial x},
%     \frac{\partial G_t(1,1,1)}{\partial y},
%     \frac{\partial G_t(1,1,1)}{\partial z}
%   \right)
%   =(2,2t,-2(1+t)).
% \]
% 因为
% \[
%   (2,2,2)(2,2t,-2(1+t))=0,
% \]
% 所以 $S$ 与 $S_t$ 在 $(1,1,1)$ 处总是正交的.
% 
% 由于 $S$ 与 $S_t$ 的交线在 $(1,1,1)$ 处的切线方程即两曲面在 $(1,1,1)$ 处的切平面的交线, 因此切线方程为
% \[
%   \begin{cases}
%     (x-1)+(y-1)+(z-1)=0,\\
%     (x-1)+t(y-1)-(1+t)(z-1)=0,
%   \end{cases}
% \]
% 化简得
% \[
%   \begin{cases}
%     x+y+z=3,\\
%     x+ty-(1+t)z=0.
%   \end{cases}
% \qedhere
% \]
% \end{solve}
% 
% 
% \subsection{多元凸函数}
% \label{subsec:ma-14-6-3}
% 
% 作为泰勒公式的应用, 我们来证明 $n$ 元函数凸性的一个结果。由于研究凸性需要考虑任意两点之间线段上的函数值与端点函数值之间的关系, 因此我们假定 $D\subset\mathbb R^n$ 是一个凸域。若 $f(x)$ 在凸域 $D$ 内具有 $k$ 阶连续偏导数, 则对于 $\forall x_0\in D$, $f(x)$ 在 $D$ 内总可以在 $x_0$ 处展成 $k$ 阶泰勒公式.
% 
% \begin{definition}\label{definition:ma-14-6-1}
% 设 $D\subset\mathbb R^n$ 是一个凸域, $f(x)$ 在 $D$ 内有定义。如果对于 $\forall x_0,x_1\in D$ 和 $\forall t\in(0,1)$, 有
% \[
%   f(tx_1+(1-t)x_2)\leqslant tf(x_1)+(1-t)f(x_2),
% \]
% 则称 $f(x)$ 在 $D$ 内是凸函数; 如果上述不等式总成立严格不等式, 则称 $f(x)$ 在 $D$ 内是严格凸函数.
% \end{definition}
% 
% 在微分学中, 我们主要利用导数 (偏导数) 来研究凸函数。回忆一下, 对一元函数 $f(x)$ 来说, 若 $f''(x)$ 存在, 且 $f''(x)>0$ 处处成立, 则 $f(x)$ 是凸的。我们前面提到过, 对于 $n$ 元函数 $f(x)$, 其导数 $f'(x)$ 在很多情形下将起着一元函数导数相似的作用。下面定理的结论说明, 在凸函数的研究中, $n$ 元二阶连续可微函数 $f(x)$ 的海色矩阵 $H_f(x)$ 起着类似于一元函数二阶导数的作用.
% 
% \begin{theorem}\label{theorem:ma-14-6-1}
% 设 $D\subset\mathbb R^n$ 是一个凸区域, 函数 $f(x)$ 在 $D$ 内具有二阶连续偏导数, 则以下结论等价:
% 
% (1) $f(x)$ 在 $D$ 内是凸函数;
% 
% (2) 对于 $\forall x_0,x\in D$, 成立 $f(x)\geqslant f(x_0)+f'(x_0)(x-x_0)^{\mathrm T}$;
% 
% (3) 对于 $\forall x_0\in D$, $f(x)$ 在 $x_0$ 处的海色矩阵 $H_f(x_0)$ 半正定.
% \end{theorem}
% 
% 
% \begin{proof}
% 先证 (1) $\Rightarrow$ (2)。由凸函数的定义知, 对于 $\forall x_0,x\in D$ 及 $\forall t\in(0,1)$, 有
% \begin{equation}
%   f(tx+(1-t)x_0)\leqslant tf(x)+(1-t)f(x_0).
%   \label{eq:ma-14-6-1}
% \end{equation}
% 由于 $f(x)$ 具有二阶连续偏导数, 我们有
% \begin{equation}
%   \begin{aligned}
%     f(tx+(1-t)x_0)
%     &=f(x_0+t(x-x_0))\\
%     &=f(x_0)+tf'(x_0)(x-x_0)^{\mathrm T}
%       +o(t|x-x_0|)\quad (t|x-x_0|\to0).
%   \end{aligned}
%   \label{eq:ma-14-6-2}
% \end{equation}
% 将式 \eqref{eq:ma-14-6-2} 代入式 \eqref{eq:ma-14-6-1} 得
% \[
%   t(f'(x_0)(x-x_0)^{\mathrm T})+o(t|x-x_0|)\leqslant t(f(x)-f(x_0)).
% \]
% 在上式中两边除以 $t$, 并注意到 $|x-x_0|$ 在 $t$ 变化的过程中是常数, 我们有
% \[
%   \lim_{t\to0+0}\frac{o(t|x-x_0|)}{t}=0,
% \]
% 因此
% \[
%   f(x)\geqslant f(x_0)+f'(x_0)(x-x_0)^{\mathrm T}.
% \]
% 
% 再证 (2) $\Rightarrow$ (3)。任取 $x_0\in D$ 及 $\Delta x\in\mathbb R^n$, $\Delta x\ne0$, 则当 $t$ 充分小时, $x_0+t\Delta x\in D$。由泰勒公式有
% \[
%   f(x_0+t\Delta x)=f(x_0)+tf'(x_0)\Delta x^{\mathrm T}
%   +\frac{t^2}{2}\Delta xH_f(x_0)\Delta x^{\mathrm T}
%   +o(t^2|\Delta x|^2)\quad (t|\Delta x|\to0).
% \]
% 在 (2) 中令 $x=x_0+t\Delta x$, 我们有
% \[
%   \frac{t^2}{2}\Delta xH_f(x_0)\Delta x^{\mathrm T}
%   +o(t^2|\Delta x|^2)\geqslant0,
% \]
% 即
% \[
%   \Delta xH_f(x_0)\Delta x^{\mathrm T}
%   +\frac{o(t^2|\Delta x|^2)}{t^2}\geqslant0.
% \]
% 令 $t\to0+0$, 即得
% \[
%   \Delta xH_f(x_0)\Delta x^{\mathrm T}\geqslant0,
% \]
% 这说明 $H_f(x_0)$ 是半正定的.
% 
% 最后证 (3) $\Rightarrow$ (1)。对于 $\forall x_1,x_2\in D$, 令 $x_0=tx_1+(1-t)x_2(t\in(0,1))$, 则有 $x_0\in D$, 从而有
% \[
%   f(x_1)=f(x_0)+f'(x_0)(x_1-x_0)^{\mathrm T}
%   +\frac12(x_1-x_0)H_f(x_0+\theta(x_1-x_0))(x_1-x_0)^{\mathrm T},
% \]
% 其中 $0<\theta<1$。由 (3) 知
% \[
%   f(x_1)\geqslant f(x_0)+f'(x_0)(x_1-x_0)^{\mathrm T}.
% \]
% 同理,
% \[
%   f(x_2)\geqslant f(x_0)+f'(x_0)(x_2-x_0)^{\mathrm T}.
% \]
% 因此有
% \[
%   \begin{aligned}
%     tf(x_1)+(1-t)f(x_2)
%     &\geqslant t[f(x_0)+f'(x_0)(x_1-x_0)^{\mathrm T}]
%       +(1-t)[f(x_0)+f'(x_0)(x_2-x_0)^{\mathrm T}]\\
%     &=f(x_0)=f(tx_1+(1-t)x_2).
%   \end{aligned}
% \qedhere
% \]
% \end{proof}
% 
% 
% \begin{remark}\label{remark:ma-14-source-122}
% 条件 (2) 说明 $u=f(x)$ 表示的 ``曲面'' 在其任一点处的切平面的上方.
% \end{remark}
% 
% 
% \begin{example}\label{example:ma-14-6-3}
% 证明: 若一元凸函数 $y=g(x)$ 在 $(a,b)$ 内具有二阶导数, 则 $f(x)=\displaystyle\sum_{i=1}^n g(x_i)$ 是
% \[
%   D=\underbrace{(a,b)\times(a,b)\times\cdots\times(a,b)}_{n\text{ 个}}
% \]
% 内的凸函数.
% \end{example}
% 
% 
% \begin{proof}
% 因为 $f'(x)=(g'(x_1),g'(x_2),\cdots,g'(x_n))$, 所以
% \[
%   H_f(x)=
%   \begin{pmatrix}
%     g''(x_1) & 0 & \cdots & 0\\
%     0 & g''(x_2) & \cdots & 0\\
%     \vdots & \vdots & & \vdots\\
%     0 & 0 & \cdots & g''(x_n)
%   \end{pmatrix}.
% \]
% 由于 $g''(x_i)\geqslant0(i=1,2,\cdots,n)$, 从代数知识知, $H_f(x)$ 是半正定矩阵。由定理~\ref{theorem:ma-14-6-1} 知 $f(x)$ 是 $D$ 内的凸函数.
% \end{proof}
% 
% 
% \section*{习题十四}
% \phantomsection
% \addcontentsline{toc}{section}{习题十四}
% \label{exercises:ma-14}
% 
% 1。设函数 $u=f(x)$ 在 $U(x_0,\delta_0)\subset\mathbb R^n(\delta_0>0)$ 内存在各个偏导数, 并且所有的偏导数在该邻域内有界, 证明 $f(x)$ 在 $x_0$ 处连续; 举例说明存在函数 $u=g(x)$, 它在 $x_0$ 的某个邻域内存在无界的各个偏导数, 但它在 $x_0$ 处连续.
% 
% 2。举例说明在 $\mathbb R^2$ 内存在函数 $z=f(x,y)$, 使得 $f(x,y)$ 在 $\mathbb R^2$ 内处处不连续, 但它在原点处存在两个偏导数.
% 
% 3。求下列函数在指定点处的偏导数:
% 
% (1) $f(x,y)=xy\ln[x^2+\sin(xy^2)+\sin(xy)]$, 求 $f_x'(1,-1),f_y'(1,-1)$;
% 
% % 整合来源：math-15.tex
% \chapter{重积分}
% \label{chap:ma-15}
% 
% 本章中, 我们将研究多元函数的重积分。多元函数的重积分理论与一元函数的定积分理论是平行的, 最主要的差别在于重积分的积分区域与闭区间相比更为复杂, 从而使得重积分的理论与计算也要复杂得多。读者在学习重积分这部分内容时, 应该对一般的重积分理论有深刻的理解, 而具体计算重积分时, 主要精力则应放在二重及三重积分上.
% 
% \section{重积分的定义}
% \label{sec:ma-15-1}
% 
% 在日常生活和科学技术中许多问题的解决都需要用到重积分。例如, 设有一物体 (如一块铁矿等), 已知其密度, 要求它的质量。若将此物体放在 $\mathbb R^3$ 中, 则从几何上看, 该物体即为空间的一个闭区域 $D$, 而该物体的密度可以通过 $D$ 上的一个函数 $\rho(x,y,z)$ 来刻画.
% 
% 当 $\rho(x,y,z)=\rho_0$ ($\rho_0$ 为常数) 时, 该物体的质量 $m=\rho_0V$, 其中 $V$ 是 $D$ 的体积; 当 $\rho(x,y,z)$ 不是常数函数时, 我们可以利用“分割--求和--取极限”的方法来求其质量。我们将 $D$ 作分割
% \[
%   \Delta=\{\Delta D_1,\Delta D_2,\cdots,\Delta D_K\},
% \]
% 使得
% \[
%   \bigcup_{k=1}^K\Delta D_k=D,
% \]
% 且每个 $\Delta D_k$ 的直径
% \[
%   \operatorname{diam}(\Delta D_k)=\sup_{x,y\in\Delta D_k}|x-y|
% \]
% 很小。这样做的原因是使得 $\rho(x,y,z)$ 在 $\Delta D_k$ 上的变化不是很大, 因此可认为 $\rho(x,y,z)$ 在 $\Delta D_k$ 上几乎是常数。在 $\Delta D_k$ 上任取 $(\xi_k,\eta_k,\zeta_k)$, 记 $\Delta V_k$ 为 $\Delta D_k$ 的体积, 作和式
% \[
%   \sum_{k=1}^K\rho(\xi_k,\eta_k,\zeta_k)\Delta V_k
% \]
% (该和式也称为 $\rho(x,y,z)$ 的黎曼和)。若当
% \[
%   \lambda(\Delta)=\max_{1\leqslant k\leqslant K}\operatorname{diam}(\Delta D_k)\to0
% \]
% 时, 上述和式存在极限 $m$, 则自然认为 $m$ 为该物体的质量.
% 
% 另外一个例子是求曲顶柱体的体积。设 $D\subset\mathbb R^2$ 是有界闭区域, 函数 $z=f(x,y)$ 在 $D$ 上连续, 且对于 $\forall(x,y)\in D$, 有 $f(x,y)>0$。因此我们有一个空间立体
% \[
%   \Omega=\{(x,y,z):(x,y)\in D,\ 0\leqslant z\leqslant f(x,y)\}.
% \]
% 通常称形如这样的立体为曲顶柱体(见图 15.1.1)。为了求 $\Omega$ 的体积 $V$, 我们可以对 $D$ 作分割
% \[
%   \Delta=\{\Delta D_1,\Delta D_2,\cdots,\Delta D_K\}.
% \]
% 在 $\Delta D_k(k=1,2,\cdots,K)$ 上取 $(\xi_k,\eta_k)$, 并记 $\Delta\sigma_k$ 为 $\Delta D_k$ 的面积。当极限
% \[
%   \lim_{\lambda(\Delta)\to0}\sum_{k=1}^K f(\xi_k,\eta_k)\Delta\sigma_k
% \]
% 存在时, 同样可以认为该极限即为 $V$ 的体积.
% 
% 上面的两个例子都归结到求一个特别的黎曼和的极限。讨论何时这种和式存在极限就是重积分理论所要研究的问题.
% 
% \subsection{\texorpdfstring{$\mathbb R^n$ 空间中集合的体积}{Rⁿ 空间中集合的体积}}
% \label{subsec:ma-15-1-1}
% 
% 在上面两个例子的黎曼和中, 我们首先面临的问题是平面区域的面积与 $\mathbb R^3$ 中区域的体积问题。在本小节中, 我们先来研究一下这个问题。在 $\mathbb R^n$ 空间中, 我们使用的是欧氏度量, 即对于
% \[
%   x=(x_1,x_2,\cdots,x_n),\quad y=(y_1,y_2,\cdots,y_n)\in\mathbb R^n,
% \]
% $x,y$ 之间的距离为
% \[
%   |x-y|=\sqrt{\sum_{j=1}^n(x_j-y_j)^2}.
% \]
% 因此我们可以定义一个 $n$ 维长方体
% \[
%   A=\{(x_1,x_2,\cdots,x_n):a_j\leqslant x_j\leqslant b_j\}
% \]
% 的体积为
% \[
%   V(A)=\prod_{j=1}^n(b_j-a_j).
% \]
% 在本章中, 如果没有特别说明, 所有的长方体都是指上述形式的长方体, 即这些长方体的各个面均是与某个坐标平面平行的.
% 
% 下面我们希望用长方体的体积来给出一般集合体积的定义。设 $E\subset\mathbb R^n$ 是一给定的非空有界集合, 我们来讨论 $E$ 的体积的定义。不妨设 $E$ 不是一个长方体。我们考虑由 $\mathbb R^n$ 中有限个长方体构成的长方体族 $\{A_k\}_{k=1}^K$, 且假定它满足条件:
% \[
%   (1)\quad A_j^\circ\cap A_l^\circ=\varnothing\quad (j,l=1,2,\cdots,K;\ j\ne l);
% \]
% \[
%   (2)\quad \bigcup_{k=1}^K A_k\supset E.
% \]
% 为了简便, 以下我们将满足上述条件的长方体族记为 $\mathcal A_E$.
% 
% 对每个给定的长方体族 $\mathcal A_E$, 记 $m(\mathcal A_E)$ 为 $\mathcal A_E$ 中完全含于 $E$ 内的那些长方体的体积之和, $M(\mathcal A_E)$ 为 $\mathcal A_E$ 中与 $E$ 的交非空的那些长方体的体积之和。由定义, 显然有 $m(\mathcal A_E)\leqslant M(\mathcal A_E)$.
% 
% 设 $\mathcal A_E=\{A_k\}_{k=1}^K,\mathcal A'_E=\{A'_j\}_{j=1}^J$ 分别为两个满足上述条件的长方体族, 称 $\mathcal A'_E$ 是 $\mathcal A_E$ 的细分, 若对每个 $A'_{j_0}\in\mathcal A'_E(j_0=1,2,\cdots,J)$, 都存在 $A_{k_0}\in\mathcal A_E$, 使得 $A'_{j_0}\subset A_{k_0}$。当 $\mathcal A'_E$ 是 $\mathcal A_E$ 的细分时, 我们有
% \[
%   m(\mathcal A_E)\leqslant m(\mathcal A'_E)\leqslant M(\mathcal A'_E)\leqslant M(\mathcal A_E).
% \]
% 因此, 类似于定积分达布理论中的上下积分, 我们定义
% \[
%   m(E)=\sup m(\mathcal A_E),\quad M(E)=\inf M(\mathcal A_E).
% \]
% 在这里 $\sup$ 和 $\inf$ 都是关于所有满足条件 (1) 和 (2) 的长方体族 $\mathcal A_E$ 取的.
% 
% \begin{definition}\label{definition:ma-15-1-1}
% 设 $E\subset\mathbb R^n$ 是有界集合, 若 $E$ 满足
% \[
%   m(E)=M(E),
% \]
% 则称 $E$ 是可求体积的, 并称 $m(E)=M(E)$ 是 $E$ 的体积, 记其为 $V(E)$.
% \end{definition}
% 
% 为了刻画可求体积的集合, 我们给出下述定理.
% 
% \begin{theorem}\label{theorem:ma-15-1-1}
% 设 $E\subset\mathbb R^n$ 是有界集合, 则 $E$ 可求体积的充分必要条件是 $V(\partial E)=0$.
% \end{theorem}
% 
% 
% \begin{proof}
% \keyterm{必要性}\quad 设 $E$ 可求体积, 由定义可知, 对于 $\forall\varepsilon>0$, 存在满足条件 (1), (2) 的长方体族 $\mathcal A_E^1$, 使得
% \[
%   M(\mathcal A_E^1)-m(\mathcal A_E^1)<\varepsilon.
% \]
% 取 $\mathcal A_E^1$ 中与 $\partial E$ 相交的长方体构成的集合为 $\mathcal A_{\partial E}^1$, 则
% \[
%   M(\mathcal A_{\partial E}^1)=M(\mathcal A_E^1)-m(\mathcal A_E^1)<\varepsilon,
% \]
% 从而有
% \[
%   \inf_{\mathcal A_{\partial E}}M(\mathcal A_{\partial E})<\varepsilon.
% \]
% 由 $\varepsilon$ 的任意性知 $V(\partial E)=0$.
% 
% \keyterm{充分性}\quad 设 $V(\partial E)=0$, 则对于 $\forall\varepsilon>0$, 存在有限个长方体构成的集合 $\mathcal A_{\partial E}^2$, 使得 $M(\mathcal A_{\partial E}^2)<\varepsilon$。任取集合 $\mathcal A_E^2$, 使得 $\mathcal A_{\partial E}^2$ 中的长方体都在 $\mathcal A_E^2$ 中, 则
% \[
%   M(\mathcal A_E^2)-m(\mathcal A_E^2)\leqslant M(\mathcal A_{\partial E}^2)<\varepsilon.
% \]
% 因此有 $M(E)=m(E)$, 从而 $E$ 可求体积.
% \end{proof}
% 
% 
% 对于 $\mathbb R^n$ 中的有界集合 $E$, 若 $V(E)=0$, 则称 $E$ 为零体积集。由体积的定义, 若 $E_1$ 与 $E_2$ 为零体积集, 显然 $E_1\cap E_2,E_1\cup E_2$ 及 $E_1\setminus E_2$ 都是零体积集。对于 $\mathbb R^2$ 中的可求体积的有界集合 $E$, 我们有时也称 $E$ 可求面积, 并记 $V(E)=\sigma(E)$.
% 
% \begin{remark}\label{remark:ma-15-source-4}
% 从体积的定义不难看出, 若集合 $E\subset\mathbb R^n$ 的体积为零, 则对于 $\forall\varepsilon>0$, 存在 $\mathbb R^n$ 中有限个长方体 $A_1,A_2,\cdots,A_K$, 使得
% \[
%   \bigcup_{k=1}^K A_k^\circ\supset E,
% \]
% 且
% \[
%   V\left(\bigcup_{k=1}^K A_k\right)<\varepsilon.
% \]
% 由有限个长方体组成的集合也称为简单集合.
% \end{remark}
% 
% 关于 $\mathbb R^n$ 的有界子集的体积, 由定理~\ref{theorem:ma-15-1-1} 我们有以下推论.
% 
% \begin{corollary}\label{corollary:ma-local-15.1-1}
% 设 $A,B$ 是 $\mathbb R^n$ 中可求体积的有界集合, 则 $A\cup B,A\cap B,A\setminus B$ 均可求体积.
% \end{corollary}
% 
% 
% \begin{proof}
% 注意到
% \[
%   \partial(A\cup B)\subset\partial(A)\cup\partial(B),\quad
%   \partial(A\cap B)\subset\partial(A)\cup\partial(B),
% \]
% \[
%   \partial(A\setminus B)\subset\partial(A)\cup\partial(B),
% \]
% 由于 $\partial(A),\partial(B)$ 的体积为零, 从而 $\partial(A\cup B),\partial(A\cap B),\partial(A-B)$ 的体积均为零。由定理~\ref{theorem:ma-15-1-1} 知推论~\ref{corollary:ma-local-15.1-1} 成立.
% \end{proof}
% 
% 
% 特别地, 对于我们经常遇到的区域, 我们容易看出下述结论成立.
% 
% \begin{corollary}\label{corollary:ma-local-15.1-2}
% 设 $D$ 是 $\mathbb R^n$ 中的有界区域, 则 $D$ 可求体积的充分必要条件是 $\overline D$ 可求体积, 并且当它们可求体积时, $V(D)=V(\overline D)$.
% \end{corollary}
% 
% 下面我们来研究一下今后常见的 $\mathbb R^2$ 中的有界区域的面积问题.
% 
% \begin{example}\label{example:ma-15-1-1}
% 设 $\mathbb R^2$ 中的有界区域 $D$ 的边界由有限条可求长曲线所组成, 证明 $D$ 可求面积.
% \end{example}
% 
% 
% \begin{proof}
% 设
% \[
%   \partial D=\bigcup_{j=1}^J\Gamma_j,
% \]
% 其中 $\Gamma_j$ 为可求长曲线。现在我们证明 $\sigma(\partial D)=0$。因此 $D$ 可求面积.
% 
% 事实上, 对于 $\forall j(1\leqslant j\leqslant J)$, 我们只要证 $\sigma(\Gamma_j)=0$ 即可。设 $\Gamma_j$ 的长度为 $l_j$, 对于 $\forall k\in\mathbb N$, 将 $\Gamma_j$ 按弧长等分成 $k$ 个小弧段。易见每一小弧段均可含于一个边长为 $4\dfrac{l_j}{k}$ 的正方形内。因此 $\Gamma_j$ 可被 $k$ 个边长为 $4\dfrac{l_j}{k}$ 的正方形所覆盖。由于
% \[
%   k\left(4\frac{l_j}{k}\right)^2=16\frac{l_j^2}{k}\to0\quad (k\to\infty),
% \]
% 因此有 $\sigma(\Gamma_j)=0$.
% 
% 由例~\ref{example:ma-15-1-1} 可知, 若平面上一个有界区域 $D$, 它是由有限条分段光滑的曲线所围, 则 $D$ 可求面积。作为习题, 请读者证明, 若曲线 $\Gamma$ 是一元连续函数 $y=f(x)(x\in[a,b])$ 的图像, 则 $\sigma(\Gamma)=0$.
% 
% 最后我们指出, 在 $\mathbb R^n$ 中是存在不可求体积的有界闭区域的。由于这类集合的构造相当复杂, 在这里我们就不对此加以介绍了.
% \end{proof}
% 
% 
% \subsection{重积分的定义}
% \label{subsec:ma-15-1-2}
% 
% 设 $D\subset\mathbb R^n$ 是可求体积的有界闭区域, 我们称
% \[
%   \Delta=\{\Delta D_1,\Delta D_2,\cdots,\Delta D_K\}
% \]
% 为 $D$ 的一个分割, 若 $\Delta$ 满足:
% 
% (1) 每个 $\Delta D_k(k=1,2,\cdots,K)$ 是 $D$ 的一个可求体积的闭子集;
% 
% (2) $\Delta D_1,\Delta D_2,\cdots,\Delta D_K$ 两两交集的体积为零, 即
% \[
%   V(\Delta D_j\cap\Delta D_k)=0\quad (j,k=1,2,\cdots,K;\ j\ne k);
% \]
% 
% (3)
% \[
%   D=\bigcup_{k=1}^K\Delta D_k.
% \]
% 
% 记 $\Delta V_k$ 为 $\Delta D_k$ 的体积, $d_k$ 为 $\Delta D_k$ 的直径, 即
% \[
%   d_k=\sup_{x,y\in\Delta D_k}|x-y|\quad (k=1,2,\cdots,K),
% \]
% 并记
% \[
%   \lambda(\Delta)=\max_{1\leqslant k\leqslant K}\{d_k\}.
% \]
% 
% \begin{definition}\label{definition:ma-15-1-2}
% 设函数 $f(x)=f(x_1,x_2,\cdots,x_n)$ 在可求体积的有界闭区域 $D\subset\mathbb R^n$ 上有定义, $\Delta$ 为 $D$ 的一个分割。在每个 $\Delta D_k$ 上任取一点 $\xi_k(k=1,2,\cdots,K)$, 作黎曼和
% \[
%   \sum_{k=1}^K f(\xi_k)\Delta V_k.
% \]
% 如果存在常数 $I$, 使得对于 $\forall\varepsilon>0$, $\exists\delta>0$, 对 $D$ 的任何分割 $\Delta$ 以及任意选取的 $\xi_k$, 当 $\lambda(\Delta)<\delta$ 时, 有
% \[
%   \left|\sum_{k=1}^K f(\xi_k)\Delta V_k-I\right|<\varepsilon,
% \]
% 则称 $f(x)$ 在 $D$ 上可积, 并且称 $I$ 为 $f(x)$ 在 $D$ 上的 $n$ 重积分, 记为
% \[
%   \idotsint_D f(x)\mathrm dV
%   \quad\text{或}\quad
%   \idotsint_D f(x_1,x_2,\cdots,x_n)\mathrm dx_1\mathrm dx_2\cdots\mathrm dx_n,
% \]
% 其中 $f(x)$ 称为被积函数, $D$ 称为积分区域, $x_1,x_2,\cdots,x_n$ 称为积分变量, $\mathrm dV=\mathrm dx_1\mathrm dx_2\cdots\mathrm dx_n$ 称为体积元素.
% \end{definition}
% 
% 在上述定义中, 特别地, 当 $n=2$ 时, 我们常常记被积函数 $f(x)=f(x,y)$, 此时 $f(x,y)$ 在 $D$ 上的二重积分记为
% \[
%   \iint_D f(x,y)\mathrm d\sigma
%   \quad\text{或}\quad
%   \iint_D f(x,y)\mathrm dx\mathrm dy.
% \]
% 这时也称 $\mathrm d\sigma=\mathrm dx\mathrm dy$ 为面积元素。当 $n=3$ 时, 记被积函数 $f(x)=f(x,y,z)$, 则 $f(x,y,z)$ 在 $D$ 上的三重积分记为
% \[
%   \iiint_D f(x,y,z)\mathrm dV
%   \quad\text{或}\quad
%   \iiint_D f(x,y,z)\mathrm dx\mathrm dy\mathrm dz.
% \]
% 在重积分中, 二重积分具有鲜明的几何意义。设 $z=f(x,y)((x,y)\in D)$, 其中 $D\subset\mathbb R^2$ 是可求面积的有界闭区域, 再设 $f(x,y)$ 在 $D$ 上连续且为正函数, 则从二重积分的定义可以看出, $\iint_D f(x,y)\mathrm dx\mathrm dy$ 是以曲面 $z=f(x,y)((x,y)\in D)$ 为顶, 以闭区域 $D$ 为底的曲顶柱体的体积 (见图 15.1.2).
% 
% 特别地, 当 $D\subset\mathbb R^2$ 是可求面积的有界闭区域时, $\iint_D\mathrm dx\mathrm dy$ 等于 $D$ 的面积; 当 $D\subset\mathbb R^3$ 是可求体积的有界闭区域时, $\iiint_D\mathrm dx\mathrm dy\mathrm dz$ 等于 $D$ 的体积.
% 
% \begin{remark}\label{remark:ma-15-source-11}
% 当 $n=1$ 时, 有界闭区域 $D$ 一定是有界闭区间。设函数 $f(x)$ 在区间 $[a,b]$ 上有定义, 上述重积分的定义与 $f(x)$ 在 $[a,b]$ 上的定积分并不一致。这主要是在定积分中我们必须考虑积分区间的有向长度的缘故, 而上述重积分不考虑区间的方向, 只考虑其长度.
% \end{remark}
% 
% 
% \section{多元函数的可积性理论与重积分的性质}
% \label{sec:ma-15-2}
% 
% \subsection{达布理论}
% \label{subsec:ma-15-2-1}
% 
% 在重积分定义中, 我们没有假定 $f(x)$ 在有界闭区域 $D\subset\mathbb R^n$ 上是有界的, 但当 $f(x)$ 可积时, 我们可以推出 $f(x)$ 在 $D$ 上必定有界。事实上, 倘若 $f(x)$ 在 $D$ 上无界, 由于 $D$ 是有界闭区域, 因此它的体积大于零。由有限覆盖定理可知, 存在 $x_0\in D$, 使得对于任何 $x_0$ 的邻域 $U(x_0,\delta)(\delta>0)$, $f(x)$ 在 $U(x_0,\delta)\cap D$ 上都无界。因此对于 $\forall\delta>0$, 令 $\Delta D_\delta$ 为 $U(x_0,\delta)\cap D$ 的闭包, 则 $\Delta D_\delta$ 的体积 $V(\Delta D_\delta)>0$。这是因为, 当 $x_0$ 是 $D$ 的内点时, 这是显然的; 当 $x_0$ 是 $D$ 的边界点时, 由于 $D$ 是 $D^\circ$ 的闭包, 因此 $U(x_0,\delta)$ 必含有 $D^\circ$ 中的点 $x_1$, 从而存在 $\delta'(0<\delta'<\delta)$ 使得 $U(x_1,\delta')\subset\Delta D_\delta$, 故这时也有 $V(\Delta D_\delta)>0$。另外, 显然存在 $D$ 的分割
% \[
%   \Delta=\{\Delta D_\delta,\Delta D_1,\cdots,\Delta D_K\},
% \]
% 且 $\lambda(\Delta)\leqslant2\delta$。如同定积分一样, 我们可以选择 $\xi_\delta\in\Delta D_\delta$, 使得 $f(x)$ 关于 $\Delta$ 的某些黎曼和的绝对值可以大于任何指定的正数。这就说明了 $f(x)$ 在 $D$ 不可积。于是我们有下面的定理.
% 
% \begin{theorem}\label{theorem:ma-15-2-1}
% 设 $D\subset\mathbb R^n$ 是可求体积的有界闭区域, 若 $f(x)$ 在 $D$ 上可积, 则 $f(x)$ 在 $D$ 上有界.
% \end{theorem}
% 
% 在本小节中, 我们总假定 $D$ 是 $\mathbb R^n$ 中可求体积的有界闭区域, $f(x)$ 是 $D$ 上的有界函数, 并记 $M$ 与 $m$ 分别为 $f(x)$ 在 $D$ 上的上、下确界.
% 
% 设
% \[
%   \Delta=\{\Delta D_1,\Delta D_2,\cdots,\Delta D_K\}
% \]
% 是 $D$ 的一个分割, 对于 $k=1,2,\cdots,K$, 记 $\Delta V_k$ 是 $\Delta D_k$ 的体积, $M_k$ 与 $m_k$ 分别为 $f(x)$ 在 $\Delta D_k(k=1,2,\cdots,K)$ 的上、下确界。我们称
% \[
%   \overline S(\Delta)=\sum_{k=1}^K M_k\Delta V_k
% \]
% 为 $f(x)$ 关于分割 $\Delta$ 的达布大和, 而称
% \[
%   \underline S(\Delta)=\sum_{k=1}^K m_k\Delta V_k
% \]
% 为 $f(x)$ 关于分割 $\Delta$ 的达布小和.
% 
% 设 $\Delta_1=\{\Delta D_1,\Delta D_2,\cdots,\Delta D_K\}$ 与 $\Delta_2=\{\Delta D'_1,\Delta D'_2,\cdots,\Delta D'_L\}$ 是 $D$ 的两个分割, 若对于 $\forall l(1\leqslant l\leqslant L)$, 都存在 $k(1\leqslant k\leqslant K)$ 使得 $\Delta D'_l\subset\Delta D_k$, 则称 $\Delta_2$ 是 $\Delta_1$ 的细分, 并记为 $\Delta_1\subset\Delta_2$。不难看出, 对 $D$ 的任意两个分割 $\Delta_1$ 与 $\Delta_2$, 总存在 $D$ 的分割 $\Delta'$, 使得 $\Delta_1\subset\Delta'$ 和 $\Delta_2\subset\Delta'$。事实上, 取
% \[
%   \Delta D_l\cap\Delta D'_k\ne\varnothing\quad (l=1,2,\cdots,L;\ k=1,2,\cdots,K)
% \]
% 组成的分割为 $\Delta'$ 即可.
% 
% 如同定积分的达布理论, 关于达布大和、达布小和, 我们有下列性质:
% 
% (1) 记 $S(\Delta)$ 为函数 $f(x)$ 关于分割 $\Delta$ 的任一黎曼和, 则有
% \[
%   \underline S(\Delta)\leqslant S(\Delta)\leqslant\overline S(\Delta).
% \]
% 
% (2) 若 $D$ 的两个分割 $\Delta_1,\Delta_2$ 满足 $\Delta_2\subset\Delta_1$, 则有
% \[
%   \underline S(\Delta_1)\geqslant\underline S(\Delta_2),\quad
%   \overline S(\Delta_1)\leqslant\overline S(\Delta_2).
% \]
% 
% (3) 对 $D$ 的任意两个分割 $\Delta_1,\Delta_2$, 总有
% \[
%   \underline S(\Delta_1)\leqslant\overline S(\Delta_2).
% \]
% 
% 现记
% \[
%   I^*=\inf_\Delta\overline S(\Delta),\quad
%   I_*=\sup_\Delta\underline S(\Delta),
% \]
% 我们分别称 $I^*,I_*$ 为 $f(x)$ 在 $D$ 上的上积分和下积分.
% 
% \begin{theorem}[{达布定理}]\label{theorem:ma-15-2-2}
% \[
%   \lim_{\lambda(\Delta)\to0}\overline S(\Delta)=I^*,\quad
%   \lim_{\lambda(\Delta)\to0}\underline S(\Delta)=I_*.
% \]
% \end{theorem}
% 
% 利用类似于定积分达布理论的讨论, 我们可以证明下面的结论.
% 
% \begin{theorem}\label{theorem:ma-15-2-3}
% 设函数 $f(x)$ 在可求体积的有界闭区域 $D\subset\mathbb R^n$ 上有界, 则 $f(x)$ 在 $D$ 上可积的充分必要条件是: 对于 $\forall\varepsilon>0$, 存在 $D$ 的分割
% \[
%   \Delta=\{\Delta D_1,\Delta D_2,\cdots,\Delta D_K\},
% \]
% 使得
% \[
%   \sum_{k=1}^K\omega_k\Delta V_k<\varepsilon,
% \]
% 其中 $\omega_k=M_k-m_k$ 为 $f(x)$ 在 $\Delta D_k(k=1,2,\cdots,K)$ 上的振幅, $\Delta V_k$ 为 $\Delta D_k$ 的体积.
% \end{theorem}
% 
% 上述定理以及性质 (1)--(3) 的证明几乎与定积分相应定理、性质的证明方法完全一样, 作为练习, 请读者自己给出.
% 
% 在重积分的定义及可积性讨论中, 当有界闭区域 $D$ 的边界很复杂时, 相应的分割 $\Delta=\{\Delta D_1,\Delta D_2,\cdots,\Delta D_K\}$ 中的每个小集合也可能是很复杂的。这使得对重积分计算公式的推导产生很大困难。值得指出的是, 当 $f(x)$ 在有界闭区域 $D$ 上可积时, 我们可以证明 $f(x)$ 在 $D$ 上的重积分 $I$ 等于一列特别的黎曼和的极限。在该序列中, 黎曼和对应的分割具有下面的形式:
% \[
%   \Delta'=\{\Delta D'_1,\Delta D'_2,\cdots,\Delta D'_J,\Delta D''_1,\Delta D''_2,\cdots,\Delta D''_L\},
% \]
% 其中, 对于 $\forall j(1\leqslant j\leqslant J)$, $\Delta D'_j$ 为完全落在 $D^\circ$ 的小闭区域 (甚至可以是小立方体), 而 $\Delta D''_l\cap\partial D\ne\varnothing(l=1,2,\cdots,L)$。在 $\Delta D'_j$ 上任取 $\xi_j$ 并记 $\Delta V'_j$ 为 $\Delta D'_j$ 的体积 $(j=1,2,\cdots,J)$, 则有
% \[
%   \lim_{\lambda(\Delta')\to0}\sum_{j=1}^J f(\xi_j)\Delta V'_j=I.
% \]
% 
% 事实上, 由 $f(x)$ 在 $D$ 上可积, 从而存在 $M>0$, 使得对于 $\forall x\in D$, 有 $|f(x)|\leqslant M$。另外, 再由 $f(x)$ 的可积性, 从而 $\exists I\in\mathbb R$, 使得对于 $\forall\varepsilon>0$, $\exists\delta>0$, 对于 $D$ 的任意分割
% \[
%   \Delta=\{\Delta D_1,\Delta D_2,\cdots,\Delta D_K\}
% \]
% 及 $\Delta D_k$ 上任意选取的 $\xi_k(k=1,2,\cdots,K)$, 当 $\lambda(\Delta)<\delta$ 时, 有
% \[
%   \left|\sum_{k=1}^K f(\xi_k)\Delta V_k-I\right|<\frac{\varepsilon}{2}
%   \quad\text{和}\quad
%   \sum_{k=1}^K\omega_kV_k<\frac{\varepsilon}{4},
% \]
% 其中 $\omega_k(k=1,2,\cdots,K)$ 为 $f(x)$ 在 $\Delta D_k$ 上的振幅.
% 
% 由 $D$ 可求体积, 总存在有限个长方体集合 $\mathcal A_D=\{A_k\}_{k=1}^{K'}$, 使得
% \[
%   M(\mathcal A_D)-m(\mathcal A_D)<\frac{\varepsilon}{4M}.
% \]
% 显然, 我们可以取得每个小长方体 $A_k$ 充分小, 使得 $\operatorname{diam}(A_k)<\delta(k=1,2,\cdots,K')$。因此
% \[
%   \Delta=\{A_k\cap D\}_{k=1}^{K'}
% \]
% 是 $D$ 的一个分割, 且 $\lambda(\Delta)<\delta$。在每个 $\Delta D_k=A_k\cap D$ 上取 $\xi_k$, 并记 $\Delta V_k$ 为 $\Delta D_k$ 的体积 $(k=1,2,\cdots,K')$, 则有
% \[
%   \left|\sum_{k=1}^{K'} f(\xi_k)\Delta V_k-I\right|<\frac{\varepsilon}{2}.
% \]
% 记
% \[
%   \Delta=\{\Delta D'_1,\Delta D'_2,\cdots,\Delta D'_J,\Delta D''_1,\Delta D''_2,\cdots,\Delta D''_L\},
% \]
% 其中 $\Delta D'_j\subset D^\circ(j=1,2,\cdots,J)$, 其余的 $\Delta D''_j$ 与 $\partial D$ 的交非空。记 $\Delta V'_j$ 为 $\Delta D'_j(j=1,2,\cdots,J)$ 的体积, $\Delta V''_l$ 为 $\Delta D''_l(l=1,2,\cdots,L)$ 的体积。在 $\Delta D'_j$ 上任取 $\xi'_j(j=1,2,\cdots,J)$, 则有
% \[
% \begin{aligned}
%   &\left|\sum_{k=1}^{K'}f(\xi_k)\Delta V_k-\sum_{j=1}^J f(\xi'_j)\Delta V'_j\right|\\
%   &\quad\leqslant
%   \sum_{j=1}^J|f(\xi_j)-f(\xi'_j)|\Delta V'_j+
%   \left|\sum_{l=1}^L f(\xi_l)\Delta V''_l\right|\\
%   &\quad\leqslant
%   \sum_{k=1}^{K'}\omega_k\Delta V_k+M\cdot(M(\mathcal A_D)-m(\mathcal A_D))
%   <\frac{\varepsilon}{4}+\frac{\varepsilon}{4}=\frac{\varepsilon}{2},
% \end{aligned}
% \]
% 从而有
% \[
%   \left|\sum_{j=1}^J f(\xi'_j)\Delta V'_j-I\right|<\varepsilon.
% \]
% 
% 注意到上述黎曼和只与 $D$ 内部的小长方体构成的集合有关。由于以后我们还将用到上面的结论, 我们将其用下面的推论形式给出.
% 
% \begin{corollary}\label{corollary:ma-15-source-15}
% 设函数 $f(x)$ 在可求体积的有界闭区域 $D$ 上可积, 且其积分值为 $I$, 则对于 $\forall\varepsilon>0$, $\exists\delta>0$, 对 $D$ 的特殊分割
% \[
%   \Delta=\{\Delta D'_1,\Delta D'_2,\cdots,\Delta D'_J,\Delta D''_1,\Delta D''_2,\cdots,\Delta D''_L\},
% \]
% 其中 $\Delta D'_j\subset D^\circ(j=1,2,\cdots,J)$, 并且 $\Delta D'_j$ 均是小长方体, 而 $\Delta D''_l\cap\partial D\ne\varnothing(l=1,2,\cdots,L)$, 当 $\lambda(\Delta)<\delta$ 时, 任取 $\xi'_j\in\Delta D'_j(j=1,2,\cdots,J)$, 则有
% \[
%   \left|\sum_{j=1}^J f(\xi'_j)\Delta V'_j-I\right|<\varepsilon,
% \]
% 其中 $\Delta V'_j$ 是 $\Delta D'_j(j=1,2,\cdots,J)$ 的体积.
% \end{corollary}
% 
% 下面我们来讨论一下可积函数类.
% 
% \begin{theorem}\label{theorem:ma-15-2-4}
% 设 $D\subset\mathbb R^n$ 是可求体积的有界闭区域, 再设函数 $f(x)$ 在 $D$ 上连续, 则 $f(x)$ 在 $D$ 上可积.
% \end{theorem}
% 
% 
% \begin{proof}
% 记 $D$ 的体积为 $V>0$。由于 $f(x)$ 在 $D$ 上连续, 且 $D$ 是紧集, 因此 $f(x)$ 在 $D$ 上一致连续, 从而对于 $\forall\varepsilon>0$, $\exists\delta>0$, 当 $D_1\subset D$ 且 $D_1$ 的直径 $\operatorname{diam}(D_1)<\delta$ 时, $f(x)$ 在 $D_1$ 上的振幅 $\omega(D_1)<\dfrac{\varepsilon}{V}$.
% 于是我们可以任取 $D$ 的一个分割 $\Delta=\{\Delta D_1,\Delta D_2,\cdots,\Delta D_K\}$, 使得 $\lambda(\Delta)<\delta$, 这样我们就有
% \[
%   \sum_{k=1}^K\omega_k\Delta V_k<\frac{\varepsilon}{V}\sum_{k=1}^K\Delta V_k=\varepsilon,
% \]
% 其中 $\omega_k$ 是 $f(x)$ 在 $\Delta D_k$ 上的振幅, $\Delta V_k$ 为 $\Delta D_k(k=1,2,\cdots,K)$ 的体积。因此 $f(x)$ 在 $D$ 上可积.
% \end{proof}
% 
% 
% \begin{remark}\label{remark:ma-15-source-18}
% 上面的证明方法可以用来证明下述结论: 设函数 $f(x)$ 在可求体积的有界闭区域 $D$ 上有界, 并且除去 $D$ 的一个零体积集外连续, 则 $f(x)$ 在 $D$ 上可积.
% \end{remark}
% 
% 
% \begin{example}\label{example:ma-15-2-1}
% 设 $\{x_k\}$ 是区间 $[0,1]$ 上所有有理数组成的序列, 定义 $D=[0,1]\times[0,1]$ 上的函数 $f(x,y)$ 如下:
% \[
%   f(x,y)=
%   \begin{cases}
%     \dfrac1k, & x=x_k(k\in\mathbb N),\ 0\leqslant y\leqslant1,\\
%     0, & \text{其他}.
%   \end{cases}
% \]
% 证明 $f(x,y)$ 在 $D$ 上可积, 且
% \[
%   \iint_D f(x,y)\mathrm dx\mathrm dy=0.
% \]
% \end{example}
% 
% 
% \begin{proof}
% 对于 $\forall\varepsilon>0$, $\exists K\in\mathbb N$, 使得 $\dfrac1K<\dfrac{\varepsilon}{2}$。任取 $k>K$, 将 $[0,1]$ 等分成 $k$ 个小区间, 相应地将 $D$ 等分成 $k^2$ 个小正方形
% \[
%   \Delta D_1,\Delta D_2,\cdots,\Delta D_{k^2}.
% \]
% 这 $k^2$ 个小正方形中最多只有 $2kK$ 个小正方形与下述线段
% \[
%   \{(x,y):x=x_1,x_2,\cdots,x_K,\ 0\leqslant y\leqslant1\}
% \]
% 的交非空。记 $f(x,y)$ 在 $\Delta D_l(1\leqslant l\leqslant k^2)$ 上的振幅为 $\omega_l$, 则有
% \[
%   \sum_{l=1}^{k^2}\omega_l\frac1{k^2}
%   \leqslant \frac{2kK}{k^2}+\frac{\varepsilon}{2}\sum_{l=1}^{k^2}\frac1{k^2}
%   =\frac{2K}{k}+\frac{\varepsilon}{2}.
% \]
% 因此取 $k>\dfrac{4K}{\varepsilon}$, 则有
% \[
%   \sum_{l=1}^{k^2}\omega_l\frac1{k^2}<\varepsilon.
% \]
% 由定理~\ref{theorem:ma-15-2-3} 知 $f(x,y)$ 在 $D$ 上可积.
% 
% 对 $D$ 的任意分割 $\Delta=\{\Delta D_1,\Delta D_2,\cdots,\Delta D_K\}$, 我们总可以取 $(\xi_k,\eta_k)\in\Delta D_k(k=1,2,\cdots,K)$, 使得 $f(\xi_k,\eta_k)=0$(除非 $\Delta D_k$ 的面积 $\Delta\sigma_k=0$)。因此
% \[
%   \sum_{k=1}^K f(\xi_k,\eta_k)\Delta\sigma_k=0,
% \]
% 所以
% \[
%   \iint_D f(x,y)\mathrm dx\mathrm dy=0.
% \]
% 
% 读者可以证明上面例子中的函数的间断点集合恰好是
% \[
%   \{(x,y):x=x_k(k\in\mathbb N),\ 0\leqslant y\leqslant1\}.
% \]
% 上例说明, 尽管 $f(x,y)$ 在无穷多条线段上不连续, 但它仍然可积.
% \end{proof}
% 
% 
% \subsection{重积分的性质}
% \label{subsec:ma-15-2-2}
% 
% 对于重积分, 我们同样有定积分所具有的一些相应性质。在下述性质中, 我们总是假定 $D\subset\mathbb R^n$ 为可求体积的有界闭区域.
% 
% \begin{property}\label{property:ma-15-2-5}
% 设函数 $f(x),g(x)$ 在 $D$ 上可积, $\alpha,\beta\in\mathbb R$ 是两常数, 则 $\alpha f(x)+\beta g(x)$ 在 $D$ 上可积, 并且有
% \[
%   \idotsint_D(\alpha f(x)+\beta g(x))\mathrm dV
%   =\alpha\idotsint_D f(x)\mathrm dV+\beta\idotsint_D g(x)\mathrm dV.
% \]
% \end{property}
% 
% 
% \begin{property}\label{property:ma-15-2-6}
% 设函数 $f(x)$ 在 $D$ 上可积, 则 $|f(x)|$ 在 $D$ 上可积, 并且有
% \[
%   \left|\idotsint_D f(x)\mathrm dV\right|
%   \leqslant \idotsint_D |f(x)|\mathrm dV.
% \]
% \end{property}
% 
% 
% \begin{property}\label{property:ma-15-2-7}
% 设 $D_1\subset\mathbb R^n,D_2\subset\mathbb R^n$ 为可求体积的有界闭区域, $D_1^\circ\cap D_2^\circ=\varnothing$, 且 $D_1\cup D_2$ 为可求体积的有界闭区域, 则 $f(x)$ 在 $D_1\cup D_2$ 上可积的充分必要条件是 $f(x)$ 在 $D_1$ 和 $D_2$ 上分别可积, 并且 $f(x)$ 在 $D_1\cup D_2$ 上可积时成立
% \[
%   \idotsint_{D_1\cup D_2}f(x)\mathrm dV
%   =\idotsint_{D_1}f(x)\mathrm dV+\idotsint_{D_2}f(x)\mathrm dV.
% \]
% 在性质~\ref{property:ma-15-2-7} 中, 当 $D_1\cap D_2=\varnothing$ 时, 我们自然地定义 $f(x)$ 在 $D_1\cup D_2$ 上的重积分为 $f(x)$ 在 $D_1$ 与 $D_2$ 上的重积分之和。容易看出这个性质可以推广到有限个两两不交可求体积的有界闭区域的情况, 我们在广义重积分中要遇到这种情况.
% \end{property}
% 
% 
% \begin{property}\label{property:ma-15-2-8}
% 设函数 $f(x),g(x)$ 在 $D$ 上可积, 且对于 $\forall x\in D$, 有 $f(x)\leqslant g(x)$, 则
% \[
%   \idotsint_D f(x)\mathrm dV\leqslant\idotsint_D g(x)\mathrm dV.
% \]
% \end{property}
% 
% 
% \begin{property}\label{property:ma-15-2-9}
% 设函数 $f(x),g(x)$ 在 $D$ 上可积, 则 $f(x)g(x)$ 在 $D$ 上可积.
% \end{property}
% 
% 
% \begin{property}[{重积分第一中值定理}]\label{property:ma-15-2-10}
% 设函数 $f(x)$ 在 $D$ 上连续, $g(x)$ 在 $D$ 上可积且不变号, 则存在 $\xi\in D$, 使得
% \[
%   \idotsint_D f(x)g(x)\mathrm dV=f(\xi)\idotsint_D g(x)\mathrm dV.
% \]
% \end{property}
% 
% 以上这些性质的证明与定积分相应性质的证明类似, 请读者自己给出.
% 
% \section{化重积分为累次积分}
% \label{sec:ma-15-3}
% 
% 我们在前面介绍了 $\mathbb R^n$ 中可求体积的有界闭区域上的重积分理论。在本节及下节中, 我们将讨论重积分的计算问题。我们主要介绍二重积分与三重积分的计算.
% 
% \subsection{化二重积分为累次积分}
% \label{subsec:ma-15-3-1}
% 
% 在讨论一般可求体积的有界闭区域上的二重积分的计算之前, 我们先来讨论一下矩形区域上的二重积分.
% 
% \begin{theorem}\label{theorem:ma-15-3-1}
% 设函数 $f(x,y)$ 在 $D=[a,b]\times[c,d]$ 上可积, 且对于 $\forall x\in[a,b]$, $I(x)=\int_c^d f(x,y)\mathrm dy$ 存在, 则定积分 $\int_a^b I(x)\mathrm dx$ 存在, 并且
% \[
%   \iint_D f(x,y)\mathrm dx\mathrm dy
%   =\int_a^b I(x)\mathrm dx
%   \triangleq\int_a^b\mathrm dx\int_c^d f(x,y)\mathrm dy.
% \]
% \end{theorem}
% 
% 
% \begin{proof}
% 对 $[a,b]$ 作分割
% \[
%   \Delta_x:a=x_0<x_1<\cdots<x_J=b,
% \]
% 并记 $\Delta x_j=x_j-x_{j-1}(j=1,2,\cdots,J)$; 对 $[c,d]$ 作分割
% \[
%   \Delta_y:c=y_0<y_1<\cdots<y_K=d,
% \]
% 并记 $\Delta y_k=y_k-y_{k-1}(k=1,2,\cdots,K)$。过点 $(x_j,0)(j=1,2,\cdots,J)$ 且平行于 $y$ 轴的直线以及过点 $(0,y_k)(k=1,2,\cdots,K)$ 且平行于 $x$ 轴的直线将 $D$ 分成了一些小矩形, 它们组成了 $D$ 的一个分割:
% \[
%   \Delta=\{\Delta D_{jk}:1\leqslant j\leqslant J,\ 1\leqslant k\leqslant K\}.
% \]
% 显然 $\lambda(\Delta)\to0$ 的充分必要条件是 $\lambda(\Delta_x)\to0$ 和 $\lambda(\Delta_y)\to0$。在 $[x_{j-1},x_j]$ 上任取 $\xi_j$, 在 $[y_{k-1},y_k]$ 上任取 $\eta_k$, 则 $(\xi_j,\eta_k)\in\Delta D_{jk}(1\leqslant j\leqslant J,1\leqslant k\leqslant K)$.
% 
% 记
% \[
%   m_{jk}=\inf_{(x,y)\in\Delta D_{jk}}f(x,y),\quad
%   M_{jk}=\sup_{(x,y)\in\Delta D_{jk}}f(x,y),
% \]
% 则有
% \[
%   \sum_{j=1}^J\sum_{k=1}^K m_{jk}\Delta x_j\Delta y_k
%   \leqslant \sum_{j=1}^J\sum_{k=1}^K f(\xi_j,\eta_k)\Delta x_j\Delta y_k
%   \leqslant \sum_{j=1}^J\sum_{k=1}^K M_{jk}\Delta x_j\Delta y_k,
% \]
% 从而有
% \[
% \begin{aligned}
%   \sum_{j=1}^J\sum_{k=1}^K m_{jk}\Delta x_j\Delta y_k
%   &\leqslant \sum_{j=1}^J\left(\int_c^d f(\xi_j,y)\mathrm dy\right)\Delta x_j\\
%   &=\sum_{j=1}^J I(\xi_j)\Delta x_j
%   \leqslant \sum_{j=1}^J\sum_{k=1}^K M_{jk}\Delta x_j\Delta y_k.
% \end{aligned}
% \]
% 由于
% \[
%   \lim_{\lambda(\Delta)\to0}\sum_{j=1}^J\sum_{k=1}^K m_{jk}\Delta x_j\Delta y_k
%   =
%   \lim_{\lambda(\Delta)\to0}\sum_{j=1}^J\sum_{k=1}^K M_{jk}\Delta x_j\Delta y_k
%   =
%   \iint_D f(x,y)\mathrm dx\mathrm dy,
% \]
% 从定积分的定义知, $I(x)$ 在 $[a,b]$ 上可积, 并且
% \[
%   \int_a^b I(x)\mathrm dx
%   =
%   \lim_{\lambda(\Delta_x)\to0}\sum_{j=1}^J I(\xi_j)\Delta x_j
%   =
%   \iint_D f(x,y)\mathrm dx\mathrm dy.
% \qedhere
% \]
% \end{proof}
% 
% 
% \paragraph*{注 1}
% 若 $f(x,y)$ 在 $D=[a,b]\times[c,d]$ 上可积, 且对于 $\forall y\in[c,d]$, $J(y)=\int_a^b f(x,y)\mathrm dx$ 存在, 则类似地可以证明 $\int_c^d J(y)\mathrm dy$ 存在, 并且
% \[
%   \iint_D f(x,y)\mathrm dx\mathrm dy
%   =
%   \int_c^d\mathrm dy\int_a^b f(x,y)\mathrm dx.
% \]
% 今后我们称形如 $\int_a^b\mathrm dx\int_c^d f(x,y)\mathrm dy$ 或 $\int_c^d\mathrm dy\int_a^b f(x,y)\mathrm dx$ 的积分为累次积分.
% 
% \paragraph*{注 2}
% 在矩形区域上的二重积分相当于一个二元函数的极限, 而累次积分相当于一个二元函数的累次极限。因此, 类似于极限与累次极限的关系, 当一个二重积分存在时, 我们不能保证该积分能用累次积分求出。这样的例子可以容易举出。事实上, 我们只要对例~\ref{example:ma-15-2-1} 稍稍加以修改即可, 例如设 $D=[0,1]\times[0,1]$, 当 $(x,y)\in D$ 时, 令
% \[
%   f(x,y)=
%   \begin{cases}
%     \dfrac1k, & x=\dfrac1k(k\in\mathbb N),\ y\text{ 为有理数},\\
%     0, & \text{其他}.
%   \end{cases}
% \]
% 我们来说明 $f(x,y)$ 在 $D$ 上的二重积分不能通过累次积分求出。用与例~\ref{example:ma-15-2-1} 完全相同的方法可以证明 $f(x,y)$ 在 $D$ 上可积且积分为零。但当 $x=\dfrac1k$ 时, $f\left(\dfrac1k,y\right)=\dfrac1k\mathrm D(y)$, 其中 $\mathrm D(y)$ 是 $[0,1]$ 上的狄利克雷函数, 从而它在 $[0,1]$ 上不可积.
% 
% \begin{example}\label{example:ma-15-3-1}
% 设 $D=[0,1]\times[0,1]$, 计算
% \[
%   I=\iint_D x(x-y)^2\mathrm dx\mathrm dy.
% \]
% \end{example}
% 
% 
% \begin{solve}
% 由于被积函数 $x(x-y)^2$ 在 $D$ 上连续, 因此
% \[
% \begin{aligned}
%   I
%   &=\int_0^1 x\mathrm dx\int_0^1(x-y)^2\mathrm dy
%     =\int_0^1 x\left[-\frac13(x-y)^3\right]_0^1\mathrm dx\\
%   &=\frac13\int_0^1x[x^3-(x-1)^3]\mathrm dx
%     =\frac1{12}.
% \end{aligned}
% \]
% 
% 下面我们讨论 $\mathbb R^2$ 中两类特殊有界闭区域上的二重积分问题。在此我们先来复习一下 X 型区域与 Y 型区域的概念。若区域
% \[
%   D=\{(x,y):a\leqslant x\leqslant b,\ \varphi_1(x)\leqslant y\leqslant\varphi_2(x)\},
% \]
% 其中 $\varphi_1(x)$ 与 $\varphi_2(x)$ 是区间 $[a,b]$ 上的连续函数, 且对于 $\forall x\in(a,b)$, 有 $\varphi_1(x)<\varphi_2(x)$, 则称 $D$ 为 X 型区域(见图 15.3.1).
% 
% 若
% \[
%   D=\{(x,y):c\leqslant y\leqslant d,\ \psi_1(y)\leqslant x\leqslant\psi_2(y)\},
% \]
% 其中 $\psi_1(y)$ 与 $\psi_2(y)$ 是区间 $[c,d]$ 上的连续函数, 且对于 $\forall y\in(c,d)$, 有 $\psi_1(y)<\psi_2(y)$, 则称 $D$ 为 Y 型区域(见图 15.3.2).
% 
% 显然任何 X 型区域与 Y 型区域都是可求面积的。对于此类区域上的二重积分, 我们有以下结论.
% \end{solve}
% 
% 
% \begin{theorem}\label{theorem:ma-15-3-2}
% 设 $D=\{(x,y):a\leqslant x\leqslant b,\varphi_1(x)\leqslant y\leqslant\varphi_2(x)\}$ 是 X 型区域, 函数 $f(x,y)$ 在 $D$ 上可积, 且对于 $\forall x\in[a,b]$,
% \[
%   I(x)=\int_{\varphi_1(x)}^{\varphi_2(x)}f(x,y)\mathrm dy
% \]
% 存在, 则定积分 $\int_a^b I(x)\mathrm dx$ 存在, 且
% \begin{equation}
%   \iint_D f(x,y)\mathrm dx\mathrm dy
%   =\int_a^b I(x)\mathrm dx
%   =\int_a^b\mathrm dx\int_{\varphi_1(x)}^{\varphi_2(x)}f(x,y)\mathrm dy.
%   \label{eq:ma-15-3-1}
% \end{equation}
% \end{theorem}
% 
% 
% \begin{proof}
% 由于函数 $\varphi_1(x)$ 与 $\varphi_2(x)$ 在 $[a,b]$ 上连续, 我们取
% \[
%   c=\min_{x\in[a,b]}\varphi_1(x)-1,\qquad
%   d=\max_{x\in[a,b]}\varphi_2(x)+1,
% \]
% 则 $D\subset D_1=[a,b]\times[c,d]$。记
% \[
%   \tilde f(x,y)=
%   \begin{cases}
%     f(x,y), & (x,y)\in D,\\
%     0, & (x,y)\in D_1\setminus D.
%   \end{cases}
% \]
% 由重积分的性质知, $\tilde f(x,y)$ 在 $D_1$ 上可积, 且对于 $\forall x\in[a,b]$,
% \[
%   \int_c^d\tilde f(x,y)\mathrm dy
%   =\int_{\varphi_1(x)}^{\varphi_2(x)}f(x,y)\mathrm dy.
% \]
% 因此
% \[
%   \iint_D f(x,y)\mathrm dx\mathrm dy
%   =\iint_{D_1}\tilde f(x,y)\mathrm dx\mathrm dy
%   =\int_a^b\mathrm dx\int_{\varphi_1(x)}^{\varphi_2(x)}f(x,y)\mathrm dy.
% \qedhere
% \]
% \end{proof}
% 
% 
% \begin{remark}\label{remark:ma-15-source-33}
% 当 $f(x,y)$ 在 Y 型区域
% \[
%   D=\{(x,y):c\leqslant y\leqslant d,\ \psi_1(y)\leqslant x\leqslant\psi_2(y)\}
% \]
% 上可积, 且对于 $\forall y\in[c,d]$, $J(y)=\int_{\psi_1(y)}^{\psi_2(y)}f(x,y)\mathrm dx$ 存在时, 则定积分 $\int_c^dJ(y)\mathrm dy$ 存在, 且
% \begin{equation}
%   \iint_D f(x,y)\mathrm dx\mathrm dy
%   =\int_c^dJ(y)\mathrm dy
%   =\int_c^d\mathrm dy\int_{\psi_1(y)}^{\psi_2(y)}f(x,y)\mathrm dx.
%   \label{eq:ma-15-3-2}
% \end{equation}
% \end{remark}
% 
% 通常我们称 \eqref{eq:ma-15-3-1} 与 \eqref{eq:ma-15-3-2} 右边的积分为重积分 $\iint_D f(x,y)\mathrm dx\mathrm dy$ 的累次积分。另外, 当 $D$ 是一般有界闭区域时, 若添加有限条光滑曲线可以将 $D$ 分成有限个 X 型或 Y 型区域, 我们也可以通过累次积分来计算 $D$ 上的二重积分.
% 
% \begin{example}\label{example:ma-15-3-2}
% 设有界闭区域 $D\subset\mathbb R^2$ 由直线 $y=0$ 及曲线 $y=x^3,x+y=2$ 所围成, $f(x,y)$ 是 $D$ 上的连续函数, 试用两种不同的积分顺序将二重积分 $\iint_D f(x,y)\mathrm dx\mathrm dy$ 化为累次积分.
% \end{example}
% 
% 
% \begin{solve}
% 参照图 15.3.3, 我们有
% \[
% \begin{aligned}
%   \iint_D f(x,y)\mathrm dx\mathrm dy
%   &=\int_0^1\mathrm dy\int_{y^{1/3}}^{2-y}f(x,y)\mathrm dx\\
%   &=\int_0^1\mathrm dx\int_0^{x^3}f(x,y)\mathrm dy
%     +\int_1^2\mathrm dx\int_0^{2-x}f(x,y)\mathrm dy.
% \end{aligned}
% \]
% 
% 从上例中可以看出, 选择不同的积分顺序, 计算的难易程度可能会有所差别.
% \end{solve}
% 
% 
% \begin{example}\label{example:ma-15-3-3}
% 计算累次积分
% \[
%   \int_0^1\mathrm dx\int_x^{\sqrt x}\frac{\sin y}{y}\mathrm dy.
% \]
% \end{example}
% 
% 
% \begin{solve}
% 由于 $\dfrac{\sin y}{y}$ 的原函数无法求出, 累次积分不能直接计算。为此考虑函数
% \[
%   f(x,y)=
%   \begin{cases}
%     \dfrac{\sin y}{y}, & y\ne0,\\
%     1, & y=0
%   \end{cases}
% \]
% 在平面区域 $D=\{(x,y):0\leqslant x\leqslant1,\ x\leqslant y\leqslant\sqrt x\}$ 上的二重积分(见图 15.3.4).
% 
% 由于 $f(x,y)$ 在 $D$ 上连续, 因此该二重积分可以化为累次积分。我们有
% \[
% \begin{aligned}
%   \int_0^1\mathrm dx\int_x^{\sqrt x}\frac{\sin y}{y}\mathrm dy
%   &=\iint_D\frac{\sin y}{y}\mathrm dx\mathrm dy
%     =\int_0^1\mathrm dy\int_{y^2}^{y}\frac{\sin y}{y}\mathrm dx\\
%   &=\int_0^1\frac{(y-y^2)\sin y}{y}\mathrm dy
%     =1-\sin1.
% \end{aligned}
% \]
% 从例~\ref{example:ma-15-3-3} 可以看出, 对于一些二重积分, 其中一个累次积分很难计算出二重积分的值, 而另一个累次积分却能容易地算出二重积分的值.
% \end{solve}
% 
% 
% \begin{example}\label{example:ma-15-3-4}
% 计算 $\mathbb R^3$ 中以柱面 $x^2+y^2=a^2$ 与 $x^2+z^2=a^2(a>0)$ 为边界且含有原点的立体体积 $V$.
% \end{example}
% 
% 
% \begin{solve}
% 由对称性, 该立体在八个卦限中都具有相等的体积。在第一卦限中, 该立体是以区域
% \[
%   D=\{(x,y):x^2+y^2\leqslant a^2,\ x\geqslant0,\ y\geqslant0\}
% \]
% 为底, 以曲面 $z=\sqrt{a^2-x^2}$ 为顶的曲顶柱体(见图 15.3.5).
% 由二重积分的几何意义知
% \[
%   V=8\iint_D z\mathrm dx\mathrm dy.
% \]
% 化二重积分为累次积分即得
% \[
% \begin{aligned}
%   V
%   &=8\iint_D z\mathrm dx\mathrm dy
%     =8\int_0^a\mathrm dx\int_0^{\sqrt{a^2-x^2}}\sqrt{a^2-x^2}\mathrm dy\\
%   &=8\int_0^a(a^2-x^2)\mathrm dx
%     =\frac{16}{3}a^3.
% \end{aligned}
% \qedhere
% \]
% \end{solve}
% 
% 
% \subsection{化三重积分为累次积分}
% \label{subsec:ma-15-3-2}
% 
% 对于三重积分, 当其积分区域具有特殊形状时, 我们可以将三重积分化为累次积分.
% 
% 设 $\Omega\subset\mathbb R^2$ 是可求面积的有界闭区域, $\varphi(x,y),\psi(x,y),(x,y)\in\Omega$ 是 $\Omega$ 上的连续函数, 且对于 $\forall(x,y)\in\Omega^\circ$, 有 $\varphi(x,y)<\psi(x,y)$, 记
% \[
%   D=\{(x,y,z):(x,y)\in\Omega,\ \varphi(x,y)\leqslant z\leqslant\psi(x,y)\}
% \]
% (见图 15.3.6), 则容易推出 $D$ 是可求体积的闭区域(见本章习题)。再设函数 $f(x,y,z)$ 在 $D$ 上连续, 则
% \[
%   I(x,y)=\int_{\varphi(x,y)}^{\psi(x,y)}f(x,y,z)\mathrm dz
% \]
% 在 $\Omega$ 上可积, 且有
% \[
%   \iiint_D f(x,y,z)\mathrm dV
%   =\iint_\Omega I(x,y)\mathrm dx\mathrm dy
%   \triangleq \iint_\Omega\mathrm dx\mathrm dy
%     \int_{\varphi(x,y)}^{\psi(x,y)}f(x,y,z)\mathrm dz.
% \]
% 若三重积分 $\iiint_D f(x,y,z)\mathrm dV$ 的积分区域 $D$ 具有如下两种形式:
% \[
%   D=\{(x,y,z):(y,z)\in\Omega,\ \varphi(y,z)\leqslant x\leqslant\psi(y,z)\},
% \]
% \[
%   D=\{(x,y,z):(x,z)\in\Omega,\ \varphi(x,z)\leqslant y\leqslant\psi(x,z)\},
% \]
% 类似地, 我们也可以将其化为累次积分.
% 
% 现在设空间有界闭区域 $D\subset\mathbb R^3$ 的边界是由有限块分片光滑曲面所组成, 并且当 $(x,y,z)\in D$ 时, 有 $a\leqslant z\leqslant b$, 又设对于 $\forall z\in[a,b]$, 过点 $(0,0,z)$ 且与 $Oxy$ 平面平行的平面截 $D$ 所得到的截面是可求面积的平面闭区域 $D_z$(见图 15.3.7), 则当函数 $f(x,y,z)$ 在 $D$ 上连续时, 有
% \[
%   \iiint_D f(x,y,z)\mathrm dV
%   =\int_a^b\left(\iint_{D_z}f(x,y,z)\mathrm dx\mathrm dy\right)\mathrm dz
%   \triangleq \int_a^b\mathrm dz\iint_{D_z}f(x,y,z)\mathrm dx\mathrm dy.
% \]
% 当用平行于 $Ozx$ 平面或 $Oyz$ 平面的平面截 $D$ 所得的截面为可求面积的平面闭区域时, 我们也有类似的公式.
% 
% \begin{example}\label{example:ma-15-3-5}
% 计算三重积分
% \[
%   I=\iiint_D\frac{\mathrm dx\mathrm dy\mathrm dz}{(1+x+y+z)^3},
% \]
% 其中 $D$ 是由平面 $x=0,y=0,z=0$ 及 $x+y+z=1$ 所围成的闭区域.
% \end{example}
% 
% 
% \begin{solve}
% 记 $\Omega\subset\mathbb R^2$ 为由直线 $x=0,y=0$ 及 $x+y=1$ 所围的闭区域, 见图 15.3.8, 则
% \[
%   D=\{(x,y,z):0\leqslant z\leqslant1-x-y,\ (x,y)\in\Omega\}.
% \]
% 于是有
% \[
% \begin{aligned}
%   I
%   &=\iint_\Omega\mathrm dx\mathrm dy\int_0^{1-x-y}\frac{\mathrm dz}{(1+x+y+z)^3}\\
%   &=\int_0^1\mathrm dx\int_0^{1-x}\mathrm dy\int_0^{1-x-y}\frac{\mathrm dz}{(1+x+y+z)^3}\\
%   &=\int_0^1\mathrm dx\int_0^{1-x}
%     \left.\frac{-1}{2(1+x+y+z)^2}\right|_0^{1-x-y}\mathrm dy\\
%   &=\int_0^1\mathrm dx\int_0^{1-x}\frac12\left[\frac1{(1+x+y)^2}-\frac14\right]\mathrm dy\\
%   &=\frac12\int_0^1\left.\left(\frac{-1}{1+x+y}-\frac y4\right)\right|_0^{1-x}\mathrm dx\\
%   &=\frac12\int_0^1\left(\frac1{1+x}-\frac{3-x}{4}\right)\mathrm dx
%     =\frac12\ln2-\frac5{16}.
% \end{aligned}
% \qedhere
% \]
% \end{solve}
% 
% 
% \begin{example}\label{example:ma-15-3-6}
% 计算三重积分
% \[
%   I=\iiint_D (x+y+z)^2\mathrm dx\mathrm dy\mathrm dz,
% \]
% 其中
% \[
%   D=\left\{(x,y,z):\frac{x^2}{a^2}+\frac{y^2}{b^2}+\frac{z^2}{c^2}\leqslant1,\ a,b,c>0\right\}.
% \]
% \end{example}
% 
% 
% \begin{solve}
% 由积分区域的对称性以及被积函数的对称性容易看出
% \[
%   \iiint_D xy\mathrm dx\mathrm dy\mathrm dz
%   =\iiint_D xz\mathrm dx\mathrm dy\mathrm dz
%   =\iiint_D yz\mathrm dx\mathrm dy\mathrm dz=0,
% \]
% 因此
% \[
%   I=\iiint_D (x^2+y^2+z^2)\mathrm dx\mathrm dy\mathrm dz.
% \]
% 对任意 $x\in[-a,a]$, 过点 $(x,0,0)$ 且平行于 $Oyz$ 平面的平面与 $D$ 的交为
% \[
%   D_x=\left\{(y,z):\frac{y^2}{b^2\left(1-\frac{x^2}{a^2}\right)}
%   +\frac{z^2}{c^2\left(1-\frac{x^2}{a^2}\right)}\leqslant1\right\},
% \]
% 其面积为 $\pi bc\left(1-\dfrac{x^2}{a^2}\right)$。因此有
% \[
% \begin{aligned}
%   \iiint_D x^2\mathrm dx\mathrm dy\mathrm dz
%   &=\int_{-a}^a x^2\mathrm dx\iint_{D_x}\mathrm dy\mathrm dz\\
%   &=\int_{-a}^a x^2\pi bc\left(1-\frac{x^2}{a^2}\right)\mathrm dx
%     =\frac4{15}\pi a^3bc.
% \end{aligned}
% \]
% 同理可计算出
% \[
%   \iiint_D y^2\mathrm dx\mathrm dy\mathrm dz=\frac4{15}\pi ab^3c,\qquad
%   \iiint_D z^2\mathrm dx\mathrm dy\mathrm dz=\frac4{15}\pi abc^3.
% \]
% 最后我们有
% \[
%   I=\frac4{15}\pi abc(a^2+b^2+c^2).
% \qedhere
% \]
% \end{solve}
% 
% 
% \section{重积分的变量替换}
% \label{sec:ma-15-4}
% 
% \subsection{重积分的变量替换公式}
% \label{subsec:ma-15-4-1}
% 
% 在本节中, 我们主要证明二重积分的变量替换公式。对于 $n(n>2)$ 重积分, 我们将不加证明地给出其变量替换公式.
% 
% 设
% \[
%   T(u,v):
%   \begin{cases}
%     x=x(u,v),\\
%     y=y(u,v)
%   \end{cases}
% \]
% 是可求面积的有界闭区域 $D$ 到可求面积的有界闭区域 $\Omega$ 的一个 $C^1$ 同胚映射, 并且假定对于 $\forall(u,v)\in D$ 有
% \[
%   \frac{\partial(x,y)}{\partial(u,v)}\ne0.
% \]
% 由此假定与偏导数的连续性知, $\dfrac{\partial(x,y)}{\partial(u,v)}$ 在 $D$ 上不改变符号。注意到向量函数的导数定义, $T(u,v)$ 在 $(u,v)\in D$ 处导数的行列式即为 $\dfrac{\partial(x,y)}{\partial(u,v)}$.
% 
% 在一元微积分中, 当 $f'(x)$ 在区间 $[a,b]$ 上连续且不变号时, $f(x)$ 在 $[a,b]$ 上单调, 从而
% \[
%   \int_a^b|f'(x)|\mathrm dx
%   =\left|\int_a^bf'(x)\mathrm dx\right|
%   =|f(b)-f(a)|.
% \]
% 这说明 $\int_a^b|f'(x)|\mathrm dx$ 是函数 $y=f(x)$ 的值域的“长度”。因此, 我们有理由猜测
% \[
%   \iint_D\left|\frac{\partial(x,y)}{\partial(u,v)}\right|\mathrm du\mathrm dv
% \]
% 应该是 $\Omega$ 的面积.
% 
% 再从微元法的观点来看, 在点 $(u_0,v_0)\in D$ 处给一增量 $(\mathrm du,\mathrm dv)$, 则以 $(u_0,v_0)$, $(u_0+\mathrm du,v_0)$, $(u_0+\mathrm du,v_0+\mathrm dv)$, $(u_0,v_0+\mathrm dv)$ 为顶点的小矩形的面积为 $\mathrm du\mathrm dv$。而这个小矩形在变换
% \[
%   T(u,v):
%   \begin{cases}
%     x=x(u,v),\\
%     y=y(u,v)
%   \end{cases}
% \]
% 下近似地映成一个平行四边形, 其四个顶点分别为
% \[
% \begin{aligned}
%   P_1&=(x(u_0,v_0),y(u_0,v_0)),\\
%   P_2&=(x(u_0+\mathrm du,v_0),y(u_0+\mathrm du,v_0)),\\
%   P_3&=(x(u_0+\mathrm du,v_0+\mathrm dv),y(u_0+\mathrm du,v_0+\mathrm dv)),\\
%   P_4&=(x(u_0,v_0+\mathrm dv),y(u_0,v_0+\mathrm dv)).
% \end{aligned}
% \]
% 由于 $T(u,v)$ 是 $C^1$ 映射, 因此有
% \[
%   |\overrightarrow{P_1P_2}|\approx
%   \left|\frac{\partial x(u_0,v_0)}{\partial u}\boldsymbol i
%   +\frac{\partial y(u_0,v_0)}{\partial u}\boldsymbol j\right|\mathrm du,
% \]
% \[
%   |\overrightarrow{P_1P_4}|\approx
%   \left|\frac{\partial x(u_0,v_0)}{\partial v}\boldsymbol i
%   +\frac{\partial y(u_0,v_0)}{\partial v}\boldsymbol j\right|\mathrm dv.
% \]
% 该近似平行四边形的面积约为
% \[
%   |\overrightarrow{P_1P_2}\times\overrightarrow{P_1P_4}|
%   \approx
%   \left|\frac{\partial(x,y)}{\partial(u,v)}\right|\mathrm du\mathrm dv.
% \]
% 这也说明了 $\dfrac{\partial(x,y)}{\partial(u,v)}$ 起着一元函数 $f(x)$ 的导数相似的作用。下面我们将给出上述结论的严格证明, 并推出相应的二重积分的变量替换公式.
% 
% \begin{theorem}\label{theorem:ma-15-4-1}
% 设 $D\subset\mathbb R^2$ 是有界闭区域, $\partial D$ 由有限条分段光滑曲线所组成, 变换
% \[
%   T(u,v):
%   \begin{cases}
%     x=x(u,v),\\
%     y=y(u,v)
%   \end{cases}
% \]
% 是 $D$ 到有界闭区域 $\Omega$ 的 $C^1$ 同胚映射, 并且在 $D$ 上处处有
% \[
%   \frac{\partial(x,y)}{\partial(u,v)}\ne0,
% \]
% 再设函数 $f(x,y)$ 在 $\Omega$ 上可积, 则
% \[
%   \iint_\Omega f(x,y)\mathrm dx\mathrm dy
%   =\iint_D f(x(u,v),y(u,v))
%   \left|\frac{\partial(x,y)}{\partial(u,v)}\right|\mathrm du\mathrm dv.
% \]
% \end{theorem}
% 
% 为了证明定理~\ref{theorem:ma-15-4-1}, 我们先做一些准备工作。设 $D_0$ 为 $D$ 内的一个正方形, 它的四个顶点为 $(u_0,v_0)$, $(u_0+\delta,v_0)$, $(u_0+\delta,v_0+\delta)$, $(u_0,v_0+\delta)$, 其中 $\delta>0$。同胚映射 $T$ 将 $D_0$ 映成了区域 $\Omega$ 内的一个由四条光滑曲线组成的一个曲边四边形 $T(D_0)$。如果我们进一步假定变换 $T$ 的四个偏导数 $x'_u,x'_v,y'_u,y'_v$ 中至少有一个在 $D_0$ 上恒大(小)于零, 取 $D_0$ 的某一条对角线将 $D_0$ 分成两个三角形, 则不难看出变换 $T$ 将这两个三角形映成两个 X 型区域(或两个 Y 型区域)。下面我们先证明一个引理.
% 
% \begin{lemma}\label{lemma:ma-15-4-2}
% 假设变换
% \[
%   T(u,v):
%   \begin{cases}
%     x=x(u,v),\\
%     y=y(u,v)
%   \end{cases}
% \]
% 在 $D$ 内具有二阶连续偏导数, $D_0$ 为 $D$ 内的一个正方形并且 $T(D_0)$ 可以分解成有限个由分段光滑曲线所围成的 X 型区域(或 Y 型区域), 则 $T(D_0)$ 的面积
% \[
%   \sigma(T(D_0))=\iint_{D_0}\left|\frac{\partial(x,y)}{\partial(u,v)}\right|\mathrm du\mathrm dv.
% \]
% \end{lemma}
% 
% 
% \begin{proof}
% 在定积分的几何应用中, 我们知道, 对于 $T(D_0)$ 的面积, 有以下的计算公式:
% \[
%   \sigma(T(D_0))=-\int_{\partial T(D_0)}y\mathrm dx
%   =\int_{\partial T(D_0)}x\mathrm dy
%   =\frac12\int_{\partial T(D_0)}x\mathrm dy-y\mathrm dx,
% \]
% 其中 $\partial T(D_0)$ 的方向取正向。曲线 $\partial T(D_0)$ 的方程可以由参数 $(u,v)$ 在 $D_0$ 四条边上的变化来描述。先假设 $T$ 将 $\partial D_0$ 的正向映成了 $\partial T(D_0)$ 的正向, 这时我们有
% \[
% \begin{aligned}
%   -\int_{\partial T(D_0)}y\mathrm dx
%   &=-\int_{u_0}^{u_0+\delta}y(u,v_0)\frac{\partial x(u,v_0)}{\partial u}\mathrm du
%     -\int_{v_0}^{v_0+\delta}y(u_0+\delta,v)\frac{\partial x(u_0+\delta,v)}{\partial v}\mathrm dv\\
%   &\quad-\int_{u_0+\delta}^{u_0}y(u,v_0+\delta)\frac{\partial x(u,v_0+\delta)}{\partial u}\mathrm du
%     -\int_{v_0+\delta}^{v_0}y(u_0,v)\frac{\partial x(u_0,v)}{\partial v}\mathrm dv\\
%   &=\int_{u_0}^{u_0+\delta}
%     \left[y(u,v_0+\delta)\frac{\partial x(u,v_0+\delta)}{\partial u}
%     -y(u,v_0)\frac{\partial x(u,v_0)}{\partial u}\right]\mathrm du\\
%   &\quad-\int_{v_0}^{v_0+\delta}
%     \left[y(u_0+\delta,v)\frac{\partial x(u_0+\delta,v)}{\partial v}
%     -y(u_0,v)\frac{\partial x(u_0,v)}{\partial v}\right]\mathrm dv\\
%   &=\iint_{D_0}\left[
%     \frac{\partial}{\partial v}\left(y\frac{\partial x}{\partial u}\right)
%     -\frac{\partial}{\partial u}\left(y\frac{\partial x}{\partial v}\right)
%     \right]\mathrm du\mathrm dv\\
%   &=\iint_{D_0}\frac{\partial(x,y)}{\partial(u,v)}\mathrm du\mathrm dv.
% \end{aligned}
% \]
% 当 $T$ 将 $\partial D_0$ 的正向映成了 $\partial T(D_0)$ 的负向时, 由上述方法同样可以证明
% \[
%   \sigma(T(D_0))=-\iint_{D_0}\frac{\partial(x,y)}{\partial(u,v)}\mathrm du\mathrm dv
%   =\iint_{D_0}\left|\frac{\partial(x,y)}{\partial(u,v)}\right|\mathrm du\mathrm dv.
% \qedhere
% \]
% \end{proof}
% 
% 
% \begin{remark}\label{remark:ma-15-source-47}
% 在引理~\ref{lemma:ma-15-4-2} 中, 我们对 $x(u,v),y(u,v)$ 假定了它们具有二阶连续偏导数。当 $x(u,v)$ 与 $y(u,v)$ 具有一阶连续偏导数时, 利用具有二阶连续偏导数的函数对它们进行逼近, 则可以证明此时引理的结论仍成立。由于这种逼近过程过于复杂, 我们在此不作介绍, 但我们还是假定引理在此稍弱的条件下仍成立.
% \end{remark}
% 
% 
% \begin{proof}[{定理~\ref{theorem:ma-15-4-1} 的证明}]
% 由于函数 $f(x,y)$ 在 $\Omega$ 上可积, 从而
% \[
%   I=\iint_\Omega f(x,y)\mathrm dx\mathrm dy
% \]
% 存在。由于 $T$ 是 $D$ 到 $\Omega$ 的同胚映射, 因此 $f(x(u,v),y(u,v))$ 在 $D$ 上可积(见本章习题)。再注意到 $\dfrac{\partial(x,y)}{\partial(u,v)}$ 在 $D$ 上连续, 从而
% \[
%   \iint_D f(x(u,v),y(u,v))
%   \left|\frac{\partial(x,y)}{\partial(u,v)}\right|\mathrm du\mathrm dv
% \]
% 存在。下面只要证明它等于 $I$ 即可.
% 
% 取 $\delta>0$ 充分小, 并在 $Ouv$ 平面上作平行于两坐标轴的两个直线族, 使得它们的间距为 $\delta$。这两个直线族对 $D$ 产生了一个分割:
% \[
%   \Delta_D=\{\Delta D_1,\Delta D_2,\cdots,\Delta D_K\},
% \]
% 即每个 $\Delta D_k(k=1,2,\cdots,K)$ 都是一个小正方形与 $D$ 的交集。利用变换 $T$, 我们得到 $\Omega$ 的一个分割
% \[
%   \Delta_\Omega=\{T(\Delta D_1),T(\Delta D_2),\cdots,T(\Delta D_K)\},
% \]
% 由于 $T$ 在 $D$ 上的一致连续性, 当 $\delta\to0$ 时, 有 $\lambda(\Delta_D)\to0$ 和 $\lambda(\Delta_\Omega)\to0$.
% 
% 现将 $\Delta_D$ 中的小闭集分成两类: 记 $\Delta'_D=\{\Delta D'_1,\Delta D'_2,\cdots,\Delta D'_J\}$ 为 $\Delta(D)$ 中所有完全落在 $D$ 内部的正方形的全体, 而记 $\Delta''_D=\{\Delta D''_1,\Delta D''_2,\cdots,\Delta D''_L\}$ 为 $\Delta_D$ 中所有含有 $\partial D$ 的点的集合。显然, $\Delta'_\Omega=\{T(\Delta D'_1),T(\Delta D'_2),\cdots,T(\Delta D'_J)\}$ 中的每个集合均落在 $\Omega$ 的内部, 而 $\Delta''_\Omega=\{T(\Delta D''_1),T(\Delta D''_2),\cdots,T(\Delta D''_L)\}$ 中的每个集合均与 $\partial\Omega$ 的交非空.
% 
% 由假设 $D$ 由有限条分段光滑曲线所围成, 从而 $D$ 可求面积。由此我们有
% \[
%   \lim_{\delta\to0}\sigma\left(\bigcup_{l=1}^L\Delta D''_l\right)=0.
% \]
% 由于变换 $T$ 是 $C^1$ 同胚映射, $\partial\Omega$ 也是由有限条分段光滑曲线所组成。因此对于 $\forall\varepsilon>0$, 存在有限个长方体组成的简单集 $A$, 使得 $\partial\Omega\subset A^\circ$ 且 $\sigma(A)<\varepsilon$。因为 $\partial\Omega$ 与 $\partial A$ 均为紧集, 从而有
% \[
%   \rho=\min_{x\in\partial\Omega,\ y\in\partial A}|x-y|>0.
% \]
% 再由 $T$ 在 $\overline D$ 的一致连续性, 当 $\delta<\dfrac\rho2$ 时, 必有 $T(\Delta D''_l)\subset A^\circ(l=1,2,\cdots,L)$。这说明了
% \[
%   \lim_{\delta\to0}\sigma\left(\bigcup_{l=1}^LT(\Delta D''_l)\right)=0.
% \]
% 注意到在 $D$ 上处处有 $\left|\dfrac{\partial(x,y)}{\partial(u,v)}\right|\ne0$, 对每个 $\Delta D'_j(j=1,2,\cdots,J)$, 由有限覆盖定理不难看出, 若将其等分成 $k^2(k\in\mathbb N)$ 个小正方形, 则当 $k$ 充分大时, 在每个小正方形上 $x'_u,x'_v,y'_u$ 及 $y'_v$ 中至少有一个不取零。因此, $T$ 将该小正方形映成一个曲边小四边形, 这个小曲边四边形必是两个由分段光滑曲线围成的 X 型或 Y 型区域的并。由引理~\ref{lemma:ma-15-4-2} 有
% \[
%   \sigma(T(\Delta D'_j))
%   =\iint_{\Delta D'_j}\left|\frac{\partial(x,y)}{\partial(u,v)}\right|\mathrm du\mathrm dv.
% \]
% 由重积分第一中值定理, 在 $\Delta D'_j(j=1,2,\cdots,J)$ 上存在 $(u'_j,v'_j)$, 使得
% \[
%   \sigma(T(\Delta D'_j))
%   =\left|\frac{\partial(x,y)}{\partial(u,v)}\right|_{(u'_j,v'_j)}
%   \sigma(\Delta D'_j).
% \]
% 
% 记 $x'_j=x(u'_j,v'_j), y'_j=y(u'_j,v'_j)(j=1,2,\cdots,J)$。在 $\Delta D''_l$ 内任取 $(u''_l,v''_l)$, 在 $T(\Delta D''_l)$ 内任取 $(x''_l,y''_l)(l=1,2,\cdots,L)$, 作和式
% \begin{equation}
%   \sum_{j=1}^J f(x'_j,y'_j)\sigma(T(\Delta D'_j))
%   +\sum_{l=1}^L f(x''_l,y''_l)\sigma(T(\Delta D''_l)),
%   \label{eq:ma-15-4-1}
% \end{equation}
% 则当 $\delta\to0$ 时, 上述和式趋于
% \[
%   \iint_\Omega f(x,y)\mathrm dx\mathrm dy.
% \]
% 另外, 我们作和式
% \begin{equation}
% \begin{aligned}
%   &\sum_{j=1}^J f(x(u'_j,v'_j),y(u'_j,v'_j))
%   \left|\frac{\partial(x,y)}{\partial(u,v)}\right|_{(u'_j,v'_j)}
%   \sigma(\Delta D'_j)\\
%   &\quad+\sum_{l=1}^L f(x(u''_l,v''_l),y(u''_l,v''_l))
%   \left|\frac{\partial(x,y)}{\partial(u,v)}\right|_{(u''_l,v''_l)}
%   \sigma(\Delta D''_l),
% \end{aligned}
%   \label{eq:ma-15-4-2}
% \end{equation}
% 则当 $\delta\to0$, 上述和式趋于
% \[
%   \iint_D f(x(u,v),y(u,v))
%   \left|\frac{\partial(x,y)}{\partial(u,v)}\right|\mathrm du\mathrm dv.
% \]
% 由于 $f(x(u,v),y(u,v))$ 和
% \[
%   \left|f(x(u,v),y(u,v))\left|\frac{\partial(x,y)}{\partial(u,v)}\right|\right|
% \]
% 在 $D$ 上可积, 从而存在 $M>0$, 使得当 $(x,y)\in\Omega$ 时, 有 $|f(x,y)|\leqslant M$, 以及当 $(u,v)\in D$ 时, 有
% \[
%   \left|f(x(u,v),y(u,v))\left|\frac{\partial(x,y)}{\partial(u,v)}\right|\right|\leqslant M.
% \]
% 由此我们有
% \[
% \begin{aligned}
%   &\left|
%   \sum_{j=1}^J f(x'_j,y'_j)\sigma(T(\Delta D'_j))
%   +\sum_{l=1}^L f(x''_l,y''_l)\sigma(T(\Delta D''_l))\right.\\
%   &\quad-\left[
%   \sum_{j=1}^J f(x(u'_j,v'_j),y(u'_j,v'_j))
%   \left|\frac{\partial(x,y)}{\partial(u,v)}\right|_{(u'_j,v'_j)}
%   \sigma(\Delta D'_j)\right.\\
%   &\qquad\left.\left.
%   +\sum_{l=1}^L f(x(u''_l,v''_l),y(u''_l,v''_l))
%   \left|\frac{\partial(x,y)}{\partial(u,v)}\right|_{(u''_l,v''_l)}
%   \sigma(\Delta D''_l)\right]\right|\\
%   &=\left|
%   \sum_{l=1}^L f(x''_l,y''_l)\sigma(T(\Delta D''_l))
%   -\sum_{l=1}^L f(x(u''_l,v''_l),y(u''_l,v''_l))
%   \left|\frac{\partial(x,y)}{\partial(u,v)}\right|_{(u''_l,v''_l)}
%   \sigma(\Delta D''_l)\right|\\
%   &\leqslant M\sigma\left(\bigcup_{l=1}^L T(\Delta D''_l)\right)
%   +M\sigma\left(\bigcup_{l=1}^L\Delta D''_l\right)\to0\quad(\delta\to0).
% \end{aligned}
% \]
% 因此有
% \[
%   \iint_\Omega f(x,y)\mathrm dx\mathrm dy
%   =\iint_D f(x(u,v),y(u,v))
%   \left|\frac{\partial(x,y)}{\partial(u,v)}\right|\mathrm du\mathrm dv.
% \qedhere
% \]
% \end{proof}
% 
% 
% 对于 $n$ 重积分, 我们不加证明地给出下述定理.
% 
% \begin{theorem}\label{theorem:ma-15-4-3}
% 设变换
% \[
%   \boldsymbol x(\boldsymbol u)=(x_1(\boldsymbol u),x_2(\boldsymbol u),\cdots,x_n(\boldsymbol u))
% \]
% 是可求体积的有界闭区域 $D\subset\mathbb R^n$ 到可求体积的有界闭区域 $\Omega\subset\mathbb R^n$ 的同胚映射, 它的各个偏导数在包含 $D$ 的区域上连续, 并且在 $D$ 上
% \[
%   \frac{\partial(x_1,\cdots,x_n)}{\partial(u_1,\cdots,u_n)}\ne0,
% \]
% 再设函数 $f(\boldsymbol x)$ 在 $\Omega$ 上可积, 则有
% \[
% \begin{aligned}
%   &\int\!\!\int\cdots\int_\Omega f(x_1,\cdots,x_n)\mathrm dx_1\cdots\mathrm dx_n\\
%   &\quad=\int\!\!\int\cdots\int_D
%   f(x_1(u_1,\cdots,u_n),\cdots,x_n(u_1,\cdots,u_n))
%   \left|\frac{\partial(x_1,\cdots,x_n)}{\partial(u_1,\cdots,u_n)}\right|
%   \mathrm du_1\cdots\mathrm du_n.
% \end{aligned}
% \]
% \end{theorem}
% 
% 
% \subsection{利用变量替换计算重积分}
% \label{subsec:ma-15-4-2}
% 
% 在本小节中, 我们举例说明如何用变量替换公式来计算重积分。在二重积分变量替换中, 极坐标变换
% \[
%   \begin{cases}
%     x=r\cos\theta,\\
%     y=r\sin\theta
%   \end{cases}
%   \quad(0\leqslant r<+\infty,0\leqslant\theta\leqslant2\pi)
% \]
% 是非常重要的变换之一。我们前面已经知道, 它的雅可比行列式为
% \[
%   \left.\frac{\partial(x,y)}{\partial(r,\theta)}\right|_{(r,\theta)}=r.
% \]
% 注意到极坐标变换在正实轴与原点处不是一一映射的, 但它是 $\{(r,\theta):0<r<+\infty,0<\theta<2\pi\}$ 到 $\mathbb R^2\setminus\{(x,y):x\geqslant0\}$ 的 $C^1$ 同胚映射.
% 
% \begin{example}\label{example:ma-15-4-1}
% 计算闭单位圆盘 $\Omega=\{(x,y):x^2+y^2\leqslant1\}$ 的面积.
% \end{example}
% 
% 
% \begin{solve}
% 极坐标变换
% \[
%   \begin{cases}
%     x=r\cos\theta,\\
%     y=r\sin\theta
%   \end{cases}
% \]
% 把 $D=\{(r,\theta):0\leqslant r\leqslant1,0\leqslant\theta\leqslant2\pi\}$ 映成 $\Omega$, 且 $\dfrac{\partial(x,y)}{\partial(r,\theta)}=r$。因此 $\Omega$ 的面积为
% \[
%   \sigma(\Omega)=\iint_\Omega\mathrm dx\mathrm dy
%   =\iint_D r\mathrm dr\mathrm d\theta
%   =\int_0^1 r\mathrm dr\int_0^{2\pi}\mathrm d\theta=\pi.
% \]
% 
% 细心的读者可以发现, 尽管极坐标变换不是 $D$ 到 $\Omega$ 的同胚。但在上例中应用极坐标变换仍然能正确计算出 $\sigma(\Omega)$。我们可以用下面的方法来说明上述计算的合理性。对于 $\forall\varepsilon>0$, 记 $\Omega_\varepsilon=\Omega\setminus\{(x,y):\varepsilon<x<1,-\varepsilon x<y<\varepsilon x\}$, 则极坐标变换是 $D_\varepsilon=\{(r,\theta):\varepsilon\leqslant r\leqslant1,2\pi-\arctan\varepsilon\leqslant\theta\leqslant\arctan\varepsilon\}$ 到 $\Omega_\varepsilon$ 的 $C^1$ 变换。因此
% \[
% \begin{aligned}
%   \sigma(\Omega)
%   &=\lim_{\varepsilon\to0+0}\iint_{\Omega_\varepsilon}\mathrm dx\mathrm dy
%   =\lim_{\varepsilon\to0+0}\iint_{D_\varepsilon}r\mathrm dr\mathrm d\theta\\
%   &=\lim_{\varepsilon\to0+0}\int_\varepsilon^1 r\mathrm dr
%   \int_{\arctan\varepsilon}^{2\pi-\arctan\varepsilon}\mathrm d\theta
%   =\int_0^1 r\mathrm dr\int_0^{2\pi}\mathrm d\theta=\pi.
% \end{aligned}
% \]
% 
% 因此, 若今后我们遇到一个映射在挖掉有限条光滑曲线后的区域之间是 $C^1$ 同胚映射时, 我们可以在原来的闭区域上直接应用定理~\ref{theorem:ma-15-4-1}.
% \end{solve}
% 
% 
% \begin{example}\label{example:ma-15-4-2}
% 计算二重积分
% \[
%   \iint_D(x^2+y^2)\mathrm dx\mathrm dy,
% \]
% 其中 $D$ 是由曲线 $(x^2+y^2)^2=a^2(x^2-y^2)$ 所围且落在第一、四象限的部分, 这里 $a>0$ 为常数.
% \end{example}
% 
% 
% \begin{solve}
% 利用极坐标变换
% \[
%   \begin{cases}
%     x=r\cos\theta,\\
%     y=r\sin\theta
%   \end{cases}
% \]
% 可把曲线方程 $(x^2+y^2)^2=a^2(x^2-y^2)$ 化为
% \[
%   r^4=a^2r^2(\cos^2\theta-\sin^2\theta)=a^2r^2\cos2\theta.
% \]
% 注意到 $x>0$ 时, 有
% \[
%   r=a\sqrt{\cos2\theta},\quad -\frac\pi4\leqslant\theta\leqslant\frac\pi4,
% \]
% 再注意到 $\dfrac{\partial(x,y)}{\partial(r,\theta)}=r$, 我们有
% \[
% \begin{aligned}
%   \iint_D(x^2+y^2)\mathrm dx\mathrm dy
%   &=\int_{-\frac\pi4}^{\frac\pi4}\mathrm d\theta
%   \int_0^{a\sqrt{\cos2\theta}}r^2\cdot r\mathrm dr\\
%   &=2\int_0^{\frac\pi4}\frac14a^4\cos^22\theta\mathrm d\theta
%   =\frac\pi8a^4.
% \end{aligned}
% \]
% 
% 下面我们介绍利用一些特殊的变量替换来计算重积分的例子.
% \end{solve}
% 
% 
% \begin{example}\label{example:ma-15-4-3}
% 计算二重积分
% \[
%   \iint_D(x^2+y^2)\mathrm dx\mathrm dy,
% \]
% 其中 $D$ 是由曲线 $x^2-y^2=1,x^2-y^2=9,xy=2,xy=4$ 所围成的闭区域.
% \end{example}
% 
% 
% \begin{solve}
% 作变换 $T$:
% \[
%   \begin{cases}
%     u=x^2-y^2,\\
%     v=2xy,
%   \end{cases}
% \]
% 则 $T^{-1}$ 将矩形
% \[
%   \Omega=\{(u,v):1\leqslant u\leqslant9,4\leqslant v\leqslant8\}
%   \]
% 一一地映成 $D$。由
% \[
%   \frac{\partial(x,y)}{\partial(u,v)}
%   =\frac1{\dfrac{\partial(u,v)}{\partial(x,y)}}
%   =\frac1{\left|\begin{matrix}2x&-2y\\2y&2x\end{matrix}\right|}
%   =\frac1{4(x^2+y^2)}
%   =\frac1{4\sqrt{u^2+v^2}},
% \]
% 得
% \[
% \begin{aligned}
%   \iint_D(x^2+y^2)\mathrm dx\mathrm dy
%   &=\iint_\Omega\sqrt{u^2+v^2}\frac1{4\sqrt{u^2+v^2}}\mathrm du\mathrm dv\\
%   &=\frac14\iint_\Omega\mathrm du\mathrm dv=8.
% \end{aligned}
% \qedhere
% \]
% \end{solve}
% 
% 
% \begin{example}\label{example:ma-15-4-4}
% 求椭球体
% \[
%   \frac{x^2}{a^2}+\frac{y^2}{b^2}+\frac{z^2}{c^2}\leqslant1\quad(a,b,c>0)
% \]
% 的体积.
% \end{example}
% 
% 
% \begin{solve}
% 取
% \[
%   D=\left\{(x,y):\frac{x^2}{a^2}+\frac{y^2}{b^2}\leqslant1\right\},
% \]
% 则由二重积分的几何意义知, 所求体积为
% \[
%   V=2\iint_D c\sqrt{1-\frac{x^2}{a^2}-\frac{y^2}{b^2}}\mathrm dx\mathrm dy.
% \]
% 作广义极坐标变换
% \[
%   \begin{cases}
%     x=ar\cos\theta,\\
%     y=br\sin\theta,
%   \end{cases}
% \]
% 则 $\dfrac{\partial(x,y)}{\partial(r,\theta)}=abr$, 且
% \[
%   \Omega=\{(r,\theta):0\leqslant r<1,0\leqslant\theta<2\pi\}
% \]
% 经变换一一地映成 $D\setminus\{(0,0)\}$, 从而
% \[
% \begin{aligned}
%   V&=2\iint_\Omega c\sqrt{1-r^2}\,abr\mathrm dr\mathrm d\theta
%   =2\int_0^{2\pi}\mathrm d\theta\int_0^1 c\sqrt{1-r^2}\,abr\mathrm dr\\
%   &=4\pi abc\cdot\frac12\cdot\frac23(1-r^2)^{\frac32}\Big|_0^1
%   =\frac43\pi abc.
% \end{aligned}
% \qedhere
% \]
% \end{solve}
% 
% 
% \begin{example}\label{example:ma-15-4-5}
% 计算二重积分
% \[
%   \iint_D\sqrt{\frac{xy}{x+y}}\mathrm dx\mathrm dy,
% \]
% 其中 $D$ 是由直线 $x=0,y=0,x+y=1$ 所围成的闭区域.
% \end{example}
% 
% 
% \begin{solve}
% 首先我们注意到
% \[
%   \lim_{D\ni(x,y)\to(0,0)}\frac{xy}{x+y}=0,
% \]
% 因此该二重积分存在。显然, 对于积分区域 $D$, 该二重积分易于化成累次积分。但由于被积函数比较复杂, 该积分仍不易求出。为此我们利用下述变换
% \[
%   \begin{cases}
%     x=r\cos^2\theta,\\
%     y=r\sin^2\theta.
%   \end{cases}
% \]
% 该变换将
% \[
%   \Omega=\left\{(r,\theta):0<r<1,0<\theta<\frac\pi2\right\}
% \]
% 一一地变到 $D$ 的内部, 并且
% \[
%   \frac{\partial(x,y)}{\partial(r,\theta)}
%   =\left|\begin{matrix}
%     \cos^2\theta&-r\sin2\theta\\
%     \sin^2\theta&r\sin2\theta
%   \end{matrix}\right|
%   =r\sin2\theta.
% \]
% 因此有
% \[
% \begin{aligned}
%   \iint_D\sqrt{\frac{xy}{x+y}}\mathrm dx\mathrm dy
%   &=\frac12\iint_\Omega r^{\frac12}\sin2\theta\cdot r\sin2\theta\,\mathrm dr\mathrm d\theta\\
%   &=\frac12\int_0^{\frac\pi2}\sin^22\theta\,\mathrm d\theta\int_0^1 r^{\frac32}\mathrm dr
%   =\frac\pi{20}.
% \end{aligned}
% \]
% 
% 对于三重积分的计算, 下面的两种变换是经常用到的.
% \end{solve}
% 
% 
% \paragraph*{柱坐标变换}
% 
% \[
%   \begin{cases}
%     x=r\cos\theta,\\
%     y=r\sin\theta,\\
%     z=z,
%   \end{cases}
% \]
% 其中 $0\leqslant r<+\infty,0\leqslant\theta<2\pi,-\infty<z<+\infty$。这种变换本质上是将空间一点投影到 $Oxy$ 平面内, 然后取它投影点的极坐标(见图 15.4.1)。因此它将
% \[
%   \{(r,\theta,z):0\leqslant r<+\infty,0\leqslant\theta<2\pi,-\infty<z<+\infty\}
% \]
% 映成 $\mathbb R^3$, 其雅可比行列式为
% \[
%   \frac{\partial(x,y,z)}{\partial(r,\theta,z)}=r.
% \]
% 
% \paragraph*{球坐标变换}
% 
% \[
%   \begin{cases}
%     x=r\sin\varphi\cos\theta,\\
%     y=r\sin\varphi\sin\theta,\\
%     z=r\cos\varphi,
%   \end{cases}
% \]
% 其中 $0\leqslant r<+\infty,0\leqslant\varphi\leqslant\pi,0\leqslant\theta<2\pi$(见图 15.4.2)。它的雅可比行列式
% \[
%   \frac{\partial(x,y,z)}{\partial(r,\varphi,\theta)}=r^2\sin\varphi.
% \]
% 显然, 在球坐标下, $r=a$ 对应于球面 $x^2+y^2+z^2=a^2$, $\varphi=b$ 对应于锥面 $z=\cot b\sqrt{x^2+y^2}$, $\theta=c$ 为以 $z$ 轴为边的半平面.
% 
% \begin{example}\label{example:ma-15-4-6}
% 设空间立体 $D$ 由曲面 $z=25-x^2-y^2$ 与平面 $z=0$ 所围成, 试求它的体积 $V$.
% \end{example}
% 
% 
% \begin{solve}
% 由三重积分的几何意义得
% \[
%   V=\iiint_D\mathrm dx\mathrm dy\mathrm dz.
% \]
% 取柱坐标变换
% \[
%   \begin{cases}
%     x=r\cos\theta,\\
%     y=r\sin\theta,\\
%     z=z,
%   \end{cases}
% \]
% 则它将
% \[
%   \Omega=\{(r,\theta,z):0\leqslant r\leqslant5,0\leqslant\theta<2\pi,0\leqslant z\leqslant25-r^2\}
% \]
% 映成 $D$。因此
% \[
% \begin{aligned}
%   V&=\iiint_\Omega\frac{\partial(x,y,z)}{\partial(r,\theta,z)}
%   \mathrm dr\mathrm d\theta\mathrm dz
%   =\int_0^{2\pi}\mathrm d\theta\int_0^5 r\mathrm dr\int_0^{25-r^2}\mathrm dz\\
%   &=2\pi\int_0^5 r(25-r^2)\mathrm dr=\frac{625}2\pi.
% \end{aligned}
% \qedhere
% \]
% \end{solve}
% 
% 
% \begin{remark}\label{remark:ma-15-source-62}
% 上例中的柱坐标变换在 $\Omega$ 和 $D$ 内分别挖掉一个光滑曲面后是 $C^1$ 同胚, 利用处理极坐标变换相似的方法, 容易验证上述的重积分的变量替换公式仍然成立.
% \end{remark}
% 
% 
% \begin{example}\label{example:ma-15-4-7}
% 设 $\mathbb R^3$ 中均匀物体 $D$ 由球面 $x^2+y^2+z^2=1$ 和锥面 $z=\sqrt{x^2+y^2}$ 所围成, 求它的质心坐标.
% \end{example}
% 
% 
% \begin{solve}
% 由对称性, 可设其质心坐标为 $(0,0,z_0)$。锥面方程 $z=\sqrt{x^2+y^2}$ 在球坐标系下为 $\varphi=\dfrac\pi4$。因此 $D$ 的体积
% \[
% \begin{aligned}
%   V&=\iiint_D\mathrm dV
%   =\int_0^{2\pi}\mathrm d\theta\int_0^{\frac\pi4}\mathrm d\varphi\int_0^1 r^2\sin\varphi\mathrm dr\\
%   &=2\pi\int_0^{\frac\pi4}\sin\varphi\frac{r^3}3\Big|_0^1\mathrm d\varphi
%   =\frac{2\pi}3\left(1-\frac{\sqrt2}2\right).
% \end{aligned}
% \]
% 下面还要计算物体到 $Oxy$ 平面的力矩 $M_{xy}$。利用微元法, 我们有
% \[
% \begin{aligned}
%   M_{xy}&=\iiint_D z\mathrm dV
%   =\int_0^{2\pi}\mathrm d\theta\int_0^{\frac\pi4}\mathrm d\varphi\int_0^1
%   r\cos\varphi\cdot r^2\sin\varphi\mathrm dr\\
%   &=2\pi\int_0^{\frac\pi4}\sin\varphi\cos\varphi\mathrm d\varphi\int_0^1r^3\mathrm dr
%   =\frac\pi8.
% \end{aligned}
% \]
% 因此有
% \[
%   z_0=\frac{M_{xy}}V=\frac{3\left(1+\dfrac{\sqrt2}2\right)}8.
% \qedhere
% \]
% \end{solve}
% 
% 
% \section{广义重积分}
% \label{sec:ma-15-5}
% 
% 类似于一元函数的无穷积分与瑕积分, 对于多元函数的重积分我们也有相应的问题: 无界区域上函数的重积分与有界区域上无界函数的重积分。前者称为无穷重积分, 而后者称为瑕重积分, 二者统称为广义重积分。与一元函数的广义积分相比, 多元函数的广义重积分最大的困难是区域的复杂性。因此, 对它有可能得到与一元函数情形具有本质差别的结果.
% 
% \subsection{无穷重积分的基本概念}
% \label{subsec:ma-15-5-1}
% 
% 我们下面来讨论无界区域上函数的重积分。我们只给出二元函数情形的详细讨论, 对于 $n(n>2)$ 元函数的情形, 相应理论可平行推出, 在此不再赘述.
% 
% 广义重积分理论可以讨论非常广泛的无界区域, 但在这里我们只对一些较为简单的区域进行讨论。设 $D\subset\mathbb R^2$ 是无界闭区域, 在本节中, 我们总是假定当 $R$ 充分大时, $D_R\cap D$ 是可求面积闭区域, 其中
% \[
%   D_R=\{(x,y):x^2+y^2\leqslant R\};
% \]
% 并且当 $R_1<R_2$, 且 $R_1$ 充分大时, $D(R_1,R_2)\cap D$ 是有限个内部不相交的可求面积的闭区域的并, 其中
% \[
%   D(R_1,R_2)=\{(x,y):R_1^2\leqslant x^2+y^2\leqslant R_2^2\}.
% \]
% 
% \begin{example}\label{example:ma-15-5-1}
% 设 $D=\{(x,y):0\leqslant x\leqslant1,y\in\mathbb R\}$, 则容易看出, 对于 $\forall R>0$ 及 $\forall R_2>R_1>1$, $D\cap D_R$ 是一个闭区域, 而 $D\cap D(R_1,R_2)$ 是两个闭区域.
% \end{example}
% 
% 
% \begin{definition}\label{definition:ma-15-5-1}
% 设 $D\subset\mathbb R^2$ 是无界闭区域, 函数 $f(x,y)$ 在 $D$ 上有定义, 并且在 $D$ 的任何可求面积的有界闭子区域上可积。若存在常数 $I$, 使得对于 $\forall\varepsilon>0$, $\exists R>0$, 对任何 $D$ 的可求面积的有界闭子区域 $\widetilde D$, 只要 $\widetilde D\supset D_R\cap D$, 就有
% \begin{equation}
%   \left|\iint_{\widetilde D}f(x,y)\mathrm dx\mathrm dy-I\right|<\varepsilon,
%   \label{eq:ma-15-5-1}
% \end{equation}
% 则称 $f(x,y)$ 在 $D$ 上的无穷重积分收敛, 并称 $I$ 为 $f(x,y)$ 在 $D$ 上的无穷重积分, 记为
% \[
%   I=\iint_D f(x,y)\mathrm dx\mathrm dy.
% \]
% 如果不存在 $I\in\mathbb R$, 使得式 \eqref{eq:ma-15-5-1} 成立, 则称无穷重积分
% \[
%   \iint_D f(x,y)\mathrm dx\mathrm dy
% \]
% 发散.
% \end{definition}
% 
% 初学者应特别注意, 上述定义中要求对任何包含 $D_R\cap D$ 的闭区域 $\widetilde D$ 成立式 \eqref{eq:ma-15-5-1}。为了强调这一点我们来看以下例子.
% 
% \begin{example}\label{example:ma-15-5-2}
% 设函数
% \[
%   f(x,y)=\frac{x}{x^2+y^2+1},
% \]
% 试讨论无穷重积分
% \[
%   \iint_{\mathbb R^2} f(x,y)\mathrm dx\mathrm dy
% \]
% 的敛散性.
% \end{example}
% 
% 
% \begin{solve}
% 对于 $\forall R>0$, 取 $r>R$, 考虑
% \[
%   \Omega_r=\{(x,y):|x|\leqslant r, |y|\leqslant r\}
% \]
% 与
% \[
%   Q_r=\{(x,y):r\leqslant x\leqslant2r, |y|\leqslant r\}.
% \]
% 容易看出 $\Omega_r\supset D_R=\{(x,y):\sqrt{x^2+y^2}\leqslant R\}$, 且 $\Omega_r\cup Q_r$ 是一个可求面积的闭区域, 显然它也包含了 $D_R$.
% 
% 下面我们来估计
% \[
%   \iint_{\Omega_r\cup Q_r} f(x,y)\mathrm dx\mathrm dy
%   =\iint_{\Omega_r}\frac{x}{x^2+y^2+1}\mathrm dx\mathrm dy
%   +\iint_{Q_r}\frac{x}{x^2+y^2+1}\mathrm dx\mathrm dy.
% \]
% 由 $\Omega_r$ 的对称性知
% \[
%   \iint_{\Omega_r}\frac{x}{x^2+y^2+1}\mathrm dx\mathrm dy=0,
% \]
% 而
% \[
%   \iint_{Q_r}\frac{x}{x^2+y^2+1}\mathrm dx\mathrm dy
%   \geqslant\iint_{Q_r}\frac{r}{(2r)^2+r^2+1}\mathrm dx\mathrm dy
%   =\frac{r}{1+5r^2}\cdot2r^2.
% \]
% 因此当 $r\to+\infty$ 时, 我们有
% \[
%   \iint_{\Omega_r\cup Q_r} f(x,y)\mathrm dx\mathrm dy\to+\infty.
% \]
% 这说明
% \[
%   \iint_{\mathbb R^2} f(x,y)\mathrm dx\mathrm dy
% \]
% 发散.
% \end{solve}
% 
% 
% \begin{remark}\label{remark:ma-15-source-69}
% 从上述例子可以看出, 尽管无穷重积分
% \[
%   \iint_{\mathbb R^2} f(x,y)\mathrm dx\mathrm dy
% \]
% 是发散的, 但对于 $\forall R>0$, 当 $r>R$ 时, 有 $D_r\supset D_R$ 且
% \[
%   \iint_{D_r} f(x,y)\mathrm dx\mathrm dy=0.
% \]
% \end{remark}
% 
% 
% \subsection{无穷重积分敛散性的判定}
% \label{subsec:ma-15-5-2}
% 
% 受函数极限与序列极限之间关系的启发, 我们对无界闭区域引进穷尽闭区域列的概念, 借助它来讨论无穷重积分的敛散性.
% 
% \begin{definition}\label{definition:ma-15-5-2}
% 设 $D\subset\mathbb R^2$ 是无界闭区域, 若可求面积的有界闭区域序列 $\{D_k\}$ 满足:
% 
% (1) $D_1\subset D_2\subset\cdots\subset D_k\subset\cdots\subset D$;
% 
% (2) 对于 $\forall R>0$, $\exists K\in\mathbb N$, 当 $k>K$ 时, $D_k\supset D\cap D_R$,
% 
% 则称 $\{D_k\}$ 是 $D$ 的一个\keyterm{穷尽闭区域列}(简称穷尽列).
% \end{definition}
% 
% 
% \begin{theorem}\label{theorem:ma-15-5-1}
% 设函数 $f(x,y)$ 在无界闭区域 $D\subset\mathbb R^2$ 上有定义, 在 $D$ 的任何可求面积的有界闭子区域上可积, 则无穷重积分
% \[
%   \iint_D f(x,y)\mathrm dx\mathrm dy
% \]
% 收敛的充分必要条件是: 对 $D$ 的任何一个穷尽列 $\{D_k\}$, 序列极限
% \[
%   \lim_{k\to\infty}\iint_{D_k}f(x,y)\mathrm dx\mathrm dy
% \]
% 存在.
% \end{theorem}
% 
% 作为练习, 请读者自己给出上述定理的证明.
% 
% \begin{theorem}\label{theorem:ma-15-5-2}
% 设 $D\subset\mathbb R^2$ 是无界闭区域, 函数 $f(x,y)$ 在 $D$ 上有定义, 在 $D$ 的任何可求面积的有界闭子区域上可积, 并且对于 $\forall(x,y)\in D$, 有 $f(x,y)\geqslant0$, 则无穷重积分
% \[
%   \iint_D f(x,y)\mathrm dx\mathrm dy
% \]
% 收敛的充分必要条件是: 存在 $D$ 的一个穷尽列 $\{D_k\}$ 及常数 $I>0$, 使得对于 $\forall k\in\mathbb N$, 有
% \[
%   \iint_{D_k}f(x,y)\mathrm dx\mathrm dy\leqslant I.
% \]
% \end{theorem}
% 
% 
% \begin{proof}
% 必要性\quad 设
% \[
%   \iint_D f(x,y)\mathrm dx\mathrm dy
% \]
% 收敛, 令
% \[
%   I=\iint_D f(x,y)\mathrm dx\mathrm dy,
% \]
% 则任取 $D$ 的一个穷尽列 $\{D_k\}$, 对于 $\forall k\in\mathbb N$, 都有
% \[
%   \iint_{D_k}f(x,y)\mathrm dx\mathrm dy\leqslant I.
% \]
% 
% 充分性\quad 设存在 $D$ 的一个穷尽列 $\{D_k\}$ 及常数 $I>0$ 满足: 对于 $\forall k\in\mathbb N$, 有
% \[
%   I_k=\iint_{D_k}f(x,y)\mathrm dx\mathrm dy\leqslant I.
% \]
% 显然 $\{I_k\}$ 是一个单调上升序列且有上界 $I$, 因此它收敛。下面设 $\{D'_j\}$ 是 $D$ 的任何一个穷尽列, 并对 $j=1,2,\cdots$, 记
% \[
%   I'_j=\iint_{D'_j}f(x,y)\mathrm dx\mathrm dy.
% \]
% 由穷尽列的定义, 对于 $\forall j\in\mathbb N$, 必存在 $k_j\in\mathbb N$, 使得有 $k_j>j$ 和 $D'_j\subset D_{k_j}$ 成立。因此有 $I'_j\leqslant I_{k_j}\leqslant I$。这说明 $\{I'_j\}$ 是单调上升有上界的序列, 从而
% \[
%   \lim_{j\to\infty}\iint_{D'_j}f(x,y)\mathrm dx\mathrm dy
% \]
% 收敛。再由定理~\ref{theorem:ma-15-5-1} 知
% \[
%   \iint_D f(x,y)\mathrm dx\mathrm dy
% \]
% 收敛.
% \end{proof}
% 
% 
% \begin{remark}\label{remark:ma-15-source-74}
% 显然, 当 $f(x,y)$ 在 $D$ 上不变号时, 对 $D$ 的任何一个穷尽列 $\{D_k\}$, 都有
% \[
%   \lim_{k\to\infty}\iint_{D_k}f(x,y)\mathrm dx\mathrm dy
%   =\iint_D f(x,y)\mathrm dx\mathrm dy.
% \]
% 上述等式可以理解为: 当
% \[
%   \iint_D f(x,y)\mathrm dx\mathrm dy
% \]
% 收敛时总成立等式, 而当
% \[
%   \iint_D f(x,y)\mathrm dx\mathrm dy
% \]
% 发散时, 等号两边都等于 $+\infty$.
% \end{remark}
% 
% 无穷重积分中与一元函数的无穷积分具有本质区别的是以下的结果.
% 
% \begin{theorem}\label{theorem:ma-15-5-3}
% 设 $D\subset\mathbb R^2$ 是无界闭区域, 函数 $f(x,y)$ 在 $D$ 上有定义, 且在 $D$ 的任何可求面积的有界闭子区域上可积, 则无穷重积分
% \[
%   \iint_D f(x,y)\mathrm dx\mathrm dy
% \]
% 收敛的充分必要条件是无穷重积分
% \[
%   \iint_D |f(x,y)|\mathrm dx\mathrm dy
% \]
% 收敛.
% \end{theorem}
% 
% 
% \begin{proof}
% 充分性\quad 设
% \[
%   \iint_D |f(x,y)|\mathrm dx\mathrm dy
% \]
% 收敛, 并记
% \[
%   \iint_D |f(x,y)|\mathrm dx\mathrm dy=I.
% \]
% 对于 $\forall(x,y)\in D$, 定义
% \[
%   f_+(x,y)=\max\{f(x,y),0\}\quad\text{和}\quad f_-(x,y)=\max\{-f(x,y),0\},
% \]
% 则 $f_+(x,y)\geqslant0$, $f_-(x,y)\geqslant0$, 且对于 $\forall(x,y)\in D$, 有
% \[
%   f(x,y)=f_+(x,y)-f_-(x,y).
% \]
% 设 $\{D_k\}$ 为 $D$ 的任意一个穷尽列, 则对于 $\forall k\in\mathbb N$, 有
% \[
%   \iint_{D_k}f_+(x,y)\mathrm dx\mathrm dy
%   \leqslant\iint_{D_k}|f(x,y)|\mathrm dx\mathrm dy\leqslant I,
% \]
% \[
%   \iint_{D_k}f_-(x,y)\mathrm dx\mathrm dy
%   \leqslant\iint_{D_k}|f(x,y)|\mathrm dx\mathrm dy\leqslant I.
% \]
% 这说明
% \[
%   \iint_D f_+(x,y)\mathrm dx\mathrm dy
% \]
% 与
% \[
%   \iint_D f_-(x,y)\mathrm dx\mathrm dy
% \]
% 均收敛, 从而
% \[
%   \iint_D f(x,y)\mathrm dx\mathrm dy
%   =\iint_D f_+(x,y)\mathrm dx\mathrm dy
%   -\iint_D f_-(x,y)\mathrm dx\mathrm dy
% \]
% 收敛.
% 
% 必要性\quad 设
% \[
%   \iint_D f(x,y)\mathrm dx\mathrm dy
% \]
% 收敛, 倘若
% \[
%   \iint_D |f(x,y)|\mathrm dx\mathrm dy
% \]
% 发散, 我们将推出矛盾.
% 
% 由于
% \[
%   \iint_D f(x,y)\mathrm dx\mathrm dy
% \]
% 收敛, 若记
% \[
%   \iint_D f(x,y)\mathrm dx\mathrm dy=I,
% \]
% 则 $\exists R_1>0$, 当 $R>R_1$ 时,
% \[
%   I-1<\iint_{D\cap D_R}f(x,y)\mathrm dx\mathrm dy<I+1.
% \]
% 由
% \[
%   \iint_D |f(x,y)|\mathrm dx\mathrm dy
% \]
% 发散有
% \[
%   \lim_{R\to+\infty}\iint_{D\cap D_R}|f(x,y)|\mathrm dx\mathrm dy=+\infty.
% \]
% 由于 $f(x,y)=f_+(x,y)-f_-(x,y)$, 因此 $f_+(x,y)$ 与 $f_-(x,y)$ 中必有一个函数, 不妨设 $f_+(x,y)$, 满足
% \[
%   \lim_{R\to+\infty}\iint_{D\cap D_R}f_+(x,y)\mathrm dx\mathrm dy=+\infty.
% \]
% 因此可选取一列 $R_k\to+\infty$, 若记 $D_k=D\cap D_{R_k}$, $D'_k=\overline{D_k\setminus D_{k-1}}$, 则当 $k$ 充分大时, 有
% \[
%   \iint_{D'_k}f_+(x,y)\mathrm dx\mathrm dy>|I|+4.
% \]
% 注意到由我们的假设, $D'_k$ 是有限个可求面积的且内部互不相交的闭区域的并。由重积分的可积性理论可知(参考定理~\ref{theorem:ma-15-2-3} 的推论), 存在 $D'_k$ 的一个分割
% \[
%   \Delta=\{\Delta D_1,\Delta D_2,\cdots,\Delta D_J\},
% \]
% 满足如下性质:
% 
% (1) 设 $\Delta$ 中与 $\partial D'_k$ 交非空的小闭区域为 $\Delta D_1,\cdots,\Delta D_l(l<J)$, 则有
% \[
%   \sum_{j=1}^l M_j\Delta\sigma_j<1,
% \]
% 其中 $M_j$ 为 $f_+(x,y)$ 在 $\Delta D_j$ 的上确界, $\Delta\sigma_j$ 为 $\Delta D_j(j=1,2,\cdots,l)$ 的面积.
% 
% (2) 记 $\Delta$ 中全部落在 $D'_k$ 内部的小闭区域为 $\Delta D_{l+1},\cdots,\Delta D_J$, 则每个 $\Delta D_j(j=l+1,\cdots,J)$ 均为一个矩形, 且满足
% \[
%   \sum_{j=l+1}^J m_j\Delta\sigma_j>|I|+3,
% \]
% 其中 $m_j$ 是 $f_+(x,y)$ 在 $\Delta D_j$ 的下确界, $\Delta\sigma_j$ 为 $\Delta D_j(j=l+1,\cdots,J)$ 的面积。不妨设 $m_j>0(j=l+1,\cdots,J')$, $m_j=0(j=J'+1,\cdots,J)$, 则
% \[
%   \sum_{j=l+1}^{J'}m_j\Delta\sigma_j>|I|+3.
% \]
% 
% 记
% \[
%   D''_k=\bigcup_{j=l+1}^{J'}\Delta D_j,
% \]
% 则当 $k$ 充分大时, 有
% \[
% \begin{aligned}
%   \iint_{D_{k-1}\cup D''_k}f(x,y)\mathrm dx\mathrm dy
%   &=\iint_{D_{k-1}}f(x,y)\mathrm dx\mathrm dy
%   +\iint_{D''_k}f(x,y)\mathrm dx\mathrm dy\\
%   &=\iint_{D_{k-1}}f(x,y)\mathrm dx\mathrm dy
%   +\iint_{D''_k}f_+(x,y)\mathrm dx\mathrm dy\\
%   &\geqslant I-1+\sum_{j=l+1}^{J'}m_j\Delta\sigma_j
%   \geqslant I-1+|I|+3\geqslant I+2.
% \end{aligned}
% \]
% 如果 $D_{k-1}\cup D''_k$ 是闭区域, 则上式与
% \[
%   \iint_D f(x,y)\mathrm dx\mathrm dy=I
% \]
% 矛盾。若 $D_{k-1}\cup D''_k$ 不连通, 注意到 $D''_k$ 是由区域 $D_k^\circ$ 内有限个小矩形组成, 因此可在 $D_k^\circ$ 中取有限个面积很小的狭窄带状闭区域组成集合 $D'''_k$, 使得 $D_{k-1}\cup D''_k\cup D'''_k$ 为有界闭区域, 且
% \[
%   \iint_{D_{k-1}\cup D''_k\cup D'''_k}f(x,y)\mathrm dx\mathrm dy>|I|+1.
% \]
% 这样我们仍然得到矛盾.
% \end{proof}
% 
% 
% 对于无穷重积分, 当区域及被积函数满足一定条件, 我们仍然可以将其化为累次积分或利用变量替换来计算它。我们不加证明地给出下面的定理.
% 
% \begin{theorem}\label{theorem:ma-15-5-4}
% 设函数 $f(x,y)$ 在 $D=[a,+\infty)\times[b,+\infty)$ 上有定义, 且在 $D$ 的任何可求面积的有界闭子区域上可积。若累次积分
% \[
%   \int_a^{+\infty}\mathrm dx\int_b^{+\infty}|f(x,y)|\mathrm dy<+\infty
% \]
% (或
% \[
%   \int_b^{+\infty}\mathrm dy\int_a^{+\infty}|f(x,y)|\mathrm dx<+\infty),
% \]
% 则
% \[
%   \iint_D f(x,y)\mathrm dx\mathrm dy
% \]
% 收敛, 且
% \[
%   \iint_D f(x,y)\mathrm dx\mathrm dy
%   =\int_a^{+\infty}\mathrm dx\int_b^{+\infty}f(x,y)\mathrm dy
% \]
% (或
% \[
%   \iint_D f(x,y)\mathrm dx\mathrm dy
%   =\int_b^{+\infty}\mathrm dy\int_a^{+\infty}f(x,y)\mathrm dx).
% \]
% 若其中的两个累次积分有一个为 $\infty$, 则
% \[
%   \iint_D f(x,y)\mathrm dx\mathrm dy
% \]
% 发散.
% \end{theorem}
% 
% 
% \begin{theorem}\label{theorem:ma-15-5-5}
% 设 $D\subset\mathbb R^2$ 是无界闭区域, 变换
% \[
%   T(u,v):
%   \begin{cases}
%     x=x(u,v),\\
%     y=y(u,v)
%   \end{cases}
%   \quad (u,v)\in D
% \]
% 是 $D$ 到 $T(D)$ 的 $C^1$ 同胚映射, 函数 $f(x,y)$ 在 $T(D)$ 的任何可求面积的有界闭子区域上可积, 则
% \begin{equation}
%   \iint_{T(D)}f(x,y)\mathrm dx\mathrm dy
%   =\iint_D f(x(u,v),y(u,v))
%   \left|\frac{\partial(x,y)}{\partial(u,v)}\right|\mathrm du\mathrm dv.
%   \label{eq:ma-15-5-2}
% \end{equation}
% \end{theorem}
% 
% 
% \begin{remark}\label{remark:ma-15-source-79}
% 等式 \eqref{eq:ma-15-5-2} 应理解为: 若等号两端有一个积分收敛, 则另一个积分也收敛, 并且积分值相等.
% \end{remark}
% 
% 
% \begin{example}\label{example:ma-15-5-3}
% 设函数
% \[
%   f(x,y)=\frac1{(x^2+y^2)^{\frac\alpha2}}=\frac1{r^\alpha},
% \]
% 其中 $r=\sqrt{x^2+y^2}$, $\alpha$ 为常数。分别就下列积分区域讨论无穷重积分
% \[
%   \iint_D f(x,y)\mathrm dx\mathrm dy
% \]
% 的敛散性:
% 
% (1) $D=\{(r,\theta):1\leqslant r<+\infty,0\leqslant\theta_1\leqslant\theta\leqslant\theta_2\leqslant2\pi\}$;
% 
% (2) $D=\{(x,y):1\leqslant x<+\infty,0\leqslant y\leqslant1\}$.
% \end{example}
% 
% 
% \begin{solve}
% (1) 由极坐标变换得
% \[
%   \iint_D\frac{\mathrm dx\mathrm dy}{(x^2+y^2)^{\frac\alpha2}}
%   =\lim_{R\to+\infty}\int_{\theta_1}^{\theta_2}\mathrm d\theta
%   \int_1^R\frac r{r^\alpha}\mathrm dr.
% \]
% 由一元函数无穷积分的收敛性知, 当 $\alpha-1>1$, 即 $\alpha>2$ 时,
% \[
%   \iint_D\frac{\mathrm dx\mathrm dy}{(x^2+y^2)^{\frac\alpha2}}
% \]
% 收敛; 而当 $\alpha\leqslant2$ 时, 该无穷重积分发散.
% 
% (2) 在对于 $\forall(x,y)\in D$ 中有
% \[
%   \frac1{(x^2+1)^{\frac\alpha2}}\leqslant
%   \frac1{(x^2+y^2)^{\frac\alpha2}}\leqslant\frac1{x^\alpha},
% \]
% 且
% \[
% \begin{aligned}
%   \iint_D\frac1{(x^2+1)^{\frac\alpha2}}\mathrm dx\mathrm dy
%   &=\lim_{X\to+\infty}\int_0^1\mathrm dy\int_1^X
%   \frac{\mathrm dx}{(x^2+1)^{\frac\alpha2}}\\
%   &=\lim_{X\to+\infty}\int_1^X
%   \frac{\mathrm dx}{(x^2+1)^{\frac\alpha2}},
% \end{aligned}
% \]
% 及
% \[
%   \iint_D\frac1{x^\alpha}\mathrm dx\mathrm dy
%   =\lim_{X\to+\infty}\int_1^X\frac{\mathrm dx}{x^\alpha}.
% \]
% 由一元非负函数无穷积分敛散性的判别法知, 当 $\alpha>1$ 时,
% \[
%   \iint_D\frac1{x^\alpha}\mathrm dx\mathrm dy
% \]
% 收敛, 从而
% \[
%   \iint_D\frac{\mathrm dx\mathrm dy}{(x^2+y^2)^{\frac\alpha2}}
% \]
% 收敛; 当 $\alpha\leqslant1$ 时,
% \[
%   \iint_D\frac{\mathrm dx\mathrm dy}{(x^2+1)^{\frac\alpha2}}
% \]
% 发散, 从而
% \[
%   \iint_D\frac{\mathrm dx\mathrm dy}{(x^2+y^2)^{\frac\alpha2}}
% \]
% 发散.
% \end{solve}
% 
% 
% \begin{remark}\label{remark:ma-15-source-82}
% 当 $1<\alpha\leqslant2$ 时,
% \[
%   \frac1{(x^2+y^2)^{\frac\alpha2}}
% \]
% 在 (1) 中 $D$ 上的无穷重积分发散, 而在 (2) 中 $D$ 上的无穷重积分收敛。这说明相同的函数在不同区域上的积分可能具有不同的敛散性。读者不难证明: 对任何一个 $\mathbb R^2$ 内的连续函数 $f(x,y)$, 总可找到一个无界闭区域 $D$, 使得
% \[
%   \iint_D f(x,y)\mathrm dx\mathrm dy
% \]
% 收敛。反过来, 对于任何一个无界闭区域 $D\subset\mathbb R^2$, 总可以找到一个 $D$ 上的正的连续函数 $g(x,y)$, 使得
% \[
%   \iint_D g(x,y)\mathrm dx\mathrm dy
% \]
% 发散.
% \end{remark}
% 
% 
% \begin{example}\label{example:ma-15-5-4}
% 讨论无穷重积分
% \[
%   \iint_D\frac{x^2-y^2}{(x^2+y^2)^2}\mathrm dx\mathrm dy
% \]
% 的敛散性, 其中
% \[
%   D=\{(x,y):1\leqslant x<+\infty,1\leqslant y<+\infty\}.
% \]
% \end{example}
% 
% 
% \begin{solve}
% 取定 $0<\theta_1<\theta_2<\pi/8$, 利用极坐标变换, 则存在 $R_0>0$, 使得
% \[
%   \Omega=\{(r,\theta):R_0\leqslant r<+\infty,\theta_1\leqslant\theta\leqslant\theta_2\}
% \]
% 在变换
% \[
%   \begin{cases}
%     x=r\cos\theta,\\
%     y=r\sin\theta
%   \end{cases}
% \]
% 下的像 $D_1$ 落在 $D$ 内。由
% \[
% \begin{aligned}
%   \iint_D\frac{|x^2-y^2|}{(x^2+y^2)^2}\mathrm dx\mathrm dy
%   &\geqslant\iint_{D_1}\frac{|x^2-y^2|}{(x^2+y^2)^2}\mathrm dx\mathrm dy\\
%   &=\iint_\Omega\frac{r^2\cos2\theta}{r^4}r\mathrm dr\mathrm d\theta\\
%   &=\lim_{R\to+\infty}\int_{\theta_1}^{\theta_2}\mathrm d\theta
%   \int_{R_0}^R\cos2\theta\frac{\mathrm dr}r=+\infty
% \end{aligned}
% \]
% 知原无穷重积分发散.
% \end{solve}
% 
% 
% \begin{example}\label{example:ma-15-5-5}
% 通过计算
% \[
%   \iint_{\mathbb R^2}e^{-(x^2+y^2)}\mathrm dx\mathrm dy
% \]
% 求
% \[
%   \int_0^{+\infty}e^{-x^2}\mathrm dx
% \]
% 的值.
% \end{example}
% 
% 
% \begin{solve}
% 取
% \[
%   D_R=\{(x,y):x^2+y^2\leqslant R^2\},
% \]
% 则由极坐标变换得
% \[
%   \lim_{R\to+\infty}\iint_{D_R}e^{-(x^2+y^2)}\mathrm dx\mathrm dy
%   =\lim_{R\to+\infty}\int_0^{2\pi}\mathrm d\theta\int_0^R e^{-r^2}r\mathrm dr=\pi.
% \]
% 因此无穷重积分
% \[
%   \iint_{\mathbb R^2}e^{-(x^2+y^2)}\mathrm dx\mathrm dy
% \]
% 收敛, 并且它的值为 $\pi$.
% 
% 现取
% \[
%   N_R=\{(x,y):-R\leqslant x\leqslant R,-R\leqslant y\leqslant R\},
% \]
% 则
% \[
% \begin{aligned}
%   \pi&=\iint_{\mathbb R^2}e^{-(x^2+y^2)}\mathrm dx\mathrm dy
%   =\lim_{R\to+\infty}\iint_{N_R}e^{-(x^2+y^2)}\mathrm dx\mathrm dy\\
%   &=\lim_{R\to+\infty}\int_{-R}^R e^{-x^2}\mathrm dx
%   \int_{-R}^R e^{-y^2}\mathrm dy
%   =\left(\int_{-\infty}^{+\infty}e^{-x^2}\mathrm dx\right)^2.
% \end{aligned}
% \]
% 所以
% \[
%   \int_0^{+\infty}e^{-x^2}\mathrm dx=\frac{\sqrt\pi}2.
% \qedhere
% \]
% \end{solve}
% 
% 
% \subsection{瑕重积分}
% \label{subsec:ma-15-5-3}
% 
% 设 $D\subset\mathbb R^2$ 是可求面积的有界闭区域, $x_0=(x_0,y_0)\in D$。在本节中, 我们对 $D$ 作如下假定: 对每个充分小的 $\delta>0$, $D\setminus U(x_0,\delta)$ 是一个闭区域; 若记
% \[
%   A(x_0,\delta_1,\delta_2)=\{x:0<\delta_1\leqslant |x-x_0|\leqslant\delta_2\},
% \]
% 则当 $\delta_2$ 充分小时, $D\cap A(x_0,\delta_1,\delta_2)$ 是有限个两两内部不交的可求面积的闭区域之并。我们引入如下瑕重积分敛散性的定义:
% 
% \begin{definition}\label{definition:ma-15-5-3}
% 设 $D\subset\mathbb R^2$ 是可求面积的有界闭区域, $x_0\in D$, 函数 $f(x,y)$ 在 $D\setminus\{x_0\}$ 上有定义且无界, 在 $D\setminus\{x_0\}$ 的任何可求面积的闭子区域上可积。若存在常数 $I$, 使得对于 $\forall\varepsilon>0$, $\exists\delta>0$, 对 $D\setminus\{x_0\}$ 的任何可求面积的闭子区域 $D'$, 只要 $D'\supset D\setminus U(x_0,\delta)$, 就有
% \begin{equation}
%   \left|\iint_{D'}f(x,y)\mathrm dx\mathrm dy-I\right|<\varepsilon,
%   \label{eq:ma-15-5-3}
% \end{equation}
% 则称瑕重积分
% \[
%   \iint_D f(x,y)\mathrm dx\mathrm dy
% \]
% 收敛, 同时称 $I$ 为 $f(x,y)$ 在 $D$ 上的瑕重积分, 并记为
% \[
%   \iint_D f(x,y)\mathrm dx\mathrm dy=I.
% \]
% 若不存在常数 $I$, 使得式 \eqref{eq:ma-15-5-3} 成立, 则称瑕重积分
% \[
%   \iint_D f(x,y)\mathrm dx\mathrm dy
% \]
% 发散.
% \end{definition}
% 
% 
% \begin{remark}\label{remark:ma-15-source-88}
% 我们把定义~\ref{definition:ma-15-5-3} 中的 $x_0=(x_0,y_0)$ 称为 $f(x,y)$ 的瑕点。如同讨论无穷重积分敛散性时引入的穷尽列, 称 $\{D_k\}$ 为可求体积的有界闭区域 $D$ 的一个穷尽列, 若 $D_1\subset D_2\subset\cdots\subset D_k\subset\cdots\subset D$, 且对于 $\forall\delta>0$, $\exists K\in\mathbb N$, 当 $k>K$ 时, 有 $D_k\supset D\setminus U(x_0,\delta)$。我们也有以下类似于判定无穷重积分敛散性的定理.
% \end{remark}
% 
% 
% \begin{theorem}\label{theorem:ma-15-5-6}
% 设 $D\subset\mathbb R^2$ 是可求面积的有界闭区域, $x_0=(x_0,y_0)\in D$, 函数 $f(x,y)$ 在 $D\setminus\{x_0\}$ 上有定义且无界, 在 $D\setminus\{x_0\}$ 的任何可求面积的闭子区域上可积, 则瑕重积分
% \[
%   \iint_D f(x,y)\mathrm dx\mathrm dy
% \]
% 收敛的充分必要条件是: 对 $D$ 的任何穷尽列 $\{D_k\}$, 序列极限
% \[
%   \lim_{k\to\infty}\iint_{D_k}f(x,y)\mathrm dx\mathrm dy
% \]
% 存在.
% \end{theorem}
% 
% 
% \begin{theorem}\label{theorem:ma-15-5-7}
% 设 $D\subset\mathbb R^2$ 是可求面积的有界闭区域, $x_0=(x_0,y_0)\in D$, 函数 $f(x,y)$ 在 $D\setminus\{x_0\}$ 上有定义、非负且无界, 在 $D\setminus\{x_0\}$ 的任何可求面积的闭子区域上可积, 则瑕重积分
% \[
%   \iint_D f(x,y)\mathrm dx\mathrm dy
% \]
% 收敛的充分必要条件是: 存在常数 $I>0$ 和 $D$ 的穷尽列 $\{D_k\}$, 使得对于 $\forall k\in\mathbb N$, 有
% \[
%   \iint_{D_k}f(x,y)\mathrm dx\mathrm dy\leqslant I.
% \]
% \end{theorem}
% 
% 
% \begin{theorem}\label{theorem:ma-15-5-8}
% 设 $D\subset\mathbb R^2$ 是可求面积的有界区域, $x_0=(x_0,y_0)\in D$, 函数 $f(x,y)$ 在 $D\setminus\{x_0\}$ 上有定义且无界, 在 $D\setminus\{x_0\}$ 的任何可求面积的闭子区域上可积, 则瑕重积分
% \[
%   \iint_D f(x,y)\mathrm dx\mathrm dy
% \]
% 收敛的充分必要条件是瑕重积分
% \[
%   \iint_D |f(x,y)|\mathrm dx\mathrm dy
% \]
% 收敛.
% \end{theorem}
% 
% 上述这些定理的证明与无穷重积分相应定理的证明相似, 请读者自证。另外, $f(x,y)$ 在 $D$ 内一点的邻域内无界的瑕重积分, 可以推广到多个点, 甚至在一些曲线上无界的情形, 请读者给出它们相应的定义与性质, 在此我们就不再赘述了.
% 
% \begin{example}\label{example:ma-15-5-6}
% 试讨论瑕重积分
% \[
%   \iint_D\frac{\mathrm dx\mathrm dy}{(x^2+y^2)^{\frac\alpha2}}
% \]
% 的敛散性, 其中
% \[
%   D=\{(x,y):x^2+y^2\leqslant1\}.
% \]
% \end{example}
% 
% 
% \begin{solve}[{解法 1}]
% 显然, 当 $\alpha>0$ 时, 原点是被积函数的瑕点。由极坐标变换得
% \[
%   \iint_D\frac{\mathrm dx\mathrm dy}{(x^2+y^2)^{\frac\alpha2}}
%   =\lim_{\delta\to0+0}\int_0^{2\pi}\mathrm d\theta\int_\delta^1\frac r{r^\alpha}\mathrm dr
%   =\lim_{\delta\to0+0}2\pi\int_\delta^1r^{1-\alpha}\mathrm dr,
% \]
% 因此, 当 $\alpha<2$ 时, 原瑕重积分收敛; 当 $\alpha\geqslant2$ 时, 原瑕重积分发散.
% \end{solve}
% 
% 
% \begin{solve}[{解法 2}]
% 记
% \[
%   D^*=\{(s,t):s^2+t^2\geqslant1\}.
% \]
% 作变换
% \[
%   \begin{cases}
%     x=\dfrac{s}{s^2+t^2},\\
%     y=\dfrac{-t}{s^2+t^2},
%   \end{cases}
%   \quad(s,t)\in D^*,
% \]
% 则变换将 $D^*$ 映成 $D\setminus\{(0,0)\}$, 且
% \[
%   \frac{\partial(x,y)}{\partial(s,t)}=\frac1{(s^2+t^2)^2}.
% \]
% 因此
% \[
%   \iint_D\frac{\mathrm dx\mathrm dy}{(x^2+y^2)^{\frac\alpha2}}
%   =\iint_{D^*}\frac{(s^2+t^2)^{\frac\alpha2}}{(s^2+t^2)^2}\mathrm ds\mathrm dt
%   =\iint_{D^*}\frac{\mathrm ds\mathrm dt}{(s^2+t^2)^{2-\frac\alpha2}}.
% \]
% 由例~\ref{example:ma-15-5-3} 知, 当 $2-\dfrac\alpha2>1$, 即 $\alpha<2$ 时, 原瑕重积分收敛; 当 $2-\dfrac\alpha2\leqslant1$, 即 $\alpha\geqslant2$ 时, 原瑕重积分发散.
% \end{solve}
% 
% 
% \paragraph*{注 1}
% 结合例~\ref{example:ma-15-5-3} 和例~\ref{example:ma-15-5-6}, 我们推知, 对任意的 $\alpha\in\mathbb R$, 广义重积分
% \[
%   \iint_{\mathbb R^2}\frac{\mathrm dx\mathrm dy}{(x^2+y^2)^{\frac\alpha2}}
% \]
% 都发散.
% 
% \paragraph*{注 2}
% 由解法 2, 我们容易将有界闭区域上具有一个瑕点的无界函数的重积分与无界区域上的重积分进行转化。例如, 从例~\ref{example:ma-15-5-3} 我们可以容易地构造以原点为心的单位圆盘 $D$ 内的一个闭区域 $\Omega$, 使得 $0\in\partial\Omega$ 且当 $2\leqslant\alpha<3$ 时, 无穷重积分
% \[
%   \iint_\Omega\frac{\mathrm dx\mathrm dy}{(x^2+y^2)^{\frac\alpha2}}
% \]
% 仍然收敛.
% 
% 事实上, 取 $(s,t)$ 平面上的闭区域
% \[
%   G=\{(s,t):1\leqslant s<+\infty,0\leqslant t\leqslant1\},
% \]
% 然后利用变换
% \[
%   \begin{cases}
%     x=\dfrac{s}{s^2+t^2},\\
%     y=\dfrac{-t}{s^2+t^2}
%   \end{cases}
% \]
% 将它变到单位圆盘内, 容易看出其像集合并上 $(0,0)$ 为一个可求面积的有界闭区域 $\Omega$。利用例~\ref{example:ma-15-5-3}, 可知
% \[
%   \iint_\Omega\frac{\mathrm dx\mathrm dy}{(x^2+y^2)^{\alpha/2}}
% \]
% 必收敛.
% 
% 最后, 我们举一个例子来说明如何判别具有多个瑕点的多元函数瑕积分的敛散性.
% 
% \begin{example}\label{example:ma-15-5-7}
% 设函数 $f(x)$ 在区间 $[0,1]$ 上连续, 且对于 $\forall x\in[0,1]$, 有 $f(x)>0$, 又设闭区域
% \[
%   D=\{(x,y):0\leqslant x\leqslant1,0\leqslant y\leqslant f(x)\}.
% \]
% 试讨论瑕重积分
% \[
%   \iint_D\frac{\mathrm dx\mathrm dy}{|y-f(x)|^\alpha}
% \]
% 的敛散性, 其中 $\alpha<1$.
% \end{example}
% 
% 
% \begin{solve}
% 设 $f(x)$ 在 $[0,1]$ 上的最大值为 $M$。取 $r>0$, 令
% \[
%   D_r=\{(x,y):0\leqslant x\leqslant1,0\leqslant y\leqslant f(x)-r\},
% \]
% 则
% \[
% \begin{aligned}
%   \iint_D\frac{\mathrm dx\mathrm dy}{|y-f(x)|^\alpha}
%   &=\lim_{r\to0+0}\iint_{D_r}\frac{\mathrm dx\mathrm dy}{|y-f(x)|^\alpha}\\
%   &=\lim_{r\to0+0}\int_0^1\mathrm dx\int_0^{f(x)-r}
%   \frac{\mathrm dy}{|y-f(x)|^\alpha}.
% \end{aligned}
% \]
% 将 $x$ 看成常数, 对定积分
% \[
%   \int_0^{f(x)-r}\frac{\mathrm dy}{|y-f(x)|^\alpha}
% \]
% 作变换 $y=f(x)-t$, 则有
% \[
% \begin{aligned}
%   &\lim_{r\to0+0}\int_0^1\mathrm dx\int_0^{f(x)-r}
%   \frac{\mathrm dy}{|y-f(x)|^\alpha}\\
%   &\quad=\lim_{r\to0+0}\int_0^1\mathrm dx\int_r^{f(x)}\frac{\mathrm dt}{t^\alpha}\\
%   &\quad\leqslant\lim_{r\to0+0}\int_0^1\mathrm dx\int_r^M\frac{\mathrm dt}{t^\alpha}
%   =\int_0^M\frac{\mathrm dt}{t^\alpha}<+\infty.
% \end{aligned}
% \]
% 因此原瑕重积分收敛.
% \end{solve}
% 
% % 整合来源：math-16.tex
% \chapter{曲线积分与曲面积分}
% \label{chap:ma-16}
% 
% 在 $\mathbb R^n(n\geqslant2)$ 中具有很多有趣的子集, 如 $\mathbb R^2$ 中的曲线, $\mathbb R^3$ 中的曲线与曲面等。当 $n>3$ 时, 除了曲线与曲面外有趣的子集就更多了。在科学技术的许多问题中, 人们会遇到这些子集上的积分问题。在本章中, 我们将研究定义在 $\mathbb R^3$ (或 $\mathbb R^2$) 中一些特殊几何体上的积分问题, 主要考虑两种简单情况, 即曲线和曲面上的积分。根据这些积分的不同背景, 人们将它们分成不同类型的积分。我们主要介绍两类积分: 一类是多元函数的积分, 这类积分称为第一型积分; 另一类是向量函数的积分, 这类积分称为第二型积分。在以下各节中, 我们将对它们作详细介绍.
% 
% \section{第一型曲线积分}
% \label{sec:ma-16-1}
% 
% \subsection{第一型曲线积分的定义}
% \label{subsec:ma-16-1-1}
% 
% 在定积分的应用中, 我们学习过曲线弧长的定义。现在我们简要回顾一下。设 $\Gamma$ 是平面内或空间 $\mathbb R^3$ 中的一条连续的简单曲线, 以 $A,B$ 为其端点 (当 $A=B$ 时, 则 $\Gamma$ 是一条约当曲线)。在 $\Gamma$ 上从 $A$ 到 $B$ 依次取点 $A=A_0,A_1,\cdots,A_K=B$, 它们构成了 $\Gamma$ 的一个分割 $T$。记 $\overline{A_{k-1}A_k}$ 为连接 $A_{k-1},A_k$ 的线段, 其长为 $s_k(k=1,2,\cdots,K)$, 则当
% \[
%   \sup_T\sum_{k=1}^K s_k<+\infty
% \]
% 时, 其中上确界是对 $\Gamma$ 所有的分割 $T$ 取的, 我们称该曲线是可求长的, 并称该上确界为 $\Gamma$ 的弧长.
% 
% 现设光滑曲线 $\Gamma$ 由参数方程
% \[
%   \begin{cases}
%     x=x(t),\\
%     y=y(t),\\
%     z=z(t)
%   \end{cases}
%   \quad(\alpha\leqslant t\leqslant \beta)
% \]
% 给出, 则其弧长 $L$ 可由下述公式计算:
% \[
%   L=\int_\alpha^\beta \sqrt{(x'(t))^2+(y'(t))^2+(z'(t))^2}\mathrm dt.
% \]
% 
% 为了引入第一型曲线积分, 我们先来研究以下物质曲线的质量问题。设 $\Gamma$ 是一条物质曲线, 它可被看成落在 $\mathbb R^3$ 中的一条曲线, 并在它上面定义了一个密度函数。现在设 $\Gamma$ 以 $A$ 为起点, 以 $B$ 为终点, 在 $\Gamma$ 上任一点 $(x,y,z)$ 处该物质曲线的线密度为 $\rho(x,y,z)$, 我们来求 $\Gamma$ 的质量。当线密度函数为常数函数, 即 $\rho(x,y,z)\equiv\rho_0$ 时, 我们可以求出该物质曲线的质量为 $\rho_0L$, 其中 $L$ 为该曲线的弧长。当 $\rho(x,y,z)$ 在 $\Gamma$ 上不是常数函数时, 依照定积分的思想, 我们可以在 $\Gamma$ 上依次取点 $A=A_0,A_1,\cdots,A_K=B$, 并称其为 $\Gamma$ 的一个分割。此分割将 $\Gamma$ 分成 $K$ 个小弧段。在以 $A_{k-1},A_k$ 为端点的弧段 $\widehat{A_{k-1}A_k}$ 上任取一点 $(\xi_k,\eta_k,\zeta_k)(k=1,2,\cdots,K)$, 作和式 $\sum\limits_{k=1}^K \rho(\xi_k,\eta_k,\zeta_k)\Delta s_k$, 其中 $\Delta s_k(k=1,2,\cdots,K)$ 是 $\widehat{A_{k-1}A_k}$ 的弧长。如同求曲边梯形的面积, 当分割越来越细时, 若该和式的极限趋向于常数 $I$, 则我们认为 $I$ 即为该物质曲线的质量。由此我们引进下面的定义.
% 
% \begin{definition}\label{definition:ma-16-1-1}
% 设 $\Gamma\subset\mathbb R^3$ 是可求长的简单曲线, 它以 $A$ 为起点, $B$ 为终点, $f(x,y,z)$ 是定义在 $\Gamma$ 上的函数。对于 $\Gamma$ 的任一分割 $T$: $A=A_0,A_1,\cdots,A_K=B$, 记 $\Delta s_k(k=1,2,\cdots,K)$ 为弧段 $\widehat{A_{k-1}A_k}$ 的弧长 (即 $\Gamma$ 介于 $A_{k-1},A_k$ 之间部分的长度) 及 $\lambda(T)=\max\limits_{1\leqslant k\leqslant K}\Delta s_k$。在 $\widehat{A_{k-1}A_k}$ 上任取一点 $(\xi_k,\eta_k,\zeta_k)(k=1,2,\cdots,K)$, 作和式
% \[
%   \sum_{k=1}^K f(\xi_k,\eta_k,\zeta_k)\Delta s_k.
% \]
% 若存在 $I\in\mathbb R$, 使得对于 $\forall\varepsilon>0$, $\exists\delta>0$, 对 $\Gamma$ 的任意分割 $T$ 及每个小弧段上任意选取的 $(\xi_k,\eta_k,\zeta_k)$, 当 $\lambda(T)<\delta$ 时, 有
% \[
%   \left|\sum_{k=1}^K f(\xi_k,\eta_k,\zeta_k)\Delta s_k-I\right|<\varepsilon,
% \]
% 即
% \[
%   \lim_{\lambda(T)\to0}\sum_{k=1}^K f(\xi_k,\eta_k,\zeta_k)\Delta s_k=I,
% \]
% 则称 $f(x,y,z)$ 沿 $\Gamma$ 的第一型曲线积分存在, 并称 $I$ 为 $f(x,y,z)$ 在 $\Gamma$ 上的第一型曲线积分, 记为
% \[
%   I=\int_\Gamma f(x,y,z)\mathrm ds
%   \quad\text{或}\quad
%   I=\int_{\widehat{AB}} f(x,y,z)\mathrm ds,
% \]
% 其中 $f(x,y,z)$ 称为被积函数, $\Gamma$ 称为积分曲线.
% \end{definition}
% 
% 从定义~\ref{definition:ma-16-1-1} 可以看出, 第一型曲线积分只与积分曲线 $\Gamma$ 的弧长及被积函数 $f(x,y,z)$ 的取值有关, 而与 $\Gamma$ 的方向无关。设 $\Gamma=\widehat{AB}$, 即 $\Gamma$ 是以 $A$ 为起点, $B$ 为终点的一条曲线。记 $\Gamma^-=\widehat{BA}$ 是 $\Gamma$ 的反向曲线, 即将 $\Gamma$ 的起点与终点对调后所得到的曲线, 则当 $f(x,y,z)$ 在 $\Gamma$ 上的第一型曲线积分存在时, 有
% \[
%   \int_\Gamma f(x,y,z)\mathrm ds=\int_{\Gamma^-} f(x,y,z)\mathrm ds.
% \]
% 
% 从定义~\ref{definition:ma-16-1-1} 还可以推知第一型曲线积分具有下述性质:
% 
% (1) 若对于 $\forall(x,y,z)\in\Gamma$, 有 $f(x,y,z)\equiv1$, 则 $\int_\Gamma f(x,y,z)\mathrm ds=\int_\Gamma \mathrm ds$ 即为 $\Gamma$ 的弧长.
% 
% (2) 若 $\Gamma$ 落在 $Oxy$ 平面上, 且对于 $\forall(x,y)\in\Gamma$ 有 $f(x,y)>0$, 则 $\int_\Gamma f(x,y)\mathrm ds$ 表示空间柱面 $S=\{(x,y,z):(x,y)\in\Gamma,0\leqslant z\leqslant f(x,y)\}$ 的面积 (见图 16.1.1, 关于曲面面积的定义将在 第~\ref{sec:ma-16-3}~节 中讨论).
% 
% (3) 第一型曲线积分具有关于被积函数的线性性, 即设 $\Gamma$ 是可求长的简单曲线, 函数 $f(x,y,z)$ 与 $g(x,y,z)$ 在 $\Gamma$ 上的第一型曲线积分存在, $\alpha,\beta\in\mathbb R$, 则 $\alpha f(x,y,z)+\beta g(x,y,z)$ 在 $\Gamma$ 上的第一型曲线积分存在, 并且
% \[
%   \int_\Gamma [\alpha f(x,y,z)+\beta g(x,y,z)]\mathrm ds
%   =\alpha\int_\Gamma f(x,y,z)\mathrm ds
%   +\beta\int_\Gamma g(x,y,z)\mathrm ds.
% \]
% 
% (4) 第一型曲线积分具有关于积分曲线的可加性, 即设 $\Gamma_1=\widehat{AB}$, $\Gamma_2=\widehat{BC}$ 均为可求长的简单曲线, 且 $\Gamma=\Gamma_1+\Gamma_2=\widehat{AC}$ 也为可求长的简单曲线, 则函数 $f(x,y,z)$ 在 $\Gamma$ 上的第一型曲线积分存在的充分必要条件是 $f(x,y,z)$ 在 $\Gamma_1$ 和 $\Gamma_2$ 上的第一型曲线积分均存在, 且当第一型曲线积分存在时, 有
% \[
%   \int_{\Gamma_1} f(x,y,z)\mathrm ds+\int_{\Gamma_2} f(x,y,z)\mathrm ds
%   =\int_\Gamma f(x,y,z)\mathrm ds.
% \]
% 以上性质的证明请读者自己给出.
% 
% \subsection{第一型曲线积分的存在性与计算公式}
% \label{subsec:ma-16-1-2}
% 
% 对于第一型曲线积分的存在性, 我们有下面的结论.
% 
% \begin{theorem}\label{theorem:ma-16-1-1}
% 设 $\Gamma=\widehat{AB}\subset\mathbb R^3$ 是可求长的简单曲线, 函数 $f(x,y,z)$ 在 $\Gamma$ 上连续, 则第一型曲线积分 $\int_\Gamma f(x,y,z)\mathrm ds$ 存在.
% \end{theorem}
% 
% 
% \begin{proof}
% 设 $\Gamma$ 以弧长 $s$ 为参数的参数方程为
% \[
%   \Gamma:
%   \begin{cases}
%     x=x(s),\\
%     y=y(s),\\
%     z=z(s)
%   \end{cases}
%   \quad(0\leqslant s\leqslant L),
% \]
% 其中 $L$ 为 $\Gamma$ 的弧长。对于 $\Gamma=\widehat{AB}$ 的一个分割 $T$: $A=A_0,A_1,\cdots,A_K=B$, 则对应于闭区间 $[0,L]$ 的一个分割 $\Delta$: $0=s_0<s_1<\cdots<s_K=L$。在 $\widehat{A_{k-1}A_k}$ 上任取一点 $(\xi_k,\eta_k,\zeta_k)$, 则必存在 $s'_k\in[s_{k-1},s_k]$, 使得 $(\xi_k,\eta_k,\zeta_k)=(x(s'_k),y(s'_k),z(s'_k))(k=1,2,\cdots,K)$。因此
% \[
%   \sum_{k=1}^K f(\xi_k,\eta_k,\zeta_k)\Delta s_k
%   =\sum_{k=1}^K f(x(s'_k),y(s'_k),z(s'_k))\Delta s_k.
% \]
% 上述等式右端是连续函数 $f(x(s),y(s),z(s))$ 在区间 $[0,L]$ 上的黎曼和。注意到 $\lambda(T)\to0$ 的充分必要条件是 $\lambda(\Delta)\to0$, 因此 $\int_\Gamma f(x,y,z)\mathrm ds$ 存在, 并且有
% \[
%   \int_\Gamma f(x,y,z)\mathrm ds=\int_0^L f(x(s),y(s),z(s))\mathrm ds.
% \qedhere
% \]
% \end{proof}
% 
% 
% 在假定 $\Gamma\subset\mathbb R^3$ 是光滑曲线时, 我们有以下关于第一型曲线积分的计算公式.
% 
% \begin{theorem}\label{theorem:ma-16-1-2}
% 设 $\Gamma$ 是光滑曲线, 其参数方程为
% \[
%   \begin{cases}
%     x=x(t),\\
%     y=y(t),\\
%     z=z(t)
%   \end{cases}
%   \quad(\alpha\leqslant t\leqslant \beta),
% \]
% 再设函数 $f(x,y,z)$ 在 $\Gamma$ 上连续, 则有
% \[
%   \int_\Gamma f(x,y,z)\mathrm ds
%   =\int_\alpha^\beta f(x(t),y(t),z(t))
%   \sqrt{[x'(t)]^2+[y'(t)]^2+[z'(t)]^2}\mathrm dt.
% \]
% \end{theorem}
% 
% 
% \begin{proof}
% 由于 $\Gamma$ 是光滑曲线, 设其弧长为 $L$。作 $[\alpha,\beta]\to[0,L]$ 的变换如下: 对于 $\forall t\in[\alpha,\beta]$, 记 $s(t)$ 为 $\Gamma$ 介于 $(x(\alpha),y(\alpha),z(\alpha))$ 与 $(x(t),y(t),z(t))$ 之间部分的弧长, 即
% \[
%   s=s(t)=\int_\alpha^t \sqrt{[x'(u)]^2+[y'(u)]^2+[z'(u)]^2}\mathrm du.
% \]
% 注意到
% \[
%   \frac{\mathrm ds}{\mathrm dt}
%   =\sqrt{[x'(t)]^2+[y'(t)]^2+[z'(t)]^2}>0,
% \]
% 记 $t=t(s)$ 为 $s(t)$ 的反函数, 因此由定积分换元公式得
% \[
% \begin{aligned}
%   \int_\Gamma f(x,y,z)\mathrm ds
%   &=\int_0^L f(x(t(s)),y(t(s)),z(t(s)))\mathrm ds\\
%   &=\int_\alpha^\beta f(x(t),y(t),z(t))
%     \sqrt{[x'(t)]^2+[y'(t)]^2+[z'(t)]^2}\mathrm dt.
% \end{aligned}
% \qedhere
% \]
% \end{proof}
% 
% 
% \begin{example}\label{example:ma-16-1-1}
% 计算第一型曲线积分 $\int_\Gamma \sqrt{1-x^2-y^2}\mathrm ds$, 其中 $\Gamma$ 是曲线 $x^2+y^2=x$.
% \end{example}
% 
% 
% \begin{solve}[{解法 1}]
% 取曲线 $\Gamma$ 的参数方程
% \[
%   \begin{cases}
%     x=\dfrac12+\dfrac12\cos t,\\
%     y=\dfrac12\sin t
%   \end{cases}
%   \quad(0\leqslant t\leqslant 2\pi),
% \]
% 则
% \[
%   \mathrm ds=\sqrt{\frac14\sin^2t+\frac14\cos^2t}\mathrm dt
%   =\frac12\mathrm dt.
% \]
% 于是
% \[
% \begin{aligned}
%   \int_\Gamma \sqrt{1-x^2-y^2}\mathrm ds
%   &=\int_\Gamma \sqrt{1-x}\mathrm ds
%     =\frac12\int_0^{2\pi}\sqrt{\frac{1-\cos t}{2}}\mathrm dt\\
%   &=\frac12\int_0^{2\pi}\left|\sin\frac t2\right|\mathrm dt
%     =\int_0^\pi\sin\frac t2\mathrm dt=2.
% \end{aligned}
% \qedhere
% \]
% \end{solve}
% 
% 
% \begin{solve}[{解法 2}]
% 由于被积函数及积分曲线的对称性, 我们只要考虑积分曲线落在第一象限的部分即可。此时曲线 $\Gamma$ 可写成如下的参数方程
% \[
%   \begin{cases}
%     x=x,\\
%     y=\sqrt{x-x^2}
%   \end{cases}
%   \quad(0\leqslant x\leqslant1),
% \]
% 从而 $\mathrm ds=\dfrac{\mathrm dx}{2\sqrt{x-x^2}}$, 因此
% \[
%   \int_\Gamma \sqrt{1-x^2-y^2}\mathrm ds
%   =2\int_0^1\sqrt{1-x}\frac{\mathrm dx}{2\sqrt{x-x^2}}
%   =\int_0^1\frac{\mathrm dx}{\sqrt x}=2.
% \qedhere
% \]
% \end{solve}
% 
% 
% \begin{remark}\label{remark:ma-16-source-9}
% 在上例中, 解法 2 中的计算更简单些, 但最后的积分是一个瑕积分。另外, 由第一型曲线积分的性质, 上例中的积分值是柱面 $x^2+y^2=x$ 被单位球面所截得的在 $Oxy$ 平面上方那一块的面积.
% \end{remark}
% 
% 
% \begin{example}\label{example:ma-16-1-2}
% 计算第一型曲线积分 $\int_\Gamma (x+y)^2\mathrm ds$, 其中 $\Gamma$ 是单位圆盘上半部分的边界, 即 $\Gamma=\Gamma_1\cup\Gamma_2$, 其中
% \[
%   \Gamma_1=\{(x,y):x^2+y^2=1,y\geqslant0\},\quad
%   \Gamma_2=\{(x,y):-1\leqslant x\leqslant1,y=0\}.
% \]
% \end{example}
% 
% 
% \begin{solve}
% 由第一型曲线积分的性质, 我们有
% \[
%   \int_\Gamma (x+y)^2\mathrm ds
%   =\int_{\Gamma_1}(x+y)^2\mathrm ds+\int_{\Gamma_2}(x+y)^2\mathrm ds.
% \]
% 由于
% \[
%   \int_{\Gamma_1}(x+y)^2\mathrm ds
%   =\int_{\Gamma_1}(x^2+y^2)\mathrm ds+2\int_{\Gamma_1}xy\mathrm ds
%   =\int_{\Gamma_1}\mathrm ds+0=\pi,
% \]
% 其中 $\int_{\Gamma_1}xy\mathrm ds=0$ 由对称性得到, 另外我们有
% \[
%   \int_{\Gamma_2}(x+y)^2\mathrm ds
%   =\int_{-1}^1 x^2\mathrm dx=\frac23,
% \]
% 所以
% \[
%   \int_\Gamma (x+y)^2\mathrm ds=\frac23+\pi.
% \qedhere
% \]
% \end{solve}
% 
% 
% \section{第二型曲线积分}
% \label{sec:ma-16-2}
% 
% \subsection{第二型曲线积分的定义}
% \label{subsec:ma-16-2-1}
% 
% 第二型曲线积分的物理模型是变力做功。设空间有一个力场, 如地球对空间中任一点都存在引力, 该引力在空间的分布就可称为地球的引力场。从数学的观点来看一个力场即是在 $\mathbb R^3$ 中定义了一个向量函数
% \[
%   \boldsymbol F(x,y,z)=(P(x,y,z),Q(x,y,z),R(x,y,z)).
% \]
% 假定在它的作用下, 一个质点沿空间曲线 $\Gamma$ 从点 $A$ 移到了点 $B$。现在我们来求该力场对质点所做的功.
% 
% 首先注意到, 若 $\boldsymbol F$ 是常力, 即它的大小与方向均是常数, 并且在它的作用下, 质点沿连接点 $A$ 和 $B$ 的直线段 $\overline{AB}$ 从 $A$ 移到 $B$, 则它所做的功为
% \[
%   W=\boldsymbol F\cdot\overrightarrow{AB}.
% \]
% 为了求解上述的一般问题, 我们用熟知的分割、求和、取极限的积分技巧。记 $\Gamma=\widehat{AB}$, 沿 $\Gamma$ 依次取分点
% \[
%   A=A_0,A_1,\cdots,A_K=B.
% \]
% 它们构成 $\Gamma$ 的一个分割 $T$, 记 $\lambda(T)=\max\limits_{1\leqslant k\leqslant K}\{\operatorname{diam}(\widehat{A_{k-1}A_k})\}$。在每一段小弧 $\widehat{A_{k-1}A_k}(k=1,2,\cdots,K)$ 上, 以线段 $\overline{A_{k-1}A_k}$ 代替 $\widehat{A_{k-1}A_k}$, 并且将 $\boldsymbol F(x,y,z)$ 在 $\widehat{A_{k-1}A_k}$ 上取常力 $\boldsymbol F(\xi_k,\eta_k,\zeta_k)((\xi_k,\eta_k,\zeta_k)\in\widehat{A_{k-1}A_k})$, 则在变力 $\boldsymbol F$ 作用下质点沿 $\Gamma$ 从 $A_{k-1}$ 移动到 $A_k$ 所做的功
% \[
%   \Delta W_k\approx \boldsymbol F(\xi_k,\eta_k,\zeta_k)\cdot\overrightarrow{A_{k-1}A_k}.
% \]
% 因此
% \[
%   W=\sum_{k=1}^K\Delta W_k
%   \approx \sum_{k=1}^K \boldsymbol F(\xi_k,\eta_k,\zeta_k)\cdot\overrightarrow{A_{k-1}A_k}.
% \]
% 当 $\lambda(T)\to0$ 时, 若上式右端和式的极限存在, 则
% \[
%   W=\lim_{\lambda(T)\to0}\sum_{k=1}^K
%   \boldsymbol F(\xi_k,\eta_k,\zeta_k)\cdot\overrightarrow{A_{k-1}A_k}.
% \]
% 记 $A_k=(x_k,y_k,z_k)(k=1,2,\cdots,K)$, 则
% \[
%   \overrightarrow{A_{k-1}A_k}
%   =(x_k-x_{k-1},y_k-y_{k-1},z_k-z_{k-1})
%   \triangleq(\Delta x_k,\Delta y_k,\Delta z_k).
% \]
% 于是我们有
% \[
% \begin{aligned}
%   &\sum_{k=1}^K \boldsymbol F(\xi_k,\eta_k,\zeta_k)\cdot\overrightarrow{A_{k-1}A_k}\\
%   &\quad=\sum_{k=1}^K\bigl(P(\xi_k,\eta_k,\zeta_k)\Delta x_k
%     +Q(\xi_k,\eta_k,\zeta_k)\Delta y_k
%     +R(\xi_k,\eta_k,\zeta_k)\Delta z_k\bigr).
% \end{aligned}
% \]
% 从上述变力做功的问题中, 我们得到一个新的和式极限。由此和式极限我们可以引进一类新的积分。另外要注意, 我们记空间一条曲线 $\Gamma=\widehat{AB}$ 是为了强调 $\Gamma$ 是以 $A$ 为起点, $B$ 为终点的。当 $A=B$ 时, 我们则必须在 $\Gamma$ 上取定一个方向.
% 
% \begin{definition}\label{definition:ma-16-2-1}
% 设 $\Gamma=\widehat{AB}$ 为空间连续曲线, 向量函数
% \[
%   \boldsymbol F(x,y,z)=(P(x,y,z),Q(x,y,z),R(x,y,z))
% \]
% 在 $\widehat{AB}$ 上有定义, 记在 $\widehat{AB}$ 上依次取分点
% \[
%   A=A_0,A_1,\cdots,A_K=B
% \]
% 组成的分割为 $T$，并记
% \[\overrightarrow{A_{k-1}A_k}=(\Delta x_k,\Delta y_k,\Delta z_k),\quad k=1,2,\cdots,K.\]
% 分割的模记为 $\lambda(T)=\max\limits_{1\leqslant k\leqslant K}\{\operatorname{diam}(\widehat{A_{k-1}A_k})\}$。在 $\widehat{A_{k-1}A_k}$ 上任取点 $(\xi_k,\eta_k,\zeta_k)(k=1,2,\cdots,K)$。若极限
% \[
%   \lim_{\lambda(T)\to0}\sum_{k=1}^K
%   [P(\xi_k,\eta_k,\zeta_k)\Delta x_k
%   +Q(\xi_k,\eta_k,\zeta_k)\Delta y_k
%   +R(\xi_k,\eta_k,\zeta_k)\Delta z_k]
% \]
% 存在并且等于 $I$, 其中常数 $I$ 不依赖 $\widehat{AB}$ 的分割 $T$ 以及 $(\xi_k,\eta_k,\zeta_k)$ 的选取, 则称该极限值 $I$ 为 $\boldsymbol F(x,y,z)$ 在 $\Gamma$ 上的第二型曲线积分, 记为
% \[
%   I=\int_\Gamma P(x,y,z)\mathrm dx+Q(x,y,z)\mathrm dy+R(x,y,z)\mathrm dz,
% \]
% 或者
% \[
%   I=\int_{\widehat{AB}}P\mathrm dx+Q\mathrm dy+R\mathrm dz,
% \]
% 其中 $\boldsymbol F(x,y,z)=(P,Q,R)$ 称为被积函数, $\Gamma$ 称为积分曲线.
% \end{definition}
% 
% 从上述定义可以看出, 第二型曲线积分是考虑一个向量函数在曲线上沿着一个方向的积分, 因此它与曲线的方向有关。显然, 第一型曲线积分是考虑一个多元函数在曲线上的积分。当然, 仅从纯数学的观点, 第二型曲线积分也可仅仅定义一个多元函数 $P(x,y,z)$ 的下述积分 $\int_{\widehat{AB}}P\mathrm dx$, $\int_{\widehat{AB}}P\mathrm dy$, $\int_{\widehat{AB}}P\mathrm dz$。但是这些积分不能反映出第二型曲线积分的物理背景.
% 
% 从定义我们可以看出第二型曲线积分具有下述性质:
% 
% (1) 若 $\int_{\widehat{AB}}P\mathrm dx+Q\mathrm dy+R\mathrm dz$ 存在, 则 $\int_{\widehat{BA}}P\mathrm dx+Q\mathrm dy+R\mathrm dz$ 存在, 并且
% \[
%   \int_{\widehat{BA}}P\mathrm dx+Q\mathrm dy+R\mathrm dz
%   =-\int_{\widehat{AB}}P\mathrm dx+Q\mathrm dy+R\mathrm dz.
% \]
% 特别地, 当 $\Gamma$ 是一条简单闭曲线时, 我们有
% \[
%   \int_{\Gamma^-}P\mathrm dx+Q\mathrm dy+R\mathrm dz
%   =-\int_\Gamma P\mathrm dx+Q\mathrm dy+R\mathrm dz.
% \]
% 
% (2) 设向量函数 $\boldsymbol f=(P_1,Q_1,R_1)$ 和 $\boldsymbol g=(P_2,Q_2,R_2)$ 在空间曲线 $\widehat{AB}$ 上的第二型曲线积分都存在, $\alpha,\beta\in\mathbb R$, 则 $\alpha\boldsymbol f+\beta\boldsymbol g$ 在 $\widehat{AB}$ 上的第二型曲线积分存在, 并且
% \[
% \begin{aligned}
%   &\int_{\widehat{AB}}(\alpha P_1\mathrm dx+\alpha Q_1\mathrm dy+\alpha R_1\mathrm dz)
%   +(\beta P_2\mathrm dx+\beta Q_2\mathrm dy+\beta R_2\mathrm dz)\\
%   &\quad=\alpha\int_{\widehat{AB}}P_1\mathrm dx+Q_1\mathrm dy+R_1\mathrm dz
%   +\beta\int_{\widehat{AB}}P_2\mathrm dx+Q_2\mathrm dy+R_2\mathrm dz.
% \end{aligned}
% \]
% 
% (3) 设 $\widehat{AB}=\widehat{AC}\cup\widehat{CB}$, 即 $\widehat{AB}$ 是由两条连续曲线 $\widehat{AC},\widehat{CB}$ 组成, 并且 $\widehat{AC}$ 与 $\widehat{CB}$ 仅在端点 $C$ 处相交, 则向量函数 $\boldsymbol F=(P,Q,R)$ 在 $\widehat{AB}$ 上的第二型曲线积分存在的充分必要条件是它在 $\widehat{AC}$ 及 $\widehat{CB}$ 上的第二型曲线积分均存在, 并且第二型曲线积分存在时, 成立等式
% \[
%   \int_{\widehat{AB}}P\mathrm dx+Q\mathrm dy+R\mathrm dz
%   =\int_{\widehat{AC}}P\mathrm dx+Q\mathrm dy+R\mathrm dz
%   +\int_{\widehat{CB}}P\mathrm dx+Q\mathrm dy+R\mathrm dz.
% \]
% 
% \subsection{第二型曲线积分的存在性与计算公式}
% \label{subsec:ma-16-2-2}
% 
% 在这里, 我们不准备讨论关于非常广泛的积分曲线和被积函数的第二型曲线积分存在性问题, 只讨论积分曲线是光滑曲线, 而被积函数是连续的情形.
% 
% \begin{theorem}\label{theorem:ma-16-2-1}
% 设 $\Gamma=\widehat{AB}$ 是以 $A$ 为起点, $B$ 为终点的光滑曲线, 其参数方程为
% \[
%   \begin{cases}
%     x=x(t),\\
%     y=y(t),\\
%     z=z(t)
%   \end{cases}
%   \quad(\alpha\leqslant t\leqslant \beta)
% \]
% 且当 $t$ 从 $\alpha$ 连续变为 $\beta$ 时, 对应曲线上的点从 $A$ 连续变为 $B$, 再假定函数 $P(x,y,z)$、$Q(x,y,z)$ 和 $R(x,y,z)$ 在 $\Gamma$ 上连续, 则向量函数 $\boldsymbol F=(P,Q,R)$ 在 $\widehat{AB}$ 上的第二型曲线积分存在, 并且
% \[
% \begin{aligned}
%   \int_{\widehat{AB}}P\mathrm dx+Q\mathrm dy+R\mathrm dz
%   &=\int_\alpha^\beta P(x(t),y(t),z(t))x'(t)\mathrm dt
%     +\int_\alpha^\beta Q(x(t),y(t),z(t))y'(t)\mathrm dt\\
%   &\quad+\int_\alpha^\beta R(x(t),y(t),z(t))z'(t)\mathrm dt.
% \end{aligned}
% \]
% \end{theorem}
% 
% 
% \begin{proof}
% 我们用微元法证之。由光滑曲线的定义知 $x'(t),y'(t),z'(t)$ 均是 $t$ 的连续函数, 且对于 $\forall t\in(\alpha,\beta)$, 有
% \[
%   x'^2(t)+y'^2(t)+z'^2(t)\ne0.
% \]
% 对于 $\forall[t,t+\mathrm dt]\subset[\alpha,\beta]$, 记 $A_t=(x(t),y(t),z(t))\in\Gamma$, $A_{t+\mathrm dt}=(x(t+\mathrm dt),y(t+\mathrm dt),z(t+\mathrm dt))\in\Gamma$, 由于
% \[
% \begin{aligned}
%   \overrightarrow{A_tA_{t+\mathrm dt}}
%   &\approx (x'(t),y'(t),z'(t))\mathrm dt\\
%   &=\frac{(x'(t),y'(t),z'(t))}
%     {\sqrt{x'^2(t)+y'^2(t)+z'^2(t)}}
%     \sqrt{x'^2(t)+y'^2(t)+z'^2(t)}\mathrm dt\\
%   &=\frac{(x'(t),y'(t),z'(t))}
%     {\sqrt{x'^2(t)+y'^2(t)+z'^2(t)}}\mathrm ds\\
%   &=(\cos\alpha,\cos\beta,\cos\gamma)\mathrm ds.
% \end{aligned}
% \]
% 其中 $(\cos\alpha,\cos\beta,\cos\gamma)$ 是曲线在点 $(x(t),y(t),z(t))$ 处的单位切向量。因此有
% \[
% \begin{aligned}
%   P\mathrm dx+Q\mathrm dy+R\mathrm dz
%   &=(P,Q,R)(\mathrm dx,\mathrm dy,\mathrm dz)\\
%   &=(P,Q,R)(\cos\alpha,\cos\beta,\cos\gamma)\mathrm ds\\
%   &=(P\cos\alpha+Q\cos\beta+R\cos\gamma)\mathrm ds.
% \end{aligned}
% \]
% 由假设, $\Gamma$ 是光滑曲线, 从而 $\cos\alpha,\cos\beta,\cos\gamma$ 是 $t$ 的连续函数。由第一型曲线积分的存在性定理知 $\int_{\widehat{AB}}(P\cos\alpha+Q\cos\beta+R\cos\gamma)\mathrm ds$ 存在, 从而 $\int_{\widehat{AB}}P\mathrm dx+Q\mathrm dy+R\mathrm dz$ 存在, 并且有公式
% \[
% \begin{aligned}
%   \int_{\widehat{AB}}P\mathrm dx+Q\mathrm dy+R\mathrm dz
%   &=\int_\alpha^\beta P(x(t),y(t),z(t))x'(t)\mathrm dt
%     +\int_\alpha^\beta Q(x(t),y(t),z(t))y'(t)\mathrm dt\\
%   &\quad+\int_\alpha^\beta R(x(t),y(t),z(t))z'(t)\mathrm dt.
% \end{aligned}
% \qedhere
% \]
% \end{proof}
% 
% 
% \begin{remark}\label{remark:ma-16-source-15}
% 在上面定理的证明过程中, 我们清楚地知道, 第二型曲线积分与第一型曲线积分有如下的关系:
% \[
%   \int_{\widehat{AB}}P\mathrm dx+Q\mathrm dy+R\mathrm dz
%   =\int_{\widehat{AB}}(P,Q,R)\cdot(\cos\alpha,\cos\beta,\cos\gamma)\mathrm ds,
% \]
% 其中 $(\cos\alpha,\cos\beta,\cos\gamma)$ 为 $\widehat{AB}$ 在 $(x,y,z)$ 处的单位切向量.
% \end{remark}
% 
% 到现在为止, 我们讨论了平面及空间的曲线积分问题。值得指出的是, 在 $\mathbb R^n(n>3)$ 中也可以讨论相应的曲线积分问题。事实上, 我们已经有了 $\mathbb R^n$ 中曲线的定义, 因此可以用完全平行于 $\mathbb R^3$ 中的理论来讨论可求长曲线、光滑曲线等问题, 从而可以定义两类不同的积分。由于这些理论与我们已经研究的 $\mathbb R^3$ 中的相应理论有很大的相似性, 我们在这里就不深入讨论了.
% 
% \begin{example}\label{example:ma-16-2-1}
% 计算第二型曲线积分 $\int_{\widehat{AB}}(-y)\mathrm dx+x\mathrm dy$, 其中 $\widehat{AB}$ 为单位圆周 $x^2+y^2=1$ 的上半部分, 方向为从点 $A(1,0)$ 到点 $B(-1,0)$.
% \end{example}
% 
% 
% \begin{solve}[{解法 1}]
% 令
% \[
%   \begin{cases}
%     x=\cos t,\\
%     y=\sin t
%   \end{cases}
%   \quad(0\leqslant t\leqslant\pi),
% \]
% 则 $t$ 从 $0$ 连续变化到 $\pi$ 时, 曲线上的点从 $A$ 连续变化到 $B$。因此
% \[
% \begin{aligned}
%   \int_{\widehat{AB}}(-y)\mathrm dx+x\mathrm dy
%   &=\int_0^\pi[-\sin t\cdot(-\sin t)+\cos t\cdot\cos t]\mathrm dt\\
%   &=\int_0^\pi\mathrm dt=\pi.
% \end{aligned}
% \qedhere
% \]
% \end{solve}
% 
% 
% \begin{solve}[{解法 2}]
% 由 $x^2+y^2=1$ 知 $x\mathrm dx+y\mathrm dy=0$, 得 $\mathrm dy=-\dfrac{x}{y}\mathrm dx$, 从而有
% \[
% \begin{aligned}
%   \int_{\widehat{AB}}(-y)\mathrm dx+x\mathrm dy
%   &=\int_1^{-1}(-y)\mathrm dx+\int_1^{-1}x\left(-\frac{x}{y}\right)\mathrm dx\\
%   &=\int_{-1}^1\left(\frac{x^2+y^2}{y}\right)\mathrm dx
%     =\int_{-1}^1\frac{\mathrm dx}{\sqrt{1-x^2}}
%     =\arcsin x\bigg|_{-1}^1=\pi.
% \end{aligned}
% \qedhere
% \]
% \end{solve}
% 
% 
% \begin{example}\label{example:ma-16-2-2}
% 计算第二型曲线积分 $\int_\Gamma (x^2-y^2)\mathrm dx-2xy\mathrm dy$, 其中 $\Gamma$ 是从点 $(0,0)$ 沿曲线 $y=x^\alpha(\alpha>0)$ 到点 $(1,1)$ 的部分.
% \end{example}
% 
% 
% \begin{solve}
% 将曲线
% \[
%   \begin{cases}
%     x=x,\\
%     y=x^\alpha
%   \end{cases}
%   \quad(x\in[0,1])
% \]
% 代入得
% \[
% \begin{aligned}
%   \int_\Gamma (x^2-y^2)\mathrm dx-2xy\mathrm dy
%   &=\int_0^1(x^2-x^{2\alpha})\mathrm dx
%     -\int_0^1 2xx^\alpha(\alpha x^{\alpha-1})\mathrm dx\\
%   &=\int_0^1[x^2-(2\alpha+1)x^{2\alpha}]\mathrm dx
%     =-\frac23.
% \end{aligned}
% \]
% 此题说明, 该第二型曲线积分与 $\alpha$ 的选取无关, 而只与这些曲线的端点有关.
% \end{solve}
% 
% 
% \begin{example}\label{example:ma-16-2-3}
% 计算第二型曲线积分 $\int_\Gamma x\mathrm dx+y\mathrm dy+z\mathrm dz$, 其中 $\Gamma$ 为球面 $x^2+y^2+z^2=1$ 与平面 $x+y+z=0$ 的交线, 从 $z$ 轴看去取逆时针方向.
% \end{example}
% 
% 
% \begin{solve}[{解法 1}]
% 作正交变换将 $\Gamma$ 变到 $O\xi\eta$ 平面:
% \[
%   \begin{cases}
%     \xi=\dfrac1{\sqrt2}(x-y),\\
%     \eta=\dfrac1{\sqrt6}(x+y-2z),\\
%     \zeta=\dfrac1{\sqrt3}(x+y+z).
%   \end{cases}
% \]
% 然后利用 $\xi=\cos t,\eta=\sin t$ 代入上述 $\Gamma$ 的方程, 即得 $\Gamma$ 的参数方程
% \[
%   \begin{cases}
%     x=\dfrac1{\sqrt2}\cos t+\dfrac1{\sqrt6}\sin t,\\
%     y=-\dfrac1{\sqrt2}\cos t+\dfrac1{\sqrt6}\sin t,\\
%     z=-\dfrac2{\sqrt6}\sin t,
%   \end{cases}
% \]
% 从而
% \[
% \begin{aligned}
%   \int_\Gamma x\mathrm dx+y\mathrm dy+z\mathrm dz
%   &=\int_0^{2\pi}\bigg[
%     \left(\frac1{\sqrt2}\cos t+\frac1{\sqrt6}\sin t\right)
%     \left(-\frac1{\sqrt2}\sin t+\frac1{\sqrt6}\cos t\right)\\
%   &\quad+
%     \left(-\frac1{\sqrt2}\cos t+\frac1{\sqrt6}\sin t\right)
%     \left(\frac1{\sqrt2}\sin t+\frac1{\sqrt6}\cos t\right)\\
%   &\quad+
%     \left(-\frac2{\sqrt6}\sin t\right)
%     \left(-\frac2{\sqrt6}\cos t\right)
%   \bigg]\mathrm dt\\
%   &=\int_0^{2\pi}0\mathrm dt=0.
% \end{aligned}
% \qedhere
% \]
% \end{solve}
% 
% 
% \begin{solve}[{解法 2}]
% 由于球面 $x^2+y^2+z^2=1$ 上每个点 $(x,y,z)$ 在 $\mathbb R^3$ 中表示的向量即是球面在该点的单位法向量 $\boldsymbol n$, 而曲线 $\Gamma$ 每个点处的单位切向量 $\boldsymbol v$ 与球面在该点的法向量正交, 因此有
% \[
%   \int_\Gamma x\mathrm dx+y\mathrm dy+z\mathrm dz
%   =\int_\Gamma \boldsymbol n\cdot\boldsymbol v\mathrm ds=0.
% \qedhere
% \]
% \end{solve}
% 
% 
% \section{第一型曲面积分}
% \label{sec:ma-16-3}
% 
% \subsection{曲面的面积}
% \label{subsec:ma-16-3-1}
% 
% 在介绍第一型曲面积分之前, 我们先来讨论一下曲面的面积问题。如何利用积分的思想来求曲面的面积呢? 受求曲线弧长的启发, 我们能否用曲面内接多边形的面积来逼近曲面的面积? 对此施瓦茨曾经有一个著名的例子, 他考虑了由圆柱面的内接全等等腰三角形构成的多面体面积的逼近: 设圆柱面 $S$ 为 $\{(x,y,z):x^2+y^2=1,0\leqslant z\leqslant1\}$。将 $S$ 的高进行了 $K$ 等分, 从而在 $S$ 上得到 $K+1$ 个平行于 $Oxy$ 平面的圆周 $\Gamma_k(k=1,2,\cdots,K+1)$。对每个圆周 $\Gamma_k(k=1,2,\cdots,K+1)$ 再 $L$ 等分, 使得 $\Gamma_{k+1}$ 上的分点在 $Oxy$ 平面的投影落在 $\Gamma_k$ 上的分点在 $Oxy$ 平面的投影的中间, 然后将每个 $\Gamma_k$ 上相邻两个分点连接, 并且它们分别与 $\Gamma_{k+1}$ 上的位于它上方中间的点连接, 这样我们就得到了一个等腰三角形。所有的这些三角形组成了 $S$ 的一个内接多边形。这个多边形的面积用初等数学的知识即可求出。容易证明, 适当选取 $K,L$ 的比例, 如 $K=L^3$, 则可以使得多边形的面积趋于 $\infty$。在上述逼近中, 若取 $K=L$, 则容易推出多边形的面积趋于 $2\pi$。因此用内接多边形的面积来逼近曲面面积的方法是行不通的.
% 
% 仔细分析一下施瓦茨的例子, 当 $K$ 与 $L$ 的比很大时, 即三角形的高与底边长的比趋于零时, 每个小三角形所在平面几乎与柱面的切平面垂直, 因此产生很多由于弯曲而增加的面积。如果该三角形所在平面越来越“平行”于切平面, 则多边形面积逼近将会成功。注意到对于圆柱面, 我们若将其摊平到平面, 则可以计算出它的面积为
% \[
%   \text{底周长}\times\text{高}=2\pi.
% \]
% 
% 为了求曲面面积, 我们的想法是用每一点处切平面上的部分来构造曲面的外切多边形进行逼近。由于要用到切平面, 我们只考虑光滑曲面.
% 
% 设曲面 $S$ 由参数方程
% \begin{equation}
%   \begin{cases}
%     x=x(u,v),\\
%     y=y(u,v),\\
%     z=z(u,v),
%   \end{cases}
%   \quad (u,v)\in D
%   \label{eq:ma-16-3-1}
% \end{equation}
% 定义, 其中 $D$ 是可求面积的有界闭区域, 函数 $x(u,v),y(u,v),z(u,v)$ 均在 $D$ 上具有连续偏导数。作为 $D$ 到 $\mathbb R^3$ 的映射, \eqref{eq:ma-16-3-1} 是单的且其雅可比矩阵
% \[
%   \begin{pmatrix}
%     x'_u & y'_u & z'_u\\
%     x'_v & y'_v & z'_v
%   \end{pmatrix}
% \]
% 在每一点 $(u,v)$ 处的秩均为 $2$。在本节中, 我们总假定曲面 $S$ 为满足上述条件的光滑曲面.
% 
% 现记 $\boldsymbol r(u,v)=(x(u,v),y(u,v),z(u,v))$ 及
% \begin{equation}
%   A=\frac{\partial(y,z)}{\partial(u,v)},\quad
%   B=\frac{\partial(z,x)}{\partial(u,v)},\quad
%   C=\frac{\partial(x,y)}{\partial(u,v)}.
%   \label{eq:ma-16-3-2}
% \end{equation}
% 则在任一点 $(u,v)$ 处行列式 $A,B,C$ 中至少有一个不为零, 并且它们均是 $(u,v)$ 的连续函数, 因此法向量 $\boldsymbol n=(A,B,C)$ 在曲面 $S$ 上连续变化.
% 
% 下面我们利用微元法来求曲面的面积。给定以 $(u_0,v_0),(u_0+\mathrm du,v_0),(u_0+\mathrm du,v_0+\mathrm dv),(u_0,v_0+\mathrm dv)$ 为顶点的一个小矩形 $D_0$。当 $(u,v)$ 在 $D_0$ 上变化时, 在 $S$ 上可得到一小块曲面, 它可近似看成一个小平行四边形, 其两邻边分别为
% \[
%   \boldsymbol r(u_0+\mathrm du,v_0)-\boldsymbol r(u_0,v_0)
%   \approx \boldsymbol r'_u(u_0,v_0)\mathrm du
% \]
% 与
% \[
%   \boldsymbol r(u_0,v_0+\mathrm dv)-\boldsymbol r(u_0,v_0)
%   \approx \boldsymbol r'_v(u_0,v_0)\mathrm dv.
% \]
% 因此它们张成的平行四边形的面积, 即面积微元为
% \[
%   \mathrm dS(u_0,v_0)=|\boldsymbol r'_u(u_0,v_0)\times\boldsymbol r'_v(u_0,v_0)|\mathrm du\mathrm dv,
% \]
% 从而曲面 $S$ 的面积 (仍用 $S$ 表示其面积) 为
% \[
%   S=\iint_D |\boldsymbol r'_u(u,v)\times\boldsymbol r'_v(u,v)|\mathrm du\mathrm dv.
% \]
% 现在我们进一步计算 $|\boldsymbol r'_u(u,v)\times\boldsymbol r'_v(u,v)|$。由
% \[
%   |\boldsymbol r'_u(u,v)\times\boldsymbol r'_v(u,v)|
%   =|(x'_u,y'_u,z'_u)\times(x'_v,y'_v,z'_v)|
%   =\sqrt{A^2+B^2+C^2},
% \]
% 若记
% \begin{equation}
%   \begin{cases}
%     E=\boldsymbol r'_u\boldsymbol r'_u=x_u'^2+y_u'^2+z_u'^2,\\
%     F=\boldsymbol r'_u\boldsymbol r'_v=x'_ux'_v+y'_uy'_v+z'_uz'_v,\\
%     G=\boldsymbol r'_v\boldsymbol r'_v=x_v'^2+y_v'^2+z_v'^2,
%   \end{cases}
%   \label{eq:ma-16-3-3}
% \end{equation}
% 则可知
% \[
%   EG-F^2=A^2+B^2+C^2.
% \]
% 因此, 我们有下面的定理.
% 
% \begin{theorem}\label{theorem:ma-16-3-1}
% 设光滑曲面 $S$ 由参数方程 \eqref{eq:ma-16-3-1} 给出, 则它的面积为
% \[
%   S=\iint_D \sqrt{A^2+B^2+C^2}\mathrm du\mathrm dv
%   =\iint_D \sqrt{EG-F^2}\mathrm du\mathrm dv,
% \]
% 其中 $A,B,C$ 和 $E,F,G$ 的定义分别见式 \eqref{eq:ma-16-3-2} 和 \eqref{eq:ma-16-3-3}.
% \end{theorem}
% 
% 特别地, 当光滑曲面 $S$ 是由 $z=f(x,y)((x,y)\in D)$ 给出时, 则有
% \[
%   S=\iint_D \sqrt{1+f_x'^2(x,y)+f_y'^2(x,y)}\mathrm dx\mathrm dy.
% \]
% 
% \begin{example}\label{example:ma-16-3-1}
% 求单位圆柱面 $x^2+y^2=1$ 在平面 $z=0$ 与 $z=1$ 之间的面积.
% \end{example}
% 
% 
% \begin{solve}
% 设所求面积为 $S$, 则 $S$ 为曲面 $y=f(z,x)=\sqrt{1-x^2}((z,x)\in D,\text{ 其中 }D=[0,1]\times[-1,1])$ 的面积的两倍。因此有
% \[
% \begin{aligned}
%   S&=2\iint_D\sqrt{1+f_x'^2+f_z'^2}\mathrm dz\mathrm dx\\
%   &=2\int_0^1\mathrm dz\int_{-1}^1\sqrt{1+\left(\frac{-x}{\sqrt{1-x^2}}\right)^2}\mathrm dx\\
%   &=2\int_{-1}^1\frac{\mathrm dx}{\sqrt{1-x^2}}=2\pi.
% \end{aligned}
% \qedhere
% \]
% \end{solve}
% 
% 
% \begin{example}\label{example:ma-16-3-2}
% 计算球面 $x^2+y^2+z^2=1$ 被柱面 $\left(x-\dfrac12\right)^2+y^2=\dfrac14$ 所截在柱面内的部分的面积.
% \end{example}
% 
% 
% \begin{solve}
% 由对称性, 所求面积 $S$ 为所截得曲面在 $Oxy$ 平面上方部分面积的两倍。设 $z=f(x,y)=\sqrt{1-x^2-y^2}((x,y)\in D)$, 其中
% \[
%   D=\left\{(x,y):\left(x-\frac12\right)^2+y^2\leqslant\frac14\right\},
% \]
% 则
% \[
%   S=2\iint_D\sqrt{1+f_x'^2+f_y'^2}\mathrm dx\mathrm dy
%   =2\iint_D\frac{\mathrm dx\mathrm dy}{\sqrt{1-x^2-y^2}}.
% \]
% 利用极坐标变换得
% \[
% \begin{aligned}
%   S&=4\int_0^{\frac\pi2}\mathrm d\theta\int_0^{\cos\theta}
%   \frac{r\mathrm dr}{\sqrt{1-r^2}}
%   =4\int_0^{\frac\pi2}\left(-\sqrt{1-r^2}\right)\bigg|_0^{\cos\theta}\mathrm d\theta\\
%   &=4\int_0^{\frac\pi2}(1-\sin\theta)\mathrm d\theta=2\pi-4.
% \end{aligned}
% \qedhere
% \]
% \end{solve}
% 
% 
% \subsection{第一型曲面积分的定义}
% \label{subsec:ma-16-3-2}
% 
% 由于曲面的复杂性, 我们在这里只讨论光滑曲面上的第一型曲面积分。具体来说, 在本小节中我们所指的曲面 $S$ 将由参数方程 \eqref{eq:ma-16-3-1} 定义, 并且它们满足所假定的光滑条件。特别地, 我们总假定闭区域 $D\subset\mathbb R^2$ 的边界由有限条光滑曲线组成。设
% \[
%   \Delta=\{\Delta D_1,\Delta D_2,\cdots,\Delta D_K\}
% \]
% 为 $D$ 的一个分割, 且假定每个 $\Delta D_k(k=1,2,\cdots,K)$ 均为由有限条光滑曲线所围成的区域。在此假定下, $\Delta S_k=\{\boldsymbol r(u,v):(u,v)\in\Delta D_k\}(k=1,2,\cdots,K)$ 是一块光滑小曲面, $T=\{\Delta S_1,\allowbreak\Delta S_2,\allowbreak\cdots,\allowbreak\Delta S_K\}$ 是 $S$ 的一个分割。记 $\lambda(\Delta)=\max\limits_{1\leqslant k\leqslant K}\operatorname{diam}(\Delta D_k)$ 与 $\lambda(T)=\max\limits_{1\leqslant k\leqslant K}\operatorname{diam}(\Delta S_k)$。由于 $\boldsymbol r(u,v)$ 在 $D$ 上一致连续且其逆映射在 $S$ 上一致连续, 因此 $\lambda(\Delta)\to0$ 的充分必要条件是 $\lambda(T)\to0$.
% 
% 对于 $S$ 的分割 $T=\{\Delta S_1,\Delta S_2,\cdots,\Delta S_K\}$, 以下我们总是假定每个 $\Delta S_k$ 均是 $S$ 上可求面积的一小块曲面, 并且存在 $D$ 的可求面积的闭子区域 $\Delta D_k$, 使得 $\boldsymbol r(u,v)$ 是 $\Delta D_k$ 到 $\Delta S_k$ 的一一对应, 从而 $\Delta=\{\Delta D_1,\Delta D_2,\cdots,\Delta D_K\}$ 为 $D$ 的一个分割.
% 
% \begin{definition}\label{definition:ma-16-3-1}
% 设 $S\subset\mathbb R^3$ 是光滑曲面, 函数 $f(x,y,z)$ 在 $S$ 上有定义, 又设 $T=\{\Delta S_1,\allowbreak\Delta S_2,\allowbreak\cdots,\allowbreak\Delta S_K\}$ 是 $S$ 的一个分割, 且仍用 $\Delta S_k(k=1,2,\cdots,K)$ 来记每块小曲面的面积, 并记
% \[
%   \lambda(T)=\max_{1\leqslant k\leqslant K}\operatorname{diam}(\Delta S_k).
% \]
% 在 $\Delta S_k$ 上任取一点 $(\xi_k,\eta_k,\zeta_k)$, 作和式 $\sum\limits_{k=1}^K f(\xi_k,\eta_k,\zeta_k)\Delta S_k$。若存在常数 $I$, 使得对于 $\forall\varepsilon>0$, $\exists\delta>0$, 对于 $S$ 的任意分割 $T$ 及任取的 $(\xi_k,\eta_k,\zeta_k)\in\Delta S_k(k=1,2,\cdots,K)$, 当 $\lambda(T)<\delta$ 时, 有
% \[
%   \left|\sum_{k=1}^K f(\xi_k,\eta_k,\zeta_k)\Delta S_k-I\right|<\varepsilon,
% \]
% 即
% \[
%   \lim_{\lambda(T)\to0}\sum_{k=1}^K f(\xi_k,\eta_k,\zeta_k)\Delta S_k=I,
% \]
% 则称 $f(x,y,z)$ 在 $S$ 上的第一型曲面积分存在, 并称 $I$ 为 $f(x,y,z)$ 在 $S$ 上的第一型曲面积分, 记为
% \[
%   I=\iint_S f(x,y,z)\mathrm dS,
% \]
% 其中 $f(x,y,z)$ 称为被积函数, $S$ 称为积分曲面.
% \end{definition}
% 
% 从定义可以推出第一型曲面积分具有以下性质:
% 
% (1) 若在 $S$ 上函数 $f(x,y,z)\equiv1$, 则 $\iint_S f(x,y,z)\mathrm dS=\iint_S\mathrm dS$ 即为曲面 $S$ 的面积.
% 
% (2) 若曲面 $S\subset\mathbb R^2$, 则函数 $f(x,y)$ 在 $S$ 上的曲面积分即为二重积分 $\iint_S f(x,y)\mathrm dx\mathrm dy$.
% 
% 另外, 第一型曲面积分同样具有关于被积函数的线性性以及积分曲面的可加性等性质, 请读者自己给出上述性质的精确描述.
% 
% \subsection{第一型曲面积分的存在性与计算公式}
% \label{subsec:ma-16-3-3}
% 
% 对于光滑曲面上连续函数的第一型曲面积分, 我们有以下的存在性定理以及计算公式.
% 
% \begin{theorem}\label{theorem:ma-16-3-2}
% 设 $S\subset\mathbb R^3$ 是光滑曲面, 其参数方程为
% \[
%   \begin{cases}
%     x=x(u,v),\\
%     y=y(u,v),\\
%     z=z(u,v),
%   \end{cases}
%   \quad (u,v)\in D,
% \]
% 其中 $D$ 是由有限条光滑曲线围成的有界闭区域, 又设函数 $f(x,y,z)$ 在 $S$ 上连续, 则 $f(x,y,z)$ 在 $S$ 上的第一型曲面积分存在, 并且
% \begin{equation}
%   \iint_S f(x,y,z)\mathrm dS
%   =\iint_D f(x(u,v),y(u,v),z(u,v))\sqrt{EG-F^2}\mathrm du\mathrm dv,
%   \label{eq:ma-16-3-4}
% \end{equation}
% 其中 $E,G,F$ 由式 \eqref{eq:ma-16-3-3} 定义.
% \end{theorem}
% 
% 
% \begin{proof}
% 考查曲面 $S$ 的任一分割 $T=\{\Delta S_1,\Delta S_2,\cdots,\Delta S_K\}$, 记其相应于 $D$ 的分割为 $\Delta=\{\Delta D_1,\Delta D_2,\cdots,\Delta D_K\}$。在每个 $\Delta S_k(k=1,2,\cdots,K)$ 上任取点 $(\xi_k,\zeta_k,\eta_k)$, 则存在 $(u_k,v_k)\in\Delta D_k$, 使得
% \[
%   (\xi_k,\zeta_k,\eta_k)=(x(u_k,v_k),y(u_k,v_k),z(u_k,v_k)).
% \]
% 因为
% \[
%   \Delta S_k=\iint_{\Delta D_k}\sqrt{EG-F^2}\mathrm du\mathrm dv,
% \]
% 所以由重积分第一中值定理得
% \[
%   \Delta S_k=\sqrt{EG-F^2}\bigg|_{(u'_k,v'_k)}\Delta\sigma_k,
% \]
% 其中 $(u'_k,v'_k)\in\Delta D_k,\Delta\sigma_k$ 为 $\Delta D_k$ 的面积。因此
% \[
% \begin{aligned}
%   &\sum_{k=1}^K f(\xi_k,\zeta_k,\eta_k)\Delta S_k\\
%   &\quad=\sum_{k=1}^K f(x(u_k,v_k),y(u_k,v_k),z(u_k,v_k))
%     \sqrt{EG-F^2}\bigg|_{(u'_k,v'_k)}\Delta\sigma_k\\
%   &\quad=\sum_{k=1}^K f(x(u_k,v_k),y(u_k,v_k),z(u_k,v_k))
%     \sqrt{EG-F^2}\bigg|_{(u_k,v_k)}\Delta\sigma_k\\
%   &\qquad+\bigg[
%     \sum_{k=1}^K f(x(u_k,v_k),y(u_k,v_k),z(u_k,v_k))
%     \sqrt{EG-F^2}\bigg|_{(u'_k,v'_k)}\Delta\sigma_k\\
%   &\qquad\quad-
%     \sum_{k=1}^K f(x(u_k,v_k),y(u_k,v_k),z(u_k,v_k))
%     \sqrt{EG-F^2}\bigg|_{(u_k,v_k)}\Delta\sigma_k
%   \bigg]\\
%   &\quad\triangleq I_1+I_2.
% \end{aligned}
% \]
% 显然
% \[
%   \lim_{\lambda(T)\to0}I_1=\lim_{\lambda(\Delta)\to0}I_1
%   =\iint_D f(x(u,v),y(u,v),z(u,v))\sqrt{EG-F^2}\mathrm du\mathrm dv.
% \]
% 记 $M$ 为 $|f(x,y,z)|$ 在 $S$ 上的最大值, $M_k,m_k$ 分别为连续函数 $\sqrt{EG-F^2}$ 在 $\Delta D_k(k=1,2,\cdots,K)$ 上的最大、最小值, 则有
% \[
%   |I_2|\leqslant M\sum_{k=1}^K(M_k-m_k)\Delta\sigma_k\to0
%   \quad(\lambda(\Delta)\to0).
% \]
% 因此
% \[
% \begin{aligned}
%   \lim_{\lambda(T)\to0}\sum_{k=1}^K f(\xi_k,\zeta_k,\eta_k)\Delta S_k
%   &=\lim_{\lambda(T)\to0}(I_1+I_2)
%     =\lim_{\lambda(\Delta)\to0}(I_1+I_2)\\
%   &=\iint_D f(x(u,v),y(u,v),z(u,v))\sqrt{EG-F^2}\mathrm du\mathrm dv.
% \end{aligned}
% \]
% 这就证明了 $f(x,y,z)$ 在 $S$ 上的第一型曲面积分存在, 并且
% \[
%   \iint_S f(x,y,z)\mathrm dS
%   =\iint_D f(x(u,v),y(u,v),z(u,v))\sqrt{EG-F^2}\mathrm du\mathrm dv.
% \qedhere
% \]
% \end{proof}
% 
% 
% 当曲面 $S$ 由方程 $z=g(x,y)((x,y)\in D)$ 给出时, 由定理~\ref{theorem:ma-16-3-2} 有
% \[
%   \iint_S f(x,y,z)\mathrm dS
%   =\iint_D f(x,y,g(x,y))\sqrt{1+g_x'^2(x,y)+g_y'^2(x,y)}\mathrm dx\mathrm dy.
% \]
% 
% \begin{example}\label{example:ma-16-3-3}
% 求第一型曲面积分 $I=\iint_S x^2z\mathrm dS$, 其中 $S$ 为以下的曲面:
% 
% (1) 上半单位球面 $z=\sqrt{1-x^2-y^2}$;
% 
% (2) 单位球面 $x^2+y^2+z^2=1$.
% \end{example}
% 
% 
% \begin{solve}
% (1) 取 $D=\{(x,y):x^2+y^2\leqslant1\}$, 则
% \[
% \begin{aligned}
%   I&=\iint_S x^2z\mathrm dS\\
%   &=\iint_D x^2\cdot\sqrt{1-x^2-y^2}
%     \sqrt{1+\frac{x^2}{1-x^2-y^2}+\frac{y^2}{1-x^2-y^2}}\mathrm dx\mathrm dy\\
%   &=\iint_D x^2\mathrm dx\mathrm dy
%     =\frac12\iint_D(x^2+y^2)\mathrm dx\mathrm dy\\
%   &=\frac12\int_0^{2\pi}\mathrm d\theta\int_0^1r^3\mathrm dr=\frac\pi4.
% \end{aligned}
% \]
% 
% (2) 分别记
% \[
%   S_1:z=\sqrt{1-x^2-y^2},\quad (x,y)\in D,
% \]
% \[
%   S_2:z=-\sqrt{1-x^2-y^2},\quad (x,y)\in D,
% \]
% 则由积分曲面及被积函数的对称性得
% \[
%   \iint_{S_1}x^2z\mathrm dS=-\iint_{S_2}x^2z\mathrm dS,
% \]
% 所以
% \[
%   \iint_S x^2z\mathrm dS=0.
% \qedhere
% \]
% \end{solve}
% 
% 
% \begin{example}\label{example:ma-16-3-4}
% 求均匀物质曲面 $S:z=f(x,y)=2-(x^2+y^2)(z\geqslant0)$ 的质心坐标.
% \end{example}
% 
% 
% \begin{solve}
% 设其质心坐标为 $(x_0,y_0,z_0)$, 由对称性有 $x_0=y_0=0$, 并且由微元法有
% \[
%   z_0=\frac{M_{xy}}{\sigma}=\frac{\iint_S z\mathrm dS}{\iint_S\mathrm dS}
% \]
% (其中 $M_{xy}$ 是曲面 $S$ 到 $Oxy$ 平面的力矩, $\sigma$ 是曲面 $S$ 的面积)。取 $D=\{(x,y):x^2+y^2\leqslant2\}$。由 $\sqrt{1+f_x'^2+f_y'^2}=\sqrt{1+4x^2+4y^2}$ 有
% \[
% \begin{aligned}
%   \iint_S\mathrm dS
%   &=\iint_D\sqrt{1+4x^2+4y^2}\mathrm dx\mathrm dy
%     =\int_0^{2\pi}\mathrm d\theta\int_0^{\sqrt2}\sqrt{1+4r^2}r\mathrm dr\\
%   &=2\pi\cdot\frac1{12}(1+4r^2)^{\frac32}\bigg|_0^{\sqrt2}
%     =\frac{13}3\pi,
% \end{aligned}
% \]
% \[
% \begin{aligned}
%   \iint_S z\mathrm dS
%   &=\iint_D(2-x^2-y^2)\sqrt{1+4x^2+4y^2}\mathrm dx\mathrm dy\\
%   &=\int_0^{2\pi}\mathrm d\theta\int_0^{\sqrt2}r(2-r^2)\sqrt{1+4r^2}\mathrm dr
%     =\frac{37}{10}\pi,
% \end{aligned}
% \]
% 所以
% \[
%   z_0=\frac{111}{130}.
% \qedhere
% \]
% \end{solve}
% 
% 
% \section{第二型曲面积分}
% \label{sec:ma-16-4}
% 
% \subsection{曲面的侧}
% \label{subsec:ma-16-4-1}
% 
% 在日常生活中, 我们常见的曲面大都具有不同的两侧, 如一个自来水管具有内侧与外侧, 当水管破裂时, 管内的水将从内侧流向外侧.
% 
% 是不是所有的曲面都能分出两侧呢? 著名的麦比乌斯 (Möbius) 带就是一个单侧曲面。该曲面的构造非常简单: 取一条长方形纸带, 记为 $ABCD$, 将纸带扭转后, 让 $A$ 与 $C$, $B$ 与 $D$ 重合后粘起来, 使纸带构成一个环带 $S$, 该环带则称为麦比乌斯带。原长方形的边 $AB$ 与 $DC$ 构成该环带 $S$ 的边界。用颜色来涂该环带, 可以不越过其边界, 而将它全部涂遍颜色。由此我们无法分出该环带不同的侧 (见图 16.4.1).
% 
% 如何用精确的语言来描述一个曲面 $S$ 是单侧还是双侧的呢? 我们有以下定义.
% 
% \begin{definition}\label{definition:ma-16-4-1}
% 设 $S$ 是光滑曲面 ($S$ 可以是封闭的, 则此时无边界, 否则它具有边界)。对 $S$ 内部的任一点 $P_0$, 取定 $S$ 在 $P_0$ 处的一个法向的朝向。若一个动点 $P$ 从 $P_0$ 出发, 不越过 $S$ 的边界, 沿 $S$ 上任何路径 $\Gamma$ 运动回到 $P_0$ 处, $S$ 在 $P$ 的法向从 $P_0$ 处选定的朝向出发连续地沿路径 $\Gamma$ 变化回到 $P_0$ 处时, 总是保持原先在 $P_0$ 处选定的朝向, 则称 $S$ 为双侧曲面, 否则称 $S$ 为单侧曲面.
% \end{definition}
% 
% 由定义, 一个单侧曲面 $S$ 上必定存在点 $P_0\in S$, 以及 $S$ 上的一条经过 $P_0$ 的闭路径 $\Gamma$, 若在 $P_0$ 处选择 $S$ 的一个法向的朝向, 然后让 $P$ 从 $P_0$ 出发沿 $\Gamma$ 运动, $P$ 处的法向连续地变化, 当 $P$ 经过 $\Gamma$ 再回到 $P_0$ 时, 法向的朝向与原来确定的朝向相反.
% 
% 现设 $S$ 是一双侧曲面, 从定义可以看出在 $S$ 上任取一点 $P_0$, 然后取定 $S$ 在 $P_0$ 处一个法向的朝向, 则该曲面一侧法向的朝向即由该处法向的朝向所取定。这个事实请读者自己来验证。利用此事实, 我们可以容易地描述由一些参数方程表示的双侧曲面的定侧问题.
% 
% 例如, 若方程
% \[
%   \begin{cases}
%     x=x(u,v),\\
%     y=y(u,v),\\
%     z=z(u,v),
%   \end{cases}\quad (u,v)\in D
% \]
% 定义了一个光滑的双侧曲面 $S$, 令
% \begin{equation}
%   A=\frac{\partial(y,z)}{\partial(u,v)},\quad
%   B=\frac{\partial(z,x)}{\partial(u,v)},\quad
%   C=\frac{\partial(x,y)}{\partial(u,v)},
%   \label{eq:ma-16-4-1}
% \end{equation}
% 则 $\boldsymbol n=\pm(A,B,C)$ 即为 $S$ 的两个不同朝向的法向量。记 $\boldsymbol n_1=(A,B,C)$, $\boldsymbol n_2=(-A,-B,-C)$。当 $P_0\in S$, 且在 $P_0$ 处 $C>0$ 时, 我们将由 $\boldsymbol n_1$ 确定的一侧称为 $S$ 的上侧, $\boldsymbol n_2$ 确定的一侧称为 $S$ 下侧。类似地, 当 $A\ne0$ 时, 我们可以说 $S$ 有前侧和后侧, 而当 $B\ne0$ 时, $S$ 有左侧和右侧.
% 
% 若光滑曲面 $S$ 的方程为
% \[
%   z=f(x,y),\quad (x,y)\in D,
% \]
% 则 $S$ 上点 $(x,y,z)$ 处的法向量为 $\boldsymbol n=\pm(-f_x',-f_y',1)$。因此, 若 $\boldsymbol n$ 取 ``+'' 号, 便得到了 $S$ 的上侧; 若 $\boldsymbol n$ 取 ``-'' 号, 则得到 $S$ 的下侧.
% 
% \subsection{第二型曲面积分的定义}
% \label{subsec:ma-16-4-2}
% 
% 第二型曲面积分的物理背景是流量的计算问题。设区域 $D\subset\mathbb R^3$ 内有某一流体在流动 (如风的运动, 海水的运动等), 对于 $\forall(x,y,z)\in D$, 其速度函数为
% \[
%   \boldsymbol v(x,y,z)=(P(x,y,z),Q(x,y,z),R(x,y,z)),
% \]
% 简记为 $\boldsymbol v=(P,Q,R)$。若不考虑流体的种类, 上述速度函数也可认为是 $D$ 内的一个流速场。进一步假设 $\boldsymbol v$ 仅仅依赖于 $D$ 内点的位置, 而不依赖于时间的变化, 这样的流速场也称为稳定流速场.
% 
% 设 $S$ 是 $D$ 内的一个光滑双侧曲面。我们的问题是: 如何来求单位时间内流体流过 $S$ 的质量? 为了使问题简单化, 我们可设该流体的密度为 1。因此, 我们只要计算该流体通过 $S$ 的流量即可.
% 
% 若 $S$ 是平面的一部分, 其面积仍由 $S$ 表示, 而 $\boldsymbol v$ 是常向量函数, 即 $P=C_1,Q=C_2,R=C_3$ 均为常数函数, 记 $S$ 的单位法向量为 $\boldsymbol n$, 则不难看出单位时间内该流体流过 $S$ 的流量
% \[
%   W=(\boldsymbol v\cdot\boldsymbol n)S,
% \]
% 即流量 $W$ 为 $S$ 的面积与 $\boldsymbol v$ 在 $\boldsymbol n$ 上的投影的乘积。有了上述的计算公式, 利用积分思想就不难推导出在一般情形下流量的计算公式了.
% 
% 我们下面仍用微元法来求解此问题。在 $S$ 上任取一小块 $\Delta S$, 其面积仍然用 $\Delta S$ 记之。任取 $(\xi,\eta,\zeta)\in\Delta S$, 设 $\Delta S$ 在 $(\xi,\eta,\zeta)$ 处的单位法向量为 $\boldsymbol n(\xi,\eta,\zeta)$, 则可以认为 $\Delta S$ 近似地是法向量为 $\boldsymbol n(\xi,\eta,\zeta)$, 面积为 $\Delta S$ 的一小块平面。再假设 $\boldsymbol v$ 在 $\Delta S$ 上也近似为 $\boldsymbol v(\xi,\eta,\zeta)$, 则由上面的分析, 我们便得到了流量微元
% \[
%   \mathrm dW=\boldsymbol v(x,y,z)\cdot\boldsymbol n(x,y,z)\mathrm dS.
% \]
% 因此, 整个流量便是
% \[
%   W=\iint_S\boldsymbol v\cdot\boldsymbol n\,\mathrm dS
%   =\iint_S(P\cos\alpha+Q\cos\beta+R\cos\gamma)\mathrm dS,
% \]
% 其中 $\cos\alpha,\cos\beta,\cos\gamma$ 分别是 $\boldsymbol n$ 的方向余弦。由于上述积分中含有曲面法向量的方向余弦, 因此它不是第一型曲面积分。它其实是下面我们要讨论的第二型曲面积分.
% 
% 对于第二型曲面积分所涉及的曲面 $S$, 我们假定它是分片光滑的双侧曲面。这类曲面由有限块光滑曲面组成, 任何两块曲面均在边界相交, 并且不存在三块曲面具有相同一段弧作为边界。如图 16.4.2 中的曲面我们将不予考虑.
% 
% 现在我们给出分片光滑双侧曲面上的第二型曲面积分的定义.
% 
% \begin{definition}\label{definition:ma-16-4-2}
% 设 $S\subset\mathbb R^3$ 是分片光滑双侧曲面, 若它有边界, 则其边界由有限条分段光滑曲线组成。给定 $S$ 的一侧, $S$ 上每点处的单位法向量记为 $\boldsymbol n=(\cos\alpha,\cos\beta,\cos\gamma)$, 向量函数 $\boldsymbol f(x,y,z)=(P(x,y,z),Q(x,y,z),R(x,y,z))$ 在 $S$ 上有定义。对 $S$ 作任意分割 $T=\{\Delta S_1,\Delta S_2,\cdots,\Delta S_K\}$, 其中 $\Delta S_k$ 是以光滑曲线为边界的小光滑曲面, 并且仍用 $\Delta S_k(k=1,2,\cdots,K)$ 来记其面积。记 $\lambda(T)=\max_{1\leqslant k\leqslant K}\operatorname{diam}(\Delta S_k)$。在 $\Delta S_k$ 上任取一点 $(\xi_k,\eta_k,\zeta_k)(k=1,2,\cdots,K)$。若存在不依赖于分割 $T$ 以及 $(\xi_k,\eta_k,\zeta_k)(k=1,2,\cdots,K)$ 的选取的常数 $I$, 使得
% \begin{equation}
% \begin{aligned}
%   \lim_{\lambda(T)\to0}\sum_{k=1}^K
%   (&P(\xi_k,\eta_k,\zeta_k)\cos\alpha
%   +Q(\xi_k,\eta_k,\zeta_k)\cos\beta\\
%   &+R(\xi_k,\eta_k,\zeta_k)\cos\gamma)\Delta S_k=I,
% \end{aligned}
% \label{eq:ma-16-4-2}
% \end{equation}
% 则称 $I$ 是 $\boldsymbol f=(P,Q,R)$ 在 $S$ 上的第二型曲面积分, 记为
% \[
%   I=\iint_S P(x,y,z)\mathrm dy\mathrm dz
%     +Q(x,y,z)\mathrm dz\mathrm dx
%     +R(x,y,z)\mathrm dx\mathrm dy,
% \]
% 其中 $\boldsymbol f=(P,Q,R)$ 称为被积函数, $S$ 称为积分曲面.
% \end{definition}
% 
% 
% \paragraph*{注 1}
% 我们来解释一下第二型曲面积分记号的合理性。由定义, 一个第二型曲面积分是式 \eqref{eq:ma-16-4-2} 定义的极限。当小曲面 $\Delta S_k$ 的直径很小时, $\Delta S_k$ 近似于 $(\xi_k,\eta_k,\zeta_k)$ 处一小块切平面, 因此它的面积与 $\cos\alpha$ 的乘积近似为该小块切平面在 $Oyz$ 平面上的投影 $\Delta\sigma^k_{yz}$ 的有向面积, 即当 $\cos\alpha>0$ 时, 取其投影的面积为正, 否则为负。因此, 若仍由 $\Delta\sigma^k_{yz}$ 记其有向面积, 则
% \[
% \begin{aligned}
%   &\lim_{\lambda(T)\to0}\sum_{k=1}^K
%     P(\xi_k,\eta_k,\zeta_k)\cos\alpha\Delta S_k\\
%   &\quad=\lim_{\lambda(T)\to0}\sum_{k=1}^K
%     P(\xi_k,\eta_k,\zeta_k)\Delta\sigma^k_{yz}.
% \end{aligned}
% \]
% 上式右边可以自然地记成 $\iint_S P\,\mathrm dy\mathrm dz$。同理我们可解释记号 $\iint_S Q\,\mathrm dz\mathrm dx$ 和 $\iint_S R\,\mathrm dx\mathrm dy$ 的合理性.
% 
% 对于第二型曲面积分, 读者不难看出它们仍然具有关于被积函数的线性性以及积分曲面的相加性。值得注意的是, 当曲面相加时要指定曲面的侧。特别需要指出的是: 双侧曲面上同一个向量函数在两个侧上的第二型曲面积分的绝对值相等, 但相差一个负号.
% 
% \subsection{第二型曲面积分的存在性与计算公式}
% \label{subsec:ma-16-4-3}
% 
% 设光滑曲面 $S$ 由参数方程
% \begin{equation}
%   \begin{cases}
%     x=x(u,v),\\
%     y=y(u,v),\\
%     z=z(u,v),
%   \end{cases}\quad (u,v)\in D
%   \label{eq:ma-16-4-3}
% \end{equation}
% 给出, 则 $S$ 上每一点处的方向余弦为
% \[
%   \cos\alpha=\pm\frac{A}{\sqrt{A^2+B^2+C^2}},\quad
%   \cos\beta=\pm\frac{B}{\sqrt{A^2+B^2+C^2}},
% \]
% \[
%   \cos\gamma=\pm\frac{C}{\sqrt{A^2+B^2+C^2}},
% \]
% 其中 $A,B,C$ 由式 \eqref{eq:ma-16-4-1} 定义。值得指出的是, 它们都是 $(u,v)\in D$ 的连续函数.
% 
% \begin{theorem}\label{theorem:ma-16-4-1}
% 设 $S$ 为光滑双侧曲面, 其参数方程由 \eqref{eq:ma-16-4-3} 给出。选定 $S$ 的一侧, 记其单位法向量为
% \[
%   \boldsymbol n=\left(
%     \frac{A}{\sqrt{A^2+B^2+C^2}},
%     \frac{B}{\sqrt{A^2+B^2+C^2}},
%     \frac{C}{\sqrt{A^2+B^2+C^2}}
%   \right),
% \]
% 再设向量函数 $\boldsymbol f(x,y,z)=(P,Q,R)$ 在 $S$ 上连续, 则 $\boldsymbol f$ 在 $S$ 上的第二型曲面积分存在, 并且
% \[
%   \iint_S P\,\mathrm dy\mathrm dz+Q\,\mathrm dz\mathrm dx+R\,\mathrm dx\mathrm dy
%   =\iint_D(PA+QB+RC)\,\mathrm du\mathrm dv.
% \]
% \end{theorem}
% 
% 
% \begin{proof}
% 由 $S$ 的光滑性及式 \eqref{eq:ma-16-4-1} 知 $\boldsymbol n$ 在 $S$ 上是连续向量函数。由
% \begin{equation}
%   \iint_S P\,\mathrm dy\mathrm dz+Q\,\mathrm dz\mathrm dx+R\,\mathrm dx\mathrm dy
%   =\iint_S(P\cos\alpha+Q\cos\beta+R\cos\gamma)\mathrm dS,
%   \label{eq:ma-16-4-4}
% \end{equation}
% 并注意到上式的右边是连续函数 $P\cos\alpha+Q\cos\beta+R\cos\gamma$ 的第一型曲面积分, 因此左边的第二型曲面积分存在。定理的可积性部分得证.
% 
% 由公式 \eqref{eq:ma-16-3-4} 得
% \[
%   \iint_S(P\cos\alpha+Q\cos\beta+R\cos\gamma)\mathrm dS
%   =\iint_D(PA+QB+RC)\,\mathrm du\mathrm dv.
% \qedhere
% \]
% \end{proof}
% 
% 
% 特别地, 若光滑曲面 $S$ 的方程为
% \[
%   z=f(x,y),\quad (x,y)\in D,
% \]
% 并取定其上侧, 其中 $D\subset\mathbb R^2$ 是由分段光滑曲线所围成的有界闭区域。由于 $S$ 上的单位法向量平行于 $(-f_x',-f_y',1)$, 容易推出
% \[
%   \iint_S P(x,y,z)\mathrm dy\mathrm dz+Q(x,y,z)\mathrm dz\mathrm dx+R(x,y,z)\mathrm dx\mathrm dy
% \]
% \[
%   =\iint_D[-P(x,y,f(x,y))f_x'(x,y)-Q(x,y,f(x,y))f_y'(x,y)
%   +R(x,y,f(x,y))]\mathrm dx\mathrm dy.
% \]
% 
% \begin{example}\label{example:ma-16-4-2}
% 计算第二型曲面积分
% \[
%   I=\iint_S x\mathrm dy\mathrm dz+y\mathrm dz\mathrm dx+z\mathrm dx\mathrm dy,
% \]
% 其中 $S$ 为由三个坐标平面及平面 $x+y+z=1$ 所围四面体的外侧.
% \end{example}
% 
% 
% \begin{solve}[{解法 1}]
% 记 $S$ 落在 $Oxy,Oyz,Ozx$ 平面上的部分分别记为 $S_z,S_x$ 及 $S_y$, 落在平面 $x+y+z=1$ 的部分记为 $S_1$ (见图 16.4.3)。在 $S_z$ 上, $z=0$, $\mathrm dy\mathrm dz=\mathrm dz\mathrm dx=0$, 从而
% \[
%   \iint_{S_z}x\mathrm dy\mathrm dz+y\mathrm dz\mathrm dx+z\mathrm dx\mathrm dy=0.
% \]
% 同理, 在 $S_y$ 与 $S_x$ 上的积分都为零。因此
% \[
%   I=\iint_{S_1}x\mathrm dy\mathrm dz+y\mathrm dz\mathrm dx+z\mathrm dx\mathrm dy.
% \]
% 记 $D=\{(x,y):x\geqslant0,y\geqslant0,x+y\leqslant1\}$, 则由对称性有
% \[
%   I=3\iint_D(1-x-y)\mathrm dx\mathrm dy
%   =3\int_0^1\mathrm dx\int_0^{1-x}(1-x-y)\mathrm dy=\frac12.
% \qedhere
% \]
% \end{solve}
% 
% 
% \begin{solve}[{解法 2}]
% 与解法 1 同理得到
% \[
%   I=\iint_{S_1}x\mathrm dy\mathrm dz+y\mathrm dz\mathrm dx+z\mathrm dx\mathrm dy.
% \]
% 由于在 $S_1$ 上点 $(x,y,z)$ 处的方向余弦 $(\cos\alpha,\cos\beta,\cos\gamma)$ 即为 $S_1$ 的单位法向量, 且
% \[
%   \cos\alpha=\cos\beta=\cos\gamma=\frac1{\sqrt3},
% \]
% 而 $S_1$ 是一个三角形, 其面积为 $\frac{\sqrt3}{2}$, 因此
% \[
%   I=\iint_{S_1}\frac{x+y+z}{\sqrt3}\mathrm dS
%   =\frac1{\sqrt3}\iint_{S_1}\mathrm dS=\frac12.
% \qedhere
% \]
% \end{solve}
% 
% 
% \section{各类积分之间的联系}
% \label{sec:ma-16-5}
% 
% 设 $f(x)\in C^1[a,b]$。众所周知, 牛顿-莱布尼茨公式
% \[
%   \int_a^b\mathrm df(x)=f(b)-f(a)
% \]
% 是一元微积分中最重要的公式。换种观点来看该公式, 如果我们定义一个函数 $f(x)$ 在 $a,b(a<b)$ 两点上的积分为 $f(b)-f(a)$, 则上述公式建立了区间 $[a,b]$ 上 $f'(x)$ 的积分与 $f(x)$ 在端点处的积分的联系。因此, 一个自然的问题是: 是否平面或空间中一个几何体的边界与内部上的某些积分总存在联系? 本节中我们将讨论这个问题.
% 
% \subsection{格林公式}
% \label{subsec:ma-16-5-1}
% 
% 格林 (Green) 公式给出的是平面区域上函数的二重积分与其相关的函数在边界上的第二型曲线积分之间的联系。在这里, 我们只讨论 $\mathbb R^2$ 中的有界闭区域 $D$, 其边界是有限条约当曲线, 且每条约当曲线都是分段光滑的。因此 $D$ 一定具有如图 16.5.1 的形状.
% 
% 对于平面 $\mathbb R^2$ 内一条约当曲线 $\Gamma$, 以 $\Gamma$ 为边界的区域有两个: 一个是有界区域, 该区域也称为 $\Gamma$ 的内部; 而另一个是无界区域, 该区域也称为 $\Gamma$ 的外部。如在图 16.5.1 中, $D$ 位于 $\Gamma_0$ 的内部, 而在 $\Gamma_k(k=1,2,\cdots,K)$ 的外部。一个区域 $D$ 称为单连通区域, 如果位于 $D$ 内的任一条约当曲线的内部总包含在 $D$ 内; 否则称 $D$ 是多连通区域。在我们假定的情形下, 若 $D$ 是一条约当曲线的内部, 则它是单连通的; 若 $D$ 由 $K(K>1)$ 条约当曲线围成, 则 $D$ 是多连通的, 此时也称 $D$ 是 $K$ 连通的.
% 
% 一条约当曲线具有两个互为相反的方向。现设 $D$ 是由有限条约当曲线所围成的区域。我们在第二册曾给出 $\partial D$ 的定向, 即若一个人沿着 $D$ 的边界曲线上的一个方向前进时, 曲线所围的区域总在他的左边, 则称该方向为 $\partial D$ 的正向。如图 16.5.1 所示, 相对 $D$ 来说, $\Gamma_0$ 的正向是逆时针方向, 而 $\Gamma_k(k=1,2,\cdots,K)$ 的正向是顺时针方向.
% 
% 有了上述准备, 我们可以证明下述定理.
% 
% \begin{theorem}[{格林公式}]\label{theorem:ma-16-5-1}
% 设 $D\subset\mathbb R^2$ 是有界闭区域, 其边界由有限条分段光滑的约当曲线组成。若函数 $P(x,y),Q(x,y)$ 在 $D$ 上具有连续偏导数, 则有
% \[
%   \int_{\partial D}P\mathrm dx+Q\mathrm dy
%   =\iint_D\left(\frac{\partial Q}{\partial x}-\frac{\partial P}{\partial y}\right)\mathrm dx\mathrm dy,
% \]
% 其中 $\partial D$ 的方向为正向.
% \end{theorem}
% 
% 
% \begin{proof}
% 我们先证 $D$ 是单连通区域的情形。证明分以下几步进行:
% 
% (1) 设 $D$ 是 $Y$ 型区域, 即
% \[
%   D=\{(x,y):\varphi(y)\leqslant x\leqslant\psi(y),c\leqslant y\leqslant d\},
% \]
% 如图 16.5.2 所示。下面我们证明
% \[
%   \int_{\partial D}Q(x,y)\mathrm dy=\iint_D\frac{\partial Q}{\partial x}\mathrm dx\mathrm dy.
% \]
% 我们有
% \[
%   \iint_D\frac{\partial Q}{\partial x}\mathrm dx\mathrm dy
%   =\int_c^d\mathrm dy\int_{\varphi(y)}^{\psi(y)}
%   \frac{\partial Q}{\partial x}\mathrm dx
%   =\int_c^d[Q(\psi(y),y)-Q(\varphi(y),y)]\mathrm dy.
% \]
% 由
% \[
% \begin{aligned}
%   \int_{\partial D}Q(x,y)\mathrm dy
%   &=\int_{\widehat{AB}}Q(x,y)\mathrm dy+\int_{\widehat{BC}}Q(x,y)\mathrm dy
%     +\int_{\widehat{CE}}Q(x,y)\mathrm dy+\int_{\widehat{EA}}Q(x,y)\mathrm dy,\\
%   \int_{\widehat{AB}}Q(x,y)\mathrm dy&=\int_{\widehat{CE}}Q(x,y)\mathrm dy=0,\\
%   \int_{\widehat{BC}}Q(x,y)\mathrm dy&=\int_c^d Q(\psi(y),y)\mathrm dy,
% \end{aligned}
% \]
% 及
% \[
%   \int_{\widehat{EA}}Q(x,y)\mathrm dy
%   =\int_d^cQ(\varphi(y),y)\mathrm dy
%   =-\int_c^dQ(\varphi(y),y)\mathrm dy,
% \]
% 因此有
% \[
%   \int_{\partial D}Q(x,y)\mathrm dy=\iint_D\frac{\partial Q}{\partial x}\mathrm dx\mathrm dy.
% \]
% 
% (2) 若 $D$ 是 $X$ 型区域, 与 (1) 同理可证
% \[
%   \int_{\partial D}P(x,y)\mathrm dx=-\iint_D\frac{\partial P}{\partial y}\mathrm dx\mathrm dy.
% \]
% 
% (3) 若 $D$ 既是 $X$ 型又是 $Y$ 型区域, 则成立
% \[
%   \int_{\partial D}P\mathrm dx+Q\mathrm dy
%   =\iint_D\left(\frac{\partial Q}{\partial x}-\frac{\partial P}{\partial y}\right)\mathrm dx\mathrm dy.
% \]
% 
% 若区域 $D$ 加上有限条光滑曲线后, 可分成有限个既是 $X$ 型又是 $Y$ 型区域, 则可在每个小区域上应用上述结果。将每个小区域边界上的曲线积分与区域上的重积分分别相加, 由于添加的曲线必然是两个小区域的公共边界, 因此在曲线积分求和中出现两次, 但因方向相反而相互抵消, 所以总有
% \[
%   \int_{\partial D}P\mathrm dx+Q\mathrm dy
%   =\iint_D\left(\frac{\partial Q}{\partial x}-\frac{\partial P}{\partial y}\right)\mathrm dx\mathrm dy.
% \]
% 对一般的单连通区域, 可以用分成有限个既是 $X$ 型又是 $Y$ 型区域的区域逼近方法证之, 但这方面内容在数学分析课程中难以介绍, 因此我们略去该证明.
% 
% 当 $D$ 是 $K(K>1)$ 连通时, 我们同样可以在区域 $D$ 内添加若干条光滑曲线, 将 $D$ 分成有限个单连通区域 (见图 16.5.3)。在每个小单连通区域上由上所证成立格林公式, 然后将每个小区域边界上的曲线积分与区域上的重积分分别相加, 注意到所添加的光滑曲线仍然是两个小区域的公共边界, 因此和式中在所添加曲线上的曲线积分将相互抵消, 所以格林公式对 $K$ 连通区域仍然成立.
% \end{proof}
% 
% 
% 在上述格林公式中, 我们假定了函数 $P$ 和 $Q$ 在 $D$ 上具有连续偏导数, 这就要求 $P,Q$ 在包含 $D$ 的某个开集内具有连续偏导数。另外, 对于有限条封闭曲线 $\Gamma$ 所围的区域 $D$, 为了强调曲线的封闭性, 今后我们常常用
% \[
%   \oint_{\partial D}P\mathrm dx+Q\mathrm dy
% \]
% 来表示沿 $\partial D$ 关于 $D$ 的正向的积分.
% 
% 对于格林公式, 我们还可以有以下形式。设光滑曲线 $\Gamma$ 由
% \[
%   \begin{cases}
%     x=x(t),\\
%     y=y(t),
%   \end{cases}\quad (t\in[\alpha,\beta])
% \]
% 给出, $t$ 的增加方向恰好是 $\Gamma$ 的正向。记该曲线在 $(x,y)$ 处的单位切向量为 $\boldsymbol v$, 则有
% \[
%   \mathrm dx=x'(t)\mathrm dt
%   =\frac{x'(t)}{\sqrt{x'^2(t)+y'^2(t)}}\sqrt{x'^2(t)+y'^2(t)}\mathrm dt
%   =\cos(\boldsymbol v,x)\mathrm ds,
% \]
% 其中 $\cos(\boldsymbol v,x)$ 为切向量 $\boldsymbol v$ 与 $x$ 轴正向的夹角的余弦。同理, 有
% \[
%   \mathrm dy=y'(t)\mathrm dt=\cos(\boldsymbol v,y)\mathrm ds,
% \]
% 其中 $\cos(\boldsymbol v,y)$ 为切向量 $\boldsymbol v$ 与 $y$ 轴正向的夹角的余弦。设函数 $P(x,y),Q(x,y)$ 在 $\Gamma$ 上连续, 则
% \[
%   \int_\Gamma P\mathrm dx+Q\mathrm dy
%   =\int_\Gamma[P\cos(\boldsymbol v,x)+Q\cos(\boldsymbol v,y)]\mathrm ds.
% \]
% 现取 $\Gamma$ 在 $(x,y)$ 处的单位法向量为 $\boldsymbol n$, 使得 $\boldsymbol n$ 与 $\boldsymbol v$ 成右手系。这样, 在封闭曲线 $\Gamma$ 围成区域 $D$ 时, 若取 $\Gamma$ 的正向, 则选取的法向量指向 $D$ 的外部 (见图 16.5.4), 从而有
% \[
%   \begin{cases}
%     \cos(\boldsymbol n,x)=\cos(\boldsymbol v,y),\\
%     \cos(\boldsymbol n,y)=-\cos(\boldsymbol v,x),
%   \end{cases}
% \]
% 其中 $\cos(\boldsymbol n,x)$ 与 $\cos(\boldsymbol n,y)$ 分别为法向 $\boldsymbol n$ 与 $x$ 轴, $y$ 轴正向的夹角的余弦。因此
% \[
% \begin{aligned}
%   &\int_\Gamma[P\cos(\boldsymbol n,x)+Q\cos(\boldsymbol n,y)]\mathrm ds\\
%   &=\int_\Gamma[-Q\cos(\boldsymbol v,x)+P\cos(\boldsymbol v,y)]\mathrm ds\\
%   &=\int_\Gamma(-Q)\mathrm dx+P\mathrm dy
%   =\iint_D\left(\frac{\partial P}{\partial x}+\frac{\partial Q}{\partial y}\right)\mathrm dx\mathrm dy.
% \end{aligned}
% \]
% 
% 另外, 在定积分的应用中, 我们曾用
% \[
%   \int_\Gamma x\mathrm dy=-\int_\Gamma y\mathrm dx
%   =\frac12\int_\Gamma x\mathrm dy-y\mathrm dx
% \]
% 来表示封闭曲线 $\Gamma$ 所围区域 $D$ 的面积。现在利用格林公式可以清楚看到上述三个积分都等于 $\iint_D\mathrm dx\mathrm dy$, 即 $D$ 的面积.
% 
% \begin{example}\label{example:ma-16-5-1}
% 设 $\Gamma=\widehat{AB}\subset\mathbb R^2$ 是从点 $A(0,0)$ 到点 $B(0,1)$ 的任一条光滑曲线, 试计算第二型曲线积分
% \[
%   \int_\Gamma 2xy\mathrm dx+(x^2-y^2)\mathrm dy.
% \]
% \end{example}
% 
% 
% \begin{solve}
% 由于 $\Gamma$ 是任一条光滑曲线, 因此直接计算将无法求出积分值。注意到若令
% \[
%   P(x,y)=2xy,\quad Q(x,y)=x^2-y^2,
% \]
% 则有
% \[
%   \frac{\partial Q}{\partial x}=2x=\frac{\partial P}{\partial y}.
% \]
% 因此对于 $\forall(x,y)\in\mathbb R^2$, 有 $\frac{\partial Q}{\partial x}-\frac{\partial P}{\partial y}=0$.
% 
% 为了应用格林公式, 设实轴上从点 $(0,0)$ 到点 $(0,1)$ 的线段为 $\Gamma_2$。我们再取一条光滑曲线 $\Gamma_1$, 使得 $\Gamma_1$ 和 $\Gamma_2$ 都以点 $(0,0),(0,1)$ 为端点, 并且 $\Gamma\cup\Gamma_1^-$ 和 $\Gamma_1\cup\Gamma_2^-$ 均为约当区域 (见图 16.5.5)。设 $\Gamma\cup\Gamma_1^-$ 所围的有界闭区域为 $D_1$, 而 $\Gamma_1\cup\Gamma_2^-$ 所围的有界闭区域为 $D_2$, 由格林公式有
% \[
% \begin{aligned}
%   \int_\Gamma P\mathrm dx+Q\mathrm dy
%   &=\int_{\Gamma_1}P\mathrm dx+Q\mathrm dy+\int_{\Gamma\cup\Gamma_1^-}P\mathrm dx+Q\mathrm dy\\
%   &=\int_{\Gamma_1}P\mathrm dx+Q\mathrm dy+\iint_{D_1}0\mathrm dx\mathrm dy\\
%   &=\int_{\Gamma_2}P\mathrm dx+Q\mathrm dy+\int_{\Gamma_1\cup\Gamma_2^-}P\mathrm dx+Q\mathrm dy\\
%   &=\int_{\Gamma_2}P\mathrm dx+Q\mathrm dy+\iint_{D_2}0\mathrm dx\mathrm dy\\
%   &=\int_0^1 0\mathrm dx=0.
% \end{aligned}
% \qedhere
% \]
% \end{solve}
% 
% 
% \begin{example}\label{example:ma-16-5-2}
% 计算第二型曲线积分
% \[
%   \oint_\Gamma\frac{x\mathrm dy-y\mathrm dx}{x^2+y^2},
% \]
% 其中 $\Gamma\subset\mathbb R^2$ 为任一条光滑约当曲线, 且 $(0,0)\notin\Gamma$.
% \end{example}
% 
% 
% \begin{solve}
% 记
% \[
%   P(x,y)=\frac{-y}{x^2+y^2},\quad Q(x,y)=\frac{x}{x^2+y^2},
% \]
% 则对任意的 $(x,y)\ne(0,0)$, 有
% \[
%   \frac{\partial Q}{\partial x}=\frac{y^2-x^2}{(x^2+y^2)^2}=\frac{\partial P}{\partial y}.
% \]
% 因此, 设 $\Gamma$ 所围成的区域为 $D$, 则当 $(0,0)\notin D$ 时, 由格林公式有
% \[
%   \oint_\Gamma\frac{x\mathrm dy-y\mathrm dx}{x^2+y^2}
%   =\iint_D0\mathrm dx\mathrm dy=0.
% \]
% 当 $(0,0)\in D$ 时, 格林公式条件不满足。对于 $R>0$, 设
% \[
%   \Gamma_R=\{(x,y):x^2+y^2=R^2\}.
% \]
% 取 $R$ 充分大, 使得 $\Gamma$ 落在 $\Gamma_R$ 的内部。如图 16.5.6, 则 $\Gamma$ 与 $\Gamma_R$ 围成一个二连通区域 $D_R$, 在此区域上应用格林公式, 有
% \[
%   0=\iint_{D_R}\left(\frac{\partial Q}{\partial x}-\frac{\partial P}{\partial y}\right)\mathrm dx\mathrm dy
%   =\int_{\Gamma_R}P\mathrm dx+Q\mathrm dy+\int_{\Gamma^-}P\mathrm dx+Q\mathrm dy.
% \]
% 因此, 我们有
% \[
% \begin{aligned}
%   \int_\Gamma P\mathrm dx+Q\mathrm dy
%   &=\int_{\Gamma_R}\frac{x\mathrm dy-y\mathrm dx}{x^2+y^2}
%     =\frac1{R^2}\int_{\Gamma_R}x\mathrm dy-y\mathrm dx\\
%   &=\frac2{R^2}\iint_{D_R}\mathrm dx\mathrm dy=2\pi.
% \end{aligned}
% \qedhere
% \]
% \end{solve}
% 
% 
% \subsection{高斯公式}
% \label{subsec:ma-16-5-2}
% 
% 在本小节中, 设 $D\subset\mathbb R^3$ 是有界闭区域, $\partial D$ 由有限块光滑双侧曲面所组成。选定 $\partial D$ 的侧为 $D$ 的外侧, 即 $D$ 的外法向量确定的侧 (此时也称外侧为 $\partial D$ 的正向)。高斯 (Gauss) 公式将给出函数在 $D$ 上的三重积分与相关函数在 $\partial D$ 上的曲面积分的关系.
% 
% \begin{theorem}[{高斯公式}]\label{theorem:ma-16-5-2}
% 设 $D\subset\mathbb R^3$ 是有界闭区域, $\partial D$ 是 $D$ 的外侧, 函数 $P(x,y,z)$、$Q(x,y,z)$ 和 $R(x,y,z)$ 在 $D$ 上具有连续偏导数, 则有
% \[
%   \iint_{\partial D}P\mathrm dy\mathrm dz+Q\mathrm dz\mathrm dx+R\mathrm dx\mathrm dy
%   =\iiint_D\left(\frac{\partial P}{\partial x}+\frac{\partial Q}{\partial y}
%   +\frac{\partial R}{\partial z}\right)\mathrm dx\mathrm dy\mathrm dz.
% \]
% \end{theorem}
% 
% 
% \begin{proof}
% 高斯公式的证明与格林公式的证明具有相似性, 基本思想是: 先设 $D$ 是一些特殊的区域, 然后再用区域逼近来证明一般区域上的高斯公式。同样, 区域逼近这一步在这里也不作介绍.
% 
% (1) 设 $D$ 是如下的闭区域:
% \[
%   D=\{(x,y,z):(y,z)\in\Omega,\varphi(y,z)\leqslant x\leqslant\psi(y,z)\}
% \]
% (通常也称为 $X$ 型区域), 其中 $\Omega\subset\mathbb R^2$ 是由有限多条光滑闭曲线所围的有界闭区域, $\varphi(y,z)$ 与 $\psi(y,z)$ 是 $\Omega$ 上的连续函数。我们来证明
% \[
%   \iint_{\partial D}P(x,y,z)\mathrm dy\mathrm dz
%   =\iiint_D\frac{\partial P}{\partial x}\mathrm dx\mathrm dy\mathrm dz.
% \]
% 其证明方法是将上式两边的积分都化为 $\Omega$ 上的二重积分。记
% \[
% \begin{aligned}
%   S_1&=\{(x,y,z):x=\psi(y,z),(y,z)\in\Omega\},\quad \text{取前侧};\\
%   S_2&=\{(x,y,z):x=\varphi(y,z),(y,z)\in\Omega\},\quad \text{取后侧};\\
%   S_3&=\{(x,y,z):\varphi(y,z)\leqslant x\leqslant\psi(y,z),(y,z)\in\Omega\},
%     \quad \text{取 $S_3$ 的侧为 $S$ 的外侧在 $S_3$ 的限制}.
% \end{aligned}
% \]
% 由于
% \[
%   \iint_{S_1}P(x,y,z)\mathrm dy\mathrm dz
%   =\iint_\Omega P(\psi(y,z),y,z)\mathrm dy\mathrm dz,
% \]
% \[
%   \iint_{S_2}P(x,y,z)\mathrm dy\mathrm dz
%   =-\iint_\Omega P(\varphi(y,z),y,z)\mathrm dy\mathrm dz,\quad
%   \iint_{S_3}P(x,y,z)\mathrm dy\mathrm dz=0,
% \]
% 因此
% \[
%   \iint_{\partial D}P(x,y,z)\mathrm dy\mathrm dz
%   =\iint_\Omega[P(\psi(y,z),y,z)-P(\varphi(y,z),y,z)]\mathrm dy\mathrm dz.
% \]
% 而
% \[
% \begin{aligned}
%   \iiint_D\frac{\partial P}{\partial x}\mathrm dx\mathrm dy\mathrm dz
%   &=\iint_\Omega\mathrm dy\mathrm dz
%     \int_{\varphi(y,z)}^{\psi(y,z)}\frac{\partial P}{\partial x}\mathrm dx\\
%   &=\iint_\Omega[P(\psi(y,z),y,z)-P(\varphi(y,z),y,z)]\mathrm dy\mathrm dz,
% \end{aligned}
% \]
% 所以
% \[
%   \iint_{\partial D}P(x,y,z)\mathrm dy\mathrm dz
%   =\iiint_D\frac{\partial P}{\partial x}\mathrm dx\mathrm dy\mathrm dz.
% \]
% 
% (2) 若 $D=\{(x,y,z):(x,y)\in\Omega,\varphi(x,y)\leqslant z\leqslant\psi(x,y)\}$ (称之为 $Z$ 型区域, 其中 $\Omega,\varphi,\psi$ 所满足的条件与 (1) 中的相同), 则我们可以证明
% \[
%   \iint_{\partial D}R(x,y,z)\mathrm dx\mathrm dy
%   =\iiint_D\frac{\partial R}{\partial z}\mathrm dx\mathrm dy\mathrm dz.
% \]
% 
% (3) 当 $D=\{(x,y,z):(z,x)\in\Omega,\varphi(z,x)\leqslant y\leqslant\psi(z,x)\}$ (也称为 $Y$ 型区域, 其中 $\Omega,\varphi,\psi$ 所满足的条件与 (1) 中的相同) 时, 我们可以证明
% \[
%   \iint_{\partial D}Q(x,y,z)\mathrm dz\mathrm dx
%   =\iiint_D\frac{\partial Q}{\partial y}\mathrm dx\mathrm dy\mathrm dz.
% \]
% 
% (4) 显然, 当 $D$ 刚好同时是这三种区域时, 定理结论成立。对于一般的中间无 ``洞'' 的有界闭区域 $D$, 当区域 $D$ 可以用有限块光滑曲面将其分成一些小闭区域, 且使得每个小闭区域同时是 $X$ 型, $Y$ 型与 $Z$ 型区域时 (如四面体区域), 高斯公式在每个小区域上成立。再将每个小区域上相应积分相加, 注意到所添加曲面必是两个小区域的公共边界, 因此在第二型曲面积分中出现两次, 但因它们的侧刚好相反, 从而相互抵消。所以对这类区域, 高斯公式依然成立。若 $D$ 不能分成有限块同时是上述的 $X$ 型, $Y$ 型和 $Z$ 型区域, 我们不加证明地指出, 利用分成同时是上述三种区域的区域逼近方法, 可以证明对于 $D$ 仍成立高斯公式.
% 
% (5) 若 $D$ 是中间有 ``洞'' 的情形。此时, 我们亦可添加若干块分片光滑的曲面, 将 $D$ 分成有限个中间无 ``洞'' 的小区域。由上所证, 高斯公式在这些小区域上成立。注意到所添加曲面必是两个小区域的公共边界, 因此对这类相当广泛的区域, 高斯公式依然成立.
% \end{proof}
% 
% 
% \begin{remark}\label{remark:ma-16-source-51}
% 对于区域有 ``洞'' 的情形, 区域 $D$ 内的 ``洞'' $D_1$ 的边界曲面的侧相对于区域 $D$ 来说, 应该是 $\partial D_1$ 的内侧 (见图 16.5.7).
% \end{remark}
% 
% 
% \begin{example}\label{example:ma-16-5-3}
% 计算第二型曲面积分
% \[
%   I=\iint_S(\sin yz+x)\mathrm dy\mathrm dz+(e^{xz}+y)\mathrm dz\mathrm dx+(xy+z)\mathrm dx\mathrm dy,
% \]
% 其中 $S$ 是单位球面上半部分的上侧, 即
% \[
%   S=\{(x,y,z):x^2+y^2+z^2=1,z\geqslant0\}.
% \]
% \end{example}
% 
% 
% \begin{solve}
% 直接计算似乎比较复杂, 因此利用高斯公式来计算。由于 $S$ 不封闭, 若我们取 $Oxy$ 平面上的单位圆盘 $S_1=\{(x,y):x^2+y^2\leqslant1\}$, 并取 $S_1$ 的下侧, 则 $S\cup S_1$ 构成了上半单位球体 $D$ 的边界, 且取外侧。由高斯公式得
% \[
% \begin{aligned}
%   &\iint_{S\cup S_1}(\sin yz+x)\mathrm dy\mathrm dz+(e^{xz}+y)\mathrm dz\mathrm dx+(xy+z)\mathrm dx\mathrm dy\\
%   &=3\iiint_D\mathrm dx\mathrm dy\mathrm dz
%     =3\cdot\frac12\cdot\frac{4\pi}{3}=2\pi,
% \end{aligned}
% \]
% 而在 $S_1$ 上 $\mathrm dy\mathrm dz=0,\mathrm dz\mathrm dx=0$, 又
% \[
%   \iint_{S_1}(xy+z)\mathrm dx\mathrm dy
%   =-\iint_{x^2+y^2\leqslant1}xy\mathrm dx\mathrm dy=0,
% \]
% 所以 $I=2\pi$.
% \end{solve}
% 
% 
% \begin{example}\label{example:ma-16-5-4}
% 设 $S\subset\mathbb R^3$ 为封闭光滑曲面, 取外侧, 以它为边界的有界闭区域是 $D$。已知点 $(\xi,\eta,\zeta)\in\mathbb R^3$ 不在 $S$ 上。计算高斯积分
% \[
%   I=\iint_S\frac{\cos(\boldsymbol r,\boldsymbol n)}{r^2}\mathrm dS,
% \]
% 其中 $\boldsymbol r=(x-\xi,y-\eta,z-\zeta)$, $r=|\boldsymbol r|$, $\boldsymbol n$ 是 $S$ 的单位外法向量.
% \end{example}
% 
% 
% \begin{solve}
% 由于
% \[
%   \cos(\boldsymbol r,\boldsymbol n)
%   =\frac{x-\xi}{r}\cos(\boldsymbol n,x)
%    +\frac{y-\eta}{r}\cos(\boldsymbol n,y)
%    +\frac{z-\zeta}{r}\cos(\boldsymbol n,z),
% \]
% 因此
% \[
%   I=\iint_S\frac{x-\xi}{r^3}\mathrm dy\mathrm dz
%     +\frac{y-\eta}{r^3}\mathrm dz\mathrm dx
%     +\frac{z-\zeta}{r^3}\mathrm dx\mathrm dy.
% \]
% 因为
% \[
%   \frac{\partial}{\partial x}\left(\frac{x-\xi}{r^3}\right)=\frac1{r^3}-\frac{3(x-\xi)^2}{r^5},\quad
%   \frac{\partial}{\partial y}\left(\frac{y-\eta}{r^3}\right)=\frac1{r^3}-\frac{3(y-\eta)^2}{r^5},
% \]
% \[
%   \frac{\partial}{\partial z}\left(\frac{z-\zeta}{r^3}\right)=\frac1{r^3}-\frac{3(z-\zeta)^2}{r^5},
% \]
% 所以
% \[
%   \frac{\partial}{\partial x}\left(\frac{x-\xi}{r^3}\right)
%   +\frac{\partial}{\partial y}\left(\frac{y-\eta}{r^3}\right)
%   +\frac{\partial}{\partial z}\left(\frac{z-\zeta}{r^3}\right)=0.
% \]
% 因此, 当 $(\xi,\eta,\zeta)\notin D$ 时, 由高斯公式得 $I=0$.
% 
% 当 $(\xi,\eta,\zeta)\in D$ 时, 我们可取 $\varepsilon$ 充分小, 使得球面
% \[
%   S_\varepsilon=\{(x,y,z):(x-\xi)^2+(y-\eta)^2+(z-\zeta)^2=\varepsilon^2\}
% \]
% 完全落在 $D$ 的内部。取 $S_\varepsilon$ 的内侧 $S_\varepsilon^-$, 设闭区域 $D_\varepsilon$ 以 $S\cup S_\varepsilon^-$ 为边界, 则
% \[
%   \iint_{S\cup S_\varepsilon^-}\frac{\cos(\boldsymbol r,\boldsymbol n)}{r^2}\mathrm dS
%   =\iiint_{D_\varepsilon}0\mathrm dx\mathrm dy\mathrm dz=0.
% \]
% 注意到在 $S_\varepsilon$ 上, $\boldsymbol r$ 与 $\boldsymbol n$ 平行, 所以
% \[
% \begin{aligned}
%   \iint_S\frac{\cos(\boldsymbol r,\boldsymbol n)}{r^2}\mathrm dS
%   &=-\iint_{S_\varepsilon^-}\frac{\cos(\boldsymbol r,\boldsymbol n)}{r^2}\mathrm dS
%     =\iint_{S_\varepsilon}\frac{\cos(\boldsymbol r,\boldsymbol n)}{r^2}\mathrm dS\\
%   &=\iint_{S_\varepsilon}\frac{\mathrm dS}{\varepsilon^2}
%     =\frac1{\varepsilon^2}4\pi\varepsilon^2=4\pi.
% \end{aligned}
% \qedhere
% \]
% \end{solve}
% 
% 
% \subsection{斯托克斯公式}
% \label{subsec:ma-16-5-3}
% 
% 斯托克斯 (Stokes) 公式是格林公式的直接推广, 它揭示了函数在空间曲面 $S$ 上的第二型曲面积分与相关函数在 $\partial S$ 上的第二型曲线积分之间的关系。为此我们先讨论一下空间曲面与其边界的定向问题.
% 
% 设 $S$ 是空间分片光滑的双侧曲面, 其边界由有限条分段光滑的空间闭曲线组成。给定 $S$ 的一侧, 设其由法向量 $\boldsymbol n$ 所确定。若一个人站立与该法向量保持一致并且沿 $\partial S$ 前进时 $S$ 在其左边, 则称他的前进方向为 $\partial S$ 关于该曲面确定的侧的正向.
% 
% 如图 16.5.8, $\Gamma_0$ 的逆时针方向及 $\Gamma_k(k=1,2,\cdots,K)$ 的顺时针方向构成了 $S$ 关于 $\boldsymbol n$ 确定的侧的正向。显然, 这样的定向与平面区域的边界定向是保持一致的.
% 
% \begin{theorem}[{斯托克斯公式}]\label{theorem:ma-16-5-3}
% 设 $S\subset\mathbb R^3$ 是光滑双侧曲面, $\partial S$ 由有限多条分段光滑曲线组成, 给定 $S$ 的一侧并取 $\partial S$ 关于该侧为正向的方向, 再设函数 $P(x,y,z)$、$Q(x,y,z)$ 和 $R(x,y,z)$ 在包含 $S$ 的某个区域内具有连续偏导数, 则
% \[
% \begin{aligned}
%   \int_{\partial S}P\mathrm dx+Q\mathrm dy+R\mathrm dz
%   &=\iint_S\left(\frac{\partial R}{\partial y}-\frac{\partial Q}{\partial z}\right)\mathrm dy\mathrm dz\\
%   &\quad+\left(\frac{\partial P}{\partial z}-\frac{\partial R}{\partial x}\right)\mathrm dz\mathrm dx
%     +\left(\frac{\partial Q}{\partial x}-\frac{\partial P}{\partial y}\right)\mathrm dx\mathrm dy.
% \end{aligned}
% \]
% \end{theorem}
% 
% 
% \begin{proof}
% 若 $\partial S$ 由有限条闭曲线所组成, 我们可以在 $S$ 上添加若干条曲线将 $S$ 分成有限块小曲面, 而每块小曲面的边界只有一条闭曲线。如同格林公式的证明, 若在小曲面上斯托克斯公式成立, 则 $S$ 上也成立。因此不妨设 $S$ 的边界是一条分段光滑的闭曲线。另外, 我们也只对几种简单的曲面进行证明, 对于一般的曲面, 可以由简单曲面逼近来证明, 但这个过程非常复杂, 在此就不作介绍了.
% 
% 具体地说, 我们先设 $S$ 由方程
% \begin{equation}
%   z=f(x,y),\quad (x,y)\in D
%   \label{eq:ma-16-5-1}
% \end{equation}
% 给出, 其中 $D\subset\mathbb R^2$ 是由一条分段光滑的约当曲线所围成的区域, $\partial D$ 的定向由 $S$ 在 $D$ 的投影所确定.
% 
% 一方面, 由第二型曲线积分的计算公式以及格林公式得
% \[
% \begin{aligned}
%   \int_{\partial S}P(x,y,z)\mathrm dx
%   &=\int_{\partial D}P(x,y,f(x,y))\mathrm dx\\
%   &=-\iint_D\left(\frac{\partial P}{\partial y}
%     +\frac{\partial P}{\partial z}\cdot\frac{\partial f}{\partial y}\right)\mathrm dx\mathrm dy.
% \end{aligned}
% \]
% 另一方面, 有
% \[
% \begin{aligned}
%   \iint_S\frac{\partial P}{\partial z}\mathrm dz\mathrm dx
%   -\frac{\partial P}{\partial y}\mathrm dx\mathrm dy
%   &=\iint_D\left[\frac{\partial P}{\partial z}\left(-\frac{\partial f}{\partial y}\right)
%     -\frac{\partial P}{\partial y}\right]\mathrm dx\mathrm dy\\
%   &=-\iint_D\left(\frac{\partial P}{\partial y}
%     +\frac{\partial P}{\partial z}\cdot\frac{\partial f}{\partial y}\right)\mathrm dx\mathrm dy.
% \end{aligned}
% \]
% 因此
% \[
%   \int_{\partial S}P\mathrm dx
%   =\iint_S\frac{\partial P}{\partial z}\mathrm dz\mathrm dx
%     -\frac{\partial P}{\partial y}\mathrm dx\mathrm dy.
% \]
% 同样的计算可以得出
% \[
%   \int_{\partial S}Q\mathrm dy
%   =\iint_S\left(-\frac{\partial Q}{\partial z}\right)\mathrm dy\mathrm dz
%     +\frac{\partial Q}{\partial x}\mathrm dx\mathrm dy.
% \]
% 
% 同理, 对于曲面 $S$ 由方程 $x=g(y,z)((y,z)\in D)$ 给出的情形, 其中 $S$ 与 $D$ 的光滑性假设如 \eqref{eq:ma-16-5-1}, 则有
% \[
%   \int_{\partial S}Q\mathrm dy+\int_{\partial S}R\mathrm dz
%   =\iint_S\left(\frac{\partial R}{\partial y}-\frac{\partial Q}{\partial z}\right)\mathrm dy\mathrm dz
%     +\frac{\partial Q}{\partial x}\mathrm dx\mathrm dy-\frac{\partial R}{\partial x}\mathrm dz\mathrm dx;
% \]
% 而对于曲面 $S$ 由方程 $y=h(z,x)((z,x)\in D)$ 给出的情形, 则有
% \[
%   \int_{\partial S}P\mathrm dx+\int_{\partial S}R\mathrm dz
%   =\iint_S\left(\frac{\partial P}{\partial z}-\frac{\partial R}{\partial x}\right)\mathrm dz\mathrm dx
%     +\frac{\partial R}{\partial y}\mathrm dy\mathrm dz-\frac{\partial P}{\partial y}\mathrm dx\mathrm dy.
% \]
% 因此, 若 $S$ 是同时能表示成上述三种形式的曲面或用光滑曲线可以将 $S$ 分成有限块这样的曲面, 则有
% \[
% \begin{aligned}
%   \int_{\partial S}P\mathrm dx+Q\mathrm dy+R\mathrm dz
%   &=\iint_S\left(\frac{\partial R}{\partial y}-\frac{\partial Q}{\partial z}\right)\mathrm dy\mathrm dz\\
%   &\quad+\left(\frac{\partial P}{\partial z}-\frac{\partial R}{\partial x}\right)\mathrm dz\mathrm dx
%     +\left(\frac{\partial Q}{\partial x}-\frac{\partial P}{\partial y}\right)\mathrm dx\mathrm dy.
% \end{aligned}
% \qedhere
% \]
% \end{proof}
% 
% 
% 在斯托克斯公式中, 曲面积分的被积函数似乎不太容易记忆, 我们可以用下述方式来记该公式:
% \[
%   \int_{\partial S}P\mathrm dx+Q\mathrm dy+R\mathrm dz
%   =\iint_S
%   \begin{vmatrix}
%     \cos\alpha&\cos\beta&\cos\gamma\\
%     \frac{\partial}{\partial x}&\frac{\partial}{\partial y}&\frac{\partial}{\partial z}\\
%     P&Q&R
%   \end{vmatrix}\mathrm dS
% \]
% \[
%   =\iint_S
%   \begin{vmatrix}
%     \mathrm dy\mathrm dz&\mathrm dz\mathrm dx&\mathrm dx\mathrm dy\\
%     \frac{\partial}{\partial x}&\frac{\partial}{\partial y}&\frac{\partial}{\partial z}\\
%     P&Q&R
%   \end{vmatrix},
% \]
% 其中 $(\cos\alpha,\cos\beta,\cos\gamma)$ 为与 $\partial S$ 的正向保持一致的法向量的方向余弦。另外, 由于此时 $\partial S$ 为封闭曲线, 我们也用 $\oint_{\partial S}P\mathrm dx+Q\mathrm dy+R\mathrm dz$ 来记 $\int_{\partial S}P\mathrm dx+Q\mathrm dy+R\mathrm dz$.
% 
% \begin{example}\label{example:ma-16-5-5}
% 计算第二型曲线积分
% \[
%   I=\oint_{\Gamma_h}(y^2-z^2)\mathrm dx+(z^2-x^2)\mathrm dy+(x^2-y^2)\mathrm dz,
% \]
% 其中 $\Gamma_h$ 是平面 $x+y+z=h(-1<h<1)$ 与球面 $x^2+y^2+z^2=1$ 的交线, 从 $z$ 轴看去取逆时针方向.
% \end{example}
% 
% 
% \begin{solve}
% 设平面 $x+y+z=h(-1<h<1)$ 被圆周 $\Gamma_h$ 所围的部分为 $S_h$, 则 $S_h$ 是半径为 $\sqrt{1-h^2}$ 的圆盘。取 $S_h$ 的法向量向上, 即选取法向量的方向余弦为
% \[
%   (\cos\alpha,\cos\beta,\cos\gamma)=\left(\frac1{\sqrt3},\frac1{\sqrt3},\frac1{\sqrt3}\right).
% \]
% 利用斯托克斯公式, 有
% \[
%   I=\iint_{S_h}
%   \begin{vmatrix}
%     \frac1{\sqrt3}&\frac1{\sqrt3}&\frac1{\sqrt3}\\
%     \frac{\partial}{\partial x}&\frac{\partial}{\partial y}&\frac{\partial}{\partial z}\\
%     y^2-z^2&z^2-x^2&x^2-y^2
%   \end{vmatrix}\mathrm dS
% \]
% \[
%   =-\frac4{\sqrt3}\iint_{S_h}(x+y+z)\mathrm dS
%   =-\frac4{\sqrt3}\iint_{S_h}h\mathrm dS
%   =-\frac{4\pi}{\sqrt3}h(1-h^2).
% \qedhere
% \]
% \end{solve}
% 
% 
% \section{微分形式简介}
% \label{sec:ma-16-6}
% 
% \subsection{微分形式}
% \label{subsec:ma-16-6-1}
% 
% 我们前面学过了许多不同的积分, 对一个多元函数来说, 笼统地说它的积分是十分不确切的事情; 就三元向量函数而言, 给了一个三元向量函数 $(P(x,y,z),Q(x,y,z),R(x,y,z))$, 对它就可以进行曲线积分与曲面积分。但如果我们写出 $P\mathrm dx+Q\mathrm dy+R\mathrm dz$ 或 $P\mathrm dy\mathrm dz+Q\mathrm dz\mathrm dx+R\mathrm dx\mathrm dy$, 就可以清楚知道前者可以在曲线上积分, 而后者可以在曲面上积分。而 $P\mathrm dx+Q\mathrm dy+R\mathrm dz$ 与 $P\mathrm dy\mathrm dz+Q\mathrm dz\mathrm dx+R\mathrm dx\mathrm dy$ 就是所谓的微分形式。当然, 微分形式的意义远不止这些, 它是现代数学中一种基本的工具, 在许多领域中有着广泛的应用。由于在后续课程中有专门课程来介绍它, 在这里我们只作一些简单介绍, 主要目的是使读者能对重积分换元公式及各种积分之间的联系有更深刻的理解.
% 
% 设 $f(\boldsymbol x)=f(x_1,x_2,\cdots,x_n)$ 是定义在 $\mathbb R^n(n\geqslant2)$ 中的一个可微函数, 则它的微分是
% \[
%   \mathrm df=\sum_{i=1}^n\frac{\partial f(\boldsymbol x)}{\partial x_i}\mathrm dx_i.
% \]
% 现在抛开微分的几何意义, 纯粹从数学的观点来看上式右边的表达式, 我们可以认为 $\sum_{i=1}^n\frac{\partial f(\boldsymbol x)}{\partial x_i}\mathrm dx_i$ 是 $\mathbb R^n$ 中点 $\boldsymbol x$ 处的一个关于 $(\mathrm dx_1,\mathrm dx_2,\cdots,\mathrm dx_n)$ 的线性组合。因此, 我们可以将 $(\mathrm dx_1,\mathrm dx_2,\cdots,\mathrm dx_n)$ 看做一组基。对给定 $\mathbb R^n$ 中的一个区域 $D$, $C^1(D)$ ($D$ 内连续可微函数的全体) 在这组基上的线性组合
% \[
%   \sum_{i=1}^n a_i(\boldsymbol x)\mathrm dx_i\quad
%   (a_i(\boldsymbol x)\in C^1(D),\ i=1,2,\cdots,n)
% \]
% 称为 $D$ 内的一个一次微分形式或 $1$-形式, 一般用 $\omega,\eta$ 等来记之。所有 $D$ 内的 $1$-形式的全体记为 $\Lambda^1(D)$。在后面的讨论中, 我们总是固定区域 $D$, 因此 $\Lambda^1(D)$ 也记为 $\Lambda^1$.
% 
% 在 $\Lambda^1$ 内, 我们可以自然地定义加法与数乘如下:
% 
% (1) 设 $\omega_j=\sum_{i=1}^n a_i^j(\boldsymbol x)\mathrm dx_i\in\Lambda^1(j=1,2)$, 定义
% \[
%   \omega_1+\omega_2=\sum_{i=1}^n(a_i^1(\boldsymbol x)+a_i^2(\boldsymbol x))\mathrm dx_i\in\Lambda^1.
% \]
% 
% (2) 设 $\omega=\sum_{i=1}^n a_i(\boldsymbol x)\mathrm dx_i\in\Lambda^1$, $\lambda\in\mathbb R$, 定义
% \[
%   \lambda\omega=\sum_{i=1}^n(\lambda a_i(\boldsymbol x))\mathrm dx_i\in\Lambda^1,
% \]
% 另外, 我们规定 $0=\sum_{i=1}^n0\mathrm dx_i$.
% 
% 读者不难验证上述两种运算满足一个线性空间所要求的各种运算定律, 从而 $\Lambda^1$ 构成一个线性空间.
% 
% 现对任意的 $k(1<k\leqslant n)$, 在 $(\mathrm dx_1,\mathrm dx_2,\cdots,\mathrm dx_n)$ 中取 $k$ 个元素, 将其形式地记成 $\mathrm dx_{i_1}\wedge\mathrm dx_{i_2}\wedge\cdots\wedge\mathrm dx_{i_k}$, 这里 $\wedge$ 暂且作为一个形式记号。然后规定
% \[
%   \mathrm dx_{i_1}\wedge\cdots\wedge\mathrm dx_{i_j}\wedge\mathrm dx_{i_{j+1}}\wedge\cdots\wedge\mathrm dx_{i_k}
%   =-\mathrm dx_{i_1}\wedge\cdots\wedge\mathrm dx_{i_{j+1}}\wedge\mathrm dx_{i_j}\wedge\cdots\wedge\mathrm dx_{i_k}.
% \]
% 在此规定下, 若 $i_j=i_l(1\leqslant j<l\leqslant k)$, 则 $\mathrm dx_{i_1}\wedge\mathrm dx_{i_2}\wedge\cdots\wedge\mathrm dx_{i_k}=0$。因此我们有 $C_n^k$ 个有序元
% \[
%   \mathrm dx_{i_1}\wedge\mathrm dx_{i_2}\wedge\cdots\wedge\mathrm dx_{i_k}
%   \quad (1\leqslant i_1<i_2<\cdots<i_k\leqslant n).
% \]
% 我们将以这 $C_n^k$ 个有序元作为基并且系数在 $C^1(D)$ 的线性空间记为 $\Lambda^k$。$\Lambda^k$ 中的元素称为 $k$ 次微分形式, 一般我们还是用 $\omega,\eta$ 等来记它们。从定义我们知道, 对于 $\forall\omega\in\Lambda^k$, 我们有下述表达式:
% \[
%   \omega=\sum_{1\leqslant i_1<i_2<\cdots<i_k\leqslant n}
%   a_{i_1i_2\cdots i_k}(\boldsymbol x)
%   \mathrm dx_{i_1}\wedge\mathrm dx_{i_2}\wedge\cdots\wedge\mathrm dx_{i_k},
% \]
% 其中 $a_{i_1i_2\cdots i_k}(\boldsymbol x)\in C^1(D)$。$\omega$ 的这种表示称为它的标准形式。特别地, 我们将 $C^1(D)$ 记为 $\Lambda^0$, 即 $\Lambda^0$ 是 $D$ 内连续可微函数的全体, 它的基是 $f(\boldsymbol x)=1$, 从而它是一维线性空间。另外, 我们记 $0$ 为任意次零微分形式.
% 
% 当 $k=n$ 时, $\Lambda^n$ 只有一个基 $\mathrm dx_1\wedge\mathrm dx_2\wedge\cdots\wedge\mathrm dx_n$, 从而 $\Lambda^n$ 也是一维线性空间.
% 
% 在 $\mathbb R^2$ 中可记 $(x_1,x_2)=(x,y)$。在 $\mathbb R^3$ 中, 我们记 $(x_1,x_2,x_3)=(x,y,z)$。例如, 在 $\mathbb R^3$ 中设 $\omega\in\Lambda^2$, 则 $\omega$ 的标准形式为
% \[
%   \omega=P\mathrm dy\wedge\mathrm dz+Q\mathrm dx\wedge\mathrm dz+R\mathrm dx\wedge\mathrm dy,
% \]
% 其中 $P,Q,R$ 是 $\mathbb R^3$ 中的连续可微函数.
% 
% \subsection{微分形式的外积}
% \label{subsec:ma-16-6-2}
% 
% 外积运算建立了不同次微分形式空间的联系。一个 $p$ 次微分形式 $\omega$ 与一个 $q$ 次微分形式 $\eta$ 的外积记为 $\omega\wedge\eta$, 它是一个 $p+q$ 次微分形式。因此只有当 $p+q\leqslant n$ 时, 外积才有意义。外积 $\omega\wedge\eta$ 的确切定义为: 设
% \[
%   \omega=\sum_{1\leqslant i_1<i_2<\cdots<i_p\leqslant n}
%   a_{i_1i_2\cdots i_p}(\boldsymbol x)
%   \mathrm dx_{i_1}\wedge\mathrm dx_{i_2}\wedge\cdots\wedge\mathrm dx_{i_p}\in\Lambda^p,
% \]
% \[
%   \eta=\sum_{1\leqslant j_1<j_2<\cdots<j_q\leqslant n}
%   a_{j_1j_2\cdots j_q}(\boldsymbol x)
%   \mathrm dx_{j_1}\wedge\mathrm dx_{j_2}\wedge\cdots\wedge\mathrm dx_{j_q}\in\Lambda^q,
% \]
% 则
% \[
%   \omega\wedge\eta=
%   \sum_{\substack{1\leqslant i_1<i_2<\cdots<i_p\leqslant n\\
%   1\leqslant j_1<j_2<\cdots<j_q\leqslant n}}
%   a_{i_1i_2\cdots i_p}(\boldsymbol x)a_{j_1j_2\cdots j_q}(\boldsymbol x)
%   \cdot\mathrm dx_{i_1}\wedge\mathrm dx_{i_2}\wedge\cdots\wedge\mathrm dx_{i_p}
%   \wedge\mathrm dx_{j_1}\wedge\mathrm dx_{j_2}\wedge\cdots\wedge\mathrm dx_{j_q}\in\Lambda^{p+q}.
% \]
% 特别地, 当 $f\in\Lambda^0,\omega\in\Lambda^k$ 时, $f\wedge\omega$ 定义为 $f\omega$.
% 
% 容易验证, 外积具有下述性质:
% 
% (1) 设 $f_1,f_2\in\Lambda^0,\omega_1,\omega_2\in\Lambda^p,\eta\in\Lambda^q$, 则有
% \[
%   (f_1\omega_1+f_2\omega_2)\wedge\eta
%   =f_1\omega_1\wedge\eta+f_2\omega_2\wedge\eta.
% \]
% 
% (2) 设 $\omega\in\Lambda^p,\eta\in\Lambda^q$, 则有 $\omega\wedge\eta=(-1)^{pq}\eta\wedge\omega$.
% 
% 上述性质 (2) 也称为外积的反对称性。此性质的证明可由一个微分形式中的基元素两两换位改变符号直接推出.
% 
% (3) 设 $\omega,\eta,\theta$ 是任意三个微分形式, 则 $\omega\wedge(\eta\wedge\theta)=(\omega\wedge\eta)\wedge\theta$.
% 
% 有了外积的概念, 微分形式 $\mathrm dx_1\wedge\mathrm dx_2$ 可以看成是 $\mathrm dx_1$ 与 $\mathrm dx_2$ 的外积, 类似地理解 $\mathrm dx_1\wedge\mathrm dx_2\wedge\cdots\wedge\mathrm dx_n$.
% 
% \begin{example}\label{example:ma-16-6-1}
% 设
% \[
%   \begin{pmatrix}x\\y\end{pmatrix}
%   =
%   \begin{pmatrix}x(u,v)\\y(u,v)\end{pmatrix}
% \]
% 是 $D\subset\mathbb R^2$ 到 $\Omega\subset\mathbb R^2$ 的 $C^1$ 同胚映射, 求 $\mathrm dx\wedge\mathrm dy$ 关于 $\mathrm du\wedge\mathrm dv$ 的表示.
% \end{example}
% 
% 
% \begin{solve}
% 由微分定义有
% \[
%   \mathrm dx=\frac{\partial x}{\partial u}\mathrm du+\frac{\partial x}{\partial v}\mathrm dv,
%   \quad
%   \mathrm dy=\frac{\partial y}{\partial u}\mathrm du+\frac{\partial y}{\partial v}\mathrm dv,
% \]
% 因此
% \[
%   \mathrm dx\wedge\mathrm dy
%   =\left(\frac{\partial x}{\partial u}\mathrm du+\frac{\partial x}{\partial v}\mathrm dv\right)
%   \wedge
%   \left(\frac{\partial y}{\partial u}\mathrm du+\frac{\partial y}{\partial v}\mathrm dv\right)
% \]
% \[
%   =\left(\frac{\partial x}{\partial u}\frac{\partial y}{\partial v}
%   -\frac{\partial x}{\partial v}\frac{\partial y}{\partial u}\right)\mathrm du\wedge\mathrm dv
%   =\frac{\partial(x,y)}{\partial(u,v)}\mathrm du\wedge\mathrm dv.
% \qedhere
% \]
% \end{solve}
% 
% 
% \begin{example}\label{example:ma-16-6-2}
% 设
% \[
%   \begin{pmatrix}x\\y\\z\end{pmatrix}
%   =
%   \begin{pmatrix}x(u,v,w)\\y(u,v,w)\\z(u,v,w)\end{pmatrix}
% \]
% 是 $D\subset\mathbb R^3$ 到 $\Omega\subset\mathbb R^3$ 的 $C^1$ 同胚映射, 求 $\mathrm dx\wedge\mathrm dy\wedge\mathrm dz$ 关于 $\mathrm du\wedge\mathrm dv\wedge\mathrm dw$ 的表示.
% \end{example}
% 
% 
% \begin{solve}
% 因为
% \[
%   \mathrm dx=\frac{\partial x}{\partial u}\mathrm du+\frac{\partial x}{\partial v}\mathrm dv+\frac{\partial x}{\partial w}\mathrm dw,
% \]
% \[
%   \mathrm dy=\frac{\partial y}{\partial u}\mathrm du+\frac{\partial y}{\partial v}\mathrm dv+\frac{\partial y}{\partial w}\mathrm dw,
% \]
% \[
%   \mathrm dz=\frac{\partial z}{\partial u}\mathrm du+\frac{\partial z}{\partial v}\mathrm dv+\frac{\partial z}{\partial w}\mathrm dw,
% \]
% 所以
% \[
%   \mathrm dx\wedge\mathrm dy
%   =\left(\frac{\partial x}{\partial u}\mathrm du+\frac{\partial x}{\partial v}\mathrm dv+\frac{\partial x}{\partial w}\mathrm dw\right)
%   \wedge
%   \left(\frac{\partial y}{\partial u}\mathrm du+\frac{\partial y}{\partial v}\mathrm dv+\frac{\partial y}{\partial w}\mathrm dw\right)
% \]
% \[
%   =\frac{\partial(x,y)}{\partial(v,w)}\mathrm dv\wedge\mathrm dw
%   +\frac{\partial(x,y)}{\partial(w,u)}\mathrm dw\wedge\mathrm du
%   +\frac{\partial(x,y)}{\partial(u,v)}\mathrm du\wedge\mathrm dv.
% \]
% \[
%   \mathrm dx\wedge\mathrm dy\wedge\mathrm dz
%   =\frac{\partial z}{\partial u}\cdot\frac{\partial(x,y)}{\partial(v,w)}
%   \mathrm dv\wedge\mathrm dw\wedge\mathrm du
%   +\frac{\partial z}{\partial v}\cdot\frac{\partial(x,y)}{\partial(w,u)}
%   \mathrm dw\wedge\mathrm du\wedge\mathrm dv
% \]
% \[
%   \quad+\frac{\partial z}{\partial w}\cdot\frac{\partial(x,y)}{\partial(u,v)}
%   \mathrm du\wedge\mathrm dv\wedge\mathrm dw
%   =\frac{\partial(x,y,z)}{\partial(u,v,w)}\mathrm du\wedge\mathrm dv\wedge\mathrm dw.
% \]
% 
% 下面我们以区域 $D\subset\mathbb R^2$ 到 $\Omega\subset\mathbb R^2$ 的 $C^1$ 同胚映射
% \[
%   \begin{pmatrix}x\\y\end{pmatrix}
%   =
%   \begin{pmatrix}x(u,v)\\y(u,v)\end{pmatrix}
% \]
% 来说明其雅可比行列式的符号与映射的保向性的联系。记 $\boldsymbol u=(u,v),\boldsymbol x=(x,y)$, 并将上述同胚映射记成 $\boldsymbol x=\boldsymbol f(\boldsymbol u)$.
% 
% 任取一点 $\boldsymbol u_0\in D$, 则由 $C^1$ 同胚映射的性质, 存在 $U(\boldsymbol u_0,\delta_0)\subset D$, 使得 $\boldsymbol f(U(\boldsymbol u_0,\delta_0))$ 是包含 $\boldsymbol f(\boldsymbol u_0)=\boldsymbol x_0$ 的一个区域。因此任取一条约当曲线 $\Gamma\subset U(\boldsymbol u_0,\delta_0)$, 则 $\gamma=\boldsymbol f(\Gamma)$ 是 $\Omega$ 内的一条约当曲线。如果 $\boldsymbol x=\boldsymbol f(\boldsymbol u)$ 将 $\Gamma$ 的正向映成 $\gamma$ 的正向 (即动点 $\boldsymbol u$ 沿 $\Gamma$ 逆时针运动时, 动点 $\boldsymbol f(\boldsymbol u)$ 沿 $\gamma$ 逆时针运动), 则称 $\boldsymbol f(\boldsymbol u)$ 在 $\boldsymbol u_0$ 处保向。若 $\boldsymbol f(\boldsymbol u)$ 在 $\boldsymbol u_0$ 处将 $\Gamma$ 的正向映成 $\gamma$ 的负向, 则称 $\boldsymbol f(\boldsymbol u)$ 在 $\boldsymbol u_0$ 处反定向。可以证明, 若 $\boldsymbol f(\boldsymbol u)$ 在 $\boldsymbol u_0$ 处保向, 则它在任意点 $\boldsymbol u\in D$ 都是保向的 (见本章习题).
% 
% 下面我们来研究保向映射的雅可比行列式的性质。任取 $\boldsymbol u_0=(u_0,v_0)\in D$, 任取过 $\boldsymbol u_0$ 的一条光滑曲线
% \[
%   \Gamma:\begin{cases}
%     u=u(t),\\
%     v=v(t),
%   \end{cases}\quad t\in(\alpha,\beta),
% \]
% 使得它满足 $u(t_0)=u_0,v(t_0)=v_0(t_0\in(\alpha,\beta))$。$\Gamma$ 在 $\boldsymbol u_0$ 处的切线斜率为
% \[
%   k_u=\frac{v'(t_0)}{u'(t_0)}.
% \]
% 曲线 $\gamma=\boldsymbol f(\Gamma)$ 的参数方程为
% \[
%   \begin{cases}
%     x=x(u(t),v(t)),\\
%     y=y(u(t),v(t)).
%   \end{cases}
% \]
% 它在 $\boldsymbol x_0=\boldsymbol f(\boldsymbol u_0)$ 处的切线的斜率
% \[
%   k_x=\frac{y'(t_0)}{x'(t_0)}
%   =\frac{y_u'(\boldsymbol x_0)u'(t_0)+y_v'(\boldsymbol x_0)v'(t_0)}
%   {x_u'(\boldsymbol x_0)u'(t_0)+x_v'(\boldsymbol x_0)v'(t_0)}
%   =\frac{y_u'(\boldsymbol x_0)+y_v'(\boldsymbol x_0)k_u}
%   {x_u'(\boldsymbol x_0)+x_v'(\boldsymbol x_0)k_u}.
% \]
% 所以
% \[
%   \frac{\mathrm dk_x}{\mathrm dk_u}
%   =
%   \frac{\left.\dfrac{\partial(x,y)}{\partial(u,v)}\right|_{\boldsymbol u_0}}
%   {[x_u'(\boldsymbol x_0)+x_v'(\boldsymbol x_0)k_u]^2}.
% \]
% 从上式可以看出, 当 $\left.\dfrac{\partial(x,y)}{\partial(u,v)}\right|_{\boldsymbol u_0}>0$ 时, $k_x$ 是 $k_u$ 的单调上升函数。由导数的几何意义知, 此时 $\boldsymbol f(\boldsymbol u)$ 在 $\boldsymbol u_0$ 处保向; 当 $\left.\dfrac{\partial(x,y)}{\partial(u,v)}\right|_{\boldsymbol u_0}<0$ 时, $\boldsymbol f(\boldsymbol u)$ 在 $\boldsymbol u_0$ 处反定向。我们已经知道, 一个变换在一点保向则处处保向, 因此 $\dfrac{\partial(x,y)}{\partial(u,v)}$ 在 $D$ 内必定不变号.
% 
% 现在回到微分形式的外积。设
% \[
%   \begin{pmatrix}x\\y\end{pmatrix}
%   =
%   \begin{pmatrix}x(u,v)\\y(u,v)\end{pmatrix}
% \]
% 是区域 $D\subset\mathbb R^2$ 到可求面积的有界闭区域 $\Omega\subset\mathbb R^2$ 的 $C^1$ 同胚映射。当它保向时, 有 $\dfrac{\partial(x,y)}{\partial(u,v)}\geqslant0$。这时我们有
% \[
%   \mathrm dx\wedge\mathrm dy
%   =\frac{\partial(x,y)}{\partial(u,v)}\mathrm du\wedge\mathrm dv.
% \]
% 当变换
% \[
%   \begin{pmatrix}x\\y\end{pmatrix}
%   =
%   \begin{pmatrix}x(u,v)\\y(u,v)\end{pmatrix}
% \]
% 不保向时, 则 $\dfrac{\partial(x,y)}{\partial(u,v)}\leqslant0$。现在我们把 $Ouv$ 平面上的 $u$ 轴改为 $v$ 轴, 而把 $v$ 轴改成 $u$ 轴。即令 $u'=v,v'=u$, 则上述变换雅可比行列式为
% \[
%   \frac{\partial(x,y)}{\partial(u',v')}
%   =
%   \begin{vmatrix}
%     \frac{\partial x}{\partial u'}&\frac{\partial x}{\partial v'}\\
%     \frac{\partial y}{\partial u'}&\frac{\partial y}{\partial v'}
%   \end{vmatrix}
%   =
%   \begin{vmatrix}
%     \frac{\partial x}{\partial v}&\frac{\partial x}{\partial u}\\
%     \frac{\partial y}{\partial v}&\frac{\partial y}{\partial u}
%   \end{vmatrix}
%   =-\frac{\partial(x,y)}{\partial(u,v)}\geqslant0,
% \]
% 从而
% \[
%   \mathrm dx\wedge\mathrm dy
%   =\frac{\partial(x,y)}{\partial(u',v')}\mathrm du'\wedge\mathrm dv'
%   =\frac{\partial(x,y)}{\partial(u,v)}\mathrm du\wedge\mathrm dv.
% \]
% 因此, 若 $f(x,y)$ 是 $\Omega$ 上的可积函数, 当映射
% \[
%   \begin{pmatrix}x\\y\end{pmatrix}
%   =
%   \begin{pmatrix}x(u,v)\\y(u,v)\end{pmatrix}
% \]
% 保向时, 有
% \[
%   \iint_\Omega f(x,y)\mathrm dx\wedge\mathrm dy
%   =
%   \iint_D f(x(u,v),y(u,v))\frac{\partial(x,y)}{\partial(u,v)}\mathrm du\wedge\mathrm dv.
% \]
% 而当
% \[
%   \begin{pmatrix}x\\y\end{pmatrix}
%   =
%   \begin{pmatrix}x(u,v)\\y(u,v)\end{pmatrix}
% \]
% 不保向时, 仍然成立
% \[
%   \iint_\Omega f(x,y)\mathrm dx\wedge\mathrm dy
%   =
%   \iint_D f(x(u,v),y(u,v))\frac{\partial(x,y)}{\partial(u,v)}\mathrm du\wedge\mathrm dv.
% \]
% 同理, 对 $n$ 重积分也有类似的结果。这说明了引入微分形式的外积后, 我们在重积分的变量替换中不必考虑变换的雅可比行列式的正负号。当 $n\geqslant3$ 时, 一个映射的保向性的刻画比 $\mathbb R^2$ 的情形会更加复杂, 在这里我们就不作介绍了.
% \end{solve}
% 
% 
% \subsection{外微分}
% \label{subsec:ma-16-6-3}
% 
% 微分形式的外微分是一个重要的运算, 利用它可以将许多我们学过的重要定理统一起来。下面我们先对 $\mathbb R^n(n\geqslant2)$ 中区域 $D$ 内的微分形式引进外微分这一概念.
% 
% \begin{definition}\label{definition:ma-16-6-1}
% 设 $\omega\in\Lambda^k(1\leqslant k\leqslant n)$, 其形式为
% \[
%   \omega=
%   \sum_{1\leqslant i_1<i_2<\cdots<i_k\leqslant n}
%   a_{i_1i_2\cdots i_k}(\boldsymbol x)
%   \mathrm dx_{i_1}\wedge\mathrm dx_{i_2}\wedge\cdots\wedge\mathrm dx_{i_k},
% \]
% 称 $k+1$ 次微分形式
% \[
%   \sum_{1\leqslant i_1<i_2<\cdots<i_k\leqslant n}
%   \left(\sum_{j=1}^n\frac{\partial a_{i_1i_2\cdots i_k}(\boldsymbol x)}{\partial x_j}\mathrm dx_j\right)
%   \wedge\mathrm dx_{i_1}\wedge\mathrm dx_{i_2}\wedge\cdots\wedge\mathrm dx_{i_k}
% \]
% 为 $\omega$ 的外微分, 记为 $\mathrm d\omega$.
% \end{definition}
% 
% 显然, 当 $\omega\in\Lambda^n$ 时, $\mathrm d\omega=0$。当 $\omega\in\Lambda^k$, 且 $\omega$ 在 $\Lambda^k$ 的基元素前面的系数函数是 $D$ 内的 $C^1$ 函数时, 称 $\omega$ 是 $C^1$ 微分形式, 这时 $\mathrm d\omega$ 的系数仅在 $D$ 内连续。若要讨论 $\mathrm d(\mathrm d\omega)\triangleq\mathrm d^2\omega$ (称之为二次外微分, 即外微分的外微分), 则要假定 $\omega$ 是 $C^2$ 微分形式, 即 $\omega$ 中在基前面的系数函数在 $D$ 内具有各个二阶连续偏导数.
% 
% 由定义可推知外微分具有下述性质:
% 
% (1) 设 $\omega_1,\omega_2\in\Lambda^k$, 则 $\mathrm d(\omega_1+\omega_2)=\mathrm d\omega_1+\mathrm d\omega_2$.
% 
% (2) 设 $\omega\in\Lambda^k,\eta\in\Lambda^l$, 则 $\mathrm d(\omega\wedge\eta)=\mathrm d\omega\wedge\eta+(-1)^k\omega\wedge\mathrm d\eta$.
% 
% (3) 当 $\omega\in\Lambda^k$, 且 $\omega$ 是 $C^2$ 微分形式时, $\mathrm d(\mathrm d\omega)=\mathrm d^2\omega=0$.
% 
% 这些性质的证明留给读者.
% 
% 下面我们来求一些微分形式的外微分.
% 
% (1) 设 $\omega=P(x,y)\mathrm dx+Q(x,y)\mathrm dy$, 则
% \[
%   \mathrm d\omega=\left(\frac{\partial Q}{\partial x}-\frac{\partial P}{\partial y}\right)\mathrm dx\wedge\mathrm dy.
% \]
% 
% (2) 设 $\eta=P(x,y,z)\mathrm dx+Q(x,y,z)\mathrm dy+R(x,y,z)\mathrm dz$, 则
% \[
%   \mathrm d\eta=
%   \left(\frac{\partial R}{\partial y}-\frac{\partial Q}{\partial z}\right)\mathrm dy\wedge\mathrm dz
%   +
%   \left(\frac{\partial P}{\partial z}-\frac{\partial R}{\partial x}\right)\mathrm dz\wedge\mathrm dx
%   +
%   \left(\frac{\partial Q}{\partial x}-\frac{\partial P}{\partial y}\right)\mathrm dx\wedge\mathrm dy.
% \]
% 
% (3) 设 $\theta=P(x,y,z)\mathrm dy\wedge\mathrm dz+Q(x,y,z)\mathrm dz\wedge\mathrm dx+R(x,y,z)\mathrm dx\wedge\mathrm dy$, 则
% \[
%   \mathrm d\theta=
%   \left(\frac{\partial P}{\partial x}+\frac{\partial Q}{\partial y}+\frac{\partial R}{\partial z}\right)
%   \mathrm dx\wedge\mathrm dy\wedge\mathrm dz.
% \]
% 
% 本节开始时我们曾指出, 给了 $\mathbb R^3$ 中的一个微分形式, 我们很清楚它在 $\mathbb R^3$ 中什么样的子集上可积分。因此, 对于一个几何体 $D$ (闭区域、曲面、曲线等), 由上述外微分的表达式, 格林公式、高斯公式以及斯托克斯公式可以统一地写成
% \[
%   \int_D\mathrm d\omega=\int_{\partial D}\omega,
% \]
% 其中 $\omega$ 是 $\overline D$ 上的一个 $C^1$ 微分形式, $\partial D$ 取正向.
% 
% 若在牛顿-莱布尼茨公式 $\int_a^b f'(x)\mathrm dx=f(x)\big|_a^b$ 中将 $f(x)\big|_a^b=f(b)-f(a)$ 看成是 $\omega=f(x)$ 在 $a,b$ 两点的积分, 上述公式也可写成 $\int_D\mathrm d\omega=\int_{\partial D}\omega$ 的形式, 其中 $D=[a,b],\partial D=\{a,b\}$.
% 
% 公式 $\int_D\mathrm d\omega=\int_{\partial D}\omega$ 也称为斯托克斯公式。尽管如此, 任何积分公式最后的计算以及它们的证明都必须要用到牛顿-莱布尼茨公式, 所以人们还是称牛顿-莱布尼茨公式为微积分基本定理.
% 
% 另外, 当 $n\geqslant4$ 时, 微分形式具有更多的形式, 因此在 $\mathbb R^n$ 中具有更多可以讨论积分问题的子集。关于这些问题在微分流形课程中将做系统的介绍, 在这里我们就不展开这方面的内容了.
% 
% \section{曲线积分与路径的无关性}
% \label{sec:ma-16-7}
% 
% 设 $D\subset\mathbb R^3$ 是一个区域, $u(x,y,z)$ 是 $D$ 内的 $C^1$ 函数, 则 $\mathrm du=u_x'\mathrm dx+u_y'\mathrm dy+u_z'\mathrm dz$ 是 $u(x,y,z)$ 在 $D$ 内的全微分。给定 $A_0\in D$, 则对于 $\forall A\in D$, $A$ 与 $A_0$ 之间均可用 $D$ 内无穷多条分段光滑的曲线连接。如果我们写出其中任意一条曲线 $\Gamma$ (以 $A_0$ 为起点, 以 $A$ 为终点) 的参数方程, 然后计算可知
% \[
%   \int_\Gamma u_x'\mathrm dx+u_y'\mathrm dy+u_z'\mathrm dz=u(A)-u(A_0).
% \]
% 换句话说, $u_x'\mathrm dx+u_y'\mathrm dy+u_z'\mathrm dz$ 在 $D$ 内沿任何光滑曲线上的积分, 只依赖于该曲线的起点与终点, 而与曲线的选取无关。对于这种现象, 我们引入以下概念.
% 
% \begin{definition}\label{definition:ma-16-7-1}
% 设 $D\subset\mathbb R^3$ 为一个区域, 函数 $P(x,y,z)$、$Q(x,y,z)$ 和 $R(x,y,z)$ 在 $D$ 内连续。如果对于 $\forall A_0,A\in D$ 及 $D$ 内任意一条以 $A_0$ 为起点, $A$ 为终点的分段光滑的曲线 $\Gamma$, 曲线积分
% \[
%   \int_\Gamma P\mathrm dx+Q\mathrm dy+R\mathrm dz
% \]
% 的值只与 $A_0,A$ 有关, 而与曲线 $\Gamma$ 的选取无关, 则称曲线积分 $\int_\Gamma P\mathrm dx+Q\mathrm dy+R\mathrm dz$ 在 $D$ 内与路径无关.
% \end{definition}
% 
% 从第二型曲线积分的物理背景可以看出, 该曲线积分与路径无关是指 $D$ 内力场 $\boldsymbol F=(P,Q,R)$ 将单位质量的粒子从点 $A_0$ 移动到点 $A$ 所做的功不依赖于粒子所走的路线, 而只依赖粒子运动的起点与终点。从直观上看, 重力做功应具有该性质, 因为重力对质点所做的功只依赖于空间这两点之间的垂直距离.
% 
% 现在, 我们来研究对于区域 $D$ 内给定的三个函数 $P,Q,R$, 何时曲线积分 $\int_\Gamma P\mathrm dx+Q\mathrm dy+R\mathrm dz$ 在 $D$ 内与路径无关。设 $P(x,y,z)$、$Q(x,y,z)$ 和 $R(x,y,z)$ 是空间 $\mathbb R^3$ 中区域 $D$ 内的三个函数, 若存在 $D$ 内的可微函数 $u=u(x,y,z)$, 使得
% \[
%   \mathrm du=P\mathrm dx+Q\mathrm dy+R\mathrm dz,
% \]
% 则称 $u$ 为 $P\mathrm dx+Q\mathrm dy+R\mathrm dz$ 的一个原函数。因此, $P\mathrm dx+Q\mathrm dy+R\mathrm dz$ 有原函数 $u$ 的等价定义是: 存在可微函数 $u$, 使得 $P\mathrm dx+Q\mathrm dy+R\mathrm dz$ 是 $u$ 的全微分.
% 
% 在本节中, 我们将主要证明以下两个定理.
% 
% \begin{theorem}\label{theorem:ma-16-7-1}
% 设 $D$ 是 $\mathbb R^3$ 中的区域, 函数 $P(x,y,z)$、$Q(x,y,z)$ 和 $R(x,y,z)$ 在 $D$ 内连续, 则以下三个命题等价:
% 
% (1) 曲线积分 $\int_\Gamma P\mathrm dx+Q\mathrm dy+R\mathrm dz$ 在 $D$ 内与路径无关;
% 
% (2) 对 $D$ 内任何分段光滑的约当曲线 $\Gamma$, 有
% \[
%   \oint_\Gamma P\mathrm dx+Q\mathrm dy+R\mathrm dz=0;
% \]
% 
% (3) $P\mathrm dx+Q\mathrm dy+R\mathrm dz$ 在 $D$ 内存在原函数.
% \end{theorem}
% 
% 
% \begin{proof}
% 先证 $(1)\Rightarrow(2)$。设 $\Gamma$ 是 $D$ 内任一条分段光滑的约当曲线。在 $\Gamma$ 上取定两点 $A_0,A$, 则 $\Gamma$ 构成了从点 $A_0$ 到点 $A$ 的两条曲线 $\Gamma_1,\Gamma_2$, 其中 $\Gamma_1$ 取曲线 $\Gamma$ 的正向, $\Gamma_2$ 则取 $\Gamma$ 的负向 (如图 16.7.1)。由 (1) 我们有
% \[
%   0=\int_{\Gamma_1}P\mathrm dx+Q\mathrm dy+R\mathrm dz
%   -\int_{\Gamma_2}P\mathrm dx+Q\mathrm dy+R\mathrm dz
% \]
% \[
%   =\int_{\Gamma_1+\Gamma_2^-}P\mathrm dx+Q\mathrm dy+R\mathrm dz
%   =\int_\Gamma P\mathrm dx+Q\mathrm dy+R\mathrm dz,
% \]
% 因此 (2) 成立.
% 
% 再证 $(2)\Rightarrow(1)$。对于 $\forall A_0,A\in D$, 任取两条以 $A_0$ 为起点, $A$ 为终点的光滑曲线 $\Gamma_1,\Gamma_2$。若 $\Gamma_1\cup\Gamma_2^-$ 组成一条约当曲线, 令 $\Gamma=\Gamma_1\cup\Gamma_2^-$, 则我们有
% \[
%   \int_{\Gamma_1}P\mathrm dx+Q\mathrm dy+R\mathrm dz
%   -\int_{\Gamma_2}P\mathrm dx+Q\mathrm dy+R\mathrm dz
% \]
% \[
%   =\int_{\Gamma_1\cup\Gamma_2^-}P\mathrm dx+Q\mathrm dy+R\mathrm dz
%   =\int_\Gamma P\mathrm dx+Q\mathrm dy+R\mathrm dz=0.
% \]
% 当 $\Gamma_1\cup\Gamma_2^-$ 不是约当曲线时, 在 $D$ 内选取一条从点 $A_0$ 到点 $A$ 的光滑曲线 $\Gamma_3$, 使得 $\Gamma_1\cup\Gamma_3^-,\Gamma_2\cup\Gamma_3^-$ 均是约当曲线, 如图 16.7.2。由上面所证有
% \[
%   \int_{\Gamma_1}P\mathrm dx+Q\mathrm dy+R\mathrm dz
%   =\int_{\Gamma_3}P\mathrm dx+Q\mathrm dy+R\mathrm dz
%   =\int_{\Gamma_2}P\mathrm dx+Q\mathrm dy+R\mathrm dz,
% \]
% 因此 (1) 成立.
% 
% 现证 $(1)\Rightarrow(3)$。取定 $A_0=(x_0,y_0,z_0)\in D$, 任取 $A=(x,y,z)\in D$, 则可记
% \[
%   u(x,y,z)=\int_{\Gamma_0}P(\xi,\eta,\zeta)\mathrm d\xi
%   +Q(\xi,\eta,\zeta)\mathrm d\eta
%   +R(\xi,\eta,\zeta)\mathrm d\zeta,
% \]
% 其中 $\Gamma_0$ 为 $D$ 中一条从点 $A_0$ 到点 $A$ 的固定的光滑曲线.
% 
% 由于 $A\in D$, 因此 $A$ 是 $D$ 的内点。对于充分小的 $\Delta x$, 有 $A_{\Delta x}=(x+\Delta x,y,z)\in D$, 于是线段 $\overline{AA_{\Delta x}}\subset D$。记 $\Gamma_{\Delta x}=\Gamma_0\cup\overline{AA_{\Delta x}}$, 如图 16.7.3, 则
% \[
%   u(x+\Delta x,y,z)-u(x,y,z)
%   =
%   \int_{\Gamma_{\Delta x}}P\mathrm d\xi+Q\mathrm d\eta+R\mathrm d\zeta
%   -\int_{\Gamma_0}P\mathrm d\xi+Q\mathrm d\eta+R\mathrm d\zeta
% \]
% \[
%   =\int_{\overline{AA_{\Delta x}}}P\mathrm d\xi+Q\mathrm d\eta+R\mathrm d\zeta
%   =\int_x^{x+\Delta x}P(\xi,\eta,\zeta)\mathrm d\xi
%   =P(x+\theta\Delta x,y,z)\Delta x\quad(0<\theta<1).
% \]
% 因此
% \[
%   \frac{\partial u}{\partial x}
%   =\lim_{\Delta x\to0}\frac{u(x+\Delta x,y,z)-u(x,y,z)}{\Delta x}
%   =\lim_{\Delta x\to0}P(x+\theta\Delta x,y,z)=P(x,y,z).
% \]
% 同理可证 $\dfrac{\partial u}{\partial y}=Q$ 及 $\dfrac{\partial u}{\partial z}=R$.
% 
% 由 $P,Q,R$ 的连续性知, $\mathrm du=P\mathrm dx+Q\mathrm dy+R\mathrm dz$ 在 $D$ 内成立, 因此 $u(x,y,z)$ 是 $P\mathrm dx+Q\mathrm dy+R\mathrm dz$ 的一个原函数.
% 
% 最后证 $(3)\Rightarrow(1)$。设 $u(x,y,z)$ 满足 $\mathrm du=P\mathrm dx+Q\mathrm dy+R\mathrm dz$, 再设 $A_0=(x_0,y_0,z_0),A_1=(x_1,y_1,z_1)\in D$ 是任意两点, 任取一条从点 $A_0$ 到点 $A$ 的分段光滑曲线
% \[
%   \Gamma:\begin{cases}
%     x=x(t),\\
%     y=y(t),\\
%     z=z(t),
%   \end{cases}\quad t\in[\alpha,\beta],
% \]
% 且 $(x(\alpha),y(\alpha),z(\alpha))=(x_0,y_0,z_0),(x(\beta),y(\beta),z(\beta))=(x_1,y_1,z_1)$, 则
% \[
%   \int_\Gamma P\mathrm dx+Q\mathrm dy+R\mathrm dz
% \]
% \[
%   =\int_\alpha^\beta
%   [P(x(t),y(t),z(t))x'(t)+Q(x(t),y(t),z(t))y'(t)
%   +R(x(t),y(t),z(t))z'(t)]\mathrm dt
% \]
% \[
%   =\int_\alpha^\beta\frac{\mathrm du(x(t),y(t),z(t))}{\mathrm dt}\mathrm dt
%   =u(x(\beta),y(\beta),z(\beta))-u(x(\alpha),y(\alpha),z(\alpha))
% \]
% \[
%   =u(x_1,y_1,z_1)-u(x_0,y_0,z_0).
% \]
% 这说明了曲线积分 $\int_\Gamma P\mathrm dx+Q\mathrm dy+R\mathrm dz$ 在 $D$ 内与路径无关.
% \end{proof}
% 
% 
% 当曲线积分 $\int_\Gamma P\mathrm dx+Q\mathrm dy+R\mathrm dz$ 在 $D$ 内与路径无关时, 对于 $A_0=(x_0,y_0,z_0),A_1=(x_1,y_1,z_1)\in D$, 我们用
% \[
%   \int_{(x_0,y_0,z_0)}^{(x_1,y_1,z_1)}P\mathrm dx+Q\mathrm dy+R\mathrm dz
% \]
% 来表示 $D$ 内从点 $A_0$ 到点 $A_1$ 的任一条路径的曲线积分.
% 
% 为了得到进一步的判定曲线积分与路径无关的准则, 我们需要对空间区域作一些限制。设 $D\subset\mathbb R^3$ 是一个区域, 如果 $D$ 内任何一条约当曲线 $\Gamma$ 总能在 $D$ 内连续地缩成 $D$ 内的一点, 则称 $D$ 是一个单连通区域。因此对于 $\mathbb R^3$ 中的单连通区域 $D$, $D$ 内部可以有“洞”, 如同心球面围成的区域就是单连通的。而对于形如轮胎的环面所围区域就不是单连通的了。空间的单连通区域 $D$ 的一个特点是: 对于 $D$ 中任何一条分段光滑的约当曲线 $\Gamma$, 我们总可以在 $D$ 内找一个分片光滑的双侧曲面 $S$, 使得 $\Gamma=\partial S$。上述事实的证明超出了数学分析课程的范围, 在这里我们就不介绍了。但是在下面定理的证明中我们将不加证明地利用该事实.
% 
% \begin{theorem}\label{theorem:ma-16-7-2}
% 设 $D\subset\mathbb R^3$ 是单连通区域, 函数 $P(x,y,z)$、$Q(x,y,z)$ 和 $R(x,y,z)$ 在 $D$ 内具有连续偏导数, 则曲线积分 $\int_\Gamma P\mathrm dx+Q\mathrm dy+R\mathrm dz$ 在 $D$ 内与路径无关的充分必要条件是: 在 $D$ 内有
% \[
%   \frac{\partial R}{\partial y}\equiv\frac{\partial Q}{\partial z},
%   \quad
%   \frac{\partial P}{\partial z}\equiv\frac{\partial R}{\partial x},
%   \quad
%   \frac{\partial Q}{\partial x}\equiv\frac{\partial P}{\partial y},
% \]
% 即向量
% \[
%   \begin{vmatrix}
%     \boldsymbol i&\boldsymbol j&\boldsymbol k\\
%     \frac{\partial}{\partial x}&\frac{\partial}{\partial y}&\frac{\partial}{\partial z}\\
%     P&Q&R
%   \end{vmatrix}
%   \equiv0.
% \]
% \end{theorem}
% 
% 
% \begin{proof}
% 必要性\quad 设 $\int_\Gamma P\mathrm dx+Q\mathrm dy+R\mathrm dz$ 在 $D$ 内与路径无关, 则由定理~\ref{theorem:ma-16-7-1} 知, 存在定义在 $D$ 内的可微函数 $u(x,y,z)$, 使得
% \[
%   P\mathrm dx+Q\mathrm dy+R\mathrm dz=u_x'\mathrm dx+u_y'\mathrm dy+u_z'\mathrm dz.
% \]
% 因此, 在 $D$ 内恒有
% \[
%   \frac{\partial R}{\partial y}
%   =\frac{\partial^2u}{\partial y\partial z}
%   =\frac{\partial^2u}{\partial z\partial y}
%   =\frac{\partial Q}{\partial z},
%   \quad
%   \frac{\partial P}{\partial z}
%   =\frac{\partial^2u}{\partial z\partial x}
%   =\frac{\partial^2u}{\partial x\partial z}
%   =\frac{\partial R}{\partial x},
% \]
% \[
%   \frac{\partial Q}{\partial x}
%   =\frac{\partial^2u}{\partial x\partial y}
%   =\frac{\partial^2u}{\partial y\partial x}
%   =\frac{\partial P}{\partial y},
% \]
% 即必要性成立.
% 
% 充分性\quad 任取 $D$ 内一条分段光滑的约当曲线 $\Gamma$, 作 $D$ 内一个分片光滑的双侧曲面 $S$, 使得 $\partial S=\Gamma$, 且选取 $S$ 的侧使得 $\Gamma$ 定向为正向。由斯托克斯公式我们有
% \[
%   \oint_\Gamma P\mathrm dx+Q\mathrm dy+R\mathrm dz
%   =
%   \iint_S
%   \begin{vmatrix}
%     \mathrm dy\mathrm dz&\mathrm dz\mathrm dx&\mathrm dx\mathrm dy\\
%     \frac{\partial}{\partial x}&\frac{\partial}{\partial y}&\frac{\partial}{\partial z}\\
%     P&Q&R
%   \end{vmatrix}
%   =0.
% \]
% 由定理~\ref{theorem:ma-16-7-1} 知, $\int_\Gamma P\mathrm dx+Q\mathrm dy+R\mathrm dz$ 在 $D$ 内与路径无关.
% \end{proof}
% 
% 
% 由于我们在很多时候要讨论平面区域内的曲线积分与路径的关系, 我们不加证明地给出下面定理.
% 
% \begin{theorem}\label{theorem:ma-16-7-3}
% 设 $D\subset\mathbb R^2$ 是一个区域, 函数 $P(x,y),Q(x,y)$ 在 $D$ 内连续, 则以下断言等价:
% 
% (1) 曲线积分 $\int_\Gamma P\mathrm dx+Q\mathrm dy$ 在 $D$ 内与路径无关;
% 
% (2) 对 $D$ 内任何分段光滑的约当曲线 $\Gamma$, 有 $\oint_\Gamma P\mathrm dx+Q\mathrm dy=0$;
% 
% (3) 在 $D$ 内存在可微函数 $u(x,y)$, 使得 $\mathrm du=P\mathrm dx+Q\mathrm dy$.
% \end{theorem}
% 
% 当 $D$ 是单连通区域且 $P,Q$ 具有连续偏导数时, 则下述断言与上述三个断言等价:
% \[
%   (4)\quad \text{在 }D\text{ 内成立 }\frac{\partial Q}{\partial x}\equiv\frac{\partial P}{\partial y}.
% \]
% 此定理的证明与定理~\ref{theorem:ma-16-7-1} 和~\ref{theorem:ma-16-7-2} 证明方法相同, 不过在证明 (4) 时, 要用到格林公式.
% 
% \begin{example}\label{example:ma-16-7-1}
% 记 $E=\{(x,y):x\in[0,1],y=0\}$, 证明曲线积分
% \[
%   \int_\Gamma\frac{x\mathrm dy-y\mathrm dx}{x^2+y^2}
%   -\frac{(x-1)\mathrm dy-y\mathrm dx}{(x-1)^2+y^2}
% \]
% 在 $D=\mathbb R^2\setminus E$ 内与路径无关.
% \end{example}
% 
% 
% \begin{proof}
% 设 $\Gamma_0\subset D$ 是任一条分段光滑的约当曲线, 记
% \[
%   I=\int_{\Gamma_0}\frac{x\mathrm dy-y\mathrm dx}{x^2+y^2}
%   -\frac{(x-1)\mathrm dy-y\mathrm dx}{(x-1)^2+y^2},
% \]
% 则
% \[
%   I=\int_{\Gamma_0}\frac{x\mathrm dy-y\mathrm dx}{x^2+y^2}
%   -\int_{\Gamma_0}\frac{(x-1)\mathrm dy-y\mathrm dx}{(x-1)^2+y^2}.
% \]
% 由于
% \[
%   \frac{\partial}{\partial x}\left(\frac{x}{x^2+y^2}\right)
%   =\frac{y^2-x^2}{(x^2+y^2)^2}
%   =\frac{\partial}{\partial y}\left(\frac{-y}{x^2+y^2}\right)
% \]
% 及
% \[
%   \frac{\partial}{\partial x}\left(\frac{x-1}{(x-1)^2+y^2}\right)
%   =\frac{y^2-(x-1)^2}{[(x-1)^2+y^2]^2}
%   =\frac{\partial}{\partial y}\left(\frac{-y}{(x-1)^2+y^2}\right),
% \]
% 因此, 对任意的光滑约当曲线 $\Gamma_0\subset D$, 设 $\Gamma_0$ 所围区域为 $D_0$, 若 $D_0\cap E=\varnothing$, 则由定理~\ref{theorem:ma-16-7-3} 知 $I=0$。当 $E\subset D_0$ 时, 点 $(0,0)$ 及 $(1,0)$ 均在 $D_0$ 内, 从而由例~\ref{example:ma-16-5-2} 知
% \[
%   \int_{\Gamma_0}\frac{x\mathrm dy-y\mathrm dx}{x^2+y^2}=2\pi,
%   \quad
%   \int_{\Gamma_0}\frac{(x-1)\mathrm dy-y\mathrm dx}{(x-1)^2+y^2}=2\pi.
% \]
% 因此同样有 $I=0$。由定理~\ref{theorem:ma-16-7-3} 知, 曲线积分
% \[
%   \int_\Gamma\frac{x\mathrm dy-y\mathrm dx}{x^2+y^2}
%   -\frac{(x-1)\mathrm dy-y\mathrm dx}{(x-1)^2+y^2}
% \]
% 在 $D$ 内与路径无关.
% \end{proof}
% 
% 
% \begin{remark}\label{remark:ma-16-source-73}
% 显然, 若在 $\mathbb R^2\setminus\{(0,0),(0,1)\}$ 内考虑曲线积分
% \[
%   \int_\Gamma\frac{x\mathrm dy-y\mathrm dx}{x^2+y^2}
%   -\frac{(x-1)\mathrm dy-y\mathrm dx}{(x-1)^2+y^2},
% \]
% 则它在 $D$ 内与路径有关。另外, 我们容易看出, 在 $\mathbb R^2\setminus\{(-\infty,0]\cup[1,+\infty)\}$ 内该曲线积分与路径无关。以后从复变函数课程中, 我们还将讨论上述的积分问题, 在那里我们可以更加清楚地知道这些曲线积分与路径无关的本质.
% \end{remark}
% 
% 
% \begin{example}\label{example:ma-16-7-2}
% 设 $f(r)$ 是定义在 $(0,+\infty)$ 内的一元连续函数, 其中 $r=\sqrt{x^2+y^2+z^2}$。证明: 当区域 $D\subset\mathbb R^3$ 不含原点时, 曲线积分
% \[
%   \int_\Gamma f(r)(x\mathrm dx+y\mathrm dy+z\mathrm dz)
% \]
% 在 $D$ 内与路径无关.
% \end{example}
% 
% 
% \begin{proof}
% 由于 $rf(r)$ 连续, 从而它存在原函数 $u(r)$。取 $u(\sqrt{x^2+y^2+z^2})$, 则当 $D$ 不含原点时, 有
% \[
%   \mathrm du
%   =\frac{\partial u}{\partial r}\cdot\frac{\partial r}{\partial x}\mathrm dx
%   +\frac{\partial u}{\partial r}\cdot\frac{\partial r}{\partial y}\mathrm dy
%   +\frac{\partial u}{\partial r}\cdot\frac{\partial r}{\partial z}\mathrm dz
% \]
% \[
%   =f(r)r\left(\frac{x}{r}\mathrm dx+\frac{y}{r}\mathrm dy+\frac{z}{r}\mathrm dz\right)
%   =f(r)(x\mathrm dx+y\mathrm dy+z\mathrm dz).
% \]
% 由定理~\ref{theorem:ma-16-7-1} 知 $\int_\Gamma f(r)(x\mathrm dx+y\mathrm dy+z\mathrm dz)$ 在 $D$ 内与路径无关.
% \end{proof}
% 
% 
% \begin{remark}\label{remark:ma-16-source-76}
% 选择不同的 $f(r)$, 我们可以得到许多在 $\mathbb R^3$ 中不含原点的区域内的函数 $P,Q,R$, 使得 $\int_\Gamma P\mathrm dx+Q\mathrm dy+R\mathrm dz$ 在 $D$ 内与路径无关.
% \end{remark}
% 
% 最后我们介绍一下如何来求原函数, 即给定微分 $P\mathrm dx+Q\mathrm dy+R\mathrm dz$, 当 $\int_\Gamma P\mathrm dx+Q\mathrm dy+R\mathrm dz$ 在区域 $D$ 内与路径无关时, 如何找到在 $D$ 内的函数 $u(x,y,z)$, 使得 $\mathrm du=P\mathrm dx+Q\mathrm dy+R\mathrm dz$.
% 
% 第一种方法是不定积分法: 由于
% \[
%   \frac{\partial u}{\partial x}=P,\quad
%   \frac{\partial u}{\partial y}=Q,\quad
%   \frac{\partial u}{\partial z}=R,
% \]
% 将 $y,z$ 视为常数, 对 $\dfrac{\partial u}{\partial x}=P$ 关于 $x$ 求不定积分得
% \[
%   u(x,y,z)=\int P\mathrm dx+\varphi(y,z),
% \]
% 其中 $\varphi(y,z)$ 是与 $x$ 无关的 $(y,z)$ 的待定函数。因此我们有
% \[
%   \frac{\partial\varphi}{\partial y}
%   =Q-\frac{\partial}{\partial y}\int P\mathrm dx.
% \]
% 在上面等式两边对 $y$ 求不定积分得
% \[
%   \varphi(y,z)=\int\left(Q-\frac{\partial}{\partial y}\int P\mathrm dx\right)\mathrm dy+\psi(z),
% \]
% 其中 $\psi(z)$ 是与 $y$ 无关的 $z$ 的待定函数。因此我们有
% \[
%   \frac{\partial\psi}{\partial z}
%   =
%   R-\frac{\partial}{\partial z}\int P\mathrm dx
%   -\frac{\partial}{\partial z}
%   \left(\int\left(Q-\frac{\partial}{\partial y}\int P\mathrm dx\right)\mathrm dy\right).
% \]
% 在上面等式两边对 $z$ 求不定积分得
% \[
%   \psi(z)=\int\left[
%   R-\frac{\partial}{\partial z}\int P\mathrm dx
%   -\frac{\partial}{\partial z}
%   \left(\int\left(Q-\frac{\partial}{\partial y}\int P\mathrm dx\right)\mathrm dy\right)
%   \right]\mathrm dz.
% \]
% 最后我们有
% \[
%   u(x,y,z)=\int P\mathrm dx
%   +\int\left(Q-\frac{\partial}{\partial y}\int P\mathrm dx\right)\mathrm dy
% \]
% \[
%   \quad+\int\left[
%   R-\frac{\partial}{\partial z}\int P\mathrm dx
%   -\frac{\partial}{\partial z}
%   \left(\int\left(Q-\frac{\partial}{\partial y}\int P\mathrm dx\right)\mathrm dy\right)
%   \right]\mathrm dz.
% \]
% 
% 第二种方法是曲线积分法: 给定 $D$ 中一点 $A_0=(x_0,y_0,z_0)$, 对于任意一点 $A=(x,y,z)\in D$, 在 $D$ 内选择一些连接 $A_0,A$ 且易于求曲线积分的路径, 然后计算曲线积分即得.
% 
% \begin{example}\label{example:ma-16-7-3}
% 在 $\mathbb R^3$ 内求 $u(x,y,z)$, 使得
% \[
%   \mathrm du(x,y,z)=(yze^{xyz}+3x^2)\mathrm dx+(xze^{xyz}+\sin y)\mathrm dy+(xye^{xyz}+2z)\mathrm dz.
% \]
% \end{example}
% 
% 
% \begin{solve}
% 记
% \[
%   P(x,y,z)=yze^{xyz}+3x^2,\quad Q(x,y,z)=xze^{xyz}+\sin y,
% \]
% \[
%   R(x,y,z)=xye^{xyz}+2z.
% \]
% 容易验证它们处处成立
% \[
%   \frac{\partial R}{\partial y}=\frac{\partial Q}{\partial z},\quad
%   \frac{\partial P}{\partial z}=\frac{\partial R}{\partial x},\quad
%   \frac{\partial Q}{\partial x}=\frac{\partial P}{\partial y},
% \]
% 因此 $\int_\Gamma P\mathrm dx+Q\mathrm dy+R\mathrm dz$ 在 $\mathbb R^3$ 内与路径无关.
% 
% 下面我们用两种方法来求 $u(x,y,z)$.
% 
% 方法 1\quad 由 $\dfrac{\partial u}{\partial x}=yze^{xyz}+3x^2$, 两边对 $x$ 求不定积分得
% \begin{equation}
%   u(x,y,z)=e^{xyz}+x^3+\varphi(y,z),
%   \label{eq:ma-16-7-1}
% \end{equation}
% 其中 $\varphi(y,z)$ 为与 $x$ 无关的 $(y,z)$ 的待定函数。上式两边对 $y$ 求偏导数得
% \begin{equation}
%   Q(x,y,z)=xze^{xyz}+\sin y
%   =\frac{\partial u}{\partial y}
%   =xze^{xyz}+\frac{\partial\varphi}{\partial y},
%   \label{eq:ma-16-7-2}
% \end{equation}
% 因此 $\dfrac{\partial\varphi}{\partial y}=\sin y$。该式两边对 $y$ 求不定积分得
% \begin{equation}
%   \varphi(y,z)=-\cos y+\psi(z),
%   \label{eq:ma-16-7-3}
% \end{equation}
% 其中 $\psi(z)$ 为与 $x,y$ 无关的 $z$ 的待定函数。将式 \eqref{eq:ma-16-7-3} 代入式 \eqref{eq:ma-16-7-1}, 再对 $z$ 求偏导数得
% \[
%   R(x,y,z)=xye^{xyz}+2z=\frac{\partial u}{\partial z}=xye^{xyz}+\psi'(z),
% \]
% 从而 $\psi'(z)=2z$。此式两边求不定积分得 $\psi(z)=z^2+C$, 其中 $C$ 为常数。因此
% \[
%   u(x,y,z)=e^{xyz}+x^3-\cos y+z^2+C.
% \]
% 
% 方法 2\quad 设 $\mathbb R^3$ 中任意一点的坐标为 $(\xi,\eta,\zeta)$, 并设 $\Gamma_1,\Gamma_2,\Gamma_3$ 分别是连接 $(0,0,0)$ 与 $(\xi,0,0)$, $(\xi,0,0)$ 与 $(\xi,\eta,0)$, $(\xi,\eta,0)$ 与 $(\xi,\eta,\zeta)$ 的线段, 则 $\Gamma_1\cup\Gamma_2\cup\Gamma_3$ 是连接 $(0,0,0)$ 与 $(\xi,\eta,\zeta)$ 的一条分段光滑曲线.
% 
% 由于曲线积分与路径无关, 我们得到 $P\mathrm dx+Q\mathrm dy+R\mathrm dz$ 的一个原函数如下:
% \[
%   u_1(\xi,\eta,\zeta)
%   =\int_{(0,0,0)}^{(\xi,\eta,\zeta)}P\mathrm dx+Q\mathrm dy+R\mathrm dz
% \]
% \[
%   =\int_0^\xi P(x,0,0)\mathrm dx
%   +\int_0^\eta Q(\xi,y,0)\mathrm dy
%   +\int_0^\zeta R(\xi,\eta,z)\mathrm dz
% \]
% \[
%   =\int_0^\xi 3x^2\mathrm dx+\int_0^\eta\sin y\mathrm dy
%   +\int_0^\zeta(\xi\eta e^{\xi\eta z}+2z)\mathrm dz
% \]
% \[
%   =\xi^3-\cos\eta+1+e^{\xi\eta\zeta}-1+\zeta^2
%   =e^{\xi\eta\zeta}+\xi^3-\cos\eta+\zeta^2.
% \]
% 将 $(\xi,\eta,\zeta)$ 换回 $(x,y,z)$, 并注意到两个原函数可以相差一个常数, 我们有
% \[
%   u(x,y,z)=e^{xyz}+x^3-\cos y+z^2+C.
% \qedhere
% \]
% \end{solve}
% 
% 
% \section{场论简介}
% \label{sec:ma-16-8}
% 
% 场的概念来源于物理学。在物理学中, 通常将一些量在空间区域的分布称为场。最常见的场有密度场、温度场、引力场、磁场等。如果一个场可以用一个函数来刻画, 则称该场为数量场; 如果一个场需要用一个向量函数来刻画, 则称该场为向量场。例如, 密度场与温度场都是数量场, 而引力场与磁场则都是向量场.
% 
% 有些场只与空间的位置有关, 而与时间的变化无关, 这些场称为稳定场, 如地球的引力场就是稳定场。而温度场等随着空间位置及时间的变化而变化, 这样的场则称为不稳定场.
% 
% \subsection{数量场的梯度}
% \label{subsec:ma-16-8-1}
% 
% 前面我们已经介绍了梯度的概念, 在这里我们来回顾一下。设 $u=u(x,y,z)$ 是区域 $D\subset\mathbb R^3$ 内的一个数量场 (我们这样来表示数量场, 实际上是建立了一个直角坐标系, 从而数量场在该坐标系下可由 $u=u(x,y,z)$ 来表示), 则当 $u$ 是可微函数时, $u$ 的梯度为
% \[
%   \operatorname{grad}u=
%   \left(\frac{\partial u}{\partial x},\frac{\partial u}{\partial y},\frac{\partial u}{\partial z}\right).
% \]
% 设 $\boldsymbol x_0=(x_0,y_0,z)\in D$ 且 $\boldsymbol v=(\cos\alpha,\cos\beta,\cos\gamma)$ 为一个单位向量, 则过 $A_0$ 沿该单位向量的方向导数为
% \[
%   \left.\frac{\partial f(\boldsymbol x_0)}{\partial\boldsymbol v}\right|_{\boldsymbol x_0}
%   =
%   \left.\left(\frac{\partial u}{\partial x}\cos\alpha+\frac{\partial u}{\partial y}\cos\beta+\frac{\partial u}{\partial z}\cos\gamma\right)\right|_{\boldsymbol x_0}
%   =
%   \left.\operatorname{grad}u\right|_{\boldsymbol x_0}\cdot\boldsymbol v.
% \]
% 因此 $\left.\operatorname{grad}u\right|_{\boldsymbol x_0}$ 与 $\boldsymbol v$ 同向时, 在 $\boldsymbol x_0$ 处沿 $\boldsymbol v$ 的方向导数达到最大。换句话说, 一个数量场的梯度是一个向量; 它是数量场改变最快的方向, 其向量的模长为数量场沿该方向的变化率。由此定义的梯度显然与表示该数量场的坐标系的选取无关.
% 
% 梯度与数量场的等位面密切相关。所谓等位面是指区域 $D$ 内满足方程 $u(x,y,z)=C(C$ 为常数) 的根的集合。当 $u(x,y,z)$ 具有较好的性质时 (如是 $C^1$ 的), 数量场的等位面一般是一个曲面。显然, 任何两个不同的等位面均不相交, 且 $D$ 内任一点都在某一等位面上。等位面可能是空集、一个点或者一个曲面。当等位面是一曲面时, 前面我们已经知道它的法向量为
% \[
%   \left(\frac{\partial u}{\partial x},\frac{\partial u}{\partial y},\frac{\partial u}{\partial z}\right).
% \]
% 由此, 等位面在点 $\boldsymbol x\in D$ 处的法向量即为 $u$ 在 $\boldsymbol x$ 处的梯度 $\left.\operatorname{grad}u\right|_{\boldsymbol x}$.
% 
% \subsection{向量场的向量线}
% \label{subsec:ma-16-8-2}
% 
% 设 $D\subset\mathbb R^3$ 为一个区域, 在 $D$ 内有一个向量场
% \[
%   \boldsymbol F(x,y,z)=(P(x,y,z),Q(x,y,z),R(x,y,z)),
% \]
% 其中 $P,Q,R$ 在 $D$ 内均具有连续偏导数。现在设 $\Gamma$ 为 $D$ 内的一条光滑曲线, 其方程为
% \[
%   \begin{cases}
%     x=x(t),\\
%     y=y(t),\\
%     z=z(t),
%   \end{cases}\quad t\in I,
% \]
% 则该曲线上任一点处的切向量为 $(x'(t),y'(t),z'(t))$。若对于 $\forall t\in I$, 该切向量与向量 $\boldsymbol F(x(t),y(t),z(t))$ 平行, 则曲线 $\Gamma$ 称为 $\boldsymbol F$ 的一条向量线。由于 $(x'(t),y'(t),z'(t))$ 平行于 $(P,Q,R)$, 我们有
% \[
%   \frac{x'(t)\mathrm dt}{P}=\frac{y'(t)\mathrm dt}{Q}=\frac{z'(t)\mathrm dt}{R},
% \]
% 即 $\Gamma$ 上的点必须满足
% \[
%   \frac{\mathrm dx}{P}=\frac{\mathrm dy}{Q}=\frac{\mathrm dz}{R}.
% \]
% 因此我们也称满足上述方程的曲线为 $\boldsymbol F$ 的向量线。对于 $D$ 内的 $C^1$ 向量函数来说, 过 $D$ 内每一点有且只有一条向量线, 并且不同的向量线互不相交, 所有向量线充满区域 $D$。关于向量线理论在后续课程中还有介绍.
% 
% \begin{example}\label{example:ma-16-8-1}
% 设向量场 $\boldsymbol F(x,y,z)=(x,y,z)$, 求 $\boldsymbol F$ 在 $\mathbb R^3\setminus\{(0,0,0)\}$ 的向量线.
% \end{example}
% 
% 
% \begin{solve}
% 设 $(x_0,y_0,z_0)\in\mathbb R^3$, 且 $(x_0,y_0,z_0)\ne(0,0,0)$, 不妨设 $x_0\ne0$。下面我们来求过点 $(x_0,y_0,z_0)$ 的向量线。由
% \[
%   \frac{\mathrm dx}{x}=\frac{\mathrm dy}{y},\quad
%   \frac{\mathrm dx}{x}=\frac{\mathrm dz}{z},
% \]
% 知 $y=c_1x,z=c_2x$, 其中 $c_1,c_2$ 为待定常数。再由该向量线过点 $(x_0,y_0,z_0)$ 知 $c_1=\dfrac{y_0}{x_0},c_2=\dfrac{z_0}{x_0}$ 即
% \[
%   \begin{cases}
%     y=\dfrac{y_0}{x_0}x,\\
%     z=\dfrac{z_0}{x_0}x,
%   \end{cases}
% \]
% 此即射线 $\dfrac{x}{x_0}=\dfrac{y}{y_0}=\dfrac{z}{z_0}$。从几何上看, 也可以简单地求出该向量线.
% \end{solve}
% 
% 
% \subsection{向量场的散度}
% \label{subsec:ma-16-8-3}
% 
% 设 $\boldsymbol F=(P(x,y,z),Q(x,y,z),R(x,y,z))$ 是空间区域 $D$ 内的一个 $C^1$ 向量场 (函数), 则 $\boldsymbol F$ 的散度是一个数量场, 其定义为
% \[
%   \operatorname{div}\boldsymbol F=
%   \frac{\partial P}{\partial x}+\frac{\partial Q}{\partial y}+\frac{\partial R}{\partial z}.
% \]
% 上述定义的散度实际上与坐标系的选取无关, 对此我们可以用高斯公式来证明之。事实上, 将 $\boldsymbol F$ 看做一个流速场。任给 $D$ 中一点 $\boldsymbol x_0$, 以 $\boldsymbol x_0$ 为心, $r$ 为半径作一小球体 $V_r$, 使得 $V_r\subset D$, 再取 $\partial V_r$ 外侧的单位法向量为 $\boldsymbol n$, 则
% \[
%   \iint_{\partial V_r}\boldsymbol F\cdot\boldsymbol n\mathrm dS
%   =
%   \iiint_{V_r}\operatorname{div}\boldsymbol F\mathrm dV.
% \]
% 由第二型曲面积分的几何意义, 上述等式的左边是该流速场单位时间从内到外通过 $\partial V_r$ 的流量, 由此它不依赖于坐标系的选取。由重积分的第一中值定理, 有
% \[
%   \operatorname{div}\boldsymbol F(\boldsymbol x_0)
%   =
%   \lim_{r\to0+0}
%   \frac{\iint_{\partial V_r}\boldsymbol F\cdot\boldsymbol n\mathrm dS}{V_r},
% \]
% 其中我们仍然用 $V_r$ 表示 $V_r$ 的体积。由此说明散度的定义也不依赖于坐标系的选取, 是一个向量场的本质属性.
% 
% 对于在空间区域 $D$ 内具有连续偏导数的向量函数 $\boldsymbol F(x,y,z)$, 若它的散度 $\operatorname{div}\boldsymbol F$ 在 $\boldsymbol x_0=(x_0,y_0,z_0)\in D$ 处不为零, 则当 $\operatorname{div}\boldsymbol F(\boldsymbol x_0)>0$ 时, 对于任何一个包含 $\boldsymbol x_0$ 的充分小光滑闭曲面 $S$ 所围的区域 $V$, 有
% \[
%   \iiint_V\operatorname{div}\boldsymbol F\mathrm dV>0.
% \]
% 由高斯公式知, 此处从内向外流出的流量为正。因此 $\boldsymbol x_0$ 也称为该向量场的源。反之, 若 $\operatorname{div}\boldsymbol F(\boldsymbol x_0)<0$, 则对任何内部含有 $\boldsymbol x_0$ 的充分小的闭曲面, 从外面里流进的流量为正。这时, $\boldsymbol x_0$ 也称该向量场的汇。当 $\operatorname{div}\boldsymbol F(\boldsymbol x_0)=0$ 时, 称向量场 $\boldsymbol F$ 在 $\boldsymbol x_0$ 处无源.
% 
% \subsection{向量场的旋度}
% \label{subsec:ma-16-8-4}
% 
% 如同上一小节, 我们仍设 $\boldsymbol F(x,y,z)=(P(x,y,z),Q(x,y,z),R(x,y,z))$ 为区域 $D\subset\mathbb R^3$ 内的一个 $C^1$ 向量场。在斯托克斯公式中, 曾出现如下与 $\boldsymbol F$ 有关的向量
% \[
%   \begin{vmatrix}
%     \boldsymbol i&\boldsymbol j&\boldsymbol k\\
%     \frac{\partial}{\partial x}&\frac{\partial}{\partial y}&\frac{\partial}{\partial z}\\
%     P&Q&R
%   \end{vmatrix}.
% \]
% 我们称此向量为 $\boldsymbol F$ 的旋度, 记为 $\operatorname{rot}\boldsymbol F$.
% 
% 设 $S$ 是一分片光滑曲面, $\partial S$ 是空间分段光滑闭曲线, 且其定向与 $S$ 的侧的单位法向量 $\boldsymbol n$ 的选定一致, 则斯托克斯公式告诉我们
% \[
%   \oint_{\partial S}P\mathrm dx+Q\mathrm dy+R\mathrm dz
%   =
%   \iint_S\operatorname{rot}\boldsymbol F\cdot\boldsymbol n\mathrm dS.
% \]
% 上式左边的线积分通常称为 $\boldsymbol F$ 沿 $\partial S$ 的环量, 它显然可记成 $\oint_{\partial S}\boldsymbol F\mathrm d\boldsymbol s$, 其中 $\mathrm d\boldsymbol s=(\mathrm dx,\mathrm dy,\mathrm dz)$。由此定义的环量与坐标系的选取无关.
% 
% 利用这一定义, 我们同样可说明旋度也与坐标系选取无关。事实上, 任取以 $\boldsymbol x_0\in D$ 为心, 半径 $r$ 充分小的一个圆盘 $S_r\subset D$, 取定 $\partial S_r$ 的方向与 $S_r$ 的一侧的单位法向量 $\boldsymbol n$ 保持一致, 则
% \[
%   \oint_{\partial S_r}\boldsymbol F\mathrm d\boldsymbol s
%   =
%   \iint_{S_r}\operatorname{rot}\boldsymbol F\cdot\boldsymbol n\mathrm dS,
% \]
% 从而
% \[
%   \lim_{r\to0}\operatorname{rot}\boldsymbol F\cdot\boldsymbol n
%   =
%   \lim_{r\to0}
%   \frac{\displaystyle\oint_{\partial S_r}\boldsymbol F\mathrm d\boldsymbol s}{\pi r^2}.
% \]
% 由于 $S_r$ 是任意一个圆盘, 从而 $\boldsymbol n$ 可取到 $\boldsymbol x_0$ 出发的任何方向。上式说明向量 $\operatorname{rot}\boldsymbol F$ 在任何方向的投影均不依赖于坐标系的选取, 从而 $\operatorname{rot}\boldsymbol F$ 与坐标系的选取无关。由斯托克斯定理, 旋度与环量是密切相关的。特别地, 曲线积分 $\int_\Gamma P\mathrm dx+Q\mathrm dy+R\mathrm dz$ 在区域 $D$ 内与路径无关的一个必要条件是 $\boldsymbol F$ 的旋度处处为零。当 $\boldsymbol F$ 的旋度处处为零时, 我们称 $\boldsymbol F$ 为一个无旋场.
% 
% \subsection{一些重要算子}
% \label{subsec:ma-16-8-5}
% 
% 所谓“算子”是作用在某些特定对象 (如函数、向量函数等) 的一些特别的运算。人们对它们赋予了一些特殊的记号。在这里我们主要介绍两个重要的算子.
% 
% \paragraph*{哈密尔顿算子}
% 
% 
% 哈密尔顿 (Hamilton) 曾在 $\mathbb R^n$ 中引进了以下算子:
% \[
%   \nabla=\left(\frac{\partial}{\partial x_1},\frac{\partial}{\partial x_2},\cdots,\frac{\partial}{\partial x_n}\right).
% \]
% 我们称该算子为哈密尔顿算子。对于区域 $D\subset\mathbb R^n$ 内的可微函数 $f(x_1,x_2,\cdots,x_n)$, 哈密尔顿算子 $\nabla$ 对 $f$ 的作用定义为
% \[
%   \nabla f=\left(\frac{\partial f}{\partial x_1},\frac{\partial f}{\partial x_2},\cdots,\frac{\partial f}{\partial x_n}\right)=\operatorname{grad}f.
% \]
% 
% 容易验证, 哈密尔顿算子 $\nabla$ 具有如下性质: 设 $c\in\mathbb R$, $f,g$ 是可微函数, 则
% \[
%   (1)\quad \nabla(cf)=c\nabla f;
% \]
% \[
%   (2)\quad \nabla(f+g)=\nabla f+\nabla g;
% \]
% \[
%   (3)\quad \nabla(fg)=f\nabla g+g\nabla f.
% \]
% 
% 另外, $\mathbb R^3$ 中的向量场 $\boldsymbol F=(P,Q,R)$ 的散度可表示为
% \[
%   \operatorname{div}\boldsymbol F
%   =
%   \frac{\partial P}{\partial x}+\frac{\partial Q}{\partial y}+\frac{\partial R}{\partial z}
%   =
%   \nabla\cdot\boldsymbol F,
% \]
% 而 $\boldsymbol F$ 的旋度表示为
% \[
%   \operatorname{rot}\boldsymbol F
%   =
%   \begin{vmatrix}
%     \boldsymbol i&\boldsymbol j&\boldsymbol k\\
%     \frac{\partial}{\partial x}&\frac{\partial}{\partial y}&\frac{\partial}{\partial z}\\
%     P&Q&R
%   \end{vmatrix}
%   =
%   \nabla\times\boldsymbol F.
% \]
% 
% \paragraph*{拉普拉斯算子}
% 
% 
% 下面我们介绍第二个重要算子, 即拉普拉斯算子。在 $\mathbb R^n$ 中, 拉普拉斯算子的定义为
% \[
%   \Delta=\sum_{i=1}^n\frac{\partial^2}{\partial x_i^2}.
% \]
% 对于在区域 $D\subset\mathbb R^n$ 内具有二阶连续偏导数的函数 $f(x_1,x_2,\cdots,x_n)$, 拉普拉斯算子 $\Delta$ 对 $f$ 的作用定义为
% \[
%   \Delta f=\sum_{i=1}^n\frac{\partial^2f}{\partial x_i^2}.
% \]
% 当 $u=f(x_1,x_2,\cdots,x_n)$ 在 $D$ 内恒满足
% \[
%   \Delta u=\sum_{i=1}^n\frac{\partial^2f}{\partial x_i^2}=0
% \]
% 时, 称 $u$ 为 $D$ 内的调和函数.
% 
% 显然一元函数 $f(x)$ 为调和函数的充分必要条件是 $f(x)$ 为线性函数。对二元函数 $u=f(x,y)$, 由定义, 当 $\Delta u=\dfrac{\partial^2f}{\partial x^2}+\dfrac{\partial^2f}{\partial y^2}=0$ 时, $u=f(x,y)$ 是调和函数。读者不难验证, $e^x\cos y,2xy$ 均是调和函数。调和函数具有许多特殊的性质, 在后续课程中将有详细介绍.
% 
% \begin{example}\label{example:ma-16-8-2}
% 设 $D\subset\mathbb R^2$ 是单连通区域, $u(x,y)$ 是 $D$ 内的调和函数。证明曲线积分
% \[
%   \int_\Gamma u_y'\mathrm dx-u_x'\mathrm dy
% \]
% 在 $D$ 内与路径无关, 且其原函数为 $D$ 内的调和函数.
% \end{example}
% 
% 
% \begin{proof}
% 由于 $P\mathrm dx+Q\mathrm dy=u_y'\mathrm dx-u_x'\mathrm dy$ 及
% \[
%   \frac{\partial Q}{\partial x}-\frac{\partial P}{\partial y}
%   =
%   \frac{\partial(-u_x')}{\partial x}-\frac{\partial u_y'}{\partial y}
%   =
%   -\left(\frac{\partial^2u}{\partial x^2}+\frac{\partial^2u}{\partial y^2}\right)
%   \equiv0,
% \]
% 再注意到 $D$ 是单连通区域, 根据定理~\ref{theorem:ma-16-7-3} 知, $\int_\Gamma u_y'\mathrm dx-u_x'\mathrm dy$ 在 $D$ 内与路径无关.
% 
% 设 $v(x,y)$ 为 $u_y'\mathrm dx-u_x'\mathrm dy$ 的一个原函数, 则
% \[
%   \frac{\partial v}{\partial x}=u_y',
%   \quad
%   \frac{\partial v}{\partial y}=-u_x'.
% \]
% 因此
% \[
%   \frac{\partial^2v}{\partial x^2}+\frac{\partial^2v}{\partial y^2}
%   =
%   \frac{\partial^2u}{\partial y\partial x}
%   -
%   \frac{\partial^2u}{\partial x\partial y}
%   \equiv0,
% \]
% 所以 $v(x,y)$ 是 $D$ 内的调和函数.
% \end{proof}
% 
% 
% \begin{remark}\label{remark:ma-16-source-83}
% 设 $u(x,y)$ 是 $D$ 内的调和函数, 在上例中得到的 $D$ 内的调和函数 $v(x,y)$ 称为 $u(x,y)$ 的共轭调和函数.
% \end{remark}
% 
% 
% \section*{习题十六}
% \phantomsection
% \addcontentsline{toc}{section}{习题十六}
% \label{exercises:ma-16}
% 
% 1。设 $\Gamma$ 是空间 $\mathbb R^3$ 中一条光滑的物质曲线, 其密度函数 $\rho(x,y,z)$ 在 $\Gamma$ 上连续, 试求 $\Gamma$ 的质心坐标.
% 
% 2。设曲线 $\Gamma$ 是单位圆周: $x^2+y^2=1$, 求下列第一型曲线积分:
% \[
%   (1)\quad \int_\Gamma x\mathrm ds;\qquad
%   (2)\quad \int_\Gamma xy\mathrm ds;\qquad
%   (3)\quad \int_\Gamma x^2\mathrm ds;\qquad
%   (4)\quad \int_\Gamma |x|\mathrm ds.
% \]
% 
% 3。设曲线 $\Gamma$ 为球面 $x^2+y^2+z^2=1$ 与平面 $x+y+z=0$ 的交线, 求下列第一型曲线积分:
% \[
%   (1)\quad \int_\Gamma x\mathrm ds;\qquad
%   (2)\quad \int_\Gamma xy\mathrm ds;\qquad
%   (3)\quad \int_\Gamma x^2\mathrm ds.
% \]
% 
% 4。求下列第一型曲线积分:
% 
% \begin{remark}
% 原稿在本题处截断，未提供完整题面。
% \end{remark}
% 
% % 整合来源：math-17.tex
% \chapter{含参变量积分}
% \label{chap:ma-17}
% 
% 含参变量积分在许多自然科学问题中常常碰到。例如, 一个不稳定流速场可以表示成 $F(x,y,z,t),(x,y,z)\in D,t\in(a,b)$, 其中 $D$ 是 $\mathbb R^3$ 中的区域, $t$ 是时间变量。一般说来, $F(x,y,z,t)$ 是 $(x,y,z,t)$ 的连续函数。现设一个光滑双侧曲面 $S$ 位于 $D$ 内, 取定 $S$ 的一侧, 则在任意时刻 $t$, $F$ 在 $S$ 上的第二型曲面积分与流量密切相关。对固定的时刻 $t$, 我们可以在 $S$ 上对 $F$ 进行第二型曲面积分。当 $t$ 变化时, 该积分便是 $t$ 的函数。一个自然的问题是: 该积分是否是 $t$ 的连续函数? 如果该积分是 $t$ 的连续函数, 则通过对 $t$ 积分就可计算出某时间段内流体通过 $S$ 的流量。以后我们还要研究: 当 $F$ 关于 $t$ 可求偏导数时, 该积分关于 $t$ 的导数是否存在? 这些问题的解决就要用到含参变量积分的理论.
% 
% 在本章中, 我们将重点讨论含有参数的函数的积分问题, 讨论的积分以定积分与广义积分为主。读者可以将这些积分推广到其他积分 (如重积分、曲线积分和曲面积分等) 的情形.
% 
% \section{含参变量定积分}
% \label{sec:ma-17-1}
% 
% 设函数 $f(x,y)$ 在平面区域 $D=[a,b]\times[c,d]$ 上有定义。若对于 $\forall x\in[a,b]$, 定积分 $\int_c^d f(x,y)\mathrm dy$ 存在, 则由此定义了区间 $[a,b]$ 上的函数
% \[
%   I(x)=\int_c^d f(x,y)\mathrm dy.
% \]
% 我们称如此定义的函数为\keyterm{含参变量定积分} (简称\keyterm{含参变量积分}), 其中 $x$ 为参变量。同理, 若对于 $\forall y\in[c,d]$, $J(y)=\int_a^b f(x,y)\mathrm dx$ 存在, 则也称 $J(y)$ 为含参变量定积分, 其中 $y$ 为参变量.
% 
% 由于定积分是一个极限过程, 因此对含参变量积分的研究有点类似于对函数序列或函数项级数的相应研究。我们主要的研究兴趣在于: 对 $f(x,y)$ 加什么条件, 可以保证 $I(x)$ 在区间 $[a,b]$ 上连续、可积以及可导? 下面我们分别来讨论这些问题.
% 
% \begin{theorem}\label{theorem:ma-17-1-1}
% 设函数 $f(x,y)$ 在区域 $D=[a,b]\times[c,d]$ 上连续, 则对于 $\forall x\in[a,b]$, 含参变量定积分 $I(x)=\int_c^d f(x,y)\mathrm dy$ 存在, 并且 $I(x)$ 在区间 $[a,b]$ 上连续.
% \end{theorem}
% 
% 
% \begin{proof}
% 由于 $f(x,y)\in C(D)$, 因此对固定的 $x$, $f(x,y)$ 在 $[c,d]$ 上连续, 从而 $I(x)=\int_c^d f(x,y)\mathrm dy$ 存在。下面证 $I(x)$ 在 $[a,b]$ 上连续。对于 $\forall x_0\in[a,b]$ 及 $\forall x\in[a,b]$, 有
% \[
%   I(x)-I(x_0)=\int_c^d [f(x,y)-f(x_0,y)]\mathrm dy.
% \]
% 由 $f(x,y)$ 在 $D$ 上的一致连续性知, 对于 $\forall\varepsilon>0,\exists\delta>0$, 当 $(x_1,y_1),(x_2,y_2)\in D$ 且满足
% \[
%   \sqrt{(x_1-x_2)^2+(y_1-y_2)^2}<\delta
% \]
% 时, 有
% \[
%   |f(x_1,y_1)-f(x_2,y_2)|<\frac{\varepsilon}{d-c}.
% \]
% 所以, 当 $|x-x_0|<\delta$ 且 $x\in[a,b]$ 时, 有
% \[
%   |I(x)-I(x_0)|
%   \leqslant \int_c^d |f(x,y)-f(x_0,y)|\mathrm dy
%   < \frac{\varepsilon}{d-c}(d-c)=\varepsilon.
% \qedhere
% \]
% \end{proof}
% 
% 
% 显然, 定理~\ref{theorem:ma-17-1-1} 中 $f(x,y)$ 在 $D$ 上连续只是 $I(x)$ 连续的充分条件。读者容易举出 $f(x,y)$ 在 $D$ 上不连续时 $I(x)$ 仍然连续的例子.
% 
% 另外, 在定理~\ref{theorem:ma-17-1-1} 的条件下, 我们有
% \[
%   \lim_{x\to x_0}\int_c^d f(x,y)\mathrm dy
%   =\lim_{x\to x_0} I(x)
%   =I(x_0)
%   =\int_c^d \lim_{x\to x_0} f(x,y)\mathrm dy.
% \]
% 因此, 函数 $I(x)$ 的连续性可以解释为极限运算与积分运算的可交换性.
% 
% \begin{remark}\label{remark:ma-17-source-3}
% 定理~\ref{theorem:ma-17-1-1} 可以作如下推广: 在定理~\ref{theorem:ma-17-1-1} 的假定下, 对于 $\forall(x,u)\in D$, 变上限含参变量积分
% \[
%   I(x,u)=\int_c^u f(x,y)\mathrm dy
% \]
% 存在, 并且二元函数 $I(x,u)$ 在 $D$ 上连续 (见本章习题)。对于变下限含参变量积分, 也有类似的结论.
% \end{remark}
% 
% 由化二重积分为累次积分的定理 (定理~\ref{theorem:ma-15-3-1}), 我们有如下关于含参变量定积分的积分定理.
% 
% \begin{theorem}\label{theorem:ma-17-1-2}
% 设函数 $f(x,y)$ 在区域 $D=[a,b]\times[c,d]$ 上连续, 则函数
% \[
%   I_1(x)=\int_c^d f(x,y)\mathrm dy
%   \quad\text{和}\quad
%   I_2(y)=\int_a^b f(x,y)\mathrm dx
% \]
% 分别在区间 $[a,b]$ 和 $[c,d]$ 上可积, 并且
% \[
%   \int_a^b I_1(x)\mathrm dx=\int_c^d I_2(y)\mathrm dy.
% \]
% 定理~\ref{theorem:ma-17-1-2} 说明, 在定理条件下, 两个求积分的顺序可以交换:
% \[
%   \int_a^b \mathrm dx\int_c^d f(x,y)\mathrm dy
%   =
%   \int_c^d \mathrm dy\int_a^b f(x,y)\mathrm dx.
% \]
% \end{theorem}
% 
% 下面我们来讨论含参变量定积分的可导性.
% 
% \begin{theorem}\label{theorem:ma-17-1-3}
% 设函数 $f(x,y)$ 及其偏导数 $f'_x(x,y)$ 在区域 $D=[a,b]\times[c,d]$ 上连续, 则函数 $I(x)=\int_c^d f(x,y)\mathrm dy$ 在区间 $[a,b]$ 上可导, 并且有
% \[
%   I'(x)=\int_c^d f'_x(x,y)\mathrm dy.
% \]
% \end{theorem}
% 
% 
% \begin{proof}
% 由于 $f'_x(x,y)$ 在 $D$ 上连续, 从而由定理~\ref{theorem:ma-17-1-1} 知函数
% \[
%   g(x)=\int_c^d f'_x(x,y)\mathrm dy
% \]
% 在 $[a,b]$ 上存在并且连续。现将 $g(x)$ 改写成
% \[
%   g(u)=\int_c^d f'_u(u,y)\mathrm dy.
% \]
% 对于 $\forall x\in[a,b]$, 对 $f'_u(u,y)$ 在 $D=[a,x]\times[c,d]$ 上应用定理~\ref{theorem:ma-17-1-2}, 有
% \[
% \begin{aligned}
%   \int_a^x g(u)\mathrm du
%   &=\int_a^x \mathrm du\int_c^d f'_u(u,y)\mathrm dy
%     =\int_c^d \mathrm dy\int_a^x f'_u(u,y)\mathrm du\\
%   &=\int_c^d [f(x,y)-f(a,y)]\mathrm dy
%     = I(x)-I(a).
% \end{aligned}
% \]
% 因此对变限定积分 $\int_a^x g(u)\mathrm du$ 求导数得
% \[
%   I'(x)=g(x)=\int_c^d f'_x(x,y)\mathrm dy.
% \qedhere
% \]
% \end{proof}
% 
% 
% 定理~\ref{theorem:ma-17-1-3} 说明了在 $f'_x(x,y)$ 连续的条件下求导数与积分的可交换性。此定理也可通过导数定义直接求出 $I'(x)$ 来加以证明.
% 
% \begin{theorem}\label{theorem:ma-17-1-4}
% 设函数 $f(x,y)$ 及其偏导数 $f'_x(x,y)$ 在区域 $D=[a,b]\times[c,d]$ 上连续, 且 $\varphi(x)\ (x\in[a,b])$ 是满足 $c\leqslant \varphi(x)\leqslant d$ 的可微函数, 则函数
% \[
%   I(x)=\int_c^{\varphi(x)} f(x,y)\mathrm dy
% \]
% 在区间 $[a,b]$ 上可导, 并且
% \[
%   I'(x)=\int_c^{\varphi(x)} f'_x(x,y)\mathrm dy
%   + f(x,\varphi(x))\varphi'(x).
% \]
% \end{theorem}
% 
% 
% \begin{proof}
% 对任意给定的 $x_0\in[a,b]$, 由 $c\leqslant\varphi(x_0)\leqslant d$ 及 $f(x,y)$ 在 $D$ 的连续性知, $f(x_0,y)$ 在 $[c,\varphi(x_0)]$ 上连续, 所以 $I(x_0)=\int_c^{\varphi(x_0)}f(x,y)\mathrm dy$ 存在。因此 $I(x)$ 在 $[a,b]$ 上有定义.
% 
% 另外, 令 $u=\varphi(x)$, 则 $I(x)$ 可以看成是 $F(x,u)=\int_c^u f(x,y)\mathrm dy$ 与 $u=\varphi(x)$ 的复合函数。由定理~\ref{theorem:ma-17-1-3} 我们知
% \[
%   \frac{\partial F(x,u)}{\partial x}=\int_c^u f'_x(x,y)\mathrm dy.
% \]
% 再由定理~\ref{theorem:ma-17-1-1} 的注知 $\frac{\partial F(x,u)}{\partial x}=\int_c^u f'_x(x,y)\mathrm dy$ 在 $D$ 上连续。对 $F(x,u)=\int_c^u f(x,y)\mathrm dy$ 两边关于 $u$ 求偏导数得
% \[
%   \frac{\partial F(x,u)}{\partial u}=f(x,u),
% \]
% 由假设它也在 $D$ 上连续。这说明 $F(x,u)$ 的两个偏导数在 $D$ 上连续。因此利用复合函数求导法得
% \[
% \begin{aligned}
%   I'(x)
%   &=\frac{\partial F(x,u)}{\partial x}
%     +\frac{\partial F(x,u)}{\partial u}\cdot\frac{\partial u}{\partial x}\\
%   &=\int_c^{\varphi(x)} f'_x(x,y)\mathrm dy
%     +f(x,\varphi(x))\varphi'(x).
% \end{aligned}
% \qedhere
% \]
% \end{proof}
% 
% 
% \begin{example}\label{example:ma-17-1-1}
% 求极限
% \[
%   \lim_{x\to0+0}\int_0^1 (1+y)^x e^{xy}\sin(x+y)\mathrm dy.
% \]
% \end{example}
% 
% 
% \begin{solve}
% 由于 $f(x,y)=(1+y)^x e^{xy}\sin(x+y)$ 在 $[0,+\infty)\times[0,1]$ 上连续, 因此 $I(x)=\int_0^1(1+y)^x e^{xy}\sin(x+y)\mathrm dy$ 在 $[0,+\infty)$ 上连续。特别地, 我们有
% \[
%   \lim_{x\to0+0}\int_0^1 (1+y)^x e^{xy}\sin(x+y)\mathrm dy
%   =\int_0^1 \sin y\,\mathrm dy
%   =1-\cos1.
% \qedhere
% \]
% \end{solve}
% 
% 
% \begin{example}\label{example:ma-17-1-2}
% 求定积分
% \[
%   \int_0^1 \frac{x^{2e-1}-x^{e-1}}{\ln x}\mathrm dx.
% \]
% \end{example}
% 
% 
% \begin{solve}
% 如果直接求此定积分, 计算将会非常困难, 因此考虑利用含参变量积分来计算。由
% \[
%   \frac{x^{2e-1}-x^{e-1}}{\ln x}
%   =
%   \int_{e-1}^{2e-1} x^y\mathrm dy
% \]
% 及 $f(x,y)=x^y$ 在 $[0,1]\times[e-1,2e-1]$ 上连续, 我们有
% \[
% \begin{aligned}
%   \int_0^1 \frac{x^{2e-1}-x^{e-1}}{\ln x}\mathrm dx
%   &=\int_0^1 \mathrm dx\int_{e-1}^{2e-1}x^y\mathrm dy\\
%   &=\int_{e-1}^{2e-1}\mathrm dy\int_0^1x^y\mathrm dx
%     =\int_{e-1}^{2e-1}\frac{\mathrm dy}{1+y}\\
%   &=\left.\ln(1+y)\right|_{e-1}^{2e-1}=\ln2.
% \end{aligned}
% \qedhere
% \]
% \end{solve}
% 
% 
% \begin{example}\label{example:ma-17-1-3}
% 设函数 $f(s,t)$ 在 $\mathbb R^2$ 上连续, 记
% \[
%   F(x)=\int_x^{x^2}\mathrm ds\int_s^x f(s,t)\mathrm dt,
% \]
% \end{example}
% 
% \noindent 求 $F'(x)$.
% 
% \begin{solve}
% 记 $g(s,x)=\int_s^x f(s,t)\mathrm dt$。由定理~\ref{theorem:ma-17-1-1} 的注知, 它在 $\mathbb R^2$ 中连续, 从而对于 $\forall x\in(-\infty,+\infty)$, 我们有
% \[
%   F(x)=\int_x^{x^2}g(s,x)\mathrm ds
%   =\int_0^{x^2}g(s,x)\mathrm ds-\int_0^x g(s,x)\mathrm ds
% \]
% 存在。另外, $\frac{\partial g(s,x)}{\partial x}=f(s,x)$ 在 $\mathbb R^2$ 上连续, 并且 $x,x^2$ 均是可导函数, 从而由定理~\ref{theorem:ma-17-1-4} 知 $F'(x)$ 存在.
% 
% 记
% \[
%   F_1(x)=\int_0^{x^2}g(s,x)\mathrm ds,\qquad
%   F_2(x)=\int_0^x g(s,x)\mathrm ds,
% \]
% 则有
% \[
% \begin{aligned}
%   F_1'(x)
%   &=\frac{\mathrm d}{\mathrm dx}\left(\int_0^{x^2}g(s,x)\mathrm ds\right)
%     =\int_0^{x^2}\frac{\partial g(s,x)}{\partial x}\mathrm ds
%       +g(x^2,x)(x^2)'\\
%   &=\int_0^{x^2}f(s,x)\mathrm ds
%     +2x\int_{x^2}^x f(x^2,t)\mathrm dt.
% \end{aligned}
% \]
% 同理,
% \[
%   F_2'(x)=\int_0^x f(s,x)\mathrm ds+\int_x^x f(x,t)\mathrm dt
%   =\int_0^x f(s,x)\mathrm ds.
% \]
% 最后我们有
% \[
%   F'(x)=F_1'(x)-F_2'(x)
%   =\int_x^{x^2}f(s,x)\mathrm ds+2x\int_{x^2}^x f(x^2,t)\mathrm dt.
% \qedhere
% \]
% \end{solve}
% 
% 
% \section{含参变量广义积分}
% \label{sec:ma-17-2}
% 
% 含参变量广义积分视积分为无穷积分或瑕积分可分为两种: 含参变量无穷积分和含参变量瑕积分。我们先来讨论含参变量无穷积分.
% 
% \subsection{含参变量无穷积分}
% \label{subsec:ma-17-2-1}
% 
% 设函数 $f(x,y)$ 在 $E\times[c,+\infty)$ 上有定义, 其中 $E\subset\mathbb R$ 为一个集合。若对于 $\forall x\in E$, 广义积分 $\int_c^{+\infty}f(x,y)\mathrm dy$ 收敛, 则可得到 $E$ 上的函数
% \[
%   I(x)=\int_c^{+\infty}f(x,y)\mathrm dy,
% \]
% 我们称该函数为\keyterm{含参变量无穷积分}.
% 
% 在这里我们仍将讨论 $I(x)$ 的连续性、可微性以及积分与无穷积分的交换性等问题。在讨论这些问题时, 一般我们假定 $E$ 是一个区间。如同函数项级数, 为了保证 $I(x)$ 能够继承 $f(x,y)$ 的一些性质, 引入含参变量无穷积分一致收敛的概念是必要的.
% 
% \begin{definition}\label{definition:ma-17-2-1}
% 设函数 $f(x,y)$ 在 $E\times[c,+\infty)$ 上有定义, 其中 $E\subset\mathbb R$ 是一个区间。若对于 $\forall\varepsilon>0,\exists A_0>0$, 当 $A>A_0$ 时, 对于 $\forall x\in E$, 有
% \[
%   \left|\int_A^{+\infty}f(x,y)\mathrm dy\right|<\varepsilon,
% \]
% 则称含参变量无穷积分 $\int_c^{+\infty}f(x,y)\mathrm dy$ 在 $E$ 上\keyterm{一致收敛}.
% \end{definition}
% 
% 在一致收敛的研究中, 柯西准则是一个非常重要的理论工具.
% 
% \begin{theorem}\label{theorem:ma-17-2-1}
% 设函数 $f(x,y)$ 在 $E\times[c,+\infty)$ 上有定义, 其中 $E\subset\mathbb R$ 是一个区间, 则含参变量无穷积分 $\int_c^{+\infty}f(x,y)\mathrm dy$ 在 $E$ 上一致收敛的充分必要条件是: 对于 $\forall\varepsilon>0,\exists A_0>0$, 当 $A,A'>A_0$ 时, 对于 $\forall x\in E$ 有
% \[
%   \left|\int_A^{A'}f(x,y)\mathrm dy\right|<\varepsilon.
% \]
% \end{theorem}
% 
% 
% \begin{proof}
% 必要性\quad 设 $\int_c^{+\infty}f(x,y)\mathrm dy$ 在 $E$ 上一致收敛, 则对于 $\forall\varepsilon>0,\exists A_0>0$, 当 $A>A_0$ 时, 对于 $\forall x\in E$ 有
% \[
%   \left|\int_A^{+\infty}f(x,y)\mathrm dy\right|<\frac{\varepsilon}{2};
% \]
% 从而当 $A,A'>A_0$ 时, 有
% \[
%   \left|\int_A^{A'}f(x,y)\mathrm dy\right|
%   \leqslant\left|\int_A^{+\infty}f(x,y)\mathrm dy\right|
%     +\left|\int_{A'}^{+\infty}f(x,y)\mathrm dy\right|
%   <\frac{\varepsilon}{2}+\frac{\varepsilon}{2}=\varepsilon.
% \]
% 
% 充分性\quad 设对于 $\forall\varepsilon>0,\exists A_0>0$, 当 $A,A'>A_0$ 时, 对于 $\forall x\in E$ 有
% \begin{equation}
%   \left|\int_A^{A'}f(x,y)\mathrm dy\right|<\frac{\varepsilon}{2}.
%   \label{eq:ma-17-2-1}
% \end{equation}
% 这说明 $\int_c^{+\infty}f(x,y)\mathrm dy$ 在每点 $x\in E$ 满足柯西准则, 从而它点点收敛。在式 \eqref{eq:ma-17-2-1} 中令 $A'\to+\infty$, 则对于 $\forall x\in E$, 都有
% \[
%   \left|\int_A^{+\infty}f(x,y)\mathrm dy\right|\leqslant\frac{\varepsilon}{2}<\varepsilon.
% \qedhere
% \]
% \end{proof}
% 
% 
% 设函数 $f(x,y)$ 在 $E\times[c,+\infty)$ 上有定义, 其中 $E\subset\mathbb R$ 是一个区间。若对于 $\forall x\in E$, $\int_c^{+\infty}|f(x,y)|\mathrm dy$ 收敛, 则称 $\int_c^{+\infty}f(x,y)\mathrm dy$ 在 $E$ 上\keyterm{绝对收敛}。显然, 若 $\int_c^{+\infty}f(x,y)\mathrm dy$ 在 $E$ 上绝对收敛, 则 $\int_c^{+\infty}f(x,y)\mathrm dy$ 在 $E$ 上收敛。另外, 若 $\int_c^{+\infty}|f(x,y)|\mathrm dy$ 在 $E$ 上一致收敛, 则称 $\int_c^{+\infty}f(x,y)\mathrm dy$ 在 $E$ 上\keyterm{绝对一致收敛}.
% 
% \begin{theorem}[{魏尔斯特拉斯定理}]\label{theorem:ma-17-2-2}
% 设函数 $f(x,y)$ 在 $E\times[c,+\infty)$ 上有定义, 其中 $E\subset\mathbb R$ 是一个区间。若存在函数 $g(y)(y\in[c,+\infty))$, 使得对于 $\forall x\in E$ 及 $\forall y\in[c,+\infty)$, 有 $|f(x,y)|\leqslant g(y)$, 并且 $\int_c^{+\infty}g(y)\mathrm dy$ 收敛, 则 $\int_c^{+\infty}f(x,y)\mathrm dy$ 在 $E$ 上绝对一致收敛.
% \end{theorem}
% 
% 
% \begin{proof}
% 对于 $\forall\varepsilon>0$, 由 $\int_c^{+\infty}g(y)\mathrm dy$ 收敛, 从而 $\exists A_0>0$, 当 $A,A'>A_0$ 时, 有
% \[
%   \left|\int_A^{A'}g(y)\mathrm dy\right|<\varepsilon.
% \]
% 我们不妨设 $A'>A$, 从而对于 $\forall x\in E$, 有
% \[
%   \left|\int_A^{A'}f(x,y)\mathrm dy\right|
%   \leqslant\int_A^{A'}|f(x,y)|\mathrm dy
%   \leqslant\int_A^{A'}g(y)\mathrm dy<\varepsilon.
% \]
% 由定理~\ref{theorem:ma-17-2-1} 知 $\int_c^{+\infty}|f(x,y)|\mathrm dy$ 在 $E$ 上一致收敛, 即 $\int_c^{+\infty}f(x,y)\mathrm dy$ 在 $E$ 上绝对一致收敛.
% \end{proof}
% 
% 
% 对于含参变量无穷积分一致收敛的判别法则, 可以参照无穷积分的收敛判别法则来加以修改得到。事实上, 只要将后者的判别法则中的条件对参数一致满足即可得到前者的判别法则。我们有下面的两个定理.
% 
% \begin{theorem}[{狄利克雷判别法}]\label{theorem:ma-17-2-3}
% 设函数 $f(x,y),g(x,y)$ 在 $E\times[c,+\infty)$ 上有定义 (其中 $E\subset\mathbb R$ 是一个区间), 并且满足:
% \[
% \begin{gathered}
%   (1)\quad \text{存在 } M>0,\ \text{对于 } \forall A>c \text{ 及 } \forall x\in E,\ \text{有 }
%   \left|\int_c^A f(x,y)\mathrm dy\right|\leqslant M;\\
%   (2)\quad \text{对任意固定的 } x\in E,\ g(x,y) \text{ 是 } y \text{ 的单调函数, 且对于 } \forall\varepsilon>0,\exists A>c,\\
%   \text{当 } y>A \text{ 时, 对一切 } x\in E,\ \text{有 } |g(x,y)|<\varepsilon,
%   \text{ 即当 } y\to+\infty \text{ 时, } g(x,y) \text{ 关于 } x \text{ 一致趋于 } 0,
% \end{gathered}
% \]
% 则含参变量无穷积分 $\int_c^{+\infty}f(x,y)g(x,y)\mathrm dy$ 在 $E$ 上一致收敛.
% \end{theorem}
% 
% 
% \begin{proof}
% 对于 $A'>A>c$, 由定积分第二中值定理, $\exists\xi\in(A,A')$, 使得
% \begin{equation}
%   \int_A^{A'}f(x,y)g(x,y)\mathrm dy
%   =
%   g(x,A)\int_A^\xi f(x,y)\mathrm dy
%   +g(x,A')\int_\xi^{A'}f(x,y)\mathrm dy.
%   \label{eq:ma-17-2-2}
% \end{equation}
% 由 (2) 知, 对于 $\forall\varepsilon>0,\exists A_0>c$, 当 $y>A_0$ 时, 对一切 $x\in E$, 有 $|g(x,y)|<\frac{\varepsilon}{4M}$。因此当 $A'>A>A_0$ 时, 由式 \eqref{eq:ma-17-2-2} 即有
% \[
% \begin{aligned}
%   \left|\int_A^{A'}f(x,y)g(x,y)\mathrm dy\right|
%   &<\frac{\varepsilon}{4M}\left|\int_c^\xi f(x,y)\mathrm dy-\int_c^A f(x,y)\mathrm dy\right|\\
%   &\quad+\frac{\varepsilon}{4M}\left|\int_c^{A'}f(x,y)\mathrm dy-\int_c^\xi f(x,y)\mathrm dy\right|\\
%   &\leqslant\frac{\varepsilon}{4M}(2M+2M)=\varepsilon.
% \end{aligned}
% \qedhere
% \]
% \end{proof}
% 
% 
% \begin{theorem}[{阿贝尔判别法}]\label{theorem:ma-17-2-4}
% 设函数 $f(x,y),g(x,y)$ 在 $E\times[c,+\infty)$ 上有定义 (其中 $E\subset\mathbb R$ 是一个区间), 并且满足:
% \[
% \begin{gathered}
%   (1)\quad \int_c^{+\infty}f(x,y)\mathrm dy \text{ 在 } E \text{ 上一致收敛};\\
%   (2)\quad \text{对任意固定的 } x\in E,\ g(x,y) \text{ 是 } y \text{ 的单调函数, 并且存在常数 } M>0,\\
%   \text{对于 } \forall x\in E \text{ 及 } \forall y\in[c,+\infty),\ \text{有 } |g(x,y)|\leqslant M,
% \end{gathered}
% \]
% 则含参变量无穷积分 $\int_c^{+\infty}f(x,y)g(x,y)\mathrm dy$ 在 $E$ 上一致收敛.
% \end{theorem}
% 
% 
% \begin{proof}
% 对于 $A'>A>c$, 我们仍然有式 \eqref{eq:ma-17-2-2}.
% 
% 由 $\int_c^{+\infty}f(x,y)\mathrm dy$ 在 $E$ 上一致收敛, 对于 $\forall\varepsilon>0,\exists A_0>c$, 当 $A'>A>A_0$ 时, 对一切 $x\in E$, 有
% \[
%   \left|\int_A^{A'}f(x,y)\mathrm dy\right|<\frac{\varepsilon}{2M},
% \]
% 从而由式 \eqref{eq:ma-17-2-2} 得
% \[
%   \left|\int_A^{A'}f(x,y)g(x,y)\mathrm dy\right|
%   <\frac{\varepsilon}{2M}\bigl(|g(x,A)|+|g(x,A')|\bigr)\leqslant\varepsilon.
% \qedhere
% \]
% \end{proof}
% 
% 
% \begin{example}\label{example:ma-17-2-1}
% 证明含参变量无穷积分 $\int_0^{+\infty}y^xe^{-y}\mathrm dy$, 满足:
% \[
% \begin{gathered}
%   (1)\quad \text{在 } [0,M]\ (0<M<+\infty) \text{ 上一致收敛};\\
%   (2)\quad \text{在 } [0,+\infty) \text{ 上不一致收敛}.
% \end{gathered}
% \]
% \end{example}
% 
% 
% \begin{proof}
% (1) 对于 $\forall x\in[0,M]$, 有
% \[
%   0\leqslant y^xe^{-y}\leqslant y^M e^{-y}.
% \]
% 由于 $\int_0^{+\infty}y^M e^{-y}\mathrm dy$ 收敛, 由定理~\ref{theorem:ma-17-2-4} 知, $\int_0^{+\infty}y^xe^{-y}\mathrm dy$ 在 $[0,M]$ 上一致收敛.
% 
% (2) 对任意的 $A_0>1$, 任取 $A_1>A_0,A_2=A_1+1$, 考查 $\int_{A_1}^{A_2}y^xe^{-y}\mathrm dy$。由于
% \[
%   \lim_{x\to+\infty}A_1^x e^{-A_2}=+\infty,
% \]
% 从而存在 $x'\in(0,+\infty)$, 使得 $A_1^{x'}e^{-A_2}\geqslant1$, 因此
% \[
%   \int_{A_1}^{A_2}y^{x'}e^{-y}\mathrm dy
%   \geqslant \int_{A_1}^{A_2}A_1^{x'}e^{-A_2}\mathrm dy\geqslant1.
% \]
% 这说明 $\int_0^{+\infty}y^xe^{-y}\mathrm dy$ 在 $[0,+\infty)$ 上不一致收敛.
% \end{proof}
% 
% 
% \begin{example}\label{example:ma-17-2-2}
% 证明含参变量无穷积分 $\int_0^{+\infty}\frac{\sin xy}{y}\mathrm dy$
% \[
% \begin{gathered}
%   (1)\quad \text{在 } [\alpha_0,+\infty)\ (\alpha_0>0) \text{ 上一致收敛};\\
%   (2)\quad \text{在 } (0,+\infty) \text{ 内不一致收敛}.
% \end{gathered}
% \]
% \end{example}
% 
% 
% \begin{proof}
% (1) 方法 1\quad 设 $f(x,y)=\sin xy,g(x,y)=\frac1y$, 则对于 $\forall x\in[\alpha_0,+\infty)$ 及 $\forall A>0$, 有
% \[
%   \left|\int_0^A f(x,y)\mathrm dy\right|
%   =
%   \left|\int_0^A\sin xy\,\mathrm dy\right|
%   =
%   \left|\frac1x(1-\cos xA)\right|
%   \leqslant\frac2{\alpha_0}.
% \]
% 而 $g(x,y)=\frac1y$ 关于 $y$ 单调, 且对 $x\in[\alpha_0,+\infty)$ 一致趋于 $0$, 由狄利克雷判别法知 $\int_0^{+\infty}\frac{\sin xy}{y}\mathrm dy$ 在 $[\alpha_0,+\infty)$ 上一致收敛.
% 
% 方法 2\quad 对于 $\forall\varepsilon>0$, 由 $\int_0^{+\infty}\frac{\sin y}{y}\mathrm dy$ 收敛, 从而存在 $A_0>0$, 使得当 $A>A_0$ 时, 有
% \[
%   \left|\int_A^{+\infty}\frac{\sin y}{y}\mathrm dy\right|<\varepsilon.
% \]
% 取 $A_1=\frac{A_0}{\alpha_0}$, 则当 $A>A_1$ 时, 对一切 $x\in[\alpha_0,+\infty)$, 有 $xA>\alpha_0 A_1=A_0$, 从而有
% \[
%   \left|\int_A^{+\infty}\frac{\sin xy}{y}\mathrm dy\right|
%   =
%   \left|\int_{xA}^{+\infty}\frac{\sin y}{y}\mathrm dy\right|
%   <\varepsilon.
% \]
% 因此 $\int_0^{+\infty}\frac{\sin xy}{y}\mathrm dy$ 在 $[\alpha_0,+\infty)$ 上一致收敛.
% 
% (2) 取 $x=\frac1{2k}\ (k\in\mathbb N)$, 我们有
% \[
% \begin{aligned}
%   \left|\int_{4k\pi}^{5k\pi}\frac{\sin xy}{y}\mathrm dy\right|
%   &=\left|\int_{4k\pi}^{5k\pi}\frac{\sin\frac{y}{2k}}{y}\mathrm dy\right|\\
%   &\geqslant \frac1{5k\pi}\left|\int_{4k\pi}^{5k\pi}\sin\frac{y}{2k}\mathrm dy\right|
%   =\frac2{5\pi}.
% \end{aligned}
% \]
% 因此存在 $\varepsilon_0=\frac1{5\pi}$, 对 $\forall A>0$, 存在正整数 $k_0>A$, 从而存在 $A'=4k_0\pi,A''=5k_0\pi>A$ 及 $x'=\frac1{2k_0}\in(0,+\infty)$, 使得
% \[
%   \left|\int_{A'}^{A''}\frac{\sin x'y}{y}\mathrm dy\right|>\varepsilon_0.
% \]
% 这说明 $\int_0^{+\infty}\frac{\sin xy}{y}\mathrm dy$ 在 $(0,+\infty)$ 内不一致收敛.
% \end{proof}
% 
% 
% \begin{example}\label{example:ma-17-2-3}
% 设函数 $f(x)$ 在 $[0,+\infty)$ 内连续, 证明无穷积分
% \[
%   \int_0^{+\infty}e^{-\alpha x}f(x)\mathrm dx
% \]
% 对 $\alpha\in(0,+\infty)$ 一致收敛的充分必要条件是无穷积分 $\int_0^{+\infty}f(x)\mathrm dx$ 收敛.
% \end{example}
% 
% 
% \begin{proof}
% 充分性。若无穷积分 $\int_0^{+\infty}f(x)\mathrm dx$ 收敛, 则它对参数 $\alpha\in(0,+\infty)$ 一致收敛。又 $e^{-\alpha x}$ 关于 $x$ 单调, 且对 $\alpha\in(0,+\infty)$ 有 $|e^{-\alpha x}|\leqslant1$, 由阿贝尔判别法知 $\int_0^{+\infty}e^{-\alpha x}f(x)\mathrm dx$ 在 $(0,+\infty)$ 上一致收敛.
% 
% 必要性。假设 $\int_0^{+\infty}f(x)\mathrm dx$ 不收敛, 则存在 $\varepsilon_0>0$, 对任意 $A_0>0$, 存在 $A_2>A_1>A_0$, 使得
% \[
%   \left|\int_{A_1}^{A_2}f(x)\mathrm dx\right|>2\varepsilon_0.
% \]
% 由于 $e^{-\alpha x}f(x)$ 在 $[A_1,A_2]\times[0,+\infty)$ 上连续, 所以
% \[
%   \lim_{\alpha\to0+0}\int_{A_1}^{A_2}e^{-\alpha x}f(x)\mathrm dx
%   =\int_{A_1}^{A_2}f(x)\mathrm dx.
% \]
% 于是存在 $\alpha'>0$, 使得
% \[
%   \left|\int_{A_1}^{A_2}e^{-\alpha'x}f(x)\mathrm dx\right|
%   \geqslant \frac12\left|\int_{A_1}^{A_2}f(x)\mathrm dx\right|
%   >\varepsilon_0.
% \]
% 这说明 $\int_0^{+\infty}e^{-\alpha x}f(x)\mathrm dx$ 在 $(0,+\infty)$ 上不一致收敛, 矛盾.
% \end{proof}
% 
% 
% \subsection{含参变量无穷积分的性质}
% \label{subsec:ma-17-2-2}
% 
% 本节讨论含参变量无穷积分的连续性、可积性与可微性.
% 
% \begin{theorem}\label{theorem:ma-17-2-5}
% 设函数 $f(x,y)$ 在 $E\times[c,+\infty)$ 上有定义, 其中 $E\subset\mathbb R$, 则含参变量无穷积分
% \[
%   \int_c^{+\infty}f(x,y)\mathrm dy
% \]
% 在 $E$ 上一致收敛的充分必要条件是: 对任意的满足条件
% \[
%   c<t_1<t_2<\cdots<t_k<\cdots,\qquad \lim_{k\to\infty}t_k=+\infty
% \]
% 的序列 $\{t_k\}$, 函数列
% \[
%   F_k(x)=\int_c^{t_k}f(x,y)\mathrm dy\quad(k=1,2,\cdots)
% \]
% 在 $E$ 上一致收敛.
% \end{theorem}
% 
% 
% \begin{proof}
% 必要性。设 $\int_c^{+\infty}f(x,y)\mathrm dy$ 在 $E$ 上一致收敛。对 $\forall\varepsilon>0$, 存在 $A_0>0$, 使得当 $A',A''>A_0$ 时, 对一切 $x\in E$, 有
% \[
%   \left|\int_{A'}^{A''}f(x,y)\mathrm dy\right|<\varepsilon.
% \]
% 任取满足条件的序列 $\{t_k\}$, 当 $k,k'$ 充分大时, 有 $t_k,t_{k'}>A_0$, 故
% \[
%   |F_k(x)-F_{k'}(x)|
%   =\left|\int_{t_{k'}}^{t_k}f(x,y)\mathrm dy\right|<\varepsilon
%   \quad(x\in E),
% \]
% 即 $\{F_k(x)\}$ 在 $E$ 上一致收敛.
% 
% 充分性。反设 $\int_c^{+\infty}f(x,y)\mathrm dy$ 在 $E$ 上不一致收敛。则存在 $\varepsilon_0>0$, 对任意 $A>0$, 存在 $A',A''>A$ 及 $x'\in E$, 使得
% \[
%   \left|\int_{A'}^{A''}f(x',y)\mathrm dy\right|\geqslant\varepsilon_0.
% \]
% 取 $A=1$, 得到 $t_1',t_1''$ 与 $x_1\in E$; 依次取 $A=t_k''$, 可得到序列
% \[
%   k<t_k'<t_k''<t_{k+1}'<t_{k+1}''\quad(k=1,2,\cdots)
% \]
% 及 $x_k\in E$, 满足
% \[
%   \left|\int_{t_k'}^{t_k''}f(x_k,y)\mathrm dy\right|\geqslant\varepsilon_0.
% \]
% 令 $\{t_k\}=\{t_1',t_1'',t_2',t_2'',\cdots\}$, 则相应的函数列 $F_k(x)=\int_c^{t_k}f(x,y)\mathrm dy$ 不在 $E$ 上一致收敛, 矛盾.
% \end{proof}
% 
% 
% 容易看出, 在定理~\ref{theorem:ma-17-2-5} 中, 若 $\int_c^{+\infty}f(x,y)\mathrm dy$ 在 $E$ 上一致收敛, 则对任意满足该定理条件的 $\{t_k\}$, 当 $k\to\infty$ 时, 都有
% \[
%   F_k(x)\rightrightarrows \int_c^{+\infty}f(x,y)\mathrm dy\quad(x\in E).
% \]
% 
% \begin{theorem}\label{theorem:ma-17-2-6}
% 设函数 $f(x,y)$ 在 $E\times[c,+\infty)$ 上连续, 其中 $E\subset\mathbb R$ 是一个区间, 并且含参变量无穷积分 $\int_c^{+\infty}f(x,y)\mathrm dy$ 在 $E$ 上一致收敛到函数 $I(x)$, 则 $I(x)$ 在 $E$ 上连续.
% \end{theorem}
% 
% 
% \begin{proof}
% 取 $t_k=c+k\ (k=1,2,\cdots)$, 令
% \[
%   F_k(x)=\int_c^{t_k}f(x,y)\mathrm dy.
% \]
% 由定理~\ref{theorem:ma-17-1-1} 知, $F_k(x)$ 在 $E$ 上连续。又由定理~\ref{theorem:ma-17-2-5} 知 $F_k(x)\rightrightarrows I(x)$, 因此 $I(x)$ 在 $E$ 上连续.
% \end{proof}
% 
% 
% \begin{theorem}\label{theorem:ma-17-2-7}
% 设函数 $f(x,y)$ 在 $[a,b]\times[c,+\infty)$ 上连续, 且含参变量无穷积分 $\int_c^{+\infty}f(x,y)\mathrm dy$ 在 $[a,b]$ 上一致收敛, 则有
% \begin{equation}
%   \int_a^b\mathrm dx\int_c^{+\infty}f(x,y)\mathrm dy
%   =
%   \int_c^{+\infty}\mathrm dy\int_a^b f(x,y)\mathrm dx.
%   \label{eq:ma-17-2-3}
% \end{equation}
% \end{theorem}
% 
% 
% \begin{proof}
% 由定理~\ref{theorem:ma-17-2-6} 知 $I(x)=\int_c^{+\infty}f(x,y)\mathrm dy$ 在 $[a,b]$ 上连续。对任意满足定理~\ref{theorem:ma-17-2-5} 条件的序列 $\{t_k\}$, 令 $F_k(x)=\int_c^{t_k}f(x,y)\mathrm dy$, 则 $F_k(x)\rightrightarrows I(x)$。因此
% \[
% \begin{aligned}
%   \int_a^b I(x)\mathrm dx
%   &=\int_a^b\lim_{k\to\infty}F_k(x)\mathrm dx\\
%   &=\lim_{k\to\infty}\int_a^b\mathrm dx\int_c^{t_k}f(x,y)\mathrm dy\\
%   &=\lim_{k\to\infty}\int_c^{t_k}\mathrm dy\int_a^b f(x,y)\mathrm dx.
% \end{aligned}
% \]
% 由于 $\{t_k\}$ 是任意的, 故 $\int_c^{+\infty}\mathrm dy\int_a^b f(x,y)\mathrm dx$ 收敛且等式 \eqref{eq:ma-17-2-3} 成立.
% \end{proof}
% 
% 
% \begin{theorem}\label{theorem:ma-17-2-8}
% 设函数 $f(x,y)$ 及其偏导数 $f'_x(x,y)$ 在 $E\times[c,+\infty)$ 上连续, 其中 $E\subset\mathbb R$ 是一个区间, 再设存在 $x_0\in E$, 使得 $\int_c^{+\infty}f(x_0,y)\mathrm dy$ 收敛, 并且 $\int_c^{+\infty}f'_x(x,y)\mathrm dy$ 在 $E$ 上一致收敛, 则
% \[
% \begin{gathered}
%   (1)\quad I(x)=\int_c^{+\infty}f(x,y)\mathrm dy \text{ 在 } E \text{ 上一致收敛};\\
%   (2)\quad I'(x)=\left(\int_c^{+\infty}f(x,y)\mathrm dy\right)'
%   =\int_c^{+\infty}f'_x(x,y)\mathrm dy.
% \end{gathered}
% \]
% \end{theorem}
% 
% 
% \begin{proof}
% 任取满足定理~\ref{theorem:ma-17-2-5} 条件的序列 $\{t_k\}$, 令
% \[
%   F_k(x)=\int_c^{t_k}f(x,y)\mathrm dy.
% \]
% 由定理~\ref{theorem:ma-17-1-3} 知
% \[
%   F'_k(x)=\int_c^{t_k}f'_x(x,y)\mathrm dy.
% \]
% 因为数列 $\{F_k(x_0)\}$ 收敛, 且函数列 $\{F'_k(x)\}$ 在 $E$ 上一致收敛, 所以 $\{F_k(x)\}$ 在 $E$ 上一致收敛, 并且可逐项求导。于是结论成立.
% \end{proof}
% 
% 
% \begin{example}\label{example:ma-17-2-5}
% 计算狄利克雷积分 $\int_0^{+\infty}\frac{\sin x}{x}\mathrm dx$.
% \end{example}
% 
% 
% \begin{solve}
% 令
% \[
%   I(\alpha)=\int_0^{+\infty}e^{-\alpha x}\frac{\sin x}{x}\mathrm dx\quad(\alpha>0).
% \]
% 由阿贝尔判别法知该积分对 $\alpha\in(0,+\infty)$ 一致收敛。对任意 $\alpha_0>0$, $\int_0^{+\infty}(-e^{-\alpha x}\sin x)\mathrm dx$ 在 $[\alpha_0,+\infty)$ 上一致收敛, 故
% \[
% \begin{aligned}
%   I'(\alpha)
%   &=-\int_0^{+\infty}e^{-\alpha x}\sin x\,\mathrm dx\\
%   &=e^{-\alpha x}\cos x\Big|_0^{+\infty}
%     +\alpha\int_0^{+\infty}e^{-\alpha x}\cos x\,\mathrm dx\\
%   &=-1+\alpha e^{-\alpha x}\sin x\Big|_0^{+\infty}
%     +\alpha^2\int_0^{+\infty}e^{-\alpha x}\sin x\,\mathrm dx.
% \end{aligned}
% \]
% 从而 $I'(\alpha)=-\frac1{1+\alpha^2}$, 所以
% \[
%   I(\alpha)=-\arctan\alpha+C.
% \]
% 又
% \[
%   |I(\alpha)|\leqslant\int_0^{+\infty}e^{-\alpha x}\mathrm dx=\frac1\alpha,
% \]
% 故 $\lim_{\alpha\to+\infty}I(\alpha)=0$, 从而 $C=\frac\pi2$。由于 $I(\alpha)$ 在 $\alpha=0$ 处右连续, 得
% \[
%   \int_0^{+\infty}\frac{\sin x}{x}\mathrm dx
%   =\lim_{\alpha\to0+0}I(\alpha)=\frac\pi2.
% \qedhere
% \]
% \end{solve}
% 
% 
% \subsection{含参变量瑕积分}
% \label{subsec:ma-17-2-3}
% 
% 设函数 $f(x,y)$ 在 $[a,b]\times(c,d]$ 上连续, 当 $x\in[a,b]$ 时, $f(x,y)$ 以 $c$ 为瑕点。若对任意 $x\in[a,b]$, 瑕积分
% \begin{equation}
%   I(x)=\int_c^d f(x,y)\mathrm dy
%   \label{eq:ma-17-2-4}
% \end{equation}
% 收敛, 则 $I(x)$ 为在 $[a,b]$ 上有定义的函数。通常称 $I(x)$ 为含参变量瑕积分, $x$ 为参变量.
% 
% 类似于含参变量无穷积分, 可定义含参变量瑕积分的一致收敛性, 并建立柯西准则、魏尔斯特拉斯判别法、狄利克雷判别法、阿贝尔判别法以及连续性、可积性和可微性等结论。例如作变量替换
% \[
%   y=c+\frac1t,
% \]
% 则
% \[
%   \int_c^d f(x,y)\mathrm dy
%   =
%   \int_{\frac1{d-c}}^{+\infty}
%   f\left(x,c+\frac1t\right)\frac{\mathrm dt}{t^2},
% \]
% 从而含参变量瑕积分可化为含参变量无穷积分来讨论.
% 
% 上述结论也可以推广到含多个参变量的积分。例如, 设 $S$ 为空间中的一片光滑曲面, $\mathbf F(x,y,z,t)$ 为流体在点 $(x,y,z)$、时刻 $t$ 的速度场, $\mathbf n(x,y,z)$ 为曲面的单位法向量。则单位时间内通过曲面 $S$ 的流量为
% \[
%   S(t)=\iint_S \mathbf F(x,y,z,t)\cdot\mathbf n(x,y,z)\mathrm dS.
% \]
% 若 $\mathbf F$ 关于 $t$ 连续, 则 $S(t)$ 在相应区间上连续。因而在时间区间 $[t_0,t_1]$ 内通过曲面 $S$ 的流量为
% \[
%   Q=\int_{t_0}^{t_1}\mathrm dt
%   \iint_S \mathbf F(x,y,z,t)\cdot\mathbf n(x,y,z)\mathrm dS.
% \]
% 特别地, 若速度场与时间无关, 则
% \[
%   Q=\int_{t_0}^{t_0+1}\mathrm dt
%   \iint_S \mathbf F(x,y,z)\cdot\mathbf n(x,y,z)\mathrm dS
%   =\iint_S \mathbf F(x,y,z)\cdot\mathbf n(x,y,z)\mathrm dS.
% \]
% 
% \section{\texorpdfstring{$\Gamma$ 函数与 B 函数}{Γ 函数与 B 函数}}
% \label{sec:ma-17-3}
% 
% \subsection{\texorpdfstring{$\Gamma$ 函数}{Γ 函数}}
% \label{subsec:ma-17-3-1}
% 
% \begin{definition}\label{definition:ma-17-3-1}
% 称含参变量积分
% \[
%   \Gamma(s)=\int_0^{+\infty}x^{s-1}e^{-x}\mathrm dx
% \]
% 为 $\Gamma$ 函数.
% \end{definition}
% 
% 当 $s<1$ 时, $x=0$ 为瑕点。将积分分成
% \[
%   \int_0^{+\infty}x^{s-1}e^{-x}\mathrm dx
%   =
%   \int_0^1x^{s-1}e^{-x}\mathrm dx+\int_1^{+\infty}x^{s-1}e^{-x}\mathrm dx.
% \]
% 第一个积分当且仅当 $s>0$ 时收敛, 第二个积分对任意实数 $s$ 都收敛。因此 $\Gamma$ 函数的定义域为 $(0,+\infty)$, 且 $\Gamma(s)>0$.
% 
% \begin{property}\label{property:ma-17-3-1}
% $\Gamma(s+1)=s\Gamma(s)\ (s>0)$。特别地, 当 $k$ 为非负整数时,
% \[
%   \Gamma(k+1)=k!.
% \]
% \end{property}
% 
% 
% \begin{proof}
% \[
% \begin{aligned}
%   \Gamma(s+1)
%   &=\int_0^{+\infty}x^s e^{-x}\mathrm dx
%   =-\int_0^{+\infty}x^s\mathrm d e^{-x}\\
%   &=-x^s e^{-x}\Big|_0^{+\infty}
%   +s\int_0^{+\infty}x^{s-1}e^{-x}\mathrm dx
%   =s\Gamma(s).
% \end{aligned}
% \]
% 由此递推即得 $\Gamma(k+1)=k!$.
% \end{proof}
% 
% 
% \begin{property}\label{property:ma-17-3-2}
% 对 $s>0$, 有
% \[
%   \Gamma(s)=2\int_0^{+\infty}x^{2s-1}e^{-x^2}\mathrm dx.
% \]
% \end{property}
% 
% 
% \begin{proof}
% 在 $\Gamma(s)$ 的定义式中令 $x=t^2$, 即得.
% \end{proof}
% 
% 
% \begin{property}\label{property:ma-17-3-3}
% $\Gamma(s)\in C^\infty(0,+\infty)$.
% \end{property}
% 
% 
% \begin{proof}
% 对任意非负整数 $k$,
% \[
%   \frac{\partial^k}{\partial s^k}\bigl(x^{s-1}e^{-x}\bigr)
%   =x^{s-1}(\ln x)^k e^{-x}.
% \]
% 任取 $s_0>0$。当 $0<x\leqslant1$ 且 $s\in[\frac{s_0}2,s_0+1]$ 时,
% \[
%   x^{s-1}|\ln x|^k e^{-x}\leqslant x^{\frac{s_0}2-1}|\ln x|^k,
% \]
% 右端在 $(0,1]$ 上可积; 当 $x\geqslant1$ 时,
% \[
%   x^{s-1}|\ln x|^k e^{-x}\leqslant x^{s_0+k}e^{-x},
% \]
% 右端在 $[1,+\infty)$ 上可积。因此
% \[
%   \int_0^{+\infty}x^{s-1}(\ln x)^k e^{-x}\mathrm dx
% \]
% 在 $[\frac{s_0}2,s_0+1]$ 上一致收敛, 从而
% \[
%   \Gamma^{(k)}(s_0)=\int_0^{+\infty}x^{s_0-1}(\ln x)^k e^{-x}\mathrm dx.
% \]
% 由于 $s_0$ 与 $k$ 任意, 得 $\Gamma(s)\in C^\infty(0,+\infty)$.
% \end{proof}
% 
% 
% \begin{property}\label{property:ma-17-3-4}
% $\Gamma(s)$ 与 $\ln\Gamma(s)$ 在 $(0,+\infty)$ 内都是严格凸函数.
% \end{property}
% 
% 
% \begin{proof}
% 由性质~\ref{property:ma-17-3-3},
% \[
%   \Gamma''(s)=\int_0^{+\infty}x^{s-1}(\ln x)^2e^{-x}\mathrm dx>0,
% \]
% 故 $\Gamma(s)$ 严格凸。又
% \[
%   (\ln\Gamma(s))''=\frac{\Gamma(s)\Gamma''(s)-(\Gamma'(s))^2}{\Gamma^2(s)}.
% \]
% 由柯西不等式,
% \[
% \begin{aligned}
%   \Gamma(s)\Gamma''(s)
%   &=
%   \int_0^{+\infty}x^{s-1}e^{-x}\mathrm dx
%   \int_0^{+\infty}x^{s-1}(\ln x)^2e^{-x}\mathrm dx\\
%   &>
%   \left(\int_0^{+\infty}|x^{s-1}\ln x\,e^{-x}|\mathrm dx\right)^2
%   \geqslant(\Gamma'(s))^2,
% \end{aligned}
% \]
% 因此 $\ln\Gamma(s)$ 严格凸.
% \end{proof}
% 
% 
% \subsection{B 函数}
% \label{subsec:ma-17-3-2}
% 
% \begin{definition}\label{definition:ma-17-3-2}
% 称含参变量积分
% \[
%   B(p,q)=\int_0^1x^{p-1}(1-x)^{q-1}\mathrm dx
% \]
% 为 B 函数, 其中 $p>0,q>0$.
% \end{definition}
% 
% 
% \begin{property}\label{property:ma-17-3-5}
% B 函数具有以下性质:
% \[
% \begin{gathered}
%   (1)\quad B(p,q)=B(q,p);\\
%   (2)\quad B(p,q)=\frac{p-1}{p+q-1}B(p-1,q)\quad(p>1,q>0).
% \end{gathered}
% \]
% \end{property}
% 
% 
% \begin{proof}
% (1) 在定义式中令 $x=1-t$ 即得.
% 
% (2) 由分部积分,
% \[
% \begin{aligned}
%   B(p,q)
%   &=\int_0^1x^{p-1}(1-x)^{q-1}\mathrm dx\\
%   &=-\frac1q\int_0^1x^{p-1}\mathrm d(1-x)^q\\
%   &=\frac{p-1}{q}\int_0^1x^{p-2}(1-x)^q\mathrm dx\\
%   &=\frac{p-1}{q}\bigl(B(p-1,q)-B(p,q)\bigr).
% \end{aligned}
% \]
% 整理即得结论。特别地, 当 $p>1,q>1$ 时,
% \[
%   B(p,q)=\frac{(p-1)(q-1)}{(p+q-1)(p+q-2)}B(p-1,q-1).
% \qedhere
% \]
% \end{proof}
% 
% 
% \begin{property}\label{property:ma-17-3-6}
% B 函数具有以下形式:
% \[
% \begin{gathered}
%   (1)\quad B(p,q)=2\int_0^{\frac\pi2}\cos^{2p-1}\theta\sin^{2q-1}\theta\,\mathrm d\theta\quad(p,q>0);\\
%   (2)\quad B(p,q)=\int_0^{+\infty}\frac{x^{q-1}}{(1+x)^{p+q}}\mathrm dx\quad(p,q>0).
% \end{gathered}
% \]
% \end{property}
% 
% 
% \begin{proof}
% (1) 令 $x=\sin^2\theta$ 作积分变换即可.
% 
% (2) 令 $x=\frac{t}{1+t}$ 作积分变换即可.
% \end{proof}
% 
% 
% \subsection{\texorpdfstring{$\Gamma$ 函数与 B 函数的关系}{Γ 函数与 B 函数的关系}}
% \label{subsec:ma-17-3-3}
% 
% \begin{theorem}\label{theorem:ma-17-3-7}
% 设 $p>0,q>0$, 则有
% \[
%   B(p,q)=\frac{\Gamma(p)\Gamma(q)}{\Gamma(p+q)}.
% \]
% \end{theorem}
% 
% 
% \begin{proof}
% 由性质~\ref{property:ma-17-3-2} 有
% \[
%   \Gamma(p)=2\int_0^{+\infty}x^{2p-1}e^{-x^2}\mathrm dx,\quad
%   \Gamma(q)=2\int_0^{+\infty}y^{2q-1}e^{-y^2}\mathrm dy.
% \]
% 若令 $D=\{(x,y):0\leqslant x<+\infty,0\leqslant y<+\infty\}$, 则有
% \[
%   \Gamma(p)\Gamma(q)
%   =4\iint_D x^{2p-1}y^{2q-1}e^{-(x^2+y^2)}\mathrm dx\mathrm dy.
% \]
% 利用极坐标变换 $x=r\cos\theta,\ y=r\sin\theta$, 并令
% \[
%   D_1=\{(r,\theta):0<r<+\infty,0\leqslant\theta\leqslant\frac\pi2\},
% \]
% 则
% \[
% \begin{aligned}
%   \Gamma(p)\Gamma(q)
%   &=4\iint_{D_1}r^{2(p+q)-1}e^{-r^2}
%     \cos^{2p-1}\theta\sin^{2q-1}\theta\,\mathrm dr\mathrm d\theta\\
%   &=
%   \left(2\int_0^{\frac\pi2}\cos^{2p-1}\theta\sin^{2q-1}\theta\,\mathrm d\theta\right)
%   \left(2\int_0^{+\infty}r^{2(p+q)-1}e^{-r^2}\mathrm dr\right)\\
%   &=B(p,q)\Gamma(p+q).
% \end{aligned}
% \qedhere
% \]
% \end{proof}
% 
% 
% \begin{corollary}\label{corollary:ma-17-source-56}
% $B(p,q)$ 在 $(0,+\infty)\times(0,+\infty)$ 内具有任意阶偏导数.
% \end{corollary}
% 
% 
% \begin{theorem}[{余元公式}]\label{theorem:ma-17-3-8}
% 设 $0<p<1$, 则有
% \[
%   B(p,1-p)=\Gamma(p)\Gamma(1-p)=\frac{\pi}{\sin p\pi}.
% \]
% \end{theorem}
% 
% 
% \begin{proof}
% 由于
% \[
%   B(p,1-p)=\int_0^{+\infty}\frac{x^{p-1}}{1+x}\mathrm dx,
% \]
% 利用变量替换 $x=\frac1t$ 有
% \[
%   \int_1^{+\infty}\frac{x^{p-1}}{1+x}\mathrm dx
%   =\int_0^1\frac{x^{-p}}{1+x}\mathrm dx.
% \]
% 因此
% \[
%   B(p,1-p)=\int_0^1\frac{x^{p-1}+x^{-p}}{1+x}\mathrm dx.
% \]
% 我们将 $\frac1{1+x}$ 展成幂级数, 从而有
% \[
% \begin{aligned}
%   B(p,1-p)
%   &=\lim_{r\to1-0}\int_0^r\frac{x^{p-1}+x^{-p}}{1+x}\mathrm dx\\
%   &=\lim_{r\to1-0}\int_0^r
%   \left[\sum_{k=0}^{+\infty}(-1)^k x^{k+p-1}
%   +\sum_{k=0}^{+\infty}(-1)^k x^{k-p}\right]\mathrm dx\\
%   &=\sum_{k=0}^{+\infty}\frac{(-1)^k}{k+p}
%   +\sum_{k=0}^{+\infty}\frac{(-1)^k}{k-p+1}\\
%   &=\frac1p+\sum_{k=1}^{+\infty}(-1)^k
%   \left(\frac1{k+p}+\frac1{p-k}\right)\\
%   &=\frac1p+\sum_{k=1}^{+\infty}(-1)^k\frac{2p}{p^2-k^2}.
% \end{aligned}
% \]
% 由于 $\cos px$ 的傅里叶级数
% \[
%   \cos px=\frac{\sin p\pi}{\pi}
%   \left[\frac1p+\sum_{k=1}^{+\infty}(-1)^k\frac{2p}{p^2-k^2}\cos kx\right]
% \]
% 在 $|x|\leqslant\pi$ 处处收敛, 令 $x=0$ 即得
% \[
%   B(p,1-p)=\frac1p+\sum_{k=1}^{+\infty}(-1)^k\frac{2p}{p^2-k^2}
%   =\frac{\pi}{\sin p\pi}.
% \qedhere
% \]
% \end{proof}
% 
% 
% 特别地, 在余元公式中令 $p=\frac12$, 得
% \[
%   B\left(\frac12,\frac12\right)=2\int_0^{\frac\pi2}\mathrm d\theta=\pi.
% \]
% 又
% \[
%   \Gamma^2\left(\frac12\right)
%   =\frac{\Gamma\left(\frac12\right)\Gamma\left(\frac12\right)}
%   {\Gamma\left(\frac12+\frac12\right)}
%   =B\left(\frac12,\frac12\right)=\pi\quad(\Gamma(1)=1),
% \]
% 即
% \[
%   \Gamma\left(\frac12\right)=\sqrt\pi.
% \]
% 由递推公式及余元公式, 若我们知道 $\Gamma(s)$ 在 $\left(0,\frac12\right)$ 中的值, 则可求出它在 $(0,+\infty)$ 的所有值。值得注意的是, 我们又一次得到了
% \[
%   \int_0^{+\infty}e^{-x^2}\mathrm dx=\frac12\Gamma\left(\frac12\right)=\frac{\sqrt\pi}{2}.
% \]
% 
% \begin{example}\label{example:ma-17-3-1}
% 计算定积分
% \[
%   I_n=\int_0^{\frac\pi2}\sin^n x\,\mathrm dx
%   =\int_0^{\frac\pi2}\cos^n x\,\mathrm dx.
% \]
% \end{example}
% 
% 
% \begin{solve}
% 由 B 函数的定义我们有
% \[
% \begin{aligned}
%   I_n
%   &=\frac12\cdot2\int_0^{\frac\pi2}\cos^{\frac12\cdot2-1}x
%     \sin^{\frac{2(n+1)}2-1}x\,\mathrm dx\\
%   &=\frac12 B\left(\frac12,\frac{n+1}{2}\right)
%   =\frac{\Gamma\left(\frac12\right)\Gamma\left(\frac{n+1}{2}\right)}
%   {2\Gamma\left(\frac12+\frac{n+1}{2}\right)}.
% \end{aligned}
% \]
% 因此
% \[
% I_n=
% \begin{cases}
%   \displaystyle \frac{(n-1)!!}{n!!}\cdot\frac\pi2, & n\text{ 为偶数},\\[1.2ex]
%   \displaystyle \frac{(n-1)!!}{n!!}, & n\text{ 为奇数}.
% \end{cases}
% \qedhere
% \]
% \end{solve}
% 
% 
% \begin{example}\label{example:ma-17-3-2}
% 求 $\mathbb R^n$ 中单位球体 $D:x_1^2+x_2^2+\cdots+x_n^2\leqslant1$ 的体积.
% \end{example}
% 
% 
% \begin{solve}
% 由 $n$ 重积分的几何意义知, 所求体积为
% \[
%   V=\idotsint_D \mathrm dx_1\mathrm dx_2\cdots\mathrm dx_n.
% \]
% 类似于三维情形的球坐标变换, 我们令变换为
% \[
% \left\{
% \begin{aligned}
%   x_1&=r\cos\theta_1,\\
%   x_2&=r\sin\theta_1\cos\theta_2,\\
%   x_3&=r\sin\theta_1\sin\theta_2\cos\theta_3,\\
%   &\cdots\cdots\cdots\\
%   x_{n-1}&=r\sin\theta_1\sin\theta_2\cdots\sin\theta_{n-2}\cos\theta_{n-1},\\
%   x_n&=r\sin\theta_1\sin\theta_2\cdots\sin\theta_{n-2}\sin\theta_{n-1},
% \end{aligned}
% \right.
% \]
% 其中 $0<r<1,\ 0<\theta_1,\theta_2,\cdots,\theta_{n-2}<\pi,\ 0<\theta_{n-1}<2\pi$, 则其雅可比行列式
% \[
%   \frac{\partial(x_1,x_2,\cdots,x_n)}
%   {\partial(r,\theta_1,\theta_2,\cdots,\theta_{n-2})}
%   =r^{n-1}\sin^{n-2}\theta_1\sin^{n-3}\theta_2\cdots\sin\theta_{n-2}.
% \]
% 由此得
% \[
% \begin{aligned}
%   V
%   &=\idotsint_D \mathrm dx_1\mathrm dx_2\cdots\mathrm dx_n\\
%   &=\int_0^{2\pi}\mathrm d\theta_{n-1}\int_0^\pi\mathrm d\theta_{n-2}
%     \cdots\int_0^\pi\mathrm d\theta_1\int_0^1
%     r^{n-1}\sin^{n-2}\theta_1\sin^{n-3}\theta_2\cdots
%     \sin\theta_{n-2}\mathrm dr\\
%   &=\frac{2\pi}{n}
%     \left(\int_0^\pi\sin^{n-2}\theta_1\mathrm d\theta_1\right)
%     \left(\int_0^\pi\sin^{n-3}\theta_2\mathrm d\theta_2\right)\cdots
%     \left(\int_0^\pi\sin\theta_{n-2}\mathrm d\theta_{n-2}\right)\\
%   &=\frac{2\pi}{n}
%     B\left(\frac12,\frac{n-1}{2}\right)
%     B\left(\frac12,\frac{n-2}{2}\right)\cdots B\left(\frac12,1\right)\\
%   &=\frac{2\pi}{n}\cdot
%     \frac{\Gamma\left(\frac12\right)\Gamma\left(\frac{n-1}{2}\right)}{\Gamma\left(\frac n2\right)}
%     \cdot
%     \frac{\Gamma\left(\frac12\right)\Gamma\left(\frac{n-2}{2}\right)}{\Gamma\left(\frac{n-1}{2}\right)}
%     \cdots
%     \frac{\Gamma\left(\frac12\right)\Gamma(1)}{\Gamma\left(\frac32\right)}\\
%   &=\frac{2\pi}{n}\cdot
%     \frac{\Gamma^{n-2}\left(\frac12\right)}{\Gamma\left(\frac n2\right)}
%   =\frac{\pi^{\frac n2}}{\Gamma\left(\frac n2+1\right)}.
% \end{aligned}
% \qedhere
% \]
% \end{solve}
% 
% 
% \section*{习题十七}
% \phantomsection
% \addcontentsline{toc}{section}{习题十七}
% \label{exercises:ma-17}
% 
% 1。举例说明在 $D=[0,1]\times[0,1]$ 上存在函数 $f(x,y)$, 使得它同时满足以下条件:
% 
% (1) $f(x,y)$ 的不连续点在 $D$ 稠密;
% 
% (2) 对于 $\forall x\in[0,1]$, $I(x)=\int_0^1f(x,y)\mathrm dy$ 存在, 并且 $I(x)$ 在 $[0,1]$ 上连续.
% 
% 2。设函数 $f(x,y)$ 在 $D=[a,b]\times[c,d]$ 上连续, 证明: 对于 $\forall(x,u)\in D$, $I(x,u)=\int_c^u f(x,y)\mathrm dy$ 存在, 并且 $I(x,u)$ 在 $D$ 上连续.
% 
% 3。设 $N(0,1)\subset\mathbb R^n$, 函数 $f(x,y)$ 在 $\overline{N(0,1)}\times\overline{N(0,1)}$ 上连续, 证明: 对于 $\forall x\in\overline{N(0,1)}$,
% \[
%   I(x)=\idotsint_{N(0,1)}f(x,y)\mathrm dy_1\mathrm dy_2\cdots\mathrm dy_n
% \]
% 存在, 并且在 $\overline{N(0,1)}$ 上连续.
% 
% 4。求下列极限:
% \[
%   (1)\quad \lim_{x\to0}\int_0^{e^x}\frac{\cos xy}{\sqrt{x^2+y^2+1}}\mathrm dy;
%   \qquad
%   (2)\quad \lim_{x\to1}\int_0^1\frac{\mathrm dy}{1+xy^2}.
% \]
% 
% 5。计算无穷积分 $\displaystyle\int_0^{+\infty}\frac{e^{-ax}-e^{-bx}}{x}\mathrm dx$，其中 $a>b>0$.
% 
% 6。计算无穷积分 $\displaystyle\int_0^{+\infty}\frac{e^{-ax^2}-e^{-bx^2}}{x^2}\mathrm dx$，其中 $b>a>0$.
% 
\end{document}
